\subsection{\Ito{} Integral}
We will work with a Wiener process $W$ on a complete probability space $(\Omega, \cH, \Pm)$ with a right-continuous filtration $\cF \coloneqq (\cF_s, s \ge 0)$ and such that $W$ is adapted to $\cF$. We also use the following definition.
\begin{defn}{Simple Processes}{SimpleProcesses}
    A process $F \coloneqq (F_s, s \in [0, t])$ is called \emph{simple} if it is of the form
    \begin{equation}\label{eq:SimpleProcessDef}
        F_s \coloneqq \xi_0 \I_{\{0\}}(s) + \sum_{k=0}^{n-1} \xi_k \I_{(s_k, s_{k+1}]}(s), \quad s \in [0, t],
    \end{equation}
    where $\{s_0, \dots, s_n : 0 = s_0 < \dots < s_n = t\} =: \Pi_n, n \in \N$, is a segmentation of $[0, t]$ and $\xi_k \in \cF_{s_k}$ for all $k$.
\end{defn}
Note that the simple process in \eqref{eq:SimpleProcessDef} is right-continuous at $s = 0$ and is left-continuous at $s \in (0, t]$. \Cref{fig:SimpleProcessRealisation} shows a realisation of a simple process on $[0, 1]$ for the segmentation
\[
\{s_0, \dots, s_{10}\} = \{0, 0.07, 0.13, 0.19, 0.41, 0.49, 0.52, 0.67, 0.71, 0.76, 1\}.
\]
\begin{figure}[H]
    \centering
    %% Creator: Matplotlib, PGF backend
%%
%% To include the figure in your LaTeX document, write
%%   \input{<filename>.pgf}
%%
%% Make sure the required packages are loaded in your preamble
%%   \usepackage{pgf}
%%
%% Also ensure that all the required font packages are loaded; for instance,
%% the lmodern package is sometimes necessary when using math font.
%%   \usepackage{lmodern}
%%
%% Figures using additional raster images can only be included by \input if
%% they are in the same directory as the main LaTeX file. For loading figures
%% from other directories you can use the `import` package
%%   \usepackage{import}
%%
%% and then include the figures with
%%   \import{<path to file>}{<filename>.pgf}
%%
%% Matplotlib used the following preamble
%%   \def\mathdefault#1{#1}
%%   \everymath=\expandafter{\the\everymath\displaystyle}
%%   \IfFileExists{scrextend.sty}{
%%     \usepackage[fontsize=10.000000pt]{scrextend}
%%   }{
%%     \renewcommand{\normalsize}{\fontsize{10.000000}{12.000000}\selectfont}
%%     \normalsize
%%   }
%%   \usepackage{amsmath}
%%   \usepackage{amssymb}
%%   \usepackage{amsfonts}
%%   \makeatletter\@ifpackageloaded{underscore}{}{\usepackage[strings]{underscore}}\makeatother
%%
\begingroup%
\makeatletter%
\begin{pgfpicture}%
\pgfpathrectangle{\pgfpointorigin}{\pgfqpoint{5.739353in}{4.503305in}}%
\pgfusepath{use as bounding box, clip}%
\begin{pgfscope}%
\pgfsetbuttcap%
\pgfsetmiterjoin%
\definecolor{currentfill}{rgb}{1.000000,1.000000,1.000000}%
\pgfsetfillcolor{currentfill}%
\pgfsetlinewidth{0.000000pt}%
\definecolor{currentstroke}{rgb}{1.000000,1.000000,1.000000}%
\pgfsetstrokecolor{currentstroke}%
\pgfsetdash{}{0pt}%
\pgfpathmoveto{\pgfqpoint{0.000000in}{0.000000in}}%
\pgfpathlineto{\pgfqpoint{5.739353in}{0.000000in}}%
\pgfpathlineto{\pgfqpoint{5.739353in}{4.503305in}}%
\pgfpathlineto{\pgfqpoint{0.000000in}{4.503305in}}%
\pgfpathlineto{\pgfqpoint{0.000000in}{0.000000in}}%
\pgfpathclose%
\pgfusepath{fill}%
\end{pgfscope}%
\begin{pgfscope}%
\pgfsetbuttcap%
\pgfsetmiterjoin%
\definecolor{currentfill}{rgb}{1.000000,1.000000,1.000000}%
\pgfsetfillcolor{currentfill}%
\pgfsetlinewidth{0.000000pt}%
\definecolor{currentstroke}{rgb}{0.000000,0.000000,0.000000}%
\pgfsetstrokecolor{currentstroke}%
\pgfsetstrokeopacity{0.000000}%
\pgfsetdash{}{0pt}%
\pgfpathmoveto{\pgfqpoint{0.563118in}{0.401080in}}%
\pgfpathlineto{\pgfqpoint{5.600618in}{0.401080in}}%
\pgfpathlineto{\pgfqpoint{5.600618in}{4.405080in}}%
\pgfpathlineto{\pgfqpoint{0.563118in}{4.405080in}}%
\pgfpathlineto{\pgfqpoint{0.563118in}{0.401080in}}%
\pgfpathclose%
\pgfusepath{fill}%
\end{pgfscope}%
\begin{pgfscope}%
\pgfpathrectangle{\pgfqpoint{0.563118in}{0.401080in}}{\pgfqpoint{5.037500in}{4.004000in}}%
\pgfusepath{clip}%
\pgfsetbuttcap%
\pgfsetroundjoin%
\pgfsetlinewidth{2.007500pt}%
\definecolor{currentstroke}{rgb}{0.000000,0.000000,0.000000}%
\pgfsetstrokecolor{currentstroke}%
\pgfsetdash{}{0pt}%
\pgfpathmoveto{\pgfqpoint{0.563118in}{4.405080in}}%
\pgfpathlineto{\pgfqpoint{0.915743in}{4.405080in}}%
\pgfusepath{stroke}%
\end{pgfscope}%
\begin{pgfscope}%
\pgfpathrectangle{\pgfqpoint{0.563118in}{0.401080in}}{\pgfqpoint{5.037500in}{4.004000in}}%
\pgfusepath{clip}%
\pgfsetbuttcap%
\pgfsetroundjoin%
\pgfsetlinewidth{2.007500pt}%
\definecolor{currentstroke}{rgb}{0.000000,0.000000,0.000000}%
\pgfsetstrokecolor{currentstroke}%
\pgfsetdash{}{0pt}%
\pgfpathmoveto{\pgfqpoint{0.915743in}{3.437447in}}%
\pgfpathlineto{\pgfqpoint{1.217993in}{3.437447in}}%
\pgfusepath{stroke}%
\end{pgfscope}%
\begin{pgfscope}%
\pgfpathrectangle{\pgfqpoint{0.563118in}{0.401080in}}{\pgfqpoint{5.037500in}{4.004000in}}%
\pgfusepath{clip}%
\pgfsetbuttcap%
\pgfsetroundjoin%
\pgfsetlinewidth{2.007500pt}%
\definecolor{currentstroke}{rgb}{0.000000,0.000000,0.000000}%
\pgfsetstrokecolor{currentstroke}%
\pgfsetdash{}{0pt}%
\pgfpathmoveto{\pgfqpoint{1.217993in}{2.770113in}}%
\pgfpathlineto{\pgfqpoint{1.520243in}{2.770113in}}%
\pgfusepath{stroke}%
\end{pgfscope}%
\begin{pgfscope}%
\pgfpathrectangle{\pgfqpoint{0.563118in}{0.401080in}}{\pgfqpoint{5.037500in}{4.004000in}}%
\pgfusepath{clip}%
\pgfsetbuttcap%
\pgfsetroundjoin%
\pgfsetlinewidth{2.007500pt}%
\definecolor{currentstroke}{rgb}{0.000000,0.000000,0.000000}%
\pgfsetstrokecolor{currentstroke}%
\pgfsetdash{}{0pt}%
\pgfpathmoveto{\pgfqpoint{1.520243in}{2.002680in}}%
\pgfpathlineto{\pgfqpoint{2.628493in}{2.002680in}}%
\pgfusepath{stroke}%
\end{pgfscope}%
\begin{pgfscope}%
\pgfpathrectangle{\pgfqpoint{0.563118in}{0.401080in}}{\pgfqpoint{5.037500in}{4.004000in}}%
\pgfusepath{clip}%
\pgfsetbuttcap%
\pgfsetroundjoin%
\pgfsetlinewidth{2.007500pt}%
\definecolor{currentstroke}{rgb}{0.000000,0.000000,0.000000}%
\pgfsetstrokecolor{currentstroke}%
\pgfsetdash{}{0pt}%
\pgfpathmoveto{\pgfqpoint{2.628493in}{0.567913in}}%
\pgfpathlineto{\pgfqpoint{3.031493in}{0.567913in}}%
\pgfusepath{stroke}%
\end{pgfscope}%
\begin{pgfscope}%
\pgfpathrectangle{\pgfqpoint{0.563118in}{0.401080in}}{\pgfqpoint{5.037500in}{4.004000in}}%
\pgfusepath{clip}%
\pgfsetbuttcap%
\pgfsetroundjoin%
\pgfsetlinewidth{2.007500pt}%
\definecolor{currentstroke}{rgb}{0.000000,0.000000,0.000000}%
\pgfsetstrokecolor{currentstroke}%
\pgfsetdash{}{0pt}%
\pgfpathmoveto{\pgfqpoint{3.031493in}{2.403080in}}%
\pgfpathlineto{\pgfqpoint{3.182618in}{2.403080in}}%
\pgfusepath{stroke}%
\end{pgfscope}%
\begin{pgfscope}%
\pgfpathrectangle{\pgfqpoint{0.563118in}{0.401080in}}{\pgfqpoint{5.037500in}{4.004000in}}%
\pgfusepath{clip}%
\pgfsetbuttcap%
\pgfsetroundjoin%
\pgfsetlinewidth{2.007500pt}%
\definecolor{currentstroke}{rgb}{0.000000,0.000000,0.000000}%
\pgfsetstrokecolor{currentstroke}%
\pgfsetdash{}{0pt}%
\pgfpathmoveto{\pgfqpoint{3.182618in}{1.602280in}}%
\pgfpathlineto{\pgfqpoint{3.938243in}{1.602280in}}%
\pgfusepath{stroke}%
\end{pgfscope}%
\begin{pgfscope}%
\pgfpathrectangle{\pgfqpoint{0.563118in}{0.401080in}}{\pgfqpoint{5.037500in}{4.004000in}}%
\pgfusepath{clip}%
\pgfsetbuttcap%
\pgfsetroundjoin%
\pgfsetlinewidth{2.007500pt}%
\definecolor{currentstroke}{rgb}{0.000000,0.000000,0.000000}%
\pgfsetstrokecolor{currentstroke}%
\pgfsetdash{}{0pt}%
\pgfpathmoveto{\pgfqpoint{3.938243in}{1.268613in}}%
\pgfpathlineto{\pgfqpoint{4.139743in}{1.268613in}}%
\pgfusepath{stroke}%
\end{pgfscope}%
\begin{pgfscope}%
\pgfpathrectangle{\pgfqpoint{0.563118in}{0.401080in}}{\pgfqpoint{5.037500in}{4.004000in}}%
\pgfusepath{clip}%
\pgfsetbuttcap%
\pgfsetroundjoin%
\pgfsetlinewidth{2.007500pt}%
\definecolor{currentstroke}{rgb}{0.000000,0.000000,0.000000}%
\pgfsetstrokecolor{currentstroke}%
\pgfsetdash{}{0pt}%
\pgfpathmoveto{\pgfqpoint{4.139743in}{0.901580in}}%
\pgfpathlineto{\pgfqpoint{4.391618in}{0.901580in}}%
\pgfusepath{stroke}%
\end{pgfscope}%
\begin{pgfscope}%
\pgfpathrectangle{\pgfqpoint{0.563118in}{0.401080in}}{\pgfqpoint{5.037500in}{4.004000in}}%
\pgfusepath{clip}%
\pgfsetbuttcap%
\pgfsetroundjoin%
\pgfsetlinewidth{2.007500pt}%
\definecolor{currentstroke}{rgb}{0.000000,0.000000,0.000000}%
\pgfsetstrokecolor{currentstroke}%
\pgfsetdash{}{0pt}%
\pgfpathmoveto{\pgfqpoint{4.391618in}{2.302980in}}%
\pgfpathlineto{\pgfqpoint{5.600618in}{2.302980in}}%
\pgfusepath{stroke}%
\end{pgfscope}%
\begin{pgfscope}%
\pgfsetbuttcap%
\pgfsetroundjoin%
\definecolor{currentfill}{rgb}{0.000000,0.000000,0.000000}%
\pgfsetfillcolor{currentfill}%
\pgfsetlinewidth{0.501875pt}%
\definecolor{currentstroke}{rgb}{0.000000,0.000000,0.000000}%
\pgfsetstrokecolor{currentstroke}%
\pgfsetdash{}{0pt}%
\pgfsys@defobject{currentmarker}{\pgfqpoint{0.000000in}{0.000000in}}{\pgfqpoint{0.000000in}{0.041667in}}{%
\pgfpathmoveto{\pgfqpoint{0.000000in}{0.000000in}}%
\pgfpathlineto{\pgfqpoint{0.000000in}{0.041667in}}%
\pgfusepath{stroke,fill}%
}%
\begin{pgfscope}%
\pgfsys@transformshift{0.563118in}{0.401080in}%
\pgfsys@useobject{currentmarker}{}%
\end{pgfscope}%
\end{pgfscope}%
\begin{pgfscope}%
\pgfsetbuttcap%
\pgfsetroundjoin%
\definecolor{currentfill}{rgb}{0.000000,0.000000,0.000000}%
\pgfsetfillcolor{currentfill}%
\pgfsetlinewidth{0.501875pt}%
\definecolor{currentstroke}{rgb}{0.000000,0.000000,0.000000}%
\pgfsetstrokecolor{currentstroke}%
\pgfsetdash{}{0pt}%
\pgfsys@defobject{currentmarker}{\pgfqpoint{0.000000in}{-0.041667in}}{\pgfqpoint{0.000000in}{0.000000in}}{%
\pgfpathmoveto{\pgfqpoint{0.000000in}{0.000000in}}%
\pgfpathlineto{\pgfqpoint{0.000000in}{-0.041667in}}%
\pgfusepath{stroke,fill}%
}%
\begin{pgfscope}%
\pgfsys@transformshift{0.563118in}{4.405080in}%
\pgfsys@useobject{currentmarker}{}%
\end{pgfscope}%
\end{pgfscope}%
\begin{pgfscope}%
\definecolor{textcolor}{rgb}{0.000000,0.000000,0.000000}%
\pgfsetstrokecolor{textcolor}%
\pgfsetfillcolor{textcolor}%
\pgftext[x=0.563118in,y=0.352469in,,top]{\color{textcolor}{\rmfamily\fontsize{10.000000}{12.000000}\selectfont\catcode`\^=\active\def^{\ifmmode\sp\else\^{}\fi}\catcode`\%=\active\def%{\%}$\mathdefault{0.0}$}}%
\end{pgfscope}%
\begin{pgfscope}%
\pgfsetbuttcap%
\pgfsetroundjoin%
\definecolor{currentfill}{rgb}{0.000000,0.000000,0.000000}%
\pgfsetfillcolor{currentfill}%
\pgfsetlinewidth{0.501875pt}%
\definecolor{currentstroke}{rgb}{0.000000,0.000000,0.000000}%
\pgfsetstrokecolor{currentstroke}%
\pgfsetdash{}{0pt}%
\pgfsys@defobject{currentmarker}{\pgfqpoint{0.000000in}{0.000000in}}{\pgfqpoint{0.000000in}{0.041667in}}{%
\pgfpathmoveto{\pgfqpoint{0.000000in}{0.000000in}}%
\pgfpathlineto{\pgfqpoint{0.000000in}{0.041667in}}%
\pgfusepath{stroke,fill}%
}%
\begin{pgfscope}%
\pgfsys@transformshift{1.570618in}{0.401080in}%
\pgfsys@useobject{currentmarker}{}%
\end{pgfscope}%
\end{pgfscope}%
\begin{pgfscope}%
\pgfsetbuttcap%
\pgfsetroundjoin%
\definecolor{currentfill}{rgb}{0.000000,0.000000,0.000000}%
\pgfsetfillcolor{currentfill}%
\pgfsetlinewidth{0.501875pt}%
\definecolor{currentstroke}{rgb}{0.000000,0.000000,0.000000}%
\pgfsetstrokecolor{currentstroke}%
\pgfsetdash{}{0pt}%
\pgfsys@defobject{currentmarker}{\pgfqpoint{0.000000in}{-0.041667in}}{\pgfqpoint{0.000000in}{0.000000in}}{%
\pgfpathmoveto{\pgfqpoint{0.000000in}{0.000000in}}%
\pgfpathlineto{\pgfqpoint{0.000000in}{-0.041667in}}%
\pgfusepath{stroke,fill}%
}%
\begin{pgfscope}%
\pgfsys@transformshift{1.570618in}{4.405080in}%
\pgfsys@useobject{currentmarker}{}%
\end{pgfscope}%
\end{pgfscope}%
\begin{pgfscope}%
\definecolor{textcolor}{rgb}{0.000000,0.000000,0.000000}%
\pgfsetstrokecolor{textcolor}%
\pgfsetfillcolor{textcolor}%
\pgftext[x=1.570618in,y=0.352469in,,top]{\color{textcolor}{\rmfamily\fontsize{10.000000}{12.000000}\selectfont\catcode`\^=\active\def^{\ifmmode\sp\else\^{}\fi}\catcode`\%=\active\def%{\%}$\mathdefault{0.2}$}}%
\end{pgfscope}%
\begin{pgfscope}%
\pgfsetbuttcap%
\pgfsetroundjoin%
\definecolor{currentfill}{rgb}{0.000000,0.000000,0.000000}%
\pgfsetfillcolor{currentfill}%
\pgfsetlinewidth{0.501875pt}%
\definecolor{currentstroke}{rgb}{0.000000,0.000000,0.000000}%
\pgfsetstrokecolor{currentstroke}%
\pgfsetdash{}{0pt}%
\pgfsys@defobject{currentmarker}{\pgfqpoint{0.000000in}{0.000000in}}{\pgfqpoint{0.000000in}{0.041667in}}{%
\pgfpathmoveto{\pgfqpoint{0.000000in}{0.000000in}}%
\pgfpathlineto{\pgfqpoint{0.000000in}{0.041667in}}%
\pgfusepath{stroke,fill}%
}%
\begin{pgfscope}%
\pgfsys@transformshift{2.578118in}{0.401080in}%
\pgfsys@useobject{currentmarker}{}%
\end{pgfscope}%
\end{pgfscope}%
\begin{pgfscope}%
\pgfsetbuttcap%
\pgfsetroundjoin%
\definecolor{currentfill}{rgb}{0.000000,0.000000,0.000000}%
\pgfsetfillcolor{currentfill}%
\pgfsetlinewidth{0.501875pt}%
\definecolor{currentstroke}{rgb}{0.000000,0.000000,0.000000}%
\pgfsetstrokecolor{currentstroke}%
\pgfsetdash{}{0pt}%
\pgfsys@defobject{currentmarker}{\pgfqpoint{0.000000in}{-0.041667in}}{\pgfqpoint{0.000000in}{0.000000in}}{%
\pgfpathmoveto{\pgfqpoint{0.000000in}{0.000000in}}%
\pgfpathlineto{\pgfqpoint{0.000000in}{-0.041667in}}%
\pgfusepath{stroke,fill}%
}%
\begin{pgfscope}%
\pgfsys@transformshift{2.578118in}{4.405080in}%
\pgfsys@useobject{currentmarker}{}%
\end{pgfscope}%
\end{pgfscope}%
\begin{pgfscope}%
\definecolor{textcolor}{rgb}{0.000000,0.000000,0.000000}%
\pgfsetstrokecolor{textcolor}%
\pgfsetfillcolor{textcolor}%
\pgftext[x=2.578118in,y=0.352469in,,top]{\color{textcolor}{\rmfamily\fontsize{10.000000}{12.000000}\selectfont\catcode`\^=\active\def^{\ifmmode\sp\else\^{}\fi}\catcode`\%=\active\def%{\%}$\mathdefault{0.4}$}}%
\end{pgfscope}%
\begin{pgfscope}%
\pgfsetbuttcap%
\pgfsetroundjoin%
\definecolor{currentfill}{rgb}{0.000000,0.000000,0.000000}%
\pgfsetfillcolor{currentfill}%
\pgfsetlinewidth{0.501875pt}%
\definecolor{currentstroke}{rgb}{0.000000,0.000000,0.000000}%
\pgfsetstrokecolor{currentstroke}%
\pgfsetdash{}{0pt}%
\pgfsys@defobject{currentmarker}{\pgfqpoint{0.000000in}{0.000000in}}{\pgfqpoint{0.000000in}{0.041667in}}{%
\pgfpathmoveto{\pgfqpoint{0.000000in}{0.000000in}}%
\pgfpathlineto{\pgfqpoint{0.000000in}{0.041667in}}%
\pgfusepath{stroke,fill}%
}%
\begin{pgfscope}%
\pgfsys@transformshift{3.585618in}{0.401080in}%
\pgfsys@useobject{currentmarker}{}%
\end{pgfscope}%
\end{pgfscope}%
\begin{pgfscope}%
\pgfsetbuttcap%
\pgfsetroundjoin%
\definecolor{currentfill}{rgb}{0.000000,0.000000,0.000000}%
\pgfsetfillcolor{currentfill}%
\pgfsetlinewidth{0.501875pt}%
\definecolor{currentstroke}{rgb}{0.000000,0.000000,0.000000}%
\pgfsetstrokecolor{currentstroke}%
\pgfsetdash{}{0pt}%
\pgfsys@defobject{currentmarker}{\pgfqpoint{0.000000in}{-0.041667in}}{\pgfqpoint{0.000000in}{0.000000in}}{%
\pgfpathmoveto{\pgfqpoint{0.000000in}{0.000000in}}%
\pgfpathlineto{\pgfqpoint{0.000000in}{-0.041667in}}%
\pgfusepath{stroke,fill}%
}%
\begin{pgfscope}%
\pgfsys@transformshift{3.585618in}{4.405080in}%
\pgfsys@useobject{currentmarker}{}%
\end{pgfscope}%
\end{pgfscope}%
\begin{pgfscope}%
\definecolor{textcolor}{rgb}{0.000000,0.000000,0.000000}%
\pgfsetstrokecolor{textcolor}%
\pgfsetfillcolor{textcolor}%
\pgftext[x=3.585618in,y=0.352469in,,top]{\color{textcolor}{\rmfamily\fontsize{10.000000}{12.000000}\selectfont\catcode`\^=\active\def^{\ifmmode\sp\else\^{}\fi}\catcode`\%=\active\def%{\%}$\mathdefault{0.6}$}}%
\end{pgfscope}%
\begin{pgfscope}%
\pgfsetbuttcap%
\pgfsetroundjoin%
\definecolor{currentfill}{rgb}{0.000000,0.000000,0.000000}%
\pgfsetfillcolor{currentfill}%
\pgfsetlinewidth{0.501875pt}%
\definecolor{currentstroke}{rgb}{0.000000,0.000000,0.000000}%
\pgfsetstrokecolor{currentstroke}%
\pgfsetdash{}{0pt}%
\pgfsys@defobject{currentmarker}{\pgfqpoint{0.000000in}{0.000000in}}{\pgfqpoint{0.000000in}{0.041667in}}{%
\pgfpathmoveto{\pgfqpoint{0.000000in}{0.000000in}}%
\pgfpathlineto{\pgfqpoint{0.000000in}{0.041667in}}%
\pgfusepath{stroke,fill}%
}%
\begin{pgfscope}%
\pgfsys@transformshift{4.593118in}{0.401080in}%
\pgfsys@useobject{currentmarker}{}%
\end{pgfscope}%
\end{pgfscope}%
\begin{pgfscope}%
\pgfsetbuttcap%
\pgfsetroundjoin%
\definecolor{currentfill}{rgb}{0.000000,0.000000,0.000000}%
\pgfsetfillcolor{currentfill}%
\pgfsetlinewidth{0.501875pt}%
\definecolor{currentstroke}{rgb}{0.000000,0.000000,0.000000}%
\pgfsetstrokecolor{currentstroke}%
\pgfsetdash{}{0pt}%
\pgfsys@defobject{currentmarker}{\pgfqpoint{0.000000in}{-0.041667in}}{\pgfqpoint{0.000000in}{0.000000in}}{%
\pgfpathmoveto{\pgfqpoint{0.000000in}{0.000000in}}%
\pgfpathlineto{\pgfqpoint{0.000000in}{-0.041667in}}%
\pgfusepath{stroke,fill}%
}%
\begin{pgfscope}%
\pgfsys@transformshift{4.593118in}{4.405080in}%
\pgfsys@useobject{currentmarker}{}%
\end{pgfscope}%
\end{pgfscope}%
\begin{pgfscope}%
\definecolor{textcolor}{rgb}{0.000000,0.000000,0.000000}%
\pgfsetstrokecolor{textcolor}%
\pgfsetfillcolor{textcolor}%
\pgftext[x=4.593118in,y=0.352469in,,top]{\color{textcolor}{\rmfamily\fontsize{10.000000}{12.000000}\selectfont\catcode`\^=\active\def^{\ifmmode\sp\else\^{}\fi}\catcode`\%=\active\def%{\%}$\mathdefault{0.8}$}}%
\end{pgfscope}%
\begin{pgfscope}%
\pgfsetbuttcap%
\pgfsetroundjoin%
\definecolor{currentfill}{rgb}{0.000000,0.000000,0.000000}%
\pgfsetfillcolor{currentfill}%
\pgfsetlinewidth{0.501875pt}%
\definecolor{currentstroke}{rgb}{0.000000,0.000000,0.000000}%
\pgfsetstrokecolor{currentstroke}%
\pgfsetdash{}{0pt}%
\pgfsys@defobject{currentmarker}{\pgfqpoint{0.000000in}{0.000000in}}{\pgfqpoint{0.000000in}{0.041667in}}{%
\pgfpathmoveto{\pgfqpoint{0.000000in}{0.000000in}}%
\pgfpathlineto{\pgfqpoint{0.000000in}{0.041667in}}%
\pgfusepath{stroke,fill}%
}%
\begin{pgfscope}%
\pgfsys@transformshift{5.600618in}{0.401080in}%
\pgfsys@useobject{currentmarker}{}%
\end{pgfscope}%
\end{pgfscope}%
\begin{pgfscope}%
\pgfsetbuttcap%
\pgfsetroundjoin%
\definecolor{currentfill}{rgb}{0.000000,0.000000,0.000000}%
\pgfsetfillcolor{currentfill}%
\pgfsetlinewidth{0.501875pt}%
\definecolor{currentstroke}{rgb}{0.000000,0.000000,0.000000}%
\pgfsetstrokecolor{currentstroke}%
\pgfsetdash{}{0pt}%
\pgfsys@defobject{currentmarker}{\pgfqpoint{0.000000in}{-0.041667in}}{\pgfqpoint{0.000000in}{0.000000in}}{%
\pgfpathmoveto{\pgfqpoint{0.000000in}{0.000000in}}%
\pgfpathlineto{\pgfqpoint{0.000000in}{-0.041667in}}%
\pgfusepath{stroke,fill}%
}%
\begin{pgfscope}%
\pgfsys@transformshift{5.600618in}{4.405080in}%
\pgfsys@useobject{currentmarker}{}%
\end{pgfscope}%
\end{pgfscope}%
\begin{pgfscope}%
\definecolor{textcolor}{rgb}{0.000000,0.000000,0.000000}%
\pgfsetstrokecolor{textcolor}%
\pgfsetfillcolor{textcolor}%
\pgftext[x=5.600618in,y=0.352469in,,top]{\color{textcolor}{\rmfamily\fontsize{10.000000}{12.000000}\selectfont\catcode`\^=\active\def^{\ifmmode\sp\else\^{}\fi}\catcode`\%=\active\def%{\%}$\mathdefault{1.0}$}}%
\end{pgfscope}%
\begin{pgfscope}%
\pgfsetbuttcap%
\pgfsetroundjoin%
\definecolor{currentfill}{rgb}{0.000000,0.000000,0.000000}%
\pgfsetfillcolor{currentfill}%
\pgfsetlinewidth{0.501875pt}%
\definecolor{currentstroke}{rgb}{0.000000,0.000000,0.000000}%
\pgfsetstrokecolor{currentstroke}%
\pgfsetdash{}{0pt}%
\pgfsys@defobject{currentmarker}{\pgfqpoint{0.000000in}{0.000000in}}{\pgfqpoint{0.000000in}{0.020833in}}{%
\pgfpathmoveto{\pgfqpoint{0.000000in}{0.000000in}}%
\pgfpathlineto{\pgfqpoint{0.000000in}{0.020833in}}%
\pgfusepath{stroke,fill}%
}%
\begin{pgfscope}%
\pgfsys@transformshift{0.814993in}{0.401080in}%
\pgfsys@useobject{currentmarker}{}%
\end{pgfscope}%
\end{pgfscope}%
\begin{pgfscope}%
\pgfsetbuttcap%
\pgfsetroundjoin%
\definecolor{currentfill}{rgb}{0.000000,0.000000,0.000000}%
\pgfsetfillcolor{currentfill}%
\pgfsetlinewidth{0.501875pt}%
\definecolor{currentstroke}{rgb}{0.000000,0.000000,0.000000}%
\pgfsetstrokecolor{currentstroke}%
\pgfsetdash{}{0pt}%
\pgfsys@defobject{currentmarker}{\pgfqpoint{0.000000in}{-0.020833in}}{\pgfqpoint{0.000000in}{0.000000in}}{%
\pgfpathmoveto{\pgfqpoint{0.000000in}{0.000000in}}%
\pgfpathlineto{\pgfqpoint{0.000000in}{-0.020833in}}%
\pgfusepath{stroke,fill}%
}%
\begin{pgfscope}%
\pgfsys@transformshift{0.814993in}{4.405080in}%
\pgfsys@useobject{currentmarker}{}%
\end{pgfscope}%
\end{pgfscope}%
\begin{pgfscope}%
\pgfsetbuttcap%
\pgfsetroundjoin%
\definecolor{currentfill}{rgb}{0.000000,0.000000,0.000000}%
\pgfsetfillcolor{currentfill}%
\pgfsetlinewidth{0.501875pt}%
\definecolor{currentstroke}{rgb}{0.000000,0.000000,0.000000}%
\pgfsetstrokecolor{currentstroke}%
\pgfsetdash{}{0pt}%
\pgfsys@defobject{currentmarker}{\pgfqpoint{0.000000in}{0.000000in}}{\pgfqpoint{0.000000in}{0.020833in}}{%
\pgfpathmoveto{\pgfqpoint{0.000000in}{0.000000in}}%
\pgfpathlineto{\pgfqpoint{0.000000in}{0.020833in}}%
\pgfusepath{stroke,fill}%
}%
\begin{pgfscope}%
\pgfsys@transformshift{1.066868in}{0.401080in}%
\pgfsys@useobject{currentmarker}{}%
\end{pgfscope}%
\end{pgfscope}%
\begin{pgfscope}%
\pgfsetbuttcap%
\pgfsetroundjoin%
\definecolor{currentfill}{rgb}{0.000000,0.000000,0.000000}%
\pgfsetfillcolor{currentfill}%
\pgfsetlinewidth{0.501875pt}%
\definecolor{currentstroke}{rgb}{0.000000,0.000000,0.000000}%
\pgfsetstrokecolor{currentstroke}%
\pgfsetdash{}{0pt}%
\pgfsys@defobject{currentmarker}{\pgfqpoint{0.000000in}{-0.020833in}}{\pgfqpoint{0.000000in}{0.000000in}}{%
\pgfpathmoveto{\pgfqpoint{0.000000in}{0.000000in}}%
\pgfpathlineto{\pgfqpoint{0.000000in}{-0.020833in}}%
\pgfusepath{stroke,fill}%
}%
\begin{pgfscope}%
\pgfsys@transformshift{1.066868in}{4.405080in}%
\pgfsys@useobject{currentmarker}{}%
\end{pgfscope}%
\end{pgfscope}%
\begin{pgfscope}%
\pgfsetbuttcap%
\pgfsetroundjoin%
\definecolor{currentfill}{rgb}{0.000000,0.000000,0.000000}%
\pgfsetfillcolor{currentfill}%
\pgfsetlinewidth{0.501875pt}%
\definecolor{currentstroke}{rgb}{0.000000,0.000000,0.000000}%
\pgfsetstrokecolor{currentstroke}%
\pgfsetdash{}{0pt}%
\pgfsys@defobject{currentmarker}{\pgfqpoint{0.000000in}{0.000000in}}{\pgfqpoint{0.000000in}{0.020833in}}{%
\pgfpathmoveto{\pgfqpoint{0.000000in}{0.000000in}}%
\pgfpathlineto{\pgfqpoint{0.000000in}{0.020833in}}%
\pgfusepath{stroke,fill}%
}%
\begin{pgfscope}%
\pgfsys@transformshift{1.318743in}{0.401080in}%
\pgfsys@useobject{currentmarker}{}%
\end{pgfscope}%
\end{pgfscope}%
\begin{pgfscope}%
\pgfsetbuttcap%
\pgfsetroundjoin%
\definecolor{currentfill}{rgb}{0.000000,0.000000,0.000000}%
\pgfsetfillcolor{currentfill}%
\pgfsetlinewidth{0.501875pt}%
\definecolor{currentstroke}{rgb}{0.000000,0.000000,0.000000}%
\pgfsetstrokecolor{currentstroke}%
\pgfsetdash{}{0pt}%
\pgfsys@defobject{currentmarker}{\pgfqpoint{0.000000in}{-0.020833in}}{\pgfqpoint{0.000000in}{0.000000in}}{%
\pgfpathmoveto{\pgfqpoint{0.000000in}{0.000000in}}%
\pgfpathlineto{\pgfqpoint{0.000000in}{-0.020833in}}%
\pgfusepath{stroke,fill}%
}%
\begin{pgfscope}%
\pgfsys@transformshift{1.318743in}{4.405080in}%
\pgfsys@useobject{currentmarker}{}%
\end{pgfscope}%
\end{pgfscope}%
\begin{pgfscope}%
\pgfsetbuttcap%
\pgfsetroundjoin%
\definecolor{currentfill}{rgb}{0.000000,0.000000,0.000000}%
\pgfsetfillcolor{currentfill}%
\pgfsetlinewidth{0.501875pt}%
\definecolor{currentstroke}{rgb}{0.000000,0.000000,0.000000}%
\pgfsetstrokecolor{currentstroke}%
\pgfsetdash{}{0pt}%
\pgfsys@defobject{currentmarker}{\pgfqpoint{0.000000in}{0.000000in}}{\pgfqpoint{0.000000in}{0.020833in}}{%
\pgfpathmoveto{\pgfqpoint{0.000000in}{0.000000in}}%
\pgfpathlineto{\pgfqpoint{0.000000in}{0.020833in}}%
\pgfusepath{stroke,fill}%
}%
\begin{pgfscope}%
\pgfsys@transformshift{1.822493in}{0.401080in}%
\pgfsys@useobject{currentmarker}{}%
\end{pgfscope}%
\end{pgfscope}%
\begin{pgfscope}%
\pgfsetbuttcap%
\pgfsetroundjoin%
\definecolor{currentfill}{rgb}{0.000000,0.000000,0.000000}%
\pgfsetfillcolor{currentfill}%
\pgfsetlinewidth{0.501875pt}%
\definecolor{currentstroke}{rgb}{0.000000,0.000000,0.000000}%
\pgfsetstrokecolor{currentstroke}%
\pgfsetdash{}{0pt}%
\pgfsys@defobject{currentmarker}{\pgfqpoint{0.000000in}{-0.020833in}}{\pgfqpoint{0.000000in}{0.000000in}}{%
\pgfpathmoveto{\pgfqpoint{0.000000in}{0.000000in}}%
\pgfpathlineto{\pgfqpoint{0.000000in}{-0.020833in}}%
\pgfusepath{stroke,fill}%
}%
\begin{pgfscope}%
\pgfsys@transformshift{1.822493in}{4.405080in}%
\pgfsys@useobject{currentmarker}{}%
\end{pgfscope}%
\end{pgfscope}%
\begin{pgfscope}%
\pgfsetbuttcap%
\pgfsetroundjoin%
\definecolor{currentfill}{rgb}{0.000000,0.000000,0.000000}%
\pgfsetfillcolor{currentfill}%
\pgfsetlinewidth{0.501875pt}%
\definecolor{currentstroke}{rgb}{0.000000,0.000000,0.000000}%
\pgfsetstrokecolor{currentstroke}%
\pgfsetdash{}{0pt}%
\pgfsys@defobject{currentmarker}{\pgfqpoint{0.000000in}{0.000000in}}{\pgfqpoint{0.000000in}{0.020833in}}{%
\pgfpathmoveto{\pgfqpoint{0.000000in}{0.000000in}}%
\pgfpathlineto{\pgfqpoint{0.000000in}{0.020833in}}%
\pgfusepath{stroke,fill}%
}%
\begin{pgfscope}%
\pgfsys@transformshift{2.074368in}{0.401080in}%
\pgfsys@useobject{currentmarker}{}%
\end{pgfscope}%
\end{pgfscope}%
\begin{pgfscope}%
\pgfsetbuttcap%
\pgfsetroundjoin%
\definecolor{currentfill}{rgb}{0.000000,0.000000,0.000000}%
\pgfsetfillcolor{currentfill}%
\pgfsetlinewidth{0.501875pt}%
\definecolor{currentstroke}{rgb}{0.000000,0.000000,0.000000}%
\pgfsetstrokecolor{currentstroke}%
\pgfsetdash{}{0pt}%
\pgfsys@defobject{currentmarker}{\pgfqpoint{0.000000in}{-0.020833in}}{\pgfqpoint{0.000000in}{0.000000in}}{%
\pgfpathmoveto{\pgfqpoint{0.000000in}{0.000000in}}%
\pgfpathlineto{\pgfqpoint{0.000000in}{-0.020833in}}%
\pgfusepath{stroke,fill}%
}%
\begin{pgfscope}%
\pgfsys@transformshift{2.074368in}{4.405080in}%
\pgfsys@useobject{currentmarker}{}%
\end{pgfscope}%
\end{pgfscope}%
\begin{pgfscope}%
\pgfsetbuttcap%
\pgfsetroundjoin%
\definecolor{currentfill}{rgb}{0.000000,0.000000,0.000000}%
\pgfsetfillcolor{currentfill}%
\pgfsetlinewidth{0.501875pt}%
\definecolor{currentstroke}{rgb}{0.000000,0.000000,0.000000}%
\pgfsetstrokecolor{currentstroke}%
\pgfsetdash{}{0pt}%
\pgfsys@defobject{currentmarker}{\pgfqpoint{0.000000in}{0.000000in}}{\pgfqpoint{0.000000in}{0.020833in}}{%
\pgfpathmoveto{\pgfqpoint{0.000000in}{0.000000in}}%
\pgfpathlineto{\pgfqpoint{0.000000in}{0.020833in}}%
\pgfusepath{stroke,fill}%
}%
\begin{pgfscope}%
\pgfsys@transformshift{2.326243in}{0.401080in}%
\pgfsys@useobject{currentmarker}{}%
\end{pgfscope}%
\end{pgfscope}%
\begin{pgfscope}%
\pgfsetbuttcap%
\pgfsetroundjoin%
\definecolor{currentfill}{rgb}{0.000000,0.000000,0.000000}%
\pgfsetfillcolor{currentfill}%
\pgfsetlinewidth{0.501875pt}%
\definecolor{currentstroke}{rgb}{0.000000,0.000000,0.000000}%
\pgfsetstrokecolor{currentstroke}%
\pgfsetdash{}{0pt}%
\pgfsys@defobject{currentmarker}{\pgfqpoint{0.000000in}{-0.020833in}}{\pgfqpoint{0.000000in}{0.000000in}}{%
\pgfpathmoveto{\pgfqpoint{0.000000in}{0.000000in}}%
\pgfpathlineto{\pgfqpoint{0.000000in}{-0.020833in}}%
\pgfusepath{stroke,fill}%
}%
\begin{pgfscope}%
\pgfsys@transformshift{2.326243in}{4.405080in}%
\pgfsys@useobject{currentmarker}{}%
\end{pgfscope}%
\end{pgfscope}%
\begin{pgfscope}%
\pgfsetbuttcap%
\pgfsetroundjoin%
\definecolor{currentfill}{rgb}{0.000000,0.000000,0.000000}%
\pgfsetfillcolor{currentfill}%
\pgfsetlinewidth{0.501875pt}%
\definecolor{currentstroke}{rgb}{0.000000,0.000000,0.000000}%
\pgfsetstrokecolor{currentstroke}%
\pgfsetdash{}{0pt}%
\pgfsys@defobject{currentmarker}{\pgfqpoint{0.000000in}{0.000000in}}{\pgfqpoint{0.000000in}{0.020833in}}{%
\pgfpathmoveto{\pgfqpoint{0.000000in}{0.000000in}}%
\pgfpathlineto{\pgfqpoint{0.000000in}{0.020833in}}%
\pgfusepath{stroke,fill}%
}%
\begin{pgfscope}%
\pgfsys@transformshift{2.829993in}{0.401080in}%
\pgfsys@useobject{currentmarker}{}%
\end{pgfscope}%
\end{pgfscope}%
\begin{pgfscope}%
\pgfsetbuttcap%
\pgfsetroundjoin%
\definecolor{currentfill}{rgb}{0.000000,0.000000,0.000000}%
\pgfsetfillcolor{currentfill}%
\pgfsetlinewidth{0.501875pt}%
\definecolor{currentstroke}{rgb}{0.000000,0.000000,0.000000}%
\pgfsetstrokecolor{currentstroke}%
\pgfsetdash{}{0pt}%
\pgfsys@defobject{currentmarker}{\pgfqpoint{0.000000in}{-0.020833in}}{\pgfqpoint{0.000000in}{0.000000in}}{%
\pgfpathmoveto{\pgfqpoint{0.000000in}{0.000000in}}%
\pgfpathlineto{\pgfqpoint{0.000000in}{-0.020833in}}%
\pgfusepath{stroke,fill}%
}%
\begin{pgfscope}%
\pgfsys@transformshift{2.829993in}{4.405080in}%
\pgfsys@useobject{currentmarker}{}%
\end{pgfscope}%
\end{pgfscope}%
\begin{pgfscope}%
\pgfsetbuttcap%
\pgfsetroundjoin%
\definecolor{currentfill}{rgb}{0.000000,0.000000,0.000000}%
\pgfsetfillcolor{currentfill}%
\pgfsetlinewidth{0.501875pt}%
\definecolor{currentstroke}{rgb}{0.000000,0.000000,0.000000}%
\pgfsetstrokecolor{currentstroke}%
\pgfsetdash{}{0pt}%
\pgfsys@defobject{currentmarker}{\pgfqpoint{0.000000in}{0.000000in}}{\pgfqpoint{0.000000in}{0.020833in}}{%
\pgfpathmoveto{\pgfqpoint{0.000000in}{0.000000in}}%
\pgfpathlineto{\pgfqpoint{0.000000in}{0.020833in}}%
\pgfusepath{stroke,fill}%
}%
\begin{pgfscope}%
\pgfsys@transformshift{3.081868in}{0.401080in}%
\pgfsys@useobject{currentmarker}{}%
\end{pgfscope}%
\end{pgfscope}%
\begin{pgfscope}%
\pgfsetbuttcap%
\pgfsetroundjoin%
\definecolor{currentfill}{rgb}{0.000000,0.000000,0.000000}%
\pgfsetfillcolor{currentfill}%
\pgfsetlinewidth{0.501875pt}%
\definecolor{currentstroke}{rgb}{0.000000,0.000000,0.000000}%
\pgfsetstrokecolor{currentstroke}%
\pgfsetdash{}{0pt}%
\pgfsys@defobject{currentmarker}{\pgfqpoint{0.000000in}{-0.020833in}}{\pgfqpoint{0.000000in}{0.000000in}}{%
\pgfpathmoveto{\pgfqpoint{0.000000in}{0.000000in}}%
\pgfpathlineto{\pgfqpoint{0.000000in}{-0.020833in}}%
\pgfusepath{stroke,fill}%
}%
\begin{pgfscope}%
\pgfsys@transformshift{3.081868in}{4.405080in}%
\pgfsys@useobject{currentmarker}{}%
\end{pgfscope}%
\end{pgfscope}%
\begin{pgfscope}%
\pgfsetbuttcap%
\pgfsetroundjoin%
\definecolor{currentfill}{rgb}{0.000000,0.000000,0.000000}%
\pgfsetfillcolor{currentfill}%
\pgfsetlinewidth{0.501875pt}%
\definecolor{currentstroke}{rgb}{0.000000,0.000000,0.000000}%
\pgfsetstrokecolor{currentstroke}%
\pgfsetdash{}{0pt}%
\pgfsys@defobject{currentmarker}{\pgfqpoint{0.000000in}{0.000000in}}{\pgfqpoint{0.000000in}{0.020833in}}{%
\pgfpathmoveto{\pgfqpoint{0.000000in}{0.000000in}}%
\pgfpathlineto{\pgfqpoint{0.000000in}{0.020833in}}%
\pgfusepath{stroke,fill}%
}%
\begin{pgfscope}%
\pgfsys@transformshift{3.333743in}{0.401080in}%
\pgfsys@useobject{currentmarker}{}%
\end{pgfscope}%
\end{pgfscope}%
\begin{pgfscope}%
\pgfsetbuttcap%
\pgfsetroundjoin%
\definecolor{currentfill}{rgb}{0.000000,0.000000,0.000000}%
\pgfsetfillcolor{currentfill}%
\pgfsetlinewidth{0.501875pt}%
\definecolor{currentstroke}{rgb}{0.000000,0.000000,0.000000}%
\pgfsetstrokecolor{currentstroke}%
\pgfsetdash{}{0pt}%
\pgfsys@defobject{currentmarker}{\pgfqpoint{0.000000in}{-0.020833in}}{\pgfqpoint{0.000000in}{0.000000in}}{%
\pgfpathmoveto{\pgfqpoint{0.000000in}{0.000000in}}%
\pgfpathlineto{\pgfqpoint{0.000000in}{-0.020833in}}%
\pgfusepath{stroke,fill}%
}%
\begin{pgfscope}%
\pgfsys@transformshift{3.333743in}{4.405080in}%
\pgfsys@useobject{currentmarker}{}%
\end{pgfscope}%
\end{pgfscope}%
\begin{pgfscope}%
\pgfsetbuttcap%
\pgfsetroundjoin%
\definecolor{currentfill}{rgb}{0.000000,0.000000,0.000000}%
\pgfsetfillcolor{currentfill}%
\pgfsetlinewidth{0.501875pt}%
\definecolor{currentstroke}{rgb}{0.000000,0.000000,0.000000}%
\pgfsetstrokecolor{currentstroke}%
\pgfsetdash{}{0pt}%
\pgfsys@defobject{currentmarker}{\pgfqpoint{0.000000in}{0.000000in}}{\pgfqpoint{0.000000in}{0.020833in}}{%
\pgfpathmoveto{\pgfqpoint{0.000000in}{0.000000in}}%
\pgfpathlineto{\pgfqpoint{0.000000in}{0.020833in}}%
\pgfusepath{stroke,fill}%
}%
\begin{pgfscope}%
\pgfsys@transformshift{3.837493in}{0.401080in}%
\pgfsys@useobject{currentmarker}{}%
\end{pgfscope}%
\end{pgfscope}%
\begin{pgfscope}%
\pgfsetbuttcap%
\pgfsetroundjoin%
\definecolor{currentfill}{rgb}{0.000000,0.000000,0.000000}%
\pgfsetfillcolor{currentfill}%
\pgfsetlinewidth{0.501875pt}%
\definecolor{currentstroke}{rgb}{0.000000,0.000000,0.000000}%
\pgfsetstrokecolor{currentstroke}%
\pgfsetdash{}{0pt}%
\pgfsys@defobject{currentmarker}{\pgfqpoint{0.000000in}{-0.020833in}}{\pgfqpoint{0.000000in}{0.000000in}}{%
\pgfpathmoveto{\pgfqpoint{0.000000in}{0.000000in}}%
\pgfpathlineto{\pgfqpoint{0.000000in}{-0.020833in}}%
\pgfusepath{stroke,fill}%
}%
\begin{pgfscope}%
\pgfsys@transformshift{3.837493in}{4.405080in}%
\pgfsys@useobject{currentmarker}{}%
\end{pgfscope}%
\end{pgfscope}%
\begin{pgfscope}%
\pgfsetbuttcap%
\pgfsetroundjoin%
\definecolor{currentfill}{rgb}{0.000000,0.000000,0.000000}%
\pgfsetfillcolor{currentfill}%
\pgfsetlinewidth{0.501875pt}%
\definecolor{currentstroke}{rgb}{0.000000,0.000000,0.000000}%
\pgfsetstrokecolor{currentstroke}%
\pgfsetdash{}{0pt}%
\pgfsys@defobject{currentmarker}{\pgfqpoint{0.000000in}{0.000000in}}{\pgfqpoint{0.000000in}{0.020833in}}{%
\pgfpathmoveto{\pgfqpoint{0.000000in}{0.000000in}}%
\pgfpathlineto{\pgfqpoint{0.000000in}{0.020833in}}%
\pgfusepath{stroke,fill}%
}%
\begin{pgfscope}%
\pgfsys@transformshift{4.089368in}{0.401080in}%
\pgfsys@useobject{currentmarker}{}%
\end{pgfscope}%
\end{pgfscope}%
\begin{pgfscope}%
\pgfsetbuttcap%
\pgfsetroundjoin%
\definecolor{currentfill}{rgb}{0.000000,0.000000,0.000000}%
\pgfsetfillcolor{currentfill}%
\pgfsetlinewidth{0.501875pt}%
\definecolor{currentstroke}{rgb}{0.000000,0.000000,0.000000}%
\pgfsetstrokecolor{currentstroke}%
\pgfsetdash{}{0pt}%
\pgfsys@defobject{currentmarker}{\pgfqpoint{0.000000in}{-0.020833in}}{\pgfqpoint{0.000000in}{0.000000in}}{%
\pgfpathmoveto{\pgfqpoint{0.000000in}{0.000000in}}%
\pgfpathlineto{\pgfqpoint{0.000000in}{-0.020833in}}%
\pgfusepath{stroke,fill}%
}%
\begin{pgfscope}%
\pgfsys@transformshift{4.089368in}{4.405080in}%
\pgfsys@useobject{currentmarker}{}%
\end{pgfscope}%
\end{pgfscope}%
\begin{pgfscope}%
\pgfsetbuttcap%
\pgfsetroundjoin%
\definecolor{currentfill}{rgb}{0.000000,0.000000,0.000000}%
\pgfsetfillcolor{currentfill}%
\pgfsetlinewidth{0.501875pt}%
\definecolor{currentstroke}{rgb}{0.000000,0.000000,0.000000}%
\pgfsetstrokecolor{currentstroke}%
\pgfsetdash{}{0pt}%
\pgfsys@defobject{currentmarker}{\pgfqpoint{0.000000in}{0.000000in}}{\pgfqpoint{0.000000in}{0.020833in}}{%
\pgfpathmoveto{\pgfqpoint{0.000000in}{0.000000in}}%
\pgfpathlineto{\pgfqpoint{0.000000in}{0.020833in}}%
\pgfusepath{stroke,fill}%
}%
\begin{pgfscope}%
\pgfsys@transformshift{4.341243in}{0.401080in}%
\pgfsys@useobject{currentmarker}{}%
\end{pgfscope}%
\end{pgfscope}%
\begin{pgfscope}%
\pgfsetbuttcap%
\pgfsetroundjoin%
\definecolor{currentfill}{rgb}{0.000000,0.000000,0.000000}%
\pgfsetfillcolor{currentfill}%
\pgfsetlinewidth{0.501875pt}%
\definecolor{currentstroke}{rgb}{0.000000,0.000000,0.000000}%
\pgfsetstrokecolor{currentstroke}%
\pgfsetdash{}{0pt}%
\pgfsys@defobject{currentmarker}{\pgfqpoint{0.000000in}{-0.020833in}}{\pgfqpoint{0.000000in}{0.000000in}}{%
\pgfpathmoveto{\pgfqpoint{0.000000in}{0.000000in}}%
\pgfpathlineto{\pgfqpoint{0.000000in}{-0.020833in}}%
\pgfusepath{stroke,fill}%
}%
\begin{pgfscope}%
\pgfsys@transformshift{4.341243in}{4.405080in}%
\pgfsys@useobject{currentmarker}{}%
\end{pgfscope}%
\end{pgfscope}%
\begin{pgfscope}%
\pgfsetbuttcap%
\pgfsetroundjoin%
\definecolor{currentfill}{rgb}{0.000000,0.000000,0.000000}%
\pgfsetfillcolor{currentfill}%
\pgfsetlinewidth{0.501875pt}%
\definecolor{currentstroke}{rgb}{0.000000,0.000000,0.000000}%
\pgfsetstrokecolor{currentstroke}%
\pgfsetdash{}{0pt}%
\pgfsys@defobject{currentmarker}{\pgfqpoint{0.000000in}{0.000000in}}{\pgfqpoint{0.000000in}{0.020833in}}{%
\pgfpathmoveto{\pgfqpoint{0.000000in}{0.000000in}}%
\pgfpathlineto{\pgfqpoint{0.000000in}{0.020833in}}%
\pgfusepath{stroke,fill}%
}%
\begin{pgfscope}%
\pgfsys@transformshift{4.844993in}{0.401080in}%
\pgfsys@useobject{currentmarker}{}%
\end{pgfscope}%
\end{pgfscope}%
\begin{pgfscope}%
\pgfsetbuttcap%
\pgfsetroundjoin%
\definecolor{currentfill}{rgb}{0.000000,0.000000,0.000000}%
\pgfsetfillcolor{currentfill}%
\pgfsetlinewidth{0.501875pt}%
\definecolor{currentstroke}{rgb}{0.000000,0.000000,0.000000}%
\pgfsetstrokecolor{currentstroke}%
\pgfsetdash{}{0pt}%
\pgfsys@defobject{currentmarker}{\pgfqpoint{0.000000in}{-0.020833in}}{\pgfqpoint{0.000000in}{0.000000in}}{%
\pgfpathmoveto{\pgfqpoint{0.000000in}{0.000000in}}%
\pgfpathlineto{\pgfqpoint{0.000000in}{-0.020833in}}%
\pgfusepath{stroke,fill}%
}%
\begin{pgfscope}%
\pgfsys@transformshift{4.844993in}{4.405080in}%
\pgfsys@useobject{currentmarker}{}%
\end{pgfscope}%
\end{pgfscope}%
\begin{pgfscope}%
\pgfsetbuttcap%
\pgfsetroundjoin%
\definecolor{currentfill}{rgb}{0.000000,0.000000,0.000000}%
\pgfsetfillcolor{currentfill}%
\pgfsetlinewidth{0.501875pt}%
\definecolor{currentstroke}{rgb}{0.000000,0.000000,0.000000}%
\pgfsetstrokecolor{currentstroke}%
\pgfsetdash{}{0pt}%
\pgfsys@defobject{currentmarker}{\pgfqpoint{0.000000in}{0.000000in}}{\pgfqpoint{0.000000in}{0.020833in}}{%
\pgfpathmoveto{\pgfqpoint{0.000000in}{0.000000in}}%
\pgfpathlineto{\pgfqpoint{0.000000in}{0.020833in}}%
\pgfusepath{stroke,fill}%
}%
\begin{pgfscope}%
\pgfsys@transformshift{5.096868in}{0.401080in}%
\pgfsys@useobject{currentmarker}{}%
\end{pgfscope}%
\end{pgfscope}%
\begin{pgfscope}%
\pgfsetbuttcap%
\pgfsetroundjoin%
\definecolor{currentfill}{rgb}{0.000000,0.000000,0.000000}%
\pgfsetfillcolor{currentfill}%
\pgfsetlinewidth{0.501875pt}%
\definecolor{currentstroke}{rgb}{0.000000,0.000000,0.000000}%
\pgfsetstrokecolor{currentstroke}%
\pgfsetdash{}{0pt}%
\pgfsys@defobject{currentmarker}{\pgfqpoint{0.000000in}{-0.020833in}}{\pgfqpoint{0.000000in}{0.000000in}}{%
\pgfpathmoveto{\pgfqpoint{0.000000in}{0.000000in}}%
\pgfpathlineto{\pgfqpoint{0.000000in}{-0.020833in}}%
\pgfusepath{stroke,fill}%
}%
\begin{pgfscope}%
\pgfsys@transformshift{5.096868in}{4.405080in}%
\pgfsys@useobject{currentmarker}{}%
\end{pgfscope}%
\end{pgfscope}%
\begin{pgfscope}%
\pgfsetbuttcap%
\pgfsetroundjoin%
\definecolor{currentfill}{rgb}{0.000000,0.000000,0.000000}%
\pgfsetfillcolor{currentfill}%
\pgfsetlinewidth{0.501875pt}%
\definecolor{currentstroke}{rgb}{0.000000,0.000000,0.000000}%
\pgfsetstrokecolor{currentstroke}%
\pgfsetdash{}{0pt}%
\pgfsys@defobject{currentmarker}{\pgfqpoint{0.000000in}{0.000000in}}{\pgfqpoint{0.000000in}{0.020833in}}{%
\pgfpathmoveto{\pgfqpoint{0.000000in}{0.000000in}}%
\pgfpathlineto{\pgfqpoint{0.000000in}{0.020833in}}%
\pgfusepath{stroke,fill}%
}%
\begin{pgfscope}%
\pgfsys@transformshift{5.348743in}{0.401080in}%
\pgfsys@useobject{currentmarker}{}%
\end{pgfscope}%
\end{pgfscope}%
\begin{pgfscope}%
\pgfsetbuttcap%
\pgfsetroundjoin%
\definecolor{currentfill}{rgb}{0.000000,0.000000,0.000000}%
\pgfsetfillcolor{currentfill}%
\pgfsetlinewidth{0.501875pt}%
\definecolor{currentstroke}{rgb}{0.000000,0.000000,0.000000}%
\pgfsetstrokecolor{currentstroke}%
\pgfsetdash{}{0pt}%
\pgfsys@defobject{currentmarker}{\pgfqpoint{0.000000in}{-0.020833in}}{\pgfqpoint{0.000000in}{0.000000in}}{%
\pgfpathmoveto{\pgfqpoint{0.000000in}{0.000000in}}%
\pgfpathlineto{\pgfqpoint{0.000000in}{-0.020833in}}%
\pgfusepath{stroke,fill}%
}%
\begin{pgfscope}%
\pgfsys@transformshift{5.348743in}{4.405080in}%
\pgfsys@useobject{currentmarker}{}%
\end{pgfscope}%
\end{pgfscope}%
\begin{pgfscope}%
\definecolor{textcolor}{rgb}{0.000000,0.000000,0.000000}%
\pgfsetstrokecolor{textcolor}%
\pgfsetfillcolor{textcolor}%
\pgftext[x=3.081868in,y=0.173457in,,top]{\color{textcolor}{\rmfamily\fontsize{10.000000}{12.000000}\selectfont\catcode`\^=\active\def^{\ifmmode\sp\else\^{}\fi}\catcode`\%=\active\def%{\%}$s$}}%
\end{pgfscope}%
\begin{pgfscope}%
\pgfsetbuttcap%
\pgfsetroundjoin%
\definecolor{currentfill}{rgb}{0.000000,0.000000,0.000000}%
\pgfsetfillcolor{currentfill}%
\pgfsetlinewidth{0.501875pt}%
\definecolor{currentstroke}{rgb}{0.000000,0.000000,0.000000}%
\pgfsetstrokecolor{currentstroke}%
\pgfsetdash{}{0pt}%
\pgfsys@defobject{currentmarker}{\pgfqpoint{0.000000in}{0.000000in}}{\pgfqpoint{0.041667in}{0.000000in}}{%
\pgfpathmoveto{\pgfqpoint{0.000000in}{0.000000in}}%
\pgfpathlineto{\pgfqpoint{0.041667in}{0.000000in}}%
\pgfusepath{stroke,fill}%
}%
\begin{pgfscope}%
\pgfsys@transformshift{0.563118in}{0.401080in}%
\pgfsys@useobject{currentmarker}{}%
\end{pgfscope}%
\end{pgfscope}%
\begin{pgfscope}%
\pgfsetbuttcap%
\pgfsetroundjoin%
\definecolor{currentfill}{rgb}{0.000000,0.000000,0.000000}%
\pgfsetfillcolor{currentfill}%
\pgfsetlinewidth{0.501875pt}%
\definecolor{currentstroke}{rgb}{0.000000,0.000000,0.000000}%
\pgfsetstrokecolor{currentstroke}%
\pgfsetdash{}{0pt}%
\pgfsys@defobject{currentmarker}{\pgfqpoint{-0.041667in}{0.000000in}}{\pgfqpoint{-0.000000in}{0.000000in}}{%
\pgfpathmoveto{\pgfqpoint{-0.000000in}{0.000000in}}%
\pgfpathlineto{\pgfqpoint{-0.041667in}{0.000000in}}%
\pgfusepath{stroke,fill}%
}%
\begin{pgfscope}%
\pgfsys@transformshift{5.600618in}{0.401080in}%
\pgfsys@useobject{currentmarker}{}%
\end{pgfscope}%
\end{pgfscope}%
\begin{pgfscope}%
\definecolor{textcolor}{rgb}{0.000000,0.000000,0.000000}%
\pgfsetstrokecolor{textcolor}%
\pgfsetfillcolor{textcolor}%
\pgftext[x=0.229012in, y=0.352855in, left, base]{\color{textcolor}{\rmfamily\fontsize{10.000000}{12.000000}\selectfont\catcode`\^=\active\def^{\ifmmode\sp\else\^{}\fi}\catcode`\%=\active\def%{\%}$\mathdefault{\ensuremath{-}1.2}$}}%
\end{pgfscope}%
\begin{pgfscope}%
\pgfsetbuttcap%
\pgfsetroundjoin%
\definecolor{currentfill}{rgb}{0.000000,0.000000,0.000000}%
\pgfsetfillcolor{currentfill}%
\pgfsetlinewidth{0.501875pt}%
\definecolor{currentstroke}{rgb}{0.000000,0.000000,0.000000}%
\pgfsetstrokecolor{currentstroke}%
\pgfsetdash{}{0pt}%
\pgfsys@defobject{currentmarker}{\pgfqpoint{0.000000in}{0.000000in}}{\pgfqpoint{0.041667in}{0.000000in}}{%
\pgfpathmoveto{\pgfqpoint{0.000000in}{0.000000in}}%
\pgfpathlineto{\pgfqpoint{0.041667in}{0.000000in}}%
\pgfusepath{stroke,fill}%
}%
\begin{pgfscope}%
\pgfsys@transformshift{0.563118in}{1.068413in}%
\pgfsys@useobject{currentmarker}{}%
\end{pgfscope}%
\end{pgfscope}%
\begin{pgfscope}%
\pgfsetbuttcap%
\pgfsetroundjoin%
\definecolor{currentfill}{rgb}{0.000000,0.000000,0.000000}%
\pgfsetfillcolor{currentfill}%
\pgfsetlinewidth{0.501875pt}%
\definecolor{currentstroke}{rgb}{0.000000,0.000000,0.000000}%
\pgfsetstrokecolor{currentstroke}%
\pgfsetdash{}{0pt}%
\pgfsys@defobject{currentmarker}{\pgfqpoint{-0.041667in}{0.000000in}}{\pgfqpoint{-0.000000in}{0.000000in}}{%
\pgfpathmoveto{\pgfqpoint{-0.000000in}{0.000000in}}%
\pgfpathlineto{\pgfqpoint{-0.041667in}{0.000000in}}%
\pgfusepath{stroke,fill}%
}%
\begin{pgfscope}%
\pgfsys@transformshift{5.600618in}{1.068413in}%
\pgfsys@useobject{currentmarker}{}%
\end{pgfscope}%
\end{pgfscope}%
\begin{pgfscope}%
\definecolor{textcolor}{rgb}{0.000000,0.000000,0.000000}%
\pgfsetstrokecolor{textcolor}%
\pgfsetfillcolor{textcolor}%
\pgftext[x=0.229012in, y=1.020188in, left, base]{\color{textcolor}{\rmfamily\fontsize{10.000000}{12.000000}\selectfont\catcode`\^=\active\def^{\ifmmode\sp\else\^{}\fi}\catcode`\%=\active\def%{\%}$\mathdefault{\ensuremath{-}1.0}$}}%
\end{pgfscope}%
\begin{pgfscope}%
\pgfsetbuttcap%
\pgfsetroundjoin%
\definecolor{currentfill}{rgb}{0.000000,0.000000,0.000000}%
\pgfsetfillcolor{currentfill}%
\pgfsetlinewidth{0.501875pt}%
\definecolor{currentstroke}{rgb}{0.000000,0.000000,0.000000}%
\pgfsetstrokecolor{currentstroke}%
\pgfsetdash{}{0pt}%
\pgfsys@defobject{currentmarker}{\pgfqpoint{0.000000in}{0.000000in}}{\pgfqpoint{0.041667in}{0.000000in}}{%
\pgfpathmoveto{\pgfqpoint{0.000000in}{0.000000in}}%
\pgfpathlineto{\pgfqpoint{0.041667in}{0.000000in}}%
\pgfusepath{stroke,fill}%
}%
\begin{pgfscope}%
\pgfsys@transformshift{0.563118in}{1.735747in}%
\pgfsys@useobject{currentmarker}{}%
\end{pgfscope}%
\end{pgfscope}%
\begin{pgfscope}%
\pgfsetbuttcap%
\pgfsetroundjoin%
\definecolor{currentfill}{rgb}{0.000000,0.000000,0.000000}%
\pgfsetfillcolor{currentfill}%
\pgfsetlinewidth{0.501875pt}%
\definecolor{currentstroke}{rgb}{0.000000,0.000000,0.000000}%
\pgfsetstrokecolor{currentstroke}%
\pgfsetdash{}{0pt}%
\pgfsys@defobject{currentmarker}{\pgfqpoint{-0.041667in}{0.000000in}}{\pgfqpoint{-0.000000in}{0.000000in}}{%
\pgfpathmoveto{\pgfqpoint{-0.000000in}{0.000000in}}%
\pgfpathlineto{\pgfqpoint{-0.041667in}{0.000000in}}%
\pgfusepath{stroke,fill}%
}%
\begin{pgfscope}%
\pgfsys@transformshift{5.600618in}{1.735747in}%
\pgfsys@useobject{currentmarker}{}%
\end{pgfscope}%
\end{pgfscope}%
\begin{pgfscope}%
\definecolor{textcolor}{rgb}{0.000000,0.000000,0.000000}%
\pgfsetstrokecolor{textcolor}%
\pgfsetfillcolor{textcolor}%
\pgftext[x=0.229012in, y=1.687521in, left, base]{\color{textcolor}{\rmfamily\fontsize{10.000000}{12.000000}\selectfont\catcode`\^=\active\def^{\ifmmode\sp\else\^{}\fi}\catcode`\%=\active\def%{\%}$\mathdefault{\ensuremath{-}0.8}$}}%
\end{pgfscope}%
\begin{pgfscope}%
\pgfsetbuttcap%
\pgfsetroundjoin%
\definecolor{currentfill}{rgb}{0.000000,0.000000,0.000000}%
\pgfsetfillcolor{currentfill}%
\pgfsetlinewidth{0.501875pt}%
\definecolor{currentstroke}{rgb}{0.000000,0.000000,0.000000}%
\pgfsetstrokecolor{currentstroke}%
\pgfsetdash{}{0pt}%
\pgfsys@defobject{currentmarker}{\pgfqpoint{0.000000in}{0.000000in}}{\pgfqpoint{0.041667in}{0.000000in}}{%
\pgfpathmoveto{\pgfqpoint{0.000000in}{0.000000in}}%
\pgfpathlineto{\pgfqpoint{0.041667in}{0.000000in}}%
\pgfusepath{stroke,fill}%
}%
\begin{pgfscope}%
\pgfsys@transformshift{0.563118in}{2.403080in}%
\pgfsys@useobject{currentmarker}{}%
\end{pgfscope}%
\end{pgfscope}%
\begin{pgfscope}%
\pgfsetbuttcap%
\pgfsetroundjoin%
\definecolor{currentfill}{rgb}{0.000000,0.000000,0.000000}%
\pgfsetfillcolor{currentfill}%
\pgfsetlinewidth{0.501875pt}%
\definecolor{currentstroke}{rgb}{0.000000,0.000000,0.000000}%
\pgfsetstrokecolor{currentstroke}%
\pgfsetdash{}{0pt}%
\pgfsys@defobject{currentmarker}{\pgfqpoint{-0.041667in}{0.000000in}}{\pgfqpoint{-0.000000in}{0.000000in}}{%
\pgfpathmoveto{\pgfqpoint{-0.000000in}{0.000000in}}%
\pgfpathlineto{\pgfqpoint{-0.041667in}{0.000000in}}%
\pgfusepath{stroke,fill}%
}%
\begin{pgfscope}%
\pgfsys@transformshift{5.600618in}{2.403080in}%
\pgfsys@useobject{currentmarker}{}%
\end{pgfscope}%
\end{pgfscope}%
\begin{pgfscope}%
\definecolor{textcolor}{rgb}{0.000000,0.000000,0.000000}%
\pgfsetstrokecolor{textcolor}%
\pgfsetfillcolor{textcolor}%
\pgftext[x=0.229012in, y=2.354855in, left, base]{\color{textcolor}{\rmfamily\fontsize{10.000000}{12.000000}\selectfont\catcode`\^=\active\def^{\ifmmode\sp\else\^{}\fi}\catcode`\%=\active\def%{\%}$\mathdefault{\ensuremath{-}0.6}$}}%
\end{pgfscope}%
\begin{pgfscope}%
\pgfsetbuttcap%
\pgfsetroundjoin%
\definecolor{currentfill}{rgb}{0.000000,0.000000,0.000000}%
\pgfsetfillcolor{currentfill}%
\pgfsetlinewidth{0.501875pt}%
\definecolor{currentstroke}{rgb}{0.000000,0.000000,0.000000}%
\pgfsetstrokecolor{currentstroke}%
\pgfsetdash{}{0pt}%
\pgfsys@defobject{currentmarker}{\pgfqpoint{0.000000in}{0.000000in}}{\pgfqpoint{0.041667in}{0.000000in}}{%
\pgfpathmoveto{\pgfqpoint{0.000000in}{0.000000in}}%
\pgfpathlineto{\pgfqpoint{0.041667in}{0.000000in}}%
\pgfusepath{stroke,fill}%
}%
\begin{pgfscope}%
\pgfsys@transformshift{0.563118in}{3.070413in}%
\pgfsys@useobject{currentmarker}{}%
\end{pgfscope}%
\end{pgfscope}%
\begin{pgfscope}%
\pgfsetbuttcap%
\pgfsetroundjoin%
\definecolor{currentfill}{rgb}{0.000000,0.000000,0.000000}%
\pgfsetfillcolor{currentfill}%
\pgfsetlinewidth{0.501875pt}%
\definecolor{currentstroke}{rgb}{0.000000,0.000000,0.000000}%
\pgfsetstrokecolor{currentstroke}%
\pgfsetdash{}{0pt}%
\pgfsys@defobject{currentmarker}{\pgfqpoint{-0.041667in}{0.000000in}}{\pgfqpoint{-0.000000in}{0.000000in}}{%
\pgfpathmoveto{\pgfqpoint{-0.000000in}{0.000000in}}%
\pgfpathlineto{\pgfqpoint{-0.041667in}{0.000000in}}%
\pgfusepath{stroke,fill}%
}%
\begin{pgfscope}%
\pgfsys@transformshift{5.600618in}{3.070413in}%
\pgfsys@useobject{currentmarker}{}%
\end{pgfscope}%
\end{pgfscope}%
\begin{pgfscope}%
\definecolor{textcolor}{rgb}{0.000000,0.000000,0.000000}%
\pgfsetstrokecolor{textcolor}%
\pgfsetfillcolor{textcolor}%
\pgftext[x=0.229012in, y=3.022188in, left, base]{\color{textcolor}{\rmfamily\fontsize{10.000000}{12.000000}\selectfont\catcode`\^=\active\def^{\ifmmode\sp\else\^{}\fi}\catcode`\%=\active\def%{\%}$\mathdefault{\ensuremath{-}0.4}$}}%
\end{pgfscope}%
\begin{pgfscope}%
\pgfsetbuttcap%
\pgfsetroundjoin%
\definecolor{currentfill}{rgb}{0.000000,0.000000,0.000000}%
\pgfsetfillcolor{currentfill}%
\pgfsetlinewidth{0.501875pt}%
\definecolor{currentstroke}{rgb}{0.000000,0.000000,0.000000}%
\pgfsetstrokecolor{currentstroke}%
\pgfsetdash{}{0pt}%
\pgfsys@defobject{currentmarker}{\pgfqpoint{0.000000in}{0.000000in}}{\pgfqpoint{0.041667in}{0.000000in}}{%
\pgfpathmoveto{\pgfqpoint{0.000000in}{0.000000in}}%
\pgfpathlineto{\pgfqpoint{0.041667in}{0.000000in}}%
\pgfusepath{stroke,fill}%
}%
\begin{pgfscope}%
\pgfsys@transformshift{0.563118in}{3.737747in}%
\pgfsys@useobject{currentmarker}{}%
\end{pgfscope}%
\end{pgfscope}%
\begin{pgfscope}%
\pgfsetbuttcap%
\pgfsetroundjoin%
\definecolor{currentfill}{rgb}{0.000000,0.000000,0.000000}%
\pgfsetfillcolor{currentfill}%
\pgfsetlinewidth{0.501875pt}%
\definecolor{currentstroke}{rgb}{0.000000,0.000000,0.000000}%
\pgfsetstrokecolor{currentstroke}%
\pgfsetdash{}{0pt}%
\pgfsys@defobject{currentmarker}{\pgfqpoint{-0.041667in}{0.000000in}}{\pgfqpoint{-0.000000in}{0.000000in}}{%
\pgfpathmoveto{\pgfqpoint{-0.000000in}{0.000000in}}%
\pgfpathlineto{\pgfqpoint{-0.041667in}{0.000000in}}%
\pgfusepath{stroke,fill}%
}%
\begin{pgfscope}%
\pgfsys@transformshift{5.600618in}{3.737747in}%
\pgfsys@useobject{currentmarker}{}%
\end{pgfscope}%
\end{pgfscope}%
\begin{pgfscope}%
\definecolor{textcolor}{rgb}{0.000000,0.000000,0.000000}%
\pgfsetstrokecolor{textcolor}%
\pgfsetfillcolor{textcolor}%
\pgftext[x=0.229012in, y=3.689521in, left, base]{\color{textcolor}{\rmfamily\fontsize{10.000000}{12.000000}\selectfont\catcode`\^=\active\def^{\ifmmode\sp\else\^{}\fi}\catcode`\%=\active\def%{\%}$\mathdefault{\ensuremath{-}0.2}$}}%
\end{pgfscope}%
\begin{pgfscope}%
\pgfsetbuttcap%
\pgfsetroundjoin%
\definecolor{currentfill}{rgb}{0.000000,0.000000,0.000000}%
\pgfsetfillcolor{currentfill}%
\pgfsetlinewidth{0.501875pt}%
\definecolor{currentstroke}{rgb}{0.000000,0.000000,0.000000}%
\pgfsetstrokecolor{currentstroke}%
\pgfsetdash{}{0pt}%
\pgfsys@defobject{currentmarker}{\pgfqpoint{0.000000in}{0.000000in}}{\pgfqpoint{0.041667in}{0.000000in}}{%
\pgfpathmoveto{\pgfqpoint{0.000000in}{0.000000in}}%
\pgfpathlineto{\pgfqpoint{0.041667in}{0.000000in}}%
\pgfusepath{stroke,fill}%
}%
\begin{pgfscope}%
\pgfsys@transformshift{0.563118in}{4.405080in}%
\pgfsys@useobject{currentmarker}{}%
\end{pgfscope}%
\end{pgfscope}%
\begin{pgfscope}%
\pgfsetbuttcap%
\pgfsetroundjoin%
\definecolor{currentfill}{rgb}{0.000000,0.000000,0.000000}%
\pgfsetfillcolor{currentfill}%
\pgfsetlinewidth{0.501875pt}%
\definecolor{currentstroke}{rgb}{0.000000,0.000000,0.000000}%
\pgfsetstrokecolor{currentstroke}%
\pgfsetdash{}{0pt}%
\pgfsys@defobject{currentmarker}{\pgfqpoint{-0.041667in}{0.000000in}}{\pgfqpoint{-0.000000in}{0.000000in}}{%
\pgfpathmoveto{\pgfqpoint{-0.000000in}{0.000000in}}%
\pgfpathlineto{\pgfqpoint{-0.041667in}{0.000000in}}%
\pgfusepath{stroke,fill}%
}%
\begin{pgfscope}%
\pgfsys@transformshift{5.600618in}{4.405080in}%
\pgfsys@useobject{currentmarker}{}%
\end{pgfscope}%
\end{pgfscope}%
\begin{pgfscope}%
\definecolor{textcolor}{rgb}{0.000000,0.000000,0.000000}%
\pgfsetstrokecolor{textcolor}%
\pgfsetfillcolor{textcolor}%
\pgftext[x=0.337037in, y=4.356855in, left, base]{\color{textcolor}{\rmfamily\fontsize{10.000000}{12.000000}\selectfont\catcode`\^=\active\def^{\ifmmode\sp\else\^{}\fi}\catcode`\%=\active\def%{\%}$\mathdefault{0.0}$}}%
\end{pgfscope}%
\begin{pgfscope}%
\pgfsetbuttcap%
\pgfsetroundjoin%
\definecolor{currentfill}{rgb}{0.000000,0.000000,0.000000}%
\pgfsetfillcolor{currentfill}%
\pgfsetlinewidth{0.501875pt}%
\definecolor{currentstroke}{rgb}{0.000000,0.000000,0.000000}%
\pgfsetstrokecolor{currentstroke}%
\pgfsetdash{}{0pt}%
\pgfsys@defobject{currentmarker}{\pgfqpoint{0.000000in}{0.000000in}}{\pgfqpoint{0.020833in}{0.000000in}}{%
\pgfpathmoveto{\pgfqpoint{0.000000in}{0.000000in}}%
\pgfpathlineto{\pgfqpoint{0.020833in}{0.000000in}}%
\pgfusepath{stroke,fill}%
}%
\begin{pgfscope}%
\pgfsys@transformshift{0.563118in}{0.567913in}%
\pgfsys@useobject{currentmarker}{}%
\end{pgfscope}%
\end{pgfscope}%
\begin{pgfscope}%
\pgfsetbuttcap%
\pgfsetroundjoin%
\definecolor{currentfill}{rgb}{0.000000,0.000000,0.000000}%
\pgfsetfillcolor{currentfill}%
\pgfsetlinewidth{0.501875pt}%
\definecolor{currentstroke}{rgb}{0.000000,0.000000,0.000000}%
\pgfsetstrokecolor{currentstroke}%
\pgfsetdash{}{0pt}%
\pgfsys@defobject{currentmarker}{\pgfqpoint{-0.020833in}{0.000000in}}{\pgfqpoint{-0.000000in}{0.000000in}}{%
\pgfpathmoveto{\pgfqpoint{-0.000000in}{0.000000in}}%
\pgfpathlineto{\pgfqpoint{-0.020833in}{0.000000in}}%
\pgfusepath{stroke,fill}%
}%
\begin{pgfscope}%
\pgfsys@transformshift{5.600618in}{0.567913in}%
\pgfsys@useobject{currentmarker}{}%
\end{pgfscope}%
\end{pgfscope}%
\begin{pgfscope}%
\pgfsetbuttcap%
\pgfsetroundjoin%
\definecolor{currentfill}{rgb}{0.000000,0.000000,0.000000}%
\pgfsetfillcolor{currentfill}%
\pgfsetlinewidth{0.501875pt}%
\definecolor{currentstroke}{rgb}{0.000000,0.000000,0.000000}%
\pgfsetstrokecolor{currentstroke}%
\pgfsetdash{}{0pt}%
\pgfsys@defobject{currentmarker}{\pgfqpoint{0.000000in}{0.000000in}}{\pgfqpoint{0.020833in}{0.000000in}}{%
\pgfpathmoveto{\pgfqpoint{0.000000in}{0.000000in}}%
\pgfpathlineto{\pgfqpoint{0.020833in}{0.000000in}}%
\pgfusepath{stroke,fill}%
}%
\begin{pgfscope}%
\pgfsys@transformshift{0.563118in}{0.734747in}%
\pgfsys@useobject{currentmarker}{}%
\end{pgfscope}%
\end{pgfscope}%
\begin{pgfscope}%
\pgfsetbuttcap%
\pgfsetroundjoin%
\definecolor{currentfill}{rgb}{0.000000,0.000000,0.000000}%
\pgfsetfillcolor{currentfill}%
\pgfsetlinewidth{0.501875pt}%
\definecolor{currentstroke}{rgb}{0.000000,0.000000,0.000000}%
\pgfsetstrokecolor{currentstroke}%
\pgfsetdash{}{0pt}%
\pgfsys@defobject{currentmarker}{\pgfqpoint{-0.020833in}{0.000000in}}{\pgfqpoint{-0.000000in}{0.000000in}}{%
\pgfpathmoveto{\pgfqpoint{-0.000000in}{0.000000in}}%
\pgfpathlineto{\pgfqpoint{-0.020833in}{0.000000in}}%
\pgfusepath{stroke,fill}%
}%
\begin{pgfscope}%
\pgfsys@transformshift{5.600618in}{0.734747in}%
\pgfsys@useobject{currentmarker}{}%
\end{pgfscope}%
\end{pgfscope}%
\begin{pgfscope}%
\pgfsetbuttcap%
\pgfsetroundjoin%
\definecolor{currentfill}{rgb}{0.000000,0.000000,0.000000}%
\pgfsetfillcolor{currentfill}%
\pgfsetlinewidth{0.501875pt}%
\definecolor{currentstroke}{rgb}{0.000000,0.000000,0.000000}%
\pgfsetstrokecolor{currentstroke}%
\pgfsetdash{}{0pt}%
\pgfsys@defobject{currentmarker}{\pgfqpoint{0.000000in}{0.000000in}}{\pgfqpoint{0.020833in}{0.000000in}}{%
\pgfpathmoveto{\pgfqpoint{0.000000in}{0.000000in}}%
\pgfpathlineto{\pgfqpoint{0.020833in}{0.000000in}}%
\pgfusepath{stroke,fill}%
}%
\begin{pgfscope}%
\pgfsys@transformshift{0.563118in}{0.901580in}%
\pgfsys@useobject{currentmarker}{}%
\end{pgfscope}%
\end{pgfscope}%
\begin{pgfscope}%
\pgfsetbuttcap%
\pgfsetroundjoin%
\definecolor{currentfill}{rgb}{0.000000,0.000000,0.000000}%
\pgfsetfillcolor{currentfill}%
\pgfsetlinewidth{0.501875pt}%
\definecolor{currentstroke}{rgb}{0.000000,0.000000,0.000000}%
\pgfsetstrokecolor{currentstroke}%
\pgfsetdash{}{0pt}%
\pgfsys@defobject{currentmarker}{\pgfqpoint{-0.020833in}{0.000000in}}{\pgfqpoint{-0.000000in}{0.000000in}}{%
\pgfpathmoveto{\pgfqpoint{-0.000000in}{0.000000in}}%
\pgfpathlineto{\pgfqpoint{-0.020833in}{0.000000in}}%
\pgfusepath{stroke,fill}%
}%
\begin{pgfscope}%
\pgfsys@transformshift{5.600618in}{0.901580in}%
\pgfsys@useobject{currentmarker}{}%
\end{pgfscope}%
\end{pgfscope}%
\begin{pgfscope}%
\pgfsetbuttcap%
\pgfsetroundjoin%
\definecolor{currentfill}{rgb}{0.000000,0.000000,0.000000}%
\pgfsetfillcolor{currentfill}%
\pgfsetlinewidth{0.501875pt}%
\definecolor{currentstroke}{rgb}{0.000000,0.000000,0.000000}%
\pgfsetstrokecolor{currentstroke}%
\pgfsetdash{}{0pt}%
\pgfsys@defobject{currentmarker}{\pgfqpoint{0.000000in}{0.000000in}}{\pgfqpoint{0.020833in}{0.000000in}}{%
\pgfpathmoveto{\pgfqpoint{0.000000in}{0.000000in}}%
\pgfpathlineto{\pgfqpoint{0.020833in}{0.000000in}}%
\pgfusepath{stroke,fill}%
}%
\begin{pgfscope}%
\pgfsys@transformshift{0.563118in}{1.235247in}%
\pgfsys@useobject{currentmarker}{}%
\end{pgfscope}%
\end{pgfscope}%
\begin{pgfscope}%
\pgfsetbuttcap%
\pgfsetroundjoin%
\definecolor{currentfill}{rgb}{0.000000,0.000000,0.000000}%
\pgfsetfillcolor{currentfill}%
\pgfsetlinewidth{0.501875pt}%
\definecolor{currentstroke}{rgb}{0.000000,0.000000,0.000000}%
\pgfsetstrokecolor{currentstroke}%
\pgfsetdash{}{0pt}%
\pgfsys@defobject{currentmarker}{\pgfqpoint{-0.020833in}{0.000000in}}{\pgfqpoint{-0.000000in}{0.000000in}}{%
\pgfpathmoveto{\pgfqpoint{-0.000000in}{0.000000in}}%
\pgfpathlineto{\pgfqpoint{-0.020833in}{0.000000in}}%
\pgfusepath{stroke,fill}%
}%
\begin{pgfscope}%
\pgfsys@transformshift{5.600618in}{1.235247in}%
\pgfsys@useobject{currentmarker}{}%
\end{pgfscope}%
\end{pgfscope}%
\begin{pgfscope}%
\pgfsetbuttcap%
\pgfsetroundjoin%
\definecolor{currentfill}{rgb}{0.000000,0.000000,0.000000}%
\pgfsetfillcolor{currentfill}%
\pgfsetlinewidth{0.501875pt}%
\definecolor{currentstroke}{rgb}{0.000000,0.000000,0.000000}%
\pgfsetstrokecolor{currentstroke}%
\pgfsetdash{}{0pt}%
\pgfsys@defobject{currentmarker}{\pgfqpoint{0.000000in}{0.000000in}}{\pgfqpoint{0.020833in}{0.000000in}}{%
\pgfpathmoveto{\pgfqpoint{0.000000in}{0.000000in}}%
\pgfpathlineto{\pgfqpoint{0.020833in}{0.000000in}}%
\pgfusepath{stroke,fill}%
}%
\begin{pgfscope}%
\pgfsys@transformshift{0.563118in}{1.402080in}%
\pgfsys@useobject{currentmarker}{}%
\end{pgfscope}%
\end{pgfscope}%
\begin{pgfscope}%
\pgfsetbuttcap%
\pgfsetroundjoin%
\definecolor{currentfill}{rgb}{0.000000,0.000000,0.000000}%
\pgfsetfillcolor{currentfill}%
\pgfsetlinewidth{0.501875pt}%
\definecolor{currentstroke}{rgb}{0.000000,0.000000,0.000000}%
\pgfsetstrokecolor{currentstroke}%
\pgfsetdash{}{0pt}%
\pgfsys@defobject{currentmarker}{\pgfqpoint{-0.020833in}{0.000000in}}{\pgfqpoint{-0.000000in}{0.000000in}}{%
\pgfpathmoveto{\pgfqpoint{-0.000000in}{0.000000in}}%
\pgfpathlineto{\pgfqpoint{-0.020833in}{0.000000in}}%
\pgfusepath{stroke,fill}%
}%
\begin{pgfscope}%
\pgfsys@transformshift{5.600618in}{1.402080in}%
\pgfsys@useobject{currentmarker}{}%
\end{pgfscope}%
\end{pgfscope}%
\begin{pgfscope}%
\pgfsetbuttcap%
\pgfsetroundjoin%
\definecolor{currentfill}{rgb}{0.000000,0.000000,0.000000}%
\pgfsetfillcolor{currentfill}%
\pgfsetlinewidth{0.501875pt}%
\definecolor{currentstroke}{rgb}{0.000000,0.000000,0.000000}%
\pgfsetstrokecolor{currentstroke}%
\pgfsetdash{}{0pt}%
\pgfsys@defobject{currentmarker}{\pgfqpoint{0.000000in}{0.000000in}}{\pgfqpoint{0.020833in}{0.000000in}}{%
\pgfpathmoveto{\pgfqpoint{0.000000in}{0.000000in}}%
\pgfpathlineto{\pgfqpoint{0.020833in}{0.000000in}}%
\pgfusepath{stroke,fill}%
}%
\begin{pgfscope}%
\pgfsys@transformshift{0.563118in}{1.568913in}%
\pgfsys@useobject{currentmarker}{}%
\end{pgfscope}%
\end{pgfscope}%
\begin{pgfscope}%
\pgfsetbuttcap%
\pgfsetroundjoin%
\definecolor{currentfill}{rgb}{0.000000,0.000000,0.000000}%
\pgfsetfillcolor{currentfill}%
\pgfsetlinewidth{0.501875pt}%
\definecolor{currentstroke}{rgb}{0.000000,0.000000,0.000000}%
\pgfsetstrokecolor{currentstroke}%
\pgfsetdash{}{0pt}%
\pgfsys@defobject{currentmarker}{\pgfqpoint{-0.020833in}{0.000000in}}{\pgfqpoint{-0.000000in}{0.000000in}}{%
\pgfpathmoveto{\pgfqpoint{-0.000000in}{0.000000in}}%
\pgfpathlineto{\pgfqpoint{-0.020833in}{0.000000in}}%
\pgfusepath{stroke,fill}%
}%
\begin{pgfscope}%
\pgfsys@transformshift{5.600618in}{1.568913in}%
\pgfsys@useobject{currentmarker}{}%
\end{pgfscope}%
\end{pgfscope}%
\begin{pgfscope}%
\pgfsetbuttcap%
\pgfsetroundjoin%
\definecolor{currentfill}{rgb}{0.000000,0.000000,0.000000}%
\pgfsetfillcolor{currentfill}%
\pgfsetlinewidth{0.501875pt}%
\definecolor{currentstroke}{rgb}{0.000000,0.000000,0.000000}%
\pgfsetstrokecolor{currentstroke}%
\pgfsetdash{}{0pt}%
\pgfsys@defobject{currentmarker}{\pgfqpoint{0.000000in}{0.000000in}}{\pgfqpoint{0.020833in}{0.000000in}}{%
\pgfpathmoveto{\pgfqpoint{0.000000in}{0.000000in}}%
\pgfpathlineto{\pgfqpoint{0.020833in}{0.000000in}}%
\pgfusepath{stroke,fill}%
}%
\begin{pgfscope}%
\pgfsys@transformshift{0.563118in}{1.902580in}%
\pgfsys@useobject{currentmarker}{}%
\end{pgfscope}%
\end{pgfscope}%
\begin{pgfscope}%
\pgfsetbuttcap%
\pgfsetroundjoin%
\definecolor{currentfill}{rgb}{0.000000,0.000000,0.000000}%
\pgfsetfillcolor{currentfill}%
\pgfsetlinewidth{0.501875pt}%
\definecolor{currentstroke}{rgb}{0.000000,0.000000,0.000000}%
\pgfsetstrokecolor{currentstroke}%
\pgfsetdash{}{0pt}%
\pgfsys@defobject{currentmarker}{\pgfqpoint{-0.020833in}{0.000000in}}{\pgfqpoint{-0.000000in}{0.000000in}}{%
\pgfpathmoveto{\pgfqpoint{-0.000000in}{0.000000in}}%
\pgfpathlineto{\pgfqpoint{-0.020833in}{0.000000in}}%
\pgfusepath{stroke,fill}%
}%
\begin{pgfscope}%
\pgfsys@transformshift{5.600618in}{1.902580in}%
\pgfsys@useobject{currentmarker}{}%
\end{pgfscope}%
\end{pgfscope}%
\begin{pgfscope}%
\pgfsetbuttcap%
\pgfsetroundjoin%
\definecolor{currentfill}{rgb}{0.000000,0.000000,0.000000}%
\pgfsetfillcolor{currentfill}%
\pgfsetlinewidth{0.501875pt}%
\definecolor{currentstroke}{rgb}{0.000000,0.000000,0.000000}%
\pgfsetstrokecolor{currentstroke}%
\pgfsetdash{}{0pt}%
\pgfsys@defobject{currentmarker}{\pgfqpoint{0.000000in}{0.000000in}}{\pgfqpoint{0.020833in}{0.000000in}}{%
\pgfpathmoveto{\pgfqpoint{0.000000in}{0.000000in}}%
\pgfpathlineto{\pgfqpoint{0.020833in}{0.000000in}}%
\pgfusepath{stroke,fill}%
}%
\begin{pgfscope}%
\pgfsys@transformshift{0.563118in}{2.069413in}%
\pgfsys@useobject{currentmarker}{}%
\end{pgfscope}%
\end{pgfscope}%
\begin{pgfscope}%
\pgfsetbuttcap%
\pgfsetroundjoin%
\definecolor{currentfill}{rgb}{0.000000,0.000000,0.000000}%
\pgfsetfillcolor{currentfill}%
\pgfsetlinewidth{0.501875pt}%
\definecolor{currentstroke}{rgb}{0.000000,0.000000,0.000000}%
\pgfsetstrokecolor{currentstroke}%
\pgfsetdash{}{0pt}%
\pgfsys@defobject{currentmarker}{\pgfqpoint{-0.020833in}{0.000000in}}{\pgfqpoint{-0.000000in}{0.000000in}}{%
\pgfpathmoveto{\pgfqpoint{-0.000000in}{0.000000in}}%
\pgfpathlineto{\pgfqpoint{-0.020833in}{0.000000in}}%
\pgfusepath{stroke,fill}%
}%
\begin{pgfscope}%
\pgfsys@transformshift{5.600618in}{2.069413in}%
\pgfsys@useobject{currentmarker}{}%
\end{pgfscope}%
\end{pgfscope}%
\begin{pgfscope}%
\pgfsetbuttcap%
\pgfsetroundjoin%
\definecolor{currentfill}{rgb}{0.000000,0.000000,0.000000}%
\pgfsetfillcolor{currentfill}%
\pgfsetlinewidth{0.501875pt}%
\definecolor{currentstroke}{rgb}{0.000000,0.000000,0.000000}%
\pgfsetstrokecolor{currentstroke}%
\pgfsetdash{}{0pt}%
\pgfsys@defobject{currentmarker}{\pgfqpoint{0.000000in}{0.000000in}}{\pgfqpoint{0.020833in}{0.000000in}}{%
\pgfpathmoveto{\pgfqpoint{0.000000in}{0.000000in}}%
\pgfpathlineto{\pgfqpoint{0.020833in}{0.000000in}}%
\pgfusepath{stroke,fill}%
}%
\begin{pgfscope}%
\pgfsys@transformshift{0.563118in}{2.236247in}%
\pgfsys@useobject{currentmarker}{}%
\end{pgfscope}%
\end{pgfscope}%
\begin{pgfscope}%
\pgfsetbuttcap%
\pgfsetroundjoin%
\definecolor{currentfill}{rgb}{0.000000,0.000000,0.000000}%
\pgfsetfillcolor{currentfill}%
\pgfsetlinewidth{0.501875pt}%
\definecolor{currentstroke}{rgb}{0.000000,0.000000,0.000000}%
\pgfsetstrokecolor{currentstroke}%
\pgfsetdash{}{0pt}%
\pgfsys@defobject{currentmarker}{\pgfqpoint{-0.020833in}{0.000000in}}{\pgfqpoint{-0.000000in}{0.000000in}}{%
\pgfpathmoveto{\pgfqpoint{-0.000000in}{0.000000in}}%
\pgfpathlineto{\pgfqpoint{-0.020833in}{0.000000in}}%
\pgfusepath{stroke,fill}%
}%
\begin{pgfscope}%
\pgfsys@transformshift{5.600618in}{2.236247in}%
\pgfsys@useobject{currentmarker}{}%
\end{pgfscope}%
\end{pgfscope}%
\begin{pgfscope}%
\pgfsetbuttcap%
\pgfsetroundjoin%
\definecolor{currentfill}{rgb}{0.000000,0.000000,0.000000}%
\pgfsetfillcolor{currentfill}%
\pgfsetlinewidth{0.501875pt}%
\definecolor{currentstroke}{rgb}{0.000000,0.000000,0.000000}%
\pgfsetstrokecolor{currentstroke}%
\pgfsetdash{}{0pt}%
\pgfsys@defobject{currentmarker}{\pgfqpoint{0.000000in}{0.000000in}}{\pgfqpoint{0.020833in}{0.000000in}}{%
\pgfpathmoveto{\pgfqpoint{0.000000in}{0.000000in}}%
\pgfpathlineto{\pgfqpoint{0.020833in}{0.000000in}}%
\pgfusepath{stroke,fill}%
}%
\begin{pgfscope}%
\pgfsys@transformshift{0.563118in}{2.569913in}%
\pgfsys@useobject{currentmarker}{}%
\end{pgfscope}%
\end{pgfscope}%
\begin{pgfscope}%
\pgfsetbuttcap%
\pgfsetroundjoin%
\definecolor{currentfill}{rgb}{0.000000,0.000000,0.000000}%
\pgfsetfillcolor{currentfill}%
\pgfsetlinewidth{0.501875pt}%
\definecolor{currentstroke}{rgb}{0.000000,0.000000,0.000000}%
\pgfsetstrokecolor{currentstroke}%
\pgfsetdash{}{0pt}%
\pgfsys@defobject{currentmarker}{\pgfqpoint{-0.020833in}{0.000000in}}{\pgfqpoint{-0.000000in}{0.000000in}}{%
\pgfpathmoveto{\pgfqpoint{-0.000000in}{0.000000in}}%
\pgfpathlineto{\pgfqpoint{-0.020833in}{0.000000in}}%
\pgfusepath{stroke,fill}%
}%
\begin{pgfscope}%
\pgfsys@transformshift{5.600618in}{2.569913in}%
\pgfsys@useobject{currentmarker}{}%
\end{pgfscope}%
\end{pgfscope}%
\begin{pgfscope}%
\pgfsetbuttcap%
\pgfsetroundjoin%
\definecolor{currentfill}{rgb}{0.000000,0.000000,0.000000}%
\pgfsetfillcolor{currentfill}%
\pgfsetlinewidth{0.501875pt}%
\definecolor{currentstroke}{rgb}{0.000000,0.000000,0.000000}%
\pgfsetstrokecolor{currentstroke}%
\pgfsetdash{}{0pt}%
\pgfsys@defobject{currentmarker}{\pgfqpoint{0.000000in}{0.000000in}}{\pgfqpoint{0.020833in}{0.000000in}}{%
\pgfpathmoveto{\pgfqpoint{0.000000in}{0.000000in}}%
\pgfpathlineto{\pgfqpoint{0.020833in}{0.000000in}}%
\pgfusepath{stroke,fill}%
}%
\begin{pgfscope}%
\pgfsys@transformshift{0.563118in}{2.736747in}%
\pgfsys@useobject{currentmarker}{}%
\end{pgfscope}%
\end{pgfscope}%
\begin{pgfscope}%
\pgfsetbuttcap%
\pgfsetroundjoin%
\definecolor{currentfill}{rgb}{0.000000,0.000000,0.000000}%
\pgfsetfillcolor{currentfill}%
\pgfsetlinewidth{0.501875pt}%
\definecolor{currentstroke}{rgb}{0.000000,0.000000,0.000000}%
\pgfsetstrokecolor{currentstroke}%
\pgfsetdash{}{0pt}%
\pgfsys@defobject{currentmarker}{\pgfqpoint{-0.020833in}{0.000000in}}{\pgfqpoint{-0.000000in}{0.000000in}}{%
\pgfpathmoveto{\pgfqpoint{-0.000000in}{0.000000in}}%
\pgfpathlineto{\pgfqpoint{-0.020833in}{0.000000in}}%
\pgfusepath{stroke,fill}%
}%
\begin{pgfscope}%
\pgfsys@transformshift{5.600618in}{2.736747in}%
\pgfsys@useobject{currentmarker}{}%
\end{pgfscope}%
\end{pgfscope}%
\begin{pgfscope}%
\pgfsetbuttcap%
\pgfsetroundjoin%
\definecolor{currentfill}{rgb}{0.000000,0.000000,0.000000}%
\pgfsetfillcolor{currentfill}%
\pgfsetlinewidth{0.501875pt}%
\definecolor{currentstroke}{rgb}{0.000000,0.000000,0.000000}%
\pgfsetstrokecolor{currentstroke}%
\pgfsetdash{}{0pt}%
\pgfsys@defobject{currentmarker}{\pgfqpoint{0.000000in}{0.000000in}}{\pgfqpoint{0.020833in}{0.000000in}}{%
\pgfpathmoveto{\pgfqpoint{0.000000in}{0.000000in}}%
\pgfpathlineto{\pgfqpoint{0.020833in}{0.000000in}}%
\pgfusepath{stroke,fill}%
}%
\begin{pgfscope}%
\pgfsys@transformshift{0.563118in}{2.903580in}%
\pgfsys@useobject{currentmarker}{}%
\end{pgfscope}%
\end{pgfscope}%
\begin{pgfscope}%
\pgfsetbuttcap%
\pgfsetroundjoin%
\definecolor{currentfill}{rgb}{0.000000,0.000000,0.000000}%
\pgfsetfillcolor{currentfill}%
\pgfsetlinewidth{0.501875pt}%
\definecolor{currentstroke}{rgb}{0.000000,0.000000,0.000000}%
\pgfsetstrokecolor{currentstroke}%
\pgfsetdash{}{0pt}%
\pgfsys@defobject{currentmarker}{\pgfqpoint{-0.020833in}{0.000000in}}{\pgfqpoint{-0.000000in}{0.000000in}}{%
\pgfpathmoveto{\pgfqpoint{-0.000000in}{0.000000in}}%
\pgfpathlineto{\pgfqpoint{-0.020833in}{0.000000in}}%
\pgfusepath{stroke,fill}%
}%
\begin{pgfscope}%
\pgfsys@transformshift{5.600618in}{2.903580in}%
\pgfsys@useobject{currentmarker}{}%
\end{pgfscope}%
\end{pgfscope}%
\begin{pgfscope}%
\pgfsetbuttcap%
\pgfsetroundjoin%
\definecolor{currentfill}{rgb}{0.000000,0.000000,0.000000}%
\pgfsetfillcolor{currentfill}%
\pgfsetlinewidth{0.501875pt}%
\definecolor{currentstroke}{rgb}{0.000000,0.000000,0.000000}%
\pgfsetstrokecolor{currentstroke}%
\pgfsetdash{}{0pt}%
\pgfsys@defobject{currentmarker}{\pgfqpoint{0.000000in}{0.000000in}}{\pgfqpoint{0.020833in}{0.000000in}}{%
\pgfpathmoveto{\pgfqpoint{0.000000in}{0.000000in}}%
\pgfpathlineto{\pgfqpoint{0.020833in}{0.000000in}}%
\pgfusepath{stroke,fill}%
}%
\begin{pgfscope}%
\pgfsys@transformshift{0.563118in}{3.237247in}%
\pgfsys@useobject{currentmarker}{}%
\end{pgfscope}%
\end{pgfscope}%
\begin{pgfscope}%
\pgfsetbuttcap%
\pgfsetroundjoin%
\definecolor{currentfill}{rgb}{0.000000,0.000000,0.000000}%
\pgfsetfillcolor{currentfill}%
\pgfsetlinewidth{0.501875pt}%
\definecolor{currentstroke}{rgb}{0.000000,0.000000,0.000000}%
\pgfsetstrokecolor{currentstroke}%
\pgfsetdash{}{0pt}%
\pgfsys@defobject{currentmarker}{\pgfqpoint{-0.020833in}{0.000000in}}{\pgfqpoint{-0.000000in}{0.000000in}}{%
\pgfpathmoveto{\pgfqpoint{-0.000000in}{0.000000in}}%
\pgfpathlineto{\pgfqpoint{-0.020833in}{0.000000in}}%
\pgfusepath{stroke,fill}%
}%
\begin{pgfscope}%
\pgfsys@transformshift{5.600618in}{3.237247in}%
\pgfsys@useobject{currentmarker}{}%
\end{pgfscope}%
\end{pgfscope}%
\begin{pgfscope}%
\pgfsetbuttcap%
\pgfsetroundjoin%
\definecolor{currentfill}{rgb}{0.000000,0.000000,0.000000}%
\pgfsetfillcolor{currentfill}%
\pgfsetlinewidth{0.501875pt}%
\definecolor{currentstroke}{rgb}{0.000000,0.000000,0.000000}%
\pgfsetstrokecolor{currentstroke}%
\pgfsetdash{}{0pt}%
\pgfsys@defobject{currentmarker}{\pgfqpoint{0.000000in}{0.000000in}}{\pgfqpoint{0.020833in}{0.000000in}}{%
\pgfpathmoveto{\pgfqpoint{0.000000in}{0.000000in}}%
\pgfpathlineto{\pgfqpoint{0.020833in}{0.000000in}}%
\pgfusepath{stroke,fill}%
}%
\begin{pgfscope}%
\pgfsys@transformshift{0.563118in}{3.404080in}%
\pgfsys@useobject{currentmarker}{}%
\end{pgfscope}%
\end{pgfscope}%
\begin{pgfscope}%
\pgfsetbuttcap%
\pgfsetroundjoin%
\definecolor{currentfill}{rgb}{0.000000,0.000000,0.000000}%
\pgfsetfillcolor{currentfill}%
\pgfsetlinewidth{0.501875pt}%
\definecolor{currentstroke}{rgb}{0.000000,0.000000,0.000000}%
\pgfsetstrokecolor{currentstroke}%
\pgfsetdash{}{0pt}%
\pgfsys@defobject{currentmarker}{\pgfqpoint{-0.020833in}{0.000000in}}{\pgfqpoint{-0.000000in}{0.000000in}}{%
\pgfpathmoveto{\pgfqpoint{-0.000000in}{0.000000in}}%
\pgfpathlineto{\pgfqpoint{-0.020833in}{0.000000in}}%
\pgfusepath{stroke,fill}%
}%
\begin{pgfscope}%
\pgfsys@transformshift{5.600618in}{3.404080in}%
\pgfsys@useobject{currentmarker}{}%
\end{pgfscope}%
\end{pgfscope}%
\begin{pgfscope}%
\pgfsetbuttcap%
\pgfsetroundjoin%
\definecolor{currentfill}{rgb}{0.000000,0.000000,0.000000}%
\pgfsetfillcolor{currentfill}%
\pgfsetlinewidth{0.501875pt}%
\definecolor{currentstroke}{rgb}{0.000000,0.000000,0.000000}%
\pgfsetstrokecolor{currentstroke}%
\pgfsetdash{}{0pt}%
\pgfsys@defobject{currentmarker}{\pgfqpoint{0.000000in}{0.000000in}}{\pgfqpoint{0.020833in}{0.000000in}}{%
\pgfpathmoveto{\pgfqpoint{0.000000in}{0.000000in}}%
\pgfpathlineto{\pgfqpoint{0.020833in}{0.000000in}}%
\pgfusepath{stroke,fill}%
}%
\begin{pgfscope}%
\pgfsys@transformshift{0.563118in}{3.570913in}%
\pgfsys@useobject{currentmarker}{}%
\end{pgfscope}%
\end{pgfscope}%
\begin{pgfscope}%
\pgfsetbuttcap%
\pgfsetroundjoin%
\definecolor{currentfill}{rgb}{0.000000,0.000000,0.000000}%
\pgfsetfillcolor{currentfill}%
\pgfsetlinewidth{0.501875pt}%
\definecolor{currentstroke}{rgb}{0.000000,0.000000,0.000000}%
\pgfsetstrokecolor{currentstroke}%
\pgfsetdash{}{0pt}%
\pgfsys@defobject{currentmarker}{\pgfqpoint{-0.020833in}{0.000000in}}{\pgfqpoint{-0.000000in}{0.000000in}}{%
\pgfpathmoveto{\pgfqpoint{-0.000000in}{0.000000in}}%
\pgfpathlineto{\pgfqpoint{-0.020833in}{0.000000in}}%
\pgfusepath{stroke,fill}%
}%
\begin{pgfscope}%
\pgfsys@transformshift{5.600618in}{3.570913in}%
\pgfsys@useobject{currentmarker}{}%
\end{pgfscope}%
\end{pgfscope}%
\begin{pgfscope}%
\pgfsetbuttcap%
\pgfsetroundjoin%
\definecolor{currentfill}{rgb}{0.000000,0.000000,0.000000}%
\pgfsetfillcolor{currentfill}%
\pgfsetlinewidth{0.501875pt}%
\definecolor{currentstroke}{rgb}{0.000000,0.000000,0.000000}%
\pgfsetstrokecolor{currentstroke}%
\pgfsetdash{}{0pt}%
\pgfsys@defobject{currentmarker}{\pgfqpoint{0.000000in}{0.000000in}}{\pgfqpoint{0.020833in}{0.000000in}}{%
\pgfpathmoveto{\pgfqpoint{0.000000in}{0.000000in}}%
\pgfpathlineto{\pgfqpoint{0.020833in}{0.000000in}}%
\pgfusepath{stroke,fill}%
}%
\begin{pgfscope}%
\pgfsys@transformshift{0.563118in}{3.904580in}%
\pgfsys@useobject{currentmarker}{}%
\end{pgfscope}%
\end{pgfscope}%
\begin{pgfscope}%
\pgfsetbuttcap%
\pgfsetroundjoin%
\definecolor{currentfill}{rgb}{0.000000,0.000000,0.000000}%
\pgfsetfillcolor{currentfill}%
\pgfsetlinewidth{0.501875pt}%
\definecolor{currentstroke}{rgb}{0.000000,0.000000,0.000000}%
\pgfsetstrokecolor{currentstroke}%
\pgfsetdash{}{0pt}%
\pgfsys@defobject{currentmarker}{\pgfqpoint{-0.020833in}{0.000000in}}{\pgfqpoint{-0.000000in}{0.000000in}}{%
\pgfpathmoveto{\pgfqpoint{-0.000000in}{0.000000in}}%
\pgfpathlineto{\pgfqpoint{-0.020833in}{0.000000in}}%
\pgfusepath{stroke,fill}%
}%
\begin{pgfscope}%
\pgfsys@transformshift{5.600618in}{3.904580in}%
\pgfsys@useobject{currentmarker}{}%
\end{pgfscope}%
\end{pgfscope}%
\begin{pgfscope}%
\pgfsetbuttcap%
\pgfsetroundjoin%
\definecolor{currentfill}{rgb}{0.000000,0.000000,0.000000}%
\pgfsetfillcolor{currentfill}%
\pgfsetlinewidth{0.501875pt}%
\definecolor{currentstroke}{rgb}{0.000000,0.000000,0.000000}%
\pgfsetstrokecolor{currentstroke}%
\pgfsetdash{}{0pt}%
\pgfsys@defobject{currentmarker}{\pgfqpoint{0.000000in}{0.000000in}}{\pgfqpoint{0.020833in}{0.000000in}}{%
\pgfpathmoveto{\pgfqpoint{0.000000in}{0.000000in}}%
\pgfpathlineto{\pgfqpoint{0.020833in}{0.000000in}}%
\pgfusepath{stroke,fill}%
}%
\begin{pgfscope}%
\pgfsys@transformshift{0.563118in}{4.071413in}%
\pgfsys@useobject{currentmarker}{}%
\end{pgfscope}%
\end{pgfscope}%
\begin{pgfscope}%
\pgfsetbuttcap%
\pgfsetroundjoin%
\definecolor{currentfill}{rgb}{0.000000,0.000000,0.000000}%
\pgfsetfillcolor{currentfill}%
\pgfsetlinewidth{0.501875pt}%
\definecolor{currentstroke}{rgb}{0.000000,0.000000,0.000000}%
\pgfsetstrokecolor{currentstroke}%
\pgfsetdash{}{0pt}%
\pgfsys@defobject{currentmarker}{\pgfqpoint{-0.020833in}{0.000000in}}{\pgfqpoint{-0.000000in}{0.000000in}}{%
\pgfpathmoveto{\pgfqpoint{-0.000000in}{0.000000in}}%
\pgfpathlineto{\pgfqpoint{-0.020833in}{0.000000in}}%
\pgfusepath{stroke,fill}%
}%
\begin{pgfscope}%
\pgfsys@transformshift{5.600618in}{4.071413in}%
\pgfsys@useobject{currentmarker}{}%
\end{pgfscope}%
\end{pgfscope}%
\begin{pgfscope}%
\pgfsetbuttcap%
\pgfsetroundjoin%
\definecolor{currentfill}{rgb}{0.000000,0.000000,0.000000}%
\pgfsetfillcolor{currentfill}%
\pgfsetlinewidth{0.501875pt}%
\definecolor{currentstroke}{rgb}{0.000000,0.000000,0.000000}%
\pgfsetstrokecolor{currentstroke}%
\pgfsetdash{}{0pt}%
\pgfsys@defobject{currentmarker}{\pgfqpoint{0.000000in}{0.000000in}}{\pgfqpoint{0.020833in}{0.000000in}}{%
\pgfpathmoveto{\pgfqpoint{0.000000in}{0.000000in}}%
\pgfpathlineto{\pgfqpoint{0.020833in}{0.000000in}}%
\pgfusepath{stroke,fill}%
}%
\begin{pgfscope}%
\pgfsys@transformshift{0.563118in}{4.238247in}%
\pgfsys@useobject{currentmarker}{}%
\end{pgfscope}%
\end{pgfscope}%
\begin{pgfscope}%
\pgfsetbuttcap%
\pgfsetroundjoin%
\definecolor{currentfill}{rgb}{0.000000,0.000000,0.000000}%
\pgfsetfillcolor{currentfill}%
\pgfsetlinewidth{0.501875pt}%
\definecolor{currentstroke}{rgb}{0.000000,0.000000,0.000000}%
\pgfsetstrokecolor{currentstroke}%
\pgfsetdash{}{0pt}%
\pgfsys@defobject{currentmarker}{\pgfqpoint{-0.020833in}{0.000000in}}{\pgfqpoint{-0.000000in}{0.000000in}}{%
\pgfpathmoveto{\pgfqpoint{-0.000000in}{0.000000in}}%
\pgfpathlineto{\pgfqpoint{-0.020833in}{0.000000in}}%
\pgfusepath{stroke,fill}%
}%
\begin{pgfscope}%
\pgfsys@transformshift{5.600618in}{4.238247in}%
\pgfsys@useobject{currentmarker}{}%
\end{pgfscope}%
\end{pgfscope}%
\begin{pgfscope}%
\definecolor{textcolor}{rgb}{0.000000,0.000000,0.000000}%
\pgfsetstrokecolor{textcolor}%
\pgfsetfillcolor{textcolor}%
\pgftext[x=0.173457in,y=2.403080in,,bottom,rotate=90.000000]{\color{textcolor}{\rmfamily\fontsize{10.000000}{12.000000}\selectfont\catcode`\^=\active\def^{\ifmmode\sp\else\^{}\fi}\catcode`\%=\active\def%{\%}$F_s$}}%
\end{pgfscope}%
\begin{pgfscope}%
\pgfpathrectangle{\pgfqpoint{0.563118in}{0.401080in}}{\pgfqpoint{5.037500in}{4.004000in}}%
\pgfusepath{clip}%
\pgfsetbuttcap%
\pgfsetroundjoin%
\definecolor{currentfill}{rgb}{0.690196,0.188235,0.376471}%
\pgfsetfillcolor{currentfill}%
\pgfsetlinewidth{1.003750pt}%
\definecolor{currentstroke}{rgb}{0.690196,0.188235,0.376471}%
\pgfsetstrokecolor{currentstroke}%
\pgfsetdash{}{0pt}%
\pgfsys@defobject{currentmarker}{\pgfqpoint{-0.034722in}{-0.034722in}}{\pgfqpoint{0.034722in}{0.034722in}}{%
\pgfpathmoveto{\pgfqpoint{0.000000in}{-0.034722in}}%
\pgfpathcurveto{\pgfqpoint{0.009208in}{-0.034722in}}{\pgfqpoint{0.018041in}{-0.031064in}}{\pgfqpoint{0.024552in}{-0.024552in}}%
\pgfpathcurveto{\pgfqpoint{0.031064in}{-0.018041in}}{\pgfqpoint{0.034722in}{-0.009208in}}{\pgfqpoint{0.034722in}{0.000000in}}%
\pgfpathcurveto{\pgfqpoint{0.034722in}{0.009208in}}{\pgfqpoint{0.031064in}{0.018041in}}{\pgfqpoint{0.024552in}{0.024552in}}%
\pgfpathcurveto{\pgfqpoint{0.018041in}{0.031064in}}{\pgfqpoint{0.009208in}{0.034722in}}{\pgfqpoint{0.000000in}{0.034722in}}%
\pgfpathcurveto{\pgfqpoint{-0.009208in}{0.034722in}}{\pgfqpoint{-0.018041in}{0.031064in}}{\pgfqpoint{-0.024552in}{0.024552in}}%
\pgfpathcurveto{\pgfqpoint{-0.031064in}{0.018041in}}{\pgfqpoint{-0.034722in}{0.009208in}}{\pgfqpoint{-0.034722in}{0.000000in}}%
\pgfpathcurveto{\pgfqpoint{-0.034722in}{-0.009208in}}{\pgfqpoint{-0.031064in}{-0.018041in}}{\pgfqpoint{-0.024552in}{-0.024552in}}%
\pgfpathcurveto{\pgfqpoint{-0.018041in}{-0.031064in}}{\pgfqpoint{-0.009208in}{-0.034722in}}{\pgfqpoint{0.000000in}{-0.034722in}}%
\pgfpathlineto{\pgfqpoint{0.000000in}{-0.034722in}}%
\pgfpathclose%
\pgfusepath{stroke,fill}%
}%
\begin{pgfscope}%
\pgfsys@transformshift{0.915743in}{4.405080in}%
\pgfsys@useobject{currentmarker}{}%
\end{pgfscope}%
\begin{pgfscope}%
\pgfsys@transformshift{1.217993in}{3.437447in}%
\pgfsys@useobject{currentmarker}{}%
\end{pgfscope}%
\begin{pgfscope}%
\pgfsys@transformshift{1.520243in}{2.770113in}%
\pgfsys@useobject{currentmarker}{}%
\end{pgfscope}%
\begin{pgfscope}%
\pgfsys@transformshift{2.628493in}{2.002680in}%
\pgfsys@useobject{currentmarker}{}%
\end{pgfscope}%
\begin{pgfscope}%
\pgfsys@transformshift{3.031493in}{0.567913in}%
\pgfsys@useobject{currentmarker}{}%
\end{pgfscope}%
\begin{pgfscope}%
\pgfsys@transformshift{3.182618in}{2.403080in}%
\pgfsys@useobject{currentmarker}{}%
\end{pgfscope}%
\begin{pgfscope}%
\pgfsys@transformshift{3.938243in}{1.602280in}%
\pgfsys@useobject{currentmarker}{}%
\end{pgfscope}%
\begin{pgfscope}%
\pgfsys@transformshift{4.139743in}{1.268613in}%
\pgfsys@useobject{currentmarker}{}%
\end{pgfscope}%
\begin{pgfscope}%
\pgfsys@transformshift{4.391618in}{0.901580in}%
\pgfsys@useobject{currentmarker}{}%
\end{pgfscope}%
\begin{pgfscope}%
\pgfsys@transformshift{5.600618in}{2.302980in}%
\pgfsys@useobject{currentmarker}{}%
\end{pgfscope}%
\end{pgfscope}%
\begin{pgfscope}%
\pgfpathrectangle{\pgfqpoint{0.563118in}{0.401080in}}{\pgfqpoint{5.037500in}{4.004000in}}%
\pgfusepath{clip}%
\pgfsetbuttcap%
\pgfsetroundjoin%
\pgfsetlinewidth{1.003750pt}%
\definecolor{currentstroke}{rgb}{0.690196,0.188235,0.376471}%
\pgfsetstrokecolor{currentstroke}%
\pgfsetdash{}{0pt}%
\pgfpathmoveto{\pgfqpoint{0.915743in}{3.402724in}}%
\pgfpathcurveto{\pgfqpoint{0.924951in}{3.402724in}}{\pgfqpoint{0.933784in}{3.406383in}}{\pgfqpoint{0.940295in}{3.412894in}}%
\pgfpathcurveto{\pgfqpoint{0.946807in}{3.419406in}}{\pgfqpoint{0.950465in}{3.428238in}}{\pgfqpoint{0.950465in}{3.437447in}}%
\pgfpathcurveto{\pgfqpoint{0.950465in}{3.446655in}}{\pgfqpoint{0.946807in}{3.455488in}}{\pgfqpoint{0.940295in}{3.461999in}}%
\pgfpathcurveto{\pgfqpoint{0.933784in}{3.468510in}}{\pgfqpoint{0.924951in}{3.472169in}}{\pgfqpoint{0.915743in}{3.472169in}}%
\pgfpathcurveto{\pgfqpoint{0.906535in}{3.472169in}}{\pgfqpoint{0.897702in}{3.468510in}}{\pgfqpoint{0.891191in}{3.461999in}}%
\pgfpathcurveto{\pgfqpoint{0.884679in}{3.455488in}}{\pgfqpoint{0.881021in}{3.446655in}}{\pgfqpoint{0.881021in}{3.437447in}}%
\pgfpathcurveto{\pgfqpoint{0.881021in}{3.428238in}}{\pgfqpoint{0.884679in}{3.419406in}}{\pgfqpoint{0.891191in}{3.412894in}}%
\pgfpathcurveto{\pgfqpoint{0.897702in}{3.406383in}}{\pgfqpoint{0.906535in}{3.402724in}}{\pgfqpoint{0.915743in}{3.402724in}}%
\pgfpathlineto{\pgfqpoint{0.915743in}{3.402724in}}%
\pgfpathclose%
\pgfusepath{stroke}%
\end{pgfscope}%
\begin{pgfscope}%
\pgfpathrectangle{\pgfqpoint{0.563118in}{0.401080in}}{\pgfqpoint{5.037500in}{4.004000in}}%
\pgfusepath{clip}%
\pgfsetbuttcap%
\pgfsetroundjoin%
\pgfsetlinewidth{1.003750pt}%
\definecolor{currentstroke}{rgb}{0.690196,0.188235,0.376471}%
\pgfsetstrokecolor{currentstroke}%
\pgfsetdash{}{0pt}%
\pgfpathmoveto{\pgfqpoint{1.217993in}{2.735391in}}%
\pgfpathcurveto{\pgfqpoint{1.227201in}{2.735391in}}{\pgfqpoint{1.236034in}{2.739050in}}{\pgfqpoint{1.242545in}{2.745561in}}%
\pgfpathcurveto{\pgfqpoint{1.249057in}{2.752072in}}{\pgfqpoint{1.252715in}{2.760905in}}{\pgfqpoint{1.252715in}{2.770113in}}%
\pgfpathcurveto{\pgfqpoint{1.252715in}{2.779322in}}{\pgfqpoint{1.249057in}{2.788154in}}{\pgfqpoint{1.242545in}{2.794666in}}%
\pgfpathcurveto{\pgfqpoint{1.236034in}{2.801177in}}{\pgfqpoint{1.227201in}{2.804836in}}{\pgfqpoint{1.217993in}{2.804836in}}%
\pgfpathcurveto{\pgfqpoint{1.208785in}{2.804836in}}{\pgfqpoint{1.199952in}{2.801177in}}{\pgfqpoint{1.193441in}{2.794666in}}%
\pgfpathcurveto{\pgfqpoint{1.186929in}{2.788154in}}{\pgfqpoint{1.183271in}{2.779322in}}{\pgfqpoint{1.183271in}{2.770113in}}%
\pgfpathcurveto{\pgfqpoint{1.183271in}{2.760905in}}{\pgfqpoint{1.186929in}{2.752072in}}{\pgfqpoint{1.193441in}{2.745561in}}%
\pgfpathcurveto{\pgfqpoint{1.199952in}{2.739050in}}{\pgfqpoint{1.208785in}{2.735391in}}{\pgfqpoint{1.217993in}{2.735391in}}%
\pgfpathlineto{\pgfqpoint{1.217993in}{2.735391in}}%
\pgfpathclose%
\pgfusepath{stroke}%
\end{pgfscope}%
\begin{pgfscope}%
\pgfpathrectangle{\pgfqpoint{0.563118in}{0.401080in}}{\pgfqpoint{5.037500in}{4.004000in}}%
\pgfusepath{clip}%
\pgfsetbuttcap%
\pgfsetroundjoin%
\pgfsetlinewidth{1.003750pt}%
\definecolor{currentstroke}{rgb}{0.690196,0.188235,0.376471}%
\pgfsetstrokecolor{currentstroke}%
\pgfsetdash{}{0pt}%
\pgfpathmoveto{\pgfqpoint{1.520243in}{1.967958in}}%
\pgfpathcurveto{\pgfqpoint{1.529451in}{1.967958in}}{\pgfqpoint{1.538284in}{1.971616in}}{\pgfqpoint{1.544795in}{1.978128in}}%
\pgfpathcurveto{\pgfqpoint{1.551307in}{1.984639in}}{\pgfqpoint{1.554965in}{1.993472in}}{\pgfqpoint{1.554965in}{2.002680in}}%
\pgfpathcurveto{\pgfqpoint{1.554965in}{2.011888in}}{\pgfqpoint{1.551307in}{2.020721in}}{\pgfqpoint{1.544795in}{2.027232in}}%
\pgfpathcurveto{\pgfqpoint{1.538284in}{2.033744in}}{\pgfqpoint{1.529451in}{2.037402in}}{\pgfqpoint{1.520243in}{2.037402in}}%
\pgfpathcurveto{\pgfqpoint{1.511035in}{2.037402in}}{\pgfqpoint{1.502202in}{2.033744in}}{\pgfqpoint{1.495691in}{2.027232in}}%
\pgfpathcurveto{\pgfqpoint{1.489179in}{2.020721in}}{\pgfqpoint{1.485521in}{2.011888in}}{\pgfqpoint{1.485521in}{2.002680in}}%
\pgfpathcurveto{\pgfqpoint{1.485521in}{1.993472in}}{\pgfqpoint{1.489179in}{1.984639in}}{\pgfqpoint{1.495691in}{1.978128in}}%
\pgfpathcurveto{\pgfqpoint{1.502202in}{1.971616in}}{\pgfqpoint{1.511035in}{1.967958in}}{\pgfqpoint{1.520243in}{1.967958in}}%
\pgfpathlineto{\pgfqpoint{1.520243in}{1.967958in}}%
\pgfpathclose%
\pgfusepath{stroke}%
\end{pgfscope}%
\begin{pgfscope}%
\pgfpathrectangle{\pgfqpoint{0.563118in}{0.401080in}}{\pgfqpoint{5.037500in}{4.004000in}}%
\pgfusepath{clip}%
\pgfsetbuttcap%
\pgfsetroundjoin%
\pgfsetlinewidth{1.003750pt}%
\definecolor{currentstroke}{rgb}{0.690196,0.188235,0.376471}%
\pgfsetstrokecolor{currentstroke}%
\pgfsetdash{}{0pt}%
\pgfpathmoveto{\pgfqpoint{2.628493in}{0.533191in}}%
\pgfpathcurveto{\pgfqpoint{2.637701in}{0.533191in}}{\pgfqpoint{2.646534in}{0.536850in}}{\pgfqpoint{2.653045in}{0.543361in}}%
\pgfpathcurveto{\pgfqpoint{2.659557in}{0.549872in}}{\pgfqpoint{2.663215in}{0.558705in}}{\pgfqpoint{2.663215in}{0.567913in}}%
\pgfpathcurveto{\pgfqpoint{2.663215in}{0.577122in}}{\pgfqpoint{2.659557in}{0.585954in}}{\pgfqpoint{2.653045in}{0.592466in}}%
\pgfpathcurveto{\pgfqpoint{2.646534in}{0.598977in}}{\pgfqpoint{2.637701in}{0.602636in}}{\pgfqpoint{2.628493in}{0.602636in}}%
\pgfpathcurveto{\pgfqpoint{2.619285in}{0.602636in}}{\pgfqpoint{2.610452in}{0.598977in}}{\pgfqpoint{2.603941in}{0.592466in}}%
\pgfpathcurveto{\pgfqpoint{2.597429in}{0.585954in}}{\pgfqpoint{2.593771in}{0.577122in}}{\pgfqpoint{2.593771in}{0.567913in}}%
\pgfpathcurveto{\pgfqpoint{2.593771in}{0.558705in}}{\pgfqpoint{2.597429in}{0.549872in}}{\pgfqpoint{2.603941in}{0.543361in}}%
\pgfpathcurveto{\pgfqpoint{2.610452in}{0.536850in}}{\pgfqpoint{2.619285in}{0.533191in}}{\pgfqpoint{2.628493in}{0.533191in}}%
\pgfpathlineto{\pgfqpoint{2.628493in}{0.533191in}}%
\pgfpathclose%
\pgfusepath{stroke}%
\end{pgfscope}%
\begin{pgfscope}%
\pgfpathrectangle{\pgfqpoint{0.563118in}{0.401080in}}{\pgfqpoint{5.037500in}{4.004000in}}%
\pgfusepath{clip}%
\pgfsetbuttcap%
\pgfsetroundjoin%
\pgfsetlinewidth{1.003750pt}%
\definecolor{currentstroke}{rgb}{0.690196,0.188235,0.376471}%
\pgfsetstrokecolor{currentstroke}%
\pgfsetdash{}{0pt}%
\pgfpathmoveto{\pgfqpoint{3.031493in}{2.368358in}}%
\pgfpathcurveto{\pgfqpoint{3.040701in}{2.368358in}}{\pgfqpoint{3.049534in}{2.372016in}}{\pgfqpoint{3.056045in}{2.378528in}}%
\pgfpathcurveto{\pgfqpoint{3.062557in}{2.385039in}}{\pgfqpoint{3.066215in}{2.393872in}}{\pgfqpoint{3.066215in}{2.403080in}}%
\pgfpathcurveto{\pgfqpoint{3.066215in}{2.412288in}}{\pgfqpoint{3.062557in}{2.421121in}}{\pgfqpoint{3.056045in}{2.427632in}}%
\pgfpathcurveto{\pgfqpoint{3.049534in}{2.434144in}}{\pgfqpoint{3.040701in}{2.437802in}}{\pgfqpoint{3.031493in}{2.437802in}}%
\pgfpathcurveto{\pgfqpoint{3.022285in}{2.437802in}}{\pgfqpoint{3.013452in}{2.434144in}}{\pgfqpoint{3.006941in}{2.427632in}}%
\pgfpathcurveto{\pgfqpoint{3.000429in}{2.421121in}}{\pgfqpoint{2.996771in}{2.412288in}}{\pgfqpoint{2.996771in}{2.403080in}}%
\pgfpathcurveto{\pgfqpoint{2.996771in}{2.393872in}}{\pgfqpoint{3.000429in}{2.385039in}}{\pgfqpoint{3.006941in}{2.378528in}}%
\pgfpathcurveto{\pgfqpoint{3.013452in}{2.372016in}}{\pgfqpoint{3.022285in}{2.368358in}}{\pgfqpoint{3.031493in}{2.368358in}}%
\pgfpathlineto{\pgfqpoint{3.031493in}{2.368358in}}%
\pgfpathclose%
\pgfusepath{stroke}%
\end{pgfscope}%
\begin{pgfscope}%
\pgfpathrectangle{\pgfqpoint{0.563118in}{0.401080in}}{\pgfqpoint{5.037500in}{4.004000in}}%
\pgfusepath{clip}%
\pgfsetbuttcap%
\pgfsetroundjoin%
\pgfsetlinewidth{1.003750pt}%
\definecolor{currentstroke}{rgb}{0.690196,0.188235,0.376471}%
\pgfsetstrokecolor{currentstroke}%
\pgfsetdash{}{0pt}%
\pgfpathmoveto{\pgfqpoint{3.182618in}{1.567558in}}%
\pgfpathcurveto{\pgfqpoint{3.191826in}{1.567558in}}{\pgfqpoint{3.200659in}{1.571216in}}{\pgfqpoint{3.207170in}{1.577728in}}%
\pgfpathcurveto{\pgfqpoint{3.213682in}{1.584239in}}{\pgfqpoint{3.217340in}{1.593072in}}{\pgfqpoint{3.217340in}{1.602280in}}%
\pgfpathcurveto{\pgfqpoint{3.217340in}{1.611488in}}{\pgfqpoint{3.213682in}{1.620321in}}{\pgfqpoint{3.207170in}{1.626832in}}%
\pgfpathcurveto{\pgfqpoint{3.200659in}{1.633344in}}{\pgfqpoint{3.191826in}{1.637002in}}{\pgfqpoint{3.182618in}{1.637002in}}%
\pgfpathcurveto{\pgfqpoint{3.173410in}{1.637002in}}{\pgfqpoint{3.164577in}{1.633344in}}{\pgfqpoint{3.158066in}{1.626832in}}%
\pgfpathcurveto{\pgfqpoint{3.151554in}{1.620321in}}{\pgfqpoint{3.147896in}{1.611488in}}{\pgfqpoint{3.147896in}{1.602280in}}%
\pgfpathcurveto{\pgfqpoint{3.147896in}{1.593072in}}{\pgfqpoint{3.151554in}{1.584239in}}{\pgfqpoint{3.158066in}{1.577728in}}%
\pgfpathcurveto{\pgfqpoint{3.164577in}{1.571216in}}{\pgfqpoint{3.173410in}{1.567558in}}{\pgfqpoint{3.182618in}{1.567558in}}%
\pgfpathlineto{\pgfqpoint{3.182618in}{1.567558in}}%
\pgfpathclose%
\pgfusepath{stroke}%
\end{pgfscope}%
\begin{pgfscope}%
\pgfpathrectangle{\pgfqpoint{0.563118in}{0.401080in}}{\pgfqpoint{5.037500in}{4.004000in}}%
\pgfusepath{clip}%
\pgfsetbuttcap%
\pgfsetroundjoin%
\pgfsetlinewidth{1.003750pt}%
\definecolor{currentstroke}{rgb}{0.690196,0.188235,0.376471}%
\pgfsetstrokecolor{currentstroke}%
\pgfsetdash{}{0pt}%
\pgfpathmoveto{\pgfqpoint{3.938243in}{1.233891in}}%
\pgfpathcurveto{\pgfqpoint{3.947451in}{1.233891in}}{\pgfqpoint{3.956284in}{1.237550in}}{\pgfqpoint{3.962795in}{1.244061in}}%
\pgfpathcurveto{\pgfqpoint{3.969307in}{1.250572in}}{\pgfqpoint{3.972965in}{1.259405in}}{\pgfqpoint{3.972965in}{1.268613in}}%
\pgfpathcurveto{\pgfqpoint{3.972965in}{1.277822in}}{\pgfqpoint{3.969307in}{1.286654in}}{\pgfqpoint{3.962795in}{1.293166in}}%
\pgfpathcurveto{\pgfqpoint{3.956284in}{1.299677in}}{\pgfqpoint{3.947451in}{1.303336in}}{\pgfqpoint{3.938243in}{1.303336in}}%
\pgfpathcurveto{\pgfqpoint{3.929035in}{1.303336in}}{\pgfqpoint{3.920202in}{1.299677in}}{\pgfqpoint{3.913691in}{1.293166in}}%
\pgfpathcurveto{\pgfqpoint{3.907179in}{1.286654in}}{\pgfqpoint{3.903521in}{1.277822in}}{\pgfqpoint{3.903521in}{1.268613in}}%
\pgfpathcurveto{\pgfqpoint{3.903521in}{1.259405in}}{\pgfqpoint{3.907179in}{1.250572in}}{\pgfqpoint{3.913691in}{1.244061in}}%
\pgfpathcurveto{\pgfqpoint{3.920202in}{1.237550in}}{\pgfqpoint{3.929035in}{1.233891in}}{\pgfqpoint{3.938243in}{1.233891in}}%
\pgfpathlineto{\pgfqpoint{3.938243in}{1.233891in}}%
\pgfpathclose%
\pgfusepath{stroke}%
\end{pgfscope}%
\begin{pgfscope}%
\pgfpathrectangle{\pgfqpoint{0.563118in}{0.401080in}}{\pgfqpoint{5.037500in}{4.004000in}}%
\pgfusepath{clip}%
\pgfsetbuttcap%
\pgfsetroundjoin%
\pgfsetlinewidth{1.003750pt}%
\definecolor{currentstroke}{rgb}{0.690196,0.188235,0.376471}%
\pgfsetstrokecolor{currentstroke}%
\pgfsetdash{}{0pt}%
\pgfpathmoveto{\pgfqpoint{4.139743in}{0.866858in}}%
\pgfpathcurveto{\pgfqpoint{4.148951in}{0.866858in}}{\pgfqpoint{4.157784in}{0.870516in}}{\pgfqpoint{4.164295in}{0.877028in}}%
\pgfpathcurveto{\pgfqpoint{4.170807in}{0.883539in}}{\pgfqpoint{4.174465in}{0.892372in}}{\pgfqpoint{4.174465in}{0.901580in}}%
\pgfpathcurveto{\pgfqpoint{4.174465in}{0.910788in}}{\pgfqpoint{4.170807in}{0.919621in}}{\pgfqpoint{4.164295in}{0.926132in}}%
\pgfpathcurveto{\pgfqpoint{4.157784in}{0.932644in}}{\pgfqpoint{4.148951in}{0.936302in}}{\pgfqpoint{4.139743in}{0.936302in}}%
\pgfpathcurveto{\pgfqpoint{4.130535in}{0.936302in}}{\pgfqpoint{4.121702in}{0.932644in}}{\pgfqpoint{4.115191in}{0.926132in}}%
\pgfpathcurveto{\pgfqpoint{4.108679in}{0.919621in}}{\pgfqpoint{4.105021in}{0.910788in}}{\pgfqpoint{4.105021in}{0.901580in}}%
\pgfpathcurveto{\pgfqpoint{4.105021in}{0.892372in}}{\pgfqpoint{4.108679in}{0.883539in}}{\pgfqpoint{4.115191in}{0.877028in}}%
\pgfpathcurveto{\pgfqpoint{4.121702in}{0.870516in}}{\pgfqpoint{4.130535in}{0.866858in}}{\pgfqpoint{4.139743in}{0.866858in}}%
\pgfpathlineto{\pgfqpoint{4.139743in}{0.866858in}}%
\pgfpathclose%
\pgfusepath{stroke}%
\end{pgfscope}%
\begin{pgfscope}%
\pgfpathrectangle{\pgfqpoint{0.563118in}{0.401080in}}{\pgfqpoint{5.037500in}{4.004000in}}%
\pgfusepath{clip}%
\pgfsetbuttcap%
\pgfsetroundjoin%
\pgfsetlinewidth{1.003750pt}%
\definecolor{currentstroke}{rgb}{0.690196,0.188235,0.376471}%
\pgfsetstrokecolor{currentstroke}%
\pgfsetdash{}{0pt}%
\pgfpathmoveto{\pgfqpoint{4.391618in}{2.268258in}}%
\pgfpathcurveto{\pgfqpoint{4.400826in}{2.268258in}}{\pgfqpoint{4.409659in}{2.271916in}}{\pgfqpoint{4.416170in}{2.278428in}}%
\pgfpathcurveto{\pgfqpoint{4.422682in}{2.284939in}}{\pgfqpoint{4.426340in}{2.293772in}}{\pgfqpoint{4.426340in}{2.302980in}}%
\pgfpathcurveto{\pgfqpoint{4.426340in}{2.312188in}}{\pgfqpoint{4.422682in}{2.321021in}}{\pgfqpoint{4.416170in}{2.327532in}}%
\pgfpathcurveto{\pgfqpoint{4.409659in}{2.334044in}}{\pgfqpoint{4.400826in}{2.337702in}}{\pgfqpoint{4.391618in}{2.337702in}}%
\pgfpathcurveto{\pgfqpoint{4.382410in}{2.337702in}}{\pgfqpoint{4.373577in}{2.334044in}}{\pgfqpoint{4.367066in}{2.327532in}}%
\pgfpathcurveto{\pgfqpoint{4.360554in}{2.321021in}}{\pgfqpoint{4.356896in}{2.312188in}}{\pgfqpoint{4.356896in}{2.302980in}}%
\pgfpathcurveto{\pgfqpoint{4.356896in}{2.293772in}}{\pgfqpoint{4.360554in}{2.284939in}}{\pgfqpoint{4.367066in}{2.278428in}}%
\pgfpathcurveto{\pgfqpoint{4.373577in}{2.271916in}}{\pgfqpoint{4.382410in}{2.268258in}}{\pgfqpoint{4.391618in}{2.268258in}}%
\pgfpathlineto{\pgfqpoint{4.391618in}{2.268258in}}%
\pgfpathclose%
\pgfusepath{stroke}%
\end{pgfscope}%
\begin{pgfscope}%
\pgfpathrectangle{\pgfqpoint{0.563118in}{0.401080in}}{\pgfqpoint{5.037500in}{4.004000in}}%
\pgfusepath{clip}%
\pgfsetbuttcap%
\pgfsetroundjoin%
\definecolor{currentfill}{rgb}{0.690196,0.188235,0.376471}%
\pgfsetfillcolor{currentfill}%
\pgfsetlinewidth{1.003750pt}%
\definecolor{currentstroke}{rgb}{0.690196,0.188235,0.376471}%
\pgfsetstrokecolor{currentstroke}%
\pgfsetdash{}{0pt}%
\pgfsys@defobject{currentmarker}{\pgfqpoint{-0.034722in}{-0.034722in}}{\pgfqpoint{0.034722in}{0.034722in}}{%
\pgfpathmoveto{\pgfqpoint{0.000000in}{-0.034722in}}%
\pgfpathcurveto{\pgfqpoint{0.009208in}{-0.034722in}}{\pgfqpoint{0.018041in}{-0.031064in}}{\pgfqpoint{0.024552in}{-0.024552in}}%
\pgfpathcurveto{\pgfqpoint{0.031064in}{-0.018041in}}{\pgfqpoint{0.034722in}{-0.009208in}}{\pgfqpoint{0.034722in}{0.000000in}}%
\pgfpathcurveto{\pgfqpoint{0.034722in}{0.009208in}}{\pgfqpoint{0.031064in}{0.018041in}}{\pgfqpoint{0.024552in}{0.024552in}}%
\pgfpathcurveto{\pgfqpoint{0.018041in}{0.031064in}}{\pgfqpoint{0.009208in}{0.034722in}}{\pgfqpoint{0.000000in}{0.034722in}}%
\pgfpathcurveto{\pgfqpoint{-0.009208in}{0.034722in}}{\pgfqpoint{-0.018041in}{0.031064in}}{\pgfqpoint{-0.024552in}{0.024552in}}%
\pgfpathcurveto{\pgfqpoint{-0.031064in}{0.018041in}}{\pgfqpoint{-0.034722in}{0.009208in}}{\pgfqpoint{-0.034722in}{0.000000in}}%
\pgfpathcurveto{\pgfqpoint{-0.034722in}{-0.009208in}}{\pgfqpoint{-0.031064in}{-0.018041in}}{\pgfqpoint{-0.024552in}{-0.024552in}}%
\pgfpathcurveto{\pgfqpoint{-0.018041in}{-0.031064in}}{\pgfqpoint{-0.009208in}{-0.034722in}}{\pgfqpoint{0.000000in}{-0.034722in}}%
\pgfpathlineto{\pgfqpoint{0.000000in}{-0.034722in}}%
\pgfpathclose%
\pgfusepath{stroke,fill}%
}%
\begin{pgfscope}%
\pgfsys@transformshift{0.563118in}{4.405080in}%
\pgfsys@useobject{currentmarker}{}%
\end{pgfscope}%
\end{pgfscope}%
\begin{pgfscope}%
\pgfsetrectcap%
\pgfsetmiterjoin%
\pgfsetlinewidth{0.501875pt}%
\definecolor{currentstroke}{rgb}{0.000000,0.000000,0.000000}%
\pgfsetstrokecolor{currentstroke}%
\pgfsetdash{}{0pt}%
\pgfpathmoveto{\pgfqpoint{0.563118in}{0.401080in}}%
\pgfpathlineto{\pgfqpoint{0.563118in}{4.405080in}}%
\pgfusepath{stroke}%
\end{pgfscope}%
\begin{pgfscope}%
\pgfsetrectcap%
\pgfsetmiterjoin%
\pgfsetlinewidth{0.501875pt}%
\definecolor{currentstroke}{rgb}{0.000000,0.000000,0.000000}%
\pgfsetstrokecolor{currentstroke}%
\pgfsetdash{}{0pt}%
\pgfpathmoveto{\pgfqpoint{5.600618in}{0.401080in}}%
\pgfpathlineto{\pgfqpoint{5.600618in}{4.405080in}}%
\pgfusepath{stroke}%
\end{pgfscope}%
\begin{pgfscope}%
\pgfsetrectcap%
\pgfsetmiterjoin%
\pgfsetlinewidth{0.501875pt}%
\definecolor{currentstroke}{rgb}{0.000000,0.000000,0.000000}%
\pgfsetstrokecolor{currentstroke}%
\pgfsetdash{}{0pt}%
\pgfpathmoveto{\pgfqpoint{0.563118in}{0.401080in}}%
\pgfpathlineto{\pgfqpoint{5.600618in}{0.401080in}}%
\pgfusepath{stroke}%
\end{pgfscope}%
\begin{pgfscope}%
\pgfsetrectcap%
\pgfsetmiterjoin%
\pgfsetlinewidth{0.501875pt}%
\definecolor{currentstroke}{rgb}{0.000000,0.000000,0.000000}%
\pgfsetstrokecolor{currentstroke}%
\pgfsetdash{}{0pt}%
\pgfpathmoveto{\pgfqpoint{0.563118in}{4.405080in}}%
\pgfpathlineto{\pgfqpoint{5.600618in}{4.405080in}}%
\pgfusepath{stroke}%
\end{pgfscope}%
\end{pgfpicture}%
\makeatother%
\endgroup%

    \caption{A realisation of a simple process.}
    \label{fig:SimpleProcessRealisation}
\end{figure}

\begin{defn}{\Ito{} Integral for Simple Processes}{ItoIntegralForSimpleProcesses}
    Suppose that $F$ is a simple process of the form \cref{eq:SimpleProcessDef}, with $\int_0^t \E F_s^2 \di s < \infty$. Then, the \emph{\Ito{} integral} of $F$ with respect to $W$ over $[0, r]$ is defined as
    \begin{equation}\label{eq:ItoIntegralSimpleDef}
        \int_0^r F_s \di W_s \coloneqq \sum_{k=0}^{n-1} \xi_k (W_{s_{k+1} \wedge r} - W_{s_k \wedge r}), \quad r \in [0, t].
    \end{equation}
\end{defn}

\begin{thm}{Properties of the \Ito{} Integral}{PropertiesOfTheItoIntegral}
    Assume that the \Ito{} integral is well-defined. Then, for $\alpha, \beta \in \R$ and $r \in [0, t]$:
    \begin{enumerate}
        \item \emph{(Linearity):} $\int_0^r [\alpha F_s + \beta G_s] \di W_s = \alpha \int_0^r F_s \di W_s + \beta \int_0^r G_s \di W_s$.
        \item \emph{(Zero mean):} $\E \int_0^r F_s \di W_s = 0$.
        \item \emph{(Isometry):} $\E \left[ \int_0^r F_s \di W_s \times \int_0^r G_s \di W_s \right] = \int_0^r \E [F_s G_s] \di s$.
        \item \emph{(Martingale):} $(\int_0^r F_s \di W_s, r \in [0, t])$ is an $L^2$-martingale.
    \end{enumerate}
\end{thm}

\begin{defn}{Canonical Approximation}{CanonicalApproximation}
    For any $\cF$-adapted process $F \coloneqq (F_s, s \in [0, t])$, we define its \emph{canonical approximation} as
    \begin{equation}\label{eq:CanonicalApproximationDef}
        F_s^{(n)} \coloneqq F_0 \I_{\{0\}}(s) + \sum_{k=0}^{n-1} F_{s_k} \I_{(s_k, s_{k+1}]}(s), \quad s \in [0, t].
    \end{equation}
\end{defn}

\begin{thm}{Properties of Left-continuous Adapted Processes}{PropertiesOfLeftContinuousAdaptedProcesses}
    Suppose that $F \coloneqq (F_s, s \in [0, t])$ is a process on $(\Omega, \cH, \Pm)$ that is adapted to the filtration $\cF$ and is left-continuous on $(0, t]$. Then, the following hold:
    \begin{enumerate}
        \item \emph{(Predictable):} The process $F$ is $\cF$-predictable (i.e., $\cF^P$-measurable).
        \item \emph{(Progressively Measurable):} For each $t$, the mapping $F : [0, t] \times \Omega \to \R$ is jointly measurable with respect to $\cB_{[0, t]} \otimes \cF_t$.
        \item \emph{(Indistinguishable):} If $\widetilde{F}$ is a modification of $F$ and almost surely left-continuous, then $\Pm(\cap_{s \le t} \{F_s = \widetilde{F}_s\}) = 1$.
    \end{enumerate}
\end{thm}

\begin{defn}{The Class of Processes $\mathcal{P}_t$}{}
Let $\mathcal{P}_t$ be the class of stochastic processes $F \coloneqq (F_s, s \in [0, t])$ that are left-continuous on $(0, t]$, $\cF$-adapted, and satisfy the norm condition:
\begin{equation*}
\|F\|_{\mathcal{P}_t}^2 \coloneqq \E \int_0^t F_s^2 \di s < \infty.
\end{equation*}
\end{defn}

\begin{thm}{Approximating the Class $\mathcal{P}_t$ via Simple Processes}{ApproximatingClassPtViaSimpleProcesses}
    Suppose that $F \in \mathcal{P}_t$. Then, there exists a sequence $(F^{(n)})$ of simple processes of the form \cref{eq:SimpleProcessDef} such that $\mathrm{mesh}(\Pi_n) \to 0$ implies that $F_s^{(n)} \xrightarrow{\text{a.s.}} F_s$ for all $s \in [0, t]$ and
    \begin{equation}\label{eq:L2ConvergenceSimple}
        \lim_{n \to \infty} \|F^{(n)} - F\|_{\mathcal{P}_t} = 0.
    \end{equation}
\end{thm}

\begin{prop}{Deterministic Integrands and Gaussian Processes}{DeterministicIntegrandsAndGaussianProcesses}
    Let $f : \R_+ \times \R_+ \to \R$ be left-continuous in the first argument and such that $\|f(\cdot, t)\|_{\mathcal{P}_t}^2 = \int_0^t f^2(s, t) \di s < \infty$. Then, the stochastic integral process
    \begin{equation}\label{eq:GaussianIntegralProcess}
        Z_t := \int_0^t f(s, t) \di W_s, \quad t \ge 0,
    \end{equation}
    is a zero-mean Gaussian process with covariance function
    \begin{equation}\label{eq:GaussianCovariance}
        \Cov(Z_{t_1}, Z_{t_2}) = \int_0^{t_1 \wedge t_2} f(s, t_1) f(s, t_2) \di s, \quad t_1, t_2 \ge 0.
    \end{equation}
\end{prop}

\begin{equation}\label{eq:UniformIntegrabilityCondition}
    \lim_{m \to \infty} \sup_{s \in [0, t]} \E F_s^2 \I_{\{F_s^2 > m\}} = 0.
\end{equation}

\begin{lem}{Quadratic Variation Integral}{QuadraticVariationIntegral}
    Suppose that $F \in \mathcal{P}_t$ satisfies the condition \cref{eq:UniformIntegrabilityCondition}. Then, as $\mathrm{mesh}(\Pi_n) \to 0$:
    \begin{equation}\label{eq:QuadraticVariationLimit}
        \sum_{k=0}^{n-1} F_{s_k} (W_{s_{k+1}} - W_{s_k})^2 \xrightarrow{L^2} \int_0^t F_s \di s.
    \end{equation}
\end{lem}

\begin{thm}{Enlarging the Class $\scP_t$ via Localization}{EnlargingClassPt}
Let $F := (F_s, s \in [0, t])$ be an $\scF$-adapted process, left-continuous on $(0, t]$, and such that (7.34) holds. Then, there exists a localizing sequence $(T_n)$ such that the stopped process $(\int_0^{r \wedge T_n} F_s \di W_s, r \in [0, t])$ is a square-integrable martingale. In addition, the \Ito\ integral of $F$ with respect to $W$ on $[0, t]$ is defined as the almost sure limit:
\begin{equation}
\int_0^t F_s \di W_s \coloneqq \lim_{n \to \infty} \int_0^{t \wedge T_n} F_s \di W_s.
\label{eq:LocalizedIntegralLimit}
\end{equation}
\end{thm}

\begin{prop}{Continuous Modification and Stopping Times}{ContinuousModificationStopping}
If $F \in \scP_t$ and $T \le t$ is a stopping time with respect to $\scF$, then almost surely:
\begin{equation}
\int_0^T F_s \di W_s = \int_0^t \I_{[0, T]}(s) F_s \di W_s.
\label{eq:StoppedIntegralIdentity}
\end{equation}
\end{prop}

\begin{thm}{Continuous Modification of the \Ito\ Integral}{ContinuousModificationIto}
Suppose that $F \in \scP_t$ and $(F^{(n)})$ is a sequence of simple processes satisfying (7.17). Denote the limit $\int_0^r F_s^{(n)} \di W_s =: Z_r^{(n)} \xrightarrow{L^2} Z_r$ as the \Ito\ integral on $[0, r]$ for $r \le t$. Then, there exists a modification $\widetilde{Z}$ of $(Z_r, r \in [0, t])$ such that the mapping $r \mapsto \widetilde{Z}_r$ is almost surely continuous. Moreover, this modification is unique up to indistinguishability.
\end{thm}

\begin{thm}{Existence of the \Ito\ Integral with Integrand $F \in \scP_t$}{ExistenceItoIntegral}
Suppose that $F \in \scP_t$ and $(F^{(n)})$ is a sequence of simple processes satisfying (7.17). Then, the following limit exists in $L^2$ norm:
\begin{equation}
\int_0^t F_s \di W_s \coloneqq \lim_{n \to \infty} \int_0^t F_s^{(n)} \di W_s, \quad r \in [0, t],
\label{eq:ItoIntegralL2Limit}
\end{equation}
and it defines the \textit{\Ito\ integral} of $F$ with respect to $W$ on $[0, t]$.
\end{thm}

\subsection{\Ito{} Calculus}
\begin{thm}{\Ito{} Formula for the Wiener Process}{}
Let $W$ be a Wiener process and $f : \R \to \R$ be twice continuously differentiable with derivatives $f'$ and $f''$. Then, almost surely
\begin{equation*}
f(W_t) - f(W_0) = \int_0^t f'(W_s) \di W_s + \frac{1}{2} \int_0^t f''(W_s) \di s.
\end{equation*}
\end{thm}

\begin{defn}{Integral With Respect To An \Ito\ Process}{IntegralItoProcess}
Let $X$ be an \Ito$(\mu, \sigma)$ process and $F$ a left-continuous $\scF$-adapted process such that almost surely $\int_{0}^{t}(|F_s\mu_s| + (F_s\sigma_s)^2)\di s < \infty$ for all $t$. Then, the integral of $F$ with respect to $X$ is defined as the \Ito$(F\mu, F\sigma)$ process with:
\begin{equation}
\int_{0}^{t} F_s \di X_s \coloneqq \int_{0}^{t} F_s \mu_s \di s + \int_{0}^{t} F_s \sigma_s \di W_s.
\end{equation}
\end{defn}

\begin{defn}{Quadratic Variation On $[0, t]$}{QuadraticVariation}
Let $\Pi_n \coloneqq \{s_k, k = 0, \dots, n\}$ be a segmentation of $[0, t]$ and let $(X^{(n)})$ be a sequence of simple processes with almost sure limit $X_s \coloneqq \lim_n X^{(n)}_s$ for all $s \in [0, t]$ as $\mathrm{mesh}(\Pi_n) \to 0$. The \textit{quadratic variation} of $X \coloneqq (X_s, s \in [0, t])$ over $[0, t]$, denoted as $\langle X \rangle_t$, is defined as the limit in probability of
\begin{equation}
\langle X^{(n)} \rangle_t \coloneqq \sum_{k=0}^{n-1} (X^{(n)}_{s_{k+1}} - X^{(n)}_{s_k})^2.
\end{equation}
\end{defn}

\begin{thm}{Quadratic Variation Of An \Ito\ Process}{QuadraticVariationIto}
If $X$ is an \Ito$(\mu, \sigma)$ process, then
\begin{equation}
\langle X^{(n)} \rangle_t \xrightarrow{L^2} \langle X \rangle_t = \int_{0}^{t} \sigma_s^2 \di s.
\end{equation}
Furthermore, if $F \in \scP_t$ satisfies the condition (7.19) and $\mathrm{mesh}(\Pi_n) \to 0$, then:
\begin{equation}
\sum_{k=0}^{n-1} F_{s_k} (X_{s_{k+1}} - X_{s_k})^2 \xrightarrow{L^2} \int_{0}^{t} F_s \sigma_s^2 \di s.
\end{equation}
\end{thm}

\begin{thm}{\Ito's Formula For \Ito\ Processes}{ItosFormula}
Let $X$ be an \Ito\ process and let $f$ be twice continuously differentiable. Then,
\begin{equation}
f(X_t) = f(X_0) + \int_{0}^{t} f'(X_s) \di X_s + \frac{1}{2} \int_{0}^{t} f''(X_s) \di \langle X \rangle_s.
\label{eq:ItoFormula}
\end{equation}
\end{thm}

\begin{defn}{Covariation Of Two \Ito\ Processes}{CovariationIto}
Let $X$ and $Y$ be two \Ito\ processes adapted to the same filtration $\scF$. Then, the \textit{covariation} between $X$ and $Y$ is defined as the quantity:
\begin{equation}
\langle X, Y \rangle_t \coloneqq \frac{\langle X + Y \rangle_t - \langle X \rangle_t - \langle Y \rangle_t}{2}.
\label{eq:CovariationDef}
\end{equation}
\end{defn}

\begin{thm}{\Ito's Formula In $\R^d$}{ItoFormulaMultiDim}
Let $(X_t)$ be a $d$-dimensional \Ito\ process, and let $f : \R^d \to \R$ be twice continuously differentiable in all variables. Then,
\begin{equation}
\di f(X_t) = \sum_{i=1}^{d} \partial_i f(X_t) \di X_{t,i} + \frac{1}{2} \sum_{i=1}^{d} \sum_{j=1}^{d} \partial_{ij} f(X_t) \di \langle X_{\cdot, i}, X_{\cdot, j} \rangle_t.
\label{eq:DifferentialItoMulti}
\end{equation}
\end{thm}

\begin{thm}{\Ito\ Processes And Martingales}{ItoMartingaleDrift}
An \Ito\ process in $\R^d$ is a local martingale if and only if its drift is 0.
\end{thm}

\clearpage
\subsection{Stochastic Differential Equations}\label{sec:StochasticDifferentialEquations}
Stochastic differential equations are based on the same principle as ordinary differential equations, relating an unknown function to its derivatives, but with the additional feature that part of the unknown function is driven by randomness.

\begin{defn}{Stochastic Differential Equation}{StochasticDifferentialEquation}
    A \textit{stochastic differential equation} (SDE) for a stochastic process $(X_t, t \ge 0)$ is an expression of the form
    \begin{equation}
        \di X_t = a(X_t, t) \di t + b(X_t, t) \di W_t,
        \label{eq:StochasticDifferentialEquationForm}
    \end{equation}
    where $(W_t, t \ge 0)$ is a Wiener process and $a$ and $b$ are deterministic functions, called the \textit{drift} and \textit{diffusion} functions. The resulting process $(X_t, t \ge 0)$ is called an \textit{(\Ito) diffusion}.
\end{defn}

Consider the autonomous SDE
\begin{equation}
    \di X_t = a(X_t) \di t + b(X_t) \di W_t.
    \label{eq:AutonomousSDE}
\end{equation}
As already mentioned, the meaning of this SDE is that $X$ is a stochastic process that satisfies the integral equation
\begin{equation}
    X_t = X_0 + \int_0^t a(X_s) \di s + \int_0^t b(X_s) \di W_s.
    \label{eq:StochasticIntegralEquation}
\end{equation}
However, it is not clear that such a process indeed exists. We show in this section that under mild conditions on $a$ and $b$ we can obtain a unique solution $X$ of the SDE via a method of successive approximations, similar to \textit{Picard iteration} for ordinary differential equations.

The condition that is placed on $a$ and $b$ is that they be \textit{Lipschitz continuous}; that is, for all $x \in \R$ there exists a constant $c$ such that
\begin{equation}
    |a(y) - a(x)| + |b(y) - b(x)| \le c|y - x|.
    \label{eq:LipschitzCondition1D}
\end{equation}

The successive approximation idea is very simple. Suppose that $X_0 = x$. Define $X_t^{(0)} \coloneqq x$ for all $t$, and let
\begin{equation}
    X_t^{(n+1)} \coloneqq x + \int_0^t a(X_s^{(n)}) \di s + \int_0^t b(X_s^{(n)}) \di W_s
    \label{eq:SuccessiveApproximation}
\end{equation}
for $n \in \N$. The following is the main existence result for SDEs:

\begin{thm}{Existence of SDE Solutions}{ExistenceOfSDESolutions}
    For any $t \in [0, v]$, the sequence $(X^{(n)})$ in \eqref{eq:SuccessiveApproximation} converges almost surely in the uniform norm on $[0, t]$ to a process $X$ that satisfies the integral equation \eqref{eq:IntegralEquation}.
\end{thm}

\begin{lem}{Integral Bound}{IntegralBound}
    Let $X^{(n)}, n = 1, 2, \dots$ be defined in \eqref{eq:SuccessiveApproximation}, with $X^{(0)} = x$. Then, for every $v > 0$ there exists a constant $\alpha$ such that
    \begin{equation*}
        \E \sup_{s \le t} \left( X_s^{(n+1)} - X_s^{(n)} \right)^2 \le \alpha \| X^{(n)} - X^{(n-1)} \|_{\mathscr{P}_t}^2, \qquad t \le v.
    \end{equation*}
\end{lem}

\begin{defn}{\Ito \; Stratonovich Integral}{ItoStratonovichIntegral}
    Let $X$ and $Y$ be \Ito processes such that $\int_0^t Y_s \di X_s$ and $\int_0^t X_s \di Y_s$ are well-defined stochastic integrals. Then, the \textit{\Ito--Stratonovich integral} of $Y$ with respect to $X$ is defined as
    \begin{equation}
        \int_0^t Y_s \circ \di X_s \coloneqq \int_0^t Y_s \di X_s + \frac{1}{2}\langle X, Y \rangle_t.
        \label{eq:ItoStratonovichIntegral}
    \end{equation}
\end{defn}

\begin{prop}{\Ito \; Stratonovich Formula}{ItoStratonovichFormula}
    Let $(X_t)$ be a $d$-dimensional \Ito process, and let $f : \R^d \to \R$ be three times continuously differentiable in all variables. Then,
    \begin{equation}
        \di f(X_t) = \sum_{i=1}^d \partial_i f(X_t) \circ \di X_{t,i}.
        \label{eq:ItoStratonovichFormula}
    \end{equation}
\end{prop}
\begin{equation}\label{eq:ChangeOfMeasureDef}
    \widetilde{\Pm}^t(A) := \E M_t \I_A, \quad A \in \cF_t.
\end{equation}
\begin{equation}\label{eq:ExponentialMartingaleDef}
    M_t = \exp\left(\int_0^t \mu_s \di W_s - \frac{1}{2} \int_0^t \mu_s^2 \di s\right) = \e^{Z_t - \frac{1}{2}\langle Z \rangle_t}, \quad t \ge 0,
\end{equation}
\begin{equation}\label{eq:KrylovCondition}
    \lim_{\varepsilon \downarrow 0} \varepsilon \ln \E \exp((1 - \varepsilon)\langle Z \rangle_t / 2) = 0.
\end{equation}
\begin{thm}{Girsanov's Change of Measure}{GirsanovsChangeOfMeasure}
    Suppose that $\mu$ satisfies the Krylov condition \cref{eq:KrylovCondition} and $M$ is the exponential martingale \cref{eq:ExponentialMartingaleDef}. Then, under the change of measure \cref{eq:ChangeOfMeasureDef} induced by $M$, the \Ito{} process
    \begin{equation*}
        \widetilde{W}_t := -\int_0^t \mu_s \di s + W_t, \quad t \ge 0
    \end{equation*}
    is a Wiener process.
\end{thm}

\begin{Pcode}{GBM.py}
N, T = 1000, 1.0
mu, sigma = 0.5, 0.2
dt = T / N
t = np.linspace(0, T, N + 1)

rng = np.random.default_rng(4)
dW = rng.normal(0, np.sqrt(dt), N)
W = np.r_[0, np.cumsum(dW)]

drift = (mu - 0.5 * sigma**2) * t
diffusion = sigma * W
S_t = np.exp(drift + diffusion)
E_St = np.exp(mu * t)
\end{Pcode}

\begin{figure}[H]
    \centering
    %% Creator: Matplotlib, PGF backend
%%
%% To include the figure in your LaTeX document, write
%%   \input{<filename>.pgf}
%%
%% Make sure the required packages are loaded in your preamble
%%   \usepackage{pgf}
%%
%% Also ensure that all the required font packages are loaded; for instance,
%% the lmodern package is sometimes necessary when using math font.
%%   \usepackage{lmodern}
%%
%% Figures using additional raster images can only be included by \input if
%% they are in the same directory as the main LaTeX file. For loading figures
%% from other directories you can use the `import` package
%%   \usepackage{import}
%%
%% and then include the figures with
%%   \import{<path to file>}{<filename>.pgf}
%%
%% Matplotlib used the following preamble
%%   \def\mathdefault#1{#1}
%%   \everymath=\expandafter{\the\everymath\displaystyle}
%%   \IfFileExists{scrextend.sty}{
%%     \usepackage[fontsize=10.000000pt]{scrextend}
%%   }{
%%     \renewcommand{\normalsize}{\fontsize{10.000000}{12.000000}\selectfont}
%%     \normalsize
%%   }
%%   
%%   \makeatletter\@ifpackageloaded{underscore}{}{\usepackage[strings]{underscore}}\makeatother
%%
\begingroup%
\makeatletter%
\begin{pgfpicture}%
\pgfpathrectangle{\pgfpointorigin}{\pgfqpoint{5.631328in}{4.455080in}}%
\pgfusepath{use as bounding box, clip}%
\begin{pgfscope}%
\pgfsetbuttcap%
\pgfsetmiterjoin%
\definecolor{currentfill}{rgb}{1.000000,1.000000,1.000000}%
\pgfsetfillcolor{currentfill}%
\pgfsetlinewidth{0.000000pt}%
\definecolor{currentstroke}{rgb}{1.000000,1.000000,1.000000}%
\pgfsetstrokecolor{currentstroke}%
\pgfsetdash{}{0pt}%
\pgfpathmoveto{\pgfqpoint{0.000000in}{0.000000in}}%
\pgfpathlineto{\pgfqpoint{5.631328in}{0.000000in}}%
\pgfpathlineto{\pgfqpoint{5.631328in}{4.455080in}}%
\pgfpathlineto{\pgfqpoint{0.000000in}{4.455080in}}%
\pgfpathlineto{\pgfqpoint{0.000000in}{0.000000in}}%
\pgfpathclose%
\pgfusepath{fill}%
\end{pgfscope}%
\begin{pgfscope}%
\pgfsetbuttcap%
\pgfsetmiterjoin%
\definecolor{currentfill}{rgb}{1.000000,1.000000,1.000000}%
\pgfsetfillcolor{currentfill}%
\pgfsetlinewidth{0.000000pt}%
\definecolor{currentstroke}{rgb}{0.000000,0.000000,0.000000}%
\pgfsetstrokecolor{currentstroke}%
\pgfsetstrokeopacity{0.000000}%
\pgfsetdash{}{0pt}%
\pgfpathmoveto{\pgfqpoint{0.455093in}{0.401080in}}%
\pgfpathlineto{\pgfqpoint{5.492593in}{0.401080in}}%
\pgfpathlineto{\pgfqpoint{5.492593in}{4.405080in}}%
\pgfpathlineto{\pgfqpoint{0.455093in}{4.405080in}}%
\pgfpathlineto{\pgfqpoint{0.455093in}{0.401080in}}%
\pgfpathclose%
\pgfusepath{fill}%
\end{pgfscope}%
\begin{pgfscope}%
\pgfsetbuttcap%
\pgfsetroundjoin%
\definecolor{currentfill}{rgb}{0.000000,0.000000,0.000000}%
\pgfsetfillcolor{currentfill}%
\pgfsetlinewidth{0.501875pt}%
\definecolor{currentstroke}{rgb}{0.000000,0.000000,0.000000}%
\pgfsetstrokecolor{currentstroke}%
\pgfsetdash{}{0pt}%
\pgfsys@defobject{currentmarker}{\pgfqpoint{0.000000in}{0.000000in}}{\pgfqpoint{0.000000in}{0.041667in}}{%
\pgfpathmoveto{\pgfqpoint{0.000000in}{0.000000in}}%
\pgfpathlineto{\pgfqpoint{0.000000in}{0.041667in}}%
\pgfusepath{stroke,fill}%
}%
\begin{pgfscope}%
\pgfsys@transformshift{0.455093in}{0.401080in}%
\pgfsys@useobject{currentmarker}{}%
\end{pgfscope}%
\end{pgfscope}%
\begin{pgfscope}%
\pgfsetbuttcap%
\pgfsetroundjoin%
\definecolor{currentfill}{rgb}{0.000000,0.000000,0.000000}%
\pgfsetfillcolor{currentfill}%
\pgfsetlinewidth{0.501875pt}%
\definecolor{currentstroke}{rgb}{0.000000,0.000000,0.000000}%
\pgfsetstrokecolor{currentstroke}%
\pgfsetdash{}{0pt}%
\pgfsys@defobject{currentmarker}{\pgfqpoint{0.000000in}{-0.041667in}}{\pgfqpoint{0.000000in}{0.000000in}}{%
\pgfpathmoveto{\pgfqpoint{0.000000in}{0.000000in}}%
\pgfpathlineto{\pgfqpoint{0.000000in}{-0.041667in}}%
\pgfusepath{stroke,fill}%
}%
\begin{pgfscope}%
\pgfsys@transformshift{0.455093in}{4.405080in}%
\pgfsys@useobject{currentmarker}{}%
\end{pgfscope}%
\end{pgfscope}%
\begin{pgfscope}%
\definecolor{textcolor}{rgb}{0.000000,0.000000,0.000000}%
\pgfsetstrokecolor{textcolor}%
\pgfsetfillcolor{textcolor}%
\pgftext[x=0.455093in,y=0.352469in,,top]{\color{textcolor}{\rmfamily\fontsize{10.000000}{12.000000}\selectfont\catcode`\^=\active\def^{\ifmmode\sp\else\^{}\fi}\catcode`\%=\active\def%{\%}$\mathdefault{0.0}$}}%
\end{pgfscope}%
\begin{pgfscope}%
\pgfsetbuttcap%
\pgfsetroundjoin%
\definecolor{currentfill}{rgb}{0.000000,0.000000,0.000000}%
\pgfsetfillcolor{currentfill}%
\pgfsetlinewidth{0.501875pt}%
\definecolor{currentstroke}{rgb}{0.000000,0.000000,0.000000}%
\pgfsetstrokecolor{currentstroke}%
\pgfsetdash{}{0pt}%
\pgfsys@defobject{currentmarker}{\pgfqpoint{0.000000in}{0.000000in}}{\pgfqpoint{0.000000in}{0.041667in}}{%
\pgfpathmoveto{\pgfqpoint{0.000000in}{0.000000in}}%
\pgfpathlineto{\pgfqpoint{0.000000in}{0.041667in}}%
\pgfusepath{stroke,fill}%
}%
\begin{pgfscope}%
\pgfsys@transformshift{1.462593in}{0.401080in}%
\pgfsys@useobject{currentmarker}{}%
\end{pgfscope}%
\end{pgfscope}%
\begin{pgfscope}%
\pgfsetbuttcap%
\pgfsetroundjoin%
\definecolor{currentfill}{rgb}{0.000000,0.000000,0.000000}%
\pgfsetfillcolor{currentfill}%
\pgfsetlinewidth{0.501875pt}%
\definecolor{currentstroke}{rgb}{0.000000,0.000000,0.000000}%
\pgfsetstrokecolor{currentstroke}%
\pgfsetdash{}{0pt}%
\pgfsys@defobject{currentmarker}{\pgfqpoint{0.000000in}{-0.041667in}}{\pgfqpoint{0.000000in}{0.000000in}}{%
\pgfpathmoveto{\pgfqpoint{0.000000in}{0.000000in}}%
\pgfpathlineto{\pgfqpoint{0.000000in}{-0.041667in}}%
\pgfusepath{stroke,fill}%
}%
\begin{pgfscope}%
\pgfsys@transformshift{1.462593in}{4.405080in}%
\pgfsys@useobject{currentmarker}{}%
\end{pgfscope}%
\end{pgfscope}%
\begin{pgfscope}%
\definecolor{textcolor}{rgb}{0.000000,0.000000,0.000000}%
\pgfsetstrokecolor{textcolor}%
\pgfsetfillcolor{textcolor}%
\pgftext[x=1.462593in,y=0.352469in,,top]{\color{textcolor}{\rmfamily\fontsize{10.000000}{12.000000}\selectfont\catcode`\^=\active\def^{\ifmmode\sp\else\^{}\fi}\catcode`\%=\active\def%{\%}$\mathdefault{0.2}$}}%
\end{pgfscope}%
\begin{pgfscope}%
\pgfsetbuttcap%
\pgfsetroundjoin%
\definecolor{currentfill}{rgb}{0.000000,0.000000,0.000000}%
\pgfsetfillcolor{currentfill}%
\pgfsetlinewidth{0.501875pt}%
\definecolor{currentstroke}{rgb}{0.000000,0.000000,0.000000}%
\pgfsetstrokecolor{currentstroke}%
\pgfsetdash{}{0pt}%
\pgfsys@defobject{currentmarker}{\pgfqpoint{0.000000in}{0.000000in}}{\pgfqpoint{0.000000in}{0.041667in}}{%
\pgfpathmoveto{\pgfqpoint{0.000000in}{0.000000in}}%
\pgfpathlineto{\pgfqpoint{0.000000in}{0.041667in}}%
\pgfusepath{stroke,fill}%
}%
\begin{pgfscope}%
\pgfsys@transformshift{2.470093in}{0.401080in}%
\pgfsys@useobject{currentmarker}{}%
\end{pgfscope}%
\end{pgfscope}%
\begin{pgfscope}%
\pgfsetbuttcap%
\pgfsetroundjoin%
\definecolor{currentfill}{rgb}{0.000000,0.000000,0.000000}%
\pgfsetfillcolor{currentfill}%
\pgfsetlinewidth{0.501875pt}%
\definecolor{currentstroke}{rgb}{0.000000,0.000000,0.000000}%
\pgfsetstrokecolor{currentstroke}%
\pgfsetdash{}{0pt}%
\pgfsys@defobject{currentmarker}{\pgfqpoint{0.000000in}{-0.041667in}}{\pgfqpoint{0.000000in}{0.000000in}}{%
\pgfpathmoveto{\pgfqpoint{0.000000in}{0.000000in}}%
\pgfpathlineto{\pgfqpoint{0.000000in}{-0.041667in}}%
\pgfusepath{stroke,fill}%
}%
\begin{pgfscope}%
\pgfsys@transformshift{2.470093in}{4.405080in}%
\pgfsys@useobject{currentmarker}{}%
\end{pgfscope}%
\end{pgfscope}%
\begin{pgfscope}%
\definecolor{textcolor}{rgb}{0.000000,0.000000,0.000000}%
\pgfsetstrokecolor{textcolor}%
\pgfsetfillcolor{textcolor}%
\pgftext[x=2.470093in,y=0.352469in,,top]{\color{textcolor}{\rmfamily\fontsize{10.000000}{12.000000}\selectfont\catcode`\^=\active\def^{\ifmmode\sp\else\^{}\fi}\catcode`\%=\active\def%{\%}$\mathdefault{0.4}$}}%
\end{pgfscope}%
\begin{pgfscope}%
\pgfsetbuttcap%
\pgfsetroundjoin%
\definecolor{currentfill}{rgb}{0.000000,0.000000,0.000000}%
\pgfsetfillcolor{currentfill}%
\pgfsetlinewidth{0.501875pt}%
\definecolor{currentstroke}{rgb}{0.000000,0.000000,0.000000}%
\pgfsetstrokecolor{currentstroke}%
\pgfsetdash{}{0pt}%
\pgfsys@defobject{currentmarker}{\pgfqpoint{0.000000in}{0.000000in}}{\pgfqpoint{0.000000in}{0.041667in}}{%
\pgfpathmoveto{\pgfqpoint{0.000000in}{0.000000in}}%
\pgfpathlineto{\pgfqpoint{0.000000in}{0.041667in}}%
\pgfusepath{stroke,fill}%
}%
\begin{pgfscope}%
\pgfsys@transformshift{3.477593in}{0.401080in}%
\pgfsys@useobject{currentmarker}{}%
\end{pgfscope}%
\end{pgfscope}%
\begin{pgfscope}%
\pgfsetbuttcap%
\pgfsetroundjoin%
\definecolor{currentfill}{rgb}{0.000000,0.000000,0.000000}%
\pgfsetfillcolor{currentfill}%
\pgfsetlinewidth{0.501875pt}%
\definecolor{currentstroke}{rgb}{0.000000,0.000000,0.000000}%
\pgfsetstrokecolor{currentstroke}%
\pgfsetdash{}{0pt}%
\pgfsys@defobject{currentmarker}{\pgfqpoint{0.000000in}{-0.041667in}}{\pgfqpoint{0.000000in}{0.000000in}}{%
\pgfpathmoveto{\pgfqpoint{0.000000in}{0.000000in}}%
\pgfpathlineto{\pgfqpoint{0.000000in}{-0.041667in}}%
\pgfusepath{stroke,fill}%
}%
\begin{pgfscope}%
\pgfsys@transformshift{3.477593in}{4.405080in}%
\pgfsys@useobject{currentmarker}{}%
\end{pgfscope}%
\end{pgfscope}%
\begin{pgfscope}%
\definecolor{textcolor}{rgb}{0.000000,0.000000,0.000000}%
\pgfsetstrokecolor{textcolor}%
\pgfsetfillcolor{textcolor}%
\pgftext[x=3.477593in,y=0.352469in,,top]{\color{textcolor}{\rmfamily\fontsize{10.000000}{12.000000}\selectfont\catcode`\^=\active\def^{\ifmmode\sp\else\^{}\fi}\catcode`\%=\active\def%{\%}$\mathdefault{0.6}$}}%
\end{pgfscope}%
\begin{pgfscope}%
\pgfsetbuttcap%
\pgfsetroundjoin%
\definecolor{currentfill}{rgb}{0.000000,0.000000,0.000000}%
\pgfsetfillcolor{currentfill}%
\pgfsetlinewidth{0.501875pt}%
\definecolor{currentstroke}{rgb}{0.000000,0.000000,0.000000}%
\pgfsetstrokecolor{currentstroke}%
\pgfsetdash{}{0pt}%
\pgfsys@defobject{currentmarker}{\pgfqpoint{0.000000in}{0.000000in}}{\pgfqpoint{0.000000in}{0.041667in}}{%
\pgfpathmoveto{\pgfqpoint{0.000000in}{0.000000in}}%
\pgfpathlineto{\pgfqpoint{0.000000in}{0.041667in}}%
\pgfusepath{stroke,fill}%
}%
\begin{pgfscope}%
\pgfsys@transformshift{4.485093in}{0.401080in}%
\pgfsys@useobject{currentmarker}{}%
\end{pgfscope}%
\end{pgfscope}%
\begin{pgfscope}%
\pgfsetbuttcap%
\pgfsetroundjoin%
\definecolor{currentfill}{rgb}{0.000000,0.000000,0.000000}%
\pgfsetfillcolor{currentfill}%
\pgfsetlinewidth{0.501875pt}%
\definecolor{currentstroke}{rgb}{0.000000,0.000000,0.000000}%
\pgfsetstrokecolor{currentstroke}%
\pgfsetdash{}{0pt}%
\pgfsys@defobject{currentmarker}{\pgfqpoint{0.000000in}{-0.041667in}}{\pgfqpoint{0.000000in}{0.000000in}}{%
\pgfpathmoveto{\pgfqpoint{0.000000in}{0.000000in}}%
\pgfpathlineto{\pgfqpoint{0.000000in}{-0.041667in}}%
\pgfusepath{stroke,fill}%
}%
\begin{pgfscope}%
\pgfsys@transformshift{4.485093in}{4.405080in}%
\pgfsys@useobject{currentmarker}{}%
\end{pgfscope}%
\end{pgfscope}%
\begin{pgfscope}%
\definecolor{textcolor}{rgb}{0.000000,0.000000,0.000000}%
\pgfsetstrokecolor{textcolor}%
\pgfsetfillcolor{textcolor}%
\pgftext[x=4.485093in,y=0.352469in,,top]{\color{textcolor}{\rmfamily\fontsize{10.000000}{12.000000}\selectfont\catcode`\^=\active\def^{\ifmmode\sp\else\^{}\fi}\catcode`\%=\active\def%{\%}$\mathdefault{0.8}$}}%
\end{pgfscope}%
\begin{pgfscope}%
\pgfsetbuttcap%
\pgfsetroundjoin%
\definecolor{currentfill}{rgb}{0.000000,0.000000,0.000000}%
\pgfsetfillcolor{currentfill}%
\pgfsetlinewidth{0.501875pt}%
\definecolor{currentstroke}{rgb}{0.000000,0.000000,0.000000}%
\pgfsetstrokecolor{currentstroke}%
\pgfsetdash{}{0pt}%
\pgfsys@defobject{currentmarker}{\pgfqpoint{0.000000in}{0.000000in}}{\pgfqpoint{0.000000in}{0.041667in}}{%
\pgfpathmoveto{\pgfqpoint{0.000000in}{0.000000in}}%
\pgfpathlineto{\pgfqpoint{0.000000in}{0.041667in}}%
\pgfusepath{stroke,fill}%
}%
\begin{pgfscope}%
\pgfsys@transformshift{5.492593in}{0.401080in}%
\pgfsys@useobject{currentmarker}{}%
\end{pgfscope}%
\end{pgfscope}%
\begin{pgfscope}%
\pgfsetbuttcap%
\pgfsetroundjoin%
\definecolor{currentfill}{rgb}{0.000000,0.000000,0.000000}%
\pgfsetfillcolor{currentfill}%
\pgfsetlinewidth{0.501875pt}%
\definecolor{currentstroke}{rgb}{0.000000,0.000000,0.000000}%
\pgfsetstrokecolor{currentstroke}%
\pgfsetdash{}{0pt}%
\pgfsys@defobject{currentmarker}{\pgfqpoint{0.000000in}{-0.041667in}}{\pgfqpoint{0.000000in}{0.000000in}}{%
\pgfpathmoveto{\pgfqpoint{0.000000in}{0.000000in}}%
\pgfpathlineto{\pgfqpoint{0.000000in}{-0.041667in}}%
\pgfusepath{stroke,fill}%
}%
\begin{pgfscope}%
\pgfsys@transformshift{5.492593in}{4.405080in}%
\pgfsys@useobject{currentmarker}{}%
\end{pgfscope}%
\end{pgfscope}%
\begin{pgfscope}%
\definecolor{textcolor}{rgb}{0.000000,0.000000,0.000000}%
\pgfsetstrokecolor{textcolor}%
\pgfsetfillcolor{textcolor}%
\pgftext[x=5.492593in,y=0.352469in,,top]{\color{textcolor}{\rmfamily\fontsize{10.000000}{12.000000}\selectfont\catcode`\^=\active\def^{\ifmmode\sp\else\^{}\fi}\catcode`\%=\active\def%{\%}$\mathdefault{1.0}$}}%
\end{pgfscope}%
\begin{pgfscope}%
\pgfsetbuttcap%
\pgfsetroundjoin%
\definecolor{currentfill}{rgb}{0.000000,0.000000,0.000000}%
\pgfsetfillcolor{currentfill}%
\pgfsetlinewidth{0.501875pt}%
\definecolor{currentstroke}{rgb}{0.000000,0.000000,0.000000}%
\pgfsetstrokecolor{currentstroke}%
\pgfsetdash{}{0pt}%
\pgfsys@defobject{currentmarker}{\pgfqpoint{0.000000in}{0.000000in}}{\pgfqpoint{0.000000in}{0.020833in}}{%
\pgfpathmoveto{\pgfqpoint{0.000000in}{0.000000in}}%
\pgfpathlineto{\pgfqpoint{0.000000in}{0.020833in}}%
\pgfusepath{stroke,fill}%
}%
\begin{pgfscope}%
\pgfsys@transformshift{0.706968in}{0.401080in}%
\pgfsys@useobject{currentmarker}{}%
\end{pgfscope}%
\end{pgfscope}%
\begin{pgfscope}%
\pgfsetbuttcap%
\pgfsetroundjoin%
\definecolor{currentfill}{rgb}{0.000000,0.000000,0.000000}%
\pgfsetfillcolor{currentfill}%
\pgfsetlinewidth{0.501875pt}%
\definecolor{currentstroke}{rgb}{0.000000,0.000000,0.000000}%
\pgfsetstrokecolor{currentstroke}%
\pgfsetdash{}{0pt}%
\pgfsys@defobject{currentmarker}{\pgfqpoint{0.000000in}{-0.020833in}}{\pgfqpoint{0.000000in}{0.000000in}}{%
\pgfpathmoveto{\pgfqpoint{0.000000in}{0.000000in}}%
\pgfpathlineto{\pgfqpoint{0.000000in}{-0.020833in}}%
\pgfusepath{stroke,fill}%
}%
\begin{pgfscope}%
\pgfsys@transformshift{0.706968in}{4.405080in}%
\pgfsys@useobject{currentmarker}{}%
\end{pgfscope}%
\end{pgfscope}%
\begin{pgfscope}%
\pgfsetbuttcap%
\pgfsetroundjoin%
\definecolor{currentfill}{rgb}{0.000000,0.000000,0.000000}%
\pgfsetfillcolor{currentfill}%
\pgfsetlinewidth{0.501875pt}%
\definecolor{currentstroke}{rgb}{0.000000,0.000000,0.000000}%
\pgfsetstrokecolor{currentstroke}%
\pgfsetdash{}{0pt}%
\pgfsys@defobject{currentmarker}{\pgfqpoint{0.000000in}{0.000000in}}{\pgfqpoint{0.000000in}{0.020833in}}{%
\pgfpathmoveto{\pgfqpoint{0.000000in}{0.000000in}}%
\pgfpathlineto{\pgfqpoint{0.000000in}{0.020833in}}%
\pgfusepath{stroke,fill}%
}%
\begin{pgfscope}%
\pgfsys@transformshift{0.958843in}{0.401080in}%
\pgfsys@useobject{currentmarker}{}%
\end{pgfscope}%
\end{pgfscope}%
\begin{pgfscope}%
\pgfsetbuttcap%
\pgfsetroundjoin%
\definecolor{currentfill}{rgb}{0.000000,0.000000,0.000000}%
\pgfsetfillcolor{currentfill}%
\pgfsetlinewidth{0.501875pt}%
\definecolor{currentstroke}{rgb}{0.000000,0.000000,0.000000}%
\pgfsetstrokecolor{currentstroke}%
\pgfsetdash{}{0pt}%
\pgfsys@defobject{currentmarker}{\pgfqpoint{0.000000in}{-0.020833in}}{\pgfqpoint{0.000000in}{0.000000in}}{%
\pgfpathmoveto{\pgfqpoint{0.000000in}{0.000000in}}%
\pgfpathlineto{\pgfqpoint{0.000000in}{-0.020833in}}%
\pgfusepath{stroke,fill}%
}%
\begin{pgfscope}%
\pgfsys@transformshift{0.958843in}{4.405080in}%
\pgfsys@useobject{currentmarker}{}%
\end{pgfscope}%
\end{pgfscope}%
\begin{pgfscope}%
\pgfsetbuttcap%
\pgfsetroundjoin%
\definecolor{currentfill}{rgb}{0.000000,0.000000,0.000000}%
\pgfsetfillcolor{currentfill}%
\pgfsetlinewidth{0.501875pt}%
\definecolor{currentstroke}{rgb}{0.000000,0.000000,0.000000}%
\pgfsetstrokecolor{currentstroke}%
\pgfsetdash{}{0pt}%
\pgfsys@defobject{currentmarker}{\pgfqpoint{0.000000in}{0.000000in}}{\pgfqpoint{0.000000in}{0.020833in}}{%
\pgfpathmoveto{\pgfqpoint{0.000000in}{0.000000in}}%
\pgfpathlineto{\pgfqpoint{0.000000in}{0.020833in}}%
\pgfusepath{stroke,fill}%
}%
\begin{pgfscope}%
\pgfsys@transformshift{1.210718in}{0.401080in}%
\pgfsys@useobject{currentmarker}{}%
\end{pgfscope}%
\end{pgfscope}%
\begin{pgfscope}%
\pgfsetbuttcap%
\pgfsetroundjoin%
\definecolor{currentfill}{rgb}{0.000000,0.000000,0.000000}%
\pgfsetfillcolor{currentfill}%
\pgfsetlinewidth{0.501875pt}%
\definecolor{currentstroke}{rgb}{0.000000,0.000000,0.000000}%
\pgfsetstrokecolor{currentstroke}%
\pgfsetdash{}{0pt}%
\pgfsys@defobject{currentmarker}{\pgfqpoint{0.000000in}{-0.020833in}}{\pgfqpoint{0.000000in}{0.000000in}}{%
\pgfpathmoveto{\pgfqpoint{0.000000in}{0.000000in}}%
\pgfpathlineto{\pgfqpoint{0.000000in}{-0.020833in}}%
\pgfusepath{stroke,fill}%
}%
\begin{pgfscope}%
\pgfsys@transformshift{1.210718in}{4.405080in}%
\pgfsys@useobject{currentmarker}{}%
\end{pgfscope}%
\end{pgfscope}%
\begin{pgfscope}%
\pgfsetbuttcap%
\pgfsetroundjoin%
\definecolor{currentfill}{rgb}{0.000000,0.000000,0.000000}%
\pgfsetfillcolor{currentfill}%
\pgfsetlinewidth{0.501875pt}%
\definecolor{currentstroke}{rgb}{0.000000,0.000000,0.000000}%
\pgfsetstrokecolor{currentstroke}%
\pgfsetdash{}{0pt}%
\pgfsys@defobject{currentmarker}{\pgfqpoint{0.000000in}{0.000000in}}{\pgfqpoint{0.000000in}{0.020833in}}{%
\pgfpathmoveto{\pgfqpoint{0.000000in}{0.000000in}}%
\pgfpathlineto{\pgfqpoint{0.000000in}{0.020833in}}%
\pgfusepath{stroke,fill}%
}%
\begin{pgfscope}%
\pgfsys@transformshift{1.714468in}{0.401080in}%
\pgfsys@useobject{currentmarker}{}%
\end{pgfscope}%
\end{pgfscope}%
\begin{pgfscope}%
\pgfsetbuttcap%
\pgfsetroundjoin%
\definecolor{currentfill}{rgb}{0.000000,0.000000,0.000000}%
\pgfsetfillcolor{currentfill}%
\pgfsetlinewidth{0.501875pt}%
\definecolor{currentstroke}{rgb}{0.000000,0.000000,0.000000}%
\pgfsetstrokecolor{currentstroke}%
\pgfsetdash{}{0pt}%
\pgfsys@defobject{currentmarker}{\pgfqpoint{0.000000in}{-0.020833in}}{\pgfqpoint{0.000000in}{0.000000in}}{%
\pgfpathmoveto{\pgfqpoint{0.000000in}{0.000000in}}%
\pgfpathlineto{\pgfqpoint{0.000000in}{-0.020833in}}%
\pgfusepath{stroke,fill}%
}%
\begin{pgfscope}%
\pgfsys@transformshift{1.714468in}{4.405080in}%
\pgfsys@useobject{currentmarker}{}%
\end{pgfscope}%
\end{pgfscope}%
\begin{pgfscope}%
\pgfsetbuttcap%
\pgfsetroundjoin%
\definecolor{currentfill}{rgb}{0.000000,0.000000,0.000000}%
\pgfsetfillcolor{currentfill}%
\pgfsetlinewidth{0.501875pt}%
\definecolor{currentstroke}{rgb}{0.000000,0.000000,0.000000}%
\pgfsetstrokecolor{currentstroke}%
\pgfsetdash{}{0pt}%
\pgfsys@defobject{currentmarker}{\pgfqpoint{0.000000in}{0.000000in}}{\pgfqpoint{0.000000in}{0.020833in}}{%
\pgfpathmoveto{\pgfqpoint{0.000000in}{0.000000in}}%
\pgfpathlineto{\pgfqpoint{0.000000in}{0.020833in}}%
\pgfusepath{stroke,fill}%
}%
\begin{pgfscope}%
\pgfsys@transformshift{1.966343in}{0.401080in}%
\pgfsys@useobject{currentmarker}{}%
\end{pgfscope}%
\end{pgfscope}%
\begin{pgfscope}%
\pgfsetbuttcap%
\pgfsetroundjoin%
\definecolor{currentfill}{rgb}{0.000000,0.000000,0.000000}%
\pgfsetfillcolor{currentfill}%
\pgfsetlinewidth{0.501875pt}%
\definecolor{currentstroke}{rgb}{0.000000,0.000000,0.000000}%
\pgfsetstrokecolor{currentstroke}%
\pgfsetdash{}{0pt}%
\pgfsys@defobject{currentmarker}{\pgfqpoint{0.000000in}{-0.020833in}}{\pgfqpoint{0.000000in}{0.000000in}}{%
\pgfpathmoveto{\pgfqpoint{0.000000in}{0.000000in}}%
\pgfpathlineto{\pgfqpoint{0.000000in}{-0.020833in}}%
\pgfusepath{stroke,fill}%
}%
\begin{pgfscope}%
\pgfsys@transformshift{1.966343in}{4.405080in}%
\pgfsys@useobject{currentmarker}{}%
\end{pgfscope}%
\end{pgfscope}%
\begin{pgfscope}%
\pgfsetbuttcap%
\pgfsetroundjoin%
\definecolor{currentfill}{rgb}{0.000000,0.000000,0.000000}%
\pgfsetfillcolor{currentfill}%
\pgfsetlinewidth{0.501875pt}%
\definecolor{currentstroke}{rgb}{0.000000,0.000000,0.000000}%
\pgfsetstrokecolor{currentstroke}%
\pgfsetdash{}{0pt}%
\pgfsys@defobject{currentmarker}{\pgfqpoint{0.000000in}{0.000000in}}{\pgfqpoint{0.000000in}{0.020833in}}{%
\pgfpathmoveto{\pgfqpoint{0.000000in}{0.000000in}}%
\pgfpathlineto{\pgfqpoint{0.000000in}{0.020833in}}%
\pgfusepath{stroke,fill}%
}%
\begin{pgfscope}%
\pgfsys@transformshift{2.218218in}{0.401080in}%
\pgfsys@useobject{currentmarker}{}%
\end{pgfscope}%
\end{pgfscope}%
\begin{pgfscope}%
\pgfsetbuttcap%
\pgfsetroundjoin%
\definecolor{currentfill}{rgb}{0.000000,0.000000,0.000000}%
\pgfsetfillcolor{currentfill}%
\pgfsetlinewidth{0.501875pt}%
\definecolor{currentstroke}{rgb}{0.000000,0.000000,0.000000}%
\pgfsetstrokecolor{currentstroke}%
\pgfsetdash{}{0pt}%
\pgfsys@defobject{currentmarker}{\pgfqpoint{0.000000in}{-0.020833in}}{\pgfqpoint{0.000000in}{0.000000in}}{%
\pgfpathmoveto{\pgfqpoint{0.000000in}{0.000000in}}%
\pgfpathlineto{\pgfqpoint{0.000000in}{-0.020833in}}%
\pgfusepath{stroke,fill}%
}%
\begin{pgfscope}%
\pgfsys@transformshift{2.218218in}{4.405080in}%
\pgfsys@useobject{currentmarker}{}%
\end{pgfscope}%
\end{pgfscope}%
\begin{pgfscope}%
\pgfsetbuttcap%
\pgfsetroundjoin%
\definecolor{currentfill}{rgb}{0.000000,0.000000,0.000000}%
\pgfsetfillcolor{currentfill}%
\pgfsetlinewidth{0.501875pt}%
\definecolor{currentstroke}{rgb}{0.000000,0.000000,0.000000}%
\pgfsetstrokecolor{currentstroke}%
\pgfsetdash{}{0pt}%
\pgfsys@defobject{currentmarker}{\pgfqpoint{0.000000in}{0.000000in}}{\pgfqpoint{0.000000in}{0.020833in}}{%
\pgfpathmoveto{\pgfqpoint{0.000000in}{0.000000in}}%
\pgfpathlineto{\pgfqpoint{0.000000in}{0.020833in}}%
\pgfusepath{stroke,fill}%
}%
\begin{pgfscope}%
\pgfsys@transformshift{2.721968in}{0.401080in}%
\pgfsys@useobject{currentmarker}{}%
\end{pgfscope}%
\end{pgfscope}%
\begin{pgfscope}%
\pgfsetbuttcap%
\pgfsetroundjoin%
\definecolor{currentfill}{rgb}{0.000000,0.000000,0.000000}%
\pgfsetfillcolor{currentfill}%
\pgfsetlinewidth{0.501875pt}%
\definecolor{currentstroke}{rgb}{0.000000,0.000000,0.000000}%
\pgfsetstrokecolor{currentstroke}%
\pgfsetdash{}{0pt}%
\pgfsys@defobject{currentmarker}{\pgfqpoint{0.000000in}{-0.020833in}}{\pgfqpoint{0.000000in}{0.000000in}}{%
\pgfpathmoveto{\pgfqpoint{0.000000in}{0.000000in}}%
\pgfpathlineto{\pgfqpoint{0.000000in}{-0.020833in}}%
\pgfusepath{stroke,fill}%
}%
\begin{pgfscope}%
\pgfsys@transformshift{2.721968in}{4.405080in}%
\pgfsys@useobject{currentmarker}{}%
\end{pgfscope}%
\end{pgfscope}%
\begin{pgfscope}%
\pgfsetbuttcap%
\pgfsetroundjoin%
\definecolor{currentfill}{rgb}{0.000000,0.000000,0.000000}%
\pgfsetfillcolor{currentfill}%
\pgfsetlinewidth{0.501875pt}%
\definecolor{currentstroke}{rgb}{0.000000,0.000000,0.000000}%
\pgfsetstrokecolor{currentstroke}%
\pgfsetdash{}{0pt}%
\pgfsys@defobject{currentmarker}{\pgfqpoint{0.000000in}{0.000000in}}{\pgfqpoint{0.000000in}{0.020833in}}{%
\pgfpathmoveto{\pgfqpoint{0.000000in}{0.000000in}}%
\pgfpathlineto{\pgfqpoint{0.000000in}{0.020833in}}%
\pgfusepath{stroke,fill}%
}%
\begin{pgfscope}%
\pgfsys@transformshift{2.973843in}{0.401080in}%
\pgfsys@useobject{currentmarker}{}%
\end{pgfscope}%
\end{pgfscope}%
\begin{pgfscope}%
\pgfsetbuttcap%
\pgfsetroundjoin%
\definecolor{currentfill}{rgb}{0.000000,0.000000,0.000000}%
\pgfsetfillcolor{currentfill}%
\pgfsetlinewidth{0.501875pt}%
\definecolor{currentstroke}{rgb}{0.000000,0.000000,0.000000}%
\pgfsetstrokecolor{currentstroke}%
\pgfsetdash{}{0pt}%
\pgfsys@defobject{currentmarker}{\pgfqpoint{0.000000in}{-0.020833in}}{\pgfqpoint{0.000000in}{0.000000in}}{%
\pgfpathmoveto{\pgfqpoint{0.000000in}{0.000000in}}%
\pgfpathlineto{\pgfqpoint{0.000000in}{-0.020833in}}%
\pgfusepath{stroke,fill}%
}%
\begin{pgfscope}%
\pgfsys@transformshift{2.973843in}{4.405080in}%
\pgfsys@useobject{currentmarker}{}%
\end{pgfscope}%
\end{pgfscope}%
\begin{pgfscope}%
\pgfsetbuttcap%
\pgfsetroundjoin%
\definecolor{currentfill}{rgb}{0.000000,0.000000,0.000000}%
\pgfsetfillcolor{currentfill}%
\pgfsetlinewidth{0.501875pt}%
\definecolor{currentstroke}{rgb}{0.000000,0.000000,0.000000}%
\pgfsetstrokecolor{currentstroke}%
\pgfsetdash{}{0pt}%
\pgfsys@defobject{currentmarker}{\pgfqpoint{0.000000in}{0.000000in}}{\pgfqpoint{0.000000in}{0.020833in}}{%
\pgfpathmoveto{\pgfqpoint{0.000000in}{0.000000in}}%
\pgfpathlineto{\pgfqpoint{0.000000in}{0.020833in}}%
\pgfusepath{stroke,fill}%
}%
\begin{pgfscope}%
\pgfsys@transformshift{3.225718in}{0.401080in}%
\pgfsys@useobject{currentmarker}{}%
\end{pgfscope}%
\end{pgfscope}%
\begin{pgfscope}%
\pgfsetbuttcap%
\pgfsetroundjoin%
\definecolor{currentfill}{rgb}{0.000000,0.000000,0.000000}%
\pgfsetfillcolor{currentfill}%
\pgfsetlinewidth{0.501875pt}%
\definecolor{currentstroke}{rgb}{0.000000,0.000000,0.000000}%
\pgfsetstrokecolor{currentstroke}%
\pgfsetdash{}{0pt}%
\pgfsys@defobject{currentmarker}{\pgfqpoint{0.000000in}{-0.020833in}}{\pgfqpoint{0.000000in}{0.000000in}}{%
\pgfpathmoveto{\pgfqpoint{0.000000in}{0.000000in}}%
\pgfpathlineto{\pgfqpoint{0.000000in}{-0.020833in}}%
\pgfusepath{stroke,fill}%
}%
\begin{pgfscope}%
\pgfsys@transformshift{3.225718in}{4.405080in}%
\pgfsys@useobject{currentmarker}{}%
\end{pgfscope}%
\end{pgfscope}%
\begin{pgfscope}%
\pgfsetbuttcap%
\pgfsetroundjoin%
\definecolor{currentfill}{rgb}{0.000000,0.000000,0.000000}%
\pgfsetfillcolor{currentfill}%
\pgfsetlinewidth{0.501875pt}%
\definecolor{currentstroke}{rgb}{0.000000,0.000000,0.000000}%
\pgfsetstrokecolor{currentstroke}%
\pgfsetdash{}{0pt}%
\pgfsys@defobject{currentmarker}{\pgfqpoint{0.000000in}{0.000000in}}{\pgfqpoint{0.000000in}{0.020833in}}{%
\pgfpathmoveto{\pgfqpoint{0.000000in}{0.000000in}}%
\pgfpathlineto{\pgfqpoint{0.000000in}{0.020833in}}%
\pgfusepath{stroke,fill}%
}%
\begin{pgfscope}%
\pgfsys@transformshift{3.729468in}{0.401080in}%
\pgfsys@useobject{currentmarker}{}%
\end{pgfscope}%
\end{pgfscope}%
\begin{pgfscope}%
\pgfsetbuttcap%
\pgfsetroundjoin%
\definecolor{currentfill}{rgb}{0.000000,0.000000,0.000000}%
\pgfsetfillcolor{currentfill}%
\pgfsetlinewidth{0.501875pt}%
\definecolor{currentstroke}{rgb}{0.000000,0.000000,0.000000}%
\pgfsetstrokecolor{currentstroke}%
\pgfsetdash{}{0pt}%
\pgfsys@defobject{currentmarker}{\pgfqpoint{0.000000in}{-0.020833in}}{\pgfqpoint{0.000000in}{0.000000in}}{%
\pgfpathmoveto{\pgfqpoint{0.000000in}{0.000000in}}%
\pgfpathlineto{\pgfqpoint{0.000000in}{-0.020833in}}%
\pgfusepath{stroke,fill}%
}%
\begin{pgfscope}%
\pgfsys@transformshift{3.729468in}{4.405080in}%
\pgfsys@useobject{currentmarker}{}%
\end{pgfscope}%
\end{pgfscope}%
\begin{pgfscope}%
\pgfsetbuttcap%
\pgfsetroundjoin%
\definecolor{currentfill}{rgb}{0.000000,0.000000,0.000000}%
\pgfsetfillcolor{currentfill}%
\pgfsetlinewidth{0.501875pt}%
\definecolor{currentstroke}{rgb}{0.000000,0.000000,0.000000}%
\pgfsetstrokecolor{currentstroke}%
\pgfsetdash{}{0pt}%
\pgfsys@defobject{currentmarker}{\pgfqpoint{0.000000in}{0.000000in}}{\pgfqpoint{0.000000in}{0.020833in}}{%
\pgfpathmoveto{\pgfqpoint{0.000000in}{0.000000in}}%
\pgfpathlineto{\pgfqpoint{0.000000in}{0.020833in}}%
\pgfusepath{stroke,fill}%
}%
\begin{pgfscope}%
\pgfsys@transformshift{3.981343in}{0.401080in}%
\pgfsys@useobject{currentmarker}{}%
\end{pgfscope}%
\end{pgfscope}%
\begin{pgfscope}%
\pgfsetbuttcap%
\pgfsetroundjoin%
\definecolor{currentfill}{rgb}{0.000000,0.000000,0.000000}%
\pgfsetfillcolor{currentfill}%
\pgfsetlinewidth{0.501875pt}%
\definecolor{currentstroke}{rgb}{0.000000,0.000000,0.000000}%
\pgfsetstrokecolor{currentstroke}%
\pgfsetdash{}{0pt}%
\pgfsys@defobject{currentmarker}{\pgfqpoint{0.000000in}{-0.020833in}}{\pgfqpoint{0.000000in}{0.000000in}}{%
\pgfpathmoveto{\pgfqpoint{0.000000in}{0.000000in}}%
\pgfpathlineto{\pgfqpoint{0.000000in}{-0.020833in}}%
\pgfusepath{stroke,fill}%
}%
\begin{pgfscope}%
\pgfsys@transformshift{3.981343in}{4.405080in}%
\pgfsys@useobject{currentmarker}{}%
\end{pgfscope}%
\end{pgfscope}%
\begin{pgfscope}%
\pgfsetbuttcap%
\pgfsetroundjoin%
\definecolor{currentfill}{rgb}{0.000000,0.000000,0.000000}%
\pgfsetfillcolor{currentfill}%
\pgfsetlinewidth{0.501875pt}%
\definecolor{currentstroke}{rgb}{0.000000,0.000000,0.000000}%
\pgfsetstrokecolor{currentstroke}%
\pgfsetdash{}{0pt}%
\pgfsys@defobject{currentmarker}{\pgfqpoint{0.000000in}{0.000000in}}{\pgfqpoint{0.000000in}{0.020833in}}{%
\pgfpathmoveto{\pgfqpoint{0.000000in}{0.000000in}}%
\pgfpathlineto{\pgfqpoint{0.000000in}{0.020833in}}%
\pgfusepath{stroke,fill}%
}%
\begin{pgfscope}%
\pgfsys@transformshift{4.233218in}{0.401080in}%
\pgfsys@useobject{currentmarker}{}%
\end{pgfscope}%
\end{pgfscope}%
\begin{pgfscope}%
\pgfsetbuttcap%
\pgfsetroundjoin%
\definecolor{currentfill}{rgb}{0.000000,0.000000,0.000000}%
\pgfsetfillcolor{currentfill}%
\pgfsetlinewidth{0.501875pt}%
\definecolor{currentstroke}{rgb}{0.000000,0.000000,0.000000}%
\pgfsetstrokecolor{currentstroke}%
\pgfsetdash{}{0pt}%
\pgfsys@defobject{currentmarker}{\pgfqpoint{0.000000in}{-0.020833in}}{\pgfqpoint{0.000000in}{0.000000in}}{%
\pgfpathmoveto{\pgfqpoint{0.000000in}{0.000000in}}%
\pgfpathlineto{\pgfqpoint{0.000000in}{-0.020833in}}%
\pgfusepath{stroke,fill}%
}%
\begin{pgfscope}%
\pgfsys@transformshift{4.233218in}{4.405080in}%
\pgfsys@useobject{currentmarker}{}%
\end{pgfscope}%
\end{pgfscope}%
\begin{pgfscope}%
\pgfsetbuttcap%
\pgfsetroundjoin%
\definecolor{currentfill}{rgb}{0.000000,0.000000,0.000000}%
\pgfsetfillcolor{currentfill}%
\pgfsetlinewidth{0.501875pt}%
\definecolor{currentstroke}{rgb}{0.000000,0.000000,0.000000}%
\pgfsetstrokecolor{currentstroke}%
\pgfsetdash{}{0pt}%
\pgfsys@defobject{currentmarker}{\pgfqpoint{0.000000in}{0.000000in}}{\pgfqpoint{0.000000in}{0.020833in}}{%
\pgfpathmoveto{\pgfqpoint{0.000000in}{0.000000in}}%
\pgfpathlineto{\pgfqpoint{0.000000in}{0.020833in}}%
\pgfusepath{stroke,fill}%
}%
\begin{pgfscope}%
\pgfsys@transformshift{4.736968in}{0.401080in}%
\pgfsys@useobject{currentmarker}{}%
\end{pgfscope}%
\end{pgfscope}%
\begin{pgfscope}%
\pgfsetbuttcap%
\pgfsetroundjoin%
\definecolor{currentfill}{rgb}{0.000000,0.000000,0.000000}%
\pgfsetfillcolor{currentfill}%
\pgfsetlinewidth{0.501875pt}%
\definecolor{currentstroke}{rgb}{0.000000,0.000000,0.000000}%
\pgfsetstrokecolor{currentstroke}%
\pgfsetdash{}{0pt}%
\pgfsys@defobject{currentmarker}{\pgfqpoint{0.000000in}{-0.020833in}}{\pgfqpoint{0.000000in}{0.000000in}}{%
\pgfpathmoveto{\pgfqpoint{0.000000in}{0.000000in}}%
\pgfpathlineto{\pgfqpoint{0.000000in}{-0.020833in}}%
\pgfusepath{stroke,fill}%
}%
\begin{pgfscope}%
\pgfsys@transformshift{4.736968in}{4.405080in}%
\pgfsys@useobject{currentmarker}{}%
\end{pgfscope}%
\end{pgfscope}%
\begin{pgfscope}%
\pgfsetbuttcap%
\pgfsetroundjoin%
\definecolor{currentfill}{rgb}{0.000000,0.000000,0.000000}%
\pgfsetfillcolor{currentfill}%
\pgfsetlinewidth{0.501875pt}%
\definecolor{currentstroke}{rgb}{0.000000,0.000000,0.000000}%
\pgfsetstrokecolor{currentstroke}%
\pgfsetdash{}{0pt}%
\pgfsys@defobject{currentmarker}{\pgfqpoint{0.000000in}{0.000000in}}{\pgfqpoint{0.000000in}{0.020833in}}{%
\pgfpathmoveto{\pgfqpoint{0.000000in}{0.000000in}}%
\pgfpathlineto{\pgfqpoint{0.000000in}{0.020833in}}%
\pgfusepath{stroke,fill}%
}%
\begin{pgfscope}%
\pgfsys@transformshift{4.988843in}{0.401080in}%
\pgfsys@useobject{currentmarker}{}%
\end{pgfscope}%
\end{pgfscope}%
\begin{pgfscope}%
\pgfsetbuttcap%
\pgfsetroundjoin%
\definecolor{currentfill}{rgb}{0.000000,0.000000,0.000000}%
\pgfsetfillcolor{currentfill}%
\pgfsetlinewidth{0.501875pt}%
\definecolor{currentstroke}{rgb}{0.000000,0.000000,0.000000}%
\pgfsetstrokecolor{currentstroke}%
\pgfsetdash{}{0pt}%
\pgfsys@defobject{currentmarker}{\pgfqpoint{0.000000in}{-0.020833in}}{\pgfqpoint{0.000000in}{0.000000in}}{%
\pgfpathmoveto{\pgfqpoint{0.000000in}{0.000000in}}%
\pgfpathlineto{\pgfqpoint{0.000000in}{-0.020833in}}%
\pgfusepath{stroke,fill}%
}%
\begin{pgfscope}%
\pgfsys@transformshift{4.988843in}{4.405080in}%
\pgfsys@useobject{currentmarker}{}%
\end{pgfscope}%
\end{pgfscope}%
\begin{pgfscope}%
\pgfsetbuttcap%
\pgfsetroundjoin%
\definecolor{currentfill}{rgb}{0.000000,0.000000,0.000000}%
\pgfsetfillcolor{currentfill}%
\pgfsetlinewidth{0.501875pt}%
\definecolor{currentstroke}{rgb}{0.000000,0.000000,0.000000}%
\pgfsetstrokecolor{currentstroke}%
\pgfsetdash{}{0pt}%
\pgfsys@defobject{currentmarker}{\pgfqpoint{0.000000in}{0.000000in}}{\pgfqpoint{0.000000in}{0.020833in}}{%
\pgfpathmoveto{\pgfqpoint{0.000000in}{0.000000in}}%
\pgfpathlineto{\pgfqpoint{0.000000in}{0.020833in}}%
\pgfusepath{stroke,fill}%
}%
\begin{pgfscope}%
\pgfsys@transformshift{5.240718in}{0.401080in}%
\pgfsys@useobject{currentmarker}{}%
\end{pgfscope}%
\end{pgfscope}%
\begin{pgfscope}%
\pgfsetbuttcap%
\pgfsetroundjoin%
\definecolor{currentfill}{rgb}{0.000000,0.000000,0.000000}%
\pgfsetfillcolor{currentfill}%
\pgfsetlinewidth{0.501875pt}%
\definecolor{currentstroke}{rgb}{0.000000,0.000000,0.000000}%
\pgfsetstrokecolor{currentstroke}%
\pgfsetdash{}{0pt}%
\pgfsys@defobject{currentmarker}{\pgfqpoint{0.000000in}{-0.020833in}}{\pgfqpoint{0.000000in}{0.000000in}}{%
\pgfpathmoveto{\pgfqpoint{0.000000in}{0.000000in}}%
\pgfpathlineto{\pgfqpoint{0.000000in}{-0.020833in}}%
\pgfusepath{stroke,fill}%
}%
\begin{pgfscope}%
\pgfsys@transformshift{5.240718in}{4.405080in}%
\pgfsys@useobject{currentmarker}{}%
\end{pgfscope}%
\end{pgfscope}%
\begin{pgfscope}%
\definecolor{textcolor}{rgb}{0.000000,0.000000,0.000000}%
\pgfsetstrokecolor{textcolor}%
\pgfsetfillcolor{textcolor}%
\pgftext[x=2.973843in,y=0.173457in,,top]{\color{textcolor}{\rmfamily\fontsize{10.000000}{12.000000}\selectfont\catcode`\^=\active\def^{\ifmmode\sp\else\^{}\fi}\catcode`\%=\active\def%{\%}Time}}%
\end{pgfscope}%
\begin{pgfscope}%
\pgfsetbuttcap%
\pgfsetroundjoin%
\definecolor{currentfill}{rgb}{0.000000,0.000000,0.000000}%
\pgfsetfillcolor{currentfill}%
\pgfsetlinewidth{0.501875pt}%
\definecolor{currentstroke}{rgb}{0.000000,0.000000,0.000000}%
\pgfsetstrokecolor{currentstroke}%
\pgfsetdash{}{0pt}%
\pgfsys@defobject{currentmarker}{\pgfqpoint{0.000000in}{0.000000in}}{\pgfqpoint{0.041667in}{0.000000in}}{%
\pgfpathmoveto{\pgfqpoint{0.000000in}{0.000000in}}%
\pgfpathlineto{\pgfqpoint{0.041667in}{0.000000in}}%
\pgfusepath{stroke,fill}%
}%
\begin{pgfscope}%
\pgfsys@transformshift{0.455093in}{0.426983in}%
\pgfsys@useobject{currentmarker}{}%
\end{pgfscope}%
\end{pgfscope}%
\begin{pgfscope}%
\pgfsetbuttcap%
\pgfsetroundjoin%
\definecolor{currentfill}{rgb}{0.000000,0.000000,0.000000}%
\pgfsetfillcolor{currentfill}%
\pgfsetlinewidth{0.501875pt}%
\definecolor{currentstroke}{rgb}{0.000000,0.000000,0.000000}%
\pgfsetstrokecolor{currentstroke}%
\pgfsetdash{}{0pt}%
\pgfsys@defobject{currentmarker}{\pgfqpoint{-0.041667in}{0.000000in}}{\pgfqpoint{-0.000000in}{0.000000in}}{%
\pgfpathmoveto{\pgfqpoint{-0.000000in}{0.000000in}}%
\pgfpathlineto{\pgfqpoint{-0.041667in}{0.000000in}}%
\pgfusepath{stroke,fill}%
}%
\begin{pgfscope}%
\pgfsys@transformshift{5.492593in}{0.426983in}%
\pgfsys@useobject{currentmarker}{}%
\end{pgfscope}%
\end{pgfscope}%
\begin{pgfscope}%
\definecolor{textcolor}{rgb}{0.000000,0.000000,0.000000}%
\pgfsetstrokecolor{textcolor}%
\pgfsetfillcolor{textcolor}%
\pgftext[x=0.229012in, y=0.378757in, left, base]{\color{textcolor}{\rmfamily\fontsize{10.000000}{12.000000}\selectfont\catcode`\^=\active\def^{\ifmmode\sp\else\^{}\fi}\catcode`\%=\active\def%{\%}$\mathdefault{0.8}$}}%
\end{pgfscope}%
\begin{pgfscope}%
\pgfsetbuttcap%
\pgfsetroundjoin%
\definecolor{currentfill}{rgb}{0.000000,0.000000,0.000000}%
\pgfsetfillcolor{currentfill}%
\pgfsetlinewidth{0.501875pt}%
\definecolor{currentstroke}{rgb}{0.000000,0.000000,0.000000}%
\pgfsetstrokecolor{currentstroke}%
\pgfsetdash{}{0pt}%
\pgfsys@defobject{currentmarker}{\pgfqpoint{0.000000in}{0.000000in}}{\pgfqpoint{0.041667in}{0.000000in}}{%
\pgfpathmoveto{\pgfqpoint{0.000000in}{0.000000in}}%
\pgfpathlineto{\pgfqpoint{0.041667in}{0.000000in}}%
\pgfusepath{stroke,fill}%
}%
\begin{pgfscope}%
\pgfsys@transformshift{0.455093in}{1.021044in}%
\pgfsys@useobject{currentmarker}{}%
\end{pgfscope}%
\end{pgfscope}%
\begin{pgfscope}%
\pgfsetbuttcap%
\pgfsetroundjoin%
\definecolor{currentfill}{rgb}{0.000000,0.000000,0.000000}%
\pgfsetfillcolor{currentfill}%
\pgfsetlinewidth{0.501875pt}%
\definecolor{currentstroke}{rgb}{0.000000,0.000000,0.000000}%
\pgfsetstrokecolor{currentstroke}%
\pgfsetdash{}{0pt}%
\pgfsys@defobject{currentmarker}{\pgfqpoint{-0.041667in}{0.000000in}}{\pgfqpoint{-0.000000in}{0.000000in}}{%
\pgfpathmoveto{\pgfqpoint{-0.000000in}{0.000000in}}%
\pgfpathlineto{\pgfqpoint{-0.041667in}{0.000000in}}%
\pgfusepath{stroke,fill}%
}%
\begin{pgfscope}%
\pgfsys@transformshift{5.492593in}{1.021044in}%
\pgfsys@useobject{currentmarker}{}%
\end{pgfscope}%
\end{pgfscope}%
\begin{pgfscope}%
\definecolor{textcolor}{rgb}{0.000000,0.000000,0.000000}%
\pgfsetstrokecolor{textcolor}%
\pgfsetfillcolor{textcolor}%
\pgftext[x=0.229012in, y=0.972819in, left, base]{\color{textcolor}{\rmfamily\fontsize{10.000000}{12.000000}\selectfont\catcode`\^=\active\def^{\ifmmode\sp\else\^{}\fi}\catcode`\%=\active\def%{\%}$\mathdefault{1.0}$}}%
\end{pgfscope}%
\begin{pgfscope}%
\pgfsetbuttcap%
\pgfsetroundjoin%
\definecolor{currentfill}{rgb}{0.000000,0.000000,0.000000}%
\pgfsetfillcolor{currentfill}%
\pgfsetlinewidth{0.501875pt}%
\definecolor{currentstroke}{rgb}{0.000000,0.000000,0.000000}%
\pgfsetstrokecolor{currentstroke}%
\pgfsetdash{}{0pt}%
\pgfsys@defobject{currentmarker}{\pgfqpoint{0.000000in}{0.000000in}}{\pgfqpoint{0.041667in}{0.000000in}}{%
\pgfpathmoveto{\pgfqpoint{0.000000in}{0.000000in}}%
\pgfpathlineto{\pgfqpoint{0.041667in}{0.000000in}}%
\pgfusepath{stroke,fill}%
}%
\begin{pgfscope}%
\pgfsys@transformshift{0.455093in}{1.615105in}%
\pgfsys@useobject{currentmarker}{}%
\end{pgfscope}%
\end{pgfscope}%
\begin{pgfscope}%
\pgfsetbuttcap%
\pgfsetroundjoin%
\definecolor{currentfill}{rgb}{0.000000,0.000000,0.000000}%
\pgfsetfillcolor{currentfill}%
\pgfsetlinewidth{0.501875pt}%
\definecolor{currentstroke}{rgb}{0.000000,0.000000,0.000000}%
\pgfsetstrokecolor{currentstroke}%
\pgfsetdash{}{0pt}%
\pgfsys@defobject{currentmarker}{\pgfqpoint{-0.041667in}{0.000000in}}{\pgfqpoint{-0.000000in}{0.000000in}}{%
\pgfpathmoveto{\pgfqpoint{-0.000000in}{0.000000in}}%
\pgfpathlineto{\pgfqpoint{-0.041667in}{0.000000in}}%
\pgfusepath{stroke,fill}%
}%
\begin{pgfscope}%
\pgfsys@transformshift{5.492593in}{1.615105in}%
\pgfsys@useobject{currentmarker}{}%
\end{pgfscope}%
\end{pgfscope}%
\begin{pgfscope}%
\definecolor{textcolor}{rgb}{0.000000,0.000000,0.000000}%
\pgfsetstrokecolor{textcolor}%
\pgfsetfillcolor{textcolor}%
\pgftext[x=0.229012in, y=1.566880in, left, base]{\color{textcolor}{\rmfamily\fontsize{10.000000}{12.000000}\selectfont\catcode`\^=\active\def^{\ifmmode\sp\else\^{}\fi}\catcode`\%=\active\def%{\%}$\mathdefault{1.2}$}}%
\end{pgfscope}%
\begin{pgfscope}%
\pgfsetbuttcap%
\pgfsetroundjoin%
\definecolor{currentfill}{rgb}{0.000000,0.000000,0.000000}%
\pgfsetfillcolor{currentfill}%
\pgfsetlinewidth{0.501875pt}%
\definecolor{currentstroke}{rgb}{0.000000,0.000000,0.000000}%
\pgfsetstrokecolor{currentstroke}%
\pgfsetdash{}{0pt}%
\pgfsys@defobject{currentmarker}{\pgfqpoint{0.000000in}{0.000000in}}{\pgfqpoint{0.041667in}{0.000000in}}{%
\pgfpathmoveto{\pgfqpoint{0.000000in}{0.000000in}}%
\pgfpathlineto{\pgfqpoint{0.041667in}{0.000000in}}%
\pgfusepath{stroke,fill}%
}%
\begin{pgfscope}%
\pgfsys@transformshift{0.455093in}{2.209166in}%
\pgfsys@useobject{currentmarker}{}%
\end{pgfscope}%
\end{pgfscope}%
\begin{pgfscope}%
\pgfsetbuttcap%
\pgfsetroundjoin%
\definecolor{currentfill}{rgb}{0.000000,0.000000,0.000000}%
\pgfsetfillcolor{currentfill}%
\pgfsetlinewidth{0.501875pt}%
\definecolor{currentstroke}{rgb}{0.000000,0.000000,0.000000}%
\pgfsetstrokecolor{currentstroke}%
\pgfsetdash{}{0pt}%
\pgfsys@defobject{currentmarker}{\pgfqpoint{-0.041667in}{0.000000in}}{\pgfqpoint{-0.000000in}{0.000000in}}{%
\pgfpathmoveto{\pgfqpoint{-0.000000in}{0.000000in}}%
\pgfpathlineto{\pgfqpoint{-0.041667in}{0.000000in}}%
\pgfusepath{stroke,fill}%
}%
\begin{pgfscope}%
\pgfsys@transformshift{5.492593in}{2.209166in}%
\pgfsys@useobject{currentmarker}{}%
\end{pgfscope}%
\end{pgfscope}%
\begin{pgfscope}%
\definecolor{textcolor}{rgb}{0.000000,0.000000,0.000000}%
\pgfsetstrokecolor{textcolor}%
\pgfsetfillcolor{textcolor}%
\pgftext[x=0.229012in, y=2.160941in, left, base]{\color{textcolor}{\rmfamily\fontsize{10.000000}{12.000000}\selectfont\catcode`\^=\active\def^{\ifmmode\sp\else\^{}\fi}\catcode`\%=\active\def%{\%}$\mathdefault{1.4}$}}%
\end{pgfscope}%
\begin{pgfscope}%
\pgfsetbuttcap%
\pgfsetroundjoin%
\definecolor{currentfill}{rgb}{0.000000,0.000000,0.000000}%
\pgfsetfillcolor{currentfill}%
\pgfsetlinewidth{0.501875pt}%
\definecolor{currentstroke}{rgb}{0.000000,0.000000,0.000000}%
\pgfsetstrokecolor{currentstroke}%
\pgfsetdash{}{0pt}%
\pgfsys@defobject{currentmarker}{\pgfqpoint{0.000000in}{0.000000in}}{\pgfqpoint{0.041667in}{0.000000in}}{%
\pgfpathmoveto{\pgfqpoint{0.000000in}{0.000000in}}%
\pgfpathlineto{\pgfqpoint{0.041667in}{0.000000in}}%
\pgfusepath{stroke,fill}%
}%
\begin{pgfscope}%
\pgfsys@transformshift{0.455093in}{2.803228in}%
\pgfsys@useobject{currentmarker}{}%
\end{pgfscope}%
\end{pgfscope}%
\begin{pgfscope}%
\pgfsetbuttcap%
\pgfsetroundjoin%
\definecolor{currentfill}{rgb}{0.000000,0.000000,0.000000}%
\pgfsetfillcolor{currentfill}%
\pgfsetlinewidth{0.501875pt}%
\definecolor{currentstroke}{rgb}{0.000000,0.000000,0.000000}%
\pgfsetstrokecolor{currentstroke}%
\pgfsetdash{}{0pt}%
\pgfsys@defobject{currentmarker}{\pgfqpoint{-0.041667in}{0.000000in}}{\pgfqpoint{-0.000000in}{0.000000in}}{%
\pgfpathmoveto{\pgfqpoint{-0.000000in}{0.000000in}}%
\pgfpathlineto{\pgfqpoint{-0.041667in}{0.000000in}}%
\pgfusepath{stroke,fill}%
}%
\begin{pgfscope}%
\pgfsys@transformshift{5.492593in}{2.803228in}%
\pgfsys@useobject{currentmarker}{}%
\end{pgfscope}%
\end{pgfscope}%
\begin{pgfscope}%
\definecolor{textcolor}{rgb}{0.000000,0.000000,0.000000}%
\pgfsetstrokecolor{textcolor}%
\pgfsetfillcolor{textcolor}%
\pgftext[x=0.229012in, y=2.755002in, left, base]{\color{textcolor}{\rmfamily\fontsize{10.000000}{12.000000}\selectfont\catcode`\^=\active\def^{\ifmmode\sp\else\^{}\fi}\catcode`\%=\active\def%{\%}$\mathdefault{1.6}$}}%
\end{pgfscope}%
\begin{pgfscope}%
\pgfsetbuttcap%
\pgfsetroundjoin%
\definecolor{currentfill}{rgb}{0.000000,0.000000,0.000000}%
\pgfsetfillcolor{currentfill}%
\pgfsetlinewidth{0.501875pt}%
\definecolor{currentstroke}{rgb}{0.000000,0.000000,0.000000}%
\pgfsetstrokecolor{currentstroke}%
\pgfsetdash{}{0pt}%
\pgfsys@defobject{currentmarker}{\pgfqpoint{0.000000in}{0.000000in}}{\pgfqpoint{0.041667in}{0.000000in}}{%
\pgfpathmoveto{\pgfqpoint{0.000000in}{0.000000in}}%
\pgfpathlineto{\pgfqpoint{0.041667in}{0.000000in}}%
\pgfusepath{stroke,fill}%
}%
\begin{pgfscope}%
\pgfsys@transformshift{0.455093in}{3.397289in}%
\pgfsys@useobject{currentmarker}{}%
\end{pgfscope}%
\end{pgfscope}%
\begin{pgfscope}%
\pgfsetbuttcap%
\pgfsetroundjoin%
\definecolor{currentfill}{rgb}{0.000000,0.000000,0.000000}%
\pgfsetfillcolor{currentfill}%
\pgfsetlinewidth{0.501875pt}%
\definecolor{currentstroke}{rgb}{0.000000,0.000000,0.000000}%
\pgfsetstrokecolor{currentstroke}%
\pgfsetdash{}{0pt}%
\pgfsys@defobject{currentmarker}{\pgfqpoint{-0.041667in}{0.000000in}}{\pgfqpoint{-0.000000in}{0.000000in}}{%
\pgfpathmoveto{\pgfqpoint{-0.000000in}{0.000000in}}%
\pgfpathlineto{\pgfqpoint{-0.041667in}{0.000000in}}%
\pgfusepath{stroke,fill}%
}%
\begin{pgfscope}%
\pgfsys@transformshift{5.492593in}{3.397289in}%
\pgfsys@useobject{currentmarker}{}%
\end{pgfscope}%
\end{pgfscope}%
\begin{pgfscope}%
\definecolor{textcolor}{rgb}{0.000000,0.000000,0.000000}%
\pgfsetstrokecolor{textcolor}%
\pgfsetfillcolor{textcolor}%
\pgftext[x=0.229012in, y=3.349064in, left, base]{\color{textcolor}{\rmfamily\fontsize{10.000000}{12.000000}\selectfont\catcode`\^=\active\def^{\ifmmode\sp\else\^{}\fi}\catcode`\%=\active\def%{\%}$\mathdefault{1.8}$}}%
\end{pgfscope}%
\begin{pgfscope}%
\pgfsetbuttcap%
\pgfsetroundjoin%
\definecolor{currentfill}{rgb}{0.000000,0.000000,0.000000}%
\pgfsetfillcolor{currentfill}%
\pgfsetlinewidth{0.501875pt}%
\definecolor{currentstroke}{rgb}{0.000000,0.000000,0.000000}%
\pgfsetstrokecolor{currentstroke}%
\pgfsetdash{}{0pt}%
\pgfsys@defobject{currentmarker}{\pgfqpoint{0.000000in}{0.000000in}}{\pgfqpoint{0.041667in}{0.000000in}}{%
\pgfpathmoveto{\pgfqpoint{0.000000in}{0.000000in}}%
\pgfpathlineto{\pgfqpoint{0.041667in}{0.000000in}}%
\pgfusepath{stroke,fill}%
}%
\begin{pgfscope}%
\pgfsys@transformshift{0.455093in}{3.991350in}%
\pgfsys@useobject{currentmarker}{}%
\end{pgfscope}%
\end{pgfscope}%
\begin{pgfscope}%
\pgfsetbuttcap%
\pgfsetroundjoin%
\definecolor{currentfill}{rgb}{0.000000,0.000000,0.000000}%
\pgfsetfillcolor{currentfill}%
\pgfsetlinewidth{0.501875pt}%
\definecolor{currentstroke}{rgb}{0.000000,0.000000,0.000000}%
\pgfsetstrokecolor{currentstroke}%
\pgfsetdash{}{0pt}%
\pgfsys@defobject{currentmarker}{\pgfqpoint{-0.041667in}{0.000000in}}{\pgfqpoint{-0.000000in}{0.000000in}}{%
\pgfpathmoveto{\pgfqpoint{-0.000000in}{0.000000in}}%
\pgfpathlineto{\pgfqpoint{-0.041667in}{0.000000in}}%
\pgfusepath{stroke,fill}%
}%
\begin{pgfscope}%
\pgfsys@transformshift{5.492593in}{3.991350in}%
\pgfsys@useobject{currentmarker}{}%
\end{pgfscope}%
\end{pgfscope}%
\begin{pgfscope}%
\definecolor{textcolor}{rgb}{0.000000,0.000000,0.000000}%
\pgfsetstrokecolor{textcolor}%
\pgfsetfillcolor{textcolor}%
\pgftext[x=0.229012in, y=3.943125in, left, base]{\color{textcolor}{\rmfamily\fontsize{10.000000}{12.000000}\selectfont\catcode`\^=\active\def^{\ifmmode\sp\else\^{}\fi}\catcode`\%=\active\def%{\%}$\mathdefault{2.0}$}}%
\end{pgfscope}%
\begin{pgfscope}%
\pgfsetbuttcap%
\pgfsetroundjoin%
\definecolor{currentfill}{rgb}{0.000000,0.000000,0.000000}%
\pgfsetfillcolor{currentfill}%
\pgfsetlinewidth{0.501875pt}%
\definecolor{currentstroke}{rgb}{0.000000,0.000000,0.000000}%
\pgfsetstrokecolor{currentstroke}%
\pgfsetdash{}{0pt}%
\pgfsys@defobject{currentmarker}{\pgfqpoint{0.000000in}{0.000000in}}{\pgfqpoint{0.020833in}{0.000000in}}{%
\pgfpathmoveto{\pgfqpoint{0.000000in}{0.000000in}}%
\pgfpathlineto{\pgfqpoint{0.020833in}{0.000000in}}%
\pgfusepath{stroke,fill}%
}%
\begin{pgfscope}%
\pgfsys@transformshift{0.455093in}{0.575498in}%
\pgfsys@useobject{currentmarker}{}%
\end{pgfscope}%
\end{pgfscope}%
\begin{pgfscope}%
\pgfsetbuttcap%
\pgfsetroundjoin%
\definecolor{currentfill}{rgb}{0.000000,0.000000,0.000000}%
\pgfsetfillcolor{currentfill}%
\pgfsetlinewidth{0.501875pt}%
\definecolor{currentstroke}{rgb}{0.000000,0.000000,0.000000}%
\pgfsetstrokecolor{currentstroke}%
\pgfsetdash{}{0pt}%
\pgfsys@defobject{currentmarker}{\pgfqpoint{-0.020833in}{0.000000in}}{\pgfqpoint{-0.000000in}{0.000000in}}{%
\pgfpathmoveto{\pgfqpoint{-0.000000in}{0.000000in}}%
\pgfpathlineto{\pgfqpoint{-0.020833in}{0.000000in}}%
\pgfusepath{stroke,fill}%
}%
\begin{pgfscope}%
\pgfsys@transformshift{5.492593in}{0.575498in}%
\pgfsys@useobject{currentmarker}{}%
\end{pgfscope}%
\end{pgfscope}%
\begin{pgfscope}%
\pgfsetbuttcap%
\pgfsetroundjoin%
\definecolor{currentfill}{rgb}{0.000000,0.000000,0.000000}%
\pgfsetfillcolor{currentfill}%
\pgfsetlinewidth{0.501875pt}%
\definecolor{currentstroke}{rgb}{0.000000,0.000000,0.000000}%
\pgfsetstrokecolor{currentstroke}%
\pgfsetdash{}{0pt}%
\pgfsys@defobject{currentmarker}{\pgfqpoint{0.000000in}{0.000000in}}{\pgfqpoint{0.020833in}{0.000000in}}{%
\pgfpathmoveto{\pgfqpoint{0.000000in}{0.000000in}}%
\pgfpathlineto{\pgfqpoint{0.020833in}{0.000000in}}%
\pgfusepath{stroke,fill}%
}%
\begin{pgfscope}%
\pgfsys@transformshift{0.455093in}{0.724013in}%
\pgfsys@useobject{currentmarker}{}%
\end{pgfscope}%
\end{pgfscope}%
\begin{pgfscope}%
\pgfsetbuttcap%
\pgfsetroundjoin%
\definecolor{currentfill}{rgb}{0.000000,0.000000,0.000000}%
\pgfsetfillcolor{currentfill}%
\pgfsetlinewidth{0.501875pt}%
\definecolor{currentstroke}{rgb}{0.000000,0.000000,0.000000}%
\pgfsetstrokecolor{currentstroke}%
\pgfsetdash{}{0pt}%
\pgfsys@defobject{currentmarker}{\pgfqpoint{-0.020833in}{0.000000in}}{\pgfqpoint{-0.000000in}{0.000000in}}{%
\pgfpathmoveto{\pgfqpoint{-0.000000in}{0.000000in}}%
\pgfpathlineto{\pgfqpoint{-0.020833in}{0.000000in}}%
\pgfusepath{stroke,fill}%
}%
\begin{pgfscope}%
\pgfsys@transformshift{5.492593in}{0.724013in}%
\pgfsys@useobject{currentmarker}{}%
\end{pgfscope}%
\end{pgfscope}%
\begin{pgfscope}%
\pgfsetbuttcap%
\pgfsetroundjoin%
\definecolor{currentfill}{rgb}{0.000000,0.000000,0.000000}%
\pgfsetfillcolor{currentfill}%
\pgfsetlinewidth{0.501875pt}%
\definecolor{currentstroke}{rgb}{0.000000,0.000000,0.000000}%
\pgfsetstrokecolor{currentstroke}%
\pgfsetdash{}{0pt}%
\pgfsys@defobject{currentmarker}{\pgfqpoint{0.000000in}{0.000000in}}{\pgfqpoint{0.020833in}{0.000000in}}{%
\pgfpathmoveto{\pgfqpoint{0.000000in}{0.000000in}}%
\pgfpathlineto{\pgfqpoint{0.020833in}{0.000000in}}%
\pgfusepath{stroke,fill}%
}%
\begin{pgfscope}%
\pgfsys@transformshift{0.455093in}{0.872529in}%
\pgfsys@useobject{currentmarker}{}%
\end{pgfscope}%
\end{pgfscope}%
\begin{pgfscope}%
\pgfsetbuttcap%
\pgfsetroundjoin%
\definecolor{currentfill}{rgb}{0.000000,0.000000,0.000000}%
\pgfsetfillcolor{currentfill}%
\pgfsetlinewidth{0.501875pt}%
\definecolor{currentstroke}{rgb}{0.000000,0.000000,0.000000}%
\pgfsetstrokecolor{currentstroke}%
\pgfsetdash{}{0pt}%
\pgfsys@defobject{currentmarker}{\pgfqpoint{-0.020833in}{0.000000in}}{\pgfqpoint{-0.000000in}{0.000000in}}{%
\pgfpathmoveto{\pgfqpoint{-0.000000in}{0.000000in}}%
\pgfpathlineto{\pgfqpoint{-0.020833in}{0.000000in}}%
\pgfusepath{stroke,fill}%
}%
\begin{pgfscope}%
\pgfsys@transformshift{5.492593in}{0.872529in}%
\pgfsys@useobject{currentmarker}{}%
\end{pgfscope}%
\end{pgfscope}%
\begin{pgfscope}%
\pgfsetbuttcap%
\pgfsetroundjoin%
\definecolor{currentfill}{rgb}{0.000000,0.000000,0.000000}%
\pgfsetfillcolor{currentfill}%
\pgfsetlinewidth{0.501875pt}%
\definecolor{currentstroke}{rgb}{0.000000,0.000000,0.000000}%
\pgfsetstrokecolor{currentstroke}%
\pgfsetdash{}{0pt}%
\pgfsys@defobject{currentmarker}{\pgfqpoint{0.000000in}{0.000000in}}{\pgfqpoint{0.020833in}{0.000000in}}{%
\pgfpathmoveto{\pgfqpoint{0.000000in}{0.000000in}}%
\pgfpathlineto{\pgfqpoint{0.020833in}{0.000000in}}%
\pgfusepath{stroke,fill}%
}%
\begin{pgfscope}%
\pgfsys@transformshift{0.455093in}{1.169559in}%
\pgfsys@useobject{currentmarker}{}%
\end{pgfscope}%
\end{pgfscope}%
\begin{pgfscope}%
\pgfsetbuttcap%
\pgfsetroundjoin%
\definecolor{currentfill}{rgb}{0.000000,0.000000,0.000000}%
\pgfsetfillcolor{currentfill}%
\pgfsetlinewidth{0.501875pt}%
\definecolor{currentstroke}{rgb}{0.000000,0.000000,0.000000}%
\pgfsetstrokecolor{currentstroke}%
\pgfsetdash{}{0pt}%
\pgfsys@defobject{currentmarker}{\pgfqpoint{-0.020833in}{0.000000in}}{\pgfqpoint{-0.000000in}{0.000000in}}{%
\pgfpathmoveto{\pgfqpoint{-0.000000in}{0.000000in}}%
\pgfpathlineto{\pgfqpoint{-0.020833in}{0.000000in}}%
\pgfusepath{stroke,fill}%
}%
\begin{pgfscope}%
\pgfsys@transformshift{5.492593in}{1.169559in}%
\pgfsys@useobject{currentmarker}{}%
\end{pgfscope}%
\end{pgfscope}%
\begin{pgfscope}%
\pgfsetbuttcap%
\pgfsetroundjoin%
\definecolor{currentfill}{rgb}{0.000000,0.000000,0.000000}%
\pgfsetfillcolor{currentfill}%
\pgfsetlinewidth{0.501875pt}%
\definecolor{currentstroke}{rgb}{0.000000,0.000000,0.000000}%
\pgfsetstrokecolor{currentstroke}%
\pgfsetdash{}{0pt}%
\pgfsys@defobject{currentmarker}{\pgfqpoint{0.000000in}{0.000000in}}{\pgfqpoint{0.020833in}{0.000000in}}{%
\pgfpathmoveto{\pgfqpoint{0.000000in}{0.000000in}}%
\pgfpathlineto{\pgfqpoint{0.020833in}{0.000000in}}%
\pgfusepath{stroke,fill}%
}%
\begin{pgfscope}%
\pgfsys@transformshift{0.455093in}{1.318074in}%
\pgfsys@useobject{currentmarker}{}%
\end{pgfscope}%
\end{pgfscope}%
\begin{pgfscope}%
\pgfsetbuttcap%
\pgfsetroundjoin%
\definecolor{currentfill}{rgb}{0.000000,0.000000,0.000000}%
\pgfsetfillcolor{currentfill}%
\pgfsetlinewidth{0.501875pt}%
\definecolor{currentstroke}{rgb}{0.000000,0.000000,0.000000}%
\pgfsetstrokecolor{currentstroke}%
\pgfsetdash{}{0pt}%
\pgfsys@defobject{currentmarker}{\pgfqpoint{-0.020833in}{0.000000in}}{\pgfqpoint{-0.000000in}{0.000000in}}{%
\pgfpathmoveto{\pgfqpoint{-0.000000in}{0.000000in}}%
\pgfpathlineto{\pgfqpoint{-0.020833in}{0.000000in}}%
\pgfusepath{stroke,fill}%
}%
\begin{pgfscope}%
\pgfsys@transformshift{5.492593in}{1.318074in}%
\pgfsys@useobject{currentmarker}{}%
\end{pgfscope}%
\end{pgfscope}%
\begin{pgfscope}%
\pgfsetbuttcap%
\pgfsetroundjoin%
\definecolor{currentfill}{rgb}{0.000000,0.000000,0.000000}%
\pgfsetfillcolor{currentfill}%
\pgfsetlinewidth{0.501875pt}%
\definecolor{currentstroke}{rgb}{0.000000,0.000000,0.000000}%
\pgfsetstrokecolor{currentstroke}%
\pgfsetdash{}{0pt}%
\pgfsys@defobject{currentmarker}{\pgfqpoint{0.000000in}{0.000000in}}{\pgfqpoint{0.020833in}{0.000000in}}{%
\pgfpathmoveto{\pgfqpoint{0.000000in}{0.000000in}}%
\pgfpathlineto{\pgfqpoint{0.020833in}{0.000000in}}%
\pgfusepath{stroke,fill}%
}%
\begin{pgfscope}%
\pgfsys@transformshift{0.455093in}{1.466590in}%
\pgfsys@useobject{currentmarker}{}%
\end{pgfscope}%
\end{pgfscope}%
\begin{pgfscope}%
\pgfsetbuttcap%
\pgfsetroundjoin%
\definecolor{currentfill}{rgb}{0.000000,0.000000,0.000000}%
\pgfsetfillcolor{currentfill}%
\pgfsetlinewidth{0.501875pt}%
\definecolor{currentstroke}{rgb}{0.000000,0.000000,0.000000}%
\pgfsetstrokecolor{currentstroke}%
\pgfsetdash{}{0pt}%
\pgfsys@defobject{currentmarker}{\pgfqpoint{-0.020833in}{0.000000in}}{\pgfqpoint{-0.000000in}{0.000000in}}{%
\pgfpathmoveto{\pgfqpoint{-0.000000in}{0.000000in}}%
\pgfpathlineto{\pgfqpoint{-0.020833in}{0.000000in}}%
\pgfusepath{stroke,fill}%
}%
\begin{pgfscope}%
\pgfsys@transformshift{5.492593in}{1.466590in}%
\pgfsys@useobject{currentmarker}{}%
\end{pgfscope}%
\end{pgfscope}%
\begin{pgfscope}%
\pgfsetbuttcap%
\pgfsetroundjoin%
\definecolor{currentfill}{rgb}{0.000000,0.000000,0.000000}%
\pgfsetfillcolor{currentfill}%
\pgfsetlinewidth{0.501875pt}%
\definecolor{currentstroke}{rgb}{0.000000,0.000000,0.000000}%
\pgfsetstrokecolor{currentstroke}%
\pgfsetdash{}{0pt}%
\pgfsys@defobject{currentmarker}{\pgfqpoint{0.000000in}{0.000000in}}{\pgfqpoint{0.020833in}{0.000000in}}{%
\pgfpathmoveto{\pgfqpoint{0.000000in}{0.000000in}}%
\pgfpathlineto{\pgfqpoint{0.020833in}{0.000000in}}%
\pgfusepath{stroke,fill}%
}%
\begin{pgfscope}%
\pgfsys@transformshift{0.455093in}{1.763620in}%
\pgfsys@useobject{currentmarker}{}%
\end{pgfscope}%
\end{pgfscope}%
\begin{pgfscope}%
\pgfsetbuttcap%
\pgfsetroundjoin%
\definecolor{currentfill}{rgb}{0.000000,0.000000,0.000000}%
\pgfsetfillcolor{currentfill}%
\pgfsetlinewidth{0.501875pt}%
\definecolor{currentstroke}{rgb}{0.000000,0.000000,0.000000}%
\pgfsetstrokecolor{currentstroke}%
\pgfsetdash{}{0pt}%
\pgfsys@defobject{currentmarker}{\pgfqpoint{-0.020833in}{0.000000in}}{\pgfqpoint{-0.000000in}{0.000000in}}{%
\pgfpathmoveto{\pgfqpoint{-0.000000in}{0.000000in}}%
\pgfpathlineto{\pgfqpoint{-0.020833in}{0.000000in}}%
\pgfusepath{stroke,fill}%
}%
\begin{pgfscope}%
\pgfsys@transformshift{5.492593in}{1.763620in}%
\pgfsys@useobject{currentmarker}{}%
\end{pgfscope}%
\end{pgfscope}%
\begin{pgfscope}%
\pgfsetbuttcap%
\pgfsetroundjoin%
\definecolor{currentfill}{rgb}{0.000000,0.000000,0.000000}%
\pgfsetfillcolor{currentfill}%
\pgfsetlinewidth{0.501875pt}%
\definecolor{currentstroke}{rgb}{0.000000,0.000000,0.000000}%
\pgfsetstrokecolor{currentstroke}%
\pgfsetdash{}{0pt}%
\pgfsys@defobject{currentmarker}{\pgfqpoint{0.000000in}{0.000000in}}{\pgfqpoint{0.020833in}{0.000000in}}{%
\pgfpathmoveto{\pgfqpoint{0.000000in}{0.000000in}}%
\pgfpathlineto{\pgfqpoint{0.020833in}{0.000000in}}%
\pgfusepath{stroke,fill}%
}%
\begin{pgfscope}%
\pgfsys@transformshift{0.455093in}{1.912136in}%
\pgfsys@useobject{currentmarker}{}%
\end{pgfscope}%
\end{pgfscope}%
\begin{pgfscope}%
\pgfsetbuttcap%
\pgfsetroundjoin%
\definecolor{currentfill}{rgb}{0.000000,0.000000,0.000000}%
\pgfsetfillcolor{currentfill}%
\pgfsetlinewidth{0.501875pt}%
\definecolor{currentstroke}{rgb}{0.000000,0.000000,0.000000}%
\pgfsetstrokecolor{currentstroke}%
\pgfsetdash{}{0pt}%
\pgfsys@defobject{currentmarker}{\pgfqpoint{-0.020833in}{0.000000in}}{\pgfqpoint{-0.000000in}{0.000000in}}{%
\pgfpathmoveto{\pgfqpoint{-0.000000in}{0.000000in}}%
\pgfpathlineto{\pgfqpoint{-0.020833in}{0.000000in}}%
\pgfusepath{stroke,fill}%
}%
\begin{pgfscope}%
\pgfsys@transformshift{5.492593in}{1.912136in}%
\pgfsys@useobject{currentmarker}{}%
\end{pgfscope}%
\end{pgfscope}%
\begin{pgfscope}%
\pgfsetbuttcap%
\pgfsetroundjoin%
\definecolor{currentfill}{rgb}{0.000000,0.000000,0.000000}%
\pgfsetfillcolor{currentfill}%
\pgfsetlinewidth{0.501875pt}%
\definecolor{currentstroke}{rgb}{0.000000,0.000000,0.000000}%
\pgfsetstrokecolor{currentstroke}%
\pgfsetdash{}{0pt}%
\pgfsys@defobject{currentmarker}{\pgfqpoint{0.000000in}{0.000000in}}{\pgfqpoint{0.020833in}{0.000000in}}{%
\pgfpathmoveto{\pgfqpoint{0.000000in}{0.000000in}}%
\pgfpathlineto{\pgfqpoint{0.020833in}{0.000000in}}%
\pgfusepath{stroke,fill}%
}%
\begin{pgfscope}%
\pgfsys@transformshift{0.455093in}{2.060651in}%
\pgfsys@useobject{currentmarker}{}%
\end{pgfscope}%
\end{pgfscope}%
\begin{pgfscope}%
\pgfsetbuttcap%
\pgfsetroundjoin%
\definecolor{currentfill}{rgb}{0.000000,0.000000,0.000000}%
\pgfsetfillcolor{currentfill}%
\pgfsetlinewidth{0.501875pt}%
\definecolor{currentstroke}{rgb}{0.000000,0.000000,0.000000}%
\pgfsetstrokecolor{currentstroke}%
\pgfsetdash{}{0pt}%
\pgfsys@defobject{currentmarker}{\pgfqpoint{-0.020833in}{0.000000in}}{\pgfqpoint{-0.000000in}{0.000000in}}{%
\pgfpathmoveto{\pgfqpoint{-0.000000in}{0.000000in}}%
\pgfpathlineto{\pgfqpoint{-0.020833in}{0.000000in}}%
\pgfusepath{stroke,fill}%
}%
\begin{pgfscope}%
\pgfsys@transformshift{5.492593in}{2.060651in}%
\pgfsys@useobject{currentmarker}{}%
\end{pgfscope}%
\end{pgfscope}%
\begin{pgfscope}%
\pgfsetbuttcap%
\pgfsetroundjoin%
\definecolor{currentfill}{rgb}{0.000000,0.000000,0.000000}%
\pgfsetfillcolor{currentfill}%
\pgfsetlinewidth{0.501875pt}%
\definecolor{currentstroke}{rgb}{0.000000,0.000000,0.000000}%
\pgfsetstrokecolor{currentstroke}%
\pgfsetdash{}{0pt}%
\pgfsys@defobject{currentmarker}{\pgfqpoint{0.000000in}{0.000000in}}{\pgfqpoint{0.020833in}{0.000000in}}{%
\pgfpathmoveto{\pgfqpoint{0.000000in}{0.000000in}}%
\pgfpathlineto{\pgfqpoint{0.020833in}{0.000000in}}%
\pgfusepath{stroke,fill}%
}%
\begin{pgfscope}%
\pgfsys@transformshift{0.455093in}{2.357682in}%
\pgfsys@useobject{currentmarker}{}%
\end{pgfscope}%
\end{pgfscope}%
\begin{pgfscope}%
\pgfsetbuttcap%
\pgfsetroundjoin%
\definecolor{currentfill}{rgb}{0.000000,0.000000,0.000000}%
\pgfsetfillcolor{currentfill}%
\pgfsetlinewidth{0.501875pt}%
\definecolor{currentstroke}{rgb}{0.000000,0.000000,0.000000}%
\pgfsetstrokecolor{currentstroke}%
\pgfsetdash{}{0pt}%
\pgfsys@defobject{currentmarker}{\pgfqpoint{-0.020833in}{0.000000in}}{\pgfqpoint{-0.000000in}{0.000000in}}{%
\pgfpathmoveto{\pgfqpoint{-0.000000in}{0.000000in}}%
\pgfpathlineto{\pgfqpoint{-0.020833in}{0.000000in}}%
\pgfusepath{stroke,fill}%
}%
\begin{pgfscope}%
\pgfsys@transformshift{5.492593in}{2.357682in}%
\pgfsys@useobject{currentmarker}{}%
\end{pgfscope}%
\end{pgfscope}%
\begin{pgfscope}%
\pgfsetbuttcap%
\pgfsetroundjoin%
\definecolor{currentfill}{rgb}{0.000000,0.000000,0.000000}%
\pgfsetfillcolor{currentfill}%
\pgfsetlinewidth{0.501875pt}%
\definecolor{currentstroke}{rgb}{0.000000,0.000000,0.000000}%
\pgfsetstrokecolor{currentstroke}%
\pgfsetdash{}{0pt}%
\pgfsys@defobject{currentmarker}{\pgfqpoint{0.000000in}{0.000000in}}{\pgfqpoint{0.020833in}{0.000000in}}{%
\pgfpathmoveto{\pgfqpoint{0.000000in}{0.000000in}}%
\pgfpathlineto{\pgfqpoint{0.020833in}{0.000000in}}%
\pgfusepath{stroke,fill}%
}%
\begin{pgfscope}%
\pgfsys@transformshift{0.455093in}{2.506197in}%
\pgfsys@useobject{currentmarker}{}%
\end{pgfscope}%
\end{pgfscope}%
\begin{pgfscope}%
\pgfsetbuttcap%
\pgfsetroundjoin%
\definecolor{currentfill}{rgb}{0.000000,0.000000,0.000000}%
\pgfsetfillcolor{currentfill}%
\pgfsetlinewidth{0.501875pt}%
\definecolor{currentstroke}{rgb}{0.000000,0.000000,0.000000}%
\pgfsetstrokecolor{currentstroke}%
\pgfsetdash{}{0pt}%
\pgfsys@defobject{currentmarker}{\pgfqpoint{-0.020833in}{0.000000in}}{\pgfqpoint{-0.000000in}{0.000000in}}{%
\pgfpathmoveto{\pgfqpoint{-0.000000in}{0.000000in}}%
\pgfpathlineto{\pgfqpoint{-0.020833in}{0.000000in}}%
\pgfusepath{stroke,fill}%
}%
\begin{pgfscope}%
\pgfsys@transformshift{5.492593in}{2.506197in}%
\pgfsys@useobject{currentmarker}{}%
\end{pgfscope}%
\end{pgfscope}%
\begin{pgfscope}%
\pgfsetbuttcap%
\pgfsetroundjoin%
\definecolor{currentfill}{rgb}{0.000000,0.000000,0.000000}%
\pgfsetfillcolor{currentfill}%
\pgfsetlinewidth{0.501875pt}%
\definecolor{currentstroke}{rgb}{0.000000,0.000000,0.000000}%
\pgfsetstrokecolor{currentstroke}%
\pgfsetdash{}{0pt}%
\pgfsys@defobject{currentmarker}{\pgfqpoint{0.000000in}{0.000000in}}{\pgfqpoint{0.020833in}{0.000000in}}{%
\pgfpathmoveto{\pgfqpoint{0.000000in}{0.000000in}}%
\pgfpathlineto{\pgfqpoint{0.020833in}{0.000000in}}%
\pgfusepath{stroke,fill}%
}%
\begin{pgfscope}%
\pgfsys@transformshift{0.455093in}{2.654712in}%
\pgfsys@useobject{currentmarker}{}%
\end{pgfscope}%
\end{pgfscope}%
\begin{pgfscope}%
\pgfsetbuttcap%
\pgfsetroundjoin%
\definecolor{currentfill}{rgb}{0.000000,0.000000,0.000000}%
\pgfsetfillcolor{currentfill}%
\pgfsetlinewidth{0.501875pt}%
\definecolor{currentstroke}{rgb}{0.000000,0.000000,0.000000}%
\pgfsetstrokecolor{currentstroke}%
\pgfsetdash{}{0pt}%
\pgfsys@defobject{currentmarker}{\pgfqpoint{-0.020833in}{0.000000in}}{\pgfqpoint{-0.000000in}{0.000000in}}{%
\pgfpathmoveto{\pgfqpoint{-0.000000in}{0.000000in}}%
\pgfpathlineto{\pgfqpoint{-0.020833in}{0.000000in}}%
\pgfusepath{stroke,fill}%
}%
\begin{pgfscope}%
\pgfsys@transformshift{5.492593in}{2.654712in}%
\pgfsys@useobject{currentmarker}{}%
\end{pgfscope}%
\end{pgfscope}%
\begin{pgfscope}%
\pgfsetbuttcap%
\pgfsetroundjoin%
\definecolor{currentfill}{rgb}{0.000000,0.000000,0.000000}%
\pgfsetfillcolor{currentfill}%
\pgfsetlinewidth{0.501875pt}%
\definecolor{currentstroke}{rgb}{0.000000,0.000000,0.000000}%
\pgfsetstrokecolor{currentstroke}%
\pgfsetdash{}{0pt}%
\pgfsys@defobject{currentmarker}{\pgfqpoint{0.000000in}{0.000000in}}{\pgfqpoint{0.020833in}{0.000000in}}{%
\pgfpathmoveto{\pgfqpoint{0.000000in}{0.000000in}}%
\pgfpathlineto{\pgfqpoint{0.020833in}{0.000000in}}%
\pgfusepath{stroke,fill}%
}%
\begin{pgfscope}%
\pgfsys@transformshift{0.455093in}{2.951743in}%
\pgfsys@useobject{currentmarker}{}%
\end{pgfscope}%
\end{pgfscope}%
\begin{pgfscope}%
\pgfsetbuttcap%
\pgfsetroundjoin%
\definecolor{currentfill}{rgb}{0.000000,0.000000,0.000000}%
\pgfsetfillcolor{currentfill}%
\pgfsetlinewidth{0.501875pt}%
\definecolor{currentstroke}{rgb}{0.000000,0.000000,0.000000}%
\pgfsetstrokecolor{currentstroke}%
\pgfsetdash{}{0pt}%
\pgfsys@defobject{currentmarker}{\pgfqpoint{-0.020833in}{0.000000in}}{\pgfqpoint{-0.000000in}{0.000000in}}{%
\pgfpathmoveto{\pgfqpoint{-0.000000in}{0.000000in}}%
\pgfpathlineto{\pgfqpoint{-0.020833in}{0.000000in}}%
\pgfusepath{stroke,fill}%
}%
\begin{pgfscope}%
\pgfsys@transformshift{5.492593in}{2.951743in}%
\pgfsys@useobject{currentmarker}{}%
\end{pgfscope}%
\end{pgfscope}%
\begin{pgfscope}%
\pgfsetbuttcap%
\pgfsetroundjoin%
\definecolor{currentfill}{rgb}{0.000000,0.000000,0.000000}%
\pgfsetfillcolor{currentfill}%
\pgfsetlinewidth{0.501875pt}%
\definecolor{currentstroke}{rgb}{0.000000,0.000000,0.000000}%
\pgfsetstrokecolor{currentstroke}%
\pgfsetdash{}{0pt}%
\pgfsys@defobject{currentmarker}{\pgfqpoint{0.000000in}{0.000000in}}{\pgfqpoint{0.020833in}{0.000000in}}{%
\pgfpathmoveto{\pgfqpoint{0.000000in}{0.000000in}}%
\pgfpathlineto{\pgfqpoint{0.020833in}{0.000000in}}%
\pgfusepath{stroke,fill}%
}%
\begin{pgfscope}%
\pgfsys@transformshift{0.455093in}{3.100258in}%
\pgfsys@useobject{currentmarker}{}%
\end{pgfscope}%
\end{pgfscope}%
\begin{pgfscope}%
\pgfsetbuttcap%
\pgfsetroundjoin%
\definecolor{currentfill}{rgb}{0.000000,0.000000,0.000000}%
\pgfsetfillcolor{currentfill}%
\pgfsetlinewidth{0.501875pt}%
\definecolor{currentstroke}{rgb}{0.000000,0.000000,0.000000}%
\pgfsetstrokecolor{currentstroke}%
\pgfsetdash{}{0pt}%
\pgfsys@defobject{currentmarker}{\pgfqpoint{-0.020833in}{0.000000in}}{\pgfqpoint{-0.000000in}{0.000000in}}{%
\pgfpathmoveto{\pgfqpoint{-0.000000in}{0.000000in}}%
\pgfpathlineto{\pgfqpoint{-0.020833in}{0.000000in}}%
\pgfusepath{stroke,fill}%
}%
\begin{pgfscope}%
\pgfsys@transformshift{5.492593in}{3.100258in}%
\pgfsys@useobject{currentmarker}{}%
\end{pgfscope}%
\end{pgfscope}%
\begin{pgfscope}%
\pgfsetbuttcap%
\pgfsetroundjoin%
\definecolor{currentfill}{rgb}{0.000000,0.000000,0.000000}%
\pgfsetfillcolor{currentfill}%
\pgfsetlinewidth{0.501875pt}%
\definecolor{currentstroke}{rgb}{0.000000,0.000000,0.000000}%
\pgfsetstrokecolor{currentstroke}%
\pgfsetdash{}{0pt}%
\pgfsys@defobject{currentmarker}{\pgfqpoint{0.000000in}{0.000000in}}{\pgfqpoint{0.020833in}{0.000000in}}{%
\pgfpathmoveto{\pgfqpoint{0.000000in}{0.000000in}}%
\pgfpathlineto{\pgfqpoint{0.020833in}{0.000000in}}%
\pgfusepath{stroke,fill}%
}%
\begin{pgfscope}%
\pgfsys@transformshift{0.455093in}{3.248774in}%
\pgfsys@useobject{currentmarker}{}%
\end{pgfscope}%
\end{pgfscope}%
\begin{pgfscope}%
\pgfsetbuttcap%
\pgfsetroundjoin%
\definecolor{currentfill}{rgb}{0.000000,0.000000,0.000000}%
\pgfsetfillcolor{currentfill}%
\pgfsetlinewidth{0.501875pt}%
\definecolor{currentstroke}{rgb}{0.000000,0.000000,0.000000}%
\pgfsetstrokecolor{currentstroke}%
\pgfsetdash{}{0pt}%
\pgfsys@defobject{currentmarker}{\pgfqpoint{-0.020833in}{0.000000in}}{\pgfqpoint{-0.000000in}{0.000000in}}{%
\pgfpathmoveto{\pgfqpoint{-0.000000in}{0.000000in}}%
\pgfpathlineto{\pgfqpoint{-0.020833in}{0.000000in}}%
\pgfusepath{stroke,fill}%
}%
\begin{pgfscope}%
\pgfsys@transformshift{5.492593in}{3.248774in}%
\pgfsys@useobject{currentmarker}{}%
\end{pgfscope}%
\end{pgfscope}%
\begin{pgfscope}%
\pgfsetbuttcap%
\pgfsetroundjoin%
\definecolor{currentfill}{rgb}{0.000000,0.000000,0.000000}%
\pgfsetfillcolor{currentfill}%
\pgfsetlinewidth{0.501875pt}%
\definecolor{currentstroke}{rgb}{0.000000,0.000000,0.000000}%
\pgfsetstrokecolor{currentstroke}%
\pgfsetdash{}{0pt}%
\pgfsys@defobject{currentmarker}{\pgfqpoint{0.000000in}{0.000000in}}{\pgfqpoint{0.020833in}{0.000000in}}{%
\pgfpathmoveto{\pgfqpoint{0.000000in}{0.000000in}}%
\pgfpathlineto{\pgfqpoint{0.020833in}{0.000000in}}%
\pgfusepath{stroke,fill}%
}%
\begin{pgfscope}%
\pgfsys@transformshift{0.455093in}{3.545804in}%
\pgfsys@useobject{currentmarker}{}%
\end{pgfscope}%
\end{pgfscope}%
\begin{pgfscope}%
\pgfsetbuttcap%
\pgfsetroundjoin%
\definecolor{currentfill}{rgb}{0.000000,0.000000,0.000000}%
\pgfsetfillcolor{currentfill}%
\pgfsetlinewidth{0.501875pt}%
\definecolor{currentstroke}{rgb}{0.000000,0.000000,0.000000}%
\pgfsetstrokecolor{currentstroke}%
\pgfsetdash{}{0pt}%
\pgfsys@defobject{currentmarker}{\pgfqpoint{-0.020833in}{0.000000in}}{\pgfqpoint{-0.000000in}{0.000000in}}{%
\pgfpathmoveto{\pgfqpoint{-0.000000in}{0.000000in}}%
\pgfpathlineto{\pgfqpoint{-0.020833in}{0.000000in}}%
\pgfusepath{stroke,fill}%
}%
\begin{pgfscope}%
\pgfsys@transformshift{5.492593in}{3.545804in}%
\pgfsys@useobject{currentmarker}{}%
\end{pgfscope}%
\end{pgfscope}%
\begin{pgfscope}%
\pgfsetbuttcap%
\pgfsetroundjoin%
\definecolor{currentfill}{rgb}{0.000000,0.000000,0.000000}%
\pgfsetfillcolor{currentfill}%
\pgfsetlinewidth{0.501875pt}%
\definecolor{currentstroke}{rgb}{0.000000,0.000000,0.000000}%
\pgfsetstrokecolor{currentstroke}%
\pgfsetdash{}{0pt}%
\pgfsys@defobject{currentmarker}{\pgfqpoint{0.000000in}{0.000000in}}{\pgfqpoint{0.020833in}{0.000000in}}{%
\pgfpathmoveto{\pgfqpoint{0.000000in}{0.000000in}}%
\pgfpathlineto{\pgfqpoint{0.020833in}{0.000000in}}%
\pgfusepath{stroke,fill}%
}%
\begin{pgfscope}%
\pgfsys@transformshift{0.455093in}{3.694319in}%
\pgfsys@useobject{currentmarker}{}%
\end{pgfscope}%
\end{pgfscope}%
\begin{pgfscope}%
\pgfsetbuttcap%
\pgfsetroundjoin%
\definecolor{currentfill}{rgb}{0.000000,0.000000,0.000000}%
\pgfsetfillcolor{currentfill}%
\pgfsetlinewidth{0.501875pt}%
\definecolor{currentstroke}{rgb}{0.000000,0.000000,0.000000}%
\pgfsetstrokecolor{currentstroke}%
\pgfsetdash{}{0pt}%
\pgfsys@defobject{currentmarker}{\pgfqpoint{-0.020833in}{0.000000in}}{\pgfqpoint{-0.000000in}{0.000000in}}{%
\pgfpathmoveto{\pgfqpoint{-0.000000in}{0.000000in}}%
\pgfpathlineto{\pgfqpoint{-0.020833in}{0.000000in}}%
\pgfusepath{stroke,fill}%
}%
\begin{pgfscope}%
\pgfsys@transformshift{5.492593in}{3.694319in}%
\pgfsys@useobject{currentmarker}{}%
\end{pgfscope}%
\end{pgfscope}%
\begin{pgfscope}%
\pgfsetbuttcap%
\pgfsetroundjoin%
\definecolor{currentfill}{rgb}{0.000000,0.000000,0.000000}%
\pgfsetfillcolor{currentfill}%
\pgfsetlinewidth{0.501875pt}%
\definecolor{currentstroke}{rgb}{0.000000,0.000000,0.000000}%
\pgfsetstrokecolor{currentstroke}%
\pgfsetdash{}{0pt}%
\pgfsys@defobject{currentmarker}{\pgfqpoint{0.000000in}{0.000000in}}{\pgfqpoint{0.020833in}{0.000000in}}{%
\pgfpathmoveto{\pgfqpoint{0.000000in}{0.000000in}}%
\pgfpathlineto{\pgfqpoint{0.020833in}{0.000000in}}%
\pgfusepath{stroke,fill}%
}%
\begin{pgfscope}%
\pgfsys@transformshift{0.455093in}{3.842835in}%
\pgfsys@useobject{currentmarker}{}%
\end{pgfscope}%
\end{pgfscope}%
\begin{pgfscope}%
\pgfsetbuttcap%
\pgfsetroundjoin%
\definecolor{currentfill}{rgb}{0.000000,0.000000,0.000000}%
\pgfsetfillcolor{currentfill}%
\pgfsetlinewidth{0.501875pt}%
\definecolor{currentstroke}{rgb}{0.000000,0.000000,0.000000}%
\pgfsetstrokecolor{currentstroke}%
\pgfsetdash{}{0pt}%
\pgfsys@defobject{currentmarker}{\pgfqpoint{-0.020833in}{0.000000in}}{\pgfqpoint{-0.000000in}{0.000000in}}{%
\pgfpathmoveto{\pgfqpoint{-0.000000in}{0.000000in}}%
\pgfpathlineto{\pgfqpoint{-0.020833in}{0.000000in}}%
\pgfusepath{stroke,fill}%
}%
\begin{pgfscope}%
\pgfsys@transformshift{5.492593in}{3.842835in}%
\pgfsys@useobject{currentmarker}{}%
\end{pgfscope}%
\end{pgfscope}%
\begin{pgfscope}%
\pgfsetbuttcap%
\pgfsetroundjoin%
\definecolor{currentfill}{rgb}{0.000000,0.000000,0.000000}%
\pgfsetfillcolor{currentfill}%
\pgfsetlinewidth{0.501875pt}%
\definecolor{currentstroke}{rgb}{0.000000,0.000000,0.000000}%
\pgfsetstrokecolor{currentstroke}%
\pgfsetdash{}{0pt}%
\pgfsys@defobject{currentmarker}{\pgfqpoint{0.000000in}{0.000000in}}{\pgfqpoint{0.020833in}{0.000000in}}{%
\pgfpathmoveto{\pgfqpoint{0.000000in}{0.000000in}}%
\pgfpathlineto{\pgfqpoint{0.020833in}{0.000000in}}%
\pgfusepath{stroke,fill}%
}%
\begin{pgfscope}%
\pgfsys@transformshift{0.455093in}{4.139865in}%
\pgfsys@useobject{currentmarker}{}%
\end{pgfscope}%
\end{pgfscope}%
\begin{pgfscope}%
\pgfsetbuttcap%
\pgfsetroundjoin%
\definecolor{currentfill}{rgb}{0.000000,0.000000,0.000000}%
\pgfsetfillcolor{currentfill}%
\pgfsetlinewidth{0.501875pt}%
\definecolor{currentstroke}{rgb}{0.000000,0.000000,0.000000}%
\pgfsetstrokecolor{currentstroke}%
\pgfsetdash{}{0pt}%
\pgfsys@defobject{currentmarker}{\pgfqpoint{-0.020833in}{0.000000in}}{\pgfqpoint{-0.000000in}{0.000000in}}{%
\pgfpathmoveto{\pgfqpoint{-0.000000in}{0.000000in}}%
\pgfpathlineto{\pgfqpoint{-0.020833in}{0.000000in}}%
\pgfusepath{stroke,fill}%
}%
\begin{pgfscope}%
\pgfsys@transformshift{5.492593in}{4.139865in}%
\pgfsys@useobject{currentmarker}{}%
\end{pgfscope}%
\end{pgfscope}%
\begin{pgfscope}%
\pgfsetbuttcap%
\pgfsetroundjoin%
\definecolor{currentfill}{rgb}{0.000000,0.000000,0.000000}%
\pgfsetfillcolor{currentfill}%
\pgfsetlinewidth{0.501875pt}%
\definecolor{currentstroke}{rgb}{0.000000,0.000000,0.000000}%
\pgfsetstrokecolor{currentstroke}%
\pgfsetdash{}{0pt}%
\pgfsys@defobject{currentmarker}{\pgfqpoint{0.000000in}{0.000000in}}{\pgfqpoint{0.020833in}{0.000000in}}{%
\pgfpathmoveto{\pgfqpoint{0.000000in}{0.000000in}}%
\pgfpathlineto{\pgfqpoint{0.020833in}{0.000000in}}%
\pgfusepath{stroke,fill}%
}%
\begin{pgfscope}%
\pgfsys@transformshift{0.455093in}{4.288381in}%
\pgfsys@useobject{currentmarker}{}%
\end{pgfscope}%
\end{pgfscope}%
\begin{pgfscope}%
\pgfsetbuttcap%
\pgfsetroundjoin%
\definecolor{currentfill}{rgb}{0.000000,0.000000,0.000000}%
\pgfsetfillcolor{currentfill}%
\pgfsetlinewidth{0.501875pt}%
\definecolor{currentstroke}{rgb}{0.000000,0.000000,0.000000}%
\pgfsetstrokecolor{currentstroke}%
\pgfsetdash{}{0pt}%
\pgfsys@defobject{currentmarker}{\pgfqpoint{-0.020833in}{0.000000in}}{\pgfqpoint{-0.000000in}{0.000000in}}{%
\pgfpathmoveto{\pgfqpoint{-0.000000in}{0.000000in}}%
\pgfpathlineto{\pgfqpoint{-0.020833in}{0.000000in}}%
\pgfusepath{stroke,fill}%
}%
\begin{pgfscope}%
\pgfsys@transformshift{5.492593in}{4.288381in}%
\pgfsys@useobject{currentmarker}{}%
\end{pgfscope}%
\end{pgfscope}%
\begin{pgfscope}%
\definecolor{textcolor}{rgb}{0.000000,0.000000,0.000000}%
\pgfsetstrokecolor{textcolor}%
\pgfsetfillcolor{textcolor}%
\pgftext[x=0.173457in,y=2.403080in,,bottom,rotate=90.000000]{\color{textcolor}{\rmfamily\fontsize{10.000000}{12.000000}\selectfont\catcode`\^=\active\def^{\ifmmode\sp\else\^{}\fi}\catcode`\%=\active\def%{\%}Geometric Brownian motion}}%
\end{pgfscope}%
\begin{pgfscope}%
\pgfpathrectangle{\pgfqpoint{0.455093in}{0.401080in}}{\pgfqpoint{5.037500in}{4.004000in}}%
\pgfusepath{clip}%
\pgfsetrectcap%
\pgfsetroundjoin%
\pgfsetlinewidth{0.501875pt}%
\definecolor{currentstroke}{rgb}{0.690196,0.188235,0.376471}%
\pgfsetstrokecolor{currentstroke}%
\pgfsetstrokeopacity{0.800000}%
\pgfsetdash{}{0pt}%
\pgfpathmoveto{\pgfqpoint{0.455093in}{1.021044in}}%
\pgfpathlineto{\pgfqpoint{0.457612in}{1.014418in}}%
\pgfpathlineto{\pgfqpoint{0.460131in}{1.005971in}}%
\pgfpathlineto{\pgfqpoint{0.462649in}{1.016936in}}%
\pgfpathlineto{\pgfqpoint{0.465168in}{1.006634in}}%
\pgfpathlineto{\pgfqpoint{0.467687in}{1.021069in}}%
\pgfpathlineto{\pgfqpoint{0.470206in}{1.026624in}}%
\pgfpathlineto{\pgfqpoint{0.472724in}{1.013023in}}%
\pgfpathlineto{\pgfqpoint{0.475243in}{0.995848in}}%
\pgfpathlineto{\pgfqpoint{0.477762in}{0.983295in}}%
\pgfpathlineto{\pgfqpoint{0.480281in}{0.996196in}}%
\pgfpathlineto{\pgfqpoint{0.485318in}{1.005158in}}%
\pgfpathlineto{\pgfqpoint{0.487837in}{1.014259in}}%
\pgfpathlineto{\pgfqpoint{0.490356in}{1.025024in}}%
\pgfpathlineto{\pgfqpoint{0.492874in}{1.045614in}}%
\pgfpathlineto{\pgfqpoint{0.495393in}{1.054364in}}%
\pgfpathlineto{\pgfqpoint{0.497912in}{1.029776in}}%
\pgfpathlineto{\pgfqpoint{0.500431in}{1.038552in}}%
\pgfpathlineto{\pgfqpoint{0.502949in}{1.016550in}}%
\pgfpathlineto{\pgfqpoint{0.505468in}{1.017824in}}%
\pgfpathlineto{\pgfqpoint{0.507987in}{1.019950in}}%
\pgfpathlineto{\pgfqpoint{0.510506in}{1.019687in}}%
\pgfpathlineto{\pgfqpoint{0.513024in}{1.019299in}}%
\pgfpathlineto{\pgfqpoint{0.515543in}{1.017753in}}%
\pgfpathlineto{\pgfqpoint{0.518062in}{1.005503in}}%
\pgfpathlineto{\pgfqpoint{0.520581in}{0.999658in}}%
\pgfpathlineto{\pgfqpoint{0.523099in}{1.007097in}}%
\pgfpathlineto{\pgfqpoint{0.525618in}{1.015797in}}%
\pgfpathlineto{\pgfqpoint{0.528137in}{1.012279in}}%
\pgfpathlineto{\pgfqpoint{0.530656in}{1.002337in}}%
\pgfpathlineto{\pgfqpoint{0.533174in}{0.986026in}}%
\pgfpathlineto{\pgfqpoint{0.535693in}{0.997091in}}%
\pgfpathlineto{\pgfqpoint{0.538212in}{0.970240in}}%
\pgfpathlineto{\pgfqpoint{0.540731in}{0.963248in}}%
\pgfpathlineto{\pgfqpoint{0.543249in}{0.952804in}}%
\pgfpathlineto{\pgfqpoint{0.545768in}{0.967743in}}%
\pgfpathlineto{\pgfqpoint{0.548287in}{0.984140in}}%
\pgfpathlineto{\pgfqpoint{0.550806in}{0.991142in}}%
\pgfpathlineto{\pgfqpoint{0.553324in}{0.990753in}}%
\pgfpathlineto{\pgfqpoint{0.555843in}{0.991172in}}%
\pgfpathlineto{\pgfqpoint{0.558362in}{0.989196in}}%
\pgfpathlineto{\pgfqpoint{0.560881in}{0.971050in}}%
\pgfpathlineto{\pgfqpoint{0.563399in}{0.964793in}}%
\pgfpathlineto{\pgfqpoint{0.565918in}{0.958108in}}%
\pgfpathlineto{\pgfqpoint{0.568437in}{0.966261in}}%
\pgfpathlineto{\pgfqpoint{0.570956in}{0.976034in}}%
\pgfpathlineto{\pgfqpoint{0.573474in}{0.989991in}}%
\pgfpathlineto{\pgfqpoint{0.575993in}{0.993912in}}%
\pgfpathlineto{\pgfqpoint{0.578512in}{1.000161in}}%
\pgfpathlineto{\pgfqpoint{0.581031in}{1.000309in}}%
\pgfpathlineto{\pgfqpoint{0.583549in}{1.004778in}}%
\pgfpathlineto{\pgfqpoint{0.586068in}{1.018335in}}%
\pgfpathlineto{\pgfqpoint{0.588587in}{1.024145in}}%
\pgfpathlineto{\pgfqpoint{0.591106in}{1.020212in}}%
\pgfpathlineto{\pgfqpoint{0.593624in}{1.025914in}}%
\pgfpathlineto{\pgfqpoint{0.596143in}{1.016417in}}%
\pgfpathlineto{\pgfqpoint{0.598662in}{1.007881in}}%
\pgfpathlineto{\pgfqpoint{0.601181in}{1.005637in}}%
\pgfpathlineto{\pgfqpoint{0.603699in}{1.013961in}}%
\pgfpathlineto{\pgfqpoint{0.606218in}{1.030334in}}%
\pgfpathlineto{\pgfqpoint{0.608737in}{1.008682in}}%
\pgfpathlineto{\pgfqpoint{0.611256in}{1.000285in}}%
\pgfpathlineto{\pgfqpoint{0.613774in}{1.010393in}}%
\pgfpathlineto{\pgfqpoint{0.616293in}{1.006205in}}%
\pgfpathlineto{\pgfqpoint{0.618812in}{1.034661in}}%
\pgfpathlineto{\pgfqpoint{0.621331in}{1.034391in}}%
\pgfpathlineto{\pgfqpoint{0.623849in}{1.024345in}}%
\pgfpathlineto{\pgfqpoint{0.626368in}{1.024724in}}%
\pgfpathlineto{\pgfqpoint{0.628887in}{1.027572in}}%
\pgfpathlineto{\pgfqpoint{0.631406in}{1.031015in}}%
\pgfpathlineto{\pgfqpoint{0.633924in}{1.022191in}}%
\pgfpathlineto{\pgfqpoint{0.636443in}{1.044936in}}%
\pgfpathlineto{\pgfqpoint{0.638962in}{1.049092in}}%
\pgfpathlineto{\pgfqpoint{0.641481in}{1.056634in}}%
\pgfpathlineto{\pgfqpoint{0.643999in}{1.052346in}}%
\pgfpathlineto{\pgfqpoint{0.646518in}{1.079079in}}%
\pgfpathlineto{\pgfqpoint{0.649037in}{1.081326in}}%
\pgfpathlineto{\pgfqpoint{0.651556in}{1.057840in}}%
\pgfpathlineto{\pgfqpoint{0.654074in}{1.057487in}}%
\pgfpathlineto{\pgfqpoint{0.656593in}{1.038956in}}%
\pgfpathlineto{\pgfqpoint{0.659112in}{1.049762in}}%
\pgfpathlineto{\pgfqpoint{0.661631in}{1.057250in}}%
\pgfpathlineto{\pgfqpoint{0.664149in}{1.053456in}}%
\pgfpathlineto{\pgfqpoint{0.666668in}{1.066612in}}%
\pgfpathlineto{\pgfqpoint{0.669187in}{1.056566in}}%
\pgfpathlineto{\pgfqpoint{0.671706in}{1.069794in}}%
\pgfpathlineto{\pgfqpoint{0.674224in}{1.085682in}}%
\pgfpathlineto{\pgfqpoint{0.676743in}{1.079065in}}%
\pgfpathlineto{\pgfqpoint{0.679262in}{1.093247in}}%
\pgfpathlineto{\pgfqpoint{0.681781in}{1.071860in}}%
\pgfpathlineto{\pgfqpoint{0.684299in}{1.074627in}}%
\pgfpathlineto{\pgfqpoint{0.686818in}{1.092746in}}%
\pgfpathlineto{\pgfqpoint{0.689337in}{1.084052in}}%
\pgfpathlineto{\pgfqpoint{0.691856in}{1.063139in}}%
\pgfpathlineto{\pgfqpoint{0.694374in}{1.075054in}}%
\pgfpathlineto{\pgfqpoint{0.696893in}{1.078998in}}%
\pgfpathlineto{\pgfqpoint{0.699412in}{1.079971in}}%
\pgfpathlineto{\pgfqpoint{0.701931in}{1.062241in}}%
\pgfpathlineto{\pgfqpoint{0.704449in}{1.039526in}}%
\pgfpathlineto{\pgfqpoint{0.706968in}{1.052371in}}%
\pgfpathlineto{\pgfqpoint{0.709487in}{1.068133in}}%
\pgfpathlineto{\pgfqpoint{0.712006in}{1.068199in}}%
\pgfpathlineto{\pgfqpoint{0.714524in}{1.062053in}}%
\pgfpathlineto{\pgfqpoint{0.717043in}{1.036047in}}%
\pgfpathlineto{\pgfqpoint{0.719562in}{1.024098in}}%
\pgfpathlineto{\pgfqpoint{0.722081in}{1.019893in}}%
\pgfpathlineto{\pgfqpoint{0.724599in}{1.009229in}}%
\pgfpathlineto{\pgfqpoint{0.727118in}{1.013845in}}%
\pgfpathlineto{\pgfqpoint{0.729637in}{1.014489in}}%
\pgfpathlineto{\pgfqpoint{0.732156in}{1.015373in}}%
\pgfpathlineto{\pgfqpoint{0.734674in}{1.011004in}}%
\pgfpathlineto{\pgfqpoint{0.737193in}{1.012192in}}%
\pgfpathlineto{\pgfqpoint{0.739712in}{1.016348in}}%
\pgfpathlineto{\pgfqpoint{0.742231in}{1.018681in}}%
\pgfpathlineto{\pgfqpoint{0.744749in}{1.029832in}}%
\pgfpathlineto{\pgfqpoint{0.749787in}{1.069344in}}%
\pgfpathlineto{\pgfqpoint{0.752306in}{1.053194in}}%
\pgfpathlineto{\pgfqpoint{0.754824in}{1.068149in}}%
\pgfpathlineto{\pgfqpoint{0.757343in}{1.068619in}}%
\pgfpathlineto{\pgfqpoint{0.759862in}{1.077958in}}%
\pgfpathlineto{\pgfqpoint{0.762381in}{1.103779in}}%
\pgfpathlineto{\pgfqpoint{0.764899in}{1.126639in}}%
\pgfpathlineto{\pgfqpoint{0.767418in}{1.108444in}}%
\pgfpathlineto{\pgfqpoint{0.769937in}{1.086440in}}%
\pgfpathlineto{\pgfqpoint{0.772456in}{1.108092in}}%
\pgfpathlineto{\pgfqpoint{0.774974in}{1.096669in}}%
\pgfpathlineto{\pgfqpoint{0.777493in}{1.109027in}}%
\pgfpathlineto{\pgfqpoint{0.780012in}{1.104184in}}%
\pgfpathlineto{\pgfqpoint{0.782531in}{1.100551in}}%
\pgfpathlineto{\pgfqpoint{0.785049in}{1.115296in}}%
\pgfpathlineto{\pgfqpoint{0.787568in}{1.110663in}}%
\pgfpathlineto{\pgfqpoint{0.790087in}{1.094121in}}%
\pgfpathlineto{\pgfqpoint{0.792606in}{1.112731in}}%
\pgfpathlineto{\pgfqpoint{0.795124in}{1.092467in}}%
\pgfpathlineto{\pgfqpoint{0.797643in}{1.099403in}}%
\pgfpathlineto{\pgfqpoint{0.800162in}{1.082272in}}%
\pgfpathlineto{\pgfqpoint{0.802681in}{1.063000in}}%
\pgfpathlineto{\pgfqpoint{0.805199in}{1.087512in}}%
\pgfpathlineto{\pgfqpoint{0.807718in}{1.079884in}}%
\pgfpathlineto{\pgfqpoint{0.810237in}{1.067202in}}%
\pgfpathlineto{\pgfqpoint{0.812756in}{1.081408in}}%
\pgfpathlineto{\pgfqpoint{0.815274in}{1.065547in}}%
\pgfpathlineto{\pgfqpoint{0.817793in}{1.046738in}}%
\pgfpathlineto{\pgfqpoint{0.820312in}{1.034389in}}%
\pgfpathlineto{\pgfqpoint{0.822831in}{1.040236in}}%
\pgfpathlineto{\pgfqpoint{0.825349in}{1.030758in}}%
\pgfpathlineto{\pgfqpoint{0.827868in}{1.025567in}}%
\pgfpathlineto{\pgfqpoint{0.830387in}{0.993501in}}%
\pgfpathlineto{\pgfqpoint{0.832906in}{1.012155in}}%
\pgfpathlineto{\pgfqpoint{0.835424in}{1.008882in}}%
\pgfpathlineto{\pgfqpoint{0.837943in}{1.025341in}}%
\pgfpathlineto{\pgfqpoint{0.840462in}{1.030689in}}%
\pgfpathlineto{\pgfqpoint{0.842981in}{1.015985in}}%
\pgfpathlineto{\pgfqpoint{0.845499in}{1.008594in}}%
\pgfpathlineto{\pgfqpoint{0.848018in}{0.978248in}}%
\pgfpathlineto{\pgfqpoint{0.850537in}{0.957872in}}%
\pgfpathlineto{\pgfqpoint{0.853056in}{0.954399in}}%
\pgfpathlineto{\pgfqpoint{0.855574in}{0.952267in}}%
\pgfpathlineto{\pgfqpoint{0.858093in}{0.939704in}}%
\pgfpathlineto{\pgfqpoint{0.860612in}{0.925907in}}%
\pgfpathlineto{\pgfqpoint{0.863131in}{0.943513in}}%
\pgfpathlineto{\pgfqpoint{0.865649in}{0.952601in}}%
\pgfpathlineto{\pgfqpoint{0.868168in}{0.941732in}}%
\pgfpathlineto{\pgfqpoint{0.870687in}{0.972939in}}%
\pgfpathlineto{\pgfqpoint{0.873206in}{0.986037in}}%
\pgfpathlineto{\pgfqpoint{0.875724in}{1.010038in}}%
\pgfpathlineto{\pgfqpoint{0.878243in}{0.997822in}}%
\pgfpathlineto{\pgfqpoint{0.880762in}{1.004033in}}%
\pgfpathlineto{\pgfqpoint{0.883281in}{1.002448in}}%
\pgfpathlineto{\pgfqpoint{0.885799in}{1.001698in}}%
\pgfpathlineto{\pgfqpoint{0.888318in}{1.002357in}}%
\pgfpathlineto{\pgfqpoint{0.890837in}{0.987891in}}%
\pgfpathlineto{\pgfqpoint{0.893356in}{1.009451in}}%
\pgfpathlineto{\pgfqpoint{0.895874in}{0.993120in}}%
\pgfpathlineto{\pgfqpoint{0.898393in}{1.001029in}}%
\pgfpathlineto{\pgfqpoint{0.900912in}{1.021415in}}%
\pgfpathlineto{\pgfqpoint{0.903431in}{1.045602in}}%
\pgfpathlineto{\pgfqpoint{0.905949in}{1.051099in}}%
\pgfpathlineto{\pgfqpoint{0.908468in}{1.051994in}}%
\pgfpathlineto{\pgfqpoint{0.910987in}{1.051341in}}%
\pgfpathlineto{\pgfqpoint{0.913506in}{1.034677in}}%
\pgfpathlineto{\pgfqpoint{0.916024in}{1.027806in}}%
\pgfpathlineto{\pgfqpoint{0.918543in}{1.026619in}}%
\pgfpathlineto{\pgfqpoint{0.921062in}{1.030776in}}%
\pgfpathlineto{\pgfqpoint{0.923581in}{1.004680in}}%
\pgfpathlineto{\pgfqpoint{0.926099in}{0.994167in}}%
\pgfpathlineto{\pgfqpoint{0.928618in}{1.002518in}}%
\pgfpathlineto{\pgfqpoint{0.931137in}{0.996061in}}%
\pgfpathlineto{\pgfqpoint{0.933656in}{0.969283in}}%
\pgfpathlineto{\pgfqpoint{0.936174in}{0.929102in}}%
\pgfpathlineto{\pgfqpoint{0.938693in}{0.927366in}}%
\pgfpathlineto{\pgfqpoint{0.941212in}{0.949395in}}%
\pgfpathlineto{\pgfqpoint{0.943731in}{0.964762in}}%
\pgfpathlineto{\pgfqpoint{0.946249in}{0.969289in}}%
\pgfpathlineto{\pgfqpoint{0.948768in}{0.978076in}}%
\pgfpathlineto{\pgfqpoint{0.951287in}{0.959469in}}%
\pgfpathlineto{\pgfqpoint{0.953806in}{0.950980in}}%
\pgfpathlineto{\pgfqpoint{0.956324in}{0.950867in}}%
\pgfpathlineto{\pgfqpoint{0.958843in}{0.946515in}}%
\pgfpathlineto{\pgfqpoint{0.961362in}{0.943137in}}%
\pgfpathlineto{\pgfqpoint{0.963881in}{0.953208in}}%
\pgfpathlineto{\pgfqpoint{0.966399in}{0.924951in}}%
\pgfpathlineto{\pgfqpoint{0.968918in}{0.931832in}}%
\pgfpathlineto{\pgfqpoint{0.971437in}{0.942488in}}%
\pgfpathlineto{\pgfqpoint{0.973956in}{0.933119in}}%
\pgfpathlineto{\pgfqpoint{0.976474in}{0.942379in}}%
\pgfpathlineto{\pgfqpoint{0.978993in}{0.952139in}}%
\pgfpathlineto{\pgfqpoint{0.981512in}{0.922482in}}%
\pgfpathlineto{\pgfqpoint{0.984031in}{0.900534in}}%
\pgfpathlineto{\pgfqpoint{0.986549in}{0.921965in}}%
\pgfpathlineto{\pgfqpoint{0.989068in}{0.910591in}}%
\pgfpathlineto{\pgfqpoint{0.991587in}{0.914260in}}%
\pgfpathlineto{\pgfqpoint{0.994106in}{0.918333in}}%
\pgfpathlineto{\pgfqpoint{0.996624in}{0.912567in}}%
\pgfpathlineto{\pgfqpoint{0.999143in}{0.905721in}}%
\pgfpathlineto{\pgfqpoint{1.001662in}{0.904096in}}%
\pgfpathlineto{\pgfqpoint{1.004181in}{0.903603in}}%
\pgfpathlineto{\pgfqpoint{1.006699in}{0.911792in}}%
\pgfpathlineto{\pgfqpoint{1.009218in}{0.933058in}}%
\pgfpathlineto{\pgfqpoint{1.011737in}{0.934070in}}%
\pgfpathlineto{\pgfqpoint{1.014256in}{0.909599in}}%
\pgfpathlineto{\pgfqpoint{1.016774in}{0.892069in}}%
\pgfpathlineto{\pgfqpoint{1.019293in}{0.909260in}}%
\pgfpathlineto{\pgfqpoint{1.021812in}{0.925370in}}%
\pgfpathlineto{\pgfqpoint{1.024331in}{0.906517in}}%
\pgfpathlineto{\pgfqpoint{1.026849in}{0.901545in}}%
\pgfpathlineto{\pgfqpoint{1.029368in}{0.908577in}}%
\pgfpathlineto{\pgfqpoint{1.031887in}{0.917706in}}%
\pgfpathlineto{\pgfqpoint{1.034406in}{0.931896in}}%
\pgfpathlineto{\pgfqpoint{1.036924in}{0.932937in}}%
\pgfpathlineto{\pgfqpoint{1.039443in}{0.937912in}}%
\pgfpathlineto{\pgfqpoint{1.041962in}{0.923594in}}%
\pgfpathlineto{\pgfqpoint{1.044481in}{0.948552in}}%
\pgfpathlineto{\pgfqpoint{1.046999in}{0.947239in}}%
\pgfpathlineto{\pgfqpoint{1.049518in}{0.944446in}}%
\pgfpathlineto{\pgfqpoint{1.052037in}{0.950116in}}%
\pgfpathlineto{\pgfqpoint{1.054556in}{0.950767in}}%
\pgfpathlineto{\pgfqpoint{1.057074in}{0.943136in}}%
\pgfpathlineto{\pgfqpoint{1.059593in}{0.961978in}}%
\pgfpathlineto{\pgfqpoint{1.062112in}{0.975734in}}%
\pgfpathlineto{\pgfqpoint{1.064631in}{0.992324in}}%
\pgfpathlineto{\pgfqpoint{1.067149in}{1.028844in}}%
\pgfpathlineto{\pgfqpoint{1.069668in}{1.020793in}}%
\pgfpathlineto{\pgfqpoint{1.072187in}{1.038227in}}%
\pgfpathlineto{\pgfqpoint{1.074706in}{1.044301in}}%
\pgfpathlineto{\pgfqpoint{1.077224in}{1.066091in}}%
\pgfpathlineto{\pgfqpoint{1.079743in}{1.075565in}}%
\pgfpathlineto{\pgfqpoint{1.082262in}{1.071606in}}%
\pgfpathlineto{\pgfqpoint{1.084781in}{1.084664in}}%
\pgfpathlineto{\pgfqpoint{1.087299in}{1.095229in}}%
\pgfpathlineto{\pgfqpoint{1.089818in}{1.105258in}}%
\pgfpathlineto{\pgfqpoint{1.092337in}{1.112383in}}%
\pgfpathlineto{\pgfqpoint{1.094856in}{1.088509in}}%
\pgfpathlineto{\pgfqpoint{1.097374in}{1.093695in}}%
\pgfpathlineto{\pgfqpoint{1.099893in}{1.101142in}}%
\pgfpathlineto{\pgfqpoint{1.102412in}{1.109209in}}%
\pgfpathlineto{\pgfqpoint{1.104931in}{1.115443in}}%
\pgfpathlineto{\pgfqpoint{1.107449in}{1.128530in}}%
\pgfpathlineto{\pgfqpoint{1.109968in}{1.129374in}}%
\pgfpathlineto{\pgfqpoint{1.112487in}{1.141233in}}%
\pgfpathlineto{\pgfqpoint{1.115006in}{1.134729in}}%
\pgfpathlineto{\pgfqpoint{1.117524in}{1.124231in}}%
\pgfpathlineto{\pgfqpoint{1.120043in}{1.105358in}}%
\pgfpathlineto{\pgfqpoint{1.122562in}{1.114623in}}%
\pgfpathlineto{\pgfqpoint{1.125081in}{1.113641in}}%
\pgfpathlineto{\pgfqpoint{1.127599in}{1.134371in}}%
\pgfpathlineto{\pgfqpoint{1.130118in}{1.143565in}}%
\pgfpathlineto{\pgfqpoint{1.132637in}{1.148774in}}%
\pgfpathlineto{\pgfqpoint{1.135156in}{1.141766in}}%
\pgfpathlineto{\pgfqpoint{1.137674in}{1.137506in}}%
\pgfpathlineto{\pgfqpoint{1.140193in}{1.142727in}}%
\pgfpathlineto{\pgfqpoint{1.142712in}{1.142075in}}%
\pgfpathlineto{\pgfqpoint{1.145231in}{1.127209in}}%
\pgfpathlineto{\pgfqpoint{1.147749in}{1.122085in}}%
\pgfpathlineto{\pgfqpoint{1.150268in}{1.109414in}}%
\pgfpathlineto{\pgfqpoint{1.152787in}{1.112192in}}%
\pgfpathlineto{\pgfqpoint{1.155306in}{1.088769in}}%
\pgfpathlineto{\pgfqpoint{1.157824in}{1.072666in}}%
\pgfpathlineto{\pgfqpoint{1.160343in}{1.069559in}}%
\pgfpathlineto{\pgfqpoint{1.162862in}{1.054857in}}%
\pgfpathlineto{\pgfqpoint{1.167899in}{1.031791in}}%
\pgfpathlineto{\pgfqpoint{1.170418in}{1.037638in}}%
\pgfpathlineto{\pgfqpoint{1.172937in}{1.030680in}}%
\pgfpathlineto{\pgfqpoint{1.175456in}{1.031743in}}%
\pgfpathlineto{\pgfqpoint{1.177974in}{1.030526in}}%
\pgfpathlineto{\pgfqpoint{1.180493in}{1.016921in}}%
\pgfpathlineto{\pgfqpoint{1.183012in}{1.030392in}}%
\pgfpathlineto{\pgfqpoint{1.185531in}{1.024037in}}%
\pgfpathlineto{\pgfqpoint{1.188049in}{1.010384in}}%
\pgfpathlineto{\pgfqpoint{1.190568in}{1.009899in}}%
\pgfpathlineto{\pgfqpoint{1.193087in}{1.025147in}}%
\pgfpathlineto{\pgfqpoint{1.195606in}{1.022512in}}%
\pgfpathlineto{\pgfqpoint{1.198124in}{1.015958in}}%
\pgfpathlineto{\pgfqpoint{1.200643in}{0.986871in}}%
\pgfpathlineto{\pgfqpoint{1.203162in}{0.988744in}}%
\pgfpathlineto{\pgfqpoint{1.205681in}{0.971316in}}%
\pgfpathlineto{\pgfqpoint{1.208199in}{0.975878in}}%
\pgfpathlineto{\pgfqpoint{1.210718in}{0.987569in}}%
\pgfpathlineto{\pgfqpoint{1.213237in}{0.979987in}}%
\pgfpathlineto{\pgfqpoint{1.215756in}{0.994486in}}%
\pgfpathlineto{\pgfqpoint{1.218274in}{0.983690in}}%
\pgfpathlineto{\pgfqpoint{1.220793in}{1.025856in}}%
\pgfpathlineto{\pgfqpoint{1.223312in}{1.036132in}}%
\pgfpathlineto{\pgfqpoint{1.225831in}{1.038855in}}%
\pgfpathlineto{\pgfqpoint{1.228349in}{1.043982in}}%
\pgfpathlineto{\pgfqpoint{1.230868in}{1.051782in}}%
\pgfpathlineto{\pgfqpoint{1.233387in}{1.056213in}}%
\pgfpathlineto{\pgfqpoint{1.235906in}{1.074916in}}%
\pgfpathlineto{\pgfqpoint{1.238424in}{1.076531in}}%
\pgfpathlineto{\pgfqpoint{1.240943in}{1.078754in}}%
\pgfpathlineto{\pgfqpoint{1.243462in}{1.087542in}}%
\pgfpathlineto{\pgfqpoint{1.245981in}{1.082757in}}%
\pgfpathlineto{\pgfqpoint{1.248499in}{1.060667in}}%
\pgfpathlineto{\pgfqpoint{1.251018in}{1.075127in}}%
\pgfpathlineto{\pgfqpoint{1.253537in}{1.069936in}}%
\pgfpathlineto{\pgfqpoint{1.256056in}{1.055699in}}%
\pgfpathlineto{\pgfqpoint{1.258574in}{1.039927in}}%
\pgfpathlineto{\pgfqpoint{1.261093in}{1.032329in}}%
\pgfpathlineto{\pgfqpoint{1.263612in}{1.030453in}}%
\pgfpathlineto{\pgfqpoint{1.266131in}{1.037965in}}%
\pgfpathlineto{\pgfqpoint{1.268649in}{1.025095in}}%
\pgfpathlineto{\pgfqpoint{1.271168in}{1.030935in}}%
\pgfpathlineto{\pgfqpoint{1.273687in}{1.024844in}}%
\pgfpathlineto{\pgfqpoint{1.276206in}{1.037030in}}%
\pgfpathlineto{\pgfqpoint{1.278724in}{1.028747in}}%
\pgfpathlineto{\pgfqpoint{1.281243in}{1.019362in}}%
\pgfpathlineto{\pgfqpoint{1.283762in}{0.997464in}}%
\pgfpathlineto{\pgfqpoint{1.286281in}{0.996026in}}%
\pgfpathlineto{\pgfqpoint{1.288799in}{1.014774in}}%
\pgfpathlineto{\pgfqpoint{1.291318in}{1.005480in}}%
\pgfpathlineto{\pgfqpoint{1.293837in}{0.990607in}}%
\pgfpathlineto{\pgfqpoint{1.296356in}{1.000254in}}%
\pgfpathlineto{\pgfqpoint{1.298874in}{0.996559in}}%
\pgfpathlineto{\pgfqpoint{1.301393in}{1.012166in}}%
\pgfpathlineto{\pgfqpoint{1.303912in}{1.023950in}}%
\pgfpathlineto{\pgfqpoint{1.306431in}{1.031693in}}%
\pgfpathlineto{\pgfqpoint{1.308949in}{1.040770in}}%
\pgfpathlineto{\pgfqpoint{1.311468in}{1.039569in}}%
\pgfpathlineto{\pgfqpoint{1.313987in}{1.056984in}}%
\pgfpathlineto{\pgfqpoint{1.316506in}{1.076606in}}%
\pgfpathlineto{\pgfqpoint{1.319024in}{1.071791in}}%
\pgfpathlineto{\pgfqpoint{1.321543in}{1.048930in}}%
\pgfpathlineto{\pgfqpoint{1.324062in}{1.049334in}}%
\pgfpathlineto{\pgfqpoint{1.326581in}{1.031803in}}%
\pgfpathlineto{\pgfqpoint{1.329099in}{1.034103in}}%
\pgfpathlineto{\pgfqpoint{1.331618in}{1.008945in}}%
\pgfpathlineto{\pgfqpoint{1.334137in}{1.016526in}}%
\pgfpathlineto{\pgfqpoint{1.336656in}{1.022361in}}%
\pgfpathlineto{\pgfqpoint{1.339174in}{1.034718in}}%
\pgfpathlineto{\pgfqpoint{1.341693in}{1.035570in}}%
\pgfpathlineto{\pgfqpoint{1.344212in}{1.020995in}}%
\pgfpathlineto{\pgfqpoint{1.346731in}{1.012192in}}%
\pgfpathlineto{\pgfqpoint{1.349249in}{1.029810in}}%
\pgfpathlineto{\pgfqpoint{1.351768in}{1.029547in}}%
\pgfpathlineto{\pgfqpoint{1.354287in}{1.033444in}}%
\pgfpathlineto{\pgfqpoint{1.356806in}{1.036805in}}%
\pgfpathlineto{\pgfqpoint{1.359324in}{1.035115in}}%
\pgfpathlineto{\pgfqpoint{1.361843in}{1.040164in}}%
\pgfpathlineto{\pgfqpoint{1.364362in}{1.039678in}}%
\pgfpathlineto{\pgfqpoint{1.366881in}{1.050025in}}%
\pgfpathlineto{\pgfqpoint{1.369399in}{1.048929in}}%
\pgfpathlineto{\pgfqpoint{1.371918in}{1.050224in}}%
\pgfpathlineto{\pgfqpoint{1.374437in}{1.074519in}}%
\pgfpathlineto{\pgfqpoint{1.376956in}{1.070489in}}%
\pgfpathlineto{\pgfqpoint{1.379474in}{1.092001in}}%
\pgfpathlineto{\pgfqpoint{1.381993in}{1.092163in}}%
\pgfpathlineto{\pgfqpoint{1.384512in}{1.098626in}}%
\pgfpathlineto{\pgfqpoint{1.387031in}{1.112187in}}%
\pgfpathlineto{\pgfqpoint{1.389549in}{1.132655in}}%
\pgfpathlineto{\pgfqpoint{1.392068in}{1.134192in}}%
\pgfpathlineto{\pgfqpoint{1.394587in}{1.125215in}}%
\pgfpathlineto{\pgfqpoint{1.397106in}{1.136024in}}%
\pgfpathlineto{\pgfqpoint{1.399624in}{1.136095in}}%
\pgfpathlineto{\pgfqpoint{1.402143in}{1.108800in}}%
\pgfpathlineto{\pgfqpoint{1.404662in}{1.101300in}}%
\pgfpathlineto{\pgfqpoint{1.407181in}{1.092247in}}%
\pgfpathlineto{\pgfqpoint{1.409699in}{1.098213in}}%
\pgfpathlineto{\pgfqpoint{1.412218in}{1.082913in}}%
\pgfpathlineto{\pgfqpoint{1.414737in}{1.094887in}}%
\pgfpathlineto{\pgfqpoint{1.417256in}{1.111731in}}%
\pgfpathlineto{\pgfqpoint{1.419774in}{1.142324in}}%
\pgfpathlineto{\pgfqpoint{1.424812in}{1.132842in}}%
\pgfpathlineto{\pgfqpoint{1.429849in}{1.106035in}}%
\pgfpathlineto{\pgfqpoint{1.432368in}{1.132456in}}%
\pgfpathlineto{\pgfqpoint{1.434887in}{1.146534in}}%
\pgfpathlineto{\pgfqpoint{1.437406in}{1.158585in}}%
\pgfpathlineto{\pgfqpoint{1.439924in}{1.147901in}}%
\pgfpathlineto{\pgfqpoint{1.442443in}{1.138491in}}%
\pgfpathlineto{\pgfqpoint{1.444962in}{1.121387in}}%
\pgfpathlineto{\pgfqpoint{1.447481in}{1.119445in}}%
\pgfpathlineto{\pgfqpoint{1.449999in}{1.141367in}}%
\pgfpathlineto{\pgfqpoint{1.452518in}{1.136198in}}%
\pgfpathlineto{\pgfqpoint{1.457556in}{1.125008in}}%
\pgfpathlineto{\pgfqpoint{1.460074in}{1.133494in}}%
\pgfpathlineto{\pgfqpoint{1.462593in}{1.138543in}}%
\pgfpathlineto{\pgfqpoint{1.465112in}{1.145497in}}%
\pgfpathlineto{\pgfqpoint{1.467631in}{1.162822in}}%
\pgfpathlineto{\pgfqpoint{1.470149in}{1.146858in}}%
\pgfpathlineto{\pgfqpoint{1.472668in}{1.151144in}}%
\pgfpathlineto{\pgfqpoint{1.475187in}{1.145347in}}%
\pgfpathlineto{\pgfqpoint{1.477706in}{1.185996in}}%
\pgfpathlineto{\pgfqpoint{1.480224in}{1.195964in}}%
\pgfpathlineto{\pgfqpoint{1.482743in}{1.183999in}}%
\pgfpathlineto{\pgfqpoint{1.485262in}{1.197351in}}%
\pgfpathlineto{\pgfqpoint{1.487781in}{1.197537in}}%
\pgfpathlineto{\pgfqpoint{1.490299in}{1.208224in}}%
\pgfpathlineto{\pgfqpoint{1.492818in}{1.191647in}}%
\pgfpathlineto{\pgfqpoint{1.495337in}{1.152683in}}%
\pgfpathlineto{\pgfqpoint{1.497856in}{1.189116in}}%
\pgfpathlineto{\pgfqpoint{1.500374in}{1.215743in}}%
\pgfpathlineto{\pgfqpoint{1.502893in}{1.226954in}}%
\pgfpathlineto{\pgfqpoint{1.505412in}{1.203786in}}%
\pgfpathlineto{\pgfqpoint{1.507931in}{1.205826in}}%
\pgfpathlineto{\pgfqpoint{1.510449in}{1.163602in}}%
\pgfpathlineto{\pgfqpoint{1.512968in}{1.141850in}}%
\pgfpathlineto{\pgfqpoint{1.515487in}{1.146764in}}%
\pgfpathlineto{\pgfqpoint{1.518006in}{1.136706in}}%
\pgfpathlineto{\pgfqpoint{1.520524in}{1.135564in}}%
\pgfpathlineto{\pgfqpoint{1.523043in}{1.137873in}}%
\pgfpathlineto{\pgfqpoint{1.525562in}{1.138349in}}%
\pgfpathlineto{\pgfqpoint{1.528081in}{1.129199in}}%
\pgfpathlineto{\pgfqpoint{1.530599in}{1.123608in}}%
\pgfpathlineto{\pgfqpoint{1.533118in}{1.154854in}}%
\pgfpathlineto{\pgfqpoint{1.535637in}{1.157307in}}%
\pgfpathlineto{\pgfqpoint{1.538156in}{1.160862in}}%
\pgfpathlineto{\pgfqpoint{1.540674in}{1.136183in}}%
\pgfpathlineto{\pgfqpoint{1.543193in}{1.119379in}}%
\pgfpathlineto{\pgfqpoint{1.545712in}{1.133578in}}%
\pgfpathlineto{\pgfqpoint{1.548231in}{1.146847in}}%
\pgfpathlineto{\pgfqpoint{1.550749in}{1.132726in}}%
\pgfpathlineto{\pgfqpoint{1.553268in}{1.136618in}}%
\pgfpathlineto{\pgfqpoint{1.555787in}{1.136757in}}%
\pgfpathlineto{\pgfqpoint{1.558306in}{1.118500in}}%
\pgfpathlineto{\pgfqpoint{1.560824in}{1.127960in}}%
\pgfpathlineto{\pgfqpoint{1.563343in}{1.145839in}}%
\pgfpathlineto{\pgfqpoint{1.565862in}{1.137930in}}%
\pgfpathlineto{\pgfqpoint{1.568381in}{1.110770in}}%
\pgfpathlineto{\pgfqpoint{1.570899in}{1.120368in}}%
\pgfpathlineto{\pgfqpoint{1.573418in}{1.139207in}}%
\pgfpathlineto{\pgfqpoint{1.575937in}{1.143063in}}%
\pgfpathlineto{\pgfqpoint{1.578456in}{1.147268in}}%
\pgfpathlineto{\pgfqpoint{1.580974in}{1.143278in}}%
\pgfpathlineto{\pgfqpoint{1.583493in}{1.132597in}}%
\pgfpathlineto{\pgfqpoint{1.586012in}{1.130288in}}%
\pgfpathlineto{\pgfqpoint{1.588531in}{1.127643in}}%
\pgfpathlineto{\pgfqpoint{1.591049in}{1.127440in}}%
\pgfpathlineto{\pgfqpoint{1.593568in}{1.118900in}}%
\pgfpathlineto{\pgfqpoint{1.596087in}{1.121108in}}%
\pgfpathlineto{\pgfqpoint{1.598606in}{1.095442in}}%
\pgfpathlineto{\pgfqpoint{1.601124in}{1.087098in}}%
\pgfpathlineto{\pgfqpoint{1.603643in}{1.084848in}}%
\pgfpathlineto{\pgfqpoint{1.606162in}{1.059193in}}%
\pgfpathlineto{\pgfqpoint{1.608681in}{1.054803in}}%
\pgfpathlineto{\pgfqpoint{1.611199in}{1.059297in}}%
\pgfpathlineto{\pgfqpoint{1.613718in}{1.084427in}}%
\pgfpathlineto{\pgfqpoint{1.616237in}{1.085866in}}%
\pgfpathlineto{\pgfqpoint{1.618756in}{1.099717in}}%
\pgfpathlineto{\pgfqpoint{1.621274in}{1.109110in}}%
\pgfpathlineto{\pgfqpoint{1.623793in}{1.120978in}}%
\pgfpathlineto{\pgfqpoint{1.626312in}{1.110936in}}%
\pgfpathlineto{\pgfqpoint{1.628831in}{1.112419in}}%
\pgfpathlineto{\pgfqpoint{1.631349in}{1.106140in}}%
\pgfpathlineto{\pgfqpoint{1.633868in}{1.114187in}}%
\pgfpathlineto{\pgfqpoint{1.636387in}{1.111723in}}%
\pgfpathlineto{\pgfqpoint{1.638906in}{1.097815in}}%
\pgfpathlineto{\pgfqpoint{1.641424in}{1.085129in}}%
\pgfpathlineto{\pgfqpoint{1.643943in}{1.102047in}}%
\pgfpathlineto{\pgfqpoint{1.646462in}{1.107428in}}%
\pgfpathlineto{\pgfqpoint{1.648981in}{1.085450in}}%
\pgfpathlineto{\pgfqpoint{1.651499in}{1.091340in}}%
\pgfpathlineto{\pgfqpoint{1.654018in}{1.109545in}}%
\pgfpathlineto{\pgfqpoint{1.656537in}{1.105255in}}%
\pgfpathlineto{\pgfqpoint{1.659056in}{1.124175in}}%
\pgfpathlineto{\pgfqpoint{1.661574in}{1.140223in}}%
\pgfpathlineto{\pgfqpoint{1.664093in}{1.147538in}}%
\pgfpathlineto{\pgfqpoint{1.666612in}{1.168414in}}%
\pgfpathlineto{\pgfqpoint{1.669131in}{1.166944in}}%
\pgfpathlineto{\pgfqpoint{1.671649in}{1.171507in}}%
\pgfpathlineto{\pgfqpoint{1.674168in}{1.148778in}}%
\pgfpathlineto{\pgfqpoint{1.676687in}{1.176370in}}%
\pgfpathlineto{\pgfqpoint{1.679206in}{1.198550in}}%
\pgfpathlineto{\pgfqpoint{1.681724in}{1.199241in}}%
\pgfpathlineto{\pgfqpoint{1.684243in}{1.198679in}}%
\pgfpathlineto{\pgfqpoint{1.686762in}{1.195175in}}%
\pgfpathlineto{\pgfqpoint{1.689281in}{1.192882in}}%
\pgfpathlineto{\pgfqpoint{1.691799in}{1.199961in}}%
\pgfpathlineto{\pgfqpoint{1.694318in}{1.222086in}}%
\pgfpathlineto{\pgfqpoint{1.696837in}{1.189388in}}%
\pgfpathlineto{\pgfqpoint{1.699356in}{1.178186in}}%
\pgfpathlineto{\pgfqpoint{1.701874in}{1.182116in}}%
\pgfpathlineto{\pgfqpoint{1.704393in}{1.194130in}}%
\pgfpathlineto{\pgfqpoint{1.706912in}{1.173184in}}%
\pgfpathlineto{\pgfqpoint{1.709431in}{1.191837in}}%
\pgfpathlineto{\pgfqpoint{1.711949in}{1.153608in}}%
\pgfpathlineto{\pgfqpoint{1.714468in}{1.129929in}}%
\pgfpathlineto{\pgfqpoint{1.716987in}{1.137547in}}%
\pgfpathlineto{\pgfqpoint{1.719506in}{1.136343in}}%
\pgfpathlineto{\pgfqpoint{1.722024in}{1.134503in}}%
\pgfpathlineto{\pgfqpoint{1.727062in}{1.118709in}}%
\pgfpathlineto{\pgfqpoint{1.729581in}{1.129521in}}%
\pgfpathlineto{\pgfqpoint{1.732099in}{1.126923in}}%
\pgfpathlineto{\pgfqpoint{1.734618in}{1.134269in}}%
\pgfpathlineto{\pgfqpoint{1.737137in}{1.130716in}}%
\pgfpathlineto{\pgfqpoint{1.739656in}{1.136403in}}%
\pgfpathlineto{\pgfqpoint{1.742174in}{1.129763in}}%
\pgfpathlineto{\pgfqpoint{1.744693in}{1.132026in}}%
\pgfpathlineto{\pgfqpoint{1.747212in}{1.140550in}}%
\pgfpathlineto{\pgfqpoint{1.749731in}{1.124603in}}%
\pgfpathlineto{\pgfqpoint{1.752249in}{1.143135in}}%
\pgfpathlineto{\pgfqpoint{1.754768in}{1.134865in}}%
\pgfpathlineto{\pgfqpoint{1.757287in}{1.125119in}}%
\pgfpathlineto{\pgfqpoint{1.759806in}{1.109040in}}%
\pgfpathlineto{\pgfqpoint{1.762324in}{1.104801in}}%
\pgfpathlineto{\pgfqpoint{1.764843in}{1.103453in}}%
\pgfpathlineto{\pgfqpoint{1.767362in}{1.122074in}}%
\pgfpathlineto{\pgfqpoint{1.769881in}{1.129050in}}%
\pgfpathlineto{\pgfqpoint{1.772399in}{1.111550in}}%
\pgfpathlineto{\pgfqpoint{1.774918in}{1.123466in}}%
\pgfpathlineto{\pgfqpoint{1.777437in}{1.109570in}}%
\pgfpathlineto{\pgfqpoint{1.779956in}{1.100449in}}%
\pgfpathlineto{\pgfqpoint{1.782474in}{1.113747in}}%
\pgfpathlineto{\pgfqpoint{1.784993in}{1.126441in}}%
\pgfpathlineto{\pgfqpoint{1.787512in}{1.105114in}}%
\pgfpathlineto{\pgfqpoint{1.790031in}{1.115027in}}%
\pgfpathlineto{\pgfqpoint{1.792549in}{1.092564in}}%
\pgfpathlineto{\pgfqpoint{1.795068in}{1.093669in}}%
\pgfpathlineto{\pgfqpoint{1.797587in}{1.091494in}}%
\pgfpathlineto{\pgfqpoint{1.800106in}{1.071802in}}%
\pgfpathlineto{\pgfqpoint{1.802624in}{1.087284in}}%
\pgfpathlineto{\pgfqpoint{1.805143in}{1.064382in}}%
\pgfpathlineto{\pgfqpoint{1.807662in}{1.057523in}}%
\pgfpathlineto{\pgfqpoint{1.810181in}{1.064509in}}%
\pgfpathlineto{\pgfqpoint{1.812699in}{1.062374in}}%
\pgfpathlineto{\pgfqpoint{1.815218in}{1.051121in}}%
\pgfpathlineto{\pgfqpoint{1.817737in}{1.060858in}}%
\pgfpathlineto{\pgfqpoint{1.820256in}{1.047872in}}%
\pgfpathlineto{\pgfqpoint{1.822774in}{1.029308in}}%
\pgfpathlineto{\pgfqpoint{1.825293in}{1.028191in}}%
\pgfpathlineto{\pgfqpoint{1.827812in}{1.001945in}}%
\pgfpathlineto{\pgfqpoint{1.830331in}{1.003361in}}%
\pgfpathlineto{\pgfqpoint{1.832849in}{0.989281in}}%
\pgfpathlineto{\pgfqpoint{1.835368in}{0.964567in}}%
\pgfpathlineto{\pgfqpoint{1.837887in}{0.963158in}}%
\pgfpathlineto{\pgfqpoint{1.840406in}{0.963395in}}%
\pgfpathlineto{\pgfqpoint{1.842924in}{0.944814in}}%
\pgfpathlineto{\pgfqpoint{1.845443in}{0.920691in}}%
\pgfpathlineto{\pgfqpoint{1.847962in}{0.926246in}}%
\pgfpathlineto{\pgfqpoint{1.850481in}{0.918050in}}%
\pgfpathlineto{\pgfqpoint{1.852999in}{0.916398in}}%
\pgfpathlineto{\pgfqpoint{1.855518in}{0.918100in}}%
\pgfpathlineto{\pgfqpoint{1.858037in}{0.914046in}}%
\pgfpathlineto{\pgfqpoint{1.860556in}{0.915011in}}%
\pgfpathlineto{\pgfqpoint{1.863074in}{0.916868in}}%
\pgfpathlineto{\pgfqpoint{1.868112in}{0.927696in}}%
\pgfpathlineto{\pgfqpoint{1.870631in}{0.922326in}}%
\pgfpathlineto{\pgfqpoint{1.873149in}{0.924934in}}%
\pgfpathlineto{\pgfqpoint{1.875668in}{0.947828in}}%
\pgfpathlineto{\pgfqpoint{1.878187in}{0.956087in}}%
\pgfpathlineto{\pgfqpoint{1.880706in}{0.950808in}}%
\pgfpathlineto{\pgfqpoint{1.883224in}{0.956364in}}%
\pgfpathlineto{\pgfqpoint{1.885743in}{0.985457in}}%
\pgfpathlineto{\pgfqpoint{1.888262in}{0.982626in}}%
\pgfpathlineto{\pgfqpoint{1.890781in}{0.976160in}}%
\pgfpathlineto{\pgfqpoint{1.893299in}{0.993898in}}%
\pgfpathlineto{\pgfqpoint{1.895818in}{1.015050in}}%
\pgfpathlineto{\pgfqpoint{1.898337in}{1.020893in}}%
\pgfpathlineto{\pgfqpoint{1.900856in}{1.033747in}}%
\pgfpathlineto{\pgfqpoint{1.903374in}{1.041782in}}%
\pgfpathlineto{\pgfqpoint{1.905893in}{1.039963in}}%
\pgfpathlineto{\pgfqpoint{1.908412in}{1.036609in}}%
\pgfpathlineto{\pgfqpoint{1.910931in}{1.064443in}}%
\pgfpathlineto{\pgfqpoint{1.913449in}{1.054036in}}%
\pgfpathlineto{\pgfqpoint{1.915968in}{1.062427in}}%
\pgfpathlineto{\pgfqpoint{1.918487in}{1.072637in}}%
\pgfpathlineto{\pgfqpoint{1.921006in}{1.076682in}}%
\pgfpathlineto{\pgfqpoint{1.923524in}{1.072669in}}%
\pgfpathlineto{\pgfqpoint{1.926043in}{1.071805in}}%
\pgfpathlineto{\pgfqpoint{1.928562in}{1.074914in}}%
\pgfpathlineto{\pgfqpoint{1.931081in}{1.071979in}}%
\pgfpathlineto{\pgfqpoint{1.933599in}{1.062358in}}%
\pgfpathlineto{\pgfqpoint{1.936118in}{1.082353in}}%
\pgfpathlineto{\pgfqpoint{1.938637in}{1.077666in}}%
\pgfpathlineto{\pgfqpoint{1.941156in}{1.088928in}}%
\pgfpathlineto{\pgfqpoint{1.943674in}{1.087355in}}%
\pgfpathlineto{\pgfqpoint{1.946193in}{1.082413in}}%
\pgfpathlineto{\pgfqpoint{1.948712in}{1.071214in}}%
\pgfpathlineto{\pgfqpoint{1.951231in}{1.079941in}}%
\pgfpathlineto{\pgfqpoint{1.953749in}{1.060226in}}%
\pgfpathlineto{\pgfqpoint{1.956268in}{1.072121in}}%
\pgfpathlineto{\pgfqpoint{1.958787in}{1.081961in}}%
\pgfpathlineto{\pgfqpoint{1.961306in}{1.057876in}}%
\pgfpathlineto{\pgfqpoint{1.963824in}{1.051535in}}%
\pgfpathlineto{\pgfqpoint{1.966343in}{1.046540in}}%
\pgfpathlineto{\pgfqpoint{1.968862in}{1.023255in}}%
\pgfpathlineto{\pgfqpoint{1.971381in}{1.048170in}}%
\pgfpathlineto{\pgfqpoint{1.973899in}{1.034017in}}%
\pgfpathlineto{\pgfqpoint{1.976418in}{1.016338in}}%
\pgfpathlineto{\pgfqpoint{1.978937in}{1.025134in}}%
\pgfpathlineto{\pgfqpoint{1.981456in}{1.047895in}}%
\pgfpathlineto{\pgfqpoint{1.983974in}{1.034279in}}%
\pgfpathlineto{\pgfqpoint{1.986493in}{1.055026in}}%
\pgfpathlineto{\pgfqpoint{1.989012in}{1.071657in}}%
\pgfpathlineto{\pgfqpoint{1.991531in}{1.073306in}}%
\pgfpathlineto{\pgfqpoint{1.994049in}{1.078320in}}%
\pgfpathlineto{\pgfqpoint{1.996568in}{1.097582in}}%
\pgfpathlineto{\pgfqpoint{1.999087in}{1.091154in}}%
\pgfpathlineto{\pgfqpoint{2.001606in}{1.081262in}}%
\pgfpathlineto{\pgfqpoint{2.004124in}{1.074582in}}%
\pgfpathlineto{\pgfqpoint{2.006643in}{1.073009in}}%
\pgfpathlineto{\pgfqpoint{2.009162in}{1.067277in}}%
\pgfpathlineto{\pgfqpoint{2.011681in}{1.057829in}}%
\pgfpathlineto{\pgfqpoint{2.014199in}{1.075398in}}%
\pgfpathlineto{\pgfqpoint{2.016718in}{1.065288in}}%
\pgfpathlineto{\pgfqpoint{2.019237in}{1.061481in}}%
\pgfpathlineto{\pgfqpoint{2.021756in}{1.042875in}}%
\pgfpathlineto{\pgfqpoint{2.024274in}{1.026342in}}%
\pgfpathlineto{\pgfqpoint{2.026793in}{1.055857in}}%
\pgfpathlineto{\pgfqpoint{2.029312in}{1.041349in}}%
\pgfpathlineto{\pgfqpoint{2.031831in}{1.029451in}}%
\pgfpathlineto{\pgfqpoint{2.034349in}{1.025299in}}%
\pgfpathlineto{\pgfqpoint{2.036868in}{1.019179in}}%
\pgfpathlineto{\pgfqpoint{2.039387in}{1.018718in}}%
\pgfpathlineto{\pgfqpoint{2.041906in}{1.019390in}}%
\pgfpathlineto{\pgfqpoint{2.044424in}{1.020558in}}%
\pgfpathlineto{\pgfqpoint{2.046943in}{1.007733in}}%
\pgfpathlineto{\pgfqpoint{2.049462in}{1.026775in}}%
\pgfpathlineto{\pgfqpoint{2.051981in}{1.013725in}}%
\pgfpathlineto{\pgfqpoint{2.054499in}{1.009412in}}%
\pgfpathlineto{\pgfqpoint{2.057018in}{0.987832in}}%
\pgfpathlineto{\pgfqpoint{2.059537in}{1.003853in}}%
\pgfpathlineto{\pgfqpoint{2.062056in}{1.021932in}}%
\pgfpathlineto{\pgfqpoint{2.064574in}{0.993455in}}%
\pgfpathlineto{\pgfqpoint{2.067093in}{0.992528in}}%
\pgfpathlineto{\pgfqpoint{2.069612in}{0.956611in}}%
\pgfpathlineto{\pgfqpoint{2.072131in}{0.951760in}}%
\pgfpathlineto{\pgfqpoint{2.074649in}{0.955166in}}%
\pgfpathlineto{\pgfqpoint{2.077168in}{0.963438in}}%
\pgfpathlineto{\pgfqpoint{2.079687in}{0.947391in}}%
\pgfpathlineto{\pgfqpoint{2.082206in}{0.943256in}}%
\pgfpathlineto{\pgfqpoint{2.084724in}{0.941142in}}%
\pgfpathlineto{\pgfqpoint{2.087243in}{0.961439in}}%
\pgfpathlineto{\pgfqpoint{2.089762in}{0.970283in}}%
\pgfpathlineto{\pgfqpoint{2.092281in}{0.977115in}}%
\pgfpathlineto{\pgfqpoint{2.094799in}{0.964098in}}%
\pgfpathlineto{\pgfqpoint{2.097318in}{0.970699in}}%
\pgfpathlineto{\pgfqpoint{2.099837in}{0.994724in}}%
\pgfpathlineto{\pgfqpoint{2.102356in}{0.988544in}}%
\pgfpathlineto{\pgfqpoint{2.104874in}{1.006572in}}%
\pgfpathlineto{\pgfqpoint{2.107393in}{0.992858in}}%
\pgfpathlineto{\pgfqpoint{2.109912in}{0.999988in}}%
\pgfpathlineto{\pgfqpoint{2.112431in}{1.015853in}}%
\pgfpathlineto{\pgfqpoint{2.114949in}{1.001151in}}%
\pgfpathlineto{\pgfqpoint{2.117468in}{0.978119in}}%
\pgfpathlineto{\pgfqpoint{2.119987in}{0.979773in}}%
\pgfpathlineto{\pgfqpoint{2.122506in}{0.966710in}}%
\pgfpathlineto{\pgfqpoint{2.125024in}{0.950767in}}%
\pgfpathlineto{\pgfqpoint{2.127543in}{0.953209in}}%
\pgfpathlineto{\pgfqpoint{2.130062in}{0.957697in}}%
\pgfpathlineto{\pgfqpoint{2.132581in}{0.921696in}}%
\pgfpathlineto{\pgfqpoint{2.135099in}{0.922819in}}%
\pgfpathlineto{\pgfqpoint{2.137618in}{0.922787in}}%
\pgfpathlineto{\pgfqpoint{2.140137in}{0.933336in}}%
\pgfpathlineto{\pgfqpoint{2.142656in}{0.949078in}}%
\pgfpathlineto{\pgfqpoint{2.145174in}{0.950858in}}%
\pgfpathlineto{\pgfqpoint{2.147693in}{0.942260in}}%
\pgfpathlineto{\pgfqpoint{2.150212in}{0.926601in}}%
\pgfpathlineto{\pgfqpoint{2.152731in}{0.918048in}}%
\pgfpathlineto{\pgfqpoint{2.155249in}{0.921933in}}%
\pgfpathlineto{\pgfqpoint{2.157768in}{0.911114in}}%
\pgfpathlineto{\pgfqpoint{2.160287in}{0.893928in}}%
\pgfpathlineto{\pgfqpoint{2.162806in}{0.883327in}}%
\pgfpathlineto{\pgfqpoint{2.165324in}{0.878288in}}%
\pgfpathlineto{\pgfqpoint{2.167843in}{0.891632in}}%
\pgfpathlineto{\pgfqpoint{2.170362in}{0.876140in}}%
\pgfpathlineto{\pgfqpoint{2.172881in}{0.853875in}}%
\pgfpathlineto{\pgfqpoint{2.175399in}{0.848838in}}%
\pgfpathlineto{\pgfqpoint{2.177918in}{0.854726in}}%
\pgfpathlineto{\pgfqpoint{2.180437in}{0.845714in}}%
\pgfpathlineto{\pgfqpoint{2.182956in}{0.828792in}}%
\pgfpathlineto{\pgfqpoint{2.185474in}{0.830556in}}%
\pgfpathlineto{\pgfqpoint{2.187993in}{0.838388in}}%
\pgfpathlineto{\pgfqpoint{2.190512in}{0.817382in}}%
\pgfpathlineto{\pgfqpoint{2.193031in}{0.798049in}}%
\pgfpathlineto{\pgfqpoint{2.195549in}{0.803008in}}%
\pgfpathlineto{\pgfqpoint{2.198068in}{0.794825in}}%
\pgfpathlineto{\pgfqpoint{2.200587in}{0.779387in}}%
\pgfpathlineto{\pgfqpoint{2.203106in}{0.783123in}}%
\pgfpathlineto{\pgfqpoint{2.205624in}{0.791920in}}%
\pgfpathlineto{\pgfqpoint{2.208143in}{0.787001in}}%
\pgfpathlineto{\pgfqpoint{2.210662in}{0.794777in}}%
\pgfpathlineto{\pgfqpoint{2.213181in}{0.786754in}}%
\pgfpathlineto{\pgfqpoint{2.215699in}{0.796802in}}%
\pgfpathlineto{\pgfqpoint{2.218218in}{0.813960in}}%
\pgfpathlineto{\pgfqpoint{2.220737in}{0.819557in}}%
\pgfpathlineto{\pgfqpoint{2.223256in}{0.808964in}}%
\pgfpathlineto{\pgfqpoint{2.225774in}{0.820397in}}%
\pgfpathlineto{\pgfqpoint{2.228293in}{0.825333in}}%
\pgfpathlineto{\pgfqpoint{2.230812in}{0.839414in}}%
\pgfpathlineto{\pgfqpoint{2.233331in}{0.823363in}}%
\pgfpathlineto{\pgfqpoint{2.235849in}{0.813928in}}%
\pgfpathlineto{\pgfqpoint{2.238368in}{0.808415in}}%
\pgfpathlineto{\pgfqpoint{2.240887in}{0.803802in}}%
\pgfpathlineto{\pgfqpoint{2.243406in}{0.801492in}}%
\pgfpathlineto{\pgfqpoint{2.245924in}{0.799431in}}%
\pgfpathlineto{\pgfqpoint{2.248443in}{0.814237in}}%
\pgfpathlineto{\pgfqpoint{2.250962in}{0.817772in}}%
\pgfpathlineto{\pgfqpoint{2.253481in}{0.802485in}}%
\pgfpathlineto{\pgfqpoint{2.255999in}{0.792015in}}%
\pgfpathlineto{\pgfqpoint{2.258518in}{0.786142in}}%
\pgfpathlineto{\pgfqpoint{2.261037in}{0.791036in}}%
\pgfpathlineto{\pgfqpoint{2.263556in}{0.792474in}}%
\pgfpathlineto{\pgfqpoint{2.266074in}{0.787676in}}%
\pgfpathlineto{\pgfqpoint{2.268593in}{0.774930in}}%
\pgfpathlineto{\pgfqpoint{2.271112in}{0.800569in}}%
\pgfpathlineto{\pgfqpoint{2.273631in}{0.805738in}}%
\pgfpathlineto{\pgfqpoint{2.276149in}{0.805633in}}%
\pgfpathlineto{\pgfqpoint{2.278668in}{0.833829in}}%
\pgfpathlineto{\pgfqpoint{2.281187in}{0.837833in}}%
\pgfpathlineto{\pgfqpoint{2.283706in}{0.830890in}}%
\pgfpathlineto{\pgfqpoint{2.286224in}{0.843790in}}%
\pgfpathlineto{\pgfqpoint{2.288743in}{0.848910in}}%
\pgfpathlineto{\pgfqpoint{2.291262in}{0.833438in}}%
\pgfpathlineto{\pgfqpoint{2.293781in}{0.850413in}}%
\pgfpathlineto{\pgfqpoint{2.296299in}{0.850224in}}%
\pgfpathlineto{\pgfqpoint{2.298818in}{0.849393in}}%
\pgfpathlineto{\pgfqpoint{2.301337in}{0.861396in}}%
\pgfpathlineto{\pgfqpoint{2.303856in}{0.841903in}}%
\pgfpathlineto{\pgfqpoint{2.306374in}{0.833141in}}%
\pgfpathlineto{\pgfqpoint{2.308893in}{0.837663in}}%
\pgfpathlineto{\pgfqpoint{2.311412in}{0.854287in}}%
\pgfpathlineto{\pgfqpoint{2.313931in}{0.868524in}}%
\pgfpathlineto{\pgfqpoint{2.316449in}{0.875608in}}%
\pgfpathlineto{\pgfqpoint{2.318968in}{0.853941in}}%
\pgfpathlineto{\pgfqpoint{2.321487in}{0.865949in}}%
\pgfpathlineto{\pgfqpoint{2.324006in}{0.862556in}}%
\pgfpathlineto{\pgfqpoint{2.326524in}{0.866185in}}%
\pgfpathlineto{\pgfqpoint{2.329043in}{0.870130in}}%
\pgfpathlineto{\pgfqpoint{2.331562in}{0.860013in}}%
\pgfpathlineto{\pgfqpoint{2.334081in}{0.867593in}}%
\pgfpathlineto{\pgfqpoint{2.336599in}{0.876040in}}%
\pgfpathlineto{\pgfqpoint{2.339118in}{0.868016in}}%
\pgfpathlineto{\pgfqpoint{2.341637in}{0.875347in}}%
\pgfpathlineto{\pgfqpoint{2.344156in}{0.859936in}}%
\pgfpathlineto{\pgfqpoint{2.346674in}{0.885046in}}%
\pgfpathlineto{\pgfqpoint{2.349193in}{0.893538in}}%
\pgfpathlineto{\pgfqpoint{2.351712in}{0.915714in}}%
\pgfpathlineto{\pgfqpoint{2.354231in}{0.921710in}}%
\pgfpathlineto{\pgfqpoint{2.359268in}{0.878114in}}%
\pgfpathlineto{\pgfqpoint{2.361787in}{0.867446in}}%
\pgfpathlineto{\pgfqpoint{2.364306in}{0.882064in}}%
\pgfpathlineto{\pgfqpoint{2.366824in}{0.894444in}}%
\pgfpathlineto{\pgfqpoint{2.369343in}{0.897394in}}%
\pgfpathlineto{\pgfqpoint{2.371862in}{0.877133in}}%
\pgfpathlineto{\pgfqpoint{2.374381in}{0.869870in}}%
\pgfpathlineto{\pgfqpoint{2.376899in}{0.867523in}}%
\pgfpathlineto{\pgfqpoint{2.379418in}{0.869458in}}%
\pgfpathlineto{\pgfqpoint{2.381937in}{0.874146in}}%
\pgfpathlineto{\pgfqpoint{2.384456in}{0.877821in}}%
\pgfpathlineto{\pgfqpoint{2.386974in}{0.886038in}}%
\pgfpathlineto{\pgfqpoint{2.389493in}{0.918817in}}%
\pgfpathlineto{\pgfqpoint{2.392012in}{0.909023in}}%
\pgfpathlineto{\pgfqpoint{2.394531in}{0.904477in}}%
\pgfpathlineto{\pgfqpoint{2.397049in}{0.924583in}}%
\pgfpathlineto{\pgfqpoint{2.399568in}{0.922441in}}%
\pgfpathlineto{\pgfqpoint{2.402087in}{0.896391in}}%
\pgfpathlineto{\pgfqpoint{2.404606in}{0.912945in}}%
\pgfpathlineto{\pgfqpoint{2.407124in}{0.908974in}}%
\pgfpathlineto{\pgfqpoint{2.409643in}{0.897081in}}%
\pgfpathlineto{\pgfqpoint{2.412162in}{0.876310in}}%
\pgfpathlineto{\pgfqpoint{2.414681in}{0.891314in}}%
\pgfpathlineto{\pgfqpoint{2.417199in}{0.868008in}}%
\pgfpathlineto{\pgfqpoint{2.419718in}{0.872250in}}%
\pgfpathlineto{\pgfqpoint{2.422237in}{0.859438in}}%
\pgfpathlineto{\pgfqpoint{2.424756in}{0.862666in}}%
\pgfpathlineto{\pgfqpoint{2.427274in}{0.869741in}}%
\pgfpathlineto{\pgfqpoint{2.429793in}{0.866781in}}%
\pgfpathlineto{\pgfqpoint{2.432312in}{0.873263in}}%
\pgfpathlineto{\pgfqpoint{2.434831in}{0.869406in}}%
\pgfpathlineto{\pgfqpoint{2.437349in}{0.869538in}}%
\pgfpathlineto{\pgfqpoint{2.439868in}{0.873633in}}%
\pgfpathlineto{\pgfqpoint{2.442387in}{0.893737in}}%
\pgfpathlineto{\pgfqpoint{2.444906in}{0.897038in}}%
\pgfpathlineto{\pgfqpoint{2.447424in}{0.908371in}}%
\pgfpathlineto{\pgfqpoint{2.449943in}{0.899558in}}%
\pgfpathlineto{\pgfqpoint{2.452462in}{0.902953in}}%
\pgfpathlineto{\pgfqpoint{2.454981in}{0.934372in}}%
\pgfpathlineto{\pgfqpoint{2.457499in}{0.935358in}}%
\pgfpathlineto{\pgfqpoint{2.460018in}{0.929955in}}%
\pgfpathlineto{\pgfqpoint{2.462537in}{0.927652in}}%
\pgfpathlineto{\pgfqpoint{2.465056in}{0.943716in}}%
\pgfpathlineto{\pgfqpoint{2.467574in}{0.939204in}}%
\pgfpathlineto{\pgfqpoint{2.470093in}{0.945598in}}%
\pgfpathlineto{\pgfqpoint{2.472612in}{0.939669in}}%
\pgfpathlineto{\pgfqpoint{2.475131in}{0.958877in}}%
\pgfpathlineto{\pgfqpoint{2.477649in}{0.930348in}}%
\pgfpathlineto{\pgfqpoint{2.480168in}{0.943770in}}%
\pgfpathlineto{\pgfqpoint{2.482687in}{0.937191in}}%
\pgfpathlineto{\pgfqpoint{2.485206in}{0.929602in}}%
\pgfpathlineto{\pgfqpoint{2.487724in}{0.947530in}}%
\pgfpathlineto{\pgfqpoint{2.490243in}{0.959864in}}%
\pgfpathlineto{\pgfqpoint{2.492762in}{0.953238in}}%
\pgfpathlineto{\pgfqpoint{2.495281in}{0.947068in}}%
\pgfpathlineto{\pgfqpoint{2.497799in}{0.932655in}}%
\pgfpathlineto{\pgfqpoint{2.500318in}{0.947821in}}%
\pgfpathlineto{\pgfqpoint{2.502837in}{0.958217in}}%
\pgfpathlineto{\pgfqpoint{2.505356in}{0.961953in}}%
\pgfpathlineto{\pgfqpoint{2.507874in}{0.956376in}}%
\pgfpathlineto{\pgfqpoint{2.510393in}{0.947747in}}%
\pgfpathlineto{\pgfqpoint{2.512912in}{0.948840in}}%
\pgfpathlineto{\pgfqpoint{2.515431in}{0.940840in}}%
\pgfpathlineto{\pgfqpoint{2.517949in}{0.923513in}}%
\pgfpathlineto{\pgfqpoint{2.520468in}{0.921169in}}%
\pgfpathlineto{\pgfqpoint{2.525506in}{0.904435in}}%
\pgfpathlineto{\pgfqpoint{2.528024in}{0.894174in}}%
\pgfpathlineto{\pgfqpoint{2.530543in}{0.881903in}}%
\pgfpathlineto{\pgfqpoint{2.533062in}{0.867337in}}%
\pgfpathlineto{\pgfqpoint{2.535581in}{0.874749in}}%
\pgfpathlineto{\pgfqpoint{2.538099in}{0.882906in}}%
\pgfpathlineto{\pgfqpoint{2.540618in}{0.887958in}}%
\pgfpathlineto{\pgfqpoint{2.543137in}{0.913161in}}%
\pgfpathlineto{\pgfqpoint{2.545656in}{0.934434in}}%
\pgfpathlineto{\pgfqpoint{2.548174in}{0.937991in}}%
\pgfpathlineto{\pgfqpoint{2.550693in}{0.945615in}}%
\pgfpathlineto{\pgfqpoint{2.553212in}{0.976568in}}%
\pgfpathlineto{\pgfqpoint{2.555731in}{0.973319in}}%
\pgfpathlineto{\pgfqpoint{2.558249in}{0.967287in}}%
\pgfpathlineto{\pgfqpoint{2.560768in}{0.976490in}}%
\pgfpathlineto{\pgfqpoint{2.563287in}{0.978605in}}%
\pgfpathlineto{\pgfqpoint{2.565806in}{0.993094in}}%
\pgfpathlineto{\pgfqpoint{2.568324in}{1.005586in}}%
\pgfpathlineto{\pgfqpoint{2.570843in}{0.991393in}}%
\pgfpathlineto{\pgfqpoint{2.573362in}{0.986062in}}%
\pgfpathlineto{\pgfqpoint{2.575881in}{1.022456in}}%
\pgfpathlineto{\pgfqpoint{2.578399in}{1.026542in}}%
\pgfpathlineto{\pgfqpoint{2.580918in}{1.008165in}}%
\pgfpathlineto{\pgfqpoint{2.583437in}{1.004277in}}%
\pgfpathlineto{\pgfqpoint{2.585956in}{0.995180in}}%
\pgfpathlineto{\pgfqpoint{2.588474in}{1.005118in}}%
\pgfpathlineto{\pgfqpoint{2.590993in}{0.995627in}}%
\pgfpathlineto{\pgfqpoint{2.593512in}{0.987767in}}%
\pgfpathlineto{\pgfqpoint{2.596031in}{0.997601in}}%
\pgfpathlineto{\pgfqpoint{2.598549in}{0.992011in}}%
\pgfpathlineto{\pgfqpoint{2.601068in}{0.986020in}}%
\pgfpathlineto{\pgfqpoint{2.603587in}{0.968665in}}%
\pgfpathlineto{\pgfqpoint{2.606106in}{0.983903in}}%
\pgfpathlineto{\pgfqpoint{2.608624in}{0.977357in}}%
\pgfpathlineto{\pgfqpoint{2.611143in}{0.990289in}}%
\pgfpathlineto{\pgfqpoint{2.613662in}{0.984748in}}%
\pgfpathlineto{\pgfqpoint{2.616181in}{0.975382in}}%
\pgfpathlineto{\pgfqpoint{2.618699in}{0.963575in}}%
\pgfpathlineto{\pgfqpoint{2.621218in}{0.943714in}}%
\pgfpathlineto{\pgfqpoint{2.623737in}{0.944050in}}%
\pgfpathlineto{\pgfqpoint{2.626256in}{0.931299in}}%
\pgfpathlineto{\pgfqpoint{2.628774in}{0.938037in}}%
\pgfpathlineto{\pgfqpoint{2.631293in}{0.940466in}}%
\pgfpathlineto{\pgfqpoint{2.633812in}{0.950550in}}%
\pgfpathlineto{\pgfqpoint{2.636331in}{0.955791in}}%
\pgfpathlineto{\pgfqpoint{2.638849in}{0.964684in}}%
\pgfpathlineto{\pgfqpoint{2.641368in}{0.963662in}}%
\pgfpathlineto{\pgfqpoint{2.643887in}{0.977190in}}%
\pgfpathlineto{\pgfqpoint{2.646406in}{0.978460in}}%
\pgfpathlineto{\pgfqpoint{2.648924in}{0.977612in}}%
\pgfpathlineto{\pgfqpoint{2.651443in}{0.985820in}}%
\pgfpathlineto{\pgfqpoint{2.653962in}{1.008993in}}%
\pgfpathlineto{\pgfqpoint{2.656481in}{0.993259in}}%
\pgfpathlineto{\pgfqpoint{2.658999in}{1.013160in}}%
\pgfpathlineto{\pgfqpoint{2.661518in}{1.030032in}}%
\pgfpathlineto{\pgfqpoint{2.664037in}{1.051831in}}%
\pgfpathlineto{\pgfqpoint{2.666556in}{1.038839in}}%
\pgfpathlineto{\pgfqpoint{2.669074in}{1.046528in}}%
\pgfpathlineto{\pgfqpoint{2.671593in}{1.041789in}}%
\pgfpathlineto{\pgfqpoint{2.674112in}{1.055226in}}%
\pgfpathlineto{\pgfqpoint{2.676631in}{1.036110in}}%
\pgfpathlineto{\pgfqpoint{2.679149in}{1.033668in}}%
\pgfpathlineto{\pgfqpoint{2.681668in}{1.043539in}}%
\pgfpathlineto{\pgfqpoint{2.684187in}{1.055854in}}%
\pgfpathlineto{\pgfqpoint{2.686706in}{1.059236in}}%
\pgfpathlineto{\pgfqpoint{2.689224in}{1.075307in}}%
\pgfpathlineto{\pgfqpoint{2.691743in}{1.059137in}}%
\pgfpathlineto{\pgfqpoint{2.694262in}{1.073932in}}%
\pgfpathlineto{\pgfqpoint{2.696781in}{1.093138in}}%
\pgfpathlineto{\pgfqpoint{2.699299in}{1.072585in}}%
\pgfpathlineto{\pgfqpoint{2.701818in}{1.080765in}}%
\pgfpathlineto{\pgfqpoint{2.704337in}{1.083464in}}%
\pgfpathlineto{\pgfqpoint{2.706856in}{1.096935in}}%
\pgfpathlineto{\pgfqpoint{2.709374in}{1.092909in}}%
\pgfpathlineto{\pgfqpoint{2.711893in}{1.098636in}}%
\pgfpathlineto{\pgfqpoint{2.714412in}{1.110275in}}%
\pgfpathlineto{\pgfqpoint{2.716931in}{1.133912in}}%
\pgfpathlineto{\pgfqpoint{2.719449in}{1.123288in}}%
\pgfpathlineto{\pgfqpoint{2.721968in}{1.114507in}}%
\pgfpathlineto{\pgfqpoint{2.724487in}{1.115001in}}%
\pgfpathlineto{\pgfqpoint{2.727006in}{1.103435in}}%
\pgfpathlineto{\pgfqpoint{2.729524in}{1.101784in}}%
\pgfpathlineto{\pgfqpoint{2.732043in}{1.124271in}}%
\pgfpathlineto{\pgfqpoint{2.734562in}{1.131537in}}%
\pgfpathlineto{\pgfqpoint{2.737081in}{1.145708in}}%
\pgfpathlineto{\pgfqpoint{2.739599in}{1.138657in}}%
\pgfpathlineto{\pgfqpoint{2.742118in}{1.148830in}}%
\pgfpathlineto{\pgfqpoint{2.744637in}{1.127861in}}%
\pgfpathlineto{\pgfqpoint{2.747156in}{1.157028in}}%
\pgfpathlineto{\pgfqpoint{2.749674in}{1.165230in}}%
\pgfpathlineto{\pgfqpoint{2.752193in}{1.174494in}}%
\pgfpathlineto{\pgfqpoint{2.754712in}{1.185646in}}%
\pgfpathlineto{\pgfqpoint{2.757231in}{1.167568in}}%
\pgfpathlineto{\pgfqpoint{2.759749in}{1.191229in}}%
\pgfpathlineto{\pgfqpoint{2.762268in}{1.206554in}}%
\pgfpathlineto{\pgfqpoint{2.764787in}{1.208490in}}%
\pgfpathlineto{\pgfqpoint{2.767306in}{1.217343in}}%
\pgfpathlineto{\pgfqpoint{2.769824in}{1.214180in}}%
\pgfpathlineto{\pgfqpoint{2.772343in}{1.221338in}}%
\pgfpathlineto{\pgfqpoint{2.774862in}{1.209756in}}%
\pgfpathlineto{\pgfqpoint{2.777381in}{1.215317in}}%
\pgfpathlineto{\pgfqpoint{2.779899in}{1.209670in}}%
\pgfpathlineto{\pgfqpoint{2.782418in}{1.203367in}}%
\pgfpathlineto{\pgfqpoint{2.784937in}{1.202934in}}%
\pgfpathlineto{\pgfqpoint{2.787456in}{1.206541in}}%
\pgfpathlineto{\pgfqpoint{2.789974in}{1.191370in}}%
\pgfpathlineto{\pgfqpoint{2.792493in}{1.182642in}}%
\pgfpathlineto{\pgfqpoint{2.795012in}{1.175171in}}%
\pgfpathlineto{\pgfqpoint{2.797531in}{1.161607in}}%
\pgfpathlineto{\pgfqpoint{2.800049in}{1.174791in}}%
\pgfpathlineto{\pgfqpoint{2.802568in}{1.161781in}}%
\pgfpathlineto{\pgfqpoint{2.805087in}{1.164041in}}%
\pgfpathlineto{\pgfqpoint{2.807606in}{1.164862in}}%
\pgfpathlineto{\pgfqpoint{2.810124in}{1.161327in}}%
\pgfpathlineto{\pgfqpoint{2.812643in}{1.163778in}}%
\pgfpathlineto{\pgfqpoint{2.815162in}{1.178839in}}%
\pgfpathlineto{\pgfqpoint{2.817681in}{1.160733in}}%
\pgfpathlineto{\pgfqpoint{2.820199in}{1.164808in}}%
\pgfpathlineto{\pgfqpoint{2.822718in}{1.172690in}}%
\pgfpathlineto{\pgfqpoint{2.825237in}{1.161622in}}%
\pgfpathlineto{\pgfqpoint{2.827756in}{1.139808in}}%
\pgfpathlineto{\pgfqpoint{2.830274in}{1.133722in}}%
\pgfpathlineto{\pgfqpoint{2.832793in}{1.155066in}}%
\pgfpathlineto{\pgfqpoint{2.835312in}{1.165560in}}%
\pgfpathlineto{\pgfqpoint{2.837831in}{1.143456in}}%
\pgfpathlineto{\pgfqpoint{2.840349in}{1.116179in}}%
\pgfpathlineto{\pgfqpoint{2.842868in}{1.131003in}}%
\pgfpathlineto{\pgfqpoint{2.845387in}{1.114159in}}%
\pgfpathlineto{\pgfqpoint{2.847906in}{1.127225in}}%
\pgfpathlineto{\pgfqpoint{2.850424in}{1.153740in}}%
\pgfpathlineto{\pgfqpoint{2.852943in}{1.124636in}}%
\pgfpathlineto{\pgfqpoint{2.855462in}{1.133138in}}%
\pgfpathlineto{\pgfqpoint{2.857981in}{1.135791in}}%
\pgfpathlineto{\pgfqpoint{2.860499in}{1.145192in}}%
\pgfpathlineto{\pgfqpoint{2.863018in}{1.141811in}}%
\pgfpathlineto{\pgfqpoint{2.865537in}{1.148986in}}%
\pgfpathlineto{\pgfqpoint{2.868056in}{1.124559in}}%
\pgfpathlineto{\pgfqpoint{2.870574in}{1.136023in}}%
\pgfpathlineto{\pgfqpoint{2.873093in}{1.157299in}}%
\pgfpathlineto{\pgfqpoint{2.875612in}{1.149573in}}%
\pgfpathlineto{\pgfqpoint{2.878131in}{1.184478in}}%
\pgfpathlineto{\pgfqpoint{2.880649in}{1.197069in}}%
\pgfpathlineto{\pgfqpoint{2.883168in}{1.211833in}}%
\pgfpathlineto{\pgfqpoint{2.885687in}{1.218448in}}%
\pgfpathlineto{\pgfqpoint{2.888206in}{1.207574in}}%
\pgfpathlineto{\pgfqpoint{2.890724in}{1.203065in}}%
\pgfpathlineto{\pgfqpoint{2.893243in}{1.218728in}}%
\pgfpathlineto{\pgfqpoint{2.895762in}{1.210229in}}%
\pgfpathlineto{\pgfqpoint{2.898281in}{1.191203in}}%
\pgfpathlineto{\pgfqpoint{2.900799in}{1.184572in}}%
\pgfpathlineto{\pgfqpoint{2.903318in}{1.161526in}}%
\pgfpathlineto{\pgfqpoint{2.905837in}{1.178488in}}%
\pgfpathlineto{\pgfqpoint{2.908356in}{1.173066in}}%
\pgfpathlineto{\pgfqpoint{2.910874in}{1.155659in}}%
\pgfpathlineto{\pgfqpoint{2.913393in}{1.171359in}}%
\pgfpathlineto{\pgfqpoint{2.915912in}{1.158196in}}%
\pgfpathlineto{\pgfqpoint{2.918431in}{1.149964in}}%
\pgfpathlineto{\pgfqpoint{2.920949in}{1.162689in}}%
\pgfpathlineto{\pgfqpoint{2.923468in}{1.167794in}}%
\pgfpathlineto{\pgfqpoint{2.925987in}{1.164894in}}%
\pgfpathlineto{\pgfqpoint{2.928506in}{1.198631in}}%
\pgfpathlineto{\pgfqpoint{2.931024in}{1.206445in}}%
\pgfpathlineto{\pgfqpoint{2.933543in}{1.209071in}}%
\pgfpathlineto{\pgfqpoint{2.936062in}{1.218248in}}%
\pgfpathlineto{\pgfqpoint{2.938581in}{1.224243in}}%
\pgfpathlineto{\pgfqpoint{2.941099in}{1.216927in}}%
\pgfpathlineto{\pgfqpoint{2.943618in}{1.216183in}}%
\pgfpathlineto{\pgfqpoint{2.946137in}{1.238269in}}%
\pgfpathlineto{\pgfqpoint{2.948656in}{1.253544in}}%
\pgfpathlineto{\pgfqpoint{2.951174in}{1.256755in}}%
\pgfpathlineto{\pgfqpoint{2.953693in}{1.271759in}}%
\pgfpathlineto{\pgfqpoint{2.956212in}{1.267723in}}%
\pgfpathlineto{\pgfqpoint{2.958731in}{1.253351in}}%
\pgfpathlineto{\pgfqpoint{2.961249in}{1.243156in}}%
\pgfpathlineto{\pgfqpoint{2.963768in}{1.272381in}}%
\pgfpathlineto{\pgfqpoint{2.966287in}{1.264347in}}%
\pgfpathlineto{\pgfqpoint{2.968806in}{1.263636in}}%
\pgfpathlineto{\pgfqpoint{2.971324in}{1.271117in}}%
\pgfpathlineto{\pgfqpoint{2.973843in}{1.259526in}}%
\pgfpathlineto{\pgfqpoint{2.976362in}{1.219252in}}%
\pgfpathlineto{\pgfqpoint{2.978881in}{1.230198in}}%
\pgfpathlineto{\pgfqpoint{2.981399in}{1.227615in}}%
\pgfpathlineto{\pgfqpoint{2.983918in}{1.206885in}}%
\pgfpathlineto{\pgfqpoint{2.986437in}{1.201650in}}%
\pgfpathlineto{\pgfqpoint{2.988956in}{1.182566in}}%
\pgfpathlineto{\pgfqpoint{2.991474in}{1.178861in}}%
\pgfpathlineto{\pgfqpoint{2.993993in}{1.198796in}}%
\pgfpathlineto{\pgfqpoint{2.996512in}{1.205672in}}%
\pgfpathlineto{\pgfqpoint{2.999031in}{1.204179in}}%
\pgfpathlineto{\pgfqpoint{3.001549in}{1.196150in}}%
\pgfpathlineto{\pgfqpoint{3.004068in}{1.196563in}}%
\pgfpathlineto{\pgfqpoint{3.006587in}{1.196281in}}%
\pgfpathlineto{\pgfqpoint{3.009106in}{1.207186in}}%
\pgfpathlineto{\pgfqpoint{3.011624in}{1.225425in}}%
\pgfpathlineto{\pgfqpoint{3.014143in}{1.218897in}}%
\pgfpathlineto{\pgfqpoint{3.016662in}{1.219628in}}%
\pgfpathlineto{\pgfqpoint{3.019181in}{1.239251in}}%
\pgfpathlineto{\pgfqpoint{3.021699in}{1.231040in}}%
\pgfpathlineto{\pgfqpoint{3.024218in}{1.225921in}}%
\pgfpathlineto{\pgfqpoint{3.026737in}{1.235110in}}%
\pgfpathlineto{\pgfqpoint{3.029256in}{1.262401in}}%
\pgfpathlineto{\pgfqpoint{3.031774in}{1.280057in}}%
\pgfpathlineto{\pgfqpoint{3.034293in}{1.281325in}}%
\pgfpathlineto{\pgfqpoint{3.036812in}{1.281160in}}%
\pgfpathlineto{\pgfqpoint{3.039331in}{1.296201in}}%
\pgfpathlineto{\pgfqpoint{3.041849in}{1.256139in}}%
\pgfpathlineto{\pgfqpoint{3.044368in}{1.258476in}}%
\pgfpathlineto{\pgfqpoint{3.046887in}{1.260272in}}%
\pgfpathlineto{\pgfqpoint{3.049406in}{1.231737in}}%
\pgfpathlineto{\pgfqpoint{3.051924in}{1.238678in}}%
\pgfpathlineto{\pgfqpoint{3.054443in}{1.240762in}}%
\pgfpathlineto{\pgfqpoint{3.056962in}{1.227820in}}%
\pgfpathlineto{\pgfqpoint{3.059481in}{1.242581in}}%
\pgfpathlineto{\pgfqpoint{3.061999in}{1.228828in}}%
\pgfpathlineto{\pgfqpoint{3.064518in}{1.218778in}}%
\pgfpathlineto{\pgfqpoint{3.067037in}{1.233000in}}%
\pgfpathlineto{\pgfqpoint{3.069556in}{1.246100in}}%
\pgfpathlineto{\pgfqpoint{3.072074in}{1.253577in}}%
\pgfpathlineto{\pgfqpoint{3.074593in}{1.226382in}}%
\pgfpathlineto{\pgfqpoint{3.077112in}{1.220581in}}%
\pgfpathlineto{\pgfqpoint{3.079631in}{1.216582in}}%
\pgfpathlineto{\pgfqpoint{3.082149in}{1.194519in}}%
\pgfpathlineto{\pgfqpoint{3.084668in}{1.185539in}}%
\pgfpathlineto{\pgfqpoint{3.087187in}{1.183679in}}%
\pgfpathlineto{\pgfqpoint{3.089706in}{1.182980in}}%
\pgfpathlineto{\pgfqpoint{3.092224in}{1.186185in}}%
\pgfpathlineto{\pgfqpoint{3.094743in}{1.194971in}}%
\pgfpathlineto{\pgfqpoint{3.097262in}{1.196722in}}%
\pgfpathlineto{\pgfqpoint{3.099781in}{1.213643in}}%
\pgfpathlineto{\pgfqpoint{3.102299in}{1.215554in}}%
\pgfpathlineto{\pgfqpoint{3.104818in}{1.235286in}}%
\pgfpathlineto{\pgfqpoint{3.107337in}{1.221432in}}%
\pgfpathlineto{\pgfqpoint{3.109856in}{1.219543in}}%
\pgfpathlineto{\pgfqpoint{3.112374in}{1.211205in}}%
\pgfpathlineto{\pgfqpoint{3.114893in}{1.215759in}}%
\pgfpathlineto{\pgfqpoint{3.117412in}{1.228448in}}%
\pgfpathlineto{\pgfqpoint{3.119931in}{1.244980in}}%
\pgfpathlineto{\pgfqpoint{3.122449in}{1.226385in}}%
\pgfpathlineto{\pgfqpoint{3.124968in}{1.219000in}}%
\pgfpathlineto{\pgfqpoint{3.127487in}{1.217334in}}%
\pgfpathlineto{\pgfqpoint{3.130006in}{1.214576in}}%
\pgfpathlineto{\pgfqpoint{3.132524in}{1.226221in}}%
\pgfpathlineto{\pgfqpoint{3.135043in}{1.243324in}}%
\pgfpathlineto{\pgfqpoint{3.137562in}{1.257216in}}%
\pgfpathlineto{\pgfqpoint{3.140081in}{1.260440in}}%
\pgfpathlineto{\pgfqpoint{3.142599in}{1.244518in}}%
\pgfpathlineto{\pgfqpoint{3.145118in}{1.220428in}}%
\pgfpathlineto{\pgfqpoint{3.147637in}{1.228683in}}%
\pgfpathlineto{\pgfqpoint{3.150156in}{1.208972in}}%
\pgfpathlineto{\pgfqpoint{3.152674in}{1.225332in}}%
\pgfpathlineto{\pgfqpoint{3.155193in}{1.246519in}}%
\pgfpathlineto{\pgfqpoint{3.157712in}{1.249934in}}%
\pgfpathlineto{\pgfqpoint{3.160231in}{1.292805in}}%
\pgfpathlineto{\pgfqpoint{3.162749in}{1.281924in}}%
\pgfpathlineto{\pgfqpoint{3.165268in}{1.267749in}}%
\pgfpathlineto{\pgfqpoint{3.167787in}{1.268358in}}%
\pgfpathlineto{\pgfqpoint{3.172824in}{1.287359in}}%
\pgfpathlineto{\pgfqpoint{3.175343in}{1.276163in}}%
\pgfpathlineto{\pgfqpoint{3.177862in}{1.267789in}}%
\pgfpathlineto{\pgfqpoint{3.180381in}{1.251716in}}%
\pgfpathlineto{\pgfqpoint{3.182899in}{1.261473in}}%
\pgfpathlineto{\pgfqpoint{3.185418in}{1.225932in}}%
\pgfpathlineto{\pgfqpoint{3.187937in}{1.224642in}}%
\pgfpathlineto{\pgfqpoint{3.190456in}{1.206821in}}%
\pgfpathlineto{\pgfqpoint{3.192974in}{1.225112in}}%
\pgfpathlineto{\pgfqpoint{3.195493in}{1.226490in}}%
\pgfpathlineto{\pgfqpoint{3.198012in}{1.220053in}}%
\pgfpathlineto{\pgfqpoint{3.200531in}{1.224465in}}%
\pgfpathlineto{\pgfqpoint{3.203049in}{1.199090in}}%
\pgfpathlineto{\pgfqpoint{3.205568in}{1.201007in}}%
\pgfpathlineto{\pgfqpoint{3.208087in}{1.210802in}}%
\pgfpathlineto{\pgfqpoint{3.210606in}{1.185561in}}%
\pgfpathlineto{\pgfqpoint{3.213124in}{1.186974in}}%
\pgfpathlineto{\pgfqpoint{3.215643in}{1.179052in}}%
\pgfpathlineto{\pgfqpoint{3.218162in}{1.194243in}}%
\pgfpathlineto{\pgfqpoint{3.220681in}{1.160118in}}%
\pgfpathlineto{\pgfqpoint{3.223199in}{1.144045in}}%
\pgfpathlineto{\pgfqpoint{3.225718in}{1.125433in}}%
\pgfpathlineto{\pgfqpoint{3.228237in}{1.131233in}}%
\pgfpathlineto{\pgfqpoint{3.230756in}{1.127429in}}%
\pgfpathlineto{\pgfqpoint{3.233274in}{1.116236in}}%
\pgfpathlineto{\pgfqpoint{3.235793in}{1.124385in}}%
\pgfpathlineto{\pgfqpoint{3.238312in}{1.160107in}}%
\pgfpathlineto{\pgfqpoint{3.240831in}{1.130615in}}%
\pgfpathlineto{\pgfqpoint{3.243349in}{1.128421in}}%
\pgfpathlineto{\pgfqpoint{3.245868in}{1.125299in}}%
\pgfpathlineto{\pgfqpoint{3.248387in}{1.111809in}}%
\pgfpathlineto{\pgfqpoint{3.250906in}{1.114489in}}%
\pgfpathlineto{\pgfqpoint{3.253424in}{1.133308in}}%
\pgfpathlineto{\pgfqpoint{3.255943in}{1.136633in}}%
\pgfpathlineto{\pgfqpoint{3.258462in}{1.148963in}}%
\pgfpathlineto{\pgfqpoint{3.260981in}{1.114197in}}%
\pgfpathlineto{\pgfqpoint{3.263499in}{1.112161in}}%
\pgfpathlineto{\pgfqpoint{3.266018in}{1.100933in}}%
\pgfpathlineto{\pgfqpoint{3.268537in}{1.086737in}}%
\pgfpathlineto{\pgfqpoint{3.271056in}{1.074294in}}%
\pgfpathlineto{\pgfqpoint{3.273574in}{1.072062in}}%
\pgfpathlineto{\pgfqpoint{3.276093in}{1.078465in}}%
\pgfpathlineto{\pgfqpoint{3.278612in}{1.084565in}}%
\pgfpathlineto{\pgfqpoint{3.281131in}{1.095936in}}%
\pgfpathlineto{\pgfqpoint{3.283649in}{1.095827in}}%
\pgfpathlineto{\pgfqpoint{3.286168in}{1.094841in}}%
\pgfpathlineto{\pgfqpoint{3.288687in}{1.091098in}}%
\pgfpathlineto{\pgfqpoint{3.291206in}{1.106958in}}%
\pgfpathlineto{\pgfqpoint{3.293724in}{1.094014in}}%
\pgfpathlineto{\pgfqpoint{3.296243in}{1.089488in}}%
\pgfpathlineto{\pgfqpoint{3.298762in}{1.090176in}}%
\pgfpathlineto{\pgfqpoint{3.301281in}{1.078629in}}%
\pgfpathlineto{\pgfqpoint{3.303799in}{1.074562in}}%
\pgfpathlineto{\pgfqpoint{3.306318in}{1.080190in}}%
\pgfpathlineto{\pgfqpoint{3.308837in}{1.088350in}}%
\pgfpathlineto{\pgfqpoint{3.311356in}{1.072502in}}%
\pgfpathlineto{\pgfqpoint{3.313874in}{1.073459in}}%
\pgfpathlineto{\pgfqpoint{3.316393in}{1.082564in}}%
\pgfpathlineto{\pgfqpoint{3.318912in}{1.065286in}}%
\pgfpathlineto{\pgfqpoint{3.321431in}{1.060534in}}%
\pgfpathlineto{\pgfqpoint{3.323949in}{1.088792in}}%
\pgfpathlineto{\pgfqpoint{3.326468in}{1.097775in}}%
\pgfpathlineto{\pgfqpoint{3.328987in}{1.128208in}}%
\pgfpathlineto{\pgfqpoint{3.331506in}{1.119982in}}%
\pgfpathlineto{\pgfqpoint{3.334024in}{1.138133in}}%
\pgfpathlineto{\pgfqpoint{3.336543in}{1.112885in}}%
\pgfpathlineto{\pgfqpoint{3.339062in}{1.125394in}}%
\pgfpathlineto{\pgfqpoint{3.341581in}{1.121820in}}%
\pgfpathlineto{\pgfqpoint{3.344099in}{1.108607in}}%
\pgfpathlineto{\pgfqpoint{3.346618in}{1.097886in}}%
\pgfpathlineto{\pgfqpoint{3.349137in}{1.088351in}}%
\pgfpathlineto{\pgfqpoint{3.351656in}{1.089047in}}%
\pgfpathlineto{\pgfqpoint{3.354174in}{1.097973in}}%
\pgfpathlineto{\pgfqpoint{3.356693in}{1.099864in}}%
\pgfpathlineto{\pgfqpoint{3.359212in}{1.099658in}}%
\pgfpathlineto{\pgfqpoint{3.361731in}{1.097020in}}%
\pgfpathlineto{\pgfqpoint{3.364249in}{1.099892in}}%
\pgfpathlineto{\pgfqpoint{3.366768in}{1.128313in}}%
\pgfpathlineto{\pgfqpoint{3.369287in}{1.110420in}}%
\pgfpathlineto{\pgfqpoint{3.371806in}{1.081787in}}%
\pgfpathlineto{\pgfqpoint{3.374324in}{1.102653in}}%
\pgfpathlineto{\pgfqpoint{3.376843in}{1.099573in}}%
\pgfpathlineto{\pgfqpoint{3.379362in}{1.099349in}}%
\pgfpathlineto{\pgfqpoint{3.381881in}{1.075407in}}%
\pgfpathlineto{\pgfqpoint{3.384399in}{1.084237in}}%
\pgfpathlineto{\pgfqpoint{3.386918in}{1.079268in}}%
\pgfpathlineto{\pgfqpoint{3.389437in}{1.080010in}}%
\pgfpathlineto{\pgfqpoint{3.391956in}{1.063732in}}%
\pgfpathlineto{\pgfqpoint{3.394474in}{1.046835in}}%
\pgfpathlineto{\pgfqpoint{3.396993in}{1.040762in}}%
\pgfpathlineto{\pgfqpoint{3.399512in}{1.041028in}}%
\pgfpathlineto{\pgfqpoint{3.402031in}{1.071591in}}%
\pgfpathlineto{\pgfqpoint{3.404549in}{1.072276in}}%
\pgfpathlineto{\pgfqpoint{3.407068in}{1.071715in}}%
\pgfpathlineto{\pgfqpoint{3.409587in}{1.070435in}}%
\pgfpathlineto{\pgfqpoint{3.412106in}{1.074179in}}%
\pgfpathlineto{\pgfqpoint{3.414624in}{1.068160in}}%
\pgfpathlineto{\pgfqpoint{3.417143in}{1.057123in}}%
\pgfpathlineto{\pgfqpoint{3.419662in}{1.058215in}}%
\pgfpathlineto{\pgfqpoint{3.422181in}{1.057365in}}%
\pgfpathlineto{\pgfqpoint{3.424699in}{1.074353in}}%
\pgfpathlineto{\pgfqpoint{3.427218in}{1.066839in}}%
\pgfpathlineto{\pgfqpoint{3.429737in}{1.052245in}}%
\pgfpathlineto{\pgfqpoint{3.432256in}{1.056586in}}%
\pgfpathlineto{\pgfqpoint{3.434774in}{1.053943in}}%
\pgfpathlineto{\pgfqpoint{3.437293in}{1.068638in}}%
\pgfpathlineto{\pgfqpoint{3.439812in}{1.063493in}}%
\pgfpathlineto{\pgfqpoint{3.442331in}{1.049207in}}%
\pgfpathlineto{\pgfqpoint{3.444849in}{1.055215in}}%
\pgfpathlineto{\pgfqpoint{3.447368in}{1.052111in}}%
\pgfpathlineto{\pgfqpoint{3.449887in}{1.055367in}}%
\pgfpathlineto{\pgfqpoint{3.452406in}{1.062839in}}%
\pgfpathlineto{\pgfqpoint{3.454924in}{1.073445in}}%
\pgfpathlineto{\pgfqpoint{3.457443in}{1.044886in}}%
\pgfpathlineto{\pgfqpoint{3.459962in}{1.050123in}}%
\pgfpathlineto{\pgfqpoint{3.462481in}{1.063266in}}%
\pgfpathlineto{\pgfqpoint{3.464999in}{1.058266in}}%
\pgfpathlineto{\pgfqpoint{3.467518in}{1.065667in}}%
\pgfpathlineto{\pgfqpoint{3.470037in}{1.057729in}}%
\pgfpathlineto{\pgfqpoint{3.472556in}{1.060532in}}%
\pgfpathlineto{\pgfqpoint{3.475074in}{1.049584in}}%
\pgfpathlineto{\pgfqpoint{3.477593in}{1.048061in}}%
\pgfpathlineto{\pgfqpoint{3.480112in}{1.074577in}}%
\pgfpathlineto{\pgfqpoint{3.482631in}{1.067770in}}%
\pgfpathlineto{\pgfqpoint{3.485149in}{1.054268in}}%
\pgfpathlineto{\pgfqpoint{3.487668in}{1.069117in}}%
\pgfpathlineto{\pgfqpoint{3.490187in}{1.073163in}}%
\pgfpathlineto{\pgfqpoint{3.492706in}{1.097778in}}%
\pgfpathlineto{\pgfqpoint{3.495224in}{1.101535in}}%
\pgfpathlineto{\pgfqpoint{3.497743in}{1.108330in}}%
\pgfpathlineto{\pgfqpoint{3.500262in}{1.103272in}}%
\pgfpathlineto{\pgfqpoint{3.502781in}{1.117696in}}%
\pgfpathlineto{\pgfqpoint{3.505299in}{1.137535in}}%
\pgfpathlineto{\pgfqpoint{3.507818in}{1.150366in}}%
\pgfpathlineto{\pgfqpoint{3.510337in}{1.136698in}}%
\pgfpathlineto{\pgfqpoint{3.512856in}{1.126818in}}%
\pgfpathlineto{\pgfqpoint{3.515374in}{1.099894in}}%
\pgfpathlineto{\pgfqpoint{3.517893in}{1.114323in}}%
\pgfpathlineto{\pgfqpoint{3.520412in}{1.135472in}}%
\pgfpathlineto{\pgfqpoint{3.522931in}{1.122995in}}%
\pgfpathlineto{\pgfqpoint{3.525449in}{1.119466in}}%
\pgfpathlineto{\pgfqpoint{3.527968in}{1.146544in}}%
\pgfpathlineto{\pgfqpoint{3.530487in}{1.169927in}}%
\pgfpathlineto{\pgfqpoint{3.533006in}{1.155945in}}%
\pgfpathlineto{\pgfqpoint{3.535524in}{1.157988in}}%
\pgfpathlineto{\pgfqpoint{3.538043in}{1.179438in}}%
\pgfpathlineto{\pgfqpoint{3.540562in}{1.179103in}}%
\pgfpathlineto{\pgfqpoint{3.543081in}{1.187985in}}%
\pgfpathlineto{\pgfqpoint{3.545599in}{1.170195in}}%
\pgfpathlineto{\pgfqpoint{3.548118in}{1.178869in}}%
\pgfpathlineto{\pgfqpoint{3.550637in}{1.181557in}}%
\pgfpathlineto{\pgfqpoint{3.553156in}{1.182416in}}%
\pgfpathlineto{\pgfqpoint{3.555674in}{1.165558in}}%
\pgfpathlineto{\pgfqpoint{3.558193in}{1.158519in}}%
\pgfpathlineto{\pgfqpoint{3.560712in}{1.168029in}}%
\pgfpathlineto{\pgfqpoint{3.563231in}{1.176355in}}%
\pgfpathlineto{\pgfqpoint{3.565749in}{1.172174in}}%
\pgfpathlineto{\pgfqpoint{3.568268in}{1.171315in}}%
\pgfpathlineto{\pgfqpoint{3.570787in}{1.170252in}}%
\pgfpathlineto{\pgfqpoint{3.573306in}{1.158539in}}%
\pgfpathlineto{\pgfqpoint{3.575824in}{1.138544in}}%
\pgfpathlineto{\pgfqpoint{3.578343in}{1.135229in}}%
\pgfpathlineto{\pgfqpoint{3.580862in}{1.145500in}}%
\pgfpathlineto{\pgfqpoint{3.583381in}{1.148122in}}%
\pgfpathlineto{\pgfqpoint{3.585899in}{1.138407in}}%
\pgfpathlineto{\pgfqpoint{3.588418in}{1.114905in}}%
\pgfpathlineto{\pgfqpoint{3.590937in}{1.130518in}}%
\pgfpathlineto{\pgfqpoint{3.593456in}{1.133254in}}%
\pgfpathlineto{\pgfqpoint{3.595974in}{1.141846in}}%
\pgfpathlineto{\pgfqpoint{3.598493in}{1.138383in}}%
\pgfpathlineto{\pgfqpoint{3.601012in}{1.128413in}}%
\pgfpathlineto{\pgfqpoint{3.603531in}{1.132549in}}%
\pgfpathlineto{\pgfqpoint{3.606049in}{1.141773in}}%
\pgfpathlineto{\pgfqpoint{3.608568in}{1.142162in}}%
\pgfpathlineto{\pgfqpoint{3.611087in}{1.142845in}}%
\pgfpathlineto{\pgfqpoint{3.613606in}{1.115198in}}%
\pgfpathlineto{\pgfqpoint{3.616124in}{1.116352in}}%
\pgfpathlineto{\pgfqpoint{3.618643in}{1.138765in}}%
\pgfpathlineto{\pgfqpoint{3.621162in}{1.147686in}}%
\pgfpathlineto{\pgfqpoint{3.623681in}{1.160000in}}%
\pgfpathlineto{\pgfqpoint{3.626199in}{1.160032in}}%
\pgfpathlineto{\pgfqpoint{3.628718in}{1.153867in}}%
\pgfpathlineto{\pgfqpoint{3.631237in}{1.142369in}}%
\pgfpathlineto{\pgfqpoint{3.633756in}{1.124455in}}%
\pgfpathlineto{\pgfqpoint{3.636274in}{1.129534in}}%
\pgfpathlineto{\pgfqpoint{3.638793in}{1.113315in}}%
\pgfpathlineto{\pgfqpoint{3.641312in}{1.101909in}}%
\pgfpathlineto{\pgfqpoint{3.643831in}{1.118855in}}%
\pgfpathlineto{\pgfqpoint{3.646349in}{1.116711in}}%
\pgfpathlineto{\pgfqpoint{3.648868in}{1.105916in}}%
\pgfpathlineto{\pgfqpoint{3.651387in}{1.103788in}}%
\pgfpathlineto{\pgfqpoint{3.653906in}{1.092089in}}%
\pgfpathlineto{\pgfqpoint{3.656424in}{1.103313in}}%
\pgfpathlineto{\pgfqpoint{3.658943in}{1.123634in}}%
\pgfpathlineto{\pgfqpoint{3.661462in}{1.141897in}}%
\pgfpathlineto{\pgfqpoint{3.663981in}{1.154138in}}%
\pgfpathlineto{\pgfqpoint{3.666499in}{1.159770in}}%
\pgfpathlineto{\pgfqpoint{3.669018in}{1.155909in}}%
\pgfpathlineto{\pgfqpoint{3.671537in}{1.180654in}}%
\pgfpathlineto{\pgfqpoint{3.674056in}{1.167190in}}%
\pgfpathlineto{\pgfqpoint{3.676574in}{1.164173in}}%
\pgfpathlineto{\pgfqpoint{3.679093in}{1.173484in}}%
\pgfpathlineto{\pgfqpoint{3.681612in}{1.166938in}}%
\pgfpathlineto{\pgfqpoint{3.684131in}{1.172017in}}%
\pgfpathlineto{\pgfqpoint{3.686649in}{1.183414in}}%
\pgfpathlineto{\pgfqpoint{3.689168in}{1.177224in}}%
\pgfpathlineto{\pgfqpoint{3.691687in}{1.163553in}}%
\pgfpathlineto{\pgfqpoint{3.694206in}{1.172541in}}%
\pgfpathlineto{\pgfqpoint{3.696724in}{1.178799in}}%
\pgfpathlineto{\pgfqpoint{3.699243in}{1.152573in}}%
\pgfpathlineto{\pgfqpoint{3.701762in}{1.138075in}}%
\pgfpathlineto{\pgfqpoint{3.704281in}{1.146439in}}%
\pgfpathlineto{\pgfqpoint{3.706799in}{1.160229in}}%
\pgfpathlineto{\pgfqpoint{3.709318in}{1.167883in}}%
\pgfpathlineto{\pgfqpoint{3.711837in}{1.166587in}}%
\pgfpathlineto{\pgfqpoint{3.714356in}{1.168364in}}%
\pgfpathlineto{\pgfqpoint{3.716874in}{1.145142in}}%
\pgfpathlineto{\pgfqpoint{3.719393in}{1.154840in}}%
\pgfpathlineto{\pgfqpoint{3.721912in}{1.175192in}}%
\pgfpathlineto{\pgfqpoint{3.724431in}{1.156964in}}%
\pgfpathlineto{\pgfqpoint{3.726949in}{1.164320in}}%
\pgfpathlineto{\pgfqpoint{3.729468in}{1.149804in}}%
\pgfpathlineto{\pgfqpoint{3.731987in}{1.168419in}}%
\pgfpathlineto{\pgfqpoint{3.734506in}{1.171379in}}%
\pgfpathlineto{\pgfqpoint{3.737024in}{1.173360in}}%
\pgfpathlineto{\pgfqpoint{3.739543in}{1.171501in}}%
\pgfpathlineto{\pgfqpoint{3.742062in}{1.165638in}}%
\pgfpathlineto{\pgfqpoint{3.744581in}{1.174134in}}%
\pgfpathlineto{\pgfqpoint{3.747099in}{1.166881in}}%
\pgfpathlineto{\pgfqpoint{3.749618in}{1.149641in}}%
\pgfpathlineto{\pgfqpoint{3.752137in}{1.151761in}}%
\pgfpathlineto{\pgfqpoint{3.754656in}{1.136089in}}%
\pgfpathlineto{\pgfqpoint{3.757174in}{1.157873in}}%
\pgfpathlineto{\pgfqpoint{3.759693in}{1.169350in}}%
\pgfpathlineto{\pgfqpoint{3.762212in}{1.179003in}}%
\pgfpathlineto{\pgfqpoint{3.764731in}{1.212397in}}%
\pgfpathlineto{\pgfqpoint{3.767249in}{1.213755in}}%
\pgfpathlineto{\pgfqpoint{3.769768in}{1.238683in}}%
\pgfpathlineto{\pgfqpoint{3.772287in}{1.218156in}}%
\pgfpathlineto{\pgfqpoint{3.774806in}{1.211070in}}%
\pgfpathlineto{\pgfqpoint{3.777324in}{1.218483in}}%
\pgfpathlineto{\pgfqpoint{3.779843in}{1.210735in}}%
\pgfpathlineto{\pgfqpoint{3.782362in}{1.231825in}}%
\pgfpathlineto{\pgfqpoint{3.784881in}{1.250678in}}%
\pgfpathlineto{\pgfqpoint{3.787399in}{1.227234in}}%
\pgfpathlineto{\pgfqpoint{3.789918in}{1.211015in}}%
\pgfpathlineto{\pgfqpoint{3.792437in}{1.211915in}}%
\pgfpathlineto{\pgfqpoint{3.794956in}{1.193925in}}%
\pgfpathlineto{\pgfqpoint{3.797474in}{1.201071in}}%
\pgfpathlineto{\pgfqpoint{3.799993in}{1.180182in}}%
\pgfpathlineto{\pgfqpoint{3.802512in}{1.168149in}}%
\pgfpathlineto{\pgfqpoint{3.805031in}{1.148481in}}%
\pgfpathlineto{\pgfqpoint{3.807549in}{1.147095in}}%
\pgfpathlineto{\pgfqpoint{3.810068in}{1.177569in}}%
\pgfpathlineto{\pgfqpoint{3.812587in}{1.191324in}}%
\pgfpathlineto{\pgfqpoint{3.815106in}{1.182810in}}%
\pgfpathlineto{\pgfqpoint{3.817624in}{1.187394in}}%
\pgfpathlineto{\pgfqpoint{3.820143in}{1.182704in}}%
\pgfpathlineto{\pgfqpoint{3.822662in}{1.200201in}}%
\pgfpathlineto{\pgfqpoint{3.825181in}{1.165999in}}%
\pgfpathlineto{\pgfqpoint{3.827699in}{1.165561in}}%
\pgfpathlineto{\pgfqpoint{3.830218in}{1.184762in}}%
\pgfpathlineto{\pgfqpoint{3.832737in}{1.184722in}}%
\pgfpathlineto{\pgfqpoint{3.835256in}{1.193228in}}%
\pgfpathlineto{\pgfqpoint{3.837774in}{1.187683in}}%
\pgfpathlineto{\pgfqpoint{3.840293in}{1.186621in}}%
\pgfpathlineto{\pgfqpoint{3.842812in}{1.173900in}}%
\pgfpathlineto{\pgfqpoint{3.845331in}{1.157280in}}%
\pgfpathlineto{\pgfqpoint{3.847849in}{1.149846in}}%
\pgfpathlineto{\pgfqpoint{3.850368in}{1.157511in}}%
\pgfpathlineto{\pgfqpoint{3.852887in}{1.138032in}}%
\pgfpathlineto{\pgfqpoint{3.855406in}{1.156150in}}%
\pgfpathlineto{\pgfqpoint{3.857924in}{1.164694in}}%
\pgfpathlineto{\pgfqpoint{3.860443in}{1.166658in}}%
\pgfpathlineto{\pgfqpoint{3.862962in}{1.145993in}}%
\pgfpathlineto{\pgfqpoint{3.865481in}{1.130310in}}%
\pgfpathlineto{\pgfqpoint{3.867999in}{1.138474in}}%
\pgfpathlineto{\pgfqpoint{3.870518in}{1.119976in}}%
\pgfpathlineto{\pgfqpoint{3.873037in}{1.115157in}}%
\pgfpathlineto{\pgfqpoint{3.875556in}{1.122650in}}%
\pgfpathlineto{\pgfqpoint{3.878074in}{1.135648in}}%
\pgfpathlineto{\pgfqpoint{3.880593in}{1.164367in}}%
\pgfpathlineto{\pgfqpoint{3.883112in}{1.150937in}}%
\pgfpathlineto{\pgfqpoint{3.885631in}{1.170793in}}%
\pgfpathlineto{\pgfqpoint{3.888149in}{1.172046in}}%
\pgfpathlineto{\pgfqpoint{3.890668in}{1.151557in}}%
\pgfpathlineto{\pgfqpoint{3.893187in}{1.134409in}}%
\pgfpathlineto{\pgfqpoint{3.895706in}{1.150091in}}%
\pgfpathlineto{\pgfqpoint{3.898224in}{1.142260in}}%
\pgfpathlineto{\pgfqpoint{3.900743in}{1.155676in}}%
\pgfpathlineto{\pgfqpoint{3.903262in}{1.165128in}}%
\pgfpathlineto{\pgfqpoint{3.905781in}{1.162572in}}%
\pgfpathlineto{\pgfqpoint{3.908299in}{1.185060in}}%
\pgfpathlineto{\pgfqpoint{3.910818in}{1.201274in}}%
\pgfpathlineto{\pgfqpoint{3.913337in}{1.212778in}}%
\pgfpathlineto{\pgfqpoint{3.915856in}{1.218517in}}%
\pgfpathlineto{\pgfqpoint{3.918374in}{1.194041in}}%
\pgfpathlineto{\pgfqpoint{3.920893in}{1.194096in}}%
\pgfpathlineto{\pgfqpoint{3.923412in}{1.199225in}}%
\pgfpathlineto{\pgfqpoint{3.925931in}{1.243530in}}%
\pgfpathlineto{\pgfqpoint{3.928449in}{1.228419in}}%
\pgfpathlineto{\pgfqpoint{3.930968in}{1.240681in}}%
\pgfpathlineto{\pgfqpoint{3.933487in}{1.224448in}}%
\pgfpathlineto{\pgfqpoint{3.936006in}{1.234090in}}%
\pgfpathlineto{\pgfqpoint{3.938524in}{1.230942in}}%
\pgfpathlineto{\pgfqpoint{3.941043in}{1.238497in}}%
\pgfpathlineto{\pgfqpoint{3.943562in}{1.209966in}}%
\pgfpathlineto{\pgfqpoint{3.946081in}{1.208314in}}%
\pgfpathlineto{\pgfqpoint{3.948599in}{1.195976in}}%
\pgfpathlineto{\pgfqpoint{3.951118in}{1.212941in}}%
\pgfpathlineto{\pgfqpoint{3.953637in}{1.210943in}}%
\pgfpathlineto{\pgfqpoint{3.956156in}{1.229613in}}%
\pgfpathlineto{\pgfqpoint{3.958674in}{1.247104in}}%
\pgfpathlineto{\pgfqpoint{3.961193in}{1.235059in}}%
\pgfpathlineto{\pgfqpoint{3.963712in}{1.231960in}}%
\pgfpathlineto{\pgfqpoint{3.966231in}{1.235360in}}%
\pgfpathlineto{\pgfqpoint{3.968749in}{1.235429in}}%
\pgfpathlineto{\pgfqpoint{3.971268in}{1.227170in}}%
\pgfpathlineto{\pgfqpoint{3.973787in}{1.237233in}}%
\pgfpathlineto{\pgfqpoint{3.976306in}{1.243288in}}%
\pgfpathlineto{\pgfqpoint{3.978824in}{1.218344in}}%
\pgfpathlineto{\pgfqpoint{3.981343in}{1.220290in}}%
\pgfpathlineto{\pgfqpoint{3.983862in}{1.218318in}}%
\pgfpathlineto{\pgfqpoint{3.986381in}{1.203207in}}%
\pgfpathlineto{\pgfqpoint{3.988899in}{1.219142in}}%
\pgfpathlineto{\pgfqpoint{3.991418in}{1.224420in}}%
\pgfpathlineto{\pgfqpoint{3.993937in}{1.216121in}}%
\pgfpathlineto{\pgfqpoint{3.996456in}{1.222773in}}%
\pgfpathlineto{\pgfqpoint{3.998974in}{1.225021in}}%
\pgfpathlineto{\pgfqpoint{4.001493in}{1.220468in}}%
\pgfpathlineto{\pgfqpoint{4.004012in}{1.237510in}}%
\pgfpathlineto{\pgfqpoint{4.006531in}{1.211027in}}%
\pgfpathlineto{\pgfqpoint{4.009049in}{1.231457in}}%
\pgfpathlineto{\pgfqpoint{4.011568in}{1.239426in}}%
\pgfpathlineto{\pgfqpoint{4.014087in}{1.237031in}}%
\pgfpathlineto{\pgfqpoint{4.016606in}{1.236411in}}%
\pgfpathlineto{\pgfqpoint{4.019124in}{1.241773in}}%
\pgfpathlineto{\pgfqpoint{4.021643in}{1.226376in}}%
\pgfpathlineto{\pgfqpoint{4.024162in}{1.197807in}}%
\pgfpathlineto{\pgfqpoint{4.026681in}{1.190935in}}%
\pgfpathlineto{\pgfqpoint{4.029199in}{1.191083in}}%
\pgfpathlineto{\pgfqpoint{4.031718in}{1.205332in}}%
\pgfpathlineto{\pgfqpoint{4.034237in}{1.170189in}}%
\pgfpathlineto{\pgfqpoint{4.036756in}{1.166503in}}%
\pgfpathlineto{\pgfqpoint{4.039274in}{1.182975in}}%
\pgfpathlineto{\pgfqpoint{4.041793in}{1.174006in}}%
\pgfpathlineto{\pgfqpoint{4.044312in}{1.161817in}}%
\pgfpathlineto{\pgfqpoint{4.046831in}{1.121055in}}%
\pgfpathlineto{\pgfqpoint{4.049349in}{1.114953in}}%
\pgfpathlineto{\pgfqpoint{4.051868in}{1.110482in}}%
\pgfpathlineto{\pgfqpoint{4.054387in}{1.110536in}}%
\pgfpathlineto{\pgfqpoint{4.056906in}{1.111451in}}%
\pgfpathlineto{\pgfqpoint{4.059424in}{1.128038in}}%
\pgfpathlineto{\pgfqpoint{4.061943in}{1.117591in}}%
\pgfpathlineto{\pgfqpoint{4.066981in}{1.104924in}}%
\pgfpathlineto{\pgfqpoint{4.069499in}{1.095198in}}%
\pgfpathlineto{\pgfqpoint{4.072018in}{1.080785in}}%
\pgfpathlineto{\pgfqpoint{4.074537in}{1.086383in}}%
\pgfpathlineto{\pgfqpoint{4.077056in}{1.065943in}}%
\pgfpathlineto{\pgfqpoint{4.079574in}{1.071763in}}%
\pgfpathlineto{\pgfqpoint{4.082093in}{1.067174in}}%
\pgfpathlineto{\pgfqpoint{4.084612in}{1.063034in}}%
\pgfpathlineto{\pgfqpoint{4.087131in}{1.065166in}}%
\pgfpathlineto{\pgfqpoint{4.089649in}{1.064467in}}%
\pgfpathlineto{\pgfqpoint{4.092168in}{1.080634in}}%
\pgfpathlineto{\pgfqpoint{4.094687in}{1.078608in}}%
\pgfpathlineto{\pgfqpoint{4.097206in}{1.087357in}}%
\pgfpathlineto{\pgfqpoint{4.099724in}{1.093169in}}%
\pgfpathlineto{\pgfqpoint{4.102243in}{1.091373in}}%
\pgfpathlineto{\pgfqpoint{4.104762in}{1.081045in}}%
\pgfpathlineto{\pgfqpoint{4.107281in}{1.097277in}}%
\pgfpathlineto{\pgfqpoint{4.109799in}{1.101652in}}%
\pgfpathlineto{\pgfqpoint{4.112318in}{1.115302in}}%
\pgfpathlineto{\pgfqpoint{4.114837in}{1.120716in}}%
\pgfpathlineto{\pgfqpoint{4.117356in}{1.118905in}}%
\pgfpathlineto{\pgfqpoint{4.119874in}{1.119590in}}%
\pgfpathlineto{\pgfqpoint{4.122393in}{1.124149in}}%
\pgfpathlineto{\pgfqpoint{4.124912in}{1.124924in}}%
\pgfpathlineto{\pgfqpoint{4.127431in}{1.117251in}}%
\pgfpathlineto{\pgfqpoint{4.129949in}{1.094631in}}%
\pgfpathlineto{\pgfqpoint{4.132468in}{1.116663in}}%
\pgfpathlineto{\pgfqpoint{4.134987in}{1.117006in}}%
\pgfpathlineto{\pgfqpoint{4.137506in}{1.109423in}}%
\pgfpathlineto{\pgfqpoint{4.140024in}{1.120941in}}%
\pgfpathlineto{\pgfqpoint{4.142543in}{1.107694in}}%
\pgfpathlineto{\pgfqpoint{4.145062in}{1.106473in}}%
\pgfpathlineto{\pgfqpoint{4.147581in}{1.091621in}}%
\pgfpathlineto{\pgfqpoint{4.150099in}{1.091275in}}%
\pgfpathlineto{\pgfqpoint{4.152618in}{1.089202in}}%
\pgfpathlineto{\pgfqpoint{4.155137in}{1.091322in}}%
\pgfpathlineto{\pgfqpoint{4.157656in}{1.118689in}}%
\pgfpathlineto{\pgfqpoint{4.160174in}{1.128668in}}%
\pgfpathlineto{\pgfqpoint{4.165212in}{1.099389in}}%
\pgfpathlineto{\pgfqpoint{4.167731in}{1.107083in}}%
\pgfpathlineto{\pgfqpoint{4.170249in}{1.131542in}}%
\pgfpathlineto{\pgfqpoint{4.172768in}{1.142117in}}%
\pgfpathlineto{\pgfqpoint{4.175287in}{1.137077in}}%
\pgfpathlineto{\pgfqpoint{4.177806in}{1.120457in}}%
\pgfpathlineto{\pgfqpoint{4.180324in}{1.159805in}}%
\pgfpathlineto{\pgfqpoint{4.182843in}{1.159653in}}%
\pgfpathlineto{\pgfqpoint{4.185362in}{1.170196in}}%
\pgfpathlineto{\pgfqpoint{4.187881in}{1.160737in}}%
\pgfpathlineto{\pgfqpoint{4.190399in}{1.164148in}}%
\pgfpathlineto{\pgfqpoint{4.192918in}{1.154257in}}%
\pgfpathlineto{\pgfqpoint{4.195437in}{1.143695in}}%
\pgfpathlineto{\pgfqpoint{4.197956in}{1.154915in}}%
\pgfpathlineto{\pgfqpoint{4.200474in}{1.166571in}}%
\pgfpathlineto{\pgfqpoint{4.202993in}{1.177192in}}%
\pgfpathlineto{\pgfqpoint{4.205512in}{1.165870in}}%
\pgfpathlineto{\pgfqpoint{4.208031in}{1.155736in}}%
\pgfpathlineto{\pgfqpoint{4.210549in}{1.164878in}}%
\pgfpathlineto{\pgfqpoint{4.213068in}{1.152600in}}%
\pgfpathlineto{\pgfqpoint{4.215587in}{1.182504in}}%
\pgfpathlineto{\pgfqpoint{4.218106in}{1.186271in}}%
\pgfpathlineto{\pgfqpoint{4.220624in}{1.173292in}}%
\pgfpathlineto{\pgfqpoint{4.223143in}{1.179801in}}%
\pgfpathlineto{\pgfqpoint{4.225662in}{1.190694in}}%
\pgfpathlineto{\pgfqpoint{4.228181in}{1.196075in}}%
\pgfpathlineto{\pgfqpoint{4.230699in}{1.194076in}}%
\pgfpathlineto{\pgfqpoint{4.233218in}{1.221129in}}%
\pgfpathlineto{\pgfqpoint{4.235737in}{1.209390in}}%
\pgfpathlineto{\pgfqpoint{4.238256in}{1.204673in}}%
\pgfpathlineto{\pgfqpoint{4.240774in}{1.218886in}}%
\pgfpathlineto{\pgfqpoint{4.243293in}{1.212265in}}%
\pgfpathlineto{\pgfqpoint{4.245812in}{1.210032in}}%
\pgfpathlineto{\pgfqpoint{4.248331in}{1.230826in}}%
\pgfpathlineto{\pgfqpoint{4.250849in}{1.207514in}}%
\pgfpathlineto{\pgfqpoint{4.253368in}{1.209944in}}%
\pgfpathlineto{\pgfqpoint{4.255887in}{1.197271in}}%
\pgfpathlineto{\pgfqpoint{4.258406in}{1.198838in}}%
\pgfpathlineto{\pgfqpoint{4.260924in}{1.208251in}}%
\pgfpathlineto{\pgfqpoint{4.263443in}{1.210940in}}%
\pgfpathlineto{\pgfqpoint{4.265962in}{1.209092in}}%
\pgfpathlineto{\pgfqpoint{4.268481in}{1.190940in}}%
\pgfpathlineto{\pgfqpoint{4.270999in}{1.201694in}}%
\pgfpathlineto{\pgfqpoint{4.273518in}{1.240004in}}%
\pgfpathlineto{\pgfqpoint{4.276037in}{1.272704in}}%
\pgfpathlineto{\pgfqpoint{4.278556in}{1.271839in}}%
\pgfpathlineto{\pgfqpoint{4.281074in}{1.255840in}}%
\pgfpathlineto{\pgfqpoint{4.283593in}{1.252835in}}%
\pgfpathlineto{\pgfqpoint{4.286112in}{1.257252in}}%
\pgfpathlineto{\pgfqpoint{4.288631in}{1.266143in}}%
\pgfpathlineto{\pgfqpoint{4.291149in}{1.278294in}}%
\pgfpathlineto{\pgfqpoint{4.293668in}{1.289832in}}%
\pgfpathlineto{\pgfqpoint{4.296187in}{1.287904in}}%
\pgfpathlineto{\pgfqpoint{4.298706in}{1.289165in}}%
\pgfpathlineto{\pgfqpoint{4.301224in}{1.298783in}}%
\pgfpathlineto{\pgfqpoint{4.303743in}{1.288674in}}%
\pgfpathlineto{\pgfqpoint{4.306262in}{1.283187in}}%
\pgfpathlineto{\pgfqpoint{4.308781in}{1.276714in}}%
\pgfpathlineto{\pgfqpoint{4.311299in}{1.254316in}}%
\pgfpathlineto{\pgfqpoint{4.313818in}{1.266778in}}%
\pgfpathlineto{\pgfqpoint{4.316337in}{1.275803in}}%
\pgfpathlineto{\pgfqpoint{4.318856in}{1.301128in}}%
\pgfpathlineto{\pgfqpoint{4.321374in}{1.294654in}}%
\pgfpathlineto{\pgfqpoint{4.323893in}{1.286707in}}%
\pgfpathlineto{\pgfqpoint{4.326412in}{1.271905in}}%
\pgfpathlineto{\pgfqpoint{4.328931in}{1.271551in}}%
\pgfpathlineto{\pgfqpoint{4.331449in}{1.262181in}}%
\pgfpathlineto{\pgfqpoint{4.333968in}{1.279189in}}%
\pgfpathlineto{\pgfqpoint{4.336487in}{1.269439in}}%
\pgfpathlineto{\pgfqpoint{4.339006in}{1.271404in}}%
\pgfpathlineto{\pgfqpoint{4.341524in}{1.243828in}}%
\pgfpathlineto{\pgfqpoint{4.344043in}{1.219722in}}%
\pgfpathlineto{\pgfqpoint{4.346562in}{1.223790in}}%
\pgfpathlineto{\pgfqpoint{4.349081in}{1.233249in}}%
\pgfpathlineto{\pgfqpoint{4.351599in}{1.246270in}}%
\pgfpathlineto{\pgfqpoint{4.354118in}{1.264167in}}%
\pgfpathlineto{\pgfqpoint{4.356637in}{1.267517in}}%
\pgfpathlineto{\pgfqpoint{4.359156in}{1.281249in}}%
\pgfpathlineto{\pgfqpoint{4.361674in}{1.290120in}}%
\pgfpathlineto{\pgfqpoint{4.364193in}{1.276135in}}%
\pgfpathlineto{\pgfqpoint{4.366712in}{1.294432in}}%
\pgfpathlineto{\pgfqpoint{4.369231in}{1.304164in}}%
\pgfpathlineto{\pgfqpoint{4.371749in}{1.296843in}}%
\pgfpathlineto{\pgfqpoint{4.374268in}{1.298661in}}%
\pgfpathlineto{\pgfqpoint{4.376787in}{1.328799in}}%
\pgfpathlineto{\pgfqpoint{4.379306in}{1.334055in}}%
\pgfpathlineto{\pgfqpoint{4.381824in}{1.323051in}}%
\pgfpathlineto{\pgfqpoint{4.384343in}{1.314084in}}%
\pgfpathlineto{\pgfqpoint{4.386862in}{1.297400in}}%
\pgfpathlineto{\pgfqpoint{4.389381in}{1.296039in}}%
\pgfpathlineto{\pgfqpoint{4.391899in}{1.297422in}}%
\pgfpathlineto{\pgfqpoint{4.394418in}{1.296440in}}%
\pgfpathlineto{\pgfqpoint{4.396937in}{1.308221in}}%
\pgfpathlineto{\pgfqpoint{4.399456in}{1.295800in}}%
\pgfpathlineto{\pgfqpoint{4.401974in}{1.286294in}}%
\pgfpathlineto{\pgfqpoint{4.404493in}{1.306714in}}%
\pgfpathlineto{\pgfqpoint{4.407012in}{1.334902in}}%
\pgfpathlineto{\pgfqpoint{4.409531in}{1.330757in}}%
\pgfpathlineto{\pgfqpoint{4.412049in}{1.346095in}}%
\pgfpathlineto{\pgfqpoint{4.414568in}{1.333044in}}%
\pgfpathlineto{\pgfqpoint{4.417087in}{1.285329in}}%
\pgfpathlineto{\pgfqpoint{4.419606in}{1.291438in}}%
\pgfpathlineto{\pgfqpoint{4.422124in}{1.299838in}}%
\pgfpathlineto{\pgfqpoint{4.424643in}{1.299374in}}%
\pgfpathlineto{\pgfqpoint{4.427162in}{1.295721in}}%
\pgfpathlineto{\pgfqpoint{4.429681in}{1.307947in}}%
\pgfpathlineto{\pgfqpoint{4.432199in}{1.256119in}}%
\pgfpathlineto{\pgfqpoint{4.434718in}{1.255131in}}%
\pgfpathlineto{\pgfqpoint{4.437237in}{1.258860in}}%
\pgfpathlineto{\pgfqpoint{4.439756in}{1.257760in}}%
\pgfpathlineto{\pgfqpoint{4.442274in}{1.243383in}}%
\pgfpathlineto{\pgfqpoint{4.444793in}{1.278036in}}%
\pgfpathlineto{\pgfqpoint{4.447312in}{1.244956in}}%
\pgfpathlineto{\pgfqpoint{4.449831in}{1.240287in}}%
\pgfpathlineto{\pgfqpoint{4.452349in}{1.250932in}}%
\pgfpathlineto{\pgfqpoint{4.454868in}{1.247164in}}%
\pgfpathlineto{\pgfqpoint{4.457387in}{1.250987in}}%
\pgfpathlineto{\pgfqpoint{4.459906in}{1.274266in}}%
\pgfpathlineto{\pgfqpoint{4.462424in}{1.278046in}}%
\pgfpathlineto{\pgfqpoint{4.464943in}{1.291546in}}%
\pgfpathlineto{\pgfqpoint{4.467462in}{1.287346in}}%
\pgfpathlineto{\pgfqpoint{4.469981in}{1.310342in}}%
\pgfpathlineto{\pgfqpoint{4.472499in}{1.323335in}}%
\pgfpathlineto{\pgfqpoint{4.475018in}{1.328136in}}%
\pgfpathlineto{\pgfqpoint{4.477537in}{1.319980in}}%
\pgfpathlineto{\pgfqpoint{4.480056in}{1.326608in}}%
\pgfpathlineto{\pgfqpoint{4.482574in}{1.321135in}}%
\pgfpathlineto{\pgfqpoint{4.485093in}{1.339846in}}%
\pgfpathlineto{\pgfqpoint{4.487612in}{1.323558in}}%
\pgfpathlineto{\pgfqpoint{4.490131in}{1.319813in}}%
\pgfpathlineto{\pgfqpoint{4.492649in}{1.342983in}}%
\pgfpathlineto{\pgfqpoint{4.495168in}{1.358504in}}%
\pgfpathlineto{\pgfqpoint{4.497687in}{1.376331in}}%
\pgfpathlineto{\pgfqpoint{4.500206in}{1.360208in}}%
\pgfpathlineto{\pgfqpoint{4.502724in}{1.354428in}}%
\pgfpathlineto{\pgfqpoint{4.505243in}{1.347390in}}%
\pgfpathlineto{\pgfqpoint{4.507762in}{1.330377in}}%
\pgfpathlineto{\pgfqpoint{4.510281in}{1.310015in}}%
\pgfpathlineto{\pgfqpoint{4.512799in}{1.310997in}}%
\pgfpathlineto{\pgfqpoint{4.515318in}{1.312117in}}%
\pgfpathlineto{\pgfqpoint{4.517837in}{1.321553in}}%
\pgfpathlineto{\pgfqpoint{4.520356in}{1.309867in}}%
\pgfpathlineto{\pgfqpoint{4.522874in}{1.318164in}}%
\pgfpathlineto{\pgfqpoint{4.525393in}{1.277909in}}%
\pgfpathlineto{\pgfqpoint{4.530431in}{1.299472in}}%
\pgfpathlineto{\pgfqpoint{4.532949in}{1.280454in}}%
\pgfpathlineto{\pgfqpoint{4.535468in}{1.271974in}}%
\pgfpathlineto{\pgfqpoint{4.537987in}{1.296600in}}%
\pgfpathlineto{\pgfqpoint{4.540506in}{1.298586in}}%
\pgfpathlineto{\pgfqpoint{4.543024in}{1.279176in}}%
\pgfpathlineto{\pgfqpoint{4.545543in}{1.296611in}}%
\pgfpathlineto{\pgfqpoint{4.548062in}{1.293422in}}%
\pgfpathlineto{\pgfqpoint{4.550581in}{1.283958in}}%
\pgfpathlineto{\pgfqpoint{4.553099in}{1.281075in}}%
\pgfpathlineto{\pgfqpoint{4.555618in}{1.299677in}}%
\pgfpathlineto{\pgfqpoint{4.558137in}{1.306917in}}%
\pgfpathlineto{\pgfqpoint{4.560656in}{1.300739in}}%
\pgfpathlineto{\pgfqpoint{4.563174in}{1.303802in}}%
\pgfpathlineto{\pgfqpoint{4.565693in}{1.286537in}}%
\pgfpathlineto{\pgfqpoint{4.568212in}{1.294361in}}%
\pgfpathlineto{\pgfqpoint{4.570731in}{1.293527in}}%
\pgfpathlineto{\pgfqpoint{4.573249in}{1.280284in}}%
\pgfpathlineto{\pgfqpoint{4.575768in}{1.266533in}}%
\pgfpathlineto{\pgfqpoint{4.578287in}{1.277561in}}%
\pgfpathlineto{\pgfqpoint{4.580806in}{1.256267in}}%
\pgfpathlineto{\pgfqpoint{4.583324in}{1.281071in}}%
\pgfpathlineto{\pgfqpoint{4.585843in}{1.281939in}}%
\pgfpathlineto{\pgfqpoint{4.588362in}{1.255550in}}%
\pgfpathlineto{\pgfqpoint{4.590881in}{1.241806in}}%
\pgfpathlineto{\pgfqpoint{4.593399in}{1.276058in}}%
\pgfpathlineto{\pgfqpoint{4.595918in}{1.284707in}}%
\pgfpathlineto{\pgfqpoint{4.598437in}{1.290594in}}%
\pgfpathlineto{\pgfqpoint{4.600956in}{1.285272in}}%
\pgfpathlineto{\pgfqpoint{4.603474in}{1.279276in}}%
\pgfpathlineto{\pgfqpoint{4.605993in}{1.257343in}}%
\pgfpathlineto{\pgfqpoint{4.608512in}{1.270855in}}%
\pgfpathlineto{\pgfqpoint{4.611031in}{1.268738in}}%
\pgfpathlineto{\pgfqpoint{4.613549in}{1.255761in}}%
\pgfpathlineto{\pgfqpoint{4.616068in}{1.276464in}}%
\pgfpathlineto{\pgfqpoint{4.618587in}{1.291038in}}%
\pgfpathlineto{\pgfqpoint{4.621106in}{1.291229in}}%
\pgfpathlineto{\pgfqpoint{4.623624in}{1.285691in}}%
\pgfpathlineto{\pgfqpoint{4.626143in}{1.293791in}}%
\pgfpathlineto{\pgfqpoint{4.628662in}{1.298281in}}%
\pgfpathlineto{\pgfqpoint{4.631181in}{1.316493in}}%
\pgfpathlineto{\pgfqpoint{4.633699in}{1.324452in}}%
\pgfpathlineto{\pgfqpoint{4.636218in}{1.324648in}}%
\pgfpathlineto{\pgfqpoint{4.638737in}{1.331604in}}%
\pgfpathlineto{\pgfqpoint{4.641256in}{1.337655in}}%
\pgfpathlineto{\pgfqpoint{4.643774in}{1.334728in}}%
\pgfpathlineto{\pgfqpoint{4.646293in}{1.316466in}}%
\pgfpathlineto{\pgfqpoint{4.648812in}{1.314518in}}%
\pgfpathlineto{\pgfqpoint{4.651331in}{1.312887in}}%
\pgfpathlineto{\pgfqpoint{4.653849in}{1.299623in}}%
\pgfpathlineto{\pgfqpoint{4.656368in}{1.303612in}}%
\pgfpathlineto{\pgfqpoint{4.658887in}{1.302869in}}%
\pgfpathlineto{\pgfqpoint{4.661406in}{1.305639in}}%
\pgfpathlineto{\pgfqpoint{4.663924in}{1.324090in}}%
\pgfpathlineto{\pgfqpoint{4.666443in}{1.299483in}}%
\pgfpathlineto{\pgfqpoint{4.668962in}{1.338279in}}%
\pgfpathlineto{\pgfqpoint{4.671481in}{1.336179in}}%
\pgfpathlineto{\pgfqpoint{4.673999in}{1.350418in}}%
\pgfpathlineto{\pgfqpoint{4.676518in}{1.360722in}}%
\pgfpathlineto{\pgfqpoint{4.679037in}{1.344450in}}%
\pgfpathlineto{\pgfqpoint{4.681556in}{1.356196in}}%
\pgfpathlineto{\pgfqpoint{4.684074in}{1.364792in}}%
\pgfpathlineto{\pgfqpoint{4.686593in}{1.371621in}}%
\pgfpathlineto{\pgfqpoint{4.689112in}{1.387303in}}%
\pgfpathlineto{\pgfqpoint{4.691631in}{1.411242in}}%
\pgfpathlineto{\pgfqpoint{4.694149in}{1.451989in}}%
\pgfpathlineto{\pgfqpoint{4.696668in}{1.465877in}}%
\pgfpathlineto{\pgfqpoint{4.699187in}{1.465797in}}%
\pgfpathlineto{\pgfqpoint{4.701706in}{1.473650in}}%
\pgfpathlineto{\pgfqpoint{4.704224in}{1.472737in}}%
\pgfpathlineto{\pgfqpoint{4.706743in}{1.476261in}}%
\pgfpathlineto{\pgfqpoint{4.709262in}{1.460315in}}%
\pgfpathlineto{\pgfqpoint{4.711781in}{1.467212in}}%
\pgfpathlineto{\pgfqpoint{4.714299in}{1.453943in}}%
\pgfpathlineto{\pgfqpoint{4.716818in}{1.474179in}}%
\pgfpathlineto{\pgfqpoint{4.719337in}{1.473204in}}%
\pgfpathlineto{\pgfqpoint{4.721856in}{1.460821in}}%
\pgfpathlineto{\pgfqpoint{4.724374in}{1.429515in}}%
\pgfpathlineto{\pgfqpoint{4.726893in}{1.438984in}}%
\pgfpathlineto{\pgfqpoint{4.729412in}{1.434077in}}%
\pgfpathlineto{\pgfqpoint{4.731931in}{1.423424in}}%
\pgfpathlineto{\pgfqpoint{4.734449in}{1.426136in}}%
\pgfpathlineto{\pgfqpoint{4.736968in}{1.439544in}}%
\pgfpathlineto{\pgfqpoint{4.739487in}{1.425659in}}%
\pgfpathlineto{\pgfqpoint{4.742006in}{1.440272in}}%
\pgfpathlineto{\pgfqpoint{4.744524in}{1.447587in}}%
\pgfpathlineto{\pgfqpoint{4.747043in}{1.444892in}}%
\pgfpathlineto{\pgfqpoint{4.749562in}{1.444537in}}%
\pgfpathlineto{\pgfqpoint{4.752081in}{1.455962in}}%
\pgfpathlineto{\pgfqpoint{4.754599in}{1.453071in}}%
\pgfpathlineto{\pgfqpoint{4.757118in}{1.471379in}}%
\pgfpathlineto{\pgfqpoint{4.759637in}{1.475335in}}%
\pgfpathlineto{\pgfqpoint{4.762156in}{1.469365in}}%
\pgfpathlineto{\pgfqpoint{4.764674in}{1.468048in}}%
\pgfpathlineto{\pgfqpoint{4.767193in}{1.462883in}}%
\pgfpathlineto{\pgfqpoint{4.769712in}{1.469455in}}%
\pgfpathlineto{\pgfqpoint{4.772231in}{1.473545in}}%
\pgfpathlineto{\pgfqpoint{4.774749in}{1.473937in}}%
\pgfpathlineto{\pgfqpoint{4.777268in}{1.478852in}}%
\pgfpathlineto{\pgfqpoint{4.779787in}{1.481037in}}%
\pgfpathlineto{\pgfqpoint{4.782306in}{1.475467in}}%
\pgfpathlineto{\pgfqpoint{4.784824in}{1.464666in}}%
\pgfpathlineto{\pgfqpoint{4.787343in}{1.451151in}}%
\pgfpathlineto{\pgfqpoint{4.789862in}{1.489130in}}%
\pgfpathlineto{\pgfqpoint{4.792381in}{1.494081in}}%
\pgfpathlineto{\pgfqpoint{4.794899in}{1.501307in}}%
\pgfpathlineto{\pgfqpoint{4.797418in}{1.516530in}}%
\pgfpathlineto{\pgfqpoint{4.799937in}{1.514059in}}%
\pgfpathlineto{\pgfqpoint{4.802456in}{1.504644in}}%
\pgfpathlineto{\pgfqpoint{4.807493in}{1.524620in}}%
\pgfpathlineto{\pgfqpoint{4.810012in}{1.539136in}}%
\pgfpathlineto{\pgfqpoint{4.812531in}{1.528804in}}%
\pgfpathlineto{\pgfqpoint{4.815049in}{1.514050in}}%
\pgfpathlineto{\pgfqpoint{4.817568in}{1.509445in}}%
\pgfpathlineto{\pgfqpoint{4.820087in}{1.525076in}}%
\pgfpathlineto{\pgfqpoint{4.822606in}{1.522357in}}%
\pgfpathlineto{\pgfqpoint{4.825124in}{1.548995in}}%
\pgfpathlineto{\pgfqpoint{4.827643in}{1.543763in}}%
\pgfpathlineto{\pgfqpoint{4.830162in}{1.540735in}}%
\pgfpathlineto{\pgfqpoint{4.832681in}{1.547462in}}%
\pgfpathlineto{\pgfqpoint{4.835199in}{1.539721in}}%
\pgfpathlineto{\pgfqpoint{4.837718in}{1.522534in}}%
\pgfpathlineto{\pgfqpoint{4.840237in}{1.524596in}}%
\pgfpathlineto{\pgfqpoint{4.842756in}{1.483252in}}%
\pgfpathlineto{\pgfqpoint{4.845274in}{1.471159in}}%
\pgfpathlineto{\pgfqpoint{4.847793in}{1.465102in}}%
\pgfpathlineto{\pgfqpoint{4.850312in}{1.452857in}}%
\pgfpathlineto{\pgfqpoint{4.852831in}{1.438824in}}%
\pgfpathlineto{\pgfqpoint{4.855349in}{1.428292in}}%
\pgfpathlineto{\pgfqpoint{4.857868in}{1.445865in}}%
\pgfpathlineto{\pgfqpoint{4.860387in}{1.434941in}}%
\pgfpathlineto{\pgfqpoint{4.862906in}{1.417863in}}%
\pgfpathlineto{\pgfqpoint{4.865424in}{1.406788in}}%
\pgfpathlineto{\pgfqpoint{4.867943in}{1.385167in}}%
\pgfpathlineto{\pgfqpoint{4.870462in}{1.377113in}}%
\pgfpathlineto{\pgfqpoint{4.872981in}{1.374263in}}%
\pgfpathlineto{\pgfqpoint{4.875499in}{1.373969in}}%
\pgfpathlineto{\pgfqpoint{4.878018in}{1.352619in}}%
\pgfpathlineto{\pgfqpoint{4.880537in}{1.359518in}}%
\pgfpathlineto{\pgfqpoint{4.885574in}{1.346976in}}%
\pgfpathlineto{\pgfqpoint{4.888093in}{1.360289in}}%
\pgfpathlineto{\pgfqpoint{4.890612in}{1.353230in}}%
\pgfpathlineto{\pgfqpoint{4.893131in}{1.341834in}}%
\pgfpathlineto{\pgfqpoint{4.895649in}{1.347852in}}%
\pgfpathlineto{\pgfqpoint{4.898168in}{1.371701in}}%
\pgfpathlineto{\pgfqpoint{4.900687in}{1.391063in}}%
\pgfpathlineto{\pgfqpoint{4.905724in}{1.376622in}}%
\pgfpathlineto{\pgfqpoint{4.908243in}{1.390901in}}%
\pgfpathlineto{\pgfqpoint{4.910762in}{1.396038in}}%
\pgfpathlineto{\pgfqpoint{4.913281in}{1.415190in}}%
\pgfpathlineto{\pgfqpoint{4.915799in}{1.433267in}}%
\pgfpathlineto{\pgfqpoint{4.918318in}{1.444096in}}%
\pgfpathlineto{\pgfqpoint{4.920837in}{1.425730in}}%
\pgfpathlineto{\pgfqpoint{4.923356in}{1.432747in}}%
\pgfpathlineto{\pgfqpoint{4.925874in}{1.426162in}}%
\pgfpathlineto{\pgfqpoint{4.928393in}{1.439293in}}%
\pgfpathlineto{\pgfqpoint{4.930912in}{1.443673in}}%
\pgfpathlineto{\pgfqpoint{4.933431in}{1.440757in}}%
\pgfpathlineto{\pgfqpoint{4.935949in}{1.387021in}}%
\pgfpathlineto{\pgfqpoint{4.938468in}{1.421545in}}%
\pgfpathlineto{\pgfqpoint{4.940987in}{1.422679in}}%
\pgfpathlineto{\pgfqpoint{4.943506in}{1.421718in}}%
\pgfpathlineto{\pgfqpoint{4.946024in}{1.418256in}}%
\pgfpathlineto{\pgfqpoint{4.948543in}{1.398531in}}%
\pgfpathlineto{\pgfqpoint{4.951062in}{1.408478in}}%
\pgfpathlineto{\pgfqpoint{4.953581in}{1.403808in}}%
\pgfpathlineto{\pgfqpoint{4.956099in}{1.403995in}}%
\pgfpathlineto{\pgfqpoint{4.958618in}{1.416639in}}%
\pgfpathlineto{\pgfqpoint{4.961137in}{1.415751in}}%
\pgfpathlineto{\pgfqpoint{4.963656in}{1.418637in}}%
\pgfpathlineto{\pgfqpoint{4.966174in}{1.388412in}}%
\pgfpathlineto{\pgfqpoint{4.968693in}{1.381213in}}%
\pgfpathlineto{\pgfqpoint{4.971212in}{1.387936in}}%
\pgfpathlineto{\pgfqpoint{4.973731in}{1.377377in}}%
\pgfpathlineto{\pgfqpoint{4.976249in}{1.379161in}}%
\pgfpathlineto{\pgfqpoint{4.978768in}{1.367014in}}%
\pgfpathlineto{\pgfqpoint{4.981287in}{1.373745in}}%
\pgfpathlineto{\pgfqpoint{4.983806in}{1.362359in}}%
\pgfpathlineto{\pgfqpoint{4.986324in}{1.360356in}}%
\pgfpathlineto{\pgfqpoint{4.988843in}{1.353826in}}%
\pgfpathlineto{\pgfqpoint{4.991362in}{1.355562in}}%
\pgfpathlineto{\pgfqpoint{4.993881in}{1.363509in}}%
\pgfpathlineto{\pgfqpoint{4.996399in}{1.386679in}}%
\pgfpathlineto{\pgfqpoint{4.998918in}{1.372652in}}%
\pgfpathlineto{\pgfqpoint{5.001437in}{1.363947in}}%
\pgfpathlineto{\pgfqpoint{5.003956in}{1.335876in}}%
\pgfpathlineto{\pgfqpoint{5.006474in}{1.345883in}}%
\pgfpathlineto{\pgfqpoint{5.008993in}{1.354445in}}%
\pgfpathlineto{\pgfqpoint{5.011512in}{1.357641in}}%
\pgfpathlineto{\pgfqpoint{5.014031in}{1.374653in}}%
\pgfpathlineto{\pgfqpoint{5.016549in}{1.368302in}}%
\pgfpathlineto{\pgfqpoint{5.019068in}{1.391112in}}%
\pgfpathlineto{\pgfqpoint{5.021587in}{1.388102in}}%
\pgfpathlineto{\pgfqpoint{5.024106in}{1.367247in}}%
\pgfpathlineto{\pgfqpoint{5.026624in}{1.358314in}}%
\pgfpathlineto{\pgfqpoint{5.029143in}{1.363294in}}%
\pgfpathlineto{\pgfqpoint{5.031662in}{1.377974in}}%
\pgfpathlineto{\pgfqpoint{5.034181in}{1.375385in}}%
\pgfpathlineto{\pgfqpoint{5.036699in}{1.379420in}}%
\pgfpathlineto{\pgfqpoint{5.039218in}{1.364640in}}%
\pgfpathlineto{\pgfqpoint{5.041737in}{1.375802in}}%
\pgfpathlineto{\pgfqpoint{5.044256in}{1.382740in}}%
\pgfpathlineto{\pgfqpoint{5.046774in}{1.395103in}}%
\pgfpathlineto{\pgfqpoint{5.049293in}{1.411343in}}%
\pgfpathlineto{\pgfqpoint{5.051812in}{1.388631in}}%
\pgfpathlineto{\pgfqpoint{5.054331in}{1.353458in}}%
\pgfpathlineto{\pgfqpoint{5.056849in}{1.352380in}}%
\pgfpathlineto{\pgfqpoint{5.059368in}{1.339518in}}%
\pgfpathlineto{\pgfqpoint{5.061887in}{1.339737in}}%
\pgfpathlineto{\pgfqpoint{5.064406in}{1.359351in}}%
\pgfpathlineto{\pgfqpoint{5.066924in}{1.340566in}}%
\pgfpathlineto{\pgfqpoint{5.069443in}{1.347839in}}%
\pgfpathlineto{\pgfqpoint{5.071962in}{1.374054in}}%
\pgfpathlineto{\pgfqpoint{5.074481in}{1.376130in}}%
\pgfpathlineto{\pgfqpoint{5.076999in}{1.384468in}}%
\pgfpathlineto{\pgfqpoint{5.079518in}{1.380830in}}%
\pgfpathlineto{\pgfqpoint{5.082037in}{1.385954in}}%
\pgfpathlineto{\pgfqpoint{5.084556in}{1.380120in}}%
\pgfpathlineto{\pgfqpoint{5.087074in}{1.370155in}}%
\pgfpathlineto{\pgfqpoint{5.089593in}{1.377099in}}%
\pgfpathlineto{\pgfqpoint{5.092112in}{1.378571in}}%
\pgfpathlineto{\pgfqpoint{5.094631in}{1.410145in}}%
\pgfpathlineto{\pgfqpoint{5.097149in}{1.400808in}}%
\pgfpathlineto{\pgfqpoint{5.099668in}{1.408344in}}%
\pgfpathlineto{\pgfqpoint{5.102187in}{1.408731in}}%
\pgfpathlineto{\pgfqpoint{5.104706in}{1.437176in}}%
\pgfpathlineto{\pgfqpoint{5.107224in}{1.430296in}}%
\pgfpathlineto{\pgfqpoint{5.109743in}{1.432321in}}%
\pgfpathlineto{\pgfqpoint{5.112262in}{1.459783in}}%
\pgfpathlineto{\pgfqpoint{5.114781in}{1.450822in}}%
\pgfpathlineto{\pgfqpoint{5.117299in}{1.466843in}}%
\pgfpathlineto{\pgfqpoint{5.119818in}{1.486197in}}%
\pgfpathlineto{\pgfqpoint{5.122337in}{1.443206in}}%
\pgfpathlineto{\pgfqpoint{5.124856in}{1.444525in}}%
\pgfpathlineto{\pgfqpoint{5.127374in}{1.419162in}}%
\pgfpathlineto{\pgfqpoint{5.129893in}{1.396839in}}%
\pgfpathlineto{\pgfqpoint{5.132412in}{1.430153in}}%
\pgfpathlineto{\pgfqpoint{5.134931in}{1.428091in}}%
\pgfpathlineto{\pgfqpoint{5.137449in}{1.425795in}}%
\pgfpathlineto{\pgfqpoint{5.139968in}{1.438654in}}%
\pgfpathlineto{\pgfqpoint{5.142487in}{1.445843in}}%
\pgfpathlineto{\pgfqpoint{5.145006in}{1.437423in}}%
\pgfpathlineto{\pgfqpoint{5.147524in}{1.448792in}}%
\pgfpathlineto{\pgfqpoint{5.150043in}{1.457734in}}%
\pgfpathlineto{\pgfqpoint{5.152562in}{1.459169in}}%
\pgfpathlineto{\pgfqpoint{5.155081in}{1.465398in}}%
\pgfpathlineto{\pgfqpoint{5.157599in}{1.458844in}}%
\pgfpathlineto{\pgfqpoint{5.160118in}{1.453195in}}%
\pgfpathlineto{\pgfqpoint{5.162637in}{1.428798in}}%
\pgfpathlineto{\pgfqpoint{5.165156in}{1.435966in}}%
\pgfpathlineto{\pgfqpoint{5.167674in}{1.426697in}}%
\pgfpathlineto{\pgfqpoint{5.170193in}{1.412347in}}%
\pgfpathlineto{\pgfqpoint{5.172712in}{1.430245in}}%
\pgfpathlineto{\pgfqpoint{5.175231in}{1.447386in}}%
\pgfpathlineto{\pgfqpoint{5.177749in}{1.461200in}}%
\pgfpathlineto{\pgfqpoint{5.180268in}{1.454187in}}%
\pgfpathlineto{\pgfqpoint{5.182787in}{1.478167in}}%
\pgfpathlineto{\pgfqpoint{5.185306in}{1.471607in}}%
\pgfpathlineto{\pgfqpoint{5.187824in}{1.482507in}}%
\pgfpathlineto{\pgfqpoint{5.190343in}{1.462254in}}%
\pgfpathlineto{\pgfqpoint{5.192862in}{1.428919in}}%
\pgfpathlineto{\pgfqpoint{5.195381in}{1.435425in}}%
\pgfpathlineto{\pgfqpoint{5.197899in}{1.443977in}}%
\pgfpathlineto{\pgfqpoint{5.200418in}{1.447535in}}%
\pgfpathlineto{\pgfqpoint{5.202937in}{1.465819in}}%
\pgfpathlineto{\pgfqpoint{5.205456in}{1.477238in}}%
\pgfpathlineto{\pgfqpoint{5.207974in}{1.474450in}}%
\pgfpathlineto{\pgfqpoint{5.210493in}{1.480349in}}%
\pgfpathlineto{\pgfqpoint{5.213012in}{1.503134in}}%
\pgfpathlineto{\pgfqpoint{5.215531in}{1.468098in}}%
\pgfpathlineto{\pgfqpoint{5.218049in}{1.475272in}}%
\pgfpathlineto{\pgfqpoint{5.220568in}{1.493858in}}%
\pgfpathlineto{\pgfqpoint{5.223087in}{1.467017in}}%
\pgfpathlineto{\pgfqpoint{5.225606in}{1.484121in}}%
\pgfpathlineto{\pgfqpoint{5.228124in}{1.494510in}}%
\pgfpathlineto{\pgfqpoint{5.230643in}{1.486230in}}%
\pgfpathlineto{\pgfqpoint{5.233162in}{1.485999in}}%
\pgfpathlineto{\pgfqpoint{5.235681in}{1.467055in}}%
\pgfpathlineto{\pgfqpoint{5.238199in}{1.475169in}}%
\pgfpathlineto{\pgfqpoint{5.240718in}{1.488012in}}%
\pgfpathlineto{\pgfqpoint{5.243237in}{1.516531in}}%
\pgfpathlineto{\pgfqpoint{5.245756in}{1.525932in}}%
\pgfpathlineto{\pgfqpoint{5.248274in}{1.524583in}}%
\pgfpathlineto{\pgfqpoint{5.250793in}{1.515196in}}%
\pgfpathlineto{\pgfqpoint{5.253312in}{1.544237in}}%
\pgfpathlineto{\pgfqpoint{5.255831in}{1.505060in}}%
\pgfpathlineto{\pgfqpoint{5.258349in}{1.531726in}}%
\pgfpathlineto{\pgfqpoint{5.260868in}{1.528993in}}%
\pgfpathlineto{\pgfqpoint{5.263387in}{1.510012in}}%
\pgfpathlineto{\pgfqpoint{5.265906in}{1.514761in}}%
\pgfpathlineto{\pgfqpoint{5.268424in}{1.493232in}}%
\pgfpathlineto{\pgfqpoint{5.270943in}{1.496562in}}%
\pgfpathlineto{\pgfqpoint{5.273462in}{1.513839in}}%
\pgfpathlineto{\pgfqpoint{5.275981in}{1.524494in}}%
\pgfpathlineto{\pgfqpoint{5.278499in}{1.544038in}}%
\pgfpathlineto{\pgfqpoint{5.281018in}{1.533763in}}%
\pgfpathlineto{\pgfqpoint{5.283537in}{1.525100in}}%
\pgfpathlineto{\pgfqpoint{5.286056in}{1.517963in}}%
\pgfpathlineto{\pgfqpoint{5.288574in}{1.520759in}}%
\pgfpathlineto{\pgfqpoint{5.291093in}{1.525707in}}%
\pgfpathlineto{\pgfqpoint{5.293612in}{1.540399in}}%
\pgfpathlineto{\pgfqpoint{5.296131in}{1.536667in}}%
\pgfpathlineto{\pgfqpoint{5.298649in}{1.550839in}}%
\pgfpathlineto{\pgfqpoint{5.301168in}{1.552028in}}%
\pgfpathlineto{\pgfqpoint{5.303687in}{1.564331in}}%
\pgfpathlineto{\pgfqpoint{5.306206in}{1.579390in}}%
\pgfpathlineto{\pgfqpoint{5.308724in}{1.573441in}}%
\pgfpathlineto{\pgfqpoint{5.311243in}{1.537837in}}%
\pgfpathlineto{\pgfqpoint{5.313762in}{1.536753in}}%
\pgfpathlineto{\pgfqpoint{5.316281in}{1.549011in}}%
\pgfpathlineto{\pgfqpoint{5.318799in}{1.570371in}}%
\pgfpathlineto{\pgfqpoint{5.321318in}{1.583976in}}%
\pgfpathlineto{\pgfqpoint{5.323837in}{1.565183in}}%
\pgfpathlineto{\pgfqpoint{5.326356in}{1.561816in}}%
\pgfpathlineto{\pgfqpoint{5.328874in}{1.565404in}}%
\pgfpathlineto{\pgfqpoint{5.331393in}{1.573611in}}%
\pgfpathlineto{\pgfqpoint{5.333912in}{1.569633in}}%
\pgfpathlineto{\pgfqpoint{5.336431in}{1.589512in}}%
\pgfpathlineto{\pgfqpoint{5.338949in}{1.593039in}}%
\pgfpathlineto{\pgfqpoint{5.341468in}{1.579903in}}%
\pgfpathlineto{\pgfqpoint{5.343987in}{1.564808in}}%
\pgfpathlineto{\pgfqpoint{5.346506in}{1.563013in}}%
\pgfpathlineto{\pgfqpoint{5.349024in}{1.566035in}}%
\pgfpathlineto{\pgfqpoint{5.351543in}{1.558146in}}%
\pgfpathlineto{\pgfqpoint{5.354062in}{1.540425in}}%
\pgfpathlineto{\pgfqpoint{5.356581in}{1.557576in}}%
\pgfpathlineto{\pgfqpoint{5.359099in}{1.587967in}}%
\pgfpathlineto{\pgfqpoint{5.361618in}{1.566119in}}%
\pgfpathlineto{\pgfqpoint{5.366656in}{1.604029in}}%
\pgfpathlineto{\pgfqpoint{5.369174in}{1.597963in}}%
\pgfpathlineto{\pgfqpoint{5.371693in}{1.571428in}}%
\pgfpathlineto{\pgfqpoint{5.374212in}{1.562661in}}%
\pgfpathlineto{\pgfqpoint{5.376731in}{1.585132in}}%
\pgfpathlineto{\pgfqpoint{5.379249in}{1.584926in}}%
\pgfpathlineto{\pgfqpoint{5.381768in}{1.595033in}}%
\pgfpathlineto{\pgfqpoint{5.384287in}{1.585065in}}%
\pgfpathlineto{\pgfqpoint{5.386806in}{1.582483in}}%
\pgfpathlineto{\pgfqpoint{5.389324in}{1.618732in}}%
\pgfpathlineto{\pgfqpoint{5.391843in}{1.599125in}}%
\pgfpathlineto{\pgfqpoint{5.394362in}{1.582420in}}%
\pgfpathlineto{\pgfqpoint{5.396881in}{1.591439in}}%
\pgfpathlineto{\pgfqpoint{5.399399in}{1.596054in}}%
\pgfpathlineto{\pgfqpoint{5.401918in}{1.601480in}}%
\pgfpathlineto{\pgfqpoint{5.404437in}{1.604317in}}%
\pgfpathlineto{\pgfqpoint{5.406956in}{1.610649in}}%
\pgfpathlineto{\pgfqpoint{5.409474in}{1.637364in}}%
\pgfpathlineto{\pgfqpoint{5.411993in}{1.629706in}}%
\pgfpathlineto{\pgfqpoint{5.414512in}{1.630104in}}%
\pgfpathlineto{\pgfqpoint{5.417031in}{1.662400in}}%
\pgfpathlineto{\pgfqpoint{5.419549in}{1.652930in}}%
\pgfpathlineto{\pgfqpoint{5.422068in}{1.654908in}}%
\pgfpathlineto{\pgfqpoint{5.424587in}{1.641528in}}%
\pgfpathlineto{\pgfqpoint{5.427106in}{1.643917in}}%
\pgfpathlineto{\pgfqpoint{5.429624in}{1.657788in}}%
\pgfpathlineto{\pgfqpoint{5.432143in}{1.650272in}}%
\pgfpathlineto{\pgfqpoint{5.434662in}{1.649650in}}%
\pgfpathlineto{\pgfqpoint{5.437181in}{1.642671in}}%
\pgfpathlineto{\pgfqpoint{5.439699in}{1.634303in}}%
\pgfpathlineto{\pgfqpoint{5.442218in}{1.636395in}}%
\pgfpathlineto{\pgfqpoint{5.444737in}{1.654713in}}%
\pgfpathlineto{\pgfqpoint{5.447256in}{1.700617in}}%
\pgfpathlineto{\pgfqpoint{5.449774in}{1.697075in}}%
\pgfpathlineto{\pgfqpoint{5.452293in}{1.681721in}}%
\pgfpathlineto{\pgfqpoint{5.454812in}{1.673022in}}%
\pgfpathlineto{\pgfqpoint{5.457331in}{1.695322in}}%
\pgfpathlineto{\pgfqpoint{5.459849in}{1.690253in}}%
\pgfpathlineto{\pgfqpoint{5.462368in}{1.711883in}}%
\pgfpathlineto{\pgfqpoint{5.464887in}{1.684562in}}%
\pgfpathlineto{\pgfqpoint{5.467406in}{1.705246in}}%
\pgfpathlineto{\pgfqpoint{5.469924in}{1.736074in}}%
\pgfpathlineto{\pgfqpoint{5.472443in}{1.742161in}}%
\pgfpathlineto{\pgfqpoint{5.474962in}{1.763277in}}%
\pgfpathlineto{\pgfqpoint{5.477481in}{1.768501in}}%
\pgfpathlineto{\pgfqpoint{5.479999in}{1.766202in}}%
\pgfpathlineto{\pgfqpoint{5.482518in}{1.790274in}}%
\pgfpathlineto{\pgfqpoint{5.485037in}{1.805423in}}%
\pgfpathlineto{\pgfqpoint{5.487556in}{1.793725in}}%
\pgfpathlineto{\pgfqpoint{5.490074in}{1.797813in}}%
\pgfpathlineto{\pgfqpoint{5.492593in}{1.794213in}}%
\pgfpathlineto{\pgfqpoint{5.492593in}{1.794213in}}%
\pgfusepath{stroke}%
\end{pgfscope}%
\begin{pgfscope}%
\pgfpathrectangle{\pgfqpoint{0.455093in}{0.401080in}}{\pgfqpoint{5.037500in}{4.004000in}}%
\pgfusepath{clip}%
\pgfsetrectcap%
\pgfsetroundjoin%
\pgfsetlinewidth{0.501875pt}%
\definecolor{currentstroke}{rgb}{0.690196,0.188235,0.376471}%
\pgfsetstrokecolor{currentstroke}%
\pgfsetstrokeopacity{0.800000}%
\pgfsetdash{}{0pt}%
\pgfpathmoveto{\pgfqpoint{0.455093in}{1.021044in}}%
\pgfpathlineto{\pgfqpoint{0.457612in}{1.001797in}}%
\pgfpathlineto{\pgfqpoint{0.460131in}{0.979241in}}%
\pgfpathlineto{\pgfqpoint{0.462649in}{0.976843in}}%
\pgfpathlineto{\pgfqpoint{0.465168in}{0.972724in}}%
\pgfpathlineto{\pgfqpoint{0.467687in}{0.987082in}}%
\pgfpathlineto{\pgfqpoint{0.470206in}{0.990064in}}%
\pgfpathlineto{\pgfqpoint{0.472724in}{0.988488in}}%
\pgfpathlineto{\pgfqpoint{0.475243in}{1.004311in}}%
\pgfpathlineto{\pgfqpoint{0.477762in}{1.005447in}}%
\pgfpathlineto{\pgfqpoint{0.480281in}{0.989507in}}%
\pgfpathlineto{\pgfqpoint{0.482799in}{0.977704in}}%
\pgfpathlineto{\pgfqpoint{0.485318in}{0.996948in}}%
\pgfpathlineto{\pgfqpoint{0.487837in}{0.992552in}}%
\pgfpathlineto{\pgfqpoint{0.490356in}{0.980412in}}%
\pgfpathlineto{\pgfqpoint{0.492874in}{0.986508in}}%
\pgfpathlineto{\pgfqpoint{0.495393in}{0.970079in}}%
\pgfpathlineto{\pgfqpoint{0.497912in}{0.987039in}}%
\pgfpathlineto{\pgfqpoint{0.500431in}{0.975455in}}%
\pgfpathlineto{\pgfqpoint{0.502949in}{0.976567in}}%
\pgfpathlineto{\pgfqpoint{0.505468in}{0.987365in}}%
\pgfpathlineto{\pgfqpoint{0.507987in}{0.992201in}}%
\pgfpathlineto{\pgfqpoint{0.510506in}{0.996657in}}%
\pgfpathlineto{\pgfqpoint{0.513024in}{0.983588in}}%
\pgfpathlineto{\pgfqpoint{0.515543in}{0.979031in}}%
\pgfpathlineto{\pgfqpoint{0.518062in}{0.991862in}}%
\pgfpathlineto{\pgfqpoint{0.520581in}{0.988363in}}%
\pgfpathlineto{\pgfqpoint{0.523099in}{0.985596in}}%
\pgfpathlineto{\pgfqpoint{0.525618in}{0.994539in}}%
\pgfpathlineto{\pgfqpoint{0.528137in}{1.013261in}}%
\pgfpathlineto{\pgfqpoint{0.530656in}{1.024878in}}%
\pgfpathlineto{\pgfqpoint{0.533174in}{1.053605in}}%
\pgfpathlineto{\pgfqpoint{0.535693in}{1.062264in}}%
\pgfpathlineto{\pgfqpoint{0.538212in}{1.037922in}}%
\pgfpathlineto{\pgfqpoint{0.540731in}{1.037771in}}%
\pgfpathlineto{\pgfqpoint{0.543249in}{1.046227in}}%
\pgfpathlineto{\pgfqpoint{0.545768in}{1.058863in}}%
\pgfpathlineto{\pgfqpoint{0.548287in}{1.061200in}}%
\pgfpathlineto{\pgfqpoint{0.550806in}{1.063178in}}%
\pgfpathlineto{\pgfqpoint{0.553324in}{1.052268in}}%
\pgfpathlineto{\pgfqpoint{0.555843in}{1.053545in}}%
\pgfpathlineto{\pgfqpoint{0.558362in}{1.048193in}}%
\pgfpathlineto{\pgfqpoint{0.560881in}{1.050023in}}%
\pgfpathlineto{\pgfqpoint{0.563399in}{1.047861in}}%
\pgfpathlineto{\pgfqpoint{0.565918in}{1.037892in}}%
\pgfpathlineto{\pgfqpoint{0.568437in}{1.068166in}}%
\pgfpathlineto{\pgfqpoint{0.570956in}{1.086533in}}%
\pgfpathlineto{\pgfqpoint{0.573474in}{1.096894in}}%
\pgfpathlineto{\pgfqpoint{0.578512in}{1.085042in}}%
\pgfpathlineto{\pgfqpoint{0.581031in}{1.078409in}}%
\pgfpathlineto{\pgfqpoint{0.583549in}{1.069805in}}%
\pgfpathlineto{\pgfqpoint{0.586068in}{1.095301in}}%
\pgfpathlineto{\pgfqpoint{0.588587in}{1.093146in}}%
\pgfpathlineto{\pgfqpoint{0.591106in}{1.081307in}}%
\pgfpathlineto{\pgfqpoint{0.593624in}{1.073759in}}%
\pgfpathlineto{\pgfqpoint{0.596143in}{1.074940in}}%
\pgfpathlineto{\pgfqpoint{0.598662in}{1.065462in}}%
\pgfpathlineto{\pgfqpoint{0.601181in}{1.054780in}}%
\pgfpathlineto{\pgfqpoint{0.603699in}{1.059163in}}%
\pgfpathlineto{\pgfqpoint{0.606218in}{1.059628in}}%
\pgfpathlineto{\pgfqpoint{0.608737in}{1.047209in}}%
\pgfpathlineto{\pgfqpoint{0.611256in}{1.055919in}}%
\pgfpathlineto{\pgfqpoint{0.613774in}{1.072510in}}%
\pgfpathlineto{\pgfqpoint{0.616293in}{1.037092in}}%
\pgfpathlineto{\pgfqpoint{0.618812in}{1.031922in}}%
\pgfpathlineto{\pgfqpoint{0.621331in}{1.042360in}}%
\pgfpathlineto{\pgfqpoint{0.623849in}{1.030942in}}%
\pgfpathlineto{\pgfqpoint{0.626368in}{1.008221in}}%
\pgfpathlineto{\pgfqpoint{0.628887in}{1.009161in}}%
\pgfpathlineto{\pgfqpoint{0.631406in}{1.006977in}}%
\pgfpathlineto{\pgfqpoint{0.633924in}{1.019756in}}%
\pgfpathlineto{\pgfqpoint{0.636443in}{1.002125in}}%
\pgfpathlineto{\pgfqpoint{0.638962in}{0.992901in}}%
\pgfpathlineto{\pgfqpoint{0.641481in}{0.977642in}}%
\pgfpathlineto{\pgfqpoint{0.643999in}{0.971786in}}%
\pgfpathlineto{\pgfqpoint{0.646518in}{0.985017in}}%
\pgfpathlineto{\pgfqpoint{0.649037in}{0.979551in}}%
\pgfpathlineto{\pgfqpoint{0.651556in}{0.980817in}}%
\pgfpathlineto{\pgfqpoint{0.654074in}{1.002612in}}%
\pgfpathlineto{\pgfqpoint{0.656593in}{1.013286in}}%
\pgfpathlineto{\pgfqpoint{0.659112in}{0.999189in}}%
\pgfpathlineto{\pgfqpoint{0.661631in}{0.976275in}}%
\pgfpathlineto{\pgfqpoint{0.664149in}{0.976530in}}%
\pgfpathlineto{\pgfqpoint{0.666668in}{0.989123in}}%
\pgfpathlineto{\pgfqpoint{0.669187in}{0.985977in}}%
\pgfpathlineto{\pgfqpoint{0.671706in}{0.977719in}}%
\pgfpathlineto{\pgfqpoint{0.674224in}{1.008980in}}%
\pgfpathlineto{\pgfqpoint{0.679262in}{1.035317in}}%
\pgfpathlineto{\pgfqpoint{0.681781in}{1.044445in}}%
\pgfpathlineto{\pgfqpoint{0.684299in}{1.072267in}}%
\pgfpathlineto{\pgfqpoint{0.686818in}{1.078963in}}%
\pgfpathlineto{\pgfqpoint{0.689337in}{1.047660in}}%
\pgfpathlineto{\pgfqpoint{0.691856in}{1.064538in}}%
\pgfpathlineto{\pgfqpoint{0.694374in}{1.066591in}}%
\pgfpathlineto{\pgfqpoint{0.696893in}{1.067822in}}%
\pgfpathlineto{\pgfqpoint{0.699412in}{1.072019in}}%
\pgfpathlineto{\pgfqpoint{0.701931in}{1.084581in}}%
\pgfpathlineto{\pgfqpoint{0.704449in}{1.073343in}}%
\pgfpathlineto{\pgfqpoint{0.706968in}{1.090396in}}%
\pgfpathlineto{\pgfqpoint{0.709487in}{1.089822in}}%
\pgfpathlineto{\pgfqpoint{0.712006in}{1.114684in}}%
\pgfpathlineto{\pgfqpoint{0.717043in}{1.126233in}}%
\pgfpathlineto{\pgfqpoint{0.719562in}{1.127391in}}%
\pgfpathlineto{\pgfqpoint{0.722081in}{1.119731in}}%
\pgfpathlineto{\pgfqpoint{0.724599in}{1.136906in}}%
\pgfpathlineto{\pgfqpoint{0.727118in}{1.166085in}}%
\pgfpathlineto{\pgfqpoint{0.729637in}{1.182521in}}%
\pgfpathlineto{\pgfqpoint{0.732156in}{1.163858in}}%
\pgfpathlineto{\pgfqpoint{0.734674in}{1.174935in}}%
\pgfpathlineto{\pgfqpoint{0.737193in}{1.177007in}}%
\pgfpathlineto{\pgfqpoint{0.739712in}{1.176630in}}%
\pgfpathlineto{\pgfqpoint{0.742231in}{1.210656in}}%
\pgfpathlineto{\pgfqpoint{0.744749in}{1.222240in}}%
\pgfpathlineto{\pgfqpoint{0.747268in}{1.219376in}}%
\pgfpathlineto{\pgfqpoint{0.752306in}{1.235560in}}%
\pgfpathlineto{\pgfqpoint{0.754824in}{1.215778in}}%
\pgfpathlineto{\pgfqpoint{0.757343in}{1.197685in}}%
\pgfpathlineto{\pgfqpoint{0.759862in}{1.191312in}}%
\pgfpathlineto{\pgfqpoint{0.762381in}{1.195402in}}%
\pgfpathlineto{\pgfqpoint{0.764899in}{1.232984in}}%
\pgfpathlineto{\pgfqpoint{0.767418in}{1.214827in}}%
\pgfpathlineto{\pgfqpoint{0.769937in}{1.214463in}}%
\pgfpathlineto{\pgfqpoint{0.772456in}{1.226217in}}%
\pgfpathlineto{\pgfqpoint{0.774974in}{1.230477in}}%
\pgfpathlineto{\pgfqpoint{0.777493in}{1.222342in}}%
\pgfpathlineto{\pgfqpoint{0.780012in}{1.223883in}}%
\pgfpathlineto{\pgfqpoint{0.782531in}{1.249438in}}%
\pgfpathlineto{\pgfqpoint{0.785049in}{1.238227in}}%
\pgfpathlineto{\pgfqpoint{0.787568in}{1.245489in}}%
\pgfpathlineto{\pgfqpoint{0.790087in}{1.260376in}}%
\pgfpathlineto{\pgfqpoint{0.792606in}{1.272573in}}%
\pgfpathlineto{\pgfqpoint{0.795124in}{1.254832in}}%
\pgfpathlineto{\pgfqpoint{0.797643in}{1.256229in}}%
\pgfpathlineto{\pgfqpoint{0.800162in}{1.271498in}}%
\pgfpathlineto{\pgfqpoint{0.802681in}{1.275552in}}%
\pgfpathlineto{\pgfqpoint{0.805199in}{1.254133in}}%
\pgfpathlineto{\pgfqpoint{0.807718in}{1.264052in}}%
\pgfpathlineto{\pgfqpoint{0.810237in}{1.256946in}}%
\pgfpathlineto{\pgfqpoint{0.812756in}{1.244694in}}%
\pgfpathlineto{\pgfqpoint{0.815274in}{1.262888in}}%
\pgfpathlineto{\pgfqpoint{0.817793in}{1.268182in}}%
\pgfpathlineto{\pgfqpoint{0.820312in}{1.269480in}}%
\pgfpathlineto{\pgfqpoint{0.822831in}{1.263251in}}%
\pgfpathlineto{\pgfqpoint{0.825349in}{1.259319in}}%
\pgfpathlineto{\pgfqpoint{0.827868in}{1.263969in}}%
\pgfpathlineto{\pgfqpoint{0.830387in}{1.258646in}}%
\pgfpathlineto{\pgfqpoint{0.832906in}{1.266628in}}%
\pgfpathlineto{\pgfqpoint{0.835424in}{1.269034in}}%
\pgfpathlineto{\pgfqpoint{0.837943in}{1.269952in}}%
\pgfpathlineto{\pgfqpoint{0.840462in}{1.269698in}}%
\pgfpathlineto{\pgfqpoint{0.842981in}{1.247423in}}%
\pgfpathlineto{\pgfqpoint{0.845499in}{1.270742in}}%
\pgfpathlineto{\pgfqpoint{0.848018in}{1.279705in}}%
\pgfpathlineto{\pgfqpoint{0.850537in}{1.262231in}}%
\pgfpathlineto{\pgfqpoint{0.853056in}{1.271862in}}%
\pgfpathlineto{\pgfqpoint{0.855574in}{1.244181in}}%
\pgfpathlineto{\pgfqpoint{0.858093in}{1.254069in}}%
\pgfpathlineto{\pgfqpoint{0.860612in}{1.266471in}}%
\pgfpathlineto{\pgfqpoint{0.863131in}{1.300916in}}%
\pgfpathlineto{\pgfqpoint{0.865649in}{1.329607in}}%
\pgfpathlineto{\pgfqpoint{0.868168in}{1.353557in}}%
\pgfpathlineto{\pgfqpoint{0.870687in}{1.357230in}}%
\pgfpathlineto{\pgfqpoint{0.873206in}{1.327668in}}%
\pgfpathlineto{\pgfqpoint{0.875724in}{1.348461in}}%
\pgfpathlineto{\pgfqpoint{0.878243in}{1.349004in}}%
\pgfpathlineto{\pgfqpoint{0.880762in}{1.366886in}}%
\pgfpathlineto{\pgfqpoint{0.883281in}{1.377054in}}%
\pgfpathlineto{\pgfqpoint{0.885799in}{1.390342in}}%
\pgfpathlineto{\pgfqpoint{0.888318in}{1.394653in}}%
\pgfpathlineto{\pgfqpoint{0.890837in}{1.424983in}}%
\pgfpathlineto{\pgfqpoint{0.893356in}{1.418028in}}%
\pgfpathlineto{\pgfqpoint{0.895874in}{1.418687in}}%
\pgfpathlineto{\pgfqpoint{0.898393in}{1.447070in}}%
\pgfpathlineto{\pgfqpoint{0.900912in}{1.473907in}}%
\pgfpathlineto{\pgfqpoint{0.903431in}{1.473450in}}%
\pgfpathlineto{\pgfqpoint{0.905949in}{1.508491in}}%
\pgfpathlineto{\pgfqpoint{0.908468in}{1.504492in}}%
\pgfpathlineto{\pgfqpoint{0.910987in}{1.498379in}}%
\pgfpathlineto{\pgfqpoint{0.913506in}{1.485750in}}%
\pgfpathlineto{\pgfqpoint{0.916024in}{1.479267in}}%
\pgfpathlineto{\pgfqpoint{0.918543in}{1.477377in}}%
\pgfpathlineto{\pgfqpoint{0.921062in}{1.494503in}}%
\pgfpathlineto{\pgfqpoint{0.923581in}{1.514220in}}%
\pgfpathlineto{\pgfqpoint{0.926099in}{1.483758in}}%
\pgfpathlineto{\pgfqpoint{0.928618in}{1.486541in}}%
\pgfpathlineto{\pgfqpoint{0.931137in}{1.462612in}}%
\pgfpathlineto{\pgfqpoint{0.933656in}{1.487243in}}%
\pgfpathlineto{\pgfqpoint{0.936174in}{1.487093in}}%
\pgfpathlineto{\pgfqpoint{0.938693in}{1.512884in}}%
\pgfpathlineto{\pgfqpoint{0.941212in}{1.506249in}}%
\pgfpathlineto{\pgfqpoint{0.943731in}{1.477985in}}%
\pgfpathlineto{\pgfqpoint{0.946249in}{1.474916in}}%
\pgfpathlineto{\pgfqpoint{0.948768in}{1.473739in}}%
\pgfpathlineto{\pgfqpoint{0.951287in}{1.452814in}}%
\pgfpathlineto{\pgfqpoint{0.953806in}{1.438915in}}%
\pgfpathlineto{\pgfqpoint{0.956324in}{1.421640in}}%
\pgfpathlineto{\pgfqpoint{0.958843in}{1.430133in}}%
\pgfpathlineto{\pgfqpoint{0.961362in}{1.409714in}}%
\pgfpathlineto{\pgfqpoint{0.963881in}{1.412960in}}%
\pgfpathlineto{\pgfqpoint{0.966399in}{1.420886in}}%
\pgfpathlineto{\pgfqpoint{0.968918in}{1.447556in}}%
\pgfpathlineto{\pgfqpoint{0.971437in}{1.472170in}}%
\pgfpathlineto{\pgfqpoint{0.973956in}{1.444792in}}%
\pgfpathlineto{\pgfqpoint{0.976474in}{1.440518in}}%
\pgfpathlineto{\pgfqpoint{0.978993in}{1.438276in}}%
\pgfpathlineto{\pgfqpoint{0.981512in}{1.452877in}}%
\pgfpathlineto{\pgfqpoint{0.984031in}{1.432841in}}%
\pgfpathlineto{\pgfqpoint{0.986549in}{1.445755in}}%
\pgfpathlineto{\pgfqpoint{0.989068in}{1.459776in}}%
\pgfpathlineto{\pgfqpoint{0.991587in}{1.445793in}}%
\pgfpathlineto{\pgfqpoint{0.994106in}{1.464772in}}%
\pgfpathlineto{\pgfqpoint{0.996624in}{1.459267in}}%
\pgfpathlineto{\pgfqpoint{0.999143in}{1.465096in}}%
\pgfpathlineto{\pgfqpoint{1.001662in}{1.439707in}}%
\pgfpathlineto{\pgfqpoint{1.004181in}{1.455032in}}%
\pgfpathlineto{\pgfqpoint{1.006699in}{1.450904in}}%
\pgfpathlineto{\pgfqpoint{1.009218in}{1.462937in}}%
\pgfpathlineto{\pgfqpoint{1.011737in}{1.442613in}}%
\pgfpathlineto{\pgfqpoint{1.014256in}{1.449329in}}%
\pgfpathlineto{\pgfqpoint{1.016774in}{1.466117in}}%
\pgfpathlineto{\pgfqpoint{1.019293in}{1.475719in}}%
\pgfpathlineto{\pgfqpoint{1.021812in}{1.454588in}}%
\pgfpathlineto{\pgfqpoint{1.024331in}{1.435115in}}%
\pgfpathlineto{\pgfqpoint{1.026849in}{1.444309in}}%
\pgfpathlineto{\pgfqpoint{1.029368in}{1.440982in}}%
\pgfpathlineto{\pgfqpoint{1.031887in}{1.424413in}}%
\pgfpathlineto{\pgfqpoint{1.034406in}{1.440267in}}%
\pgfpathlineto{\pgfqpoint{1.036924in}{1.431296in}}%
\pgfpathlineto{\pgfqpoint{1.039443in}{1.460329in}}%
\pgfpathlineto{\pgfqpoint{1.041962in}{1.469005in}}%
\pgfpathlineto{\pgfqpoint{1.044481in}{1.496928in}}%
\pgfpathlineto{\pgfqpoint{1.046999in}{1.502715in}}%
\pgfpathlineto{\pgfqpoint{1.049518in}{1.526461in}}%
\pgfpathlineto{\pgfqpoint{1.052037in}{1.526285in}}%
\pgfpathlineto{\pgfqpoint{1.054556in}{1.532154in}}%
\pgfpathlineto{\pgfqpoint{1.057074in}{1.574579in}}%
\pgfpathlineto{\pgfqpoint{1.059593in}{1.557730in}}%
\pgfpathlineto{\pgfqpoint{1.062112in}{1.563480in}}%
\pgfpathlineto{\pgfqpoint{1.064631in}{1.566894in}}%
\pgfpathlineto{\pgfqpoint{1.067149in}{1.553629in}}%
\pgfpathlineto{\pgfqpoint{1.069668in}{1.553367in}}%
\pgfpathlineto{\pgfqpoint{1.072187in}{1.564568in}}%
\pgfpathlineto{\pgfqpoint{1.074706in}{1.568047in}}%
\pgfpathlineto{\pgfqpoint{1.077224in}{1.584925in}}%
\pgfpathlineto{\pgfqpoint{1.079743in}{1.551913in}}%
\pgfpathlineto{\pgfqpoint{1.082262in}{1.535683in}}%
\pgfpathlineto{\pgfqpoint{1.084781in}{1.536640in}}%
\pgfpathlineto{\pgfqpoint{1.087299in}{1.549238in}}%
\pgfpathlineto{\pgfqpoint{1.089818in}{1.569416in}}%
\pgfpathlineto{\pgfqpoint{1.092337in}{1.555513in}}%
\pgfpathlineto{\pgfqpoint{1.094856in}{1.558626in}}%
\pgfpathlineto{\pgfqpoint{1.097374in}{1.548715in}}%
\pgfpathlineto{\pgfqpoint{1.099893in}{1.543105in}}%
\pgfpathlineto{\pgfqpoint{1.102412in}{1.578205in}}%
\pgfpathlineto{\pgfqpoint{1.104931in}{1.572445in}}%
\pgfpathlineto{\pgfqpoint{1.107449in}{1.574004in}}%
\pgfpathlineto{\pgfqpoint{1.109968in}{1.576142in}}%
\pgfpathlineto{\pgfqpoint{1.112487in}{1.579755in}}%
\pgfpathlineto{\pgfqpoint{1.115006in}{1.577499in}}%
\pgfpathlineto{\pgfqpoint{1.117524in}{1.548334in}}%
\pgfpathlineto{\pgfqpoint{1.120043in}{1.546947in}}%
\pgfpathlineto{\pgfqpoint{1.122562in}{1.533551in}}%
\pgfpathlineto{\pgfqpoint{1.125081in}{1.556108in}}%
\pgfpathlineto{\pgfqpoint{1.127599in}{1.569744in}}%
\pgfpathlineto{\pgfqpoint{1.130118in}{1.578675in}}%
\pgfpathlineto{\pgfqpoint{1.135156in}{1.579030in}}%
\pgfpathlineto{\pgfqpoint{1.137674in}{1.545498in}}%
\pgfpathlineto{\pgfqpoint{1.140193in}{1.557872in}}%
\pgfpathlineto{\pgfqpoint{1.142712in}{1.557542in}}%
\pgfpathlineto{\pgfqpoint{1.145231in}{1.589308in}}%
\pgfpathlineto{\pgfqpoint{1.147749in}{1.597690in}}%
\pgfpathlineto{\pgfqpoint{1.150268in}{1.594954in}}%
\pgfpathlineto{\pgfqpoint{1.152787in}{1.614622in}}%
\pgfpathlineto{\pgfqpoint{1.155306in}{1.595271in}}%
\pgfpathlineto{\pgfqpoint{1.157824in}{1.639737in}}%
\pgfpathlineto{\pgfqpoint{1.160343in}{1.625354in}}%
\pgfpathlineto{\pgfqpoint{1.162862in}{1.631762in}}%
\pgfpathlineto{\pgfqpoint{1.165381in}{1.643793in}}%
\pgfpathlineto{\pgfqpoint{1.167899in}{1.633910in}}%
\pgfpathlineto{\pgfqpoint{1.170418in}{1.636790in}}%
\pgfpathlineto{\pgfqpoint{1.172937in}{1.653198in}}%
\pgfpathlineto{\pgfqpoint{1.175456in}{1.678132in}}%
\pgfpathlineto{\pgfqpoint{1.177974in}{1.668631in}}%
\pgfpathlineto{\pgfqpoint{1.180493in}{1.661859in}}%
\pgfpathlineto{\pgfqpoint{1.183012in}{1.660542in}}%
\pgfpathlineto{\pgfqpoint{1.185531in}{1.682762in}}%
\pgfpathlineto{\pgfqpoint{1.188049in}{1.669885in}}%
\pgfpathlineto{\pgfqpoint{1.190568in}{1.681767in}}%
\pgfpathlineto{\pgfqpoint{1.193087in}{1.701875in}}%
\pgfpathlineto{\pgfqpoint{1.195606in}{1.690935in}}%
\pgfpathlineto{\pgfqpoint{1.198124in}{1.663277in}}%
\pgfpathlineto{\pgfqpoint{1.200643in}{1.649009in}}%
\pgfpathlineto{\pgfqpoint{1.203162in}{1.653011in}}%
\pgfpathlineto{\pgfqpoint{1.205681in}{1.660617in}}%
\pgfpathlineto{\pgfqpoint{1.208199in}{1.680381in}}%
\pgfpathlineto{\pgfqpoint{1.210718in}{1.673612in}}%
\pgfpathlineto{\pgfqpoint{1.213237in}{1.660688in}}%
\pgfpathlineto{\pgfqpoint{1.215756in}{1.624665in}}%
\pgfpathlineto{\pgfqpoint{1.218274in}{1.619873in}}%
\pgfpathlineto{\pgfqpoint{1.220793in}{1.630818in}}%
\pgfpathlineto{\pgfqpoint{1.223312in}{1.637631in}}%
\pgfpathlineto{\pgfqpoint{1.225831in}{1.648280in}}%
\pgfpathlineto{\pgfqpoint{1.228349in}{1.660205in}}%
\pgfpathlineto{\pgfqpoint{1.230868in}{1.638474in}}%
\pgfpathlineto{\pgfqpoint{1.233387in}{1.640432in}}%
\pgfpathlineto{\pgfqpoint{1.235906in}{1.646928in}}%
\pgfpathlineto{\pgfqpoint{1.238424in}{1.624627in}}%
\pgfpathlineto{\pgfqpoint{1.240943in}{1.633199in}}%
\pgfpathlineto{\pgfqpoint{1.243462in}{1.623804in}}%
\pgfpathlineto{\pgfqpoint{1.245981in}{1.647444in}}%
\pgfpathlineto{\pgfqpoint{1.248499in}{1.663657in}}%
\pgfpathlineto{\pgfqpoint{1.251018in}{1.673473in}}%
\pgfpathlineto{\pgfqpoint{1.253537in}{1.690406in}}%
\pgfpathlineto{\pgfqpoint{1.256056in}{1.718689in}}%
\pgfpathlineto{\pgfqpoint{1.258574in}{1.716359in}}%
\pgfpathlineto{\pgfqpoint{1.261093in}{1.712839in}}%
\pgfpathlineto{\pgfqpoint{1.263612in}{1.743236in}}%
\pgfpathlineto{\pgfqpoint{1.266131in}{1.728896in}}%
\pgfpathlineto{\pgfqpoint{1.268649in}{1.723228in}}%
\pgfpathlineto{\pgfqpoint{1.271168in}{1.696594in}}%
\pgfpathlineto{\pgfqpoint{1.273687in}{1.696287in}}%
\pgfpathlineto{\pgfqpoint{1.276206in}{1.704249in}}%
\pgfpathlineto{\pgfqpoint{1.278724in}{1.747100in}}%
\pgfpathlineto{\pgfqpoint{1.281243in}{1.724267in}}%
\pgfpathlineto{\pgfqpoint{1.283762in}{1.718014in}}%
\pgfpathlineto{\pgfqpoint{1.286281in}{1.733800in}}%
\pgfpathlineto{\pgfqpoint{1.288799in}{1.747000in}}%
\pgfpathlineto{\pgfqpoint{1.291318in}{1.746131in}}%
\pgfpathlineto{\pgfqpoint{1.293837in}{1.748954in}}%
\pgfpathlineto{\pgfqpoint{1.296356in}{1.769299in}}%
\pgfpathlineto{\pgfqpoint{1.298874in}{1.753202in}}%
\pgfpathlineto{\pgfqpoint{1.301393in}{1.748634in}}%
\pgfpathlineto{\pgfqpoint{1.303912in}{1.736145in}}%
\pgfpathlineto{\pgfqpoint{1.306431in}{1.745544in}}%
\pgfpathlineto{\pgfqpoint{1.308949in}{1.757729in}}%
\pgfpathlineto{\pgfqpoint{1.311468in}{1.776241in}}%
\pgfpathlineto{\pgfqpoint{1.313987in}{1.782768in}}%
\pgfpathlineto{\pgfqpoint{1.316506in}{1.784794in}}%
\pgfpathlineto{\pgfqpoint{1.319024in}{1.770124in}}%
\pgfpathlineto{\pgfqpoint{1.321543in}{1.792246in}}%
\pgfpathlineto{\pgfqpoint{1.324062in}{1.817574in}}%
\pgfpathlineto{\pgfqpoint{1.326581in}{1.807437in}}%
\pgfpathlineto{\pgfqpoint{1.329099in}{1.796747in}}%
\pgfpathlineto{\pgfqpoint{1.331618in}{1.822063in}}%
\pgfpathlineto{\pgfqpoint{1.334137in}{1.842160in}}%
\pgfpathlineto{\pgfqpoint{1.336656in}{1.834629in}}%
\pgfpathlineto{\pgfqpoint{1.339174in}{1.838425in}}%
\pgfpathlineto{\pgfqpoint{1.341693in}{1.840319in}}%
\pgfpathlineto{\pgfqpoint{1.344212in}{1.833687in}}%
\pgfpathlineto{\pgfqpoint{1.346731in}{1.811477in}}%
\pgfpathlineto{\pgfqpoint{1.349249in}{1.806161in}}%
\pgfpathlineto{\pgfqpoint{1.351768in}{1.803373in}}%
\pgfpathlineto{\pgfqpoint{1.354287in}{1.813397in}}%
\pgfpathlineto{\pgfqpoint{1.356806in}{1.812754in}}%
\pgfpathlineto{\pgfqpoint{1.359324in}{1.814148in}}%
\pgfpathlineto{\pgfqpoint{1.361843in}{1.799806in}}%
\pgfpathlineto{\pgfqpoint{1.364362in}{1.814379in}}%
\pgfpathlineto{\pgfqpoint{1.366881in}{1.807307in}}%
\pgfpathlineto{\pgfqpoint{1.369399in}{1.788089in}}%
\pgfpathlineto{\pgfqpoint{1.371918in}{1.790965in}}%
\pgfpathlineto{\pgfqpoint{1.374437in}{1.743197in}}%
\pgfpathlineto{\pgfqpoint{1.376956in}{1.753102in}}%
\pgfpathlineto{\pgfqpoint{1.379474in}{1.770733in}}%
\pgfpathlineto{\pgfqpoint{1.381993in}{1.785343in}}%
\pgfpathlineto{\pgfqpoint{1.384512in}{1.760007in}}%
\pgfpathlineto{\pgfqpoint{1.387031in}{1.755867in}}%
\pgfpathlineto{\pgfqpoint{1.389549in}{1.752301in}}%
\pgfpathlineto{\pgfqpoint{1.392068in}{1.752188in}}%
\pgfpathlineto{\pgfqpoint{1.394587in}{1.734749in}}%
\pgfpathlineto{\pgfqpoint{1.397106in}{1.725887in}}%
\pgfpathlineto{\pgfqpoint{1.399624in}{1.744620in}}%
\pgfpathlineto{\pgfqpoint{1.402143in}{1.739671in}}%
\pgfpathlineto{\pgfqpoint{1.404662in}{1.711549in}}%
\pgfpathlineto{\pgfqpoint{1.407181in}{1.698028in}}%
\pgfpathlineto{\pgfqpoint{1.409699in}{1.677723in}}%
\pgfpathlineto{\pgfqpoint{1.412218in}{1.698266in}}%
\pgfpathlineto{\pgfqpoint{1.414737in}{1.675084in}}%
\pgfpathlineto{\pgfqpoint{1.419774in}{1.651175in}}%
\pgfpathlineto{\pgfqpoint{1.422293in}{1.678685in}}%
\pgfpathlineto{\pgfqpoint{1.427331in}{1.711342in}}%
\pgfpathlineto{\pgfqpoint{1.429849in}{1.735728in}}%
\pgfpathlineto{\pgfqpoint{1.432368in}{1.742959in}}%
\pgfpathlineto{\pgfqpoint{1.434887in}{1.764861in}}%
\pgfpathlineto{\pgfqpoint{1.437406in}{1.747899in}}%
\pgfpathlineto{\pgfqpoint{1.439924in}{1.762785in}}%
\pgfpathlineto{\pgfqpoint{1.442443in}{1.778523in}}%
\pgfpathlineto{\pgfqpoint{1.444962in}{1.779612in}}%
\pgfpathlineto{\pgfqpoint{1.447481in}{1.759621in}}%
\pgfpathlineto{\pgfqpoint{1.449999in}{1.751739in}}%
\pgfpathlineto{\pgfqpoint{1.452518in}{1.763322in}}%
\pgfpathlineto{\pgfqpoint{1.455037in}{1.765292in}}%
\pgfpathlineto{\pgfqpoint{1.460074in}{1.797943in}}%
\pgfpathlineto{\pgfqpoint{1.462593in}{1.828119in}}%
\pgfpathlineto{\pgfqpoint{1.465112in}{1.866906in}}%
\pgfpathlineto{\pgfqpoint{1.467631in}{1.857545in}}%
\pgfpathlineto{\pgfqpoint{1.470149in}{1.846690in}}%
\pgfpathlineto{\pgfqpoint{1.472668in}{1.858905in}}%
\pgfpathlineto{\pgfqpoint{1.475187in}{1.871837in}}%
\pgfpathlineto{\pgfqpoint{1.477706in}{1.875524in}}%
\pgfpathlineto{\pgfqpoint{1.480224in}{1.903253in}}%
\pgfpathlineto{\pgfqpoint{1.482743in}{1.933659in}}%
\pgfpathlineto{\pgfqpoint{1.485262in}{1.939111in}}%
\pgfpathlineto{\pgfqpoint{1.487781in}{1.984519in}}%
\pgfpathlineto{\pgfqpoint{1.490299in}{1.982150in}}%
\pgfpathlineto{\pgfqpoint{1.492818in}{1.961398in}}%
\pgfpathlineto{\pgfqpoint{1.495337in}{1.930623in}}%
\pgfpathlineto{\pgfqpoint{1.497856in}{1.951321in}}%
\pgfpathlineto{\pgfqpoint{1.500374in}{1.966906in}}%
\pgfpathlineto{\pgfqpoint{1.502893in}{1.985513in}}%
\pgfpathlineto{\pgfqpoint{1.505412in}{2.012625in}}%
\pgfpathlineto{\pgfqpoint{1.507931in}{1.982858in}}%
\pgfpathlineto{\pgfqpoint{1.510449in}{1.996918in}}%
\pgfpathlineto{\pgfqpoint{1.512968in}{2.018204in}}%
\pgfpathlineto{\pgfqpoint{1.515487in}{2.009137in}}%
\pgfpathlineto{\pgfqpoint{1.518006in}{1.974245in}}%
\pgfpathlineto{\pgfqpoint{1.520524in}{1.984784in}}%
\pgfpathlineto{\pgfqpoint{1.523043in}{1.988712in}}%
\pgfpathlineto{\pgfqpoint{1.525562in}{1.988237in}}%
\pgfpathlineto{\pgfqpoint{1.528081in}{1.992943in}}%
\pgfpathlineto{\pgfqpoint{1.530599in}{2.011801in}}%
\pgfpathlineto{\pgfqpoint{1.533118in}{2.002336in}}%
\pgfpathlineto{\pgfqpoint{1.535637in}{1.987255in}}%
\pgfpathlineto{\pgfqpoint{1.538156in}{1.966720in}}%
\pgfpathlineto{\pgfqpoint{1.540674in}{1.983077in}}%
\pgfpathlineto{\pgfqpoint{1.543193in}{1.981582in}}%
\pgfpathlineto{\pgfqpoint{1.545712in}{1.989385in}}%
\pgfpathlineto{\pgfqpoint{1.548231in}{1.978517in}}%
\pgfpathlineto{\pgfqpoint{1.550749in}{1.987903in}}%
\pgfpathlineto{\pgfqpoint{1.553268in}{1.973983in}}%
\pgfpathlineto{\pgfqpoint{1.555787in}{1.989244in}}%
\pgfpathlineto{\pgfqpoint{1.558306in}{2.000980in}}%
\pgfpathlineto{\pgfqpoint{1.560824in}{1.992521in}}%
\pgfpathlineto{\pgfqpoint{1.563343in}{1.974091in}}%
\pgfpathlineto{\pgfqpoint{1.565862in}{1.975743in}}%
\pgfpathlineto{\pgfqpoint{1.568381in}{1.942512in}}%
\pgfpathlineto{\pgfqpoint{1.570899in}{1.944457in}}%
\pgfpathlineto{\pgfqpoint{1.573418in}{1.942508in}}%
\pgfpathlineto{\pgfqpoint{1.575937in}{1.942012in}}%
\pgfpathlineto{\pgfqpoint{1.578456in}{1.970039in}}%
\pgfpathlineto{\pgfqpoint{1.580974in}{1.957111in}}%
\pgfpathlineto{\pgfqpoint{1.583493in}{1.948455in}}%
\pgfpathlineto{\pgfqpoint{1.586012in}{1.932537in}}%
\pgfpathlineto{\pgfqpoint{1.588531in}{1.931471in}}%
\pgfpathlineto{\pgfqpoint{1.591049in}{1.920355in}}%
\pgfpathlineto{\pgfqpoint{1.593568in}{1.937547in}}%
\pgfpathlineto{\pgfqpoint{1.596087in}{1.929654in}}%
\pgfpathlineto{\pgfqpoint{1.598606in}{1.922511in}}%
\pgfpathlineto{\pgfqpoint{1.601124in}{1.918033in}}%
\pgfpathlineto{\pgfqpoint{1.603643in}{1.878857in}}%
\pgfpathlineto{\pgfqpoint{1.606162in}{1.848095in}}%
\pgfpathlineto{\pgfqpoint{1.608681in}{1.825079in}}%
\pgfpathlineto{\pgfqpoint{1.611199in}{1.817758in}}%
\pgfpathlineto{\pgfqpoint{1.613718in}{1.826545in}}%
\pgfpathlineto{\pgfqpoint{1.616237in}{1.799962in}}%
\pgfpathlineto{\pgfqpoint{1.618756in}{1.798719in}}%
\pgfpathlineto{\pgfqpoint{1.621274in}{1.836826in}}%
\pgfpathlineto{\pgfqpoint{1.623793in}{1.826291in}}%
\pgfpathlineto{\pgfqpoint{1.626312in}{1.818909in}}%
\pgfpathlineto{\pgfqpoint{1.628831in}{1.829370in}}%
\pgfpathlineto{\pgfqpoint{1.631349in}{1.844610in}}%
\pgfpathlineto{\pgfqpoint{1.633868in}{1.855928in}}%
\pgfpathlineto{\pgfqpoint{1.636387in}{1.852994in}}%
\pgfpathlineto{\pgfqpoint{1.638906in}{1.843498in}}%
\pgfpathlineto{\pgfqpoint{1.641424in}{1.828155in}}%
\pgfpathlineto{\pgfqpoint{1.643943in}{1.839899in}}%
\pgfpathlineto{\pgfqpoint{1.646462in}{1.848250in}}%
\pgfpathlineto{\pgfqpoint{1.648981in}{1.860856in}}%
\pgfpathlineto{\pgfqpoint{1.651499in}{1.856585in}}%
\pgfpathlineto{\pgfqpoint{1.654018in}{1.862837in}}%
\pgfpathlineto{\pgfqpoint{1.656537in}{1.850206in}}%
\pgfpathlineto{\pgfqpoint{1.659056in}{1.845390in}}%
\pgfpathlineto{\pgfqpoint{1.661574in}{1.819942in}}%
\pgfpathlineto{\pgfqpoint{1.664093in}{1.770523in}}%
\pgfpathlineto{\pgfqpoint{1.666612in}{1.779701in}}%
\pgfpathlineto{\pgfqpoint{1.669131in}{1.789414in}}%
\pgfpathlineto{\pgfqpoint{1.671649in}{1.786980in}}%
\pgfpathlineto{\pgfqpoint{1.674168in}{1.773101in}}%
\pgfpathlineto{\pgfqpoint{1.676687in}{1.773200in}}%
\pgfpathlineto{\pgfqpoint{1.679206in}{1.766370in}}%
\pgfpathlineto{\pgfqpoint{1.681724in}{1.780167in}}%
\pgfpathlineto{\pgfqpoint{1.684243in}{1.795336in}}%
\pgfpathlineto{\pgfqpoint{1.686762in}{1.776881in}}%
\pgfpathlineto{\pgfqpoint{1.689281in}{1.744944in}}%
\pgfpathlineto{\pgfqpoint{1.691799in}{1.762394in}}%
\pgfpathlineto{\pgfqpoint{1.694318in}{1.772846in}}%
\pgfpathlineto{\pgfqpoint{1.696837in}{1.775859in}}%
\pgfpathlineto{\pgfqpoint{1.699356in}{1.782918in}}%
\pgfpathlineto{\pgfqpoint{1.701874in}{1.800734in}}%
\pgfpathlineto{\pgfqpoint{1.704393in}{1.789979in}}%
\pgfpathlineto{\pgfqpoint{1.706912in}{1.790767in}}%
\pgfpathlineto{\pgfqpoint{1.709431in}{1.792299in}}%
\pgfpathlineto{\pgfqpoint{1.711949in}{1.791454in}}%
\pgfpathlineto{\pgfqpoint{1.714468in}{1.799786in}}%
\pgfpathlineto{\pgfqpoint{1.716987in}{1.796888in}}%
\pgfpathlineto{\pgfqpoint{1.719506in}{1.834202in}}%
\pgfpathlineto{\pgfqpoint{1.722024in}{1.806029in}}%
\pgfpathlineto{\pgfqpoint{1.724543in}{1.787526in}}%
\pgfpathlineto{\pgfqpoint{1.727062in}{1.786579in}}%
\pgfpathlineto{\pgfqpoint{1.729581in}{1.793288in}}%
\pgfpathlineto{\pgfqpoint{1.732099in}{1.795603in}}%
\pgfpathlineto{\pgfqpoint{1.734618in}{1.782615in}}%
\pgfpathlineto{\pgfqpoint{1.737137in}{1.797696in}}%
\pgfpathlineto{\pgfqpoint{1.739656in}{1.792463in}}%
\pgfpathlineto{\pgfqpoint{1.742174in}{1.826395in}}%
\pgfpathlineto{\pgfqpoint{1.744693in}{1.839483in}}%
\pgfpathlineto{\pgfqpoint{1.747212in}{1.840213in}}%
\pgfpathlineto{\pgfqpoint{1.749731in}{1.841922in}}%
\pgfpathlineto{\pgfqpoint{1.752249in}{1.853114in}}%
\pgfpathlineto{\pgfqpoint{1.754768in}{1.858284in}}%
\pgfpathlineto{\pgfqpoint{1.757287in}{1.862651in}}%
\pgfpathlineto{\pgfqpoint{1.759806in}{1.861425in}}%
\pgfpathlineto{\pgfqpoint{1.762324in}{1.854223in}}%
\pgfpathlineto{\pgfqpoint{1.764843in}{1.839320in}}%
\pgfpathlineto{\pgfqpoint{1.767362in}{1.847224in}}%
\pgfpathlineto{\pgfqpoint{1.769881in}{1.865288in}}%
\pgfpathlineto{\pgfqpoint{1.772399in}{1.854115in}}%
\pgfpathlineto{\pgfqpoint{1.774918in}{1.832465in}}%
\pgfpathlineto{\pgfqpoint{1.777437in}{1.857781in}}%
\pgfpathlineto{\pgfqpoint{1.779956in}{1.861814in}}%
\pgfpathlineto{\pgfqpoint{1.784993in}{1.866727in}}%
\pgfpathlineto{\pgfqpoint{1.787512in}{1.894221in}}%
\pgfpathlineto{\pgfqpoint{1.790031in}{1.893537in}}%
\pgfpathlineto{\pgfqpoint{1.792549in}{1.909599in}}%
\pgfpathlineto{\pgfqpoint{1.795068in}{1.939459in}}%
\pgfpathlineto{\pgfqpoint{1.797587in}{1.949909in}}%
\pgfpathlineto{\pgfqpoint{1.800106in}{1.940253in}}%
\pgfpathlineto{\pgfqpoint{1.802624in}{1.938823in}}%
\pgfpathlineto{\pgfqpoint{1.805143in}{1.933465in}}%
\pgfpathlineto{\pgfqpoint{1.807662in}{1.916976in}}%
\pgfpathlineto{\pgfqpoint{1.810181in}{1.927547in}}%
\pgfpathlineto{\pgfqpoint{1.812699in}{1.926447in}}%
\pgfpathlineto{\pgfqpoint{1.815218in}{1.905386in}}%
\pgfpathlineto{\pgfqpoint{1.817737in}{1.911563in}}%
\pgfpathlineto{\pgfqpoint{1.820256in}{1.940449in}}%
\pgfpathlineto{\pgfqpoint{1.822774in}{1.946986in}}%
\pgfpathlineto{\pgfqpoint{1.825293in}{1.979805in}}%
\pgfpathlineto{\pgfqpoint{1.827812in}{1.963839in}}%
\pgfpathlineto{\pgfqpoint{1.830331in}{1.937599in}}%
\pgfpathlineto{\pgfqpoint{1.832849in}{1.929949in}}%
\pgfpathlineto{\pgfqpoint{1.835368in}{1.953796in}}%
\pgfpathlineto{\pgfqpoint{1.837887in}{1.961213in}}%
\pgfpathlineto{\pgfqpoint{1.840406in}{1.957890in}}%
\pgfpathlineto{\pgfqpoint{1.842924in}{1.940968in}}%
\pgfpathlineto{\pgfqpoint{1.845443in}{1.930311in}}%
\pgfpathlineto{\pgfqpoint{1.847962in}{1.956907in}}%
\pgfpathlineto{\pgfqpoint{1.850481in}{1.974233in}}%
\pgfpathlineto{\pgfqpoint{1.855518in}{2.019343in}}%
\pgfpathlineto{\pgfqpoint{1.858037in}{2.035773in}}%
\pgfpathlineto{\pgfqpoint{1.860556in}{2.022637in}}%
\pgfpathlineto{\pgfqpoint{1.863074in}{2.034045in}}%
\pgfpathlineto{\pgfqpoint{1.865593in}{2.006463in}}%
\pgfpathlineto{\pgfqpoint{1.868112in}{2.014785in}}%
\pgfpathlineto{\pgfqpoint{1.870631in}{2.045787in}}%
\pgfpathlineto{\pgfqpoint{1.873149in}{2.030538in}}%
\pgfpathlineto{\pgfqpoint{1.875668in}{1.999069in}}%
\pgfpathlineto{\pgfqpoint{1.878187in}{1.987730in}}%
\pgfpathlineto{\pgfqpoint{1.880706in}{2.000943in}}%
\pgfpathlineto{\pgfqpoint{1.883224in}{2.016001in}}%
\pgfpathlineto{\pgfqpoint{1.885743in}{2.039604in}}%
\pgfpathlineto{\pgfqpoint{1.888262in}{2.020867in}}%
\pgfpathlineto{\pgfqpoint{1.890781in}{2.020994in}}%
\pgfpathlineto{\pgfqpoint{1.893299in}{2.049604in}}%
\pgfpathlineto{\pgfqpoint{1.895818in}{2.021881in}}%
\pgfpathlineto{\pgfqpoint{1.898337in}{2.029155in}}%
\pgfpathlineto{\pgfqpoint{1.900856in}{2.022729in}}%
\pgfpathlineto{\pgfqpoint{1.903374in}{2.010506in}}%
\pgfpathlineto{\pgfqpoint{1.905893in}{2.037797in}}%
\pgfpathlineto{\pgfqpoint{1.908412in}{2.028915in}}%
\pgfpathlineto{\pgfqpoint{1.910931in}{2.025145in}}%
\pgfpathlineto{\pgfqpoint{1.913449in}{2.036095in}}%
\pgfpathlineto{\pgfqpoint{1.915968in}{2.064118in}}%
\pgfpathlineto{\pgfqpoint{1.918487in}{2.071002in}}%
\pgfpathlineto{\pgfqpoint{1.921006in}{2.069875in}}%
\pgfpathlineto{\pgfqpoint{1.923524in}{2.076668in}}%
\pgfpathlineto{\pgfqpoint{1.926043in}{2.067378in}}%
\pgfpathlineto{\pgfqpoint{1.928562in}{2.068310in}}%
\pgfpathlineto{\pgfqpoint{1.931081in}{2.081453in}}%
\pgfpathlineto{\pgfqpoint{1.933599in}{2.087589in}}%
\pgfpathlineto{\pgfqpoint{1.936118in}{2.098694in}}%
\pgfpathlineto{\pgfqpoint{1.938637in}{2.106766in}}%
\pgfpathlineto{\pgfqpoint{1.941156in}{2.091723in}}%
\pgfpathlineto{\pgfqpoint{1.943674in}{2.097584in}}%
\pgfpathlineto{\pgfqpoint{1.946193in}{2.072465in}}%
\pgfpathlineto{\pgfqpoint{1.948712in}{2.058332in}}%
\pgfpathlineto{\pgfqpoint{1.951231in}{2.040128in}}%
\pgfpathlineto{\pgfqpoint{1.953749in}{2.034325in}}%
\pgfpathlineto{\pgfqpoint{1.956268in}{2.005211in}}%
\pgfpathlineto{\pgfqpoint{1.958787in}{2.018699in}}%
\pgfpathlineto{\pgfqpoint{1.961306in}{2.025565in}}%
\pgfpathlineto{\pgfqpoint{1.963824in}{2.013904in}}%
\pgfpathlineto{\pgfqpoint{1.966343in}{1.983045in}}%
\pgfpathlineto{\pgfqpoint{1.968862in}{1.991987in}}%
\pgfpathlineto{\pgfqpoint{1.971381in}{2.001846in}}%
\pgfpathlineto{\pgfqpoint{1.976418in}{2.023705in}}%
\pgfpathlineto{\pgfqpoint{1.978937in}{2.031636in}}%
\pgfpathlineto{\pgfqpoint{1.981456in}{2.060128in}}%
\pgfpathlineto{\pgfqpoint{1.983974in}{2.046355in}}%
\pgfpathlineto{\pgfqpoint{1.986493in}{2.065471in}}%
\pgfpathlineto{\pgfqpoint{1.989012in}{2.042736in}}%
\pgfpathlineto{\pgfqpoint{1.991531in}{2.033727in}}%
\pgfpathlineto{\pgfqpoint{1.994049in}{2.016814in}}%
\pgfpathlineto{\pgfqpoint{1.996568in}{2.030011in}}%
\pgfpathlineto{\pgfqpoint{1.999087in}{2.032585in}}%
\pgfpathlineto{\pgfqpoint{2.001606in}{2.014905in}}%
\pgfpathlineto{\pgfqpoint{2.004124in}{2.025884in}}%
\pgfpathlineto{\pgfqpoint{2.006643in}{2.021325in}}%
\pgfpathlineto{\pgfqpoint{2.009162in}{2.035325in}}%
\pgfpathlineto{\pgfqpoint{2.011681in}{2.039519in}}%
\pgfpathlineto{\pgfqpoint{2.014199in}{2.035333in}}%
\pgfpathlineto{\pgfqpoint{2.016718in}{2.048287in}}%
\pgfpathlineto{\pgfqpoint{2.019237in}{2.013950in}}%
\pgfpathlineto{\pgfqpoint{2.021756in}{2.042551in}}%
\pgfpathlineto{\pgfqpoint{2.024274in}{2.035462in}}%
\pgfpathlineto{\pgfqpoint{2.026793in}{2.042673in}}%
\pgfpathlineto{\pgfqpoint{2.031831in}{2.089373in}}%
\pgfpathlineto{\pgfqpoint{2.034349in}{2.100314in}}%
\pgfpathlineto{\pgfqpoint{2.036868in}{2.101748in}}%
\pgfpathlineto{\pgfqpoint{2.039387in}{2.106994in}}%
\pgfpathlineto{\pgfqpoint{2.041906in}{2.091343in}}%
\pgfpathlineto{\pgfqpoint{2.044424in}{2.095180in}}%
\pgfpathlineto{\pgfqpoint{2.046943in}{2.077429in}}%
\pgfpathlineto{\pgfqpoint{2.049462in}{2.118265in}}%
\pgfpathlineto{\pgfqpoint{2.051981in}{2.091688in}}%
\pgfpathlineto{\pgfqpoint{2.054499in}{2.097757in}}%
\pgfpathlineto{\pgfqpoint{2.057018in}{2.114691in}}%
\pgfpathlineto{\pgfqpoint{2.059537in}{2.106827in}}%
\pgfpathlineto{\pgfqpoint{2.062056in}{2.100074in}}%
\pgfpathlineto{\pgfqpoint{2.064574in}{2.095593in}}%
\pgfpathlineto{\pgfqpoint{2.067093in}{2.096485in}}%
\pgfpathlineto{\pgfqpoint{2.069612in}{2.081643in}}%
\pgfpathlineto{\pgfqpoint{2.072131in}{2.114427in}}%
\pgfpathlineto{\pgfqpoint{2.074649in}{2.130239in}}%
\pgfpathlineto{\pgfqpoint{2.077168in}{2.141322in}}%
\pgfpathlineto{\pgfqpoint{2.079687in}{2.164988in}}%
\pgfpathlineto{\pgfqpoint{2.082206in}{2.179337in}}%
\pgfpathlineto{\pgfqpoint{2.084724in}{2.235006in}}%
\pgfpathlineto{\pgfqpoint{2.089762in}{2.201163in}}%
\pgfpathlineto{\pgfqpoint{2.092281in}{2.206784in}}%
\pgfpathlineto{\pgfqpoint{2.094799in}{2.204498in}}%
\pgfpathlineto{\pgfqpoint{2.097318in}{2.213766in}}%
\pgfpathlineto{\pgfqpoint{2.099837in}{2.223885in}}%
\pgfpathlineto{\pgfqpoint{2.102356in}{2.243169in}}%
\pgfpathlineto{\pgfqpoint{2.104874in}{2.255340in}}%
\pgfpathlineto{\pgfqpoint{2.107393in}{2.244731in}}%
\pgfpathlineto{\pgfqpoint{2.109912in}{2.224879in}}%
\pgfpathlineto{\pgfqpoint{2.112431in}{2.206105in}}%
\pgfpathlineto{\pgfqpoint{2.114949in}{2.194446in}}%
\pgfpathlineto{\pgfqpoint{2.117468in}{2.177223in}}%
\pgfpathlineto{\pgfqpoint{2.119987in}{2.167726in}}%
\pgfpathlineto{\pgfqpoint{2.122506in}{2.179329in}}%
\pgfpathlineto{\pgfqpoint{2.125024in}{2.170732in}}%
\pgfpathlineto{\pgfqpoint{2.127543in}{2.168563in}}%
\pgfpathlineto{\pgfqpoint{2.130062in}{2.163855in}}%
\pgfpathlineto{\pgfqpoint{2.132581in}{2.157407in}}%
\pgfpathlineto{\pgfqpoint{2.135099in}{2.159059in}}%
\pgfpathlineto{\pgfqpoint{2.137618in}{2.150912in}}%
\pgfpathlineto{\pgfqpoint{2.140137in}{2.170167in}}%
\pgfpathlineto{\pgfqpoint{2.142656in}{2.136523in}}%
\pgfpathlineto{\pgfqpoint{2.145174in}{2.150969in}}%
\pgfpathlineto{\pgfqpoint{2.147693in}{2.121426in}}%
\pgfpathlineto{\pgfqpoint{2.150212in}{2.124693in}}%
\pgfpathlineto{\pgfqpoint{2.152731in}{2.144844in}}%
\pgfpathlineto{\pgfqpoint{2.155249in}{2.141973in}}%
\pgfpathlineto{\pgfqpoint{2.157768in}{2.202095in}}%
\pgfpathlineto{\pgfqpoint{2.160287in}{2.181266in}}%
\pgfpathlineto{\pgfqpoint{2.162806in}{2.180331in}}%
\pgfpathlineto{\pgfqpoint{2.165324in}{2.150923in}}%
\pgfpathlineto{\pgfqpoint{2.167843in}{2.176567in}}%
\pgfpathlineto{\pgfqpoint{2.170362in}{2.187931in}}%
\pgfpathlineto{\pgfqpoint{2.172881in}{2.205273in}}%
\pgfpathlineto{\pgfqpoint{2.175399in}{2.190769in}}%
\pgfpathlineto{\pgfqpoint{2.177918in}{2.170809in}}%
\pgfpathlineto{\pgfqpoint{2.180437in}{2.127999in}}%
\pgfpathlineto{\pgfqpoint{2.182956in}{2.154362in}}%
\pgfpathlineto{\pgfqpoint{2.185474in}{2.152520in}}%
\pgfpathlineto{\pgfqpoint{2.187993in}{2.160848in}}%
\pgfpathlineto{\pgfqpoint{2.190512in}{2.190905in}}%
\pgfpathlineto{\pgfqpoint{2.193031in}{2.196783in}}%
\pgfpathlineto{\pgfqpoint{2.195549in}{2.181503in}}%
\pgfpathlineto{\pgfqpoint{2.198068in}{2.182195in}}%
\pgfpathlineto{\pgfqpoint{2.200587in}{2.185114in}}%
\pgfpathlineto{\pgfqpoint{2.203106in}{2.150994in}}%
\pgfpathlineto{\pgfqpoint{2.205624in}{2.141029in}}%
\pgfpathlineto{\pgfqpoint{2.208143in}{2.156437in}}%
\pgfpathlineto{\pgfqpoint{2.210662in}{2.151656in}}%
\pgfpathlineto{\pgfqpoint{2.213181in}{2.140288in}}%
\pgfpathlineto{\pgfqpoint{2.215699in}{2.152309in}}%
\pgfpathlineto{\pgfqpoint{2.218218in}{2.168356in}}%
\pgfpathlineto{\pgfqpoint{2.220737in}{2.195275in}}%
\pgfpathlineto{\pgfqpoint{2.223256in}{2.182437in}}%
\pgfpathlineto{\pgfqpoint{2.225774in}{2.171071in}}%
\pgfpathlineto{\pgfqpoint{2.228293in}{2.195587in}}%
\pgfpathlineto{\pgfqpoint{2.230812in}{2.228489in}}%
\pgfpathlineto{\pgfqpoint{2.233331in}{2.217343in}}%
\pgfpathlineto{\pgfqpoint{2.235849in}{2.226964in}}%
\pgfpathlineto{\pgfqpoint{2.238368in}{2.226606in}}%
\pgfpathlineto{\pgfqpoint{2.240887in}{2.204237in}}%
\pgfpathlineto{\pgfqpoint{2.243406in}{2.225541in}}%
\pgfpathlineto{\pgfqpoint{2.245924in}{2.243306in}}%
\pgfpathlineto{\pgfqpoint{2.248443in}{2.250169in}}%
\pgfpathlineto{\pgfqpoint{2.250962in}{2.269996in}}%
\pgfpathlineto{\pgfqpoint{2.253481in}{2.267461in}}%
\pgfpathlineto{\pgfqpoint{2.255999in}{2.277879in}}%
\pgfpathlineto{\pgfqpoint{2.258518in}{2.278143in}}%
\pgfpathlineto{\pgfqpoint{2.261037in}{2.293988in}}%
\pgfpathlineto{\pgfqpoint{2.263556in}{2.271716in}}%
\pgfpathlineto{\pgfqpoint{2.266074in}{2.283158in}}%
\pgfpathlineto{\pgfqpoint{2.268593in}{2.279301in}}%
\pgfpathlineto{\pgfqpoint{2.271112in}{2.266304in}}%
\pgfpathlineto{\pgfqpoint{2.273631in}{2.228451in}}%
\pgfpathlineto{\pgfqpoint{2.276149in}{2.226168in}}%
\pgfpathlineto{\pgfqpoint{2.278668in}{2.214228in}}%
\pgfpathlineto{\pgfqpoint{2.281187in}{2.255257in}}%
\pgfpathlineto{\pgfqpoint{2.283706in}{2.213899in}}%
\pgfpathlineto{\pgfqpoint{2.286224in}{2.235115in}}%
\pgfpathlineto{\pgfqpoint{2.288743in}{2.257559in}}%
\pgfpathlineto{\pgfqpoint{2.291262in}{2.289306in}}%
\pgfpathlineto{\pgfqpoint{2.293781in}{2.316415in}}%
\pgfpathlineto{\pgfqpoint{2.296299in}{2.304291in}}%
\pgfpathlineto{\pgfqpoint{2.298818in}{2.284315in}}%
\pgfpathlineto{\pgfqpoint{2.301337in}{2.287264in}}%
\pgfpathlineto{\pgfqpoint{2.303856in}{2.279314in}}%
\pgfpathlineto{\pgfqpoint{2.306374in}{2.295938in}}%
\pgfpathlineto{\pgfqpoint{2.308893in}{2.303253in}}%
\pgfpathlineto{\pgfqpoint{2.311412in}{2.319652in}}%
\pgfpathlineto{\pgfqpoint{2.313931in}{2.337196in}}%
\pgfpathlineto{\pgfqpoint{2.316449in}{2.332982in}}%
\pgfpathlineto{\pgfqpoint{2.318968in}{2.328065in}}%
\pgfpathlineto{\pgfqpoint{2.321487in}{2.380249in}}%
\pgfpathlineto{\pgfqpoint{2.324006in}{2.368400in}}%
\pgfpathlineto{\pgfqpoint{2.326524in}{2.336905in}}%
\pgfpathlineto{\pgfqpoint{2.329043in}{2.323357in}}%
\pgfpathlineto{\pgfqpoint{2.331562in}{2.342923in}}%
\pgfpathlineto{\pgfqpoint{2.334081in}{2.330603in}}%
\pgfpathlineto{\pgfqpoint{2.336599in}{2.348173in}}%
\pgfpathlineto{\pgfqpoint{2.339118in}{2.380831in}}%
\pgfpathlineto{\pgfqpoint{2.341637in}{2.368302in}}%
\pgfpathlineto{\pgfqpoint{2.344156in}{2.368381in}}%
\pgfpathlineto{\pgfqpoint{2.346674in}{2.360107in}}%
\pgfpathlineto{\pgfqpoint{2.349193in}{2.348369in}}%
\pgfpathlineto{\pgfqpoint{2.351712in}{2.346282in}}%
\pgfpathlineto{\pgfqpoint{2.354231in}{2.354557in}}%
\pgfpathlineto{\pgfqpoint{2.356749in}{2.378848in}}%
\pgfpathlineto{\pgfqpoint{2.359268in}{2.384869in}}%
\pgfpathlineto{\pgfqpoint{2.361787in}{2.382404in}}%
\pgfpathlineto{\pgfqpoint{2.364306in}{2.370211in}}%
\pgfpathlineto{\pgfqpoint{2.366824in}{2.343794in}}%
\pgfpathlineto{\pgfqpoint{2.369343in}{2.380044in}}%
\pgfpathlineto{\pgfqpoint{2.371862in}{2.367797in}}%
\pgfpathlineto{\pgfqpoint{2.374381in}{2.350255in}}%
\pgfpathlineto{\pgfqpoint{2.376899in}{2.375293in}}%
\pgfpathlineto{\pgfqpoint{2.379418in}{2.353425in}}%
\pgfpathlineto{\pgfqpoint{2.381937in}{2.359057in}}%
\pgfpathlineto{\pgfqpoint{2.384456in}{2.378471in}}%
\pgfpathlineto{\pgfqpoint{2.386974in}{2.343706in}}%
\pgfpathlineto{\pgfqpoint{2.389493in}{2.348257in}}%
\pgfpathlineto{\pgfqpoint{2.392012in}{2.367595in}}%
\pgfpathlineto{\pgfqpoint{2.394531in}{2.366338in}}%
\pgfpathlineto{\pgfqpoint{2.397049in}{2.362925in}}%
\pgfpathlineto{\pgfqpoint{2.399568in}{2.370157in}}%
\pgfpathlineto{\pgfqpoint{2.402087in}{2.436350in}}%
\pgfpathlineto{\pgfqpoint{2.404606in}{2.433632in}}%
\pgfpathlineto{\pgfqpoint{2.407124in}{2.417089in}}%
\pgfpathlineto{\pgfqpoint{2.409643in}{2.467752in}}%
\pgfpathlineto{\pgfqpoint{2.412162in}{2.481389in}}%
\pgfpathlineto{\pgfqpoint{2.414681in}{2.486305in}}%
\pgfpathlineto{\pgfqpoint{2.417199in}{2.524006in}}%
\pgfpathlineto{\pgfqpoint{2.419718in}{2.542366in}}%
\pgfpathlineto{\pgfqpoint{2.422237in}{2.557740in}}%
\pgfpathlineto{\pgfqpoint{2.424756in}{2.598820in}}%
\pgfpathlineto{\pgfqpoint{2.427274in}{2.612611in}}%
\pgfpathlineto{\pgfqpoint{2.429793in}{2.607561in}}%
\pgfpathlineto{\pgfqpoint{2.432312in}{2.563741in}}%
\pgfpathlineto{\pgfqpoint{2.434831in}{2.585157in}}%
\pgfpathlineto{\pgfqpoint{2.437349in}{2.566338in}}%
\pgfpathlineto{\pgfqpoint{2.439868in}{2.577843in}}%
\pgfpathlineto{\pgfqpoint{2.442387in}{2.576780in}}%
\pgfpathlineto{\pgfqpoint{2.444906in}{2.579519in}}%
\pgfpathlineto{\pgfqpoint{2.447424in}{2.587369in}}%
\pgfpathlineto{\pgfqpoint{2.449943in}{2.612623in}}%
\pgfpathlineto{\pgfqpoint{2.452462in}{2.630634in}}%
\pgfpathlineto{\pgfqpoint{2.454981in}{2.623507in}}%
\pgfpathlineto{\pgfqpoint{2.457499in}{2.640594in}}%
\pgfpathlineto{\pgfqpoint{2.460018in}{2.629113in}}%
\pgfpathlineto{\pgfqpoint{2.462537in}{2.631552in}}%
\pgfpathlineto{\pgfqpoint{2.465056in}{2.655140in}}%
\pgfpathlineto{\pgfqpoint{2.467574in}{2.618563in}}%
\pgfpathlineto{\pgfqpoint{2.470093in}{2.575633in}}%
\pgfpathlineto{\pgfqpoint{2.472612in}{2.575587in}}%
\pgfpathlineto{\pgfqpoint{2.475131in}{2.622337in}}%
\pgfpathlineto{\pgfqpoint{2.477649in}{2.611263in}}%
\pgfpathlineto{\pgfqpoint{2.480168in}{2.633218in}}%
\pgfpathlineto{\pgfqpoint{2.482687in}{2.595442in}}%
\pgfpathlineto{\pgfqpoint{2.485206in}{2.571679in}}%
\pgfpathlineto{\pgfqpoint{2.487724in}{2.611124in}}%
\pgfpathlineto{\pgfqpoint{2.490243in}{2.615989in}}%
\pgfpathlineto{\pgfqpoint{2.492762in}{2.617892in}}%
\pgfpathlineto{\pgfqpoint{2.495281in}{2.622258in}}%
\pgfpathlineto{\pgfqpoint{2.497799in}{2.633824in}}%
\pgfpathlineto{\pgfqpoint{2.500318in}{2.666053in}}%
\pgfpathlineto{\pgfqpoint{2.502837in}{2.624862in}}%
\pgfpathlineto{\pgfqpoint{2.505356in}{2.605357in}}%
\pgfpathlineto{\pgfqpoint{2.507874in}{2.610860in}}%
\pgfpathlineto{\pgfqpoint{2.510393in}{2.610598in}}%
\pgfpathlineto{\pgfqpoint{2.512912in}{2.583978in}}%
\pgfpathlineto{\pgfqpoint{2.515431in}{2.592801in}}%
\pgfpathlineto{\pgfqpoint{2.517949in}{2.586350in}}%
\pgfpathlineto{\pgfqpoint{2.520468in}{2.597985in}}%
\pgfpathlineto{\pgfqpoint{2.522987in}{2.595442in}}%
\pgfpathlineto{\pgfqpoint{2.525506in}{2.586055in}}%
\pgfpathlineto{\pgfqpoint{2.528024in}{2.595307in}}%
\pgfpathlineto{\pgfqpoint{2.530543in}{2.576200in}}%
\pgfpathlineto{\pgfqpoint{2.533062in}{2.551350in}}%
\pgfpathlineto{\pgfqpoint{2.535581in}{2.558380in}}%
\pgfpathlineto{\pgfqpoint{2.538099in}{2.563411in}}%
\pgfpathlineto{\pgfqpoint{2.540618in}{2.577849in}}%
\pgfpathlineto{\pgfqpoint{2.543137in}{2.635925in}}%
\pgfpathlineto{\pgfqpoint{2.545656in}{2.611097in}}%
\pgfpathlineto{\pgfqpoint{2.548174in}{2.634774in}}%
\pgfpathlineto{\pgfqpoint{2.550693in}{2.645109in}}%
\pgfpathlineto{\pgfqpoint{2.553212in}{2.644708in}}%
\pgfpathlineto{\pgfqpoint{2.555731in}{2.658484in}}%
\pgfpathlineto{\pgfqpoint{2.558249in}{2.624990in}}%
\pgfpathlineto{\pgfqpoint{2.560768in}{2.608322in}}%
\pgfpathlineto{\pgfqpoint{2.563287in}{2.584103in}}%
\pgfpathlineto{\pgfqpoint{2.565806in}{2.582290in}}%
\pgfpathlineto{\pgfqpoint{2.568324in}{2.594952in}}%
\pgfpathlineto{\pgfqpoint{2.570843in}{2.552748in}}%
\pgfpathlineto{\pgfqpoint{2.573362in}{2.572897in}}%
\pgfpathlineto{\pgfqpoint{2.575881in}{2.572133in}}%
\pgfpathlineto{\pgfqpoint{2.578399in}{2.584737in}}%
\pgfpathlineto{\pgfqpoint{2.580918in}{2.555866in}}%
\pgfpathlineto{\pgfqpoint{2.583437in}{2.563336in}}%
\pgfpathlineto{\pgfqpoint{2.585956in}{2.576348in}}%
\pgfpathlineto{\pgfqpoint{2.588474in}{2.605012in}}%
\pgfpathlineto{\pgfqpoint{2.590993in}{2.599664in}}%
\pgfpathlineto{\pgfqpoint{2.593512in}{2.583222in}}%
\pgfpathlineto{\pgfqpoint{2.596031in}{2.615873in}}%
\pgfpathlineto{\pgfqpoint{2.598549in}{2.637376in}}%
\pgfpathlineto{\pgfqpoint{2.601068in}{2.630601in}}%
\pgfpathlineto{\pgfqpoint{2.603587in}{2.644063in}}%
\pgfpathlineto{\pgfqpoint{2.606106in}{2.673881in}}%
\pgfpathlineto{\pgfqpoint{2.608624in}{2.668152in}}%
\pgfpathlineto{\pgfqpoint{2.611143in}{2.699014in}}%
\pgfpathlineto{\pgfqpoint{2.613662in}{2.691614in}}%
\pgfpathlineto{\pgfqpoint{2.616181in}{2.664477in}}%
\pgfpathlineto{\pgfqpoint{2.618699in}{2.656000in}}%
\pgfpathlineto{\pgfqpoint{2.621218in}{2.633930in}}%
\pgfpathlineto{\pgfqpoint{2.623737in}{2.639120in}}%
\pgfpathlineto{\pgfqpoint{2.626256in}{2.665118in}}%
\pgfpathlineto{\pgfqpoint{2.628774in}{2.673582in}}%
\pgfpathlineto{\pgfqpoint{2.631293in}{2.662129in}}%
\pgfpathlineto{\pgfqpoint{2.633812in}{2.657097in}}%
\pgfpathlineto{\pgfqpoint{2.636331in}{2.660541in}}%
\pgfpathlineto{\pgfqpoint{2.638849in}{2.631683in}}%
\pgfpathlineto{\pgfqpoint{2.641368in}{2.639271in}}%
\pgfpathlineto{\pgfqpoint{2.643887in}{2.649818in}}%
\pgfpathlineto{\pgfqpoint{2.646406in}{2.655431in}}%
\pgfpathlineto{\pgfqpoint{2.648924in}{2.640245in}}%
\pgfpathlineto{\pgfqpoint{2.651443in}{2.637234in}}%
\pgfpathlineto{\pgfqpoint{2.653962in}{2.647268in}}%
\pgfpathlineto{\pgfqpoint{2.656481in}{2.675661in}}%
\pgfpathlineto{\pgfqpoint{2.658999in}{2.677520in}}%
\pgfpathlineto{\pgfqpoint{2.661518in}{2.677799in}}%
\pgfpathlineto{\pgfqpoint{2.664037in}{2.658574in}}%
\pgfpathlineto{\pgfqpoint{2.669074in}{2.681667in}}%
\pgfpathlineto{\pgfqpoint{2.671593in}{2.662408in}}%
\pgfpathlineto{\pgfqpoint{2.674112in}{2.666129in}}%
\pgfpathlineto{\pgfqpoint{2.676631in}{2.669187in}}%
\pgfpathlineto{\pgfqpoint{2.679149in}{2.685645in}}%
\pgfpathlineto{\pgfqpoint{2.681668in}{2.698686in}}%
\pgfpathlineto{\pgfqpoint{2.684187in}{2.710196in}}%
\pgfpathlineto{\pgfqpoint{2.686706in}{2.680953in}}%
\pgfpathlineto{\pgfqpoint{2.689224in}{2.723105in}}%
\pgfpathlineto{\pgfqpoint{2.691743in}{2.725935in}}%
\pgfpathlineto{\pgfqpoint{2.694262in}{2.692772in}}%
\pgfpathlineto{\pgfqpoint{2.696781in}{2.713191in}}%
\pgfpathlineto{\pgfqpoint{2.699299in}{2.716117in}}%
\pgfpathlineto{\pgfqpoint{2.701818in}{2.743563in}}%
\pgfpathlineto{\pgfqpoint{2.704337in}{2.739056in}}%
\pgfpathlineto{\pgfqpoint{2.706856in}{2.727236in}}%
\pgfpathlineto{\pgfqpoint{2.709374in}{2.752522in}}%
\pgfpathlineto{\pgfqpoint{2.711893in}{2.782958in}}%
\pgfpathlineto{\pgfqpoint{2.714412in}{2.781251in}}%
\pgfpathlineto{\pgfqpoint{2.716931in}{2.799258in}}%
\pgfpathlineto{\pgfqpoint{2.719449in}{2.748280in}}%
\pgfpathlineto{\pgfqpoint{2.721968in}{2.743478in}}%
\pgfpathlineto{\pgfqpoint{2.724487in}{2.755868in}}%
\pgfpathlineto{\pgfqpoint{2.727006in}{2.792007in}}%
\pgfpathlineto{\pgfqpoint{2.729524in}{2.784659in}}%
\pgfpathlineto{\pgfqpoint{2.732043in}{2.800412in}}%
\pgfpathlineto{\pgfqpoint{2.734562in}{2.807792in}}%
\pgfpathlineto{\pgfqpoint{2.737081in}{2.803829in}}%
\pgfpathlineto{\pgfqpoint{2.739599in}{2.817196in}}%
\pgfpathlineto{\pgfqpoint{2.742118in}{2.845065in}}%
\pgfpathlineto{\pgfqpoint{2.744637in}{2.864549in}}%
\pgfpathlineto{\pgfqpoint{2.747156in}{2.832070in}}%
\pgfpathlineto{\pgfqpoint{2.749674in}{2.853520in}}%
\pgfpathlineto{\pgfqpoint{2.752193in}{2.854960in}}%
\pgfpathlineto{\pgfqpoint{2.754712in}{2.848049in}}%
\pgfpathlineto{\pgfqpoint{2.757231in}{2.850945in}}%
\pgfpathlineto{\pgfqpoint{2.759749in}{2.832379in}}%
\pgfpathlineto{\pgfqpoint{2.762268in}{2.777719in}}%
\pgfpathlineto{\pgfqpoint{2.764787in}{2.764420in}}%
\pgfpathlineto{\pgfqpoint{2.767306in}{2.756569in}}%
\pgfpathlineto{\pgfqpoint{2.769824in}{2.752497in}}%
\pgfpathlineto{\pgfqpoint{2.772343in}{2.750520in}}%
\pgfpathlineto{\pgfqpoint{2.774862in}{2.742717in}}%
\pgfpathlineto{\pgfqpoint{2.777381in}{2.743815in}}%
\pgfpathlineto{\pgfqpoint{2.779899in}{2.748880in}}%
\pgfpathlineto{\pgfqpoint{2.782418in}{2.746235in}}%
\pgfpathlineto{\pgfqpoint{2.784937in}{2.783800in}}%
\pgfpathlineto{\pgfqpoint{2.787456in}{2.770985in}}%
\pgfpathlineto{\pgfqpoint{2.789974in}{2.760056in}}%
\pgfpathlineto{\pgfqpoint{2.792493in}{2.778857in}}%
\pgfpathlineto{\pgfqpoint{2.795012in}{2.805528in}}%
\pgfpathlineto{\pgfqpoint{2.797531in}{2.802093in}}%
\pgfpathlineto{\pgfqpoint{2.800049in}{2.790852in}}%
\pgfpathlineto{\pgfqpoint{2.802568in}{2.794844in}}%
\pgfpathlineto{\pgfqpoint{2.805087in}{2.786371in}}%
\pgfpathlineto{\pgfqpoint{2.807606in}{2.751187in}}%
\pgfpathlineto{\pgfqpoint{2.810124in}{2.754210in}}%
\pgfpathlineto{\pgfqpoint{2.812643in}{2.744681in}}%
\pgfpathlineto{\pgfqpoint{2.815162in}{2.728907in}}%
\pgfpathlineto{\pgfqpoint{2.817681in}{2.728193in}}%
\pgfpathlineto{\pgfqpoint{2.820199in}{2.759686in}}%
\pgfpathlineto{\pgfqpoint{2.822718in}{2.753979in}}%
\pgfpathlineto{\pgfqpoint{2.825237in}{2.770349in}}%
\pgfpathlineto{\pgfqpoint{2.827756in}{2.791031in}}%
\pgfpathlineto{\pgfqpoint{2.830274in}{2.793780in}}%
\pgfpathlineto{\pgfqpoint{2.832793in}{2.808748in}}%
\pgfpathlineto{\pgfqpoint{2.835312in}{2.801365in}}%
\pgfpathlineto{\pgfqpoint{2.837831in}{2.801047in}}%
\pgfpathlineto{\pgfqpoint{2.840349in}{2.808619in}}%
\pgfpathlineto{\pgfqpoint{2.842868in}{2.815819in}}%
\pgfpathlineto{\pgfqpoint{2.845387in}{2.837547in}}%
\pgfpathlineto{\pgfqpoint{2.847906in}{2.834105in}}%
\pgfpathlineto{\pgfqpoint{2.850424in}{2.839820in}}%
\pgfpathlineto{\pgfqpoint{2.852943in}{2.797748in}}%
\pgfpathlineto{\pgfqpoint{2.855462in}{2.793768in}}%
\pgfpathlineto{\pgfqpoint{2.857981in}{2.778669in}}%
\pgfpathlineto{\pgfqpoint{2.860499in}{2.776037in}}%
\pgfpathlineto{\pgfqpoint{2.863018in}{2.735423in}}%
\pgfpathlineto{\pgfqpoint{2.865537in}{2.721335in}}%
\pgfpathlineto{\pgfqpoint{2.868056in}{2.720466in}}%
\pgfpathlineto{\pgfqpoint{2.870574in}{2.694518in}}%
\pgfpathlineto{\pgfqpoint{2.873093in}{2.681202in}}%
\pgfpathlineto{\pgfqpoint{2.875612in}{2.670360in}}%
\pgfpathlineto{\pgfqpoint{2.878131in}{2.649590in}}%
\pgfpathlineto{\pgfqpoint{2.880649in}{2.647904in}}%
\pgfpathlineto{\pgfqpoint{2.883168in}{2.643279in}}%
\pgfpathlineto{\pgfqpoint{2.885687in}{2.682181in}}%
\pgfpathlineto{\pgfqpoint{2.888206in}{2.681363in}}%
\pgfpathlineto{\pgfqpoint{2.890724in}{2.654587in}}%
\pgfpathlineto{\pgfqpoint{2.893243in}{2.680451in}}%
\pgfpathlineto{\pgfqpoint{2.895762in}{2.678710in}}%
\pgfpathlineto{\pgfqpoint{2.898281in}{2.670847in}}%
\pgfpathlineto{\pgfqpoint{2.903318in}{2.623570in}}%
\pgfpathlineto{\pgfqpoint{2.905837in}{2.634114in}}%
\pgfpathlineto{\pgfqpoint{2.908356in}{2.652098in}}%
\pgfpathlineto{\pgfqpoint{2.910874in}{2.667724in}}%
\pgfpathlineto{\pgfqpoint{2.913393in}{2.647605in}}%
\pgfpathlineto{\pgfqpoint{2.915912in}{2.674031in}}%
\pgfpathlineto{\pgfqpoint{2.918431in}{2.657388in}}%
\pgfpathlineto{\pgfqpoint{2.920949in}{2.621450in}}%
\pgfpathlineto{\pgfqpoint{2.923468in}{2.650802in}}%
\pgfpathlineto{\pgfqpoint{2.925987in}{2.666624in}}%
\pgfpathlineto{\pgfqpoint{2.928506in}{2.668774in}}%
\pgfpathlineto{\pgfqpoint{2.931024in}{2.678131in}}%
\pgfpathlineto{\pgfqpoint{2.933543in}{2.695099in}}%
\pgfpathlineto{\pgfqpoint{2.936062in}{2.681901in}}%
\pgfpathlineto{\pgfqpoint{2.938581in}{2.696634in}}%
\pgfpathlineto{\pgfqpoint{2.941099in}{2.682228in}}%
\pgfpathlineto{\pgfqpoint{2.943618in}{2.705863in}}%
\pgfpathlineto{\pgfqpoint{2.946137in}{2.698425in}}%
\pgfpathlineto{\pgfqpoint{2.948656in}{2.705559in}}%
\pgfpathlineto{\pgfqpoint{2.951174in}{2.723475in}}%
\pgfpathlineto{\pgfqpoint{2.953693in}{2.736757in}}%
\pgfpathlineto{\pgfqpoint{2.956212in}{2.759308in}}%
\pgfpathlineto{\pgfqpoint{2.958731in}{2.740876in}}%
\pgfpathlineto{\pgfqpoint{2.961249in}{2.734921in}}%
\pgfpathlineto{\pgfqpoint{2.963768in}{2.750294in}}%
\pgfpathlineto{\pgfqpoint{2.966287in}{2.736159in}}%
\pgfpathlineto{\pgfqpoint{2.968806in}{2.716434in}}%
\pgfpathlineto{\pgfqpoint{2.971324in}{2.725460in}}%
\pgfpathlineto{\pgfqpoint{2.973843in}{2.755130in}}%
\pgfpathlineto{\pgfqpoint{2.976362in}{2.763939in}}%
\pgfpathlineto{\pgfqpoint{2.978881in}{2.782973in}}%
\pgfpathlineto{\pgfqpoint{2.981399in}{2.764329in}}%
\pgfpathlineto{\pgfqpoint{2.983918in}{2.776591in}}%
\pgfpathlineto{\pgfqpoint{2.986437in}{2.791502in}}%
\pgfpathlineto{\pgfqpoint{2.988956in}{2.794523in}}%
\pgfpathlineto{\pgfqpoint{2.991474in}{2.805033in}}%
\pgfpathlineto{\pgfqpoint{2.993993in}{2.809852in}}%
\pgfpathlineto{\pgfqpoint{2.996512in}{2.825837in}}%
\pgfpathlineto{\pgfqpoint{2.999031in}{2.798347in}}%
\pgfpathlineto{\pgfqpoint{3.001549in}{2.809586in}}%
\pgfpathlineto{\pgfqpoint{3.004068in}{2.826431in}}%
\pgfpathlineto{\pgfqpoint{3.006587in}{2.838366in}}%
\pgfpathlineto{\pgfqpoint{3.009106in}{2.868913in}}%
\pgfpathlineto{\pgfqpoint{3.011624in}{2.849757in}}%
\pgfpathlineto{\pgfqpoint{3.014143in}{2.850070in}}%
\pgfpathlineto{\pgfqpoint{3.016662in}{2.847259in}}%
\pgfpathlineto{\pgfqpoint{3.019181in}{2.873342in}}%
\pgfpathlineto{\pgfqpoint{3.021699in}{2.877086in}}%
\pgfpathlineto{\pgfqpoint{3.024218in}{2.868690in}}%
\pgfpathlineto{\pgfqpoint{3.026737in}{2.866943in}}%
\pgfpathlineto{\pgfqpoint{3.029256in}{2.855360in}}%
\pgfpathlineto{\pgfqpoint{3.031774in}{2.866314in}}%
\pgfpathlineto{\pgfqpoint{3.034293in}{2.861802in}}%
\pgfpathlineto{\pgfqpoint{3.036812in}{2.865882in}}%
\pgfpathlineto{\pgfqpoint{3.039331in}{2.854878in}}%
\pgfpathlineto{\pgfqpoint{3.041849in}{2.863493in}}%
\pgfpathlineto{\pgfqpoint{3.044368in}{2.886334in}}%
\pgfpathlineto{\pgfqpoint{3.046887in}{2.867731in}}%
\pgfpathlineto{\pgfqpoint{3.051924in}{2.852352in}}%
\pgfpathlineto{\pgfqpoint{3.054443in}{2.848734in}}%
\pgfpathlineto{\pgfqpoint{3.056962in}{2.831379in}}%
\pgfpathlineto{\pgfqpoint{3.059481in}{2.821527in}}%
\pgfpathlineto{\pgfqpoint{3.061999in}{2.786739in}}%
\pgfpathlineto{\pgfqpoint{3.064518in}{2.785533in}}%
\pgfpathlineto{\pgfqpoint{3.067037in}{2.802857in}}%
\pgfpathlineto{\pgfqpoint{3.069556in}{2.787253in}}%
\pgfpathlineto{\pgfqpoint{3.072074in}{2.783971in}}%
\pgfpathlineto{\pgfqpoint{3.074593in}{2.764868in}}%
\pgfpathlineto{\pgfqpoint{3.077112in}{2.724929in}}%
\pgfpathlineto{\pgfqpoint{3.079631in}{2.728509in}}%
\pgfpathlineto{\pgfqpoint{3.082149in}{2.729170in}}%
\pgfpathlineto{\pgfqpoint{3.084668in}{2.730702in}}%
\pgfpathlineto{\pgfqpoint{3.087187in}{2.743398in}}%
\pgfpathlineto{\pgfqpoint{3.089706in}{2.761883in}}%
\pgfpathlineto{\pgfqpoint{3.092224in}{2.749535in}}%
\pgfpathlineto{\pgfqpoint{3.094743in}{2.758160in}}%
\pgfpathlineto{\pgfqpoint{3.097262in}{2.775126in}}%
\pgfpathlineto{\pgfqpoint{3.099781in}{2.778229in}}%
\pgfpathlineto{\pgfqpoint{3.102299in}{2.741695in}}%
\pgfpathlineto{\pgfqpoint{3.104818in}{2.771814in}}%
\pgfpathlineto{\pgfqpoint{3.107337in}{2.808513in}}%
\pgfpathlineto{\pgfqpoint{3.109856in}{2.853393in}}%
\pgfpathlineto{\pgfqpoint{3.112374in}{2.868399in}}%
\pgfpathlineto{\pgfqpoint{3.114893in}{2.822526in}}%
\pgfpathlineto{\pgfqpoint{3.117412in}{2.793975in}}%
\pgfpathlineto{\pgfqpoint{3.119931in}{2.753122in}}%
\pgfpathlineto{\pgfqpoint{3.122449in}{2.731072in}}%
\pgfpathlineto{\pgfqpoint{3.124968in}{2.737353in}}%
\pgfpathlineto{\pgfqpoint{3.127487in}{2.753120in}}%
\pgfpathlineto{\pgfqpoint{3.130006in}{2.746115in}}%
\pgfpathlineto{\pgfqpoint{3.132524in}{2.767604in}}%
\pgfpathlineto{\pgfqpoint{3.135043in}{2.783354in}}%
\pgfpathlineto{\pgfqpoint{3.137562in}{2.796168in}}%
\pgfpathlineto{\pgfqpoint{3.140081in}{2.797482in}}%
\pgfpathlineto{\pgfqpoint{3.142599in}{2.772462in}}%
\pgfpathlineto{\pgfqpoint{3.145118in}{2.801746in}}%
\pgfpathlineto{\pgfqpoint{3.147637in}{2.789629in}}%
\pgfpathlineto{\pgfqpoint{3.150156in}{2.783982in}}%
\pgfpathlineto{\pgfqpoint{3.152674in}{2.777033in}}%
\pgfpathlineto{\pgfqpoint{3.155193in}{2.755787in}}%
\pgfpathlineto{\pgfqpoint{3.157712in}{2.751679in}}%
\pgfpathlineto{\pgfqpoint{3.160231in}{2.739049in}}%
\pgfpathlineto{\pgfqpoint{3.162749in}{2.755453in}}%
\pgfpathlineto{\pgfqpoint{3.165268in}{2.743588in}}%
\pgfpathlineto{\pgfqpoint{3.167787in}{2.750324in}}%
\pgfpathlineto{\pgfqpoint{3.170306in}{2.738801in}}%
\pgfpathlineto{\pgfqpoint{3.172824in}{2.738159in}}%
\pgfpathlineto{\pgfqpoint{3.175343in}{2.737292in}}%
\pgfpathlineto{\pgfqpoint{3.177862in}{2.744288in}}%
\pgfpathlineto{\pgfqpoint{3.180381in}{2.778304in}}%
\pgfpathlineto{\pgfqpoint{3.182899in}{2.758390in}}%
\pgfpathlineto{\pgfqpoint{3.185418in}{2.771364in}}%
\pgfpathlineto{\pgfqpoint{3.187937in}{2.774553in}}%
\pgfpathlineto{\pgfqpoint{3.190456in}{2.765472in}}%
\pgfpathlineto{\pgfqpoint{3.192974in}{2.774708in}}%
\pgfpathlineto{\pgfqpoint{3.195493in}{2.743634in}}%
\pgfpathlineto{\pgfqpoint{3.198012in}{2.749862in}}%
\pgfpathlineto{\pgfqpoint{3.200531in}{2.731513in}}%
\pgfpathlineto{\pgfqpoint{3.203049in}{2.741702in}}%
\pgfpathlineto{\pgfqpoint{3.205568in}{2.737126in}}%
\pgfpathlineto{\pgfqpoint{3.208087in}{2.753714in}}%
\pgfpathlineto{\pgfqpoint{3.210606in}{2.775565in}}%
\pgfpathlineto{\pgfqpoint{3.213124in}{2.768311in}}%
\pgfpathlineto{\pgfqpoint{3.215643in}{2.778518in}}%
\pgfpathlineto{\pgfqpoint{3.218162in}{2.802731in}}%
\pgfpathlineto{\pgfqpoint{3.220681in}{2.790537in}}%
\pgfpathlineto{\pgfqpoint{3.223199in}{2.751918in}}%
\pgfpathlineto{\pgfqpoint{3.225718in}{2.747424in}}%
\pgfpathlineto{\pgfqpoint{3.228237in}{2.758775in}}%
\pgfpathlineto{\pgfqpoint{3.230756in}{2.733907in}}%
\pgfpathlineto{\pgfqpoint{3.233274in}{2.749022in}}%
\pgfpathlineto{\pgfqpoint{3.235793in}{2.738539in}}%
\pgfpathlineto{\pgfqpoint{3.238312in}{2.725991in}}%
\pgfpathlineto{\pgfqpoint{3.240831in}{2.719874in}}%
\pgfpathlineto{\pgfqpoint{3.243349in}{2.721172in}}%
\pgfpathlineto{\pgfqpoint{3.245868in}{2.694296in}}%
\pgfpathlineto{\pgfqpoint{3.248387in}{2.677934in}}%
\pgfpathlineto{\pgfqpoint{3.250906in}{2.696573in}}%
\pgfpathlineto{\pgfqpoint{3.253424in}{2.699235in}}%
\pgfpathlineto{\pgfqpoint{3.255943in}{2.692151in}}%
\pgfpathlineto{\pgfqpoint{3.258462in}{2.674641in}}%
\pgfpathlineto{\pgfqpoint{3.260981in}{2.667985in}}%
\pgfpathlineto{\pgfqpoint{3.263499in}{2.666939in}}%
\pgfpathlineto{\pgfqpoint{3.266018in}{2.691302in}}%
\pgfpathlineto{\pgfqpoint{3.268537in}{2.670342in}}%
\pgfpathlineto{\pgfqpoint{3.271056in}{2.674350in}}%
\pgfpathlineto{\pgfqpoint{3.273574in}{2.679545in}}%
\pgfpathlineto{\pgfqpoint{3.276093in}{2.686583in}}%
\pgfpathlineto{\pgfqpoint{3.278612in}{2.703187in}}%
\pgfpathlineto{\pgfqpoint{3.281131in}{2.687251in}}%
\pgfpathlineto{\pgfqpoint{3.283649in}{2.707764in}}%
\pgfpathlineto{\pgfqpoint{3.286168in}{2.731631in}}%
\pgfpathlineto{\pgfqpoint{3.288687in}{2.719130in}}%
\pgfpathlineto{\pgfqpoint{3.291206in}{2.677461in}}%
\pgfpathlineto{\pgfqpoint{3.293724in}{2.668093in}}%
\pgfpathlineto{\pgfqpoint{3.296243in}{2.707319in}}%
\pgfpathlineto{\pgfqpoint{3.298762in}{2.737660in}}%
\pgfpathlineto{\pgfqpoint{3.301281in}{2.725333in}}%
\pgfpathlineto{\pgfqpoint{3.303799in}{2.738880in}}%
\pgfpathlineto{\pgfqpoint{3.306318in}{2.755385in}}%
\pgfpathlineto{\pgfqpoint{3.308837in}{2.766094in}}%
\pgfpathlineto{\pgfqpoint{3.311356in}{2.736753in}}%
\pgfpathlineto{\pgfqpoint{3.313874in}{2.725587in}}%
\pgfpathlineto{\pgfqpoint{3.316393in}{2.739612in}}%
\pgfpathlineto{\pgfqpoint{3.318912in}{2.756872in}}%
\pgfpathlineto{\pgfqpoint{3.321431in}{2.722953in}}%
\pgfpathlineto{\pgfqpoint{3.323949in}{2.731281in}}%
\pgfpathlineto{\pgfqpoint{3.326468in}{2.741210in}}%
\pgfpathlineto{\pgfqpoint{3.328987in}{2.738623in}}%
\pgfpathlineto{\pgfqpoint{3.331506in}{2.746030in}}%
\pgfpathlineto{\pgfqpoint{3.334024in}{2.749148in}}%
\pgfpathlineto{\pgfqpoint{3.336543in}{2.746956in}}%
\pgfpathlineto{\pgfqpoint{3.339062in}{2.732248in}}%
\pgfpathlineto{\pgfqpoint{3.341581in}{2.741138in}}%
\pgfpathlineto{\pgfqpoint{3.344099in}{2.727440in}}%
\pgfpathlineto{\pgfqpoint{3.346618in}{2.701550in}}%
\pgfpathlineto{\pgfqpoint{3.349137in}{2.672328in}}%
\pgfpathlineto{\pgfqpoint{3.351656in}{2.690196in}}%
\pgfpathlineto{\pgfqpoint{3.354174in}{2.686664in}}%
\pgfpathlineto{\pgfqpoint{3.356693in}{2.679327in}}%
\pgfpathlineto{\pgfqpoint{3.359212in}{2.654760in}}%
\pgfpathlineto{\pgfqpoint{3.361731in}{2.649594in}}%
\pgfpathlineto{\pgfqpoint{3.364249in}{2.661926in}}%
\pgfpathlineto{\pgfqpoint{3.366768in}{2.666166in}}%
\pgfpathlineto{\pgfqpoint{3.369287in}{2.667628in}}%
\pgfpathlineto{\pgfqpoint{3.371806in}{2.649101in}}%
\pgfpathlineto{\pgfqpoint{3.374324in}{2.654089in}}%
\pgfpathlineto{\pgfqpoint{3.376843in}{2.628795in}}%
\pgfpathlineto{\pgfqpoint{3.379362in}{2.645823in}}%
\pgfpathlineto{\pgfqpoint{3.381881in}{2.644646in}}%
\pgfpathlineto{\pgfqpoint{3.384399in}{2.630931in}}%
\pgfpathlineto{\pgfqpoint{3.386918in}{2.671281in}}%
\pgfpathlineto{\pgfqpoint{3.389437in}{2.664020in}}%
\pgfpathlineto{\pgfqpoint{3.391956in}{2.673001in}}%
\pgfpathlineto{\pgfqpoint{3.394474in}{2.670839in}}%
\pgfpathlineto{\pgfqpoint{3.399512in}{2.669205in}}%
\pgfpathlineto{\pgfqpoint{3.402031in}{2.667128in}}%
\pgfpathlineto{\pgfqpoint{3.404549in}{2.676535in}}%
\pgfpathlineto{\pgfqpoint{3.407068in}{2.703520in}}%
\pgfpathlineto{\pgfqpoint{3.409587in}{2.704446in}}%
\pgfpathlineto{\pgfqpoint{3.412106in}{2.695379in}}%
\pgfpathlineto{\pgfqpoint{3.414624in}{2.699389in}}%
\pgfpathlineto{\pgfqpoint{3.417143in}{2.677584in}}%
\pgfpathlineto{\pgfqpoint{3.419662in}{2.708824in}}%
\pgfpathlineto{\pgfqpoint{3.422181in}{2.762565in}}%
\pgfpathlineto{\pgfqpoint{3.424699in}{2.751305in}}%
\pgfpathlineto{\pgfqpoint{3.427218in}{2.733894in}}%
\pgfpathlineto{\pgfqpoint{3.429737in}{2.757209in}}%
\pgfpathlineto{\pgfqpoint{3.432256in}{2.793354in}}%
\pgfpathlineto{\pgfqpoint{3.434774in}{2.781041in}}%
\pgfpathlineto{\pgfqpoint{3.437293in}{2.762065in}}%
\pgfpathlineto{\pgfqpoint{3.439812in}{2.757760in}}%
\pgfpathlineto{\pgfqpoint{3.442331in}{2.726874in}}%
\pgfpathlineto{\pgfqpoint{3.444849in}{2.738383in}}%
\pgfpathlineto{\pgfqpoint{3.447368in}{2.760799in}}%
\pgfpathlineto{\pgfqpoint{3.449887in}{2.758285in}}%
\pgfpathlineto{\pgfqpoint{3.452406in}{2.783281in}}%
\pgfpathlineto{\pgfqpoint{3.454924in}{2.800890in}}%
\pgfpathlineto{\pgfqpoint{3.457443in}{2.768027in}}%
\pgfpathlineto{\pgfqpoint{3.459962in}{2.759758in}}%
\pgfpathlineto{\pgfqpoint{3.462481in}{2.755317in}}%
\pgfpathlineto{\pgfqpoint{3.464999in}{2.755347in}}%
\pgfpathlineto{\pgfqpoint{3.467518in}{2.764799in}}%
\pgfpathlineto{\pgfqpoint{3.470037in}{2.743406in}}%
\pgfpathlineto{\pgfqpoint{3.472556in}{2.774238in}}%
\pgfpathlineto{\pgfqpoint{3.475074in}{2.789073in}}%
\pgfpathlineto{\pgfqpoint{3.477593in}{2.839803in}}%
\pgfpathlineto{\pgfqpoint{3.480112in}{2.822404in}}%
\pgfpathlineto{\pgfqpoint{3.482631in}{2.829224in}}%
\pgfpathlineto{\pgfqpoint{3.485149in}{2.826308in}}%
\pgfpathlineto{\pgfqpoint{3.487668in}{2.839808in}}%
\pgfpathlineto{\pgfqpoint{3.490187in}{2.821888in}}%
\pgfpathlineto{\pgfqpoint{3.492706in}{2.840519in}}%
\pgfpathlineto{\pgfqpoint{3.495224in}{2.841358in}}%
\pgfpathlineto{\pgfqpoint{3.500262in}{2.878352in}}%
\pgfpathlineto{\pgfqpoint{3.502781in}{2.868812in}}%
\pgfpathlineto{\pgfqpoint{3.505299in}{2.844432in}}%
\pgfpathlineto{\pgfqpoint{3.507818in}{2.875608in}}%
\pgfpathlineto{\pgfqpoint{3.510337in}{2.905226in}}%
\pgfpathlineto{\pgfqpoint{3.512856in}{2.899535in}}%
\pgfpathlineto{\pgfqpoint{3.515374in}{2.879792in}}%
\pgfpathlineto{\pgfqpoint{3.517893in}{2.916719in}}%
\pgfpathlineto{\pgfqpoint{3.520412in}{2.912973in}}%
\pgfpathlineto{\pgfqpoint{3.522931in}{2.916992in}}%
\pgfpathlineto{\pgfqpoint{3.525449in}{2.890502in}}%
\pgfpathlineto{\pgfqpoint{3.527968in}{2.934006in}}%
\pgfpathlineto{\pgfqpoint{3.530487in}{2.950268in}}%
\pgfpathlineto{\pgfqpoint{3.533006in}{2.923434in}}%
\pgfpathlineto{\pgfqpoint{3.535524in}{2.912542in}}%
\pgfpathlineto{\pgfqpoint{3.538043in}{2.930342in}}%
\pgfpathlineto{\pgfqpoint{3.540562in}{2.898552in}}%
\pgfpathlineto{\pgfqpoint{3.543081in}{2.862836in}}%
\pgfpathlineto{\pgfqpoint{3.545599in}{2.865870in}}%
\pgfpathlineto{\pgfqpoint{3.548118in}{2.861899in}}%
\pgfpathlineto{\pgfqpoint{3.550637in}{2.850325in}}%
\pgfpathlineto{\pgfqpoint{3.553156in}{2.852219in}}%
\pgfpathlineto{\pgfqpoint{3.555674in}{2.847953in}}%
\pgfpathlineto{\pgfqpoint{3.558193in}{2.838648in}}%
\pgfpathlineto{\pgfqpoint{3.560712in}{2.852853in}}%
\pgfpathlineto{\pgfqpoint{3.563231in}{2.839239in}}%
\pgfpathlineto{\pgfqpoint{3.565749in}{2.850334in}}%
\pgfpathlineto{\pgfqpoint{3.568268in}{2.828313in}}%
\pgfpathlineto{\pgfqpoint{3.570787in}{2.801818in}}%
\pgfpathlineto{\pgfqpoint{3.573306in}{2.828149in}}%
\pgfpathlineto{\pgfqpoint{3.575824in}{2.834371in}}%
\pgfpathlineto{\pgfqpoint{3.580862in}{2.799731in}}%
\pgfpathlineto{\pgfqpoint{3.583381in}{2.818805in}}%
\pgfpathlineto{\pgfqpoint{3.585899in}{2.822681in}}%
\pgfpathlineto{\pgfqpoint{3.588418in}{2.802377in}}%
\pgfpathlineto{\pgfqpoint{3.590937in}{2.828612in}}%
\pgfpathlineto{\pgfqpoint{3.593456in}{2.842475in}}%
\pgfpathlineto{\pgfqpoint{3.595974in}{2.841680in}}%
\pgfpathlineto{\pgfqpoint{3.598493in}{2.834594in}}%
\pgfpathlineto{\pgfqpoint{3.601012in}{2.804731in}}%
\pgfpathlineto{\pgfqpoint{3.603531in}{2.816082in}}%
\pgfpathlineto{\pgfqpoint{3.606049in}{2.845285in}}%
\pgfpathlineto{\pgfqpoint{3.608568in}{2.837482in}}%
\pgfpathlineto{\pgfqpoint{3.611087in}{2.881692in}}%
\pgfpathlineto{\pgfqpoint{3.613606in}{2.886832in}}%
\pgfpathlineto{\pgfqpoint{3.616124in}{2.925493in}}%
\pgfpathlineto{\pgfqpoint{3.618643in}{2.943326in}}%
\pgfpathlineto{\pgfqpoint{3.621162in}{2.926672in}}%
\pgfpathlineto{\pgfqpoint{3.623681in}{2.935806in}}%
\pgfpathlineto{\pgfqpoint{3.626199in}{2.957212in}}%
\pgfpathlineto{\pgfqpoint{3.628718in}{2.938547in}}%
\pgfpathlineto{\pgfqpoint{3.631237in}{2.939436in}}%
\pgfpathlineto{\pgfqpoint{3.633756in}{2.935677in}}%
\pgfpathlineto{\pgfqpoint{3.636274in}{2.916044in}}%
\pgfpathlineto{\pgfqpoint{3.638793in}{2.936742in}}%
\pgfpathlineto{\pgfqpoint{3.641312in}{2.942730in}}%
\pgfpathlineto{\pgfqpoint{3.643831in}{2.930327in}}%
\pgfpathlineto{\pgfqpoint{3.646349in}{2.928415in}}%
\pgfpathlineto{\pgfqpoint{3.648868in}{2.936500in}}%
\pgfpathlineto{\pgfqpoint{3.651387in}{2.929988in}}%
\pgfpathlineto{\pgfqpoint{3.653906in}{2.908067in}}%
\pgfpathlineto{\pgfqpoint{3.656424in}{2.875555in}}%
\pgfpathlineto{\pgfqpoint{3.658943in}{2.885387in}}%
\pgfpathlineto{\pgfqpoint{3.661462in}{2.911215in}}%
\pgfpathlineto{\pgfqpoint{3.663981in}{2.927133in}}%
\pgfpathlineto{\pgfqpoint{3.666499in}{2.913875in}}%
\pgfpathlineto{\pgfqpoint{3.669018in}{2.905384in}}%
\pgfpathlineto{\pgfqpoint{3.671537in}{2.900624in}}%
\pgfpathlineto{\pgfqpoint{3.674056in}{2.901253in}}%
\pgfpathlineto{\pgfqpoint{3.676574in}{2.926770in}}%
\pgfpathlineto{\pgfqpoint{3.679093in}{2.936850in}}%
\pgfpathlineto{\pgfqpoint{3.681612in}{2.976604in}}%
\pgfpathlineto{\pgfqpoint{3.684131in}{2.960069in}}%
\pgfpathlineto{\pgfqpoint{3.686649in}{2.931316in}}%
\pgfpathlineto{\pgfqpoint{3.689168in}{2.960693in}}%
\pgfpathlineto{\pgfqpoint{3.691687in}{2.975088in}}%
\pgfpathlineto{\pgfqpoint{3.694206in}{2.973888in}}%
\pgfpathlineto{\pgfqpoint{3.696724in}{2.964346in}}%
\pgfpathlineto{\pgfqpoint{3.699243in}{2.986259in}}%
\pgfpathlineto{\pgfqpoint{3.701762in}{2.985972in}}%
\pgfpathlineto{\pgfqpoint{3.704281in}{2.996283in}}%
\pgfpathlineto{\pgfqpoint{3.706799in}{2.971972in}}%
\pgfpathlineto{\pgfqpoint{3.709318in}{3.002772in}}%
\pgfpathlineto{\pgfqpoint{3.711837in}{2.985433in}}%
\pgfpathlineto{\pgfqpoint{3.714356in}{2.969484in}}%
\pgfpathlineto{\pgfqpoint{3.716874in}{2.983323in}}%
\pgfpathlineto{\pgfqpoint{3.719393in}{2.968463in}}%
\pgfpathlineto{\pgfqpoint{3.721912in}{2.981427in}}%
\pgfpathlineto{\pgfqpoint{3.724431in}{2.941584in}}%
\pgfpathlineto{\pgfqpoint{3.726949in}{2.946780in}}%
\pgfpathlineto{\pgfqpoint{3.729468in}{2.947545in}}%
\pgfpathlineto{\pgfqpoint{3.731987in}{2.913294in}}%
\pgfpathlineto{\pgfqpoint{3.734506in}{2.946759in}}%
\pgfpathlineto{\pgfqpoint{3.737024in}{2.948658in}}%
\pgfpathlineto{\pgfqpoint{3.739543in}{2.944662in}}%
\pgfpathlineto{\pgfqpoint{3.742062in}{2.948729in}}%
\pgfpathlineto{\pgfqpoint{3.744581in}{2.938949in}}%
\pgfpathlineto{\pgfqpoint{3.747099in}{2.963784in}}%
\pgfpathlineto{\pgfqpoint{3.749618in}{2.940122in}}%
\pgfpathlineto{\pgfqpoint{3.752137in}{2.909670in}}%
\pgfpathlineto{\pgfqpoint{3.754656in}{2.888766in}}%
\pgfpathlineto{\pgfqpoint{3.757174in}{2.865480in}}%
\pgfpathlineto{\pgfqpoint{3.759693in}{2.862313in}}%
\pgfpathlineto{\pgfqpoint{3.762212in}{2.861182in}}%
\pgfpathlineto{\pgfqpoint{3.764731in}{2.862226in}}%
\pgfpathlineto{\pgfqpoint{3.767249in}{2.865314in}}%
\pgfpathlineto{\pgfqpoint{3.769768in}{2.845690in}}%
\pgfpathlineto{\pgfqpoint{3.772287in}{2.843066in}}%
\pgfpathlineto{\pgfqpoint{3.774806in}{2.864865in}}%
\pgfpathlineto{\pgfqpoint{3.777324in}{2.842262in}}%
\pgfpathlineto{\pgfqpoint{3.779843in}{2.858439in}}%
\pgfpathlineto{\pgfqpoint{3.782362in}{2.849410in}}%
\pgfpathlineto{\pgfqpoint{3.784881in}{2.830235in}}%
\pgfpathlineto{\pgfqpoint{3.787399in}{2.793456in}}%
\pgfpathlineto{\pgfqpoint{3.789918in}{2.771582in}}%
\pgfpathlineto{\pgfqpoint{3.792437in}{2.742505in}}%
\pgfpathlineto{\pgfqpoint{3.797474in}{2.770033in}}%
\pgfpathlineto{\pgfqpoint{3.799993in}{2.786181in}}%
\pgfpathlineto{\pgfqpoint{3.802512in}{2.781714in}}%
\pgfpathlineto{\pgfqpoint{3.805031in}{2.780842in}}%
\pgfpathlineto{\pgfqpoint{3.807549in}{2.772302in}}%
\pgfpathlineto{\pgfqpoint{3.810068in}{2.787895in}}%
\pgfpathlineto{\pgfqpoint{3.812587in}{2.782921in}}%
\pgfpathlineto{\pgfqpoint{3.815106in}{2.776814in}}%
\pgfpathlineto{\pgfqpoint{3.817624in}{2.727600in}}%
\pgfpathlineto{\pgfqpoint{3.820143in}{2.735912in}}%
\pgfpathlineto{\pgfqpoint{3.822662in}{2.745253in}}%
\pgfpathlineto{\pgfqpoint{3.825181in}{2.722923in}}%
\pgfpathlineto{\pgfqpoint{3.827699in}{2.730362in}}%
\pgfpathlineto{\pgfqpoint{3.830218in}{2.723871in}}%
\pgfpathlineto{\pgfqpoint{3.832737in}{2.748483in}}%
\pgfpathlineto{\pgfqpoint{3.835256in}{2.787319in}}%
\pgfpathlineto{\pgfqpoint{3.837774in}{2.811174in}}%
\pgfpathlineto{\pgfqpoint{3.840293in}{2.773143in}}%
\pgfpathlineto{\pgfqpoint{3.842812in}{2.732795in}}%
\pgfpathlineto{\pgfqpoint{3.845331in}{2.717601in}}%
\pgfpathlineto{\pgfqpoint{3.847849in}{2.748251in}}%
\pgfpathlineto{\pgfqpoint{3.850368in}{2.744122in}}%
\pgfpathlineto{\pgfqpoint{3.852887in}{2.696586in}}%
\pgfpathlineto{\pgfqpoint{3.855406in}{2.678401in}}%
\pgfpathlineto{\pgfqpoint{3.857924in}{2.690190in}}%
\pgfpathlineto{\pgfqpoint{3.860443in}{2.697203in}}%
\pgfpathlineto{\pgfqpoint{3.862962in}{2.677664in}}%
\pgfpathlineto{\pgfqpoint{3.865481in}{2.670753in}}%
\pgfpathlineto{\pgfqpoint{3.867999in}{2.681106in}}%
\pgfpathlineto{\pgfqpoint{3.870518in}{2.733722in}}%
\pgfpathlineto{\pgfqpoint{3.873037in}{2.710180in}}%
\pgfpathlineto{\pgfqpoint{3.875556in}{2.695847in}}%
\pgfpathlineto{\pgfqpoint{3.878074in}{2.703910in}}%
\pgfpathlineto{\pgfqpoint{3.880593in}{2.680215in}}%
\pgfpathlineto{\pgfqpoint{3.883112in}{2.704316in}}%
\pgfpathlineto{\pgfqpoint{3.885631in}{2.685300in}}%
\pgfpathlineto{\pgfqpoint{3.888149in}{2.667333in}}%
\pgfpathlineto{\pgfqpoint{3.890668in}{2.676877in}}%
\pgfpathlineto{\pgfqpoint{3.893187in}{2.684519in}}%
\pgfpathlineto{\pgfqpoint{3.895706in}{2.716135in}}%
\pgfpathlineto{\pgfqpoint{3.898224in}{2.704384in}}%
\pgfpathlineto{\pgfqpoint{3.900743in}{2.675421in}}%
\pgfpathlineto{\pgfqpoint{3.903262in}{2.640052in}}%
\pgfpathlineto{\pgfqpoint{3.905781in}{2.617904in}}%
\pgfpathlineto{\pgfqpoint{3.908299in}{2.638019in}}%
\pgfpathlineto{\pgfqpoint{3.910818in}{2.618514in}}%
\pgfpathlineto{\pgfqpoint{3.913337in}{2.587423in}}%
\pgfpathlineto{\pgfqpoint{3.918374in}{2.611496in}}%
\pgfpathlineto{\pgfqpoint{3.920893in}{2.602647in}}%
\pgfpathlineto{\pgfqpoint{3.923412in}{2.591074in}}%
\pgfpathlineto{\pgfqpoint{3.925931in}{2.602760in}}%
\pgfpathlineto{\pgfqpoint{3.928449in}{2.548733in}}%
\pgfpathlineto{\pgfqpoint{3.930968in}{2.531421in}}%
\pgfpathlineto{\pgfqpoint{3.933487in}{2.536701in}}%
\pgfpathlineto{\pgfqpoint{3.936006in}{2.533552in}}%
\pgfpathlineto{\pgfqpoint{3.938524in}{2.584147in}}%
\pgfpathlineto{\pgfqpoint{3.941043in}{2.592536in}}%
\pgfpathlineto{\pgfqpoint{3.943562in}{2.600231in}}%
\pgfpathlineto{\pgfqpoint{3.946081in}{2.627108in}}%
\pgfpathlineto{\pgfqpoint{3.948599in}{2.638132in}}%
\pgfpathlineto{\pgfqpoint{3.951118in}{2.616461in}}%
\pgfpathlineto{\pgfqpoint{3.953637in}{2.619422in}}%
\pgfpathlineto{\pgfqpoint{3.956156in}{2.649062in}}%
\pgfpathlineto{\pgfqpoint{3.958674in}{2.654614in}}%
\pgfpathlineto{\pgfqpoint{3.961193in}{2.674399in}}%
\pgfpathlineto{\pgfqpoint{3.963712in}{2.625501in}}%
\pgfpathlineto{\pgfqpoint{3.966231in}{2.607788in}}%
\pgfpathlineto{\pgfqpoint{3.968749in}{2.637310in}}%
\pgfpathlineto{\pgfqpoint{3.971268in}{2.618381in}}%
\pgfpathlineto{\pgfqpoint{3.973787in}{2.585582in}}%
\pgfpathlineto{\pgfqpoint{3.976306in}{2.589935in}}%
\pgfpathlineto{\pgfqpoint{3.978824in}{2.608200in}}%
\pgfpathlineto{\pgfqpoint{3.981343in}{2.659140in}}%
\pgfpathlineto{\pgfqpoint{3.983862in}{2.631672in}}%
\pgfpathlineto{\pgfqpoint{3.986381in}{2.646414in}}%
\pgfpathlineto{\pgfqpoint{3.988899in}{2.638138in}}%
\pgfpathlineto{\pgfqpoint{3.991418in}{2.614280in}}%
\pgfpathlineto{\pgfqpoint{3.993937in}{2.591624in}}%
\pgfpathlineto{\pgfqpoint{3.996456in}{2.590530in}}%
\pgfpathlineto{\pgfqpoint{3.998974in}{2.637872in}}%
\pgfpathlineto{\pgfqpoint{4.001493in}{2.643444in}}%
\pgfpathlineto{\pgfqpoint{4.004012in}{2.642372in}}%
\pgfpathlineto{\pgfqpoint{4.006531in}{2.636605in}}%
\pgfpathlineto{\pgfqpoint{4.009049in}{2.608506in}}%
\pgfpathlineto{\pgfqpoint{4.011568in}{2.618256in}}%
\pgfpathlineto{\pgfqpoint{4.014087in}{2.612145in}}%
\pgfpathlineto{\pgfqpoint{4.016606in}{2.611504in}}%
\pgfpathlineto{\pgfqpoint{4.019124in}{2.622401in}}%
\pgfpathlineto{\pgfqpoint{4.021643in}{2.632276in}}%
\pgfpathlineto{\pgfqpoint{4.024162in}{2.613750in}}%
\pgfpathlineto{\pgfqpoint{4.026681in}{2.609947in}}%
\pgfpathlineto{\pgfqpoint{4.029199in}{2.622941in}}%
\pgfpathlineto{\pgfqpoint{4.031718in}{2.650469in}}%
\pgfpathlineto{\pgfqpoint{4.034237in}{2.666060in}}%
\pgfpathlineto{\pgfqpoint{4.036756in}{2.698575in}}%
\pgfpathlineto{\pgfqpoint{4.039274in}{2.722493in}}%
\pgfpathlineto{\pgfqpoint{4.041793in}{2.706514in}}%
\pgfpathlineto{\pgfqpoint{4.044312in}{2.695095in}}%
\pgfpathlineto{\pgfqpoint{4.046831in}{2.744660in}}%
\pgfpathlineto{\pgfqpoint{4.049349in}{2.737539in}}%
\pgfpathlineto{\pgfqpoint{4.051868in}{2.711611in}}%
\pgfpathlineto{\pgfqpoint{4.054387in}{2.748222in}}%
\pgfpathlineto{\pgfqpoint{4.056906in}{2.768986in}}%
\pgfpathlineto{\pgfqpoint{4.059424in}{2.763371in}}%
\pgfpathlineto{\pgfqpoint{4.061943in}{2.739305in}}%
\pgfpathlineto{\pgfqpoint{4.064462in}{2.737735in}}%
\pgfpathlineto{\pgfqpoint{4.066981in}{2.728293in}}%
\pgfpathlineto{\pgfqpoint{4.069499in}{2.754420in}}%
\pgfpathlineto{\pgfqpoint{4.072018in}{2.731743in}}%
\pgfpathlineto{\pgfqpoint{4.074537in}{2.750981in}}%
\pgfpathlineto{\pgfqpoint{4.077056in}{2.746482in}}%
\pgfpathlineto{\pgfqpoint{4.079574in}{2.739589in}}%
\pgfpathlineto{\pgfqpoint{4.082093in}{2.713988in}}%
\pgfpathlineto{\pgfqpoint{4.084612in}{2.720720in}}%
\pgfpathlineto{\pgfqpoint{4.087131in}{2.715143in}}%
\pgfpathlineto{\pgfqpoint{4.089649in}{2.737174in}}%
\pgfpathlineto{\pgfqpoint{4.092168in}{2.743322in}}%
\pgfpathlineto{\pgfqpoint{4.094687in}{2.760142in}}%
\pgfpathlineto{\pgfqpoint{4.097206in}{2.739801in}}%
\pgfpathlineto{\pgfqpoint{4.099724in}{2.760918in}}%
\pgfpathlineto{\pgfqpoint{4.102243in}{2.788610in}}%
\pgfpathlineto{\pgfqpoint{4.104762in}{2.796234in}}%
\pgfpathlineto{\pgfqpoint{4.107281in}{2.831445in}}%
\pgfpathlineto{\pgfqpoint{4.109799in}{2.835577in}}%
\pgfpathlineto{\pgfqpoint{4.112318in}{2.851443in}}%
\pgfpathlineto{\pgfqpoint{4.114837in}{2.845039in}}%
\pgfpathlineto{\pgfqpoint{4.117356in}{2.840022in}}%
\pgfpathlineto{\pgfqpoint{4.119874in}{2.824034in}}%
\pgfpathlineto{\pgfqpoint{4.122393in}{2.827130in}}%
\pgfpathlineto{\pgfqpoint{4.124912in}{2.810690in}}%
\pgfpathlineto{\pgfqpoint{4.127431in}{2.786301in}}%
\pgfpathlineto{\pgfqpoint{4.129949in}{2.788212in}}%
\pgfpathlineto{\pgfqpoint{4.132468in}{2.806913in}}%
\pgfpathlineto{\pgfqpoint{4.134987in}{2.829139in}}%
\pgfpathlineto{\pgfqpoint{4.137506in}{2.814326in}}%
\pgfpathlineto{\pgfqpoint{4.140024in}{2.790154in}}%
\pgfpathlineto{\pgfqpoint{4.142543in}{2.791720in}}%
\pgfpathlineto{\pgfqpoint{4.145062in}{2.783221in}}%
\pgfpathlineto{\pgfqpoint{4.150099in}{2.792193in}}%
\pgfpathlineto{\pgfqpoint{4.152618in}{2.774985in}}%
\pgfpathlineto{\pgfqpoint{4.155137in}{2.795353in}}%
\pgfpathlineto{\pgfqpoint{4.157656in}{2.836318in}}%
\pgfpathlineto{\pgfqpoint{4.160174in}{2.846519in}}%
\pgfpathlineto{\pgfqpoint{4.162693in}{2.872451in}}%
\pgfpathlineto{\pgfqpoint{4.165212in}{2.840507in}}%
\pgfpathlineto{\pgfqpoint{4.167731in}{2.862850in}}%
\pgfpathlineto{\pgfqpoint{4.170249in}{2.869978in}}%
\pgfpathlineto{\pgfqpoint{4.172768in}{2.875912in}}%
\pgfpathlineto{\pgfqpoint{4.175287in}{2.878562in}}%
\pgfpathlineto{\pgfqpoint{4.177806in}{2.882790in}}%
\pgfpathlineto{\pgfqpoint{4.180324in}{2.846527in}}%
\pgfpathlineto{\pgfqpoint{4.182843in}{2.851045in}}%
\pgfpathlineto{\pgfqpoint{4.185362in}{2.830093in}}%
\pgfpathlineto{\pgfqpoint{4.187881in}{2.818711in}}%
\pgfpathlineto{\pgfqpoint{4.190399in}{2.827209in}}%
\pgfpathlineto{\pgfqpoint{4.192918in}{2.804027in}}%
\pgfpathlineto{\pgfqpoint{4.195437in}{2.778156in}}%
\pgfpathlineto{\pgfqpoint{4.197956in}{2.767287in}}%
\pgfpathlineto{\pgfqpoint{4.200474in}{2.751327in}}%
\pgfpathlineto{\pgfqpoint{4.202993in}{2.762091in}}%
\pgfpathlineto{\pgfqpoint{4.205512in}{2.793548in}}%
\pgfpathlineto{\pgfqpoint{4.208031in}{2.817839in}}%
\pgfpathlineto{\pgfqpoint{4.210549in}{2.856069in}}%
\pgfpathlineto{\pgfqpoint{4.213068in}{2.862760in}}%
\pgfpathlineto{\pgfqpoint{4.215587in}{2.844808in}}%
\pgfpathlineto{\pgfqpoint{4.218106in}{2.851487in}}%
\pgfpathlineto{\pgfqpoint{4.220624in}{2.846056in}}%
\pgfpathlineto{\pgfqpoint{4.223143in}{2.874380in}}%
\pgfpathlineto{\pgfqpoint{4.225662in}{2.881052in}}%
\pgfpathlineto{\pgfqpoint{4.228181in}{2.904612in}}%
\pgfpathlineto{\pgfqpoint{4.230699in}{2.890259in}}%
\pgfpathlineto{\pgfqpoint{4.233218in}{2.851790in}}%
\pgfpathlineto{\pgfqpoint{4.235737in}{2.873399in}}%
\pgfpathlineto{\pgfqpoint{4.238256in}{2.875522in}}%
\pgfpathlineto{\pgfqpoint{4.240774in}{2.876631in}}%
\pgfpathlineto{\pgfqpoint{4.243293in}{2.865062in}}%
\pgfpathlineto{\pgfqpoint{4.245812in}{2.838166in}}%
\pgfpathlineto{\pgfqpoint{4.248331in}{2.843713in}}%
\pgfpathlineto{\pgfqpoint{4.250849in}{2.846163in}}%
\pgfpathlineto{\pgfqpoint{4.253368in}{2.807794in}}%
\pgfpathlineto{\pgfqpoint{4.255887in}{2.809306in}}%
\pgfpathlineto{\pgfqpoint{4.258406in}{2.800458in}}%
\pgfpathlineto{\pgfqpoint{4.260924in}{2.779153in}}%
\pgfpathlineto{\pgfqpoint{4.263443in}{2.775788in}}%
\pgfpathlineto{\pgfqpoint{4.265962in}{2.756411in}}%
\pgfpathlineto{\pgfqpoint{4.268481in}{2.758066in}}%
\pgfpathlineto{\pgfqpoint{4.270999in}{2.776734in}}%
\pgfpathlineto{\pgfqpoint{4.273518in}{2.796230in}}%
\pgfpathlineto{\pgfqpoint{4.276037in}{2.798089in}}%
\pgfpathlineto{\pgfqpoint{4.278556in}{2.831371in}}%
\pgfpathlineto{\pgfqpoint{4.281074in}{2.844624in}}%
\pgfpathlineto{\pgfqpoint{4.283593in}{2.855483in}}%
\pgfpathlineto{\pgfqpoint{4.286112in}{2.845461in}}%
\pgfpathlineto{\pgfqpoint{4.288631in}{2.806854in}}%
\pgfpathlineto{\pgfqpoint{4.291149in}{2.819998in}}%
\pgfpathlineto{\pgfqpoint{4.293668in}{2.829697in}}%
\pgfpathlineto{\pgfqpoint{4.296187in}{2.815305in}}%
\pgfpathlineto{\pgfqpoint{4.298706in}{2.812918in}}%
\pgfpathlineto{\pgfqpoint{4.301224in}{2.826333in}}%
\pgfpathlineto{\pgfqpoint{4.303743in}{2.854264in}}%
\pgfpathlineto{\pgfqpoint{4.306262in}{2.860552in}}%
\pgfpathlineto{\pgfqpoint{4.308781in}{2.890806in}}%
\pgfpathlineto{\pgfqpoint{4.311299in}{2.881301in}}%
\pgfpathlineto{\pgfqpoint{4.313818in}{2.904034in}}%
\pgfpathlineto{\pgfqpoint{4.316337in}{2.891679in}}%
\pgfpathlineto{\pgfqpoint{4.318856in}{2.914156in}}%
\pgfpathlineto{\pgfqpoint{4.321374in}{2.952202in}}%
\pgfpathlineto{\pgfqpoint{4.323893in}{2.975122in}}%
\pgfpathlineto{\pgfqpoint{4.326412in}{2.962895in}}%
\pgfpathlineto{\pgfqpoint{4.328931in}{2.959597in}}%
\pgfpathlineto{\pgfqpoint{4.331449in}{2.984631in}}%
\pgfpathlineto{\pgfqpoint{4.333968in}{3.019870in}}%
\pgfpathlineto{\pgfqpoint{4.336487in}{3.029904in}}%
\pgfpathlineto{\pgfqpoint{4.339006in}{3.019169in}}%
\pgfpathlineto{\pgfqpoint{4.341524in}{3.016430in}}%
\pgfpathlineto{\pgfqpoint{4.344043in}{3.027268in}}%
\pgfpathlineto{\pgfqpoint{4.346562in}{3.055076in}}%
\pgfpathlineto{\pgfqpoint{4.349081in}{3.056811in}}%
\pgfpathlineto{\pgfqpoint{4.351599in}{3.077319in}}%
\pgfpathlineto{\pgfqpoint{4.354118in}{3.080793in}}%
\pgfpathlineto{\pgfqpoint{4.356637in}{3.043909in}}%
\pgfpathlineto{\pgfqpoint{4.359156in}{3.068642in}}%
\pgfpathlineto{\pgfqpoint{4.361674in}{3.053705in}}%
\pgfpathlineto{\pgfqpoint{4.364193in}{3.062083in}}%
\pgfpathlineto{\pgfqpoint{4.366712in}{3.050239in}}%
\pgfpathlineto{\pgfqpoint{4.369231in}{3.043787in}}%
\pgfpathlineto{\pgfqpoint{4.371749in}{3.043328in}}%
\pgfpathlineto{\pgfqpoint{4.374268in}{3.079282in}}%
\pgfpathlineto{\pgfqpoint{4.376787in}{3.059428in}}%
\pgfpathlineto{\pgfqpoint{4.379306in}{3.053352in}}%
\pgfpathlineto{\pgfqpoint{4.381824in}{3.068436in}}%
\pgfpathlineto{\pgfqpoint{4.384343in}{3.077965in}}%
\pgfpathlineto{\pgfqpoint{4.386862in}{3.057401in}}%
\pgfpathlineto{\pgfqpoint{4.389381in}{3.055477in}}%
\pgfpathlineto{\pgfqpoint{4.391899in}{3.075654in}}%
\pgfpathlineto{\pgfqpoint{4.394418in}{3.084453in}}%
\pgfpathlineto{\pgfqpoint{4.396937in}{3.082759in}}%
\pgfpathlineto{\pgfqpoint{4.399456in}{3.086201in}}%
\pgfpathlineto{\pgfqpoint{4.401974in}{3.096684in}}%
\pgfpathlineto{\pgfqpoint{4.404493in}{3.066261in}}%
\pgfpathlineto{\pgfqpoint{4.407012in}{3.063528in}}%
\pgfpathlineto{\pgfqpoint{4.409531in}{3.087153in}}%
\pgfpathlineto{\pgfqpoint{4.412049in}{3.063337in}}%
\pgfpathlineto{\pgfqpoint{4.414568in}{3.055763in}}%
\pgfpathlineto{\pgfqpoint{4.417087in}{3.066910in}}%
\pgfpathlineto{\pgfqpoint{4.419606in}{3.086294in}}%
\pgfpathlineto{\pgfqpoint{4.422124in}{3.102913in}}%
\pgfpathlineto{\pgfqpoint{4.424643in}{3.126935in}}%
\pgfpathlineto{\pgfqpoint{4.427162in}{3.128276in}}%
\pgfpathlineto{\pgfqpoint{4.429681in}{3.147112in}}%
\pgfpathlineto{\pgfqpoint{4.432199in}{3.124390in}}%
\pgfpathlineto{\pgfqpoint{4.434718in}{3.154372in}}%
\pgfpathlineto{\pgfqpoint{4.437237in}{3.158439in}}%
\pgfpathlineto{\pgfqpoint{4.439756in}{3.149482in}}%
\pgfpathlineto{\pgfqpoint{4.442274in}{3.147782in}}%
\pgfpathlineto{\pgfqpoint{4.447312in}{3.195887in}}%
\pgfpathlineto{\pgfqpoint{4.449831in}{3.206997in}}%
\pgfpathlineto{\pgfqpoint{4.452349in}{3.232036in}}%
\pgfpathlineto{\pgfqpoint{4.454868in}{3.269790in}}%
\pgfpathlineto{\pgfqpoint{4.457387in}{3.273462in}}%
\pgfpathlineto{\pgfqpoint{4.459906in}{3.290161in}}%
\pgfpathlineto{\pgfqpoint{4.462424in}{3.294732in}}%
\pgfpathlineto{\pgfqpoint{4.464943in}{3.292968in}}%
\pgfpathlineto{\pgfqpoint{4.467462in}{3.315669in}}%
\pgfpathlineto{\pgfqpoint{4.469981in}{3.332942in}}%
\pgfpathlineto{\pgfqpoint{4.472499in}{3.334825in}}%
\pgfpathlineto{\pgfqpoint{4.475018in}{3.330585in}}%
\pgfpathlineto{\pgfqpoint{4.477537in}{3.360073in}}%
\pgfpathlineto{\pgfqpoint{4.480056in}{3.379640in}}%
\pgfpathlineto{\pgfqpoint{4.482574in}{3.368770in}}%
\pgfpathlineto{\pgfqpoint{4.485093in}{3.383937in}}%
\pgfpathlineto{\pgfqpoint{4.487612in}{3.377508in}}%
\pgfpathlineto{\pgfqpoint{4.490131in}{3.360610in}}%
\pgfpathlineto{\pgfqpoint{4.492649in}{3.365489in}}%
\pgfpathlineto{\pgfqpoint{4.495168in}{3.371966in}}%
\pgfpathlineto{\pgfqpoint{4.497687in}{3.378144in}}%
\pgfpathlineto{\pgfqpoint{4.500206in}{3.397269in}}%
\pgfpathlineto{\pgfqpoint{4.502724in}{3.393313in}}%
\pgfpathlineto{\pgfqpoint{4.505243in}{3.378196in}}%
\pgfpathlineto{\pgfqpoint{4.507762in}{3.368426in}}%
\pgfpathlineto{\pgfqpoint{4.510281in}{3.414475in}}%
\pgfpathlineto{\pgfqpoint{4.512799in}{3.442861in}}%
\pgfpathlineto{\pgfqpoint{4.515318in}{3.448243in}}%
\pgfpathlineto{\pgfqpoint{4.517837in}{3.438154in}}%
\pgfpathlineto{\pgfqpoint{4.520356in}{3.449898in}}%
\pgfpathlineto{\pgfqpoint{4.522874in}{3.492725in}}%
\pgfpathlineto{\pgfqpoint{4.525393in}{3.506142in}}%
\pgfpathlineto{\pgfqpoint{4.527912in}{3.500673in}}%
\pgfpathlineto{\pgfqpoint{4.530431in}{3.501188in}}%
\pgfpathlineto{\pgfqpoint{4.532949in}{3.487084in}}%
\pgfpathlineto{\pgfqpoint{4.535468in}{3.510661in}}%
\pgfpathlineto{\pgfqpoint{4.537987in}{3.507217in}}%
\pgfpathlineto{\pgfqpoint{4.540506in}{3.552724in}}%
\pgfpathlineto{\pgfqpoint{4.543024in}{3.546486in}}%
\pgfpathlineto{\pgfqpoint{4.545543in}{3.589523in}}%
\pgfpathlineto{\pgfqpoint{4.548062in}{3.587330in}}%
\pgfpathlineto{\pgfqpoint{4.550581in}{3.607595in}}%
\pgfpathlineto{\pgfqpoint{4.553099in}{3.634661in}}%
\pgfpathlineto{\pgfqpoint{4.555618in}{3.626648in}}%
\pgfpathlineto{\pgfqpoint{4.558137in}{3.632209in}}%
\pgfpathlineto{\pgfqpoint{4.560656in}{3.607537in}}%
\pgfpathlineto{\pgfqpoint{4.563174in}{3.643519in}}%
\pgfpathlineto{\pgfqpoint{4.565693in}{3.654212in}}%
\pgfpathlineto{\pgfqpoint{4.568212in}{3.641314in}}%
\pgfpathlineto{\pgfqpoint{4.570731in}{3.660728in}}%
\pgfpathlineto{\pgfqpoint{4.573249in}{3.654406in}}%
\pgfpathlineto{\pgfqpoint{4.575768in}{3.715036in}}%
\pgfpathlineto{\pgfqpoint{4.578287in}{3.703201in}}%
\pgfpathlineto{\pgfqpoint{4.580806in}{3.753887in}}%
\pgfpathlineto{\pgfqpoint{4.583324in}{3.732566in}}%
\pgfpathlineto{\pgfqpoint{4.585843in}{3.734556in}}%
\pgfpathlineto{\pgfqpoint{4.588362in}{3.779775in}}%
\pgfpathlineto{\pgfqpoint{4.590881in}{3.797819in}}%
\pgfpathlineto{\pgfqpoint{4.593399in}{3.774114in}}%
\pgfpathlineto{\pgfqpoint{4.595918in}{3.791643in}}%
\pgfpathlineto{\pgfqpoint{4.598437in}{3.758947in}}%
\pgfpathlineto{\pgfqpoint{4.600956in}{3.708013in}}%
\pgfpathlineto{\pgfqpoint{4.603474in}{3.725619in}}%
\pgfpathlineto{\pgfqpoint{4.605993in}{3.774223in}}%
\pgfpathlineto{\pgfqpoint{4.608512in}{3.749132in}}%
\pgfpathlineto{\pgfqpoint{4.611031in}{3.710371in}}%
\pgfpathlineto{\pgfqpoint{4.613549in}{3.664524in}}%
\pgfpathlineto{\pgfqpoint{4.616068in}{3.638520in}}%
\pgfpathlineto{\pgfqpoint{4.618587in}{3.647764in}}%
\pgfpathlineto{\pgfqpoint{4.621106in}{3.651620in}}%
\pgfpathlineto{\pgfqpoint{4.623624in}{3.685776in}}%
\pgfpathlineto{\pgfqpoint{4.626143in}{3.679877in}}%
\pgfpathlineto{\pgfqpoint{4.628662in}{3.705804in}}%
\pgfpathlineto{\pgfqpoint{4.631181in}{3.718093in}}%
\pgfpathlineto{\pgfqpoint{4.633699in}{3.679057in}}%
\pgfpathlineto{\pgfqpoint{4.636218in}{3.724335in}}%
\pgfpathlineto{\pgfqpoint{4.638737in}{3.721712in}}%
\pgfpathlineto{\pgfqpoint{4.641256in}{3.722953in}}%
\pgfpathlineto{\pgfqpoint{4.643774in}{3.701799in}}%
\pgfpathlineto{\pgfqpoint{4.646293in}{3.671208in}}%
\pgfpathlineto{\pgfqpoint{4.648812in}{3.687948in}}%
\pgfpathlineto{\pgfqpoint{4.651331in}{3.666205in}}%
\pgfpathlineto{\pgfqpoint{4.653849in}{3.676655in}}%
\pgfpathlineto{\pgfqpoint{4.656368in}{3.639286in}}%
\pgfpathlineto{\pgfqpoint{4.658887in}{3.636406in}}%
\pgfpathlineto{\pgfqpoint{4.661406in}{3.633135in}}%
\pgfpathlineto{\pgfqpoint{4.663924in}{3.654057in}}%
\pgfpathlineto{\pgfqpoint{4.666443in}{3.660769in}}%
\pgfpathlineto{\pgfqpoint{4.668962in}{3.631731in}}%
\pgfpathlineto{\pgfqpoint{4.671481in}{3.656290in}}%
\pgfpathlineto{\pgfqpoint{4.673999in}{3.685985in}}%
\pgfpathlineto{\pgfqpoint{4.676518in}{3.723132in}}%
\pgfpathlineto{\pgfqpoint{4.679037in}{3.696289in}}%
\pgfpathlineto{\pgfqpoint{4.681556in}{3.656510in}}%
\pgfpathlineto{\pgfqpoint{4.684074in}{3.687246in}}%
\pgfpathlineto{\pgfqpoint{4.686593in}{3.674548in}}%
\pgfpathlineto{\pgfqpoint{4.689112in}{3.690588in}}%
\pgfpathlineto{\pgfqpoint{4.691631in}{3.701824in}}%
\pgfpathlineto{\pgfqpoint{4.694149in}{3.698940in}}%
\pgfpathlineto{\pgfqpoint{4.696668in}{3.674130in}}%
\pgfpathlineto{\pgfqpoint{4.699187in}{3.659767in}}%
\pgfpathlineto{\pgfqpoint{4.701706in}{3.617618in}}%
\pgfpathlineto{\pgfqpoint{4.704224in}{3.619502in}}%
\pgfpathlineto{\pgfqpoint{4.706743in}{3.645694in}}%
\pgfpathlineto{\pgfqpoint{4.709262in}{3.625285in}}%
\pgfpathlineto{\pgfqpoint{4.711781in}{3.694125in}}%
\pgfpathlineto{\pgfqpoint{4.714299in}{3.738637in}}%
\pgfpathlineto{\pgfqpoint{4.716818in}{3.739864in}}%
\pgfpathlineto{\pgfqpoint{4.719337in}{3.724330in}}%
\pgfpathlineto{\pgfqpoint{4.721856in}{3.748570in}}%
\pgfpathlineto{\pgfqpoint{4.724374in}{3.736401in}}%
\pgfpathlineto{\pgfqpoint{4.726893in}{3.739930in}}%
\pgfpathlineto{\pgfqpoint{4.729412in}{3.729405in}}%
\pgfpathlineto{\pgfqpoint{4.731931in}{3.709057in}}%
\pgfpathlineto{\pgfqpoint{4.734449in}{3.701345in}}%
\pgfpathlineto{\pgfqpoint{4.736968in}{3.701441in}}%
\pgfpathlineto{\pgfqpoint{4.739487in}{3.686703in}}%
\pgfpathlineto{\pgfqpoint{4.742006in}{3.692838in}}%
\pgfpathlineto{\pgfqpoint{4.744524in}{3.711924in}}%
\pgfpathlineto{\pgfqpoint{4.747043in}{3.708151in}}%
\pgfpathlineto{\pgfqpoint{4.749562in}{3.735922in}}%
\pgfpathlineto{\pgfqpoint{4.752081in}{3.788743in}}%
\pgfpathlineto{\pgfqpoint{4.754599in}{3.742124in}}%
\pgfpathlineto{\pgfqpoint{4.757118in}{3.748333in}}%
\pgfpathlineto{\pgfqpoint{4.759637in}{3.746947in}}%
\pgfpathlineto{\pgfqpoint{4.762156in}{3.720441in}}%
\pgfpathlineto{\pgfqpoint{4.764674in}{3.736714in}}%
\pgfpathlineto{\pgfqpoint{4.767193in}{3.733438in}}%
\pgfpathlineto{\pgfqpoint{4.769712in}{3.771026in}}%
\pgfpathlineto{\pgfqpoint{4.772231in}{3.779647in}}%
\pgfpathlineto{\pgfqpoint{4.774749in}{3.820161in}}%
\pgfpathlineto{\pgfqpoint{4.777268in}{3.812485in}}%
\pgfpathlineto{\pgfqpoint{4.779787in}{3.796535in}}%
\pgfpathlineto{\pgfqpoint{4.782306in}{3.867516in}}%
\pgfpathlineto{\pgfqpoint{4.784824in}{3.902353in}}%
\pgfpathlineto{\pgfqpoint{4.787343in}{3.886080in}}%
\pgfpathlineto{\pgfqpoint{4.789862in}{3.854718in}}%
\pgfpathlineto{\pgfqpoint{4.792381in}{3.863023in}}%
\pgfpathlineto{\pgfqpoint{4.794899in}{3.870065in}}%
\pgfpathlineto{\pgfqpoint{4.797418in}{3.893555in}}%
\pgfpathlineto{\pgfqpoint{4.799937in}{3.884881in}}%
\pgfpathlineto{\pgfqpoint{4.802456in}{3.898585in}}%
\pgfpathlineto{\pgfqpoint{4.804974in}{3.927394in}}%
\pgfpathlineto{\pgfqpoint{4.807493in}{3.882403in}}%
\pgfpathlineto{\pgfqpoint{4.810012in}{3.811198in}}%
\pgfpathlineto{\pgfqpoint{4.812531in}{3.817327in}}%
\pgfpathlineto{\pgfqpoint{4.815049in}{3.792362in}}%
\pgfpathlineto{\pgfqpoint{4.817568in}{3.785270in}}%
\pgfpathlineto{\pgfqpoint{4.820087in}{3.777106in}}%
\pgfpathlineto{\pgfqpoint{4.822606in}{3.777091in}}%
\pgfpathlineto{\pgfqpoint{4.825124in}{3.774394in}}%
\pgfpathlineto{\pgfqpoint{4.827643in}{3.715807in}}%
\pgfpathlineto{\pgfqpoint{4.830162in}{3.695086in}}%
\pgfpathlineto{\pgfqpoint{4.832681in}{3.711710in}}%
\pgfpathlineto{\pgfqpoint{4.835199in}{3.718486in}}%
\pgfpathlineto{\pgfqpoint{4.837718in}{3.704175in}}%
\pgfpathlineto{\pgfqpoint{4.840237in}{3.683161in}}%
\pgfpathlineto{\pgfqpoint{4.842756in}{3.679787in}}%
\pgfpathlineto{\pgfqpoint{4.845274in}{3.649815in}}%
\pgfpathlineto{\pgfqpoint{4.847793in}{3.666176in}}%
\pgfpathlineto{\pgfqpoint{4.850312in}{3.637228in}}%
\pgfpathlineto{\pgfqpoint{4.852831in}{3.653808in}}%
\pgfpathlineto{\pgfqpoint{4.855349in}{3.663028in}}%
\pgfpathlineto{\pgfqpoint{4.857868in}{3.648385in}}%
\pgfpathlineto{\pgfqpoint{4.860387in}{3.616823in}}%
\pgfpathlineto{\pgfqpoint{4.862906in}{3.616041in}}%
\pgfpathlineto{\pgfqpoint{4.865424in}{3.591330in}}%
\pgfpathlineto{\pgfqpoint{4.867943in}{3.600181in}}%
\pgfpathlineto{\pgfqpoint{4.870462in}{3.550747in}}%
\pgfpathlineto{\pgfqpoint{4.872981in}{3.536660in}}%
\pgfpathlineto{\pgfqpoint{4.875499in}{3.507078in}}%
\pgfpathlineto{\pgfqpoint{4.878018in}{3.489985in}}%
\pgfpathlineto{\pgfqpoint{4.880537in}{3.479380in}}%
\pgfpathlineto{\pgfqpoint{4.883056in}{3.491197in}}%
\pgfpathlineto{\pgfqpoint{4.885574in}{3.479600in}}%
\pgfpathlineto{\pgfqpoint{4.888093in}{3.489753in}}%
\pgfpathlineto{\pgfqpoint{4.890612in}{3.463068in}}%
\pgfpathlineto{\pgfqpoint{4.893131in}{3.445831in}}%
\pgfpathlineto{\pgfqpoint{4.895649in}{3.432928in}}%
\pgfpathlineto{\pgfqpoint{4.898168in}{3.427042in}}%
\pgfpathlineto{\pgfqpoint{4.900687in}{3.467083in}}%
\pgfpathlineto{\pgfqpoint{4.903206in}{3.491945in}}%
\pgfpathlineto{\pgfqpoint{4.905724in}{3.499645in}}%
\pgfpathlineto{\pgfqpoint{4.908243in}{3.509784in}}%
\pgfpathlineto{\pgfqpoint{4.910762in}{3.499343in}}%
\pgfpathlineto{\pgfqpoint{4.913281in}{3.463319in}}%
\pgfpathlineto{\pgfqpoint{4.915799in}{3.497179in}}%
\pgfpathlineto{\pgfqpoint{4.918318in}{3.529649in}}%
\pgfpathlineto{\pgfqpoint{4.920837in}{3.534533in}}%
\pgfpathlineto{\pgfqpoint{4.923356in}{3.516608in}}%
\pgfpathlineto{\pgfqpoint{4.925874in}{3.517890in}}%
\pgfpathlineto{\pgfqpoint{4.928393in}{3.497517in}}%
\pgfpathlineto{\pgfqpoint{4.930912in}{3.490323in}}%
\pgfpathlineto{\pgfqpoint{4.933431in}{3.477928in}}%
\pgfpathlineto{\pgfqpoint{4.935949in}{3.480608in}}%
\pgfpathlineto{\pgfqpoint{4.938468in}{3.491644in}}%
\pgfpathlineto{\pgfqpoint{4.940987in}{3.461993in}}%
\pgfpathlineto{\pgfqpoint{4.943506in}{3.492169in}}%
\pgfpathlineto{\pgfqpoint{4.946024in}{3.507772in}}%
\pgfpathlineto{\pgfqpoint{4.948543in}{3.542741in}}%
\pgfpathlineto{\pgfqpoint{4.951062in}{3.532061in}}%
\pgfpathlineto{\pgfqpoint{4.953581in}{3.551380in}}%
\pgfpathlineto{\pgfqpoint{4.956099in}{3.547270in}}%
\pgfpathlineto{\pgfqpoint{4.958618in}{3.538109in}}%
\pgfpathlineto{\pgfqpoint{4.961137in}{3.568112in}}%
\pgfpathlineto{\pgfqpoint{4.963656in}{3.549239in}}%
\pgfpathlineto{\pgfqpoint{4.966174in}{3.549581in}}%
\pgfpathlineto{\pgfqpoint{4.968693in}{3.578603in}}%
\pgfpathlineto{\pgfqpoint{4.971212in}{3.598299in}}%
\pgfpathlineto{\pgfqpoint{4.973731in}{3.669188in}}%
\pgfpathlineto{\pgfqpoint{4.976249in}{3.653581in}}%
\pgfpathlineto{\pgfqpoint{4.978768in}{3.663752in}}%
\pgfpathlineto{\pgfqpoint{4.981287in}{3.665048in}}%
\pgfpathlineto{\pgfqpoint{4.983806in}{3.650874in}}%
\pgfpathlineto{\pgfqpoint{4.986324in}{3.702703in}}%
\pgfpathlineto{\pgfqpoint{4.988843in}{3.694574in}}%
\pgfpathlineto{\pgfqpoint{4.991362in}{3.751145in}}%
\pgfpathlineto{\pgfqpoint{4.993881in}{3.769759in}}%
\pgfpathlineto{\pgfqpoint{4.996399in}{3.780343in}}%
\pgfpathlineto{\pgfqpoint{4.998918in}{3.798126in}}%
\pgfpathlineto{\pgfqpoint{5.001437in}{3.783894in}}%
\pgfpathlineto{\pgfqpoint{5.003956in}{3.762435in}}%
\pgfpathlineto{\pgfqpoint{5.006474in}{3.785996in}}%
\pgfpathlineto{\pgfqpoint{5.008993in}{3.773798in}}%
\pgfpathlineto{\pgfqpoint{5.011512in}{3.754074in}}%
\pgfpathlineto{\pgfqpoint{5.014031in}{3.749468in}}%
\pgfpathlineto{\pgfqpoint{5.016549in}{3.763419in}}%
\pgfpathlineto{\pgfqpoint{5.019068in}{3.751108in}}%
\pgfpathlineto{\pgfqpoint{5.021587in}{3.764279in}}%
\pgfpathlineto{\pgfqpoint{5.024106in}{3.794398in}}%
\pgfpathlineto{\pgfqpoint{5.026624in}{3.807768in}}%
\pgfpathlineto{\pgfqpoint{5.029143in}{3.809914in}}%
\pgfpathlineto{\pgfqpoint{5.031662in}{3.800476in}}%
\pgfpathlineto{\pgfqpoint{5.034181in}{3.775148in}}%
\pgfpathlineto{\pgfqpoint{5.036699in}{3.766186in}}%
\pgfpathlineto{\pgfqpoint{5.039218in}{3.767858in}}%
\pgfpathlineto{\pgfqpoint{5.041737in}{3.849707in}}%
\pgfpathlineto{\pgfqpoint{5.044256in}{3.842075in}}%
\pgfpathlineto{\pgfqpoint{5.046774in}{3.836185in}}%
\pgfpathlineto{\pgfqpoint{5.049293in}{3.870539in}}%
\pgfpathlineto{\pgfqpoint{5.051812in}{3.885505in}}%
\pgfpathlineto{\pgfqpoint{5.054331in}{3.885054in}}%
\pgfpathlineto{\pgfqpoint{5.056849in}{3.915102in}}%
\pgfpathlineto{\pgfqpoint{5.059368in}{3.935729in}}%
\pgfpathlineto{\pgfqpoint{5.061887in}{3.947692in}}%
\pgfpathlineto{\pgfqpoint{5.064406in}{3.919775in}}%
\pgfpathlineto{\pgfqpoint{5.066924in}{3.872470in}}%
\pgfpathlineto{\pgfqpoint{5.069443in}{3.868195in}}%
\pgfpathlineto{\pgfqpoint{5.071962in}{3.814594in}}%
\pgfpathlineto{\pgfqpoint{5.074481in}{3.841808in}}%
\pgfpathlineto{\pgfqpoint{5.076999in}{3.820517in}}%
\pgfpathlineto{\pgfqpoint{5.079518in}{3.825350in}}%
\pgfpathlineto{\pgfqpoint{5.082037in}{3.831731in}}%
\pgfpathlineto{\pgfqpoint{5.084556in}{3.798388in}}%
\pgfpathlineto{\pgfqpoint{5.087074in}{3.822611in}}%
\pgfpathlineto{\pgfqpoint{5.089593in}{3.854200in}}%
\pgfpathlineto{\pgfqpoint{5.092112in}{3.856643in}}%
\pgfpathlineto{\pgfqpoint{5.094631in}{3.896800in}}%
\pgfpathlineto{\pgfqpoint{5.097149in}{3.830810in}}%
\pgfpathlineto{\pgfqpoint{5.099668in}{3.868636in}}%
\pgfpathlineto{\pgfqpoint{5.102187in}{3.892952in}}%
\pgfpathlineto{\pgfqpoint{5.104706in}{3.896269in}}%
\pgfpathlineto{\pgfqpoint{5.107224in}{3.894938in}}%
\pgfpathlineto{\pgfqpoint{5.109743in}{3.922039in}}%
\pgfpathlineto{\pgfqpoint{5.112262in}{3.893705in}}%
\pgfpathlineto{\pgfqpoint{5.114781in}{3.899364in}}%
\pgfpathlineto{\pgfqpoint{5.117299in}{3.888688in}}%
\pgfpathlineto{\pgfqpoint{5.119818in}{3.874489in}}%
\pgfpathlineto{\pgfqpoint{5.122337in}{3.917818in}}%
\pgfpathlineto{\pgfqpoint{5.124856in}{3.934341in}}%
\pgfpathlineto{\pgfqpoint{5.127374in}{3.916002in}}%
\pgfpathlineto{\pgfqpoint{5.129893in}{3.881055in}}%
\pgfpathlineto{\pgfqpoint{5.132412in}{3.888363in}}%
\pgfpathlineto{\pgfqpoint{5.134931in}{3.879197in}}%
\pgfpathlineto{\pgfqpoint{5.137449in}{3.888064in}}%
\pgfpathlineto{\pgfqpoint{5.139968in}{3.893647in}}%
\pgfpathlineto{\pgfqpoint{5.142487in}{3.956882in}}%
\pgfpathlineto{\pgfqpoint{5.145006in}{3.974132in}}%
\pgfpathlineto{\pgfqpoint{5.147524in}{3.948136in}}%
\pgfpathlineto{\pgfqpoint{5.150043in}{3.908334in}}%
\pgfpathlineto{\pgfqpoint{5.152562in}{3.910156in}}%
\pgfpathlineto{\pgfqpoint{5.155081in}{3.867709in}}%
\pgfpathlineto{\pgfqpoint{5.157599in}{3.829630in}}%
\pgfpathlineto{\pgfqpoint{5.160118in}{3.816903in}}%
\pgfpathlineto{\pgfqpoint{5.162637in}{3.822804in}}%
\pgfpathlineto{\pgfqpoint{5.165156in}{3.820517in}}%
\pgfpathlineto{\pgfqpoint{5.167674in}{3.805115in}}%
\pgfpathlineto{\pgfqpoint{5.170193in}{3.817941in}}%
\pgfpathlineto{\pgfqpoint{5.172712in}{3.840537in}}%
\pgfpathlineto{\pgfqpoint{5.175231in}{3.810229in}}%
\pgfpathlineto{\pgfqpoint{5.177749in}{3.787994in}}%
\pgfpathlineto{\pgfqpoint{5.180268in}{3.819351in}}%
\pgfpathlineto{\pgfqpoint{5.182787in}{3.779951in}}%
\pgfpathlineto{\pgfqpoint{5.185306in}{3.774044in}}%
\pgfpathlineto{\pgfqpoint{5.187824in}{3.759441in}}%
\pgfpathlineto{\pgfqpoint{5.190343in}{3.755005in}}%
\pgfpathlineto{\pgfqpoint{5.192862in}{3.796711in}}%
\pgfpathlineto{\pgfqpoint{5.195381in}{3.845573in}}%
\pgfpathlineto{\pgfqpoint{5.197899in}{3.848551in}}%
\pgfpathlineto{\pgfqpoint{5.200418in}{3.861130in}}%
\pgfpathlineto{\pgfqpoint{5.202937in}{3.894560in}}%
\pgfpathlineto{\pgfqpoint{5.205456in}{3.902890in}}%
\pgfpathlineto{\pgfqpoint{5.207974in}{3.926857in}}%
\pgfpathlineto{\pgfqpoint{5.210493in}{3.924701in}}%
\pgfpathlineto{\pgfqpoint{5.213012in}{3.936530in}}%
\pgfpathlineto{\pgfqpoint{5.215531in}{3.910539in}}%
\pgfpathlineto{\pgfqpoint{5.218049in}{3.859876in}}%
\pgfpathlineto{\pgfqpoint{5.220568in}{3.849491in}}%
\pgfpathlineto{\pgfqpoint{5.223087in}{3.862745in}}%
\pgfpathlineto{\pgfqpoint{5.225606in}{3.871343in}}%
\pgfpathlineto{\pgfqpoint{5.228124in}{3.887867in}}%
\pgfpathlineto{\pgfqpoint{5.230643in}{3.845640in}}%
\pgfpathlineto{\pgfqpoint{5.233162in}{3.784369in}}%
\pgfpathlineto{\pgfqpoint{5.235681in}{3.775982in}}%
\pgfpathlineto{\pgfqpoint{5.238199in}{3.738848in}}%
\pgfpathlineto{\pgfqpoint{5.240718in}{3.742285in}}%
\pgfpathlineto{\pgfqpoint{5.243237in}{3.742812in}}%
\pgfpathlineto{\pgfqpoint{5.245756in}{3.775899in}}%
\pgfpathlineto{\pgfqpoint{5.248274in}{3.743340in}}%
\pgfpathlineto{\pgfqpoint{5.250793in}{3.730007in}}%
\pgfpathlineto{\pgfqpoint{5.253312in}{3.696183in}}%
\pgfpathlineto{\pgfqpoint{5.255831in}{3.730899in}}%
\pgfpathlineto{\pgfqpoint{5.258349in}{3.743767in}}%
\pgfpathlineto{\pgfqpoint{5.260868in}{3.743712in}}%
\pgfpathlineto{\pgfqpoint{5.263387in}{3.711135in}}%
\pgfpathlineto{\pgfqpoint{5.265906in}{3.652517in}}%
\pgfpathlineto{\pgfqpoint{5.268424in}{3.649388in}}%
\pgfpathlineto{\pgfqpoint{5.270943in}{3.667409in}}%
\pgfpathlineto{\pgfqpoint{5.273462in}{3.684595in}}%
\pgfpathlineto{\pgfqpoint{5.275981in}{3.698106in}}%
\pgfpathlineto{\pgfqpoint{5.278499in}{3.723728in}}%
\pgfpathlineto{\pgfqpoint{5.281018in}{3.725195in}}%
\pgfpathlineto{\pgfqpoint{5.283537in}{3.746753in}}%
\pgfpathlineto{\pgfqpoint{5.286056in}{3.698058in}}%
\pgfpathlineto{\pgfqpoint{5.288574in}{3.716643in}}%
\pgfpathlineto{\pgfqpoint{5.291093in}{3.772772in}}%
\pgfpathlineto{\pgfqpoint{5.293612in}{3.745315in}}%
\pgfpathlineto{\pgfqpoint{5.296131in}{3.744031in}}%
\pgfpathlineto{\pgfqpoint{5.298649in}{3.794587in}}%
\pgfpathlineto{\pgfqpoint{5.301168in}{3.806706in}}%
\pgfpathlineto{\pgfqpoint{5.303687in}{3.822323in}}%
\pgfpathlineto{\pgfqpoint{5.306206in}{3.833146in}}%
\pgfpathlineto{\pgfqpoint{5.308724in}{3.823715in}}%
\pgfpathlineto{\pgfqpoint{5.311243in}{3.827331in}}%
\pgfpathlineto{\pgfqpoint{5.313762in}{3.845513in}}%
\pgfpathlineto{\pgfqpoint{5.316281in}{3.818826in}}%
\pgfpathlineto{\pgfqpoint{5.318799in}{3.801484in}}%
\pgfpathlineto{\pgfqpoint{5.321318in}{3.792459in}}%
\pgfpathlineto{\pgfqpoint{5.323837in}{3.764797in}}%
\pgfpathlineto{\pgfqpoint{5.326356in}{3.814518in}}%
\pgfpathlineto{\pgfqpoint{5.328874in}{3.825759in}}%
\pgfpathlineto{\pgfqpoint{5.331393in}{3.811512in}}%
\pgfpathlineto{\pgfqpoint{5.333912in}{3.819302in}}%
\pgfpathlineto{\pgfqpoint{5.336431in}{3.839903in}}%
\pgfpathlineto{\pgfqpoint{5.338949in}{3.851418in}}%
\pgfpathlineto{\pgfqpoint{5.341468in}{3.867816in}}%
\pgfpathlineto{\pgfqpoint{5.343987in}{3.887597in}}%
\pgfpathlineto{\pgfqpoint{5.346506in}{3.911311in}}%
\pgfpathlineto{\pgfqpoint{5.349024in}{3.905220in}}%
\pgfpathlineto{\pgfqpoint{5.351543in}{3.928909in}}%
\pgfpathlineto{\pgfqpoint{5.354062in}{3.928067in}}%
\pgfpathlineto{\pgfqpoint{5.356581in}{3.954100in}}%
\pgfpathlineto{\pgfqpoint{5.359099in}{3.957330in}}%
\pgfpathlineto{\pgfqpoint{5.361618in}{3.938362in}}%
\pgfpathlineto{\pgfqpoint{5.364137in}{3.940369in}}%
\pgfpathlineto{\pgfqpoint{5.366656in}{3.916304in}}%
\pgfpathlineto{\pgfqpoint{5.369174in}{3.984902in}}%
\pgfpathlineto{\pgfqpoint{5.371693in}{4.026830in}}%
\pgfpathlineto{\pgfqpoint{5.374212in}{3.980405in}}%
\pgfpathlineto{\pgfqpoint{5.376731in}{3.936125in}}%
\pgfpathlineto{\pgfqpoint{5.379249in}{3.886332in}}%
\pgfpathlineto{\pgfqpoint{5.381768in}{3.893899in}}%
\pgfpathlineto{\pgfqpoint{5.384287in}{3.885685in}}%
\pgfpathlineto{\pgfqpoint{5.386806in}{3.879891in}}%
\pgfpathlineto{\pgfqpoint{5.389324in}{3.908536in}}%
\pgfpathlineto{\pgfqpoint{5.391843in}{3.910798in}}%
\pgfpathlineto{\pgfqpoint{5.394362in}{3.932169in}}%
\pgfpathlineto{\pgfqpoint{5.396881in}{3.875223in}}%
\pgfpathlineto{\pgfqpoint{5.399399in}{3.886475in}}%
\pgfpathlineto{\pgfqpoint{5.401918in}{3.856099in}}%
\pgfpathlineto{\pgfqpoint{5.404437in}{3.889504in}}%
\pgfpathlineto{\pgfqpoint{5.406956in}{3.861798in}}%
\pgfpathlineto{\pgfqpoint{5.409474in}{3.880230in}}%
\pgfpathlineto{\pgfqpoint{5.411993in}{3.843415in}}%
\pgfpathlineto{\pgfqpoint{5.414512in}{3.848638in}}%
\pgfpathlineto{\pgfqpoint{5.417031in}{3.854519in}}%
\pgfpathlineto{\pgfqpoint{5.419549in}{3.858144in}}%
\pgfpathlineto{\pgfqpoint{5.422068in}{3.807707in}}%
\pgfpathlineto{\pgfqpoint{5.424587in}{3.798612in}}%
\pgfpathlineto{\pgfqpoint{5.427106in}{3.792718in}}%
\pgfpathlineto{\pgfqpoint{5.429624in}{3.736624in}}%
\pgfpathlineto{\pgfqpoint{5.432143in}{3.779518in}}%
\pgfpathlineto{\pgfqpoint{5.434662in}{3.825037in}}%
\pgfpathlineto{\pgfqpoint{5.437181in}{3.829263in}}%
\pgfpathlineto{\pgfqpoint{5.439699in}{3.857756in}}%
\pgfpathlineto{\pgfqpoint{5.442218in}{3.840354in}}%
\pgfpathlineto{\pgfqpoint{5.444737in}{3.841653in}}%
\pgfpathlineto{\pgfqpoint{5.447256in}{3.871312in}}%
\pgfpathlineto{\pgfqpoint{5.449774in}{3.887926in}}%
\pgfpathlineto{\pgfqpoint{5.452293in}{3.884677in}}%
\pgfpathlineto{\pgfqpoint{5.454812in}{3.901968in}}%
\pgfpathlineto{\pgfqpoint{5.457331in}{3.897441in}}%
\pgfpathlineto{\pgfqpoint{5.459849in}{3.878969in}}%
\pgfpathlineto{\pgfqpoint{5.462368in}{3.905587in}}%
\pgfpathlineto{\pgfqpoint{5.464887in}{3.923279in}}%
\pgfpathlineto{\pgfqpoint{5.467406in}{3.913388in}}%
\pgfpathlineto{\pgfqpoint{5.469924in}{3.887401in}}%
\pgfpathlineto{\pgfqpoint{5.472443in}{3.881118in}}%
\pgfpathlineto{\pgfqpoint{5.474962in}{3.847039in}}%
\pgfpathlineto{\pgfqpoint{5.477481in}{3.838231in}}%
\pgfpathlineto{\pgfqpoint{5.479999in}{3.876657in}}%
\pgfpathlineto{\pgfqpoint{5.482518in}{3.862043in}}%
\pgfpathlineto{\pgfqpoint{5.485037in}{3.828491in}}%
\pgfpathlineto{\pgfqpoint{5.487556in}{3.832634in}}%
\pgfpathlineto{\pgfqpoint{5.490074in}{3.828248in}}%
\pgfpathlineto{\pgfqpoint{5.492593in}{3.833652in}}%
\pgfpathlineto{\pgfqpoint{5.492593in}{3.833652in}}%
\pgfusepath{stroke}%
\end{pgfscope}%
\begin{pgfscope}%
\pgfpathrectangle{\pgfqpoint{0.455093in}{0.401080in}}{\pgfqpoint{5.037500in}{4.004000in}}%
\pgfusepath{clip}%
\pgfsetrectcap%
\pgfsetroundjoin%
\pgfsetlinewidth{0.501875pt}%
\definecolor{currentstroke}{rgb}{0.690196,0.188235,0.376471}%
\pgfsetstrokecolor{currentstroke}%
\pgfsetstrokeopacity{0.800000}%
\pgfsetdash{}{0pt}%
\pgfpathmoveto{\pgfqpoint{0.455093in}{1.021044in}}%
\pgfpathlineto{\pgfqpoint{0.457612in}{1.007323in}}%
\pgfpathlineto{\pgfqpoint{0.460131in}{1.010835in}}%
\pgfpathlineto{\pgfqpoint{0.462649in}{1.003061in}}%
\pgfpathlineto{\pgfqpoint{0.465168in}{1.012585in}}%
\pgfpathlineto{\pgfqpoint{0.467687in}{1.021328in}}%
\pgfpathlineto{\pgfqpoint{0.470206in}{1.037456in}}%
\pgfpathlineto{\pgfqpoint{0.472724in}{1.039821in}}%
\pgfpathlineto{\pgfqpoint{0.475243in}{1.046143in}}%
\pgfpathlineto{\pgfqpoint{0.477762in}{1.044429in}}%
\pgfpathlineto{\pgfqpoint{0.480281in}{1.020105in}}%
\pgfpathlineto{\pgfqpoint{0.482799in}{1.010417in}}%
\pgfpathlineto{\pgfqpoint{0.485318in}{0.992850in}}%
\pgfpathlineto{\pgfqpoint{0.487837in}{0.977766in}}%
\pgfpathlineto{\pgfqpoint{0.490356in}{0.983328in}}%
\pgfpathlineto{\pgfqpoint{0.492874in}{0.981744in}}%
\pgfpathlineto{\pgfqpoint{0.495393in}{0.969759in}}%
\pgfpathlineto{\pgfqpoint{0.497912in}{0.981944in}}%
\pgfpathlineto{\pgfqpoint{0.500431in}{1.004476in}}%
\pgfpathlineto{\pgfqpoint{0.502949in}{1.022610in}}%
\pgfpathlineto{\pgfqpoint{0.505468in}{1.047471in}}%
\pgfpathlineto{\pgfqpoint{0.507987in}{1.064700in}}%
\pgfpathlineto{\pgfqpoint{0.510506in}{1.056260in}}%
\pgfpathlineto{\pgfqpoint{0.513024in}{1.043788in}}%
\pgfpathlineto{\pgfqpoint{0.515543in}{1.026414in}}%
\pgfpathlineto{\pgfqpoint{0.518062in}{1.030728in}}%
\pgfpathlineto{\pgfqpoint{0.520581in}{1.011027in}}%
\pgfpathlineto{\pgfqpoint{0.523099in}{1.020436in}}%
\pgfpathlineto{\pgfqpoint{0.525618in}{1.036187in}}%
\pgfpathlineto{\pgfqpoint{0.528137in}{1.026984in}}%
\pgfpathlineto{\pgfqpoint{0.530656in}{1.004286in}}%
\pgfpathlineto{\pgfqpoint{0.533174in}{0.998327in}}%
\pgfpathlineto{\pgfqpoint{0.535693in}{1.001396in}}%
\pgfpathlineto{\pgfqpoint{0.538212in}{1.000182in}}%
\pgfpathlineto{\pgfqpoint{0.540731in}{1.016728in}}%
\pgfpathlineto{\pgfqpoint{0.543249in}{1.030252in}}%
\pgfpathlineto{\pgfqpoint{0.545768in}{1.050386in}}%
\pgfpathlineto{\pgfqpoint{0.548287in}{1.060672in}}%
\pgfpathlineto{\pgfqpoint{0.550806in}{1.057677in}}%
\pgfpathlineto{\pgfqpoint{0.553324in}{1.090691in}}%
\pgfpathlineto{\pgfqpoint{0.555843in}{1.100253in}}%
\pgfpathlineto{\pgfqpoint{0.558362in}{1.112381in}}%
\pgfpathlineto{\pgfqpoint{0.560881in}{1.113740in}}%
\pgfpathlineto{\pgfqpoint{0.563399in}{1.145553in}}%
\pgfpathlineto{\pgfqpoint{0.565918in}{1.146345in}}%
\pgfpathlineto{\pgfqpoint{0.568437in}{1.127167in}}%
\pgfpathlineto{\pgfqpoint{0.570956in}{1.127795in}}%
\pgfpathlineto{\pgfqpoint{0.573474in}{1.123218in}}%
\pgfpathlineto{\pgfqpoint{0.575993in}{1.119030in}}%
\pgfpathlineto{\pgfqpoint{0.578512in}{1.126833in}}%
\pgfpathlineto{\pgfqpoint{0.581031in}{1.139665in}}%
\pgfpathlineto{\pgfqpoint{0.583549in}{1.136129in}}%
\pgfpathlineto{\pgfqpoint{0.586068in}{1.154037in}}%
\pgfpathlineto{\pgfqpoint{0.588587in}{1.131362in}}%
\pgfpathlineto{\pgfqpoint{0.591106in}{1.117478in}}%
\pgfpathlineto{\pgfqpoint{0.593624in}{1.131714in}}%
\pgfpathlineto{\pgfqpoint{0.596143in}{1.144460in}}%
\pgfpathlineto{\pgfqpoint{0.598662in}{1.170921in}}%
\pgfpathlineto{\pgfqpoint{0.601181in}{1.170732in}}%
\pgfpathlineto{\pgfqpoint{0.603699in}{1.163947in}}%
\pgfpathlineto{\pgfqpoint{0.606218in}{1.142167in}}%
\pgfpathlineto{\pgfqpoint{0.608737in}{1.136585in}}%
\pgfpathlineto{\pgfqpoint{0.611256in}{1.151203in}}%
\pgfpathlineto{\pgfqpoint{0.613774in}{1.143905in}}%
\pgfpathlineto{\pgfqpoint{0.616293in}{1.154094in}}%
\pgfpathlineto{\pgfqpoint{0.618812in}{1.146076in}}%
\pgfpathlineto{\pgfqpoint{0.621331in}{1.146623in}}%
\pgfpathlineto{\pgfqpoint{0.623849in}{1.150343in}}%
\pgfpathlineto{\pgfqpoint{0.626368in}{1.135593in}}%
\pgfpathlineto{\pgfqpoint{0.628887in}{1.116369in}}%
\pgfpathlineto{\pgfqpoint{0.631406in}{1.109538in}}%
\pgfpathlineto{\pgfqpoint{0.633924in}{1.111210in}}%
\pgfpathlineto{\pgfqpoint{0.636443in}{1.120801in}}%
\pgfpathlineto{\pgfqpoint{0.638962in}{1.129489in}}%
\pgfpathlineto{\pgfqpoint{0.641481in}{1.119962in}}%
\pgfpathlineto{\pgfqpoint{0.643999in}{1.125525in}}%
\pgfpathlineto{\pgfqpoint{0.646518in}{1.114407in}}%
\pgfpathlineto{\pgfqpoint{0.649037in}{1.118241in}}%
\pgfpathlineto{\pgfqpoint{0.651556in}{1.111710in}}%
\pgfpathlineto{\pgfqpoint{0.654074in}{1.101975in}}%
\pgfpathlineto{\pgfqpoint{0.656593in}{1.098016in}}%
\pgfpathlineto{\pgfqpoint{0.659112in}{1.105188in}}%
\pgfpathlineto{\pgfqpoint{0.661631in}{1.097489in}}%
\pgfpathlineto{\pgfqpoint{0.664149in}{1.109838in}}%
\pgfpathlineto{\pgfqpoint{0.666668in}{1.097639in}}%
\pgfpathlineto{\pgfqpoint{0.669187in}{1.096099in}}%
\pgfpathlineto{\pgfqpoint{0.671706in}{1.106041in}}%
\pgfpathlineto{\pgfqpoint{0.674224in}{1.103828in}}%
\pgfpathlineto{\pgfqpoint{0.676743in}{1.111660in}}%
\pgfpathlineto{\pgfqpoint{0.679262in}{1.092787in}}%
\pgfpathlineto{\pgfqpoint{0.681781in}{1.087038in}}%
\pgfpathlineto{\pgfqpoint{0.684299in}{1.082061in}}%
\pgfpathlineto{\pgfqpoint{0.686818in}{1.065859in}}%
\pgfpathlineto{\pgfqpoint{0.689337in}{1.071209in}}%
\pgfpathlineto{\pgfqpoint{0.691856in}{1.075285in}}%
\pgfpathlineto{\pgfqpoint{0.694374in}{1.052627in}}%
\pgfpathlineto{\pgfqpoint{0.696893in}{1.050571in}}%
\pgfpathlineto{\pgfqpoint{0.699412in}{1.045567in}}%
\pgfpathlineto{\pgfqpoint{0.701931in}{1.079088in}}%
\pgfpathlineto{\pgfqpoint{0.704449in}{1.073629in}}%
\pgfpathlineto{\pgfqpoint{0.706968in}{1.060453in}}%
\pgfpathlineto{\pgfqpoint{0.709487in}{1.049102in}}%
\pgfpathlineto{\pgfqpoint{0.712006in}{1.040252in}}%
\pgfpathlineto{\pgfqpoint{0.714524in}{1.042431in}}%
\pgfpathlineto{\pgfqpoint{0.717043in}{1.047583in}}%
\pgfpathlineto{\pgfqpoint{0.719562in}{1.042038in}}%
\pgfpathlineto{\pgfqpoint{0.722081in}{1.037919in}}%
\pgfpathlineto{\pgfqpoint{0.724599in}{1.048998in}}%
\pgfpathlineto{\pgfqpoint{0.727118in}{1.064891in}}%
\pgfpathlineto{\pgfqpoint{0.729637in}{1.067888in}}%
\pgfpathlineto{\pgfqpoint{0.732156in}{1.087327in}}%
\pgfpathlineto{\pgfqpoint{0.734674in}{1.091645in}}%
\pgfpathlineto{\pgfqpoint{0.737193in}{1.090686in}}%
\pgfpathlineto{\pgfqpoint{0.739712in}{1.098746in}}%
\pgfpathlineto{\pgfqpoint{0.742231in}{1.105978in}}%
\pgfpathlineto{\pgfqpoint{0.744749in}{1.104105in}}%
\pgfpathlineto{\pgfqpoint{0.747268in}{1.112118in}}%
\pgfpathlineto{\pgfqpoint{0.749787in}{1.118500in}}%
\pgfpathlineto{\pgfqpoint{0.752306in}{1.107355in}}%
\pgfpathlineto{\pgfqpoint{0.754824in}{1.099563in}}%
\pgfpathlineto{\pgfqpoint{0.757343in}{1.079947in}}%
\pgfpathlineto{\pgfqpoint{0.759862in}{1.080729in}}%
\pgfpathlineto{\pgfqpoint{0.762381in}{1.110285in}}%
\pgfpathlineto{\pgfqpoint{0.764899in}{1.126410in}}%
\pgfpathlineto{\pgfqpoint{0.767418in}{1.109323in}}%
\pgfpathlineto{\pgfqpoint{0.769937in}{1.111659in}}%
\pgfpathlineto{\pgfqpoint{0.772456in}{1.109307in}}%
\pgfpathlineto{\pgfqpoint{0.774974in}{1.131081in}}%
\pgfpathlineto{\pgfqpoint{0.777493in}{1.130934in}}%
\pgfpathlineto{\pgfqpoint{0.780012in}{1.160025in}}%
\pgfpathlineto{\pgfqpoint{0.782531in}{1.167592in}}%
\pgfpathlineto{\pgfqpoint{0.785049in}{1.181618in}}%
\pgfpathlineto{\pgfqpoint{0.787568in}{1.181139in}}%
\pgfpathlineto{\pgfqpoint{0.790087in}{1.164967in}}%
\pgfpathlineto{\pgfqpoint{0.792606in}{1.138383in}}%
\pgfpathlineto{\pgfqpoint{0.795124in}{1.144656in}}%
\pgfpathlineto{\pgfqpoint{0.797643in}{1.133519in}}%
\pgfpathlineto{\pgfqpoint{0.800162in}{1.128644in}}%
\pgfpathlineto{\pgfqpoint{0.802681in}{1.140266in}}%
\pgfpathlineto{\pgfqpoint{0.805199in}{1.157920in}}%
\pgfpathlineto{\pgfqpoint{0.807718in}{1.133183in}}%
\pgfpathlineto{\pgfqpoint{0.810237in}{1.145283in}}%
\pgfpathlineto{\pgfqpoint{0.812756in}{1.164600in}}%
\pgfpathlineto{\pgfqpoint{0.815274in}{1.165007in}}%
\pgfpathlineto{\pgfqpoint{0.817793in}{1.162429in}}%
\pgfpathlineto{\pgfqpoint{0.820312in}{1.153125in}}%
\pgfpathlineto{\pgfqpoint{0.822831in}{1.154456in}}%
\pgfpathlineto{\pgfqpoint{0.825349in}{1.148027in}}%
\pgfpathlineto{\pgfqpoint{0.827868in}{1.136422in}}%
\pgfpathlineto{\pgfqpoint{0.830387in}{1.150288in}}%
\pgfpathlineto{\pgfqpoint{0.832906in}{1.143852in}}%
\pgfpathlineto{\pgfqpoint{0.835424in}{1.139620in}}%
\pgfpathlineto{\pgfqpoint{0.837943in}{1.138652in}}%
\pgfpathlineto{\pgfqpoint{0.840462in}{1.134668in}}%
\pgfpathlineto{\pgfqpoint{0.842981in}{1.133773in}}%
\pgfpathlineto{\pgfqpoint{0.845499in}{1.162260in}}%
\pgfpathlineto{\pgfqpoint{0.848018in}{1.167379in}}%
\pgfpathlineto{\pgfqpoint{0.850537in}{1.159096in}}%
\pgfpathlineto{\pgfqpoint{0.853056in}{1.134791in}}%
\pgfpathlineto{\pgfqpoint{0.855574in}{1.133972in}}%
\pgfpathlineto{\pgfqpoint{0.858093in}{1.156215in}}%
\pgfpathlineto{\pgfqpoint{0.860612in}{1.159551in}}%
\pgfpathlineto{\pgfqpoint{0.863131in}{1.149530in}}%
\pgfpathlineto{\pgfqpoint{0.865649in}{1.142116in}}%
\pgfpathlineto{\pgfqpoint{0.868168in}{1.141879in}}%
\pgfpathlineto{\pgfqpoint{0.870687in}{1.161544in}}%
\pgfpathlineto{\pgfqpoint{0.873206in}{1.145764in}}%
\pgfpathlineto{\pgfqpoint{0.875724in}{1.150130in}}%
\pgfpathlineto{\pgfqpoint{0.878243in}{1.151643in}}%
\pgfpathlineto{\pgfqpoint{0.880762in}{1.145500in}}%
\pgfpathlineto{\pgfqpoint{0.883281in}{1.132806in}}%
\pgfpathlineto{\pgfqpoint{0.885799in}{1.117500in}}%
\pgfpathlineto{\pgfqpoint{0.888318in}{1.105135in}}%
\pgfpathlineto{\pgfqpoint{0.890837in}{1.095900in}}%
\pgfpathlineto{\pgfqpoint{0.893356in}{1.089532in}}%
\pgfpathlineto{\pgfqpoint{0.895874in}{1.085421in}}%
\pgfpathlineto{\pgfqpoint{0.898393in}{1.075732in}}%
\pgfpathlineto{\pgfqpoint{0.900912in}{1.078252in}}%
\pgfpathlineto{\pgfqpoint{0.903431in}{1.053570in}}%
\pgfpathlineto{\pgfqpoint{0.905949in}{1.025056in}}%
\pgfpathlineto{\pgfqpoint{0.908468in}{1.020687in}}%
\pgfpathlineto{\pgfqpoint{0.910987in}{1.024977in}}%
\pgfpathlineto{\pgfqpoint{0.913506in}{1.037530in}}%
\pgfpathlineto{\pgfqpoint{0.916024in}{1.031985in}}%
\pgfpathlineto{\pgfqpoint{0.918543in}{1.037340in}}%
\pgfpathlineto{\pgfqpoint{0.921062in}{1.071197in}}%
\pgfpathlineto{\pgfqpoint{0.923581in}{1.070398in}}%
\pgfpathlineto{\pgfqpoint{0.926099in}{1.094381in}}%
\pgfpathlineto{\pgfqpoint{0.928618in}{1.083030in}}%
\pgfpathlineto{\pgfqpoint{0.931137in}{1.077156in}}%
\pgfpathlineto{\pgfqpoint{0.933656in}{1.077775in}}%
\pgfpathlineto{\pgfqpoint{0.936174in}{1.100872in}}%
\pgfpathlineto{\pgfqpoint{0.938693in}{1.101636in}}%
\pgfpathlineto{\pgfqpoint{0.941212in}{1.103665in}}%
\pgfpathlineto{\pgfqpoint{0.943731in}{1.084706in}}%
\pgfpathlineto{\pgfqpoint{0.946249in}{1.088910in}}%
\pgfpathlineto{\pgfqpoint{0.948768in}{1.099927in}}%
\pgfpathlineto{\pgfqpoint{0.951287in}{1.105056in}}%
\pgfpathlineto{\pgfqpoint{0.953806in}{1.106584in}}%
\pgfpathlineto{\pgfqpoint{0.956324in}{1.125837in}}%
\pgfpathlineto{\pgfqpoint{0.958843in}{1.118116in}}%
\pgfpathlineto{\pgfqpoint{0.961362in}{1.112112in}}%
\pgfpathlineto{\pgfqpoint{0.963881in}{1.127353in}}%
\pgfpathlineto{\pgfqpoint{0.966399in}{1.138438in}}%
\pgfpathlineto{\pgfqpoint{0.968918in}{1.159473in}}%
\pgfpathlineto{\pgfqpoint{0.971437in}{1.157487in}}%
\pgfpathlineto{\pgfqpoint{0.973956in}{1.165295in}}%
\pgfpathlineto{\pgfqpoint{0.976474in}{1.184467in}}%
\pgfpathlineto{\pgfqpoint{0.978993in}{1.174946in}}%
\pgfpathlineto{\pgfqpoint{0.981512in}{1.193212in}}%
\pgfpathlineto{\pgfqpoint{0.984031in}{1.172194in}}%
\pgfpathlineto{\pgfqpoint{0.986549in}{1.168120in}}%
\pgfpathlineto{\pgfqpoint{0.989068in}{1.179245in}}%
\pgfpathlineto{\pgfqpoint{0.991587in}{1.152898in}}%
\pgfpathlineto{\pgfqpoint{0.994106in}{1.157251in}}%
\pgfpathlineto{\pgfqpoint{0.996624in}{1.140388in}}%
\pgfpathlineto{\pgfqpoint{0.999143in}{1.129238in}}%
\pgfpathlineto{\pgfqpoint{1.001662in}{1.119726in}}%
\pgfpathlineto{\pgfqpoint{1.004181in}{1.116002in}}%
\pgfpathlineto{\pgfqpoint{1.006699in}{1.133059in}}%
\pgfpathlineto{\pgfqpoint{1.009218in}{1.130326in}}%
\pgfpathlineto{\pgfqpoint{1.011737in}{1.149840in}}%
\pgfpathlineto{\pgfqpoint{1.014256in}{1.178651in}}%
\pgfpathlineto{\pgfqpoint{1.016774in}{1.191015in}}%
\pgfpathlineto{\pgfqpoint{1.019293in}{1.194596in}}%
\pgfpathlineto{\pgfqpoint{1.021812in}{1.199874in}}%
\pgfpathlineto{\pgfqpoint{1.024331in}{1.199548in}}%
\pgfpathlineto{\pgfqpoint{1.026849in}{1.203868in}}%
\pgfpathlineto{\pgfqpoint{1.029368in}{1.205507in}}%
\pgfpathlineto{\pgfqpoint{1.031887in}{1.181117in}}%
\pgfpathlineto{\pgfqpoint{1.034406in}{1.175332in}}%
\pgfpathlineto{\pgfqpoint{1.036924in}{1.181515in}}%
\pgfpathlineto{\pgfqpoint{1.039443in}{1.199955in}}%
\pgfpathlineto{\pgfqpoint{1.041962in}{1.224694in}}%
\pgfpathlineto{\pgfqpoint{1.044481in}{1.233105in}}%
\pgfpathlineto{\pgfqpoint{1.046999in}{1.229387in}}%
\pgfpathlineto{\pgfqpoint{1.049518in}{1.219839in}}%
\pgfpathlineto{\pgfqpoint{1.052037in}{1.219970in}}%
\pgfpathlineto{\pgfqpoint{1.054556in}{1.197494in}}%
\pgfpathlineto{\pgfqpoint{1.057074in}{1.199132in}}%
\pgfpathlineto{\pgfqpoint{1.059593in}{1.227036in}}%
\pgfpathlineto{\pgfqpoint{1.062112in}{1.225285in}}%
\pgfpathlineto{\pgfqpoint{1.064631in}{1.207669in}}%
\pgfpathlineto{\pgfqpoint{1.067149in}{1.196862in}}%
\pgfpathlineto{\pgfqpoint{1.069668in}{1.189555in}}%
\pgfpathlineto{\pgfqpoint{1.072187in}{1.214840in}}%
\pgfpathlineto{\pgfqpoint{1.074706in}{1.209999in}}%
\pgfpathlineto{\pgfqpoint{1.077224in}{1.231344in}}%
\pgfpathlineto{\pgfqpoint{1.079743in}{1.245080in}}%
\pgfpathlineto{\pgfqpoint{1.082262in}{1.237889in}}%
\pgfpathlineto{\pgfqpoint{1.084781in}{1.246268in}}%
\pgfpathlineto{\pgfqpoint{1.087299in}{1.242343in}}%
\pgfpathlineto{\pgfqpoint{1.089818in}{1.258834in}}%
\pgfpathlineto{\pgfqpoint{1.092337in}{1.280271in}}%
\pgfpathlineto{\pgfqpoint{1.094856in}{1.273108in}}%
\pgfpathlineto{\pgfqpoint{1.097374in}{1.266761in}}%
\pgfpathlineto{\pgfqpoint{1.099893in}{1.288270in}}%
\pgfpathlineto{\pgfqpoint{1.102412in}{1.288438in}}%
\pgfpathlineto{\pgfqpoint{1.104931in}{1.318104in}}%
\pgfpathlineto{\pgfqpoint{1.107449in}{1.306306in}}%
\pgfpathlineto{\pgfqpoint{1.109968in}{1.319095in}}%
\pgfpathlineto{\pgfqpoint{1.112487in}{1.314027in}}%
\pgfpathlineto{\pgfqpoint{1.115006in}{1.302351in}}%
\pgfpathlineto{\pgfqpoint{1.117524in}{1.286917in}}%
\pgfpathlineto{\pgfqpoint{1.120043in}{1.267166in}}%
\pgfpathlineto{\pgfqpoint{1.122562in}{1.241199in}}%
\pgfpathlineto{\pgfqpoint{1.125081in}{1.246040in}}%
\pgfpathlineto{\pgfqpoint{1.127599in}{1.250349in}}%
\pgfpathlineto{\pgfqpoint{1.130118in}{1.255850in}}%
\pgfpathlineto{\pgfqpoint{1.132637in}{1.229360in}}%
\pgfpathlineto{\pgfqpoint{1.135156in}{1.248666in}}%
\pgfpathlineto{\pgfqpoint{1.137674in}{1.262348in}}%
\pgfpathlineto{\pgfqpoint{1.140193in}{1.292863in}}%
\pgfpathlineto{\pgfqpoint{1.142712in}{1.295001in}}%
\pgfpathlineto{\pgfqpoint{1.145231in}{1.302771in}}%
\pgfpathlineto{\pgfqpoint{1.147749in}{1.306721in}}%
\pgfpathlineto{\pgfqpoint{1.150268in}{1.315275in}}%
\pgfpathlineto{\pgfqpoint{1.152787in}{1.357987in}}%
\pgfpathlineto{\pgfqpoint{1.155306in}{1.339113in}}%
\pgfpathlineto{\pgfqpoint{1.157824in}{1.331944in}}%
\pgfpathlineto{\pgfqpoint{1.160343in}{1.327291in}}%
\pgfpathlineto{\pgfqpoint{1.162862in}{1.306951in}}%
\pgfpathlineto{\pgfqpoint{1.165381in}{1.314133in}}%
\pgfpathlineto{\pgfqpoint{1.167899in}{1.283166in}}%
\pgfpathlineto{\pgfqpoint{1.170418in}{1.282857in}}%
\pgfpathlineto{\pgfqpoint{1.172937in}{1.251477in}}%
\pgfpathlineto{\pgfqpoint{1.175456in}{1.244279in}}%
\pgfpathlineto{\pgfqpoint{1.177974in}{1.262185in}}%
\pgfpathlineto{\pgfqpoint{1.180493in}{1.289792in}}%
\pgfpathlineto{\pgfqpoint{1.183012in}{1.273214in}}%
\pgfpathlineto{\pgfqpoint{1.185531in}{1.267149in}}%
\pgfpathlineto{\pgfqpoint{1.188049in}{1.267723in}}%
\pgfpathlineto{\pgfqpoint{1.190568in}{1.266771in}}%
\pgfpathlineto{\pgfqpoint{1.193087in}{1.250256in}}%
\pgfpathlineto{\pgfqpoint{1.195606in}{1.257984in}}%
\pgfpathlineto{\pgfqpoint{1.198124in}{1.258979in}}%
\pgfpathlineto{\pgfqpoint{1.200643in}{1.277010in}}%
\pgfpathlineto{\pgfqpoint{1.203162in}{1.289291in}}%
\pgfpathlineto{\pgfqpoint{1.205681in}{1.322513in}}%
\pgfpathlineto{\pgfqpoint{1.208199in}{1.316309in}}%
\pgfpathlineto{\pgfqpoint{1.210718in}{1.320228in}}%
\pgfpathlineto{\pgfqpoint{1.213237in}{1.320179in}}%
\pgfpathlineto{\pgfqpoint{1.215756in}{1.309924in}}%
\pgfpathlineto{\pgfqpoint{1.218274in}{1.327282in}}%
\pgfpathlineto{\pgfqpoint{1.220793in}{1.328929in}}%
\pgfpathlineto{\pgfqpoint{1.223312in}{1.318617in}}%
\pgfpathlineto{\pgfqpoint{1.225831in}{1.315940in}}%
\pgfpathlineto{\pgfqpoint{1.228349in}{1.321229in}}%
\pgfpathlineto{\pgfqpoint{1.230868in}{1.331770in}}%
\pgfpathlineto{\pgfqpoint{1.233387in}{1.340134in}}%
\pgfpathlineto{\pgfqpoint{1.235906in}{1.367961in}}%
\pgfpathlineto{\pgfqpoint{1.238424in}{1.379262in}}%
\pgfpathlineto{\pgfqpoint{1.240943in}{1.382132in}}%
\pgfpathlineto{\pgfqpoint{1.243462in}{1.391349in}}%
\pgfpathlineto{\pgfqpoint{1.245981in}{1.394886in}}%
\pgfpathlineto{\pgfqpoint{1.248499in}{1.400242in}}%
\pgfpathlineto{\pgfqpoint{1.251018in}{1.392372in}}%
\pgfpathlineto{\pgfqpoint{1.253537in}{1.385868in}}%
\pgfpathlineto{\pgfqpoint{1.256056in}{1.389800in}}%
\pgfpathlineto{\pgfqpoint{1.258574in}{1.409202in}}%
\pgfpathlineto{\pgfqpoint{1.261093in}{1.427754in}}%
\pgfpathlineto{\pgfqpoint{1.263612in}{1.412572in}}%
\pgfpathlineto{\pgfqpoint{1.266131in}{1.403595in}}%
\pgfpathlineto{\pgfqpoint{1.268649in}{1.398837in}}%
\pgfpathlineto{\pgfqpoint{1.271168in}{1.399952in}}%
\pgfpathlineto{\pgfqpoint{1.273687in}{1.411516in}}%
\pgfpathlineto{\pgfqpoint{1.276206in}{1.416306in}}%
\pgfpathlineto{\pgfqpoint{1.278724in}{1.432279in}}%
\pgfpathlineto{\pgfqpoint{1.281243in}{1.429997in}}%
\pgfpathlineto{\pgfqpoint{1.283762in}{1.462353in}}%
\pgfpathlineto{\pgfqpoint{1.286281in}{1.448964in}}%
\pgfpathlineto{\pgfqpoint{1.288799in}{1.455930in}}%
\pgfpathlineto{\pgfqpoint{1.291318in}{1.460997in}}%
\pgfpathlineto{\pgfqpoint{1.293837in}{1.478052in}}%
\pgfpathlineto{\pgfqpoint{1.296356in}{1.458438in}}%
\pgfpathlineto{\pgfqpoint{1.298874in}{1.460430in}}%
\pgfpathlineto{\pgfqpoint{1.301393in}{1.457328in}}%
\pgfpathlineto{\pgfqpoint{1.303912in}{1.456150in}}%
\pgfpathlineto{\pgfqpoint{1.306431in}{1.449960in}}%
\pgfpathlineto{\pgfqpoint{1.308949in}{1.432164in}}%
\pgfpathlineto{\pgfqpoint{1.311468in}{1.435168in}}%
\pgfpathlineto{\pgfqpoint{1.313987in}{1.429411in}}%
\pgfpathlineto{\pgfqpoint{1.316506in}{1.429156in}}%
\pgfpathlineto{\pgfqpoint{1.319024in}{1.462113in}}%
\pgfpathlineto{\pgfqpoint{1.321543in}{1.445907in}}%
\pgfpathlineto{\pgfqpoint{1.324062in}{1.443392in}}%
\pgfpathlineto{\pgfqpoint{1.326581in}{1.438406in}}%
\pgfpathlineto{\pgfqpoint{1.329099in}{1.437947in}}%
\pgfpathlineto{\pgfqpoint{1.331618in}{1.430396in}}%
\pgfpathlineto{\pgfqpoint{1.334137in}{1.459917in}}%
\pgfpathlineto{\pgfqpoint{1.336656in}{1.467497in}}%
\pgfpathlineto{\pgfqpoint{1.339174in}{1.447976in}}%
\pgfpathlineto{\pgfqpoint{1.341693in}{1.454816in}}%
\pgfpathlineto{\pgfqpoint{1.344212in}{1.472777in}}%
\pgfpathlineto{\pgfqpoint{1.346731in}{1.459104in}}%
\pgfpathlineto{\pgfqpoint{1.349249in}{1.465694in}}%
\pgfpathlineto{\pgfqpoint{1.351768in}{1.483170in}}%
\pgfpathlineto{\pgfqpoint{1.354287in}{1.508399in}}%
\pgfpathlineto{\pgfqpoint{1.356806in}{1.505762in}}%
\pgfpathlineto{\pgfqpoint{1.364362in}{1.499507in}}%
\pgfpathlineto{\pgfqpoint{1.366881in}{1.494437in}}%
\pgfpathlineto{\pgfqpoint{1.369399in}{1.493882in}}%
\pgfpathlineto{\pgfqpoint{1.371918in}{1.472012in}}%
\pgfpathlineto{\pgfqpoint{1.374437in}{1.455227in}}%
\pgfpathlineto{\pgfqpoint{1.376956in}{1.476464in}}%
\pgfpathlineto{\pgfqpoint{1.379474in}{1.464235in}}%
\pgfpathlineto{\pgfqpoint{1.381993in}{1.458702in}}%
\pgfpathlineto{\pgfqpoint{1.384512in}{1.457637in}}%
\pgfpathlineto{\pgfqpoint{1.387031in}{1.479836in}}%
\pgfpathlineto{\pgfqpoint{1.389549in}{1.479004in}}%
\pgfpathlineto{\pgfqpoint{1.392068in}{1.471063in}}%
\pgfpathlineto{\pgfqpoint{1.394587in}{1.484654in}}%
\pgfpathlineto{\pgfqpoint{1.397106in}{1.475830in}}%
\pgfpathlineto{\pgfqpoint{1.399624in}{1.502219in}}%
\pgfpathlineto{\pgfqpoint{1.402143in}{1.494364in}}%
\pgfpathlineto{\pgfqpoint{1.404662in}{1.476891in}}%
\pgfpathlineto{\pgfqpoint{1.407181in}{1.485884in}}%
\pgfpathlineto{\pgfqpoint{1.409699in}{1.461603in}}%
\pgfpathlineto{\pgfqpoint{1.412218in}{1.451545in}}%
\pgfpathlineto{\pgfqpoint{1.414737in}{1.453356in}}%
\pgfpathlineto{\pgfqpoint{1.417256in}{1.468161in}}%
\pgfpathlineto{\pgfqpoint{1.419774in}{1.464795in}}%
\pgfpathlineto{\pgfqpoint{1.422293in}{1.462578in}}%
\pgfpathlineto{\pgfqpoint{1.424812in}{1.465003in}}%
\pgfpathlineto{\pgfqpoint{1.427331in}{1.451893in}}%
\pgfpathlineto{\pgfqpoint{1.429849in}{1.482197in}}%
\pgfpathlineto{\pgfqpoint{1.432368in}{1.491252in}}%
\pgfpathlineto{\pgfqpoint{1.434887in}{1.489601in}}%
\pgfpathlineto{\pgfqpoint{1.437406in}{1.473985in}}%
\pgfpathlineto{\pgfqpoint{1.439924in}{1.464866in}}%
\pgfpathlineto{\pgfqpoint{1.442443in}{1.461531in}}%
\pgfpathlineto{\pgfqpoint{1.444962in}{1.460444in}}%
\pgfpathlineto{\pgfqpoint{1.447481in}{1.435036in}}%
\pgfpathlineto{\pgfqpoint{1.449999in}{1.429875in}}%
\pgfpathlineto{\pgfqpoint{1.452518in}{1.443980in}}%
\pgfpathlineto{\pgfqpoint{1.455037in}{1.455715in}}%
\pgfpathlineto{\pgfqpoint{1.457556in}{1.441166in}}%
\pgfpathlineto{\pgfqpoint{1.460074in}{1.446285in}}%
\pgfpathlineto{\pgfqpoint{1.462593in}{1.439679in}}%
\pgfpathlineto{\pgfqpoint{1.465112in}{1.431968in}}%
\pgfpathlineto{\pgfqpoint{1.467631in}{1.434942in}}%
\pgfpathlineto{\pgfqpoint{1.470149in}{1.420700in}}%
\pgfpathlineto{\pgfqpoint{1.472668in}{1.420112in}}%
\pgfpathlineto{\pgfqpoint{1.475187in}{1.418526in}}%
\pgfpathlineto{\pgfqpoint{1.477706in}{1.419944in}}%
\pgfpathlineto{\pgfqpoint{1.480224in}{1.409725in}}%
\pgfpathlineto{\pgfqpoint{1.482743in}{1.411234in}}%
\pgfpathlineto{\pgfqpoint{1.485262in}{1.411161in}}%
\pgfpathlineto{\pgfqpoint{1.487781in}{1.439656in}}%
\pgfpathlineto{\pgfqpoint{1.490299in}{1.427834in}}%
\pgfpathlineto{\pgfqpoint{1.492818in}{1.401060in}}%
\pgfpathlineto{\pgfqpoint{1.495337in}{1.396536in}}%
\pgfpathlineto{\pgfqpoint{1.497856in}{1.387242in}}%
\pgfpathlineto{\pgfqpoint{1.500374in}{1.393157in}}%
\pgfpathlineto{\pgfqpoint{1.502893in}{1.394788in}}%
\pgfpathlineto{\pgfqpoint{1.505412in}{1.395773in}}%
\pgfpathlineto{\pgfqpoint{1.507931in}{1.421155in}}%
\pgfpathlineto{\pgfqpoint{1.510449in}{1.402619in}}%
\pgfpathlineto{\pgfqpoint{1.512968in}{1.421166in}}%
\pgfpathlineto{\pgfqpoint{1.515487in}{1.411430in}}%
\pgfpathlineto{\pgfqpoint{1.518006in}{1.416025in}}%
\pgfpathlineto{\pgfqpoint{1.520524in}{1.451399in}}%
\pgfpathlineto{\pgfqpoint{1.523043in}{1.438192in}}%
\pgfpathlineto{\pgfqpoint{1.525562in}{1.446731in}}%
\pgfpathlineto{\pgfqpoint{1.528081in}{1.444681in}}%
\pgfpathlineto{\pgfqpoint{1.530599in}{1.441148in}}%
\pgfpathlineto{\pgfqpoint{1.533118in}{1.397371in}}%
\pgfpathlineto{\pgfqpoint{1.535637in}{1.404211in}}%
\pgfpathlineto{\pgfqpoint{1.538156in}{1.410636in}}%
\pgfpathlineto{\pgfqpoint{1.540674in}{1.411899in}}%
\pgfpathlineto{\pgfqpoint{1.543193in}{1.423037in}}%
\pgfpathlineto{\pgfqpoint{1.545712in}{1.423936in}}%
\pgfpathlineto{\pgfqpoint{1.548231in}{1.415325in}}%
\pgfpathlineto{\pgfqpoint{1.550749in}{1.426860in}}%
\pgfpathlineto{\pgfqpoint{1.553268in}{1.419347in}}%
\pgfpathlineto{\pgfqpoint{1.555787in}{1.404241in}}%
\pgfpathlineto{\pgfqpoint{1.558306in}{1.406299in}}%
\pgfpathlineto{\pgfqpoint{1.560824in}{1.395443in}}%
\pgfpathlineto{\pgfqpoint{1.563343in}{1.377247in}}%
\pgfpathlineto{\pgfqpoint{1.565862in}{1.369433in}}%
\pgfpathlineto{\pgfqpoint{1.568381in}{1.359509in}}%
\pgfpathlineto{\pgfqpoint{1.570899in}{1.366116in}}%
\pgfpathlineto{\pgfqpoint{1.575937in}{1.332478in}}%
\pgfpathlineto{\pgfqpoint{1.578456in}{1.339698in}}%
\pgfpathlineto{\pgfqpoint{1.580974in}{1.324923in}}%
\pgfpathlineto{\pgfqpoint{1.583493in}{1.326705in}}%
\pgfpathlineto{\pgfqpoint{1.586012in}{1.328058in}}%
\pgfpathlineto{\pgfqpoint{1.588531in}{1.320387in}}%
\pgfpathlineto{\pgfqpoint{1.591049in}{1.327318in}}%
\pgfpathlineto{\pgfqpoint{1.593568in}{1.331372in}}%
\pgfpathlineto{\pgfqpoint{1.596087in}{1.335974in}}%
\pgfpathlineto{\pgfqpoint{1.598606in}{1.324436in}}%
\pgfpathlineto{\pgfqpoint{1.601124in}{1.340547in}}%
\pgfpathlineto{\pgfqpoint{1.603643in}{1.335626in}}%
\pgfpathlineto{\pgfqpoint{1.606162in}{1.338803in}}%
\pgfpathlineto{\pgfqpoint{1.608681in}{1.363135in}}%
\pgfpathlineto{\pgfqpoint{1.611199in}{1.334464in}}%
\pgfpathlineto{\pgfqpoint{1.613718in}{1.337958in}}%
\pgfpathlineto{\pgfqpoint{1.616237in}{1.320467in}}%
\pgfpathlineto{\pgfqpoint{1.618756in}{1.308324in}}%
\pgfpathlineto{\pgfqpoint{1.621274in}{1.284241in}}%
\pgfpathlineto{\pgfqpoint{1.623793in}{1.262410in}}%
\pgfpathlineto{\pgfqpoint{1.626312in}{1.258715in}}%
\pgfpathlineto{\pgfqpoint{1.628831in}{1.273628in}}%
\pgfpathlineto{\pgfqpoint{1.631349in}{1.286983in}}%
\pgfpathlineto{\pgfqpoint{1.633868in}{1.294970in}}%
\pgfpathlineto{\pgfqpoint{1.636387in}{1.292862in}}%
\pgfpathlineto{\pgfqpoint{1.638906in}{1.321142in}}%
\pgfpathlineto{\pgfqpoint{1.641424in}{1.311055in}}%
\pgfpathlineto{\pgfqpoint{1.643943in}{1.293860in}}%
\pgfpathlineto{\pgfqpoint{1.646462in}{1.312714in}}%
\pgfpathlineto{\pgfqpoint{1.648981in}{1.312379in}}%
\pgfpathlineto{\pgfqpoint{1.651499in}{1.311227in}}%
\pgfpathlineto{\pgfqpoint{1.654018in}{1.323937in}}%
\pgfpathlineto{\pgfqpoint{1.656537in}{1.328554in}}%
\pgfpathlineto{\pgfqpoint{1.659056in}{1.326966in}}%
\pgfpathlineto{\pgfqpoint{1.661574in}{1.323158in}}%
\pgfpathlineto{\pgfqpoint{1.664093in}{1.330586in}}%
\pgfpathlineto{\pgfqpoint{1.666612in}{1.354469in}}%
\pgfpathlineto{\pgfqpoint{1.669131in}{1.363752in}}%
\pgfpathlineto{\pgfqpoint{1.671649in}{1.370976in}}%
\pgfpathlineto{\pgfqpoint{1.674168in}{1.371379in}}%
\pgfpathlineto{\pgfqpoint{1.676687in}{1.349351in}}%
\pgfpathlineto{\pgfqpoint{1.679206in}{1.330653in}}%
\pgfpathlineto{\pgfqpoint{1.681724in}{1.344129in}}%
\pgfpathlineto{\pgfqpoint{1.684243in}{1.335891in}}%
\pgfpathlineto{\pgfqpoint{1.686762in}{1.318526in}}%
\pgfpathlineto{\pgfqpoint{1.689281in}{1.325520in}}%
\pgfpathlineto{\pgfqpoint{1.691799in}{1.317196in}}%
\pgfpathlineto{\pgfqpoint{1.694318in}{1.318998in}}%
\pgfpathlineto{\pgfqpoint{1.696837in}{1.319380in}}%
\pgfpathlineto{\pgfqpoint{1.699356in}{1.333793in}}%
\pgfpathlineto{\pgfqpoint{1.701874in}{1.320545in}}%
\pgfpathlineto{\pgfqpoint{1.704393in}{1.310830in}}%
\pgfpathlineto{\pgfqpoint{1.706912in}{1.299839in}}%
\pgfpathlineto{\pgfqpoint{1.709431in}{1.321854in}}%
\pgfpathlineto{\pgfqpoint{1.711949in}{1.341317in}}%
\pgfpathlineto{\pgfqpoint{1.714468in}{1.325619in}}%
\pgfpathlineto{\pgfqpoint{1.716987in}{1.317888in}}%
\pgfpathlineto{\pgfqpoint{1.719506in}{1.328581in}}%
\pgfpathlineto{\pgfqpoint{1.722024in}{1.333437in}}%
\pgfpathlineto{\pgfqpoint{1.724543in}{1.348653in}}%
\pgfpathlineto{\pgfqpoint{1.727062in}{1.349117in}}%
\pgfpathlineto{\pgfqpoint{1.729581in}{1.352391in}}%
\pgfpathlineto{\pgfqpoint{1.732099in}{1.356659in}}%
\pgfpathlineto{\pgfqpoint{1.734618in}{1.361173in}}%
\pgfpathlineto{\pgfqpoint{1.737137in}{1.377429in}}%
\pgfpathlineto{\pgfqpoint{1.739656in}{1.376032in}}%
\pgfpathlineto{\pgfqpoint{1.742174in}{1.357401in}}%
\pgfpathlineto{\pgfqpoint{1.744693in}{1.349360in}}%
\pgfpathlineto{\pgfqpoint{1.747212in}{1.354392in}}%
\pgfpathlineto{\pgfqpoint{1.749731in}{1.356825in}}%
\pgfpathlineto{\pgfqpoint{1.752249in}{1.343311in}}%
\pgfpathlineto{\pgfqpoint{1.754768in}{1.339007in}}%
\pgfpathlineto{\pgfqpoint{1.757287in}{1.350279in}}%
\pgfpathlineto{\pgfqpoint{1.759806in}{1.329958in}}%
\pgfpathlineto{\pgfqpoint{1.762324in}{1.317758in}}%
\pgfpathlineto{\pgfqpoint{1.764843in}{1.324641in}}%
\pgfpathlineto{\pgfqpoint{1.767362in}{1.330243in}}%
\pgfpathlineto{\pgfqpoint{1.769881in}{1.346701in}}%
\pgfpathlineto{\pgfqpoint{1.772399in}{1.375345in}}%
\pgfpathlineto{\pgfqpoint{1.774918in}{1.406412in}}%
\pgfpathlineto{\pgfqpoint{1.777437in}{1.421837in}}%
\pgfpathlineto{\pgfqpoint{1.779956in}{1.435331in}}%
\pgfpathlineto{\pgfqpoint{1.782474in}{1.438164in}}%
\pgfpathlineto{\pgfqpoint{1.784993in}{1.447746in}}%
\pgfpathlineto{\pgfqpoint{1.787512in}{1.432030in}}%
\pgfpathlineto{\pgfqpoint{1.790031in}{1.428362in}}%
\pgfpathlineto{\pgfqpoint{1.792549in}{1.444111in}}%
\pgfpathlineto{\pgfqpoint{1.795068in}{1.426989in}}%
\pgfpathlineto{\pgfqpoint{1.797587in}{1.434174in}}%
\pgfpathlineto{\pgfqpoint{1.800106in}{1.446796in}}%
\pgfpathlineto{\pgfqpoint{1.802624in}{1.455632in}}%
\pgfpathlineto{\pgfqpoint{1.805143in}{1.461739in}}%
\pgfpathlineto{\pgfqpoint{1.807662in}{1.473087in}}%
\pgfpathlineto{\pgfqpoint{1.810181in}{1.448464in}}%
\pgfpathlineto{\pgfqpoint{1.812699in}{1.472484in}}%
\pgfpathlineto{\pgfqpoint{1.815218in}{1.489916in}}%
\pgfpathlineto{\pgfqpoint{1.817737in}{1.479745in}}%
\pgfpathlineto{\pgfqpoint{1.820256in}{1.509508in}}%
\pgfpathlineto{\pgfqpoint{1.822774in}{1.531946in}}%
\pgfpathlineto{\pgfqpoint{1.825293in}{1.532105in}}%
\pgfpathlineto{\pgfqpoint{1.827812in}{1.536195in}}%
\pgfpathlineto{\pgfqpoint{1.830331in}{1.523018in}}%
\pgfpathlineto{\pgfqpoint{1.832849in}{1.534205in}}%
\pgfpathlineto{\pgfqpoint{1.835368in}{1.530575in}}%
\pgfpathlineto{\pgfqpoint{1.837887in}{1.531737in}}%
\pgfpathlineto{\pgfqpoint{1.840406in}{1.539399in}}%
\pgfpathlineto{\pgfqpoint{1.842924in}{1.530532in}}%
\pgfpathlineto{\pgfqpoint{1.845443in}{1.552439in}}%
\pgfpathlineto{\pgfqpoint{1.847962in}{1.540182in}}%
\pgfpathlineto{\pgfqpoint{1.850481in}{1.571920in}}%
\pgfpathlineto{\pgfqpoint{1.852999in}{1.598231in}}%
\pgfpathlineto{\pgfqpoint{1.855518in}{1.586108in}}%
\pgfpathlineto{\pgfqpoint{1.858037in}{1.578708in}}%
\pgfpathlineto{\pgfqpoint{1.860556in}{1.577548in}}%
\pgfpathlineto{\pgfqpoint{1.863074in}{1.584270in}}%
\pgfpathlineto{\pgfqpoint{1.865593in}{1.569067in}}%
\pgfpathlineto{\pgfqpoint{1.868112in}{1.582831in}}%
\pgfpathlineto{\pgfqpoint{1.870631in}{1.573521in}}%
\pgfpathlineto{\pgfqpoint{1.873149in}{1.545834in}}%
\pgfpathlineto{\pgfqpoint{1.875668in}{1.534690in}}%
\pgfpathlineto{\pgfqpoint{1.878187in}{1.512722in}}%
\pgfpathlineto{\pgfqpoint{1.880706in}{1.481441in}}%
\pgfpathlineto{\pgfqpoint{1.883224in}{1.481850in}}%
\pgfpathlineto{\pgfqpoint{1.885743in}{1.482136in}}%
\pgfpathlineto{\pgfqpoint{1.888262in}{1.477894in}}%
\pgfpathlineto{\pgfqpoint{1.890781in}{1.484646in}}%
\pgfpathlineto{\pgfqpoint{1.893299in}{1.509567in}}%
\pgfpathlineto{\pgfqpoint{1.895818in}{1.544597in}}%
\pgfpathlineto{\pgfqpoint{1.898337in}{1.547902in}}%
\pgfpathlineto{\pgfqpoint{1.903374in}{1.514169in}}%
\pgfpathlineto{\pgfqpoint{1.905893in}{1.533351in}}%
\pgfpathlineto{\pgfqpoint{1.908412in}{1.550634in}}%
\pgfpathlineto{\pgfqpoint{1.910931in}{1.537979in}}%
\pgfpathlineto{\pgfqpoint{1.913449in}{1.521662in}}%
\pgfpathlineto{\pgfqpoint{1.915968in}{1.544386in}}%
\pgfpathlineto{\pgfqpoint{1.918487in}{1.533388in}}%
\pgfpathlineto{\pgfqpoint{1.921006in}{1.541279in}}%
\pgfpathlineto{\pgfqpoint{1.923524in}{1.544540in}}%
\pgfpathlineto{\pgfqpoint{1.926043in}{1.529980in}}%
\pgfpathlineto{\pgfqpoint{1.928562in}{1.488052in}}%
\pgfpathlineto{\pgfqpoint{1.931081in}{1.470285in}}%
\pgfpathlineto{\pgfqpoint{1.933599in}{1.490870in}}%
\pgfpathlineto{\pgfqpoint{1.936118in}{1.496522in}}%
\pgfpathlineto{\pgfqpoint{1.938637in}{1.499086in}}%
\pgfpathlineto{\pgfqpoint{1.941156in}{1.503255in}}%
\pgfpathlineto{\pgfqpoint{1.943674in}{1.493693in}}%
\pgfpathlineto{\pgfqpoint{1.946193in}{1.504849in}}%
\pgfpathlineto{\pgfqpoint{1.948712in}{1.509344in}}%
\pgfpathlineto{\pgfqpoint{1.951231in}{1.520682in}}%
\pgfpathlineto{\pgfqpoint{1.953749in}{1.534683in}}%
\pgfpathlineto{\pgfqpoint{1.956268in}{1.543269in}}%
\pgfpathlineto{\pgfqpoint{1.958787in}{1.544693in}}%
\pgfpathlineto{\pgfqpoint{1.961306in}{1.532107in}}%
\pgfpathlineto{\pgfqpoint{1.963824in}{1.515531in}}%
\pgfpathlineto{\pgfqpoint{1.966343in}{1.531962in}}%
\pgfpathlineto{\pgfqpoint{1.968862in}{1.541967in}}%
\pgfpathlineto{\pgfqpoint{1.971381in}{1.548415in}}%
\pgfpathlineto{\pgfqpoint{1.973899in}{1.532059in}}%
\pgfpathlineto{\pgfqpoint{1.976418in}{1.541585in}}%
\pgfpathlineto{\pgfqpoint{1.978937in}{1.518976in}}%
\pgfpathlineto{\pgfqpoint{1.981456in}{1.516656in}}%
\pgfpathlineto{\pgfqpoint{1.983974in}{1.519788in}}%
\pgfpathlineto{\pgfqpoint{1.986493in}{1.504383in}}%
\pgfpathlineto{\pgfqpoint{1.989012in}{1.496296in}}%
\pgfpathlineto{\pgfqpoint{1.991531in}{1.463075in}}%
\pgfpathlineto{\pgfqpoint{1.994049in}{1.461541in}}%
\pgfpathlineto{\pgfqpoint{1.996568in}{1.487493in}}%
\pgfpathlineto{\pgfqpoint{1.999087in}{1.490052in}}%
\pgfpathlineto{\pgfqpoint{2.001606in}{1.513280in}}%
\pgfpathlineto{\pgfqpoint{2.004124in}{1.522508in}}%
\pgfpathlineto{\pgfqpoint{2.006643in}{1.508274in}}%
\pgfpathlineto{\pgfqpoint{2.009162in}{1.526195in}}%
\pgfpathlineto{\pgfqpoint{2.011681in}{1.541768in}}%
\pgfpathlineto{\pgfqpoint{2.014199in}{1.533345in}}%
\pgfpathlineto{\pgfqpoint{2.016718in}{1.541995in}}%
\pgfpathlineto{\pgfqpoint{2.019237in}{1.540930in}}%
\pgfpathlineto{\pgfqpoint{2.021756in}{1.548682in}}%
\pgfpathlineto{\pgfqpoint{2.024274in}{1.549276in}}%
\pgfpathlineto{\pgfqpoint{2.026793in}{1.537508in}}%
\pgfpathlineto{\pgfqpoint{2.029312in}{1.550283in}}%
\pgfpathlineto{\pgfqpoint{2.031831in}{1.540164in}}%
\pgfpathlineto{\pgfqpoint{2.034349in}{1.542439in}}%
\pgfpathlineto{\pgfqpoint{2.036868in}{1.541142in}}%
\pgfpathlineto{\pgfqpoint{2.039387in}{1.555879in}}%
\pgfpathlineto{\pgfqpoint{2.041906in}{1.551366in}}%
\pgfpathlineto{\pgfqpoint{2.044424in}{1.547733in}}%
\pgfpathlineto{\pgfqpoint{2.046943in}{1.556888in}}%
\pgfpathlineto{\pgfqpoint{2.049462in}{1.563530in}}%
\pgfpathlineto{\pgfqpoint{2.051981in}{1.586182in}}%
\pgfpathlineto{\pgfqpoint{2.054499in}{1.589444in}}%
\pgfpathlineto{\pgfqpoint{2.057018in}{1.624693in}}%
\pgfpathlineto{\pgfqpoint{2.059537in}{1.620874in}}%
\pgfpathlineto{\pgfqpoint{2.062056in}{1.641878in}}%
\pgfpathlineto{\pgfqpoint{2.064574in}{1.615320in}}%
\pgfpathlineto{\pgfqpoint{2.067093in}{1.616564in}}%
\pgfpathlineto{\pgfqpoint{2.069612in}{1.603195in}}%
\pgfpathlineto{\pgfqpoint{2.072131in}{1.598806in}}%
\pgfpathlineto{\pgfqpoint{2.074649in}{1.603840in}}%
\pgfpathlineto{\pgfqpoint{2.077168in}{1.618754in}}%
\pgfpathlineto{\pgfqpoint{2.079687in}{1.608988in}}%
\pgfpathlineto{\pgfqpoint{2.082206in}{1.581744in}}%
\pgfpathlineto{\pgfqpoint{2.084724in}{1.590533in}}%
\pgfpathlineto{\pgfqpoint{2.087243in}{1.601465in}}%
\pgfpathlineto{\pgfqpoint{2.089762in}{1.577675in}}%
\pgfpathlineto{\pgfqpoint{2.092281in}{1.577234in}}%
\pgfpathlineto{\pgfqpoint{2.094799in}{1.572718in}}%
\pgfpathlineto{\pgfqpoint{2.097318in}{1.581842in}}%
\pgfpathlineto{\pgfqpoint{2.099837in}{1.565996in}}%
\pgfpathlineto{\pgfqpoint{2.102356in}{1.559138in}}%
\pgfpathlineto{\pgfqpoint{2.104874in}{1.561971in}}%
\pgfpathlineto{\pgfqpoint{2.107393in}{1.575546in}}%
\pgfpathlineto{\pgfqpoint{2.109912in}{1.577361in}}%
\pgfpathlineto{\pgfqpoint{2.112431in}{1.581279in}}%
\pgfpathlineto{\pgfqpoint{2.114949in}{1.594395in}}%
\pgfpathlineto{\pgfqpoint{2.117468in}{1.612861in}}%
\pgfpathlineto{\pgfqpoint{2.119987in}{1.625115in}}%
\pgfpathlineto{\pgfqpoint{2.122506in}{1.640296in}}%
\pgfpathlineto{\pgfqpoint{2.125024in}{1.641034in}}%
\pgfpathlineto{\pgfqpoint{2.127543in}{1.636153in}}%
\pgfpathlineto{\pgfqpoint{2.130062in}{1.620113in}}%
\pgfpathlineto{\pgfqpoint{2.132581in}{1.649621in}}%
\pgfpathlineto{\pgfqpoint{2.135099in}{1.630996in}}%
\pgfpathlineto{\pgfqpoint{2.137618in}{1.660714in}}%
\pgfpathlineto{\pgfqpoint{2.140137in}{1.645326in}}%
\pgfpathlineto{\pgfqpoint{2.142656in}{1.647183in}}%
\pgfpathlineto{\pgfqpoint{2.145174in}{1.637682in}}%
\pgfpathlineto{\pgfqpoint{2.147693in}{1.634433in}}%
\pgfpathlineto{\pgfqpoint{2.150212in}{1.635472in}}%
\pgfpathlineto{\pgfqpoint{2.152731in}{1.632830in}}%
\pgfpathlineto{\pgfqpoint{2.155249in}{1.653444in}}%
\pgfpathlineto{\pgfqpoint{2.157768in}{1.681122in}}%
\pgfpathlineto{\pgfqpoint{2.160287in}{1.669897in}}%
\pgfpathlineto{\pgfqpoint{2.162806in}{1.662148in}}%
\pgfpathlineto{\pgfqpoint{2.165324in}{1.682085in}}%
\pgfpathlineto{\pgfqpoint{2.167843in}{1.704800in}}%
\pgfpathlineto{\pgfqpoint{2.170362in}{1.695320in}}%
\pgfpathlineto{\pgfqpoint{2.172881in}{1.704069in}}%
\pgfpathlineto{\pgfqpoint{2.175399in}{1.679903in}}%
\pgfpathlineto{\pgfqpoint{2.177918in}{1.690055in}}%
\pgfpathlineto{\pgfqpoint{2.180437in}{1.669230in}}%
\pgfpathlineto{\pgfqpoint{2.182956in}{1.701316in}}%
\pgfpathlineto{\pgfqpoint{2.185474in}{1.706347in}}%
\pgfpathlineto{\pgfqpoint{2.187993in}{1.726002in}}%
\pgfpathlineto{\pgfqpoint{2.190512in}{1.722042in}}%
\pgfpathlineto{\pgfqpoint{2.195549in}{1.777961in}}%
\pgfpathlineto{\pgfqpoint{2.198068in}{1.767374in}}%
\pgfpathlineto{\pgfqpoint{2.200587in}{1.773477in}}%
\pgfpathlineto{\pgfqpoint{2.203106in}{1.767713in}}%
\pgfpathlineto{\pgfqpoint{2.205624in}{1.747669in}}%
\pgfpathlineto{\pgfqpoint{2.208143in}{1.742494in}}%
\pgfpathlineto{\pgfqpoint{2.213181in}{1.752091in}}%
\pgfpathlineto{\pgfqpoint{2.215699in}{1.719865in}}%
\pgfpathlineto{\pgfqpoint{2.218218in}{1.735562in}}%
\pgfpathlineto{\pgfqpoint{2.220737in}{1.738991in}}%
\pgfpathlineto{\pgfqpoint{2.223256in}{1.754775in}}%
\pgfpathlineto{\pgfqpoint{2.225774in}{1.772649in}}%
\pgfpathlineto{\pgfqpoint{2.228293in}{1.787151in}}%
\pgfpathlineto{\pgfqpoint{2.230812in}{1.802386in}}%
\pgfpathlineto{\pgfqpoint{2.233331in}{1.805489in}}%
\pgfpathlineto{\pgfqpoint{2.235849in}{1.810972in}}%
\pgfpathlineto{\pgfqpoint{2.238368in}{1.828149in}}%
\pgfpathlineto{\pgfqpoint{2.240887in}{1.835045in}}%
\pgfpathlineto{\pgfqpoint{2.243406in}{1.815324in}}%
\pgfpathlineto{\pgfqpoint{2.245924in}{1.796818in}}%
\pgfpathlineto{\pgfqpoint{2.248443in}{1.802522in}}%
\pgfpathlineto{\pgfqpoint{2.250962in}{1.800231in}}%
\pgfpathlineto{\pgfqpoint{2.253481in}{1.789939in}}%
\pgfpathlineto{\pgfqpoint{2.255999in}{1.796830in}}%
\pgfpathlineto{\pgfqpoint{2.258518in}{1.810529in}}%
\pgfpathlineto{\pgfqpoint{2.261037in}{1.819179in}}%
\pgfpathlineto{\pgfqpoint{2.263556in}{1.801408in}}%
\pgfpathlineto{\pgfqpoint{2.266074in}{1.822481in}}%
\pgfpathlineto{\pgfqpoint{2.268593in}{1.852388in}}%
\pgfpathlineto{\pgfqpoint{2.271112in}{1.848605in}}%
\pgfpathlineto{\pgfqpoint{2.273631in}{1.868533in}}%
\pgfpathlineto{\pgfqpoint{2.276149in}{1.871264in}}%
\pgfpathlineto{\pgfqpoint{2.278668in}{1.856632in}}%
\pgfpathlineto{\pgfqpoint{2.281187in}{1.866442in}}%
\pgfpathlineto{\pgfqpoint{2.283706in}{1.889648in}}%
\pgfpathlineto{\pgfqpoint{2.286224in}{1.929034in}}%
\pgfpathlineto{\pgfqpoint{2.288743in}{1.917363in}}%
\pgfpathlineto{\pgfqpoint{2.291262in}{1.915134in}}%
\pgfpathlineto{\pgfqpoint{2.293781in}{1.908399in}}%
\pgfpathlineto{\pgfqpoint{2.296299in}{1.893910in}}%
\pgfpathlineto{\pgfqpoint{2.301337in}{1.845436in}}%
\pgfpathlineto{\pgfqpoint{2.303856in}{1.849064in}}%
\pgfpathlineto{\pgfqpoint{2.306374in}{1.864333in}}%
\pgfpathlineto{\pgfqpoint{2.308893in}{1.895047in}}%
\pgfpathlineto{\pgfqpoint{2.311412in}{1.894032in}}%
\pgfpathlineto{\pgfqpoint{2.313931in}{1.897920in}}%
\pgfpathlineto{\pgfqpoint{2.316449in}{1.892586in}}%
\pgfpathlineto{\pgfqpoint{2.321487in}{1.909718in}}%
\pgfpathlineto{\pgfqpoint{2.324006in}{1.956806in}}%
\pgfpathlineto{\pgfqpoint{2.326524in}{1.959297in}}%
\pgfpathlineto{\pgfqpoint{2.329043in}{1.950478in}}%
\pgfpathlineto{\pgfqpoint{2.331562in}{1.958238in}}%
\pgfpathlineto{\pgfqpoint{2.334081in}{1.992248in}}%
\pgfpathlineto{\pgfqpoint{2.336599in}{1.959849in}}%
\pgfpathlineto{\pgfqpoint{2.339118in}{1.955353in}}%
\pgfpathlineto{\pgfqpoint{2.341637in}{1.977421in}}%
\pgfpathlineto{\pgfqpoint{2.344156in}{1.972360in}}%
\pgfpathlineto{\pgfqpoint{2.346674in}{2.004335in}}%
\pgfpathlineto{\pgfqpoint{2.349193in}{1.970134in}}%
\pgfpathlineto{\pgfqpoint{2.351712in}{1.971843in}}%
\pgfpathlineto{\pgfqpoint{2.354231in}{1.993181in}}%
\pgfpathlineto{\pgfqpoint{2.356749in}{1.965337in}}%
\pgfpathlineto{\pgfqpoint{2.359268in}{1.985483in}}%
\pgfpathlineto{\pgfqpoint{2.361787in}{1.982592in}}%
\pgfpathlineto{\pgfqpoint{2.364306in}{1.976814in}}%
\pgfpathlineto{\pgfqpoint{2.366824in}{1.990565in}}%
\pgfpathlineto{\pgfqpoint{2.369343in}{1.998368in}}%
\pgfpathlineto{\pgfqpoint{2.371862in}{1.969652in}}%
\pgfpathlineto{\pgfqpoint{2.374381in}{1.949509in}}%
\pgfpathlineto{\pgfqpoint{2.376899in}{1.970030in}}%
\pgfpathlineto{\pgfqpoint{2.379418in}{1.969482in}}%
\pgfpathlineto{\pgfqpoint{2.381937in}{1.940759in}}%
\pgfpathlineto{\pgfqpoint{2.384456in}{1.963537in}}%
\pgfpathlineto{\pgfqpoint{2.386974in}{1.964658in}}%
\pgfpathlineto{\pgfqpoint{2.389493in}{1.979072in}}%
\pgfpathlineto{\pgfqpoint{2.392012in}{1.969171in}}%
\pgfpathlineto{\pgfqpoint{2.394531in}{1.955011in}}%
\pgfpathlineto{\pgfqpoint{2.397049in}{1.957512in}}%
\pgfpathlineto{\pgfqpoint{2.399568in}{1.971290in}}%
\pgfpathlineto{\pgfqpoint{2.402087in}{1.988838in}}%
\pgfpathlineto{\pgfqpoint{2.404606in}{1.992367in}}%
\pgfpathlineto{\pgfqpoint{2.407124in}{1.957587in}}%
\pgfpathlineto{\pgfqpoint{2.409643in}{1.958505in}}%
\pgfpathlineto{\pgfqpoint{2.412162in}{1.954410in}}%
\pgfpathlineto{\pgfqpoint{2.414681in}{1.954088in}}%
\pgfpathlineto{\pgfqpoint{2.417199in}{1.956576in}}%
\pgfpathlineto{\pgfqpoint{2.419718in}{1.931653in}}%
\pgfpathlineto{\pgfqpoint{2.422237in}{1.929323in}}%
\pgfpathlineto{\pgfqpoint{2.424756in}{1.906426in}}%
\pgfpathlineto{\pgfqpoint{2.427274in}{1.894507in}}%
\pgfpathlineto{\pgfqpoint{2.429793in}{1.915358in}}%
\pgfpathlineto{\pgfqpoint{2.432312in}{1.934987in}}%
\pgfpathlineto{\pgfqpoint{2.434831in}{1.936718in}}%
\pgfpathlineto{\pgfqpoint{2.437349in}{1.920016in}}%
\pgfpathlineto{\pgfqpoint{2.439868in}{1.933373in}}%
\pgfpathlineto{\pgfqpoint{2.442387in}{1.919754in}}%
\pgfpathlineto{\pgfqpoint{2.444906in}{1.941967in}}%
\pgfpathlineto{\pgfqpoint{2.449943in}{1.925717in}}%
\pgfpathlineto{\pgfqpoint{2.452462in}{1.961860in}}%
\pgfpathlineto{\pgfqpoint{2.454981in}{1.967621in}}%
\pgfpathlineto{\pgfqpoint{2.457499in}{1.981162in}}%
\pgfpathlineto{\pgfqpoint{2.460018in}{1.969418in}}%
\pgfpathlineto{\pgfqpoint{2.462537in}{1.978309in}}%
\pgfpathlineto{\pgfqpoint{2.465056in}{1.983454in}}%
\pgfpathlineto{\pgfqpoint{2.467574in}{1.996817in}}%
\pgfpathlineto{\pgfqpoint{2.470093in}{1.972577in}}%
\pgfpathlineto{\pgfqpoint{2.472612in}{1.990681in}}%
\pgfpathlineto{\pgfqpoint{2.475131in}{1.994080in}}%
\pgfpathlineto{\pgfqpoint{2.477649in}{1.995869in}}%
\pgfpathlineto{\pgfqpoint{2.480168in}{1.983869in}}%
\pgfpathlineto{\pgfqpoint{2.482687in}{1.975445in}}%
\pgfpathlineto{\pgfqpoint{2.485206in}{1.963984in}}%
\pgfpathlineto{\pgfqpoint{2.487724in}{1.971248in}}%
\pgfpathlineto{\pgfqpoint{2.490243in}{1.989652in}}%
\pgfpathlineto{\pgfqpoint{2.492762in}{1.996445in}}%
\pgfpathlineto{\pgfqpoint{2.495281in}{2.000280in}}%
\pgfpathlineto{\pgfqpoint{2.497799in}{2.007688in}}%
\pgfpathlineto{\pgfqpoint{2.500318in}{1.974639in}}%
\pgfpathlineto{\pgfqpoint{2.502837in}{1.918343in}}%
\pgfpathlineto{\pgfqpoint{2.505356in}{1.928952in}}%
\pgfpathlineto{\pgfqpoint{2.507874in}{1.937894in}}%
\pgfpathlineto{\pgfqpoint{2.510393in}{1.931758in}}%
\pgfpathlineto{\pgfqpoint{2.512912in}{1.923624in}}%
\pgfpathlineto{\pgfqpoint{2.515431in}{1.936238in}}%
\pgfpathlineto{\pgfqpoint{2.517949in}{1.933554in}}%
\pgfpathlineto{\pgfqpoint{2.520468in}{1.939862in}}%
\pgfpathlineto{\pgfqpoint{2.522987in}{1.935845in}}%
\pgfpathlineto{\pgfqpoint{2.525506in}{1.947283in}}%
\pgfpathlineto{\pgfqpoint{2.528024in}{1.945753in}}%
\pgfpathlineto{\pgfqpoint{2.530543in}{1.964209in}}%
\pgfpathlineto{\pgfqpoint{2.533062in}{1.951496in}}%
\pgfpathlineto{\pgfqpoint{2.535581in}{1.990477in}}%
\pgfpathlineto{\pgfqpoint{2.538099in}{1.996972in}}%
\pgfpathlineto{\pgfqpoint{2.540618in}{1.970982in}}%
\pgfpathlineto{\pgfqpoint{2.543137in}{1.976273in}}%
\pgfpathlineto{\pgfqpoint{2.545656in}{1.979298in}}%
\pgfpathlineto{\pgfqpoint{2.548174in}{1.993155in}}%
\pgfpathlineto{\pgfqpoint{2.550693in}{1.984353in}}%
\pgfpathlineto{\pgfqpoint{2.553212in}{1.978169in}}%
\pgfpathlineto{\pgfqpoint{2.555731in}{1.989049in}}%
\pgfpathlineto{\pgfqpoint{2.558249in}{1.986410in}}%
\pgfpathlineto{\pgfqpoint{2.560768in}{2.005110in}}%
\pgfpathlineto{\pgfqpoint{2.563287in}{2.022652in}}%
\pgfpathlineto{\pgfqpoint{2.565806in}{2.003838in}}%
\pgfpathlineto{\pgfqpoint{2.568324in}{1.990681in}}%
\pgfpathlineto{\pgfqpoint{2.570843in}{1.990402in}}%
\pgfpathlineto{\pgfqpoint{2.573362in}{2.017357in}}%
\pgfpathlineto{\pgfqpoint{2.575881in}{2.009312in}}%
\pgfpathlineto{\pgfqpoint{2.578399in}{2.010179in}}%
\pgfpathlineto{\pgfqpoint{2.580918in}{2.032619in}}%
\pgfpathlineto{\pgfqpoint{2.583437in}{2.042246in}}%
\pgfpathlineto{\pgfqpoint{2.585956in}{2.036718in}}%
\pgfpathlineto{\pgfqpoint{2.588474in}{2.019829in}}%
\pgfpathlineto{\pgfqpoint{2.590993in}{2.017007in}}%
\pgfpathlineto{\pgfqpoint{2.593512in}{2.000183in}}%
\pgfpathlineto{\pgfqpoint{2.596031in}{2.039482in}}%
\pgfpathlineto{\pgfqpoint{2.598549in}{2.063286in}}%
\pgfpathlineto{\pgfqpoint{2.601068in}{2.072626in}}%
\pgfpathlineto{\pgfqpoint{2.603587in}{2.088367in}}%
\pgfpathlineto{\pgfqpoint{2.606106in}{2.081122in}}%
\pgfpathlineto{\pgfqpoint{2.608624in}{2.075357in}}%
\pgfpathlineto{\pgfqpoint{2.611143in}{2.058326in}}%
\pgfpathlineto{\pgfqpoint{2.613662in}{2.067277in}}%
\pgfpathlineto{\pgfqpoint{2.616181in}{2.112846in}}%
\pgfpathlineto{\pgfqpoint{2.618699in}{2.126624in}}%
\pgfpathlineto{\pgfqpoint{2.621218in}{2.121306in}}%
\pgfpathlineto{\pgfqpoint{2.623737in}{2.122092in}}%
\pgfpathlineto{\pgfqpoint{2.626256in}{2.138595in}}%
\pgfpathlineto{\pgfqpoint{2.628774in}{2.144623in}}%
\pgfpathlineto{\pgfqpoint{2.631293in}{2.169521in}}%
\pgfpathlineto{\pgfqpoint{2.633812in}{2.142777in}}%
\pgfpathlineto{\pgfqpoint{2.636331in}{2.131577in}}%
\pgfpathlineto{\pgfqpoint{2.638849in}{2.110935in}}%
\pgfpathlineto{\pgfqpoint{2.641368in}{2.144165in}}%
\pgfpathlineto{\pgfqpoint{2.643887in}{2.154198in}}%
\pgfpathlineto{\pgfqpoint{2.646406in}{2.144363in}}%
\pgfpathlineto{\pgfqpoint{2.648924in}{2.137507in}}%
\pgfpathlineto{\pgfqpoint{2.651443in}{2.148591in}}%
\pgfpathlineto{\pgfqpoint{2.653962in}{2.153632in}}%
\pgfpathlineto{\pgfqpoint{2.656481in}{2.155804in}}%
\pgfpathlineto{\pgfqpoint{2.658999in}{2.210447in}}%
\pgfpathlineto{\pgfqpoint{2.661518in}{2.239111in}}%
\pgfpathlineto{\pgfqpoint{2.664037in}{2.235275in}}%
\pgfpathlineto{\pgfqpoint{2.666556in}{2.232850in}}%
\pgfpathlineto{\pgfqpoint{2.669074in}{2.229508in}}%
\pgfpathlineto{\pgfqpoint{2.671593in}{2.235833in}}%
\pgfpathlineto{\pgfqpoint{2.674112in}{2.242974in}}%
\pgfpathlineto{\pgfqpoint{2.676631in}{2.263867in}}%
\pgfpathlineto{\pgfqpoint{2.679149in}{2.261522in}}%
\pgfpathlineto{\pgfqpoint{2.681668in}{2.269974in}}%
\pgfpathlineto{\pgfqpoint{2.684187in}{2.268264in}}%
\pgfpathlineto{\pgfqpoint{2.686706in}{2.255277in}}%
\pgfpathlineto{\pgfqpoint{2.689224in}{2.262634in}}%
\pgfpathlineto{\pgfqpoint{2.691743in}{2.281863in}}%
\pgfpathlineto{\pgfqpoint{2.694262in}{2.297754in}}%
\pgfpathlineto{\pgfqpoint{2.696781in}{2.308139in}}%
\pgfpathlineto{\pgfqpoint{2.699299in}{2.284761in}}%
\pgfpathlineto{\pgfqpoint{2.701818in}{2.266190in}}%
\pgfpathlineto{\pgfqpoint{2.704337in}{2.259764in}}%
\pgfpathlineto{\pgfqpoint{2.706856in}{2.273601in}}%
\pgfpathlineto{\pgfqpoint{2.709374in}{2.294529in}}%
\pgfpathlineto{\pgfqpoint{2.711893in}{2.277950in}}%
\pgfpathlineto{\pgfqpoint{2.714412in}{2.278260in}}%
\pgfpathlineto{\pgfqpoint{2.716931in}{2.280463in}}%
\pgfpathlineto{\pgfqpoint{2.719449in}{2.276701in}}%
\pgfpathlineto{\pgfqpoint{2.721968in}{2.288839in}}%
\pgfpathlineto{\pgfqpoint{2.724487in}{2.285239in}}%
\pgfpathlineto{\pgfqpoint{2.727006in}{2.268972in}}%
\pgfpathlineto{\pgfqpoint{2.729524in}{2.272061in}}%
\pgfpathlineto{\pgfqpoint{2.732043in}{2.286448in}}%
\pgfpathlineto{\pgfqpoint{2.734562in}{2.291758in}}%
\pgfpathlineto{\pgfqpoint{2.737081in}{2.269045in}}%
\pgfpathlineto{\pgfqpoint{2.739599in}{2.270779in}}%
\pgfpathlineto{\pgfqpoint{2.742118in}{2.282542in}}%
\pgfpathlineto{\pgfqpoint{2.744637in}{2.311720in}}%
\pgfpathlineto{\pgfqpoint{2.747156in}{2.317469in}}%
\pgfpathlineto{\pgfqpoint{2.749674in}{2.281739in}}%
\pgfpathlineto{\pgfqpoint{2.752193in}{2.270516in}}%
\pgfpathlineto{\pgfqpoint{2.754712in}{2.278294in}}%
\pgfpathlineto{\pgfqpoint{2.757231in}{2.289411in}}%
\pgfpathlineto{\pgfqpoint{2.759749in}{2.314411in}}%
\pgfpathlineto{\pgfqpoint{2.762268in}{2.309154in}}%
\pgfpathlineto{\pgfqpoint{2.764787in}{2.314083in}}%
\pgfpathlineto{\pgfqpoint{2.767306in}{2.316747in}}%
\pgfpathlineto{\pgfqpoint{2.769824in}{2.315963in}}%
\pgfpathlineto{\pgfqpoint{2.772343in}{2.328314in}}%
\pgfpathlineto{\pgfqpoint{2.774862in}{2.337628in}}%
\pgfpathlineto{\pgfqpoint{2.777381in}{2.315070in}}%
\pgfpathlineto{\pgfqpoint{2.779899in}{2.307958in}}%
\pgfpathlineto{\pgfqpoint{2.782418in}{2.282149in}}%
\pgfpathlineto{\pgfqpoint{2.784937in}{2.293242in}}%
\pgfpathlineto{\pgfqpoint{2.787456in}{2.276191in}}%
\pgfpathlineto{\pgfqpoint{2.789974in}{2.278489in}}%
\pgfpathlineto{\pgfqpoint{2.792493in}{2.319675in}}%
\pgfpathlineto{\pgfqpoint{2.795012in}{2.303008in}}%
\pgfpathlineto{\pgfqpoint{2.797531in}{2.325838in}}%
\pgfpathlineto{\pgfqpoint{2.800049in}{2.299626in}}%
\pgfpathlineto{\pgfqpoint{2.802568in}{2.308450in}}%
\pgfpathlineto{\pgfqpoint{2.805087in}{2.305355in}}%
\pgfpathlineto{\pgfqpoint{2.807606in}{2.299496in}}%
\pgfpathlineto{\pgfqpoint{2.810124in}{2.279414in}}%
\pgfpathlineto{\pgfqpoint{2.812643in}{2.291978in}}%
\pgfpathlineto{\pgfqpoint{2.815162in}{2.288924in}}%
\pgfpathlineto{\pgfqpoint{2.817681in}{2.285517in}}%
\pgfpathlineto{\pgfqpoint{2.820199in}{2.308263in}}%
\pgfpathlineto{\pgfqpoint{2.822718in}{2.312523in}}%
\pgfpathlineto{\pgfqpoint{2.825237in}{2.326157in}}%
\pgfpathlineto{\pgfqpoint{2.827756in}{2.334061in}}%
\pgfpathlineto{\pgfqpoint{2.830274in}{2.312186in}}%
\pgfpathlineto{\pgfqpoint{2.832793in}{2.302355in}}%
\pgfpathlineto{\pgfqpoint{2.835312in}{2.293581in}}%
\pgfpathlineto{\pgfqpoint{2.837831in}{2.293751in}}%
\pgfpathlineto{\pgfqpoint{2.840349in}{2.290310in}}%
\pgfpathlineto{\pgfqpoint{2.842868in}{2.310744in}}%
\pgfpathlineto{\pgfqpoint{2.845387in}{2.294993in}}%
\pgfpathlineto{\pgfqpoint{2.847906in}{2.307533in}}%
\pgfpathlineto{\pgfqpoint{2.850424in}{2.290283in}}%
\pgfpathlineto{\pgfqpoint{2.852943in}{2.306115in}}%
\pgfpathlineto{\pgfqpoint{2.855462in}{2.305421in}}%
\pgfpathlineto{\pgfqpoint{2.857981in}{2.288148in}}%
\pgfpathlineto{\pgfqpoint{2.860499in}{2.288513in}}%
\pgfpathlineto{\pgfqpoint{2.863018in}{2.274445in}}%
\pgfpathlineto{\pgfqpoint{2.865537in}{2.293346in}}%
\pgfpathlineto{\pgfqpoint{2.868056in}{2.335084in}}%
\pgfpathlineto{\pgfqpoint{2.870574in}{2.322048in}}%
\pgfpathlineto{\pgfqpoint{2.873093in}{2.321590in}}%
\pgfpathlineto{\pgfqpoint{2.875612in}{2.363029in}}%
\pgfpathlineto{\pgfqpoint{2.878131in}{2.377435in}}%
\pgfpathlineto{\pgfqpoint{2.880649in}{2.380050in}}%
\pgfpathlineto{\pgfqpoint{2.883168in}{2.371733in}}%
\pgfpathlineto{\pgfqpoint{2.885687in}{2.352000in}}%
\pgfpathlineto{\pgfqpoint{2.888206in}{2.367527in}}%
\pgfpathlineto{\pgfqpoint{2.890724in}{2.394753in}}%
\pgfpathlineto{\pgfqpoint{2.893243in}{2.391171in}}%
\pgfpathlineto{\pgfqpoint{2.895762in}{2.383698in}}%
\pgfpathlineto{\pgfqpoint{2.898281in}{2.372391in}}%
\pgfpathlineto{\pgfqpoint{2.900799in}{2.359248in}}%
\pgfpathlineto{\pgfqpoint{2.903318in}{2.356900in}}%
\pgfpathlineto{\pgfqpoint{2.908356in}{2.392399in}}%
\pgfpathlineto{\pgfqpoint{2.910874in}{2.367414in}}%
\pgfpathlineto{\pgfqpoint{2.913393in}{2.351853in}}%
\pgfpathlineto{\pgfqpoint{2.915912in}{2.359942in}}%
\pgfpathlineto{\pgfqpoint{2.918431in}{2.352883in}}%
\pgfpathlineto{\pgfqpoint{2.920949in}{2.361713in}}%
\pgfpathlineto{\pgfqpoint{2.923468in}{2.343977in}}%
\pgfpathlineto{\pgfqpoint{2.925987in}{2.338824in}}%
\pgfpathlineto{\pgfqpoint{2.928506in}{2.333129in}}%
\pgfpathlineto{\pgfqpoint{2.931024in}{2.309084in}}%
\pgfpathlineto{\pgfqpoint{2.933543in}{2.309938in}}%
\pgfpathlineto{\pgfqpoint{2.936062in}{2.325084in}}%
\pgfpathlineto{\pgfqpoint{2.938581in}{2.325338in}}%
\pgfpathlineto{\pgfqpoint{2.941099in}{2.321571in}}%
\pgfpathlineto{\pgfqpoint{2.943618in}{2.301875in}}%
\pgfpathlineto{\pgfqpoint{2.946137in}{2.291048in}}%
\pgfpathlineto{\pgfqpoint{2.948656in}{2.285974in}}%
\pgfpathlineto{\pgfqpoint{2.951174in}{2.310702in}}%
\pgfpathlineto{\pgfqpoint{2.953693in}{2.322185in}}%
\pgfpathlineto{\pgfqpoint{2.956212in}{2.309996in}}%
\pgfpathlineto{\pgfqpoint{2.958731in}{2.313439in}}%
\pgfpathlineto{\pgfqpoint{2.961249in}{2.322434in}}%
\pgfpathlineto{\pgfqpoint{2.963768in}{2.304707in}}%
\pgfpathlineto{\pgfqpoint{2.966287in}{2.306504in}}%
\pgfpathlineto{\pgfqpoint{2.968806in}{2.353083in}}%
\pgfpathlineto{\pgfqpoint{2.971324in}{2.371789in}}%
\pgfpathlineto{\pgfqpoint{2.973843in}{2.386942in}}%
\pgfpathlineto{\pgfqpoint{2.976362in}{2.374592in}}%
\pgfpathlineto{\pgfqpoint{2.978881in}{2.391803in}}%
\pgfpathlineto{\pgfqpoint{2.981399in}{2.401110in}}%
\pgfpathlineto{\pgfqpoint{2.983918in}{2.432913in}}%
\pgfpathlineto{\pgfqpoint{2.986437in}{2.481903in}}%
\pgfpathlineto{\pgfqpoint{2.988956in}{2.474416in}}%
\pgfpathlineto{\pgfqpoint{2.991474in}{2.478189in}}%
\pgfpathlineto{\pgfqpoint{2.993993in}{2.469609in}}%
\pgfpathlineto{\pgfqpoint{2.996512in}{2.514619in}}%
\pgfpathlineto{\pgfqpoint{2.999031in}{2.483888in}}%
\pgfpathlineto{\pgfqpoint{3.001549in}{2.517004in}}%
\pgfpathlineto{\pgfqpoint{3.004068in}{2.495852in}}%
\pgfpathlineto{\pgfqpoint{3.006587in}{2.518210in}}%
\pgfpathlineto{\pgfqpoint{3.009106in}{2.533707in}}%
\pgfpathlineto{\pgfqpoint{3.011624in}{2.543041in}}%
\pgfpathlineto{\pgfqpoint{3.014143in}{2.545661in}}%
\pgfpathlineto{\pgfqpoint{3.016662in}{2.541127in}}%
\pgfpathlineto{\pgfqpoint{3.019181in}{2.539700in}}%
\pgfpathlineto{\pgfqpoint{3.021699in}{2.538549in}}%
\pgfpathlineto{\pgfqpoint{3.024218in}{2.530218in}}%
\pgfpathlineto{\pgfqpoint{3.026737in}{2.522384in}}%
\pgfpathlineto{\pgfqpoint{3.029256in}{2.547176in}}%
\pgfpathlineto{\pgfqpoint{3.031774in}{2.538488in}}%
\pgfpathlineto{\pgfqpoint{3.034293in}{2.538133in}}%
\pgfpathlineto{\pgfqpoint{3.036812in}{2.527331in}}%
\pgfpathlineto{\pgfqpoint{3.039331in}{2.548325in}}%
\pgfpathlineto{\pgfqpoint{3.041849in}{2.535830in}}%
\pgfpathlineto{\pgfqpoint{3.044368in}{2.515426in}}%
\pgfpathlineto{\pgfqpoint{3.046887in}{2.520025in}}%
\pgfpathlineto{\pgfqpoint{3.049406in}{2.531088in}}%
\pgfpathlineto{\pgfqpoint{3.051924in}{2.544837in}}%
\pgfpathlineto{\pgfqpoint{3.054443in}{2.582025in}}%
\pgfpathlineto{\pgfqpoint{3.056962in}{2.600797in}}%
\pgfpathlineto{\pgfqpoint{3.059481in}{2.581055in}}%
\pgfpathlineto{\pgfqpoint{3.061999in}{2.583822in}}%
\pgfpathlineto{\pgfqpoint{3.064518in}{2.563484in}}%
\pgfpathlineto{\pgfqpoint{3.067037in}{2.562959in}}%
\pgfpathlineto{\pgfqpoint{3.069556in}{2.572390in}}%
\pgfpathlineto{\pgfqpoint{3.072074in}{2.556182in}}%
\pgfpathlineto{\pgfqpoint{3.074593in}{2.548308in}}%
\pgfpathlineto{\pgfqpoint{3.077112in}{2.547173in}}%
\pgfpathlineto{\pgfqpoint{3.079631in}{2.552977in}}%
\pgfpathlineto{\pgfqpoint{3.082149in}{2.541900in}}%
\pgfpathlineto{\pgfqpoint{3.084668in}{2.561177in}}%
\pgfpathlineto{\pgfqpoint{3.087187in}{2.552358in}}%
\pgfpathlineto{\pgfqpoint{3.089706in}{2.559275in}}%
\pgfpathlineto{\pgfqpoint{3.092224in}{2.581286in}}%
\pgfpathlineto{\pgfqpoint{3.094743in}{2.589233in}}%
\pgfpathlineto{\pgfqpoint{3.097262in}{2.630143in}}%
\pgfpathlineto{\pgfqpoint{3.099781in}{2.665146in}}%
\pgfpathlineto{\pgfqpoint{3.102299in}{2.672429in}}%
\pgfpathlineto{\pgfqpoint{3.104818in}{2.677267in}}%
\pgfpathlineto{\pgfqpoint{3.107337in}{2.703843in}}%
\pgfpathlineto{\pgfqpoint{3.109856in}{2.712276in}}%
\pgfpathlineto{\pgfqpoint{3.112374in}{2.716235in}}%
\pgfpathlineto{\pgfqpoint{3.114893in}{2.712548in}}%
\pgfpathlineto{\pgfqpoint{3.117412in}{2.756876in}}%
\pgfpathlineto{\pgfqpoint{3.119931in}{2.717656in}}%
\pgfpathlineto{\pgfqpoint{3.122449in}{2.692806in}}%
\pgfpathlineto{\pgfqpoint{3.124968in}{2.705767in}}%
\pgfpathlineto{\pgfqpoint{3.127487in}{2.665335in}}%
\pgfpathlineto{\pgfqpoint{3.130006in}{2.644940in}}%
\pgfpathlineto{\pgfqpoint{3.132524in}{2.623218in}}%
\pgfpathlineto{\pgfqpoint{3.135043in}{2.616378in}}%
\pgfpathlineto{\pgfqpoint{3.137562in}{2.628201in}}%
\pgfpathlineto{\pgfqpoint{3.140081in}{2.624807in}}%
\pgfpathlineto{\pgfqpoint{3.142599in}{2.669777in}}%
\pgfpathlineto{\pgfqpoint{3.145118in}{2.662431in}}%
\pgfpathlineto{\pgfqpoint{3.147637in}{2.643209in}}%
\pgfpathlineto{\pgfqpoint{3.150156in}{2.608275in}}%
\pgfpathlineto{\pgfqpoint{3.152674in}{2.611772in}}%
\pgfpathlineto{\pgfqpoint{3.155193in}{2.629136in}}%
\pgfpathlineto{\pgfqpoint{3.157712in}{2.653806in}}%
\pgfpathlineto{\pgfqpoint{3.160231in}{2.624945in}}%
\pgfpathlineto{\pgfqpoint{3.162749in}{2.614976in}}%
\pgfpathlineto{\pgfqpoint{3.165268in}{2.601481in}}%
\pgfpathlineto{\pgfqpoint{3.167787in}{2.587053in}}%
\pgfpathlineto{\pgfqpoint{3.170306in}{2.606867in}}%
\pgfpathlineto{\pgfqpoint{3.172824in}{2.596179in}}%
\pgfpathlineto{\pgfqpoint{3.175343in}{2.590330in}}%
\pgfpathlineto{\pgfqpoint{3.177862in}{2.626017in}}%
\pgfpathlineto{\pgfqpoint{3.180381in}{2.605507in}}%
\pgfpathlineto{\pgfqpoint{3.182899in}{2.582620in}}%
\pgfpathlineto{\pgfqpoint{3.185418in}{2.620070in}}%
\pgfpathlineto{\pgfqpoint{3.187937in}{2.634221in}}%
\pgfpathlineto{\pgfqpoint{3.190456in}{2.611776in}}%
\pgfpathlineto{\pgfqpoint{3.192974in}{2.557572in}}%
\pgfpathlineto{\pgfqpoint{3.195493in}{2.558450in}}%
\pgfpathlineto{\pgfqpoint{3.198012in}{2.573050in}}%
\pgfpathlineto{\pgfqpoint{3.200531in}{2.568259in}}%
\pgfpathlineto{\pgfqpoint{3.203049in}{2.567259in}}%
\pgfpathlineto{\pgfqpoint{3.205568in}{2.580399in}}%
\pgfpathlineto{\pgfqpoint{3.208087in}{2.539007in}}%
\pgfpathlineto{\pgfqpoint{3.210606in}{2.502669in}}%
\pgfpathlineto{\pgfqpoint{3.213124in}{2.490855in}}%
\pgfpathlineto{\pgfqpoint{3.215643in}{2.514258in}}%
\pgfpathlineto{\pgfqpoint{3.218162in}{2.533501in}}%
\pgfpathlineto{\pgfqpoint{3.220681in}{2.577325in}}%
\pgfpathlineto{\pgfqpoint{3.223199in}{2.587789in}}%
\pgfpathlineto{\pgfqpoint{3.225718in}{2.579210in}}%
\pgfpathlineto{\pgfqpoint{3.228237in}{2.567907in}}%
\pgfpathlineto{\pgfqpoint{3.230756in}{2.578385in}}%
\pgfpathlineto{\pgfqpoint{3.233274in}{2.603569in}}%
\pgfpathlineto{\pgfqpoint{3.235793in}{2.586617in}}%
\pgfpathlineto{\pgfqpoint{3.238312in}{2.592652in}}%
\pgfpathlineto{\pgfqpoint{3.240831in}{2.626741in}}%
\pgfpathlineto{\pgfqpoint{3.243349in}{2.659229in}}%
\pgfpathlineto{\pgfqpoint{3.245868in}{2.648180in}}%
\pgfpathlineto{\pgfqpoint{3.248387in}{2.657185in}}%
\pgfpathlineto{\pgfqpoint{3.250906in}{2.607038in}}%
\pgfpathlineto{\pgfqpoint{3.253424in}{2.616603in}}%
\pgfpathlineto{\pgfqpoint{3.255943in}{2.578571in}}%
\pgfpathlineto{\pgfqpoint{3.258462in}{2.579896in}}%
\pgfpathlineto{\pgfqpoint{3.260981in}{2.605190in}}%
\pgfpathlineto{\pgfqpoint{3.263499in}{2.612504in}}%
\pgfpathlineto{\pgfqpoint{3.266018in}{2.610017in}}%
\pgfpathlineto{\pgfqpoint{3.268537in}{2.626697in}}%
\pgfpathlineto{\pgfqpoint{3.271056in}{2.632137in}}%
\pgfpathlineto{\pgfqpoint{3.273574in}{2.614798in}}%
\pgfpathlineto{\pgfqpoint{3.276093in}{2.613524in}}%
\pgfpathlineto{\pgfqpoint{3.278612in}{2.612890in}}%
\pgfpathlineto{\pgfqpoint{3.281131in}{2.645864in}}%
\pgfpathlineto{\pgfqpoint{3.283649in}{2.652777in}}%
\pgfpathlineto{\pgfqpoint{3.286168in}{2.624347in}}%
\pgfpathlineto{\pgfqpoint{3.288687in}{2.617736in}}%
\pgfpathlineto{\pgfqpoint{3.291206in}{2.623262in}}%
\pgfpathlineto{\pgfqpoint{3.293724in}{2.642170in}}%
\pgfpathlineto{\pgfqpoint{3.296243in}{2.642478in}}%
\pgfpathlineto{\pgfqpoint{3.298762in}{2.637535in}}%
\pgfpathlineto{\pgfqpoint{3.301281in}{2.644875in}}%
\pgfpathlineto{\pgfqpoint{3.303799in}{2.637282in}}%
\pgfpathlineto{\pgfqpoint{3.306318in}{2.667401in}}%
\pgfpathlineto{\pgfqpoint{3.308837in}{2.690313in}}%
\pgfpathlineto{\pgfqpoint{3.311356in}{2.669462in}}%
\pgfpathlineto{\pgfqpoint{3.313874in}{2.703887in}}%
\pgfpathlineto{\pgfqpoint{3.316393in}{2.729516in}}%
\pgfpathlineto{\pgfqpoint{3.318912in}{2.729486in}}%
\pgfpathlineto{\pgfqpoint{3.321431in}{2.710801in}}%
\pgfpathlineto{\pgfqpoint{3.323949in}{2.694338in}}%
\pgfpathlineto{\pgfqpoint{3.326468in}{2.656390in}}%
\pgfpathlineto{\pgfqpoint{3.328987in}{2.613883in}}%
\pgfpathlineto{\pgfqpoint{3.331506in}{2.676974in}}%
\pgfpathlineto{\pgfqpoint{3.334024in}{2.682836in}}%
\pgfpathlineto{\pgfqpoint{3.336543in}{2.682801in}}%
\pgfpathlineto{\pgfqpoint{3.339062in}{2.702854in}}%
\pgfpathlineto{\pgfqpoint{3.341581in}{2.689624in}}%
\pgfpathlineto{\pgfqpoint{3.344099in}{2.700078in}}%
\pgfpathlineto{\pgfqpoint{3.346618in}{2.700756in}}%
\pgfpathlineto{\pgfqpoint{3.349137in}{2.717629in}}%
\pgfpathlineto{\pgfqpoint{3.351656in}{2.721475in}}%
\pgfpathlineto{\pgfqpoint{3.354174in}{2.695809in}}%
\pgfpathlineto{\pgfqpoint{3.356693in}{2.756754in}}%
\pgfpathlineto{\pgfqpoint{3.359212in}{2.778169in}}%
\pgfpathlineto{\pgfqpoint{3.361731in}{2.786937in}}%
\pgfpathlineto{\pgfqpoint{3.364249in}{2.770866in}}%
\pgfpathlineto{\pgfqpoint{3.366768in}{2.772038in}}%
\pgfpathlineto{\pgfqpoint{3.369287in}{2.768532in}}%
\pgfpathlineto{\pgfqpoint{3.371806in}{2.750531in}}%
\pgfpathlineto{\pgfqpoint{3.374324in}{2.743847in}}%
\pgfpathlineto{\pgfqpoint{3.376843in}{2.785206in}}%
\pgfpathlineto{\pgfqpoint{3.379362in}{2.763032in}}%
\pgfpathlineto{\pgfqpoint{3.381881in}{2.753670in}}%
\pgfpathlineto{\pgfqpoint{3.384399in}{2.782707in}}%
\pgfpathlineto{\pgfqpoint{3.386918in}{2.736891in}}%
\pgfpathlineto{\pgfqpoint{3.389437in}{2.746078in}}%
\pgfpathlineto{\pgfqpoint{3.391956in}{2.747736in}}%
\pgfpathlineto{\pgfqpoint{3.394474in}{2.749646in}}%
\pgfpathlineto{\pgfqpoint{3.396993in}{2.730674in}}%
\pgfpathlineto{\pgfqpoint{3.399512in}{2.762869in}}%
\pgfpathlineto{\pgfqpoint{3.402031in}{2.737602in}}%
\pgfpathlineto{\pgfqpoint{3.404549in}{2.741362in}}%
\pgfpathlineto{\pgfqpoint{3.407068in}{2.741887in}}%
\pgfpathlineto{\pgfqpoint{3.409587in}{2.728936in}}%
\pgfpathlineto{\pgfqpoint{3.412106in}{2.723590in}}%
\pgfpathlineto{\pgfqpoint{3.414624in}{2.716967in}}%
\pgfpathlineto{\pgfqpoint{3.417143in}{2.712844in}}%
\pgfpathlineto{\pgfqpoint{3.419662in}{2.725642in}}%
\pgfpathlineto{\pgfqpoint{3.422181in}{2.709572in}}%
\pgfpathlineto{\pgfqpoint{3.424699in}{2.724159in}}%
\pgfpathlineto{\pgfqpoint{3.427218in}{2.726024in}}%
\pgfpathlineto{\pgfqpoint{3.429737in}{2.725703in}}%
\pgfpathlineto{\pgfqpoint{3.432256in}{2.734603in}}%
\pgfpathlineto{\pgfqpoint{3.434774in}{2.745513in}}%
\pgfpathlineto{\pgfqpoint{3.437293in}{2.745915in}}%
\pgfpathlineto{\pgfqpoint{3.439812in}{2.738147in}}%
\pgfpathlineto{\pgfqpoint{3.442331in}{2.728461in}}%
\pgfpathlineto{\pgfqpoint{3.444849in}{2.720085in}}%
\pgfpathlineto{\pgfqpoint{3.447368in}{2.743638in}}%
\pgfpathlineto{\pgfqpoint{3.449887in}{2.747084in}}%
\pgfpathlineto{\pgfqpoint{3.452406in}{2.767745in}}%
\pgfpathlineto{\pgfqpoint{3.454924in}{2.759478in}}%
\pgfpathlineto{\pgfqpoint{3.457443in}{2.715793in}}%
\pgfpathlineto{\pgfqpoint{3.459962in}{2.701252in}}%
\pgfpathlineto{\pgfqpoint{3.462481in}{2.698791in}}%
\pgfpathlineto{\pgfqpoint{3.464999in}{2.718775in}}%
\pgfpathlineto{\pgfqpoint{3.467518in}{2.697816in}}%
\pgfpathlineto{\pgfqpoint{3.470037in}{2.683092in}}%
\pgfpathlineto{\pgfqpoint{3.472556in}{2.675017in}}%
\pgfpathlineto{\pgfqpoint{3.475074in}{2.681214in}}%
\pgfpathlineto{\pgfqpoint{3.477593in}{2.689705in}}%
\pgfpathlineto{\pgfqpoint{3.480112in}{2.674659in}}%
\pgfpathlineto{\pgfqpoint{3.482631in}{2.696542in}}%
\pgfpathlineto{\pgfqpoint{3.485149in}{2.674490in}}%
\pgfpathlineto{\pgfqpoint{3.487668in}{2.697194in}}%
\pgfpathlineto{\pgfqpoint{3.490187in}{2.648178in}}%
\pgfpathlineto{\pgfqpoint{3.492706in}{2.679163in}}%
\pgfpathlineto{\pgfqpoint{3.495224in}{2.683087in}}%
\pgfpathlineto{\pgfqpoint{3.497743in}{2.665734in}}%
\pgfpathlineto{\pgfqpoint{3.500262in}{2.674129in}}%
\pgfpathlineto{\pgfqpoint{3.502781in}{2.679697in}}%
\pgfpathlineto{\pgfqpoint{3.505299in}{2.670917in}}%
\pgfpathlineto{\pgfqpoint{3.507818in}{2.645032in}}%
\pgfpathlineto{\pgfqpoint{3.510337in}{2.634911in}}%
\pgfpathlineto{\pgfqpoint{3.512856in}{2.631795in}}%
\pgfpathlineto{\pgfqpoint{3.515374in}{2.670316in}}%
\pgfpathlineto{\pgfqpoint{3.517893in}{2.646452in}}%
\pgfpathlineto{\pgfqpoint{3.520412in}{2.658934in}}%
\pgfpathlineto{\pgfqpoint{3.522931in}{2.696495in}}%
\pgfpathlineto{\pgfqpoint{3.525449in}{2.694517in}}%
\pgfpathlineto{\pgfqpoint{3.527968in}{2.678724in}}%
\pgfpathlineto{\pgfqpoint{3.530487in}{2.683175in}}%
\pgfpathlineto{\pgfqpoint{3.533006in}{2.696198in}}%
\pgfpathlineto{\pgfqpoint{3.535524in}{2.691985in}}%
\pgfpathlineto{\pgfqpoint{3.538043in}{2.715690in}}%
\pgfpathlineto{\pgfqpoint{3.540562in}{2.710861in}}%
\pgfpathlineto{\pgfqpoint{3.543081in}{2.696995in}}%
\pgfpathlineto{\pgfqpoint{3.545599in}{2.723030in}}%
\pgfpathlineto{\pgfqpoint{3.548118in}{2.710728in}}%
\pgfpathlineto{\pgfqpoint{3.550637in}{2.709299in}}%
\pgfpathlineto{\pgfqpoint{3.553156in}{2.711917in}}%
\pgfpathlineto{\pgfqpoint{3.555674in}{2.730089in}}%
\pgfpathlineto{\pgfqpoint{3.558193in}{2.738219in}}%
\pgfpathlineto{\pgfqpoint{3.560712in}{2.713129in}}%
\pgfpathlineto{\pgfqpoint{3.563231in}{2.758302in}}%
\pgfpathlineto{\pgfqpoint{3.565749in}{2.753334in}}%
\pgfpathlineto{\pgfqpoint{3.568268in}{2.778396in}}%
\pgfpathlineto{\pgfqpoint{3.570787in}{2.772432in}}%
\pgfpathlineto{\pgfqpoint{3.573306in}{2.799522in}}%
\pgfpathlineto{\pgfqpoint{3.575824in}{2.786934in}}%
\pgfpathlineto{\pgfqpoint{3.578343in}{2.797229in}}%
\pgfpathlineto{\pgfqpoint{3.580862in}{2.836724in}}%
\pgfpathlineto{\pgfqpoint{3.583381in}{2.816324in}}%
\pgfpathlineto{\pgfqpoint{3.585899in}{2.855545in}}%
\pgfpathlineto{\pgfqpoint{3.588418in}{2.858001in}}%
\pgfpathlineto{\pgfqpoint{3.590937in}{2.840856in}}%
\pgfpathlineto{\pgfqpoint{3.593456in}{2.828775in}}%
\pgfpathlineto{\pgfqpoint{3.595974in}{2.878932in}}%
\pgfpathlineto{\pgfqpoint{3.601012in}{2.846917in}}%
\pgfpathlineto{\pgfqpoint{3.603531in}{2.837656in}}%
\pgfpathlineto{\pgfqpoint{3.606049in}{2.850529in}}%
\pgfpathlineto{\pgfqpoint{3.608568in}{2.888881in}}%
\pgfpathlineto{\pgfqpoint{3.611087in}{2.856169in}}%
\pgfpathlineto{\pgfqpoint{3.613606in}{2.831458in}}%
\pgfpathlineto{\pgfqpoint{3.616124in}{2.796725in}}%
\pgfpathlineto{\pgfqpoint{3.618643in}{2.779505in}}%
\pgfpathlineto{\pgfqpoint{3.621162in}{2.769613in}}%
\pgfpathlineto{\pgfqpoint{3.623681in}{2.795711in}}%
\pgfpathlineto{\pgfqpoint{3.626199in}{2.787043in}}%
\pgfpathlineto{\pgfqpoint{3.628718in}{2.802821in}}%
\pgfpathlineto{\pgfqpoint{3.631237in}{2.799445in}}%
\pgfpathlineto{\pgfqpoint{3.633756in}{2.823629in}}%
\pgfpathlineto{\pgfqpoint{3.636274in}{2.807909in}}%
\pgfpathlineto{\pgfqpoint{3.638793in}{2.839742in}}%
\pgfpathlineto{\pgfqpoint{3.641312in}{2.771534in}}%
\pgfpathlineto{\pgfqpoint{3.643831in}{2.799767in}}%
\pgfpathlineto{\pgfqpoint{3.646349in}{2.847906in}}%
\pgfpathlineto{\pgfqpoint{3.648868in}{2.838185in}}%
\pgfpathlineto{\pgfqpoint{3.651387in}{2.822675in}}%
\pgfpathlineto{\pgfqpoint{3.653906in}{2.834507in}}%
\pgfpathlineto{\pgfqpoint{3.656424in}{2.839329in}}%
\pgfpathlineto{\pgfqpoint{3.658943in}{2.861705in}}%
\pgfpathlineto{\pgfqpoint{3.661462in}{2.842981in}}%
\pgfpathlineto{\pgfqpoint{3.663981in}{2.837636in}}%
\pgfpathlineto{\pgfqpoint{3.666499in}{2.819168in}}%
\pgfpathlineto{\pgfqpoint{3.669018in}{2.824007in}}%
\pgfpathlineto{\pgfqpoint{3.671537in}{2.825011in}}%
\pgfpathlineto{\pgfqpoint{3.674056in}{2.830723in}}%
\pgfpathlineto{\pgfqpoint{3.676574in}{2.792406in}}%
\pgfpathlineto{\pgfqpoint{3.679093in}{2.796987in}}%
\pgfpathlineto{\pgfqpoint{3.681612in}{2.792865in}}%
\pgfpathlineto{\pgfqpoint{3.684131in}{2.813940in}}%
\pgfpathlineto{\pgfqpoint{3.686649in}{2.786821in}}%
\pgfpathlineto{\pgfqpoint{3.689168in}{2.788305in}}%
\pgfpathlineto{\pgfqpoint{3.691687in}{2.766406in}}%
\pgfpathlineto{\pgfqpoint{3.694206in}{2.733070in}}%
\pgfpathlineto{\pgfqpoint{3.696724in}{2.718775in}}%
\pgfpathlineto{\pgfqpoint{3.699243in}{2.746445in}}%
\pgfpathlineto{\pgfqpoint{3.701762in}{2.761452in}}%
\pgfpathlineto{\pgfqpoint{3.704281in}{2.794141in}}%
\pgfpathlineto{\pgfqpoint{3.706799in}{2.795815in}}%
\pgfpathlineto{\pgfqpoint{3.709318in}{2.812978in}}%
\pgfpathlineto{\pgfqpoint{3.711837in}{2.821359in}}%
\pgfpathlineto{\pgfqpoint{3.714356in}{2.814430in}}%
\pgfpathlineto{\pgfqpoint{3.716874in}{2.838127in}}%
\pgfpathlineto{\pgfqpoint{3.719393in}{2.860636in}}%
\pgfpathlineto{\pgfqpoint{3.721912in}{2.878648in}}%
\pgfpathlineto{\pgfqpoint{3.724431in}{2.919967in}}%
\pgfpathlineto{\pgfqpoint{3.726949in}{2.902727in}}%
\pgfpathlineto{\pgfqpoint{3.729468in}{2.894533in}}%
\pgfpathlineto{\pgfqpoint{3.731987in}{2.871548in}}%
\pgfpathlineto{\pgfqpoint{3.734506in}{2.883007in}}%
\pgfpathlineto{\pgfqpoint{3.737024in}{2.871995in}}%
\pgfpathlineto{\pgfqpoint{3.739543in}{2.880064in}}%
\pgfpathlineto{\pgfqpoint{3.742062in}{2.881608in}}%
\pgfpathlineto{\pgfqpoint{3.744581in}{2.891102in}}%
\pgfpathlineto{\pgfqpoint{3.747099in}{2.912402in}}%
\pgfpathlineto{\pgfqpoint{3.749618in}{2.891889in}}%
\pgfpathlineto{\pgfqpoint{3.752137in}{2.876918in}}%
\pgfpathlineto{\pgfqpoint{3.754656in}{2.865854in}}%
\pgfpathlineto{\pgfqpoint{3.757174in}{2.849831in}}%
\pgfpathlineto{\pgfqpoint{3.759693in}{2.847499in}}%
\pgfpathlineto{\pgfqpoint{3.762212in}{2.808192in}}%
\pgfpathlineto{\pgfqpoint{3.764731in}{2.820751in}}%
\pgfpathlineto{\pgfqpoint{3.767249in}{2.798180in}}%
\pgfpathlineto{\pgfqpoint{3.769768in}{2.801328in}}%
\pgfpathlineto{\pgfqpoint{3.772287in}{2.769910in}}%
\pgfpathlineto{\pgfqpoint{3.774806in}{2.819804in}}%
\pgfpathlineto{\pgfqpoint{3.777324in}{2.832666in}}%
\pgfpathlineto{\pgfqpoint{3.779843in}{2.837767in}}%
\pgfpathlineto{\pgfqpoint{3.782362in}{2.857268in}}%
\pgfpathlineto{\pgfqpoint{3.784881in}{2.849818in}}%
\pgfpathlineto{\pgfqpoint{3.787399in}{2.837785in}}%
\pgfpathlineto{\pgfqpoint{3.789918in}{2.855772in}}%
\pgfpathlineto{\pgfqpoint{3.792437in}{2.854625in}}%
\pgfpathlineto{\pgfqpoint{3.794956in}{2.859532in}}%
\pgfpathlineto{\pgfqpoint{3.797474in}{2.852173in}}%
\pgfpathlineto{\pgfqpoint{3.799993in}{2.837218in}}%
\pgfpathlineto{\pgfqpoint{3.802512in}{2.819590in}}%
\pgfpathlineto{\pgfqpoint{3.805031in}{2.833649in}}%
\pgfpathlineto{\pgfqpoint{3.807549in}{2.832942in}}%
\pgfpathlineto{\pgfqpoint{3.810068in}{2.823483in}}%
\pgfpathlineto{\pgfqpoint{3.812587in}{2.874710in}}%
\pgfpathlineto{\pgfqpoint{3.815106in}{2.886883in}}%
\pgfpathlineto{\pgfqpoint{3.817624in}{2.907797in}}%
\pgfpathlineto{\pgfqpoint{3.820143in}{2.921724in}}%
\pgfpathlineto{\pgfqpoint{3.822662in}{2.963978in}}%
\pgfpathlineto{\pgfqpoint{3.825181in}{2.963961in}}%
\pgfpathlineto{\pgfqpoint{3.827699in}{3.001007in}}%
\pgfpathlineto{\pgfqpoint{3.830218in}{3.011795in}}%
\pgfpathlineto{\pgfqpoint{3.832737in}{2.986172in}}%
\pgfpathlineto{\pgfqpoint{3.835256in}{3.027759in}}%
\pgfpathlineto{\pgfqpoint{3.837774in}{3.028730in}}%
\pgfpathlineto{\pgfqpoint{3.840293in}{3.016566in}}%
\pgfpathlineto{\pgfqpoint{3.842812in}{3.047003in}}%
\pgfpathlineto{\pgfqpoint{3.845331in}{3.082093in}}%
\pgfpathlineto{\pgfqpoint{3.847849in}{3.081722in}}%
\pgfpathlineto{\pgfqpoint{3.850368in}{3.059262in}}%
\pgfpathlineto{\pgfqpoint{3.852887in}{3.027140in}}%
\pgfpathlineto{\pgfqpoint{3.855406in}{3.032426in}}%
\pgfpathlineto{\pgfqpoint{3.857924in}{3.016929in}}%
\pgfpathlineto{\pgfqpoint{3.860443in}{2.978833in}}%
\pgfpathlineto{\pgfqpoint{3.862962in}{2.989361in}}%
\pgfpathlineto{\pgfqpoint{3.865481in}{2.957264in}}%
\pgfpathlineto{\pgfqpoint{3.867999in}{2.948451in}}%
\pgfpathlineto{\pgfqpoint{3.870518in}{2.920808in}}%
\pgfpathlineto{\pgfqpoint{3.873037in}{2.908262in}}%
\pgfpathlineto{\pgfqpoint{3.875556in}{2.917378in}}%
\pgfpathlineto{\pgfqpoint{3.878074in}{2.917993in}}%
\pgfpathlineto{\pgfqpoint{3.880593in}{2.938101in}}%
\pgfpathlineto{\pgfqpoint{3.883112in}{2.927367in}}%
\pgfpathlineto{\pgfqpoint{3.885631in}{2.937975in}}%
\pgfpathlineto{\pgfqpoint{3.888149in}{2.958024in}}%
\pgfpathlineto{\pgfqpoint{3.890668in}{2.928894in}}%
\pgfpathlineto{\pgfqpoint{3.893187in}{2.890363in}}%
\pgfpathlineto{\pgfqpoint{3.895706in}{2.930486in}}%
\pgfpathlineto{\pgfqpoint{3.898224in}{2.927797in}}%
\pgfpathlineto{\pgfqpoint{3.900743in}{2.941290in}}%
\pgfpathlineto{\pgfqpoint{3.903262in}{2.933410in}}%
\pgfpathlineto{\pgfqpoint{3.905781in}{2.954926in}}%
\pgfpathlineto{\pgfqpoint{3.908299in}{2.915104in}}%
\pgfpathlineto{\pgfqpoint{3.910818in}{2.886193in}}%
\pgfpathlineto{\pgfqpoint{3.913337in}{2.804252in}}%
\pgfpathlineto{\pgfqpoint{3.915856in}{2.825690in}}%
\pgfpathlineto{\pgfqpoint{3.918374in}{2.808136in}}%
\pgfpathlineto{\pgfqpoint{3.923412in}{2.729974in}}%
\pgfpathlineto{\pgfqpoint{3.925931in}{2.733139in}}%
\pgfpathlineto{\pgfqpoint{3.928449in}{2.752979in}}%
\pgfpathlineto{\pgfqpoint{3.930968in}{2.728359in}}%
\pgfpathlineto{\pgfqpoint{3.933487in}{2.709649in}}%
\pgfpathlineto{\pgfqpoint{3.936006in}{2.732390in}}%
\pgfpathlineto{\pgfqpoint{3.938524in}{2.744409in}}%
\pgfpathlineto{\pgfqpoint{3.941043in}{2.777442in}}%
\pgfpathlineto{\pgfqpoint{3.943562in}{2.820079in}}%
\pgfpathlineto{\pgfqpoint{3.946081in}{2.791690in}}%
\pgfpathlineto{\pgfqpoint{3.948599in}{2.824823in}}%
\pgfpathlineto{\pgfqpoint{3.951118in}{2.796437in}}%
\pgfpathlineto{\pgfqpoint{3.953637in}{2.806282in}}%
\pgfpathlineto{\pgfqpoint{3.956156in}{2.768347in}}%
\pgfpathlineto{\pgfqpoint{3.958674in}{2.721933in}}%
\pgfpathlineto{\pgfqpoint{3.961193in}{2.715079in}}%
\pgfpathlineto{\pgfqpoint{3.963712in}{2.693194in}}%
\pgfpathlineto{\pgfqpoint{3.966231in}{2.687109in}}%
\pgfpathlineto{\pgfqpoint{3.971268in}{2.653144in}}%
\pgfpathlineto{\pgfqpoint{3.973787in}{2.647364in}}%
\pgfpathlineto{\pgfqpoint{3.976306in}{2.632846in}}%
\pgfpathlineto{\pgfqpoint{3.978824in}{2.632739in}}%
\pgfpathlineto{\pgfqpoint{3.981343in}{2.609877in}}%
\pgfpathlineto{\pgfqpoint{3.983862in}{2.619919in}}%
\pgfpathlineto{\pgfqpoint{3.986381in}{2.595070in}}%
\pgfpathlineto{\pgfqpoint{3.988899in}{2.616546in}}%
\pgfpathlineto{\pgfqpoint{3.991418in}{2.591514in}}%
\pgfpathlineto{\pgfqpoint{3.993937in}{2.623447in}}%
\pgfpathlineto{\pgfqpoint{3.996456in}{2.626720in}}%
\pgfpathlineto{\pgfqpoint{3.998974in}{2.617336in}}%
\pgfpathlineto{\pgfqpoint{4.001493in}{2.625479in}}%
\pgfpathlineto{\pgfqpoint{4.004012in}{2.643374in}}%
\pgfpathlineto{\pgfqpoint{4.006531in}{2.651546in}}%
\pgfpathlineto{\pgfqpoint{4.009049in}{2.661779in}}%
\pgfpathlineto{\pgfqpoint{4.011568in}{2.667207in}}%
\pgfpathlineto{\pgfqpoint{4.014087in}{2.690670in}}%
\pgfpathlineto{\pgfqpoint{4.016606in}{2.658947in}}%
\pgfpathlineto{\pgfqpoint{4.019124in}{2.651111in}}%
\pgfpathlineto{\pgfqpoint{4.021643in}{2.645285in}}%
\pgfpathlineto{\pgfqpoint{4.024162in}{2.614250in}}%
\pgfpathlineto{\pgfqpoint{4.026681in}{2.645513in}}%
\pgfpathlineto{\pgfqpoint{4.029199in}{2.659423in}}%
\pgfpathlineto{\pgfqpoint{4.031718in}{2.680493in}}%
\pgfpathlineto{\pgfqpoint{4.034237in}{2.654711in}}%
\pgfpathlineto{\pgfqpoint{4.036756in}{2.665890in}}%
\pgfpathlineto{\pgfqpoint{4.039274in}{2.660213in}}%
\pgfpathlineto{\pgfqpoint{4.041793in}{2.660016in}}%
\pgfpathlineto{\pgfqpoint{4.044312in}{2.686014in}}%
\pgfpathlineto{\pgfqpoint{4.046831in}{2.681896in}}%
\pgfpathlineto{\pgfqpoint{4.049349in}{2.701479in}}%
\pgfpathlineto{\pgfqpoint{4.051868in}{2.731947in}}%
\pgfpathlineto{\pgfqpoint{4.054387in}{2.720830in}}%
\pgfpathlineto{\pgfqpoint{4.056906in}{2.712525in}}%
\pgfpathlineto{\pgfqpoint{4.059424in}{2.710681in}}%
\pgfpathlineto{\pgfqpoint{4.061943in}{2.714489in}}%
\pgfpathlineto{\pgfqpoint{4.064462in}{2.698449in}}%
\pgfpathlineto{\pgfqpoint{4.066981in}{2.692932in}}%
\pgfpathlineto{\pgfqpoint{4.069499in}{2.705348in}}%
\pgfpathlineto{\pgfqpoint{4.072018in}{2.723064in}}%
\pgfpathlineto{\pgfqpoint{4.074537in}{2.717102in}}%
\pgfpathlineto{\pgfqpoint{4.077056in}{2.715073in}}%
\pgfpathlineto{\pgfqpoint{4.079574in}{2.674709in}}%
\pgfpathlineto{\pgfqpoint{4.082093in}{2.710865in}}%
\pgfpathlineto{\pgfqpoint{4.084612in}{2.695160in}}%
\pgfpathlineto{\pgfqpoint{4.087131in}{2.668041in}}%
\pgfpathlineto{\pgfqpoint{4.089649in}{2.673174in}}%
\pgfpathlineto{\pgfqpoint{4.092168in}{2.658018in}}%
\pgfpathlineto{\pgfqpoint{4.094687in}{2.662304in}}%
\pgfpathlineto{\pgfqpoint{4.097206in}{2.639761in}}%
\pgfpathlineto{\pgfqpoint{4.099724in}{2.656729in}}%
\pgfpathlineto{\pgfqpoint{4.102243in}{2.681464in}}%
\pgfpathlineto{\pgfqpoint{4.104762in}{2.709565in}}%
\pgfpathlineto{\pgfqpoint{4.107281in}{2.692399in}}%
\pgfpathlineto{\pgfqpoint{4.109799in}{2.645043in}}%
\pgfpathlineto{\pgfqpoint{4.112318in}{2.624148in}}%
\pgfpathlineto{\pgfqpoint{4.114837in}{2.607158in}}%
\pgfpathlineto{\pgfqpoint{4.117356in}{2.611263in}}%
\pgfpathlineto{\pgfqpoint{4.119874in}{2.618794in}}%
\pgfpathlineto{\pgfqpoint{4.122393in}{2.598919in}}%
\pgfpathlineto{\pgfqpoint{4.124912in}{2.644728in}}%
\pgfpathlineto{\pgfqpoint{4.127431in}{2.614482in}}%
\pgfpathlineto{\pgfqpoint{4.129949in}{2.642662in}}%
\pgfpathlineto{\pgfqpoint{4.132468in}{2.663465in}}%
\pgfpathlineto{\pgfqpoint{4.134987in}{2.674571in}}%
\pgfpathlineto{\pgfqpoint{4.137506in}{2.662278in}}%
\pgfpathlineto{\pgfqpoint{4.140024in}{2.678821in}}%
\pgfpathlineto{\pgfqpoint{4.142543in}{2.658679in}}%
\pgfpathlineto{\pgfqpoint{4.145062in}{2.675417in}}%
\pgfpathlineto{\pgfqpoint{4.147581in}{2.658756in}}%
\pgfpathlineto{\pgfqpoint{4.150099in}{2.665871in}}%
\pgfpathlineto{\pgfqpoint{4.152618in}{2.656150in}}%
\pgfpathlineto{\pgfqpoint{4.155137in}{2.664630in}}%
\pgfpathlineto{\pgfqpoint{4.157656in}{2.636368in}}%
\pgfpathlineto{\pgfqpoint{4.160174in}{2.623512in}}%
\pgfpathlineto{\pgfqpoint{4.162693in}{2.616489in}}%
\pgfpathlineto{\pgfqpoint{4.165212in}{2.612718in}}%
\pgfpathlineto{\pgfqpoint{4.167731in}{2.567019in}}%
\pgfpathlineto{\pgfqpoint{4.170249in}{2.577338in}}%
\pgfpathlineto{\pgfqpoint{4.172768in}{2.559937in}}%
\pgfpathlineto{\pgfqpoint{4.175287in}{2.563467in}}%
\pgfpathlineto{\pgfqpoint{4.177806in}{2.585846in}}%
\pgfpathlineto{\pgfqpoint{4.180324in}{2.625040in}}%
\pgfpathlineto{\pgfqpoint{4.182843in}{2.593747in}}%
\pgfpathlineto{\pgfqpoint{4.185362in}{2.595099in}}%
\pgfpathlineto{\pgfqpoint{4.187881in}{2.596962in}}%
\pgfpathlineto{\pgfqpoint{4.190399in}{2.590408in}}%
\pgfpathlineto{\pgfqpoint{4.192918in}{2.591201in}}%
\pgfpathlineto{\pgfqpoint{4.195437in}{2.621438in}}%
\pgfpathlineto{\pgfqpoint{4.197956in}{2.640519in}}%
\pgfpathlineto{\pgfqpoint{4.200474in}{2.650726in}}%
\pgfpathlineto{\pgfqpoint{4.202993in}{2.639216in}}%
\pgfpathlineto{\pgfqpoint{4.205512in}{2.669829in}}%
\pgfpathlineto{\pgfqpoint{4.208031in}{2.673459in}}%
\pgfpathlineto{\pgfqpoint{4.210549in}{2.673258in}}%
\pgfpathlineto{\pgfqpoint{4.213068in}{2.660916in}}%
\pgfpathlineto{\pgfqpoint{4.215587in}{2.655293in}}%
\pgfpathlineto{\pgfqpoint{4.218106in}{2.661044in}}%
\pgfpathlineto{\pgfqpoint{4.220624in}{2.659759in}}%
\pgfpathlineto{\pgfqpoint{4.223143in}{2.640706in}}%
\pgfpathlineto{\pgfqpoint{4.225662in}{2.662281in}}%
\pgfpathlineto{\pgfqpoint{4.228181in}{2.704601in}}%
\pgfpathlineto{\pgfqpoint{4.230699in}{2.711636in}}%
\pgfpathlineto{\pgfqpoint{4.233218in}{2.683985in}}%
\pgfpathlineto{\pgfqpoint{4.235737in}{2.705204in}}%
\pgfpathlineto{\pgfqpoint{4.238256in}{2.696696in}}%
\pgfpathlineto{\pgfqpoint{4.240774in}{2.702954in}}%
\pgfpathlineto{\pgfqpoint{4.243293in}{2.725589in}}%
\pgfpathlineto{\pgfqpoint{4.245812in}{2.718668in}}%
\pgfpathlineto{\pgfqpoint{4.248331in}{2.746680in}}%
\pgfpathlineto{\pgfqpoint{4.250849in}{2.754055in}}%
\pgfpathlineto{\pgfqpoint{4.253368in}{2.735069in}}%
\pgfpathlineto{\pgfqpoint{4.255887in}{2.732608in}}%
\pgfpathlineto{\pgfqpoint{4.258406in}{2.745344in}}%
\pgfpathlineto{\pgfqpoint{4.260924in}{2.737026in}}%
\pgfpathlineto{\pgfqpoint{4.263443in}{2.740704in}}%
\pgfpathlineto{\pgfqpoint{4.265962in}{2.753002in}}%
\pgfpathlineto{\pgfqpoint{4.268481in}{2.745494in}}%
\pgfpathlineto{\pgfqpoint{4.270999in}{2.739220in}}%
\pgfpathlineto{\pgfqpoint{4.273518in}{2.735173in}}%
\pgfpathlineto{\pgfqpoint{4.276037in}{2.755452in}}%
\pgfpathlineto{\pgfqpoint{4.278556in}{2.753636in}}%
\pgfpathlineto{\pgfqpoint{4.281074in}{2.713638in}}%
\pgfpathlineto{\pgfqpoint{4.283593in}{2.708667in}}%
\pgfpathlineto{\pgfqpoint{4.286112in}{2.707778in}}%
\pgfpathlineto{\pgfqpoint{4.288631in}{2.684667in}}%
\pgfpathlineto{\pgfqpoint{4.291149in}{2.691313in}}%
\pgfpathlineto{\pgfqpoint{4.293668in}{2.680856in}}%
\pgfpathlineto{\pgfqpoint{4.296187in}{2.686262in}}%
\pgfpathlineto{\pgfqpoint{4.298706in}{2.690934in}}%
\pgfpathlineto{\pgfqpoint{4.301224in}{2.690996in}}%
\pgfpathlineto{\pgfqpoint{4.303743in}{2.702852in}}%
\pgfpathlineto{\pgfqpoint{4.306262in}{2.692980in}}%
\pgfpathlineto{\pgfqpoint{4.313818in}{2.631383in}}%
\pgfpathlineto{\pgfqpoint{4.316337in}{2.657574in}}%
\pgfpathlineto{\pgfqpoint{4.318856in}{2.694091in}}%
\pgfpathlineto{\pgfqpoint{4.321374in}{2.682019in}}%
\pgfpathlineto{\pgfqpoint{4.323893in}{2.661998in}}%
\pgfpathlineto{\pgfqpoint{4.326412in}{2.657577in}}%
\pgfpathlineto{\pgfqpoint{4.328931in}{2.674054in}}%
\pgfpathlineto{\pgfqpoint{4.331449in}{2.654553in}}%
\pgfpathlineto{\pgfqpoint{4.333968in}{2.641924in}}%
\pgfpathlineto{\pgfqpoint{4.336487in}{2.618006in}}%
\pgfpathlineto{\pgfqpoint{4.339006in}{2.614419in}}%
\pgfpathlineto{\pgfqpoint{4.341524in}{2.631775in}}%
\pgfpathlineto{\pgfqpoint{4.344043in}{2.669447in}}%
\pgfpathlineto{\pgfqpoint{4.346562in}{2.687103in}}%
\pgfpathlineto{\pgfqpoint{4.349081in}{2.701378in}}%
\pgfpathlineto{\pgfqpoint{4.351599in}{2.693426in}}%
\pgfpathlineto{\pgfqpoint{4.354118in}{2.674505in}}%
\pgfpathlineto{\pgfqpoint{4.356637in}{2.684771in}}%
\pgfpathlineto{\pgfqpoint{4.359156in}{2.688409in}}%
\pgfpathlineto{\pgfqpoint{4.361674in}{2.721478in}}%
\pgfpathlineto{\pgfqpoint{4.364193in}{2.691677in}}%
\pgfpathlineto{\pgfqpoint{4.366712in}{2.690176in}}%
\pgfpathlineto{\pgfqpoint{4.369231in}{2.687834in}}%
\pgfpathlineto{\pgfqpoint{4.371749in}{2.727597in}}%
\pgfpathlineto{\pgfqpoint{4.374268in}{2.702305in}}%
\pgfpathlineto{\pgfqpoint{4.376787in}{2.696718in}}%
\pgfpathlineto{\pgfqpoint{4.379306in}{2.684136in}}%
\pgfpathlineto{\pgfqpoint{4.381824in}{2.683501in}}%
\pgfpathlineto{\pgfqpoint{4.384343in}{2.677762in}}%
\pgfpathlineto{\pgfqpoint{4.386862in}{2.666821in}}%
\pgfpathlineto{\pgfqpoint{4.389381in}{2.635874in}}%
\pgfpathlineto{\pgfqpoint{4.391899in}{2.634117in}}%
\pgfpathlineto{\pgfqpoint{4.394418in}{2.644341in}}%
\pgfpathlineto{\pgfqpoint{4.396937in}{2.636960in}}%
\pgfpathlineto{\pgfqpoint{4.399456in}{2.612695in}}%
\pgfpathlineto{\pgfqpoint{4.401974in}{2.612900in}}%
\pgfpathlineto{\pgfqpoint{4.404493in}{2.593458in}}%
\pgfpathlineto{\pgfqpoint{4.407012in}{2.578290in}}%
\pgfpathlineto{\pgfqpoint{4.409531in}{2.538823in}}%
\pgfpathlineto{\pgfqpoint{4.412049in}{2.559988in}}%
\pgfpathlineto{\pgfqpoint{4.414568in}{2.566605in}}%
\pgfpathlineto{\pgfqpoint{4.417087in}{2.539159in}}%
\pgfpathlineto{\pgfqpoint{4.419606in}{2.555203in}}%
\pgfpathlineto{\pgfqpoint{4.422124in}{2.552497in}}%
\pgfpathlineto{\pgfqpoint{4.424643in}{2.572210in}}%
\pgfpathlineto{\pgfqpoint{4.427162in}{2.536388in}}%
\pgfpathlineto{\pgfqpoint{4.429681in}{2.532940in}}%
\pgfpathlineto{\pgfqpoint{4.432199in}{2.565268in}}%
\pgfpathlineto{\pgfqpoint{4.434718in}{2.562082in}}%
\pgfpathlineto{\pgfqpoint{4.437237in}{2.569960in}}%
\pgfpathlineto{\pgfqpoint{4.439756in}{2.581531in}}%
\pgfpathlineto{\pgfqpoint{4.442274in}{2.589714in}}%
\pgfpathlineto{\pgfqpoint{4.444793in}{2.580573in}}%
\pgfpathlineto{\pgfqpoint{4.447312in}{2.578176in}}%
\pgfpathlineto{\pgfqpoint{4.449831in}{2.554101in}}%
\pgfpathlineto{\pgfqpoint{4.452349in}{2.588931in}}%
\pgfpathlineto{\pgfqpoint{4.454868in}{2.584736in}}%
\pgfpathlineto{\pgfqpoint{4.457387in}{2.613726in}}%
\pgfpathlineto{\pgfqpoint{4.459906in}{2.618218in}}%
\pgfpathlineto{\pgfqpoint{4.462424in}{2.615887in}}%
\pgfpathlineto{\pgfqpoint{4.464943in}{2.653217in}}%
\pgfpathlineto{\pgfqpoint{4.467462in}{2.674144in}}%
\pgfpathlineto{\pgfqpoint{4.469981in}{2.690459in}}%
\pgfpathlineto{\pgfqpoint{4.472499in}{2.648527in}}%
\pgfpathlineto{\pgfqpoint{4.475018in}{2.684665in}}%
\pgfpathlineto{\pgfqpoint{4.477537in}{2.674709in}}%
\pgfpathlineto{\pgfqpoint{4.480056in}{2.646478in}}%
\pgfpathlineto{\pgfqpoint{4.482574in}{2.645606in}}%
\pgfpathlineto{\pgfqpoint{4.485093in}{2.635819in}}%
\pgfpathlineto{\pgfqpoint{4.487612in}{2.631694in}}%
\pgfpathlineto{\pgfqpoint{4.490131in}{2.652559in}}%
\pgfpathlineto{\pgfqpoint{4.492649in}{2.651756in}}%
\pgfpathlineto{\pgfqpoint{4.495168in}{2.668718in}}%
\pgfpathlineto{\pgfqpoint{4.497687in}{2.657345in}}%
\pgfpathlineto{\pgfqpoint{4.500206in}{2.626895in}}%
\pgfpathlineto{\pgfqpoint{4.502724in}{2.624847in}}%
\pgfpathlineto{\pgfqpoint{4.505243in}{2.628902in}}%
\pgfpathlineto{\pgfqpoint{4.507762in}{2.610933in}}%
\pgfpathlineto{\pgfqpoint{4.510281in}{2.602536in}}%
\pgfpathlineto{\pgfqpoint{4.512799in}{2.616792in}}%
\pgfpathlineto{\pgfqpoint{4.515318in}{2.641146in}}%
\pgfpathlineto{\pgfqpoint{4.517837in}{2.644120in}}%
\pgfpathlineto{\pgfqpoint{4.520356in}{2.600436in}}%
\pgfpathlineto{\pgfqpoint{4.522874in}{2.569009in}}%
\pgfpathlineto{\pgfqpoint{4.525393in}{2.557442in}}%
\pgfpathlineto{\pgfqpoint{4.527912in}{2.530808in}}%
\pgfpathlineto{\pgfqpoint{4.530431in}{2.548551in}}%
\pgfpathlineto{\pgfqpoint{4.532949in}{2.538513in}}%
\pgfpathlineto{\pgfqpoint{4.535468in}{2.506233in}}%
\pgfpathlineto{\pgfqpoint{4.537987in}{2.506333in}}%
\pgfpathlineto{\pgfqpoint{4.540506in}{2.529361in}}%
\pgfpathlineto{\pgfqpoint{4.543024in}{2.514113in}}%
\pgfpathlineto{\pgfqpoint{4.545543in}{2.482848in}}%
\pgfpathlineto{\pgfqpoint{4.548062in}{2.475961in}}%
\pgfpathlineto{\pgfqpoint{4.550581in}{2.476095in}}%
\pgfpathlineto{\pgfqpoint{4.553099in}{2.460249in}}%
\pgfpathlineto{\pgfqpoint{4.555618in}{2.467996in}}%
\pgfpathlineto{\pgfqpoint{4.558137in}{2.499503in}}%
\pgfpathlineto{\pgfqpoint{4.560656in}{2.487942in}}%
\pgfpathlineto{\pgfqpoint{4.563174in}{2.480487in}}%
\pgfpathlineto{\pgfqpoint{4.565693in}{2.499435in}}%
\pgfpathlineto{\pgfqpoint{4.568212in}{2.506479in}}%
\pgfpathlineto{\pgfqpoint{4.570731in}{2.501157in}}%
\pgfpathlineto{\pgfqpoint{4.573249in}{2.503104in}}%
\pgfpathlineto{\pgfqpoint{4.575768in}{2.515795in}}%
\pgfpathlineto{\pgfqpoint{4.578287in}{2.539047in}}%
\pgfpathlineto{\pgfqpoint{4.580806in}{2.542472in}}%
\pgfpathlineto{\pgfqpoint{4.583324in}{2.521193in}}%
\pgfpathlineto{\pgfqpoint{4.585843in}{2.526325in}}%
\pgfpathlineto{\pgfqpoint{4.588362in}{2.508611in}}%
\pgfpathlineto{\pgfqpoint{4.590881in}{2.453706in}}%
\pgfpathlineto{\pgfqpoint{4.593399in}{2.473725in}}%
\pgfpathlineto{\pgfqpoint{4.595918in}{2.467578in}}%
\pgfpathlineto{\pgfqpoint{4.598437in}{2.469961in}}%
\pgfpathlineto{\pgfqpoint{4.600956in}{2.470997in}}%
\pgfpathlineto{\pgfqpoint{4.603474in}{2.467347in}}%
\pgfpathlineto{\pgfqpoint{4.605993in}{2.487996in}}%
\pgfpathlineto{\pgfqpoint{4.608512in}{2.492198in}}%
\pgfpathlineto{\pgfqpoint{4.611031in}{2.482499in}}%
\pgfpathlineto{\pgfqpoint{4.613549in}{2.477523in}}%
\pgfpathlineto{\pgfqpoint{4.616068in}{2.476571in}}%
\pgfpathlineto{\pgfqpoint{4.618587in}{2.503415in}}%
\pgfpathlineto{\pgfqpoint{4.621106in}{2.541874in}}%
\pgfpathlineto{\pgfqpoint{4.623624in}{2.556942in}}%
\pgfpathlineto{\pgfqpoint{4.626143in}{2.594361in}}%
\pgfpathlineto{\pgfqpoint{4.628662in}{2.598356in}}%
\pgfpathlineto{\pgfqpoint{4.631181in}{2.578793in}}%
\pgfpathlineto{\pgfqpoint{4.633699in}{2.580034in}}%
\pgfpathlineto{\pgfqpoint{4.636218in}{2.577156in}}%
\pgfpathlineto{\pgfqpoint{4.638737in}{2.590624in}}%
\pgfpathlineto{\pgfqpoint{4.641256in}{2.605909in}}%
\pgfpathlineto{\pgfqpoint{4.643774in}{2.627388in}}%
\pgfpathlineto{\pgfqpoint{4.646293in}{2.618611in}}%
\pgfpathlineto{\pgfqpoint{4.648812in}{2.610430in}}%
\pgfpathlineto{\pgfqpoint{4.651331in}{2.628329in}}%
\pgfpathlineto{\pgfqpoint{4.653849in}{2.658034in}}%
\pgfpathlineto{\pgfqpoint{4.656368in}{2.637622in}}%
\pgfpathlineto{\pgfqpoint{4.658887in}{2.667476in}}%
\pgfpathlineto{\pgfqpoint{4.661406in}{2.681541in}}%
\pgfpathlineto{\pgfqpoint{4.663924in}{2.679880in}}%
\pgfpathlineto{\pgfqpoint{4.666443in}{2.687033in}}%
\pgfpathlineto{\pgfqpoint{4.668962in}{2.694623in}}%
\pgfpathlineto{\pgfqpoint{4.671481in}{2.732581in}}%
\pgfpathlineto{\pgfqpoint{4.673999in}{2.712131in}}%
\pgfpathlineto{\pgfqpoint{4.676518in}{2.710205in}}%
\pgfpathlineto{\pgfqpoint{4.679037in}{2.721817in}}%
\pgfpathlineto{\pgfqpoint{4.681556in}{2.678855in}}%
\pgfpathlineto{\pgfqpoint{4.684074in}{2.665001in}}%
\pgfpathlineto{\pgfqpoint{4.686593in}{2.684464in}}%
\pgfpathlineto{\pgfqpoint{4.689112in}{2.676440in}}%
\pgfpathlineto{\pgfqpoint{4.691631in}{2.714050in}}%
\pgfpathlineto{\pgfqpoint{4.694149in}{2.703538in}}%
\pgfpathlineto{\pgfqpoint{4.696668in}{2.643794in}}%
\pgfpathlineto{\pgfqpoint{4.699187in}{2.645688in}}%
\pgfpathlineto{\pgfqpoint{4.701706in}{2.667564in}}%
\pgfpathlineto{\pgfqpoint{4.704224in}{2.680853in}}%
\pgfpathlineto{\pgfqpoint{4.706743in}{2.703134in}}%
\pgfpathlineto{\pgfqpoint{4.709262in}{2.681599in}}%
\pgfpathlineto{\pgfqpoint{4.711781in}{2.715090in}}%
\pgfpathlineto{\pgfqpoint{4.714299in}{2.710874in}}%
\pgfpathlineto{\pgfqpoint{4.716818in}{2.695343in}}%
\pgfpathlineto{\pgfqpoint{4.719337in}{2.681539in}}%
\pgfpathlineto{\pgfqpoint{4.721856in}{2.658686in}}%
\pgfpathlineto{\pgfqpoint{4.724374in}{2.696213in}}%
\pgfpathlineto{\pgfqpoint{4.726893in}{2.694171in}}%
\pgfpathlineto{\pgfqpoint{4.729412in}{2.675176in}}%
\pgfpathlineto{\pgfqpoint{4.731931in}{2.697220in}}%
\pgfpathlineto{\pgfqpoint{4.734449in}{2.696717in}}%
\pgfpathlineto{\pgfqpoint{4.736968in}{2.692535in}}%
\pgfpathlineto{\pgfqpoint{4.739487in}{2.707713in}}%
\pgfpathlineto{\pgfqpoint{4.742006in}{2.726508in}}%
\pgfpathlineto{\pgfqpoint{4.744524in}{2.724058in}}%
\pgfpathlineto{\pgfqpoint{4.747043in}{2.733343in}}%
\pgfpathlineto{\pgfqpoint{4.749562in}{2.709479in}}%
\pgfpathlineto{\pgfqpoint{4.752081in}{2.728785in}}%
\pgfpathlineto{\pgfqpoint{4.754599in}{2.770332in}}%
\pgfpathlineto{\pgfqpoint{4.757118in}{2.756688in}}%
\pgfpathlineto{\pgfqpoint{4.759637in}{2.747122in}}%
\pgfpathlineto{\pgfqpoint{4.762156in}{2.759340in}}%
\pgfpathlineto{\pgfqpoint{4.764674in}{2.777433in}}%
\pgfpathlineto{\pgfqpoint{4.767193in}{2.741834in}}%
\pgfpathlineto{\pgfqpoint{4.769712in}{2.772007in}}%
\pgfpathlineto{\pgfqpoint{4.772231in}{2.790315in}}%
\pgfpathlineto{\pgfqpoint{4.774749in}{2.758545in}}%
\pgfpathlineto{\pgfqpoint{4.777268in}{2.761717in}}%
\pgfpathlineto{\pgfqpoint{4.779787in}{2.744571in}}%
\pgfpathlineto{\pgfqpoint{4.782306in}{2.749134in}}%
\pgfpathlineto{\pgfqpoint{4.784824in}{2.734967in}}%
\pgfpathlineto{\pgfqpoint{4.787343in}{2.742989in}}%
\pgfpathlineto{\pgfqpoint{4.789862in}{2.771609in}}%
\pgfpathlineto{\pgfqpoint{4.792381in}{2.755613in}}%
\pgfpathlineto{\pgfqpoint{4.794899in}{2.797748in}}%
\pgfpathlineto{\pgfqpoint{4.797418in}{2.805306in}}%
\pgfpathlineto{\pgfqpoint{4.799937in}{2.834243in}}%
\pgfpathlineto{\pgfqpoint{4.802456in}{2.860284in}}%
\pgfpathlineto{\pgfqpoint{4.804974in}{2.864845in}}%
\pgfpathlineto{\pgfqpoint{4.807493in}{2.849381in}}%
\pgfpathlineto{\pgfqpoint{4.810012in}{2.865460in}}%
\pgfpathlineto{\pgfqpoint{4.812531in}{2.835801in}}%
\pgfpathlineto{\pgfqpoint{4.815049in}{2.817331in}}%
\pgfpathlineto{\pgfqpoint{4.817568in}{2.812963in}}%
\pgfpathlineto{\pgfqpoint{4.820087in}{2.809970in}}%
\pgfpathlineto{\pgfqpoint{4.822606in}{2.842030in}}%
\pgfpathlineto{\pgfqpoint{4.825124in}{2.834198in}}%
\pgfpathlineto{\pgfqpoint{4.827643in}{2.791840in}}%
\pgfpathlineto{\pgfqpoint{4.830162in}{2.805183in}}%
\pgfpathlineto{\pgfqpoint{4.832681in}{2.774143in}}%
\pgfpathlineto{\pgfqpoint{4.835199in}{2.779281in}}%
\pgfpathlineto{\pgfqpoint{4.837718in}{2.793691in}}%
\pgfpathlineto{\pgfqpoint{4.840237in}{2.754654in}}%
\pgfpathlineto{\pgfqpoint{4.842756in}{2.722290in}}%
\pgfpathlineto{\pgfqpoint{4.845274in}{2.754188in}}%
\pgfpathlineto{\pgfqpoint{4.847793in}{2.761142in}}%
\pgfpathlineto{\pgfqpoint{4.850312in}{2.718727in}}%
\pgfpathlineto{\pgfqpoint{4.855349in}{2.674187in}}%
\pgfpathlineto{\pgfqpoint{4.857868in}{2.684334in}}%
\pgfpathlineto{\pgfqpoint{4.860387in}{2.718658in}}%
\pgfpathlineto{\pgfqpoint{4.862906in}{2.738913in}}%
\pgfpathlineto{\pgfqpoint{4.865424in}{2.711382in}}%
\pgfpathlineto{\pgfqpoint{4.867943in}{2.693028in}}%
\pgfpathlineto{\pgfqpoint{4.870462in}{2.678993in}}%
\pgfpathlineto{\pgfqpoint{4.872981in}{2.682995in}}%
\pgfpathlineto{\pgfqpoint{4.875499in}{2.662794in}}%
\pgfpathlineto{\pgfqpoint{4.878018in}{2.684179in}}%
\pgfpathlineto{\pgfqpoint{4.880537in}{2.651584in}}%
\pgfpathlineto{\pgfqpoint{4.883056in}{2.649558in}}%
\pgfpathlineto{\pgfqpoint{4.885574in}{2.667126in}}%
\pgfpathlineto{\pgfqpoint{4.888093in}{2.703498in}}%
\pgfpathlineto{\pgfqpoint{4.890612in}{2.726060in}}%
\pgfpathlineto{\pgfqpoint{4.893131in}{2.765506in}}%
\pgfpathlineto{\pgfqpoint{4.895649in}{2.754250in}}%
\pgfpathlineto{\pgfqpoint{4.898168in}{2.792419in}}%
\pgfpathlineto{\pgfqpoint{4.900687in}{2.787674in}}%
\pgfpathlineto{\pgfqpoint{4.903206in}{2.803718in}}%
\pgfpathlineto{\pgfqpoint{4.905724in}{2.797805in}}%
\pgfpathlineto{\pgfqpoint{4.908243in}{2.773136in}}%
\pgfpathlineto{\pgfqpoint{4.910762in}{2.768258in}}%
\pgfpathlineto{\pgfqpoint{4.913281in}{2.775916in}}%
\pgfpathlineto{\pgfqpoint{4.915799in}{2.773635in}}%
\pgfpathlineto{\pgfqpoint{4.918318in}{2.773495in}}%
\pgfpathlineto{\pgfqpoint{4.920837in}{2.745245in}}%
\pgfpathlineto{\pgfqpoint{4.923356in}{2.753068in}}%
\pgfpathlineto{\pgfqpoint{4.925874in}{2.740790in}}%
\pgfpathlineto{\pgfqpoint{4.928393in}{2.750919in}}%
\pgfpathlineto{\pgfqpoint{4.930912in}{2.703843in}}%
\pgfpathlineto{\pgfqpoint{4.933431in}{2.699982in}}%
\pgfpathlineto{\pgfqpoint{4.935949in}{2.692736in}}%
\pgfpathlineto{\pgfqpoint{4.938468in}{2.709755in}}%
\pgfpathlineto{\pgfqpoint{4.940987in}{2.706652in}}%
\pgfpathlineto{\pgfqpoint{4.943506in}{2.707365in}}%
\pgfpathlineto{\pgfqpoint{4.946024in}{2.718574in}}%
\pgfpathlineto{\pgfqpoint{4.948543in}{2.776312in}}%
\pgfpathlineto{\pgfqpoint{4.951062in}{2.781740in}}%
\pgfpathlineto{\pgfqpoint{4.953581in}{2.781328in}}%
\pgfpathlineto{\pgfqpoint{4.956099in}{2.803164in}}%
\pgfpathlineto{\pgfqpoint{4.958618in}{2.769928in}}%
\pgfpathlineto{\pgfqpoint{4.961137in}{2.753813in}}%
\pgfpathlineto{\pgfqpoint{4.963656in}{2.731807in}}%
\pgfpathlineto{\pgfqpoint{4.966174in}{2.670465in}}%
\pgfpathlineto{\pgfqpoint{4.968693in}{2.662862in}}%
\pgfpathlineto{\pgfqpoint{4.971212in}{2.690556in}}%
\pgfpathlineto{\pgfqpoint{4.973731in}{2.681515in}}%
\pgfpathlineto{\pgfqpoint{4.976249in}{2.720740in}}%
\pgfpathlineto{\pgfqpoint{4.978768in}{2.682661in}}%
\pgfpathlineto{\pgfqpoint{4.981287in}{2.691825in}}%
\pgfpathlineto{\pgfqpoint{4.983806in}{2.699633in}}%
\pgfpathlineto{\pgfqpoint{4.986324in}{2.718380in}}%
\pgfpathlineto{\pgfqpoint{4.988843in}{2.691314in}}%
\pgfpathlineto{\pgfqpoint{4.991362in}{2.660050in}}%
\pgfpathlineto{\pgfqpoint{4.993881in}{2.685304in}}%
\pgfpathlineto{\pgfqpoint{4.996399in}{2.679439in}}%
\pgfpathlineto{\pgfqpoint{4.998918in}{2.699494in}}%
\pgfpathlineto{\pgfqpoint{5.001437in}{2.686769in}}%
\pgfpathlineto{\pgfqpoint{5.003956in}{2.659809in}}%
\pgfpathlineto{\pgfqpoint{5.006474in}{2.677066in}}%
\pgfpathlineto{\pgfqpoint{5.008993in}{2.698486in}}%
\pgfpathlineto{\pgfqpoint{5.011512in}{2.709938in}}%
\pgfpathlineto{\pgfqpoint{5.014031in}{2.725762in}}%
\pgfpathlineto{\pgfqpoint{5.016549in}{2.716728in}}%
\pgfpathlineto{\pgfqpoint{5.019068in}{2.721269in}}%
\pgfpathlineto{\pgfqpoint{5.021587in}{2.730433in}}%
\pgfpathlineto{\pgfqpoint{5.024106in}{2.747582in}}%
\pgfpathlineto{\pgfqpoint{5.026624in}{2.725848in}}%
\pgfpathlineto{\pgfqpoint{5.029143in}{2.742109in}}%
\pgfpathlineto{\pgfqpoint{5.031662in}{2.736046in}}%
\pgfpathlineto{\pgfqpoint{5.034181in}{2.720625in}}%
\pgfpathlineto{\pgfqpoint{5.036699in}{2.694970in}}%
\pgfpathlineto{\pgfqpoint{5.039218in}{2.648464in}}%
\pgfpathlineto{\pgfqpoint{5.041737in}{2.643908in}}%
\pgfpathlineto{\pgfqpoint{5.044256in}{2.632020in}}%
\pgfpathlineto{\pgfqpoint{5.046774in}{2.646873in}}%
\pgfpathlineto{\pgfqpoint{5.049293in}{2.650002in}}%
\pgfpathlineto{\pgfqpoint{5.051812in}{2.676311in}}%
\pgfpathlineto{\pgfqpoint{5.054331in}{2.693504in}}%
\pgfpathlineto{\pgfqpoint{5.056849in}{2.644021in}}%
\pgfpathlineto{\pgfqpoint{5.059368in}{2.650484in}}%
\pgfpathlineto{\pgfqpoint{5.061887in}{2.657413in}}%
\pgfpathlineto{\pgfqpoint{5.064406in}{2.670668in}}%
\pgfpathlineto{\pgfqpoint{5.066924in}{2.675843in}}%
\pgfpathlineto{\pgfqpoint{5.069443in}{2.672294in}}%
\pgfpathlineto{\pgfqpoint{5.071962in}{2.696134in}}%
\pgfpathlineto{\pgfqpoint{5.074481in}{2.684899in}}%
\pgfpathlineto{\pgfqpoint{5.076999in}{2.691590in}}%
\pgfpathlineto{\pgfqpoint{5.079518in}{2.735409in}}%
\pgfpathlineto{\pgfqpoint{5.082037in}{2.760095in}}%
\pgfpathlineto{\pgfqpoint{5.084556in}{2.781358in}}%
\pgfpathlineto{\pgfqpoint{5.087074in}{2.759364in}}%
\pgfpathlineto{\pgfqpoint{5.089593in}{2.761041in}}%
\pgfpathlineto{\pgfqpoint{5.092112in}{2.773548in}}%
\pgfpathlineto{\pgfqpoint{5.094631in}{2.783150in}}%
\pgfpathlineto{\pgfqpoint{5.097149in}{2.776446in}}%
\pgfpathlineto{\pgfqpoint{5.102187in}{2.815114in}}%
\pgfpathlineto{\pgfqpoint{5.104706in}{2.765851in}}%
\pgfpathlineto{\pgfqpoint{5.107224in}{2.757354in}}%
\pgfpathlineto{\pgfqpoint{5.109743in}{2.770837in}}%
\pgfpathlineto{\pgfqpoint{5.112262in}{2.792221in}}%
\pgfpathlineto{\pgfqpoint{5.114781in}{2.803141in}}%
\pgfpathlineto{\pgfqpoint{5.117299in}{2.778716in}}%
\pgfpathlineto{\pgfqpoint{5.119818in}{2.785442in}}%
\pgfpathlineto{\pgfqpoint{5.122337in}{2.759561in}}%
\pgfpathlineto{\pgfqpoint{5.124856in}{2.756070in}}%
\pgfpathlineto{\pgfqpoint{5.127374in}{2.712282in}}%
\pgfpathlineto{\pgfqpoint{5.129893in}{2.708632in}}%
\pgfpathlineto{\pgfqpoint{5.132412in}{2.702160in}}%
\pgfpathlineto{\pgfqpoint{5.134931in}{2.686565in}}%
\pgfpathlineto{\pgfqpoint{5.137449in}{2.697528in}}%
\pgfpathlineto{\pgfqpoint{5.139968in}{2.669437in}}%
\pgfpathlineto{\pgfqpoint{5.142487in}{2.662998in}}%
\pgfpathlineto{\pgfqpoint{5.145006in}{2.664924in}}%
\pgfpathlineto{\pgfqpoint{5.147524in}{2.652288in}}%
\pgfpathlineto{\pgfqpoint{5.150043in}{2.698911in}}%
\pgfpathlineto{\pgfqpoint{5.152562in}{2.715160in}}%
\pgfpathlineto{\pgfqpoint{5.155081in}{2.726311in}}%
\pgfpathlineto{\pgfqpoint{5.157599in}{2.708851in}}%
\pgfpathlineto{\pgfqpoint{5.160118in}{2.695007in}}%
\pgfpathlineto{\pgfqpoint{5.162637in}{2.699300in}}%
\pgfpathlineto{\pgfqpoint{5.165156in}{2.663578in}}%
\pgfpathlineto{\pgfqpoint{5.167674in}{2.640466in}}%
\pgfpathlineto{\pgfqpoint{5.170193in}{2.650039in}}%
\pgfpathlineto{\pgfqpoint{5.172712in}{2.640806in}}%
\pgfpathlineto{\pgfqpoint{5.175231in}{2.600757in}}%
\pgfpathlineto{\pgfqpoint{5.177749in}{2.579955in}}%
\pgfpathlineto{\pgfqpoint{5.180268in}{2.567670in}}%
\pgfpathlineto{\pgfqpoint{5.182787in}{2.572505in}}%
\pgfpathlineto{\pgfqpoint{5.185306in}{2.546833in}}%
\pgfpathlineto{\pgfqpoint{5.187824in}{2.549920in}}%
\pgfpathlineto{\pgfqpoint{5.190343in}{2.570337in}}%
\pgfpathlineto{\pgfqpoint{5.192862in}{2.580713in}}%
\pgfpathlineto{\pgfqpoint{5.197899in}{2.652352in}}%
\pgfpathlineto{\pgfqpoint{5.200418in}{2.643638in}}%
\pgfpathlineto{\pgfqpoint{5.202937in}{2.648236in}}%
\pgfpathlineto{\pgfqpoint{5.205456in}{2.662772in}}%
\pgfpathlineto{\pgfqpoint{5.207974in}{2.649064in}}%
\pgfpathlineto{\pgfqpoint{5.210493in}{2.680575in}}%
\pgfpathlineto{\pgfqpoint{5.213012in}{2.704466in}}%
\pgfpathlineto{\pgfqpoint{5.215531in}{2.695601in}}%
\pgfpathlineto{\pgfqpoint{5.218049in}{2.712393in}}%
\pgfpathlineto{\pgfqpoint{5.220568in}{2.716060in}}%
\pgfpathlineto{\pgfqpoint{5.223087in}{2.680800in}}%
\pgfpathlineto{\pgfqpoint{5.225606in}{2.662402in}}%
\pgfpathlineto{\pgfqpoint{5.228124in}{2.651078in}}%
\pgfpathlineto{\pgfqpoint{5.230643in}{2.628333in}}%
\pgfpathlineto{\pgfqpoint{5.233162in}{2.646289in}}%
\pgfpathlineto{\pgfqpoint{5.235681in}{2.622742in}}%
\pgfpathlineto{\pgfqpoint{5.238199in}{2.595912in}}%
\pgfpathlineto{\pgfqpoint{5.240718in}{2.573737in}}%
\pgfpathlineto{\pgfqpoint{5.243237in}{2.558461in}}%
\pgfpathlineto{\pgfqpoint{5.245756in}{2.541482in}}%
\pgfpathlineto{\pgfqpoint{5.248274in}{2.579151in}}%
\pgfpathlineto{\pgfqpoint{5.250793in}{2.596533in}}%
\pgfpathlineto{\pgfqpoint{5.253312in}{2.612391in}}%
\pgfpathlineto{\pgfqpoint{5.255831in}{2.589992in}}%
\pgfpathlineto{\pgfqpoint{5.258349in}{2.619432in}}%
\pgfpathlineto{\pgfqpoint{5.260868in}{2.616843in}}%
\pgfpathlineto{\pgfqpoint{5.263387in}{2.625920in}}%
\pgfpathlineto{\pgfqpoint{5.265906in}{2.642060in}}%
\pgfpathlineto{\pgfqpoint{5.268424in}{2.659920in}}%
\pgfpathlineto{\pgfqpoint{5.270943in}{2.664272in}}%
\pgfpathlineto{\pgfqpoint{5.273462in}{2.644097in}}%
\pgfpathlineto{\pgfqpoint{5.275981in}{2.650367in}}%
\pgfpathlineto{\pgfqpoint{5.278499in}{2.639659in}}%
\pgfpathlineto{\pgfqpoint{5.281018in}{2.656529in}}%
\pgfpathlineto{\pgfqpoint{5.283537in}{2.638547in}}%
\pgfpathlineto{\pgfqpoint{5.286056in}{2.644147in}}%
\pgfpathlineto{\pgfqpoint{5.288574in}{2.650942in}}%
\pgfpathlineto{\pgfqpoint{5.291093in}{2.685568in}}%
\pgfpathlineto{\pgfqpoint{5.293612in}{2.704007in}}%
\pgfpathlineto{\pgfqpoint{5.296131in}{2.704826in}}%
\pgfpathlineto{\pgfqpoint{5.298649in}{2.728152in}}%
\pgfpathlineto{\pgfqpoint{5.301168in}{2.717274in}}%
\pgfpathlineto{\pgfqpoint{5.303687in}{2.711724in}}%
\pgfpathlineto{\pgfqpoint{5.306206in}{2.690759in}}%
\pgfpathlineto{\pgfqpoint{5.308724in}{2.696664in}}%
\pgfpathlineto{\pgfqpoint{5.311243in}{2.699242in}}%
\pgfpathlineto{\pgfqpoint{5.313762in}{2.712279in}}%
\pgfpathlineto{\pgfqpoint{5.316281in}{2.681230in}}%
\pgfpathlineto{\pgfqpoint{5.318799in}{2.711227in}}%
\pgfpathlineto{\pgfqpoint{5.321318in}{2.720699in}}%
\pgfpathlineto{\pgfqpoint{5.323837in}{2.724179in}}%
\pgfpathlineto{\pgfqpoint{5.326356in}{2.698000in}}%
\pgfpathlineto{\pgfqpoint{5.328874in}{2.668681in}}%
\pgfpathlineto{\pgfqpoint{5.331393in}{2.669864in}}%
\pgfpathlineto{\pgfqpoint{5.333912in}{2.692167in}}%
\pgfpathlineto{\pgfqpoint{5.336431in}{2.677801in}}%
\pgfpathlineto{\pgfqpoint{5.338949in}{2.669889in}}%
\pgfpathlineto{\pgfqpoint{5.341468in}{2.664456in}}%
\pgfpathlineto{\pgfqpoint{5.343987in}{2.649696in}}%
\pgfpathlineto{\pgfqpoint{5.346506in}{2.675693in}}%
\pgfpathlineto{\pgfqpoint{5.349024in}{2.716406in}}%
\pgfpathlineto{\pgfqpoint{5.351543in}{2.719635in}}%
\pgfpathlineto{\pgfqpoint{5.354062in}{2.758652in}}%
\pgfpathlineto{\pgfqpoint{5.356581in}{2.786078in}}%
\pgfpathlineto{\pgfqpoint{5.359099in}{2.785230in}}%
\pgfpathlineto{\pgfqpoint{5.361618in}{2.819694in}}%
\pgfpathlineto{\pgfqpoint{5.364137in}{2.808960in}}%
\pgfpathlineto{\pgfqpoint{5.366656in}{2.820394in}}%
\pgfpathlineto{\pgfqpoint{5.369174in}{2.781702in}}%
\pgfpathlineto{\pgfqpoint{5.371693in}{2.788098in}}%
\pgfpathlineto{\pgfqpoint{5.374212in}{2.788157in}}%
\pgfpathlineto{\pgfqpoint{5.376731in}{2.779536in}}%
\pgfpathlineto{\pgfqpoint{5.379249in}{2.768645in}}%
\pgfpathlineto{\pgfqpoint{5.381768in}{2.776729in}}%
\pgfpathlineto{\pgfqpoint{5.384287in}{2.802520in}}%
\pgfpathlineto{\pgfqpoint{5.386806in}{2.781858in}}%
\pgfpathlineto{\pgfqpoint{5.389324in}{2.796473in}}%
\pgfpathlineto{\pgfqpoint{5.391843in}{2.803497in}}%
\pgfpathlineto{\pgfqpoint{5.394362in}{2.755393in}}%
\pgfpathlineto{\pgfqpoint{5.396881in}{2.755803in}}%
\pgfpathlineto{\pgfqpoint{5.399399in}{2.734536in}}%
\pgfpathlineto{\pgfqpoint{5.401918in}{2.748981in}}%
\pgfpathlineto{\pgfqpoint{5.404437in}{2.795843in}}%
\pgfpathlineto{\pgfqpoint{5.406956in}{2.764289in}}%
\pgfpathlineto{\pgfqpoint{5.409474in}{2.753237in}}%
\pgfpathlineto{\pgfqpoint{5.411993in}{2.768769in}}%
\pgfpathlineto{\pgfqpoint{5.414512in}{2.739442in}}%
\pgfpathlineto{\pgfqpoint{5.417031in}{2.750419in}}%
\pgfpathlineto{\pgfqpoint{5.419549in}{2.704434in}}%
\pgfpathlineto{\pgfqpoint{5.422068in}{2.698386in}}%
\pgfpathlineto{\pgfqpoint{5.424587in}{2.733884in}}%
\pgfpathlineto{\pgfqpoint{5.427106in}{2.742372in}}%
\pgfpathlineto{\pgfqpoint{5.429624in}{2.743087in}}%
\pgfpathlineto{\pgfqpoint{5.432143in}{2.755182in}}%
\pgfpathlineto{\pgfqpoint{5.434662in}{2.743480in}}%
\pgfpathlineto{\pgfqpoint{5.437181in}{2.753332in}}%
\pgfpathlineto{\pgfqpoint{5.439699in}{2.748766in}}%
\pgfpathlineto{\pgfqpoint{5.442218in}{2.719251in}}%
\pgfpathlineto{\pgfqpoint{5.444737in}{2.721271in}}%
\pgfpathlineto{\pgfqpoint{5.447256in}{2.739705in}}%
\pgfpathlineto{\pgfqpoint{5.449774in}{2.700543in}}%
\pgfpathlineto{\pgfqpoint{5.452293in}{2.687361in}}%
\pgfpathlineto{\pgfqpoint{5.454812in}{2.677115in}}%
\pgfpathlineto{\pgfqpoint{5.457331in}{2.668056in}}%
\pgfpathlineto{\pgfqpoint{5.459849in}{2.671886in}}%
\pgfpathlineto{\pgfqpoint{5.462368in}{2.667248in}}%
\pgfpathlineto{\pgfqpoint{5.464887in}{2.677617in}}%
\pgfpathlineto{\pgfqpoint{5.467406in}{2.660062in}}%
\pgfpathlineto{\pgfqpoint{5.469924in}{2.670246in}}%
\pgfpathlineto{\pgfqpoint{5.472443in}{2.649421in}}%
\pgfpathlineto{\pgfqpoint{5.474962in}{2.623704in}}%
\pgfpathlineto{\pgfqpoint{5.477481in}{2.618189in}}%
\pgfpathlineto{\pgfqpoint{5.479999in}{2.618641in}}%
\pgfpathlineto{\pgfqpoint{5.482518in}{2.582760in}}%
\pgfpathlineto{\pgfqpoint{5.485037in}{2.574602in}}%
\pgfpathlineto{\pgfqpoint{5.487556in}{2.580959in}}%
\pgfpathlineto{\pgfqpoint{5.490074in}{2.579578in}}%
\pgfpathlineto{\pgfqpoint{5.492593in}{2.570102in}}%
\pgfpathlineto{\pgfqpoint{5.492593in}{2.570102in}}%
\pgfusepath{stroke}%
\end{pgfscope}%
\begin{pgfscope}%
\pgfpathrectangle{\pgfqpoint{0.455093in}{0.401080in}}{\pgfqpoint{5.037500in}{4.004000in}}%
\pgfusepath{clip}%
\pgfsetrectcap%
\pgfsetroundjoin%
\pgfsetlinewidth{0.501875pt}%
\definecolor{currentstroke}{rgb}{0.690196,0.188235,0.376471}%
\pgfsetstrokecolor{currentstroke}%
\pgfsetstrokeopacity{0.800000}%
\pgfsetdash{}{0pt}%
\pgfpathmoveto{\pgfqpoint{0.455093in}{1.021044in}}%
\pgfpathlineto{\pgfqpoint{0.457612in}{0.999728in}}%
\pgfpathlineto{\pgfqpoint{0.460131in}{1.022761in}}%
\pgfpathlineto{\pgfqpoint{0.462649in}{1.027242in}}%
\pgfpathlineto{\pgfqpoint{0.465168in}{1.030462in}}%
\pgfpathlineto{\pgfqpoint{0.467687in}{1.029030in}}%
\pgfpathlineto{\pgfqpoint{0.470206in}{1.027787in}}%
\pgfpathlineto{\pgfqpoint{0.472724in}{1.064650in}}%
\pgfpathlineto{\pgfqpoint{0.475243in}{1.056600in}}%
\pgfpathlineto{\pgfqpoint{0.477762in}{1.067002in}}%
\pgfpathlineto{\pgfqpoint{0.480281in}{1.063525in}}%
\pgfpathlineto{\pgfqpoint{0.482799in}{1.051935in}}%
\pgfpathlineto{\pgfqpoint{0.485318in}{1.059404in}}%
\pgfpathlineto{\pgfqpoint{0.487837in}{1.059710in}}%
\pgfpathlineto{\pgfqpoint{0.490356in}{1.066618in}}%
\pgfpathlineto{\pgfqpoint{0.492874in}{1.087709in}}%
\pgfpathlineto{\pgfqpoint{0.495393in}{1.090119in}}%
\pgfpathlineto{\pgfqpoint{0.497912in}{1.091895in}}%
\pgfpathlineto{\pgfqpoint{0.500431in}{1.089669in}}%
\pgfpathlineto{\pgfqpoint{0.502949in}{1.087618in}}%
\pgfpathlineto{\pgfqpoint{0.505468in}{1.053611in}}%
\pgfpathlineto{\pgfqpoint{0.510506in}{1.075831in}}%
\pgfpathlineto{\pgfqpoint{0.513024in}{1.080860in}}%
\pgfpathlineto{\pgfqpoint{0.515543in}{1.082458in}}%
\pgfpathlineto{\pgfqpoint{0.518062in}{1.106682in}}%
\pgfpathlineto{\pgfqpoint{0.520581in}{1.107396in}}%
\pgfpathlineto{\pgfqpoint{0.523099in}{1.126046in}}%
\pgfpathlineto{\pgfqpoint{0.525618in}{1.116167in}}%
\pgfpathlineto{\pgfqpoint{0.528137in}{1.113376in}}%
\pgfpathlineto{\pgfqpoint{0.530656in}{1.133116in}}%
\pgfpathlineto{\pgfqpoint{0.533174in}{1.134896in}}%
\pgfpathlineto{\pgfqpoint{0.535693in}{1.130280in}}%
\pgfpathlineto{\pgfqpoint{0.538212in}{1.141530in}}%
\pgfpathlineto{\pgfqpoint{0.540731in}{1.133041in}}%
\pgfpathlineto{\pgfqpoint{0.543249in}{1.126497in}}%
\pgfpathlineto{\pgfqpoint{0.545768in}{1.116294in}}%
\pgfpathlineto{\pgfqpoint{0.548287in}{1.112306in}}%
\pgfpathlineto{\pgfqpoint{0.550806in}{1.115375in}}%
\pgfpathlineto{\pgfqpoint{0.553324in}{1.143367in}}%
\pgfpathlineto{\pgfqpoint{0.555843in}{1.156105in}}%
\pgfpathlineto{\pgfqpoint{0.558362in}{1.166094in}}%
\pgfpathlineto{\pgfqpoint{0.560881in}{1.152436in}}%
\pgfpathlineto{\pgfqpoint{0.563399in}{1.165967in}}%
\pgfpathlineto{\pgfqpoint{0.565918in}{1.173367in}}%
\pgfpathlineto{\pgfqpoint{0.568437in}{1.168901in}}%
\pgfpathlineto{\pgfqpoint{0.570956in}{1.175681in}}%
\pgfpathlineto{\pgfqpoint{0.573474in}{1.186072in}}%
\pgfpathlineto{\pgfqpoint{0.575993in}{1.182952in}}%
\pgfpathlineto{\pgfqpoint{0.578512in}{1.173179in}}%
\pgfpathlineto{\pgfqpoint{0.581031in}{1.172958in}}%
\pgfpathlineto{\pgfqpoint{0.583549in}{1.174772in}}%
\pgfpathlineto{\pgfqpoint{0.586068in}{1.185984in}}%
\pgfpathlineto{\pgfqpoint{0.588587in}{1.198135in}}%
\pgfpathlineto{\pgfqpoint{0.591106in}{1.201252in}}%
\pgfpathlineto{\pgfqpoint{0.593624in}{1.206486in}}%
\pgfpathlineto{\pgfqpoint{0.596143in}{1.215147in}}%
\pgfpathlineto{\pgfqpoint{0.598662in}{1.209438in}}%
\pgfpathlineto{\pgfqpoint{0.601181in}{1.221980in}}%
\pgfpathlineto{\pgfqpoint{0.603699in}{1.224981in}}%
\pgfpathlineto{\pgfqpoint{0.606218in}{1.237678in}}%
\pgfpathlineto{\pgfqpoint{0.608737in}{1.251590in}}%
\pgfpathlineto{\pgfqpoint{0.611256in}{1.272125in}}%
\pgfpathlineto{\pgfqpoint{0.613774in}{1.265262in}}%
\pgfpathlineto{\pgfqpoint{0.616293in}{1.272676in}}%
\pgfpathlineto{\pgfqpoint{0.618812in}{1.264804in}}%
\pgfpathlineto{\pgfqpoint{0.621331in}{1.279988in}}%
\pgfpathlineto{\pgfqpoint{0.623849in}{1.302546in}}%
\pgfpathlineto{\pgfqpoint{0.626368in}{1.274338in}}%
\pgfpathlineto{\pgfqpoint{0.628887in}{1.276560in}}%
\pgfpathlineto{\pgfqpoint{0.631406in}{1.284820in}}%
\pgfpathlineto{\pgfqpoint{0.633924in}{1.295143in}}%
\pgfpathlineto{\pgfqpoint{0.636443in}{1.287173in}}%
\pgfpathlineto{\pgfqpoint{0.638962in}{1.282327in}}%
\pgfpathlineto{\pgfqpoint{0.641481in}{1.300550in}}%
\pgfpathlineto{\pgfqpoint{0.643999in}{1.301620in}}%
\pgfpathlineto{\pgfqpoint{0.646518in}{1.302587in}}%
\pgfpathlineto{\pgfqpoint{0.651556in}{1.250035in}}%
\pgfpathlineto{\pgfqpoint{0.654074in}{1.256283in}}%
\pgfpathlineto{\pgfqpoint{0.656593in}{1.254709in}}%
\pgfpathlineto{\pgfqpoint{0.659112in}{1.220618in}}%
\pgfpathlineto{\pgfqpoint{0.661631in}{1.228175in}}%
\pgfpathlineto{\pgfqpoint{0.664149in}{1.233243in}}%
\pgfpathlineto{\pgfqpoint{0.666668in}{1.225614in}}%
\pgfpathlineto{\pgfqpoint{0.669187in}{1.225190in}}%
\pgfpathlineto{\pgfqpoint{0.671706in}{1.262186in}}%
\pgfpathlineto{\pgfqpoint{0.674224in}{1.257991in}}%
\pgfpathlineto{\pgfqpoint{0.676743in}{1.259823in}}%
\pgfpathlineto{\pgfqpoint{0.679262in}{1.260978in}}%
\pgfpathlineto{\pgfqpoint{0.681781in}{1.264195in}}%
\pgfpathlineto{\pgfqpoint{0.684299in}{1.262610in}}%
\pgfpathlineto{\pgfqpoint{0.686818in}{1.263110in}}%
\pgfpathlineto{\pgfqpoint{0.689337in}{1.277915in}}%
\pgfpathlineto{\pgfqpoint{0.691856in}{1.259393in}}%
\pgfpathlineto{\pgfqpoint{0.694374in}{1.272200in}}%
\pgfpathlineto{\pgfqpoint{0.696893in}{1.294380in}}%
\pgfpathlineto{\pgfqpoint{0.699412in}{1.300070in}}%
\pgfpathlineto{\pgfqpoint{0.701931in}{1.302653in}}%
\pgfpathlineto{\pgfqpoint{0.704449in}{1.303039in}}%
\pgfpathlineto{\pgfqpoint{0.706968in}{1.306716in}}%
\pgfpathlineto{\pgfqpoint{0.709487in}{1.310885in}}%
\pgfpathlineto{\pgfqpoint{0.712006in}{1.304048in}}%
\pgfpathlineto{\pgfqpoint{0.714524in}{1.310863in}}%
\pgfpathlineto{\pgfqpoint{0.717043in}{1.332827in}}%
\pgfpathlineto{\pgfqpoint{0.719562in}{1.327978in}}%
\pgfpathlineto{\pgfqpoint{0.722081in}{1.322716in}}%
\pgfpathlineto{\pgfqpoint{0.724599in}{1.324917in}}%
\pgfpathlineto{\pgfqpoint{0.727118in}{1.340507in}}%
\pgfpathlineto{\pgfqpoint{0.729637in}{1.343776in}}%
\pgfpathlineto{\pgfqpoint{0.732156in}{1.364463in}}%
\pgfpathlineto{\pgfqpoint{0.734674in}{1.359950in}}%
\pgfpathlineto{\pgfqpoint{0.737193in}{1.395443in}}%
\pgfpathlineto{\pgfqpoint{0.739712in}{1.397254in}}%
\pgfpathlineto{\pgfqpoint{0.742231in}{1.408822in}}%
\pgfpathlineto{\pgfqpoint{0.744749in}{1.397280in}}%
\pgfpathlineto{\pgfqpoint{0.747268in}{1.389906in}}%
\pgfpathlineto{\pgfqpoint{0.749787in}{1.391566in}}%
\pgfpathlineto{\pgfqpoint{0.752306in}{1.404714in}}%
\pgfpathlineto{\pgfqpoint{0.754824in}{1.410736in}}%
\pgfpathlineto{\pgfqpoint{0.757343in}{1.370996in}}%
\pgfpathlineto{\pgfqpoint{0.762381in}{1.432378in}}%
\pgfpathlineto{\pgfqpoint{0.764899in}{1.423032in}}%
\pgfpathlineto{\pgfqpoint{0.767418in}{1.442214in}}%
\pgfpathlineto{\pgfqpoint{0.769937in}{1.470150in}}%
\pgfpathlineto{\pgfqpoint{0.772456in}{1.488896in}}%
\pgfpathlineto{\pgfqpoint{0.774974in}{1.469828in}}%
\pgfpathlineto{\pgfqpoint{0.777493in}{1.503281in}}%
\pgfpathlineto{\pgfqpoint{0.780012in}{1.517018in}}%
\pgfpathlineto{\pgfqpoint{0.782531in}{1.526585in}}%
\pgfpathlineto{\pgfqpoint{0.785049in}{1.525036in}}%
\pgfpathlineto{\pgfqpoint{0.787568in}{1.522024in}}%
\pgfpathlineto{\pgfqpoint{0.790087in}{1.522113in}}%
\pgfpathlineto{\pgfqpoint{0.792606in}{1.516572in}}%
\pgfpathlineto{\pgfqpoint{0.795124in}{1.537615in}}%
\pgfpathlineto{\pgfqpoint{0.797643in}{1.510931in}}%
\pgfpathlineto{\pgfqpoint{0.800162in}{1.489338in}}%
\pgfpathlineto{\pgfqpoint{0.802681in}{1.486306in}}%
\pgfpathlineto{\pgfqpoint{0.805199in}{1.483835in}}%
\pgfpathlineto{\pgfqpoint{0.807718in}{1.486606in}}%
\pgfpathlineto{\pgfqpoint{0.810237in}{1.504856in}}%
\pgfpathlineto{\pgfqpoint{0.812756in}{1.497372in}}%
\pgfpathlineto{\pgfqpoint{0.815274in}{1.499688in}}%
\pgfpathlineto{\pgfqpoint{0.817793in}{1.526021in}}%
\pgfpathlineto{\pgfqpoint{0.820312in}{1.527299in}}%
\pgfpathlineto{\pgfqpoint{0.822831in}{1.535090in}}%
\pgfpathlineto{\pgfqpoint{0.825349in}{1.548590in}}%
\pgfpathlineto{\pgfqpoint{0.827868in}{1.544645in}}%
\pgfpathlineto{\pgfqpoint{0.830387in}{1.567292in}}%
\pgfpathlineto{\pgfqpoint{0.832906in}{1.575583in}}%
\pgfpathlineto{\pgfqpoint{0.835424in}{1.548251in}}%
\pgfpathlineto{\pgfqpoint{0.837943in}{1.571265in}}%
\pgfpathlineto{\pgfqpoint{0.840462in}{1.574573in}}%
\pgfpathlineto{\pgfqpoint{0.842981in}{1.576774in}}%
\pgfpathlineto{\pgfqpoint{0.845499in}{1.579212in}}%
\pgfpathlineto{\pgfqpoint{0.848018in}{1.584158in}}%
\pgfpathlineto{\pgfqpoint{0.850537in}{1.587596in}}%
\pgfpathlineto{\pgfqpoint{0.853056in}{1.600700in}}%
\pgfpathlineto{\pgfqpoint{0.855574in}{1.610311in}}%
\pgfpathlineto{\pgfqpoint{0.858093in}{1.600588in}}%
\pgfpathlineto{\pgfqpoint{0.860612in}{1.639191in}}%
\pgfpathlineto{\pgfqpoint{0.863131in}{1.619405in}}%
\pgfpathlineto{\pgfqpoint{0.865649in}{1.615906in}}%
\pgfpathlineto{\pgfqpoint{0.868168in}{1.599654in}}%
\pgfpathlineto{\pgfqpoint{0.870687in}{1.594694in}}%
\pgfpathlineto{\pgfqpoint{0.873206in}{1.602772in}}%
\pgfpathlineto{\pgfqpoint{0.875724in}{1.579840in}}%
\pgfpathlineto{\pgfqpoint{0.878243in}{1.590282in}}%
\pgfpathlineto{\pgfqpoint{0.880762in}{1.615214in}}%
\pgfpathlineto{\pgfqpoint{0.883281in}{1.602006in}}%
\pgfpathlineto{\pgfqpoint{0.885799in}{1.600702in}}%
\pgfpathlineto{\pgfqpoint{0.888318in}{1.598415in}}%
\pgfpathlineto{\pgfqpoint{0.890837in}{1.605276in}}%
\pgfpathlineto{\pgfqpoint{0.893356in}{1.579907in}}%
\pgfpathlineto{\pgfqpoint{0.895874in}{1.596290in}}%
\pgfpathlineto{\pgfqpoint{0.898393in}{1.598778in}}%
\pgfpathlineto{\pgfqpoint{0.900912in}{1.597158in}}%
\pgfpathlineto{\pgfqpoint{0.903431in}{1.625520in}}%
\pgfpathlineto{\pgfqpoint{0.905949in}{1.625788in}}%
\pgfpathlineto{\pgfqpoint{0.908468in}{1.638063in}}%
\pgfpathlineto{\pgfqpoint{0.910987in}{1.604863in}}%
\pgfpathlineto{\pgfqpoint{0.913506in}{1.619876in}}%
\pgfpathlineto{\pgfqpoint{0.916024in}{1.603113in}}%
\pgfpathlineto{\pgfqpoint{0.918543in}{1.610702in}}%
\pgfpathlineto{\pgfqpoint{0.921062in}{1.607752in}}%
\pgfpathlineto{\pgfqpoint{0.923581in}{1.581704in}}%
\pgfpathlineto{\pgfqpoint{0.926099in}{1.587693in}}%
\pgfpathlineto{\pgfqpoint{0.928618in}{1.585186in}}%
\pgfpathlineto{\pgfqpoint{0.931137in}{1.585933in}}%
\pgfpathlineto{\pgfqpoint{0.933656in}{1.598021in}}%
\pgfpathlineto{\pgfqpoint{0.936174in}{1.588700in}}%
\pgfpathlineto{\pgfqpoint{0.938693in}{1.593570in}}%
\pgfpathlineto{\pgfqpoint{0.941212in}{1.578970in}}%
\pgfpathlineto{\pgfqpoint{0.943731in}{1.572423in}}%
\pgfpathlineto{\pgfqpoint{0.946249in}{1.559494in}}%
\pgfpathlineto{\pgfqpoint{0.948768in}{1.587131in}}%
\pgfpathlineto{\pgfqpoint{0.951287in}{1.574418in}}%
\pgfpathlineto{\pgfqpoint{0.953806in}{1.565766in}}%
\pgfpathlineto{\pgfqpoint{0.956324in}{1.550768in}}%
\pgfpathlineto{\pgfqpoint{0.958843in}{1.552909in}}%
\pgfpathlineto{\pgfqpoint{0.961362in}{1.558456in}}%
\pgfpathlineto{\pgfqpoint{0.963881in}{1.585879in}}%
\pgfpathlineto{\pgfqpoint{0.966399in}{1.584954in}}%
\pgfpathlineto{\pgfqpoint{0.968918in}{1.589423in}}%
\pgfpathlineto{\pgfqpoint{0.971437in}{1.600657in}}%
\pgfpathlineto{\pgfqpoint{0.973956in}{1.562145in}}%
\pgfpathlineto{\pgfqpoint{0.976474in}{1.546815in}}%
\pgfpathlineto{\pgfqpoint{0.978993in}{1.577511in}}%
\pgfpathlineto{\pgfqpoint{0.981512in}{1.564925in}}%
\pgfpathlineto{\pgfqpoint{0.984031in}{1.525862in}}%
\pgfpathlineto{\pgfqpoint{0.986549in}{1.542562in}}%
\pgfpathlineto{\pgfqpoint{0.989068in}{1.543675in}}%
\pgfpathlineto{\pgfqpoint{0.991587in}{1.537933in}}%
\pgfpathlineto{\pgfqpoint{0.994106in}{1.537524in}}%
\pgfpathlineto{\pgfqpoint{0.996624in}{1.541869in}}%
\pgfpathlineto{\pgfqpoint{0.999143in}{1.552117in}}%
\pgfpathlineto{\pgfqpoint{1.001662in}{1.539724in}}%
\pgfpathlineto{\pgfqpoint{1.004181in}{1.510263in}}%
\pgfpathlineto{\pgfqpoint{1.006699in}{1.524416in}}%
\pgfpathlineto{\pgfqpoint{1.009218in}{1.526256in}}%
\pgfpathlineto{\pgfqpoint{1.011737in}{1.568935in}}%
\pgfpathlineto{\pgfqpoint{1.014256in}{1.566105in}}%
\pgfpathlineto{\pgfqpoint{1.016774in}{1.569214in}}%
\pgfpathlineto{\pgfqpoint{1.019293in}{1.550942in}}%
\pgfpathlineto{\pgfqpoint{1.021812in}{1.563935in}}%
\pgfpathlineto{\pgfqpoint{1.024331in}{1.572564in}}%
\pgfpathlineto{\pgfqpoint{1.026849in}{1.556328in}}%
\pgfpathlineto{\pgfqpoint{1.029368in}{1.548765in}}%
\pgfpathlineto{\pgfqpoint{1.031887in}{1.550344in}}%
\pgfpathlineto{\pgfqpoint{1.034406in}{1.573773in}}%
\pgfpathlineto{\pgfqpoint{1.036924in}{1.594298in}}%
\pgfpathlineto{\pgfqpoint{1.039443in}{1.603610in}}%
\pgfpathlineto{\pgfqpoint{1.041962in}{1.609693in}}%
\pgfpathlineto{\pgfqpoint{1.044481in}{1.603654in}}%
\pgfpathlineto{\pgfqpoint{1.046999in}{1.606372in}}%
\pgfpathlineto{\pgfqpoint{1.049518in}{1.611780in}}%
\pgfpathlineto{\pgfqpoint{1.052037in}{1.598583in}}%
\pgfpathlineto{\pgfqpoint{1.054556in}{1.575555in}}%
\pgfpathlineto{\pgfqpoint{1.057074in}{1.575236in}}%
\pgfpathlineto{\pgfqpoint{1.059593in}{1.584687in}}%
\pgfpathlineto{\pgfqpoint{1.062112in}{1.597843in}}%
\pgfpathlineto{\pgfqpoint{1.064631in}{1.608417in}}%
\pgfpathlineto{\pgfqpoint{1.067149in}{1.604761in}}%
\pgfpathlineto{\pgfqpoint{1.069668in}{1.610656in}}%
\pgfpathlineto{\pgfqpoint{1.074706in}{1.625950in}}%
\pgfpathlineto{\pgfqpoint{1.077224in}{1.607059in}}%
\pgfpathlineto{\pgfqpoint{1.079743in}{1.609420in}}%
\pgfpathlineto{\pgfqpoint{1.082262in}{1.612898in}}%
\pgfpathlineto{\pgfqpoint{1.084781in}{1.630435in}}%
\pgfpathlineto{\pgfqpoint{1.087299in}{1.619287in}}%
\pgfpathlineto{\pgfqpoint{1.089818in}{1.613161in}}%
\pgfpathlineto{\pgfqpoint{1.092337in}{1.621469in}}%
\pgfpathlineto{\pgfqpoint{1.094856in}{1.621751in}}%
\pgfpathlineto{\pgfqpoint{1.097374in}{1.638016in}}%
\pgfpathlineto{\pgfqpoint{1.099893in}{1.640610in}}%
\pgfpathlineto{\pgfqpoint{1.102412in}{1.638065in}}%
\pgfpathlineto{\pgfqpoint{1.104931in}{1.616742in}}%
\pgfpathlineto{\pgfqpoint{1.107449in}{1.616828in}}%
\pgfpathlineto{\pgfqpoint{1.109968in}{1.624997in}}%
\pgfpathlineto{\pgfqpoint{1.112487in}{1.610978in}}%
\pgfpathlineto{\pgfqpoint{1.115006in}{1.616629in}}%
\pgfpathlineto{\pgfqpoint{1.117524in}{1.619066in}}%
\pgfpathlineto{\pgfqpoint{1.122562in}{1.595630in}}%
\pgfpathlineto{\pgfqpoint{1.125081in}{1.607152in}}%
\pgfpathlineto{\pgfqpoint{1.127599in}{1.613337in}}%
\pgfpathlineto{\pgfqpoint{1.130118in}{1.604299in}}%
\pgfpathlineto{\pgfqpoint{1.132637in}{1.599058in}}%
\pgfpathlineto{\pgfqpoint{1.135156in}{1.584425in}}%
\pgfpathlineto{\pgfqpoint{1.137674in}{1.584901in}}%
\pgfpathlineto{\pgfqpoint{1.140193in}{1.589223in}}%
\pgfpathlineto{\pgfqpoint{1.142712in}{1.587438in}}%
\pgfpathlineto{\pgfqpoint{1.145231in}{1.599938in}}%
\pgfpathlineto{\pgfqpoint{1.147749in}{1.623344in}}%
\pgfpathlineto{\pgfqpoint{1.150268in}{1.641660in}}%
\pgfpathlineto{\pgfqpoint{1.152787in}{1.629802in}}%
\pgfpathlineto{\pgfqpoint{1.155306in}{1.628104in}}%
\pgfpathlineto{\pgfqpoint{1.157824in}{1.620614in}}%
\pgfpathlineto{\pgfqpoint{1.160343in}{1.614799in}}%
\pgfpathlineto{\pgfqpoint{1.162862in}{1.609656in}}%
\pgfpathlineto{\pgfqpoint{1.165381in}{1.591626in}}%
\pgfpathlineto{\pgfqpoint{1.167899in}{1.583872in}}%
\pgfpathlineto{\pgfqpoint{1.170418in}{1.575777in}}%
\pgfpathlineto{\pgfqpoint{1.172937in}{1.568358in}}%
\pgfpathlineto{\pgfqpoint{1.175456in}{1.547231in}}%
\pgfpathlineto{\pgfqpoint{1.177974in}{1.553445in}}%
\pgfpathlineto{\pgfqpoint{1.180493in}{1.568764in}}%
\pgfpathlineto{\pgfqpoint{1.183012in}{1.597539in}}%
\pgfpathlineto{\pgfqpoint{1.185531in}{1.601284in}}%
\pgfpathlineto{\pgfqpoint{1.188049in}{1.585659in}}%
\pgfpathlineto{\pgfqpoint{1.190568in}{1.573106in}}%
\pgfpathlineto{\pgfqpoint{1.193087in}{1.549508in}}%
\pgfpathlineto{\pgfqpoint{1.195606in}{1.574893in}}%
\pgfpathlineto{\pgfqpoint{1.198124in}{1.537953in}}%
\pgfpathlineto{\pgfqpoint{1.200643in}{1.528684in}}%
\pgfpathlineto{\pgfqpoint{1.203162in}{1.541173in}}%
\pgfpathlineto{\pgfqpoint{1.205681in}{1.535536in}}%
\pgfpathlineto{\pgfqpoint{1.208199in}{1.536243in}}%
\pgfpathlineto{\pgfqpoint{1.210718in}{1.528574in}}%
\pgfpathlineto{\pgfqpoint{1.213237in}{1.517248in}}%
\pgfpathlineto{\pgfqpoint{1.215756in}{1.487939in}}%
\pgfpathlineto{\pgfqpoint{1.218274in}{1.510165in}}%
\pgfpathlineto{\pgfqpoint{1.220793in}{1.495596in}}%
\pgfpathlineto{\pgfqpoint{1.223312in}{1.486550in}}%
\pgfpathlineto{\pgfqpoint{1.225831in}{1.488759in}}%
\pgfpathlineto{\pgfqpoint{1.228349in}{1.488321in}}%
\pgfpathlineto{\pgfqpoint{1.230868in}{1.467807in}}%
\pgfpathlineto{\pgfqpoint{1.233387in}{1.490852in}}%
\pgfpathlineto{\pgfqpoint{1.235906in}{1.465893in}}%
\pgfpathlineto{\pgfqpoint{1.238424in}{1.442050in}}%
\pgfpathlineto{\pgfqpoint{1.240943in}{1.431886in}}%
\pgfpathlineto{\pgfqpoint{1.243462in}{1.465579in}}%
\pgfpathlineto{\pgfqpoint{1.245981in}{1.472562in}}%
\pgfpathlineto{\pgfqpoint{1.248499in}{1.464815in}}%
\pgfpathlineto{\pgfqpoint{1.251018in}{1.464201in}}%
\pgfpathlineto{\pgfqpoint{1.253537in}{1.462860in}}%
\pgfpathlineto{\pgfqpoint{1.256056in}{1.486280in}}%
\pgfpathlineto{\pgfqpoint{1.258574in}{1.469978in}}%
\pgfpathlineto{\pgfqpoint{1.261093in}{1.454767in}}%
\pgfpathlineto{\pgfqpoint{1.263612in}{1.461893in}}%
\pgfpathlineto{\pgfqpoint{1.266131in}{1.452028in}}%
\pgfpathlineto{\pgfqpoint{1.268649in}{1.478323in}}%
\pgfpathlineto{\pgfqpoint{1.271168in}{1.453767in}}%
\pgfpathlineto{\pgfqpoint{1.273687in}{1.436854in}}%
\pgfpathlineto{\pgfqpoint{1.276206in}{1.466776in}}%
\pgfpathlineto{\pgfqpoint{1.278724in}{1.476836in}}%
\pgfpathlineto{\pgfqpoint{1.281243in}{1.468108in}}%
\pgfpathlineto{\pgfqpoint{1.283762in}{1.498324in}}%
\pgfpathlineto{\pgfqpoint{1.286281in}{1.477090in}}%
\pgfpathlineto{\pgfqpoint{1.288799in}{1.480065in}}%
\pgfpathlineto{\pgfqpoint{1.291318in}{1.455540in}}%
\pgfpathlineto{\pgfqpoint{1.293837in}{1.459232in}}%
\pgfpathlineto{\pgfqpoint{1.296356in}{1.439167in}}%
\pgfpathlineto{\pgfqpoint{1.298874in}{1.428522in}}%
\pgfpathlineto{\pgfqpoint{1.301393in}{1.454463in}}%
\pgfpathlineto{\pgfqpoint{1.303912in}{1.464729in}}%
\pgfpathlineto{\pgfqpoint{1.306431in}{1.446820in}}%
\pgfpathlineto{\pgfqpoint{1.308949in}{1.425498in}}%
\pgfpathlineto{\pgfqpoint{1.311468in}{1.420180in}}%
\pgfpathlineto{\pgfqpoint{1.313987in}{1.427734in}}%
\pgfpathlineto{\pgfqpoint{1.316506in}{1.441926in}}%
\pgfpathlineto{\pgfqpoint{1.319024in}{1.437018in}}%
\pgfpathlineto{\pgfqpoint{1.321543in}{1.432435in}}%
\pgfpathlineto{\pgfqpoint{1.324062in}{1.437977in}}%
\pgfpathlineto{\pgfqpoint{1.326581in}{1.457035in}}%
\pgfpathlineto{\pgfqpoint{1.329099in}{1.447787in}}%
\pgfpathlineto{\pgfqpoint{1.331618in}{1.459011in}}%
\pgfpathlineto{\pgfqpoint{1.334137in}{1.428225in}}%
\pgfpathlineto{\pgfqpoint{1.336656in}{1.432626in}}%
\pgfpathlineto{\pgfqpoint{1.339174in}{1.438488in}}%
\pgfpathlineto{\pgfqpoint{1.341693in}{1.435069in}}%
\pgfpathlineto{\pgfqpoint{1.344212in}{1.459525in}}%
\pgfpathlineto{\pgfqpoint{1.346731in}{1.466306in}}%
\pgfpathlineto{\pgfqpoint{1.349249in}{1.463703in}}%
\pgfpathlineto{\pgfqpoint{1.351768in}{1.456204in}}%
\pgfpathlineto{\pgfqpoint{1.354287in}{1.447542in}}%
\pgfpathlineto{\pgfqpoint{1.356806in}{1.453399in}}%
\pgfpathlineto{\pgfqpoint{1.359324in}{1.475289in}}%
\pgfpathlineto{\pgfqpoint{1.361843in}{1.490023in}}%
\pgfpathlineto{\pgfqpoint{1.364362in}{1.506785in}}%
\pgfpathlineto{\pgfqpoint{1.366881in}{1.521741in}}%
\pgfpathlineto{\pgfqpoint{1.369399in}{1.515617in}}%
\pgfpathlineto{\pgfqpoint{1.371918in}{1.543252in}}%
\pgfpathlineto{\pgfqpoint{1.374437in}{1.561349in}}%
\pgfpathlineto{\pgfqpoint{1.376956in}{1.554938in}}%
\pgfpathlineto{\pgfqpoint{1.379474in}{1.528500in}}%
\pgfpathlineto{\pgfqpoint{1.381993in}{1.536048in}}%
\pgfpathlineto{\pgfqpoint{1.384512in}{1.517855in}}%
\pgfpathlineto{\pgfqpoint{1.387031in}{1.520328in}}%
\pgfpathlineto{\pgfqpoint{1.389549in}{1.484586in}}%
\pgfpathlineto{\pgfqpoint{1.392068in}{1.497832in}}%
\pgfpathlineto{\pgfqpoint{1.394587in}{1.486321in}}%
\pgfpathlineto{\pgfqpoint{1.397106in}{1.492689in}}%
\pgfpathlineto{\pgfqpoint{1.399624in}{1.523699in}}%
\pgfpathlineto{\pgfqpoint{1.402143in}{1.540758in}}%
\pgfpathlineto{\pgfqpoint{1.404662in}{1.577936in}}%
\pgfpathlineto{\pgfqpoint{1.407181in}{1.545178in}}%
\pgfpathlineto{\pgfqpoint{1.409699in}{1.522619in}}%
\pgfpathlineto{\pgfqpoint{1.412218in}{1.509325in}}%
\pgfpathlineto{\pgfqpoint{1.414737in}{1.499851in}}%
\pgfpathlineto{\pgfqpoint{1.417256in}{1.497718in}}%
\pgfpathlineto{\pgfqpoint{1.419774in}{1.502660in}}%
\pgfpathlineto{\pgfqpoint{1.422293in}{1.513841in}}%
\pgfpathlineto{\pgfqpoint{1.424812in}{1.524183in}}%
\pgfpathlineto{\pgfqpoint{1.427331in}{1.520400in}}%
\pgfpathlineto{\pgfqpoint{1.429849in}{1.495084in}}%
\pgfpathlineto{\pgfqpoint{1.432368in}{1.511844in}}%
\pgfpathlineto{\pgfqpoint{1.434887in}{1.498540in}}%
\pgfpathlineto{\pgfqpoint{1.437406in}{1.470006in}}%
\pgfpathlineto{\pgfqpoint{1.439924in}{1.484766in}}%
\pgfpathlineto{\pgfqpoint{1.442443in}{1.466732in}}%
\pgfpathlineto{\pgfqpoint{1.444962in}{1.500317in}}%
\pgfpathlineto{\pgfqpoint{1.447481in}{1.522837in}}%
\pgfpathlineto{\pgfqpoint{1.449999in}{1.540727in}}%
\pgfpathlineto{\pgfqpoint{1.452518in}{1.546595in}}%
\pgfpathlineto{\pgfqpoint{1.455037in}{1.526487in}}%
\pgfpathlineto{\pgfqpoint{1.457556in}{1.522286in}}%
\pgfpathlineto{\pgfqpoint{1.460074in}{1.513041in}}%
\pgfpathlineto{\pgfqpoint{1.462593in}{1.513507in}}%
\pgfpathlineto{\pgfqpoint{1.465112in}{1.498225in}}%
\pgfpathlineto{\pgfqpoint{1.467631in}{1.495092in}}%
\pgfpathlineto{\pgfqpoint{1.470149in}{1.526833in}}%
\pgfpathlineto{\pgfqpoint{1.472668in}{1.530295in}}%
\pgfpathlineto{\pgfqpoint{1.475187in}{1.523938in}}%
\pgfpathlineto{\pgfqpoint{1.477706in}{1.527994in}}%
\pgfpathlineto{\pgfqpoint{1.480224in}{1.515828in}}%
\pgfpathlineto{\pgfqpoint{1.482743in}{1.540098in}}%
\pgfpathlineto{\pgfqpoint{1.485262in}{1.560018in}}%
\pgfpathlineto{\pgfqpoint{1.487781in}{1.548616in}}%
\pgfpathlineto{\pgfqpoint{1.490299in}{1.554467in}}%
\pgfpathlineto{\pgfqpoint{1.492818in}{1.568803in}}%
\pgfpathlineto{\pgfqpoint{1.495337in}{1.577157in}}%
\pgfpathlineto{\pgfqpoint{1.500374in}{1.608960in}}%
\pgfpathlineto{\pgfqpoint{1.502893in}{1.602547in}}%
\pgfpathlineto{\pgfqpoint{1.505412in}{1.596542in}}%
\pgfpathlineto{\pgfqpoint{1.507931in}{1.610352in}}%
\pgfpathlineto{\pgfqpoint{1.510449in}{1.607337in}}%
\pgfpathlineto{\pgfqpoint{1.512968in}{1.623606in}}%
\pgfpathlineto{\pgfqpoint{1.515487in}{1.607191in}}%
\pgfpathlineto{\pgfqpoint{1.518006in}{1.587631in}}%
\pgfpathlineto{\pgfqpoint{1.520524in}{1.570145in}}%
\pgfpathlineto{\pgfqpoint{1.523043in}{1.566703in}}%
\pgfpathlineto{\pgfqpoint{1.525562in}{1.547077in}}%
\pgfpathlineto{\pgfqpoint{1.528081in}{1.557116in}}%
\pgfpathlineto{\pgfqpoint{1.530599in}{1.580390in}}%
\pgfpathlineto{\pgfqpoint{1.533118in}{1.577527in}}%
\pgfpathlineto{\pgfqpoint{1.535637in}{1.588021in}}%
\pgfpathlineto{\pgfqpoint{1.538156in}{1.605093in}}%
\pgfpathlineto{\pgfqpoint{1.540674in}{1.617261in}}%
\pgfpathlineto{\pgfqpoint{1.543193in}{1.602795in}}%
\pgfpathlineto{\pgfqpoint{1.545712in}{1.604558in}}%
\pgfpathlineto{\pgfqpoint{1.548231in}{1.638549in}}%
\pgfpathlineto{\pgfqpoint{1.550749in}{1.646469in}}%
\pgfpathlineto{\pgfqpoint{1.553268in}{1.666020in}}%
\pgfpathlineto{\pgfqpoint{1.555787in}{1.652079in}}%
\pgfpathlineto{\pgfqpoint{1.558306in}{1.653089in}}%
\pgfpathlineto{\pgfqpoint{1.560824in}{1.655127in}}%
\pgfpathlineto{\pgfqpoint{1.563343in}{1.651203in}}%
\pgfpathlineto{\pgfqpoint{1.565862in}{1.655733in}}%
\pgfpathlineto{\pgfqpoint{1.568381in}{1.673330in}}%
\pgfpathlineto{\pgfqpoint{1.570899in}{1.678563in}}%
\pgfpathlineto{\pgfqpoint{1.573418in}{1.703691in}}%
\pgfpathlineto{\pgfqpoint{1.575937in}{1.691328in}}%
\pgfpathlineto{\pgfqpoint{1.578456in}{1.718062in}}%
\pgfpathlineto{\pgfqpoint{1.580974in}{1.700065in}}%
\pgfpathlineto{\pgfqpoint{1.583493in}{1.707534in}}%
\pgfpathlineto{\pgfqpoint{1.586012in}{1.718596in}}%
\pgfpathlineto{\pgfqpoint{1.588531in}{1.747672in}}%
\pgfpathlineto{\pgfqpoint{1.591049in}{1.722017in}}%
\pgfpathlineto{\pgfqpoint{1.593568in}{1.726539in}}%
\pgfpathlineto{\pgfqpoint{1.596087in}{1.731760in}}%
\pgfpathlineto{\pgfqpoint{1.598606in}{1.744071in}}%
\pgfpathlineto{\pgfqpoint{1.601124in}{1.729333in}}%
\pgfpathlineto{\pgfqpoint{1.603643in}{1.777340in}}%
\pgfpathlineto{\pgfqpoint{1.606162in}{1.778069in}}%
\pgfpathlineto{\pgfqpoint{1.608681in}{1.754106in}}%
\pgfpathlineto{\pgfqpoint{1.611199in}{1.749271in}}%
\pgfpathlineto{\pgfqpoint{1.613718in}{1.738692in}}%
\pgfpathlineto{\pgfqpoint{1.616237in}{1.762734in}}%
\pgfpathlineto{\pgfqpoint{1.618756in}{1.767028in}}%
\pgfpathlineto{\pgfqpoint{1.621274in}{1.755759in}}%
\pgfpathlineto{\pgfqpoint{1.623793in}{1.770041in}}%
\pgfpathlineto{\pgfqpoint{1.626312in}{1.774658in}}%
\pgfpathlineto{\pgfqpoint{1.628831in}{1.765067in}}%
\pgfpathlineto{\pgfqpoint{1.631349in}{1.740994in}}%
\pgfpathlineto{\pgfqpoint{1.633868in}{1.729694in}}%
\pgfpathlineto{\pgfqpoint{1.636387in}{1.695278in}}%
\pgfpathlineto{\pgfqpoint{1.638906in}{1.695313in}}%
\pgfpathlineto{\pgfqpoint{1.641424in}{1.684404in}}%
\pgfpathlineto{\pgfqpoint{1.643943in}{1.691056in}}%
\pgfpathlineto{\pgfqpoint{1.646462in}{1.686889in}}%
\pgfpathlineto{\pgfqpoint{1.648981in}{1.665903in}}%
\pgfpathlineto{\pgfqpoint{1.651499in}{1.671080in}}%
\pgfpathlineto{\pgfqpoint{1.654018in}{1.671811in}}%
\pgfpathlineto{\pgfqpoint{1.656537in}{1.686124in}}%
\pgfpathlineto{\pgfqpoint{1.659056in}{1.711931in}}%
\pgfpathlineto{\pgfqpoint{1.661574in}{1.731975in}}%
\pgfpathlineto{\pgfqpoint{1.664093in}{1.718293in}}%
\pgfpathlineto{\pgfqpoint{1.666612in}{1.716623in}}%
\pgfpathlineto{\pgfqpoint{1.669131in}{1.709458in}}%
\pgfpathlineto{\pgfqpoint{1.671649in}{1.738936in}}%
\pgfpathlineto{\pgfqpoint{1.674168in}{1.737485in}}%
\pgfpathlineto{\pgfqpoint{1.676687in}{1.748203in}}%
\pgfpathlineto{\pgfqpoint{1.679206in}{1.757484in}}%
\pgfpathlineto{\pgfqpoint{1.681724in}{1.773561in}}%
\pgfpathlineto{\pgfqpoint{1.684243in}{1.783632in}}%
\pgfpathlineto{\pgfqpoint{1.686762in}{1.816947in}}%
\pgfpathlineto{\pgfqpoint{1.689281in}{1.819760in}}%
\pgfpathlineto{\pgfqpoint{1.691799in}{1.825424in}}%
\pgfpathlineto{\pgfqpoint{1.694318in}{1.823420in}}%
\pgfpathlineto{\pgfqpoint{1.696837in}{1.842613in}}%
\pgfpathlineto{\pgfqpoint{1.699356in}{1.866215in}}%
\pgfpathlineto{\pgfqpoint{1.701874in}{1.894491in}}%
\pgfpathlineto{\pgfqpoint{1.704393in}{1.896408in}}%
\pgfpathlineto{\pgfqpoint{1.706912in}{1.918181in}}%
\pgfpathlineto{\pgfqpoint{1.709431in}{1.923304in}}%
\pgfpathlineto{\pgfqpoint{1.711949in}{1.902916in}}%
\pgfpathlineto{\pgfqpoint{1.714468in}{1.909734in}}%
\pgfpathlineto{\pgfqpoint{1.716987in}{1.926928in}}%
\pgfpathlineto{\pgfqpoint{1.719506in}{1.947888in}}%
\pgfpathlineto{\pgfqpoint{1.722024in}{1.961522in}}%
\pgfpathlineto{\pgfqpoint{1.724543in}{1.951995in}}%
\pgfpathlineto{\pgfqpoint{1.727062in}{1.940015in}}%
\pgfpathlineto{\pgfqpoint{1.729581in}{1.962917in}}%
\pgfpathlineto{\pgfqpoint{1.732099in}{1.943687in}}%
\pgfpathlineto{\pgfqpoint{1.734618in}{1.951256in}}%
\pgfpathlineto{\pgfqpoint{1.737137in}{1.920812in}}%
\pgfpathlineto{\pgfqpoint{1.739656in}{1.951191in}}%
\pgfpathlineto{\pgfqpoint{1.742174in}{1.970073in}}%
\pgfpathlineto{\pgfqpoint{1.747212in}{1.958994in}}%
\pgfpathlineto{\pgfqpoint{1.749731in}{1.955218in}}%
\pgfpathlineto{\pgfqpoint{1.752249in}{1.937334in}}%
\pgfpathlineto{\pgfqpoint{1.754768in}{1.923127in}}%
\pgfpathlineto{\pgfqpoint{1.757287in}{1.932492in}}%
\pgfpathlineto{\pgfqpoint{1.759806in}{1.939146in}}%
\pgfpathlineto{\pgfqpoint{1.762324in}{1.949937in}}%
\pgfpathlineto{\pgfqpoint{1.764843in}{1.977192in}}%
\pgfpathlineto{\pgfqpoint{1.767362in}{1.975519in}}%
\pgfpathlineto{\pgfqpoint{1.769881in}{2.008326in}}%
\pgfpathlineto{\pgfqpoint{1.772399in}{2.000551in}}%
\pgfpathlineto{\pgfqpoint{1.774918in}{2.014731in}}%
\pgfpathlineto{\pgfqpoint{1.777437in}{2.013942in}}%
\pgfpathlineto{\pgfqpoint{1.779956in}{2.015031in}}%
\pgfpathlineto{\pgfqpoint{1.782474in}{2.031680in}}%
\pgfpathlineto{\pgfqpoint{1.784993in}{2.047561in}}%
\pgfpathlineto{\pgfqpoint{1.787512in}{2.050692in}}%
\pgfpathlineto{\pgfqpoint{1.790031in}{2.044248in}}%
\pgfpathlineto{\pgfqpoint{1.792549in}{2.056385in}}%
\pgfpathlineto{\pgfqpoint{1.795068in}{2.076696in}}%
\pgfpathlineto{\pgfqpoint{1.797587in}{2.072532in}}%
\pgfpathlineto{\pgfqpoint{1.800106in}{2.055417in}}%
\pgfpathlineto{\pgfqpoint{1.802624in}{2.061395in}}%
\pgfpathlineto{\pgfqpoint{1.805143in}{2.064782in}}%
\pgfpathlineto{\pgfqpoint{1.807662in}{2.051147in}}%
\pgfpathlineto{\pgfqpoint{1.810181in}{2.071790in}}%
\pgfpathlineto{\pgfqpoint{1.812699in}{2.087043in}}%
\pgfpathlineto{\pgfqpoint{1.815218in}{2.087073in}}%
\pgfpathlineto{\pgfqpoint{1.817737in}{2.065275in}}%
\pgfpathlineto{\pgfqpoint{1.820256in}{2.052125in}}%
\pgfpathlineto{\pgfqpoint{1.822774in}{2.054286in}}%
\pgfpathlineto{\pgfqpoint{1.825293in}{2.052218in}}%
\pgfpathlineto{\pgfqpoint{1.827812in}{2.068780in}}%
\pgfpathlineto{\pgfqpoint{1.830331in}{2.068056in}}%
\pgfpathlineto{\pgfqpoint{1.832849in}{2.046182in}}%
\pgfpathlineto{\pgfqpoint{1.835368in}{2.028689in}}%
\pgfpathlineto{\pgfqpoint{1.837887in}{2.040130in}}%
\pgfpathlineto{\pgfqpoint{1.840406in}{2.044861in}}%
\pgfpathlineto{\pgfqpoint{1.842924in}{2.063237in}}%
\pgfpathlineto{\pgfqpoint{1.845443in}{2.062848in}}%
\pgfpathlineto{\pgfqpoint{1.847962in}{2.057513in}}%
\pgfpathlineto{\pgfqpoint{1.850481in}{2.047137in}}%
\pgfpathlineto{\pgfqpoint{1.852999in}{2.025610in}}%
\pgfpathlineto{\pgfqpoint{1.855518in}{2.017148in}}%
\pgfpathlineto{\pgfqpoint{1.858037in}{2.043796in}}%
\pgfpathlineto{\pgfqpoint{1.860556in}{2.055155in}}%
\pgfpathlineto{\pgfqpoint{1.863074in}{2.045598in}}%
\pgfpathlineto{\pgfqpoint{1.865593in}{2.048500in}}%
\pgfpathlineto{\pgfqpoint{1.868112in}{2.047681in}}%
\pgfpathlineto{\pgfqpoint{1.870631in}{2.032964in}}%
\pgfpathlineto{\pgfqpoint{1.873149in}{2.052529in}}%
\pgfpathlineto{\pgfqpoint{1.875668in}{2.062721in}}%
\pgfpathlineto{\pgfqpoint{1.878187in}{2.035086in}}%
\pgfpathlineto{\pgfqpoint{1.880706in}{2.051481in}}%
\pgfpathlineto{\pgfqpoint{1.883224in}{2.053680in}}%
\pgfpathlineto{\pgfqpoint{1.885743in}{2.047197in}}%
\pgfpathlineto{\pgfqpoint{1.888262in}{2.016211in}}%
\pgfpathlineto{\pgfqpoint{1.890781in}{2.020429in}}%
\pgfpathlineto{\pgfqpoint{1.893299in}{2.023288in}}%
\pgfpathlineto{\pgfqpoint{1.895818in}{2.010362in}}%
\pgfpathlineto{\pgfqpoint{1.898337in}{1.987583in}}%
\pgfpathlineto{\pgfqpoint{1.900856in}{1.974743in}}%
\pgfpathlineto{\pgfqpoint{1.903374in}{1.968043in}}%
\pgfpathlineto{\pgfqpoint{1.905893in}{1.978459in}}%
\pgfpathlineto{\pgfqpoint{1.908412in}{1.964998in}}%
\pgfpathlineto{\pgfqpoint{1.910931in}{1.960886in}}%
\pgfpathlineto{\pgfqpoint{1.913449in}{1.970563in}}%
\pgfpathlineto{\pgfqpoint{1.915968in}{1.948037in}}%
\pgfpathlineto{\pgfqpoint{1.918487in}{1.946557in}}%
\pgfpathlineto{\pgfqpoint{1.921006in}{1.968648in}}%
\pgfpathlineto{\pgfqpoint{1.923524in}{1.989643in}}%
\pgfpathlineto{\pgfqpoint{1.926043in}{2.012247in}}%
\pgfpathlineto{\pgfqpoint{1.928562in}{2.012642in}}%
\pgfpathlineto{\pgfqpoint{1.931081in}{2.007829in}}%
\pgfpathlineto{\pgfqpoint{1.933599in}{2.010185in}}%
\pgfpathlineto{\pgfqpoint{1.938637in}{1.968988in}}%
\pgfpathlineto{\pgfqpoint{1.941156in}{2.002131in}}%
\pgfpathlineto{\pgfqpoint{1.943674in}{2.013402in}}%
\pgfpathlineto{\pgfqpoint{1.946193in}{2.000599in}}%
\pgfpathlineto{\pgfqpoint{1.948712in}{2.006750in}}%
\pgfpathlineto{\pgfqpoint{1.951231in}{1.995742in}}%
\pgfpathlineto{\pgfqpoint{1.953749in}{2.005586in}}%
\pgfpathlineto{\pgfqpoint{1.956268in}{1.992650in}}%
\pgfpathlineto{\pgfqpoint{1.958787in}{1.991046in}}%
\pgfpathlineto{\pgfqpoint{1.961306in}{1.992159in}}%
\pgfpathlineto{\pgfqpoint{1.963824in}{2.015366in}}%
\pgfpathlineto{\pgfqpoint{1.966343in}{2.029535in}}%
\pgfpathlineto{\pgfqpoint{1.968862in}{2.047923in}}%
\pgfpathlineto{\pgfqpoint{1.971381in}{2.057375in}}%
\pgfpathlineto{\pgfqpoint{1.973899in}{2.057237in}}%
\pgfpathlineto{\pgfqpoint{1.976418in}{2.075280in}}%
\pgfpathlineto{\pgfqpoint{1.978937in}{2.095355in}}%
\pgfpathlineto{\pgfqpoint{1.981456in}{2.114633in}}%
\pgfpathlineto{\pgfqpoint{1.983974in}{2.113752in}}%
\pgfpathlineto{\pgfqpoint{1.986493in}{2.124398in}}%
\pgfpathlineto{\pgfqpoint{1.989012in}{2.124535in}}%
\pgfpathlineto{\pgfqpoint{1.991531in}{2.127786in}}%
\pgfpathlineto{\pgfqpoint{1.994049in}{2.100103in}}%
\pgfpathlineto{\pgfqpoint{1.996568in}{2.094539in}}%
\pgfpathlineto{\pgfqpoint{1.999087in}{2.099446in}}%
\pgfpathlineto{\pgfqpoint{2.001606in}{2.090867in}}%
\pgfpathlineto{\pgfqpoint{2.004124in}{2.076844in}}%
\pgfpathlineto{\pgfqpoint{2.006643in}{2.097822in}}%
\pgfpathlineto{\pgfqpoint{2.009162in}{2.130942in}}%
\pgfpathlineto{\pgfqpoint{2.011681in}{2.110921in}}%
\pgfpathlineto{\pgfqpoint{2.014199in}{2.098987in}}%
\pgfpathlineto{\pgfqpoint{2.016718in}{2.085439in}}%
\pgfpathlineto{\pgfqpoint{2.019237in}{2.095891in}}%
\pgfpathlineto{\pgfqpoint{2.021756in}{2.087297in}}%
\pgfpathlineto{\pgfqpoint{2.024274in}{2.085380in}}%
\pgfpathlineto{\pgfqpoint{2.026793in}{2.074484in}}%
\pgfpathlineto{\pgfqpoint{2.029312in}{2.083189in}}%
\pgfpathlineto{\pgfqpoint{2.031831in}{2.081178in}}%
\pgfpathlineto{\pgfqpoint{2.034349in}{2.107821in}}%
\pgfpathlineto{\pgfqpoint{2.036868in}{2.101777in}}%
\pgfpathlineto{\pgfqpoint{2.039387in}{2.107685in}}%
\pgfpathlineto{\pgfqpoint{2.041906in}{2.126809in}}%
\pgfpathlineto{\pgfqpoint{2.044424in}{2.108918in}}%
\pgfpathlineto{\pgfqpoint{2.046943in}{2.107485in}}%
\pgfpathlineto{\pgfqpoint{2.049462in}{2.107128in}}%
\pgfpathlineto{\pgfqpoint{2.051981in}{2.111938in}}%
\pgfpathlineto{\pgfqpoint{2.054499in}{2.113531in}}%
\pgfpathlineto{\pgfqpoint{2.057018in}{2.133766in}}%
\pgfpathlineto{\pgfqpoint{2.059537in}{2.151265in}}%
\pgfpathlineto{\pgfqpoint{2.062056in}{2.132957in}}%
\pgfpathlineto{\pgfqpoint{2.064574in}{2.117532in}}%
\pgfpathlineto{\pgfqpoint{2.067093in}{2.154187in}}%
\pgfpathlineto{\pgfqpoint{2.069612in}{2.145806in}}%
\pgfpathlineto{\pgfqpoint{2.072131in}{2.135981in}}%
\pgfpathlineto{\pgfqpoint{2.074649in}{2.128058in}}%
\pgfpathlineto{\pgfqpoint{2.077168in}{2.088797in}}%
\pgfpathlineto{\pgfqpoint{2.079687in}{2.103251in}}%
\pgfpathlineto{\pgfqpoint{2.082206in}{2.113202in}}%
\pgfpathlineto{\pgfqpoint{2.084724in}{2.114091in}}%
\pgfpathlineto{\pgfqpoint{2.087243in}{2.102463in}}%
\pgfpathlineto{\pgfqpoint{2.089762in}{2.094354in}}%
\pgfpathlineto{\pgfqpoint{2.092281in}{2.117092in}}%
\pgfpathlineto{\pgfqpoint{2.094799in}{2.134507in}}%
\pgfpathlineto{\pgfqpoint{2.097318in}{2.140936in}}%
\pgfpathlineto{\pgfqpoint{2.099837in}{2.176115in}}%
\pgfpathlineto{\pgfqpoint{2.102356in}{2.180828in}}%
\pgfpathlineto{\pgfqpoint{2.104874in}{2.176884in}}%
\pgfpathlineto{\pgfqpoint{2.107393in}{2.156303in}}%
\pgfpathlineto{\pgfqpoint{2.109912in}{2.156311in}}%
\pgfpathlineto{\pgfqpoint{2.112431in}{2.177655in}}%
\pgfpathlineto{\pgfqpoint{2.114949in}{2.172628in}}%
\pgfpathlineto{\pgfqpoint{2.117468in}{2.150443in}}%
\pgfpathlineto{\pgfqpoint{2.119987in}{2.140006in}}%
\pgfpathlineto{\pgfqpoint{2.122506in}{2.114438in}}%
\pgfpathlineto{\pgfqpoint{2.125024in}{2.063272in}}%
\pgfpathlineto{\pgfqpoint{2.127543in}{2.076281in}}%
\pgfpathlineto{\pgfqpoint{2.130062in}{2.074806in}}%
\pgfpathlineto{\pgfqpoint{2.132581in}{2.082935in}}%
\pgfpathlineto{\pgfqpoint{2.135099in}{2.083519in}}%
\pgfpathlineto{\pgfqpoint{2.137618in}{2.093411in}}%
\pgfpathlineto{\pgfqpoint{2.140137in}{2.107627in}}%
\pgfpathlineto{\pgfqpoint{2.142656in}{2.096819in}}%
\pgfpathlineto{\pgfqpoint{2.145174in}{2.099172in}}%
\pgfpathlineto{\pgfqpoint{2.147693in}{2.081976in}}%
\pgfpathlineto{\pgfqpoint{2.150212in}{2.072696in}}%
\pgfpathlineto{\pgfqpoint{2.152731in}{2.074521in}}%
\pgfpathlineto{\pgfqpoint{2.155249in}{2.088275in}}%
\pgfpathlineto{\pgfqpoint{2.157768in}{2.124258in}}%
\pgfpathlineto{\pgfqpoint{2.160287in}{2.126442in}}%
\pgfpathlineto{\pgfqpoint{2.162806in}{2.136781in}}%
\pgfpathlineto{\pgfqpoint{2.165324in}{2.162690in}}%
\pgfpathlineto{\pgfqpoint{2.167843in}{2.161534in}}%
\pgfpathlineto{\pgfqpoint{2.170362in}{2.142473in}}%
\pgfpathlineto{\pgfqpoint{2.172881in}{2.145112in}}%
\pgfpathlineto{\pgfqpoint{2.175399in}{2.160173in}}%
\pgfpathlineto{\pgfqpoint{2.177918in}{2.152708in}}%
\pgfpathlineto{\pgfqpoint{2.180437in}{2.150426in}}%
\pgfpathlineto{\pgfqpoint{2.182956in}{2.175940in}}%
\pgfpathlineto{\pgfqpoint{2.185474in}{2.170758in}}%
\pgfpathlineto{\pgfqpoint{2.187993in}{2.171335in}}%
\pgfpathlineto{\pgfqpoint{2.190512in}{2.148577in}}%
\pgfpathlineto{\pgfqpoint{2.193031in}{2.147142in}}%
\pgfpathlineto{\pgfqpoint{2.195549in}{2.142296in}}%
\pgfpathlineto{\pgfqpoint{2.198068in}{2.134938in}}%
\pgfpathlineto{\pgfqpoint{2.200587in}{2.138631in}}%
\pgfpathlineto{\pgfqpoint{2.203106in}{2.147467in}}%
\pgfpathlineto{\pgfqpoint{2.205624in}{2.151032in}}%
\pgfpathlineto{\pgfqpoint{2.208143in}{2.176486in}}%
\pgfpathlineto{\pgfqpoint{2.210662in}{2.191608in}}%
\pgfpathlineto{\pgfqpoint{2.213181in}{2.207861in}}%
\pgfpathlineto{\pgfqpoint{2.215699in}{2.214547in}}%
\pgfpathlineto{\pgfqpoint{2.218218in}{2.183840in}}%
\pgfpathlineto{\pgfqpoint{2.220737in}{2.198804in}}%
\pgfpathlineto{\pgfqpoint{2.223256in}{2.194937in}}%
\pgfpathlineto{\pgfqpoint{2.225774in}{2.207340in}}%
\pgfpathlineto{\pgfqpoint{2.228293in}{2.214598in}}%
\pgfpathlineto{\pgfqpoint{2.230812in}{2.175972in}}%
\pgfpathlineto{\pgfqpoint{2.233331in}{2.188459in}}%
\pgfpathlineto{\pgfqpoint{2.235849in}{2.189983in}}%
\pgfpathlineto{\pgfqpoint{2.238368in}{2.218070in}}%
\pgfpathlineto{\pgfqpoint{2.240887in}{2.226569in}}%
\pgfpathlineto{\pgfqpoint{2.243406in}{2.233683in}}%
\pgfpathlineto{\pgfqpoint{2.245924in}{2.213302in}}%
\pgfpathlineto{\pgfqpoint{2.248443in}{2.229749in}}%
\pgfpathlineto{\pgfqpoint{2.253481in}{2.253190in}}%
\pgfpathlineto{\pgfqpoint{2.255999in}{2.257488in}}%
\pgfpathlineto{\pgfqpoint{2.258518in}{2.260600in}}%
\pgfpathlineto{\pgfqpoint{2.261037in}{2.222863in}}%
\pgfpathlineto{\pgfqpoint{2.263556in}{2.229038in}}%
\pgfpathlineto{\pgfqpoint{2.266074in}{2.264433in}}%
\pgfpathlineto{\pgfqpoint{2.268593in}{2.211204in}}%
\pgfpathlineto{\pgfqpoint{2.271112in}{2.219215in}}%
\pgfpathlineto{\pgfqpoint{2.273631in}{2.168901in}}%
\pgfpathlineto{\pgfqpoint{2.276149in}{2.189439in}}%
\pgfpathlineto{\pgfqpoint{2.278668in}{2.157914in}}%
\pgfpathlineto{\pgfqpoint{2.281187in}{2.161900in}}%
\pgfpathlineto{\pgfqpoint{2.283706in}{2.156339in}}%
\pgfpathlineto{\pgfqpoint{2.286224in}{2.138132in}}%
\pgfpathlineto{\pgfqpoint{2.288743in}{2.130037in}}%
\pgfpathlineto{\pgfqpoint{2.291262in}{2.143242in}}%
\pgfpathlineto{\pgfqpoint{2.293781in}{2.164640in}}%
\pgfpathlineto{\pgfqpoint{2.296299in}{2.158719in}}%
\pgfpathlineto{\pgfqpoint{2.298818in}{2.171517in}}%
\pgfpathlineto{\pgfqpoint{2.301337in}{2.174871in}}%
\pgfpathlineto{\pgfqpoint{2.303856in}{2.156808in}}%
\pgfpathlineto{\pgfqpoint{2.306374in}{2.143719in}}%
\pgfpathlineto{\pgfqpoint{2.308893in}{2.135108in}}%
\pgfpathlineto{\pgfqpoint{2.311412in}{2.148750in}}%
\pgfpathlineto{\pgfqpoint{2.313931in}{2.138244in}}%
\pgfpathlineto{\pgfqpoint{2.316449in}{2.164518in}}%
\pgfpathlineto{\pgfqpoint{2.318968in}{2.171411in}}%
\pgfpathlineto{\pgfqpoint{2.321487in}{2.169281in}}%
\pgfpathlineto{\pgfqpoint{2.324006in}{2.144659in}}%
\pgfpathlineto{\pgfqpoint{2.326524in}{2.131171in}}%
\pgfpathlineto{\pgfqpoint{2.329043in}{2.104242in}}%
\pgfpathlineto{\pgfqpoint{2.331562in}{2.056528in}}%
\pgfpathlineto{\pgfqpoint{2.334081in}{2.066347in}}%
\pgfpathlineto{\pgfqpoint{2.336599in}{2.071189in}}%
\pgfpathlineto{\pgfqpoint{2.339118in}{2.090006in}}%
\pgfpathlineto{\pgfqpoint{2.344156in}{2.043065in}}%
\pgfpathlineto{\pgfqpoint{2.346674in}{2.031569in}}%
\pgfpathlineto{\pgfqpoint{2.349193in}{2.051635in}}%
\pgfpathlineto{\pgfqpoint{2.351712in}{2.033587in}}%
\pgfpathlineto{\pgfqpoint{2.354231in}{2.020173in}}%
\pgfpathlineto{\pgfqpoint{2.356749in}{2.021517in}}%
\pgfpathlineto{\pgfqpoint{2.359268in}{2.012746in}}%
\pgfpathlineto{\pgfqpoint{2.361787in}{2.026598in}}%
\pgfpathlineto{\pgfqpoint{2.364306in}{2.022413in}}%
\pgfpathlineto{\pgfqpoint{2.366824in}{2.052730in}}%
\pgfpathlineto{\pgfqpoint{2.369343in}{2.013436in}}%
\pgfpathlineto{\pgfqpoint{2.371862in}{1.989728in}}%
\pgfpathlineto{\pgfqpoint{2.374381in}{2.030673in}}%
\pgfpathlineto{\pgfqpoint{2.376899in}{2.059590in}}%
\pgfpathlineto{\pgfqpoint{2.379418in}{2.023713in}}%
\pgfpathlineto{\pgfqpoint{2.381937in}{2.017077in}}%
\pgfpathlineto{\pgfqpoint{2.384456in}{2.027550in}}%
\pgfpathlineto{\pgfqpoint{2.386974in}{2.043493in}}%
\pgfpathlineto{\pgfqpoint{2.389493in}{2.035353in}}%
\pgfpathlineto{\pgfqpoint{2.392012in}{2.071109in}}%
\pgfpathlineto{\pgfqpoint{2.394531in}{2.089584in}}%
\pgfpathlineto{\pgfqpoint{2.397049in}{2.063703in}}%
\pgfpathlineto{\pgfqpoint{2.399568in}{2.046994in}}%
\pgfpathlineto{\pgfqpoint{2.402087in}{2.037053in}}%
\pgfpathlineto{\pgfqpoint{2.404606in}{2.024215in}}%
\pgfpathlineto{\pgfqpoint{2.407124in}{2.019186in}}%
\pgfpathlineto{\pgfqpoint{2.409643in}{2.033204in}}%
\pgfpathlineto{\pgfqpoint{2.412162in}{2.059625in}}%
\pgfpathlineto{\pgfqpoint{2.414681in}{2.061433in}}%
\pgfpathlineto{\pgfqpoint{2.417199in}{2.031690in}}%
\pgfpathlineto{\pgfqpoint{2.419718in}{2.029742in}}%
\pgfpathlineto{\pgfqpoint{2.422237in}{1.994164in}}%
\pgfpathlineto{\pgfqpoint{2.424756in}{2.004381in}}%
\pgfpathlineto{\pgfqpoint{2.427274in}{2.021862in}}%
\pgfpathlineto{\pgfqpoint{2.429793in}{2.018250in}}%
\pgfpathlineto{\pgfqpoint{2.432312in}{2.031708in}}%
\pgfpathlineto{\pgfqpoint{2.434831in}{2.031799in}}%
\pgfpathlineto{\pgfqpoint{2.437349in}{2.036595in}}%
\pgfpathlineto{\pgfqpoint{2.439868in}{2.037793in}}%
\pgfpathlineto{\pgfqpoint{2.442387in}{2.066636in}}%
\pgfpathlineto{\pgfqpoint{2.444906in}{2.068240in}}%
\pgfpathlineto{\pgfqpoint{2.447424in}{2.058308in}}%
\pgfpathlineto{\pgfqpoint{2.449943in}{2.071584in}}%
\pgfpathlineto{\pgfqpoint{2.452462in}{2.062637in}}%
\pgfpathlineto{\pgfqpoint{2.454981in}{2.038863in}}%
\pgfpathlineto{\pgfqpoint{2.457499in}{2.023826in}}%
\pgfpathlineto{\pgfqpoint{2.460018in}{2.015088in}}%
\pgfpathlineto{\pgfqpoint{2.462537in}{2.020392in}}%
\pgfpathlineto{\pgfqpoint{2.465056in}{2.012207in}}%
\pgfpathlineto{\pgfqpoint{2.467574in}{2.027378in}}%
\pgfpathlineto{\pgfqpoint{2.470093in}{2.023904in}}%
\pgfpathlineto{\pgfqpoint{2.472612in}{2.034415in}}%
\pgfpathlineto{\pgfqpoint{2.475131in}{2.050302in}}%
\pgfpathlineto{\pgfqpoint{2.477649in}{2.037362in}}%
\pgfpathlineto{\pgfqpoint{2.480168in}{2.009546in}}%
\pgfpathlineto{\pgfqpoint{2.482687in}{2.011926in}}%
\pgfpathlineto{\pgfqpoint{2.485206in}{1.992526in}}%
\pgfpathlineto{\pgfqpoint{2.487724in}{2.012884in}}%
\pgfpathlineto{\pgfqpoint{2.490243in}{1.980332in}}%
\pgfpathlineto{\pgfqpoint{2.492762in}{2.009288in}}%
\pgfpathlineto{\pgfqpoint{2.495281in}{1.999426in}}%
\pgfpathlineto{\pgfqpoint{2.497799in}{1.985283in}}%
\pgfpathlineto{\pgfqpoint{2.500318in}{1.993738in}}%
\pgfpathlineto{\pgfqpoint{2.502837in}{1.941770in}}%
\pgfpathlineto{\pgfqpoint{2.505356in}{1.939673in}}%
\pgfpathlineto{\pgfqpoint{2.507874in}{1.949109in}}%
\pgfpathlineto{\pgfqpoint{2.512912in}{1.981768in}}%
\pgfpathlineto{\pgfqpoint{2.515431in}{1.982015in}}%
\pgfpathlineto{\pgfqpoint{2.517949in}{2.008322in}}%
\pgfpathlineto{\pgfqpoint{2.520468in}{2.020693in}}%
\pgfpathlineto{\pgfqpoint{2.522987in}{2.028566in}}%
\pgfpathlineto{\pgfqpoint{2.525506in}{2.026324in}}%
\pgfpathlineto{\pgfqpoint{2.528024in}{2.012386in}}%
\pgfpathlineto{\pgfqpoint{2.530543in}{2.008666in}}%
\pgfpathlineto{\pgfqpoint{2.533062in}{2.000508in}}%
\pgfpathlineto{\pgfqpoint{2.535581in}{2.003445in}}%
\pgfpathlineto{\pgfqpoint{2.538099in}{2.014655in}}%
\pgfpathlineto{\pgfqpoint{2.540618in}{2.003882in}}%
\pgfpathlineto{\pgfqpoint{2.543137in}{2.005868in}}%
\pgfpathlineto{\pgfqpoint{2.545656in}{2.011940in}}%
\pgfpathlineto{\pgfqpoint{2.548174in}{2.003227in}}%
\pgfpathlineto{\pgfqpoint{2.550693in}{1.991155in}}%
\pgfpathlineto{\pgfqpoint{2.553212in}{1.989901in}}%
\pgfpathlineto{\pgfqpoint{2.555731in}{2.002970in}}%
\pgfpathlineto{\pgfqpoint{2.558249in}{2.025238in}}%
\pgfpathlineto{\pgfqpoint{2.560768in}{2.049547in}}%
\pgfpathlineto{\pgfqpoint{2.563287in}{2.061347in}}%
\pgfpathlineto{\pgfqpoint{2.565806in}{2.050519in}}%
\pgfpathlineto{\pgfqpoint{2.568324in}{2.057398in}}%
\pgfpathlineto{\pgfqpoint{2.570843in}{2.025809in}}%
\pgfpathlineto{\pgfqpoint{2.573362in}{2.073035in}}%
\pgfpathlineto{\pgfqpoint{2.575881in}{2.086669in}}%
\pgfpathlineto{\pgfqpoint{2.578399in}{2.095701in}}%
\pgfpathlineto{\pgfqpoint{2.580918in}{2.116445in}}%
\pgfpathlineto{\pgfqpoint{2.583437in}{2.105697in}}%
\pgfpathlineto{\pgfqpoint{2.585956in}{2.058673in}}%
\pgfpathlineto{\pgfqpoint{2.588474in}{2.056421in}}%
\pgfpathlineto{\pgfqpoint{2.590993in}{2.058206in}}%
\pgfpathlineto{\pgfqpoint{2.593512in}{2.043808in}}%
\pgfpathlineto{\pgfqpoint{2.596031in}{2.051905in}}%
\pgfpathlineto{\pgfqpoint{2.598549in}{2.084513in}}%
\pgfpathlineto{\pgfqpoint{2.601068in}{2.085232in}}%
\pgfpathlineto{\pgfqpoint{2.603587in}{2.095425in}}%
\pgfpathlineto{\pgfqpoint{2.606106in}{2.104796in}}%
\pgfpathlineto{\pgfqpoint{2.608624in}{2.139526in}}%
\pgfpathlineto{\pgfqpoint{2.611143in}{2.127444in}}%
\pgfpathlineto{\pgfqpoint{2.613662in}{2.137552in}}%
\pgfpathlineto{\pgfqpoint{2.616181in}{2.152862in}}%
\pgfpathlineto{\pgfqpoint{2.618699in}{2.151306in}}%
\pgfpathlineto{\pgfqpoint{2.621218in}{2.151337in}}%
\pgfpathlineto{\pgfqpoint{2.623737in}{2.122166in}}%
\pgfpathlineto{\pgfqpoint{2.626256in}{2.137741in}}%
\pgfpathlineto{\pgfqpoint{2.628774in}{2.131504in}}%
\pgfpathlineto{\pgfqpoint{2.631293in}{2.117348in}}%
\pgfpathlineto{\pgfqpoint{2.633812in}{2.141342in}}%
\pgfpathlineto{\pgfqpoint{2.636331in}{2.119418in}}%
\pgfpathlineto{\pgfqpoint{2.638849in}{2.112808in}}%
\pgfpathlineto{\pgfqpoint{2.641368in}{2.132363in}}%
\pgfpathlineto{\pgfqpoint{2.643887in}{2.129559in}}%
\pgfpathlineto{\pgfqpoint{2.646406in}{2.145156in}}%
\pgfpathlineto{\pgfqpoint{2.648924in}{2.131414in}}%
\pgfpathlineto{\pgfqpoint{2.651443in}{2.105265in}}%
\pgfpathlineto{\pgfqpoint{2.653962in}{2.142678in}}%
\pgfpathlineto{\pgfqpoint{2.656481in}{2.149463in}}%
\pgfpathlineto{\pgfqpoint{2.658999in}{2.147369in}}%
\pgfpathlineto{\pgfqpoint{2.661518in}{2.146609in}}%
\pgfpathlineto{\pgfqpoint{2.664037in}{2.156091in}}%
\pgfpathlineto{\pgfqpoint{2.666556in}{2.117532in}}%
\pgfpathlineto{\pgfqpoint{2.669074in}{2.125395in}}%
\pgfpathlineto{\pgfqpoint{2.671593in}{2.129300in}}%
\pgfpathlineto{\pgfqpoint{2.674112in}{2.101013in}}%
\pgfpathlineto{\pgfqpoint{2.676631in}{2.100100in}}%
\pgfpathlineto{\pgfqpoint{2.679149in}{2.064403in}}%
\pgfpathlineto{\pgfqpoint{2.681668in}{2.086084in}}%
\pgfpathlineto{\pgfqpoint{2.684187in}{2.084909in}}%
\pgfpathlineto{\pgfqpoint{2.686706in}{2.093145in}}%
\pgfpathlineto{\pgfqpoint{2.689224in}{2.080523in}}%
\pgfpathlineto{\pgfqpoint{2.691743in}{2.108159in}}%
\pgfpathlineto{\pgfqpoint{2.694262in}{2.141534in}}%
\pgfpathlineto{\pgfqpoint{2.696781in}{2.139152in}}%
\pgfpathlineto{\pgfqpoint{2.699299in}{2.158032in}}%
\pgfpathlineto{\pgfqpoint{2.701818in}{2.162665in}}%
\pgfpathlineto{\pgfqpoint{2.704337in}{2.184952in}}%
\pgfpathlineto{\pgfqpoint{2.706856in}{2.198585in}}%
\pgfpathlineto{\pgfqpoint{2.709374in}{2.216353in}}%
\pgfpathlineto{\pgfqpoint{2.711893in}{2.233083in}}%
\pgfpathlineto{\pgfqpoint{2.714412in}{2.210107in}}%
\pgfpathlineto{\pgfqpoint{2.716931in}{2.224874in}}%
\pgfpathlineto{\pgfqpoint{2.719449in}{2.260825in}}%
\pgfpathlineto{\pgfqpoint{2.721968in}{2.266089in}}%
\pgfpathlineto{\pgfqpoint{2.724487in}{2.268079in}}%
\pgfpathlineto{\pgfqpoint{2.727006in}{2.280766in}}%
\pgfpathlineto{\pgfqpoint{2.729524in}{2.261686in}}%
\pgfpathlineto{\pgfqpoint{2.732043in}{2.256789in}}%
\pgfpathlineto{\pgfqpoint{2.734562in}{2.264524in}}%
\pgfpathlineto{\pgfqpoint{2.737081in}{2.286416in}}%
\pgfpathlineto{\pgfqpoint{2.739599in}{2.290658in}}%
\pgfpathlineto{\pgfqpoint{2.742118in}{2.267912in}}%
\pgfpathlineto{\pgfqpoint{2.744637in}{2.295379in}}%
\pgfpathlineto{\pgfqpoint{2.747156in}{2.273946in}}%
\pgfpathlineto{\pgfqpoint{2.749674in}{2.250937in}}%
\pgfpathlineto{\pgfqpoint{2.754712in}{2.304708in}}%
\pgfpathlineto{\pgfqpoint{2.757231in}{2.317790in}}%
\pgfpathlineto{\pgfqpoint{2.759749in}{2.349758in}}%
\pgfpathlineto{\pgfqpoint{2.762268in}{2.349323in}}%
\pgfpathlineto{\pgfqpoint{2.764787in}{2.332466in}}%
\pgfpathlineto{\pgfqpoint{2.767306in}{2.310761in}}%
\pgfpathlineto{\pgfqpoint{2.769824in}{2.308413in}}%
\pgfpathlineto{\pgfqpoint{2.772343in}{2.321769in}}%
\pgfpathlineto{\pgfqpoint{2.774862in}{2.354604in}}%
\pgfpathlineto{\pgfqpoint{2.777381in}{2.346573in}}%
\pgfpathlineto{\pgfqpoint{2.779899in}{2.368399in}}%
\pgfpathlineto{\pgfqpoint{2.782418in}{2.355784in}}%
\pgfpathlineto{\pgfqpoint{2.784937in}{2.353871in}}%
\pgfpathlineto{\pgfqpoint{2.787456in}{2.368785in}}%
\pgfpathlineto{\pgfqpoint{2.789974in}{2.361139in}}%
\pgfpathlineto{\pgfqpoint{2.792493in}{2.378087in}}%
\pgfpathlineto{\pgfqpoint{2.795012in}{2.362402in}}%
\pgfpathlineto{\pgfqpoint{2.797531in}{2.371037in}}%
\pgfpathlineto{\pgfqpoint{2.800049in}{2.374298in}}%
\pgfpathlineto{\pgfqpoint{2.802568in}{2.399064in}}%
\pgfpathlineto{\pgfqpoint{2.805087in}{2.380853in}}%
\pgfpathlineto{\pgfqpoint{2.807606in}{2.397817in}}%
\pgfpathlineto{\pgfqpoint{2.810124in}{2.370439in}}%
\pgfpathlineto{\pgfqpoint{2.815162in}{2.357150in}}%
\pgfpathlineto{\pgfqpoint{2.817681in}{2.360583in}}%
\pgfpathlineto{\pgfqpoint{2.820199in}{2.340121in}}%
\pgfpathlineto{\pgfqpoint{2.822718in}{2.363455in}}%
\pgfpathlineto{\pgfqpoint{2.825237in}{2.388640in}}%
\pgfpathlineto{\pgfqpoint{2.827756in}{2.401103in}}%
\pgfpathlineto{\pgfqpoint{2.830274in}{2.355034in}}%
\pgfpathlineto{\pgfqpoint{2.832793in}{2.355031in}}%
\pgfpathlineto{\pgfqpoint{2.835312in}{2.359734in}}%
\pgfpathlineto{\pgfqpoint{2.837831in}{2.378219in}}%
\pgfpathlineto{\pgfqpoint{2.840349in}{2.392301in}}%
\pgfpathlineto{\pgfqpoint{2.842868in}{2.380720in}}%
\pgfpathlineto{\pgfqpoint{2.845387in}{2.355024in}}%
\pgfpathlineto{\pgfqpoint{2.847906in}{2.387367in}}%
\pgfpathlineto{\pgfqpoint{2.850424in}{2.388206in}}%
\pgfpathlineto{\pgfqpoint{2.852943in}{2.369945in}}%
\pgfpathlineto{\pgfqpoint{2.855462in}{2.341319in}}%
\pgfpathlineto{\pgfqpoint{2.857981in}{2.359404in}}%
\pgfpathlineto{\pgfqpoint{2.860499in}{2.352608in}}%
\pgfpathlineto{\pgfqpoint{2.863018in}{2.343831in}}%
\pgfpathlineto{\pgfqpoint{2.865537in}{2.345592in}}%
\pgfpathlineto{\pgfqpoint{2.868056in}{2.325010in}}%
\pgfpathlineto{\pgfqpoint{2.870574in}{2.341939in}}%
\pgfpathlineto{\pgfqpoint{2.873093in}{2.327476in}}%
\pgfpathlineto{\pgfqpoint{2.875612in}{2.329407in}}%
\pgfpathlineto{\pgfqpoint{2.878131in}{2.310954in}}%
\pgfpathlineto{\pgfqpoint{2.880649in}{2.309052in}}%
\pgfpathlineto{\pgfqpoint{2.883168in}{2.309443in}}%
\pgfpathlineto{\pgfqpoint{2.885687in}{2.312254in}}%
\pgfpathlineto{\pgfqpoint{2.888206in}{2.295967in}}%
\pgfpathlineto{\pgfqpoint{2.890724in}{2.282388in}}%
\pgfpathlineto{\pgfqpoint{2.893243in}{2.313401in}}%
\pgfpathlineto{\pgfqpoint{2.895762in}{2.308958in}}%
\pgfpathlineto{\pgfqpoint{2.898281in}{2.308774in}}%
\pgfpathlineto{\pgfqpoint{2.900799in}{2.351396in}}%
\pgfpathlineto{\pgfqpoint{2.903318in}{2.322670in}}%
\pgfpathlineto{\pgfqpoint{2.905837in}{2.357913in}}%
\pgfpathlineto{\pgfqpoint{2.908356in}{2.352308in}}%
\pgfpathlineto{\pgfqpoint{2.910874in}{2.370930in}}%
\pgfpathlineto{\pgfqpoint{2.913393in}{2.376237in}}%
\pgfpathlineto{\pgfqpoint{2.915912in}{2.375353in}}%
\pgfpathlineto{\pgfqpoint{2.918431in}{2.390471in}}%
\pgfpathlineto{\pgfqpoint{2.920949in}{2.384010in}}%
\pgfpathlineto{\pgfqpoint{2.923468in}{2.420066in}}%
\pgfpathlineto{\pgfqpoint{2.925987in}{2.382688in}}%
\pgfpathlineto{\pgfqpoint{2.928506in}{2.383649in}}%
\pgfpathlineto{\pgfqpoint{2.931024in}{2.368000in}}%
\pgfpathlineto{\pgfqpoint{2.933543in}{2.377165in}}%
\pgfpathlineto{\pgfqpoint{2.936062in}{2.370277in}}%
\pgfpathlineto{\pgfqpoint{2.938581in}{2.350397in}}%
\pgfpathlineto{\pgfqpoint{2.941099in}{2.348373in}}%
\pgfpathlineto{\pgfqpoint{2.943618in}{2.403018in}}%
\pgfpathlineto{\pgfqpoint{2.946137in}{2.417934in}}%
\pgfpathlineto{\pgfqpoint{2.948656in}{2.423287in}}%
\pgfpathlineto{\pgfqpoint{2.953693in}{2.396324in}}%
\pgfpathlineto{\pgfqpoint{2.956212in}{2.366330in}}%
\pgfpathlineto{\pgfqpoint{2.958731in}{2.343892in}}%
\pgfpathlineto{\pgfqpoint{2.961249in}{2.345410in}}%
\pgfpathlineto{\pgfqpoint{2.963768in}{2.350365in}}%
\pgfpathlineto{\pgfqpoint{2.966287in}{2.372680in}}%
\pgfpathlineto{\pgfqpoint{2.968806in}{2.357464in}}%
\pgfpathlineto{\pgfqpoint{2.971324in}{2.358947in}}%
\pgfpathlineto{\pgfqpoint{2.973843in}{2.386211in}}%
\pgfpathlineto{\pgfqpoint{2.976362in}{2.388464in}}%
\pgfpathlineto{\pgfqpoint{2.978881in}{2.396344in}}%
\pgfpathlineto{\pgfqpoint{2.981399in}{2.406169in}}%
\pgfpathlineto{\pgfqpoint{2.983918in}{2.467997in}}%
\pgfpathlineto{\pgfqpoint{2.986437in}{2.448659in}}%
\pgfpathlineto{\pgfqpoint{2.988956in}{2.461739in}}%
\pgfpathlineto{\pgfqpoint{2.991474in}{2.463268in}}%
\pgfpathlineto{\pgfqpoint{2.993993in}{2.447258in}}%
\pgfpathlineto{\pgfqpoint{2.996512in}{2.472530in}}%
\pgfpathlineto{\pgfqpoint{2.999031in}{2.467443in}}%
\pgfpathlineto{\pgfqpoint{3.001549in}{2.501485in}}%
\pgfpathlineto{\pgfqpoint{3.004068in}{2.503415in}}%
\pgfpathlineto{\pgfqpoint{3.006587in}{2.538149in}}%
\pgfpathlineto{\pgfqpoint{3.009106in}{2.503646in}}%
\pgfpathlineto{\pgfqpoint{3.011624in}{2.525430in}}%
\pgfpathlineto{\pgfqpoint{3.014143in}{2.559139in}}%
\pgfpathlineto{\pgfqpoint{3.016662in}{2.572827in}}%
\pgfpathlineto{\pgfqpoint{3.019181in}{2.550329in}}%
\pgfpathlineto{\pgfqpoint{3.021699in}{2.554712in}}%
\pgfpathlineto{\pgfqpoint{3.024218in}{2.558760in}}%
\pgfpathlineto{\pgfqpoint{3.029256in}{2.613528in}}%
\pgfpathlineto{\pgfqpoint{3.031774in}{2.629873in}}%
\pgfpathlineto{\pgfqpoint{3.034293in}{2.660747in}}%
\pgfpathlineto{\pgfqpoint{3.036812in}{2.641871in}}%
\pgfpathlineto{\pgfqpoint{3.039331in}{2.670089in}}%
\pgfpathlineto{\pgfqpoint{3.041849in}{2.655654in}}%
\pgfpathlineto{\pgfqpoint{3.044368in}{2.636974in}}%
\pgfpathlineto{\pgfqpoint{3.046887in}{2.649354in}}%
\pgfpathlineto{\pgfqpoint{3.049406in}{2.646804in}}%
\pgfpathlineto{\pgfqpoint{3.054443in}{2.631660in}}%
\pgfpathlineto{\pgfqpoint{3.056962in}{2.592398in}}%
\pgfpathlineto{\pgfqpoint{3.059481in}{2.615910in}}%
\pgfpathlineto{\pgfqpoint{3.061999in}{2.633850in}}%
\pgfpathlineto{\pgfqpoint{3.064518in}{2.586246in}}%
\pgfpathlineto{\pgfqpoint{3.067037in}{2.618207in}}%
\pgfpathlineto{\pgfqpoint{3.069556in}{2.636884in}}%
\pgfpathlineto{\pgfqpoint{3.072074in}{2.642821in}}%
\pgfpathlineto{\pgfqpoint{3.074593in}{2.644216in}}%
\pgfpathlineto{\pgfqpoint{3.077112in}{2.675204in}}%
\pgfpathlineto{\pgfqpoint{3.079631in}{2.678547in}}%
\pgfpathlineto{\pgfqpoint{3.082149in}{2.721281in}}%
\pgfpathlineto{\pgfqpoint{3.084668in}{2.746916in}}%
\pgfpathlineto{\pgfqpoint{3.087187in}{2.721416in}}%
\pgfpathlineto{\pgfqpoint{3.089706in}{2.691464in}}%
\pgfpathlineto{\pgfqpoint{3.092224in}{2.691199in}}%
\pgfpathlineto{\pgfqpoint{3.094743in}{2.685132in}}%
\pgfpathlineto{\pgfqpoint{3.097262in}{2.656528in}}%
\pgfpathlineto{\pgfqpoint{3.099781in}{2.689221in}}%
\pgfpathlineto{\pgfqpoint{3.102299in}{2.719927in}}%
\pgfpathlineto{\pgfqpoint{3.104818in}{2.710146in}}%
\pgfpathlineto{\pgfqpoint{3.107337in}{2.743277in}}%
\pgfpathlineto{\pgfqpoint{3.109856in}{2.705709in}}%
\pgfpathlineto{\pgfqpoint{3.112374in}{2.676840in}}%
\pgfpathlineto{\pgfqpoint{3.114893in}{2.661672in}}%
\pgfpathlineto{\pgfqpoint{3.117412in}{2.697307in}}%
\pgfpathlineto{\pgfqpoint{3.119931in}{2.734433in}}%
\pgfpathlineto{\pgfqpoint{3.122449in}{2.751115in}}%
\pgfpathlineto{\pgfqpoint{3.124968in}{2.727806in}}%
\pgfpathlineto{\pgfqpoint{3.127487in}{2.693348in}}%
\pgfpathlineto{\pgfqpoint{3.130006in}{2.714384in}}%
\pgfpathlineto{\pgfqpoint{3.135043in}{2.752786in}}%
\pgfpathlineto{\pgfqpoint{3.137562in}{2.737736in}}%
\pgfpathlineto{\pgfqpoint{3.140081in}{2.767592in}}%
\pgfpathlineto{\pgfqpoint{3.142599in}{2.762860in}}%
\pgfpathlineto{\pgfqpoint{3.145118in}{2.758933in}}%
\pgfpathlineto{\pgfqpoint{3.147637in}{2.783247in}}%
\pgfpathlineto{\pgfqpoint{3.150156in}{2.793664in}}%
\pgfpathlineto{\pgfqpoint{3.152674in}{2.804817in}}%
\pgfpathlineto{\pgfqpoint{3.155193in}{2.796683in}}%
\pgfpathlineto{\pgfqpoint{3.157712in}{2.770648in}}%
\pgfpathlineto{\pgfqpoint{3.160231in}{2.783577in}}%
\pgfpathlineto{\pgfqpoint{3.162749in}{2.784579in}}%
\pgfpathlineto{\pgfqpoint{3.165268in}{2.787144in}}%
\pgfpathlineto{\pgfqpoint{3.167787in}{2.800271in}}%
\pgfpathlineto{\pgfqpoint{3.170306in}{2.790641in}}%
\pgfpathlineto{\pgfqpoint{3.172824in}{2.853694in}}%
\pgfpathlineto{\pgfqpoint{3.175343in}{2.826407in}}%
\pgfpathlineto{\pgfqpoint{3.177862in}{2.811682in}}%
\pgfpathlineto{\pgfqpoint{3.180381in}{2.781275in}}%
\pgfpathlineto{\pgfqpoint{3.182899in}{2.809557in}}%
\pgfpathlineto{\pgfqpoint{3.185418in}{2.823246in}}%
\pgfpathlineto{\pgfqpoint{3.187937in}{2.821479in}}%
\pgfpathlineto{\pgfqpoint{3.190456in}{2.870618in}}%
\pgfpathlineto{\pgfqpoint{3.192974in}{2.895696in}}%
\pgfpathlineto{\pgfqpoint{3.195493in}{2.903344in}}%
\pgfpathlineto{\pgfqpoint{3.198012in}{2.950428in}}%
\pgfpathlineto{\pgfqpoint{3.200531in}{2.963752in}}%
\pgfpathlineto{\pgfqpoint{3.203049in}{2.974509in}}%
\pgfpathlineto{\pgfqpoint{3.205568in}{2.972985in}}%
\pgfpathlineto{\pgfqpoint{3.208087in}{2.969960in}}%
\pgfpathlineto{\pgfqpoint{3.210606in}{2.961605in}}%
\pgfpathlineto{\pgfqpoint{3.213124in}{2.939579in}}%
\pgfpathlineto{\pgfqpoint{3.215643in}{2.904749in}}%
\pgfpathlineto{\pgfqpoint{3.218162in}{2.908426in}}%
\pgfpathlineto{\pgfqpoint{3.220681in}{2.876683in}}%
\pgfpathlineto{\pgfqpoint{3.223199in}{2.870499in}}%
\pgfpathlineto{\pgfqpoint{3.225718in}{2.876855in}}%
\pgfpathlineto{\pgfqpoint{3.228237in}{2.831373in}}%
\pgfpathlineto{\pgfqpoint{3.230756in}{2.826439in}}%
\pgfpathlineto{\pgfqpoint{3.233274in}{2.858661in}}%
\pgfpathlineto{\pgfqpoint{3.235793in}{2.837178in}}%
\pgfpathlineto{\pgfqpoint{3.238312in}{2.841717in}}%
\pgfpathlineto{\pgfqpoint{3.240831in}{2.873123in}}%
\pgfpathlineto{\pgfqpoint{3.245868in}{2.909972in}}%
\pgfpathlineto{\pgfqpoint{3.248387in}{2.933386in}}%
\pgfpathlineto{\pgfqpoint{3.250906in}{2.922468in}}%
\pgfpathlineto{\pgfqpoint{3.253424in}{2.953526in}}%
\pgfpathlineto{\pgfqpoint{3.255943in}{2.963137in}}%
\pgfpathlineto{\pgfqpoint{3.258462in}{2.902482in}}%
\pgfpathlineto{\pgfqpoint{3.260981in}{2.895354in}}%
\pgfpathlineto{\pgfqpoint{3.263499in}{2.881150in}}%
\pgfpathlineto{\pgfqpoint{3.266018in}{2.827207in}}%
\pgfpathlineto{\pgfqpoint{3.268537in}{2.798785in}}%
\pgfpathlineto{\pgfqpoint{3.271056in}{2.801887in}}%
\pgfpathlineto{\pgfqpoint{3.273574in}{2.789774in}}%
\pgfpathlineto{\pgfqpoint{3.276093in}{2.768690in}}%
\pgfpathlineto{\pgfqpoint{3.278612in}{2.757445in}}%
\pgfpathlineto{\pgfqpoint{3.281131in}{2.767819in}}%
\pgfpathlineto{\pgfqpoint{3.283649in}{2.766701in}}%
\pgfpathlineto{\pgfqpoint{3.286168in}{2.770351in}}%
\pgfpathlineto{\pgfqpoint{3.288687in}{2.765291in}}%
\pgfpathlineto{\pgfqpoint{3.291206in}{2.754550in}}%
\pgfpathlineto{\pgfqpoint{3.293724in}{2.754458in}}%
\pgfpathlineto{\pgfqpoint{3.296243in}{2.756396in}}%
\pgfpathlineto{\pgfqpoint{3.298762in}{2.744760in}}%
\pgfpathlineto{\pgfqpoint{3.301281in}{2.751683in}}%
\pgfpathlineto{\pgfqpoint{3.303799in}{2.738083in}}%
\pgfpathlineto{\pgfqpoint{3.306318in}{2.729464in}}%
\pgfpathlineto{\pgfqpoint{3.308837in}{2.707063in}}%
\pgfpathlineto{\pgfqpoint{3.311356in}{2.724872in}}%
\pgfpathlineto{\pgfqpoint{3.313874in}{2.706672in}}%
\pgfpathlineto{\pgfqpoint{3.316393in}{2.742178in}}%
\pgfpathlineto{\pgfqpoint{3.318912in}{2.797526in}}%
\pgfpathlineto{\pgfqpoint{3.321431in}{2.776994in}}%
\pgfpathlineto{\pgfqpoint{3.323949in}{2.743165in}}%
\pgfpathlineto{\pgfqpoint{3.326468in}{2.765601in}}%
\pgfpathlineto{\pgfqpoint{3.328987in}{2.740547in}}%
\pgfpathlineto{\pgfqpoint{3.331506in}{2.771406in}}%
\pgfpathlineto{\pgfqpoint{3.334024in}{2.769104in}}%
\pgfpathlineto{\pgfqpoint{3.336543in}{2.764812in}}%
\pgfpathlineto{\pgfqpoint{3.339062in}{2.762677in}}%
\pgfpathlineto{\pgfqpoint{3.341581in}{2.755213in}}%
\pgfpathlineto{\pgfqpoint{3.344099in}{2.767934in}}%
\pgfpathlineto{\pgfqpoint{3.346618in}{2.784923in}}%
\pgfpathlineto{\pgfqpoint{3.349137in}{2.796493in}}%
\pgfpathlineto{\pgfqpoint{3.351656in}{2.790553in}}%
\pgfpathlineto{\pgfqpoint{3.354174in}{2.793841in}}%
\pgfpathlineto{\pgfqpoint{3.356693in}{2.825205in}}%
\pgfpathlineto{\pgfqpoint{3.359212in}{2.790581in}}%
\pgfpathlineto{\pgfqpoint{3.361731in}{2.801311in}}%
\pgfpathlineto{\pgfqpoint{3.364249in}{2.795521in}}%
\pgfpathlineto{\pgfqpoint{3.366768in}{2.795541in}}%
\pgfpathlineto{\pgfqpoint{3.369287in}{2.802862in}}%
\pgfpathlineto{\pgfqpoint{3.371806in}{2.782138in}}%
\pgfpathlineto{\pgfqpoint{3.374324in}{2.793910in}}%
\pgfpathlineto{\pgfqpoint{3.376843in}{2.832947in}}%
\pgfpathlineto{\pgfqpoint{3.379362in}{2.808145in}}%
\pgfpathlineto{\pgfqpoint{3.381881in}{2.807188in}}%
\pgfpathlineto{\pgfqpoint{3.384399in}{2.842743in}}%
\pgfpathlineto{\pgfqpoint{3.386918in}{2.842805in}}%
\pgfpathlineto{\pgfqpoint{3.389437in}{2.881528in}}%
\pgfpathlineto{\pgfqpoint{3.391956in}{2.882703in}}%
\pgfpathlineto{\pgfqpoint{3.394474in}{2.892530in}}%
\pgfpathlineto{\pgfqpoint{3.399512in}{2.899404in}}%
\pgfpathlineto{\pgfqpoint{3.402031in}{2.912145in}}%
\pgfpathlineto{\pgfqpoint{3.404549in}{2.926502in}}%
\pgfpathlineto{\pgfqpoint{3.407068in}{2.932326in}}%
\pgfpathlineto{\pgfqpoint{3.409587in}{2.923281in}}%
\pgfpathlineto{\pgfqpoint{3.412106in}{2.940191in}}%
\pgfpathlineto{\pgfqpoint{3.414624in}{2.932999in}}%
\pgfpathlineto{\pgfqpoint{3.417143in}{2.961812in}}%
\pgfpathlineto{\pgfqpoint{3.419662in}{2.960744in}}%
\pgfpathlineto{\pgfqpoint{3.422181in}{2.933747in}}%
\pgfpathlineto{\pgfqpoint{3.424699in}{2.976281in}}%
\pgfpathlineto{\pgfqpoint{3.427218in}{2.975758in}}%
\pgfpathlineto{\pgfqpoint{3.429737in}{2.986539in}}%
\pgfpathlineto{\pgfqpoint{3.432256in}{2.953790in}}%
\pgfpathlineto{\pgfqpoint{3.434774in}{2.938922in}}%
\pgfpathlineto{\pgfqpoint{3.437293in}{2.914875in}}%
\pgfpathlineto{\pgfqpoint{3.439812in}{2.885540in}}%
\pgfpathlineto{\pgfqpoint{3.442331in}{2.860386in}}%
\pgfpathlineto{\pgfqpoint{3.444849in}{2.858228in}}%
\pgfpathlineto{\pgfqpoint{3.447368in}{2.849899in}}%
\pgfpathlineto{\pgfqpoint{3.449887in}{2.855456in}}%
\pgfpathlineto{\pgfqpoint{3.452406in}{2.837148in}}%
\pgfpathlineto{\pgfqpoint{3.454924in}{2.830351in}}%
\pgfpathlineto{\pgfqpoint{3.457443in}{2.796209in}}%
\pgfpathlineto{\pgfqpoint{3.459962in}{2.807467in}}%
\pgfpathlineto{\pgfqpoint{3.462481in}{2.806707in}}%
\pgfpathlineto{\pgfqpoint{3.464999in}{2.820957in}}%
\pgfpathlineto{\pgfqpoint{3.467518in}{2.834094in}}%
\pgfpathlineto{\pgfqpoint{3.470037in}{2.817289in}}%
\pgfpathlineto{\pgfqpoint{3.472556in}{2.832298in}}%
\pgfpathlineto{\pgfqpoint{3.475074in}{2.832263in}}%
\pgfpathlineto{\pgfqpoint{3.477593in}{2.851448in}}%
\pgfpathlineto{\pgfqpoint{3.480112in}{2.851757in}}%
\pgfpathlineto{\pgfqpoint{3.482631in}{2.831195in}}%
\pgfpathlineto{\pgfqpoint{3.485149in}{2.798527in}}%
\pgfpathlineto{\pgfqpoint{3.487668in}{2.822252in}}%
\pgfpathlineto{\pgfqpoint{3.490187in}{2.824795in}}%
\pgfpathlineto{\pgfqpoint{3.492706in}{2.825808in}}%
\pgfpathlineto{\pgfqpoint{3.495224in}{2.847536in}}%
\pgfpathlineto{\pgfqpoint{3.497743in}{2.887038in}}%
\pgfpathlineto{\pgfqpoint{3.500262in}{2.887459in}}%
\pgfpathlineto{\pgfqpoint{3.502781in}{2.874932in}}%
\pgfpathlineto{\pgfqpoint{3.505299in}{2.885914in}}%
\pgfpathlineto{\pgfqpoint{3.507818in}{2.916141in}}%
\pgfpathlineto{\pgfqpoint{3.510337in}{2.900882in}}%
\pgfpathlineto{\pgfqpoint{3.512856in}{2.920378in}}%
\pgfpathlineto{\pgfqpoint{3.515374in}{2.896036in}}%
\pgfpathlineto{\pgfqpoint{3.517893in}{2.913523in}}%
\pgfpathlineto{\pgfqpoint{3.520412in}{2.947634in}}%
\pgfpathlineto{\pgfqpoint{3.522931in}{2.985066in}}%
\pgfpathlineto{\pgfqpoint{3.525449in}{2.996491in}}%
\pgfpathlineto{\pgfqpoint{3.527968in}{3.019114in}}%
\pgfpathlineto{\pgfqpoint{3.530487in}{3.023478in}}%
\pgfpathlineto{\pgfqpoint{3.533006in}{2.990967in}}%
\pgfpathlineto{\pgfqpoint{3.535524in}{2.996164in}}%
\pgfpathlineto{\pgfqpoint{3.538043in}{2.988241in}}%
\pgfpathlineto{\pgfqpoint{3.540562in}{2.958401in}}%
\pgfpathlineto{\pgfqpoint{3.543081in}{2.979680in}}%
\pgfpathlineto{\pgfqpoint{3.545599in}{2.976537in}}%
\pgfpathlineto{\pgfqpoint{3.548118in}{2.987352in}}%
\pgfpathlineto{\pgfqpoint{3.550637in}{2.988945in}}%
\pgfpathlineto{\pgfqpoint{3.553156in}{2.972001in}}%
\pgfpathlineto{\pgfqpoint{3.555674in}{3.003516in}}%
\pgfpathlineto{\pgfqpoint{3.558193in}{3.003694in}}%
\pgfpathlineto{\pgfqpoint{3.560712in}{3.015534in}}%
\pgfpathlineto{\pgfqpoint{3.563231in}{3.005006in}}%
\pgfpathlineto{\pgfqpoint{3.565749in}{3.037756in}}%
\pgfpathlineto{\pgfqpoint{3.568268in}{3.067089in}}%
\pgfpathlineto{\pgfqpoint{3.570787in}{3.073592in}}%
\pgfpathlineto{\pgfqpoint{3.573306in}{3.046650in}}%
\pgfpathlineto{\pgfqpoint{3.575824in}{3.032577in}}%
\pgfpathlineto{\pgfqpoint{3.578343in}{3.045994in}}%
\pgfpathlineto{\pgfqpoint{3.580862in}{3.067487in}}%
\pgfpathlineto{\pgfqpoint{3.583381in}{3.040082in}}%
\pgfpathlineto{\pgfqpoint{3.585899in}{3.045004in}}%
\pgfpathlineto{\pgfqpoint{3.588418in}{3.085303in}}%
\pgfpathlineto{\pgfqpoint{3.590937in}{3.105015in}}%
\pgfpathlineto{\pgfqpoint{3.593456in}{3.091597in}}%
\pgfpathlineto{\pgfqpoint{3.595974in}{3.067150in}}%
\pgfpathlineto{\pgfqpoint{3.598493in}{3.074339in}}%
\pgfpathlineto{\pgfqpoint{3.601012in}{3.042524in}}%
\pgfpathlineto{\pgfqpoint{3.603531in}{3.023462in}}%
\pgfpathlineto{\pgfqpoint{3.606049in}{3.018853in}}%
\pgfpathlineto{\pgfqpoint{3.608568in}{3.037764in}}%
\pgfpathlineto{\pgfqpoint{3.611087in}{3.009700in}}%
\pgfpathlineto{\pgfqpoint{3.616124in}{3.020539in}}%
\pgfpathlineto{\pgfqpoint{3.618643in}{3.022057in}}%
\pgfpathlineto{\pgfqpoint{3.621162in}{3.011566in}}%
\pgfpathlineto{\pgfqpoint{3.623681in}{3.049183in}}%
\pgfpathlineto{\pgfqpoint{3.626199in}{3.025811in}}%
\pgfpathlineto{\pgfqpoint{3.628718in}{2.998272in}}%
\pgfpathlineto{\pgfqpoint{3.631237in}{2.977693in}}%
\pgfpathlineto{\pgfqpoint{3.633756in}{3.001976in}}%
\pgfpathlineto{\pgfqpoint{3.636274in}{3.037897in}}%
\pgfpathlineto{\pgfqpoint{3.638793in}{3.040893in}}%
\pgfpathlineto{\pgfqpoint{3.641312in}{3.036743in}}%
\pgfpathlineto{\pgfqpoint{3.643831in}{3.004446in}}%
\pgfpathlineto{\pgfqpoint{3.646349in}{2.959508in}}%
\pgfpathlineto{\pgfqpoint{3.648868in}{2.952815in}}%
\pgfpathlineto{\pgfqpoint{3.651387in}{2.959028in}}%
\pgfpathlineto{\pgfqpoint{3.653906in}{2.963920in}}%
\pgfpathlineto{\pgfqpoint{3.656424in}{2.920940in}}%
\pgfpathlineto{\pgfqpoint{3.658943in}{2.966033in}}%
\pgfpathlineto{\pgfqpoint{3.661462in}{3.024219in}}%
\pgfpathlineto{\pgfqpoint{3.663981in}{2.982933in}}%
\pgfpathlineto{\pgfqpoint{3.666499in}{2.992467in}}%
\pgfpathlineto{\pgfqpoint{3.669018in}{3.000508in}}%
\pgfpathlineto{\pgfqpoint{3.674056in}{3.033413in}}%
\pgfpathlineto{\pgfqpoint{3.676574in}{3.041222in}}%
\pgfpathlineto{\pgfqpoint{3.679093in}{3.040138in}}%
\pgfpathlineto{\pgfqpoint{3.681612in}{3.041314in}}%
\pgfpathlineto{\pgfqpoint{3.684131in}{3.046909in}}%
\pgfpathlineto{\pgfqpoint{3.686649in}{3.019908in}}%
\pgfpathlineto{\pgfqpoint{3.689168in}{3.038324in}}%
\pgfpathlineto{\pgfqpoint{3.691687in}{3.061802in}}%
\pgfpathlineto{\pgfqpoint{3.694206in}{3.095936in}}%
\pgfpathlineto{\pgfqpoint{3.696724in}{3.098320in}}%
\pgfpathlineto{\pgfqpoint{3.699243in}{3.101221in}}%
\pgfpathlineto{\pgfqpoint{3.701762in}{3.104774in}}%
\pgfpathlineto{\pgfqpoint{3.704281in}{3.084341in}}%
\pgfpathlineto{\pgfqpoint{3.706799in}{3.102648in}}%
\pgfpathlineto{\pgfqpoint{3.709318in}{3.074597in}}%
\pgfpathlineto{\pgfqpoint{3.711837in}{3.070763in}}%
\pgfpathlineto{\pgfqpoint{3.714356in}{3.078946in}}%
\pgfpathlineto{\pgfqpoint{3.716874in}{3.088890in}}%
\pgfpathlineto{\pgfqpoint{3.719393in}{3.080441in}}%
\pgfpathlineto{\pgfqpoint{3.721912in}{3.058195in}}%
\pgfpathlineto{\pgfqpoint{3.724431in}{3.058837in}}%
\pgfpathlineto{\pgfqpoint{3.726949in}{3.084601in}}%
\pgfpathlineto{\pgfqpoint{3.729468in}{3.069646in}}%
\pgfpathlineto{\pgfqpoint{3.731987in}{3.074697in}}%
\pgfpathlineto{\pgfqpoint{3.734506in}{3.069304in}}%
\pgfpathlineto{\pgfqpoint{3.737024in}{3.046384in}}%
\pgfpathlineto{\pgfqpoint{3.739543in}{3.028589in}}%
\pgfpathlineto{\pgfqpoint{3.742062in}{3.041620in}}%
\pgfpathlineto{\pgfqpoint{3.744581in}{3.089775in}}%
\pgfpathlineto{\pgfqpoint{3.747099in}{3.085961in}}%
\pgfpathlineto{\pgfqpoint{3.749618in}{3.093264in}}%
\pgfpathlineto{\pgfqpoint{3.752137in}{3.081237in}}%
\pgfpathlineto{\pgfqpoint{3.754656in}{3.067659in}}%
\pgfpathlineto{\pgfqpoint{3.757174in}{3.045513in}}%
\pgfpathlineto{\pgfqpoint{3.759693in}{3.038355in}}%
\pgfpathlineto{\pgfqpoint{3.762212in}{3.040936in}}%
\pgfpathlineto{\pgfqpoint{3.764731in}{2.996417in}}%
\pgfpathlineto{\pgfqpoint{3.767249in}{2.997635in}}%
\pgfpathlineto{\pgfqpoint{3.769768in}{2.944456in}}%
\pgfpathlineto{\pgfqpoint{3.772287in}{2.933014in}}%
\pgfpathlineto{\pgfqpoint{3.774806in}{2.908265in}}%
\pgfpathlineto{\pgfqpoint{3.777324in}{2.946620in}}%
\pgfpathlineto{\pgfqpoint{3.779843in}{2.940075in}}%
\pgfpathlineto{\pgfqpoint{3.782362in}{2.905455in}}%
\pgfpathlineto{\pgfqpoint{3.784881in}{2.898138in}}%
\pgfpathlineto{\pgfqpoint{3.787399in}{2.904742in}}%
\pgfpathlineto{\pgfqpoint{3.789918in}{2.900375in}}%
\pgfpathlineto{\pgfqpoint{3.792437in}{2.888176in}}%
\pgfpathlineto{\pgfqpoint{3.794956in}{2.859169in}}%
\pgfpathlineto{\pgfqpoint{3.797474in}{2.882759in}}%
\pgfpathlineto{\pgfqpoint{3.799993in}{2.859282in}}%
\pgfpathlineto{\pgfqpoint{3.802512in}{2.834490in}}%
\pgfpathlineto{\pgfqpoint{3.805031in}{2.831689in}}%
\pgfpathlineto{\pgfqpoint{3.807549in}{2.893500in}}%
\pgfpathlineto{\pgfqpoint{3.810068in}{2.891413in}}%
\pgfpathlineto{\pgfqpoint{3.812587in}{2.911162in}}%
\pgfpathlineto{\pgfqpoint{3.815106in}{2.886566in}}%
\pgfpathlineto{\pgfqpoint{3.817624in}{2.899103in}}%
\pgfpathlineto{\pgfqpoint{3.820143in}{2.933884in}}%
\pgfpathlineto{\pgfqpoint{3.822662in}{2.905251in}}%
\pgfpathlineto{\pgfqpoint{3.825181in}{2.899866in}}%
\pgfpathlineto{\pgfqpoint{3.827699in}{2.903635in}}%
\pgfpathlineto{\pgfqpoint{3.830218in}{2.879189in}}%
\pgfpathlineto{\pgfqpoint{3.832737in}{2.876487in}}%
\pgfpathlineto{\pgfqpoint{3.835256in}{2.828262in}}%
\pgfpathlineto{\pgfqpoint{3.837774in}{2.813985in}}%
\pgfpathlineto{\pgfqpoint{3.840293in}{2.829987in}}%
\pgfpathlineto{\pgfqpoint{3.842812in}{2.841163in}}%
\pgfpathlineto{\pgfqpoint{3.845331in}{2.818268in}}%
\pgfpathlineto{\pgfqpoint{3.847849in}{2.828331in}}%
\pgfpathlineto{\pgfqpoint{3.850368in}{2.831899in}}%
\pgfpathlineto{\pgfqpoint{3.852887in}{2.833408in}}%
\pgfpathlineto{\pgfqpoint{3.855406in}{2.838435in}}%
\pgfpathlineto{\pgfqpoint{3.857924in}{2.860926in}}%
\pgfpathlineto{\pgfqpoint{3.860443in}{2.843535in}}%
\pgfpathlineto{\pgfqpoint{3.862962in}{2.837651in}}%
\pgfpathlineto{\pgfqpoint{3.865481in}{2.815762in}}%
\pgfpathlineto{\pgfqpoint{3.867999in}{2.823954in}}%
\pgfpathlineto{\pgfqpoint{3.870518in}{2.811648in}}%
\pgfpathlineto{\pgfqpoint{3.873037in}{2.797693in}}%
\pgfpathlineto{\pgfqpoint{3.875556in}{2.794206in}}%
\pgfpathlineto{\pgfqpoint{3.878074in}{2.788792in}}%
\pgfpathlineto{\pgfqpoint{3.880593in}{2.778286in}}%
\pgfpathlineto{\pgfqpoint{3.883112in}{2.769041in}}%
\pgfpathlineto{\pgfqpoint{3.885631in}{2.784738in}}%
\pgfpathlineto{\pgfqpoint{3.888149in}{2.739574in}}%
\pgfpathlineto{\pgfqpoint{3.890668in}{2.731722in}}%
\pgfpathlineto{\pgfqpoint{3.893187in}{2.714221in}}%
\pgfpathlineto{\pgfqpoint{3.895706in}{2.748670in}}%
\pgfpathlineto{\pgfqpoint{3.898224in}{2.753014in}}%
\pgfpathlineto{\pgfqpoint{3.900743in}{2.774453in}}%
\pgfpathlineto{\pgfqpoint{3.903262in}{2.760576in}}%
\pgfpathlineto{\pgfqpoint{3.905781in}{2.776074in}}%
\pgfpathlineto{\pgfqpoint{3.908299in}{2.751692in}}%
\pgfpathlineto{\pgfqpoint{3.910818in}{2.770560in}}%
\pgfpathlineto{\pgfqpoint{3.913337in}{2.801366in}}%
\pgfpathlineto{\pgfqpoint{3.915856in}{2.788411in}}%
\pgfpathlineto{\pgfqpoint{3.918374in}{2.801235in}}%
\pgfpathlineto{\pgfqpoint{3.920893in}{2.820318in}}%
\pgfpathlineto{\pgfqpoint{3.923412in}{2.816855in}}%
\pgfpathlineto{\pgfqpoint{3.925931in}{2.793394in}}%
\pgfpathlineto{\pgfqpoint{3.928449in}{2.806133in}}%
\pgfpathlineto{\pgfqpoint{3.930968in}{2.838905in}}%
\pgfpathlineto{\pgfqpoint{3.933487in}{2.826684in}}%
\pgfpathlineto{\pgfqpoint{3.936006in}{2.813317in}}%
\pgfpathlineto{\pgfqpoint{3.938524in}{2.843774in}}%
\pgfpathlineto{\pgfqpoint{3.941043in}{2.849937in}}%
\pgfpathlineto{\pgfqpoint{3.943562in}{2.871041in}}%
\pgfpathlineto{\pgfqpoint{3.946081in}{2.826849in}}%
\pgfpathlineto{\pgfqpoint{3.948599in}{2.871343in}}%
\pgfpathlineto{\pgfqpoint{3.951118in}{2.889109in}}%
\pgfpathlineto{\pgfqpoint{3.953637in}{2.897453in}}%
\pgfpathlineto{\pgfqpoint{3.956156in}{2.898682in}}%
\pgfpathlineto{\pgfqpoint{3.958674in}{2.883889in}}%
\pgfpathlineto{\pgfqpoint{3.961193in}{2.923804in}}%
\pgfpathlineto{\pgfqpoint{3.963712in}{2.932478in}}%
\pgfpathlineto{\pgfqpoint{3.966231in}{2.945920in}}%
\pgfpathlineto{\pgfqpoint{3.968749in}{2.936848in}}%
\pgfpathlineto{\pgfqpoint{3.971268in}{2.975464in}}%
\pgfpathlineto{\pgfqpoint{3.973787in}{2.958301in}}%
\pgfpathlineto{\pgfqpoint{3.976306in}{2.933080in}}%
\pgfpathlineto{\pgfqpoint{3.978824in}{2.950545in}}%
\pgfpathlineto{\pgfqpoint{3.981343in}{2.944815in}}%
\pgfpathlineto{\pgfqpoint{3.983862in}{2.944437in}}%
\pgfpathlineto{\pgfqpoint{3.986381in}{2.925042in}}%
\pgfpathlineto{\pgfqpoint{3.988899in}{2.931413in}}%
\pgfpathlineto{\pgfqpoint{3.991418in}{2.909615in}}%
\pgfpathlineto{\pgfqpoint{3.993937in}{2.904585in}}%
\pgfpathlineto{\pgfqpoint{3.996456in}{2.918569in}}%
\pgfpathlineto{\pgfqpoint{3.998974in}{2.936192in}}%
\pgfpathlineto{\pgfqpoint{4.001493in}{2.944982in}}%
\pgfpathlineto{\pgfqpoint{4.004012in}{2.955296in}}%
\pgfpathlineto{\pgfqpoint{4.006531in}{2.985531in}}%
\pgfpathlineto{\pgfqpoint{4.009049in}{3.011361in}}%
\pgfpathlineto{\pgfqpoint{4.011568in}{2.991966in}}%
\pgfpathlineto{\pgfqpoint{4.014087in}{3.036102in}}%
\pgfpathlineto{\pgfqpoint{4.016606in}{3.027855in}}%
\pgfpathlineto{\pgfqpoint{4.019124in}{3.095088in}}%
\pgfpathlineto{\pgfqpoint{4.021643in}{3.108728in}}%
\pgfpathlineto{\pgfqpoint{4.024162in}{3.091524in}}%
\pgfpathlineto{\pgfqpoint{4.026681in}{3.076303in}}%
\pgfpathlineto{\pgfqpoint{4.029199in}{3.068114in}}%
\pgfpathlineto{\pgfqpoint{4.031718in}{3.057772in}}%
\pgfpathlineto{\pgfqpoint{4.034237in}{3.054266in}}%
\pgfpathlineto{\pgfqpoint{4.036756in}{3.022311in}}%
\pgfpathlineto{\pgfqpoint{4.039274in}{2.983816in}}%
\pgfpathlineto{\pgfqpoint{4.041793in}{3.020378in}}%
\pgfpathlineto{\pgfqpoint{4.044312in}{3.035566in}}%
\pgfpathlineto{\pgfqpoint{4.046831in}{3.016481in}}%
\pgfpathlineto{\pgfqpoint{4.049349in}{3.038884in}}%
\pgfpathlineto{\pgfqpoint{4.051868in}{3.058860in}}%
\pgfpathlineto{\pgfqpoint{4.054387in}{3.041841in}}%
\pgfpathlineto{\pgfqpoint{4.056906in}{3.055747in}}%
\pgfpathlineto{\pgfqpoint{4.059424in}{3.094286in}}%
\pgfpathlineto{\pgfqpoint{4.061943in}{3.119090in}}%
\pgfpathlineto{\pgfqpoint{4.064462in}{3.137706in}}%
\pgfpathlineto{\pgfqpoint{4.066981in}{3.122176in}}%
\pgfpathlineto{\pgfqpoint{4.069499in}{3.143378in}}%
\pgfpathlineto{\pgfqpoint{4.072018in}{3.138282in}}%
\pgfpathlineto{\pgfqpoint{4.074537in}{3.151095in}}%
\pgfpathlineto{\pgfqpoint{4.077056in}{3.126172in}}%
\pgfpathlineto{\pgfqpoint{4.079574in}{3.164605in}}%
\pgfpathlineto{\pgfqpoint{4.082093in}{3.140416in}}%
\pgfpathlineto{\pgfqpoint{4.084612in}{3.130810in}}%
\pgfpathlineto{\pgfqpoint{4.087131in}{3.099984in}}%
\pgfpathlineto{\pgfqpoint{4.089649in}{3.094328in}}%
\pgfpathlineto{\pgfqpoint{4.092168in}{3.066623in}}%
\pgfpathlineto{\pgfqpoint{4.094687in}{3.091412in}}%
\pgfpathlineto{\pgfqpoint{4.097206in}{3.089563in}}%
\pgfpathlineto{\pgfqpoint{4.099724in}{3.084293in}}%
\pgfpathlineto{\pgfqpoint{4.102243in}{3.083852in}}%
\pgfpathlineto{\pgfqpoint{4.104762in}{3.067905in}}%
\pgfpathlineto{\pgfqpoint{4.107281in}{3.069727in}}%
\pgfpathlineto{\pgfqpoint{4.109799in}{3.044016in}}%
\pgfpathlineto{\pgfqpoint{4.112318in}{3.001603in}}%
\pgfpathlineto{\pgfqpoint{4.114837in}{3.028034in}}%
\pgfpathlineto{\pgfqpoint{4.117356in}{2.994683in}}%
\pgfpathlineto{\pgfqpoint{4.119874in}{3.017117in}}%
\pgfpathlineto{\pgfqpoint{4.122393in}{3.017948in}}%
\pgfpathlineto{\pgfqpoint{4.124912in}{3.015620in}}%
\pgfpathlineto{\pgfqpoint{4.127431in}{3.014169in}}%
\pgfpathlineto{\pgfqpoint{4.129949in}{3.002293in}}%
\pgfpathlineto{\pgfqpoint{4.132468in}{2.995250in}}%
\pgfpathlineto{\pgfqpoint{4.134987in}{3.033449in}}%
\pgfpathlineto{\pgfqpoint{4.137506in}{3.033626in}}%
\pgfpathlineto{\pgfqpoint{4.140024in}{3.039043in}}%
\pgfpathlineto{\pgfqpoint{4.142543in}{3.028993in}}%
\pgfpathlineto{\pgfqpoint{4.145062in}{3.030280in}}%
\pgfpathlineto{\pgfqpoint{4.147581in}{3.045992in}}%
\pgfpathlineto{\pgfqpoint{4.150099in}{3.034342in}}%
\pgfpathlineto{\pgfqpoint{4.152618in}{3.023551in}}%
\pgfpathlineto{\pgfqpoint{4.155137in}{3.026987in}}%
\pgfpathlineto{\pgfqpoint{4.157656in}{3.073624in}}%
\pgfpathlineto{\pgfqpoint{4.160174in}{3.060819in}}%
\pgfpathlineto{\pgfqpoint{4.162693in}{3.066811in}}%
\pgfpathlineto{\pgfqpoint{4.165212in}{3.055590in}}%
\pgfpathlineto{\pgfqpoint{4.167731in}{3.037472in}}%
\pgfpathlineto{\pgfqpoint{4.170249in}{3.059400in}}%
\pgfpathlineto{\pgfqpoint{4.172768in}{3.031924in}}%
\pgfpathlineto{\pgfqpoint{4.175287in}{2.990151in}}%
\pgfpathlineto{\pgfqpoint{4.177806in}{3.007904in}}%
\pgfpathlineto{\pgfqpoint{4.180324in}{3.028867in}}%
\pgfpathlineto{\pgfqpoint{4.182843in}{3.074279in}}%
\pgfpathlineto{\pgfqpoint{4.185362in}{3.082387in}}%
\pgfpathlineto{\pgfqpoint{4.187881in}{3.094522in}}%
\pgfpathlineto{\pgfqpoint{4.190399in}{3.096996in}}%
\pgfpathlineto{\pgfqpoint{4.192918in}{3.076929in}}%
\pgfpathlineto{\pgfqpoint{4.195437in}{3.037420in}}%
\pgfpathlineto{\pgfqpoint{4.197956in}{3.037492in}}%
\pgfpathlineto{\pgfqpoint{4.200474in}{3.064817in}}%
\pgfpathlineto{\pgfqpoint{4.202993in}{3.067177in}}%
\pgfpathlineto{\pgfqpoint{4.205512in}{3.074542in}}%
\pgfpathlineto{\pgfqpoint{4.208031in}{3.059407in}}%
\pgfpathlineto{\pgfqpoint{4.210549in}{3.090734in}}%
\pgfpathlineto{\pgfqpoint{4.215587in}{3.080774in}}%
\pgfpathlineto{\pgfqpoint{4.218106in}{3.071707in}}%
\pgfpathlineto{\pgfqpoint{4.220624in}{3.066513in}}%
\pgfpathlineto{\pgfqpoint{4.223143in}{3.085064in}}%
\pgfpathlineto{\pgfqpoint{4.225662in}{3.112244in}}%
\pgfpathlineto{\pgfqpoint{4.228181in}{3.101219in}}%
\pgfpathlineto{\pgfqpoint{4.230699in}{3.098594in}}%
\pgfpathlineto{\pgfqpoint{4.233218in}{3.086806in}}%
\pgfpathlineto{\pgfqpoint{4.235737in}{3.057924in}}%
\pgfpathlineto{\pgfqpoint{4.238256in}{3.023032in}}%
\pgfpathlineto{\pgfqpoint{4.240774in}{3.020320in}}%
\pgfpathlineto{\pgfqpoint{4.243293in}{3.029984in}}%
\pgfpathlineto{\pgfqpoint{4.245812in}{3.002629in}}%
\pgfpathlineto{\pgfqpoint{4.248331in}{3.011062in}}%
\pgfpathlineto{\pgfqpoint{4.250849in}{3.006199in}}%
\pgfpathlineto{\pgfqpoint{4.253368in}{3.008614in}}%
\pgfpathlineto{\pgfqpoint{4.258406in}{3.052732in}}%
\pgfpathlineto{\pgfqpoint{4.260924in}{3.065198in}}%
\pgfpathlineto{\pgfqpoint{4.263443in}{3.049085in}}%
\pgfpathlineto{\pgfqpoint{4.265962in}{3.015137in}}%
\pgfpathlineto{\pgfqpoint{4.268481in}{2.987592in}}%
\pgfpathlineto{\pgfqpoint{4.270999in}{2.971401in}}%
\pgfpathlineto{\pgfqpoint{4.273518in}{2.948469in}}%
\pgfpathlineto{\pgfqpoint{4.276037in}{2.926817in}}%
\pgfpathlineto{\pgfqpoint{4.278556in}{2.912577in}}%
\pgfpathlineto{\pgfqpoint{4.281074in}{2.883768in}}%
\pgfpathlineto{\pgfqpoint{4.283593in}{2.889889in}}%
\pgfpathlineto{\pgfqpoint{4.286112in}{2.890114in}}%
\pgfpathlineto{\pgfqpoint{4.288631in}{2.915978in}}%
\pgfpathlineto{\pgfqpoint{4.291149in}{2.930339in}}%
\pgfpathlineto{\pgfqpoint{4.293668in}{2.940045in}}%
\pgfpathlineto{\pgfqpoint{4.296187in}{2.979744in}}%
\pgfpathlineto{\pgfqpoint{4.298706in}{2.947472in}}%
\pgfpathlineto{\pgfqpoint{4.301224in}{2.976268in}}%
\pgfpathlineto{\pgfqpoint{4.306262in}{2.992378in}}%
\pgfpathlineto{\pgfqpoint{4.308781in}{3.021964in}}%
\pgfpathlineto{\pgfqpoint{4.311299in}{3.031600in}}%
\pgfpathlineto{\pgfqpoint{4.313818in}{3.031577in}}%
\pgfpathlineto{\pgfqpoint{4.316337in}{3.041068in}}%
\pgfpathlineto{\pgfqpoint{4.318856in}{3.070237in}}%
\pgfpathlineto{\pgfqpoint{4.321374in}{3.042047in}}%
\pgfpathlineto{\pgfqpoint{4.323893in}{3.012061in}}%
\pgfpathlineto{\pgfqpoint{4.326412in}{3.005824in}}%
\pgfpathlineto{\pgfqpoint{4.328931in}{3.018203in}}%
\pgfpathlineto{\pgfqpoint{4.331449in}{3.035262in}}%
\pgfpathlineto{\pgfqpoint{4.333968in}{2.989135in}}%
\pgfpathlineto{\pgfqpoint{4.339006in}{2.996696in}}%
\pgfpathlineto{\pgfqpoint{4.341524in}{2.992029in}}%
\pgfpathlineto{\pgfqpoint{4.344043in}{2.983641in}}%
\pgfpathlineto{\pgfqpoint{4.346562in}{2.948315in}}%
\pgfpathlineto{\pgfqpoint{4.349081in}{2.910728in}}%
\pgfpathlineto{\pgfqpoint{4.351599in}{2.930428in}}%
\pgfpathlineto{\pgfqpoint{4.354118in}{2.954872in}}%
\pgfpathlineto{\pgfqpoint{4.356637in}{2.924910in}}%
\pgfpathlineto{\pgfqpoint{4.359156in}{2.916063in}}%
\pgfpathlineto{\pgfqpoint{4.361674in}{2.942530in}}%
\pgfpathlineto{\pgfqpoint{4.364193in}{2.945249in}}%
\pgfpathlineto{\pgfqpoint{4.366712in}{2.960292in}}%
\pgfpathlineto{\pgfqpoint{4.369231in}{2.910602in}}%
\pgfpathlineto{\pgfqpoint{4.371749in}{2.928894in}}%
\pgfpathlineto{\pgfqpoint{4.374268in}{2.909881in}}%
\pgfpathlineto{\pgfqpoint{4.376787in}{2.876727in}}%
\pgfpathlineto{\pgfqpoint{4.379306in}{2.917151in}}%
\pgfpathlineto{\pgfqpoint{4.381824in}{2.939404in}}%
\pgfpathlineto{\pgfqpoint{4.384343in}{2.920501in}}%
\pgfpathlineto{\pgfqpoint{4.386862in}{2.905881in}}%
\pgfpathlineto{\pgfqpoint{4.389381in}{2.894027in}}%
\pgfpathlineto{\pgfqpoint{4.391899in}{2.890435in}}%
\pgfpathlineto{\pgfqpoint{4.394418in}{2.857905in}}%
\pgfpathlineto{\pgfqpoint{4.396937in}{2.817092in}}%
\pgfpathlineto{\pgfqpoint{4.399456in}{2.845141in}}%
\pgfpathlineto{\pgfqpoint{4.401974in}{2.856843in}}%
\pgfpathlineto{\pgfqpoint{4.404493in}{2.859780in}}%
\pgfpathlineto{\pgfqpoint{4.407012in}{2.913447in}}%
\pgfpathlineto{\pgfqpoint{4.409531in}{2.922461in}}%
\pgfpathlineto{\pgfqpoint{4.412049in}{2.904847in}}%
\pgfpathlineto{\pgfqpoint{4.414568in}{2.873335in}}%
\pgfpathlineto{\pgfqpoint{4.417087in}{2.880881in}}%
\pgfpathlineto{\pgfqpoint{4.419606in}{2.855795in}}%
\pgfpathlineto{\pgfqpoint{4.422124in}{2.806698in}}%
\pgfpathlineto{\pgfqpoint{4.424643in}{2.774678in}}%
\pgfpathlineto{\pgfqpoint{4.427162in}{2.781053in}}%
\pgfpathlineto{\pgfqpoint{4.429681in}{2.805611in}}%
\pgfpathlineto{\pgfqpoint{4.432199in}{2.793452in}}%
\pgfpathlineto{\pgfqpoint{4.434718in}{2.802315in}}%
\pgfpathlineto{\pgfqpoint{4.437237in}{2.821543in}}%
\pgfpathlineto{\pgfqpoint{4.439756in}{2.809635in}}%
\pgfpathlineto{\pgfqpoint{4.442274in}{2.809965in}}%
\pgfpathlineto{\pgfqpoint{4.444793in}{2.840169in}}%
\pgfpathlineto{\pgfqpoint{4.447312in}{2.848929in}}%
\pgfpathlineto{\pgfqpoint{4.449831in}{2.897381in}}%
\pgfpathlineto{\pgfqpoint{4.452349in}{2.906231in}}%
\pgfpathlineto{\pgfqpoint{4.454868in}{2.908161in}}%
\pgfpathlineto{\pgfqpoint{4.457387in}{2.854123in}}%
\pgfpathlineto{\pgfqpoint{4.459906in}{2.835653in}}%
\pgfpathlineto{\pgfqpoint{4.462424in}{2.836324in}}%
\pgfpathlineto{\pgfqpoint{4.464943in}{2.823312in}}%
\pgfpathlineto{\pgfqpoint{4.467462in}{2.827603in}}%
\pgfpathlineto{\pgfqpoint{4.469981in}{2.792911in}}%
\pgfpathlineto{\pgfqpoint{4.472499in}{2.827873in}}%
\pgfpathlineto{\pgfqpoint{4.475018in}{2.817541in}}%
\pgfpathlineto{\pgfqpoint{4.477537in}{2.817742in}}%
\pgfpathlineto{\pgfqpoint{4.480056in}{2.826359in}}%
\pgfpathlineto{\pgfqpoint{4.482574in}{2.821280in}}%
\pgfpathlineto{\pgfqpoint{4.485093in}{2.798832in}}%
\pgfpathlineto{\pgfqpoint{4.487612in}{2.778691in}}%
\pgfpathlineto{\pgfqpoint{4.490131in}{2.757124in}}%
\pgfpathlineto{\pgfqpoint{4.492649in}{2.758003in}}%
\pgfpathlineto{\pgfqpoint{4.495168in}{2.782140in}}%
\pgfpathlineto{\pgfqpoint{4.497687in}{2.759700in}}%
\pgfpathlineto{\pgfqpoint{4.500206in}{2.767996in}}%
\pgfpathlineto{\pgfqpoint{4.502724in}{2.767329in}}%
\pgfpathlineto{\pgfqpoint{4.505243in}{2.770643in}}%
\pgfpathlineto{\pgfqpoint{4.507762in}{2.763143in}}%
\pgfpathlineto{\pgfqpoint{4.510281in}{2.763872in}}%
\pgfpathlineto{\pgfqpoint{4.512799in}{2.750899in}}%
\pgfpathlineto{\pgfqpoint{4.515318in}{2.775760in}}%
\pgfpathlineto{\pgfqpoint{4.517837in}{2.793028in}}%
\pgfpathlineto{\pgfqpoint{4.520356in}{2.820743in}}%
\pgfpathlineto{\pgfqpoint{4.522874in}{2.823272in}}%
\pgfpathlineto{\pgfqpoint{4.525393in}{2.863820in}}%
\pgfpathlineto{\pgfqpoint{4.527912in}{2.854110in}}%
\pgfpathlineto{\pgfqpoint{4.530431in}{2.859144in}}%
\pgfpathlineto{\pgfqpoint{4.532949in}{2.841721in}}%
\pgfpathlineto{\pgfqpoint{4.535468in}{2.829640in}}%
\pgfpathlineto{\pgfqpoint{4.537987in}{2.812775in}}%
\pgfpathlineto{\pgfqpoint{4.540506in}{2.815599in}}%
\pgfpathlineto{\pgfqpoint{4.543024in}{2.793760in}}%
\pgfpathlineto{\pgfqpoint{4.545543in}{2.763837in}}%
\pgfpathlineto{\pgfqpoint{4.548062in}{2.772790in}}%
\pgfpathlineto{\pgfqpoint{4.550581in}{2.740067in}}%
\pgfpathlineto{\pgfqpoint{4.553099in}{2.760514in}}%
\pgfpathlineto{\pgfqpoint{4.555618in}{2.794670in}}%
\pgfpathlineto{\pgfqpoint{4.558137in}{2.777448in}}%
\pgfpathlineto{\pgfqpoint{4.560656in}{2.802316in}}%
\pgfpathlineto{\pgfqpoint{4.563174in}{2.787088in}}%
\pgfpathlineto{\pgfqpoint{4.565693in}{2.787123in}}%
\pgfpathlineto{\pgfqpoint{4.568212in}{2.808488in}}%
\pgfpathlineto{\pgfqpoint{4.570731in}{2.805303in}}%
\pgfpathlineto{\pgfqpoint{4.573249in}{2.805985in}}%
\pgfpathlineto{\pgfqpoint{4.575768in}{2.769076in}}%
\pgfpathlineto{\pgfqpoint{4.578287in}{2.784973in}}%
\pgfpathlineto{\pgfqpoint{4.580806in}{2.789829in}}%
\pgfpathlineto{\pgfqpoint{4.583324in}{2.796216in}}%
\pgfpathlineto{\pgfqpoint{4.585843in}{2.809151in}}%
\pgfpathlineto{\pgfqpoint{4.588362in}{2.782682in}}%
\pgfpathlineto{\pgfqpoint{4.590881in}{2.782883in}}%
\pgfpathlineto{\pgfqpoint{4.593399in}{2.819989in}}%
\pgfpathlineto{\pgfqpoint{4.595918in}{2.834269in}}%
\pgfpathlineto{\pgfqpoint{4.598437in}{2.804481in}}%
\pgfpathlineto{\pgfqpoint{4.600956in}{2.800114in}}%
\pgfpathlineto{\pgfqpoint{4.603474in}{2.757518in}}%
\pgfpathlineto{\pgfqpoint{4.605993in}{2.772634in}}%
\pgfpathlineto{\pgfqpoint{4.608512in}{2.784962in}}%
\pgfpathlineto{\pgfqpoint{4.611031in}{2.771792in}}%
\pgfpathlineto{\pgfqpoint{4.613549in}{2.744358in}}%
\pgfpathlineto{\pgfqpoint{4.616068in}{2.750041in}}%
\pgfpathlineto{\pgfqpoint{4.618587in}{2.757540in}}%
\pgfpathlineto{\pgfqpoint{4.621106in}{2.808161in}}%
\pgfpathlineto{\pgfqpoint{4.623624in}{2.817364in}}%
\pgfpathlineto{\pgfqpoint{4.626143in}{2.819112in}}%
\pgfpathlineto{\pgfqpoint{4.628662in}{2.830490in}}%
\pgfpathlineto{\pgfqpoint{4.631181in}{2.850864in}}%
\pgfpathlineto{\pgfqpoint{4.633699in}{2.832130in}}%
\pgfpathlineto{\pgfqpoint{4.636218in}{2.779561in}}%
\pgfpathlineto{\pgfqpoint{4.638737in}{2.762282in}}%
\pgfpathlineto{\pgfqpoint{4.641256in}{2.783701in}}%
\pgfpathlineto{\pgfqpoint{4.643774in}{2.745233in}}%
\pgfpathlineto{\pgfqpoint{4.646293in}{2.746152in}}%
\pgfpathlineto{\pgfqpoint{4.648812in}{2.749475in}}%
\pgfpathlineto{\pgfqpoint{4.651331in}{2.725093in}}%
\pgfpathlineto{\pgfqpoint{4.653849in}{2.694315in}}%
\pgfpathlineto{\pgfqpoint{4.656368in}{2.660690in}}%
\pgfpathlineto{\pgfqpoint{4.658887in}{2.659182in}}%
\pgfpathlineto{\pgfqpoint{4.661406in}{2.679937in}}%
\pgfpathlineto{\pgfqpoint{4.663924in}{2.666760in}}%
\pgfpathlineto{\pgfqpoint{4.666443in}{2.691929in}}%
\pgfpathlineto{\pgfqpoint{4.668962in}{2.676203in}}%
\pgfpathlineto{\pgfqpoint{4.671481in}{2.729889in}}%
\pgfpathlineto{\pgfqpoint{4.673999in}{2.751779in}}%
\pgfpathlineto{\pgfqpoint{4.676518in}{2.770053in}}%
\pgfpathlineto{\pgfqpoint{4.679037in}{2.738268in}}%
\pgfpathlineto{\pgfqpoint{4.681556in}{2.709669in}}%
\pgfpathlineto{\pgfqpoint{4.684074in}{2.754697in}}%
\pgfpathlineto{\pgfqpoint{4.686593in}{2.711494in}}%
\pgfpathlineto{\pgfqpoint{4.689112in}{2.695089in}}%
\pgfpathlineto{\pgfqpoint{4.691631in}{2.669735in}}%
\pgfpathlineto{\pgfqpoint{4.694149in}{2.651787in}}%
\pgfpathlineto{\pgfqpoint{4.696668in}{2.643233in}}%
\pgfpathlineto{\pgfqpoint{4.699187in}{2.665330in}}%
\pgfpathlineto{\pgfqpoint{4.701706in}{2.664978in}}%
\pgfpathlineto{\pgfqpoint{4.704224in}{2.657640in}}%
\pgfpathlineto{\pgfqpoint{4.706743in}{2.625414in}}%
\pgfpathlineto{\pgfqpoint{4.709262in}{2.616274in}}%
\pgfpathlineto{\pgfqpoint{4.711781in}{2.602619in}}%
\pgfpathlineto{\pgfqpoint{4.714299in}{2.647164in}}%
\pgfpathlineto{\pgfqpoint{4.716818in}{2.611069in}}%
\pgfpathlineto{\pgfqpoint{4.719337in}{2.640729in}}%
\pgfpathlineto{\pgfqpoint{4.721856in}{2.667505in}}%
\pgfpathlineto{\pgfqpoint{4.724374in}{2.648364in}}%
\pgfpathlineto{\pgfqpoint{4.726893in}{2.682383in}}%
\pgfpathlineto{\pgfqpoint{4.731931in}{2.704625in}}%
\pgfpathlineto{\pgfqpoint{4.734449in}{2.697704in}}%
\pgfpathlineto{\pgfqpoint{4.736968in}{2.673239in}}%
\pgfpathlineto{\pgfqpoint{4.739487in}{2.698920in}}%
\pgfpathlineto{\pgfqpoint{4.742006in}{2.714113in}}%
\pgfpathlineto{\pgfqpoint{4.744524in}{2.730294in}}%
\pgfpathlineto{\pgfqpoint{4.747043in}{2.703426in}}%
\pgfpathlineto{\pgfqpoint{4.749562in}{2.715229in}}%
\pgfpathlineto{\pgfqpoint{4.752081in}{2.715874in}}%
\pgfpathlineto{\pgfqpoint{4.754599in}{2.712684in}}%
\pgfpathlineto{\pgfqpoint{4.757118in}{2.694207in}}%
\pgfpathlineto{\pgfqpoint{4.759637in}{2.668573in}}%
\pgfpathlineto{\pgfqpoint{4.762156in}{2.688905in}}%
\pgfpathlineto{\pgfqpoint{4.764674in}{2.701987in}}%
\pgfpathlineto{\pgfqpoint{4.767193in}{2.673196in}}%
\pgfpathlineto{\pgfqpoint{4.769712in}{2.693185in}}%
\pgfpathlineto{\pgfqpoint{4.772231in}{2.686560in}}%
\pgfpathlineto{\pgfqpoint{4.774749in}{2.695445in}}%
\pgfpathlineto{\pgfqpoint{4.777268in}{2.691121in}}%
\pgfpathlineto{\pgfqpoint{4.779787in}{2.696740in}}%
\pgfpathlineto{\pgfqpoint{4.782306in}{2.679888in}}%
\pgfpathlineto{\pgfqpoint{4.784824in}{2.711095in}}%
\pgfpathlineto{\pgfqpoint{4.787343in}{2.733292in}}%
\pgfpathlineto{\pgfqpoint{4.789862in}{2.718065in}}%
\pgfpathlineto{\pgfqpoint{4.792381in}{2.697669in}}%
\pgfpathlineto{\pgfqpoint{4.794899in}{2.724032in}}%
\pgfpathlineto{\pgfqpoint{4.797418in}{2.742970in}}%
\pgfpathlineto{\pgfqpoint{4.799937in}{2.747409in}}%
\pgfpathlineto{\pgfqpoint{4.802456in}{2.728048in}}%
\pgfpathlineto{\pgfqpoint{4.804974in}{2.730640in}}%
\pgfpathlineto{\pgfqpoint{4.807493in}{2.774512in}}%
\pgfpathlineto{\pgfqpoint{4.810012in}{2.788647in}}%
\pgfpathlineto{\pgfqpoint{4.812531in}{2.757657in}}%
\pgfpathlineto{\pgfqpoint{4.815049in}{2.746268in}}%
\pgfpathlineto{\pgfqpoint{4.817568in}{2.750419in}}%
\pgfpathlineto{\pgfqpoint{4.820087in}{2.740224in}}%
\pgfpathlineto{\pgfqpoint{4.822606in}{2.742825in}}%
\pgfpathlineto{\pgfqpoint{4.825124in}{2.746739in}}%
\pgfpathlineto{\pgfqpoint{4.827643in}{2.758012in}}%
\pgfpathlineto{\pgfqpoint{4.830162in}{2.778297in}}%
\pgfpathlineto{\pgfqpoint{4.832681in}{2.784375in}}%
\pgfpathlineto{\pgfqpoint{4.835199in}{2.765837in}}%
\pgfpathlineto{\pgfqpoint{4.837718in}{2.772141in}}%
\pgfpathlineto{\pgfqpoint{4.840237in}{2.790855in}}%
\pgfpathlineto{\pgfqpoint{4.842756in}{2.768328in}}%
\pgfpathlineto{\pgfqpoint{4.845274in}{2.771007in}}%
\pgfpathlineto{\pgfqpoint{4.847793in}{2.791340in}}%
\pgfpathlineto{\pgfqpoint{4.850312in}{2.819524in}}%
\pgfpathlineto{\pgfqpoint{4.852831in}{2.820240in}}%
\pgfpathlineto{\pgfqpoint{4.855349in}{2.814023in}}%
\pgfpathlineto{\pgfqpoint{4.857868in}{2.819057in}}%
\pgfpathlineto{\pgfqpoint{4.860387in}{2.822824in}}%
\pgfpathlineto{\pgfqpoint{4.862906in}{2.831977in}}%
\pgfpathlineto{\pgfqpoint{4.867943in}{2.789414in}}%
\pgfpathlineto{\pgfqpoint{4.870462in}{2.814852in}}%
\pgfpathlineto{\pgfqpoint{4.872981in}{2.859489in}}%
\pgfpathlineto{\pgfqpoint{4.875499in}{2.852855in}}%
\pgfpathlineto{\pgfqpoint{4.878018in}{2.856037in}}%
\pgfpathlineto{\pgfqpoint{4.880537in}{2.847773in}}%
\pgfpathlineto{\pgfqpoint{4.883056in}{2.829323in}}%
\pgfpathlineto{\pgfqpoint{4.885574in}{2.790843in}}%
\pgfpathlineto{\pgfqpoint{4.888093in}{2.779910in}}%
\pgfpathlineto{\pgfqpoint{4.890612in}{2.769576in}}%
\pgfpathlineto{\pgfqpoint{4.893131in}{2.732950in}}%
\pgfpathlineto{\pgfqpoint{4.895649in}{2.751159in}}%
\pgfpathlineto{\pgfqpoint{4.898168in}{2.778174in}}%
\pgfpathlineto{\pgfqpoint{4.900687in}{2.740668in}}%
\pgfpathlineto{\pgfqpoint{4.903206in}{2.731027in}}%
\pgfpathlineto{\pgfqpoint{4.905724in}{2.742484in}}%
\pgfpathlineto{\pgfqpoint{4.908243in}{2.720094in}}%
\pgfpathlineto{\pgfqpoint{4.910762in}{2.682273in}}%
\pgfpathlineto{\pgfqpoint{4.913281in}{2.704593in}}%
\pgfpathlineto{\pgfqpoint{4.915799in}{2.676249in}}%
\pgfpathlineto{\pgfqpoint{4.918318in}{2.714999in}}%
\pgfpathlineto{\pgfqpoint{4.920837in}{2.737697in}}%
\pgfpathlineto{\pgfqpoint{4.923356in}{2.786109in}}%
\pgfpathlineto{\pgfqpoint{4.925874in}{2.813238in}}%
\pgfpathlineto{\pgfqpoint{4.928393in}{2.804859in}}%
\pgfpathlineto{\pgfqpoint{4.930912in}{2.800770in}}%
\pgfpathlineto{\pgfqpoint{4.933431in}{2.790934in}}%
\pgfpathlineto{\pgfqpoint{4.935949in}{2.815904in}}%
\pgfpathlineto{\pgfqpoint{4.938468in}{2.874556in}}%
\pgfpathlineto{\pgfqpoint{4.940987in}{2.876793in}}%
\pgfpathlineto{\pgfqpoint{4.943506in}{2.903815in}}%
\pgfpathlineto{\pgfqpoint{4.946024in}{2.953061in}}%
\pgfpathlineto{\pgfqpoint{4.948543in}{2.968803in}}%
\pgfpathlineto{\pgfqpoint{4.951062in}{2.938180in}}%
\pgfpathlineto{\pgfqpoint{4.953581in}{2.966238in}}%
\pgfpathlineto{\pgfqpoint{4.956099in}{2.992155in}}%
\pgfpathlineto{\pgfqpoint{4.958618in}{2.959674in}}%
\pgfpathlineto{\pgfqpoint{4.961137in}{2.950837in}}%
\pgfpathlineto{\pgfqpoint{4.963656in}{2.942648in}}%
\pgfpathlineto{\pgfqpoint{4.966174in}{2.939232in}}%
\pgfpathlineto{\pgfqpoint{4.968693in}{2.943967in}}%
\pgfpathlineto{\pgfqpoint{4.971212in}{2.912826in}}%
\pgfpathlineto{\pgfqpoint{4.973731in}{2.901339in}}%
\pgfpathlineto{\pgfqpoint{4.976249in}{2.938299in}}%
\pgfpathlineto{\pgfqpoint{4.978768in}{2.939675in}}%
\pgfpathlineto{\pgfqpoint{4.981287in}{2.925262in}}%
\pgfpathlineto{\pgfqpoint{4.983806in}{2.974763in}}%
\pgfpathlineto{\pgfqpoint{4.986324in}{2.993220in}}%
\pgfpathlineto{\pgfqpoint{4.988843in}{2.989345in}}%
\pgfpathlineto{\pgfqpoint{4.991362in}{2.988771in}}%
\pgfpathlineto{\pgfqpoint{4.993881in}{3.000582in}}%
\pgfpathlineto{\pgfqpoint{4.996399in}{3.021774in}}%
\pgfpathlineto{\pgfqpoint{4.998918in}{3.030112in}}%
\pgfpathlineto{\pgfqpoint{5.001437in}{3.010365in}}%
\pgfpathlineto{\pgfqpoint{5.003956in}{3.055263in}}%
\pgfpathlineto{\pgfqpoint{5.006474in}{3.053850in}}%
\pgfpathlineto{\pgfqpoint{5.008993in}{3.049185in}}%
\pgfpathlineto{\pgfqpoint{5.011512in}{3.047603in}}%
\pgfpathlineto{\pgfqpoint{5.014031in}{3.034211in}}%
\pgfpathlineto{\pgfqpoint{5.016549in}{2.988864in}}%
\pgfpathlineto{\pgfqpoint{5.019068in}{3.008539in}}%
\pgfpathlineto{\pgfqpoint{5.021587in}{3.037289in}}%
\pgfpathlineto{\pgfqpoint{5.024106in}{3.062397in}}%
\pgfpathlineto{\pgfqpoint{5.026624in}{3.056785in}}%
\pgfpathlineto{\pgfqpoint{5.029143in}{3.061908in}}%
\pgfpathlineto{\pgfqpoint{5.031662in}{3.040998in}}%
\pgfpathlineto{\pgfqpoint{5.034181in}{3.072450in}}%
\pgfpathlineto{\pgfqpoint{5.036699in}{3.082563in}}%
\pgfpathlineto{\pgfqpoint{5.039218in}{3.079282in}}%
\pgfpathlineto{\pgfqpoint{5.041737in}{3.022949in}}%
\pgfpathlineto{\pgfqpoint{5.044256in}{2.956675in}}%
\pgfpathlineto{\pgfqpoint{5.046774in}{2.936988in}}%
\pgfpathlineto{\pgfqpoint{5.049293in}{2.927650in}}%
\pgfpathlineto{\pgfqpoint{5.051812in}{2.882947in}}%
\pgfpathlineto{\pgfqpoint{5.054331in}{2.908892in}}%
\pgfpathlineto{\pgfqpoint{5.056849in}{2.944528in}}%
\pgfpathlineto{\pgfqpoint{5.059368in}{2.886590in}}%
\pgfpathlineto{\pgfqpoint{5.061887in}{2.899576in}}%
\pgfpathlineto{\pgfqpoint{5.064406in}{2.897173in}}%
\pgfpathlineto{\pgfqpoint{5.066924in}{2.897157in}}%
\pgfpathlineto{\pgfqpoint{5.069443in}{2.913392in}}%
\pgfpathlineto{\pgfqpoint{5.071962in}{2.902589in}}%
\pgfpathlineto{\pgfqpoint{5.074481in}{2.892515in}}%
\pgfpathlineto{\pgfqpoint{5.076999in}{2.886996in}}%
\pgfpathlineto{\pgfqpoint{5.079518in}{2.863807in}}%
\pgfpathlineto{\pgfqpoint{5.082037in}{2.853356in}}%
\pgfpathlineto{\pgfqpoint{5.084556in}{2.835992in}}%
\pgfpathlineto{\pgfqpoint{5.087074in}{2.820497in}}%
\pgfpathlineto{\pgfqpoint{5.089593in}{2.824320in}}%
\pgfpathlineto{\pgfqpoint{5.092112in}{2.809516in}}%
\pgfpathlineto{\pgfqpoint{5.094631in}{2.803263in}}%
\pgfpathlineto{\pgfqpoint{5.097149in}{2.813111in}}%
\pgfpathlineto{\pgfqpoint{5.099668in}{2.807945in}}%
\pgfpathlineto{\pgfqpoint{5.102187in}{2.814054in}}%
\pgfpathlineto{\pgfqpoint{5.104706in}{2.823472in}}%
\pgfpathlineto{\pgfqpoint{5.107224in}{2.828282in}}%
\pgfpathlineto{\pgfqpoint{5.109743in}{2.787738in}}%
\pgfpathlineto{\pgfqpoint{5.112262in}{2.801655in}}%
\pgfpathlineto{\pgfqpoint{5.114781in}{2.832869in}}%
\pgfpathlineto{\pgfqpoint{5.117299in}{2.836805in}}%
\pgfpathlineto{\pgfqpoint{5.119818in}{2.844524in}}%
\pgfpathlineto{\pgfqpoint{5.124856in}{2.874208in}}%
\pgfpathlineto{\pgfqpoint{5.127374in}{2.879585in}}%
\pgfpathlineto{\pgfqpoint{5.129893in}{2.851789in}}%
\pgfpathlineto{\pgfqpoint{5.132412in}{2.874006in}}%
\pgfpathlineto{\pgfqpoint{5.134931in}{2.878153in}}%
\pgfpathlineto{\pgfqpoint{5.137449in}{2.868474in}}%
\pgfpathlineto{\pgfqpoint{5.139968in}{2.881116in}}%
\pgfpathlineto{\pgfqpoint{5.142487in}{2.897152in}}%
\pgfpathlineto{\pgfqpoint{5.145006in}{2.898361in}}%
\pgfpathlineto{\pgfqpoint{5.147524in}{2.919173in}}%
\pgfpathlineto{\pgfqpoint{5.150043in}{2.957353in}}%
\pgfpathlineto{\pgfqpoint{5.152562in}{2.971117in}}%
\pgfpathlineto{\pgfqpoint{5.155081in}{3.002534in}}%
\pgfpathlineto{\pgfqpoint{5.157599in}{2.980567in}}%
\pgfpathlineto{\pgfqpoint{5.160118in}{2.997541in}}%
\pgfpathlineto{\pgfqpoint{5.162637in}{2.998926in}}%
\pgfpathlineto{\pgfqpoint{5.165156in}{2.993240in}}%
\pgfpathlineto{\pgfqpoint{5.167674in}{2.975827in}}%
\pgfpathlineto{\pgfqpoint{5.170193in}{3.018798in}}%
\pgfpathlineto{\pgfqpoint{5.172712in}{3.024997in}}%
\pgfpathlineto{\pgfqpoint{5.175231in}{3.042978in}}%
\pgfpathlineto{\pgfqpoint{5.177749in}{3.048799in}}%
\pgfpathlineto{\pgfqpoint{5.180268in}{3.051362in}}%
\pgfpathlineto{\pgfqpoint{5.182787in}{3.055047in}}%
\pgfpathlineto{\pgfqpoint{5.185306in}{3.036603in}}%
\pgfpathlineto{\pgfqpoint{5.187824in}{3.021674in}}%
\pgfpathlineto{\pgfqpoint{5.190343in}{3.094166in}}%
\pgfpathlineto{\pgfqpoint{5.192862in}{3.059681in}}%
\pgfpathlineto{\pgfqpoint{5.195381in}{3.063834in}}%
\pgfpathlineto{\pgfqpoint{5.197899in}{3.029400in}}%
\pgfpathlineto{\pgfqpoint{5.200418in}{3.006751in}}%
\pgfpathlineto{\pgfqpoint{5.202937in}{2.995656in}}%
\pgfpathlineto{\pgfqpoint{5.205456in}{3.014581in}}%
\pgfpathlineto{\pgfqpoint{5.207974in}{2.997667in}}%
\pgfpathlineto{\pgfqpoint{5.210493in}{3.002168in}}%
\pgfpathlineto{\pgfqpoint{5.213012in}{2.989377in}}%
\pgfpathlineto{\pgfqpoint{5.215531in}{2.988064in}}%
\pgfpathlineto{\pgfqpoint{5.218049in}{3.043365in}}%
\pgfpathlineto{\pgfqpoint{5.220568in}{3.034941in}}%
\pgfpathlineto{\pgfqpoint{5.223087in}{3.044901in}}%
\pgfpathlineto{\pgfqpoint{5.225606in}{3.055315in}}%
\pgfpathlineto{\pgfqpoint{5.228124in}{3.063143in}}%
\pgfpathlineto{\pgfqpoint{5.230643in}{3.049574in}}%
\pgfpathlineto{\pgfqpoint{5.233162in}{3.097473in}}%
\pgfpathlineto{\pgfqpoint{5.235681in}{3.113335in}}%
\pgfpathlineto{\pgfqpoint{5.238199in}{3.076969in}}%
\pgfpathlineto{\pgfqpoint{5.240718in}{3.069874in}}%
\pgfpathlineto{\pgfqpoint{5.243237in}{3.086576in}}%
\pgfpathlineto{\pgfqpoint{5.245756in}{3.120815in}}%
\pgfpathlineto{\pgfqpoint{5.248274in}{3.093871in}}%
\pgfpathlineto{\pgfqpoint{5.250793in}{3.069231in}}%
\pgfpathlineto{\pgfqpoint{5.253312in}{3.068349in}}%
\pgfpathlineto{\pgfqpoint{5.255831in}{3.040926in}}%
\pgfpathlineto{\pgfqpoint{5.258349in}{3.055917in}}%
\pgfpathlineto{\pgfqpoint{5.260868in}{3.077919in}}%
\pgfpathlineto{\pgfqpoint{5.263387in}{3.063878in}}%
\pgfpathlineto{\pgfqpoint{5.265906in}{3.080112in}}%
\pgfpathlineto{\pgfqpoint{5.268424in}{3.039970in}}%
\pgfpathlineto{\pgfqpoint{5.270943in}{3.067065in}}%
\pgfpathlineto{\pgfqpoint{5.273462in}{3.046967in}}%
\pgfpathlineto{\pgfqpoint{5.275981in}{3.037550in}}%
\pgfpathlineto{\pgfqpoint{5.278499in}{3.053837in}}%
\pgfpathlineto{\pgfqpoint{5.281018in}{3.041795in}}%
\pgfpathlineto{\pgfqpoint{5.283537in}{3.009836in}}%
\pgfpathlineto{\pgfqpoint{5.286056in}{3.012714in}}%
\pgfpathlineto{\pgfqpoint{5.288574in}{2.979249in}}%
\pgfpathlineto{\pgfqpoint{5.293612in}{3.016847in}}%
\pgfpathlineto{\pgfqpoint{5.296131in}{3.039123in}}%
\pgfpathlineto{\pgfqpoint{5.298649in}{3.023515in}}%
\pgfpathlineto{\pgfqpoint{5.301168in}{2.982101in}}%
\pgfpathlineto{\pgfqpoint{5.303687in}{2.985965in}}%
\pgfpathlineto{\pgfqpoint{5.306206in}{3.011219in}}%
\pgfpathlineto{\pgfqpoint{5.308724in}{3.025716in}}%
\pgfpathlineto{\pgfqpoint{5.311243in}{3.015692in}}%
\pgfpathlineto{\pgfqpoint{5.313762in}{3.014410in}}%
\pgfpathlineto{\pgfqpoint{5.316281in}{2.987151in}}%
\pgfpathlineto{\pgfqpoint{5.318799in}{2.958723in}}%
\pgfpathlineto{\pgfqpoint{5.321318in}{2.938182in}}%
\pgfpathlineto{\pgfqpoint{5.323837in}{2.936295in}}%
\pgfpathlineto{\pgfqpoint{5.326356in}{2.944827in}}%
\pgfpathlineto{\pgfqpoint{5.328874in}{2.920855in}}%
\pgfpathlineto{\pgfqpoint{5.331393in}{2.888911in}}%
\pgfpathlineto{\pgfqpoint{5.333912in}{2.925587in}}%
\pgfpathlineto{\pgfqpoint{5.336431in}{2.912542in}}%
\pgfpathlineto{\pgfqpoint{5.338949in}{2.938623in}}%
\pgfpathlineto{\pgfqpoint{5.341468in}{2.929443in}}%
\pgfpathlineto{\pgfqpoint{5.343987in}{2.935240in}}%
\pgfpathlineto{\pgfqpoint{5.346506in}{2.921372in}}%
\pgfpathlineto{\pgfqpoint{5.349024in}{2.934453in}}%
\pgfpathlineto{\pgfqpoint{5.351543in}{2.937667in}}%
\pgfpathlineto{\pgfqpoint{5.354062in}{2.906959in}}%
\pgfpathlineto{\pgfqpoint{5.356581in}{2.904065in}}%
\pgfpathlineto{\pgfqpoint{5.359099in}{2.915532in}}%
\pgfpathlineto{\pgfqpoint{5.361618in}{2.909697in}}%
\pgfpathlineto{\pgfqpoint{5.364137in}{2.913949in}}%
\pgfpathlineto{\pgfqpoint{5.366656in}{2.956206in}}%
\pgfpathlineto{\pgfqpoint{5.369174in}{2.971991in}}%
\pgfpathlineto{\pgfqpoint{5.371693in}{2.980748in}}%
\pgfpathlineto{\pgfqpoint{5.374212in}{3.017378in}}%
\pgfpathlineto{\pgfqpoint{5.376731in}{3.016808in}}%
\pgfpathlineto{\pgfqpoint{5.379249in}{3.007091in}}%
\pgfpathlineto{\pgfqpoint{5.381768in}{3.003413in}}%
\pgfpathlineto{\pgfqpoint{5.384287in}{2.994398in}}%
\pgfpathlineto{\pgfqpoint{5.389324in}{3.059336in}}%
\pgfpathlineto{\pgfqpoint{5.391843in}{3.072530in}}%
\pgfpathlineto{\pgfqpoint{5.394362in}{3.045170in}}%
\pgfpathlineto{\pgfqpoint{5.396881in}{3.063621in}}%
\pgfpathlineto{\pgfqpoint{5.399399in}{3.049527in}}%
\pgfpathlineto{\pgfqpoint{5.401918in}{3.029821in}}%
\pgfpathlineto{\pgfqpoint{5.404437in}{3.011062in}}%
\pgfpathlineto{\pgfqpoint{5.406956in}{3.030763in}}%
\pgfpathlineto{\pgfqpoint{5.409474in}{3.047409in}}%
\pgfpathlineto{\pgfqpoint{5.411993in}{3.056409in}}%
\pgfpathlineto{\pgfqpoint{5.414512in}{3.057731in}}%
\pgfpathlineto{\pgfqpoint{5.417031in}{3.003884in}}%
\pgfpathlineto{\pgfqpoint{5.419549in}{3.014266in}}%
\pgfpathlineto{\pgfqpoint{5.422068in}{3.050383in}}%
\pgfpathlineto{\pgfqpoint{5.424587in}{3.062975in}}%
\pgfpathlineto{\pgfqpoint{5.427106in}{3.050349in}}%
\pgfpathlineto{\pgfqpoint{5.429624in}{3.061361in}}%
\pgfpathlineto{\pgfqpoint{5.432143in}{3.043888in}}%
\pgfpathlineto{\pgfqpoint{5.434662in}{3.064054in}}%
\pgfpathlineto{\pgfqpoint{5.437181in}{3.049273in}}%
\pgfpathlineto{\pgfqpoint{5.439699in}{3.029762in}}%
\pgfpathlineto{\pgfqpoint{5.442218in}{3.021942in}}%
\pgfpathlineto{\pgfqpoint{5.444737in}{3.016562in}}%
\pgfpathlineto{\pgfqpoint{5.447256in}{2.973348in}}%
\pgfpathlineto{\pgfqpoint{5.449774in}{2.925473in}}%
\pgfpathlineto{\pgfqpoint{5.452293in}{2.894296in}}%
\pgfpathlineto{\pgfqpoint{5.454812in}{2.911706in}}%
\pgfpathlineto{\pgfqpoint{5.457331in}{2.886222in}}%
\pgfpathlineto{\pgfqpoint{5.459849in}{2.884000in}}%
\pgfpathlineto{\pgfqpoint{5.462368in}{2.887033in}}%
\pgfpathlineto{\pgfqpoint{5.464887in}{2.855141in}}%
\pgfpathlineto{\pgfqpoint{5.467406in}{2.847259in}}%
\pgfpathlineto{\pgfqpoint{5.469924in}{2.876410in}}%
\pgfpathlineto{\pgfqpoint{5.472443in}{2.880689in}}%
\pgfpathlineto{\pgfqpoint{5.474962in}{2.858269in}}%
\pgfpathlineto{\pgfqpoint{5.477481in}{2.900326in}}%
\pgfpathlineto{\pgfqpoint{5.482518in}{2.938846in}}%
\pgfpathlineto{\pgfqpoint{5.485037in}{2.925911in}}%
\pgfpathlineto{\pgfqpoint{5.487556in}{2.925283in}}%
\pgfpathlineto{\pgfqpoint{5.490074in}{2.923434in}}%
\pgfpathlineto{\pgfqpoint{5.492593in}{2.957338in}}%
\pgfpathlineto{\pgfqpoint{5.492593in}{2.957338in}}%
\pgfusepath{stroke}%
\end{pgfscope}%
\begin{pgfscope}%
\pgfpathrectangle{\pgfqpoint{0.455093in}{0.401080in}}{\pgfqpoint{5.037500in}{4.004000in}}%
\pgfusepath{clip}%
\pgfsetrectcap%
\pgfsetroundjoin%
\pgfsetlinewidth{0.501875pt}%
\definecolor{currentstroke}{rgb}{0.690196,0.188235,0.376471}%
\pgfsetstrokecolor{currentstroke}%
\pgfsetstrokeopacity{0.800000}%
\pgfsetdash{}{0pt}%
\pgfpathmoveto{\pgfqpoint{0.455093in}{1.021044in}}%
\pgfpathlineto{\pgfqpoint{0.457612in}{1.036174in}}%
\pgfpathlineto{\pgfqpoint{0.460131in}{1.032788in}}%
\pgfpathlineto{\pgfqpoint{0.462649in}{1.024945in}}%
\pgfpathlineto{\pgfqpoint{0.465168in}{1.020092in}}%
\pgfpathlineto{\pgfqpoint{0.467687in}{1.009745in}}%
\pgfpathlineto{\pgfqpoint{0.470206in}{1.017567in}}%
\pgfpathlineto{\pgfqpoint{0.472724in}{0.978569in}}%
\pgfpathlineto{\pgfqpoint{0.475243in}{0.991197in}}%
\pgfpathlineto{\pgfqpoint{0.477762in}{0.995508in}}%
\pgfpathlineto{\pgfqpoint{0.480281in}{0.997539in}}%
\pgfpathlineto{\pgfqpoint{0.482799in}{1.002801in}}%
\pgfpathlineto{\pgfqpoint{0.485318in}{0.994185in}}%
\pgfpathlineto{\pgfqpoint{0.487837in}{0.993389in}}%
\pgfpathlineto{\pgfqpoint{0.490356in}{1.002421in}}%
\pgfpathlineto{\pgfqpoint{0.492874in}{1.023488in}}%
\pgfpathlineto{\pgfqpoint{0.495393in}{1.022758in}}%
\pgfpathlineto{\pgfqpoint{0.497912in}{1.021409in}}%
\pgfpathlineto{\pgfqpoint{0.500431in}{0.997096in}}%
\pgfpathlineto{\pgfqpoint{0.502949in}{1.007965in}}%
\pgfpathlineto{\pgfqpoint{0.505468in}{1.020181in}}%
\pgfpathlineto{\pgfqpoint{0.507987in}{1.039127in}}%
\pgfpathlineto{\pgfqpoint{0.510506in}{1.026540in}}%
\pgfpathlineto{\pgfqpoint{0.513024in}{1.031552in}}%
\pgfpathlineto{\pgfqpoint{0.515543in}{1.013486in}}%
\pgfpathlineto{\pgfqpoint{0.518062in}{1.014442in}}%
\pgfpathlineto{\pgfqpoint{0.520581in}{1.019362in}}%
\pgfpathlineto{\pgfqpoint{0.523099in}{1.016259in}}%
\pgfpathlineto{\pgfqpoint{0.525618in}{1.015475in}}%
\pgfpathlineto{\pgfqpoint{0.528137in}{1.013846in}}%
\pgfpathlineto{\pgfqpoint{0.530656in}{1.000797in}}%
\pgfpathlineto{\pgfqpoint{0.533174in}{0.978512in}}%
\pgfpathlineto{\pgfqpoint{0.535693in}{0.966687in}}%
\pgfpathlineto{\pgfqpoint{0.538212in}{0.962055in}}%
\pgfpathlineto{\pgfqpoint{0.540731in}{0.962665in}}%
\pgfpathlineto{\pgfqpoint{0.543249in}{0.975263in}}%
\pgfpathlineto{\pgfqpoint{0.545768in}{0.971105in}}%
\pgfpathlineto{\pgfqpoint{0.548287in}{0.958923in}}%
\pgfpathlineto{\pgfqpoint{0.550806in}{0.988699in}}%
\pgfpathlineto{\pgfqpoint{0.553324in}{0.984841in}}%
\pgfpathlineto{\pgfqpoint{0.555843in}{0.994134in}}%
\pgfpathlineto{\pgfqpoint{0.558362in}{0.988289in}}%
\pgfpathlineto{\pgfqpoint{0.560881in}{0.996699in}}%
\pgfpathlineto{\pgfqpoint{0.563399in}{0.986396in}}%
\pgfpathlineto{\pgfqpoint{0.565918in}{0.968876in}}%
\pgfpathlineto{\pgfqpoint{0.568437in}{0.965366in}}%
\pgfpathlineto{\pgfqpoint{0.570956in}{0.971250in}}%
\pgfpathlineto{\pgfqpoint{0.573474in}{0.967971in}}%
\pgfpathlineto{\pgfqpoint{0.575993in}{0.949539in}}%
\pgfpathlineto{\pgfqpoint{0.578512in}{0.937532in}}%
\pgfpathlineto{\pgfqpoint{0.581031in}{0.915721in}}%
\pgfpathlineto{\pgfqpoint{0.583549in}{0.902406in}}%
\pgfpathlineto{\pgfqpoint{0.586068in}{0.925532in}}%
\pgfpathlineto{\pgfqpoint{0.588587in}{0.929227in}}%
\pgfpathlineto{\pgfqpoint{0.591106in}{0.942984in}}%
\pgfpathlineto{\pgfqpoint{0.593624in}{0.926959in}}%
\pgfpathlineto{\pgfqpoint{0.596143in}{0.933529in}}%
\pgfpathlineto{\pgfqpoint{0.598662in}{0.944140in}}%
\pgfpathlineto{\pgfqpoint{0.601181in}{0.948966in}}%
\pgfpathlineto{\pgfqpoint{0.603699in}{0.937308in}}%
\pgfpathlineto{\pgfqpoint{0.606218in}{0.949935in}}%
\pgfpathlineto{\pgfqpoint{0.608737in}{0.962004in}}%
\pgfpathlineto{\pgfqpoint{0.611256in}{0.950838in}}%
\pgfpathlineto{\pgfqpoint{0.613774in}{0.938953in}}%
\pgfpathlineto{\pgfqpoint{0.616293in}{0.954036in}}%
\pgfpathlineto{\pgfqpoint{0.618812in}{0.956429in}}%
\pgfpathlineto{\pgfqpoint{0.621331in}{0.944640in}}%
\pgfpathlineto{\pgfqpoint{0.623849in}{0.947956in}}%
\pgfpathlineto{\pgfqpoint{0.626368in}{0.967607in}}%
\pgfpathlineto{\pgfqpoint{0.628887in}{0.985834in}}%
\pgfpathlineto{\pgfqpoint{0.631406in}{0.994149in}}%
\pgfpathlineto{\pgfqpoint{0.633924in}{0.998597in}}%
\pgfpathlineto{\pgfqpoint{0.636443in}{0.989478in}}%
\pgfpathlineto{\pgfqpoint{0.638962in}{0.994497in}}%
\pgfpathlineto{\pgfqpoint{0.641481in}{1.008655in}}%
\pgfpathlineto{\pgfqpoint{0.643999in}{1.002900in}}%
\pgfpathlineto{\pgfqpoint{0.646518in}{0.999942in}}%
\pgfpathlineto{\pgfqpoint{0.649037in}{1.005751in}}%
\pgfpathlineto{\pgfqpoint{0.651556in}{1.015855in}}%
\pgfpathlineto{\pgfqpoint{0.654074in}{1.042886in}}%
\pgfpathlineto{\pgfqpoint{0.656593in}{1.055096in}}%
\pgfpathlineto{\pgfqpoint{0.659112in}{1.073902in}}%
\pgfpathlineto{\pgfqpoint{0.661631in}{1.072497in}}%
\pgfpathlineto{\pgfqpoint{0.664149in}{1.052601in}}%
\pgfpathlineto{\pgfqpoint{0.666668in}{1.056339in}}%
\pgfpathlineto{\pgfqpoint{0.669187in}{1.071196in}}%
\pgfpathlineto{\pgfqpoint{0.671706in}{1.065273in}}%
\pgfpathlineto{\pgfqpoint{0.674224in}{1.068721in}}%
\pgfpathlineto{\pgfqpoint{0.676743in}{1.058214in}}%
\pgfpathlineto{\pgfqpoint{0.679262in}{1.068642in}}%
\pgfpathlineto{\pgfqpoint{0.681781in}{1.074522in}}%
\pgfpathlineto{\pgfqpoint{0.686818in}{1.059960in}}%
\pgfpathlineto{\pgfqpoint{0.689337in}{1.085150in}}%
\pgfpathlineto{\pgfqpoint{0.691856in}{1.077158in}}%
\pgfpathlineto{\pgfqpoint{0.694374in}{1.074241in}}%
\pgfpathlineto{\pgfqpoint{0.696893in}{1.067223in}}%
\pgfpathlineto{\pgfqpoint{0.699412in}{1.073232in}}%
\pgfpathlineto{\pgfqpoint{0.701931in}{1.075468in}}%
\pgfpathlineto{\pgfqpoint{0.704449in}{1.074343in}}%
\pgfpathlineto{\pgfqpoint{0.706968in}{1.085056in}}%
\pgfpathlineto{\pgfqpoint{0.709487in}{1.089222in}}%
\pgfpathlineto{\pgfqpoint{0.712006in}{1.106813in}}%
\pgfpathlineto{\pgfqpoint{0.714524in}{1.110624in}}%
\pgfpathlineto{\pgfqpoint{0.717043in}{1.121848in}}%
\pgfpathlineto{\pgfqpoint{0.719562in}{1.122359in}}%
\pgfpathlineto{\pgfqpoint{0.722081in}{1.109101in}}%
\pgfpathlineto{\pgfqpoint{0.724599in}{1.153194in}}%
\pgfpathlineto{\pgfqpoint{0.727118in}{1.146811in}}%
\pgfpathlineto{\pgfqpoint{0.729637in}{1.125860in}}%
\pgfpathlineto{\pgfqpoint{0.732156in}{1.087872in}}%
\pgfpathlineto{\pgfqpoint{0.734674in}{1.086134in}}%
\pgfpathlineto{\pgfqpoint{0.737193in}{1.083357in}}%
\pgfpathlineto{\pgfqpoint{0.739712in}{1.100749in}}%
\pgfpathlineto{\pgfqpoint{0.742231in}{1.111143in}}%
\pgfpathlineto{\pgfqpoint{0.744749in}{1.105591in}}%
\pgfpathlineto{\pgfqpoint{0.747268in}{1.086368in}}%
\pgfpathlineto{\pgfqpoint{0.749787in}{1.106019in}}%
\pgfpathlineto{\pgfqpoint{0.752306in}{1.120644in}}%
\pgfpathlineto{\pgfqpoint{0.754824in}{1.118336in}}%
\pgfpathlineto{\pgfqpoint{0.757343in}{1.134281in}}%
\pgfpathlineto{\pgfqpoint{0.759862in}{1.129695in}}%
\pgfpathlineto{\pgfqpoint{0.762381in}{1.114767in}}%
\pgfpathlineto{\pgfqpoint{0.764899in}{1.117541in}}%
\pgfpathlineto{\pgfqpoint{0.767418in}{1.132065in}}%
\pgfpathlineto{\pgfqpoint{0.769937in}{1.137071in}}%
\pgfpathlineto{\pgfqpoint{0.772456in}{1.131714in}}%
\pgfpathlineto{\pgfqpoint{0.774974in}{1.139221in}}%
\pgfpathlineto{\pgfqpoint{0.777493in}{1.122183in}}%
\pgfpathlineto{\pgfqpoint{0.780012in}{1.145288in}}%
\pgfpathlineto{\pgfqpoint{0.782531in}{1.147964in}}%
\pgfpathlineto{\pgfqpoint{0.785049in}{1.151468in}}%
\pgfpathlineto{\pgfqpoint{0.787568in}{1.151408in}}%
\pgfpathlineto{\pgfqpoint{0.790087in}{1.162332in}}%
\pgfpathlineto{\pgfqpoint{0.792606in}{1.156748in}}%
\pgfpathlineto{\pgfqpoint{0.795124in}{1.174774in}}%
\pgfpathlineto{\pgfqpoint{0.797643in}{1.180341in}}%
\pgfpathlineto{\pgfqpoint{0.800162in}{1.175714in}}%
\pgfpathlineto{\pgfqpoint{0.802681in}{1.172461in}}%
\pgfpathlineto{\pgfqpoint{0.805199in}{1.180697in}}%
\pgfpathlineto{\pgfqpoint{0.807718in}{1.157967in}}%
\pgfpathlineto{\pgfqpoint{0.810237in}{1.174932in}}%
\pgfpathlineto{\pgfqpoint{0.812756in}{1.175216in}}%
\pgfpathlineto{\pgfqpoint{0.815274in}{1.171620in}}%
\pgfpathlineto{\pgfqpoint{0.817793in}{1.171269in}}%
\pgfpathlineto{\pgfqpoint{0.820312in}{1.199981in}}%
\pgfpathlineto{\pgfqpoint{0.822831in}{1.182976in}}%
\pgfpathlineto{\pgfqpoint{0.825349in}{1.183329in}}%
\pgfpathlineto{\pgfqpoint{0.827868in}{1.197529in}}%
\pgfpathlineto{\pgfqpoint{0.830387in}{1.194402in}}%
\pgfpathlineto{\pgfqpoint{0.832906in}{1.226596in}}%
\pgfpathlineto{\pgfqpoint{0.835424in}{1.225250in}}%
\pgfpathlineto{\pgfqpoint{0.837943in}{1.228477in}}%
\pgfpathlineto{\pgfqpoint{0.840462in}{1.231034in}}%
\pgfpathlineto{\pgfqpoint{0.842981in}{1.251746in}}%
\pgfpathlineto{\pgfqpoint{0.845499in}{1.220507in}}%
\pgfpathlineto{\pgfqpoint{0.848018in}{1.247954in}}%
\pgfpathlineto{\pgfqpoint{0.850537in}{1.227883in}}%
\pgfpathlineto{\pgfqpoint{0.853056in}{1.239942in}}%
\pgfpathlineto{\pgfqpoint{0.855574in}{1.237062in}}%
\pgfpathlineto{\pgfqpoint{0.858093in}{1.235296in}}%
\pgfpathlineto{\pgfqpoint{0.860612in}{1.228452in}}%
\pgfpathlineto{\pgfqpoint{0.863131in}{1.225742in}}%
\pgfpathlineto{\pgfqpoint{0.865649in}{1.220160in}}%
\pgfpathlineto{\pgfqpoint{0.868168in}{1.226157in}}%
\pgfpathlineto{\pgfqpoint{0.870687in}{1.245282in}}%
\pgfpathlineto{\pgfqpoint{0.873206in}{1.232339in}}%
\pgfpathlineto{\pgfqpoint{0.875724in}{1.264448in}}%
\pgfpathlineto{\pgfqpoint{0.878243in}{1.268264in}}%
\pgfpathlineto{\pgfqpoint{0.880762in}{1.277098in}}%
\pgfpathlineto{\pgfqpoint{0.883281in}{1.256678in}}%
\pgfpathlineto{\pgfqpoint{0.885799in}{1.255983in}}%
\pgfpathlineto{\pgfqpoint{0.888318in}{1.227418in}}%
\pgfpathlineto{\pgfqpoint{0.890837in}{1.227756in}}%
\pgfpathlineto{\pgfqpoint{0.893356in}{1.246694in}}%
\pgfpathlineto{\pgfqpoint{0.895874in}{1.270520in}}%
\pgfpathlineto{\pgfqpoint{0.898393in}{1.276109in}}%
\pgfpathlineto{\pgfqpoint{0.900912in}{1.261649in}}%
\pgfpathlineto{\pgfqpoint{0.903431in}{1.243395in}}%
\pgfpathlineto{\pgfqpoint{0.905949in}{1.242027in}}%
\pgfpathlineto{\pgfqpoint{0.908468in}{1.215628in}}%
\pgfpathlineto{\pgfqpoint{0.910987in}{1.204636in}}%
\pgfpathlineto{\pgfqpoint{0.913506in}{1.209742in}}%
\pgfpathlineto{\pgfqpoint{0.916024in}{1.250756in}}%
\pgfpathlineto{\pgfqpoint{0.918543in}{1.263815in}}%
\pgfpathlineto{\pgfqpoint{0.921062in}{1.259913in}}%
\pgfpathlineto{\pgfqpoint{0.923581in}{1.241551in}}%
\pgfpathlineto{\pgfqpoint{0.926099in}{1.247909in}}%
\pgfpathlineto{\pgfqpoint{0.928618in}{1.257196in}}%
\pgfpathlineto{\pgfqpoint{0.931137in}{1.251160in}}%
\pgfpathlineto{\pgfqpoint{0.933656in}{1.268773in}}%
\pgfpathlineto{\pgfqpoint{0.936174in}{1.273531in}}%
\pgfpathlineto{\pgfqpoint{0.938693in}{1.282717in}}%
\pgfpathlineto{\pgfqpoint{0.941212in}{1.272180in}}%
\pgfpathlineto{\pgfqpoint{0.943731in}{1.256718in}}%
\pgfpathlineto{\pgfqpoint{0.946249in}{1.235300in}}%
\pgfpathlineto{\pgfqpoint{0.948768in}{1.283997in}}%
\pgfpathlineto{\pgfqpoint{0.951287in}{1.303015in}}%
\pgfpathlineto{\pgfqpoint{0.953806in}{1.298759in}}%
\pgfpathlineto{\pgfqpoint{0.956324in}{1.282698in}}%
\pgfpathlineto{\pgfqpoint{0.958843in}{1.311941in}}%
\pgfpathlineto{\pgfqpoint{0.961362in}{1.332489in}}%
\pgfpathlineto{\pgfqpoint{0.963881in}{1.339187in}}%
\pgfpathlineto{\pgfqpoint{0.966399in}{1.350500in}}%
\pgfpathlineto{\pgfqpoint{0.968918in}{1.343902in}}%
\pgfpathlineto{\pgfqpoint{0.971437in}{1.323373in}}%
\pgfpathlineto{\pgfqpoint{0.973956in}{1.348216in}}%
\pgfpathlineto{\pgfqpoint{0.976474in}{1.348243in}}%
\pgfpathlineto{\pgfqpoint{0.978993in}{1.341006in}}%
\pgfpathlineto{\pgfqpoint{0.981512in}{1.316288in}}%
\pgfpathlineto{\pgfqpoint{0.984031in}{1.327324in}}%
\pgfpathlineto{\pgfqpoint{0.986549in}{1.315757in}}%
\pgfpathlineto{\pgfqpoint{0.989068in}{1.301123in}}%
\pgfpathlineto{\pgfqpoint{0.991587in}{1.309310in}}%
\pgfpathlineto{\pgfqpoint{0.994106in}{1.297151in}}%
\pgfpathlineto{\pgfqpoint{0.996624in}{1.289950in}}%
\pgfpathlineto{\pgfqpoint{0.999143in}{1.267189in}}%
\pgfpathlineto{\pgfqpoint{1.001662in}{1.258047in}}%
\pgfpathlineto{\pgfqpoint{1.004181in}{1.262314in}}%
\pgfpathlineto{\pgfqpoint{1.006699in}{1.263854in}}%
\pgfpathlineto{\pgfqpoint{1.009218in}{1.239695in}}%
\pgfpathlineto{\pgfqpoint{1.011737in}{1.232222in}}%
\pgfpathlineto{\pgfqpoint{1.014256in}{1.247458in}}%
\pgfpathlineto{\pgfqpoint{1.016774in}{1.251037in}}%
\pgfpathlineto{\pgfqpoint{1.019293in}{1.251338in}}%
\pgfpathlineto{\pgfqpoint{1.021812in}{1.236387in}}%
\pgfpathlineto{\pgfqpoint{1.024331in}{1.207866in}}%
\pgfpathlineto{\pgfqpoint{1.026849in}{1.164697in}}%
\pgfpathlineto{\pgfqpoint{1.029368in}{1.162406in}}%
\pgfpathlineto{\pgfqpoint{1.031887in}{1.152892in}}%
\pgfpathlineto{\pgfqpoint{1.034406in}{1.178029in}}%
\pgfpathlineto{\pgfqpoint{1.036924in}{1.168311in}}%
\pgfpathlineto{\pgfqpoint{1.039443in}{1.165623in}}%
\pgfpathlineto{\pgfqpoint{1.041962in}{1.176643in}}%
\pgfpathlineto{\pgfqpoint{1.044481in}{1.188651in}}%
\pgfpathlineto{\pgfqpoint{1.046999in}{1.164032in}}%
\pgfpathlineto{\pgfqpoint{1.049518in}{1.155593in}}%
\pgfpathlineto{\pgfqpoint{1.052037in}{1.157080in}}%
\pgfpathlineto{\pgfqpoint{1.054556in}{1.156736in}}%
\pgfpathlineto{\pgfqpoint{1.057074in}{1.140885in}}%
\pgfpathlineto{\pgfqpoint{1.059593in}{1.148523in}}%
\pgfpathlineto{\pgfqpoint{1.062112in}{1.176920in}}%
\pgfpathlineto{\pgfqpoint{1.064631in}{1.170088in}}%
\pgfpathlineto{\pgfqpoint{1.067149in}{1.168616in}}%
\pgfpathlineto{\pgfqpoint{1.069668in}{1.141790in}}%
\pgfpathlineto{\pgfqpoint{1.072187in}{1.142482in}}%
\pgfpathlineto{\pgfqpoint{1.074706in}{1.143321in}}%
\pgfpathlineto{\pgfqpoint{1.077224in}{1.148955in}}%
\pgfpathlineto{\pgfqpoint{1.079743in}{1.166901in}}%
\pgfpathlineto{\pgfqpoint{1.082262in}{1.170408in}}%
\pgfpathlineto{\pgfqpoint{1.084781in}{1.165787in}}%
\pgfpathlineto{\pgfqpoint{1.087299in}{1.167525in}}%
\pgfpathlineto{\pgfqpoint{1.089818in}{1.181980in}}%
\pgfpathlineto{\pgfqpoint{1.092337in}{1.168452in}}%
\pgfpathlineto{\pgfqpoint{1.094856in}{1.176432in}}%
\pgfpathlineto{\pgfqpoint{1.097374in}{1.169380in}}%
\pgfpathlineto{\pgfqpoint{1.099893in}{1.173058in}}%
\pgfpathlineto{\pgfqpoint{1.102412in}{1.172218in}}%
\pgfpathlineto{\pgfqpoint{1.104931in}{1.190288in}}%
\pgfpathlineto{\pgfqpoint{1.107449in}{1.187742in}}%
\pgfpathlineto{\pgfqpoint{1.109968in}{1.191316in}}%
\pgfpathlineto{\pgfqpoint{1.112487in}{1.170244in}}%
\pgfpathlineto{\pgfqpoint{1.115006in}{1.160507in}}%
\pgfpathlineto{\pgfqpoint{1.117524in}{1.153525in}}%
\pgfpathlineto{\pgfqpoint{1.120043in}{1.141446in}}%
\pgfpathlineto{\pgfqpoint{1.122562in}{1.155155in}}%
\pgfpathlineto{\pgfqpoint{1.125081in}{1.160399in}}%
\pgfpathlineto{\pgfqpoint{1.127599in}{1.167474in}}%
\pgfpathlineto{\pgfqpoint{1.132637in}{1.153938in}}%
\pgfpathlineto{\pgfqpoint{1.135156in}{1.175162in}}%
\pgfpathlineto{\pgfqpoint{1.137674in}{1.195324in}}%
\pgfpathlineto{\pgfqpoint{1.140193in}{1.180719in}}%
\pgfpathlineto{\pgfqpoint{1.142712in}{1.198206in}}%
\pgfpathlineto{\pgfqpoint{1.145231in}{1.191946in}}%
\pgfpathlineto{\pgfqpoint{1.147749in}{1.191512in}}%
\pgfpathlineto{\pgfqpoint{1.150268in}{1.208299in}}%
\pgfpathlineto{\pgfqpoint{1.152787in}{1.214785in}}%
\pgfpathlineto{\pgfqpoint{1.155306in}{1.206679in}}%
\pgfpathlineto{\pgfqpoint{1.157824in}{1.222175in}}%
\pgfpathlineto{\pgfqpoint{1.160343in}{1.240886in}}%
\pgfpathlineto{\pgfqpoint{1.162862in}{1.223142in}}%
\pgfpathlineto{\pgfqpoint{1.165381in}{1.215207in}}%
\pgfpathlineto{\pgfqpoint{1.167899in}{1.199263in}}%
\pgfpathlineto{\pgfqpoint{1.170418in}{1.192611in}}%
\pgfpathlineto{\pgfqpoint{1.172937in}{1.203858in}}%
\pgfpathlineto{\pgfqpoint{1.175456in}{1.190360in}}%
\pgfpathlineto{\pgfqpoint{1.177974in}{1.201054in}}%
\pgfpathlineto{\pgfqpoint{1.180493in}{1.204347in}}%
\pgfpathlineto{\pgfqpoint{1.183012in}{1.219958in}}%
\pgfpathlineto{\pgfqpoint{1.185531in}{1.261171in}}%
\pgfpathlineto{\pgfqpoint{1.188049in}{1.264465in}}%
\pgfpathlineto{\pgfqpoint{1.190568in}{1.258580in}}%
\pgfpathlineto{\pgfqpoint{1.193087in}{1.287551in}}%
\pgfpathlineto{\pgfqpoint{1.195606in}{1.266769in}}%
\pgfpathlineto{\pgfqpoint{1.198124in}{1.230783in}}%
\pgfpathlineto{\pgfqpoint{1.200643in}{1.258476in}}%
\pgfpathlineto{\pgfqpoint{1.203162in}{1.249907in}}%
\pgfpathlineto{\pgfqpoint{1.205681in}{1.248053in}}%
\pgfpathlineto{\pgfqpoint{1.208199in}{1.221076in}}%
\pgfpathlineto{\pgfqpoint{1.210718in}{1.205463in}}%
\pgfpathlineto{\pgfqpoint{1.213237in}{1.201349in}}%
\pgfpathlineto{\pgfqpoint{1.215756in}{1.192248in}}%
\pgfpathlineto{\pgfqpoint{1.218274in}{1.169268in}}%
\pgfpathlineto{\pgfqpoint{1.220793in}{1.187606in}}%
\pgfpathlineto{\pgfqpoint{1.223312in}{1.175194in}}%
\pgfpathlineto{\pgfqpoint{1.225831in}{1.160166in}}%
\pgfpathlineto{\pgfqpoint{1.228349in}{1.167299in}}%
\pgfpathlineto{\pgfqpoint{1.230868in}{1.159069in}}%
\pgfpathlineto{\pgfqpoint{1.233387in}{1.129805in}}%
\pgfpathlineto{\pgfqpoint{1.235906in}{1.118200in}}%
\pgfpathlineto{\pgfqpoint{1.238424in}{1.145128in}}%
\pgfpathlineto{\pgfqpoint{1.240943in}{1.160286in}}%
\pgfpathlineto{\pgfqpoint{1.243462in}{1.161927in}}%
\pgfpathlineto{\pgfqpoint{1.245981in}{1.171835in}}%
\pgfpathlineto{\pgfqpoint{1.248499in}{1.167990in}}%
\pgfpathlineto{\pgfqpoint{1.251018in}{1.185208in}}%
\pgfpathlineto{\pgfqpoint{1.253537in}{1.170080in}}%
\pgfpathlineto{\pgfqpoint{1.256056in}{1.143964in}}%
\pgfpathlineto{\pgfqpoint{1.258574in}{1.112210in}}%
\pgfpathlineto{\pgfqpoint{1.261093in}{1.108183in}}%
\pgfpathlineto{\pgfqpoint{1.263612in}{1.128759in}}%
\pgfpathlineto{\pgfqpoint{1.266131in}{1.143317in}}%
\pgfpathlineto{\pgfqpoint{1.268649in}{1.136491in}}%
\pgfpathlineto{\pgfqpoint{1.271168in}{1.115068in}}%
\pgfpathlineto{\pgfqpoint{1.273687in}{1.111815in}}%
\pgfpathlineto{\pgfqpoint{1.276206in}{1.110499in}}%
\pgfpathlineto{\pgfqpoint{1.278724in}{1.096585in}}%
\pgfpathlineto{\pgfqpoint{1.281243in}{1.095602in}}%
\pgfpathlineto{\pgfqpoint{1.283762in}{1.097132in}}%
\pgfpathlineto{\pgfqpoint{1.286281in}{1.105915in}}%
\pgfpathlineto{\pgfqpoint{1.288799in}{1.122791in}}%
\pgfpathlineto{\pgfqpoint{1.291318in}{1.097135in}}%
\pgfpathlineto{\pgfqpoint{1.293837in}{1.095934in}}%
\pgfpathlineto{\pgfqpoint{1.296356in}{1.106755in}}%
\pgfpathlineto{\pgfqpoint{1.298874in}{1.110943in}}%
\pgfpathlineto{\pgfqpoint{1.301393in}{1.109039in}}%
\pgfpathlineto{\pgfqpoint{1.303912in}{1.124363in}}%
\pgfpathlineto{\pgfqpoint{1.306431in}{1.134365in}}%
\pgfpathlineto{\pgfqpoint{1.308949in}{1.150140in}}%
\pgfpathlineto{\pgfqpoint{1.311468in}{1.155246in}}%
\pgfpathlineto{\pgfqpoint{1.313987in}{1.143469in}}%
\pgfpathlineto{\pgfqpoint{1.319024in}{1.137590in}}%
\pgfpathlineto{\pgfqpoint{1.321543in}{1.116630in}}%
\pgfpathlineto{\pgfqpoint{1.324062in}{1.103895in}}%
\pgfpathlineto{\pgfqpoint{1.326581in}{1.080100in}}%
\pgfpathlineto{\pgfqpoint{1.329099in}{1.077519in}}%
\pgfpathlineto{\pgfqpoint{1.331618in}{1.075536in}}%
\pgfpathlineto{\pgfqpoint{1.334137in}{1.074334in}}%
\pgfpathlineto{\pgfqpoint{1.336656in}{1.066517in}}%
\pgfpathlineto{\pgfqpoint{1.339174in}{1.080574in}}%
\pgfpathlineto{\pgfqpoint{1.341693in}{1.114328in}}%
\pgfpathlineto{\pgfqpoint{1.344212in}{1.135696in}}%
\pgfpathlineto{\pgfqpoint{1.346731in}{1.161097in}}%
\pgfpathlineto{\pgfqpoint{1.349249in}{1.149944in}}%
\pgfpathlineto{\pgfqpoint{1.351768in}{1.165207in}}%
\pgfpathlineto{\pgfqpoint{1.354287in}{1.158961in}}%
\pgfpathlineto{\pgfqpoint{1.356806in}{1.143884in}}%
\pgfpathlineto{\pgfqpoint{1.359324in}{1.152785in}}%
\pgfpathlineto{\pgfqpoint{1.361843in}{1.169512in}}%
\pgfpathlineto{\pgfqpoint{1.364362in}{1.162609in}}%
\pgfpathlineto{\pgfqpoint{1.366881in}{1.147877in}}%
\pgfpathlineto{\pgfqpoint{1.369399in}{1.155149in}}%
\pgfpathlineto{\pgfqpoint{1.371918in}{1.152480in}}%
\pgfpathlineto{\pgfqpoint{1.374437in}{1.139927in}}%
\pgfpathlineto{\pgfqpoint{1.376956in}{1.117182in}}%
\pgfpathlineto{\pgfqpoint{1.379474in}{1.101965in}}%
\pgfpathlineto{\pgfqpoint{1.381993in}{1.107227in}}%
\pgfpathlineto{\pgfqpoint{1.384512in}{1.116239in}}%
\pgfpathlineto{\pgfqpoint{1.387031in}{1.117873in}}%
\pgfpathlineto{\pgfqpoint{1.389549in}{1.104147in}}%
\pgfpathlineto{\pgfqpoint{1.392068in}{1.111657in}}%
\pgfpathlineto{\pgfqpoint{1.394587in}{1.104767in}}%
\pgfpathlineto{\pgfqpoint{1.397106in}{1.093969in}}%
\pgfpathlineto{\pgfqpoint{1.399624in}{1.067939in}}%
\pgfpathlineto{\pgfqpoint{1.402143in}{1.070285in}}%
\pgfpathlineto{\pgfqpoint{1.404662in}{1.065619in}}%
\pgfpathlineto{\pgfqpoint{1.407181in}{1.071000in}}%
\pgfpathlineto{\pgfqpoint{1.409699in}{1.075858in}}%
\pgfpathlineto{\pgfqpoint{1.412218in}{1.079937in}}%
\pgfpathlineto{\pgfqpoint{1.414737in}{1.078041in}}%
\pgfpathlineto{\pgfqpoint{1.417256in}{1.088119in}}%
\pgfpathlineto{\pgfqpoint{1.419774in}{1.100351in}}%
\pgfpathlineto{\pgfqpoint{1.422293in}{1.125755in}}%
\pgfpathlineto{\pgfqpoint{1.424812in}{1.133276in}}%
\pgfpathlineto{\pgfqpoint{1.427331in}{1.130359in}}%
\pgfpathlineto{\pgfqpoint{1.429849in}{1.140181in}}%
\pgfpathlineto{\pgfqpoint{1.432368in}{1.165951in}}%
\pgfpathlineto{\pgfqpoint{1.434887in}{1.154161in}}%
\pgfpathlineto{\pgfqpoint{1.437406in}{1.159398in}}%
\pgfpathlineto{\pgfqpoint{1.439924in}{1.152607in}}%
\pgfpathlineto{\pgfqpoint{1.442443in}{1.156792in}}%
\pgfpathlineto{\pgfqpoint{1.444962in}{1.159866in}}%
\pgfpathlineto{\pgfqpoint{1.447481in}{1.164729in}}%
\pgfpathlineto{\pgfqpoint{1.449999in}{1.177174in}}%
\pgfpathlineto{\pgfqpoint{1.452518in}{1.186377in}}%
\pgfpathlineto{\pgfqpoint{1.455037in}{1.196161in}}%
\pgfpathlineto{\pgfqpoint{1.457556in}{1.179540in}}%
\pgfpathlineto{\pgfqpoint{1.460074in}{1.197977in}}%
\pgfpathlineto{\pgfqpoint{1.462593in}{1.179312in}}%
\pgfpathlineto{\pgfqpoint{1.465112in}{1.181664in}}%
\pgfpathlineto{\pgfqpoint{1.470149in}{1.187491in}}%
\pgfpathlineto{\pgfqpoint{1.472668in}{1.154981in}}%
\pgfpathlineto{\pgfqpoint{1.475187in}{1.150154in}}%
\pgfpathlineto{\pgfqpoint{1.477706in}{1.135049in}}%
\pgfpathlineto{\pgfqpoint{1.480224in}{1.152804in}}%
\pgfpathlineto{\pgfqpoint{1.482743in}{1.162986in}}%
\pgfpathlineto{\pgfqpoint{1.485262in}{1.169833in}}%
\pgfpathlineto{\pgfqpoint{1.487781in}{1.166941in}}%
\pgfpathlineto{\pgfqpoint{1.490299in}{1.149580in}}%
\pgfpathlineto{\pgfqpoint{1.495337in}{1.182312in}}%
\pgfpathlineto{\pgfqpoint{1.497856in}{1.164547in}}%
\pgfpathlineto{\pgfqpoint{1.500374in}{1.196762in}}%
\pgfpathlineto{\pgfqpoint{1.502893in}{1.214033in}}%
\pgfpathlineto{\pgfqpoint{1.505412in}{1.223621in}}%
\pgfpathlineto{\pgfqpoint{1.507931in}{1.220611in}}%
\pgfpathlineto{\pgfqpoint{1.510449in}{1.213159in}}%
\pgfpathlineto{\pgfqpoint{1.512968in}{1.212656in}}%
\pgfpathlineto{\pgfqpoint{1.515487in}{1.225174in}}%
\pgfpathlineto{\pgfqpoint{1.518006in}{1.222804in}}%
\pgfpathlineto{\pgfqpoint{1.520524in}{1.214853in}}%
\pgfpathlineto{\pgfqpoint{1.525562in}{1.191628in}}%
\pgfpathlineto{\pgfqpoint{1.528081in}{1.199665in}}%
\pgfpathlineto{\pgfqpoint{1.530599in}{1.176908in}}%
\pgfpathlineto{\pgfqpoint{1.533118in}{1.175589in}}%
\pgfpathlineto{\pgfqpoint{1.535637in}{1.181921in}}%
\pgfpathlineto{\pgfqpoint{1.538156in}{1.183148in}}%
\pgfpathlineto{\pgfqpoint{1.540674in}{1.173979in}}%
\pgfpathlineto{\pgfqpoint{1.543193in}{1.175745in}}%
\pgfpathlineto{\pgfqpoint{1.545712in}{1.196910in}}%
\pgfpathlineto{\pgfqpoint{1.548231in}{1.208721in}}%
\pgfpathlineto{\pgfqpoint{1.550749in}{1.201683in}}%
\pgfpathlineto{\pgfqpoint{1.553268in}{1.191622in}}%
\pgfpathlineto{\pgfqpoint{1.555787in}{1.178076in}}%
\pgfpathlineto{\pgfqpoint{1.558306in}{1.178187in}}%
\pgfpathlineto{\pgfqpoint{1.560824in}{1.195347in}}%
\pgfpathlineto{\pgfqpoint{1.563343in}{1.211414in}}%
\pgfpathlineto{\pgfqpoint{1.565862in}{1.214292in}}%
\pgfpathlineto{\pgfqpoint{1.568381in}{1.186432in}}%
\pgfpathlineto{\pgfqpoint{1.570899in}{1.211494in}}%
\pgfpathlineto{\pgfqpoint{1.573418in}{1.230574in}}%
\pgfpathlineto{\pgfqpoint{1.575937in}{1.223951in}}%
\pgfpathlineto{\pgfqpoint{1.578456in}{1.241825in}}%
\pgfpathlineto{\pgfqpoint{1.580974in}{1.248889in}}%
\pgfpathlineto{\pgfqpoint{1.583493in}{1.225372in}}%
\pgfpathlineto{\pgfqpoint{1.586012in}{1.232032in}}%
\pgfpathlineto{\pgfqpoint{1.588531in}{1.220199in}}%
\pgfpathlineto{\pgfqpoint{1.591049in}{1.214080in}}%
\pgfpathlineto{\pgfqpoint{1.593568in}{1.225404in}}%
\pgfpathlineto{\pgfqpoint{1.596087in}{1.232881in}}%
\pgfpathlineto{\pgfqpoint{1.598606in}{1.230601in}}%
\pgfpathlineto{\pgfqpoint{1.601124in}{1.232416in}}%
\pgfpathlineto{\pgfqpoint{1.603643in}{1.239274in}}%
\pgfpathlineto{\pgfqpoint{1.606162in}{1.231200in}}%
\pgfpathlineto{\pgfqpoint{1.608681in}{1.246080in}}%
\pgfpathlineto{\pgfqpoint{1.611199in}{1.259147in}}%
\pgfpathlineto{\pgfqpoint{1.613718in}{1.243834in}}%
\pgfpathlineto{\pgfqpoint{1.616237in}{1.275601in}}%
\pgfpathlineto{\pgfqpoint{1.618756in}{1.288377in}}%
\pgfpathlineto{\pgfqpoint{1.621274in}{1.283359in}}%
\pgfpathlineto{\pgfqpoint{1.623793in}{1.275138in}}%
\pgfpathlineto{\pgfqpoint{1.626312in}{1.265660in}}%
\pgfpathlineto{\pgfqpoint{1.628831in}{1.235901in}}%
\pgfpathlineto{\pgfqpoint{1.631349in}{1.240793in}}%
\pgfpathlineto{\pgfqpoint{1.633868in}{1.241318in}}%
\pgfpathlineto{\pgfqpoint{1.636387in}{1.250688in}}%
\pgfpathlineto{\pgfqpoint{1.638906in}{1.238952in}}%
\pgfpathlineto{\pgfqpoint{1.641424in}{1.227862in}}%
\pgfpathlineto{\pgfqpoint{1.643943in}{1.234941in}}%
\pgfpathlineto{\pgfqpoint{1.646462in}{1.264649in}}%
\pgfpathlineto{\pgfqpoint{1.648981in}{1.273646in}}%
\pgfpathlineto{\pgfqpoint{1.651499in}{1.290420in}}%
\pgfpathlineto{\pgfqpoint{1.654018in}{1.300417in}}%
\pgfpathlineto{\pgfqpoint{1.656537in}{1.308968in}}%
\pgfpathlineto{\pgfqpoint{1.659056in}{1.322588in}}%
\pgfpathlineto{\pgfqpoint{1.661574in}{1.346626in}}%
\pgfpathlineto{\pgfqpoint{1.664093in}{1.354649in}}%
\pgfpathlineto{\pgfqpoint{1.666612in}{1.352035in}}%
\pgfpathlineto{\pgfqpoint{1.669131in}{1.329268in}}%
\pgfpathlineto{\pgfqpoint{1.671649in}{1.336628in}}%
\pgfpathlineto{\pgfqpoint{1.674168in}{1.357908in}}%
\pgfpathlineto{\pgfqpoint{1.676687in}{1.372380in}}%
\pgfpathlineto{\pgfqpoint{1.679206in}{1.372228in}}%
\pgfpathlineto{\pgfqpoint{1.681724in}{1.369073in}}%
\pgfpathlineto{\pgfqpoint{1.684243in}{1.374634in}}%
\pgfpathlineto{\pgfqpoint{1.686762in}{1.368996in}}%
\pgfpathlineto{\pgfqpoint{1.689281in}{1.385181in}}%
\pgfpathlineto{\pgfqpoint{1.691799in}{1.400522in}}%
\pgfpathlineto{\pgfqpoint{1.694318in}{1.429582in}}%
\pgfpathlineto{\pgfqpoint{1.696837in}{1.414311in}}%
\pgfpathlineto{\pgfqpoint{1.699356in}{1.413799in}}%
\pgfpathlineto{\pgfqpoint{1.701874in}{1.416294in}}%
\pgfpathlineto{\pgfqpoint{1.704393in}{1.427245in}}%
\pgfpathlineto{\pgfqpoint{1.706912in}{1.421431in}}%
\pgfpathlineto{\pgfqpoint{1.709431in}{1.445187in}}%
\pgfpathlineto{\pgfqpoint{1.711949in}{1.412976in}}%
\pgfpathlineto{\pgfqpoint{1.714468in}{1.436262in}}%
\pgfpathlineto{\pgfqpoint{1.719506in}{1.495131in}}%
\pgfpathlineto{\pgfqpoint{1.722024in}{1.513306in}}%
\pgfpathlineto{\pgfqpoint{1.724543in}{1.525903in}}%
\pgfpathlineto{\pgfqpoint{1.727062in}{1.521890in}}%
\pgfpathlineto{\pgfqpoint{1.729581in}{1.543491in}}%
\pgfpathlineto{\pgfqpoint{1.732099in}{1.545562in}}%
\pgfpathlineto{\pgfqpoint{1.734618in}{1.563685in}}%
\pgfpathlineto{\pgfqpoint{1.737137in}{1.549635in}}%
\pgfpathlineto{\pgfqpoint{1.739656in}{1.559924in}}%
\pgfpathlineto{\pgfqpoint{1.742174in}{1.543671in}}%
\pgfpathlineto{\pgfqpoint{1.744693in}{1.524654in}}%
\pgfpathlineto{\pgfqpoint{1.747212in}{1.517680in}}%
\pgfpathlineto{\pgfqpoint{1.749731in}{1.519859in}}%
\pgfpathlineto{\pgfqpoint{1.752249in}{1.472792in}}%
\pgfpathlineto{\pgfqpoint{1.754768in}{1.457075in}}%
\pgfpathlineto{\pgfqpoint{1.757287in}{1.484011in}}%
\pgfpathlineto{\pgfqpoint{1.759806in}{1.513636in}}%
\pgfpathlineto{\pgfqpoint{1.762324in}{1.502309in}}%
\pgfpathlineto{\pgfqpoint{1.764843in}{1.504877in}}%
\pgfpathlineto{\pgfqpoint{1.767362in}{1.525571in}}%
\pgfpathlineto{\pgfqpoint{1.769881in}{1.494378in}}%
\pgfpathlineto{\pgfqpoint{1.772399in}{1.516232in}}%
\pgfpathlineto{\pgfqpoint{1.774918in}{1.513324in}}%
\pgfpathlineto{\pgfqpoint{1.777437in}{1.518609in}}%
\pgfpathlineto{\pgfqpoint{1.779956in}{1.519749in}}%
\pgfpathlineto{\pgfqpoint{1.782474in}{1.524197in}}%
\pgfpathlineto{\pgfqpoint{1.784993in}{1.554420in}}%
\pgfpathlineto{\pgfqpoint{1.787512in}{1.556721in}}%
\pgfpathlineto{\pgfqpoint{1.790031in}{1.581766in}}%
\pgfpathlineto{\pgfqpoint{1.792549in}{1.582755in}}%
\pgfpathlineto{\pgfqpoint{1.795068in}{1.587503in}}%
\pgfpathlineto{\pgfqpoint{1.797587in}{1.572119in}}%
\pgfpathlineto{\pgfqpoint{1.800106in}{1.567686in}}%
\pgfpathlineto{\pgfqpoint{1.802624in}{1.569154in}}%
\pgfpathlineto{\pgfqpoint{1.805143in}{1.571423in}}%
\pgfpathlineto{\pgfqpoint{1.807662in}{1.571178in}}%
\pgfpathlineto{\pgfqpoint{1.810181in}{1.586184in}}%
\pgfpathlineto{\pgfqpoint{1.812699in}{1.593500in}}%
\pgfpathlineto{\pgfqpoint{1.815218in}{1.612652in}}%
\pgfpathlineto{\pgfqpoint{1.817737in}{1.622629in}}%
\pgfpathlineto{\pgfqpoint{1.820256in}{1.618043in}}%
\pgfpathlineto{\pgfqpoint{1.822774in}{1.614776in}}%
\pgfpathlineto{\pgfqpoint{1.825293in}{1.596060in}}%
\pgfpathlineto{\pgfqpoint{1.827812in}{1.595753in}}%
\pgfpathlineto{\pgfqpoint{1.830331in}{1.584720in}}%
\pgfpathlineto{\pgfqpoint{1.832849in}{1.606287in}}%
\pgfpathlineto{\pgfqpoint{1.835368in}{1.593771in}}%
\pgfpathlineto{\pgfqpoint{1.837887in}{1.596629in}}%
\pgfpathlineto{\pgfqpoint{1.840406in}{1.592695in}}%
\pgfpathlineto{\pgfqpoint{1.842924in}{1.610369in}}%
\pgfpathlineto{\pgfqpoint{1.845443in}{1.591656in}}%
\pgfpathlineto{\pgfqpoint{1.847962in}{1.604181in}}%
\pgfpathlineto{\pgfqpoint{1.850481in}{1.604468in}}%
\pgfpathlineto{\pgfqpoint{1.852999in}{1.600114in}}%
\pgfpathlineto{\pgfqpoint{1.855518in}{1.578650in}}%
\pgfpathlineto{\pgfqpoint{1.858037in}{1.580890in}}%
\pgfpathlineto{\pgfqpoint{1.863074in}{1.542347in}}%
\pgfpathlineto{\pgfqpoint{1.865593in}{1.538094in}}%
\pgfpathlineto{\pgfqpoint{1.868112in}{1.544130in}}%
\pgfpathlineto{\pgfqpoint{1.870631in}{1.571306in}}%
\pgfpathlineto{\pgfqpoint{1.873149in}{1.552349in}}%
\pgfpathlineto{\pgfqpoint{1.875668in}{1.565670in}}%
\pgfpathlineto{\pgfqpoint{1.878187in}{1.571509in}}%
\pgfpathlineto{\pgfqpoint{1.880706in}{1.568508in}}%
\pgfpathlineto{\pgfqpoint{1.883224in}{1.579712in}}%
\pgfpathlineto{\pgfqpoint{1.885743in}{1.585501in}}%
\pgfpathlineto{\pgfqpoint{1.888262in}{1.573501in}}%
\pgfpathlineto{\pgfqpoint{1.890781in}{1.593669in}}%
\pgfpathlineto{\pgfqpoint{1.893299in}{1.573839in}}%
\pgfpathlineto{\pgfqpoint{1.895818in}{1.570640in}}%
\pgfpathlineto{\pgfqpoint{1.900856in}{1.614200in}}%
\pgfpathlineto{\pgfqpoint{1.903374in}{1.618796in}}%
\pgfpathlineto{\pgfqpoint{1.905893in}{1.617217in}}%
\pgfpathlineto{\pgfqpoint{1.908412in}{1.629504in}}%
\pgfpathlineto{\pgfqpoint{1.910931in}{1.624092in}}%
\pgfpathlineto{\pgfqpoint{1.913449in}{1.629900in}}%
\pgfpathlineto{\pgfqpoint{1.915968in}{1.632053in}}%
\pgfpathlineto{\pgfqpoint{1.918487in}{1.627684in}}%
\pgfpathlineto{\pgfqpoint{1.921006in}{1.630715in}}%
\pgfpathlineto{\pgfqpoint{1.923524in}{1.603750in}}%
\pgfpathlineto{\pgfqpoint{1.926043in}{1.574545in}}%
\pgfpathlineto{\pgfqpoint{1.928562in}{1.578615in}}%
\pgfpathlineto{\pgfqpoint{1.931081in}{1.583930in}}%
\pgfpathlineto{\pgfqpoint{1.933599in}{1.591718in}}%
\pgfpathlineto{\pgfqpoint{1.936118in}{1.587179in}}%
\pgfpathlineto{\pgfqpoint{1.938637in}{1.581851in}}%
\pgfpathlineto{\pgfqpoint{1.941156in}{1.563180in}}%
\pgfpathlineto{\pgfqpoint{1.943674in}{1.572101in}}%
\pgfpathlineto{\pgfqpoint{1.946193in}{1.564155in}}%
\pgfpathlineto{\pgfqpoint{1.948712in}{1.567537in}}%
\pgfpathlineto{\pgfqpoint{1.951231in}{1.592711in}}%
\pgfpathlineto{\pgfqpoint{1.953749in}{1.584069in}}%
\pgfpathlineto{\pgfqpoint{1.956268in}{1.584872in}}%
\pgfpathlineto{\pgfqpoint{1.958787in}{1.571897in}}%
\pgfpathlineto{\pgfqpoint{1.961306in}{1.579348in}}%
\pgfpathlineto{\pgfqpoint{1.963824in}{1.579657in}}%
\pgfpathlineto{\pgfqpoint{1.966343in}{1.594259in}}%
\pgfpathlineto{\pgfqpoint{1.968862in}{1.607038in}}%
\pgfpathlineto{\pgfqpoint{1.971381in}{1.593765in}}%
\pgfpathlineto{\pgfqpoint{1.973899in}{1.595231in}}%
\pgfpathlineto{\pgfqpoint{1.976418in}{1.607948in}}%
\pgfpathlineto{\pgfqpoint{1.978937in}{1.612772in}}%
\pgfpathlineto{\pgfqpoint{1.981456in}{1.613540in}}%
\pgfpathlineto{\pgfqpoint{1.983974in}{1.624287in}}%
\pgfpathlineto{\pgfqpoint{1.986493in}{1.599474in}}%
\pgfpathlineto{\pgfqpoint{1.989012in}{1.585082in}}%
\pgfpathlineto{\pgfqpoint{1.991531in}{1.565250in}}%
\pgfpathlineto{\pgfqpoint{1.994049in}{1.552171in}}%
\pgfpathlineto{\pgfqpoint{1.996568in}{1.546550in}}%
\pgfpathlineto{\pgfqpoint{1.999087in}{1.543610in}}%
\pgfpathlineto{\pgfqpoint{2.001606in}{1.545440in}}%
\pgfpathlineto{\pgfqpoint{2.004124in}{1.554315in}}%
\pgfpathlineto{\pgfqpoint{2.006643in}{1.552200in}}%
\pgfpathlineto{\pgfqpoint{2.009162in}{1.547265in}}%
\pgfpathlineto{\pgfqpoint{2.011681in}{1.537340in}}%
\pgfpathlineto{\pgfqpoint{2.014199in}{1.536661in}}%
\pgfpathlineto{\pgfqpoint{2.016718in}{1.512035in}}%
\pgfpathlineto{\pgfqpoint{2.019237in}{1.525401in}}%
\pgfpathlineto{\pgfqpoint{2.021756in}{1.482402in}}%
\pgfpathlineto{\pgfqpoint{2.024274in}{1.488529in}}%
\pgfpathlineto{\pgfqpoint{2.026793in}{1.500786in}}%
\pgfpathlineto{\pgfqpoint{2.029312in}{1.483272in}}%
\pgfpathlineto{\pgfqpoint{2.031831in}{1.469028in}}%
\pgfpathlineto{\pgfqpoint{2.034349in}{1.482467in}}%
\pgfpathlineto{\pgfqpoint{2.036868in}{1.477854in}}%
\pgfpathlineto{\pgfqpoint{2.039387in}{1.476977in}}%
\pgfpathlineto{\pgfqpoint{2.041906in}{1.464934in}}%
\pgfpathlineto{\pgfqpoint{2.044424in}{1.453995in}}%
\pgfpathlineto{\pgfqpoint{2.046943in}{1.458689in}}%
\pgfpathlineto{\pgfqpoint{2.049462in}{1.462507in}}%
\pgfpathlineto{\pgfqpoint{2.051981in}{1.468112in}}%
\pgfpathlineto{\pgfqpoint{2.054499in}{1.429158in}}%
\pgfpathlineto{\pgfqpoint{2.057018in}{1.413659in}}%
\pgfpathlineto{\pgfqpoint{2.059537in}{1.416901in}}%
\pgfpathlineto{\pgfqpoint{2.062056in}{1.405071in}}%
\pgfpathlineto{\pgfqpoint{2.064574in}{1.429453in}}%
\pgfpathlineto{\pgfqpoint{2.067093in}{1.417525in}}%
\pgfpathlineto{\pgfqpoint{2.069612in}{1.402451in}}%
\pgfpathlineto{\pgfqpoint{2.072131in}{1.413812in}}%
\pgfpathlineto{\pgfqpoint{2.074649in}{1.439693in}}%
\pgfpathlineto{\pgfqpoint{2.077168in}{1.436486in}}%
\pgfpathlineto{\pgfqpoint{2.079687in}{1.432811in}}%
\pgfpathlineto{\pgfqpoint{2.082206in}{1.427953in}}%
\pgfpathlineto{\pgfqpoint{2.084724in}{1.458353in}}%
\pgfpathlineto{\pgfqpoint{2.087243in}{1.457114in}}%
\pgfpathlineto{\pgfqpoint{2.089762in}{1.452474in}}%
\pgfpathlineto{\pgfqpoint{2.092281in}{1.469493in}}%
\pgfpathlineto{\pgfqpoint{2.094799in}{1.450329in}}%
\pgfpathlineto{\pgfqpoint{2.097318in}{1.470714in}}%
\pgfpathlineto{\pgfqpoint{2.099837in}{1.480234in}}%
\pgfpathlineto{\pgfqpoint{2.102356in}{1.475891in}}%
\pgfpathlineto{\pgfqpoint{2.104874in}{1.482890in}}%
\pgfpathlineto{\pgfqpoint{2.107393in}{1.509913in}}%
\pgfpathlineto{\pgfqpoint{2.109912in}{1.507522in}}%
\pgfpathlineto{\pgfqpoint{2.112431in}{1.486405in}}%
\pgfpathlineto{\pgfqpoint{2.114949in}{1.475534in}}%
\pgfpathlineto{\pgfqpoint{2.117468in}{1.469884in}}%
\pgfpathlineto{\pgfqpoint{2.119987in}{1.483722in}}%
\pgfpathlineto{\pgfqpoint{2.122506in}{1.507226in}}%
\pgfpathlineto{\pgfqpoint{2.125024in}{1.515567in}}%
\pgfpathlineto{\pgfqpoint{2.127543in}{1.505709in}}%
\pgfpathlineto{\pgfqpoint{2.130062in}{1.500219in}}%
\pgfpathlineto{\pgfqpoint{2.132581in}{1.513136in}}%
\pgfpathlineto{\pgfqpoint{2.135099in}{1.513308in}}%
\pgfpathlineto{\pgfqpoint{2.137618in}{1.550513in}}%
\pgfpathlineto{\pgfqpoint{2.140137in}{1.579753in}}%
\pgfpathlineto{\pgfqpoint{2.142656in}{1.573336in}}%
\pgfpathlineto{\pgfqpoint{2.145174in}{1.571555in}}%
\pgfpathlineto{\pgfqpoint{2.147693in}{1.583831in}}%
\pgfpathlineto{\pgfqpoint{2.150212in}{1.564683in}}%
\pgfpathlineto{\pgfqpoint{2.152731in}{1.546696in}}%
\pgfpathlineto{\pgfqpoint{2.155249in}{1.568778in}}%
\pgfpathlineto{\pgfqpoint{2.157768in}{1.563884in}}%
\pgfpathlineto{\pgfqpoint{2.160287in}{1.567744in}}%
\pgfpathlineto{\pgfqpoint{2.162806in}{1.584799in}}%
\pgfpathlineto{\pgfqpoint{2.165324in}{1.571243in}}%
\pgfpathlineto{\pgfqpoint{2.170362in}{1.565267in}}%
\pgfpathlineto{\pgfqpoint{2.172881in}{1.571975in}}%
\pgfpathlineto{\pgfqpoint{2.175399in}{1.551188in}}%
\pgfpathlineto{\pgfqpoint{2.177918in}{1.558913in}}%
\pgfpathlineto{\pgfqpoint{2.180437in}{1.533837in}}%
\pgfpathlineto{\pgfqpoint{2.182956in}{1.541416in}}%
\pgfpathlineto{\pgfqpoint{2.185474in}{1.530097in}}%
\pgfpathlineto{\pgfqpoint{2.187993in}{1.525067in}}%
\pgfpathlineto{\pgfqpoint{2.190512in}{1.540165in}}%
\pgfpathlineto{\pgfqpoint{2.193031in}{1.524468in}}%
\pgfpathlineto{\pgfqpoint{2.195549in}{1.515558in}}%
\pgfpathlineto{\pgfqpoint{2.198068in}{1.508167in}}%
\pgfpathlineto{\pgfqpoint{2.200587in}{1.498830in}}%
\pgfpathlineto{\pgfqpoint{2.203106in}{1.502039in}}%
\pgfpathlineto{\pgfqpoint{2.205624in}{1.500521in}}%
\pgfpathlineto{\pgfqpoint{2.208143in}{1.521658in}}%
\pgfpathlineto{\pgfqpoint{2.210662in}{1.550721in}}%
\pgfpathlineto{\pgfqpoint{2.213181in}{1.569027in}}%
\pgfpathlineto{\pgfqpoint{2.215699in}{1.573593in}}%
\pgfpathlineto{\pgfqpoint{2.218218in}{1.572229in}}%
\pgfpathlineto{\pgfqpoint{2.220737in}{1.570035in}}%
\pgfpathlineto{\pgfqpoint{2.223256in}{1.560725in}}%
\pgfpathlineto{\pgfqpoint{2.225774in}{1.549900in}}%
\pgfpathlineto{\pgfqpoint{2.228293in}{1.556903in}}%
\pgfpathlineto{\pgfqpoint{2.230812in}{1.528683in}}%
\pgfpathlineto{\pgfqpoint{2.233331in}{1.555419in}}%
\pgfpathlineto{\pgfqpoint{2.235849in}{1.557510in}}%
\pgfpathlineto{\pgfqpoint{2.238368in}{1.569851in}}%
\pgfpathlineto{\pgfqpoint{2.240887in}{1.594364in}}%
\pgfpathlineto{\pgfqpoint{2.243406in}{1.598842in}}%
\pgfpathlineto{\pgfqpoint{2.245924in}{1.587983in}}%
\pgfpathlineto{\pgfqpoint{2.248443in}{1.599571in}}%
\pgfpathlineto{\pgfqpoint{2.250962in}{1.615210in}}%
\pgfpathlineto{\pgfqpoint{2.253481in}{1.585486in}}%
\pgfpathlineto{\pgfqpoint{2.255999in}{1.594860in}}%
\pgfpathlineto{\pgfqpoint{2.258518in}{1.585141in}}%
\pgfpathlineto{\pgfqpoint{2.261037in}{1.584023in}}%
\pgfpathlineto{\pgfqpoint{2.263556in}{1.588101in}}%
\pgfpathlineto{\pgfqpoint{2.266074in}{1.596606in}}%
\pgfpathlineto{\pgfqpoint{2.268593in}{1.598493in}}%
\pgfpathlineto{\pgfqpoint{2.273631in}{1.609634in}}%
\pgfpathlineto{\pgfqpoint{2.276149in}{1.633917in}}%
\pgfpathlineto{\pgfqpoint{2.278668in}{1.635185in}}%
\pgfpathlineto{\pgfqpoint{2.281187in}{1.656182in}}%
\pgfpathlineto{\pgfqpoint{2.283706in}{1.634444in}}%
\pgfpathlineto{\pgfqpoint{2.286224in}{1.634242in}}%
\pgfpathlineto{\pgfqpoint{2.288743in}{1.623205in}}%
\pgfpathlineto{\pgfqpoint{2.291262in}{1.644293in}}%
\pgfpathlineto{\pgfqpoint{2.293781in}{1.635015in}}%
\pgfpathlineto{\pgfqpoint{2.296299in}{1.664782in}}%
\pgfpathlineto{\pgfqpoint{2.298818in}{1.644475in}}%
\pgfpathlineto{\pgfqpoint{2.301337in}{1.640884in}}%
\pgfpathlineto{\pgfqpoint{2.303856in}{1.631843in}}%
\pgfpathlineto{\pgfqpoint{2.306374in}{1.631972in}}%
\pgfpathlineto{\pgfqpoint{2.308893in}{1.595501in}}%
\pgfpathlineto{\pgfqpoint{2.311412in}{1.606483in}}%
\pgfpathlineto{\pgfqpoint{2.313931in}{1.627564in}}%
\pgfpathlineto{\pgfqpoint{2.316449in}{1.644962in}}%
\pgfpathlineto{\pgfqpoint{2.318968in}{1.626096in}}%
\pgfpathlineto{\pgfqpoint{2.321487in}{1.610784in}}%
\pgfpathlineto{\pgfqpoint{2.324006in}{1.618208in}}%
\pgfpathlineto{\pgfqpoint{2.326524in}{1.624363in}}%
\pgfpathlineto{\pgfqpoint{2.329043in}{1.654257in}}%
\pgfpathlineto{\pgfqpoint{2.331562in}{1.660314in}}%
\pgfpathlineto{\pgfqpoint{2.334081in}{1.680018in}}%
\pgfpathlineto{\pgfqpoint{2.336599in}{1.682349in}}%
\pgfpathlineto{\pgfqpoint{2.339118in}{1.712212in}}%
\pgfpathlineto{\pgfqpoint{2.341637in}{1.698707in}}%
\pgfpathlineto{\pgfqpoint{2.344156in}{1.661001in}}%
\pgfpathlineto{\pgfqpoint{2.346674in}{1.659054in}}%
\pgfpathlineto{\pgfqpoint{2.349193in}{1.657329in}}%
\pgfpathlineto{\pgfqpoint{2.351712in}{1.659520in}}%
\pgfpathlineto{\pgfqpoint{2.354231in}{1.684189in}}%
\pgfpathlineto{\pgfqpoint{2.356749in}{1.690628in}}%
\pgfpathlineto{\pgfqpoint{2.359268in}{1.691454in}}%
\pgfpathlineto{\pgfqpoint{2.361787in}{1.669055in}}%
\pgfpathlineto{\pgfqpoint{2.364306in}{1.677272in}}%
\pgfpathlineto{\pgfqpoint{2.366824in}{1.652655in}}%
\pgfpathlineto{\pgfqpoint{2.369343in}{1.663652in}}%
\pgfpathlineto{\pgfqpoint{2.371862in}{1.681407in}}%
\pgfpathlineto{\pgfqpoint{2.374381in}{1.665556in}}%
\pgfpathlineto{\pgfqpoint{2.376899in}{1.672373in}}%
\pgfpathlineto{\pgfqpoint{2.379418in}{1.694431in}}%
\pgfpathlineto{\pgfqpoint{2.381937in}{1.707365in}}%
\pgfpathlineto{\pgfqpoint{2.384456in}{1.665954in}}%
\pgfpathlineto{\pgfqpoint{2.386974in}{1.664312in}}%
\pgfpathlineto{\pgfqpoint{2.389493in}{1.676514in}}%
\pgfpathlineto{\pgfqpoint{2.392012in}{1.680729in}}%
\pgfpathlineto{\pgfqpoint{2.394531in}{1.682224in}}%
\pgfpathlineto{\pgfqpoint{2.397049in}{1.692644in}}%
\pgfpathlineto{\pgfqpoint{2.399568in}{1.688427in}}%
\pgfpathlineto{\pgfqpoint{2.402087in}{1.686875in}}%
\pgfpathlineto{\pgfqpoint{2.404606in}{1.680903in}}%
\pgfpathlineto{\pgfqpoint{2.407124in}{1.682456in}}%
\pgfpathlineto{\pgfqpoint{2.409643in}{1.706607in}}%
\pgfpathlineto{\pgfqpoint{2.412162in}{1.703535in}}%
\pgfpathlineto{\pgfqpoint{2.414681in}{1.671722in}}%
\pgfpathlineto{\pgfqpoint{2.417199in}{1.678605in}}%
\pgfpathlineto{\pgfqpoint{2.419718in}{1.678342in}}%
\pgfpathlineto{\pgfqpoint{2.422237in}{1.711071in}}%
\pgfpathlineto{\pgfqpoint{2.424756in}{1.738949in}}%
\pgfpathlineto{\pgfqpoint{2.427274in}{1.743563in}}%
\pgfpathlineto{\pgfqpoint{2.429793in}{1.734782in}}%
\pgfpathlineto{\pgfqpoint{2.432312in}{1.737667in}}%
\pgfpathlineto{\pgfqpoint{2.434831in}{1.752481in}}%
\pgfpathlineto{\pgfqpoint{2.437349in}{1.768544in}}%
\pgfpathlineto{\pgfqpoint{2.439868in}{1.780590in}}%
\pgfpathlineto{\pgfqpoint{2.442387in}{1.764343in}}%
\pgfpathlineto{\pgfqpoint{2.444906in}{1.744877in}}%
\pgfpathlineto{\pgfqpoint{2.447424in}{1.707378in}}%
\pgfpathlineto{\pgfqpoint{2.449943in}{1.715269in}}%
\pgfpathlineto{\pgfqpoint{2.452462in}{1.716938in}}%
\pgfpathlineto{\pgfqpoint{2.454981in}{1.722189in}}%
\pgfpathlineto{\pgfqpoint{2.457499in}{1.719901in}}%
\pgfpathlineto{\pgfqpoint{2.460018in}{1.714814in}}%
\pgfpathlineto{\pgfqpoint{2.462537in}{1.696762in}}%
\pgfpathlineto{\pgfqpoint{2.465056in}{1.674097in}}%
\pgfpathlineto{\pgfqpoint{2.467574in}{1.675166in}}%
\pgfpathlineto{\pgfqpoint{2.470093in}{1.679344in}}%
\pgfpathlineto{\pgfqpoint{2.472612in}{1.646482in}}%
\pgfpathlineto{\pgfqpoint{2.475131in}{1.663061in}}%
\pgfpathlineto{\pgfqpoint{2.477649in}{1.675708in}}%
\pgfpathlineto{\pgfqpoint{2.480168in}{1.675590in}}%
\pgfpathlineto{\pgfqpoint{2.482687in}{1.708760in}}%
\pgfpathlineto{\pgfqpoint{2.485206in}{1.733930in}}%
\pgfpathlineto{\pgfqpoint{2.487724in}{1.750168in}}%
\pgfpathlineto{\pgfqpoint{2.490243in}{1.764379in}}%
\pgfpathlineto{\pgfqpoint{2.492762in}{1.761441in}}%
\pgfpathlineto{\pgfqpoint{2.495281in}{1.771708in}}%
\pgfpathlineto{\pgfqpoint{2.497799in}{1.762717in}}%
\pgfpathlineto{\pgfqpoint{2.500318in}{1.773121in}}%
\pgfpathlineto{\pgfqpoint{2.505356in}{1.754314in}}%
\pgfpathlineto{\pgfqpoint{2.507874in}{1.770135in}}%
\pgfpathlineto{\pgfqpoint{2.510393in}{1.758019in}}%
\pgfpathlineto{\pgfqpoint{2.512912in}{1.763331in}}%
\pgfpathlineto{\pgfqpoint{2.515431in}{1.738810in}}%
\pgfpathlineto{\pgfqpoint{2.517949in}{1.736316in}}%
\pgfpathlineto{\pgfqpoint{2.520468in}{1.730823in}}%
\pgfpathlineto{\pgfqpoint{2.522987in}{1.720008in}}%
\pgfpathlineto{\pgfqpoint{2.525506in}{1.731736in}}%
\pgfpathlineto{\pgfqpoint{2.528024in}{1.738799in}}%
\pgfpathlineto{\pgfqpoint{2.530543in}{1.710759in}}%
\pgfpathlineto{\pgfqpoint{2.533062in}{1.700531in}}%
\pgfpathlineto{\pgfqpoint{2.535581in}{1.702463in}}%
\pgfpathlineto{\pgfqpoint{2.538099in}{1.695777in}}%
\pgfpathlineto{\pgfqpoint{2.540618in}{1.711498in}}%
\pgfpathlineto{\pgfqpoint{2.543137in}{1.724938in}}%
\pgfpathlineto{\pgfqpoint{2.545656in}{1.697082in}}%
\pgfpathlineto{\pgfqpoint{2.548174in}{1.688674in}}%
\pgfpathlineto{\pgfqpoint{2.550693in}{1.657507in}}%
\pgfpathlineto{\pgfqpoint{2.553212in}{1.670374in}}%
\pgfpathlineto{\pgfqpoint{2.555731in}{1.668220in}}%
\pgfpathlineto{\pgfqpoint{2.558249in}{1.670864in}}%
\pgfpathlineto{\pgfqpoint{2.560768in}{1.664330in}}%
\pgfpathlineto{\pgfqpoint{2.563287in}{1.680353in}}%
\pgfpathlineto{\pgfqpoint{2.565806in}{1.701313in}}%
\pgfpathlineto{\pgfqpoint{2.568324in}{1.694023in}}%
\pgfpathlineto{\pgfqpoint{2.570843in}{1.681736in}}%
\pgfpathlineto{\pgfqpoint{2.573362in}{1.679573in}}%
\pgfpathlineto{\pgfqpoint{2.575881in}{1.654475in}}%
\pgfpathlineto{\pgfqpoint{2.578399in}{1.660583in}}%
\pgfpathlineto{\pgfqpoint{2.580918in}{1.640937in}}%
\pgfpathlineto{\pgfqpoint{2.583437in}{1.646069in}}%
\pgfpathlineto{\pgfqpoint{2.585956in}{1.653990in}}%
\pgfpathlineto{\pgfqpoint{2.588474in}{1.657817in}}%
\pgfpathlineto{\pgfqpoint{2.590993in}{1.630542in}}%
\pgfpathlineto{\pgfqpoint{2.593512in}{1.669884in}}%
\pgfpathlineto{\pgfqpoint{2.596031in}{1.683616in}}%
\pgfpathlineto{\pgfqpoint{2.598549in}{1.701147in}}%
\pgfpathlineto{\pgfqpoint{2.601068in}{1.700202in}}%
\pgfpathlineto{\pgfqpoint{2.603587in}{1.726113in}}%
\pgfpathlineto{\pgfqpoint{2.606106in}{1.744457in}}%
\pgfpathlineto{\pgfqpoint{2.608624in}{1.745537in}}%
\pgfpathlineto{\pgfqpoint{2.611143in}{1.723536in}}%
\pgfpathlineto{\pgfqpoint{2.613662in}{1.727918in}}%
\pgfpathlineto{\pgfqpoint{2.616181in}{1.724989in}}%
\pgfpathlineto{\pgfqpoint{2.618699in}{1.700814in}}%
\pgfpathlineto{\pgfqpoint{2.621218in}{1.707452in}}%
\pgfpathlineto{\pgfqpoint{2.623737in}{1.698172in}}%
\pgfpathlineto{\pgfqpoint{2.626256in}{1.704752in}}%
\pgfpathlineto{\pgfqpoint{2.628774in}{1.685140in}}%
\pgfpathlineto{\pgfqpoint{2.631293in}{1.674022in}}%
\pgfpathlineto{\pgfqpoint{2.633812in}{1.617696in}}%
\pgfpathlineto{\pgfqpoint{2.636331in}{1.628391in}}%
\pgfpathlineto{\pgfqpoint{2.638849in}{1.609584in}}%
\pgfpathlineto{\pgfqpoint{2.641368in}{1.616627in}}%
\pgfpathlineto{\pgfqpoint{2.643887in}{1.614232in}}%
\pgfpathlineto{\pgfqpoint{2.646406in}{1.615876in}}%
\pgfpathlineto{\pgfqpoint{2.648924in}{1.641457in}}%
\pgfpathlineto{\pgfqpoint{2.651443in}{1.617299in}}%
\pgfpathlineto{\pgfqpoint{2.653962in}{1.590188in}}%
\pgfpathlineto{\pgfqpoint{2.656481in}{1.596319in}}%
\pgfpathlineto{\pgfqpoint{2.658999in}{1.594907in}}%
\pgfpathlineto{\pgfqpoint{2.661518in}{1.584397in}}%
\pgfpathlineto{\pgfqpoint{2.664037in}{1.575126in}}%
\pgfpathlineto{\pgfqpoint{2.666556in}{1.594450in}}%
\pgfpathlineto{\pgfqpoint{2.671593in}{1.629432in}}%
\pgfpathlineto{\pgfqpoint{2.674112in}{1.662097in}}%
\pgfpathlineto{\pgfqpoint{2.676631in}{1.686061in}}%
\pgfpathlineto{\pgfqpoint{2.679149in}{1.695434in}}%
\pgfpathlineto{\pgfqpoint{2.681668in}{1.692376in}}%
\pgfpathlineto{\pgfqpoint{2.684187in}{1.690644in}}%
\pgfpathlineto{\pgfqpoint{2.686706in}{1.674472in}}%
\pgfpathlineto{\pgfqpoint{2.689224in}{1.690351in}}%
\pgfpathlineto{\pgfqpoint{2.691743in}{1.672679in}}%
\pgfpathlineto{\pgfqpoint{2.694262in}{1.667154in}}%
\pgfpathlineto{\pgfqpoint{2.696781in}{1.688411in}}%
\pgfpathlineto{\pgfqpoint{2.699299in}{1.683286in}}%
\pgfpathlineto{\pgfqpoint{2.701818in}{1.706506in}}%
\pgfpathlineto{\pgfqpoint{2.704337in}{1.749366in}}%
\pgfpathlineto{\pgfqpoint{2.706856in}{1.750246in}}%
\pgfpathlineto{\pgfqpoint{2.709374in}{1.722323in}}%
\pgfpathlineto{\pgfqpoint{2.711893in}{1.722965in}}%
\pgfpathlineto{\pgfqpoint{2.714412in}{1.691895in}}%
\pgfpathlineto{\pgfqpoint{2.716931in}{1.673608in}}%
\pgfpathlineto{\pgfqpoint{2.719449in}{1.705048in}}%
\pgfpathlineto{\pgfqpoint{2.721968in}{1.707545in}}%
\pgfpathlineto{\pgfqpoint{2.724487in}{1.727790in}}%
\pgfpathlineto{\pgfqpoint{2.727006in}{1.710791in}}%
\pgfpathlineto{\pgfqpoint{2.729524in}{1.699161in}}%
\pgfpathlineto{\pgfqpoint{2.732043in}{1.693637in}}%
\pgfpathlineto{\pgfqpoint{2.734562in}{1.702912in}}%
\pgfpathlineto{\pgfqpoint{2.737081in}{1.723440in}}%
\pgfpathlineto{\pgfqpoint{2.739599in}{1.709569in}}%
\pgfpathlineto{\pgfqpoint{2.742118in}{1.712633in}}%
\pgfpathlineto{\pgfqpoint{2.744637in}{1.713784in}}%
\pgfpathlineto{\pgfqpoint{2.747156in}{1.708551in}}%
\pgfpathlineto{\pgfqpoint{2.749674in}{1.698107in}}%
\pgfpathlineto{\pgfqpoint{2.752193in}{1.674766in}}%
\pgfpathlineto{\pgfqpoint{2.754712in}{1.665922in}}%
\pgfpathlineto{\pgfqpoint{2.757231in}{1.670442in}}%
\pgfpathlineto{\pgfqpoint{2.759749in}{1.666583in}}%
\pgfpathlineto{\pgfqpoint{2.762268in}{1.635550in}}%
\pgfpathlineto{\pgfqpoint{2.764787in}{1.649773in}}%
\pgfpathlineto{\pgfqpoint{2.767306in}{1.636493in}}%
\pgfpathlineto{\pgfqpoint{2.769824in}{1.671220in}}%
\pgfpathlineto{\pgfqpoint{2.772343in}{1.671089in}}%
\pgfpathlineto{\pgfqpoint{2.774862in}{1.691305in}}%
\pgfpathlineto{\pgfqpoint{2.777381in}{1.723666in}}%
\pgfpathlineto{\pgfqpoint{2.779899in}{1.740397in}}%
\pgfpathlineto{\pgfqpoint{2.782418in}{1.729722in}}%
\pgfpathlineto{\pgfqpoint{2.784937in}{1.733223in}}%
\pgfpathlineto{\pgfqpoint{2.787456in}{1.731263in}}%
\pgfpathlineto{\pgfqpoint{2.789974in}{1.760778in}}%
\pgfpathlineto{\pgfqpoint{2.792493in}{1.747593in}}%
\pgfpathlineto{\pgfqpoint{2.795012in}{1.757672in}}%
\pgfpathlineto{\pgfqpoint{2.797531in}{1.719921in}}%
\pgfpathlineto{\pgfqpoint{2.800049in}{1.740332in}}%
\pgfpathlineto{\pgfqpoint{2.802568in}{1.745712in}}%
\pgfpathlineto{\pgfqpoint{2.805087in}{1.742638in}}%
\pgfpathlineto{\pgfqpoint{2.810124in}{1.777391in}}%
\pgfpathlineto{\pgfqpoint{2.812643in}{1.770308in}}%
\pgfpathlineto{\pgfqpoint{2.815162in}{1.769364in}}%
\pgfpathlineto{\pgfqpoint{2.817681in}{1.758442in}}%
\pgfpathlineto{\pgfqpoint{2.820199in}{1.723649in}}%
\pgfpathlineto{\pgfqpoint{2.822718in}{1.703799in}}%
\pgfpathlineto{\pgfqpoint{2.825237in}{1.693073in}}%
\pgfpathlineto{\pgfqpoint{2.827756in}{1.720897in}}%
\pgfpathlineto{\pgfqpoint{2.830274in}{1.700040in}}%
\pgfpathlineto{\pgfqpoint{2.832793in}{1.693414in}}%
\pgfpathlineto{\pgfqpoint{2.835312in}{1.690837in}}%
\pgfpathlineto{\pgfqpoint{2.837831in}{1.707154in}}%
\pgfpathlineto{\pgfqpoint{2.840349in}{1.714548in}}%
\pgfpathlineto{\pgfqpoint{2.842868in}{1.751809in}}%
\pgfpathlineto{\pgfqpoint{2.845387in}{1.735201in}}%
\pgfpathlineto{\pgfqpoint{2.847906in}{1.733135in}}%
\pgfpathlineto{\pgfqpoint{2.852943in}{1.785641in}}%
\pgfpathlineto{\pgfqpoint{2.855462in}{1.798559in}}%
\pgfpathlineto{\pgfqpoint{2.857981in}{1.794136in}}%
\pgfpathlineto{\pgfqpoint{2.860499in}{1.786752in}}%
\pgfpathlineto{\pgfqpoint{2.863018in}{1.792059in}}%
\pgfpathlineto{\pgfqpoint{2.865537in}{1.804389in}}%
\pgfpathlineto{\pgfqpoint{2.868056in}{1.787744in}}%
\pgfpathlineto{\pgfqpoint{2.870574in}{1.768828in}}%
\pgfpathlineto{\pgfqpoint{2.873093in}{1.778616in}}%
\pgfpathlineto{\pgfqpoint{2.875612in}{1.785764in}}%
\pgfpathlineto{\pgfqpoint{2.878131in}{1.720198in}}%
\pgfpathlineto{\pgfqpoint{2.880649in}{1.705023in}}%
\pgfpathlineto{\pgfqpoint{2.883168in}{1.726922in}}%
\pgfpathlineto{\pgfqpoint{2.885687in}{1.709663in}}%
\pgfpathlineto{\pgfqpoint{2.888206in}{1.715711in}}%
\pgfpathlineto{\pgfqpoint{2.890724in}{1.701989in}}%
\pgfpathlineto{\pgfqpoint{2.893243in}{1.735815in}}%
\pgfpathlineto{\pgfqpoint{2.895762in}{1.734593in}}%
\pgfpathlineto{\pgfqpoint{2.898281in}{1.724544in}}%
\pgfpathlineto{\pgfqpoint{2.900799in}{1.732528in}}%
\pgfpathlineto{\pgfqpoint{2.903318in}{1.758001in}}%
\pgfpathlineto{\pgfqpoint{2.905837in}{1.781484in}}%
\pgfpathlineto{\pgfqpoint{2.908356in}{1.799341in}}%
\pgfpathlineto{\pgfqpoint{2.910874in}{1.798912in}}%
\pgfpathlineto{\pgfqpoint{2.913393in}{1.806370in}}%
\pgfpathlineto{\pgfqpoint{2.915912in}{1.796946in}}%
\pgfpathlineto{\pgfqpoint{2.918431in}{1.795660in}}%
\pgfpathlineto{\pgfqpoint{2.920949in}{1.805222in}}%
\pgfpathlineto{\pgfqpoint{2.923468in}{1.821957in}}%
\pgfpathlineto{\pgfqpoint{2.925987in}{1.808552in}}%
\pgfpathlineto{\pgfqpoint{2.928506in}{1.829665in}}%
\pgfpathlineto{\pgfqpoint{2.931024in}{1.825589in}}%
\pgfpathlineto{\pgfqpoint{2.933543in}{1.830299in}}%
\pgfpathlineto{\pgfqpoint{2.936062in}{1.825748in}}%
\pgfpathlineto{\pgfqpoint{2.938581in}{1.813538in}}%
\pgfpathlineto{\pgfqpoint{2.941099in}{1.837126in}}%
\pgfpathlineto{\pgfqpoint{2.943618in}{1.851015in}}%
\pgfpathlineto{\pgfqpoint{2.946137in}{1.842311in}}%
\pgfpathlineto{\pgfqpoint{2.948656in}{1.816877in}}%
\pgfpathlineto{\pgfqpoint{2.951174in}{1.827051in}}%
\pgfpathlineto{\pgfqpoint{2.953693in}{1.815030in}}%
\pgfpathlineto{\pgfqpoint{2.956212in}{1.861516in}}%
\pgfpathlineto{\pgfqpoint{2.958731in}{1.884796in}}%
\pgfpathlineto{\pgfqpoint{2.961249in}{1.902044in}}%
\pgfpathlineto{\pgfqpoint{2.963768in}{1.904162in}}%
\pgfpathlineto{\pgfqpoint{2.966287in}{1.885196in}}%
\pgfpathlineto{\pgfqpoint{2.968806in}{1.881696in}}%
\pgfpathlineto{\pgfqpoint{2.971324in}{1.886225in}}%
\pgfpathlineto{\pgfqpoint{2.973843in}{1.884194in}}%
\pgfpathlineto{\pgfqpoint{2.976362in}{1.855499in}}%
\pgfpathlineto{\pgfqpoint{2.978881in}{1.839646in}}%
\pgfpathlineto{\pgfqpoint{2.981399in}{1.852414in}}%
\pgfpathlineto{\pgfqpoint{2.983918in}{1.868019in}}%
\pgfpathlineto{\pgfqpoint{2.986437in}{1.868381in}}%
\pgfpathlineto{\pgfqpoint{2.988956in}{1.886424in}}%
\pgfpathlineto{\pgfqpoint{2.991474in}{1.888541in}}%
\pgfpathlineto{\pgfqpoint{2.993993in}{1.861348in}}%
\pgfpathlineto{\pgfqpoint{2.996512in}{1.835599in}}%
\pgfpathlineto{\pgfqpoint{2.999031in}{1.837406in}}%
\pgfpathlineto{\pgfqpoint{3.001549in}{1.857882in}}%
\pgfpathlineto{\pgfqpoint{3.004068in}{1.903362in}}%
\pgfpathlineto{\pgfqpoint{3.006587in}{1.919868in}}%
\pgfpathlineto{\pgfqpoint{3.009106in}{1.892323in}}%
\pgfpathlineto{\pgfqpoint{3.011624in}{1.906853in}}%
\pgfpathlineto{\pgfqpoint{3.014143in}{1.926877in}}%
\pgfpathlineto{\pgfqpoint{3.016662in}{1.922833in}}%
\pgfpathlineto{\pgfqpoint{3.019181in}{1.919225in}}%
\pgfpathlineto{\pgfqpoint{3.021699in}{1.915257in}}%
\pgfpathlineto{\pgfqpoint{3.024218in}{1.898527in}}%
\pgfpathlineto{\pgfqpoint{3.026737in}{1.880433in}}%
\pgfpathlineto{\pgfqpoint{3.029256in}{1.881666in}}%
\pgfpathlineto{\pgfqpoint{3.031774in}{1.861275in}}%
\pgfpathlineto{\pgfqpoint{3.034293in}{1.815781in}}%
\pgfpathlineto{\pgfqpoint{3.036812in}{1.802706in}}%
\pgfpathlineto{\pgfqpoint{3.039331in}{1.791417in}}%
\pgfpathlineto{\pgfqpoint{3.041849in}{1.796954in}}%
\pgfpathlineto{\pgfqpoint{3.044368in}{1.790852in}}%
\pgfpathlineto{\pgfqpoint{3.046887in}{1.773343in}}%
\pgfpathlineto{\pgfqpoint{3.049406in}{1.797953in}}%
\pgfpathlineto{\pgfqpoint{3.051924in}{1.824361in}}%
\pgfpathlineto{\pgfqpoint{3.054443in}{1.809528in}}%
\pgfpathlineto{\pgfqpoint{3.056962in}{1.783586in}}%
\pgfpathlineto{\pgfqpoint{3.059481in}{1.783314in}}%
\pgfpathlineto{\pgfqpoint{3.061999in}{1.744771in}}%
\pgfpathlineto{\pgfqpoint{3.064518in}{1.748403in}}%
\pgfpathlineto{\pgfqpoint{3.067037in}{1.767350in}}%
\pgfpathlineto{\pgfqpoint{3.069556in}{1.762903in}}%
\pgfpathlineto{\pgfqpoint{3.072074in}{1.771646in}}%
\pgfpathlineto{\pgfqpoint{3.074593in}{1.754000in}}%
\pgfpathlineto{\pgfqpoint{3.077112in}{1.763811in}}%
\pgfpathlineto{\pgfqpoint{3.079631in}{1.771194in}}%
\pgfpathlineto{\pgfqpoint{3.082149in}{1.773655in}}%
\pgfpathlineto{\pgfqpoint{3.084668in}{1.769275in}}%
\pgfpathlineto{\pgfqpoint{3.087187in}{1.785250in}}%
\pgfpathlineto{\pgfqpoint{3.089706in}{1.818668in}}%
\pgfpathlineto{\pgfqpoint{3.092224in}{1.822031in}}%
\pgfpathlineto{\pgfqpoint{3.094743in}{1.785392in}}%
\pgfpathlineto{\pgfqpoint{3.097262in}{1.774438in}}%
\pgfpathlineto{\pgfqpoint{3.099781in}{1.767854in}}%
\pgfpathlineto{\pgfqpoint{3.102299in}{1.804989in}}%
\pgfpathlineto{\pgfqpoint{3.104818in}{1.817392in}}%
\pgfpathlineto{\pgfqpoint{3.107337in}{1.845860in}}%
\pgfpathlineto{\pgfqpoint{3.109856in}{1.838333in}}%
\pgfpathlineto{\pgfqpoint{3.112374in}{1.853502in}}%
\pgfpathlineto{\pgfqpoint{3.114893in}{1.862282in}}%
\pgfpathlineto{\pgfqpoint{3.117412in}{1.851149in}}%
\pgfpathlineto{\pgfqpoint{3.119931in}{1.842981in}}%
\pgfpathlineto{\pgfqpoint{3.122449in}{1.822869in}}%
\pgfpathlineto{\pgfqpoint{3.124968in}{1.810576in}}%
\pgfpathlineto{\pgfqpoint{3.127487in}{1.822986in}}%
\pgfpathlineto{\pgfqpoint{3.130006in}{1.815213in}}%
\pgfpathlineto{\pgfqpoint{3.132524in}{1.842030in}}%
\pgfpathlineto{\pgfqpoint{3.135043in}{1.840423in}}%
\pgfpathlineto{\pgfqpoint{3.137562in}{1.842189in}}%
\pgfpathlineto{\pgfqpoint{3.140081in}{1.827566in}}%
\pgfpathlineto{\pgfqpoint{3.142599in}{1.825983in}}%
\pgfpathlineto{\pgfqpoint{3.145118in}{1.818312in}}%
\pgfpathlineto{\pgfqpoint{3.147637in}{1.813122in}}%
\pgfpathlineto{\pgfqpoint{3.150156in}{1.813192in}}%
\pgfpathlineto{\pgfqpoint{3.152674in}{1.810956in}}%
\pgfpathlineto{\pgfqpoint{3.155193in}{1.794156in}}%
\pgfpathlineto{\pgfqpoint{3.157712in}{1.783102in}}%
\pgfpathlineto{\pgfqpoint{3.160231in}{1.805050in}}%
\pgfpathlineto{\pgfqpoint{3.162749in}{1.817892in}}%
\pgfpathlineto{\pgfqpoint{3.165268in}{1.801733in}}%
\pgfpathlineto{\pgfqpoint{3.167787in}{1.790737in}}%
\pgfpathlineto{\pgfqpoint{3.170306in}{1.787563in}}%
\pgfpathlineto{\pgfqpoint{3.172824in}{1.771930in}}%
\pgfpathlineto{\pgfqpoint{3.175343in}{1.777626in}}%
\pgfpathlineto{\pgfqpoint{3.177862in}{1.759174in}}%
\pgfpathlineto{\pgfqpoint{3.180381in}{1.754138in}}%
\pgfpathlineto{\pgfqpoint{3.182899in}{1.762465in}}%
\pgfpathlineto{\pgfqpoint{3.185418in}{1.795545in}}%
\pgfpathlineto{\pgfqpoint{3.187937in}{1.819687in}}%
\pgfpathlineto{\pgfqpoint{3.190456in}{1.812145in}}%
\pgfpathlineto{\pgfqpoint{3.192974in}{1.837681in}}%
\pgfpathlineto{\pgfqpoint{3.195493in}{1.847044in}}%
\pgfpathlineto{\pgfqpoint{3.198012in}{1.887961in}}%
\pgfpathlineto{\pgfqpoint{3.200531in}{1.873021in}}%
\pgfpathlineto{\pgfqpoint{3.203049in}{1.861881in}}%
\pgfpathlineto{\pgfqpoint{3.205568in}{1.848914in}}%
\pgfpathlineto{\pgfqpoint{3.208087in}{1.834882in}}%
\pgfpathlineto{\pgfqpoint{3.210606in}{1.854669in}}%
\pgfpathlineto{\pgfqpoint{3.213124in}{1.850344in}}%
\pgfpathlineto{\pgfqpoint{3.215643in}{1.841440in}}%
\pgfpathlineto{\pgfqpoint{3.218162in}{1.828452in}}%
\pgfpathlineto{\pgfqpoint{3.220681in}{1.812507in}}%
\pgfpathlineto{\pgfqpoint{3.223199in}{1.824459in}}%
\pgfpathlineto{\pgfqpoint{3.225718in}{1.809614in}}%
\pgfpathlineto{\pgfqpoint{3.228237in}{1.797087in}}%
\pgfpathlineto{\pgfqpoint{3.230756in}{1.789592in}}%
\pgfpathlineto{\pgfqpoint{3.233274in}{1.755231in}}%
\pgfpathlineto{\pgfqpoint{3.235793in}{1.753203in}}%
\pgfpathlineto{\pgfqpoint{3.238312in}{1.780405in}}%
\pgfpathlineto{\pgfqpoint{3.240831in}{1.815553in}}%
\pgfpathlineto{\pgfqpoint{3.243349in}{1.821142in}}%
\pgfpathlineto{\pgfqpoint{3.245868in}{1.829243in}}%
\pgfpathlineto{\pgfqpoint{3.250906in}{1.852956in}}%
\pgfpathlineto{\pgfqpoint{3.253424in}{1.852959in}}%
\pgfpathlineto{\pgfqpoint{3.255943in}{1.878832in}}%
\pgfpathlineto{\pgfqpoint{3.258462in}{1.875332in}}%
\pgfpathlineto{\pgfqpoint{3.260981in}{1.908278in}}%
\pgfpathlineto{\pgfqpoint{3.263499in}{1.904923in}}%
\pgfpathlineto{\pgfqpoint{3.266018in}{1.896630in}}%
\pgfpathlineto{\pgfqpoint{3.268537in}{1.885059in}}%
\pgfpathlineto{\pgfqpoint{3.271056in}{1.906241in}}%
\pgfpathlineto{\pgfqpoint{3.273574in}{1.920927in}}%
\pgfpathlineto{\pgfqpoint{3.276093in}{1.887854in}}%
\pgfpathlineto{\pgfqpoint{3.278612in}{1.853375in}}%
\pgfpathlineto{\pgfqpoint{3.281131in}{1.844872in}}%
\pgfpathlineto{\pgfqpoint{3.283649in}{1.828574in}}%
\pgfpathlineto{\pgfqpoint{3.286168in}{1.833182in}}%
\pgfpathlineto{\pgfqpoint{3.288687in}{1.851041in}}%
\pgfpathlineto{\pgfqpoint{3.291206in}{1.827645in}}%
\pgfpathlineto{\pgfqpoint{3.293724in}{1.850069in}}%
\pgfpathlineto{\pgfqpoint{3.296243in}{1.835261in}}%
\pgfpathlineto{\pgfqpoint{3.298762in}{1.831622in}}%
\pgfpathlineto{\pgfqpoint{3.301281in}{1.819199in}}%
\pgfpathlineto{\pgfqpoint{3.303799in}{1.818084in}}%
\pgfpathlineto{\pgfqpoint{3.306318in}{1.803291in}}%
\pgfpathlineto{\pgfqpoint{3.308837in}{1.795706in}}%
\pgfpathlineto{\pgfqpoint{3.311356in}{1.775056in}}%
\pgfpathlineto{\pgfqpoint{3.313874in}{1.784430in}}%
\pgfpathlineto{\pgfqpoint{3.316393in}{1.768370in}}%
\pgfpathlineto{\pgfqpoint{3.318912in}{1.759925in}}%
\pgfpathlineto{\pgfqpoint{3.321431in}{1.759983in}}%
\pgfpathlineto{\pgfqpoint{3.323949in}{1.782293in}}%
\pgfpathlineto{\pgfqpoint{3.326468in}{1.784005in}}%
\pgfpathlineto{\pgfqpoint{3.328987in}{1.771598in}}%
\pgfpathlineto{\pgfqpoint{3.331506in}{1.782287in}}%
\pgfpathlineto{\pgfqpoint{3.334024in}{1.772314in}}%
\pgfpathlineto{\pgfqpoint{3.336543in}{1.789137in}}%
\pgfpathlineto{\pgfqpoint{3.339062in}{1.799747in}}%
\pgfpathlineto{\pgfqpoint{3.341581in}{1.789453in}}%
\pgfpathlineto{\pgfqpoint{3.346618in}{1.813616in}}%
\pgfpathlineto{\pgfqpoint{3.349137in}{1.801564in}}%
\pgfpathlineto{\pgfqpoint{3.351656in}{1.800343in}}%
\pgfpathlineto{\pgfqpoint{3.354174in}{1.823114in}}%
\pgfpathlineto{\pgfqpoint{3.356693in}{1.818180in}}%
\pgfpathlineto{\pgfqpoint{3.359212in}{1.811870in}}%
\pgfpathlineto{\pgfqpoint{3.361731in}{1.807142in}}%
\pgfpathlineto{\pgfqpoint{3.364249in}{1.819982in}}%
\pgfpathlineto{\pgfqpoint{3.366768in}{1.830837in}}%
\pgfpathlineto{\pgfqpoint{3.369287in}{1.826490in}}%
\pgfpathlineto{\pgfqpoint{3.371806in}{1.828971in}}%
\pgfpathlineto{\pgfqpoint{3.374324in}{1.835858in}}%
\pgfpathlineto{\pgfqpoint{3.376843in}{1.839399in}}%
\pgfpathlineto{\pgfqpoint{3.379362in}{1.860514in}}%
\pgfpathlineto{\pgfqpoint{3.381881in}{1.863944in}}%
\pgfpathlineto{\pgfqpoint{3.384399in}{1.887176in}}%
\pgfpathlineto{\pgfqpoint{3.386918in}{1.869265in}}%
\pgfpathlineto{\pgfqpoint{3.389437in}{1.890054in}}%
\pgfpathlineto{\pgfqpoint{3.391956in}{1.913395in}}%
\pgfpathlineto{\pgfqpoint{3.394474in}{1.926987in}}%
\pgfpathlineto{\pgfqpoint{3.396993in}{1.909776in}}%
\pgfpathlineto{\pgfqpoint{3.399512in}{1.912207in}}%
\pgfpathlineto{\pgfqpoint{3.402031in}{1.917034in}}%
\pgfpathlineto{\pgfqpoint{3.404549in}{1.910752in}}%
\pgfpathlineto{\pgfqpoint{3.407068in}{1.885538in}}%
\pgfpathlineto{\pgfqpoint{3.409587in}{1.863136in}}%
\pgfpathlineto{\pgfqpoint{3.412106in}{1.865714in}}%
\pgfpathlineto{\pgfqpoint{3.414624in}{1.874179in}}%
\pgfpathlineto{\pgfqpoint{3.417143in}{1.856260in}}%
\pgfpathlineto{\pgfqpoint{3.419662in}{1.870239in}}%
\pgfpathlineto{\pgfqpoint{3.422181in}{1.868315in}}%
\pgfpathlineto{\pgfqpoint{3.424699in}{1.861698in}}%
\pgfpathlineto{\pgfqpoint{3.427218in}{1.852608in}}%
\pgfpathlineto{\pgfqpoint{3.429737in}{1.864401in}}%
\pgfpathlineto{\pgfqpoint{3.432256in}{1.836050in}}%
\pgfpathlineto{\pgfqpoint{3.434774in}{1.825299in}}%
\pgfpathlineto{\pgfqpoint{3.437293in}{1.815325in}}%
\pgfpathlineto{\pgfqpoint{3.439812in}{1.848305in}}%
\pgfpathlineto{\pgfqpoint{3.442331in}{1.842223in}}%
\pgfpathlineto{\pgfqpoint{3.444849in}{1.844207in}}%
\pgfpathlineto{\pgfqpoint{3.449887in}{1.821578in}}%
\pgfpathlineto{\pgfqpoint{3.452406in}{1.849231in}}%
\pgfpathlineto{\pgfqpoint{3.454924in}{1.853089in}}%
\pgfpathlineto{\pgfqpoint{3.457443in}{1.844006in}}%
\pgfpathlineto{\pgfqpoint{3.459962in}{1.844039in}}%
\pgfpathlineto{\pgfqpoint{3.462481in}{1.866383in}}%
\pgfpathlineto{\pgfqpoint{3.464999in}{1.844112in}}%
\pgfpathlineto{\pgfqpoint{3.467518in}{1.845959in}}%
\pgfpathlineto{\pgfqpoint{3.470037in}{1.842465in}}%
\pgfpathlineto{\pgfqpoint{3.472556in}{1.848262in}}%
\pgfpathlineto{\pgfqpoint{3.475074in}{1.840213in}}%
\pgfpathlineto{\pgfqpoint{3.477593in}{1.810540in}}%
\pgfpathlineto{\pgfqpoint{3.480112in}{1.811449in}}%
\pgfpathlineto{\pgfqpoint{3.482631in}{1.855361in}}%
\pgfpathlineto{\pgfqpoint{3.485149in}{1.852706in}}%
\pgfpathlineto{\pgfqpoint{3.487668in}{1.873747in}}%
\pgfpathlineto{\pgfqpoint{3.490187in}{1.869413in}}%
\pgfpathlineto{\pgfqpoint{3.492706in}{1.883526in}}%
\pgfpathlineto{\pgfqpoint{3.495224in}{1.903910in}}%
\pgfpathlineto{\pgfqpoint{3.497743in}{1.892196in}}%
\pgfpathlineto{\pgfqpoint{3.500262in}{1.920958in}}%
\pgfpathlineto{\pgfqpoint{3.502781in}{1.941045in}}%
\pgfpathlineto{\pgfqpoint{3.505299in}{1.934570in}}%
\pgfpathlineto{\pgfqpoint{3.507818in}{1.966743in}}%
\pgfpathlineto{\pgfqpoint{3.510337in}{2.009487in}}%
\pgfpathlineto{\pgfqpoint{3.512856in}{1.999675in}}%
\pgfpathlineto{\pgfqpoint{3.517893in}{1.936329in}}%
\pgfpathlineto{\pgfqpoint{3.520412in}{1.978255in}}%
\pgfpathlineto{\pgfqpoint{3.522931in}{1.980214in}}%
\pgfpathlineto{\pgfqpoint{3.525449in}{2.001017in}}%
\pgfpathlineto{\pgfqpoint{3.527968in}{2.029065in}}%
\pgfpathlineto{\pgfqpoint{3.530487in}{2.032069in}}%
\pgfpathlineto{\pgfqpoint{3.533006in}{2.041732in}}%
\pgfpathlineto{\pgfqpoint{3.535524in}{2.036846in}}%
\pgfpathlineto{\pgfqpoint{3.538043in}{2.045388in}}%
\pgfpathlineto{\pgfqpoint{3.540562in}{2.065708in}}%
\pgfpathlineto{\pgfqpoint{3.543081in}{2.060691in}}%
\pgfpathlineto{\pgfqpoint{3.545599in}{2.076651in}}%
\pgfpathlineto{\pgfqpoint{3.548118in}{2.084013in}}%
\pgfpathlineto{\pgfqpoint{3.550637in}{2.076147in}}%
\pgfpathlineto{\pgfqpoint{3.553156in}{2.077485in}}%
\pgfpathlineto{\pgfqpoint{3.555674in}{2.092798in}}%
\pgfpathlineto{\pgfqpoint{3.558193in}{2.096520in}}%
\pgfpathlineto{\pgfqpoint{3.560712in}{2.086095in}}%
\pgfpathlineto{\pgfqpoint{3.563231in}{2.095231in}}%
\pgfpathlineto{\pgfqpoint{3.565749in}{2.103425in}}%
\pgfpathlineto{\pgfqpoint{3.568268in}{2.084802in}}%
\pgfpathlineto{\pgfqpoint{3.570787in}{2.095320in}}%
\pgfpathlineto{\pgfqpoint{3.573306in}{2.106247in}}%
\pgfpathlineto{\pgfqpoint{3.575824in}{2.121320in}}%
\pgfpathlineto{\pgfqpoint{3.578343in}{2.120684in}}%
\pgfpathlineto{\pgfqpoint{3.580862in}{2.166745in}}%
\pgfpathlineto{\pgfqpoint{3.583381in}{2.159713in}}%
\pgfpathlineto{\pgfqpoint{3.585899in}{2.161161in}}%
\pgfpathlineto{\pgfqpoint{3.588418in}{2.154919in}}%
\pgfpathlineto{\pgfqpoint{3.590937in}{2.177178in}}%
\pgfpathlineto{\pgfqpoint{3.593456in}{2.151041in}}%
\pgfpathlineto{\pgfqpoint{3.595974in}{2.145325in}}%
\pgfpathlineto{\pgfqpoint{3.598493in}{2.157417in}}%
\pgfpathlineto{\pgfqpoint{3.601012in}{2.161747in}}%
\pgfpathlineto{\pgfqpoint{3.603531in}{2.153130in}}%
\pgfpathlineto{\pgfqpoint{3.606049in}{2.124866in}}%
\pgfpathlineto{\pgfqpoint{3.608568in}{2.119254in}}%
\pgfpathlineto{\pgfqpoint{3.611087in}{2.130391in}}%
\pgfpathlineto{\pgfqpoint{3.613606in}{2.128301in}}%
\pgfpathlineto{\pgfqpoint{3.616124in}{2.159245in}}%
\pgfpathlineto{\pgfqpoint{3.618643in}{2.177624in}}%
\pgfpathlineto{\pgfqpoint{3.621162in}{2.207193in}}%
\pgfpathlineto{\pgfqpoint{3.623681in}{2.189664in}}%
\pgfpathlineto{\pgfqpoint{3.626199in}{2.183281in}}%
\pgfpathlineto{\pgfqpoint{3.628718in}{2.205440in}}%
\pgfpathlineto{\pgfqpoint{3.631237in}{2.192069in}}%
\pgfpathlineto{\pgfqpoint{3.633756in}{2.184670in}}%
\pgfpathlineto{\pgfqpoint{3.636274in}{2.181741in}}%
\pgfpathlineto{\pgfqpoint{3.638793in}{2.168420in}}%
\pgfpathlineto{\pgfqpoint{3.641312in}{2.181716in}}%
\pgfpathlineto{\pgfqpoint{3.643831in}{2.192351in}}%
\pgfpathlineto{\pgfqpoint{3.646349in}{2.194341in}}%
\pgfpathlineto{\pgfqpoint{3.648868in}{2.159321in}}%
\pgfpathlineto{\pgfqpoint{3.651387in}{2.166154in}}%
\pgfpathlineto{\pgfqpoint{3.653906in}{2.149891in}}%
\pgfpathlineto{\pgfqpoint{3.656424in}{2.125871in}}%
\pgfpathlineto{\pgfqpoint{3.658943in}{2.142741in}}%
\pgfpathlineto{\pgfqpoint{3.661462in}{2.127953in}}%
\pgfpathlineto{\pgfqpoint{3.663981in}{2.111701in}}%
\pgfpathlineto{\pgfqpoint{3.666499in}{2.100542in}}%
\pgfpathlineto{\pgfqpoint{3.669018in}{2.102031in}}%
\pgfpathlineto{\pgfqpoint{3.671537in}{2.129901in}}%
\pgfpathlineto{\pgfqpoint{3.674056in}{2.131237in}}%
\pgfpathlineto{\pgfqpoint{3.676574in}{2.106591in}}%
\pgfpathlineto{\pgfqpoint{3.679093in}{2.103342in}}%
\pgfpathlineto{\pgfqpoint{3.681612in}{2.109427in}}%
\pgfpathlineto{\pgfqpoint{3.684131in}{2.110912in}}%
\pgfpathlineto{\pgfqpoint{3.686649in}{2.142609in}}%
\pgfpathlineto{\pgfqpoint{3.689168in}{2.159815in}}%
\pgfpathlineto{\pgfqpoint{3.691687in}{2.189286in}}%
\pgfpathlineto{\pgfqpoint{3.694206in}{2.192868in}}%
\pgfpathlineto{\pgfqpoint{3.696724in}{2.171807in}}%
\pgfpathlineto{\pgfqpoint{3.699243in}{2.170490in}}%
\pgfpathlineto{\pgfqpoint{3.701762in}{2.210586in}}%
\pgfpathlineto{\pgfqpoint{3.704281in}{2.212284in}}%
\pgfpathlineto{\pgfqpoint{3.706799in}{2.217798in}}%
\pgfpathlineto{\pgfqpoint{3.709318in}{2.193626in}}%
\pgfpathlineto{\pgfqpoint{3.711837in}{2.199894in}}%
\pgfpathlineto{\pgfqpoint{3.714356in}{2.196031in}}%
\pgfpathlineto{\pgfqpoint{3.716874in}{2.217689in}}%
\pgfpathlineto{\pgfqpoint{3.719393in}{2.212466in}}%
\pgfpathlineto{\pgfqpoint{3.721912in}{2.252427in}}%
\pgfpathlineto{\pgfqpoint{3.724431in}{2.256671in}}%
\pgfpathlineto{\pgfqpoint{3.726949in}{2.282919in}}%
\pgfpathlineto{\pgfqpoint{3.729468in}{2.282432in}}%
\pgfpathlineto{\pgfqpoint{3.731987in}{2.288751in}}%
\pgfpathlineto{\pgfqpoint{3.734506in}{2.333304in}}%
\pgfpathlineto{\pgfqpoint{3.737024in}{2.333348in}}%
\pgfpathlineto{\pgfqpoint{3.739543in}{2.304852in}}%
\pgfpathlineto{\pgfqpoint{3.742062in}{2.297708in}}%
\pgfpathlineto{\pgfqpoint{3.744581in}{2.289025in}}%
\pgfpathlineto{\pgfqpoint{3.747099in}{2.288779in}}%
\pgfpathlineto{\pgfqpoint{3.749618in}{2.284715in}}%
\pgfpathlineto{\pgfqpoint{3.752137in}{2.262620in}}%
\pgfpathlineto{\pgfqpoint{3.754656in}{2.253906in}}%
\pgfpathlineto{\pgfqpoint{3.757174in}{2.278225in}}%
\pgfpathlineto{\pgfqpoint{3.759693in}{2.278519in}}%
\pgfpathlineto{\pgfqpoint{3.762212in}{2.265909in}}%
\pgfpathlineto{\pgfqpoint{3.764731in}{2.285678in}}%
\pgfpathlineto{\pgfqpoint{3.767249in}{2.284138in}}%
\pgfpathlineto{\pgfqpoint{3.769768in}{2.292882in}}%
\pgfpathlineto{\pgfqpoint{3.772287in}{2.307464in}}%
\pgfpathlineto{\pgfqpoint{3.774806in}{2.319285in}}%
\pgfpathlineto{\pgfqpoint{3.777324in}{2.324395in}}%
\pgfpathlineto{\pgfqpoint{3.779843in}{2.325139in}}%
\pgfpathlineto{\pgfqpoint{3.782362in}{2.296536in}}%
\pgfpathlineto{\pgfqpoint{3.784881in}{2.294325in}}%
\pgfpathlineto{\pgfqpoint{3.787399in}{2.286280in}}%
\pgfpathlineto{\pgfqpoint{3.789918in}{2.294068in}}%
\pgfpathlineto{\pgfqpoint{3.792437in}{2.288978in}}%
\pgfpathlineto{\pgfqpoint{3.794956in}{2.302883in}}%
\pgfpathlineto{\pgfqpoint{3.797474in}{2.304870in}}%
\pgfpathlineto{\pgfqpoint{3.799993in}{2.283827in}}%
\pgfpathlineto{\pgfqpoint{3.802512in}{2.294032in}}%
\pgfpathlineto{\pgfqpoint{3.805031in}{2.312520in}}%
\pgfpathlineto{\pgfqpoint{3.807549in}{2.310672in}}%
\pgfpathlineto{\pgfqpoint{3.810068in}{2.330849in}}%
\pgfpathlineto{\pgfqpoint{3.812587in}{2.360846in}}%
\pgfpathlineto{\pgfqpoint{3.815106in}{2.353847in}}%
\pgfpathlineto{\pgfqpoint{3.817624in}{2.362749in}}%
\pgfpathlineto{\pgfqpoint{3.820143in}{2.354740in}}%
\pgfpathlineto{\pgfqpoint{3.822662in}{2.344420in}}%
\pgfpathlineto{\pgfqpoint{3.825181in}{2.408691in}}%
\pgfpathlineto{\pgfqpoint{3.827699in}{2.396091in}}%
\pgfpathlineto{\pgfqpoint{3.830218in}{2.404201in}}%
\pgfpathlineto{\pgfqpoint{3.832737in}{2.430516in}}%
\pgfpathlineto{\pgfqpoint{3.835256in}{2.428720in}}%
\pgfpathlineto{\pgfqpoint{3.837774in}{2.419457in}}%
\pgfpathlineto{\pgfqpoint{3.840293in}{2.388604in}}%
\pgfpathlineto{\pgfqpoint{3.842812in}{2.389740in}}%
\pgfpathlineto{\pgfqpoint{3.845331in}{2.411189in}}%
\pgfpathlineto{\pgfqpoint{3.847849in}{2.415174in}}%
\pgfpathlineto{\pgfqpoint{3.850368in}{2.407182in}}%
\pgfpathlineto{\pgfqpoint{3.852887in}{2.400091in}}%
\pgfpathlineto{\pgfqpoint{3.855406in}{2.413554in}}%
\pgfpathlineto{\pgfqpoint{3.857924in}{2.404556in}}%
\pgfpathlineto{\pgfqpoint{3.860443in}{2.412125in}}%
\pgfpathlineto{\pgfqpoint{3.862962in}{2.441393in}}%
\pgfpathlineto{\pgfqpoint{3.865481in}{2.412026in}}%
\pgfpathlineto{\pgfqpoint{3.867999in}{2.420411in}}%
\pgfpathlineto{\pgfqpoint{3.870518in}{2.408216in}}%
\pgfpathlineto{\pgfqpoint{3.873037in}{2.421403in}}%
\pgfpathlineto{\pgfqpoint{3.875556in}{2.403719in}}%
\pgfpathlineto{\pgfqpoint{3.878074in}{2.394549in}}%
\pgfpathlineto{\pgfqpoint{3.880593in}{2.381711in}}%
\pgfpathlineto{\pgfqpoint{3.883112in}{2.394974in}}%
\pgfpathlineto{\pgfqpoint{3.885631in}{2.395924in}}%
\pgfpathlineto{\pgfqpoint{3.888149in}{2.404357in}}%
\pgfpathlineto{\pgfqpoint{3.890668in}{2.387339in}}%
\pgfpathlineto{\pgfqpoint{3.893187in}{2.368261in}}%
\pgfpathlineto{\pgfqpoint{3.895706in}{2.367703in}}%
\pgfpathlineto{\pgfqpoint{3.898224in}{2.350485in}}%
\pgfpathlineto{\pgfqpoint{3.900743in}{2.378862in}}%
\pgfpathlineto{\pgfqpoint{3.903262in}{2.383891in}}%
\pgfpathlineto{\pgfqpoint{3.905781in}{2.414128in}}%
\pgfpathlineto{\pgfqpoint{3.908299in}{2.429131in}}%
\pgfpathlineto{\pgfqpoint{3.910818in}{2.422416in}}%
\pgfpathlineto{\pgfqpoint{3.913337in}{2.403983in}}%
\pgfpathlineto{\pgfqpoint{3.915856in}{2.392036in}}%
\pgfpathlineto{\pgfqpoint{3.918374in}{2.382656in}}%
\pgfpathlineto{\pgfqpoint{3.920893in}{2.380883in}}%
\pgfpathlineto{\pgfqpoint{3.923412in}{2.390388in}}%
\pgfpathlineto{\pgfqpoint{3.925931in}{2.367205in}}%
\pgfpathlineto{\pgfqpoint{3.928449in}{2.392606in}}%
\pgfpathlineto{\pgfqpoint{3.930968in}{2.423751in}}%
\pgfpathlineto{\pgfqpoint{3.933487in}{2.414371in}}%
\pgfpathlineto{\pgfqpoint{3.936006in}{2.457786in}}%
\pgfpathlineto{\pgfqpoint{3.938524in}{2.451487in}}%
\pgfpathlineto{\pgfqpoint{3.941043in}{2.457258in}}%
\pgfpathlineto{\pgfqpoint{3.943562in}{2.462452in}}%
\pgfpathlineto{\pgfqpoint{3.946081in}{2.477322in}}%
\pgfpathlineto{\pgfqpoint{3.948599in}{2.480124in}}%
\pgfpathlineto{\pgfqpoint{3.951118in}{2.494860in}}%
\pgfpathlineto{\pgfqpoint{3.953637in}{2.483533in}}%
\pgfpathlineto{\pgfqpoint{3.956156in}{2.467209in}}%
\pgfpathlineto{\pgfqpoint{3.958674in}{2.484328in}}%
\pgfpathlineto{\pgfqpoint{3.961193in}{2.440264in}}%
\pgfpathlineto{\pgfqpoint{3.963712in}{2.434490in}}%
\pgfpathlineto{\pgfqpoint{3.966231in}{2.460960in}}%
\pgfpathlineto{\pgfqpoint{3.968749in}{2.456976in}}%
\pgfpathlineto{\pgfqpoint{3.971268in}{2.468315in}}%
\pgfpathlineto{\pgfqpoint{3.973787in}{2.472262in}}%
\pgfpathlineto{\pgfqpoint{3.976306in}{2.465099in}}%
\pgfpathlineto{\pgfqpoint{3.978824in}{2.481550in}}%
\pgfpathlineto{\pgfqpoint{3.981343in}{2.480256in}}%
\pgfpathlineto{\pgfqpoint{3.983862in}{2.477014in}}%
\pgfpathlineto{\pgfqpoint{3.986381in}{2.496927in}}%
\pgfpathlineto{\pgfqpoint{3.988899in}{2.492139in}}%
\pgfpathlineto{\pgfqpoint{3.991418in}{2.503735in}}%
\pgfpathlineto{\pgfqpoint{3.993937in}{2.489062in}}%
\pgfpathlineto{\pgfqpoint{3.996456in}{2.482818in}}%
\pgfpathlineto{\pgfqpoint{3.998974in}{2.478025in}}%
\pgfpathlineto{\pgfqpoint{4.001493in}{2.487862in}}%
\pgfpathlineto{\pgfqpoint{4.004012in}{2.485908in}}%
\pgfpathlineto{\pgfqpoint{4.006531in}{2.464581in}}%
\pgfpathlineto{\pgfqpoint{4.009049in}{2.478701in}}%
\pgfpathlineto{\pgfqpoint{4.011568in}{2.464637in}}%
\pgfpathlineto{\pgfqpoint{4.014087in}{2.472122in}}%
\pgfpathlineto{\pgfqpoint{4.016606in}{2.504617in}}%
\pgfpathlineto{\pgfqpoint{4.019124in}{2.500508in}}%
\pgfpathlineto{\pgfqpoint{4.021643in}{2.438211in}}%
\pgfpathlineto{\pgfqpoint{4.024162in}{2.456099in}}%
\pgfpathlineto{\pgfqpoint{4.026681in}{2.478337in}}%
\pgfpathlineto{\pgfqpoint{4.029199in}{2.510182in}}%
\pgfpathlineto{\pgfqpoint{4.031718in}{2.505729in}}%
\pgfpathlineto{\pgfqpoint{4.034237in}{2.500326in}}%
\pgfpathlineto{\pgfqpoint{4.036756in}{2.496771in}}%
\pgfpathlineto{\pgfqpoint{4.039274in}{2.505994in}}%
\pgfpathlineto{\pgfqpoint{4.041793in}{2.484022in}}%
\pgfpathlineto{\pgfqpoint{4.044312in}{2.503497in}}%
\pgfpathlineto{\pgfqpoint{4.046831in}{2.521373in}}%
\pgfpathlineto{\pgfqpoint{4.049349in}{2.532646in}}%
\pgfpathlineto{\pgfqpoint{4.051868in}{2.561281in}}%
\pgfpathlineto{\pgfqpoint{4.054387in}{2.560344in}}%
\pgfpathlineto{\pgfqpoint{4.056906in}{2.548168in}}%
\pgfpathlineto{\pgfqpoint{4.059424in}{2.564974in}}%
\pgfpathlineto{\pgfqpoint{4.061943in}{2.585530in}}%
\pgfpathlineto{\pgfqpoint{4.064462in}{2.587310in}}%
\pgfpathlineto{\pgfqpoint{4.066981in}{2.575363in}}%
\pgfpathlineto{\pgfqpoint{4.069499in}{2.583378in}}%
\pgfpathlineto{\pgfqpoint{4.072018in}{2.616362in}}%
\pgfpathlineto{\pgfqpoint{4.074537in}{2.636763in}}%
\pgfpathlineto{\pgfqpoint{4.077056in}{2.658157in}}%
\pgfpathlineto{\pgfqpoint{4.079574in}{2.635852in}}%
\pgfpathlineto{\pgfqpoint{4.082093in}{2.641553in}}%
\pgfpathlineto{\pgfqpoint{4.084612in}{2.636967in}}%
\pgfpathlineto{\pgfqpoint{4.087131in}{2.666967in}}%
\pgfpathlineto{\pgfqpoint{4.089649in}{2.671523in}}%
\pgfpathlineto{\pgfqpoint{4.092168in}{2.659874in}}%
\pgfpathlineto{\pgfqpoint{4.094687in}{2.663657in}}%
\pgfpathlineto{\pgfqpoint{4.097206in}{2.690814in}}%
\pgfpathlineto{\pgfqpoint{4.099724in}{2.700692in}}%
\pgfpathlineto{\pgfqpoint{4.102243in}{2.713016in}}%
\pgfpathlineto{\pgfqpoint{4.104762in}{2.741455in}}%
\pgfpathlineto{\pgfqpoint{4.107281in}{2.733653in}}%
\pgfpathlineto{\pgfqpoint{4.109799in}{2.710700in}}%
\pgfpathlineto{\pgfqpoint{4.112318in}{2.742105in}}%
\pgfpathlineto{\pgfqpoint{4.114837in}{2.762006in}}%
\pgfpathlineto{\pgfqpoint{4.117356in}{2.739265in}}%
\pgfpathlineto{\pgfqpoint{4.119874in}{2.736337in}}%
\pgfpathlineto{\pgfqpoint{4.122393in}{2.725123in}}%
\pgfpathlineto{\pgfqpoint{4.124912in}{2.738560in}}%
\pgfpathlineto{\pgfqpoint{4.127431in}{2.756747in}}%
\pgfpathlineto{\pgfqpoint{4.129949in}{2.780430in}}%
\pgfpathlineto{\pgfqpoint{4.132468in}{2.792648in}}%
\pgfpathlineto{\pgfqpoint{4.134987in}{2.833183in}}%
\pgfpathlineto{\pgfqpoint{4.137506in}{2.854051in}}%
\pgfpathlineto{\pgfqpoint{4.140024in}{2.822655in}}%
\pgfpathlineto{\pgfqpoint{4.142543in}{2.789016in}}%
\pgfpathlineto{\pgfqpoint{4.145062in}{2.814961in}}%
\pgfpathlineto{\pgfqpoint{4.147581in}{2.834289in}}%
\pgfpathlineto{\pgfqpoint{4.150099in}{2.798869in}}%
\pgfpathlineto{\pgfqpoint{4.152618in}{2.807065in}}%
\pgfpathlineto{\pgfqpoint{4.155137in}{2.820098in}}%
\pgfpathlineto{\pgfqpoint{4.157656in}{2.836550in}}%
\pgfpathlineto{\pgfqpoint{4.160174in}{2.846073in}}%
\pgfpathlineto{\pgfqpoint{4.162693in}{2.849453in}}%
\pgfpathlineto{\pgfqpoint{4.165212in}{2.873928in}}%
\pgfpathlineto{\pgfqpoint{4.167731in}{2.843728in}}%
\pgfpathlineto{\pgfqpoint{4.170249in}{2.834133in}}%
\pgfpathlineto{\pgfqpoint{4.172768in}{2.848694in}}%
\pgfpathlineto{\pgfqpoint{4.175287in}{2.841185in}}%
\pgfpathlineto{\pgfqpoint{4.177806in}{2.876880in}}%
\pgfpathlineto{\pgfqpoint{4.180324in}{2.855248in}}%
\pgfpathlineto{\pgfqpoint{4.182843in}{2.862853in}}%
\pgfpathlineto{\pgfqpoint{4.185362in}{2.874281in}}%
\pgfpathlineto{\pgfqpoint{4.187881in}{2.876648in}}%
\pgfpathlineto{\pgfqpoint{4.190399in}{2.884183in}}%
\pgfpathlineto{\pgfqpoint{4.192918in}{2.911512in}}%
\pgfpathlineto{\pgfqpoint{4.195437in}{2.920797in}}%
\pgfpathlineto{\pgfqpoint{4.197956in}{2.909398in}}%
\pgfpathlineto{\pgfqpoint{4.200474in}{2.918100in}}%
\pgfpathlineto{\pgfqpoint{4.202993in}{2.922378in}}%
\pgfpathlineto{\pgfqpoint{4.208031in}{2.912840in}}%
\pgfpathlineto{\pgfqpoint{4.210549in}{2.950960in}}%
\pgfpathlineto{\pgfqpoint{4.213068in}{2.966517in}}%
\pgfpathlineto{\pgfqpoint{4.215587in}{2.948870in}}%
\pgfpathlineto{\pgfqpoint{4.218106in}{2.945188in}}%
\pgfpathlineto{\pgfqpoint{4.220624in}{2.959896in}}%
\pgfpathlineto{\pgfqpoint{4.223143in}{2.911566in}}%
\pgfpathlineto{\pgfqpoint{4.225662in}{2.898352in}}%
\pgfpathlineto{\pgfqpoint{4.228181in}{2.914118in}}%
\pgfpathlineto{\pgfqpoint{4.230699in}{2.925808in}}%
\pgfpathlineto{\pgfqpoint{4.233218in}{2.900027in}}%
\pgfpathlineto{\pgfqpoint{4.235737in}{2.934068in}}%
\pgfpathlineto{\pgfqpoint{4.238256in}{2.953443in}}%
\pgfpathlineto{\pgfqpoint{4.240774in}{2.938779in}}%
\pgfpathlineto{\pgfqpoint{4.243293in}{2.938495in}}%
\pgfpathlineto{\pgfqpoint{4.245812in}{2.930447in}}%
\pgfpathlineto{\pgfqpoint{4.248331in}{2.966514in}}%
\pgfpathlineto{\pgfqpoint{4.250849in}{2.969052in}}%
\pgfpathlineto{\pgfqpoint{4.253368in}{2.960921in}}%
\pgfpathlineto{\pgfqpoint{4.255887in}{2.959580in}}%
\pgfpathlineto{\pgfqpoint{4.258406in}{2.922102in}}%
\pgfpathlineto{\pgfqpoint{4.260924in}{2.917609in}}%
\pgfpathlineto{\pgfqpoint{4.263443in}{2.931153in}}%
\pgfpathlineto{\pgfqpoint{4.265962in}{2.887026in}}%
\pgfpathlineto{\pgfqpoint{4.268481in}{2.884288in}}%
\pgfpathlineto{\pgfqpoint{4.270999in}{2.859026in}}%
\pgfpathlineto{\pgfqpoint{4.273518in}{2.871706in}}%
\pgfpathlineto{\pgfqpoint{4.276037in}{2.866036in}}%
\pgfpathlineto{\pgfqpoint{4.278556in}{2.865607in}}%
\pgfpathlineto{\pgfqpoint{4.281074in}{2.801928in}}%
\pgfpathlineto{\pgfqpoint{4.283593in}{2.812457in}}%
\pgfpathlineto{\pgfqpoint{4.286112in}{2.836449in}}%
\pgfpathlineto{\pgfqpoint{4.288631in}{2.859164in}}%
\pgfpathlineto{\pgfqpoint{4.291149in}{2.885089in}}%
\pgfpathlineto{\pgfqpoint{4.293668in}{2.897272in}}%
\pgfpathlineto{\pgfqpoint{4.296187in}{2.890944in}}%
\pgfpathlineto{\pgfqpoint{4.298706in}{2.883435in}}%
\pgfpathlineto{\pgfqpoint{4.301224in}{2.889085in}}%
\pgfpathlineto{\pgfqpoint{4.303743in}{2.912060in}}%
\pgfpathlineto{\pgfqpoint{4.306262in}{2.944318in}}%
\pgfpathlineto{\pgfqpoint{4.308781in}{2.978446in}}%
\pgfpathlineto{\pgfqpoint{4.311299in}{3.005926in}}%
\pgfpathlineto{\pgfqpoint{4.313818in}{2.996348in}}%
\pgfpathlineto{\pgfqpoint{4.316337in}{2.978925in}}%
\pgfpathlineto{\pgfqpoint{4.318856in}{2.975598in}}%
\pgfpathlineto{\pgfqpoint{4.321374in}{2.940966in}}%
\pgfpathlineto{\pgfqpoint{4.323893in}{2.959318in}}%
\pgfpathlineto{\pgfqpoint{4.326412in}{2.931812in}}%
\pgfpathlineto{\pgfqpoint{4.328931in}{2.923681in}}%
\pgfpathlineto{\pgfqpoint{4.331449in}{2.922518in}}%
\pgfpathlineto{\pgfqpoint{4.333968in}{2.970520in}}%
\pgfpathlineto{\pgfqpoint{4.336487in}{2.963152in}}%
\pgfpathlineto{\pgfqpoint{4.341524in}{3.003157in}}%
\pgfpathlineto{\pgfqpoint{4.344043in}{3.029088in}}%
\pgfpathlineto{\pgfqpoint{4.346562in}{3.063236in}}%
\pgfpathlineto{\pgfqpoint{4.349081in}{3.061119in}}%
\pgfpathlineto{\pgfqpoint{4.351599in}{3.066201in}}%
\pgfpathlineto{\pgfqpoint{4.354118in}{3.066100in}}%
\pgfpathlineto{\pgfqpoint{4.356637in}{3.077096in}}%
\pgfpathlineto{\pgfqpoint{4.359156in}{3.124193in}}%
\pgfpathlineto{\pgfqpoint{4.361674in}{3.119825in}}%
\pgfpathlineto{\pgfqpoint{4.364193in}{3.094970in}}%
\pgfpathlineto{\pgfqpoint{4.366712in}{3.090961in}}%
\pgfpathlineto{\pgfqpoint{4.369231in}{3.074495in}}%
\pgfpathlineto{\pgfqpoint{4.371749in}{3.075919in}}%
\pgfpathlineto{\pgfqpoint{4.374268in}{3.051298in}}%
\pgfpathlineto{\pgfqpoint{4.376787in}{3.033495in}}%
\pgfpathlineto{\pgfqpoint{4.379306in}{3.022675in}}%
\pgfpathlineto{\pgfqpoint{4.381824in}{3.014190in}}%
\pgfpathlineto{\pgfqpoint{4.384343in}{3.008747in}}%
\pgfpathlineto{\pgfqpoint{4.386862in}{3.023589in}}%
\pgfpathlineto{\pgfqpoint{4.389381in}{2.996523in}}%
\pgfpathlineto{\pgfqpoint{4.391899in}{2.982582in}}%
\pgfpathlineto{\pgfqpoint{4.394418in}{3.009897in}}%
\pgfpathlineto{\pgfqpoint{4.396937in}{3.003039in}}%
\pgfpathlineto{\pgfqpoint{4.399456in}{3.020362in}}%
\pgfpathlineto{\pgfqpoint{4.401974in}{3.029946in}}%
\pgfpathlineto{\pgfqpoint{4.404493in}{3.035651in}}%
\pgfpathlineto{\pgfqpoint{4.407012in}{3.029517in}}%
\pgfpathlineto{\pgfqpoint{4.409531in}{3.048396in}}%
\pgfpathlineto{\pgfqpoint{4.412049in}{2.999652in}}%
\pgfpathlineto{\pgfqpoint{4.414568in}{2.991443in}}%
\pgfpathlineto{\pgfqpoint{4.417087in}{2.990803in}}%
\pgfpathlineto{\pgfqpoint{4.422124in}{2.930302in}}%
\pgfpathlineto{\pgfqpoint{4.424643in}{2.946679in}}%
\pgfpathlineto{\pgfqpoint{4.427162in}{2.941947in}}%
\pgfpathlineto{\pgfqpoint{4.429681in}{2.955507in}}%
\pgfpathlineto{\pgfqpoint{4.432199in}{2.942092in}}%
\pgfpathlineto{\pgfqpoint{4.434718in}{2.964091in}}%
\pgfpathlineto{\pgfqpoint{4.437237in}{2.940080in}}%
\pgfpathlineto{\pgfqpoint{4.439756in}{2.929926in}}%
\pgfpathlineto{\pgfqpoint{4.442274in}{2.937961in}}%
\pgfpathlineto{\pgfqpoint{4.444793in}{2.972546in}}%
\pgfpathlineto{\pgfqpoint{4.447312in}{2.982377in}}%
\pgfpathlineto{\pgfqpoint{4.449831in}{2.971825in}}%
\pgfpathlineto{\pgfqpoint{4.452349in}{2.987710in}}%
\pgfpathlineto{\pgfqpoint{4.454868in}{2.980190in}}%
\pgfpathlineto{\pgfqpoint{4.457387in}{2.963797in}}%
\pgfpathlineto{\pgfqpoint{4.459906in}{2.989007in}}%
\pgfpathlineto{\pgfqpoint{4.462424in}{3.026625in}}%
\pgfpathlineto{\pgfqpoint{4.464943in}{3.013172in}}%
\pgfpathlineto{\pgfqpoint{4.467462in}{3.025796in}}%
\pgfpathlineto{\pgfqpoint{4.469981in}{3.040657in}}%
\pgfpathlineto{\pgfqpoint{4.472499in}{3.086499in}}%
\pgfpathlineto{\pgfqpoint{4.475018in}{3.074421in}}%
\pgfpathlineto{\pgfqpoint{4.477537in}{3.088496in}}%
\pgfpathlineto{\pgfqpoint{4.480056in}{3.055449in}}%
\pgfpathlineto{\pgfqpoint{4.482574in}{3.073681in}}%
\pgfpathlineto{\pgfqpoint{4.485093in}{3.094946in}}%
\pgfpathlineto{\pgfqpoint{4.487612in}{3.086649in}}%
\pgfpathlineto{\pgfqpoint{4.490131in}{3.055841in}}%
\pgfpathlineto{\pgfqpoint{4.492649in}{3.059402in}}%
\pgfpathlineto{\pgfqpoint{4.495168in}{3.056241in}}%
\pgfpathlineto{\pgfqpoint{4.497687in}{3.042437in}}%
\pgfpathlineto{\pgfqpoint{4.500206in}{3.005025in}}%
\pgfpathlineto{\pgfqpoint{4.502724in}{2.970268in}}%
\pgfpathlineto{\pgfqpoint{4.505243in}{2.975976in}}%
\pgfpathlineto{\pgfqpoint{4.507762in}{3.003883in}}%
\pgfpathlineto{\pgfqpoint{4.510281in}{2.945498in}}%
\pgfpathlineto{\pgfqpoint{4.512799in}{2.983343in}}%
\pgfpathlineto{\pgfqpoint{4.515318in}{2.998047in}}%
\pgfpathlineto{\pgfqpoint{4.517837in}{2.997025in}}%
\pgfpathlineto{\pgfqpoint{4.520356in}{2.969283in}}%
\pgfpathlineto{\pgfqpoint{4.522874in}{2.964774in}}%
\pgfpathlineto{\pgfqpoint{4.525393in}{2.996766in}}%
\pgfpathlineto{\pgfqpoint{4.527912in}{3.043664in}}%
\pgfpathlineto{\pgfqpoint{4.530431in}{3.076145in}}%
\pgfpathlineto{\pgfqpoint{4.532949in}{3.075488in}}%
\pgfpathlineto{\pgfqpoint{4.535468in}{3.056169in}}%
\pgfpathlineto{\pgfqpoint{4.537987in}{3.099740in}}%
\pgfpathlineto{\pgfqpoint{4.540506in}{3.124551in}}%
\pgfpathlineto{\pgfqpoint{4.543024in}{3.129775in}}%
\pgfpathlineto{\pgfqpoint{4.545543in}{3.156664in}}%
\pgfpathlineto{\pgfqpoint{4.548062in}{3.174492in}}%
\pgfpathlineto{\pgfqpoint{4.550581in}{3.193555in}}%
\pgfpathlineto{\pgfqpoint{4.553099in}{3.185255in}}%
\pgfpathlineto{\pgfqpoint{4.555618in}{3.239113in}}%
\pgfpathlineto{\pgfqpoint{4.558137in}{3.209086in}}%
\pgfpathlineto{\pgfqpoint{4.560656in}{3.215698in}}%
\pgfpathlineto{\pgfqpoint{4.563174in}{3.183737in}}%
\pgfpathlineto{\pgfqpoint{4.565693in}{3.184650in}}%
\pgfpathlineto{\pgfqpoint{4.568212in}{3.171923in}}%
\pgfpathlineto{\pgfqpoint{4.570731in}{3.191222in}}%
\pgfpathlineto{\pgfqpoint{4.573249in}{3.185710in}}%
\pgfpathlineto{\pgfqpoint{4.575768in}{3.169026in}}%
\pgfpathlineto{\pgfqpoint{4.578287in}{3.185351in}}%
\pgfpathlineto{\pgfqpoint{4.580806in}{3.186845in}}%
\pgfpathlineto{\pgfqpoint{4.583324in}{3.198858in}}%
\pgfpathlineto{\pgfqpoint{4.585843in}{3.246728in}}%
\pgfpathlineto{\pgfqpoint{4.588362in}{3.217661in}}%
\pgfpathlineto{\pgfqpoint{4.590881in}{3.222690in}}%
\pgfpathlineto{\pgfqpoint{4.593399in}{3.240132in}}%
\pgfpathlineto{\pgfqpoint{4.595918in}{3.212604in}}%
\pgfpathlineto{\pgfqpoint{4.598437in}{3.186722in}}%
\pgfpathlineto{\pgfqpoint{4.600956in}{3.188661in}}%
\pgfpathlineto{\pgfqpoint{4.603474in}{3.207547in}}%
\pgfpathlineto{\pgfqpoint{4.605993in}{3.217187in}}%
\pgfpathlineto{\pgfqpoint{4.608512in}{3.247156in}}%
\pgfpathlineto{\pgfqpoint{4.611031in}{3.242566in}}%
\pgfpathlineto{\pgfqpoint{4.613549in}{3.198410in}}%
\pgfpathlineto{\pgfqpoint{4.616068in}{3.182733in}}%
\pgfpathlineto{\pgfqpoint{4.618587in}{3.183368in}}%
\pgfpathlineto{\pgfqpoint{4.621106in}{3.177588in}}%
\pgfpathlineto{\pgfqpoint{4.623624in}{3.180854in}}%
\pgfpathlineto{\pgfqpoint{4.626143in}{3.139720in}}%
\pgfpathlineto{\pgfqpoint{4.628662in}{3.161266in}}%
\pgfpathlineto{\pgfqpoint{4.631181in}{3.198622in}}%
\pgfpathlineto{\pgfqpoint{4.633699in}{3.219707in}}%
\pgfpathlineto{\pgfqpoint{4.636218in}{3.189788in}}%
\pgfpathlineto{\pgfqpoint{4.638737in}{3.223806in}}%
\pgfpathlineto{\pgfqpoint{4.641256in}{3.225529in}}%
\pgfpathlineto{\pgfqpoint{4.643774in}{3.235407in}}%
\pgfpathlineto{\pgfqpoint{4.646293in}{3.246391in}}%
\pgfpathlineto{\pgfqpoint{4.648812in}{3.269197in}}%
\pgfpathlineto{\pgfqpoint{4.651331in}{3.272023in}}%
\pgfpathlineto{\pgfqpoint{4.653849in}{3.310649in}}%
\pgfpathlineto{\pgfqpoint{4.656368in}{3.329986in}}%
\pgfpathlineto{\pgfqpoint{4.658887in}{3.320395in}}%
\pgfpathlineto{\pgfqpoint{4.661406in}{3.302485in}}%
\pgfpathlineto{\pgfqpoint{4.663924in}{3.321889in}}%
\pgfpathlineto{\pgfqpoint{4.666443in}{3.340377in}}%
\pgfpathlineto{\pgfqpoint{4.668962in}{3.368155in}}%
\pgfpathlineto{\pgfqpoint{4.671481in}{3.370210in}}%
\pgfpathlineto{\pgfqpoint{4.673999in}{3.393469in}}%
\pgfpathlineto{\pgfqpoint{4.676518in}{3.394511in}}%
\pgfpathlineto{\pgfqpoint{4.679037in}{3.425349in}}%
\pgfpathlineto{\pgfqpoint{4.681556in}{3.420160in}}%
\pgfpathlineto{\pgfqpoint{4.684074in}{3.453838in}}%
\pgfpathlineto{\pgfqpoint{4.686593in}{3.453223in}}%
\pgfpathlineto{\pgfqpoint{4.689112in}{3.463499in}}%
\pgfpathlineto{\pgfqpoint{4.691631in}{3.406293in}}%
\pgfpathlineto{\pgfqpoint{4.694149in}{3.413831in}}%
\pgfpathlineto{\pgfqpoint{4.696668in}{3.445005in}}%
\pgfpathlineto{\pgfqpoint{4.699187in}{3.482961in}}%
\pgfpathlineto{\pgfqpoint{4.701706in}{3.499511in}}%
\pgfpathlineto{\pgfqpoint{4.704224in}{3.493452in}}%
\pgfpathlineto{\pgfqpoint{4.706743in}{3.469645in}}%
\pgfpathlineto{\pgfqpoint{4.709262in}{3.485603in}}%
\pgfpathlineto{\pgfqpoint{4.711781in}{3.486444in}}%
\pgfpathlineto{\pgfqpoint{4.714299in}{3.488342in}}%
\pgfpathlineto{\pgfqpoint{4.716818in}{3.510553in}}%
\pgfpathlineto{\pgfqpoint{4.719337in}{3.537859in}}%
\pgfpathlineto{\pgfqpoint{4.721856in}{3.484404in}}%
\pgfpathlineto{\pgfqpoint{4.724374in}{3.497521in}}%
\pgfpathlineto{\pgfqpoint{4.726893in}{3.514716in}}%
\pgfpathlineto{\pgfqpoint{4.729412in}{3.543687in}}%
\pgfpathlineto{\pgfqpoint{4.731931in}{3.609713in}}%
\pgfpathlineto{\pgfqpoint{4.734449in}{3.621454in}}%
\pgfpathlineto{\pgfqpoint{4.736968in}{3.587738in}}%
\pgfpathlineto{\pgfqpoint{4.739487in}{3.597840in}}%
\pgfpathlineto{\pgfqpoint{4.742006in}{3.572365in}}%
\pgfpathlineto{\pgfqpoint{4.744524in}{3.562583in}}%
\pgfpathlineto{\pgfqpoint{4.747043in}{3.555831in}}%
\pgfpathlineto{\pgfqpoint{4.749562in}{3.546836in}}%
\pgfpathlineto{\pgfqpoint{4.752081in}{3.509287in}}%
\pgfpathlineto{\pgfqpoint{4.754599in}{3.537832in}}%
\pgfpathlineto{\pgfqpoint{4.757118in}{3.546525in}}%
\pgfpathlineto{\pgfqpoint{4.759637in}{3.509216in}}%
\pgfpathlineto{\pgfqpoint{4.762156in}{3.531999in}}%
\pgfpathlineto{\pgfqpoint{4.764674in}{3.588494in}}%
\pgfpathlineto{\pgfqpoint{4.767193in}{3.550245in}}%
\pgfpathlineto{\pgfqpoint{4.769712in}{3.526732in}}%
\pgfpathlineto{\pgfqpoint{4.772231in}{3.494172in}}%
\pgfpathlineto{\pgfqpoint{4.774749in}{3.552885in}}%
\pgfpathlineto{\pgfqpoint{4.777268in}{3.570184in}}%
\pgfpathlineto{\pgfqpoint{4.779787in}{3.554283in}}%
\pgfpathlineto{\pgfqpoint{4.782306in}{3.509194in}}%
\pgfpathlineto{\pgfqpoint{4.784824in}{3.534584in}}%
\pgfpathlineto{\pgfqpoint{4.787343in}{3.566476in}}%
\pgfpathlineto{\pgfqpoint{4.789862in}{3.535922in}}%
\pgfpathlineto{\pgfqpoint{4.792381in}{3.550039in}}%
\pgfpathlineto{\pgfqpoint{4.794899in}{3.573744in}}%
\pgfpathlineto{\pgfqpoint{4.797418in}{3.561126in}}%
\pgfpathlineto{\pgfqpoint{4.799937in}{3.620427in}}%
\pgfpathlineto{\pgfqpoint{4.802456in}{3.612067in}}%
\pgfpathlineto{\pgfqpoint{4.804974in}{3.593301in}}%
\pgfpathlineto{\pgfqpoint{4.807493in}{3.553967in}}%
\pgfpathlineto{\pgfqpoint{4.810012in}{3.525000in}}%
\pgfpathlineto{\pgfqpoint{4.812531in}{3.498888in}}%
\pgfpathlineto{\pgfqpoint{4.817568in}{3.543778in}}%
\pgfpathlineto{\pgfqpoint{4.820087in}{3.562098in}}%
\pgfpathlineto{\pgfqpoint{4.822606in}{3.560669in}}%
\pgfpathlineto{\pgfqpoint{4.825124in}{3.589146in}}%
\pgfpathlineto{\pgfqpoint{4.827643in}{3.644330in}}%
\pgfpathlineto{\pgfqpoint{4.830162in}{3.643833in}}%
\pgfpathlineto{\pgfqpoint{4.832681in}{3.603789in}}%
\pgfpathlineto{\pgfqpoint{4.835199in}{3.609666in}}%
\pgfpathlineto{\pgfqpoint{4.837718in}{3.591065in}}%
\pgfpathlineto{\pgfqpoint{4.840237in}{3.553441in}}%
\pgfpathlineto{\pgfqpoint{4.842756in}{3.569632in}}%
\pgfpathlineto{\pgfqpoint{4.845274in}{3.594754in}}%
\pgfpathlineto{\pgfqpoint{4.847793in}{3.601886in}}%
\pgfpathlineto{\pgfqpoint{4.850312in}{3.587962in}}%
\pgfpathlineto{\pgfqpoint{4.852831in}{3.554366in}}%
\pgfpathlineto{\pgfqpoint{4.855349in}{3.528389in}}%
\pgfpathlineto{\pgfqpoint{4.857868in}{3.547145in}}%
\pgfpathlineto{\pgfqpoint{4.860387in}{3.572493in}}%
\pgfpathlineto{\pgfqpoint{4.862906in}{3.551517in}}%
\pgfpathlineto{\pgfqpoint{4.865424in}{3.550913in}}%
\pgfpathlineto{\pgfqpoint{4.867943in}{3.535088in}}%
\pgfpathlineto{\pgfqpoint{4.870462in}{3.556770in}}%
\pgfpathlineto{\pgfqpoint{4.872981in}{3.581148in}}%
\pgfpathlineto{\pgfqpoint{4.875499in}{3.633335in}}%
\pgfpathlineto{\pgfqpoint{4.878018in}{3.622214in}}%
\pgfpathlineto{\pgfqpoint{4.880537in}{3.586288in}}%
\pgfpathlineto{\pgfqpoint{4.883056in}{3.584717in}}%
\pgfpathlineto{\pgfqpoint{4.885574in}{3.585889in}}%
\pgfpathlineto{\pgfqpoint{4.888093in}{3.560834in}}%
\pgfpathlineto{\pgfqpoint{4.890612in}{3.542354in}}%
\pgfpathlineto{\pgfqpoint{4.893131in}{3.584098in}}%
\pgfpathlineto{\pgfqpoint{4.895649in}{3.580218in}}%
\pgfpathlineto{\pgfqpoint{4.898168in}{3.561792in}}%
\pgfpathlineto{\pgfqpoint{4.900687in}{3.619274in}}%
\pgfpathlineto{\pgfqpoint{4.903206in}{3.609656in}}%
\pgfpathlineto{\pgfqpoint{4.905724in}{3.595666in}}%
\pgfpathlineto{\pgfqpoint{4.908243in}{3.614514in}}%
\pgfpathlineto{\pgfqpoint{4.910762in}{3.623044in}}%
\pgfpathlineto{\pgfqpoint{4.913281in}{3.608761in}}%
\pgfpathlineto{\pgfqpoint{4.915799in}{3.641096in}}%
\pgfpathlineto{\pgfqpoint{4.918318in}{3.627322in}}%
\pgfpathlineto{\pgfqpoint{4.920837in}{3.659564in}}%
\pgfpathlineto{\pgfqpoint{4.923356in}{3.700022in}}%
\pgfpathlineto{\pgfqpoint{4.925874in}{3.723368in}}%
\pgfpathlineto{\pgfqpoint{4.928393in}{3.698497in}}%
\pgfpathlineto{\pgfqpoint{4.930912in}{3.703739in}}%
\pgfpathlineto{\pgfqpoint{4.933431in}{3.734654in}}%
\pgfpathlineto{\pgfqpoint{4.935949in}{3.725493in}}%
\pgfpathlineto{\pgfqpoint{4.938468in}{3.731251in}}%
\pgfpathlineto{\pgfqpoint{4.940987in}{3.734468in}}%
\pgfpathlineto{\pgfqpoint{4.943506in}{3.701266in}}%
\pgfpathlineto{\pgfqpoint{4.946024in}{3.706722in}}%
\pgfpathlineto{\pgfqpoint{4.948543in}{3.759102in}}%
\pgfpathlineto{\pgfqpoint{4.951062in}{3.758024in}}%
\pgfpathlineto{\pgfqpoint{4.953581in}{3.781098in}}%
\pgfpathlineto{\pgfqpoint{4.956099in}{3.824840in}}%
\pgfpathlineto{\pgfqpoint{4.958618in}{3.809220in}}%
\pgfpathlineto{\pgfqpoint{4.961137in}{3.811013in}}%
\pgfpathlineto{\pgfqpoint{4.963656in}{3.835460in}}%
\pgfpathlineto{\pgfqpoint{4.966174in}{3.854831in}}%
\pgfpathlineto{\pgfqpoint{4.968693in}{3.882879in}}%
\pgfpathlineto{\pgfqpoint{4.971212in}{3.887300in}}%
\pgfpathlineto{\pgfqpoint{4.973731in}{3.874776in}}%
\pgfpathlineto{\pgfqpoint{4.976249in}{3.868799in}}%
\pgfpathlineto{\pgfqpoint{4.978768in}{3.896524in}}%
\pgfpathlineto{\pgfqpoint{4.981287in}{3.909953in}}%
\pgfpathlineto{\pgfqpoint{4.983806in}{3.922438in}}%
\pgfpathlineto{\pgfqpoint{4.986324in}{3.936794in}}%
\pgfpathlineto{\pgfqpoint{4.988843in}{3.966426in}}%
\pgfpathlineto{\pgfqpoint{4.991362in}{3.948445in}}%
\pgfpathlineto{\pgfqpoint{4.993881in}{3.919789in}}%
\pgfpathlineto{\pgfqpoint{4.996399in}{3.932943in}}%
\pgfpathlineto{\pgfqpoint{4.998918in}{3.968358in}}%
\pgfpathlineto{\pgfqpoint{5.001437in}{3.985052in}}%
\pgfpathlineto{\pgfqpoint{5.003956in}{3.970820in}}%
\pgfpathlineto{\pgfqpoint{5.006474in}{4.016517in}}%
\pgfpathlineto{\pgfqpoint{5.008993in}{4.001273in}}%
\pgfpathlineto{\pgfqpoint{5.011512in}{3.998406in}}%
\pgfpathlineto{\pgfqpoint{5.014031in}{3.977532in}}%
\pgfpathlineto{\pgfqpoint{5.016549in}{3.976084in}}%
\pgfpathlineto{\pgfqpoint{5.019068in}{3.952382in}}%
\pgfpathlineto{\pgfqpoint{5.021587in}{3.901504in}}%
\pgfpathlineto{\pgfqpoint{5.024106in}{3.875343in}}%
\pgfpathlineto{\pgfqpoint{5.026624in}{3.921594in}}%
\pgfpathlineto{\pgfqpoint{5.029143in}{3.854292in}}%
\pgfpathlineto{\pgfqpoint{5.031662in}{3.836841in}}%
\pgfpathlineto{\pgfqpoint{5.034181in}{3.855545in}}%
\pgfpathlineto{\pgfqpoint{5.036699in}{3.850238in}}%
\pgfpathlineto{\pgfqpoint{5.039218in}{3.805769in}}%
\pgfpathlineto{\pgfqpoint{5.041737in}{3.771986in}}%
\pgfpathlineto{\pgfqpoint{5.044256in}{3.763644in}}%
\pgfpathlineto{\pgfqpoint{5.046774in}{3.769234in}}%
\pgfpathlineto{\pgfqpoint{5.049293in}{3.757770in}}%
\pgfpathlineto{\pgfqpoint{5.051812in}{3.785024in}}%
\pgfpathlineto{\pgfqpoint{5.054331in}{3.809383in}}%
\pgfpathlineto{\pgfqpoint{5.056849in}{3.803442in}}%
\pgfpathlineto{\pgfqpoint{5.059368in}{3.792715in}}%
\pgfpathlineto{\pgfqpoint{5.061887in}{3.781184in}}%
\pgfpathlineto{\pgfqpoint{5.064406in}{3.784999in}}%
\pgfpathlineto{\pgfqpoint{5.066924in}{3.788168in}}%
\pgfpathlineto{\pgfqpoint{5.069443in}{3.781469in}}%
\pgfpathlineto{\pgfqpoint{5.071962in}{3.742168in}}%
\pgfpathlineto{\pgfqpoint{5.074481in}{3.767205in}}%
\pgfpathlineto{\pgfqpoint{5.076999in}{3.764379in}}%
\pgfpathlineto{\pgfqpoint{5.079518in}{3.746155in}}%
\pgfpathlineto{\pgfqpoint{5.082037in}{3.754585in}}%
\pgfpathlineto{\pgfqpoint{5.084556in}{3.760513in}}%
\pgfpathlineto{\pgfqpoint{5.087074in}{3.772892in}}%
\pgfpathlineto{\pgfqpoint{5.089593in}{3.741561in}}%
\pgfpathlineto{\pgfqpoint{5.092112in}{3.704053in}}%
\pgfpathlineto{\pgfqpoint{5.094631in}{3.717386in}}%
\pgfpathlineto{\pgfqpoint{5.097149in}{3.702601in}}%
\pgfpathlineto{\pgfqpoint{5.099668in}{3.689728in}}%
\pgfpathlineto{\pgfqpoint{5.102187in}{3.696472in}}%
\pgfpathlineto{\pgfqpoint{5.104706in}{3.675640in}}%
\pgfpathlineto{\pgfqpoint{5.107224in}{3.713386in}}%
\pgfpathlineto{\pgfqpoint{5.109743in}{3.717544in}}%
\pgfpathlineto{\pgfqpoint{5.112262in}{3.750756in}}%
\pgfpathlineto{\pgfqpoint{5.114781in}{3.734488in}}%
\pgfpathlineto{\pgfqpoint{5.117299in}{3.745793in}}%
\pgfpathlineto{\pgfqpoint{5.119818in}{3.745973in}}%
\pgfpathlineto{\pgfqpoint{5.122337in}{3.766062in}}%
\pgfpathlineto{\pgfqpoint{5.124856in}{3.770650in}}%
\pgfpathlineto{\pgfqpoint{5.127374in}{3.820390in}}%
\pgfpathlineto{\pgfqpoint{5.129893in}{3.829279in}}%
\pgfpathlineto{\pgfqpoint{5.132412in}{3.875390in}}%
\pgfpathlineto{\pgfqpoint{5.134931in}{3.864892in}}%
\pgfpathlineto{\pgfqpoint{5.137449in}{3.882892in}}%
\pgfpathlineto{\pgfqpoint{5.139968in}{3.884715in}}%
\pgfpathlineto{\pgfqpoint{5.142487in}{3.834899in}}%
\pgfpathlineto{\pgfqpoint{5.145006in}{3.811138in}}%
\pgfpathlineto{\pgfqpoint{5.147524in}{3.772178in}}%
\pgfpathlineto{\pgfqpoint{5.150043in}{3.746889in}}%
\pgfpathlineto{\pgfqpoint{5.152562in}{3.785310in}}%
\pgfpathlineto{\pgfqpoint{5.155081in}{3.779145in}}%
\pgfpathlineto{\pgfqpoint{5.157599in}{3.816871in}}%
\pgfpathlineto{\pgfqpoint{5.160118in}{3.817597in}}%
\pgfpathlineto{\pgfqpoint{5.162637in}{3.822608in}}%
\pgfpathlineto{\pgfqpoint{5.165156in}{3.830012in}}%
\pgfpathlineto{\pgfqpoint{5.167674in}{3.830593in}}%
\pgfpathlineto{\pgfqpoint{5.170193in}{3.828282in}}%
\pgfpathlineto{\pgfqpoint{5.172712in}{3.795396in}}%
\pgfpathlineto{\pgfqpoint{5.175231in}{3.790798in}}%
\pgfpathlineto{\pgfqpoint{5.177749in}{3.771790in}}%
\pgfpathlineto{\pgfqpoint{5.180268in}{3.785824in}}%
\pgfpathlineto{\pgfqpoint{5.182787in}{3.787843in}}%
\pgfpathlineto{\pgfqpoint{5.185306in}{3.800469in}}%
\pgfpathlineto{\pgfqpoint{5.187824in}{3.814499in}}%
\pgfpathlineto{\pgfqpoint{5.192862in}{3.823283in}}%
\pgfpathlineto{\pgfqpoint{5.195381in}{3.869345in}}%
\pgfpathlineto{\pgfqpoint{5.197899in}{3.878884in}}%
\pgfpathlineto{\pgfqpoint{5.200418in}{3.889363in}}%
\pgfpathlineto{\pgfqpoint{5.202937in}{3.906879in}}%
\pgfpathlineto{\pgfqpoint{5.205456in}{3.886714in}}%
\pgfpathlineto{\pgfqpoint{5.207974in}{3.833258in}}%
\pgfpathlineto{\pgfqpoint{5.210493in}{3.817979in}}%
\pgfpathlineto{\pgfqpoint{5.213012in}{3.819605in}}%
\pgfpathlineto{\pgfqpoint{5.215531in}{3.799963in}}%
\pgfpathlineto{\pgfqpoint{5.218049in}{3.835673in}}%
\pgfpathlineto{\pgfqpoint{5.220568in}{3.855703in}}%
\pgfpathlineto{\pgfqpoint{5.223087in}{3.874820in}}%
\pgfpathlineto{\pgfqpoint{5.225606in}{3.910009in}}%
\pgfpathlineto{\pgfqpoint{5.228124in}{3.929681in}}%
\pgfpathlineto{\pgfqpoint{5.230643in}{3.915406in}}%
\pgfpathlineto{\pgfqpoint{5.233162in}{3.904142in}}%
\pgfpathlineto{\pgfqpoint{5.235681in}{3.934736in}}%
\pgfpathlineto{\pgfqpoint{5.238199in}{3.950506in}}%
\pgfpathlineto{\pgfqpoint{5.240718in}{3.942084in}}%
\pgfpathlineto{\pgfqpoint{5.243237in}{3.938686in}}%
\pgfpathlineto{\pgfqpoint{5.245756in}{3.939431in}}%
\pgfpathlineto{\pgfqpoint{5.248274in}{4.009553in}}%
\pgfpathlineto{\pgfqpoint{5.253312in}{4.043213in}}%
\pgfpathlineto{\pgfqpoint{5.255831in}{4.055965in}}%
\pgfpathlineto{\pgfqpoint{5.258349in}{4.057071in}}%
\pgfpathlineto{\pgfqpoint{5.260868in}{4.036396in}}%
\pgfpathlineto{\pgfqpoint{5.263387in}{3.995994in}}%
\pgfpathlineto{\pgfqpoint{5.265906in}{4.016193in}}%
\pgfpathlineto{\pgfqpoint{5.268424in}{4.019051in}}%
\pgfpathlineto{\pgfqpoint{5.270943in}{4.030712in}}%
\pgfpathlineto{\pgfqpoint{5.273462in}{4.022579in}}%
\pgfpathlineto{\pgfqpoint{5.275981in}{4.063485in}}%
\pgfpathlineto{\pgfqpoint{5.278499in}{4.032187in}}%
\pgfpathlineto{\pgfqpoint{5.281018in}{4.052006in}}%
\pgfpathlineto{\pgfqpoint{5.283537in}{4.041305in}}%
\pgfpathlineto{\pgfqpoint{5.286056in}{4.058043in}}%
\pgfpathlineto{\pgfqpoint{5.288574in}{3.998114in}}%
\pgfpathlineto{\pgfqpoint{5.291093in}{3.979782in}}%
\pgfpathlineto{\pgfqpoint{5.293612in}{3.973425in}}%
\pgfpathlineto{\pgfqpoint{5.296131in}{3.982914in}}%
\pgfpathlineto{\pgfqpoint{5.298649in}{3.964832in}}%
\pgfpathlineto{\pgfqpoint{5.301168in}{3.959193in}}%
\pgfpathlineto{\pgfqpoint{5.303687in}{3.939618in}}%
\pgfpathlineto{\pgfqpoint{5.306206in}{3.951801in}}%
\pgfpathlineto{\pgfqpoint{5.308724in}{3.982081in}}%
\pgfpathlineto{\pgfqpoint{5.311243in}{3.998989in}}%
\pgfpathlineto{\pgfqpoint{5.313762in}{4.074015in}}%
\pgfpathlineto{\pgfqpoint{5.316281in}{4.067725in}}%
\pgfpathlineto{\pgfqpoint{5.318799in}{4.043562in}}%
\pgfpathlineto{\pgfqpoint{5.321318in}{4.085369in}}%
\pgfpathlineto{\pgfqpoint{5.323837in}{4.074176in}}%
\pgfpathlineto{\pgfqpoint{5.326356in}{4.060631in}}%
\pgfpathlineto{\pgfqpoint{5.328874in}{4.084526in}}%
\pgfpathlineto{\pgfqpoint{5.331393in}{4.079223in}}%
\pgfpathlineto{\pgfqpoint{5.333912in}{4.066517in}}%
\pgfpathlineto{\pgfqpoint{5.336431in}{4.034172in}}%
\pgfpathlineto{\pgfqpoint{5.338949in}{4.033452in}}%
\pgfpathlineto{\pgfqpoint{5.341468in}{4.036171in}}%
\pgfpathlineto{\pgfqpoint{5.343987in}{4.028027in}}%
\pgfpathlineto{\pgfqpoint{5.346506in}{3.996521in}}%
\pgfpathlineto{\pgfqpoint{5.349024in}{3.977580in}}%
\pgfpathlineto{\pgfqpoint{5.351543in}{3.995486in}}%
\pgfpathlineto{\pgfqpoint{5.354062in}{3.992136in}}%
\pgfpathlineto{\pgfqpoint{5.356581in}{4.038894in}}%
\pgfpathlineto{\pgfqpoint{5.359099in}{4.064254in}}%
\pgfpathlineto{\pgfqpoint{5.361618in}{4.040506in}}%
\pgfpathlineto{\pgfqpoint{5.364137in}{4.007868in}}%
\pgfpathlineto{\pgfqpoint{5.366656in}{4.001684in}}%
\pgfpathlineto{\pgfqpoint{5.369174in}{3.989127in}}%
\pgfpathlineto{\pgfqpoint{5.371693in}{4.010971in}}%
\pgfpathlineto{\pgfqpoint{5.374212in}{3.981600in}}%
\pgfpathlineto{\pgfqpoint{5.376731in}{3.972156in}}%
\pgfpathlineto{\pgfqpoint{5.379249in}{4.036283in}}%
\pgfpathlineto{\pgfqpoint{5.381768in}{4.049478in}}%
\pgfpathlineto{\pgfqpoint{5.384287in}{3.993140in}}%
\pgfpathlineto{\pgfqpoint{5.386806in}{4.006063in}}%
\pgfpathlineto{\pgfqpoint{5.389324in}{4.045057in}}%
\pgfpathlineto{\pgfqpoint{5.391843in}{4.061877in}}%
\pgfpathlineto{\pgfqpoint{5.394362in}{4.041706in}}%
\pgfpathlineto{\pgfqpoint{5.396881in}{4.074054in}}%
\pgfpathlineto{\pgfqpoint{5.399399in}{4.093512in}}%
\pgfpathlineto{\pgfqpoint{5.401918in}{4.070958in}}%
\pgfpathlineto{\pgfqpoint{5.404437in}{4.059607in}}%
\pgfpathlineto{\pgfqpoint{5.406956in}{4.057166in}}%
\pgfpathlineto{\pgfqpoint{5.409474in}{4.031814in}}%
\pgfpathlineto{\pgfqpoint{5.411993in}{4.019893in}}%
\pgfpathlineto{\pgfqpoint{5.414512in}{4.055750in}}%
\pgfpathlineto{\pgfqpoint{5.417031in}{4.003043in}}%
\pgfpathlineto{\pgfqpoint{5.419549in}{4.000208in}}%
\pgfpathlineto{\pgfqpoint{5.422068in}{3.990361in}}%
\pgfpathlineto{\pgfqpoint{5.424587in}{3.993795in}}%
\pgfpathlineto{\pgfqpoint{5.427106in}{4.022397in}}%
\pgfpathlineto{\pgfqpoint{5.429624in}{4.035990in}}%
\pgfpathlineto{\pgfqpoint{5.432143in}{4.021319in}}%
\pgfpathlineto{\pgfqpoint{5.434662in}{3.989062in}}%
\pgfpathlineto{\pgfqpoint{5.437181in}{4.027407in}}%
\pgfpathlineto{\pgfqpoint{5.439699in}{3.974349in}}%
\pgfpathlineto{\pgfqpoint{5.442218in}{3.947431in}}%
\pgfpathlineto{\pgfqpoint{5.444737in}{3.982370in}}%
\pgfpathlineto{\pgfqpoint{5.447256in}{3.977303in}}%
\pgfpathlineto{\pgfqpoint{5.449774in}{3.991557in}}%
\pgfpathlineto{\pgfqpoint{5.452293in}{3.991906in}}%
\pgfpathlineto{\pgfqpoint{5.454812in}{3.962639in}}%
\pgfpathlineto{\pgfqpoint{5.457331in}{3.973038in}}%
\pgfpathlineto{\pgfqpoint{5.459849in}{3.972019in}}%
\pgfpathlineto{\pgfqpoint{5.462368in}{3.946382in}}%
\pgfpathlineto{\pgfqpoint{5.464887in}{3.961370in}}%
\pgfpathlineto{\pgfqpoint{5.467406in}{3.964386in}}%
\pgfpathlineto{\pgfqpoint{5.469924in}{3.930229in}}%
\pgfpathlineto{\pgfqpoint{5.472443in}{3.928179in}}%
\pgfpathlineto{\pgfqpoint{5.474962in}{3.930957in}}%
\pgfpathlineto{\pgfqpoint{5.477481in}{3.921266in}}%
\pgfpathlineto{\pgfqpoint{5.479999in}{3.963536in}}%
\pgfpathlineto{\pgfqpoint{5.482518in}{3.989740in}}%
\pgfpathlineto{\pgfqpoint{5.485037in}{3.995496in}}%
\pgfpathlineto{\pgfqpoint{5.487556in}{3.991060in}}%
\pgfpathlineto{\pgfqpoint{5.490074in}{3.965184in}}%
\pgfpathlineto{\pgfqpoint{5.492593in}{3.930402in}}%
\pgfpathlineto{\pgfqpoint{5.492593in}{3.930402in}}%
\pgfusepath{stroke}%
\end{pgfscope}%
\begin{pgfscope}%
\pgfpathrectangle{\pgfqpoint{0.455093in}{0.401080in}}{\pgfqpoint{5.037500in}{4.004000in}}%
\pgfusepath{clip}%
\pgfsetrectcap%
\pgfsetroundjoin%
\pgfsetlinewidth{0.501875pt}%
\definecolor{currentstroke}{rgb}{0.690196,0.188235,0.376471}%
\pgfsetstrokecolor{currentstroke}%
\pgfsetstrokeopacity{0.800000}%
\pgfsetdash{}{0pt}%
\pgfpathmoveto{\pgfqpoint{0.455093in}{1.021044in}}%
\pgfpathlineto{\pgfqpoint{0.457612in}{1.025855in}}%
\pgfpathlineto{\pgfqpoint{0.460131in}{1.015650in}}%
\pgfpathlineto{\pgfqpoint{0.465168in}{1.003969in}}%
\pgfpathlineto{\pgfqpoint{0.467687in}{1.006777in}}%
\pgfpathlineto{\pgfqpoint{0.470206in}{1.017811in}}%
\pgfpathlineto{\pgfqpoint{0.472724in}{1.012590in}}%
\pgfpathlineto{\pgfqpoint{0.475243in}{0.991381in}}%
\pgfpathlineto{\pgfqpoint{0.477762in}{0.984893in}}%
\pgfpathlineto{\pgfqpoint{0.480281in}{0.982454in}}%
\pgfpathlineto{\pgfqpoint{0.482799in}{0.965587in}}%
\pgfpathlineto{\pgfqpoint{0.485318in}{0.996122in}}%
\pgfpathlineto{\pgfqpoint{0.487837in}{0.991385in}}%
\pgfpathlineto{\pgfqpoint{0.490356in}{0.960261in}}%
\pgfpathlineto{\pgfqpoint{0.492874in}{0.953064in}}%
\pgfpathlineto{\pgfqpoint{0.495393in}{0.984363in}}%
\pgfpathlineto{\pgfqpoint{0.497912in}{0.964836in}}%
\pgfpathlineto{\pgfqpoint{0.500431in}{0.971418in}}%
\pgfpathlineto{\pgfqpoint{0.502949in}{0.991796in}}%
\pgfpathlineto{\pgfqpoint{0.505468in}{1.000596in}}%
\pgfpathlineto{\pgfqpoint{0.507987in}{0.998640in}}%
\pgfpathlineto{\pgfqpoint{0.510506in}{1.001984in}}%
\pgfpathlineto{\pgfqpoint{0.513024in}{0.998218in}}%
\pgfpathlineto{\pgfqpoint{0.515543in}{0.999161in}}%
\pgfpathlineto{\pgfqpoint{0.518062in}{0.989098in}}%
\pgfpathlineto{\pgfqpoint{0.520581in}{0.992292in}}%
\pgfpathlineto{\pgfqpoint{0.523099in}{0.987462in}}%
\pgfpathlineto{\pgfqpoint{0.525618in}{0.988303in}}%
\pgfpathlineto{\pgfqpoint{0.528137in}{0.966197in}}%
\pgfpathlineto{\pgfqpoint{0.530656in}{0.973560in}}%
\pgfpathlineto{\pgfqpoint{0.533174in}{0.960917in}}%
\pgfpathlineto{\pgfqpoint{0.535693in}{0.962948in}}%
\pgfpathlineto{\pgfqpoint{0.538212in}{0.967752in}}%
\pgfpathlineto{\pgfqpoint{0.540731in}{0.958894in}}%
\pgfpathlineto{\pgfqpoint{0.543249in}{0.947236in}}%
\pgfpathlineto{\pgfqpoint{0.545768in}{0.941848in}}%
\pgfpathlineto{\pgfqpoint{0.548287in}{0.932772in}}%
\pgfpathlineto{\pgfqpoint{0.550806in}{0.931812in}}%
\pgfpathlineto{\pgfqpoint{0.553324in}{0.942114in}}%
\pgfpathlineto{\pgfqpoint{0.555843in}{0.945722in}}%
\pgfpathlineto{\pgfqpoint{0.558362in}{0.946366in}}%
\pgfpathlineto{\pgfqpoint{0.560881in}{0.941405in}}%
\pgfpathlineto{\pgfqpoint{0.565918in}{0.946613in}}%
\pgfpathlineto{\pgfqpoint{0.568437in}{0.944028in}}%
\pgfpathlineto{\pgfqpoint{0.570956in}{0.960878in}}%
\pgfpathlineto{\pgfqpoint{0.573474in}{0.962077in}}%
\pgfpathlineto{\pgfqpoint{0.575993in}{0.968962in}}%
\pgfpathlineto{\pgfqpoint{0.578512in}{0.982217in}}%
\pgfpathlineto{\pgfqpoint{0.581031in}{0.973539in}}%
\pgfpathlineto{\pgfqpoint{0.583549in}{0.972623in}}%
\pgfpathlineto{\pgfqpoint{0.586068in}{0.963141in}}%
\pgfpathlineto{\pgfqpoint{0.588587in}{0.974265in}}%
\pgfpathlineto{\pgfqpoint{0.591106in}{0.972092in}}%
\pgfpathlineto{\pgfqpoint{0.593624in}{0.983110in}}%
\pgfpathlineto{\pgfqpoint{0.596143in}{0.998345in}}%
\pgfpathlineto{\pgfqpoint{0.598662in}{0.999592in}}%
\pgfpathlineto{\pgfqpoint{0.601181in}{0.972701in}}%
\pgfpathlineto{\pgfqpoint{0.603699in}{0.973426in}}%
\pgfpathlineto{\pgfqpoint{0.606218in}{0.979633in}}%
\pgfpathlineto{\pgfqpoint{0.608737in}{0.976681in}}%
\pgfpathlineto{\pgfqpoint{0.611256in}{0.969297in}}%
\pgfpathlineto{\pgfqpoint{0.613774in}{0.990688in}}%
\pgfpathlineto{\pgfqpoint{0.616293in}{0.990862in}}%
\pgfpathlineto{\pgfqpoint{0.618812in}{0.995184in}}%
\pgfpathlineto{\pgfqpoint{0.621331in}{0.991294in}}%
\pgfpathlineto{\pgfqpoint{0.623849in}{0.992654in}}%
\pgfpathlineto{\pgfqpoint{0.626368in}{0.989740in}}%
\pgfpathlineto{\pgfqpoint{0.628887in}{0.966587in}}%
\pgfpathlineto{\pgfqpoint{0.631406in}{0.959877in}}%
\pgfpathlineto{\pgfqpoint{0.633924in}{0.969402in}}%
\pgfpathlineto{\pgfqpoint{0.636443in}{0.991471in}}%
\pgfpathlineto{\pgfqpoint{0.638962in}{0.994876in}}%
\pgfpathlineto{\pgfqpoint{0.641481in}{0.972590in}}%
\pgfpathlineto{\pgfqpoint{0.643999in}{0.966622in}}%
\pgfpathlineto{\pgfqpoint{0.646518in}{0.967503in}}%
\pgfpathlineto{\pgfqpoint{0.649037in}{0.966621in}}%
\pgfpathlineto{\pgfqpoint{0.651556in}{0.972405in}}%
\pgfpathlineto{\pgfqpoint{0.654074in}{0.966301in}}%
\pgfpathlineto{\pgfqpoint{0.656593in}{0.974008in}}%
\pgfpathlineto{\pgfqpoint{0.659112in}{0.972564in}}%
\pgfpathlineto{\pgfqpoint{0.661631in}{0.965749in}}%
\pgfpathlineto{\pgfqpoint{0.664149in}{0.969750in}}%
\pgfpathlineto{\pgfqpoint{0.666668in}{0.977132in}}%
\pgfpathlineto{\pgfqpoint{0.669187in}{0.949790in}}%
\pgfpathlineto{\pgfqpoint{0.671706in}{0.943952in}}%
\pgfpathlineto{\pgfqpoint{0.674224in}{0.937107in}}%
\pgfpathlineto{\pgfqpoint{0.676743in}{0.920092in}}%
\pgfpathlineto{\pgfqpoint{0.679262in}{0.904735in}}%
\pgfpathlineto{\pgfqpoint{0.684299in}{0.885421in}}%
\pgfpathlineto{\pgfqpoint{0.686818in}{0.876392in}}%
\pgfpathlineto{\pgfqpoint{0.689337in}{0.876205in}}%
\pgfpathlineto{\pgfqpoint{0.691856in}{0.872287in}}%
\pgfpathlineto{\pgfqpoint{0.694374in}{0.879498in}}%
\pgfpathlineto{\pgfqpoint{0.696893in}{0.887089in}}%
\pgfpathlineto{\pgfqpoint{0.699412in}{0.877233in}}%
\pgfpathlineto{\pgfqpoint{0.701931in}{0.885804in}}%
\pgfpathlineto{\pgfqpoint{0.704449in}{0.879246in}}%
\pgfpathlineto{\pgfqpoint{0.706968in}{0.871565in}}%
\pgfpathlineto{\pgfqpoint{0.709487in}{0.858428in}}%
\pgfpathlineto{\pgfqpoint{0.712006in}{0.862168in}}%
\pgfpathlineto{\pgfqpoint{0.714524in}{0.856689in}}%
\pgfpathlineto{\pgfqpoint{0.717043in}{0.869329in}}%
\pgfpathlineto{\pgfqpoint{0.719562in}{0.872660in}}%
\pgfpathlineto{\pgfqpoint{0.722081in}{0.874361in}}%
\pgfpathlineto{\pgfqpoint{0.724599in}{0.851798in}}%
\pgfpathlineto{\pgfqpoint{0.727118in}{0.838750in}}%
\pgfpathlineto{\pgfqpoint{0.729637in}{0.856768in}}%
\pgfpathlineto{\pgfqpoint{0.732156in}{0.856947in}}%
\pgfpathlineto{\pgfqpoint{0.734674in}{0.833799in}}%
\pgfpathlineto{\pgfqpoint{0.737193in}{0.820999in}}%
\pgfpathlineto{\pgfqpoint{0.739712in}{0.816066in}}%
\pgfpathlineto{\pgfqpoint{0.742231in}{0.812498in}}%
\pgfpathlineto{\pgfqpoint{0.744749in}{0.803193in}}%
\pgfpathlineto{\pgfqpoint{0.747268in}{0.778388in}}%
\pgfpathlineto{\pgfqpoint{0.749787in}{0.781843in}}%
\pgfpathlineto{\pgfqpoint{0.752306in}{0.783232in}}%
\pgfpathlineto{\pgfqpoint{0.754824in}{0.770082in}}%
\pgfpathlineto{\pgfqpoint{0.757343in}{0.777528in}}%
\pgfpathlineto{\pgfqpoint{0.759862in}{0.767782in}}%
\pgfpathlineto{\pgfqpoint{0.762381in}{0.760968in}}%
\pgfpathlineto{\pgfqpoint{0.764899in}{0.775024in}}%
\pgfpathlineto{\pgfqpoint{0.767418in}{0.757075in}}%
\pgfpathlineto{\pgfqpoint{0.769937in}{0.753137in}}%
\pgfpathlineto{\pgfqpoint{0.772456in}{0.762380in}}%
\pgfpathlineto{\pgfqpoint{0.774974in}{0.784198in}}%
\pgfpathlineto{\pgfqpoint{0.777493in}{0.791057in}}%
\pgfpathlineto{\pgfqpoint{0.780012in}{0.791513in}}%
\pgfpathlineto{\pgfqpoint{0.782531in}{0.803328in}}%
\pgfpathlineto{\pgfqpoint{0.785049in}{0.805437in}}%
\pgfpathlineto{\pgfqpoint{0.787568in}{0.787832in}}%
\pgfpathlineto{\pgfqpoint{0.790087in}{0.801301in}}%
\pgfpathlineto{\pgfqpoint{0.792606in}{0.818807in}}%
\pgfpathlineto{\pgfqpoint{0.795124in}{0.831796in}}%
\pgfpathlineto{\pgfqpoint{0.797643in}{0.826760in}}%
\pgfpathlineto{\pgfqpoint{0.800162in}{0.814502in}}%
\pgfpathlineto{\pgfqpoint{0.802681in}{0.796528in}}%
\pgfpathlineto{\pgfqpoint{0.805199in}{0.810593in}}%
\pgfpathlineto{\pgfqpoint{0.807718in}{0.813644in}}%
\pgfpathlineto{\pgfqpoint{0.810237in}{0.797797in}}%
\pgfpathlineto{\pgfqpoint{0.812756in}{0.819328in}}%
\pgfpathlineto{\pgfqpoint{0.815274in}{0.833839in}}%
\pgfpathlineto{\pgfqpoint{0.817793in}{0.836704in}}%
\pgfpathlineto{\pgfqpoint{0.820312in}{0.821592in}}%
\pgfpathlineto{\pgfqpoint{0.822831in}{0.828217in}}%
\pgfpathlineto{\pgfqpoint{0.825349in}{0.835315in}}%
\pgfpathlineto{\pgfqpoint{0.827868in}{0.832251in}}%
\pgfpathlineto{\pgfqpoint{0.832906in}{0.800830in}}%
\pgfpathlineto{\pgfqpoint{0.835424in}{0.790463in}}%
\pgfpathlineto{\pgfqpoint{0.837943in}{0.772822in}}%
\pgfpathlineto{\pgfqpoint{0.840462in}{0.784568in}}%
\pgfpathlineto{\pgfqpoint{0.842981in}{0.784712in}}%
\pgfpathlineto{\pgfqpoint{0.845499in}{0.744246in}}%
\pgfpathlineto{\pgfqpoint{0.848018in}{0.741832in}}%
\pgfpathlineto{\pgfqpoint{0.850537in}{0.734756in}}%
\pgfpathlineto{\pgfqpoint{0.853056in}{0.738741in}}%
\pgfpathlineto{\pgfqpoint{0.855574in}{0.759012in}}%
\pgfpathlineto{\pgfqpoint{0.858093in}{0.769522in}}%
\pgfpathlineto{\pgfqpoint{0.860612in}{0.742958in}}%
\pgfpathlineto{\pgfqpoint{0.863131in}{0.728473in}}%
\pgfpathlineto{\pgfqpoint{0.865649in}{0.731695in}}%
\pgfpathlineto{\pgfqpoint{0.868168in}{0.723861in}}%
\pgfpathlineto{\pgfqpoint{0.870687in}{0.718166in}}%
\pgfpathlineto{\pgfqpoint{0.873206in}{0.699157in}}%
\pgfpathlineto{\pgfqpoint{0.875724in}{0.703378in}}%
\pgfpathlineto{\pgfqpoint{0.878243in}{0.719641in}}%
\pgfpathlineto{\pgfqpoint{0.880762in}{0.723212in}}%
\pgfpathlineto{\pgfqpoint{0.883281in}{0.710206in}}%
\pgfpathlineto{\pgfqpoint{0.885799in}{0.720861in}}%
\pgfpathlineto{\pgfqpoint{0.888318in}{0.706497in}}%
\pgfpathlineto{\pgfqpoint{0.890837in}{0.720940in}}%
\pgfpathlineto{\pgfqpoint{0.893356in}{0.714757in}}%
\pgfpathlineto{\pgfqpoint{0.895874in}{0.697891in}}%
\pgfpathlineto{\pgfqpoint{0.898393in}{0.691352in}}%
\pgfpathlineto{\pgfqpoint{0.900912in}{0.669121in}}%
\pgfpathlineto{\pgfqpoint{0.903431in}{0.667896in}}%
\pgfpathlineto{\pgfqpoint{0.905949in}{0.695096in}}%
\pgfpathlineto{\pgfqpoint{0.908468in}{0.694136in}}%
\pgfpathlineto{\pgfqpoint{0.910987in}{0.704669in}}%
\pgfpathlineto{\pgfqpoint{0.913506in}{0.684807in}}%
\pgfpathlineto{\pgfqpoint{0.916024in}{0.671372in}}%
\pgfpathlineto{\pgfqpoint{0.918543in}{0.682319in}}%
\pgfpathlineto{\pgfqpoint{0.921062in}{0.683882in}}%
\pgfpathlineto{\pgfqpoint{0.923581in}{0.658598in}}%
\pgfpathlineto{\pgfqpoint{0.926099in}{0.668740in}}%
\pgfpathlineto{\pgfqpoint{0.928618in}{0.670582in}}%
\pgfpathlineto{\pgfqpoint{0.933656in}{0.628079in}}%
\pgfpathlineto{\pgfqpoint{0.936174in}{0.611054in}}%
\pgfpathlineto{\pgfqpoint{0.938693in}{0.626583in}}%
\pgfpathlineto{\pgfqpoint{0.941212in}{0.625220in}}%
\pgfpathlineto{\pgfqpoint{0.943731in}{0.625579in}}%
\pgfpathlineto{\pgfqpoint{0.946249in}{0.616545in}}%
\pgfpathlineto{\pgfqpoint{0.948768in}{0.594709in}}%
\pgfpathlineto{\pgfqpoint{0.951287in}{0.583080in}}%
\pgfpathlineto{\pgfqpoint{0.953806in}{0.595068in}}%
\pgfpathlineto{\pgfqpoint{0.956324in}{0.613318in}}%
\pgfpathlineto{\pgfqpoint{0.958843in}{0.609482in}}%
\pgfpathlineto{\pgfqpoint{0.961362in}{0.604725in}}%
\pgfpathlineto{\pgfqpoint{0.963881in}{0.587257in}}%
\pgfpathlineto{\pgfqpoint{0.966399in}{0.589781in}}%
\pgfpathlineto{\pgfqpoint{0.968918in}{0.592598in}}%
\pgfpathlineto{\pgfqpoint{0.971437in}{0.585193in}}%
\pgfpathlineto{\pgfqpoint{0.973956in}{0.589425in}}%
\pgfpathlineto{\pgfqpoint{0.976474in}{0.590236in}}%
\pgfpathlineto{\pgfqpoint{0.978993in}{0.602298in}}%
\pgfpathlineto{\pgfqpoint{0.981512in}{0.601537in}}%
\pgfpathlineto{\pgfqpoint{0.984031in}{0.602062in}}%
\pgfpathlineto{\pgfqpoint{0.986549in}{0.602384in}}%
\pgfpathlineto{\pgfqpoint{0.989068in}{0.585483in}}%
\pgfpathlineto{\pgfqpoint{0.991587in}{0.598217in}}%
\pgfpathlineto{\pgfqpoint{0.994106in}{0.611707in}}%
\pgfpathlineto{\pgfqpoint{0.996624in}{0.611372in}}%
\pgfpathlineto{\pgfqpoint{0.999143in}{0.620424in}}%
\pgfpathlineto{\pgfqpoint{1.001662in}{0.619561in}}%
\pgfpathlineto{\pgfqpoint{1.004181in}{0.613597in}}%
\pgfpathlineto{\pgfqpoint{1.006699in}{0.642791in}}%
\pgfpathlineto{\pgfqpoint{1.009218in}{0.660551in}}%
\pgfpathlineto{\pgfqpoint{1.011737in}{0.663365in}}%
\pgfpathlineto{\pgfqpoint{1.014256in}{0.678916in}}%
\pgfpathlineto{\pgfqpoint{1.016774in}{0.690373in}}%
\pgfpathlineto{\pgfqpoint{1.019293in}{0.680677in}}%
\pgfpathlineto{\pgfqpoint{1.021812in}{0.678815in}}%
\pgfpathlineto{\pgfqpoint{1.024331in}{0.674951in}}%
\pgfpathlineto{\pgfqpoint{1.026849in}{0.673555in}}%
\pgfpathlineto{\pgfqpoint{1.029368in}{0.681875in}}%
\pgfpathlineto{\pgfqpoint{1.031887in}{0.700761in}}%
\pgfpathlineto{\pgfqpoint{1.034406in}{0.704441in}}%
\pgfpathlineto{\pgfqpoint{1.036924in}{0.694485in}}%
\pgfpathlineto{\pgfqpoint{1.039443in}{0.703568in}}%
\pgfpathlineto{\pgfqpoint{1.041962in}{0.695084in}}%
\pgfpathlineto{\pgfqpoint{1.044481in}{0.699276in}}%
\pgfpathlineto{\pgfqpoint{1.046999in}{0.714718in}}%
\pgfpathlineto{\pgfqpoint{1.049518in}{0.720881in}}%
\pgfpathlineto{\pgfqpoint{1.052037in}{0.708728in}}%
\pgfpathlineto{\pgfqpoint{1.054556in}{0.714376in}}%
\pgfpathlineto{\pgfqpoint{1.057074in}{0.691885in}}%
\pgfpathlineto{\pgfqpoint{1.059593in}{0.720178in}}%
\pgfpathlineto{\pgfqpoint{1.062112in}{0.731547in}}%
\pgfpathlineto{\pgfqpoint{1.064631in}{0.718639in}}%
\pgfpathlineto{\pgfqpoint{1.067149in}{0.700022in}}%
\pgfpathlineto{\pgfqpoint{1.069668in}{0.702100in}}%
\pgfpathlineto{\pgfqpoint{1.072187in}{0.711872in}}%
\pgfpathlineto{\pgfqpoint{1.074706in}{0.711917in}}%
\pgfpathlineto{\pgfqpoint{1.077224in}{0.699132in}}%
\pgfpathlineto{\pgfqpoint{1.082262in}{0.713510in}}%
\pgfpathlineto{\pgfqpoint{1.084781in}{0.721225in}}%
\pgfpathlineto{\pgfqpoint{1.087299in}{0.734436in}}%
\pgfpathlineto{\pgfqpoint{1.089818in}{0.714406in}}%
\pgfpathlineto{\pgfqpoint{1.092337in}{0.708144in}}%
\pgfpathlineto{\pgfqpoint{1.094856in}{0.715965in}}%
\pgfpathlineto{\pgfqpoint{1.097374in}{0.721668in}}%
\pgfpathlineto{\pgfqpoint{1.099893in}{0.729569in}}%
\pgfpathlineto{\pgfqpoint{1.102412in}{0.712516in}}%
\pgfpathlineto{\pgfqpoint{1.104931in}{0.712783in}}%
\pgfpathlineto{\pgfqpoint{1.107449in}{0.703165in}}%
\pgfpathlineto{\pgfqpoint{1.109968in}{0.704038in}}%
\pgfpathlineto{\pgfqpoint{1.112487in}{0.710185in}}%
\pgfpathlineto{\pgfqpoint{1.115006in}{0.716841in}}%
\pgfpathlineto{\pgfqpoint{1.117524in}{0.725456in}}%
\pgfpathlineto{\pgfqpoint{1.120043in}{0.706754in}}%
\pgfpathlineto{\pgfqpoint{1.122562in}{0.697432in}}%
\pgfpathlineto{\pgfqpoint{1.125081in}{0.719787in}}%
\pgfpathlineto{\pgfqpoint{1.127599in}{0.728044in}}%
\pgfpathlineto{\pgfqpoint{1.130118in}{0.727701in}}%
\pgfpathlineto{\pgfqpoint{1.132637in}{0.732030in}}%
\pgfpathlineto{\pgfqpoint{1.135156in}{0.741630in}}%
\pgfpathlineto{\pgfqpoint{1.137674in}{0.754858in}}%
\pgfpathlineto{\pgfqpoint{1.140193in}{0.759929in}}%
\pgfpathlineto{\pgfqpoint{1.142712in}{0.770425in}}%
\pgfpathlineto{\pgfqpoint{1.145231in}{0.791690in}}%
\pgfpathlineto{\pgfqpoint{1.147749in}{0.808432in}}%
\pgfpathlineto{\pgfqpoint{1.150268in}{0.813378in}}%
\pgfpathlineto{\pgfqpoint{1.152787in}{0.783044in}}%
\pgfpathlineto{\pgfqpoint{1.155306in}{0.778321in}}%
\pgfpathlineto{\pgfqpoint{1.157824in}{0.772148in}}%
\pgfpathlineto{\pgfqpoint{1.160343in}{0.758296in}}%
\pgfpathlineto{\pgfqpoint{1.162862in}{0.760285in}}%
\pgfpathlineto{\pgfqpoint{1.165381in}{0.781459in}}%
\pgfpathlineto{\pgfqpoint{1.167899in}{0.767835in}}%
\pgfpathlineto{\pgfqpoint{1.170418in}{0.775069in}}%
\pgfpathlineto{\pgfqpoint{1.172937in}{0.769421in}}%
\pgfpathlineto{\pgfqpoint{1.175456in}{0.779392in}}%
\pgfpathlineto{\pgfqpoint{1.177974in}{0.767320in}}%
\pgfpathlineto{\pgfqpoint{1.180493in}{0.747261in}}%
\pgfpathlineto{\pgfqpoint{1.183012in}{0.759784in}}%
\pgfpathlineto{\pgfqpoint{1.185531in}{0.738329in}}%
\pgfpathlineto{\pgfqpoint{1.188049in}{0.735526in}}%
\pgfpathlineto{\pgfqpoint{1.190568in}{0.732169in}}%
\pgfpathlineto{\pgfqpoint{1.193087in}{0.731090in}}%
\pgfpathlineto{\pgfqpoint{1.195606in}{0.709256in}}%
\pgfpathlineto{\pgfqpoint{1.198124in}{0.719825in}}%
\pgfpathlineto{\pgfqpoint{1.200643in}{0.726267in}}%
\pgfpathlineto{\pgfqpoint{1.203162in}{0.738506in}}%
\pgfpathlineto{\pgfqpoint{1.205681in}{0.747027in}}%
\pgfpathlineto{\pgfqpoint{1.208199in}{0.732513in}}%
\pgfpathlineto{\pgfqpoint{1.210718in}{0.710652in}}%
\pgfpathlineto{\pgfqpoint{1.213237in}{0.705808in}}%
\pgfpathlineto{\pgfqpoint{1.215756in}{0.697012in}}%
\pgfpathlineto{\pgfqpoint{1.218274in}{0.694160in}}%
\pgfpathlineto{\pgfqpoint{1.220793in}{0.699449in}}%
\pgfpathlineto{\pgfqpoint{1.223312in}{0.696343in}}%
\pgfpathlineto{\pgfqpoint{1.225831in}{0.717994in}}%
\pgfpathlineto{\pgfqpoint{1.228349in}{0.715669in}}%
\pgfpathlineto{\pgfqpoint{1.230868in}{0.723558in}}%
\pgfpathlineto{\pgfqpoint{1.233387in}{0.732723in}}%
\pgfpathlineto{\pgfqpoint{1.235906in}{0.718735in}}%
\pgfpathlineto{\pgfqpoint{1.243462in}{0.700948in}}%
\pgfpathlineto{\pgfqpoint{1.245981in}{0.703150in}}%
\pgfpathlineto{\pgfqpoint{1.248499in}{0.690436in}}%
\pgfpathlineto{\pgfqpoint{1.251018in}{0.675412in}}%
\pgfpathlineto{\pgfqpoint{1.253537in}{0.665075in}}%
\pgfpathlineto{\pgfqpoint{1.256056in}{0.682202in}}%
\pgfpathlineto{\pgfqpoint{1.258574in}{0.663820in}}%
\pgfpathlineto{\pgfqpoint{1.261093in}{0.678880in}}%
\pgfpathlineto{\pgfqpoint{1.263612in}{0.685066in}}%
\pgfpathlineto{\pgfqpoint{1.266131in}{0.682823in}}%
\pgfpathlineto{\pgfqpoint{1.268649in}{0.693472in}}%
\pgfpathlineto{\pgfqpoint{1.271168in}{0.674517in}}%
\pgfpathlineto{\pgfqpoint{1.273687in}{0.683113in}}%
\pgfpathlineto{\pgfqpoint{1.276206in}{0.676992in}}%
\pgfpathlineto{\pgfqpoint{1.278724in}{0.648061in}}%
\pgfpathlineto{\pgfqpoint{1.281243in}{0.657295in}}%
\pgfpathlineto{\pgfqpoint{1.283762in}{0.655385in}}%
\pgfpathlineto{\pgfqpoint{1.286281in}{0.645025in}}%
\pgfpathlineto{\pgfqpoint{1.288799in}{0.647290in}}%
\pgfpathlineto{\pgfqpoint{1.291318in}{0.625019in}}%
\pgfpathlineto{\pgfqpoint{1.293837in}{0.606480in}}%
\pgfpathlineto{\pgfqpoint{1.296356in}{0.609844in}}%
\pgfpathlineto{\pgfqpoint{1.298874in}{0.610712in}}%
\pgfpathlineto{\pgfqpoint{1.301393in}{0.608140in}}%
\pgfpathlineto{\pgfqpoint{1.303912in}{0.621725in}}%
\pgfpathlineto{\pgfqpoint{1.306431in}{0.613665in}}%
\pgfpathlineto{\pgfqpoint{1.308949in}{0.611700in}}%
\pgfpathlineto{\pgfqpoint{1.311468in}{0.610504in}}%
\pgfpathlineto{\pgfqpoint{1.313987in}{0.640305in}}%
\pgfpathlineto{\pgfqpoint{1.316506in}{0.657149in}}%
\pgfpathlineto{\pgfqpoint{1.319024in}{0.649047in}}%
\pgfpathlineto{\pgfqpoint{1.321543in}{0.633687in}}%
\pgfpathlineto{\pgfqpoint{1.324062in}{0.647865in}}%
\pgfpathlineto{\pgfqpoint{1.326581in}{0.639760in}}%
\pgfpathlineto{\pgfqpoint{1.329099in}{0.654285in}}%
\pgfpathlineto{\pgfqpoint{1.331618in}{0.675022in}}%
\pgfpathlineto{\pgfqpoint{1.334137in}{0.679842in}}%
\pgfpathlineto{\pgfqpoint{1.336656in}{0.684232in}}%
\pgfpathlineto{\pgfqpoint{1.339174in}{0.716107in}}%
\pgfpathlineto{\pgfqpoint{1.341693in}{0.737932in}}%
\pgfpathlineto{\pgfqpoint{1.344212in}{0.741282in}}%
\pgfpathlineto{\pgfqpoint{1.346731in}{0.749287in}}%
\pgfpathlineto{\pgfqpoint{1.349249in}{0.750482in}}%
\pgfpathlineto{\pgfqpoint{1.351768in}{0.753882in}}%
\pgfpathlineto{\pgfqpoint{1.354287in}{0.771792in}}%
\pgfpathlineto{\pgfqpoint{1.356806in}{0.776152in}}%
\pgfpathlineto{\pgfqpoint{1.359324in}{0.799736in}}%
\pgfpathlineto{\pgfqpoint{1.361843in}{0.807114in}}%
\pgfpathlineto{\pgfqpoint{1.364362in}{0.821980in}}%
\pgfpathlineto{\pgfqpoint{1.366881in}{0.809820in}}%
\pgfpathlineto{\pgfqpoint{1.369399in}{0.786985in}}%
\pgfpathlineto{\pgfqpoint{1.371918in}{0.780733in}}%
\pgfpathlineto{\pgfqpoint{1.374437in}{0.766660in}}%
\pgfpathlineto{\pgfqpoint{1.376956in}{0.769348in}}%
\pgfpathlineto{\pgfqpoint{1.379474in}{0.783177in}}%
\pgfpathlineto{\pgfqpoint{1.381993in}{0.791986in}}%
\pgfpathlineto{\pgfqpoint{1.384512in}{0.791221in}}%
\pgfpathlineto{\pgfqpoint{1.387031in}{0.774586in}}%
\pgfpathlineto{\pgfqpoint{1.389549in}{0.762847in}}%
\pgfpathlineto{\pgfqpoint{1.392068in}{0.769137in}}%
\pgfpathlineto{\pgfqpoint{1.394587in}{0.768382in}}%
\pgfpathlineto{\pgfqpoint{1.397106in}{0.780246in}}%
\pgfpathlineto{\pgfqpoint{1.399624in}{0.778391in}}%
\pgfpathlineto{\pgfqpoint{1.402143in}{0.772622in}}%
\pgfpathlineto{\pgfqpoint{1.404662in}{0.780828in}}%
\pgfpathlineto{\pgfqpoint{1.407181in}{0.794239in}}%
\pgfpathlineto{\pgfqpoint{1.409699in}{0.794329in}}%
\pgfpathlineto{\pgfqpoint{1.412218in}{0.784144in}}%
\pgfpathlineto{\pgfqpoint{1.414737in}{0.766895in}}%
\pgfpathlineto{\pgfqpoint{1.417256in}{0.762155in}}%
\pgfpathlineto{\pgfqpoint{1.419774in}{0.767146in}}%
\pgfpathlineto{\pgfqpoint{1.422293in}{0.769706in}}%
\pgfpathlineto{\pgfqpoint{1.424812in}{0.758562in}}%
\pgfpathlineto{\pgfqpoint{1.427331in}{0.774197in}}%
\pgfpathlineto{\pgfqpoint{1.429849in}{0.779627in}}%
\pgfpathlineto{\pgfqpoint{1.432368in}{0.791968in}}%
\pgfpathlineto{\pgfqpoint{1.434887in}{0.788907in}}%
\pgfpathlineto{\pgfqpoint{1.437406in}{0.793325in}}%
\pgfpathlineto{\pgfqpoint{1.439924in}{0.812426in}}%
\pgfpathlineto{\pgfqpoint{1.442443in}{0.813402in}}%
\pgfpathlineto{\pgfqpoint{1.444962in}{0.780587in}}%
\pgfpathlineto{\pgfqpoint{1.447481in}{0.766196in}}%
\pgfpathlineto{\pgfqpoint{1.449999in}{0.756500in}}%
\pgfpathlineto{\pgfqpoint{1.452518in}{0.743485in}}%
\pgfpathlineto{\pgfqpoint{1.455037in}{0.727285in}}%
\pgfpathlineto{\pgfqpoint{1.457556in}{0.745745in}}%
\pgfpathlineto{\pgfqpoint{1.460074in}{0.756593in}}%
\pgfpathlineto{\pgfqpoint{1.462593in}{0.753763in}}%
\pgfpathlineto{\pgfqpoint{1.465112in}{0.744329in}}%
\pgfpathlineto{\pgfqpoint{1.467631in}{0.745500in}}%
\pgfpathlineto{\pgfqpoint{1.470149in}{0.765364in}}%
\pgfpathlineto{\pgfqpoint{1.472668in}{0.772747in}}%
\pgfpathlineto{\pgfqpoint{1.475187in}{0.764188in}}%
\pgfpathlineto{\pgfqpoint{1.477706in}{0.762936in}}%
\pgfpathlineto{\pgfqpoint{1.480224in}{0.752583in}}%
\pgfpathlineto{\pgfqpoint{1.482743in}{0.760512in}}%
\pgfpathlineto{\pgfqpoint{1.485262in}{0.745871in}}%
\pgfpathlineto{\pgfqpoint{1.487781in}{0.754601in}}%
\pgfpathlineto{\pgfqpoint{1.490299in}{0.765937in}}%
\pgfpathlineto{\pgfqpoint{1.492818in}{0.768385in}}%
\pgfpathlineto{\pgfqpoint{1.495337in}{0.760962in}}%
\pgfpathlineto{\pgfqpoint{1.497856in}{0.757007in}}%
\pgfpathlineto{\pgfqpoint{1.500374in}{0.771749in}}%
\pgfpathlineto{\pgfqpoint{1.502893in}{0.771101in}}%
\pgfpathlineto{\pgfqpoint{1.505412in}{0.771140in}}%
\pgfpathlineto{\pgfqpoint{1.507931in}{0.781488in}}%
\pgfpathlineto{\pgfqpoint{1.510449in}{0.794992in}}%
\pgfpathlineto{\pgfqpoint{1.512968in}{0.796004in}}%
\pgfpathlineto{\pgfqpoint{1.515487in}{0.808067in}}%
\pgfpathlineto{\pgfqpoint{1.518006in}{0.811790in}}%
\pgfpathlineto{\pgfqpoint{1.520524in}{0.819484in}}%
\pgfpathlineto{\pgfqpoint{1.523043in}{0.799041in}}%
\pgfpathlineto{\pgfqpoint{1.525562in}{0.786171in}}%
\pgfpathlineto{\pgfqpoint{1.528081in}{0.775227in}}%
\pgfpathlineto{\pgfqpoint{1.530599in}{0.759192in}}%
\pgfpathlineto{\pgfqpoint{1.533118in}{0.759757in}}%
\pgfpathlineto{\pgfqpoint{1.535637in}{0.777928in}}%
\pgfpathlineto{\pgfqpoint{1.538156in}{0.783620in}}%
\pgfpathlineto{\pgfqpoint{1.540674in}{0.793758in}}%
\pgfpathlineto{\pgfqpoint{1.543193in}{0.785377in}}%
\pgfpathlineto{\pgfqpoint{1.545712in}{0.773770in}}%
\pgfpathlineto{\pgfqpoint{1.548231in}{0.797681in}}%
\pgfpathlineto{\pgfqpoint{1.550749in}{0.804792in}}%
\pgfpathlineto{\pgfqpoint{1.553268in}{0.807810in}}%
\pgfpathlineto{\pgfqpoint{1.555787in}{0.817179in}}%
\pgfpathlineto{\pgfqpoint{1.558306in}{0.791364in}}%
\pgfpathlineto{\pgfqpoint{1.560824in}{0.777576in}}%
\pgfpathlineto{\pgfqpoint{1.563343in}{0.765829in}}%
\pgfpathlineto{\pgfqpoint{1.565862in}{0.783585in}}%
\pgfpathlineto{\pgfqpoint{1.568381in}{0.792790in}}%
\pgfpathlineto{\pgfqpoint{1.570899in}{0.766433in}}%
\pgfpathlineto{\pgfqpoint{1.573418in}{0.811931in}}%
\pgfpathlineto{\pgfqpoint{1.575937in}{0.816102in}}%
\pgfpathlineto{\pgfqpoint{1.578456in}{0.831575in}}%
\pgfpathlineto{\pgfqpoint{1.580974in}{0.837906in}}%
\pgfpathlineto{\pgfqpoint{1.583493in}{0.831889in}}%
\pgfpathlineto{\pgfqpoint{1.586012in}{0.847474in}}%
\pgfpathlineto{\pgfqpoint{1.591049in}{0.821857in}}%
\pgfpathlineto{\pgfqpoint{1.593568in}{0.829385in}}%
\pgfpathlineto{\pgfqpoint{1.596087in}{0.841348in}}%
\pgfpathlineto{\pgfqpoint{1.598606in}{0.853967in}}%
\pgfpathlineto{\pgfqpoint{1.601124in}{0.829828in}}%
\pgfpathlineto{\pgfqpoint{1.603643in}{0.815564in}}%
\pgfpathlineto{\pgfqpoint{1.606162in}{0.836720in}}%
\pgfpathlineto{\pgfqpoint{1.608681in}{0.862364in}}%
\pgfpathlineto{\pgfqpoint{1.611199in}{0.848022in}}%
\pgfpathlineto{\pgfqpoint{1.613718in}{0.845444in}}%
\pgfpathlineto{\pgfqpoint{1.616237in}{0.858021in}}%
\pgfpathlineto{\pgfqpoint{1.618756in}{0.866483in}}%
\pgfpathlineto{\pgfqpoint{1.621274in}{0.894822in}}%
\pgfpathlineto{\pgfqpoint{1.623793in}{0.905975in}}%
\pgfpathlineto{\pgfqpoint{1.626312in}{0.912915in}}%
\pgfpathlineto{\pgfqpoint{1.628831in}{0.916418in}}%
\pgfpathlineto{\pgfqpoint{1.631349in}{0.923041in}}%
\pgfpathlineto{\pgfqpoint{1.633868in}{0.925524in}}%
\pgfpathlineto{\pgfqpoint{1.636387in}{0.923441in}}%
\pgfpathlineto{\pgfqpoint{1.638906in}{0.922326in}}%
\pgfpathlineto{\pgfqpoint{1.641424in}{0.927219in}}%
\pgfpathlineto{\pgfqpoint{1.643943in}{0.934665in}}%
\pgfpathlineto{\pgfqpoint{1.646462in}{0.902293in}}%
\pgfpathlineto{\pgfqpoint{1.648981in}{0.904828in}}%
\pgfpathlineto{\pgfqpoint{1.651499in}{0.929143in}}%
\pgfpathlineto{\pgfqpoint{1.654018in}{0.935683in}}%
\pgfpathlineto{\pgfqpoint{1.656537in}{0.923315in}}%
\pgfpathlineto{\pgfqpoint{1.659056in}{0.906023in}}%
\pgfpathlineto{\pgfqpoint{1.661574in}{0.929626in}}%
\pgfpathlineto{\pgfqpoint{1.664093in}{0.901990in}}%
\pgfpathlineto{\pgfqpoint{1.666612in}{0.904779in}}%
\pgfpathlineto{\pgfqpoint{1.669131in}{0.920486in}}%
\pgfpathlineto{\pgfqpoint{1.671649in}{0.922275in}}%
\pgfpathlineto{\pgfqpoint{1.674168in}{0.943194in}}%
\pgfpathlineto{\pgfqpoint{1.676687in}{0.954083in}}%
\pgfpathlineto{\pgfqpoint{1.679206in}{0.949635in}}%
\pgfpathlineto{\pgfqpoint{1.681724in}{0.938472in}}%
\pgfpathlineto{\pgfqpoint{1.684243in}{0.919222in}}%
\pgfpathlineto{\pgfqpoint{1.686762in}{0.908986in}}%
\pgfpathlineto{\pgfqpoint{1.689281in}{0.919830in}}%
\pgfpathlineto{\pgfqpoint{1.691799in}{0.927315in}}%
\pgfpathlineto{\pgfqpoint{1.694318in}{0.928528in}}%
\pgfpathlineto{\pgfqpoint{1.696837in}{0.934726in}}%
\pgfpathlineto{\pgfqpoint{1.699356in}{0.950822in}}%
\pgfpathlineto{\pgfqpoint{1.701874in}{0.931261in}}%
\pgfpathlineto{\pgfqpoint{1.704393in}{0.915831in}}%
\pgfpathlineto{\pgfqpoint{1.706912in}{0.915369in}}%
\pgfpathlineto{\pgfqpoint{1.709431in}{0.919058in}}%
\pgfpathlineto{\pgfqpoint{1.711949in}{0.925112in}}%
\pgfpathlineto{\pgfqpoint{1.714468in}{0.903369in}}%
\pgfpathlineto{\pgfqpoint{1.716987in}{0.919216in}}%
\pgfpathlineto{\pgfqpoint{1.719506in}{0.921906in}}%
\pgfpathlineto{\pgfqpoint{1.722024in}{0.923554in}}%
\pgfpathlineto{\pgfqpoint{1.724543in}{0.961862in}}%
\pgfpathlineto{\pgfqpoint{1.727062in}{0.983109in}}%
\pgfpathlineto{\pgfqpoint{1.729581in}{0.996145in}}%
\pgfpathlineto{\pgfqpoint{1.732099in}{0.993642in}}%
\pgfpathlineto{\pgfqpoint{1.734618in}{1.001124in}}%
\pgfpathlineto{\pgfqpoint{1.737137in}{1.003050in}}%
\pgfpathlineto{\pgfqpoint{1.739656in}{0.999693in}}%
\pgfpathlineto{\pgfqpoint{1.742174in}{0.997146in}}%
\pgfpathlineto{\pgfqpoint{1.744693in}{0.991459in}}%
\pgfpathlineto{\pgfqpoint{1.747212in}{0.994747in}}%
\pgfpathlineto{\pgfqpoint{1.749731in}{0.994861in}}%
\pgfpathlineto{\pgfqpoint{1.752249in}{1.008391in}}%
\pgfpathlineto{\pgfqpoint{1.754768in}{1.011570in}}%
\pgfpathlineto{\pgfqpoint{1.757287in}{1.025659in}}%
\pgfpathlineto{\pgfqpoint{1.759806in}{1.021048in}}%
\pgfpathlineto{\pgfqpoint{1.762324in}{1.025038in}}%
\pgfpathlineto{\pgfqpoint{1.764843in}{1.031279in}}%
\pgfpathlineto{\pgfqpoint{1.767362in}{1.038964in}}%
\pgfpathlineto{\pgfqpoint{1.769881in}{1.039688in}}%
\pgfpathlineto{\pgfqpoint{1.772399in}{1.047460in}}%
\pgfpathlineto{\pgfqpoint{1.774918in}{1.041457in}}%
\pgfpathlineto{\pgfqpoint{1.777437in}{1.066632in}}%
\pgfpathlineto{\pgfqpoint{1.779956in}{1.065649in}}%
\pgfpathlineto{\pgfqpoint{1.782474in}{1.076106in}}%
\pgfpathlineto{\pgfqpoint{1.784993in}{1.103189in}}%
\pgfpathlineto{\pgfqpoint{1.787512in}{1.094251in}}%
\pgfpathlineto{\pgfqpoint{1.790031in}{1.096177in}}%
\pgfpathlineto{\pgfqpoint{1.792549in}{1.090555in}}%
\pgfpathlineto{\pgfqpoint{1.795068in}{1.103153in}}%
\pgfpathlineto{\pgfqpoint{1.797587in}{1.113552in}}%
\pgfpathlineto{\pgfqpoint{1.800106in}{1.083240in}}%
\pgfpathlineto{\pgfqpoint{1.802624in}{1.077651in}}%
\pgfpathlineto{\pgfqpoint{1.805143in}{1.065587in}}%
\pgfpathlineto{\pgfqpoint{1.807662in}{1.048250in}}%
\pgfpathlineto{\pgfqpoint{1.810181in}{1.045913in}}%
\pgfpathlineto{\pgfqpoint{1.812699in}{1.032900in}}%
\pgfpathlineto{\pgfqpoint{1.815218in}{1.068188in}}%
\pgfpathlineto{\pgfqpoint{1.817737in}{1.060937in}}%
\pgfpathlineto{\pgfqpoint{1.820256in}{1.050906in}}%
\pgfpathlineto{\pgfqpoint{1.822774in}{1.064141in}}%
\pgfpathlineto{\pgfqpoint{1.825293in}{1.073329in}}%
\pgfpathlineto{\pgfqpoint{1.827812in}{1.080049in}}%
\pgfpathlineto{\pgfqpoint{1.830331in}{1.061412in}}%
\pgfpathlineto{\pgfqpoint{1.832849in}{1.037785in}}%
\pgfpathlineto{\pgfqpoint{1.835368in}{1.029258in}}%
\pgfpathlineto{\pgfqpoint{1.837887in}{1.054867in}}%
\pgfpathlineto{\pgfqpoint{1.840406in}{1.074624in}}%
\pgfpathlineto{\pgfqpoint{1.842924in}{1.065905in}}%
\pgfpathlineto{\pgfqpoint{1.845443in}{1.078769in}}%
\pgfpathlineto{\pgfqpoint{1.847962in}{1.065987in}}%
\pgfpathlineto{\pgfqpoint{1.850481in}{1.085141in}}%
\pgfpathlineto{\pgfqpoint{1.852999in}{1.108291in}}%
\pgfpathlineto{\pgfqpoint{1.855518in}{1.109855in}}%
\pgfpathlineto{\pgfqpoint{1.858037in}{1.100890in}}%
\pgfpathlineto{\pgfqpoint{1.860556in}{1.114601in}}%
\pgfpathlineto{\pgfqpoint{1.863074in}{1.129957in}}%
\pgfpathlineto{\pgfqpoint{1.865593in}{1.129719in}}%
\pgfpathlineto{\pgfqpoint{1.868112in}{1.156111in}}%
\pgfpathlineto{\pgfqpoint{1.870631in}{1.144730in}}%
\pgfpathlineto{\pgfqpoint{1.873149in}{1.135655in}}%
\pgfpathlineto{\pgfqpoint{1.875668in}{1.148615in}}%
\pgfpathlineto{\pgfqpoint{1.878187in}{1.153974in}}%
\pgfpathlineto{\pgfqpoint{1.880706in}{1.165660in}}%
\pgfpathlineto{\pgfqpoint{1.883224in}{1.159874in}}%
\pgfpathlineto{\pgfqpoint{1.885743in}{1.167074in}}%
\pgfpathlineto{\pgfqpoint{1.888262in}{1.147530in}}%
\pgfpathlineto{\pgfqpoint{1.890781in}{1.163405in}}%
\pgfpathlineto{\pgfqpoint{1.893299in}{1.156099in}}%
\pgfpathlineto{\pgfqpoint{1.895818in}{1.159282in}}%
\pgfpathlineto{\pgfqpoint{1.898337in}{1.165218in}}%
\pgfpathlineto{\pgfqpoint{1.900856in}{1.160605in}}%
\pgfpathlineto{\pgfqpoint{1.903374in}{1.165986in}}%
\pgfpathlineto{\pgfqpoint{1.905893in}{1.167460in}}%
\pgfpathlineto{\pgfqpoint{1.908412in}{1.159464in}}%
\pgfpathlineto{\pgfqpoint{1.910931in}{1.147136in}}%
\pgfpathlineto{\pgfqpoint{1.913449in}{1.160023in}}%
\pgfpathlineto{\pgfqpoint{1.915968in}{1.180988in}}%
\pgfpathlineto{\pgfqpoint{1.918487in}{1.194539in}}%
\pgfpathlineto{\pgfqpoint{1.921006in}{1.202242in}}%
\pgfpathlineto{\pgfqpoint{1.923524in}{1.205750in}}%
\pgfpathlineto{\pgfqpoint{1.926043in}{1.211594in}}%
\pgfpathlineto{\pgfqpoint{1.928562in}{1.201234in}}%
\pgfpathlineto{\pgfqpoint{1.931081in}{1.208978in}}%
\pgfpathlineto{\pgfqpoint{1.933599in}{1.226434in}}%
\pgfpathlineto{\pgfqpoint{1.936118in}{1.216634in}}%
\pgfpathlineto{\pgfqpoint{1.938637in}{1.227989in}}%
\pgfpathlineto{\pgfqpoint{1.941156in}{1.228924in}}%
\pgfpathlineto{\pgfqpoint{1.943674in}{1.225063in}}%
\pgfpathlineto{\pgfqpoint{1.946193in}{1.271464in}}%
\pgfpathlineto{\pgfqpoint{1.948712in}{1.289681in}}%
\pgfpathlineto{\pgfqpoint{1.951231in}{1.271743in}}%
\pgfpathlineto{\pgfqpoint{1.953749in}{1.277121in}}%
\pgfpathlineto{\pgfqpoint{1.956268in}{1.268782in}}%
\pgfpathlineto{\pgfqpoint{1.958787in}{1.291947in}}%
\pgfpathlineto{\pgfqpoint{1.961306in}{1.285871in}}%
\pgfpathlineto{\pgfqpoint{1.963824in}{1.275443in}}%
\pgfpathlineto{\pgfqpoint{1.966343in}{1.264233in}}%
\pgfpathlineto{\pgfqpoint{1.968862in}{1.274139in}}%
\pgfpathlineto{\pgfqpoint{1.971381in}{1.299968in}}%
\pgfpathlineto{\pgfqpoint{1.973899in}{1.291104in}}%
\pgfpathlineto{\pgfqpoint{1.976418in}{1.293166in}}%
\pgfpathlineto{\pgfqpoint{1.978937in}{1.268744in}}%
\pgfpathlineto{\pgfqpoint{1.981456in}{1.263291in}}%
\pgfpathlineto{\pgfqpoint{1.983974in}{1.242247in}}%
\pgfpathlineto{\pgfqpoint{1.986493in}{1.235581in}}%
\pgfpathlineto{\pgfqpoint{1.989012in}{1.225048in}}%
\pgfpathlineto{\pgfqpoint{1.991531in}{1.175072in}}%
\pgfpathlineto{\pgfqpoint{1.994049in}{1.171903in}}%
\pgfpathlineto{\pgfqpoint{1.999087in}{1.134690in}}%
\pgfpathlineto{\pgfqpoint{2.001606in}{1.141102in}}%
\pgfpathlineto{\pgfqpoint{2.004124in}{1.137438in}}%
\pgfpathlineto{\pgfqpoint{2.006643in}{1.141476in}}%
\pgfpathlineto{\pgfqpoint{2.009162in}{1.148946in}}%
\pgfpathlineto{\pgfqpoint{2.011681in}{1.139930in}}%
\pgfpathlineto{\pgfqpoint{2.014199in}{1.153044in}}%
\pgfpathlineto{\pgfqpoint{2.016718in}{1.156655in}}%
\pgfpathlineto{\pgfqpoint{2.019237in}{1.138191in}}%
\pgfpathlineto{\pgfqpoint{2.021756in}{1.130173in}}%
\pgfpathlineto{\pgfqpoint{2.024274in}{1.127525in}}%
\pgfpathlineto{\pgfqpoint{2.026793in}{1.102533in}}%
\pgfpathlineto{\pgfqpoint{2.029312in}{1.103637in}}%
\pgfpathlineto{\pgfqpoint{2.031831in}{1.134199in}}%
\pgfpathlineto{\pgfqpoint{2.034349in}{1.124593in}}%
\pgfpathlineto{\pgfqpoint{2.036868in}{1.133591in}}%
\pgfpathlineto{\pgfqpoint{2.039387in}{1.132580in}}%
\pgfpathlineto{\pgfqpoint{2.041906in}{1.119583in}}%
\pgfpathlineto{\pgfqpoint{2.044424in}{1.117110in}}%
\pgfpathlineto{\pgfqpoint{2.046943in}{1.112777in}}%
\pgfpathlineto{\pgfqpoint{2.049462in}{1.088882in}}%
\pgfpathlineto{\pgfqpoint{2.051981in}{1.099065in}}%
\pgfpathlineto{\pgfqpoint{2.054499in}{1.099802in}}%
\pgfpathlineto{\pgfqpoint{2.057018in}{1.127587in}}%
\pgfpathlineto{\pgfqpoint{2.059537in}{1.158866in}}%
\pgfpathlineto{\pgfqpoint{2.062056in}{1.160433in}}%
\pgfpathlineto{\pgfqpoint{2.064574in}{1.172159in}}%
\pgfpathlineto{\pgfqpoint{2.067093in}{1.196919in}}%
\pgfpathlineto{\pgfqpoint{2.069612in}{1.184972in}}%
\pgfpathlineto{\pgfqpoint{2.072131in}{1.165673in}}%
\pgfpathlineto{\pgfqpoint{2.074649in}{1.169530in}}%
\pgfpathlineto{\pgfqpoint{2.077168in}{1.160158in}}%
\pgfpathlineto{\pgfqpoint{2.079687in}{1.172709in}}%
\pgfpathlineto{\pgfqpoint{2.082206in}{1.190876in}}%
\pgfpathlineto{\pgfqpoint{2.084724in}{1.181549in}}%
\pgfpathlineto{\pgfqpoint{2.087243in}{1.179467in}}%
\pgfpathlineto{\pgfqpoint{2.089762in}{1.160866in}}%
\pgfpathlineto{\pgfqpoint{2.094799in}{1.176496in}}%
\pgfpathlineto{\pgfqpoint{2.097318in}{1.190899in}}%
\pgfpathlineto{\pgfqpoint{2.099837in}{1.176603in}}%
\pgfpathlineto{\pgfqpoint{2.102356in}{1.185884in}}%
\pgfpathlineto{\pgfqpoint{2.104874in}{1.178290in}}%
\pgfpathlineto{\pgfqpoint{2.107393in}{1.174520in}}%
\pgfpathlineto{\pgfqpoint{2.109912in}{1.213323in}}%
\pgfpathlineto{\pgfqpoint{2.112431in}{1.218033in}}%
\pgfpathlineto{\pgfqpoint{2.114949in}{1.199412in}}%
\pgfpathlineto{\pgfqpoint{2.117468in}{1.222787in}}%
\pgfpathlineto{\pgfqpoint{2.119987in}{1.245176in}}%
\pgfpathlineto{\pgfqpoint{2.122506in}{1.224597in}}%
\pgfpathlineto{\pgfqpoint{2.125024in}{1.213691in}}%
\pgfpathlineto{\pgfqpoint{2.127543in}{1.217208in}}%
\pgfpathlineto{\pgfqpoint{2.130062in}{1.215168in}}%
\pgfpathlineto{\pgfqpoint{2.132581in}{1.220024in}}%
\pgfpathlineto{\pgfqpoint{2.135099in}{1.190000in}}%
\pgfpathlineto{\pgfqpoint{2.137618in}{1.200438in}}%
\pgfpathlineto{\pgfqpoint{2.140137in}{1.195854in}}%
\pgfpathlineto{\pgfqpoint{2.142656in}{1.202985in}}%
\pgfpathlineto{\pgfqpoint{2.145174in}{1.222006in}}%
\pgfpathlineto{\pgfqpoint{2.147693in}{1.228310in}}%
\pgfpathlineto{\pgfqpoint{2.150212in}{1.235538in}}%
\pgfpathlineto{\pgfqpoint{2.152731in}{1.219049in}}%
\pgfpathlineto{\pgfqpoint{2.155249in}{1.217972in}}%
\pgfpathlineto{\pgfqpoint{2.157768in}{1.227010in}}%
\pgfpathlineto{\pgfqpoint{2.160287in}{1.232016in}}%
\pgfpathlineto{\pgfqpoint{2.162806in}{1.242592in}}%
\pgfpathlineto{\pgfqpoint{2.165324in}{1.254898in}}%
\pgfpathlineto{\pgfqpoint{2.167843in}{1.255499in}}%
\pgfpathlineto{\pgfqpoint{2.170362in}{1.243177in}}%
\pgfpathlineto{\pgfqpoint{2.172881in}{1.255660in}}%
\pgfpathlineto{\pgfqpoint{2.175399in}{1.253744in}}%
\pgfpathlineto{\pgfqpoint{2.177918in}{1.273151in}}%
\pgfpathlineto{\pgfqpoint{2.180437in}{1.267659in}}%
\pgfpathlineto{\pgfqpoint{2.182956in}{1.255051in}}%
\pgfpathlineto{\pgfqpoint{2.185474in}{1.255748in}}%
\pgfpathlineto{\pgfqpoint{2.187993in}{1.254749in}}%
\pgfpathlineto{\pgfqpoint{2.190512in}{1.257186in}}%
\pgfpathlineto{\pgfqpoint{2.193031in}{1.245150in}}%
\pgfpathlineto{\pgfqpoint{2.195549in}{1.260455in}}%
\pgfpathlineto{\pgfqpoint{2.198068in}{1.256693in}}%
\pgfpathlineto{\pgfqpoint{2.200587in}{1.271963in}}%
\pgfpathlineto{\pgfqpoint{2.203106in}{1.282866in}}%
\pgfpathlineto{\pgfqpoint{2.205624in}{1.295266in}}%
\pgfpathlineto{\pgfqpoint{2.208143in}{1.300667in}}%
\pgfpathlineto{\pgfqpoint{2.210662in}{1.283104in}}%
\pgfpathlineto{\pgfqpoint{2.213181in}{1.297284in}}%
\pgfpathlineto{\pgfqpoint{2.215699in}{1.289665in}}%
\pgfpathlineto{\pgfqpoint{2.218218in}{1.287113in}}%
\pgfpathlineto{\pgfqpoint{2.220737in}{1.270104in}}%
\pgfpathlineto{\pgfqpoint{2.223256in}{1.272561in}}%
\pgfpathlineto{\pgfqpoint{2.225774in}{1.293497in}}%
\pgfpathlineto{\pgfqpoint{2.228293in}{1.308057in}}%
\pgfpathlineto{\pgfqpoint{2.230812in}{1.304358in}}%
\pgfpathlineto{\pgfqpoint{2.233331in}{1.302967in}}%
\pgfpathlineto{\pgfqpoint{2.235849in}{1.282578in}}%
\pgfpathlineto{\pgfqpoint{2.238368in}{1.302646in}}%
\pgfpathlineto{\pgfqpoint{2.240887in}{1.287454in}}%
\pgfpathlineto{\pgfqpoint{2.243406in}{1.322676in}}%
\pgfpathlineto{\pgfqpoint{2.245924in}{1.327298in}}%
\pgfpathlineto{\pgfqpoint{2.248443in}{1.329198in}}%
\pgfpathlineto{\pgfqpoint{2.250962in}{1.324516in}}%
\pgfpathlineto{\pgfqpoint{2.253481in}{1.335358in}}%
\pgfpathlineto{\pgfqpoint{2.255999in}{1.294118in}}%
\pgfpathlineto{\pgfqpoint{2.258518in}{1.296428in}}%
\pgfpathlineto{\pgfqpoint{2.261037in}{1.286295in}}%
\pgfpathlineto{\pgfqpoint{2.263556in}{1.267942in}}%
\pgfpathlineto{\pgfqpoint{2.266074in}{1.247787in}}%
\pgfpathlineto{\pgfqpoint{2.268593in}{1.244275in}}%
\pgfpathlineto{\pgfqpoint{2.271112in}{1.213554in}}%
\pgfpathlineto{\pgfqpoint{2.273631in}{1.236024in}}%
\pgfpathlineto{\pgfqpoint{2.276149in}{1.246937in}}%
\pgfpathlineto{\pgfqpoint{2.278668in}{1.226295in}}%
\pgfpathlineto{\pgfqpoint{2.281187in}{1.211233in}}%
\pgfpathlineto{\pgfqpoint{2.283706in}{1.226587in}}%
\pgfpathlineto{\pgfqpoint{2.286224in}{1.256457in}}%
\pgfpathlineto{\pgfqpoint{2.288743in}{1.258759in}}%
\pgfpathlineto{\pgfqpoint{2.291262in}{1.229157in}}%
\pgfpathlineto{\pgfqpoint{2.293781in}{1.228922in}}%
\pgfpathlineto{\pgfqpoint{2.296299in}{1.217924in}}%
\pgfpathlineto{\pgfqpoint{2.298818in}{1.215221in}}%
\pgfpathlineto{\pgfqpoint{2.301337in}{1.241168in}}%
\pgfpathlineto{\pgfqpoint{2.303856in}{1.226939in}}%
\pgfpathlineto{\pgfqpoint{2.306374in}{1.231800in}}%
\pgfpathlineto{\pgfqpoint{2.308893in}{1.232707in}}%
\pgfpathlineto{\pgfqpoint{2.311412in}{1.249316in}}%
\pgfpathlineto{\pgfqpoint{2.313931in}{1.240082in}}%
\pgfpathlineto{\pgfqpoint{2.316449in}{1.225887in}}%
\pgfpathlineto{\pgfqpoint{2.318968in}{1.213682in}}%
\pgfpathlineto{\pgfqpoint{2.321487in}{1.220062in}}%
\pgfpathlineto{\pgfqpoint{2.324006in}{1.230290in}}%
\pgfpathlineto{\pgfqpoint{2.326524in}{1.231892in}}%
\pgfpathlineto{\pgfqpoint{2.329043in}{1.225403in}}%
\pgfpathlineto{\pgfqpoint{2.331562in}{1.203734in}}%
\pgfpathlineto{\pgfqpoint{2.334081in}{1.238447in}}%
\pgfpathlineto{\pgfqpoint{2.336599in}{1.240399in}}%
\pgfpathlineto{\pgfqpoint{2.339118in}{1.245178in}}%
\pgfpathlineto{\pgfqpoint{2.341637in}{1.244151in}}%
\pgfpathlineto{\pgfqpoint{2.344156in}{1.256028in}}%
\pgfpathlineto{\pgfqpoint{2.346674in}{1.269319in}}%
\pgfpathlineto{\pgfqpoint{2.349193in}{1.296950in}}%
\pgfpathlineto{\pgfqpoint{2.351712in}{1.298928in}}%
\pgfpathlineto{\pgfqpoint{2.354231in}{1.296065in}}%
\pgfpathlineto{\pgfqpoint{2.356749in}{1.296903in}}%
\pgfpathlineto{\pgfqpoint{2.359268in}{1.303556in}}%
\pgfpathlineto{\pgfqpoint{2.361787in}{1.317768in}}%
\pgfpathlineto{\pgfqpoint{2.364306in}{1.310727in}}%
\pgfpathlineto{\pgfqpoint{2.366824in}{1.287250in}}%
\pgfpathlineto{\pgfqpoint{2.369343in}{1.295512in}}%
\pgfpathlineto{\pgfqpoint{2.371862in}{1.325878in}}%
\pgfpathlineto{\pgfqpoint{2.374381in}{1.346651in}}%
\pgfpathlineto{\pgfqpoint{2.376899in}{1.362310in}}%
\pgfpathlineto{\pgfqpoint{2.379418in}{1.364827in}}%
\pgfpathlineto{\pgfqpoint{2.381937in}{1.367898in}}%
\pgfpathlineto{\pgfqpoint{2.384456in}{1.373552in}}%
\pgfpathlineto{\pgfqpoint{2.386974in}{1.374947in}}%
\pgfpathlineto{\pgfqpoint{2.389493in}{1.382005in}}%
\pgfpathlineto{\pgfqpoint{2.392012in}{1.401658in}}%
\pgfpathlineto{\pgfqpoint{2.394531in}{1.392273in}}%
\pgfpathlineto{\pgfqpoint{2.397049in}{1.366940in}}%
\pgfpathlineto{\pgfqpoint{2.399568in}{1.365406in}}%
\pgfpathlineto{\pgfqpoint{2.402087in}{1.379640in}}%
\pgfpathlineto{\pgfqpoint{2.404606in}{1.380943in}}%
\pgfpathlineto{\pgfqpoint{2.407124in}{1.371156in}}%
\pgfpathlineto{\pgfqpoint{2.409643in}{1.401138in}}%
\pgfpathlineto{\pgfqpoint{2.412162in}{1.381592in}}%
\pgfpathlineto{\pgfqpoint{2.414681in}{1.384386in}}%
\pgfpathlineto{\pgfqpoint{2.417199in}{1.355160in}}%
\pgfpathlineto{\pgfqpoint{2.419718in}{1.373639in}}%
\pgfpathlineto{\pgfqpoint{2.422237in}{1.349674in}}%
\pgfpathlineto{\pgfqpoint{2.424756in}{1.362325in}}%
\pgfpathlineto{\pgfqpoint{2.427274in}{1.386945in}}%
\pgfpathlineto{\pgfqpoint{2.429793in}{1.364216in}}%
\pgfpathlineto{\pgfqpoint{2.432312in}{1.363753in}}%
\pgfpathlineto{\pgfqpoint{2.434831in}{1.382176in}}%
\pgfpathlineto{\pgfqpoint{2.437349in}{1.399807in}}%
\pgfpathlineto{\pgfqpoint{2.439868in}{1.396347in}}%
\pgfpathlineto{\pgfqpoint{2.442387in}{1.364749in}}%
\pgfpathlineto{\pgfqpoint{2.444906in}{1.365385in}}%
\pgfpathlineto{\pgfqpoint{2.447424in}{1.367212in}}%
\pgfpathlineto{\pgfqpoint{2.449943in}{1.370829in}}%
\pgfpathlineto{\pgfqpoint{2.452462in}{1.348069in}}%
\pgfpathlineto{\pgfqpoint{2.454981in}{1.317282in}}%
\pgfpathlineto{\pgfqpoint{2.457499in}{1.323801in}}%
\pgfpathlineto{\pgfqpoint{2.460018in}{1.320830in}}%
\pgfpathlineto{\pgfqpoint{2.462537in}{1.327896in}}%
\pgfpathlineto{\pgfqpoint{2.465056in}{1.336345in}}%
\pgfpathlineto{\pgfqpoint{2.467574in}{1.345186in}}%
\pgfpathlineto{\pgfqpoint{2.470093in}{1.343883in}}%
\pgfpathlineto{\pgfqpoint{2.472612in}{1.343067in}}%
\pgfpathlineto{\pgfqpoint{2.475131in}{1.335473in}}%
\pgfpathlineto{\pgfqpoint{2.477649in}{1.339749in}}%
\pgfpathlineto{\pgfqpoint{2.480168in}{1.351025in}}%
\pgfpathlineto{\pgfqpoint{2.482687in}{1.355935in}}%
\pgfpathlineto{\pgfqpoint{2.485206in}{1.356905in}}%
\pgfpathlineto{\pgfqpoint{2.487724in}{1.338513in}}%
\pgfpathlineto{\pgfqpoint{2.490243in}{1.313595in}}%
\pgfpathlineto{\pgfqpoint{2.492762in}{1.307366in}}%
\pgfpathlineto{\pgfqpoint{2.495281in}{1.297484in}}%
\pgfpathlineto{\pgfqpoint{2.497799in}{1.293484in}}%
\pgfpathlineto{\pgfqpoint{2.500318in}{1.293743in}}%
\pgfpathlineto{\pgfqpoint{2.502837in}{1.296756in}}%
\pgfpathlineto{\pgfqpoint{2.505356in}{1.329592in}}%
\pgfpathlineto{\pgfqpoint{2.507874in}{1.337609in}}%
\pgfpathlineto{\pgfqpoint{2.510393in}{1.350495in}}%
\pgfpathlineto{\pgfqpoint{2.512912in}{1.351116in}}%
\pgfpathlineto{\pgfqpoint{2.515431in}{1.323911in}}%
\pgfpathlineto{\pgfqpoint{2.517949in}{1.327488in}}%
\pgfpathlineto{\pgfqpoint{2.520468in}{1.353280in}}%
\pgfpathlineto{\pgfqpoint{2.522987in}{1.336838in}}%
\pgfpathlineto{\pgfqpoint{2.525506in}{1.337224in}}%
\pgfpathlineto{\pgfqpoint{2.528024in}{1.346036in}}%
\pgfpathlineto{\pgfqpoint{2.530543in}{1.341465in}}%
\pgfpathlineto{\pgfqpoint{2.533062in}{1.360267in}}%
\pgfpathlineto{\pgfqpoint{2.535581in}{1.357676in}}%
\pgfpathlineto{\pgfqpoint{2.538099in}{1.359836in}}%
\pgfpathlineto{\pgfqpoint{2.540618in}{1.349855in}}%
\pgfpathlineto{\pgfqpoint{2.543137in}{1.348474in}}%
\pgfpathlineto{\pgfqpoint{2.545656in}{1.353578in}}%
\pgfpathlineto{\pgfqpoint{2.548174in}{1.352618in}}%
\pgfpathlineto{\pgfqpoint{2.550693in}{1.351173in}}%
\pgfpathlineto{\pgfqpoint{2.553212in}{1.352296in}}%
\pgfpathlineto{\pgfqpoint{2.555731in}{1.366097in}}%
\pgfpathlineto{\pgfqpoint{2.558249in}{1.347636in}}%
\pgfpathlineto{\pgfqpoint{2.560768in}{1.338051in}}%
\pgfpathlineto{\pgfqpoint{2.563287in}{1.329772in}}%
\pgfpathlineto{\pgfqpoint{2.565806in}{1.326135in}}%
\pgfpathlineto{\pgfqpoint{2.568324in}{1.298026in}}%
\pgfpathlineto{\pgfqpoint{2.570843in}{1.309134in}}%
\pgfpathlineto{\pgfqpoint{2.575881in}{1.265852in}}%
\pgfpathlineto{\pgfqpoint{2.578399in}{1.280762in}}%
\pgfpathlineto{\pgfqpoint{2.580918in}{1.275358in}}%
\pgfpathlineto{\pgfqpoint{2.583437in}{1.262241in}}%
\pgfpathlineto{\pgfqpoint{2.585956in}{1.272624in}}%
\pgfpathlineto{\pgfqpoint{2.588474in}{1.254130in}}%
\pgfpathlineto{\pgfqpoint{2.590993in}{1.279220in}}%
\pgfpathlineto{\pgfqpoint{2.593512in}{1.277378in}}%
\pgfpathlineto{\pgfqpoint{2.596031in}{1.276407in}}%
\pgfpathlineto{\pgfqpoint{2.598549in}{1.287348in}}%
\pgfpathlineto{\pgfqpoint{2.601068in}{1.289384in}}%
\pgfpathlineto{\pgfqpoint{2.603587in}{1.293950in}}%
\pgfpathlineto{\pgfqpoint{2.606106in}{1.286488in}}%
\pgfpathlineto{\pgfqpoint{2.608624in}{1.262785in}}%
\pgfpathlineto{\pgfqpoint{2.611143in}{1.267106in}}%
\pgfpathlineto{\pgfqpoint{2.613662in}{1.265571in}}%
\pgfpathlineto{\pgfqpoint{2.616181in}{1.250733in}}%
\pgfpathlineto{\pgfqpoint{2.618699in}{1.255432in}}%
\pgfpathlineto{\pgfqpoint{2.621218in}{1.241138in}}%
\pgfpathlineto{\pgfqpoint{2.623737in}{1.257628in}}%
\pgfpathlineto{\pgfqpoint{2.626256in}{1.239211in}}%
\pgfpathlineto{\pgfqpoint{2.628774in}{1.209065in}}%
\pgfpathlineto{\pgfqpoint{2.631293in}{1.180283in}}%
\pgfpathlineto{\pgfqpoint{2.633812in}{1.177257in}}%
\pgfpathlineto{\pgfqpoint{2.636331in}{1.182307in}}%
\pgfpathlineto{\pgfqpoint{2.638849in}{1.200019in}}%
\pgfpathlineto{\pgfqpoint{2.641368in}{1.183858in}}%
\pgfpathlineto{\pgfqpoint{2.643887in}{1.191307in}}%
\pgfpathlineto{\pgfqpoint{2.646406in}{1.213916in}}%
\pgfpathlineto{\pgfqpoint{2.648924in}{1.196238in}}%
\pgfpathlineto{\pgfqpoint{2.651443in}{1.182856in}}%
\pgfpathlineto{\pgfqpoint{2.653962in}{1.167767in}}%
\pgfpathlineto{\pgfqpoint{2.656481in}{1.183927in}}%
\pgfpathlineto{\pgfqpoint{2.658999in}{1.184111in}}%
\pgfpathlineto{\pgfqpoint{2.661518in}{1.195086in}}%
\pgfpathlineto{\pgfqpoint{2.664037in}{1.181008in}}%
\pgfpathlineto{\pgfqpoint{2.666556in}{1.196944in}}%
\pgfpathlineto{\pgfqpoint{2.669074in}{1.194102in}}%
\pgfpathlineto{\pgfqpoint{2.671593in}{1.191577in}}%
\pgfpathlineto{\pgfqpoint{2.674112in}{1.202395in}}%
\pgfpathlineto{\pgfqpoint{2.676631in}{1.200964in}}%
\pgfpathlineto{\pgfqpoint{2.679149in}{1.205778in}}%
\pgfpathlineto{\pgfqpoint{2.681668in}{1.203835in}}%
\pgfpathlineto{\pgfqpoint{2.684187in}{1.194178in}}%
\pgfpathlineto{\pgfqpoint{2.686706in}{1.206078in}}%
\pgfpathlineto{\pgfqpoint{2.689224in}{1.207349in}}%
\pgfpathlineto{\pgfqpoint{2.691743in}{1.232755in}}%
\pgfpathlineto{\pgfqpoint{2.694262in}{1.246214in}}%
\pgfpathlineto{\pgfqpoint{2.696781in}{1.253444in}}%
\pgfpathlineto{\pgfqpoint{2.699299in}{1.245496in}}%
\pgfpathlineto{\pgfqpoint{2.701818in}{1.246205in}}%
\pgfpathlineto{\pgfqpoint{2.704337in}{1.261076in}}%
\pgfpathlineto{\pgfqpoint{2.706856in}{1.284729in}}%
\pgfpathlineto{\pgfqpoint{2.709374in}{1.304380in}}%
\pgfpathlineto{\pgfqpoint{2.711893in}{1.309651in}}%
\pgfpathlineto{\pgfqpoint{2.714412in}{1.311241in}}%
\pgfpathlineto{\pgfqpoint{2.716931in}{1.328717in}}%
\pgfpathlineto{\pgfqpoint{2.719449in}{1.326199in}}%
\pgfpathlineto{\pgfqpoint{2.721968in}{1.362690in}}%
\pgfpathlineto{\pgfqpoint{2.724487in}{1.359885in}}%
\pgfpathlineto{\pgfqpoint{2.729524in}{1.394064in}}%
\pgfpathlineto{\pgfqpoint{2.732043in}{1.377943in}}%
\pgfpathlineto{\pgfqpoint{2.734562in}{1.387623in}}%
\pgfpathlineto{\pgfqpoint{2.737081in}{1.404148in}}%
\pgfpathlineto{\pgfqpoint{2.739599in}{1.409981in}}%
\pgfpathlineto{\pgfqpoint{2.742118in}{1.417764in}}%
\pgfpathlineto{\pgfqpoint{2.744637in}{1.406503in}}%
\pgfpathlineto{\pgfqpoint{2.747156in}{1.396991in}}%
\pgfpathlineto{\pgfqpoint{2.749674in}{1.396019in}}%
\pgfpathlineto{\pgfqpoint{2.752193in}{1.401177in}}%
\pgfpathlineto{\pgfqpoint{2.754712in}{1.399854in}}%
\pgfpathlineto{\pgfqpoint{2.757231in}{1.403750in}}%
\pgfpathlineto{\pgfqpoint{2.759749in}{1.379414in}}%
\pgfpathlineto{\pgfqpoint{2.762268in}{1.382150in}}%
\pgfpathlineto{\pgfqpoint{2.764787in}{1.381782in}}%
\pgfpathlineto{\pgfqpoint{2.767306in}{1.395925in}}%
\pgfpathlineto{\pgfqpoint{2.769824in}{1.389856in}}%
\pgfpathlineto{\pgfqpoint{2.772343in}{1.408596in}}%
\pgfpathlineto{\pgfqpoint{2.774862in}{1.411227in}}%
\pgfpathlineto{\pgfqpoint{2.777381in}{1.408374in}}%
\pgfpathlineto{\pgfqpoint{2.779899in}{1.390788in}}%
\pgfpathlineto{\pgfqpoint{2.782418in}{1.363413in}}%
\pgfpathlineto{\pgfqpoint{2.784937in}{1.357172in}}%
\pgfpathlineto{\pgfqpoint{2.787456in}{1.363789in}}%
\pgfpathlineto{\pgfqpoint{2.789974in}{1.351495in}}%
\pgfpathlineto{\pgfqpoint{2.792493in}{1.343573in}}%
\pgfpathlineto{\pgfqpoint{2.795012in}{1.323018in}}%
\pgfpathlineto{\pgfqpoint{2.797531in}{1.317791in}}%
\pgfpathlineto{\pgfqpoint{2.800049in}{1.309867in}}%
\pgfpathlineto{\pgfqpoint{2.802568in}{1.328089in}}%
\pgfpathlineto{\pgfqpoint{2.805087in}{1.328366in}}%
\pgfpathlineto{\pgfqpoint{2.807606in}{1.357373in}}%
\pgfpathlineto{\pgfqpoint{2.810124in}{1.361595in}}%
\pgfpathlineto{\pgfqpoint{2.812643in}{1.367888in}}%
\pgfpathlineto{\pgfqpoint{2.815162in}{1.348251in}}%
\pgfpathlineto{\pgfqpoint{2.817681in}{1.343184in}}%
\pgfpathlineto{\pgfqpoint{2.820199in}{1.339468in}}%
\pgfpathlineto{\pgfqpoint{2.822718in}{1.367346in}}%
\pgfpathlineto{\pgfqpoint{2.825237in}{1.401980in}}%
\pgfpathlineto{\pgfqpoint{2.827756in}{1.386668in}}%
\pgfpathlineto{\pgfqpoint{2.830274in}{1.403453in}}%
\pgfpathlineto{\pgfqpoint{2.832793in}{1.389107in}}%
\pgfpathlineto{\pgfqpoint{2.835312in}{1.383517in}}%
\pgfpathlineto{\pgfqpoint{2.837831in}{1.400660in}}%
\pgfpathlineto{\pgfqpoint{2.840349in}{1.378903in}}%
\pgfpathlineto{\pgfqpoint{2.842868in}{1.417996in}}%
\pgfpathlineto{\pgfqpoint{2.845387in}{1.422232in}}%
\pgfpathlineto{\pgfqpoint{2.847906in}{1.428867in}}%
\pgfpathlineto{\pgfqpoint{2.850424in}{1.454424in}}%
\pgfpathlineto{\pgfqpoint{2.852943in}{1.446288in}}%
\pgfpathlineto{\pgfqpoint{2.855462in}{1.477289in}}%
\pgfpathlineto{\pgfqpoint{2.857981in}{1.487860in}}%
\pgfpathlineto{\pgfqpoint{2.860499in}{1.477307in}}%
\pgfpathlineto{\pgfqpoint{2.863018in}{1.457740in}}%
\pgfpathlineto{\pgfqpoint{2.865537in}{1.457349in}}%
\pgfpathlineto{\pgfqpoint{2.868056in}{1.455295in}}%
\pgfpathlineto{\pgfqpoint{2.870574in}{1.454175in}}%
\pgfpathlineto{\pgfqpoint{2.873093in}{1.472622in}}%
\pgfpathlineto{\pgfqpoint{2.875612in}{1.455402in}}%
\pgfpathlineto{\pgfqpoint{2.878131in}{1.450841in}}%
\pgfpathlineto{\pgfqpoint{2.880649in}{1.436985in}}%
\pgfpathlineto{\pgfqpoint{2.883168in}{1.438124in}}%
\pgfpathlineto{\pgfqpoint{2.885687in}{1.425645in}}%
\pgfpathlineto{\pgfqpoint{2.888206in}{1.414360in}}%
\pgfpathlineto{\pgfqpoint{2.890724in}{1.453184in}}%
\pgfpathlineto{\pgfqpoint{2.893243in}{1.435128in}}%
\pgfpathlineto{\pgfqpoint{2.895762in}{1.431079in}}%
\pgfpathlineto{\pgfqpoint{2.898281in}{1.428326in}}%
\pgfpathlineto{\pgfqpoint{2.900799in}{1.416132in}}%
\pgfpathlineto{\pgfqpoint{2.903318in}{1.399534in}}%
\pgfpathlineto{\pgfqpoint{2.905837in}{1.400857in}}%
\pgfpathlineto{\pgfqpoint{2.908356in}{1.442194in}}%
\pgfpathlineto{\pgfqpoint{2.910874in}{1.446505in}}%
\pgfpathlineto{\pgfqpoint{2.913393in}{1.460422in}}%
\pgfpathlineto{\pgfqpoint{2.915912in}{1.441275in}}%
\pgfpathlineto{\pgfqpoint{2.918431in}{1.464652in}}%
\pgfpathlineto{\pgfqpoint{2.920949in}{1.463193in}}%
\pgfpathlineto{\pgfqpoint{2.923468in}{1.470819in}}%
\pgfpathlineto{\pgfqpoint{2.925987in}{1.450426in}}%
\pgfpathlineto{\pgfqpoint{2.928506in}{1.436031in}}%
\pgfpathlineto{\pgfqpoint{2.931024in}{1.431151in}}%
\pgfpathlineto{\pgfqpoint{2.933543in}{1.405730in}}%
\pgfpathlineto{\pgfqpoint{2.936062in}{1.413770in}}%
\pgfpathlineto{\pgfqpoint{2.938581in}{1.404533in}}%
\pgfpathlineto{\pgfqpoint{2.941099in}{1.414333in}}%
\pgfpathlineto{\pgfqpoint{2.943618in}{1.402690in}}%
\pgfpathlineto{\pgfqpoint{2.946137in}{1.384767in}}%
\pgfpathlineto{\pgfqpoint{2.948656in}{1.381903in}}%
\pgfpathlineto{\pgfqpoint{2.951174in}{1.379871in}}%
\pgfpathlineto{\pgfqpoint{2.953693in}{1.390050in}}%
\pgfpathlineto{\pgfqpoint{2.956212in}{1.381143in}}%
\pgfpathlineto{\pgfqpoint{2.958731in}{1.399571in}}%
\pgfpathlineto{\pgfqpoint{2.961249in}{1.400389in}}%
\pgfpathlineto{\pgfqpoint{2.963768in}{1.377737in}}%
\pgfpathlineto{\pgfqpoint{2.966287in}{1.385886in}}%
\pgfpathlineto{\pgfqpoint{2.968806in}{1.389729in}}%
\pgfpathlineto{\pgfqpoint{2.971324in}{1.372054in}}%
\pgfpathlineto{\pgfqpoint{2.973843in}{1.382371in}}%
\pgfpathlineto{\pgfqpoint{2.976362in}{1.381362in}}%
\pgfpathlineto{\pgfqpoint{2.978881in}{1.386013in}}%
\pgfpathlineto{\pgfqpoint{2.981399in}{1.372409in}}%
\pgfpathlineto{\pgfqpoint{2.983918in}{1.368963in}}%
\pgfpathlineto{\pgfqpoint{2.986437in}{1.367479in}}%
\pgfpathlineto{\pgfqpoint{2.988956in}{1.372571in}}%
\pgfpathlineto{\pgfqpoint{2.991474in}{1.375238in}}%
\pgfpathlineto{\pgfqpoint{2.993993in}{1.370134in}}%
\pgfpathlineto{\pgfqpoint{2.996512in}{1.374396in}}%
\pgfpathlineto{\pgfqpoint{2.999031in}{1.356821in}}%
\pgfpathlineto{\pgfqpoint{3.001549in}{1.376014in}}%
\pgfpathlineto{\pgfqpoint{3.004068in}{1.390480in}}%
\pgfpathlineto{\pgfqpoint{3.006587in}{1.408716in}}%
\pgfpathlineto{\pgfqpoint{3.009106in}{1.391412in}}%
\pgfpathlineto{\pgfqpoint{3.011624in}{1.406236in}}%
\pgfpathlineto{\pgfqpoint{3.014143in}{1.417026in}}%
\pgfpathlineto{\pgfqpoint{3.016662in}{1.430991in}}%
\pgfpathlineto{\pgfqpoint{3.019181in}{1.427808in}}%
\pgfpathlineto{\pgfqpoint{3.021699in}{1.426409in}}%
\pgfpathlineto{\pgfqpoint{3.024218in}{1.429468in}}%
\pgfpathlineto{\pgfqpoint{3.026737in}{1.429248in}}%
\pgfpathlineto{\pgfqpoint{3.029256in}{1.418897in}}%
\pgfpathlineto{\pgfqpoint{3.031774in}{1.419572in}}%
\pgfpathlineto{\pgfqpoint{3.034293in}{1.403958in}}%
\pgfpathlineto{\pgfqpoint{3.036812in}{1.418694in}}%
\pgfpathlineto{\pgfqpoint{3.039331in}{1.413781in}}%
\pgfpathlineto{\pgfqpoint{3.041849in}{1.390253in}}%
\pgfpathlineto{\pgfqpoint{3.044368in}{1.407639in}}%
\pgfpathlineto{\pgfqpoint{3.046887in}{1.411801in}}%
\pgfpathlineto{\pgfqpoint{3.049406in}{1.435453in}}%
\pgfpathlineto{\pgfqpoint{3.051924in}{1.408834in}}%
\pgfpathlineto{\pgfqpoint{3.054443in}{1.421613in}}%
\pgfpathlineto{\pgfqpoint{3.056962in}{1.426788in}}%
\pgfpathlineto{\pgfqpoint{3.059481in}{1.410703in}}%
\pgfpathlineto{\pgfqpoint{3.061999in}{1.409600in}}%
\pgfpathlineto{\pgfqpoint{3.064518in}{1.412698in}}%
\pgfpathlineto{\pgfqpoint{3.067037in}{1.421319in}}%
\pgfpathlineto{\pgfqpoint{3.069556in}{1.443060in}}%
\pgfpathlineto{\pgfqpoint{3.072074in}{1.452223in}}%
\pgfpathlineto{\pgfqpoint{3.074593in}{1.442594in}}%
\pgfpathlineto{\pgfqpoint{3.077112in}{1.431693in}}%
\pgfpathlineto{\pgfqpoint{3.079631in}{1.447805in}}%
\pgfpathlineto{\pgfqpoint{3.082149in}{1.452551in}}%
\pgfpathlineto{\pgfqpoint{3.084668in}{1.458837in}}%
\pgfpathlineto{\pgfqpoint{3.087187in}{1.456066in}}%
\pgfpathlineto{\pgfqpoint{3.092224in}{1.466685in}}%
\pgfpathlineto{\pgfqpoint{3.094743in}{1.475755in}}%
\pgfpathlineto{\pgfqpoint{3.097262in}{1.469165in}}%
\pgfpathlineto{\pgfqpoint{3.102299in}{1.498337in}}%
\pgfpathlineto{\pgfqpoint{3.104818in}{1.504935in}}%
\pgfpathlineto{\pgfqpoint{3.107337in}{1.515229in}}%
\pgfpathlineto{\pgfqpoint{3.109856in}{1.494394in}}%
\pgfpathlineto{\pgfqpoint{3.112374in}{1.494229in}}%
\pgfpathlineto{\pgfqpoint{3.114893in}{1.488999in}}%
\pgfpathlineto{\pgfqpoint{3.117412in}{1.488934in}}%
\pgfpathlineto{\pgfqpoint{3.119931in}{1.476838in}}%
\pgfpathlineto{\pgfqpoint{3.122449in}{1.481236in}}%
\pgfpathlineto{\pgfqpoint{3.124968in}{1.473267in}}%
\pgfpathlineto{\pgfqpoint{3.127487in}{1.464633in}}%
\pgfpathlineto{\pgfqpoint{3.130006in}{1.497013in}}%
\pgfpathlineto{\pgfqpoint{3.132524in}{1.509962in}}%
\pgfpathlineto{\pgfqpoint{3.135043in}{1.529264in}}%
\pgfpathlineto{\pgfqpoint{3.137562in}{1.520414in}}%
\pgfpathlineto{\pgfqpoint{3.140081in}{1.520335in}}%
\pgfpathlineto{\pgfqpoint{3.142599in}{1.546657in}}%
\pgfpathlineto{\pgfqpoint{3.145118in}{1.563823in}}%
\pgfpathlineto{\pgfqpoint{3.147637in}{1.552346in}}%
\pgfpathlineto{\pgfqpoint{3.150156in}{1.561254in}}%
\pgfpathlineto{\pgfqpoint{3.152674in}{1.576265in}}%
\pgfpathlineto{\pgfqpoint{3.155193in}{1.568488in}}%
\pgfpathlineto{\pgfqpoint{3.157712in}{1.574125in}}%
\pgfpathlineto{\pgfqpoint{3.160231in}{1.575030in}}%
\pgfpathlineto{\pgfqpoint{3.162749in}{1.572686in}}%
\pgfpathlineto{\pgfqpoint{3.165268in}{1.579759in}}%
\pgfpathlineto{\pgfqpoint{3.167787in}{1.576873in}}%
\pgfpathlineto{\pgfqpoint{3.170306in}{1.588159in}}%
\pgfpathlineto{\pgfqpoint{3.172824in}{1.601191in}}%
\pgfpathlineto{\pgfqpoint{3.175343in}{1.584361in}}%
\pgfpathlineto{\pgfqpoint{3.177862in}{1.594640in}}%
\pgfpathlineto{\pgfqpoint{3.180381in}{1.586280in}}%
\pgfpathlineto{\pgfqpoint{3.182899in}{1.589700in}}%
\pgfpathlineto{\pgfqpoint{3.185418in}{1.573587in}}%
\pgfpathlineto{\pgfqpoint{3.187937in}{1.571060in}}%
\pgfpathlineto{\pgfqpoint{3.190456in}{1.555862in}}%
\pgfpathlineto{\pgfqpoint{3.192974in}{1.558963in}}%
\pgfpathlineto{\pgfqpoint{3.195493in}{1.550043in}}%
\pgfpathlineto{\pgfqpoint{3.198012in}{1.527638in}}%
\pgfpathlineto{\pgfqpoint{3.200531in}{1.553745in}}%
\pgfpathlineto{\pgfqpoint{3.203049in}{1.547333in}}%
\pgfpathlineto{\pgfqpoint{3.205568in}{1.543724in}}%
\pgfpathlineto{\pgfqpoint{3.208087in}{1.551196in}}%
\pgfpathlineto{\pgfqpoint{3.210606in}{1.550937in}}%
\pgfpathlineto{\pgfqpoint{3.213124in}{1.536669in}}%
\pgfpathlineto{\pgfqpoint{3.215643in}{1.506367in}}%
\pgfpathlineto{\pgfqpoint{3.218162in}{1.490519in}}%
\pgfpathlineto{\pgfqpoint{3.220681in}{1.485729in}}%
\pgfpathlineto{\pgfqpoint{3.223199in}{1.484715in}}%
\pgfpathlineto{\pgfqpoint{3.225718in}{1.473068in}}%
\pgfpathlineto{\pgfqpoint{3.228237in}{1.460306in}}%
\pgfpathlineto{\pgfqpoint{3.230756in}{1.452284in}}%
\pgfpathlineto{\pgfqpoint{3.233274in}{1.456385in}}%
\pgfpathlineto{\pgfqpoint{3.235793in}{1.452531in}}%
\pgfpathlineto{\pgfqpoint{3.238312in}{1.484159in}}%
\pgfpathlineto{\pgfqpoint{3.240831in}{1.486206in}}%
\pgfpathlineto{\pgfqpoint{3.243349in}{1.473225in}}%
\pgfpathlineto{\pgfqpoint{3.245868in}{1.477556in}}%
\pgfpathlineto{\pgfqpoint{3.248387in}{1.480006in}}%
\pgfpathlineto{\pgfqpoint{3.250906in}{1.500833in}}%
\pgfpathlineto{\pgfqpoint{3.253424in}{1.499648in}}%
\pgfpathlineto{\pgfqpoint{3.255943in}{1.492438in}}%
\pgfpathlineto{\pgfqpoint{3.258462in}{1.499885in}}%
\pgfpathlineto{\pgfqpoint{3.260981in}{1.476945in}}%
\pgfpathlineto{\pgfqpoint{3.263499in}{1.476779in}}%
\pgfpathlineto{\pgfqpoint{3.266018in}{1.474225in}}%
\pgfpathlineto{\pgfqpoint{3.268537in}{1.485268in}}%
\pgfpathlineto{\pgfqpoint{3.271056in}{1.480772in}}%
\pgfpathlineto{\pgfqpoint{3.273574in}{1.490427in}}%
\pgfpathlineto{\pgfqpoint{3.276093in}{1.495023in}}%
\pgfpathlineto{\pgfqpoint{3.278612in}{1.506385in}}%
\pgfpathlineto{\pgfqpoint{3.281131in}{1.504981in}}%
\pgfpathlineto{\pgfqpoint{3.283649in}{1.530453in}}%
\pgfpathlineto{\pgfqpoint{3.286168in}{1.523835in}}%
\pgfpathlineto{\pgfqpoint{3.288687in}{1.515728in}}%
\pgfpathlineto{\pgfqpoint{3.291206in}{1.524656in}}%
\pgfpathlineto{\pgfqpoint{3.293724in}{1.556338in}}%
\pgfpathlineto{\pgfqpoint{3.296243in}{1.570997in}}%
\pgfpathlineto{\pgfqpoint{3.298762in}{1.575507in}}%
\pgfpathlineto{\pgfqpoint{3.301281in}{1.598443in}}%
\pgfpathlineto{\pgfqpoint{3.303799in}{1.606460in}}%
\pgfpathlineto{\pgfqpoint{3.306318in}{1.601769in}}%
\pgfpathlineto{\pgfqpoint{3.308837in}{1.605565in}}%
\pgfpathlineto{\pgfqpoint{3.311356in}{1.600278in}}%
\pgfpathlineto{\pgfqpoint{3.313874in}{1.583476in}}%
\pgfpathlineto{\pgfqpoint{3.316393in}{1.557632in}}%
\pgfpathlineto{\pgfqpoint{3.318912in}{1.562311in}}%
\pgfpathlineto{\pgfqpoint{3.321431in}{1.532823in}}%
\pgfpathlineto{\pgfqpoint{3.323949in}{1.507446in}}%
\pgfpathlineto{\pgfqpoint{3.326468in}{1.505631in}}%
\pgfpathlineto{\pgfqpoint{3.328987in}{1.504533in}}%
\pgfpathlineto{\pgfqpoint{3.331506in}{1.472369in}}%
\pgfpathlineto{\pgfqpoint{3.334024in}{1.456120in}}%
\pgfpathlineto{\pgfqpoint{3.336543in}{1.451579in}}%
\pgfpathlineto{\pgfqpoint{3.339062in}{1.435459in}}%
\pgfpathlineto{\pgfqpoint{3.341581in}{1.422329in}}%
\pgfpathlineto{\pgfqpoint{3.344099in}{1.427054in}}%
\pgfpathlineto{\pgfqpoint{3.346618in}{1.454418in}}%
\pgfpathlineto{\pgfqpoint{3.349137in}{1.454095in}}%
\pgfpathlineto{\pgfqpoint{3.351656in}{1.451822in}}%
\pgfpathlineto{\pgfqpoint{3.354174in}{1.465153in}}%
\pgfpathlineto{\pgfqpoint{3.356693in}{1.468675in}}%
\pgfpathlineto{\pgfqpoint{3.359212in}{1.478349in}}%
\pgfpathlineto{\pgfqpoint{3.361731in}{1.480153in}}%
\pgfpathlineto{\pgfqpoint{3.364249in}{1.472328in}}%
\pgfpathlineto{\pgfqpoint{3.366768in}{1.460817in}}%
\pgfpathlineto{\pgfqpoint{3.369287in}{1.456225in}}%
\pgfpathlineto{\pgfqpoint{3.371806in}{1.465015in}}%
\pgfpathlineto{\pgfqpoint{3.374324in}{1.460902in}}%
\pgfpathlineto{\pgfqpoint{3.376843in}{1.472493in}}%
\pgfpathlineto{\pgfqpoint{3.379362in}{1.478367in}}%
\pgfpathlineto{\pgfqpoint{3.381881in}{1.499539in}}%
\pgfpathlineto{\pgfqpoint{3.384399in}{1.498695in}}%
\pgfpathlineto{\pgfqpoint{3.386918in}{1.511422in}}%
\pgfpathlineto{\pgfqpoint{3.389437in}{1.499439in}}%
\pgfpathlineto{\pgfqpoint{3.391956in}{1.483377in}}%
\pgfpathlineto{\pgfqpoint{3.394474in}{1.471190in}}%
\pgfpathlineto{\pgfqpoint{3.396993in}{1.452027in}}%
\pgfpathlineto{\pgfqpoint{3.399512in}{1.453357in}}%
\pgfpathlineto{\pgfqpoint{3.402031in}{1.452344in}}%
\pgfpathlineto{\pgfqpoint{3.404549in}{1.437847in}}%
\pgfpathlineto{\pgfqpoint{3.407068in}{1.450998in}}%
\pgfpathlineto{\pgfqpoint{3.409587in}{1.463368in}}%
\pgfpathlineto{\pgfqpoint{3.412106in}{1.429920in}}%
\pgfpathlineto{\pgfqpoint{3.414624in}{1.409152in}}%
\pgfpathlineto{\pgfqpoint{3.417143in}{1.419078in}}%
\pgfpathlineto{\pgfqpoint{3.419662in}{1.419820in}}%
\pgfpathlineto{\pgfqpoint{3.422181in}{1.414129in}}%
\pgfpathlineto{\pgfqpoint{3.424699in}{1.417859in}}%
\pgfpathlineto{\pgfqpoint{3.427218in}{1.437903in}}%
\pgfpathlineto{\pgfqpoint{3.429737in}{1.421865in}}%
\pgfpathlineto{\pgfqpoint{3.432256in}{1.424777in}}%
\pgfpathlineto{\pgfqpoint{3.434774in}{1.415713in}}%
\pgfpathlineto{\pgfqpoint{3.437293in}{1.418326in}}%
\pgfpathlineto{\pgfqpoint{3.439812in}{1.395727in}}%
\pgfpathlineto{\pgfqpoint{3.442331in}{1.390720in}}%
\pgfpathlineto{\pgfqpoint{3.444849in}{1.400892in}}%
\pgfpathlineto{\pgfqpoint{3.447368in}{1.390356in}}%
\pgfpathlineto{\pgfqpoint{3.449887in}{1.386171in}}%
\pgfpathlineto{\pgfqpoint{3.452406in}{1.372444in}}%
\pgfpathlineto{\pgfqpoint{3.454924in}{1.377278in}}%
\pgfpathlineto{\pgfqpoint{3.457443in}{1.388017in}}%
\pgfpathlineto{\pgfqpoint{3.459962in}{1.396055in}}%
\pgfpathlineto{\pgfqpoint{3.462481in}{1.373218in}}%
\pgfpathlineto{\pgfqpoint{3.464999in}{1.374804in}}%
\pgfpathlineto{\pgfqpoint{3.467518in}{1.372420in}}%
\pgfpathlineto{\pgfqpoint{3.470037in}{1.376739in}}%
\pgfpathlineto{\pgfqpoint{3.472556in}{1.389502in}}%
\pgfpathlineto{\pgfqpoint{3.475074in}{1.420630in}}%
\pgfpathlineto{\pgfqpoint{3.477593in}{1.396127in}}%
\pgfpathlineto{\pgfqpoint{3.480112in}{1.408175in}}%
\pgfpathlineto{\pgfqpoint{3.482631in}{1.395771in}}%
\pgfpathlineto{\pgfqpoint{3.485149in}{1.384658in}}%
\pgfpathlineto{\pgfqpoint{3.487668in}{1.395776in}}%
\pgfpathlineto{\pgfqpoint{3.490187in}{1.418730in}}%
\pgfpathlineto{\pgfqpoint{3.492706in}{1.411879in}}%
\pgfpathlineto{\pgfqpoint{3.495224in}{1.399404in}}%
\pgfpathlineto{\pgfqpoint{3.497743in}{1.404073in}}%
\pgfpathlineto{\pgfqpoint{3.500262in}{1.410222in}}%
\pgfpathlineto{\pgfqpoint{3.502781in}{1.403821in}}%
\pgfpathlineto{\pgfqpoint{3.505299in}{1.419654in}}%
\pgfpathlineto{\pgfqpoint{3.507818in}{1.405650in}}%
\pgfpathlineto{\pgfqpoint{3.510337in}{1.412219in}}%
\pgfpathlineto{\pgfqpoint{3.512856in}{1.389412in}}%
\pgfpathlineto{\pgfqpoint{3.515374in}{1.396873in}}%
\pgfpathlineto{\pgfqpoint{3.517893in}{1.398279in}}%
\pgfpathlineto{\pgfqpoint{3.520412in}{1.404715in}}%
\pgfpathlineto{\pgfqpoint{3.522931in}{1.404665in}}%
\pgfpathlineto{\pgfqpoint{3.527968in}{1.390674in}}%
\pgfpathlineto{\pgfqpoint{3.530487in}{1.405672in}}%
\pgfpathlineto{\pgfqpoint{3.533006in}{1.402529in}}%
\pgfpathlineto{\pgfqpoint{3.535524in}{1.412487in}}%
\pgfpathlineto{\pgfqpoint{3.538043in}{1.405405in}}%
\pgfpathlineto{\pgfqpoint{3.540562in}{1.418318in}}%
\pgfpathlineto{\pgfqpoint{3.543081in}{1.419685in}}%
\pgfpathlineto{\pgfqpoint{3.545599in}{1.429593in}}%
\pgfpathlineto{\pgfqpoint{3.548118in}{1.396714in}}%
\pgfpathlineto{\pgfqpoint{3.550637in}{1.419519in}}%
\pgfpathlineto{\pgfqpoint{3.553156in}{1.425969in}}%
\pgfpathlineto{\pgfqpoint{3.555674in}{1.420688in}}%
\pgfpathlineto{\pgfqpoint{3.558193in}{1.389287in}}%
\pgfpathlineto{\pgfqpoint{3.560712in}{1.417385in}}%
\pgfpathlineto{\pgfqpoint{3.563231in}{1.428128in}}%
\pgfpathlineto{\pgfqpoint{3.565749in}{1.456171in}}%
\pgfpathlineto{\pgfqpoint{3.568268in}{1.441180in}}%
\pgfpathlineto{\pgfqpoint{3.570787in}{1.462875in}}%
\pgfpathlineto{\pgfqpoint{3.573306in}{1.511052in}}%
\pgfpathlineto{\pgfqpoint{3.575824in}{1.529966in}}%
\pgfpathlineto{\pgfqpoint{3.578343in}{1.535218in}}%
\pgfpathlineto{\pgfqpoint{3.580862in}{1.538846in}}%
\pgfpathlineto{\pgfqpoint{3.583381in}{1.525339in}}%
\pgfpathlineto{\pgfqpoint{3.585899in}{1.519051in}}%
\pgfpathlineto{\pgfqpoint{3.588418in}{1.489887in}}%
\pgfpathlineto{\pgfqpoint{3.590937in}{1.482150in}}%
\pgfpathlineto{\pgfqpoint{3.593456in}{1.482074in}}%
\pgfpathlineto{\pgfqpoint{3.595974in}{1.513338in}}%
\pgfpathlineto{\pgfqpoint{3.598493in}{1.517568in}}%
\pgfpathlineto{\pgfqpoint{3.601012in}{1.543659in}}%
\pgfpathlineto{\pgfqpoint{3.603531in}{1.544476in}}%
\pgfpathlineto{\pgfqpoint{3.606049in}{1.537344in}}%
\pgfpathlineto{\pgfqpoint{3.608568in}{1.535693in}}%
\pgfpathlineto{\pgfqpoint{3.611087in}{1.527357in}}%
\pgfpathlineto{\pgfqpoint{3.613606in}{1.536852in}}%
\pgfpathlineto{\pgfqpoint{3.616124in}{1.543950in}}%
\pgfpathlineto{\pgfqpoint{3.618643in}{1.547324in}}%
\pgfpathlineto{\pgfqpoint{3.621162in}{1.535742in}}%
\pgfpathlineto{\pgfqpoint{3.623681in}{1.512983in}}%
\pgfpathlineto{\pgfqpoint{3.626199in}{1.496275in}}%
\pgfpathlineto{\pgfqpoint{3.628718in}{1.514460in}}%
\pgfpathlineto{\pgfqpoint{3.631237in}{1.505016in}}%
\pgfpathlineto{\pgfqpoint{3.633756in}{1.489035in}}%
\pgfpathlineto{\pgfqpoint{3.636274in}{1.487327in}}%
\pgfpathlineto{\pgfqpoint{3.638793in}{1.489559in}}%
\pgfpathlineto{\pgfqpoint{3.641312in}{1.496636in}}%
\pgfpathlineto{\pgfqpoint{3.643831in}{1.520359in}}%
\pgfpathlineto{\pgfqpoint{3.646349in}{1.511512in}}%
\pgfpathlineto{\pgfqpoint{3.648868in}{1.501576in}}%
\pgfpathlineto{\pgfqpoint{3.651387in}{1.525864in}}%
\pgfpathlineto{\pgfqpoint{3.653906in}{1.557986in}}%
\pgfpathlineto{\pgfqpoint{3.656424in}{1.570376in}}%
\pgfpathlineto{\pgfqpoint{3.658943in}{1.578846in}}%
\pgfpathlineto{\pgfqpoint{3.661462in}{1.557855in}}%
\pgfpathlineto{\pgfqpoint{3.663981in}{1.575610in}}%
\pgfpathlineto{\pgfqpoint{3.666499in}{1.576739in}}%
\pgfpathlineto{\pgfqpoint{3.669018in}{1.621583in}}%
\pgfpathlineto{\pgfqpoint{3.671537in}{1.593477in}}%
\pgfpathlineto{\pgfqpoint{3.674056in}{1.581243in}}%
\pgfpathlineto{\pgfqpoint{3.676574in}{1.574652in}}%
\pgfpathlineto{\pgfqpoint{3.679093in}{1.583560in}}%
\pgfpathlineto{\pgfqpoint{3.681612in}{1.572481in}}%
\pgfpathlineto{\pgfqpoint{3.684131in}{1.574592in}}%
\pgfpathlineto{\pgfqpoint{3.686649in}{1.586117in}}%
\pgfpathlineto{\pgfqpoint{3.689168in}{1.595801in}}%
\pgfpathlineto{\pgfqpoint{3.691687in}{1.586604in}}%
\pgfpathlineto{\pgfqpoint{3.694206in}{1.575312in}}%
\pgfpathlineto{\pgfqpoint{3.696724in}{1.576652in}}%
\pgfpathlineto{\pgfqpoint{3.699243in}{1.568210in}}%
\pgfpathlineto{\pgfqpoint{3.701762in}{1.551687in}}%
\pgfpathlineto{\pgfqpoint{3.704281in}{1.522592in}}%
\pgfpathlineto{\pgfqpoint{3.706799in}{1.534123in}}%
\pgfpathlineto{\pgfqpoint{3.709318in}{1.513479in}}%
\pgfpathlineto{\pgfqpoint{3.711837in}{1.491535in}}%
\pgfpathlineto{\pgfqpoint{3.714356in}{1.465648in}}%
\pgfpathlineto{\pgfqpoint{3.716874in}{1.478504in}}%
\pgfpathlineto{\pgfqpoint{3.719393in}{1.456960in}}%
\pgfpathlineto{\pgfqpoint{3.721912in}{1.457590in}}%
\pgfpathlineto{\pgfqpoint{3.724431in}{1.483885in}}%
\pgfpathlineto{\pgfqpoint{3.726949in}{1.469037in}}%
\pgfpathlineto{\pgfqpoint{3.729468in}{1.423825in}}%
\pgfpathlineto{\pgfqpoint{3.731987in}{1.449278in}}%
\pgfpathlineto{\pgfqpoint{3.734506in}{1.440744in}}%
\pgfpathlineto{\pgfqpoint{3.737024in}{1.453206in}}%
\pgfpathlineto{\pgfqpoint{3.739543in}{1.430795in}}%
\pgfpathlineto{\pgfqpoint{3.742062in}{1.440642in}}%
\pgfpathlineto{\pgfqpoint{3.744581in}{1.449831in}}%
\pgfpathlineto{\pgfqpoint{3.747099in}{1.462774in}}%
\pgfpathlineto{\pgfqpoint{3.749618in}{1.458339in}}%
\pgfpathlineto{\pgfqpoint{3.752137in}{1.447296in}}%
\pgfpathlineto{\pgfqpoint{3.754656in}{1.465340in}}%
\pgfpathlineto{\pgfqpoint{3.757174in}{1.450587in}}%
\pgfpathlineto{\pgfqpoint{3.759693in}{1.452022in}}%
\pgfpathlineto{\pgfqpoint{3.762212in}{1.447618in}}%
\pgfpathlineto{\pgfqpoint{3.764731in}{1.441848in}}%
\pgfpathlineto{\pgfqpoint{3.767249in}{1.447785in}}%
\pgfpathlineto{\pgfqpoint{3.769768in}{1.412658in}}%
\pgfpathlineto{\pgfqpoint{3.772287in}{1.425045in}}%
\pgfpathlineto{\pgfqpoint{3.774806in}{1.429425in}}%
\pgfpathlineto{\pgfqpoint{3.777324in}{1.428028in}}%
\pgfpathlineto{\pgfqpoint{3.779843in}{1.432974in}}%
\pgfpathlineto{\pgfqpoint{3.782362in}{1.434484in}}%
\pgfpathlineto{\pgfqpoint{3.784881in}{1.465564in}}%
\pgfpathlineto{\pgfqpoint{3.787399in}{1.452079in}}%
\pgfpathlineto{\pgfqpoint{3.789918in}{1.436282in}}%
\pgfpathlineto{\pgfqpoint{3.792437in}{1.436972in}}%
\pgfpathlineto{\pgfqpoint{3.794956in}{1.424327in}}%
\pgfpathlineto{\pgfqpoint{3.797474in}{1.445992in}}%
\pgfpathlineto{\pgfqpoint{3.799993in}{1.471826in}}%
\pgfpathlineto{\pgfqpoint{3.802512in}{1.455325in}}%
\pgfpathlineto{\pgfqpoint{3.805031in}{1.448228in}}%
\pgfpathlineto{\pgfqpoint{3.807549in}{1.468616in}}%
\pgfpathlineto{\pgfqpoint{3.810068in}{1.473287in}}%
\pgfpathlineto{\pgfqpoint{3.812587in}{1.488597in}}%
\pgfpathlineto{\pgfqpoint{3.815106in}{1.476813in}}%
\pgfpathlineto{\pgfqpoint{3.817624in}{1.470838in}}%
\pgfpathlineto{\pgfqpoint{3.820143in}{1.486045in}}%
\pgfpathlineto{\pgfqpoint{3.822662in}{1.485715in}}%
\pgfpathlineto{\pgfqpoint{3.825181in}{1.482509in}}%
\pgfpathlineto{\pgfqpoint{3.827699in}{1.478143in}}%
\pgfpathlineto{\pgfqpoint{3.830218in}{1.450871in}}%
\pgfpathlineto{\pgfqpoint{3.832737in}{1.448787in}}%
\pgfpathlineto{\pgfqpoint{3.835256in}{1.453593in}}%
\pgfpathlineto{\pgfqpoint{3.837774in}{1.468195in}}%
\pgfpathlineto{\pgfqpoint{3.840293in}{1.466069in}}%
\pgfpathlineto{\pgfqpoint{3.842812in}{1.488673in}}%
\pgfpathlineto{\pgfqpoint{3.845331in}{1.480587in}}%
\pgfpathlineto{\pgfqpoint{3.847849in}{1.478577in}}%
\pgfpathlineto{\pgfqpoint{3.850368in}{1.478922in}}%
\pgfpathlineto{\pgfqpoint{3.852887in}{1.468904in}}%
\pgfpathlineto{\pgfqpoint{3.855406in}{1.470895in}}%
\pgfpathlineto{\pgfqpoint{3.857924in}{1.464615in}}%
\pgfpathlineto{\pgfqpoint{3.860443in}{1.477637in}}%
\pgfpathlineto{\pgfqpoint{3.862962in}{1.504758in}}%
\pgfpathlineto{\pgfqpoint{3.865481in}{1.515202in}}%
\pgfpathlineto{\pgfqpoint{3.867999in}{1.490361in}}%
\pgfpathlineto{\pgfqpoint{3.870518in}{1.476647in}}%
\pgfpathlineto{\pgfqpoint{3.873037in}{1.451444in}}%
\pgfpathlineto{\pgfqpoint{3.875556in}{1.432223in}}%
\pgfpathlineto{\pgfqpoint{3.878074in}{1.438679in}}%
\pgfpathlineto{\pgfqpoint{3.880593in}{1.435622in}}%
\pgfpathlineto{\pgfqpoint{3.883112in}{1.447102in}}%
\pgfpathlineto{\pgfqpoint{3.885631in}{1.406005in}}%
\pgfpathlineto{\pgfqpoint{3.888149in}{1.402985in}}%
\pgfpathlineto{\pgfqpoint{3.890668in}{1.417420in}}%
\pgfpathlineto{\pgfqpoint{3.893187in}{1.419121in}}%
\pgfpathlineto{\pgfqpoint{3.895706in}{1.438118in}}%
\pgfpathlineto{\pgfqpoint{3.898224in}{1.419730in}}%
\pgfpathlineto{\pgfqpoint{3.900743in}{1.413558in}}%
\pgfpathlineto{\pgfqpoint{3.903262in}{1.420773in}}%
\pgfpathlineto{\pgfqpoint{3.905781in}{1.422455in}}%
\pgfpathlineto{\pgfqpoint{3.908299in}{1.433952in}}%
\pgfpathlineto{\pgfqpoint{3.910818in}{1.434483in}}%
\pgfpathlineto{\pgfqpoint{3.913337in}{1.461029in}}%
\pgfpathlineto{\pgfqpoint{3.915856in}{1.471153in}}%
\pgfpathlineto{\pgfqpoint{3.918374in}{1.485970in}}%
\pgfpathlineto{\pgfqpoint{3.920893in}{1.481669in}}%
\pgfpathlineto{\pgfqpoint{3.923412in}{1.469353in}}%
\pgfpathlineto{\pgfqpoint{3.925931in}{1.462704in}}%
\pgfpathlineto{\pgfqpoint{3.928449in}{1.493594in}}%
\pgfpathlineto{\pgfqpoint{3.930968in}{1.501815in}}%
\pgfpathlineto{\pgfqpoint{3.933487in}{1.507544in}}%
\pgfpathlineto{\pgfqpoint{3.936006in}{1.517769in}}%
\pgfpathlineto{\pgfqpoint{3.938524in}{1.490137in}}%
\pgfpathlineto{\pgfqpoint{3.941043in}{1.516118in}}%
\pgfpathlineto{\pgfqpoint{3.943562in}{1.508542in}}%
\pgfpathlineto{\pgfqpoint{3.946081in}{1.502804in}}%
\pgfpathlineto{\pgfqpoint{3.948599in}{1.502608in}}%
\pgfpathlineto{\pgfqpoint{3.951118in}{1.511240in}}%
\pgfpathlineto{\pgfqpoint{3.953637in}{1.498942in}}%
\pgfpathlineto{\pgfqpoint{3.956156in}{1.511987in}}%
\pgfpathlineto{\pgfqpoint{3.958674in}{1.499644in}}%
\pgfpathlineto{\pgfqpoint{3.961193in}{1.490509in}}%
\pgfpathlineto{\pgfqpoint{3.963712in}{1.505130in}}%
\pgfpathlineto{\pgfqpoint{3.966231in}{1.482640in}}%
\pgfpathlineto{\pgfqpoint{3.968749in}{1.487195in}}%
\pgfpathlineto{\pgfqpoint{3.971268in}{1.493803in}}%
\pgfpathlineto{\pgfqpoint{3.973787in}{1.527861in}}%
\pgfpathlineto{\pgfqpoint{3.976306in}{1.555653in}}%
\pgfpathlineto{\pgfqpoint{3.978824in}{1.557611in}}%
\pgfpathlineto{\pgfqpoint{3.981343in}{1.526132in}}%
\pgfpathlineto{\pgfqpoint{3.983862in}{1.542297in}}%
\pgfpathlineto{\pgfqpoint{3.986381in}{1.565628in}}%
\pgfpathlineto{\pgfqpoint{3.988899in}{1.554441in}}%
\pgfpathlineto{\pgfqpoint{3.991418in}{1.545775in}}%
\pgfpathlineto{\pgfqpoint{3.993937in}{1.550432in}}%
\pgfpathlineto{\pgfqpoint{3.996456in}{1.538044in}}%
\pgfpathlineto{\pgfqpoint{3.998974in}{1.539384in}}%
\pgfpathlineto{\pgfqpoint{4.001493in}{1.565533in}}%
\pgfpathlineto{\pgfqpoint{4.004012in}{1.593791in}}%
\pgfpathlineto{\pgfqpoint{4.006531in}{1.606326in}}%
\pgfpathlineto{\pgfqpoint{4.009049in}{1.597904in}}%
\pgfpathlineto{\pgfqpoint{4.011568in}{1.619193in}}%
\pgfpathlineto{\pgfqpoint{4.014087in}{1.596979in}}%
\pgfpathlineto{\pgfqpoint{4.016606in}{1.588819in}}%
\pgfpathlineto{\pgfqpoint{4.019124in}{1.592428in}}%
\pgfpathlineto{\pgfqpoint{4.021643in}{1.591824in}}%
\pgfpathlineto{\pgfqpoint{4.024162in}{1.599162in}}%
\pgfpathlineto{\pgfqpoint{4.026681in}{1.613296in}}%
\pgfpathlineto{\pgfqpoint{4.029199in}{1.608884in}}%
\pgfpathlineto{\pgfqpoint{4.031718in}{1.644926in}}%
\pgfpathlineto{\pgfqpoint{4.034237in}{1.652089in}}%
\pgfpathlineto{\pgfqpoint{4.036756in}{1.660942in}}%
\pgfpathlineto{\pgfqpoint{4.039274in}{1.670426in}}%
\pgfpathlineto{\pgfqpoint{4.041793in}{1.682551in}}%
\pgfpathlineto{\pgfqpoint{4.044312in}{1.685592in}}%
\pgfpathlineto{\pgfqpoint{4.046831in}{1.686281in}}%
\pgfpathlineto{\pgfqpoint{4.049349in}{1.636889in}}%
\pgfpathlineto{\pgfqpoint{4.051868in}{1.634300in}}%
\pgfpathlineto{\pgfqpoint{4.054387in}{1.638258in}}%
\pgfpathlineto{\pgfqpoint{4.056906in}{1.619985in}}%
\pgfpathlineto{\pgfqpoint{4.059424in}{1.610038in}}%
\pgfpathlineto{\pgfqpoint{4.064462in}{1.584367in}}%
\pgfpathlineto{\pgfqpoint{4.066981in}{1.589562in}}%
\pgfpathlineto{\pgfqpoint{4.069499in}{1.563874in}}%
\pgfpathlineto{\pgfqpoint{4.072018in}{1.593548in}}%
\pgfpathlineto{\pgfqpoint{4.074537in}{1.613919in}}%
\pgfpathlineto{\pgfqpoint{4.077056in}{1.593778in}}%
\pgfpathlineto{\pgfqpoint{4.079574in}{1.565339in}}%
\pgfpathlineto{\pgfqpoint{4.084612in}{1.586367in}}%
\pgfpathlineto{\pgfqpoint{4.087131in}{1.594117in}}%
\pgfpathlineto{\pgfqpoint{4.089649in}{1.591536in}}%
\pgfpathlineto{\pgfqpoint{4.092168in}{1.585079in}}%
\pgfpathlineto{\pgfqpoint{4.094687in}{1.589090in}}%
\pgfpathlineto{\pgfqpoint{4.097206in}{1.605343in}}%
\pgfpathlineto{\pgfqpoint{4.099724in}{1.586597in}}%
\pgfpathlineto{\pgfqpoint{4.102243in}{1.595631in}}%
\pgfpathlineto{\pgfqpoint{4.104762in}{1.574470in}}%
\pgfpathlineto{\pgfqpoint{4.107281in}{1.597985in}}%
\pgfpathlineto{\pgfqpoint{4.109799in}{1.631396in}}%
\pgfpathlineto{\pgfqpoint{4.112318in}{1.618741in}}%
\pgfpathlineto{\pgfqpoint{4.114837in}{1.622466in}}%
\pgfpathlineto{\pgfqpoint{4.117356in}{1.608352in}}%
\pgfpathlineto{\pgfqpoint{4.119874in}{1.611697in}}%
\pgfpathlineto{\pgfqpoint{4.122393in}{1.618596in}}%
\pgfpathlineto{\pgfqpoint{4.124912in}{1.624743in}}%
\pgfpathlineto{\pgfqpoint{4.127431in}{1.634162in}}%
\pgfpathlineto{\pgfqpoint{4.129949in}{1.639194in}}%
\pgfpathlineto{\pgfqpoint{4.132468in}{1.626854in}}%
\pgfpathlineto{\pgfqpoint{4.134987in}{1.636474in}}%
\pgfpathlineto{\pgfqpoint{4.137506in}{1.642701in}}%
\pgfpathlineto{\pgfqpoint{4.140024in}{1.634010in}}%
\pgfpathlineto{\pgfqpoint{4.142543in}{1.627264in}}%
\pgfpathlineto{\pgfqpoint{4.145062in}{1.639557in}}%
\pgfpathlineto{\pgfqpoint{4.147581in}{1.648543in}}%
\pgfpathlineto{\pgfqpoint{4.150099in}{1.623131in}}%
\pgfpathlineto{\pgfqpoint{4.152618in}{1.632630in}}%
\pgfpathlineto{\pgfqpoint{4.155137in}{1.621742in}}%
\pgfpathlineto{\pgfqpoint{4.157656in}{1.641298in}}%
\pgfpathlineto{\pgfqpoint{4.160174in}{1.622714in}}%
\pgfpathlineto{\pgfqpoint{4.162693in}{1.651335in}}%
\pgfpathlineto{\pgfqpoint{4.165212in}{1.657997in}}%
\pgfpathlineto{\pgfqpoint{4.167731in}{1.670760in}}%
\pgfpathlineto{\pgfqpoint{4.170249in}{1.657133in}}%
\pgfpathlineto{\pgfqpoint{4.172768in}{1.678940in}}%
\pgfpathlineto{\pgfqpoint{4.175287in}{1.683081in}}%
\pgfpathlineto{\pgfqpoint{4.177806in}{1.668377in}}%
\pgfpathlineto{\pgfqpoint{4.180324in}{1.678261in}}%
\pgfpathlineto{\pgfqpoint{4.182843in}{1.709219in}}%
\pgfpathlineto{\pgfqpoint{4.185362in}{1.705905in}}%
\pgfpathlineto{\pgfqpoint{4.190399in}{1.673163in}}%
\pgfpathlineto{\pgfqpoint{4.192918in}{1.697328in}}%
\pgfpathlineto{\pgfqpoint{4.195437in}{1.728804in}}%
\pgfpathlineto{\pgfqpoint{4.197956in}{1.731845in}}%
\pgfpathlineto{\pgfqpoint{4.200474in}{1.714682in}}%
\pgfpathlineto{\pgfqpoint{4.202993in}{1.719650in}}%
\pgfpathlineto{\pgfqpoint{4.205512in}{1.734928in}}%
\pgfpathlineto{\pgfqpoint{4.208031in}{1.747279in}}%
\pgfpathlineto{\pgfqpoint{4.210549in}{1.756701in}}%
\pgfpathlineto{\pgfqpoint{4.213068in}{1.759119in}}%
\pgfpathlineto{\pgfqpoint{4.215587in}{1.777124in}}%
\pgfpathlineto{\pgfqpoint{4.218106in}{1.762160in}}%
\pgfpathlineto{\pgfqpoint{4.220624in}{1.782369in}}%
\pgfpathlineto{\pgfqpoint{4.223143in}{1.786190in}}%
\pgfpathlineto{\pgfqpoint{4.225662in}{1.791045in}}%
\pgfpathlineto{\pgfqpoint{4.228181in}{1.773291in}}%
\pgfpathlineto{\pgfqpoint{4.230699in}{1.745145in}}%
\pgfpathlineto{\pgfqpoint{4.233218in}{1.731377in}}%
\pgfpathlineto{\pgfqpoint{4.235737in}{1.739674in}}%
\pgfpathlineto{\pgfqpoint{4.238256in}{1.748617in}}%
\pgfpathlineto{\pgfqpoint{4.240774in}{1.739111in}}%
\pgfpathlineto{\pgfqpoint{4.243293in}{1.732929in}}%
\pgfpathlineto{\pgfqpoint{4.245812in}{1.732310in}}%
\pgfpathlineto{\pgfqpoint{4.248331in}{1.734775in}}%
\pgfpathlineto{\pgfqpoint{4.250849in}{1.749083in}}%
\pgfpathlineto{\pgfqpoint{4.253368in}{1.752569in}}%
\pgfpathlineto{\pgfqpoint{4.255887in}{1.757476in}}%
\pgfpathlineto{\pgfqpoint{4.258406in}{1.768788in}}%
\pgfpathlineto{\pgfqpoint{4.260924in}{1.776257in}}%
\pgfpathlineto{\pgfqpoint{4.263443in}{1.807753in}}%
\pgfpathlineto{\pgfqpoint{4.265962in}{1.805121in}}%
\pgfpathlineto{\pgfqpoint{4.268481in}{1.838496in}}%
\pgfpathlineto{\pgfqpoint{4.270999in}{1.820329in}}%
\pgfpathlineto{\pgfqpoint{4.273518in}{1.838725in}}%
\pgfpathlineto{\pgfqpoint{4.276037in}{1.833115in}}%
\pgfpathlineto{\pgfqpoint{4.278556in}{1.810212in}}%
\pgfpathlineto{\pgfqpoint{4.281074in}{1.825658in}}%
\pgfpathlineto{\pgfqpoint{4.283593in}{1.813731in}}%
\pgfpathlineto{\pgfqpoint{4.286112in}{1.808567in}}%
\pgfpathlineto{\pgfqpoint{4.288631in}{1.772085in}}%
\pgfpathlineto{\pgfqpoint{4.291149in}{1.780545in}}%
\pgfpathlineto{\pgfqpoint{4.293668in}{1.790632in}}%
\pgfpathlineto{\pgfqpoint{4.296187in}{1.797529in}}%
\pgfpathlineto{\pgfqpoint{4.298706in}{1.798753in}}%
\pgfpathlineto{\pgfqpoint{4.301224in}{1.807873in}}%
\pgfpathlineto{\pgfqpoint{4.303743in}{1.798138in}}%
\pgfpathlineto{\pgfqpoint{4.306262in}{1.782437in}}%
\pgfpathlineto{\pgfqpoint{4.311299in}{1.818256in}}%
\pgfpathlineto{\pgfqpoint{4.313818in}{1.819199in}}%
\pgfpathlineto{\pgfqpoint{4.316337in}{1.803705in}}%
\pgfpathlineto{\pgfqpoint{4.321374in}{1.746802in}}%
\pgfpathlineto{\pgfqpoint{4.323893in}{1.772235in}}%
\pgfpathlineto{\pgfqpoint{4.326412in}{1.763377in}}%
\pgfpathlineto{\pgfqpoint{4.328931in}{1.752315in}}%
\pgfpathlineto{\pgfqpoint{4.331449in}{1.726039in}}%
\pgfpathlineto{\pgfqpoint{4.333968in}{1.737652in}}%
\pgfpathlineto{\pgfqpoint{4.336487in}{1.738356in}}%
\pgfpathlineto{\pgfqpoint{4.339006in}{1.721898in}}%
\pgfpathlineto{\pgfqpoint{4.341524in}{1.728916in}}%
\pgfpathlineto{\pgfqpoint{4.344043in}{1.734670in}}%
\pgfpathlineto{\pgfqpoint{4.346562in}{1.722269in}}%
\pgfpathlineto{\pgfqpoint{4.349081in}{1.721217in}}%
\pgfpathlineto{\pgfqpoint{4.351599in}{1.743894in}}%
\pgfpathlineto{\pgfqpoint{4.354118in}{1.734397in}}%
\pgfpathlineto{\pgfqpoint{4.356637in}{1.732157in}}%
\pgfpathlineto{\pgfqpoint{4.359156in}{1.757272in}}%
\pgfpathlineto{\pgfqpoint{4.361674in}{1.749004in}}%
\pgfpathlineto{\pgfqpoint{4.364193in}{1.746402in}}%
\pgfpathlineto{\pgfqpoint{4.366712in}{1.728104in}}%
\pgfpathlineto{\pgfqpoint{4.369231in}{1.723558in}}%
\pgfpathlineto{\pgfqpoint{4.371749in}{1.727123in}}%
\pgfpathlineto{\pgfqpoint{4.374268in}{1.729984in}}%
\pgfpathlineto{\pgfqpoint{4.376787in}{1.728226in}}%
\pgfpathlineto{\pgfqpoint{4.379306in}{1.722888in}}%
\pgfpathlineto{\pgfqpoint{4.381824in}{1.729761in}}%
\pgfpathlineto{\pgfqpoint{4.384343in}{1.729809in}}%
\pgfpathlineto{\pgfqpoint{4.386862in}{1.732901in}}%
\pgfpathlineto{\pgfqpoint{4.389381in}{1.729846in}}%
\pgfpathlineto{\pgfqpoint{4.391899in}{1.719585in}}%
\pgfpathlineto{\pgfqpoint{4.394418in}{1.716520in}}%
\pgfpathlineto{\pgfqpoint{4.396937in}{1.739536in}}%
\pgfpathlineto{\pgfqpoint{4.399456in}{1.735139in}}%
\pgfpathlineto{\pgfqpoint{4.401974in}{1.752151in}}%
\pgfpathlineto{\pgfqpoint{4.404493in}{1.736380in}}%
\pgfpathlineto{\pgfqpoint{4.407012in}{1.735819in}}%
\pgfpathlineto{\pgfqpoint{4.409531in}{1.715626in}}%
\pgfpathlineto{\pgfqpoint{4.412049in}{1.711556in}}%
\pgfpathlineto{\pgfqpoint{4.414568in}{1.711710in}}%
\pgfpathlineto{\pgfqpoint{4.417087in}{1.707016in}}%
\pgfpathlineto{\pgfqpoint{4.419606in}{1.732505in}}%
\pgfpathlineto{\pgfqpoint{4.422124in}{1.760221in}}%
\pgfpathlineto{\pgfqpoint{4.424643in}{1.751962in}}%
\pgfpathlineto{\pgfqpoint{4.427162in}{1.763260in}}%
\pgfpathlineto{\pgfqpoint{4.429681in}{1.755784in}}%
\pgfpathlineto{\pgfqpoint{4.432199in}{1.752319in}}%
\pgfpathlineto{\pgfqpoint{4.434718in}{1.771884in}}%
\pgfpathlineto{\pgfqpoint{4.437237in}{1.768856in}}%
\pgfpathlineto{\pgfqpoint{4.439756in}{1.798664in}}%
\pgfpathlineto{\pgfqpoint{4.442274in}{1.773695in}}%
\pgfpathlineto{\pgfqpoint{4.444793in}{1.780963in}}%
\pgfpathlineto{\pgfqpoint{4.447312in}{1.785447in}}%
\pgfpathlineto{\pgfqpoint{4.449831in}{1.786665in}}%
\pgfpathlineto{\pgfqpoint{4.452349in}{1.764673in}}%
\pgfpathlineto{\pgfqpoint{4.454868in}{1.778260in}}%
\pgfpathlineto{\pgfqpoint{4.457387in}{1.778771in}}%
\pgfpathlineto{\pgfqpoint{4.459906in}{1.778379in}}%
\pgfpathlineto{\pgfqpoint{4.462424in}{1.787546in}}%
\pgfpathlineto{\pgfqpoint{4.464943in}{1.789413in}}%
\pgfpathlineto{\pgfqpoint{4.467462in}{1.753701in}}%
\pgfpathlineto{\pgfqpoint{4.469981in}{1.764408in}}%
\pgfpathlineto{\pgfqpoint{4.472499in}{1.767761in}}%
\pgfpathlineto{\pgfqpoint{4.475018in}{1.750658in}}%
\pgfpathlineto{\pgfqpoint{4.477537in}{1.766878in}}%
\pgfpathlineto{\pgfqpoint{4.480056in}{1.767303in}}%
\pgfpathlineto{\pgfqpoint{4.482574in}{1.776976in}}%
\pgfpathlineto{\pgfqpoint{4.485093in}{1.814269in}}%
\pgfpathlineto{\pgfqpoint{4.487612in}{1.830800in}}%
\pgfpathlineto{\pgfqpoint{4.490131in}{1.816088in}}%
\pgfpathlineto{\pgfqpoint{4.492649in}{1.830545in}}%
\pgfpathlineto{\pgfqpoint{4.495168in}{1.833771in}}%
\pgfpathlineto{\pgfqpoint{4.497687in}{1.853649in}}%
\pgfpathlineto{\pgfqpoint{4.500206in}{1.849257in}}%
\pgfpathlineto{\pgfqpoint{4.502724in}{1.833715in}}%
\pgfpathlineto{\pgfqpoint{4.505243in}{1.832014in}}%
\pgfpathlineto{\pgfqpoint{4.507762in}{1.842307in}}%
\pgfpathlineto{\pgfqpoint{4.510281in}{1.821243in}}%
\pgfpathlineto{\pgfqpoint{4.512799in}{1.820462in}}%
\pgfpathlineto{\pgfqpoint{4.515318in}{1.824899in}}%
\pgfpathlineto{\pgfqpoint{4.517837in}{1.784061in}}%
\pgfpathlineto{\pgfqpoint{4.520356in}{1.790691in}}%
\pgfpathlineto{\pgfqpoint{4.522874in}{1.780016in}}%
\pgfpathlineto{\pgfqpoint{4.525393in}{1.783232in}}%
\pgfpathlineto{\pgfqpoint{4.527912in}{1.802731in}}%
\pgfpathlineto{\pgfqpoint{4.530431in}{1.799254in}}%
\pgfpathlineto{\pgfqpoint{4.532949in}{1.805442in}}%
\pgfpathlineto{\pgfqpoint{4.535468in}{1.820037in}}%
\pgfpathlineto{\pgfqpoint{4.537987in}{1.804756in}}%
\pgfpathlineto{\pgfqpoint{4.540506in}{1.794142in}}%
\pgfpathlineto{\pgfqpoint{4.543024in}{1.781606in}}%
\pgfpathlineto{\pgfqpoint{4.545543in}{1.763141in}}%
\pgfpathlineto{\pgfqpoint{4.548062in}{1.770600in}}%
\pgfpathlineto{\pgfqpoint{4.550581in}{1.758442in}}%
\pgfpathlineto{\pgfqpoint{4.553099in}{1.758881in}}%
\pgfpathlineto{\pgfqpoint{4.555618in}{1.742327in}}%
\pgfpathlineto{\pgfqpoint{4.558137in}{1.752111in}}%
\pgfpathlineto{\pgfqpoint{4.560656in}{1.746325in}}%
\pgfpathlineto{\pgfqpoint{4.563174in}{1.717378in}}%
\pgfpathlineto{\pgfqpoint{4.565693in}{1.728421in}}%
\pgfpathlineto{\pgfqpoint{4.568212in}{1.722817in}}%
\pgfpathlineto{\pgfqpoint{4.570731in}{1.731202in}}%
\pgfpathlineto{\pgfqpoint{4.573249in}{1.743353in}}%
\pgfpathlineto{\pgfqpoint{4.575768in}{1.727796in}}%
\pgfpathlineto{\pgfqpoint{4.578287in}{1.726277in}}%
\pgfpathlineto{\pgfqpoint{4.580806in}{1.723618in}}%
\pgfpathlineto{\pgfqpoint{4.583324in}{1.727720in}}%
\pgfpathlineto{\pgfqpoint{4.585843in}{1.708730in}}%
\pgfpathlineto{\pgfqpoint{4.588362in}{1.716289in}}%
\pgfpathlineto{\pgfqpoint{4.590881in}{1.754984in}}%
\pgfpathlineto{\pgfqpoint{4.593399in}{1.795633in}}%
\pgfpathlineto{\pgfqpoint{4.595918in}{1.807974in}}%
\pgfpathlineto{\pgfqpoint{4.598437in}{1.805105in}}%
\pgfpathlineto{\pgfqpoint{4.600956in}{1.806160in}}%
\pgfpathlineto{\pgfqpoint{4.603474in}{1.806442in}}%
\pgfpathlineto{\pgfqpoint{4.605993in}{1.823560in}}%
\pgfpathlineto{\pgfqpoint{4.608512in}{1.814386in}}%
\pgfpathlineto{\pgfqpoint{4.611031in}{1.809139in}}%
\pgfpathlineto{\pgfqpoint{4.613549in}{1.792728in}}%
\pgfpathlineto{\pgfqpoint{4.616068in}{1.786221in}}%
\pgfpathlineto{\pgfqpoint{4.618587in}{1.788774in}}%
\pgfpathlineto{\pgfqpoint{4.621106in}{1.801462in}}%
\pgfpathlineto{\pgfqpoint{4.623624in}{1.785927in}}%
\pgfpathlineto{\pgfqpoint{4.626143in}{1.767549in}}%
\pgfpathlineto{\pgfqpoint{4.628662in}{1.766579in}}%
\pgfpathlineto{\pgfqpoint{4.631181in}{1.771738in}}%
\pgfpathlineto{\pgfqpoint{4.633699in}{1.763776in}}%
\pgfpathlineto{\pgfqpoint{4.636218in}{1.763083in}}%
\pgfpathlineto{\pgfqpoint{4.638737in}{1.758352in}}%
\pgfpathlineto{\pgfqpoint{4.641256in}{1.779582in}}%
\pgfpathlineto{\pgfqpoint{4.643774in}{1.810384in}}%
\pgfpathlineto{\pgfqpoint{4.646293in}{1.836082in}}%
\pgfpathlineto{\pgfqpoint{4.648812in}{1.828119in}}%
\pgfpathlineto{\pgfqpoint{4.651331in}{1.869262in}}%
\pgfpathlineto{\pgfqpoint{4.653849in}{1.865159in}}%
\pgfpathlineto{\pgfqpoint{4.656368in}{1.865888in}}%
\pgfpathlineto{\pgfqpoint{4.658887in}{1.864827in}}%
\pgfpathlineto{\pgfqpoint{4.661406in}{1.887026in}}%
\pgfpathlineto{\pgfqpoint{4.663924in}{1.879903in}}%
\pgfpathlineto{\pgfqpoint{4.666443in}{1.885687in}}%
\pgfpathlineto{\pgfqpoint{4.668962in}{1.880295in}}%
\pgfpathlineto{\pgfqpoint{4.671481in}{1.855921in}}%
\pgfpathlineto{\pgfqpoint{4.673999in}{1.844541in}}%
\pgfpathlineto{\pgfqpoint{4.676518in}{1.842260in}}%
\pgfpathlineto{\pgfqpoint{4.679037in}{1.814150in}}%
\pgfpathlineto{\pgfqpoint{4.681556in}{1.821566in}}%
\pgfpathlineto{\pgfqpoint{4.684074in}{1.835305in}}%
\pgfpathlineto{\pgfqpoint{4.686593in}{1.846110in}}%
\pgfpathlineto{\pgfqpoint{4.689112in}{1.828265in}}%
\pgfpathlineto{\pgfqpoint{4.691631in}{1.827496in}}%
\pgfpathlineto{\pgfqpoint{4.694149in}{1.809194in}}%
\pgfpathlineto{\pgfqpoint{4.696668in}{1.808910in}}%
\pgfpathlineto{\pgfqpoint{4.699187in}{1.830464in}}%
\pgfpathlineto{\pgfqpoint{4.701706in}{1.807184in}}%
\pgfpathlineto{\pgfqpoint{4.704224in}{1.811439in}}%
\pgfpathlineto{\pgfqpoint{4.706743in}{1.796227in}}%
\pgfpathlineto{\pgfqpoint{4.709262in}{1.809184in}}%
\pgfpathlineto{\pgfqpoint{4.711781in}{1.805261in}}%
\pgfpathlineto{\pgfqpoint{4.714299in}{1.818315in}}%
\pgfpathlineto{\pgfqpoint{4.716818in}{1.850280in}}%
\pgfpathlineto{\pgfqpoint{4.719337in}{1.858977in}}%
\pgfpathlineto{\pgfqpoint{4.721856in}{1.910947in}}%
\pgfpathlineto{\pgfqpoint{4.724374in}{1.876462in}}%
\pgfpathlineto{\pgfqpoint{4.726893in}{1.853874in}}%
\pgfpathlineto{\pgfqpoint{4.729412in}{1.856900in}}%
\pgfpathlineto{\pgfqpoint{4.731931in}{1.856601in}}%
\pgfpathlineto{\pgfqpoint{4.734449in}{1.857065in}}%
\pgfpathlineto{\pgfqpoint{4.736968in}{1.861850in}}%
\pgfpathlineto{\pgfqpoint{4.739487in}{1.869653in}}%
\pgfpathlineto{\pgfqpoint{4.742006in}{1.902035in}}%
\pgfpathlineto{\pgfqpoint{4.744524in}{1.910135in}}%
\pgfpathlineto{\pgfqpoint{4.747043in}{1.905316in}}%
\pgfpathlineto{\pgfqpoint{4.749562in}{1.936984in}}%
\pgfpathlineto{\pgfqpoint{4.752081in}{1.964123in}}%
\pgfpathlineto{\pgfqpoint{4.754599in}{1.965706in}}%
\pgfpathlineto{\pgfqpoint{4.757118in}{1.966119in}}%
\pgfpathlineto{\pgfqpoint{4.759637in}{1.950532in}}%
\pgfpathlineto{\pgfqpoint{4.762156in}{1.924801in}}%
\pgfpathlineto{\pgfqpoint{4.764674in}{1.943252in}}%
\pgfpathlineto{\pgfqpoint{4.767193in}{1.924325in}}%
\pgfpathlineto{\pgfqpoint{4.769712in}{1.907494in}}%
\pgfpathlineto{\pgfqpoint{4.772231in}{1.903350in}}%
\pgfpathlineto{\pgfqpoint{4.774749in}{1.882924in}}%
\pgfpathlineto{\pgfqpoint{4.777268in}{1.894625in}}%
\pgfpathlineto{\pgfqpoint{4.779787in}{1.893466in}}%
\pgfpathlineto{\pgfqpoint{4.782306in}{1.886356in}}%
\pgfpathlineto{\pgfqpoint{4.784824in}{1.889946in}}%
\pgfpathlineto{\pgfqpoint{4.787343in}{1.898478in}}%
\pgfpathlineto{\pgfqpoint{4.789862in}{1.926101in}}%
\pgfpathlineto{\pgfqpoint{4.792381in}{1.920613in}}%
\pgfpathlineto{\pgfqpoint{4.794899in}{1.906574in}}%
\pgfpathlineto{\pgfqpoint{4.797418in}{1.925040in}}%
\pgfpathlineto{\pgfqpoint{4.799937in}{1.958320in}}%
\pgfpathlineto{\pgfqpoint{4.802456in}{1.961510in}}%
\pgfpathlineto{\pgfqpoint{4.804974in}{1.970949in}}%
\pgfpathlineto{\pgfqpoint{4.807493in}{1.956794in}}%
\pgfpathlineto{\pgfqpoint{4.810012in}{1.970995in}}%
\pgfpathlineto{\pgfqpoint{4.812531in}{1.971446in}}%
\pgfpathlineto{\pgfqpoint{4.815049in}{1.947460in}}%
\pgfpathlineto{\pgfqpoint{4.817568in}{1.961439in}}%
\pgfpathlineto{\pgfqpoint{4.820087in}{1.967823in}}%
\pgfpathlineto{\pgfqpoint{4.822606in}{1.975894in}}%
\pgfpathlineto{\pgfqpoint{4.825124in}{1.977899in}}%
\pgfpathlineto{\pgfqpoint{4.827643in}{1.971192in}}%
\pgfpathlineto{\pgfqpoint{4.830162in}{1.965693in}}%
\pgfpathlineto{\pgfqpoint{4.832681in}{1.983852in}}%
\pgfpathlineto{\pgfqpoint{4.835199in}{1.978139in}}%
\pgfpathlineto{\pgfqpoint{4.837718in}{1.978650in}}%
\pgfpathlineto{\pgfqpoint{4.840237in}{1.967272in}}%
\pgfpathlineto{\pgfqpoint{4.842756in}{1.974908in}}%
\pgfpathlineto{\pgfqpoint{4.845274in}{1.969061in}}%
\pgfpathlineto{\pgfqpoint{4.847793in}{1.961214in}}%
\pgfpathlineto{\pgfqpoint{4.850312in}{1.958752in}}%
\pgfpathlineto{\pgfqpoint{4.852831in}{1.980250in}}%
\pgfpathlineto{\pgfqpoint{4.855349in}{1.963361in}}%
\pgfpathlineto{\pgfqpoint{4.857868in}{1.971916in}}%
\pgfpathlineto{\pgfqpoint{4.860387in}{1.966280in}}%
\pgfpathlineto{\pgfqpoint{4.862906in}{1.966467in}}%
\pgfpathlineto{\pgfqpoint{4.865424in}{1.967309in}}%
\pgfpathlineto{\pgfqpoint{4.867943in}{1.955779in}}%
\pgfpathlineto{\pgfqpoint{4.870462in}{1.974013in}}%
\pgfpathlineto{\pgfqpoint{4.872981in}{2.026023in}}%
\pgfpathlineto{\pgfqpoint{4.875499in}{2.047299in}}%
\pgfpathlineto{\pgfqpoint{4.878018in}{2.048379in}}%
\pgfpathlineto{\pgfqpoint{4.880537in}{2.057432in}}%
\pgfpathlineto{\pgfqpoint{4.883056in}{2.065573in}}%
\pgfpathlineto{\pgfqpoint{4.885574in}{2.062457in}}%
\pgfpathlineto{\pgfqpoint{4.888093in}{2.052146in}}%
\pgfpathlineto{\pgfqpoint{4.890612in}{2.038172in}}%
\pgfpathlineto{\pgfqpoint{4.893131in}{2.026803in}}%
\pgfpathlineto{\pgfqpoint{4.895649in}{2.029914in}}%
\pgfpathlineto{\pgfqpoint{4.898168in}{1.989540in}}%
\pgfpathlineto{\pgfqpoint{4.900687in}{1.983825in}}%
\pgfpathlineto{\pgfqpoint{4.903206in}{1.976387in}}%
\pgfpathlineto{\pgfqpoint{4.905724in}{1.975511in}}%
\pgfpathlineto{\pgfqpoint{4.908243in}{1.991773in}}%
\pgfpathlineto{\pgfqpoint{4.910762in}{1.979003in}}%
\pgfpathlineto{\pgfqpoint{4.913281in}{1.977866in}}%
\pgfpathlineto{\pgfqpoint{4.915799in}{1.987521in}}%
\pgfpathlineto{\pgfqpoint{4.918318in}{1.977396in}}%
\pgfpathlineto{\pgfqpoint{4.920837in}{1.985589in}}%
\pgfpathlineto{\pgfqpoint{4.923356in}{1.983736in}}%
\pgfpathlineto{\pgfqpoint{4.925874in}{2.016408in}}%
\pgfpathlineto{\pgfqpoint{4.928393in}{1.971687in}}%
\pgfpathlineto{\pgfqpoint{4.930912in}{1.934362in}}%
\pgfpathlineto{\pgfqpoint{4.933431in}{1.941624in}}%
\pgfpathlineto{\pgfqpoint{4.935949in}{1.967709in}}%
\pgfpathlineto{\pgfqpoint{4.938468in}{1.971704in}}%
\pgfpathlineto{\pgfqpoint{4.940987in}{1.991132in}}%
\pgfpathlineto{\pgfqpoint{4.943506in}{2.001581in}}%
\pgfpathlineto{\pgfqpoint{4.946024in}{1.983890in}}%
\pgfpathlineto{\pgfqpoint{4.948543in}{1.941272in}}%
\pgfpathlineto{\pgfqpoint{4.951062in}{1.937680in}}%
\pgfpathlineto{\pgfqpoint{4.956099in}{1.917032in}}%
\pgfpathlineto{\pgfqpoint{4.958618in}{1.939243in}}%
\pgfpathlineto{\pgfqpoint{4.961137in}{1.941578in}}%
\pgfpathlineto{\pgfqpoint{4.963656in}{1.978228in}}%
\pgfpathlineto{\pgfqpoint{4.966174in}{1.966758in}}%
\pgfpathlineto{\pgfqpoint{4.968693in}{1.953436in}}%
\pgfpathlineto{\pgfqpoint{4.971212in}{1.965748in}}%
\pgfpathlineto{\pgfqpoint{4.973731in}{1.975991in}}%
\pgfpathlineto{\pgfqpoint{4.976249in}{1.963839in}}%
\pgfpathlineto{\pgfqpoint{4.978768in}{1.937944in}}%
\pgfpathlineto{\pgfqpoint{4.981287in}{1.954032in}}%
\pgfpathlineto{\pgfqpoint{4.983806in}{1.961398in}}%
\pgfpathlineto{\pgfqpoint{4.986324in}{1.948342in}}%
\pgfpathlineto{\pgfqpoint{4.988843in}{1.960817in}}%
\pgfpathlineto{\pgfqpoint{4.991362in}{1.962842in}}%
\pgfpathlineto{\pgfqpoint{4.993881in}{1.971894in}}%
\pgfpathlineto{\pgfqpoint{4.996399in}{1.961525in}}%
\pgfpathlineto{\pgfqpoint{4.998918in}{1.955017in}}%
\pgfpathlineto{\pgfqpoint{5.001437in}{1.967290in}}%
\pgfpathlineto{\pgfqpoint{5.003956in}{1.994831in}}%
\pgfpathlineto{\pgfqpoint{5.006474in}{1.963824in}}%
\pgfpathlineto{\pgfqpoint{5.008993in}{1.968042in}}%
\pgfpathlineto{\pgfqpoint{5.011512in}{1.989454in}}%
\pgfpathlineto{\pgfqpoint{5.014031in}{1.991564in}}%
\pgfpathlineto{\pgfqpoint{5.016549in}{2.019794in}}%
\pgfpathlineto{\pgfqpoint{5.019068in}{2.065137in}}%
\pgfpathlineto{\pgfqpoint{5.021587in}{2.063524in}}%
\pgfpathlineto{\pgfqpoint{5.024106in}{2.095227in}}%
\pgfpathlineto{\pgfqpoint{5.026624in}{2.093613in}}%
\pgfpathlineto{\pgfqpoint{5.029143in}{2.132359in}}%
\pgfpathlineto{\pgfqpoint{5.031662in}{2.137129in}}%
\pgfpathlineto{\pgfqpoint{5.034181in}{2.148538in}}%
\pgfpathlineto{\pgfqpoint{5.036699in}{2.132217in}}%
\pgfpathlineto{\pgfqpoint{5.039218in}{2.151614in}}%
\pgfpathlineto{\pgfqpoint{5.041737in}{2.162264in}}%
\pgfpathlineto{\pgfqpoint{5.044256in}{2.159905in}}%
\pgfpathlineto{\pgfqpoint{5.046774in}{2.173189in}}%
\pgfpathlineto{\pgfqpoint{5.049293in}{2.183626in}}%
\pgfpathlineto{\pgfqpoint{5.051812in}{2.200476in}}%
\pgfpathlineto{\pgfqpoint{5.054331in}{2.174053in}}%
\pgfpathlineto{\pgfqpoint{5.056849in}{2.209523in}}%
\pgfpathlineto{\pgfqpoint{5.059368in}{2.210001in}}%
\pgfpathlineto{\pgfqpoint{5.061887in}{2.210618in}}%
\pgfpathlineto{\pgfqpoint{5.064406in}{2.208955in}}%
\pgfpathlineto{\pgfqpoint{5.066924in}{2.183519in}}%
\pgfpathlineto{\pgfqpoint{5.071962in}{2.114005in}}%
\pgfpathlineto{\pgfqpoint{5.074481in}{2.103076in}}%
\pgfpathlineto{\pgfqpoint{5.076999in}{2.128635in}}%
\pgfpathlineto{\pgfqpoint{5.079518in}{2.135710in}}%
\pgfpathlineto{\pgfqpoint{5.082037in}{2.141081in}}%
\pgfpathlineto{\pgfqpoint{5.084556in}{2.122916in}}%
\pgfpathlineto{\pgfqpoint{5.087074in}{2.138259in}}%
\pgfpathlineto{\pgfqpoint{5.089593in}{2.135283in}}%
\pgfpathlineto{\pgfqpoint{5.092112in}{2.151337in}}%
\pgfpathlineto{\pgfqpoint{5.097149in}{2.186119in}}%
\pgfpathlineto{\pgfqpoint{5.099668in}{2.163475in}}%
\pgfpathlineto{\pgfqpoint{5.102187in}{2.162735in}}%
\pgfpathlineto{\pgfqpoint{5.104706in}{2.176168in}}%
\pgfpathlineto{\pgfqpoint{5.107224in}{2.154473in}}%
\pgfpathlineto{\pgfqpoint{5.109743in}{2.140076in}}%
\pgfpathlineto{\pgfqpoint{5.112262in}{2.160257in}}%
\pgfpathlineto{\pgfqpoint{5.114781in}{2.144793in}}%
\pgfpathlineto{\pgfqpoint{5.117299in}{2.140237in}}%
\pgfpathlineto{\pgfqpoint{5.119818in}{2.146591in}}%
\pgfpathlineto{\pgfqpoint{5.122337in}{2.145434in}}%
\pgfpathlineto{\pgfqpoint{5.124856in}{2.161838in}}%
\pgfpathlineto{\pgfqpoint{5.127374in}{2.175253in}}%
\pgfpathlineto{\pgfqpoint{5.129893in}{2.148206in}}%
\pgfpathlineto{\pgfqpoint{5.132412in}{2.151863in}}%
\pgfpathlineto{\pgfqpoint{5.134931in}{2.188390in}}%
\pgfpathlineto{\pgfqpoint{5.137449in}{2.172087in}}%
\pgfpathlineto{\pgfqpoint{5.139968in}{2.182013in}}%
\pgfpathlineto{\pgfqpoint{5.142487in}{2.154879in}}%
\pgfpathlineto{\pgfqpoint{5.145006in}{2.179593in}}%
\pgfpathlineto{\pgfqpoint{5.147524in}{2.151444in}}%
\pgfpathlineto{\pgfqpoint{5.150043in}{2.172648in}}%
\pgfpathlineto{\pgfqpoint{5.152562in}{2.206650in}}%
\pgfpathlineto{\pgfqpoint{5.155081in}{2.205288in}}%
\pgfpathlineto{\pgfqpoint{5.157599in}{2.211922in}}%
\pgfpathlineto{\pgfqpoint{5.160118in}{2.212648in}}%
\pgfpathlineto{\pgfqpoint{5.162637in}{2.200908in}}%
\pgfpathlineto{\pgfqpoint{5.167674in}{2.250508in}}%
\pgfpathlineto{\pgfqpoint{5.170193in}{2.279062in}}%
\pgfpathlineto{\pgfqpoint{5.175231in}{2.234131in}}%
\pgfpathlineto{\pgfqpoint{5.177749in}{2.229585in}}%
\pgfpathlineto{\pgfqpoint{5.180268in}{2.212039in}}%
\pgfpathlineto{\pgfqpoint{5.182787in}{2.237334in}}%
\pgfpathlineto{\pgfqpoint{5.185306in}{2.249560in}}%
\pgfpathlineto{\pgfqpoint{5.187824in}{2.294121in}}%
\pgfpathlineto{\pgfqpoint{5.190343in}{2.271708in}}%
\pgfpathlineto{\pgfqpoint{5.192862in}{2.282533in}}%
\pgfpathlineto{\pgfqpoint{5.195381in}{2.308508in}}%
\pgfpathlineto{\pgfqpoint{5.200418in}{2.319208in}}%
\pgfpathlineto{\pgfqpoint{5.202937in}{2.320662in}}%
\pgfpathlineto{\pgfqpoint{5.205456in}{2.355658in}}%
\pgfpathlineto{\pgfqpoint{5.207974in}{2.396506in}}%
\pgfpathlineto{\pgfqpoint{5.210493in}{2.391635in}}%
\pgfpathlineto{\pgfqpoint{5.213012in}{2.377676in}}%
\pgfpathlineto{\pgfqpoint{5.215531in}{2.364241in}}%
\pgfpathlineto{\pgfqpoint{5.218049in}{2.359285in}}%
\pgfpathlineto{\pgfqpoint{5.220568in}{2.358378in}}%
\pgfpathlineto{\pgfqpoint{5.223087in}{2.363333in}}%
\pgfpathlineto{\pgfqpoint{5.225606in}{2.358142in}}%
\pgfpathlineto{\pgfqpoint{5.228124in}{2.359823in}}%
\pgfpathlineto{\pgfqpoint{5.230643in}{2.333941in}}%
\pgfpathlineto{\pgfqpoint{5.233162in}{2.296704in}}%
\pgfpathlineto{\pgfqpoint{5.235681in}{2.307146in}}%
\pgfpathlineto{\pgfqpoint{5.238199in}{2.328362in}}%
\pgfpathlineto{\pgfqpoint{5.240718in}{2.326244in}}%
\pgfpathlineto{\pgfqpoint{5.243237in}{2.309139in}}%
\pgfpathlineto{\pgfqpoint{5.245756in}{2.305900in}}%
\pgfpathlineto{\pgfqpoint{5.248274in}{2.298038in}}%
\pgfpathlineto{\pgfqpoint{5.250793in}{2.335988in}}%
\pgfpathlineto{\pgfqpoint{5.253312in}{2.316698in}}%
\pgfpathlineto{\pgfqpoint{5.255831in}{2.363104in}}%
\pgfpathlineto{\pgfqpoint{5.258349in}{2.382621in}}%
\pgfpathlineto{\pgfqpoint{5.260868in}{2.390904in}}%
\pgfpathlineto{\pgfqpoint{5.263387in}{2.375058in}}%
\pgfpathlineto{\pgfqpoint{5.265906in}{2.365549in}}%
\pgfpathlineto{\pgfqpoint{5.268424in}{2.354543in}}%
\pgfpathlineto{\pgfqpoint{5.270943in}{2.346110in}}%
\pgfpathlineto{\pgfqpoint{5.273462in}{2.370223in}}%
\pgfpathlineto{\pgfqpoint{5.275981in}{2.384462in}}%
\pgfpathlineto{\pgfqpoint{5.278499in}{2.390348in}}%
\pgfpathlineto{\pgfqpoint{5.281018in}{2.405613in}}%
\pgfpathlineto{\pgfqpoint{5.283537in}{2.406196in}}%
\pgfpathlineto{\pgfqpoint{5.286056in}{2.405318in}}%
\pgfpathlineto{\pgfqpoint{5.288574in}{2.409583in}}%
\pgfpathlineto{\pgfqpoint{5.291093in}{2.453597in}}%
\pgfpathlineto{\pgfqpoint{5.293612in}{2.476034in}}%
\pgfpathlineto{\pgfqpoint{5.296131in}{2.464528in}}%
\pgfpathlineto{\pgfqpoint{5.298649in}{2.475481in}}%
\pgfpathlineto{\pgfqpoint{5.301168in}{2.487662in}}%
\pgfpathlineto{\pgfqpoint{5.303687in}{2.467038in}}%
\pgfpathlineto{\pgfqpoint{5.306206in}{2.473885in}}%
\pgfpathlineto{\pgfqpoint{5.308724in}{2.456056in}}%
\pgfpathlineto{\pgfqpoint{5.311243in}{2.450677in}}%
\pgfpathlineto{\pgfqpoint{5.313762in}{2.443405in}}%
\pgfpathlineto{\pgfqpoint{5.316281in}{2.443838in}}%
\pgfpathlineto{\pgfqpoint{5.318799in}{2.441399in}}%
\pgfpathlineto{\pgfqpoint{5.321318in}{2.432793in}}%
\pgfpathlineto{\pgfqpoint{5.323837in}{2.414631in}}%
\pgfpathlineto{\pgfqpoint{5.326356in}{2.400982in}}%
\pgfpathlineto{\pgfqpoint{5.328874in}{2.402096in}}%
\pgfpathlineto{\pgfqpoint{5.331393in}{2.398391in}}%
\pgfpathlineto{\pgfqpoint{5.333912in}{2.403157in}}%
\pgfpathlineto{\pgfqpoint{5.336431in}{2.387196in}}%
\pgfpathlineto{\pgfqpoint{5.338949in}{2.415752in}}%
\pgfpathlineto{\pgfqpoint{5.341468in}{2.455510in}}%
\pgfpathlineto{\pgfqpoint{5.343987in}{2.395556in}}%
\pgfpathlineto{\pgfqpoint{5.346506in}{2.416141in}}%
\pgfpathlineto{\pgfqpoint{5.349024in}{2.430098in}}%
\pgfpathlineto{\pgfqpoint{5.351543in}{2.456615in}}%
\pgfpathlineto{\pgfqpoint{5.354062in}{2.496754in}}%
\pgfpathlineto{\pgfqpoint{5.356581in}{2.492820in}}%
\pgfpathlineto{\pgfqpoint{5.359099in}{2.458527in}}%
\pgfpathlineto{\pgfqpoint{5.361618in}{2.449760in}}%
\pgfpathlineto{\pgfqpoint{5.364137in}{2.425541in}}%
\pgfpathlineto{\pgfqpoint{5.366656in}{2.423987in}}%
\pgfpathlineto{\pgfqpoint{5.369174in}{2.393441in}}%
\pgfpathlineto{\pgfqpoint{5.371693in}{2.368898in}}%
\pgfpathlineto{\pgfqpoint{5.374212in}{2.389823in}}%
\pgfpathlineto{\pgfqpoint{5.376731in}{2.379971in}}%
\pgfpathlineto{\pgfqpoint{5.379249in}{2.361691in}}%
\pgfpathlineto{\pgfqpoint{5.381768in}{2.351119in}}%
\pgfpathlineto{\pgfqpoint{5.384287in}{2.341986in}}%
\pgfpathlineto{\pgfqpoint{5.386806in}{2.353855in}}%
\pgfpathlineto{\pgfqpoint{5.389324in}{2.318653in}}%
\pgfpathlineto{\pgfqpoint{5.391843in}{2.317187in}}%
\pgfpathlineto{\pgfqpoint{5.394362in}{2.299728in}}%
\pgfpathlineto{\pgfqpoint{5.396881in}{2.344599in}}%
\pgfpathlineto{\pgfqpoint{5.399399in}{2.372670in}}%
\pgfpathlineto{\pgfqpoint{5.401918in}{2.369283in}}%
\pgfpathlineto{\pgfqpoint{5.404437in}{2.384258in}}%
\pgfpathlineto{\pgfqpoint{5.406956in}{2.372479in}}%
\pgfpathlineto{\pgfqpoint{5.409474in}{2.384257in}}%
\pgfpathlineto{\pgfqpoint{5.411993in}{2.391313in}}%
\pgfpathlineto{\pgfqpoint{5.417031in}{2.423137in}}%
\pgfpathlineto{\pgfqpoint{5.419549in}{2.459322in}}%
\pgfpathlineto{\pgfqpoint{5.422068in}{2.464205in}}%
\pgfpathlineto{\pgfqpoint{5.424587in}{2.492922in}}%
\pgfpathlineto{\pgfqpoint{5.427106in}{2.461112in}}%
\pgfpathlineto{\pgfqpoint{5.429624in}{2.484215in}}%
\pgfpathlineto{\pgfqpoint{5.432143in}{2.486172in}}%
\pgfpathlineto{\pgfqpoint{5.434662in}{2.454642in}}%
\pgfpathlineto{\pgfqpoint{5.437181in}{2.443085in}}%
\pgfpathlineto{\pgfqpoint{5.439699in}{2.437559in}}%
\pgfpathlineto{\pgfqpoint{5.442218in}{2.418882in}}%
\pgfpathlineto{\pgfqpoint{5.444737in}{2.417525in}}%
\pgfpathlineto{\pgfqpoint{5.447256in}{2.429446in}}%
\pgfpathlineto{\pgfqpoint{5.449774in}{2.416203in}}%
\pgfpathlineto{\pgfqpoint{5.452293in}{2.406836in}}%
\pgfpathlineto{\pgfqpoint{5.454812in}{2.409071in}}%
\pgfpathlineto{\pgfqpoint{5.457331in}{2.395586in}}%
\pgfpathlineto{\pgfqpoint{5.459849in}{2.386061in}}%
\pgfpathlineto{\pgfqpoint{5.462368in}{2.384789in}}%
\pgfpathlineto{\pgfqpoint{5.464887in}{2.366005in}}%
\pgfpathlineto{\pgfqpoint{5.467406in}{2.402114in}}%
\pgfpathlineto{\pgfqpoint{5.469924in}{2.424177in}}%
\pgfpathlineto{\pgfqpoint{5.472443in}{2.412689in}}%
\pgfpathlineto{\pgfqpoint{5.474962in}{2.423249in}}%
\pgfpathlineto{\pgfqpoint{5.477481in}{2.407536in}}%
\pgfpathlineto{\pgfqpoint{5.479999in}{2.436614in}}%
\pgfpathlineto{\pgfqpoint{5.482518in}{2.418093in}}%
\pgfpathlineto{\pgfqpoint{5.485037in}{2.414817in}}%
\pgfpathlineto{\pgfqpoint{5.487556in}{2.440272in}}%
\pgfpathlineto{\pgfqpoint{5.490074in}{2.440240in}}%
\pgfpathlineto{\pgfqpoint{5.492593in}{2.446490in}}%
\pgfpathlineto{\pgfqpoint{5.492593in}{2.446490in}}%
\pgfusepath{stroke}%
\end{pgfscope}%
\begin{pgfscope}%
\pgfpathrectangle{\pgfqpoint{0.455093in}{0.401080in}}{\pgfqpoint{5.037500in}{4.004000in}}%
\pgfusepath{clip}%
\pgfsetrectcap%
\pgfsetroundjoin%
\pgfsetlinewidth{0.501875pt}%
\definecolor{currentstroke}{rgb}{0.690196,0.188235,0.376471}%
\pgfsetstrokecolor{currentstroke}%
\pgfsetstrokeopacity{0.800000}%
\pgfsetdash{}{0pt}%
\pgfpathmoveto{\pgfqpoint{0.455093in}{1.021044in}}%
\pgfpathlineto{\pgfqpoint{0.457612in}{1.016438in}}%
\pgfpathlineto{\pgfqpoint{0.460131in}{1.001725in}}%
\pgfpathlineto{\pgfqpoint{0.462649in}{0.998914in}}%
\pgfpathlineto{\pgfqpoint{0.465168in}{0.984165in}}%
\pgfpathlineto{\pgfqpoint{0.467687in}{0.987865in}}%
\pgfpathlineto{\pgfqpoint{0.470206in}{0.969547in}}%
\pgfpathlineto{\pgfqpoint{0.472724in}{0.963853in}}%
\pgfpathlineto{\pgfqpoint{0.475243in}{0.978017in}}%
\pgfpathlineto{\pgfqpoint{0.477762in}{0.979730in}}%
\pgfpathlineto{\pgfqpoint{0.480281in}{0.969525in}}%
\pgfpathlineto{\pgfqpoint{0.482799in}{0.982030in}}%
\pgfpathlineto{\pgfqpoint{0.485318in}{0.977594in}}%
\pgfpathlineto{\pgfqpoint{0.487837in}{0.973982in}}%
\pgfpathlineto{\pgfqpoint{0.490356in}{0.987556in}}%
\pgfpathlineto{\pgfqpoint{0.492874in}{1.002458in}}%
\pgfpathlineto{\pgfqpoint{0.495393in}{0.988572in}}%
\pgfpathlineto{\pgfqpoint{0.497912in}{0.983496in}}%
\pgfpathlineto{\pgfqpoint{0.500431in}{0.993115in}}%
\pgfpathlineto{\pgfqpoint{0.502949in}{0.971880in}}%
\pgfpathlineto{\pgfqpoint{0.505468in}{0.968000in}}%
\pgfpathlineto{\pgfqpoint{0.507987in}{0.995723in}}%
\pgfpathlineto{\pgfqpoint{0.510506in}{0.980777in}}%
\pgfpathlineto{\pgfqpoint{0.513024in}{0.992191in}}%
\pgfpathlineto{\pgfqpoint{0.515543in}{0.980407in}}%
\pgfpathlineto{\pgfqpoint{0.518062in}{0.975316in}}%
\pgfpathlineto{\pgfqpoint{0.520581in}{0.983859in}}%
\pgfpathlineto{\pgfqpoint{0.523099in}{0.984896in}}%
\pgfpathlineto{\pgfqpoint{0.525618in}{0.977053in}}%
\pgfpathlineto{\pgfqpoint{0.528137in}{0.986933in}}%
\pgfpathlineto{\pgfqpoint{0.530656in}{0.986260in}}%
\pgfpathlineto{\pgfqpoint{0.533174in}{0.976499in}}%
\pgfpathlineto{\pgfqpoint{0.535693in}{0.975392in}}%
\pgfpathlineto{\pgfqpoint{0.538212in}{0.973604in}}%
\pgfpathlineto{\pgfqpoint{0.540731in}{0.968394in}}%
\pgfpathlineto{\pgfqpoint{0.543249in}{0.986034in}}%
\pgfpathlineto{\pgfqpoint{0.545768in}{0.986868in}}%
\pgfpathlineto{\pgfqpoint{0.548287in}{0.996017in}}%
\pgfpathlineto{\pgfqpoint{0.550806in}{0.998887in}}%
\pgfpathlineto{\pgfqpoint{0.553324in}{0.962272in}}%
\pgfpathlineto{\pgfqpoint{0.555843in}{0.970845in}}%
\pgfpathlineto{\pgfqpoint{0.558362in}{0.982698in}}%
\pgfpathlineto{\pgfqpoint{0.560881in}{0.973098in}}%
\pgfpathlineto{\pgfqpoint{0.563399in}{0.979772in}}%
\pgfpathlineto{\pgfqpoint{0.565918in}{0.973065in}}%
\pgfpathlineto{\pgfqpoint{0.568437in}{0.980412in}}%
\pgfpathlineto{\pgfqpoint{0.570956in}{0.965864in}}%
\pgfpathlineto{\pgfqpoint{0.573474in}{0.972464in}}%
\pgfpathlineto{\pgfqpoint{0.575993in}{0.980017in}}%
\pgfpathlineto{\pgfqpoint{0.578512in}{0.980623in}}%
\pgfpathlineto{\pgfqpoint{0.581031in}{0.998784in}}%
\pgfpathlineto{\pgfqpoint{0.583549in}{0.978894in}}%
\pgfpathlineto{\pgfqpoint{0.586068in}{0.971171in}}%
\pgfpathlineto{\pgfqpoint{0.588587in}{0.969868in}}%
\pgfpathlineto{\pgfqpoint{0.591106in}{0.989674in}}%
\pgfpathlineto{\pgfqpoint{0.593624in}{0.975268in}}%
\pgfpathlineto{\pgfqpoint{0.596143in}{0.956836in}}%
\pgfpathlineto{\pgfqpoint{0.598662in}{0.942342in}}%
\pgfpathlineto{\pgfqpoint{0.601181in}{0.960024in}}%
\pgfpathlineto{\pgfqpoint{0.603699in}{0.970898in}}%
\pgfpathlineto{\pgfqpoint{0.606218in}{0.968444in}}%
\pgfpathlineto{\pgfqpoint{0.608737in}{0.952694in}}%
\pgfpathlineto{\pgfqpoint{0.611256in}{0.950825in}}%
\pgfpathlineto{\pgfqpoint{0.613774in}{0.958369in}}%
\pgfpathlineto{\pgfqpoint{0.616293in}{0.957450in}}%
\pgfpathlineto{\pgfqpoint{0.618812in}{0.968508in}}%
\pgfpathlineto{\pgfqpoint{0.621331in}{0.967949in}}%
\pgfpathlineto{\pgfqpoint{0.623849in}{0.978062in}}%
\pgfpathlineto{\pgfqpoint{0.626368in}{0.964016in}}%
\pgfpathlineto{\pgfqpoint{0.628887in}{0.994252in}}%
\pgfpathlineto{\pgfqpoint{0.631406in}{0.988465in}}%
\pgfpathlineto{\pgfqpoint{0.633924in}{0.975036in}}%
\pgfpathlineto{\pgfqpoint{0.636443in}{0.999022in}}%
\pgfpathlineto{\pgfqpoint{0.638962in}{1.004632in}}%
\pgfpathlineto{\pgfqpoint{0.641481in}{1.005013in}}%
\pgfpathlineto{\pgfqpoint{0.643999in}{0.994168in}}%
\pgfpathlineto{\pgfqpoint{0.646518in}{0.985391in}}%
\pgfpathlineto{\pgfqpoint{0.649037in}{0.998162in}}%
\pgfpathlineto{\pgfqpoint{0.651556in}{1.002708in}}%
\pgfpathlineto{\pgfqpoint{0.654074in}{1.015655in}}%
\pgfpathlineto{\pgfqpoint{0.656593in}{1.018190in}}%
\pgfpathlineto{\pgfqpoint{0.659112in}{1.005288in}}%
\pgfpathlineto{\pgfqpoint{0.661631in}{1.004517in}}%
\pgfpathlineto{\pgfqpoint{0.664149in}{1.034270in}}%
\pgfpathlineto{\pgfqpoint{0.666668in}{1.031334in}}%
\pgfpathlineto{\pgfqpoint{0.669187in}{1.048359in}}%
\pgfpathlineto{\pgfqpoint{0.671706in}{1.032939in}}%
\pgfpathlineto{\pgfqpoint{0.674224in}{1.023635in}}%
\pgfpathlineto{\pgfqpoint{0.676743in}{1.023752in}}%
\pgfpathlineto{\pgfqpoint{0.679262in}{1.023122in}}%
\pgfpathlineto{\pgfqpoint{0.681781in}{1.033622in}}%
\pgfpathlineto{\pgfqpoint{0.684299in}{1.027519in}}%
\pgfpathlineto{\pgfqpoint{0.686818in}{1.036065in}}%
\pgfpathlineto{\pgfqpoint{0.689337in}{1.059572in}}%
\pgfpathlineto{\pgfqpoint{0.691856in}{1.073873in}}%
\pgfpathlineto{\pgfqpoint{0.694374in}{1.076007in}}%
\pgfpathlineto{\pgfqpoint{0.696893in}{1.064871in}}%
\pgfpathlineto{\pgfqpoint{0.699412in}{1.080195in}}%
\pgfpathlineto{\pgfqpoint{0.701931in}{1.066642in}}%
\pgfpathlineto{\pgfqpoint{0.704449in}{1.075381in}}%
\pgfpathlineto{\pgfqpoint{0.706968in}{1.051387in}}%
\pgfpathlineto{\pgfqpoint{0.709487in}{1.071534in}}%
\pgfpathlineto{\pgfqpoint{0.712006in}{1.052462in}}%
\pgfpathlineto{\pgfqpoint{0.714524in}{1.076044in}}%
\pgfpathlineto{\pgfqpoint{0.717043in}{1.070497in}}%
\pgfpathlineto{\pgfqpoint{0.719562in}{1.046857in}}%
\pgfpathlineto{\pgfqpoint{0.722081in}{1.054409in}}%
\pgfpathlineto{\pgfqpoint{0.724599in}{1.074908in}}%
\pgfpathlineto{\pgfqpoint{0.727118in}{1.074407in}}%
\pgfpathlineto{\pgfqpoint{0.729637in}{1.091662in}}%
\pgfpathlineto{\pgfqpoint{0.732156in}{1.090829in}}%
\pgfpathlineto{\pgfqpoint{0.734674in}{1.078047in}}%
\pgfpathlineto{\pgfqpoint{0.737193in}{1.099509in}}%
\pgfpathlineto{\pgfqpoint{0.739712in}{1.104161in}}%
\pgfpathlineto{\pgfqpoint{0.742231in}{1.085486in}}%
\pgfpathlineto{\pgfqpoint{0.744749in}{1.071084in}}%
\pgfpathlineto{\pgfqpoint{0.747268in}{1.058080in}}%
\pgfpathlineto{\pgfqpoint{0.749787in}{1.068859in}}%
\pgfpathlineto{\pgfqpoint{0.752306in}{1.064376in}}%
\pgfpathlineto{\pgfqpoint{0.754824in}{1.076547in}}%
\pgfpathlineto{\pgfqpoint{0.757343in}{1.101051in}}%
\pgfpathlineto{\pgfqpoint{0.759862in}{1.097447in}}%
\pgfpathlineto{\pgfqpoint{0.762381in}{1.092756in}}%
\pgfpathlineto{\pgfqpoint{0.764899in}{1.093501in}}%
\pgfpathlineto{\pgfqpoint{0.767418in}{1.112869in}}%
\pgfpathlineto{\pgfqpoint{0.769937in}{1.100350in}}%
\pgfpathlineto{\pgfqpoint{0.772456in}{1.115125in}}%
\pgfpathlineto{\pgfqpoint{0.774974in}{1.107478in}}%
\pgfpathlineto{\pgfqpoint{0.777493in}{1.100915in}}%
\pgfpathlineto{\pgfqpoint{0.780012in}{1.085095in}}%
\pgfpathlineto{\pgfqpoint{0.782531in}{1.084844in}}%
\pgfpathlineto{\pgfqpoint{0.785049in}{1.085663in}}%
\pgfpathlineto{\pgfqpoint{0.787568in}{1.114931in}}%
\pgfpathlineto{\pgfqpoint{0.790087in}{1.114670in}}%
\pgfpathlineto{\pgfqpoint{0.792606in}{1.120364in}}%
\pgfpathlineto{\pgfqpoint{0.795124in}{1.113089in}}%
\pgfpathlineto{\pgfqpoint{0.797643in}{1.124111in}}%
\pgfpathlineto{\pgfqpoint{0.800162in}{1.149071in}}%
\pgfpathlineto{\pgfqpoint{0.802681in}{1.162864in}}%
\pgfpathlineto{\pgfqpoint{0.805199in}{1.164318in}}%
\pgfpathlineto{\pgfqpoint{0.807718in}{1.162397in}}%
\pgfpathlineto{\pgfqpoint{0.810237in}{1.174273in}}%
\pgfpathlineto{\pgfqpoint{0.812756in}{1.189991in}}%
\pgfpathlineto{\pgfqpoint{0.815274in}{1.177006in}}%
\pgfpathlineto{\pgfqpoint{0.817793in}{1.175595in}}%
\pgfpathlineto{\pgfqpoint{0.820312in}{1.201686in}}%
\pgfpathlineto{\pgfqpoint{0.822831in}{1.204060in}}%
\pgfpathlineto{\pgfqpoint{0.825349in}{1.201165in}}%
\pgfpathlineto{\pgfqpoint{0.827868in}{1.198773in}}%
\pgfpathlineto{\pgfqpoint{0.830387in}{1.168229in}}%
\pgfpathlineto{\pgfqpoint{0.832906in}{1.160106in}}%
\pgfpathlineto{\pgfqpoint{0.835424in}{1.154665in}}%
\pgfpathlineto{\pgfqpoint{0.837943in}{1.142795in}}%
\pgfpathlineto{\pgfqpoint{0.840462in}{1.161045in}}%
\pgfpathlineto{\pgfqpoint{0.842981in}{1.177423in}}%
\pgfpathlineto{\pgfqpoint{0.845499in}{1.185884in}}%
\pgfpathlineto{\pgfqpoint{0.848018in}{1.178343in}}%
\pgfpathlineto{\pgfqpoint{0.850537in}{1.177401in}}%
\pgfpathlineto{\pgfqpoint{0.853056in}{1.171492in}}%
\pgfpathlineto{\pgfqpoint{0.855574in}{1.172133in}}%
\pgfpathlineto{\pgfqpoint{0.858093in}{1.161123in}}%
\pgfpathlineto{\pgfqpoint{0.860612in}{1.165413in}}%
\pgfpathlineto{\pgfqpoint{0.863131in}{1.139092in}}%
\pgfpathlineto{\pgfqpoint{0.865649in}{1.161575in}}%
\pgfpathlineto{\pgfqpoint{0.868168in}{1.158189in}}%
\pgfpathlineto{\pgfqpoint{0.870687in}{1.169650in}}%
\pgfpathlineto{\pgfqpoint{0.873206in}{1.166235in}}%
\pgfpathlineto{\pgfqpoint{0.875724in}{1.188669in}}%
\pgfpathlineto{\pgfqpoint{0.878243in}{1.200114in}}%
\pgfpathlineto{\pgfqpoint{0.880762in}{1.202948in}}%
\pgfpathlineto{\pgfqpoint{0.883281in}{1.212955in}}%
\pgfpathlineto{\pgfqpoint{0.885799in}{1.237658in}}%
\pgfpathlineto{\pgfqpoint{0.888318in}{1.253813in}}%
\pgfpathlineto{\pgfqpoint{0.890837in}{1.262338in}}%
\pgfpathlineto{\pgfqpoint{0.893356in}{1.256129in}}%
\pgfpathlineto{\pgfqpoint{0.898393in}{1.235362in}}%
\pgfpathlineto{\pgfqpoint{0.900912in}{1.248290in}}%
\pgfpathlineto{\pgfqpoint{0.903431in}{1.226200in}}%
\pgfpathlineto{\pgfqpoint{0.905949in}{1.244016in}}%
\pgfpathlineto{\pgfqpoint{0.908468in}{1.255877in}}%
\pgfpathlineto{\pgfqpoint{0.910987in}{1.255784in}}%
\pgfpathlineto{\pgfqpoint{0.913506in}{1.259490in}}%
\pgfpathlineto{\pgfqpoint{0.916024in}{1.268756in}}%
\pgfpathlineto{\pgfqpoint{0.918543in}{1.265624in}}%
\pgfpathlineto{\pgfqpoint{0.921062in}{1.257902in}}%
\pgfpathlineto{\pgfqpoint{0.923581in}{1.287797in}}%
\pgfpathlineto{\pgfqpoint{0.928618in}{1.280597in}}%
\pgfpathlineto{\pgfqpoint{0.931137in}{1.278969in}}%
\pgfpathlineto{\pgfqpoint{0.933656in}{1.279059in}}%
\pgfpathlineto{\pgfqpoint{0.936174in}{1.297926in}}%
\pgfpathlineto{\pgfqpoint{0.938693in}{1.276915in}}%
\pgfpathlineto{\pgfqpoint{0.941212in}{1.274146in}}%
\pgfpathlineto{\pgfqpoint{0.943731in}{1.260114in}}%
\pgfpathlineto{\pgfqpoint{0.946249in}{1.230420in}}%
\pgfpathlineto{\pgfqpoint{0.948768in}{1.215436in}}%
\pgfpathlineto{\pgfqpoint{0.953806in}{1.229013in}}%
\pgfpathlineto{\pgfqpoint{0.956324in}{1.219168in}}%
\pgfpathlineto{\pgfqpoint{0.958843in}{1.224109in}}%
\pgfpathlineto{\pgfqpoint{0.961362in}{1.205399in}}%
\pgfpathlineto{\pgfqpoint{0.963881in}{1.221656in}}%
\pgfpathlineto{\pgfqpoint{0.966399in}{1.201162in}}%
\pgfpathlineto{\pgfqpoint{0.968918in}{1.217699in}}%
\pgfpathlineto{\pgfqpoint{0.971437in}{1.208166in}}%
\pgfpathlineto{\pgfqpoint{0.973956in}{1.212540in}}%
\pgfpathlineto{\pgfqpoint{0.976474in}{1.226509in}}%
\pgfpathlineto{\pgfqpoint{0.978993in}{1.236571in}}%
\pgfpathlineto{\pgfqpoint{0.981512in}{1.219986in}}%
\pgfpathlineto{\pgfqpoint{0.984031in}{1.248084in}}%
\pgfpathlineto{\pgfqpoint{0.986549in}{1.240879in}}%
\pgfpathlineto{\pgfqpoint{0.989068in}{1.245496in}}%
\pgfpathlineto{\pgfqpoint{0.991587in}{1.255795in}}%
\pgfpathlineto{\pgfqpoint{0.994106in}{1.282846in}}%
\pgfpathlineto{\pgfqpoint{0.996624in}{1.298432in}}%
\pgfpathlineto{\pgfqpoint{0.999143in}{1.300924in}}%
\pgfpathlineto{\pgfqpoint{1.001662in}{1.297822in}}%
\pgfpathlineto{\pgfqpoint{1.004181in}{1.297430in}}%
\pgfpathlineto{\pgfqpoint{1.006699in}{1.297703in}}%
\pgfpathlineto{\pgfqpoint{1.009218in}{1.294532in}}%
\pgfpathlineto{\pgfqpoint{1.011737in}{1.329578in}}%
\pgfpathlineto{\pgfqpoint{1.014256in}{1.334893in}}%
\pgfpathlineto{\pgfqpoint{1.016774in}{1.312554in}}%
\pgfpathlineto{\pgfqpoint{1.019293in}{1.327816in}}%
\pgfpathlineto{\pgfqpoint{1.021812in}{1.314974in}}%
\pgfpathlineto{\pgfqpoint{1.024331in}{1.300303in}}%
\pgfpathlineto{\pgfqpoint{1.026849in}{1.301142in}}%
\pgfpathlineto{\pgfqpoint{1.029368in}{1.327883in}}%
\pgfpathlineto{\pgfqpoint{1.031887in}{1.341801in}}%
\pgfpathlineto{\pgfqpoint{1.034406in}{1.334470in}}%
\pgfpathlineto{\pgfqpoint{1.036924in}{1.353053in}}%
\pgfpathlineto{\pgfqpoint{1.039443in}{1.378674in}}%
\pgfpathlineto{\pgfqpoint{1.041962in}{1.396541in}}%
\pgfpathlineto{\pgfqpoint{1.044481in}{1.375699in}}%
\pgfpathlineto{\pgfqpoint{1.046999in}{1.401851in}}%
\pgfpathlineto{\pgfqpoint{1.049518in}{1.417970in}}%
\pgfpathlineto{\pgfqpoint{1.052037in}{1.400448in}}%
\pgfpathlineto{\pgfqpoint{1.054556in}{1.377032in}}%
\pgfpathlineto{\pgfqpoint{1.057074in}{1.361954in}}%
\pgfpathlineto{\pgfqpoint{1.059593in}{1.387952in}}%
\pgfpathlineto{\pgfqpoint{1.062112in}{1.410629in}}%
\pgfpathlineto{\pgfqpoint{1.064631in}{1.408965in}}%
\pgfpathlineto{\pgfqpoint{1.067149in}{1.408279in}}%
\pgfpathlineto{\pgfqpoint{1.069668in}{1.414858in}}%
\pgfpathlineto{\pgfqpoint{1.072187in}{1.413398in}}%
\pgfpathlineto{\pgfqpoint{1.074706in}{1.410058in}}%
\pgfpathlineto{\pgfqpoint{1.077224in}{1.375254in}}%
\pgfpathlineto{\pgfqpoint{1.079743in}{1.362415in}}%
\pgfpathlineto{\pgfqpoint{1.082262in}{1.368988in}}%
\pgfpathlineto{\pgfqpoint{1.084781in}{1.379778in}}%
\pgfpathlineto{\pgfqpoint{1.087299in}{1.365878in}}%
\pgfpathlineto{\pgfqpoint{1.089818in}{1.389150in}}%
\pgfpathlineto{\pgfqpoint{1.092337in}{1.397850in}}%
\pgfpathlineto{\pgfqpoint{1.094856in}{1.387065in}}%
\pgfpathlineto{\pgfqpoint{1.097374in}{1.378338in}}%
\pgfpathlineto{\pgfqpoint{1.099893in}{1.358376in}}%
\pgfpathlineto{\pgfqpoint{1.102412in}{1.366105in}}%
\pgfpathlineto{\pgfqpoint{1.104931in}{1.350902in}}%
\pgfpathlineto{\pgfqpoint{1.107449in}{1.360087in}}%
\pgfpathlineto{\pgfqpoint{1.109968in}{1.384319in}}%
\pgfpathlineto{\pgfqpoint{1.112487in}{1.369238in}}%
\pgfpathlineto{\pgfqpoint{1.115006in}{1.332861in}}%
\pgfpathlineto{\pgfqpoint{1.117524in}{1.342805in}}%
\pgfpathlineto{\pgfqpoint{1.120043in}{1.343702in}}%
\pgfpathlineto{\pgfqpoint{1.122562in}{1.337889in}}%
\pgfpathlineto{\pgfqpoint{1.127599in}{1.307513in}}%
\pgfpathlineto{\pgfqpoint{1.130118in}{1.334126in}}%
\pgfpathlineto{\pgfqpoint{1.132637in}{1.355353in}}%
\pgfpathlineto{\pgfqpoint{1.135156in}{1.336783in}}%
\pgfpathlineto{\pgfqpoint{1.137674in}{1.320566in}}%
\pgfpathlineto{\pgfqpoint{1.140193in}{1.346025in}}%
\pgfpathlineto{\pgfqpoint{1.142712in}{1.323026in}}%
\pgfpathlineto{\pgfqpoint{1.145231in}{1.346520in}}%
\pgfpathlineto{\pgfqpoint{1.147749in}{1.366703in}}%
\pgfpathlineto{\pgfqpoint{1.150268in}{1.360932in}}%
\pgfpathlineto{\pgfqpoint{1.152787in}{1.374072in}}%
\pgfpathlineto{\pgfqpoint{1.155306in}{1.358327in}}%
\pgfpathlineto{\pgfqpoint{1.157824in}{1.375914in}}%
\pgfpathlineto{\pgfqpoint{1.160343in}{1.384214in}}%
\pgfpathlineto{\pgfqpoint{1.162862in}{1.376155in}}%
\pgfpathlineto{\pgfqpoint{1.165381in}{1.388309in}}%
\pgfpathlineto{\pgfqpoint{1.167899in}{1.382009in}}%
\pgfpathlineto{\pgfqpoint{1.170418in}{1.396421in}}%
\pgfpathlineto{\pgfqpoint{1.172937in}{1.435097in}}%
\pgfpathlineto{\pgfqpoint{1.175456in}{1.413040in}}%
\pgfpathlineto{\pgfqpoint{1.177974in}{1.426414in}}%
\pgfpathlineto{\pgfqpoint{1.180493in}{1.432147in}}%
\pgfpathlineto{\pgfqpoint{1.183012in}{1.433646in}}%
\pgfpathlineto{\pgfqpoint{1.185531in}{1.442764in}}%
\pgfpathlineto{\pgfqpoint{1.188049in}{1.443689in}}%
\pgfpathlineto{\pgfqpoint{1.190568in}{1.437335in}}%
\pgfpathlineto{\pgfqpoint{1.193087in}{1.432501in}}%
\pgfpathlineto{\pgfqpoint{1.195606in}{1.445530in}}%
\pgfpathlineto{\pgfqpoint{1.198124in}{1.438895in}}%
\pgfpathlineto{\pgfqpoint{1.200643in}{1.461948in}}%
\pgfpathlineto{\pgfqpoint{1.203162in}{1.476645in}}%
\pgfpathlineto{\pgfqpoint{1.205681in}{1.490435in}}%
\pgfpathlineto{\pgfqpoint{1.208199in}{1.493970in}}%
\pgfpathlineto{\pgfqpoint{1.210718in}{1.498661in}}%
\pgfpathlineto{\pgfqpoint{1.213237in}{1.513299in}}%
\pgfpathlineto{\pgfqpoint{1.215756in}{1.498239in}}%
\pgfpathlineto{\pgfqpoint{1.218274in}{1.490647in}}%
\pgfpathlineto{\pgfqpoint{1.220793in}{1.481482in}}%
\pgfpathlineto{\pgfqpoint{1.223312in}{1.495622in}}%
\pgfpathlineto{\pgfqpoint{1.225831in}{1.471796in}}%
\pgfpathlineto{\pgfqpoint{1.228349in}{1.470255in}}%
\pgfpathlineto{\pgfqpoint{1.230868in}{1.492310in}}%
\pgfpathlineto{\pgfqpoint{1.233387in}{1.476274in}}%
\pgfpathlineto{\pgfqpoint{1.235906in}{1.479959in}}%
\pgfpathlineto{\pgfqpoint{1.238424in}{1.486026in}}%
\pgfpathlineto{\pgfqpoint{1.243462in}{1.470184in}}%
\pgfpathlineto{\pgfqpoint{1.245981in}{1.479119in}}%
\pgfpathlineto{\pgfqpoint{1.248499in}{1.470492in}}%
\pgfpathlineto{\pgfqpoint{1.253537in}{1.467347in}}%
\pgfpathlineto{\pgfqpoint{1.256056in}{1.456975in}}%
\pgfpathlineto{\pgfqpoint{1.258574in}{1.453156in}}%
\pgfpathlineto{\pgfqpoint{1.261093in}{1.423780in}}%
\pgfpathlineto{\pgfqpoint{1.263612in}{1.396242in}}%
\pgfpathlineto{\pgfqpoint{1.266131in}{1.387736in}}%
\pgfpathlineto{\pgfqpoint{1.268649in}{1.352879in}}%
\pgfpathlineto{\pgfqpoint{1.273687in}{1.389518in}}%
\pgfpathlineto{\pgfqpoint{1.276206in}{1.406792in}}%
\pgfpathlineto{\pgfqpoint{1.278724in}{1.405977in}}%
\pgfpathlineto{\pgfqpoint{1.281243in}{1.425356in}}%
\pgfpathlineto{\pgfqpoint{1.283762in}{1.423727in}}%
\pgfpathlineto{\pgfqpoint{1.286281in}{1.429083in}}%
\pgfpathlineto{\pgfqpoint{1.288799in}{1.434813in}}%
\pgfpathlineto{\pgfqpoint{1.291318in}{1.442917in}}%
\pgfpathlineto{\pgfqpoint{1.293837in}{1.448558in}}%
\pgfpathlineto{\pgfqpoint{1.296356in}{1.445323in}}%
\pgfpathlineto{\pgfqpoint{1.298874in}{1.444852in}}%
\pgfpathlineto{\pgfqpoint{1.301393in}{1.443058in}}%
\pgfpathlineto{\pgfqpoint{1.306431in}{1.428222in}}%
\pgfpathlineto{\pgfqpoint{1.308949in}{1.437413in}}%
\pgfpathlineto{\pgfqpoint{1.311468in}{1.451601in}}%
\pgfpathlineto{\pgfqpoint{1.313987in}{1.464365in}}%
\pgfpathlineto{\pgfqpoint{1.316506in}{1.489758in}}%
\pgfpathlineto{\pgfqpoint{1.319024in}{1.470244in}}%
\pgfpathlineto{\pgfqpoint{1.321543in}{1.462614in}}%
\pgfpathlineto{\pgfqpoint{1.324062in}{1.428686in}}%
\pgfpathlineto{\pgfqpoint{1.326581in}{1.428051in}}%
\pgfpathlineto{\pgfqpoint{1.329099in}{1.450039in}}%
\pgfpathlineto{\pgfqpoint{1.331618in}{1.436511in}}%
\pgfpathlineto{\pgfqpoint{1.334137in}{1.439002in}}%
\pgfpathlineto{\pgfqpoint{1.336656in}{1.455781in}}%
\pgfpathlineto{\pgfqpoint{1.339174in}{1.468866in}}%
\pgfpathlineto{\pgfqpoint{1.341693in}{1.489907in}}%
\pgfpathlineto{\pgfqpoint{1.344212in}{1.500906in}}%
\pgfpathlineto{\pgfqpoint{1.346731in}{1.491052in}}%
\pgfpathlineto{\pgfqpoint{1.349249in}{1.478452in}}%
\pgfpathlineto{\pgfqpoint{1.351768in}{1.478427in}}%
\pgfpathlineto{\pgfqpoint{1.354287in}{1.462617in}}%
\pgfpathlineto{\pgfqpoint{1.356806in}{1.455361in}}%
\pgfpathlineto{\pgfqpoint{1.359324in}{1.473961in}}%
\pgfpathlineto{\pgfqpoint{1.361843in}{1.467210in}}%
\pgfpathlineto{\pgfqpoint{1.364362in}{1.463625in}}%
\pgfpathlineto{\pgfqpoint{1.366881in}{1.448621in}}%
\pgfpathlineto{\pgfqpoint{1.369399in}{1.436607in}}%
\pgfpathlineto{\pgfqpoint{1.371918in}{1.450523in}}%
\pgfpathlineto{\pgfqpoint{1.374437in}{1.448979in}}%
\pgfpathlineto{\pgfqpoint{1.376956in}{1.442800in}}%
\pgfpathlineto{\pgfqpoint{1.379474in}{1.447000in}}%
\pgfpathlineto{\pgfqpoint{1.381993in}{1.453404in}}%
\pgfpathlineto{\pgfqpoint{1.384512in}{1.444159in}}%
\pgfpathlineto{\pgfqpoint{1.387031in}{1.405044in}}%
\pgfpathlineto{\pgfqpoint{1.389549in}{1.406686in}}%
\pgfpathlineto{\pgfqpoint{1.392068in}{1.390487in}}%
\pgfpathlineto{\pgfqpoint{1.394587in}{1.389932in}}%
\pgfpathlineto{\pgfqpoint{1.397106in}{1.429018in}}%
\pgfpathlineto{\pgfqpoint{1.399624in}{1.400297in}}%
\pgfpathlineto{\pgfqpoint{1.402143in}{1.400097in}}%
\pgfpathlineto{\pgfqpoint{1.404662in}{1.386281in}}%
\pgfpathlineto{\pgfqpoint{1.407181in}{1.406270in}}%
\pgfpathlineto{\pgfqpoint{1.409699in}{1.400700in}}%
\pgfpathlineto{\pgfqpoint{1.412218in}{1.408380in}}%
\pgfpathlineto{\pgfqpoint{1.414737in}{1.405546in}}%
\pgfpathlineto{\pgfqpoint{1.417256in}{1.398342in}}%
\pgfpathlineto{\pgfqpoint{1.419774in}{1.416213in}}%
\pgfpathlineto{\pgfqpoint{1.422293in}{1.420486in}}%
\pgfpathlineto{\pgfqpoint{1.424812in}{1.438062in}}%
\pgfpathlineto{\pgfqpoint{1.427331in}{1.437335in}}%
\pgfpathlineto{\pgfqpoint{1.429849in}{1.451154in}}%
\pgfpathlineto{\pgfqpoint{1.432368in}{1.444063in}}%
\pgfpathlineto{\pgfqpoint{1.434887in}{1.465271in}}%
\pgfpathlineto{\pgfqpoint{1.437406in}{1.491537in}}%
\pgfpathlineto{\pgfqpoint{1.439924in}{1.470518in}}%
\pgfpathlineto{\pgfqpoint{1.442443in}{1.492249in}}%
\pgfpathlineto{\pgfqpoint{1.444962in}{1.492200in}}%
\pgfpathlineto{\pgfqpoint{1.447481in}{1.499899in}}%
\pgfpathlineto{\pgfqpoint{1.449999in}{1.487348in}}%
\pgfpathlineto{\pgfqpoint{1.452518in}{1.494581in}}%
\pgfpathlineto{\pgfqpoint{1.455037in}{1.473717in}}%
\pgfpathlineto{\pgfqpoint{1.457556in}{1.502080in}}%
\pgfpathlineto{\pgfqpoint{1.462593in}{1.486657in}}%
\pgfpathlineto{\pgfqpoint{1.465112in}{1.462726in}}%
\pgfpathlineto{\pgfqpoint{1.467631in}{1.452681in}}%
\pgfpathlineto{\pgfqpoint{1.470149in}{1.447842in}}%
\pgfpathlineto{\pgfqpoint{1.472668in}{1.460697in}}%
\pgfpathlineto{\pgfqpoint{1.475187in}{1.451107in}}%
\pgfpathlineto{\pgfqpoint{1.477706in}{1.437080in}}%
\pgfpathlineto{\pgfqpoint{1.480224in}{1.446998in}}%
\pgfpathlineto{\pgfqpoint{1.482743in}{1.443819in}}%
\pgfpathlineto{\pgfqpoint{1.485262in}{1.458628in}}%
\pgfpathlineto{\pgfqpoint{1.487781in}{1.474008in}}%
\pgfpathlineto{\pgfqpoint{1.490299in}{1.479463in}}%
\pgfpathlineto{\pgfqpoint{1.492818in}{1.475746in}}%
\pgfpathlineto{\pgfqpoint{1.495337in}{1.481039in}}%
\pgfpathlineto{\pgfqpoint{1.497856in}{1.501823in}}%
\pgfpathlineto{\pgfqpoint{1.500374in}{1.488734in}}%
\pgfpathlineto{\pgfqpoint{1.502893in}{1.473914in}}%
\pgfpathlineto{\pgfqpoint{1.505412in}{1.478567in}}%
\pgfpathlineto{\pgfqpoint{1.507931in}{1.465194in}}%
\pgfpathlineto{\pgfqpoint{1.510449in}{1.463687in}}%
\pgfpathlineto{\pgfqpoint{1.512968in}{1.458095in}}%
\pgfpathlineto{\pgfqpoint{1.515487in}{1.464142in}}%
\pgfpathlineto{\pgfqpoint{1.518006in}{1.468116in}}%
\pgfpathlineto{\pgfqpoint{1.520524in}{1.484667in}}%
\pgfpathlineto{\pgfqpoint{1.523043in}{1.476444in}}%
\pgfpathlineto{\pgfqpoint{1.525562in}{1.480036in}}%
\pgfpathlineto{\pgfqpoint{1.528081in}{1.471887in}}%
\pgfpathlineto{\pgfqpoint{1.530599in}{1.479114in}}%
\pgfpathlineto{\pgfqpoint{1.533118in}{1.471789in}}%
\pgfpathlineto{\pgfqpoint{1.535637in}{1.481777in}}%
\pgfpathlineto{\pgfqpoint{1.538156in}{1.485853in}}%
\pgfpathlineto{\pgfqpoint{1.540674in}{1.478469in}}%
\pgfpathlineto{\pgfqpoint{1.543193in}{1.472424in}}%
\pgfpathlineto{\pgfqpoint{1.545712in}{1.459646in}}%
\pgfpathlineto{\pgfqpoint{1.550749in}{1.432343in}}%
\pgfpathlineto{\pgfqpoint{1.553268in}{1.449215in}}%
\pgfpathlineto{\pgfqpoint{1.555787in}{1.456723in}}%
\pgfpathlineto{\pgfqpoint{1.558306in}{1.455329in}}%
\pgfpathlineto{\pgfqpoint{1.560824in}{1.459659in}}%
\pgfpathlineto{\pgfqpoint{1.563343in}{1.447134in}}%
\pgfpathlineto{\pgfqpoint{1.565862in}{1.443665in}}%
\pgfpathlineto{\pgfqpoint{1.568381in}{1.456115in}}%
\pgfpathlineto{\pgfqpoint{1.570899in}{1.434507in}}%
\pgfpathlineto{\pgfqpoint{1.573418in}{1.455678in}}%
\pgfpathlineto{\pgfqpoint{1.575937in}{1.448078in}}%
\pgfpathlineto{\pgfqpoint{1.578456in}{1.447788in}}%
\pgfpathlineto{\pgfqpoint{1.580974in}{1.449445in}}%
\pgfpathlineto{\pgfqpoint{1.583493in}{1.443956in}}%
\pgfpathlineto{\pgfqpoint{1.586012in}{1.433055in}}%
\pgfpathlineto{\pgfqpoint{1.588531in}{1.435195in}}%
\pgfpathlineto{\pgfqpoint{1.591049in}{1.443878in}}%
\pgfpathlineto{\pgfqpoint{1.593568in}{1.477625in}}%
\pgfpathlineto{\pgfqpoint{1.596087in}{1.484531in}}%
\pgfpathlineto{\pgfqpoint{1.598606in}{1.519428in}}%
\pgfpathlineto{\pgfqpoint{1.601124in}{1.540136in}}%
\pgfpathlineto{\pgfqpoint{1.603643in}{1.517261in}}%
\pgfpathlineto{\pgfqpoint{1.606162in}{1.506861in}}%
\pgfpathlineto{\pgfqpoint{1.608681in}{1.514972in}}%
\pgfpathlineto{\pgfqpoint{1.611199in}{1.514408in}}%
\pgfpathlineto{\pgfqpoint{1.613718in}{1.518884in}}%
\pgfpathlineto{\pgfqpoint{1.616237in}{1.514689in}}%
\pgfpathlineto{\pgfqpoint{1.618756in}{1.504553in}}%
\pgfpathlineto{\pgfqpoint{1.621274in}{1.501704in}}%
\pgfpathlineto{\pgfqpoint{1.623793in}{1.509856in}}%
\pgfpathlineto{\pgfqpoint{1.626312in}{1.521202in}}%
\pgfpathlineto{\pgfqpoint{1.628831in}{1.502818in}}%
\pgfpathlineto{\pgfqpoint{1.631349in}{1.494360in}}%
\pgfpathlineto{\pgfqpoint{1.633868in}{1.498653in}}%
\pgfpathlineto{\pgfqpoint{1.636387in}{1.495682in}}%
\pgfpathlineto{\pgfqpoint{1.638906in}{1.529841in}}%
\pgfpathlineto{\pgfqpoint{1.641424in}{1.529970in}}%
\pgfpathlineto{\pgfqpoint{1.643943in}{1.551139in}}%
\pgfpathlineto{\pgfqpoint{1.646462in}{1.555176in}}%
\pgfpathlineto{\pgfqpoint{1.648981in}{1.567503in}}%
\pgfpathlineto{\pgfqpoint{1.651499in}{1.578056in}}%
\pgfpathlineto{\pgfqpoint{1.654018in}{1.583609in}}%
\pgfpathlineto{\pgfqpoint{1.656537in}{1.548334in}}%
\pgfpathlineto{\pgfqpoint{1.659056in}{1.575681in}}%
\pgfpathlineto{\pgfqpoint{1.661574in}{1.586950in}}%
\pgfpathlineto{\pgfqpoint{1.664093in}{1.595862in}}%
\pgfpathlineto{\pgfqpoint{1.666612in}{1.595385in}}%
\pgfpathlineto{\pgfqpoint{1.669131in}{1.588870in}}%
\pgfpathlineto{\pgfqpoint{1.671649in}{1.598362in}}%
\pgfpathlineto{\pgfqpoint{1.674168in}{1.597328in}}%
\pgfpathlineto{\pgfqpoint{1.676687in}{1.579586in}}%
\pgfpathlineto{\pgfqpoint{1.679206in}{1.586739in}}%
\pgfpathlineto{\pgfqpoint{1.681724in}{1.580744in}}%
\pgfpathlineto{\pgfqpoint{1.684243in}{1.589964in}}%
\pgfpathlineto{\pgfqpoint{1.691799in}{1.628751in}}%
\pgfpathlineto{\pgfqpoint{1.694318in}{1.647154in}}%
\pgfpathlineto{\pgfqpoint{1.696837in}{1.634615in}}%
\pgfpathlineto{\pgfqpoint{1.699356in}{1.620916in}}%
\pgfpathlineto{\pgfqpoint{1.701874in}{1.620521in}}%
\pgfpathlineto{\pgfqpoint{1.704393in}{1.590791in}}%
\pgfpathlineto{\pgfqpoint{1.706912in}{1.609825in}}%
\pgfpathlineto{\pgfqpoint{1.709431in}{1.629710in}}%
\pgfpathlineto{\pgfqpoint{1.711949in}{1.616569in}}%
\pgfpathlineto{\pgfqpoint{1.714468in}{1.611095in}}%
\pgfpathlineto{\pgfqpoint{1.716987in}{1.625011in}}%
\pgfpathlineto{\pgfqpoint{1.719506in}{1.625195in}}%
\pgfpathlineto{\pgfqpoint{1.722024in}{1.607111in}}%
\pgfpathlineto{\pgfqpoint{1.724543in}{1.576429in}}%
\pgfpathlineto{\pgfqpoint{1.727062in}{1.567154in}}%
\pgfpathlineto{\pgfqpoint{1.729581in}{1.542078in}}%
\pgfpathlineto{\pgfqpoint{1.732099in}{1.569080in}}%
\pgfpathlineto{\pgfqpoint{1.734618in}{1.533839in}}%
\pgfpathlineto{\pgfqpoint{1.737137in}{1.566573in}}%
\pgfpathlineto{\pgfqpoint{1.739656in}{1.572898in}}%
\pgfpathlineto{\pgfqpoint{1.742174in}{1.577741in}}%
\pgfpathlineto{\pgfqpoint{1.744693in}{1.597719in}}%
\pgfpathlineto{\pgfqpoint{1.747212in}{1.611817in}}%
\pgfpathlineto{\pgfqpoint{1.749731in}{1.617589in}}%
\pgfpathlineto{\pgfqpoint{1.752249in}{1.600803in}}%
\pgfpathlineto{\pgfqpoint{1.754768in}{1.611778in}}%
\pgfpathlineto{\pgfqpoint{1.757287in}{1.618867in}}%
\pgfpathlineto{\pgfqpoint{1.759806in}{1.611599in}}%
\pgfpathlineto{\pgfqpoint{1.762324in}{1.590696in}}%
\pgfpathlineto{\pgfqpoint{1.764843in}{1.565368in}}%
\pgfpathlineto{\pgfqpoint{1.767362in}{1.570305in}}%
\pgfpathlineto{\pgfqpoint{1.769881in}{1.589029in}}%
\pgfpathlineto{\pgfqpoint{1.772399in}{1.585020in}}%
\pgfpathlineto{\pgfqpoint{1.774918in}{1.560873in}}%
\pgfpathlineto{\pgfqpoint{1.777437in}{1.574852in}}%
\pgfpathlineto{\pgfqpoint{1.779956in}{1.544732in}}%
\pgfpathlineto{\pgfqpoint{1.782474in}{1.558010in}}%
\pgfpathlineto{\pgfqpoint{1.784993in}{1.558034in}}%
\pgfpathlineto{\pgfqpoint{1.787512in}{1.581462in}}%
\pgfpathlineto{\pgfqpoint{1.792549in}{1.618664in}}%
\pgfpathlineto{\pgfqpoint{1.795068in}{1.606419in}}%
\pgfpathlineto{\pgfqpoint{1.797587in}{1.596136in}}%
\pgfpathlineto{\pgfqpoint{1.800106in}{1.575070in}}%
\pgfpathlineto{\pgfqpoint{1.802624in}{1.544746in}}%
\pgfpathlineto{\pgfqpoint{1.805143in}{1.531998in}}%
\pgfpathlineto{\pgfqpoint{1.807662in}{1.549874in}}%
\pgfpathlineto{\pgfqpoint{1.810181in}{1.541704in}}%
\pgfpathlineto{\pgfqpoint{1.812699in}{1.534199in}}%
\pgfpathlineto{\pgfqpoint{1.815218in}{1.553047in}}%
\pgfpathlineto{\pgfqpoint{1.817737in}{1.539904in}}%
\pgfpathlineto{\pgfqpoint{1.820256in}{1.526093in}}%
\pgfpathlineto{\pgfqpoint{1.822774in}{1.548161in}}%
\pgfpathlineto{\pgfqpoint{1.825293in}{1.562546in}}%
\pgfpathlineto{\pgfqpoint{1.827812in}{1.546810in}}%
\pgfpathlineto{\pgfqpoint{1.830331in}{1.543159in}}%
\pgfpathlineto{\pgfqpoint{1.832849in}{1.559591in}}%
\pgfpathlineto{\pgfqpoint{1.835368in}{1.578871in}}%
\pgfpathlineto{\pgfqpoint{1.837887in}{1.575523in}}%
\pgfpathlineto{\pgfqpoint{1.840406in}{1.581624in}}%
\pgfpathlineto{\pgfqpoint{1.842924in}{1.590132in}}%
\pgfpathlineto{\pgfqpoint{1.845443in}{1.623550in}}%
\pgfpathlineto{\pgfqpoint{1.847962in}{1.606650in}}%
\pgfpathlineto{\pgfqpoint{1.850481in}{1.608569in}}%
\pgfpathlineto{\pgfqpoint{1.855518in}{1.540925in}}%
\pgfpathlineto{\pgfqpoint{1.858037in}{1.569011in}}%
\pgfpathlineto{\pgfqpoint{1.860556in}{1.589929in}}%
\pgfpathlineto{\pgfqpoint{1.863074in}{1.596488in}}%
\pgfpathlineto{\pgfqpoint{1.865593in}{1.640284in}}%
\pgfpathlineto{\pgfqpoint{1.868112in}{1.656728in}}%
\pgfpathlineto{\pgfqpoint{1.870631in}{1.640638in}}%
\pgfpathlineto{\pgfqpoint{1.873149in}{1.648680in}}%
\pgfpathlineto{\pgfqpoint{1.875668in}{1.661640in}}%
\pgfpathlineto{\pgfqpoint{1.878187in}{1.666025in}}%
\pgfpathlineto{\pgfqpoint{1.880706in}{1.692909in}}%
\pgfpathlineto{\pgfqpoint{1.883224in}{1.707583in}}%
\pgfpathlineto{\pgfqpoint{1.885743in}{1.726571in}}%
\pgfpathlineto{\pgfqpoint{1.888262in}{1.741801in}}%
\pgfpathlineto{\pgfqpoint{1.890781in}{1.733651in}}%
\pgfpathlineto{\pgfqpoint{1.893299in}{1.738033in}}%
\pgfpathlineto{\pgfqpoint{1.895818in}{1.734560in}}%
\pgfpathlineto{\pgfqpoint{1.898337in}{1.758873in}}%
\pgfpathlineto{\pgfqpoint{1.900856in}{1.752372in}}%
\pgfpathlineto{\pgfqpoint{1.903374in}{1.768029in}}%
\pgfpathlineto{\pgfqpoint{1.905893in}{1.782886in}}%
\pgfpathlineto{\pgfqpoint{1.908412in}{1.778498in}}%
\pgfpathlineto{\pgfqpoint{1.910931in}{1.735369in}}%
\pgfpathlineto{\pgfqpoint{1.913449in}{1.747200in}}%
\pgfpathlineto{\pgfqpoint{1.915968in}{1.788767in}}%
\pgfpathlineto{\pgfqpoint{1.918487in}{1.798704in}}%
\pgfpathlineto{\pgfqpoint{1.921006in}{1.773229in}}%
\pgfpathlineto{\pgfqpoint{1.923524in}{1.779684in}}%
\pgfpathlineto{\pgfqpoint{1.926043in}{1.791742in}}%
\pgfpathlineto{\pgfqpoint{1.928562in}{1.807941in}}%
\pgfpathlineto{\pgfqpoint{1.931081in}{1.780108in}}%
\pgfpathlineto{\pgfqpoint{1.933599in}{1.805143in}}%
\pgfpathlineto{\pgfqpoint{1.936118in}{1.840901in}}%
\pgfpathlineto{\pgfqpoint{1.938637in}{1.874343in}}%
\pgfpathlineto{\pgfqpoint{1.941156in}{1.856620in}}%
\pgfpathlineto{\pgfqpoint{1.943674in}{1.864795in}}%
\pgfpathlineto{\pgfqpoint{1.946193in}{1.832363in}}%
\pgfpathlineto{\pgfqpoint{1.951231in}{1.842862in}}%
\pgfpathlineto{\pgfqpoint{1.953749in}{1.812788in}}%
\pgfpathlineto{\pgfqpoint{1.956268in}{1.807472in}}%
\pgfpathlineto{\pgfqpoint{1.958787in}{1.802525in}}%
\pgfpathlineto{\pgfqpoint{1.963824in}{1.837179in}}%
\pgfpathlineto{\pgfqpoint{1.966343in}{1.830629in}}%
\pgfpathlineto{\pgfqpoint{1.968862in}{1.839974in}}%
\pgfpathlineto{\pgfqpoint{1.971381in}{1.829559in}}%
\pgfpathlineto{\pgfqpoint{1.973899in}{1.809716in}}%
\pgfpathlineto{\pgfqpoint{1.978937in}{1.756043in}}%
\pgfpathlineto{\pgfqpoint{1.981456in}{1.737900in}}%
\pgfpathlineto{\pgfqpoint{1.983974in}{1.756420in}}%
\pgfpathlineto{\pgfqpoint{1.986493in}{1.742823in}}%
\pgfpathlineto{\pgfqpoint{1.989012in}{1.696428in}}%
\pgfpathlineto{\pgfqpoint{1.991531in}{1.710035in}}%
\pgfpathlineto{\pgfqpoint{1.994049in}{1.709012in}}%
\pgfpathlineto{\pgfqpoint{1.996568in}{1.720643in}}%
\pgfpathlineto{\pgfqpoint{1.999087in}{1.707220in}}%
\pgfpathlineto{\pgfqpoint{2.001606in}{1.732993in}}%
\pgfpathlineto{\pgfqpoint{2.004124in}{1.736117in}}%
\pgfpathlineto{\pgfqpoint{2.006643in}{1.746800in}}%
\pgfpathlineto{\pgfqpoint{2.009162in}{1.758230in}}%
\pgfpathlineto{\pgfqpoint{2.011681in}{1.787270in}}%
\pgfpathlineto{\pgfqpoint{2.014199in}{1.778997in}}%
\pgfpathlineto{\pgfqpoint{2.016718in}{1.784231in}}%
\pgfpathlineto{\pgfqpoint{2.019237in}{1.800266in}}%
\pgfpathlineto{\pgfqpoint{2.021756in}{1.803666in}}%
\pgfpathlineto{\pgfqpoint{2.024274in}{1.835852in}}%
\pgfpathlineto{\pgfqpoint{2.026793in}{1.843119in}}%
\pgfpathlineto{\pgfqpoint{2.029312in}{1.814118in}}%
\pgfpathlineto{\pgfqpoint{2.031831in}{1.818293in}}%
\pgfpathlineto{\pgfqpoint{2.034349in}{1.849148in}}%
\pgfpathlineto{\pgfqpoint{2.036868in}{1.813388in}}%
\pgfpathlineto{\pgfqpoint{2.039387in}{1.830699in}}%
\pgfpathlineto{\pgfqpoint{2.041906in}{1.829034in}}%
\pgfpathlineto{\pgfqpoint{2.044424in}{1.831790in}}%
\pgfpathlineto{\pgfqpoint{2.046943in}{1.823111in}}%
\pgfpathlineto{\pgfqpoint{2.049462in}{1.845127in}}%
\pgfpathlineto{\pgfqpoint{2.051981in}{1.864111in}}%
\pgfpathlineto{\pgfqpoint{2.054499in}{1.867404in}}%
\pgfpathlineto{\pgfqpoint{2.057018in}{1.888964in}}%
\pgfpathlineto{\pgfqpoint{2.059537in}{1.880887in}}%
\pgfpathlineto{\pgfqpoint{2.062056in}{1.881531in}}%
\pgfpathlineto{\pgfqpoint{2.064574in}{1.855877in}}%
\pgfpathlineto{\pgfqpoint{2.067093in}{1.858675in}}%
\pgfpathlineto{\pgfqpoint{2.069612in}{1.864811in}}%
\pgfpathlineto{\pgfqpoint{2.072131in}{1.868107in}}%
\pgfpathlineto{\pgfqpoint{2.074649in}{1.868104in}}%
\pgfpathlineto{\pgfqpoint{2.077168in}{1.866741in}}%
\pgfpathlineto{\pgfqpoint{2.079687in}{1.883614in}}%
\pgfpathlineto{\pgfqpoint{2.082206in}{1.887988in}}%
\pgfpathlineto{\pgfqpoint{2.084724in}{1.865681in}}%
\pgfpathlineto{\pgfqpoint{2.087243in}{1.902679in}}%
\pgfpathlineto{\pgfqpoint{2.089762in}{1.920679in}}%
\pgfpathlineto{\pgfqpoint{2.092281in}{1.924999in}}%
\pgfpathlineto{\pgfqpoint{2.094799in}{1.917828in}}%
\pgfpathlineto{\pgfqpoint{2.097318in}{1.944946in}}%
\pgfpathlineto{\pgfqpoint{2.099837in}{1.962093in}}%
\pgfpathlineto{\pgfqpoint{2.102356in}{1.985449in}}%
\pgfpathlineto{\pgfqpoint{2.104874in}{1.993524in}}%
\pgfpathlineto{\pgfqpoint{2.107393in}{1.997134in}}%
\pgfpathlineto{\pgfqpoint{2.109912in}{2.009751in}}%
\pgfpathlineto{\pgfqpoint{2.112431in}{2.005315in}}%
\pgfpathlineto{\pgfqpoint{2.114949in}{2.012343in}}%
\pgfpathlineto{\pgfqpoint{2.117468in}{1.990192in}}%
\pgfpathlineto{\pgfqpoint{2.119987in}{1.977387in}}%
\pgfpathlineto{\pgfqpoint{2.122506in}{1.973915in}}%
\pgfpathlineto{\pgfqpoint{2.125024in}{1.969570in}}%
\pgfpathlineto{\pgfqpoint{2.127543in}{1.964250in}}%
\pgfpathlineto{\pgfqpoint{2.130062in}{1.977050in}}%
\pgfpathlineto{\pgfqpoint{2.132581in}{1.938809in}}%
\pgfpathlineto{\pgfqpoint{2.135099in}{1.974276in}}%
\pgfpathlineto{\pgfqpoint{2.137618in}{1.959369in}}%
\pgfpathlineto{\pgfqpoint{2.140137in}{1.978840in}}%
\pgfpathlineto{\pgfqpoint{2.142656in}{1.942268in}}%
\pgfpathlineto{\pgfqpoint{2.145174in}{1.936597in}}%
\pgfpathlineto{\pgfqpoint{2.147693in}{1.925487in}}%
\pgfpathlineto{\pgfqpoint{2.150212in}{1.945025in}}%
\pgfpathlineto{\pgfqpoint{2.152731in}{1.976805in}}%
\pgfpathlineto{\pgfqpoint{2.155249in}{1.990545in}}%
\pgfpathlineto{\pgfqpoint{2.157768in}{1.967979in}}%
\pgfpathlineto{\pgfqpoint{2.160287in}{1.980490in}}%
\pgfpathlineto{\pgfqpoint{2.162806in}{1.994353in}}%
\pgfpathlineto{\pgfqpoint{2.165324in}{1.984090in}}%
\pgfpathlineto{\pgfqpoint{2.167843in}{2.002903in}}%
\pgfpathlineto{\pgfqpoint{2.170362in}{1.993100in}}%
\pgfpathlineto{\pgfqpoint{2.172881in}{2.001648in}}%
\pgfpathlineto{\pgfqpoint{2.175399in}{2.008869in}}%
\pgfpathlineto{\pgfqpoint{2.177918in}{1.992021in}}%
\pgfpathlineto{\pgfqpoint{2.180437in}{1.983863in}}%
\pgfpathlineto{\pgfqpoint{2.182956in}{1.997174in}}%
\pgfpathlineto{\pgfqpoint{2.185474in}{2.016877in}}%
\pgfpathlineto{\pgfqpoint{2.187993in}{2.017708in}}%
\pgfpathlineto{\pgfqpoint{2.190512in}{2.011447in}}%
\pgfpathlineto{\pgfqpoint{2.193031in}{2.028346in}}%
\pgfpathlineto{\pgfqpoint{2.195549in}{2.030497in}}%
\pgfpathlineto{\pgfqpoint{2.198068in}{2.038298in}}%
\pgfpathlineto{\pgfqpoint{2.200587in}{2.050246in}}%
\pgfpathlineto{\pgfqpoint{2.203106in}{2.061029in}}%
\pgfpathlineto{\pgfqpoint{2.205624in}{2.042498in}}%
\pgfpathlineto{\pgfqpoint{2.208143in}{2.044238in}}%
\pgfpathlineto{\pgfqpoint{2.210662in}{2.027180in}}%
\pgfpathlineto{\pgfqpoint{2.213181in}{2.013771in}}%
\pgfpathlineto{\pgfqpoint{2.215699in}{2.024031in}}%
\pgfpathlineto{\pgfqpoint{2.218218in}{2.014964in}}%
\pgfpathlineto{\pgfqpoint{2.220737in}{2.007046in}}%
\pgfpathlineto{\pgfqpoint{2.223256in}{2.007328in}}%
\pgfpathlineto{\pgfqpoint{2.225774in}{1.994308in}}%
\pgfpathlineto{\pgfqpoint{2.228293in}{2.014376in}}%
\pgfpathlineto{\pgfqpoint{2.230812in}{2.048515in}}%
\pgfpathlineto{\pgfqpoint{2.233331in}{2.009364in}}%
\pgfpathlineto{\pgfqpoint{2.235849in}{1.975405in}}%
\pgfpathlineto{\pgfqpoint{2.238368in}{1.982720in}}%
\pgfpathlineto{\pgfqpoint{2.240887in}{1.961029in}}%
\pgfpathlineto{\pgfqpoint{2.243406in}{1.981996in}}%
\pgfpathlineto{\pgfqpoint{2.245924in}{1.971669in}}%
\pgfpathlineto{\pgfqpoint{2.248443in}{1.974112in}}%
\pgfpathlineto{\pgfqpoint{2.250962in}{2.001882in}}%
\pgfpathlineto{\pgfqpoint{2.253481in}{1.966009in}}%
\pgfpathlineto{\pgfqpoint{2.255999in}{1.978655in}}%
\pgfpathlineto{\pgfqpoint{2.258518in}{1.995504in}}%
\pgfpathlineto{\pgfqpoint{2.261037in}{1.976917in}}%
\pgfpathlineto{\pgfqpoint{2.263556in}{1.940434in}}%
\pgfpathlineto{\pgfqpoint{2.266074in}{1.943380in}}%
\pgfpathlineto{\pgfqpoint{2.268593in}{1.938424in}}%
\pgfpathlineto{\pgfqpoint{2.271112in}{1.945357in}}%
\pgfpathlineto{\pgfqpoint{2.273631in}{1.943204in}}%
\pgfpathlineto{\pgfqpoint{2.276149in}{1.954167in}}%
\pgfpathlineto{\pgfqpoint{2.278668in}{1.974305in}}%
\pgfpathlineto{\pgfqpoint{2.281187in}{2.015069in}}%
\pgfpathlineto{\pgfqpoint{2.283706in}{1.980204in}}%
\pgfpathlineto{\pgfqpoint{2.286224in}{1.973386in}}%
\pgfpathlineto{\pgfqpoint{2.288743in}{1.984981in}}%
\pgfpathlineto{\pgfqpoint{2.291262in}{1.993718in}}%
\pgfpathlineto{\pgfqpoint{2.293781in}{2.006029in}}%
\pgfpathlineto{\pgfqpoint{2.296299in}{2.028907in}}%
\pgfpathlineto{\pgfqpoint{2.298818in}{2.021130in}}%
\pgfpathlineto{\pgfqpoint{2.301337in}{2.029399in}}%
\pgfpathlineto{\pgfqpoint{2.303856in}{2.021790in}}%
\pgfpathlineto{\pgfqpoint{2.306374in}{2.026199in}}%
\pgfpathlineto{\pgfqpoint{2.308893in}{2.028627in}}%
\pgfpathlineto{\pgfqpoint{2.311412in}{2.034030in}}%
\pgfpathlineto{\pgfqpoint{2.313931in}{2.015911in}}%
\pgfpathlineto{\pgfqpoint{2.316449in}{2.012635in}}%
\pgfpathlineto{\pgfqpoint{2.318968in}{2.033744in}}%
\pgfpathlineto{\pgfqpoint{2.321487in}{2.051762in}}%
\pgfpathlineto{\pgfqpoint{2.324006in}{2.049777in}}%
\pgfpathlineto{\pgfqpoint{2.326524in}{2.062341in}}%
\pgfpathlineto{\pgfqpoint{2.329043in}{2.057141in}}%
\pgfpathlineto{\pgfqpoint{2.331562in}{2.061695in}}%
\pgfpathlineto{\pgfqpoint{2.334081in}{2.090070in}}%
\pgfpathlineto{\pgfqpoint{2.336599in}{2.098790in}}%
\pgfpathlineto{\pgfqpoint{2.339118in}{2.098057in}}%
\pgfpathlineto{\pgfqpoint{2.341637in}{2.093133in}}%
\pgfpathlineto{\pgfqpoint{2.344156in}{2.093257in}}%
\pgfpathlineto{\pgfqpoint{2.346674in}{2.094388in}}%
\pgfpathlineto{\pgfqpoint{2.349193in}{2.076605in}}%
\pgfpathlineto{\pgfqpoint{2.351712in}{2.096840in}}%
\pgfpathlineto{\pgfqpoint{2.354231in}{2.101686in}}%
\pgfpathlineto{\pgfqpoint{2.356749in}{2.098709in}}%
\pgfpathlineto{\pgfqpoint{2.359268in}{2.088772in}}%
\pgfpathlineto{\pgfqpoint{2.361787in}{2.085385in}}%
\pgfpathlineto{\pgfqpoint{2.364306in}{2.105052in}}%
\pgfpathlineto{\pgfqpoint{2.366824in}{2.089058in}}%
\pgfpathlineto{\pgfqpoint{2.369343in}{2.077412in}}%
\pgfpathlineto{\pgfqpoint{2.371862in}{2.077691in}}%
\pgfpathlineto{\pgfqpoint{2.374381in}{2.049508in}}%
\pgfpathlineto{\pgfqpoint{2.376899in}{2.060150in}}%
\pgfpathlineto{\pgfqpoint{2.379418in}{2.064211in}}%
\pgfpathlineto{\pgfqpoint{2.381937in}{2.034920in}}%
\pgfpathlineto{\pgfqpoint{2.384456in}{2.038999in}}%
\pgfpathlineto{\pgfqpoint{2.386974in}{2.037096in}}%
\pgfpathlineto{\pgfqpoint{2.389493in}{2.012091in}}%
\pgfpathlineto{\pgfqpoint{2.392012in}{2.045301in}}%
\pgfpathlineto{\pgfqpoint{2.394531in}{2.099110in}}%
\pgfpathlineto{\pgfqpoint{2.397049in}{2.137315in}}%
\pgfpathlineto{\pgfqpoint{2.399568in}{2.102012in}}%
\pgfpathlineto{\pgfqpoint{2.402087in}{2.114875in}}%
\pgfpathlineto{\pgfqpoint{2.404606in}{2.124709in}}%
\pgfpathlineto{\pgfqpoint{2.407124in}{2.120763in}}%
\pgfpathlineto{\pgfqpoint{2.409643in}{2.095587in}}%
\pgfpathlineto{\pgfqpoint{2.412162in}{2.118926in}}%
\pgfpathlineto{\pgfqpoint{2.414681in}{2.129136in}}%
\pgfpathlineto{\pgfqpoint{2.417199in}{2.156714in}}%
\pgfpathlineto{\pgfqpoint{2.419718in}{2.133033in}}%
\pgfpathlineto{\pgfqpoint{2.422237in}{2.140364in}}%
\pgfpathlineto{\pgfqpoint{2.424756in}{2.121087in}}%
\pgfpathlineto{\pgfqpoint{2.427274in}{2.083285in}}%
\pgfpathlineto{\pgfqpoint{2.429793in}{2.061272in}}%
\pgfpathlineto{\pgfqpoint{2.432312in}{2.091241in}}%
\pgfpathlineto{\pgfqpoint{2.434831in}{2.099792in}}%
\pgfpathlineto{\pgfqpoint{2.437349in}{2.103382in}}%
\pgfpathlineto{\pgfqpoint{2.439868in}{2.139153in}}%
\pgfpathlineto{\pgfqpoint{2.442387in}{2.131112in}}%
\pgfpathlineto{\pgfqpoint{2.444906in}{2.140693in}}%
\pgfpathlineto{\pgfqpoint{2.447424in}{2.127132in}}%
\pgfpathlineto{\pgfqpoint{2.449943in}{2.116868in}}%
\pgfpathlineto{\pgfqpoint{2.452462in}{2.084085in}}%
\pgfpathlineto{\pgfqpoint{2.454981in}{2.107067in}}%
\pgfpathlineto{\pgfqpoint{2.457499in}{2.108195in}}%
\pgfpathlineto{\pgfqpoint{2.460018in}{2.101943in}}%
\pgfpathlineto{\pgfqpoint{2.462537in}{2.133095in}}%
\pgfpathlineto{\pgfqpoint{2.465056in}{2.179949in}}%
\pgfpathlineto{\pgfqpoint{2.467574in}{2.198712in}}%
\pgfpathlineto{\pgfqpoint{2.470093in}{2.178932in}}%
\pgfpathlineto{\pgfqpoint{2.472612in}{2.170844in}}%
\pgfpathlineto{\pgfqpoint{2.475131in}{2.198399in}}%
\pgfpathlineto{\pgfqpoint{2.477649in}{2.216849in}}%
\pgfpathlineto{\pgfqpoint{2.480168in}{2.215748in}}%
\pgfpathlineto{\pgfqpoint{2.482687in}{2.185910in}}%
\pgfpathlineto{\pgfqpoint{2.485206in}{2.209562in}}%
\pgfpathlineto{\pgfqpoint{2.487724in}{2.214106in}}%
\pgfpathlineto{\pgfqpoint{2.490243in}{2.199999in}}%
\pgfpathlineto{\pgfqpoint{2.492762in}{2.209043in}}%
\pgfpathlineto{\pgfqpoint{2.495281in}{2.204592in}}%
\pgfpathlineto{\pgfqpoint{2.497799in}{2.212033in}}%
\pgfpathlineto{\pgfqpoint{2.500318in}{2.205519in}}%
\pgfpathlineto{\pgfqpoint{2.502837in}{2.209550in}}%
\pgfpathlineto{\pgfqpoint{2.505356in}{2.212899in}}%
\pgfpathlineto{\pgfqpoint{2.507874in}{2.207492in}}%
\pgfpathlineto{\pgfqpoint{2.510393in}{2.232627in}}%
\pgfpathlineto{\pgfqpoint{2.512912in}{2.241566in}}%
\pgfpathlineto{\pgfqpoint{2.515431in}{2.249564in}}%
\pgfpathlineto{\pgfqpoint{2.517949in}{2.277527in}}%
\pgfpathlineto{\pgfqpoint{2.520468in}{2.297843in}}%
\pgfpathlineto{\pgfqpoint{2.522987in}{2.305761in}}%
\pgfpathlineto{\pgfqpoint{2.525506in}{2.323812in}}%
\pgfpathlineto{\pgfqpoint{2.528024in}{2.299742in}}%
\pgfpathlineto{\pgfqpoint{2.530543in}{2.306046in}}%
\pgfpathlineto{\pgfqpoint{2.533062in}{2.303554in}}%
\pgfpathlineto{\pgfqpoint{2.535581in}{2.276520in}}%
\pgfpathlineto{\pgfqpoint{2.538099in}{2.269257in}}%
\pgfpathlineto{\pgfqpoint{2.540618in}{2.283183in}}%
\pgfpathlineto{\pgfqpoint{2.543137in}{2.265207in}}%
\pgfpathlineto{\pgfqpoint{2.545656in}{2.275909in}}%
\pgfpathlineto{\pgfqpoint{2.548174in}{2.235715in}}%
\pgfpathlineto{\pgfqpoint{2.550693in}{2.255214in}}%
\pgfpathlineto{\pgfqpoint{2.553212in}{2.240850in}}%
\pgfpathlineto{\pgfqpoint{2.555731in}{2.213239in}}%
\pgfpathlineto{\pgfqpoint{2.558249in}{2.211525in}}%
\pgfpathlineto{\pgfqpoint{2.560768in}{2.235742in}}%
\pgfpathlineto{\pgfqpoint{2.563287in}{2.251249in}}%
\pgfpathlineto{\pgfqpoint{2.568324in}{2.210916in}}%
\pgfpathlineto{\pgfqpoint{2.570843in}{2.208885in}}%
\pgfpathlineto{\pgfqpoint{2.573362in}{2.203000in}}%
\pgfpathlineto{\pgfqpoint{2.575881in}{2.206403in}}%
\pgfpathlineto{\pgfqpoint{2.578399in}{2.195118in}}%
\pgfpathlineto{\pgfqpoint{2.580918in}{2.177291in}}%
\pgfpathlineto{\pgfqpoint{2.583437in}{2.175752in}}%
\pgfpathlineto{\pgfqpoint{2.585956in}{2.143386in}}%
\pgfpathlineto{\pgfqpoint{2.588474in}{2.129484in}}%
\pgfpathlineto{\pgfqpoint{2.590993in}{2.123974in}}%
\pgfpathlineto{\pgfqpoint{2.593512in}{2.132661in}}%
\pgfpathlineto{\pgfqpoint{2.596031in}{2.122730in}}%
\pgfpathlineto{\pgfqpoint{2.598549in}{2.142989in}}%
\pgfpathlineto{\pgfqpoint{2.601068in}{2.111239in}}%
\pgfpathlineto{\pgfqpoint{2.603587in}{2.115172in}}%
\pgfpathlineto{\pgfqpoint{2.606106in}{2.112223in}}%
\pgfpathlineto{\pgfqpoint{2.608624in}{2.129289in}}%
\pgfpathlineto{\pgfqpoint{2.611143in}{2.114411in}}%
\pgfpathlineto{\pgfqpoint{2.613662in}{2.133847in}}%
\pgfpathlineto{\pgfqpoint{2.616181in}{2.149058in}}%
\pgfpathlineto{\pgfqpoint{2.618699in}{2.141869in}}%
\pgfpathlineto{\pgfqpoint{2.621218in}{2.141999in}}%
\pgfpathlineto{\pgfqpoint{2.623737in}{2.136596in}}%
\pgfpathlineto{\pgfqpoint{2.626256in}{2.137018in}}%
\pgfpathlineto{\pgfqpoint{2.628774in}{2.135373in}}%
\pgfpathlineto{\pgfqpoint{2.631293in}{2.156606in}}%
\pgfpathlineto{\pgfqpoint{2.633812in}{2.162779in}}%
\pgfpathlineto{\pgfqpoint{2.636331in}{2.169233in}}%
\pgfpathlineto{\pgfqpoint{2.638849in}{2.183475in}}%
\pgfpathlineto{\pgfqpoint{2.641368in}{2.177351in}}%
\pgfpathlineto{\pgfqpoint{2.643887in}{2.191016in}}%
\pgfpathlineto{\pgfqpoint{2.646406in}{2.204060in}}%
\pgfpathlineto{\pgfqpoint{2.648924in}{2.174175in}}%
\pgfpathlineto{\pgfqpoint{2.651443in}{2.131680in}}%
\pgfpathlineto{\pgfqpoint{2.653962in}{2.114637in}}%
\pgfpathlineto{\pgfqpoint{2.656481in}{2.112398in}}%
\pgfpathlineto{\pgfqpoint{2.658999in}{2.103911in}}%
\pgfpathlineto{\pgfqpoint{2.661518in}{2.113561in}}%
\pgfpathlineto{\pgfqpoint{2.664037in}{2.124369in}}%
\pgfpathlineto{\pgfqpoint{2.666556in}{2.113777in}}%
\pgfpathlineto{\pgfqpoint{2.669074in}{2.118784in}}%
\pgfpathlineto{\pgfqpoint{2.671593in}{2.098019in}}%
\pgfpathlineto{\pgfqpoint{2.674112in}{2.121852in}}%
\pgfpathlineto{\pgfqpoint{2.676631in}{2.105666in}}%
\pgfpathlineto{\pgfqpoint{2.679149in}{2.077243in}}%
\pgfpathlineto{\pgfqpoint{2.681668in}{2.072045in}}%
\pgfpathlineto{\pgfqpoint{2.684187in}{2.073687in}}%
\pgfpathlineto{\pgfqpoint{2.686706in}{2.048787in}}%
\pgfpathlineto{\pgfqpoint{2.689224in}{2.037916in}}%
\pgfpathlineto{\pgfqpoint{2.694262in}{1.989959in}}%
\pgfpathlineto{\pgfqpoint{2.699299in}{2.007961in}}%
\pgfpathlineto{\pgfqpoint{2.701818in}{1.986429in}}%
\pgfpathlineto{\pgfqpoint{2.704337in}{1.981224in}}%
\pgfpathlineto{\pgfqpoint{2.706856in}{1.985696in}}%
\pgfpathlineto{\pgfqpoint{2.709374in}{1.980626in}}%
\pgfpathlineto{\pgfqpoint{2.711893in}{1.993567in}}%
\pgfpathlineto{\pgfqpoint{2.714412in}{2.015199in}}%
\pgfpathlineto{\pgfqpoint{2.716931in}{2.010779in}}%
\pgfpathlineto{\pgfqpoint{2.719449in}{2.011940in}}%
\pgfpathlineto{\pgfqpoint{2.721968in}{1.989384in}}%
\pgfpathlineto{\pgfqpoint{2.724487in}{2.016448in}}%
\pgfpathlineto{\pgfqpoint{2.727006in}{2.007047in}}%
\pgfpathlineto{\pgfqpoint{2.729524in}{2.005812in}}%
\pgfpathlineto{\pgfqpoint{2.732043in}{1.986020in}}%
\pgfpathlineto{\pgfqpoint{2.734562in}{2.010451in}}%
\pgfpathlineto{\pgfqpoint{2.737081in}{2.006290in}}%
\pgfpathlineto{\pgfqpoint{2.739599in}{2.024446in}}%
\pgfpathlineto{\pgfqpoint{2.742118in}{1.981123in}}%
\pgfpathlineto{\pgfqpoint{2.744637in}{1.992941in}}%
\pgfpathlineto{\pgfqpoint{2.747156in}{1.963431in}}%
\pgfpathlineto{\pgfqpoint{2.749674in}{1.958920in}}%
\pgfpathlineto{\pgfqpoint{2.752193in}{1.962848in}}%
\pgfpathlineto{\pgfqpoint{2.754712in}{1.943418in}}%
\pgfpathlineto{\pgfqpoint{2.757231in}{1.934387in}}%
\pgfpathlineto{\pgfqpoint{2.759749in}{1.976964in}}%
\pgfpathlineto{\pgfqpoint{2.762268in}{1.971507in}}%
\pgfpathlineto{\pgfqpoint{2.764787in}{1.958622in}}%
\pgfpathlineto{\pgfqpoint{2.767306in}{1.993536in}}%
\pgfpathlineto{\pgfqpoint{2.769824in}{2.000555in}}%
\pgfpathlineto{\pgfqpoint{2.772343in}{1.983960in}}%
\pgfpathlineto{\pgfqpoint{2.774862in}{1.996064in}}%
\pgfpathlineto{\pgfqpoint{2.777381in}{2.027403in}}%
\pgfpathlineto{\pgfqpoint{2.779899in}{2.034471in}}%
\pgfpathlineto{\pgfqpoint{2.782418in}{2.038412in}}%
\pgfpathlineto{\pgfqpoint{2.784937in}{2.019445in}}%
\pgfpathlineto{\pgfqpoint{2.787456in}{2.042297in}}%
\pgfpathlineto{\pgfqpoint{2.789974in}{2.067071in}}%
\pgfpathlineto{\pgfqpoint{2.792493in}{2.058225in}}%
\pgfpathlineto{\pgfqpoint{2.795012in}{2.058883in}}%
\pgfpathlineto{\pgfqpoint{2.797531in}{2.045430in}}%
\pgfpathlineto{\pgfqpoint{2.800049in}{2.056324in}}%
\pgfpathlineto{\pgfqpoint{2.802568in}{2.046119in}}%
\pgfpathlineto{\pgfqpoint{2.805087in}{2.026538in}}%
\pgfpathlineto{\pgfqpoint{2.807606in}{2.013714in}}%
\pgfpathlineto{\pgfqpoint{2.810124in}{2.025317in}}%
\pgfpathlineto{\pgfqpoint{2.812643in}{2.042046in}}%
\pgfpathlineto{\pgfqpoint{2.815162in}{2.050907in}}%
\pgfpathlineto{\pgfqpoint{2.817681in}{2.028830in}}%
\pgfpathlineto{\pgfqpoint{2.820199in}{2.021368in}}%
\pgfpathlineto{\pgfqpoint{2.822718in}{2.032894in}}%
\pgfpathlineto{\pgfqpoint{2.825237in}{2.040941in}}%
\pgfpathlineto{\pgfqpoint{2.827756in}{2.038796in}}%
\pgfpathlineto{\pgfqpoint{2.830274in}{2.011963in}}%
\pgfpathlineto{\pgfqpoint{2.832793in}{2.023895in}}%
\pgfpathlineto{\pgfqpoint{2.835312in}{2.018131in}}%
\pgfpathlineto{\pgfqpoint{2.837831in}{2.029199in}}%
\pgfpathlineto{\pgfqpoint{2.840349in}{2.022637in}}%
\pgfpathlineto{\pgfqpoint{2.842868in}{2.030842in}}%
\pgfpathlineto{\pgfqpoint{2.845387in}{2.007251in}}%
\pgfpathlineto{\pgfqpoint{2.847906in}{2.011443in}}%
\pgfpathlineto{\pgfqpoint{2.850424in}{1.993452in}}%
\pgfpathlineto{\pgfqpoint{2.852943in}{1.984821in}}%
\pgfpathlineto{\pgfqpoint{2.855462in}{1.963393in}}%
\pgfpathlineto{\pgfqpoint{2.857981in}{1.979510in}}%
\pgfpathlineto{\pgfqpoint{2.860499in}{1.989339in}}%
\pgfpathlineto{\pgfqpoint{2.863018in}{1.994387in}}%
\pgfpathlineto{\pgfqpoint{2.865537in}{2.005351in}}%
\pgfpathlineto{\pgfqpoint{2.868056in}{2.024623in}}%
\pgfpathlineto{\pgfqpoint{2.870574in}{2.026612in}}%
\pgfpathlineto{\pgfqpoint{2.873093in}{2.038207in}}%
\pgfpathlineto{\pgfqpoint{2.875612in}{2.021483in}}%
\pgfpathlineto{\pgfqpoint{2.878131in}{1.985725in}}%
\pgfpathlineto{\pgfqpoint{2.880649in}{1.996322in}}%
\pgfpathlineto{\pgfqpoint{2.883168in}{2.014874in}}%
\pgfpathlineto{\pgfqpoint{2.888206in}{2.054409in}}%
\pgfpathlineto{\pgfqpoint{2.890724in}{2.066393in}}%
\pgfpathlineto{\pgfqpoint{2.893243in}{2.095141in}}%
\pgfpathlineto{\pgfqpoint{2.895762in}{2.100214in}}%
\pgfpathlineto{\pgfqpoint{2.898281in}{2.138398in}}%
\pgfpathlineto{\pgfqpoint{2.900799in}{2.132728in}}%
\pgfpathlineto{\pgfqpoint{2.903318in}{2.119622in}}%
\pgfpathlineto{\pgfqpoint{2.905837in}{2.124556in}}%
\pgfpathlineto{\pgfqpoint{2.908356in}{2.128074in}}%
\pgfpathlineto{\pgfqpoint{2.910874in}{2.104149in}}%
\pgfpathlineto{\pgfqpoint{2.913393in}{2.122757in}}%
\pgfpathlineto{\pgfqpoint{2.915912in}{2.122881in}}%
\pgfpathlineto{\pgfqpoint{2.918431in}{2.123666in}}%
\pgfpathlineto{\pgfqpoint{2.920949in}{2.117966in}}%
\pgfpathlineto{\pgfqpoint{2.923468in}{2.144983in}}%
\pgfpathlineto{\pgfqpoint{2.925987in}{2.150125in}}%
\pgfpathlineto{\pgfqpoint{2.928506in}{2.124109in}}%
\pgfpathlineto{\pgfqpoint{2.931024in}{2.132685in}}%
\pgfpathlineto{\pgfqpoint{2.933543in}{2.112574in}}%
\pgfpathlineto{\pgfqpoint{2.936062in}{2.104989in}}%
\pgfpathlineto{\pgfqpoint{2.938581in}{2.094179in}}%
\pgfpathlineto{\pgfqpoint{2.941099in}{2.086912in}}%
\pgfpathlineto{\pgfqpoint{2.943618in}{2.091196in}}%
\pgfpathlineto{\pgfqpoint{2.946137in}{2.094891in}}%
\pgfpathlineto{\pgfqpoint{2.948656in}{2.085453in}}%
\pgfpathlineto{\pgfqpoint{2.951174in}{2.057495in}}%
\pgfpathlineto{\pgfqpoint{2.953693in}{2.045827in}}%
\pgfpathlineto{\pgfqpoint{2.956212in}{2.049102in}}%
\pgfpathlineto{\pgfqpoint{2.958731in}{2.032860in}}%
\pgfpathlineto{\pgfqpoint{2.961249in}{2.073704in}}%
\pgfpathlineto{\pgfqpoint{2.963768in}{2.061560in}}%
\pgfpathlineto{\pgfqpoint{2.966287in}{2.076625in}}%
\pgfpathlineto{\pgfqpoint{2.968806in}{2.054268in}}%
\pgfpathlineto{\pgfqpoint{2.971324in}{2.018412in}}%
\pgfpathlineto{\pgfqpoint{2.973843in}{2.015489in}}%
\pgfpathlineto{\pgfqpoint{2.976362in}{2.049169in}}%
\pgfpathlineto{\pgfqpoint{2.978881in}{2.036706in}}%
\pgfpathlineto{\pgfqpoint{2.981399in}{2.026000in}}%
\pgfpathlineto{\pgfqpoint{2.983918in}{2.068367in}}%
\pgfpathlineto{\pgfqpoint{2.986437in}{2.033006in}}%
\pgfpathlineto{\pgfqpoint{2.991474in}{2.035694in}}%
\pgfpathlineto{\pgfqpoint{2.993993in}{2.017810in}}%
\pgfpathlineto{\pgfqpoint{2.996512in}{2.023249in}}%
\pgfpathlineto{\pgfqpoint{2.999031in}{2.004316in}}%
\pgfpathlineto{\pgfqpoint{3.001549in}{2.020403in}}%
\pgfpathlineto{\pgfqpoint{3.004068in}{1.995286in}}%
\pgfpathlineto{\pgfqpoint{3.006587in}{1.996490in}}%
\pgfpathlineto{\pgfqpoint{3.009106in}{2.007083in}}%
\pgfpathlineto{\pgfqpoint{3.011624in}{2.012493in}}%
\pgfpathlineto{\pgfqpoint{3.014143in}{1.982956in}}%
\pgfpathlineto{\pgfqpoint{3.016662in}{1.991076in}}%
\pgfpathlineto{\pgfqpoint{3.019181in}{1.982984in}}%
\pgfpathlineto{\pgfqpoint{3.021699in}{1.964210in}}%
\pgfpathlineto{\pgfqpoint{3.024218in}{1.963075in}}%
\pgfpathlineto{\pgfqpoint{3.026737in}{1.981819in}}%
\pgfpathlineto{\pgfqpoint{3.029256in}{1.969916in}}%
\pgfpathlineto{\pgfqpoint{3.031774in}{1.976162in}}%
\pgfpathlineto{\pgfqpoint{3.034293in}{2.009190in}}%
\pgfpathlineto{\pgfqpoint{3.036812in}{1.998144in}}%
\pgfpathlineto{\pgfqpoint{3.039331in}{1.981664in}}%
\pgfpathlineto{\pgfqpoint{3.041849in}{2.013056in}}%
\pgfpathlineto{\pgfqpoint{3.044368in}{2.006330in}}%
\pgfpathlineto{\pgfqpoint{3.046887in}{2.035532in}}%
\pgfpathlineto{\pgfqpoint{3.049406in}{2.029068in}}%
\pgfpathlineto{\pgfqpoint{3.051924in}{2.057023in}}%
\pgfpathlineto{\pgfqpoint{3.054443in}{2.067931in}}%
\pgfpathlineto{\pgfqpoint{3.056962in}{2.077398in}}%
\pgfpathlineto{\pgfqpoint{3.059481in}{2.113328in}}%
\pgfpathlineto{\pgfqpoint{3.061999in}{2.110754in}}%
\pgfpathlineto{\pgfqpoint{3.064518in}{2.126161in}}%
\pgfpathlineto{\pgfqpoint{3.067037in}{2.126667in}}%
\pgfpathlineto{\pgfqpoint{3.069556in}{2.140175in}}%
\pgfpathlineto{\pgfqpoint{3.072074in}{2.131636in}}%
\pgfpathlineto{\pgfqpoint{3.074593in}{2.142541in}}%
\pgfpathlineto{\pgfqpoint{3.077112in}{2.180362in}}%
\pgfpathlineto{\pgfqpoint{3.079631in}{2.139730in}}%
\pgfpathlineto{\pgfqpoint{3.082149in}{2.103403in}}%
\pgfpathlineto{\pgfqpoint{3.084668in}{2.061486in}}%
\pgfpathlineto{\pgfqpoint{3.087187in}{2.081861in}}%
\pgfpathlineto{\pgfqpoint{3.089706in}{2.075889in}}%
\pgfpathlineto{\pgfqpoint{3.092224in}{2.055318in}}%
\pgfpathlineto{\pgfqpoint{3.094743in}{2.053554in}}%
\pgfpathlineto{\pgfqpoint{3.097262in}{2.034881in}}%
\pgfpathlineto{\pgfqpoint{3.099781in}{2.027315in}}%
\pgfpathlineto{\pgfqpoint{3.102299in}{2.011965in}}%
\pgfpathlineto{\pgfqpoint{3.104818in}{2.026886in}}%
\pgfpathlineto{\pgfqpoint{3.107337in}{2.015056in}}%
\pgfpathlineto{\pgfqpoint{3.109856in}{2.037122in}}%
\pgfpathlineto{\pgfqpoint{3.112374in}{2.009110in}}%
\pgfpathlineto{\pgfqpoint{3.114893in}{2.018910in}}%
\pgfpathlineto{\pgfqpoint{3.117412in}{2.034613in}}%
\pgfpathlineto{\pgfqpoint{3.119931in}{2.008259in}}%
\pgfpathlineto{\pgfqpoint{3.122449in}{2.016425in}}%
\pgfpathlineto{\pgfqpoint{3.124968in}{2.005315in}}%
\pgfpathlineto{\pgfqpoint{3.127487in}{2.029816in}}%
\pgfpathlineto{\pgfqpoint{3.130006in}{2.043585in}}%
\pgfpathlineto{\pgfqpoint{3.132524in}{2.065630in}}%
\pgfpathlineto{\pgfqpoint{3.135043in}{2.066640in}}%
\pgfpathlineto{\pgfqpoint{3.137562in}{2.082980in}}%
\pgfpathlineto{\pgfqpoint{3.140081in}{2.080411in}}%
\pgfpathlineto{\pgfqpoint{3.142599in}{2.089633in}}%
\pgfpathlineto{\pgfqpoint{3.145118in}{2.085384in}}%
\pgfpathlineto{\pgfqpoint{3.147637in}{2.097704in}}%
\pgfpathlineto{\pgfqpoint{3.150156in}{2.115656in}}%
\pgfpathlineto{\pgfqpoint{3.152674in}{2.103717in}}%
\pgfpathlineto{\pgfqpoint{3.155193in}{2.095700in}}%
\pgfpathlineto{\pgfqpoint{3.157712in}{2.086497in}}%
\pgfpathlineto{\pgfqpoint{3.160231in}{2.057361in}}%
\pgfpathlineto{\pgfqpoint{3.162749in}{2.070218in}}%
\pgfpathlineto{\pgfqpoint{3.165268in}{2.064885in}}%
\pgfpathlineto{\pgfqpoint{3.167787in}{2.084763in}}%
\pgfpathlineto{\pgfqpoint{3.170306in}{2.091088in}}%
\pgfpathlineto{\pgfqpoint{3.172824in}{2.099042in}}%
\pgfpathlineto{\pgfqpoint{3.175343in}{2.104499in}}%
\pgfpathlineto{\pgfqpoint{3.177862in}{2.114346in}}%
\pgfpathlineto{\pgfqpoint{3.180381in}{2.101553in}}%
\pgfpathlineto{\pgfqpoint{3.182899in}{2.116721in}}%
\pgfpathlineto{\pgfqpoint{3.185418in}{2.098451in}}%
\pgfpathlineto{\pgfqpoint{3.187937in}{2.122456in}}%
\pgfpathlineto{\pgfqpoint{3.190456in}{2.120620in}}%
\pgfpathlineto{\pgfqpoint{3.192974in}{2.110930in}}%
\pgfpathlineto{\pgfqpoint{3.195493in}{2.147103in}}%
\pgfpathlineto{\pgfqpoint{3.198012in}{2.084770in}}%
\pgfpathlineto{\pgfqpoint{3.200531in}{2.096803in}}%
\pgfpathlineto{\pgfqpoint{3.203049in}{2.114195in}}%
\pgfpathlineto{\pgfqpoint{3.205568in}{2.097372in}}%
\pgfpathlineto{\pgfqpoint{3.208087in}{2.073459in}}%
\pgfpathlineto{\pgfqpoint{3.210606in}{2.059626in}}%
\pgfpathlineto{\pgfqpoint{3.213124in}{2.048804in}}%
\pgfpathlineto{\pgfqpoint{3.215643in}{2.005319in}}%
\pgfpathlineto{\pgfqpoint{3.218162in}{2.036963in}}%
\pgfpathlineto{\pgfqpoint{3.220681in}{2.037975in}}%
\pgfpathlineto{\pgfqpoint{3.223199in}{2.057143in}}%
\pgfpathlineto{\pgfqpoint{3.225718in}{2.069411in}}%
\pgfpathlineto{\pgfqpoint{3.228237in}{2.084307in}}%
\pgfpathlineto{\pgfqpoint{3.230756in}{2.102925in}}%
\pgfpathlineto{\pgfqpoint{3.233274in}{2.124752in}}%
\pgfpathlineto{\pgfqpoint{3.235793in}{2.137131in}}%
\pgfpathlineto{\pgfqpoint{3.238312in}{2.136942in}}%
\pgfpathlineto{\pgfqpoint{3.240831in}{2.140185in}}%
\pgfpathlineto{\pgfqpoint{3.243349in}{2.141951in}}%
\pgfpathlineto{\pgfqpoint{3.245868in}{2.149633in}}%
\pgfpathlineto{\pgfqpoint{3.248387in}{2.158954in}}%
\pgfpathlineto{\pgfqpoint{3.250906in}{2.169838in}}%
\pgfpathlineto{\pgfqpoint{3.253424in}{2.206577in}}%
\pgfpathlineto{\pgfqpoint{3.255943in}{2.198289in}}%
\pgfpathlineto{\pgfqpoint{3.258462in}{2.185937in}}%
\pgfpathlineto{\pgfqpoint{3.260981in}{2.178520in}}%
\pgfpathlineto{\pgfqpoint{3.263499in}{2.187725in}}%
\pgfpathlineto{\pgfqpoint{3.266018in}{2.171381in}}%
\pgfpathlineto{\pgfqpoint{3.268537in}{2.161395in}}%
\pgfpathlineto{\pgfqpoint{3.271056in}{2.177589in}}%
\pgfpathlineto{\pgfqpoint{3.273574in}{2.179261in}}%
\pgfpathlineto{\pgfqpoint{3.276093in}{2.228935in}}%
\pgfpathlineto{\pgfqpoint{3.278612in}{2.226711in}}%
\pgfpathlineto{\pgfqpoint{3.281131in}{2.221779in}}%
\pgfpathlineto{\pgfqpoint{3.283649in}{2.223756in}}%
\pgfpathlineto{\pgfqpoint{3.286168in}{2.222867in}}%
\pgfpathlineto{\pgfqpoint{3.288687in}{2.208184in}}%
\pgfpathlineto{\pgfqpoint{3.291206in}{2.221509in}}%
\pgfpathlineto{\pgfqpoint{3.293724in}{2.215549in}}%
\pgfpathlineto{\pgfqpoint{3.296243in}{2.189262in}}%
\pgfpathlineto{\pgfqpoint{3.298762in}{2.203784in}}%
\pgfpathlineto{\pgfqpoint{3.301281in}{2.226912in}}%
\pgfpathlineto{\pgfqpoint{3.303799in}{2.185629in}}%
\pgfpathlineto{\pgfqpoint{3.306318in}{2.189944in}}%
\pgfpathlineto{\pgfqpoint{3.308837in}{2.151302in}}%
\pgfpathlineto{\pgfqpoint{3.311356in}{2.139321in}}%
\pgfpathlineto{\pgfqpoint{3.313874in}{2.098542in}}%
\pgfpathlineto{\pgfqpoint{3.316393in}{2.085527in}}%
\pgfpathlineto{\pgfqpoint{3.318912in}{2.077653in}}%
\pgfpathlineto{\pgfqpoint{3.321431in}{2.094240in}}%
\pgfpathlineto{\pgfqpoint{3.323949in}{2.096426in}}%
\pgfpathlineto{\pgfqpoint{3.326468in}{2.090332in}}%
\pgfpathlineto{\pgfqpoint{3.328987in}{2.094251in}}%
\pgfpathlineto{\pgfqpoint{3.331506in}{2.103685in}}%
\pgfpathlineto{\pgfqpoint{3.334024in}{2.104995in}}%
\pgfpathlineto{\pgfqpoint{3.336543in}{2.095114in}}%
\pgfpathlineto{\pgfqpoint{3.339062in}{2.116645in}}%
\pgfpathlineto{\pgfqpoint{3.341581in}{2.150990in}}%
\pgfpathlineto{\pgfqpoint{3.344099in}{2.124708in}}%
\pgfpathlineto{\pgfqpoint{3.346618in}{2.113442in}}%
\pgfpathlineto{\pgfqpoint{3.349137in}{2.123025in}}%
\pgfpathlineto{\pgfqpoint{3.351656in}{2.119291in}}%
\pgfpathlineto{\pgfqpoint{3.354174in}{2.153775in}}%
\pgfpathlineto{\pgfqpoint{3.356693in}{2.128594in}}%
\pgfpathlineto{\pgfqpoint{3.359212in}{2.117515in}}%
\pgfpathlineto{\pgfqpoint{3.361731in}{2.105593in}}%
\pgfpathlineto{\pgfqpoint{3.364249in}{2.145097in}}%
\pgfpathlineto{\pgfqpoint{3.366768in}{2.152632in}}%
\pgfpathlineto{\pgfqpoint{3.369287in}{2.156113in}}%
\pgfpathlineto{\pgfqpoint{3.371806in}{2.197023in}}%
\pgfpathlineto{\pgfqpoint{3.374324in}{2.183233in}}%
\pgfpathlineto{\pgfqpoint{3.376843in}{2.156506in}}%
\pgfpathlineto{\pgfqpoint{3.379362in}{2.174411in}}%
\pgfpathlineto{\pgfqpoint{3.381881in}{2.153635in}}%
\pgfpathlineto{\pgfqpoint{3.384399in}{2.173293in}}%
\pgfpathlineto{\pgfqpoint{3.386918in}{2.188782in}}%
\pgfpathlineto{\pgfqpoint{3.389437in}{2.191042in}}%
\pgfpathlineto{\pgfqpoint{3.391956in}{2.217132in}}%
\pgfpathlineto{\pgfqpoint{3.394474in}{2.205993in}}%
\pgfpathlineto{\pgfqpoint{3.396993in}{2.205873in}}%
\pgfpathlineto{\pgfqpoint{3.399512in}{2.160802in}}%
\pgfpathlineto{\pgfqpoint{3.402031in}{2.150647in}}%
\pgfpathlineto{\pgfqpoint{3.404549in}{2.118057in}}%
\pgfpathlineto{\pgfqpoint{3.407068in}{2.102194in}}%
\pgfpathlineto{\pgfqpoint{3.409587in}{2.106169in}}%
\pgfpathlineto{\pgfqpoint{3.412106in}{2.129859in}}%
\pgfpathlineto{\pgfqpoint{3.414624in}{2.145575in}}%
\pgfpathlineto{\pgfqpoint{3.417143in}{2.113402in}}%
\pgfpathlineto{\pgfqpoint{3.419662in}{2.131840in}}%
\pgfpathlineto{\pgfqpoint{3.422181in}{2.114800in}}%
\pgfpathlineto{\pgfqpoint{3.424699in}{2.119420in}}%
\pgfpathlineto{\pgfqpoint{3.427218in}{2.078979in}}%
\pgfpathlineto{\pgfqpoint{3.429737in}{2.072867in}}%
\pgfpathlineto{\pgfqpoint{3.434774in}{2.039386in}}%
\pgfpathlineto{\pgfqpoint{3.437293in}{2.027266in}}%
\pgfpathlineto{\pgfqpoint{3.439812in}{2.001117in}}%
\pgfpathlineto{\pgfqpoint{3.442331in}{2.006795in}}%
\pgfpathlineto{\pgfqpoint{3.444849in}{2.032885in}}%
\pgfpathlineto{\pgfqpoint{3.447368in}{2.034596in}}%
\pgfpathlineto{\pgfqpoint{3.449887in}{2.035869in}}%
\pgfpathlineto{\pgfqpoint{3.452406in}{2.044709in}}%
\pgfpathlineto{\pgfqpoint{3.454924in}{2.040367in}}%
\pgfpathlineto{\pgfqpoint{3.457443in}{2.037181in}}%
\pgfpathlineto{\pgfqpoint{3.459962in}{2.110666in}}%
\pgfpathlineto{\pgfqpoint{3.462481in}{2.129016in}}%
\pgfpathlineto{\pgfqpoint{3.464999in}{2.184629in}}%
\pgfpathlineto{\pgfqpoint{3.467518in}{2.176902in}}%
\pgfpathlineto{\pgfqpoint{3.470037in}{2.186557in}}%
\pgfpathlineto{\pgfqpoint{3.472556in}{2.219057in}}%
\pgfpathlineto{\pgfqpoint{3.475074in}{2.214656in}}%
\pgfpathlineto{\pgfqpoint{3.477593in}{2.203726in}}%
\pgfpathlineto{\pgfqpoint{3.480112in}{2.227714in}}%
\pgfpathlineto{\pgfqpoint{3.482631in}{2.216619in}}%
\pgfpathlineto{\pgfqpoint{3.485149in}{2.232140in}}%
\pgfpathlineto{\pgfqpoint{3.487668in}{2.238844in}}%
\pgfpathlineto{\pgfqpoint{3.490187in}{2.223744in}}%
\pgfpathlineto{\pgfqpoint{3.492706in}{2.254933in}}%
\pgfpathlineto{\pgfqpoint{3.495224in}{2.231446in}}%
\pgfpathlineto{\pgfqpoint{3.497743in}{2.255244in}}%
\pgfpathlineto{\pgfqpoint{3.500262in}{2.250596in}}%
\pgfpathlineto{\pgfqpoint{3.502781in}{2.254247in}}%
\pgfpathlineto{\pgfqpoint{3.505299in}{2.266456in}}%
\pgfpathlineto{\pgfqpoint{3.507818in}{2.263995in}}%
\pgfpathlineto{\pgfqpoint{3.510337in}{2.281832in}}%
\pgfpathlineto{\pgfqpoint{3.512856in}{2.294395in}}%
\pgfpathlineto{\pgfqpoint{3.515374in}{2.290352in}}%
\pgfpathlineto{\pgfqpoint{3.517893in}{2.295609in}}%
\pgfpathlineto{\pgfqpoint{3.522931in}{2.324239in}}%
\pgfpathlineto{\pgfqpoint{3.525449in}{2.383365in}}%
\pgfpathlineto{\pgfqpoint{3.527968in}{2.389714in}}%
\pgfpathlineto{\pgfqpoint{3.530487in}{2.377688in}}%
\pgfpathlineto{\pgfqpoint{3.533006in}{2.381183in}}%
\pgfpathlineto{\pgfqpoint{3.535524in}{2.402146in}}%
\pgfpathlineto{\pgfqpoint{3.538043in}{2.403511in}}%
\pgfpathlineto{\pgfqpoint{3.540562in}{2.395445in}}%
\pgfpathlineto{\pgfqpoint{3.543081in}{2.405880in}}%
\pgfpathlineto{\pgfqpoint{3.545599in}{2.438986in}}%
\pgfpathlineto{\pgfqpoint{3.548118in}{2.485215in}}%
\pgfpathlineto{\pgfqpoint{3.550637in}{2.477846in}}%
\pgfpathlineto{\pgfqpoint{3.553156in}{2.487404in}}%
\pgfpathlineto{\pgfqpoint{3.555674in}{2.492104in}}%
\pgfpathlineto{\pgfqpoint{3.558193in}{2.477171in}}%
\pgfpathlineto{\pgfqpoint{3.560712in}{2.488851in}}%
\pgfpathlineto{\pgfqpoint{3.563231in}{2.547946in}}%
\pgfpathlineto{\pgfqpoint{3.565749in}{2.582019in}}%
\pgfpathlineto{\pgfqpoint{3.568268in}{2.569214in}}%
\pgfpathlineto{\pgfqpoint{3.570787in}{2.572669in}}%
\pgfpathlineto{\pgfqpoint{3.573306in}{2.568102in}}%
\pgfpathlineto{\pgfqpoint{3.575824in}{2.589895in}}%
\pgfpathlineto{\pgfqpoint{3.578343in}{2.555992in}}%
\pgfpathlineto{\pgfqpoint{3.580862in}{2.583656in}}%
\pgfpathlineto{\pgfqpoint{3.583381in}{2.578050in}}%
\pgfpathlineto{\pgfqpoint{3.585899in}{2.566904in}}%
\pgfpathlineto{\pgfqpoint{3.588418in}{2.608740in}}%
\pgfpathlineto{\pgfqpoint{3.590937in}{2.621387in}}%
\pgfpathlineto{\pgfqpoint{3.593456in}{2.658716in}}%
\pgfpathlineto{\pgfqpoint{3.598493in}{2.722138in}}%
\pgfpathlineto{\pgfqpoint{3.601012in}{2.711477in}}%
\pgfpathlineto{\pgfqpoint{3.603531in}{2.713226in}}%
\pgfpathlineto{\pgfqpoint{3.606049in}{2.718207in}}%
\pgfpathlineto{\pgfqpoint{3.608568in}{2.708071in}}%
\pgfpathlineto{\pgfqpoint{3.611087in}{2.693598in}}%
\pgfpathlineto{\pgfqpoint{3.613606in}{2.679880in}}%
\pgfpathlineto{\pgfqpoint{3.616124in}{2.694583in}}%
\pgfpathlineto{\pgfqpoint{3.618643in}{2.691685in}}%
\pgfpathlineto{\pgfqpoint{3.621162in}{2.707951in}}%
\pgfpathlineto{\pgfqpoint{3.623681in}{2.719601in}}%
\pgfpathlineto{\pgfqpoint{3.626199in}{2.715823in}}%
\pgfpathlineto{\pgfqpoint{3.628718in}{2.749079in}}%
\pgfpathlineto{\pgfqpoint{3.631237in}{2.777508in}}%
\pgfpathlineto{\pgfqpoint{3.633756in}{2.799065in}}%
\pgfpathlineto{\pgfqpoint{3.636274in}{2.758815in}}%
\pgfpathlineto{\pgfqpoint{3.638793in}{2.762092in}}%
\pgfpathlineto{\pgfqpoint{3.641312in}{2.738720in}}%
\pgfpathlineto{\pgfqpoint{3.643831in}{2.738753in}}%
\pgfpathlineto{\pgfqpoint{3.646349in}{2.729110in}}%
\pgfpathlineto{\pgfqpoint{3.648868in}{2.712362in}}%
\pgfpathlineto{\pgfqpoint{3.651387in}{2.705276in}}%
\pgfpathlineto{\pgfqpoint{3.653906in}{2.685904in}}%
\pgfpathlineto{\pgfqpoint{3.656424in}{2.683712in}}%
\pgfpathlineto{\pgfqpoint{3.658943in}{2.699620in}}%
\pgfpathlineto{\pgfqpoint{3.661462in}{2.706900in}}%
\pgfpathlineto{\pgfqpoint{3.663981in}{2.721111in}}%
\pgfpathlineto{\pgfqpoint{3.666499in}{2.727495in}}%
\pgfpathlineto{\pgfqpoint{3.669018in}{2.753066in}}%
\pgfpathlineto{\pgfqpoint{3.671537in}{2.740282in}}%
\pgfpathlineto{\pgfqpoint{3.674056in}{2.729809in}}%
\pgfpathlineto{\pgfqpoint{3.676574in}{2.778452in}}%
\pgfpathlineto{\pgfqpoint{3.679093in}{2.782869in}}%
\pgfpathlineto{\pgfqpoint{3.681612in}{2.796429in}}%
\pgfpathlineto{\pgfqpoint{3.684131in}{2.796561in}}%
\pgfpathlineto{\pgfqpoint{3.686649in}{2.815443in}}%
\pgfpathlineto{\pgfqpoint{3.689168in}{2.781402in}}%
\pgfpathlineto{\pgfqpoint{3.691687in}{2.771990in}}%
\pgfpathlineto{\pgfqpoint{3.694206in}{2.741215in}}%
\pgfpathlineto{\pgfqpoint{3.699243in}{2.806907in}}%
\pgfpathlineto{\pgfqpoint{3.701762in}{2.781133in}}%
\pgfpathlineto{\pgfqpoint{3.704281in}{2.751172in}}%
\pgfpathlineto{\pgfqpoint{3.706799in}{2.735577in}}%
\pgfpathlineto{\pgfqpoint{3.709318in}{2.758371in}}%
\pgfpathlineto{\pgfqpoint{3.711837in}{2.749952in}}%
\pgfpathlineto{\pgfqpoint{3.714356in}{2.775036in}}%
\pgfpathlineto{\pgfqpoint{3.716874in}{2.760784in}}%
\pgfpathlineto{\pgfqpoint{3.719393in}{2.775172in}}%
\pgfpathlineto{\pgfqpoint{3.721912in}{2.775784in}}%
\pgfpathlineto{\pgfqpoint{3.724431in}{2.759901in}}%
\pgfpathlineto{\pgfqpoint{3.726949in}{2.752804in}}%
\pgfpathlineto{\pgfqpoint{3.729468in}{2.775550in}}%
\pgfpathlineto{\pgfqpoint{3.731987in}{2.824278in}}%
\pgfpathlineto{\pgfqpoint{3.734506in}{2.811363in}}%
\pgfpathlineto{\pgfqpoint{3.737024in}{2.808074in}}%
\pgfpathlineto{\pgfqpoint{3.739543in}{2.818610in}}%
\pgfpathlineto{\pgfqpoint{3.742062in}{2.835962in}}%
\pgfpathlineto{\pgfqpoint{3.744581in}{2.816744in}}%
\pgfpathlineto{\pgfqpoint{3.747099in}{2.864660in}}%
\pgfpathlineto{\pgfqpoint{3.749618in}{2.850769in}}%
\pgfpathlineto{\pgfqpoint{3.752137in}{2.827098in}}%
\pgfpathlineto{\pgfqpoint{3.754656in}{2.800762in}}%
\pgfpathlineto{\pgfqpoint{3.757174in}{2.802800in}}%
\pgfpathlineto{\pgfqpoint{3.759693in}{2.795386in}}%
\pgfpathlineto{\pgfqpoint{3.762212in}{2.797543in}}%
\pgfpathlineto{\pgfqpoint{3.764731in}{2.803796in}}%
\pgfpathlineto{\pgfqpoint{3.767249in}{2.775161in}}%
\pgfpathlineto{\pgfqpoint{3.769768in}{2.760078in}}%
\pgfpathlineto{\pgfqpoint{3.772287in}{2.749996in}}%
\pgfpathlineto{\pgfqpoint{3.774806in}{2.783046in}}%
\pgfpathlineto{\pgfqpoint{3.777324in}{2.763937in}}%
\pgfpathlineto{\pgfqpoint{3.779843in}{2.745817in}}%
\pgfpathlineto{\pgfqpoint{3.782362in}{2.781469in}}%
\pgfpathlineto{\pgfqpoint{3.784881in}{2.754213in}}%
\pgfpathlineto{\pgfqpoint{3.787399in}{2.728221in}}%
\pgfpathlineto{\pgfqpoint{3.789918in}{2.721131in}}%
\pgfpathlineto{\pgfqpoint{3.792437in}{2.776993in}}%
\pgfpathlineto{\pgfqpoint{3.794956in}{2.776027in}}%
\pgfpathlineto{\pgfqpoint{3.797474in}{2.810781in}}%
\pgfpathlineto{\pgfqpoint{3.799993in}{2.787087in}}%
\pgfpathlineto{\pgfqpoint{3.802512in}{2.847211in}}%
\pgfpathlineto{\pgfqpoint{3.805031in}{2.838792in}}%
\pgfpathlineto{\pgfqpoint{3.807549in}{2.860461in}}%
\pgfpathlineto{\pgfqpoint{3.810068in}{2.857059in}}%
\pgfpathlineto{\pgfqpoint{3.812587in}{2.855255in}}%
\pgfpathlineto{\pgfqpoint{3.815106in}{2.847015in}}%
\pgfpathlineto{\pgfqpoint{3.817624in}{2.883011in}}%
\pgfpathlineto{\pgfqpoint{3.820143in}{2.903524in}}%
\pgfpathlineto{\pgfqpoint{3.822662in}{2.912158in}}%
\pgfpathlineto{\pgfqpoint{3.825181in}{2.919653in}}%
\pgfpathlineto{\pgfqpoint{3.827699in}{2.890721in}}%
\pgfpathlineto{\pgfqpoint{3.830218in}{2.865549in}}%
\pgfpathlineto{\pgfqpoint{3.832737in}{2.863251in}}%
\pgfpathlineto{\pgfqpoint{3.835256in}{2.871623in}}%
\pgfpathlineto{\pgfqpoint{3.837774in}{2.870071in}}%
\pgfpathlineto{\pgfqpoint{3.840293in}{2.874909in}}%
\pgfpathlineto{\pgfqpoint{3.842812in}{2.865883in}}%
\pgfpathlineto{\pgfqpoint{3.845331in}{2.879559in}}%
\pgfpathlineto{\pgfqpoint{3.847849in}{2.891852in}}%
\pgfpathlineto{\pgfqpoint{3.850368in}{2.883671in}}%
\pgfpathlineto{\pgfqpoint{3.852887in}{2.898526in}}%
\pgfpathlineto{\pgfqpoint{3.855406in}{2.906601in}}%
\pgfpathlineto{\pgfqpoint{3.857924in}{2.928815in}}%
\pgfpathlineto{\pgfqpoint{3.860443in}{2.945416in}}%
\pgfpathlineto{\pgfqpoint{3.862962in}{2.936066in}}%
\pgfpathlineto{\pgfqpoint{3.865481in}{2.929225in}}%
\pgfpathlineto{\pgfqpoint{3.867999in}{2.918026in}}%
\pgfpathlineto{\pgfqpoint{3.870518in}{2.933869in}}%
\pgfpathlineto{\pgfqpoint{3.873037in}{2.943501in}}%
\pgfpathlineto{\pgfqpoint{3.875556in}{2.984511in}}%
\pgfpathlineto{\pgfqpoint{3.878074in}{2.971637in}}%
\pgfpathlineto{\pgfqpoint{3.880593in}{2.989317in}}%
\pgfpathlineto{\pgfqpoint{3.883112in}{2.954321in}}%
\pgfpathlineto{\pgfqpoint{3.885631in}{2.931228in}}%
\pgfpathlineto{\pgfqpoint{3.888149in}{2.951581in}}%
\pgfpathlineto{\pgfqpoint{3.890668in}{2.970640in}}%
\pgfpathlineto{\pgfqpoint{3.893187in}{2.972036in}}%
\pgfpathlineto{\pgfqpoint{3.895706in}{2.969387in}}%
\pgfpathlineto{\pgfqpoint{3.898224in}{2.969515in}}%
\pgfpathlineto{\pgfqpoint{3.900743in}{2.946240in}}%
\pgfpathlineto{\pgfqpoint{3.903262in}{2.966550in}}%
\pgfpathlineto{\pgfqpoint{3.905781in}{2.994972in}}%
\pgfpathlineto{\pgfqpoint{3.908299in}{3.006093in}}%
\pgfpathlineto{\pgfqpoint{3.910818in}{3.026204in}}%
\pgfpathlineto{\pgfqpoint{3.913337in}{3.062906in}}%
\pgfpathlineto{\pgfqpoint{3.915856in}{3.036085in}}%
\pgfpathlineto{\pgfqpoint{3.918374in}{2.994924in}}%
\pgfpathlineto{\pgfqpoint{3.920893in}{3.039573in}}%
\pgfpathlineto{\pgfqpoint{3.923412in}{3.061196in}}%
\pgfpathlineto{\pgfqpoint{3.925931in}{3.057893in}}%
\pgfpathlineto{\pgfqpoint{3.928449in}{3.065094in}}%
\pgfpathlineto{\pgfqpoint{3.930968in}{3.056173in}}%
\pgfpathlineto{\pgfqpoint{3.933487in}{3.028503in}}%
\pgfpathlineto{\pgfqpoint{3.936006in}{3.043587in}}%
\pgfpathlineto{\pgfqpoint{3.938524in}{3.019394in}}%
\pgfpathlineto{\pgfqpoint{3.941043in}{3.034677in}}%
\pgfpathlineto{\pgfqpoint{3.943562in}{3.036727in}}%
\pgfpathlineto{\pgfqpoint{3.946081in}{3.026542in}}%
\pgfpathlineto{\pgfqpoint{3.948599in}{3.015550in}}%
\pgfpathlineto{\pgfqpoint{3.951118in}{3.047052in}}%
\pgfpathlineto{\pgfqpoint{3.953637in}{3.024732in}}%
\pgfpathlineto{\pgfqpoint{3.956156in}{2.997722in}}%
\pgfpathlineto{\pgfqpoint{3.958674in}{3.028122in}}%
\pgfpathlineto{\pgfqpoint{3.961193in}{3.034611in}}%
\pgfpathlineto{\pgfqpoint{3.963712in}{3.062847in}}%
\pgfpathlineto{\pgfqpoint{3.966231in}{3.043717in}}%
\pgfpathlineto{\pgfqpoint{3.968749in}{3.040632in}}%
\pgfpathlineto{\pgfqpoint{3.971268in}{3.012599in}}%
\pgfpathlineto{\pgfqpoint{3.973787in}{3.033945in}}%
\pgfpathlineto{\pgfqpoint{3.978824in}{3.129853in}}%
\pgfpathlineto{\pgfqpoint{3.981343in}{3.137931in}}%
\pgfpathlineto{\pgfqpoint{3.983862in}{3.153396in}}%
\pgfpathlineto{\pgfqpoint{3.986381in}{3.161875in}}%
\pgfpathlineto{\pgfqpoint{3.988899in}{3.180151in}}%
\pgfpathlineto{\pgfqpoint{3.991418in}{3.173173in}}%
\pgfpathlineto{\pgfqpoint{3.993937in}{3.178363in}}%
\pgfpathlineto{\pgfqpoint{3.996456in}{3.155508in}}%
\pgfpathlineto{\pgfqpoint{3.998974in}{3.168589in}}%
\pgfpathlineto{\pgfqpoint{4.001493in}{3.179607in}}%
\pgfpathlineto{\pgfqpoint{4.004012in}{3.208258in}}%
\pgfpathlineto{\pgfqpoint{4.006531in}{3.204104in}}%
\pgfpathlineto{\pgfqpoint{4.009049in}{3.225158in}}%
\pgfpathlineto{\pgfqpoint{4.011568in}{3.213062in}}%
\pgfpathlineto{\pgfqpoint{4.014087in}{3.206413in}}%
\pgfpathlineto{\pgfqpoint{4.016606in}{3.204148in}}%
\pgfpathlineto{\pgfqpoint{4.019124in}{3.199599in}}%
\pgfpathlineto{\pgfqpoint{4.021643in}{3.176940in}}%
\pgfpathlineto{\pgfqpoint{4.024162in}{3.132441in}}%
\pgfpathlineto{\pgfqpoint{4.026681in}{3.142260in}}%
\pgfpathlineto{\pgfqpoint{4.029199in}{3.147441in}}%
\pgfpathlineto{\pgfqpoint{4.031718in}{3.139131in}}%
\pgfpathlineto{\pgfqpoint{4.034237in}{3.168153in}}%
\pgfpathlineto{\pgfqpoint{4.036756in}{3.175478in}}%
\pgfpathlineto{\pgfqpoint{4.039274in}{3.189670in}}%
\pgfpathlineto{\pgfqpoint{4.041793in}{3.156390in}}%
\pgfpathlineto{\pgfqpoint{4.044312in}{3.177591in}}%
\pgfpathlineto{\pgfqpoint{4.046831in}{3.217494in}}%
\pgfpathlineto{\pgfqpoint{4.049349in}{3.196854in}}%
\pgfpathlineto{\pgfqpoint{4.051868in}{3.168564in}}%
\pgfpathlineto{\pgfqpoint{4.054387in}{3.176805in}}%
\pgfpathlineto{\pgfqpoint{4.056906in}{3.207180in}}%
\pgfpathlineto{\pgfqpoint{4.059424in}{3.208479in}}%
\pgfpathlineto{\pgfqpoint{4.061943in}{3.215179in}}%
\pgfpathlineto{\pgfqpoint{4.064462in}{3.227002in}}%
\pgfpathlineto{\pgfqpoint{4.066981in}{3.222508in}}%
\pgfpathlineto{\pgfqpoint{4.069499in}{3.222572in}}%
\pgfpathlineto{\pgfqpoint{4.072018in}{3.210375in}}%
\pgfpathlineto{\pgfqpoint{4.074537in}{3.222694in}}%
\pgfpathlineto{\pgfqpoint{4.077056in}{3.192542in}}%
\pgfpathlineto{\pgfqpoint{4.079574in}{3.201765in}}%
\pgfpathlineto{\pgfqpoint{4.082093in}{3.196832in}}%
\pgfpathlineto{\pgfqpoint{4.084612in}{3.206476in}}%
\pgfpathlineto{\pgfqpoint{4.087131in}{3.193257in}}%
\pgfpathlineto{\pgfqpoint{4.089649in}{3.228552in}}%
\pgfpathlineto{\pgfqpoint{4.092168in}{3.281174in}}%
\pgfpathlineto{\pgfqpoint{4.094687in}{3.287368in}}%
\pgfpathlineto{\pgfqpoint{4.097206in}{3.277933in}}%
\pgfpathlineto{\pgfqpoint{4.099724in}{3.296475in}}%
\pgfpathlineto{\pgfqpoint{4.102243in}{3.321250in}}%
\pgfpathlineto{\pgfqpoint{4.104762in}{3.326889in}}%
\pgfpathlineto{\pgfqpoint{4.107281in}{3.330059in}}%
\pgfpathlineto{\pgfqpoint{4.109799in}{3.363335in}}%
\pgfpathlineto{\pgfqpoint{4.112318in}{3.367598in}}%
\pgfpathlineto{\pgfqpoint{4.114837in}{3.350200in}}%
\pgfpathlineto{\pgfqpoint{4.117356in}{3.377859in}}%
\pgfpathlineto{\pgfqpoint{4.119874in}{3.391681in}}%
\pgfpathlineto{\pgfqpoint{4.122393in}{3.417100in}}%
\pgfpathlineto{\pgfqpoint{4.124912in}{3.433399in}}%
\pgfpathlineto{\pgfqpoint{4.127431in}{3.423777in}}%
\pgfpathlineto{\pgfqpoint{4.129949in}{3.445248in}}%
\pgfpathlineto{\pgfqpoint{4.132468in}{3.464741in}}%
\pgfpathlineto{\pgfqpoint{4.134987in}{3.483038in}}%
\pgfpathlineto{\pgfqpoint{4.137506in}{3.475151in}}%
\pgfpathlineto{\pgfqpoint{4.140024in}{3.481638in}}%
\pgfpathlineto{\pgfqpoint{4.142543in}{3.462813in}}%
\pgfpathlineto{\pgfqpoint{4.145062in}{3.486213in}}%
\pgfpathlineto{\pgfqpoint{4.147581in}{3.458166in}}%
\pgfpathlineto{\pgfqpoint{4.150099in}{3.452970in}}%
\pgfpathlineto{\pgfqpoint{4.152618in}{3.454373in}}%
\pgfpathlineto{\pgfqpoint{4.155137in}{3.465791in}}%
\pgfpathlineto{\pgfqpoint{4.157656in}{3.445587in}}%
\pgfpathlineto{\pgfqpoint{4.160174in}{3.417391in}}%
\pgfpathlineto{\pgfqpoint{4.162693in}{3.462258in}}%
\pgfpathlineto{\pgfqpoint{4.165212in}{3.429205in}}%
\pgfpathlineto{\pgfqpoint{4.167731in}{3.448304in}}%
\pgfpathlineto{\pgfqpoint{4.170249in}{3.477266in}}%
\pgfpathlineto{\pgfqpoint{4.172768in}{3.427844in}}%
\pgfpathlineto{\pgfqpoint{4.175287in}{3.430744in}}%
\pgfpathlineto{\pgfqpoint{4.177806in}{3.456681in}}%
\pgfpathlineto{\pgfqpoint{4.180324in}{3.422561in}}%
\pgfpathlineto{\pgfqpoint{4.182843in}{3.391705in}}%
\pgfpathlineto{\pgfqpoint{4.185362in}{3.392284in}}%
\pgfpathlineto{\pgfqpoint{4.187881in}{3.365081in}}%
\pgfpathlineto{\pgfqpoint{4.190399in}{3.372059in}}%
\pgfpathlineto{\pgfqpoint{4.192918in}{3.410682in}}%
\pgfpathlineto{\pgfqpoint{4.195437in}{3.385705in}}%
\pgfpathlineto{\pgfqpoint{4.197956in}{3.385086in}}%
\pgfpathlineto{\pgfqpoint{4.200474in}{3.399406in}}%
\pgfpathlineto{\pgfqpoint{4.202993in}{3.394291in}}%
\pgfpathlineto{\pgfqpoint{4.205512in}{3.407928in}}%
\pgfpathlineto{\pgfqpoint{4.208031in}{3.411647in}}%
\pgfpathlineto{\pgfqpoint{4.210549in}{3.386790in}}%
\pgfpathlineto{\pgfqpoint{4.213068in}{3.393921in}}%
\pgfpathlineto{\pgfqpoint{4.215587in}{3.359643in}}%
\pgfpathlineto{\pgfqpoint{4.218106in}{3.332805in}}%
\pgfpathlineto{\pgfqpoint{4.220624in}{3.333295in}}%
\pgfpathlineto{\pgfqpoint{4.225662in}{3.359270in}}%
\pgfpathlineto{\pgfqpoint{4.228181in}{3.352579in}}%
\pgfpathlineto{\pgfqpoint{4.230699in}{3.340289in}}%
\pgfpathlineto{\pgfqpoint{4.233218in}{3.376622in}}%
\pgfpathlineto{\pgfqpoint{4.235737in}{3.394347in}}%
\pgfpathlineto{\pgfqpoint{4.238256in}{3.386409in}}%
\pgfpathlineto{\pgfqpoint{4.240774in}{3.385910in}}%
\pgfpathlineto{\pgfqpoint{4.243293in}{3.400411in}}%
\pgfpathlineto{\pgfqpoint{4.245812in}{3.379042in}}%
\pgfpathlineto{\pgfqpoint{4.248331in}{3.303735in}}%
\pgfpathlineto{\pgfqpoint{4.250849in}{3.258598in}}%
\pgfpathlineto{\pgfqpoint{4.253368in}{3.306617in}}%
\pgfpathlineto{\pgfqpoint{4.255887in}{3.292780in}}%
\pgfpathlineto{\pgfqpoint{4.258406in}{3.298682in}}%
\pgfpathlineto{\pgfqpoint{4.260924in}{3.291818in}}%
\pgfpathlineto{\pgfqpoint{4.263443in}{3.280842in}}%
\pgfpathlineto{\pgfqpoint{4.265962in}{3.283012in}}%
\pgfpathlineto{\pgfqpoint{4.268481in}{3.278058in}}%
\pgfpathlineto{\pgfqpoint{4.270999in}{3.294271in}}%
\pgfpathlineto{\pgfqpoint{4.273518in}{3.285284in}}%
\pgfpathlineto{\pgfqpoint{4.276037in}{3.263867in}}%
\pgfpathlineto{\pgfqpoint{4.278556in}{3.259153in}}%
\pgfpathlineto{\pgfqpoint{4.281074in}{3.274489in}}%
\pgfpathlineto{\pgfqpoint{4.283593in}{3.250456in}}%
\pgfpathlineto{\pgfqpoint{4.286112in}{3.253708in}}%
\pgfpathlineto{\pgfqpoint{4.288631in}{3.252445in}}%
\pgfpathlineto{\pgfqpoint{4.291149in}{3.244947in}}%
\pgfpathlineto{\pgfqpoint{4.293668in}{3.203874in}}%
\pgfpathlineto{\pgfqpoint{4.296187in}{3.232913in}}%
\pgfpathlineto{\pgfqpoint{4.298706in}{3.233122in}}%
\pgfpathlineto{\pgfqpoint{4.301224in}{3.279224in}}%
\pgfpathlineto{\pgfqpoint{4.303743in}{3.245102in}}%
\pgfpathlineto{\pgfqpoint{4.306262in}{3.269704in}}%
\pgfpathlineto{\pgfqpoint{4.308781in}{3.313626in}}%
\pgfpathlineto{\pgfqpoint{4.311299in}{3.303781in}}%
\pgfpathlineto{\pgfqpoint{4.313818in}{3.301277in}}%
\pgfpathlineto{\pgfqpoint{4.316337in}{3.271210in}}%
\pgfpathlineto{\pgfqpoint{4.318856in}{3.259354in}}%
\pgfpathlineto{\pgfqpoint{4.321374in}{3.269451in}}%
\pgfpathlineto{\pgfqpoint{4.323893in}{3.252548in}}%
\pgfpathlineto{\pgfqpoint{4.326412in}{3.279294in}}%
\pgfpathlineto{\pgfqpoint{4.328931in}{3.321794in}}%
\pgfpathlineto{\pgfqpoint{4.331449in}{3.319895in}}%
\pgfpathlineto{\pgfqpoint{4.333968in}{3.308466in}}%
\pgfpathlineto{\pgfqpoint{4.336487in}{3.322546in}}%
\pgfpathlineto{\pgfqpoint{4.339006in}{3.313101in}}%
\pgfpathlineto{\pgfqpoint{4.341524in}{3.323174in}}%
\pgfpathlineto{\pgfqpoint{4.344043in}{3.342432in}}%
\pgfpathlineto{\pgfqpoint{4.346562in}{3.369044in}}%
\pgfpathlineto{\pgfqpoint{4.349081in}{3.348436in}}%
\pgfpathlineto{\pgfqpoint{4.351599in}{3.359932in}}%
\pgfpathlineto{\pgfqpoint{4.354118in}{3.370623in}}%
\pgfpathlineto{\pgfqpoint{4.356637in}{3.411215in}}%
\pgfpathlineto{\pgfqpoint{4.359156in}{3.472024in}}%
\pgfpathlineto{\pgfqpoint{4.361674in}{3.487972in}}%
\pgfpathlineto{\pgfqpoint{4.364193in}{3.530855in}}%
\pgfpathlineto{\pgfqpoint{4.366712in}{3.580648in}}%
\pgfpathlineto{\pgfqpoint{4.369231in}{3.583893in}}%
\pgfpathlineto{\pgfqpoint{4.371749in}{3.581051in}}%
\pgfpathlineto{\pgfqpoint{4.374268in}{3.567424in}}%
\pgfpathlineto{\pgfqpoint{4.376787in}{3.542518in}}%
\pgfpathlineto{\pgfqpoint{4.379306in}{3.569524in}}%
\pgfpathlineto{\pgfqpoint{4.381824in}{3.555257in}}%
\pgfpathlineto{\pgfqpoint{4.384343in}{3.580421in}}%
\pgfpathlineto{\pgfqpoint{4.386862in}{3.581611in}}%
\pgfpathlineto{\pgfqpoint{4.389381in}{3.555006in}}%
\pgfpathlineto{\pgfqpoint{4.391899in}{3.590686in}}%
\pgfpathlineto{\pgfqpoint{4.394418in}{3.574920in}}%
\pgfpathlineto{\pgfqpoint{4.396937in}{3.563499in}}%
\pgfpathlineto{\pgfqpoint{4.399456in}{3.571685in}}%
\pgfpathlineto{\pgfqpoint{4.401974in}{3.577033in}}%
\pgfpathlineto{\pgfqpoint{4.404493in}{3.613108in}}%
\pgfpathlineto{\pgfqpoint{4.407012in}{3.686393in}}%
\pgfpathlineto{\pgfqpoint{4.409531in}{3.650865in}}%
\pgfpathlineto{\pgfqpoint{4.412049in}{3.613547in}}%
\pgfpathlineto{\pgfqpoint{4.414568in}{3.607313in}}%
\pgfpathlineto{\pgfqpoint{4.417087in}{3.659020in}}%
\pgfpathlineto{\pgfqpoint{4.419606in}{3.623106in}}%
\pgfpathlineto{\pgfqpoint{4.422124in}{3.572899in}}%
\pgfpathlineto{\pgfqpoint{4.424643in}{3.542438in}}%
\pgfpathlineto{\pgfqpoint{4.427162in}{3.530991in}}%
\pgfpathlineto{\pgfqpoint{4.429681in}{3.552707in}}%
\pgfpathlineto{\pgfqpoint{4.432199in}{3.619770in}}%
\pgfpathlineto{\pgfqpoint{4.434718in}{3.620808in}}%
\pgfpathlineto{\pgfqpoint{4.437237in}{3.561149in}}%
\pgfpathlineto{\pgfqpoint{4.439756in}{3.579139in}}%
\pgfpathlineto{\pgfqpoint{4.442274in}{3.579664in}}%
\pgfpathlineto{\pgfqpoint{4.444793in}{3.596289in}}%
\pgfpathlineto{\pgfqpoint{4.447312in}{3.600173in}}%
\pgfpathlineto{\pgfqpoint{4.449831in}{3.611985in}}%
\pgfpathlineto{\pgfqpoint{4.452349in}{3.619906in}}%
\pgfpathlineto{\pgfqpoint{4.454868in}{3.554230in}}%
\pgfpathlineto{\pgfqpoint{4.457387in}{3.530756in}}%
\pgfpathlineto{\pgfqpoint{4.459906in}{3.538006in}}%
\pgfpathlineto{\pgfqpoint{4.462424in}{3.525740in}}%
\pgfpathlineto{\pgfqpoint{4.464943in}{3.507850in}}%
\pgfpathlineto{\pgfqpoint{4.467462in}{3.470796in}}%
\pgfpathlineto{\pgfqpoint{4.469981in}{3.497291in}}%
\pgfpathlineto{\pgfqpoint{4.475018in}{3.513945in}}%
\pgfpathlineto{\pgfqpoint{4.477537in}{3.494219in}}%
\pgfpathlineto{\pgfqpoint{4.480056in}{3.538064in}}%
\pgfpathlineto{\pgfqpoint{4.482574in}{3.557182in}}%
\pgfpathlineto{\pgfqpoint{4.485093in}{3.625666in}}%
\pgfpathlineto{\pgfqpoint{4.487612in}{3.596914in}}%
\pgfpathlineto{\pgfqpoint{4.490131in}{3.593618in}}%
\pgfpathlineto{\pgfqpoint{4.492649in}{3.610067in}}%
\pgfpathlineto{\pgfqpoint{4.495168in}{3.606608in}}%
\pgfpathlineto{\pgfqpoint{4.497687in}{3.617970in}}%
\pgfpathlineto{\pgfqpoint{4.500206in}{3.623672in}}%
\pgfpathlineto{\pgfqpoint{4.502724in}{3.621209in}}%
\pgfpathlineto{\pgfqpoint{4.505243in}{3.611012in}}%
\pgfpathlineto{\pgfqpoint{4.507762in}{3.590366in}}%
\pgfpathlineto{\pgfqpoint{4.510281in}{3.577285in}}%
\pgfpathlineto{\pgfqpoint{4.512799in}{3.516078in}}%
\pgfpathlineto{\pgfqpoint{4.515318in}{3.519081in}}%
\pgfpathlineto{\pgfqpoint{4.517837in}{3.499731in}}%
\pgfpathlineto{\pgfqpoint{4.522874in}{3.457160in}}%
\pgfpathlineto{\pgfqpoint{4.525393in}{3.479104in}}%
\pgfpathlineto{\pgfqpoint{4.527912in}{3.465025in}}%
\pgfpathlineto{\pgfqpoint{4.530431in}{3.468233in}}%
\pgfpathlineto{\pgfqpoint{4.532949in}{3.461865in}}%
\pgfpathlineto{\pgfqpoint{4.535468in}{3.425614in}}%
\pgfpathlineto{\pgfqpoint{4.537987in}{3.435510in}}%
\pgfpathlineto{\pgfqpoint{4.540506in}{3.405528in}}%
\pgfpathlineto{\pgfqpoint{4.543024in}{3.406329in}}%
\pgfpathlineto{\pgfqpoint{4.545543in}{3.394173in}}%
\pgfpathlineto{\pgfqpoint{4.548062in}{3.408877in}}%
\pgfpathlineto{\pgfqpoint{4.550581in}{3.435454in}}%
\pgfpathlineto{\pgfqpoint{4.553099in}{3.433543in}}%
\pgfpathlineto{\pgfqpoint{4.555618in}{3.459800in}}%
\pgfpathlineto{\pgfqpoint{4.558137in}{3.419699in}}%
\pgfpathlineto{\pgfqpoint{4.560656in}{3.358344in}}%
\pgfpathlineto{\pgfqpoint{4.563174in}{3.372497in}}%
\pgfpathlineto{\pgfqpoint{4.565693in}{3.370998in}}%
\pgfpathlineto{\pgfqpoint{4.568212in}{3.389478in}}%
\pgfpathlineto{\pgfqpoint{4.570731in}{3.383731in}}%
\pgfpathlineto{\pgfqpoint{4.575768in}{3.435588in}}%
\pgfpathlineto{\pgfqpoint{4.578287in}{3.445159in}}%
\pgfpathlineto{\pgfqpoint{4.580806in}{3.432264in}}%
\pgfpathlineto{\pgfqpoint{4.583324in}{3.454141in}}%
\pgfpathlineto{\pgfqpoint{4.585843in}{3.440775in}}%
\pgfpathlineto{\pgfqpoint{4.588362in}{3.458147in}}%
\pgfpathlineto{\pgfqpoint{4.590881in}{3.469115in}}%
\pgfpathlineto{\pgfqpoint{4.593399in}{3.441105in}}%
\pgfpathlineto{\pgfqpoint{4.595918in}{3.452193in}}%
\pgfpathlineto{\pgfqpoint{4.598437in}{3.436904in}}%
\pgfpathlineto{\pgfqpoint{4.600956in}{3.440934in}}%
\pgfpathlineto{\pgfqpoint{4.603474in}{3.422429in}}%
\pgfpathlineto{\pgfqpoint{4.605993in}{3.417304in}}%
\pgfpathlineto{\pgfqpoint{4.608512in}{3.395030in}}%
\pgfpathlineto{\pgfqpoint{4.611031in}{3.399175in}}%
\pgfpathlineto{\pgfqpoint{4.613549in}{3.393397in}}%
\pgfpathlineto{\pgfqpoint{4.616068in}{3.374872in}}%
\pgfpathlineto{\pgfqpoint{4.618587in}{3.400283in}}%
\pgfpathlineto{\pgfqpoint{4.621106in}{3.424429in}}%
\pgfpathlineto{\pgfqpoint{4.623624in}{3.422695in}}%
\pgfpathlineto{\pgfqpoint{4.626143in}{3.461825in}}%
\pgfpathlineto{\pgfqpoint{4.628662in}{3.435673in}}%
\pgfpathlineto{\pgfqpoint{4.631181in}{3.441918in}}%
\pgfpathlineto{\pgfqpoint{4.633699in}{3.418682in}}%
\pgfpathlineto{\pgfqpoint{4.636218in}{3.445116in}}%
\pgfpathlineto{\pgfqpoint{4.638737in}{3.464303in}}%
\pgfpathlineto{\pgfqpoint{4.641256in}{3.464904in}}%
\pgfpathlineto{\pgfqpoint{4.643774in}{3.448726in}}%
\pgfpathlineto{\pgfqpoint{4.646293in}{3.435132in}}%
\pgfpathlineto{\pgfqpoint{4.648812in}{3.424842in}}%
\pgfpathlineto{\pgfqpoint{4.651331in}{3.440734in}}%
\pgfpathlineto{\pgfqpoint{4.653849in}{3.450495in}}%
\pgfpathlineto{\pgfqpoint{4.656368in}{3.462015in}}%
\pgfpathlineto{\pgfqpoint{4.658887in}{3.448301in}}%
\pgfpathlineto{\pgfqpoint{4.661406in}{3.399272in}}%
\pgfpathlineto{\pgfqpoint{4.663924in}{3.415588in}}%
\pgfpathlineto{\pgfqpoint{4.666443in}{3.394448in}}%
\pgfpathlineto{\pgfqpoint{4.668962in}{3.367743in}}%
\pgfpathlineto{\pgfqpoint{4.671481in}{3.377600in}}%
\pgfpathlineto{\pgfqpoint{4.673999in}{3.345278in}}%
\pgfpathlineto{\pgfqpoint{4.676518in}{3.306999in}}%
\pgfpathlineto{\pgfqpoint{4.679037in}{3.306114in}}%
\pgfpathlineto{\pgfqpoint{4.681556in}{3.323662in}}%
\pgfpathlineto{\pgfqpoint{4.684074in}{3.292845in}}%
\pgfpathlineto{\pgfqpoint{4.686593in}{3.284089in}}%
\pgfpathlineto{\pgfqpoint{4.689112in}{3.272959in}}%
\pgfpathlineto{\pgfqpoint{4.691631in}{3.279582in}}%
\pgfpathlineto{\pgfqpoint{4.694149in}{3.313899in}}%
\pgfpathlineto{\pgfqpoint{4.696668in}{3.329991in}}%
\pgfpathlineto{\pgfqpoint{4.699187in}{3.359369in}}%
\pgfpathlineto{\pgfqpoint{4.701706in}{3.359342in}}%
\pgfpathlineto{\pgfqpoint{4.704224in}{3.360784in}}%
\pgfpathlineto{\pgfqpoint{4.706743in}{3.354883in}}%
\pgfpathlineto{\pgfqpoint{4.709262in}{3.326052in}}%
\pgfpathlineto{\pgfqpoint{4.711781in}{3.353902in}}%
\pgfpathlineto{\pgfqpoint{4.714299in}{3.382968in}}%
\pgfpathlineto{\pgfqpoint{4.716818in}{3.373220in}}%
\pgfpathlineto{\pgfqpoint{4.719337in}{3.365382in}}%
\pgfpathlineto{\pgfqpoint{4.721856in}{3.377533in}}%
\pgfpathlineto{\pgfqpoint{4.724374in}{3.400518in}}%
\pgfpathlineto{\pgfqpoint{4.726893in}{3.392229in}}%
\pgfpathlineto{\pgfqpoint{4.729412in}{3.414093in}}%
\pgfpathlineto{\pgfqpoint{4.731931in}{3.407374in}}%
\pgfpathlineto{\pgfqpoint{4.734449in}{3.394630in}}%
\pgfpathlineto{\pgfqpoint{4.736968in}{3.363389in}}%
\pgfpathlineto{\pgfqpoint{4.739487in}{3.363682in}}%
\pgfpathlineto{\pgfqpoint{4.742006in}{3.364942in}}%
\pgfpathlineto{\pgfqpoint{4.744524in}{3.361755in}}%
\pgfpathlineto{\pgfqpoint{4.747043in}{3.348017in}}%
\pgfpathlineto{\pgfqpoint{4.749562in}{3.329409in}}%
\pgfpathlineto{\pgfqpoint{4.752081in}{3.339552in}}%
\pgfpathlineto{\pgfqpoint{4.754599in}{3.364055in}}%
\pgfpathlineto{\pgfqpoint{4.757118in}{3.360981in}}%
\pgfpathlineto{\pgfqpoint{4.759637in}{3.342555in}}%
\pgfpathlineto{\pgfqpoint{4.762156in}{3.339118in}}%
\pgfpathlineto{\pgfqpoint{4.764674in}{3.348408in}}%
\pgfpathlineto{\pgfqpoint{4.767193in}{3.355383in}}%
\pgfpathlineto{\pgfqpoint{4.769712in}{3.366890in}}%
\pgfpathlineto{\pgfqpoint{4.772231in}{3.381401in}}%
\pgfpathlineto{\pgfqpoint{4.774749in}{3.363569in}}%
\pgfpathlineto{\pgfqpoint{4.777268in}{3.363058in}}%
\pgfpathlineto{\pgfqpoint{4.779787in}{3.334728in}}%
\pgfpathlineto{\pgfqpoint{4.782306in}{3.336307in}}%
\pgfpathlineto{\pgfqpoint{4.784824in}{3.345432in}}%
\pgfpathlineto{\pgfqpoint{4.787343in}{3.327817in}}%
\pgfpathlineto{\pgfqpoint{4.789862in}{3.320038in}}%
\pgfpathlineto{\pgfqpoint{4.792381in}{3.346166in}}%
\pgfpathlineto{\pgfqpoint{4.794899in}{3.330650in}}%
\pgfpathlineto{\pgfqpoint{4.797418in}{3.336071in}}%
\pgfpathlineto{\pgfqpoint{4.799937in}{3.372678in}}%
\pgfpathlineto{\pgfqpoint{4.802456in}{3.395266in}}%
\pgfpathlineto{\pgfqpoint{4.804974in}{3.424936in}}%
\pgfpathlineto{\pgfqpoint{4.807493in}{3.442772in}}%
\pgfpathlineto{\pgfqpoint{4.810012in}{3.495586in}}%
\pgfpathlineto{\pgfqpoint{4.812531in}{3.493978in}}%
\pgfpathlineto{\pgfqpoint{4.815049in}{3.493088in}}%
\pgfpathlineto{\pgfqpoint{4.817568in}{3.480110in}}%
\pgfpathlineto{\pgfqpoint{4.820087in}{3.488409in}}%
\pgfpathlineto{\pgfqpoint{4.822606in}{3.504456in}}%
\pgfpathlineto{\pgfqpoint{4.825124in}{3.513439in}}%
\pgfpathlineto{\pgfqpoint{4.827643in}{3.489001in}}%
\pgfpathlineto{\pgfqpoint{4.830162in}{3.480477in}}%
\pgfpathlineto{\pgfqpoint{4.832681in}{3.521059in}}%
\pgfpathlineto{\pgfqpoint{4.835199in}{3.521435in}}%
\pgfpathlineto{\pgfqpoint{4.837718in}{3.503079in}}%
\pgfpathlineto{\pgfqpoint{4.840237in}{3.491922in}}%
\pgfpathlineto{\pgfqpoint{4.842756in}{3.491539in}}%
\pgfpathlineto{\pgfqpoint{4.845274in}{3.477482in}}%
\pgfpathlineto{\pgfqpoint{4.847793in}{3.457851in}}%
\pgfpathlineto{\pgfqpoint{4.850312in}{3.466115in}}%
\pgfpathlineto{\pgfqpoint{4.852831in}{3.454485in}}%
\pgfpathlineto{\pgfqpoint{4.855349in}{3.424860in}}%
\pgfpathlineto{\pgfqpoint{4.857868in}{3.415018in}}%
\pgfpathlineto{\pgfqpoint{4.860387in}{3.417877in}}%
\pgfpathlineto{\pgfqpoint{4.862906in}{3.403689in}}%
\pgfpathlineto{\pgfqpoint{4.865424in}{3.408911in}}%
\pgfpathlineto{\pgfqpoint{4.867943in}{3.412110in}}%
\pgfpathlineto{\pgfqpoint{4.870462in}{3.427172in}}%
\pgfpathlineto{\pgfqpoint{4.872981in}{3.448090in}}%
\pgfpathlineto{\pgfqpoint{4.875499in}{3.453870in}}%
\pgfpathlineto{\pgfqpoint{4.878018in}{3.475231in}}%
\pgfpathlineto{\pgfqpoint{4.880537in}{3.491846in}}%
\pgfpathlineto{\pgfqpoint{4.883056in}{3.482309in}}%
\pgfpathlineto{\pgfqpoint{4.885574in}{3.503981in}}%
\pgfpathlineto{\pgfqpoint{4.888093in}{3.469180in}}%
\pgfpathlineto{\pgfqpoint{4.890612in}{3.435731in}}%
\pgfpathlineto{\pgfqpoint{4.893131in}{3.432860in}}%
\pgfpathlineto{\pgfqpoint{4.895649in}{3.421424in}}%
\pgfpathlineto{\pgfqpoint{4.898168in}{3.400859in}}%
\pgfpathlineto{\pgfqpoint{4.900687in}{3.375540in}}%
\pgfpathlineto{\pgfqpoint{4.903206in}{3.388333in}}%
\pgfpathlineto{\pgfqpoint{4.905724in}{3.418185in}}%
\pgfpathlineto{\pgfqpoint{4.908243in}{3.434394in}}%
\pgfpathlineto{\pgfqpoint{4.910762in}{3.418222in}}%
\pgfpathlineto{\pgfqpoint{4.913281in}{3.389655in}}%
\pgfpathlineto{\pgfqpoint{4.915799in}{3.402616in}}%
\pgfpathlineto{\pgfqpoint{4.918318in}{3.443224in}}%
\pgfpathlineto{\pgfqpoint{4.920837in}{3.457060in}}%
\pgfpathlineto{\pgfqpoint{4.923356in}{3.491037in}}%
\pgfpathlineto{\pgfqpoint{4.925874in}{3.475452in}}%
\pgfpathlineto{\pgfqpoint{4.928393in}{3.521280in}}%
\pgfpathlineto{\pgfqpoint{4.930912in}{3.515154in}}%
\pgfpathlineto{\pgfqpoint{4.933431in}{3.544358in}}%
\pgfpathlineto{\pgfqpoint{4.935949in}{3.555083in}}%
\pgfpathlineto{\pgfqpoint{4.938468in}{3.548190in}}%
\pgfpathlineto{\pgfqpoint{4.943506in}{3.587173in}}%
\pgfpathlineto{\pgfqpoint{4.946024in}{3.598694in}}%
\pgfpathlineto{\pgfqpoint{4.948543in}{3.571157in}}%
\pgfpathlineto{\pgfqpoint{4.951062in}{3.589378in}}%
\pgfpathlineto{\pgfqpoint{4.953581in}{3.610895in}}%
\pgfpathlineto{\pgfqpoint{4.956099in}{3.612975in}}%
\pgfpathlineto{\pgfqpoint{4.958618in}{3.592542in}}%
\pgfpathlineto{\pgfqpoint{4.961137in}{3.585812in}}%
\pgfpathlineto{\pgfqpoint{4.963656in}{3.590936in}}%
\pgfpathlineto{\pgfqpoint{4.966174in}{3.599755in}}%
\pgfpathlineto{\pgfqpoint{4.968693in}{3.612292in}}%
\pgfpathlineto{\pgfqpoint{4.971212in}{3.608976in}}%
\pgfpathlineto{\pgfqpoint{4.973731in}{3.679655in}}%
\pgfpathlineto{\pgfqpoint{4.976249in}{3.608625in}}%
\pgfpathlineto{\pgfqpoint{4.978768in}{3.574542in}}%
\pgfpathlineto{\pgfqpoint{4.981287in}{3.556349in}}%
\pgfpathlineto{\pgfqpoint{4.983806in}{3.566057in}}%
\pgfpathlineto{\pgfqpoint{4.986324in}{3.629070in}}%
\pgfpathlineto{\pgfqpoint{4.988843in}{3.644790in}}%
\pgfpathlineto{\pgfqpoint{4.991362in}{3.619482in}}%
\pgfpathlineto{\pgfqpoint{4.993881in}{3.624317in}}%
\pgfpathlineto{\pgfqpoint{4.996399in}{3.595407in}}%
\pgfpathlineto{\pgfqpoint{4.998918in}{3.579691in}}%
\pgfpathlineto{\pgfqpoint{5.001437in}{3.551876in}}%
\pgfpathlineto{\pgfqpoint{5.003956in}{3.496690in}}%
\pgfpathlineto{\pgfqpoint{5.006474in}{3.507498in}}%
\pgfpathlineto{\pgfqpoint{5.008993in}{3.505734in}}%
\pgfpathlineto{\pgfqpoint{5.011512in}{3.526007in}}%
\pgfpathlineto{\pgfqpoint{5.014031in}{3.483093in}}%
\pgfpathlineto{\pgfqpoint{5.016549in}{3.471480in}}%
\pgfpathlineto{\pgfqpoint{5.019068in}{3.482066in}}%
\pgfpathlineto{\pgfqpoint{5.024106in}{3.505998in}}%
\pgfpathlineto{\pgfqpoint{5.026624in}{3.496115in}}%
\pgfpathlineto{\pgfqpoint{5.029143in}{3.482085in}}%
\pgfpathlineto{\pgfqpoint{5.031662in}{3.478237in}}%
\pgfpathlineto{\pgfqpoint{5.034181in}{3.458331in}}%
\pgfpathlineto{\pgfqpoint{5.036699in}{3.437337in}}%
\pgfpathlineto{\pgfqpoint{5.039218in}{3.454396in}}%
\pgfpathlineto{\pgfqpoint{5.041737in}{3.467340in}}%
\pgfpathlineto{\pgfqpoint{5.044256in}{3.482553in}}%
\pgfpathlineto{\pgfqpoint{5.046774in}{3.483910in}}%
\pgfpathlineto{\pgfqpoint{5.049293in}{3.490590in}}%
\pgfpathlineto{\pgfqpoint{5.051812in}{3.460967in}}%
\pgfpathlineto{\pgfqpoint{5.054331in}{3.423700in}}%
\pgfpathlineto{\pgfqpoint{5.056849in}{3.388031in}}%
\pgfpathlineto{\pgfqpoint{5.059368in}{3.385999in}}%
\pgfpathlineto{\pgfqpoint{5.061887in}{3.345596in}}%
\pgfpathlineto{\pgfqpoint{5.064406in}{3.326205in}}%
\pgfpathlineto{\pgfqpoint{5.066924in}{3.370347in}}%
\pgfpathlineto{\pgfqpoint{5.069443in}{3.389625in}}%
\pgfpathlineto{\pgfqpoint{5.071962in}{3.398505in}}%
\pgfpathlineto{\pgfqpoint{5.074481in}{3.375709in}}%
\pgfpathlineto{\pgfqpoint{5.076999in}{3.376077in}}%
\pgfpathlineto{\pgfqpoint{5.079518in}{3.350548in}}%
\pgfpathlineto{\pgfqpoint{5.082037in}{3.352823in}}%
\pgfpathlineto{\pgfqpoint{5.084556in}{3.355832in}}%
\pgfpathlineto{\pgfqpoint{5.087074in}{3.359594in}}%
\pgfpathlineto{\pgfqpoint{5.089593in}{3.372210in}}%
\pgfpathlineto{\pgfqpoint{5.092112in}{3.357288in}}%
\pgfpathlineto{\pgfqpoint{5.094631in}{3.389878in}}%
\pgfpathlineto{\pgfqpoint{5.097149in}{3.384351in}}%
\pgfpathlineto{\pgfqpoint{5.099668in}{3.375187in}}%
\pgfpathlineto{\pgfqpoint{5.102187in}{3.335986in}}%
\pgfpathlineto{\pgfqpoint{5.104706in}{3.312441in}}%
\pgfpathlineto{\pgfqpoint{5.107224in}{3.268222in}}%
\pgfpathlineto{\pgfqpoint{5.109743in}{3.257749in}}%
\pgfpathlineto{\pgfqpoint{5.112262in}{3.237786in}}%
\pgfpathlineto{\pgfqpoint{5.114781in}{3.213797in}}%
\pgfpathlineto{\pgfqpoint{5.117299in}{3.230039in}}%
\pgfpathlineto{\pgfqpoint{5.119818in}{3.190773in}}%
\pgfpathlineto{\pgfqpoint{5.122337in}{3.229771in}}%
\pgfpathlineto{\pgfqpoint{5.124856in}{3.257501in}}%
\pgfpathlineto{\pgfqpoint{5.127374in}{3.299181in}}%
\pgfpathlineto{\pgfqpoint{5.129893in}{3.283028in}}%
\pgfpathlineto{\pgfqpoint{5.132412in}{3.313728in}}%
\pgfpathlineto{\pgfqpoint{5.134931in}{3.338503in}}%
\pgfpathlineto{\pgfqpoint{5.137449in}{3.346174in}}%
\pgfpathlineto{\pgfqpoint{5.139968in}{3.338115in}}%
\pgfpathlineto{\pgfqpoint{5.142487in}{3.342041in}}%
\pgfpathlineto{\pgfqpoint{5.145006in}{3.301183in}}%
\pgfpathlineto{\pgfqpoint{5.147524in}{3.310678in}}%
\pgfpathlineto{\pgfqpoint{5.150043in}{3.285136in}}%
\pgfpathlineto{\pgfqpoint{5.152562in}{3.312673in}}%
\pgfpathlineto{\pgfqpoint{5.155081in}{3.296076in}}%
\pgfpathlineto{\pgfqpoint{5.157599in}{3.322240in}}%
\pgfpathlineto{\pgfqpoint{5.160118in}{3.333057in}}%
\pgfpathlineto{\pgfqpoint{5.162637in}{3.308932in}}%
\pgfpathlineto{\pgfqpoint{5.165156in}{3.309633in}}%
\pgfpathlineto{\pgfqpoint{5.167674in}{3.331896in}}%
\pgfpathlineto{\pgfqpoint{5.170193in}{3.319706in}}%
\pgfpathlineto{\pgfqpoint{5.172712in}{3.364524in}}%
\pgfpathlineto{\pgfqpoint{5.175231in}{3.334420in}}%
\pgfpathlineto{\pgfqpoint{5.177749in}{3.395077in}}%
\pgfpathlineto{\pgfqpoint{5.180268in}{3.411518in}}%
\pgfpathlineto{\pgfqpoint{5.182787in}{3.415369in}}%
\pgfpathlineto{\pgfqpoint{5.185306in}{3.400053in}}%
\pgfpathlineto{\pgfqpoint{5.187824in}{3.402047in}}%
\pgfpathlineto{\pgfqpoint{5.190343in}{3.400194in}}%
\pgfpathlineto{\pgfqpoint{5.192862in}{3.392154in}}%
\pgfpathlineto{\pgfqpoint{5.195381in}{3.378980in}}%
\pgfpathlineto{\pgfqpoint{5.197899in}{3.399478in}}%
\pgfpathlineto{\pgfqpoint{5.200418in}{3.413000in}}%
\pgfpathlineto{\pgfqpoint{5.202937in}{3.423328in}}%
\pgfpathlineto{\pgfqpoint{5.205456in}{3.438665in}}%
\pgfpathlineto{\pgfqpoint{5.207974in}{3.436050in}}%
\pgfpathlineto{\pgfqpoint{5.210493in}{3.415350in}}%
\pgfpathlineto{\pgfqpoint{5.213012in}{3.403823in}}%
\pgfpathlineto{\pgfqpoint{5.215531in}{3.373317in}}%
\pgfpathlineto{\pgfqpoint{5.218049in}{3.390125in}}%
\pgfpathlineto{\pgfqpoint{5.220568in}{3.386018in}}%
\pgfpathlineto{\pgfqpoint{5.223087in}{3.385567in}}%
\pgfpathlineto{\pgfqpoint{5.225606in}{3.373492in}}%
\pgfpathlineto{\pgfqpoint{5.228124in}{3.404352in}}%
\pgfpathlineto{\pgfqpoint{5.230643in}{3.407324in}}%
\pgfpathlineto{\pgfqpoint{5.233162in}{3.402972in}}%
\pgfpathlineto{\pgfqpoint{5.235681in}{3.386377in}}%
\pgfpathlineto{\pgfqpoint{5.238199in}{3.402492in}}%
\pgfpathlineto{\pgfqpoint{5.240718in}{3.437766in}}%
\pgfpathlineto{\pgfqpoint{5.243237in}{3.411104in}}%
\pgfpathlineto{\pgfqpoint{5.248274in}{3.475506in}}%
\pgfpathlineto{\pgfqpoint{5.250793in}{3.456352in}}%
\pgfpathlineto{\pgfqpoint{5.253312in}{3.510014in}}%
\pgfpathlineto{\pgfqpoint{5.255831in}{3.513195in}}%
\pgfpathlineto{\pgfqpoint{5.258349in}{3.542208in}}%
\pgfpathlineto{\pgfqpoint{5.260868in}{3.527063in}}%
\pgfpathlineto{\pgfqpoint{5.263387in}{3.567838in}}%
\pgfpathlineto{\pgfqpoint{5.265906in}{3.523167in}}%
\pgfpathlineto{\pgfqpoint{5.268424in}{3.488360in}}%
\pgfpathlineto{\pgfqpoint{5.270943in}{3.528737in}}%
\pgfpathlineto{\pgfqpoint{5.273462in}{3.553731in}}%
\pgfpathlineto{\pgfqpoint{5.275981in}{3.566901in}}%
\pgfpathlineto{\pgfqpoint{5.278499in}{3.557807in}}%
\pgfpathlineto{\pgfqpoint{5.281018in}{3.607449in}}%
\pgfpathlineto{\pgfqpoint{5.283537in}{3.593625in}}%
\pgfpathlineto{\pgfqpoint{5.286056in}{3.614991in}}%
\pgfpathlineto{\pgfqpoint{5.288574in}{3.610145in}}%
\pgfpathlineto{\pgfqpoint{5.291093in}{3.615117in}}%
\pgfpathlineto{\pgfqpoint{5.293612in}{3.637166in}}%
\pgfpathlineto{\pgfqpoint{5.296131in}{3.647331in}}%
\pgfpathlineto{\pgfqpoint{5.298649in}{3.672668in}}%
\pgfpathlineto{\pgfqpoint{5.301168in}{3.641251in}}%
\pgfpathlineto{\pgfqpoint{5.303687in}{3.635400in}}%
\pgfpathlineto{\pgfqpoint{5.306206in}{3.659616in}}%
\pgfpathlineto{\pgfqpoint{5.308724in}{3.645198in}}%
\pgfpathlineto{\pgfqpoint{5.311243in}{3.662451in}}%
\pgfpathlineto{\pgfqpoint{5.316281in}{3.619963in}}%
\pgfpathlineto{\pgfqpoint{5.318799in}{3.621034in}}%
\pgfpathlineto{\pgfqpoint{5.321318in}{3.597219in}}%
\pgfpathlineto{\pgfqpoint{5.323837in}{3.620792in}}%
\pgfpathlineto{\pgfqpoint{5.326356in}{3.650698in}}%
\pgfpathlineto{\pgfqpoint{5.328874in}{3.657940in}}%
\pgfpathlineto{\pgfqpoint{5.331393in}{3.647839in}}%
\pgfpathlineto{\pgfqpoint{5.333912in}{3.622355in}}%
\pgfpathlineto{\pgfqpoint{5.336431in}{3.629209in}}%
\pgfpathlineto{\pgfqpoint{5.338949in}{3.661185in}}%
\pgfpathlineto{\pgfqpoint{5.341468in}{3.642628in}}%
\pgfpathlineto{\pgfqpoint{5.343987in}{3.612925in}}%
\pgfpathlineto{\pgfqpoint{5.346506in}{3.618896in}}%
\pgfpathlineto{\pgfqpoint{5.349024in}{3.609900in}}%
\pgfpathlineto{\pgfqpoint{5.351543in}{3.653109in}}%
\pgfpathlineto{\pgfqpoint{5.354062in}{3.662973in}}%
\pgfpathlineto{\pgfqpoint{5.356581in}{3.679118in}}%
\pgfpathlineto{\pgfqpoint{5.359099in}{3.670158in}}%
\pgfpathlineto{\pgfqpoint{5.361618in}{3.658191in}}%
\pgfpathlineto{\pgfqpoint{5.364137in}{3.645491in}}%
\pgfpathlineto{\pgfqpoint{5.366656in}{3.738872in}}%
\pgfpathlineto{\pgfqpoint{5.369174in}{3.717249in}}%
\pgfpathlineto{\pgfqpoint{5.371693in}{3.756524in}}%
\pgfpathlineto{\pgfqpoint{5.374212in}{3.744829in}}%
\pgfpathlineto{\pgfqpoint{5.376731in}{3.758978in}}%
\pgfpathlineto{\pgfqpoint{5.379249in}{3.751763in}}%
\pgfpathlineto{\pgfqpoint{5.381768in}{3.780206in}}%
\pgfpathlineto{\pgfqpoint{5.384287in}{3.780156in}}%
\pgfpathlineto{\pgfqpoint{5.386806in}{3.778487in}}%
\pgfpathlineto{\pgfqpoint{5.389324in}{3.801409in}}%
\pgfpathlineto{\pgfqpoint{5.391843in}{3.793819in}}%
\pgfpathlineto{\pgfqpoint{5.394362in}{3.769720in}}%
\pgfpathlineto{\pgfqpoint{5.396881in}{3.782061in}}%
\pgfpathlineto{\pgfqpoint{5.399399in}{3.761410in}}%
\pgfpathlineto{\pgfqpoint{5.401918in}{3.725651in}}%
\pgfpathlineto{\pgfqpoint{5.404437in}{3.752285in}}%
\pgfpathlineto{\pgfqpoint{5.406956in}{3.762986in}}%
\pgfpathlineto{\pgfqpoint{5.409474in}{3.759200in}}%
\pgfpathlineto{\pgfqpoint{5.411993in}{3.750326in}}%
\pgfpathlineto{\pgfqpoint{5.414512in}{3.739399in}}%
\pgfpathlineto{\pgfqpoint{5.417031in}{3.768081in}}%
\pgfpathlineto{\pgfqpoint{5.419549in}{3.789281in}}%
\pgfpathlineto{\pgfqpoint{5.422068in}{3.759289in}}%
\pgfpathlineto{\pgfqpoint{5.424587in}{3.783522in}}%
\pgfpathlineto{\pgfqpoint{5.427106in}{3.815308in}}%
\pgfpathlineto{\pgfqpoint{5.429624in}{3.850706in}}%
\pgfpathlineto{\pgfqpoint{5.432143in}{3.810348in}}%
\pgfpathlineto{\pgfqpoint{5.434662in}{3.777706in}}%
\pgfpathlineto{\pgfqpoint{5.437181in}{3.799598in}}%
\pgfpathlineto{\pgfqpoint{5.439699in}{3.836156in}}%
\pgfpathlineto{\pgfqpoint{5.442218in}{3.818848in}}%
\pgfpathlineto{\pgfqpoint{5.444737in}{3.837582in}}%
\pgfpathlineto{\pgfqpoint{5.447256in}{3.877105in}}%
\pgfpathlineto{\pgfqpoint{5.449774in}{3.922552in}}%
\pgfpathlineto{\pgfqpoint{5.452293in}{3.984381in}}%
\pgfpathlineto{\pgfqpoint{5.454812in}{3.982736in}}%
\pgfpathlineto{\pgfqpoint{5.457331in}{3.997750in}}%
\pgfpathlineto{\pgfqpoint{5.459849in}{4.033086in}}%
\pgfpathlineto{\pgfqpoint{5.462368in}{4.026097in}}%
\pgfpathlineto{\pgfqpoint{5.464887in}{4.025591in}}%
\pgfpathlineto{\pgfqpoint{5.467406in}{4.050603in}}%
\pgfpathlineto{\pgfqpoint{5.469924in}{4.052440in}}%
\pgfpathlineto{\pgfqpoint{5.472443in}{4.065786in}}%
\pgfpathlineto{\pgfqpoint{5.474962in}{4.058182in}}%
\pgfpathlineto{\pgfqpoint{5.477481in}{4.066546in}}%
\pgfpathlineto{\pgfqpoint{5.479999in}{4.065454in}}%
\pgfpathlineto{\pgfqpoint{5.482518in}{4.091054in}}%
\pgfpathlineto{\pgfqpoint{5.485037in}{4.187576in}}%
\pgfpathlineto{\pgfqpoint{5.487556in}{4.223080in}}%
\pgfpathlineto{\pgfqpoint{5.490074in}{4.186817in}}%
\pgfpathlineto{\pgfqpoint{5.492593in}{4.197245in}}%
\pgfpathlineto{\pgfqpoint{5.492593in}{4.197245in}}%
\pgfusepath{stroke}%
\end{pgfscope}%
\begin{pgfscope}%
\pgfpathrectangle{\pgfqpoint{0.455093in}{0.401080in}}{\pgfqpoint{5.037500in}{4.004000in}}%
\pgfusepath{clip}%
\pgfsetrectcap%
\pgfsetroundjoin%
\pgfsetlinewidth{0.501875pt}%
\definecolor{currentstroke}{rgb}{0.690196,0.188235,0.376471}%
\pgfsetstrokecolor{currentstroke}%
\pgfsetstrokeopacity{0.800000}%
\pgfsetdash{}{0pt}%
\pgfpathmoveto{\pgfqpoint{0.455093in}{1.021044in}}%
\pgfpathlineto{\pgfqpoint{0.457612in}{1.024732in}}%
\pgfpathlineto{\pgfqpoint{0.460131in}{1.041031in}}%
\pgfpathlineto{\pgfqpoint{0.462649in}{1.059641in}}%
\pgfpathlineto{\pgfqpoint{0.465168in}{1.046487in}}%
\pgfpathlineto{\pgfqpoint{0.467687in}{1.048519in}}%
\pgfpathlineto{\pgfqpoint{0.470206in}{1.049205in}}%
\pgfpathlineto{\pgfqpoint{0.472724in}{1.055314in}}%
\pgfpathlineto{\pgfqpoint{0.475243in}{1.029632in}}%
\pgfpathlineto{\pgfqpoint{0.477762in}{1.020312in}}%
\pgfpathlineto{\pgfqpoint{0.480281in}{1.026258in}}%
\pgfpathlineto{\pgfqpoint{0.482799in}{1.022649in}}%
\pgfpathlineto{\pgfqpoint{0.485318in}{1.019423in}}%
\pgfpathlineto{\pgfqpoint{0.487837in}{1.019875in}}%
\pgfpathlineto{\pgfqpoint{0.490356in}{1.033255in}}%
\pgfpathlineto{\pgfqpoint{0.492874in}{1.037525in}}%
\pgfpathlineto{\pgfqpoint{0.495393in}{1.016140in}}%
\pgfpathlineto{\pgfqpoint{0.497912in}{1.026466in}}%
\pgfpathlineto{\pgfqpoint{0.500431in}{1.032186in}}%
\pgfpathlineto{\pgfqpoint{0.502949in}{1.010714in}}%
\pgfpathlineto{\pgfqpoint{0.505468in}{1.008395in}}%
\pgfpathlineto{\pgfqpoint{0.507987in}{0.989276in}}%
\pgfpathlineto{\pgfqpoint{0.510506in}{0.976857in}}%
\pgfpathlineto{\pgfqpoint{0.513024in}{0.977643in}}%
\pgfpathlineto{\pgfqpoint{0.515543in}{0.964489in}}%
\pgfpathlineto{\pgfqpoint{0.518062in}{0.972091in}}%
\pgfpathlineto{\pgfqpoint{0.520581in}{0.962214in}}%
\pgfpathlineto{\pgfqpoint{0.523099in}{0.961969in}}%
\pgfpathlineto{\pgfqpoint{0.525618in}{0.977871in}}%
\pgfpathlineto{\pgfqpoint{0.528137in}{0.958388in}}%
\pgfpathlineto{\pgfqpoint{0.530656in}{0.961909in}}%
\pgfpathlineto{\pgfqpoint{0.533174in}{0.981606in}}%
\pgfpathlineto{\pgfqpoint{0.535693in}{0.977967in}}%
\pgfpathlineto{\pgfqpoint{0.538212in}{0.990446in}}%
\pgfpathlineto{\pgfqpoint{0.540731in}{0.997299in}}%
\pgfpathlineto{\pgfqpoint{0.543249in}{1.005630in}}%
\pgfpathlineto{\pgfqpoint{0.545768in}{0.998055in}}%
\pgfpathlineto{\pgfqpoint{0.548287in}{1.002138in}}%
\pgfpathlineto{\pgfqpoint{0.553324in}{1.028479in}}%
\pgfpathlineto{\pgfqpoint{0.555843in}{1.021881in}}%
\pgfpathlineto{\pgfqpoint{0.558362in}{1.009693in}}%
\pgfpathlineto{\pgfqpoint{0.560881in}{1.012875in}}%
\pgfpathlineto{\pgfqpoint{0.563399in}{0.995836in}}%
\pgfpathlineto{\pgfqpoint{0.565918in}{0.996475in}}%
\pgfpathlineto{\pgfqpoint{0.568437in}{1.001637in}}%
\pgfpathlineto{\pgfqpoint{0.570956in}{1.001579in}}%
\pgfpathlineto{\pgfqpoint{0.573474in}{1.022269in}}%
\pgfpathlineto{\pgfqpoint{0.575993in}{1.007175in}}%
\pgfpathlineto{\pgfqpoint{0.578512in}{1.020304in}}%
\pgfpathlineto{\pgfqpoint{0.581031in}{0.997971in}}%
\pgfpathlineto{\pgfqpoint{0.583549in}{1.001901in}}%
\pgfpathlineto{\pgfqpoint{0.586068in}{1.017112in}}%
\pgfpathlineto{\pgfqpoint{0.588587in}{1.029547in}}%
\pgfpathlineto{\pgfqpoint{0.591106in}{1.019452in}}%
\pgfpathlineto{\pgfqpoint{0.593624in}{1.024519in}}%
\pgfpathlineto{\pgfqpoint{0.596143in}{1.012019in}}%
\pgfpathlineto{\pgfqpoint{0.598662in}{1.015905in}}%
\pgfpathlineto{\pgfqpoint{0.601181in}{1.033285in}}%
\pgfpathlineto{\pgfqpoint{0.603699in}{1.018646in}}%
\pgfpathlineto{\pgfqpoint{0.606218in}{1.012222in}}%
\pgfpathlineto{\pgfqpoint{0.608737in}{1.011254in}}%
\pgfpathlineto{\pgfqpoint{0.611256in}{1.017867in}}%
\pgfpathlineto{\pgfqpoint{0.613774in}{1.016308in}}%
\pgfpathlineto{\pgfqpoint{0.616293in}{1.033924in}}%
\pgfpathlineto{\pgfqpoint{0.618812in}{1.042071in}}%
\pgfpathlineto{\pgfqpoint{0.621331in}{1.046690in}}%
\pgfpathlineto{\pgfqpoint{0.623849in}{1.047835in}}%
\pgfpathlineto{\pgfqpoint{0.626368in}{1.050730in}}%
\pgfpathlineto{\pgfqpoint{0.628887in}{1.045445in}}%
\pgfpathlineto{\pgfqpoint{0.631406in}{1.016423in}}%
\pgfpathlineto{\pgfqpoint{0.633924in}{1.008074in}}%
\pgfpathlineto{\pgfqpoint{0.636443in}{0.981946in}}%
\pgfpathlineto{\pgfqpoint{0.638962in}{0.974519in}}%
\pgfpathlineto{\pgfqpoint{0.641481in}{0.971221in}}%
\pgfpathlineto{\pgfqpoint{0.643999in}{0.968718in}}%
\pgfpathlineto{\pgfqpoint{0.646518in}{0.957180in}}%
\pgfpathlineto{\pgfqpoint{0.649037in}{0.948973in}}%
\pgfpathlineto{\pgfqpoint{0.651556in}{0.938354in}}%
\pgfpathlineto{\pgfqpoint{0.654074in}{0.939665in}}%
\pgfpathlineto{\pgfqpoint{0.656593in}{0.952364in}}%
\pgfpathlineto{\pgfqpoint{0.659112in}{0.974655in}}%
\pgfpathlineto{\pgfqpoint{0.661631in}{0.984658in}}%
\pgfpathlineto{\pgfqpoint{0.664149in}{0.986025in}}%
\pgfpathlineto{\pgfqpoint{0.666668in}{0.987881in}}%
\pgfpathlineto{\pgfqpoint{0.669187in}{0.989161in}}%
\pgfpathlineto{\pgfqpoint{0.671706in}{0.982065in}}%
\pgfpathlineto{\pgfqpoint{0.674224in}{1.026162in}}%
\pgfpathlineto{\pgfqpoint{0.676743in}{1.019637in}}%
\pgfpathlineto{\pgfqpoint{0.679262in}{1.012654in}}%
\pgfpathlineto{\pgfqpoint{0.681781in}{1.008388in}}%
\pgfpathlineto{\pgfqpoint{0.684299in}{1.019489in}}%
\pgfpathlineto{\pgfqpoint{0.686818in}{1.024164in}}%
\pgfpathlineto{\pgfqpoint{0.689337in}{1.024676in}}%
\pgfpathlineto{\pgfqpoint{0.691856in}{1.027996in}}%
\pgfpathlineto{\pgfqpoint{0.694374in}{1.038689in}}%
\pgfpathlineto{\pgfqpoint{0.696893in}{1.034799in}}%
\pgfpathlineto{\pgfqpoint{0.699412in}{1.045405in}}%
\pgfpathlineto{\pgfqpoint{0.701931in}{1.083400in}}%
\pgfpathlineto{\pgfqpoint{0.704449in}{1.079101in}}%
\pgfpathlineto{\pgfqpoint{0.706968in}{1.073736in}}%
\pgfpathlineto{\pgfqpoint{0.709487in}{1.073321in}}%
\pgfpathlineto{\pgfqpoint{0.712006in}{1.059341in}}%
\pgfpathlineto{\pgfqpoint{0.714524in}{1.073113in}}%
\pgfpathlineto{\pgfqpoint{0.717043in}{1.077167in}}%
\pgfpathlineto{\pgfqpoint{0.719562in}{1.077864in}}%
\pgfpathlineto{\pgfqpoint{0.722081in}{1.079795in}}%
\pgfpathlineto{\pgfqpoint{0.724599in}{1.067455in}}%
\pgfpathlineto{\pgfqpoint{0.727118in}{1.084158in}}%
\pgfpathlineto{\pgfqpoint{0.729637in}{1.088735in}}%
\pgfpathlineto{\pgfqpoint{0.732156in}{1.093789in}}%
\pgfpathlineto{\pgfqpoint{0.734674in}{1.112027in}}%
\pgfpathlineto{\pgfqpoint{0.737193in}{1.126771in}}%
\pgfpathlineto{\pgfqpoint{0.739712in}{1.125561in}}%
\pgfpathlineto{\pgfqpoint{0.742231in}{1.107065in}}%
\pgfpathlineto{\pgfqpoint{0.744749in}{1.115209in}}%
\pgfpathlineto{\pgfqpoint{0.747268in}{1.118627in}}%
\pgfpathlineto{\pgfqpoint{0.749787in}{1.112424in}}%
\pgfpathlineto{\pgfqpoint{0.752306in}{1.113507in}}%
\pgfpathlineto{\pgfqpoint{0.754824in}{1.116188in}}%
\pgfpathlineto{\pgfqpoint{0.757343in}{1.132156in}}%
\pgfpathlineto{\pgfqpoint{0.759862in}{1.156862in}}%
\pgfpathlineto{\pgfqpoint{0.762381in}{1.126708in}}%
\pgfpathlineto{\pgfqpoint{0.764899in}{1.141705in}}%
\pgfpathlineto{\pgfqpoint{0.767418in}{1.123397in}}%
\pgfpathlineto{\pgfqpoint{0.769937in}{1.136473in}}%
\pgfpathlineto{\pgfqpoint{0.772456in}{1.139898in}}%
\pgfpathlineto{\pgfqpoint{0.774974in}{1.151870in}}%
\pgfpathlineto{\pgfqpoint{0.777493in}{1.138405in}}%
\pgfpathlineto{\pgfqpoint{0.780012in}{1.159026in}}%
\pgfpathlineto{\pgfqpoint{0.782531in}{1.156608in}}%
\pgfpathlineto{\pgfqpoint{0.785049in}{1.176423in}}%
\pgfpathlineto{\pgfqpoint{0.787568in}{1.153213in}}%
\pgfpathlineto{\pgfqpoint{0.790087in}{1.125343in}}%
\pgfpathlineto{\pgfqpoint{0.792606in}{1.095355in}}%
\pgfpathlineto{\pgfqpoint{0.795124in}{1.113529in}}%
\pgfpathlineto{\pgfqpoint{0.797643in}{1.096952in}}%
\pgfpathlineto{\pgfqpoint{0.800162in}{1.106589in}}%
\pgfpathlineto{\pgfqpoint{0.802681in}{1.129815in}}%
\pgfpathlineto{\pgfqpoint{0.805199in}{1.151817in}}%
\pgfpathlineto{\pgfqpoint{0.807718in}{1.166821in}}%
\pgfpathlineto{\pgfqpoint{0.810237in}{1.177508in}}%
\pgfpathlineto{\pgfqpoint{0.812756in}{1.197114in}}%
\pgfpathlineto{\pgfqpoint{0.815274in}{1.184730in}}%
\pgfpathlineto{\pgfqpoint{0.817793in}{1.175531in}}%
\pgfpathlineto{\pgfqpoint{0.820312in}{1.174427in}}%
\pgfpathlineto{\pgfqpoint{0.822831in}{1.187712in}}%
\pgfpathlineto{\pgfqpoint{0.825349in}{1.218786in}}%
\pgfpathlineto{\pgfqpoint{0.827868in}{1.235110in}}%
\pgfpathlineto{\pgfqpoint{0.830387in}{1.228867in}}%
\pgfpathlineto{\pgfqpoint{0.832906in}{1.226262in}}%
\pgfpathlineto{\pgfqpoint{0.835424in}{1.242174in}}%
\pgfpathlineto{\pgfqpoint{0.837943in}{1.265832in}}%
\pgfpathlineto{\pgfqpoint{0.840462in}{1.264629in}}%
\pgfpathlineto{\pgfqpoint{0.842981in}{1.243830in}}%
\pgfpathlineto{\pgfqpoint{0.845499in}{1.249605in}}%
\pgfpathlineto{\pgfqpoint{0.848018in}{1.251326in}}%
\pgfpathlineto{\pgfqpoint{0.850537in}{1.251350in}}%
\pgfpathlineto{\pgfqpoint{0.853056in}{1.240147in}}%
\pgfpathlineto{\pgfqpoint{0.855574in}{1.218673in}}%
\pgfpathlineto{\pgfqpoint{0.858093in}{1.225584in}}%
\pgfpathlineto{\pgfqpoint{0.860612in}{1.220936in}}%
\pgfpathlineto{\pgfqpoint{0.863131in}{1.236770in}}%
\pgfpathlineto{\pgfqpoint{0.865649in}{1.240659in}}%
\pgfpathlineto{\pgfqpoint{0.868168in}{1.255882in}}%
\pgfpathlineto{\pgfqpoint{0.870687in}{1.251282in}}%
\pgfpathlineto{\pgfqpoint{0.873206in}{1.252496in}}%
\pgfpathlineto{\pgfqpoint{0.875724in}{1.228890in}}%
\pgfpathlineto{\pgfqpoint{0.878243in}{1.235526in}}%
\pgfpathlineto{\pgfqpoint{0.880762in}{1.230491in}}%
\pgfpathlineto{\pgfqpoint{0.883281in}{1.244742in}}%
\pgfpathlineto{\pgfqpoint{0.885799in}{1.225826in}}%
\pgfpathlineto{\pgfqpoint{0.888318in}{1.231918in}}%
\pgfpathlineto{\pgfqpoint{0.890837in}{1.250864in}}%
\pgfpathlineto{\pgfqpoint{0.893356in}{1.241633in}}%
\pgfpathlineto{\pgfqpoint{0.895874in}{1.255970in}}%
\pgfpathlineto{\pgfqpoint{0.898393in}{1.293297in}}%
\pgfpathlineto{\pgfqpoint{0.900912in}{1.321983in}}%
\pgfpathlineto{\pgfqpoint{0.903431in}{1.353332in}}%
\pgfpathlineto{\pgfqpoint{0.905949in}{1.363219in}}%
\pgfpathlineto{\pgfqpoint{0.908468in}{1.375591in}}%
\pgfpathlineto{\pgfqpoint{0.910987in}{1.369367in}}%
\pgfpathlineto{\pgfqpoint{0.913506in}{1.345642in}}%
\pgfpathlineto{\pgfqpoint{0.916024in}{1.358639in}}%
\pgfpathlineto{\pgfqpoint{0.918543in}{1.345687in}}%
\pgfpathlineto{\pgfqpoint{0.921062in}{1.354575in}}%
\pgfpathlineto{\pgfqpoint{0.923581in}{1.346823in}}%
\pgfpathlineto{\pgfqpoint{0.926099in}{1.357275in}}%
\pgfpathlineto{\pgfqpoint{0.928618in}{1.342358in}}%
\pgfpathlineto{\pgfqpoint{0.931137in}{1.354369in}}%
\pgfpathlineto{\pgfqpoint{0.933656in}{1.338143in}}%
\pgfpathlineto{\pgfqpoint{0.936174in}{1.351812in}}%
\pgfpathlineto{\pgfqpoint{0.938693in}{1.334467in}}%
\pgfpathlineto{\pgfqpoint{0.941212in}{1.344260in}}%
\pgfpathlineto{\pgfqpoint{0.943731in}{1.338631in}}%
\pgfpathlineto{\pgfqpoint{0.946249in}{1.345849in}}%
\pgfpathlineto{\pgfqpoint{0.948768in}{1.333967in}}%
\pgfpathlineto{\pgfqpoint{0.951287in}{1.335577in}}%
\pgfpathlineto{\pgfqpoint{0.953806in}{1.345637in}}%
\pgfpathlineto{\pgfqpoint{0.956324in}{1.331446in}}%
\pgfpathlineto{\pgfqpoint{0.958843in}{1.339332in}}%
\pgfpathlineto{\pgfqpoint{0.961362in}{1.350135in}}%
\pgfpathlineto{\pgfqpoint{0.963881in}{1.347507in}}%
\pgfpathlineto{\pgfqpoint{0.966399in}{1.370099in}}%
\pgfpathlineto{\pgfqpoint{0.968918in}{1.358402in}}%
\pgfpathlineto{\pgfqpoint{0.971437in}{1.360162in}}%
\pgfpathlineto{\pgfqpoint{0.973956in}{1.348941in}}%
\pgfpathlineto{\pgfqpoint{0.976474in}{1.345658in}}%
\pgfpathlineto{\pgfqpoint{0.978993in}{1.339903in}}%
\pgfpathlineto{\pgfqpoint{0.981512in}{1.339313in}}%
\pgfpathlineto{\pgfqpoint{0.984031in}{1.348210in}}%
\pgfpathlineto{\pgfqpoint{0.986549in}{1.350284in}}%
\pgfpathlineto{\pgfqpoint{0.989068in}{1.381726in}}%
\pgfpathlineto{\pgfqpoint{0.991587in}{1.367009in}}%
\pgfpathlineto{\pgfqpoint{0.994106in}{1.365846in}}%
\pgfpathlineto{\pgfqpoint{0.996624in}{1.334292in}}%
\pgfpathlineto{\pgfqpoint{0.999143in}{1.324731in}}%
\pgfpathlineto{\pgfqpoint{1.001662in}{1.325962in}}%
\pgfpathlineto{\pgfqpoint{1.006699in}{1.330875in}}%
\pgfpathlineto{\pgfqpoint{1.009218in}{1.325309in}}%
\pgfpathlineto{\pgfqpoint{1.011737in}{1.303860in}}%
\pgfpathlineto{\pgfqpoint{1.014256in}{1.295348in}}%
\pgfpathlineto{\pgfqpoint{1.016774in}{1.284746in}}%
\pgfpathlineto{\pgfqpoint{1.021812in}{1.310461in}}%
\pgfpathlineto{\pgfqpoint{1.024331in}{1.291607in}}%
\pgfpathlineto{\pgfqpoint{1.026849in}{1.299411in}}%
\pgfpathlineto{\pgfqpoint{1.029368in}{1.315310in}}%
\pgfpathlineto{\pgfqpoint{1.031887in}{1.336167in}}%
\pgfpathlineto{\pgfqpoint{1.034406in}{1.341519in}}%
\pgfpathlineto{\pgfqpoint{1.036924in}{1.333216in}}%
\pgfpathlineto{\pgfqpoint{1.039443in}{1.347439in}}%
\pgfpathlineto{\pgfqpoint{1.041962in}{1.362594in}}%
\pgfpathlineto{\pgfqpoint{1.044481in}{1.348420in}}%
\pgfpathlineto{\pgfqpoint{1.046999in}{1.336225in}}%
\pgfpathlineto{\pgfqpoint{1.049518in}{1.334410in}}%
\pgfpathlineto{\pgfqpoint{1.052037in}{1.324698in}}%
\pgfpathlineto{\pgfqpoint{1.054556in}{1.309411in}}%
\pgfpathlineto{\pgfqpoint{1.057074in}{1.307208in}}%
\pgfpathlineto{\pgfqpoint{1.059593in}{1.328261in}}%
\pgfpathlineto{\pgfqpoint{1.062112in}{1.314193in}}%
\pgfpathlineto{\pgfqpoint{1.064631in}{1.298500in}}%
\pgfpathlineto{\pgfqpoint{1.067149in}{1.306248in}}%
\pgfpathlineto{\pgfqpoint{1.069668in}{1.286189in}}%
\pgfpathlineto{\pgfqpoint{1.072187in}{1.285435in}}%
\pgfpathlineto{\pgfqpoint{1.074706in}{1.304015in}}%
\pgfpathlineto{\pgfqpoint{1.077224in}{1.296884in}}%
\pgfpathlineto{\pgfqpoint{1.079743in}{1.295061in}}%
\pgfpathlineto{\pgfqpoint{1.082262in}{1.306034in}}%
\pgfpathlineto{\pgfqpoint{1.084781in}{1.302535in}}%
\pgfpathlineto{\pgfqpoint{1.087299in}{1.299797in}}%
\pgfpathlineto{\pgfqpoint{1.089818in}{1.316847in}}%
\pgfpathlineto{\pgfqpoint{1.092337in}{1.282492in}}%
\pgfpathlineto{\pgfqpoint{1.094856in}{1.279996in}}%
\pgfpathlineto{\pgfqpoint{1.097374in}{1.273015in}}%
\pgfpathlineto{\pgfqpoint{1.099893in}{1.278233in}}%
\pgfpathlineto{\pgfqpoint{1.102412in}{1.258644in}}%
\pgfpathlineto{\pgfqpoint{1.104931in}{1.252332in}}%
\pgfpathlineto{\pgfqpoint{1.107449in}{1.232512in}}%
\pgfpathlineto{\pgfqpoint{1.109968in}{1.222271in}}%
\pgfpathlineto{\pgfqpoint{1.112487in}{1.219127in}}%
\pgfpathlineto{\pgfqpoint{1.115006in}{1.232486in}}%
\pgfpathlineto{\pgfqpoint{1.117524in}{1.216299in}}%
\pgfpathlineto{\pgfqpoint{1.120043in}{1.218409in}}%
\pgfpathlineto{\pgfqpoint{1.122562in}{1.218200in}}%
\pgfpathlineto{\pgfqpoint{1.125081in}{1.200106in}}%
\pgfpathlineto{\pgfqpoint{1.127599in}{1.199837in}}%
\pgfpathlineto{\pgfqpoint{1.130118in}{1.221445in}}%
\pgfpathlineto{\pgfqpoint{1.132637in}{1.226921in}}%
\pgfpathlineto{\pgfqpoint{1.135156in}{1.210580in}}%
\pgfpathlineto{\pgfqpoint{1.137674in}{1.234716in}}%
\pgfpathlineto{\pgfqpoint{1.140193in}{1.255025in}}%
\pgfpathlineto{\pgfqpoint{1.142712in}{1.262778in}}%
\pgfpathlineto{\pgfqpoint{1.145231in}{1.289091in}}%
\pgfpathlineto{\pgfqpoint{1.147749in}{1.288771in}}%
\pgfpathlineto{\pgfqpoint{1.150268in}{1.278668in}}%
\pgfpathlineto{\pgfqpoint{1.152787in}{1.293137in}}%
\pgfpathlineto{\pgfqpoint{1.155306in}{1.288715in}}%
\pgfpathlineto{\pgfqpoint{1.157824in}{1.308734in}}%
\pgfpathlineto{\pgfqpoint{1.160343in}{1.318409in}}%
\pgfpathlineto{\pgfqpoint{1.162862in}{1.314537in}}%
\pgfpathlineto{\pgfqpoint{1.165381in}{1.315449in}}%
\pgfpathlineto{\pgfqpoint{1.167899in}{1.323827in}}%
\pgfpathlineto{\pgfqpoint{1.170418in}{1.311782in}}%
\pgfpathlineto{\pgfqpoint{1.172937in}{1.329064in}}%
\pgfpathlineto{\pgfqpoint{1.175456in}{1.313349in}}%
\pgfpathlineto{\pgfqpoint{1.177974in}{1.302264in}}%
\pgfpathlineto{\pgfqpoint{1.180493in}{1.293497in}}%
\pgfpathlineto{\pgfqpoint{1.183012in}{1.281797in}}%
\pgfpathlineto{\pgfqpoint{1.185531in}{1.329918in}}%
\pgfpathlineto{\pgfqpoint{1.188049in}{1.336803in}}%
\pgfpathlineto{\pgfqpoint{1.190568in}{1.335069in}}%
\pgfpathlineto{\pgfqpoint{1.193087in}{1.348516in}}%
\pgfpathlineto{\pgfqpoint{1.195606in}{1.343472in}}%
\pgfpathlineto{\pgfqpoint{1.198124in}{1.369700in}}%
\pgfpathlineto{\pgfqpoint{1.200643in}{1.372097in}}%
\pgfpathlineto{\pgfqpoint{1.203162in}{1.371521in}}%
\pgfpathlineto{\pgfqpoint{1.205681in}{1.389220in}}%
\pgfpathlineto{\pgfqpoint{1.208199in}{1.416550in}}%
\pgfpathlineto{\pgfqpoint{1.210718in}{1.420120in}}%
\pgfpathlineto{\pgfqpoint{1.213237in}{1.428863in}}%
\pgfpathlineto{\pgfqpoint{1.215756in}{1.435611in}}%
\pgfpathlineto{\pgfqpoint{1.218274in}{1.430574in}}%
\pgfpathlineto{\pgfqpoint{1.220793in}{1.453994in}}%
\pgfpathlineto{\pgfqpoint{1.223312in}{1.457991in}}%
\pgfpathlineto{\pgfqpoint{1.225831in}{1.470567in}}%
\pgfpathlineto{\pgfqpoint{1.228349in}{1.460808in}}%
\pgfpathlineto{\pgfqpoint{1.230868in}{1.473775in}}%
\pgfpathlineto{\pgfqpoint{1.233387in}{1.482141in}}%
\pgfpathlineto{\pgfqpoint{1.235906in}{1.473169in}}%
\pgfpathlineto{\pgfqpoint{1.238424in}{1.488767in}}%
\pgfpathlineto{\pgfqpoint{1.240943in}{1.457106in}}%
\pgfpathlineto{\pgfqpoint{1.243462in}{1.441320in}}%
\pgfpathlineto{\pgfqpoint{1.245981in}{1.424684in}}%
\pgfpathlineto{\pgfqpoint{1.248499in}{1.417153in}}%
\pgfpathlineto{\pgfqpoint{1.251018in}{1.431005in}}%
\pgfpathlineto{\pgfqpoint{1.253537in}{1.427681in}}%
\pgfpathlineto{\pgfqpoint{1.256056in}{1.437755in}}%
\pgfpathlineto{\pgfqpoint{1.258574in}{1.432251in}}%
\pgfpathlineto{\pgfqpoint{1.261093in}{1.442158in}}%
\pgfpathlineto{\pgfqpoint{1.263612in}{1.435825in}}%
\pgfpathlineto{\pgfqpoint{1.266131in}{1.446519in}}%
\pgfpathlineto{\pgfqpoint{1.268649in}{1.456216in}}%
\pgfpathlineto{\pgfqpoint{1.271168in}{1.469055in}}%
\pgfpathlineto{\pgfqpoint{1.273687in}{1.440914in}}%
\pgfpathlineto{\pgfqpoint{1.276206in}{1.425788in}}%
\pgfpathlineto{\pgfqpoint{1.278724in}{1.416357in}}%
\pgfpathlineto{\pgfqpoint{1.281243in}{1.413756in}}%
\pgfpathlineto{\pgfqpoint{1.283762in}{1.386676in}}%
\pgfpathlineto{\pgfqpoint{1.286281in}{1.372248in}}%
\pgfpathlineto{\pgfqpoint{1.288799in}{1.368053in}}%
\pgfpathlineto{\pgfqpoint{1.291318in}{1.363199in}}%
\pgfpathlineto{\pgfqpoint{1.293837in}{1.379263in}}%
\pgfpathlineto{\pgfqpoint{1.296356in}{1.363809in}}%
\pgfpathlineto{\pgfqpoint{1.298874in}{1.359405in}}%
\pgfpathlineto{\pgfqpoint{1.301393in}{1.348555in}}%
\pgfpathlineto{\pgfqpoint{1.303912in}{1.356970in}}%
\pgfpathlineto{\pgfqpoint{1.306431in}{1.349774in}}%
\pgfpathlineto{\pgfqpoint{1.311468in}{1.314601in}}%
\pgfpathlineto{\pgfqpoint{1.313987in}{1.315256in}}%
\pgfpathlineto{\pgfqpoint{1.316506in}{1.309971in}}%
\pgfpathlineto{\pgfqpoint{1.319024in}{1.326117in}}%
\pgfpathlineto{\pgfqpoint{1.321543in}{1.312856in}}%
\pgfpathlineto{\pgfqpoint{1.324062in}{1.333846in}}%
\pgfpathlineto{\pgfqpoint{1.326581in}{1.349030in}}%
\pgfpathlineto{\pgfqpoint{1.329099in}{1.340841in}}%
\pgfpathlineto{\pgfqpoint{1.331618in}{1.361866in}}%
\pgfpathlineto{\pgfqpoint{1.334137in}{1.357062in}}%
\pgfpathlineto{\pgfqpoint{1.336656in}{1.356269in}}%
\pgfpathlineto{\pgfqpoint{1.339174in}{1.382799in}}%
\pgfpathlineto{\pgfqpoint{1.341693in}{1.381871in}}%
\pgfpathlineto{\pgfqpoint{1.344212in}{1.387814in}}%
\pgfpathlineto{\pgfqpoint{1.346731in}{1.398770in}}%
\pgfpathlineto{\pgfqpoint{1.349249in}{1.397777in}}%
\pgfpathlineto{\pgfqpoint{1.351768in}{1.380354in}}%
\pgfpathlineto{\pgfqpoint{1.354287in}{1.367994in}}%
\pgfpathlineto{\pgfqpoint{1.356806in}{1.376266in}}%
\pgfpathlineto{\pgfqpoint{1.359324in}{1.378493in}}%
\pgfpathlineto{\pgfqpoint{1.361843in}{1.387518in}}%
\pgfpathlineto{\pgfqpoint{1.364362in}{1.382923in}}%
\pgfpathlineto{\pgfqpoint{1.366881in}{1.386143in}}%
\pgfpathlineto{\pgfqpoint{1.369399in}{1.386020in}}%
\pgfpathlineto{\pgfqpoint{1.371918in}{1.374811in}}%
\pgfpathlineto{\pgfqpoint{1.374437in}{1.369531in}}%
\pgfpathlineto{\pgfqpoint{1.376956in}{1.363835in}}%
\pgfpathlineto{\pgfqpoint{1.379474in}{1.371963in}}%
\pgfpathlineto{\pgfqpoint{1.381993in}{1.389924in}}%
\pgfpathlineto{\pgfqpoint{1.384512in}{1.387427in}}%
\pgfpathlineto{\pgfqpoint{1.387031in}{1.392074in}}%
\pgfpathlineto{\pgfqpoint{1.389549in}{1.394865in}}%
\pgfpathlineto{\pgfqpoint{1.392068in}{1.399139in}}%
\pgfpathlineto{\pgfqpoint{1.394587in}{1.408735in}}%
\pgfpathlineto{\pgfqpoint{1.397106in}{1.404925in}}%
\pgfpathlineto{\pgfqpoint{1.399624in}{1.389089in}}%
\pgfpathlineto{\pgfqpoint{1.402143in}{1.401369in}}%
\pgfpathlineto{\pgfqpoint{1.404662in}{1.387390in}}%
\pgfpathlineto{\pgfqpoint{1.407181in}{1.419490in}}%
\pgfpathlineto{\pgfqpoint{1.409699in}{1.398107in}}%
\pgfpathlineto{\pgfqpoint{1.412218in}{1.410669in}}%
\pgfpathlineto{\pgfqpoint{1.414737in}{1.405965in}}%
\pgfpathlineto{\pgfqpoint{1.417256in}{1.395179in}}%
\pgfpathlineto{\pgfqpoint{1.419774in}{1.397536in}}%
\pgfpathlineto{\pgfqpoint{1.422293in}{1.415683in}}%
\pgfpathlineto{\pgfqpoint{1.424812in}{1.417110in}}%
\pgfpathlineto{\pgfqpoint{1.427331in}{1.440155in}}%
\pgfpathlineto{\pgfqpoint{1.429849in}{1.431678in}}%
\pgfpathlineto{\pgfqpoint{1.432368in}{1.436288in}}%
\pgfpathlineto{\pgfqpoint{1.434887in}{1.448530in}}%
\pgfpathlineto{\pgfqpoint{1.437406in}{1.466925in}}%
\pgfpathlineto{\pgfqpoint{1.439924in}{1.473875in}}%
\pgfpathlineto{\pgfqpoint{1.442443in}{1.473305in}}%
\pgfpathlineto{\pgfqpoint{1.444962in}{1.496855in}}%
\pgfpathlineto{\pgfqpoint{1.447481in}{1.514947in}}%
\pgfpathlineto{\pgfqpoint{1.449999in}{1.521069in}}%
\pgfpathlineto{\pgfqpoint{1.452518in}{1.498081in}}%
\pgfpathlineto{\pgfqpoint{1.455037in}{1.528485in}}%
\pgfpathlineto{\pgfqpoint{1.457556in}{1.548127in}}%
\pgfpathlineto{\pgfqpoint{1.460074in}{1.557012in}}%
\pgfpathlineto{\pgfqpoint{1.462593in}{1.561210in}}%
\pgfpathlineto{\pgfqpoint{1.465112in}{1.553498in}}%
\pgfpathlineto{\pgfqpoint{1.467631in}{1.581371in}}%
\pgfpathlineto{\pgfqpoint{1.470149in}{1.607893in}}%
\pgfpathlineto{\pgfqpoint{1.475187in}{1.575426in}}%
\pgfpathlineto{\pgfqpoint{1.477706in}{1.538466in}}%
\pgfpathlineto{\pgfqpoint{1.480224in}{1.514907in}}%
\pgfpathlineto{\pgfqpoint{1.482743in}{1.520489in}}%
\pgfpathlineto{\pgfqpoint{1.485262in}{1.515512in}}%
\pgfpathlineto{\pgfqpoint{1.487781in}{1.513124in}}%
\pgfpathlineto{\pgfqpoint{1.490299in}{1.504990in}}%
\pgfpathlineto{\pgfqpoint{1.492818in}{1.491294in}}%
\pgfpathlineto{\pgfqpoint{1.495337in}{1.496190in}}%
\pgfpathlineto{\pgfqpoint{1.497856in}{1.477491in}}%
\pgfpathlineto{\pgfqpoint{1.500374in}{1.478884in}}%
\pgfpathlineto{\pgfqpoint{1.502893in}{1.494400in}}%
\pgfpathlineto{\pgfqpoint{1.505412in}{1.483899in}}%
\pgfpathlineto{\pgfqpoint{1.507931in}{1.475299in}}%
\pgfpathlineto{\pgfqpoint{1.510449in}{1.475737in}}%
\pgfpathlineto{\pgfqpoint{1.512968in}{1.481999in}}%
\pgfpathlineto{\pgfqpoint{1.515487in}{1.473136in}}%
\pgfpathlineto{\pgfqpoint{1.518006in}{1.475047in}}%
\pgfpathlineto{\pgfqpoint{1.520524in}{1.485288in}}%
\pgfpathlineto{\pgfqpoint{1.523043in}{1.462109in}}%
\pgfpathlineto{\pgfqpoint{1.525562in}{1.463032in}}%
\pgfpathlineto{\pgfqpoint{1.528081in}{1.451231in}}%
\pgfpathlineto{\pgfqpoint{1.530599in}{1.448432in}}%
\pgfpathlineto{\pgfqpoint{1.533118in}{1.434656in}}%
\pgfpathlineto{\pgfqpoint{1.535637in}{1.444047in}}%
\pgfpathlineto{\pgfqpoint{1.538156in}{1.430609in}}%
\pgfpathlineto{\pgfqpoint{1.540674in}{1.402515in}}%
\pgfpathlineto{\pgfqpoint{1.543193in}{1.401185in}}%
\pgfpathlineto{\pgfqpoint{1.545712in}{1.393885in}}%
\pgfpathlineto{\pgfqpoint{1.548231in}{1.390342in}}%
\pgfpathlineto{\pgfqpoint{1.550749in}{1.387360in}}%
\pgfpathlineto{\pgfqpoint{1.553268in}{1.380940in}}%
\pgfpathlineto{\pgfqpoint{1.555787in}{1.373306in}}%
\pgfpathlineto{\pgfqpoint{1.558306in}{1.393748in}}%
\pgfpathlineto{\pgfqpoint{1.560824in}{1.402016in}}%
\pgfpathlineto{\pgfqpoint{1.563343in}{1.393597in}}%
\pgfpathlineto{\pgfqpoint{1.565862in}{1.394118in}}%
\pgfpathlineto{\pgfqpoint{1.568381in}{1.388786in}}%
\pgfpathlineto{\pgfqpoint{1.570899in}{1.390060in}}%
\pgfpathlineto{\pgfqpoint{1.573418in}{1.367946in}}%
\pgfpathlineto{\pgfqpoint{1.575937in}{1.386453in}}%
\pgfpathlineto{\pgfqpoint{1.578456in}{1.379865in}}%
\pgfpathlineto{\pgfqpoint{1.580974in}{1.408780in}}%
\pgfpathlineto{\pgfqpoint{1.583493in}{1.424047in}}%
\pgfpathlineto{\pgfqpoint{1.586012in}{1.403051in}}%
\pgfpathlineto{\pgfqpoint{1.588531in}{1.401465in}}%
\pgfpathlineto{\pgfqpoint{1.591049in}{1.405025in}}%
\pgfpathlineto{\pgfqpoint{1.593568in}{1.419285in}}%
\pgfpathlineto{\pgfqpoint{1.596087in}{1.413493in}}%
\pgfpathlineto{\pgfqpoint{1.598606in}{1.401518in}}%
\pgfpathlineto{\pgfqpoint{1.601124in}{1.396102in}}%
\pgfpathlineto{\pgfqpoint{1.603643in}{1.398030in}}%
\pgfpathlineto{\pgfqpoint{1.606162in}{1.382751in}}%
\pgfpathlineto{\pgfqpoint{1.608681in}{1.379904in}}%
\pgfpathlineto{\pgfqpoint{1.611199in}{1.391302in}}%
\pgfpathlineto{\pgfqpoint{1.613718in}{1.386226in}}%
\pgfpathlineto{\pgfqpoint{1.616237in}{1.389521in}}%
\pgfpathlineto{\pgfqpoint{1.618756in}{1.407786in}}%
\pgfpathlineto{\pgfqpoint{1.621274in}{1.408686in}}%
\pgfpathlineto{\pgfqpoint{1.623793in}{1.390909in}}%
\pgfpathlineto{\pgfqpoint{1.626312in}{1.399160in}}%
\pgfpathlineto{\pgfqpoint{1.628831in}{1.375284in}}%
\pgfpathlineto{\pgfqpoint{1.631349in}{1.352919in}}%
\pgfpathlineto{\pgfqpoint{1.633868in}{1.365311in}}%
\pgfpathlineto{\pgfqpoint{1.636387in}{1.352495in}}%
\pgfpathlineto{\pgfqpoint{1.638906in}{1.328577in}}%
\pgfpathlineto{\pgfqpoint{1.641424in}{1.323655in}}%
\pgfpathlineto{\pgfqpoint{1.643943in}{1.328180in}}%
\pgfpathlineto{\pgfqpoint{1.646462in}{1.313010in}}%
\pgfpathlineto{\pgfqpoint{1.648981in}{1.314361in}}%
\pgfpathlineto{\pgfqpoint{1.651499in}{1.316682in}}%
\pgfpathlineto{\pgfqpoint{1.654018in}{1.314421in}}%
\pgfpathlineto{\pgfqpoint{1.656537in}{1.305074in}}%
\pgfpathlineto{\pgfqpoint{1.659056in}{1.307928in}}%
\pgfpathlineto{\pgfqpoint{1.661574in}{1.336775in}}%
\pgfpathlineto{\pgfqpoint{1.664093in}{1.347985in}}%
\pgfpathlineto{\pgfqpoint{1.666612in}{1.339998in}}%
\pgfpathlineto{\pgfqpoint{1.669131in}{1.346432in}}%
\pgfpathlineto{\pgfqpoint{1.671649in}{1.341087in}}%
\pgfpathlineto{\pgfqpoint{1.674168in}{1.357533in}}%
\pgfpathlineto{\pgfqpoint{1.676687in}{1.364636in}}%
\pgfpathlineto{\pgfqpoint{1.679206in}{1.360846in}}%
\pgfpathlineto{\pgfqpoint{1.681724in}{1.360422in}}%
\pgfpathlineto{\pgfqpoint{1.684243in}{1.371808in}}%
\pgfpathlineto{\pgfqpoint{1.686762in}{1.385504in}}%
\pgfpathlineto{\pgfqpoint{1.689281in}{1.366206in}}%
\pgfpathlineto{\pgfqpoint{1.691799in}{1.363488in}}%
\pgfpathlineto{\pgfqpoint{1.694318in}{1.341879in}}%
\pgfpathlineto{\pgfqpoint{1.696837in}{1.304793in}}%
\pgfpathlineto{\pgfqpoint{1.699356in}{1.284203in}}%
\pgfpathlineto{\pgfqpoint{1.701874in}{1.272946in}}%
\pgfpathlineto{\pgfqpoint{1.704393in}{1.274329in}}%
\pgfpathlineto{\pgfqpoint{1.706912in}{1.259146in}}%
\pgfpathlineto{\pgfqpoint{1.709431in}{1.249591in}}%
\pgfpathlineto{\pgfqpoint{1.711949in}{1.239220in}}%
\pgfpathlineto{\pgfqpoint{1.714468in}{1.258766in}}%
\pgfpathlineto{\pgfqpoint{1.716987in}{1.258635in}}%
\pgfpathlineto{\pgfqpoint{1.719506in}{1.247250in}}%
\pgfpathlineto{\pgfqpoint{1.722024in}{1.245472in}}%
\pgfpathlineto{\pgfqpoint{1.724543in}{1.244253in}}%
\pgfpathlineto{\pgfqpoint{1.727062in}{1.234931in}}%
\pgfpathlineto{\pgfqpoint{1.729581in}{1.238692in}}%
\pgfpathlineto{\pgfqpoint{1.732099in}{1.243847in}}%
\pgfpathlineto{\pgfqpoint{1.734618in}{1.251073in}}%
\pgfpathlineto{\pgfqpoint{1.737137in}{1.252900in}}%
\pgfpathlineto{\pgfqpoint{1.739656in}{1.257722in}}%
\pgfpathlineto{\pgfqpoint{1.742174in}{1.258015in}}%
\pgfpathlineto{\pgfqpoint{1.744693in}{1.256965in}}%
\pgfpathlineto{\pgfqpoint{1.747212in}{1.255588in}}%
\pgfpathlineto{\pgfqpoint{1.749731in}{1.233404in}}%
\pgfpathlineto{\pgfqpoint{1.752249in}{1.230903in}}%
\pgfpathlineto{\pgfqpoint{1.754768in}{1.238361in}}%
\pgfpathlineto{\pgfqpoint{1.757287in}{1.243603in}}%
\pgfpathlineto{\pgfqpoint{1.759806in}{1.257674in}}%
\pgfpathlineto{\pgfqpoint{1.762324in}{1.259306in}}%
\pgfpathlineto{\pgfqpoint{1.764843in}{1.263862in}}%
\pgfpathlineto{\pgfqpoint{1.767362in}{1.233589in}}%
\pgfpathlineto{\pgfqpoint{1.769881in}{1.225382in}}%
\pgfpathlineto{\pgfqpoint{1.772399in}{1.227948in}}%
\pgfpathlineto{\pgfqpoint{1.774918in}{1.251751in}}%
\pgfpathlineto{\pgfqpoint{1.777437in}{1.282275in}}%
\pgfpathlineto{\pgfqpoint{1.779956in}{1.311417in}}%
\pgfpathlineto{\pgfqpoint{1.782474in}{1.323680in}}%
\pgfpathlineto{\pgfqpoint{1.784993in}{1.333523in}}%
\pgfpathlineto{\pgfqpoint{1.790031in}{1.361565in}}%
\pgfpathlineto{\pgfqpoint{1.792549in}{1.385302in}}%
\pgfpathlineto{\pgfqpoint{1.795068in}{1.427181in}}%
\pgfpathlineto{\pgfqpoint{1.797587in}{1.435953in}}%
\pgfpathlineto{\pgfqpoint{1.800106in}{1.460275in}}%
\pgfpathlineto{\pgfqpoint{1.802624in}{1.463403in}}%
\pgfpathlineto{\pgfqpoint{1.805143in}{1.461796in}}%
\pgfpathlineto{\pgfqpoint{1.807662in}{1.471272in}}%
\pgfpathlineto{\pgfqpoint{1.810181in}{1.486835in}}%
\pgfpathlineto{\pgfqpoint{1.812699in}{1.508332in}}%
\pgfpathlineto{\pgfqpoint{1.815218in}{1.507054in}}%
\pgfpathlineto{\pgfqpoint{1.817737in}{1.505936in}}%
\pgfpathlineto{\pgfqpoint{1.820256in}{1.490614in}}%
\pgfpathlineto{\pgfqpoint{1.822774in}{1.473778in}}%
\pgfpathlineto{\pgfqpoint{1.825293in}{1.472303in}}%
\pgfpathlineto{\pgfqpoint{1.827812in}{1.492532in}}%
\pgfpathlineto{\pgfqpoint{1.830331in}{1.484959in}}%
\pgfpathlineto{\pgfqpoint{1.832849in}{1.458626in}}%
\pgfpathlineto{\pgfqpoint{1.835368in}{1.461430in}}%
\pgfpathlineto{\pgfqpoint{1.837887in}{1.450338in}}%
\pgfpathlineto{\pgfqpoint{1.840406in}{1.447761in}}%
\pgfpathlineto{\pgfqpoint{1.842924in}{1.442321in}}%
\pgfpathlineto{\pgfqpoint{1.845443in}{1.451790in}}%
\pgfpathlineto{\pgfqpoint{1.847962in}{1.477799in}}%
\pgfpathlineto{\pgfqpoint{1.850481in}{1.502275in}}%
\pgfpathlineto{\pgfqpoint{1.852999in}{1.467623in}}%
\pgfpathlineto{\pgfqpoint{1.855518in}{1.465400in}}%
\pgfpathlineto{\pgfqpoint{1.858037in}{1.458740in}}%
\pgfpathlineto{\pgfqpoint{1.860556in}{1.448908in}}%
\pgfpathlineto{\pgfqpoint{1.863074in}{1.435111in}}%
\pgfpathlineto{\pgfqpoint{1.865593in}{1.431542in}}%
\pgfpathlineto{\pgfqpoint{1.868112in}{1.435698in}}%
\pgfpathlineto{\pgfqpoint{1.870631in}{1.472852in}}%
\pgfpathlineto{\pgfqpoint{1.873149in}{1.467586in}}%
\pgfpathlineto{\pgfqpoint{1.875668in}{1.479933in}}%
\pgfpathlineto{\pgfqpoint{1.878187in}{1.465940in}}%
\pgfpathlineto{\pgfqpoint{1.880706in}{1.447913in}}%
\pgfpathlineto{\pgfqpoint{1.883224in}{1.434449in}}%
\pgfpathlineto{\pgfqpoint{1.885743in}{1.450520in}}%
\pgfpathlineto{\pgfqpoint{1.888262in}{1.429536in}}%
\pgfpathlineto{\pgfqpoint{1.890781in}{1.465129in}}%
\pgfpathlineto{\pgfqpoint{1.893299in}{1.474952in}}%
\pgfpathlineto{\pgfqpoint{1.895818in}{1.468906in}}%
\pgfpathlineto{\pgfqpoint{1.898337in}{1.450478in}}%
\pgfpathlineto{\pgfqpoint{1.900856in}{1.441704in}}%
\pgfpathlineto{\pgfqpoint{1.903374in}{1.442262in}}%
\pgfpathlineto{\pgfqpoint{1.905893in}{1.439211in}}%
\pgfpathlineto{\pgfqpoint{1.908412in}{1.461896in}}%
\pgfpathlineto{\pgfqpoint{1.910931in}{1.459962in}}%
\pgfpathlineto{\pgfqpoint{1.913449in}{1.470379in}}%
\pgfpathlineto{\pgfqpoint{1.915968in}{1.463360in}}%
\pgfpathlineto{\pgfqpoint{1.918487in}{1.468447in}}%
\pgfpathlineto{\pgfqpoint{1.921006in}{1.473791in}}%
\pgfpathlineto{\pgfqpoint{1.923524in}{1.495401in}}%
\pgfpathlineto{\pgfqpoint{1.926043in}{1.471217in}}%
\pgfpathlineto{\pgfqpoint{1.928562in}{1.484282in}}%
\pgfpathlineto{\pgfqpoint{1.931081in}{1.495819in}}%
\pgfpathlineto{\pgfqpoint{1.933599in}{1.486199in}}%
\pgfpathlineto{\pgfqpoint{1.936118in}{1.493693in}}%
\pgfpathlineto{\pgfqpoint{1.938637in}{1.483379in}}%
\pgfpathlineto{\pgfqpoint{1.941156in}{1.472215in}}%
\pgfpathlineto{\pgfqpoint{1.943674in}{1.493838in}}%
\pgfpathlineto{\pgfqpoint{1.946193in}{1.512152in}}%
\pgfpathlineto{\pgfqpoint{1.948712in}{1.527155in}}%
\pgfpathlineto{\pgfqpoint{1.951231in}{1.526372in}}%
\pgfpathlineto{\pgfqpoint{1.953749in}{1.542850in}}%
\pgfpathlineto{\pgfqpoint{1.956268in}{1.557209in}}%
\pgfpathlineto{\pgfqpoint{1.958787in}{1.572893in}}%
\pgfpathlineto{\pgfqpoint{1.961306in}{1.603582in}}%
\pgfpathlineto{\pgfqpoint{1.963824in}{1.612993in}}%
\pgfpathlineto{\pgfqpoint{1.966343in}{1.578016in}}%
\pgfpathlineto{\pgfqpoint{1.968862in}{1.577456in}}%
\pgfpathlineto{\pgfqpoint{1.971381in}{1.620097in}}%
\pgfpathlineto{\pgfqpoint{1.973899in}{1.617175in}}%
\pgfpathlineto{\pgfqpoint{1.976418in}{1.639468in}}%
\pgfpathlineto{\pgfqpoint{1.978937in}{1.648025in}}%
\pgfpathlineto{\pgfqpoint{1.981456in}{1.634817in}}%
\pgfpathlineto{\pgfqpoint{1.983974in}{1.642456in}}%
\pgfpathlineto{\pgfqpoint{1.986493in}{1.629337in}}%
\pgfpathlineto{\pgfqpoint{1.989012in}{1.653363in}}%
\pgfpathlineto{\pgfqpoint{1.991531in}{1.643660in}}%
\pgfpathlineto{\pgfqpoint{1.994049in}{1.644694in}}%
\pgfpathlineto{\pgfqpoint{1.996568in}{1.662643in}}%
\pgfpathlineto{\pgfqpoint{1.999087in}{1.659654in}}%
\pgfpathlineto{\pgfqpoint{2.001606in}{1.658887in}}%
\pgfpathlineto{\pgfqpoint{2.004124in}{1.679902in}}%
\pgfpathlineto{\pgfqpoint{2.006643in}{1.690666in}}%
\pgfpathlineto{\pgfqpoint{2.009162in}{1.709423in}}%
\pgfpathlineto{\pgfqpoint{2.011681in}{1.693099in}}%
\pgfpathlineto{\pgfqpoint{2.014199in}{1.688284in}}%
\pgfpathlineto{\pgfqpoint{2.016718in}{1.669817in}}%
\pgfpathlineto{\pgfqpoint{2.019237in}{1.674790in}}%
\pgfpathlineto{\pgfqpoint{2.021756in}{1.672351in}}%
\pgfpathlineto{\pgfqpoint{2.024274in}{1.682055in}}%
\pgfpathlineto{\pgfqpoint{2.026793in}{1.686814in}}%
\pgfpathlineto{\pgfqpoint{2.029312in}{1.676264in}}%
\pgfpathlineto{\pgfqpoint{2.031831in}{1.662614in}}%
\pgfpathlineto{\pgfqpoint{2.034349in}{1.657663in}}%
\pgfpathlineto{\pgfqpoint{2.036868in}{1.666311in}}%
\pgfpathlineto{\pgfqpoint{2.039387in}{1.691168in}}%
\pgfpathlineto{\pgfqpoint{2.041906in}{1.691898in}}%
\pgfpathlineto{\pgfqpoint{2.044424in}{1.692426in}}%
\pgfpathlineto{\pgfqpoint{2.046943in}{1.715714in}}%
\pgfpathlineto{\pgfqpoint{2.049462in}{1.705125in}}%
\pgfpathlineto{\pgfqpoint{2.051981in}{1.720979in}}%
\pgfpathlineto{\pgfqpoint{2.054499in}{1.734410in}}%
\pgfpathlineto{\pgfqpoint{2.057018in}{1.751025in}}%
\pgfpathlineto{\pgfqpoint{2.059537in}{1.731886in}}%
\pgfpathlineto{\pgfqpoint{2.062056in}{1.729897in}}%
\pgfpathlineto{\pgfqpoint{2.064574in}{1.725957in}}%
\pgfpathlineto{\pgfqpoint{2.067093in}{1.721469in}}%
\pgfpathlineto{\pgfqpoint{2.069612in}{1.724019in}}%
\pgfpathlineto{\pgfqpoint{2.072131in}{1.715644in}}%
\pgfpathlineto{\pgfqpoint{2.074649in}{1.760278in}}%
\pgfpathlineto{\pgfqpoint{2.077168in}{1.772870in}}%
\pgfpathlineto{\pgfqpoint{2.079687in}{1.751226in}}%
\pgfpathlineto{\pgfqpoint{2.082206in}{1.768721in}}%
\pgfpathlineto{\pgfqpoint{2.084724in}{1.758706in}}%
\pgfpathlineto{\pgfqpoint{2.087243in}{1.750457in}}%
\pgfpathlineto{\pgfqpoint{2.089762in}{1.743950in}}%
\pgfpathlineto{\pgfqpoint{2.092281in}{1.764020in}}%
\pgfpathlineto{\pgfqpoint{2.094799in}{1.777392in}}%
\pgfpathlineto{\pgfqpoint{2.097318in}{1.793364in}}%
\pgfpathlineto{\pgfqpoint{2.099837in}{1.781333in}}%
\pgfpathlineto{\pgfqpoint{2.102356in}{1.783414in}}%
\pgfpathlineto{\pgfqpoint{2.104874in}{1.773016in}}%
\pgfpathlineto{\pgfqpoint{2.107393in}{1.787907in}}%
\pgfpathlineto{\pgfqpoint{2.109912in}{1.801033in}}%
\pgfpathlineto{\pgfqpoint{2.112431in}{1.798012in}}%
\pgfpathlineto{\pgfqpoint{2.114949in}{1.769663in}}%
\pgfpathlineto{\pgfqpoint{2.117468in}{1.786177in}}%
\pgfpathlineto{\pgfqpoint{2.119987in}{1.787945in}}%
\pgfpathlineto{\pgfqpoint{2.122506in}{1.812145in}}%
\pgfpathlineto{\pgfqpoint{2.125024in}{1.806564in}}%
\pgfpathlineto{\pgfqpoint{2.127543in}{1.829595in}}%
\pgfpathlineto{\pgfqpoint{2.130062in}{1.810316in}}%
\pgfpathlineto{\pgfqpoint{2.132581in}{1.799539in}}%
\pgfpathlineto{\pgfqpoint{2.135099in}{1.777483in}}%
\pgfpathlineto{\pgfqpoint{2.137618in}{1.785413in}}%
\pgfpathlineto{\pgfqpoint{2.140137in}{1.731715in}}%
\pgfpathlineto{\pgfqpoint{2.142656in}{1.732172in}}%
\pgfpathlineto{\pgfqpoint{2.145174in}{1.712899in}}%
\pgfpathlineto{\pgfqpoint{2.147693in}{1.712519in}}%
\pgfpathlineto{\pgfqpoint{2.150212in}{1.723856in}}%
\pgfpathlineto{\pgfqpoint{2.152731in}{1.736492in}}%
\pgfpathlineto{\pgfqpoint{2.155249in}{1.762520in}}%
\pgfpathlineto{\pgfqpoint{2.157768in}{1.747682in}}%
\pgfpathlineto{\pgfqpoint{2.160287in}{1.762200in}}%
\pgfpathlineto{\pgfqpoint{2.162806in}{1.761009in}}%
\pgfpathlineto{\pgfqpoint{2.165324in}{1.739963in}}%
\pgfpathlineto{\pgfqpoint{2.167843in}{1.729730in}}%
\pgfpathlineto{\pgfqpoint{2.170362in}{1.720256in}}%
\pgfpathlineto{\pgfqpoint{2.172881in}{1.677500in}}%
\pgfpathlineto{\pgfqpoint{2.175399in}{1.662565in}}%
\pgfpathlineto{\pgfqpoint{2.177918in}{1.681343in}}%
\pgfpathlineto{\pgfqpoint{2.180437in}{1.675161in}}%
\pgfpathlineto{\pgfqpoint{2.182956in}{1.679536in}}%
\pgfpathlineto{\pgfqpoint{2.185474in}{1.688178in}}%
\pgfpathlineto{\pgfqpoint{2.187993in}{1.695289in}}%
\pgfpathlineto{\pgfqpoint{2.190512in}{1.669742in}}%
\pgfpathlineto{\pgfqpoint{2.193031in}{1.695030in}}%
\pgfpathlineto{\pgfqpoint{2.195549in}{1.702447in}}%
\pgfpathlineto{\pgfqpoint{2.198068in}{1.702122in}}%
\pgfpathlineto{\pgfqpoint{2.200587in}{1.720042in}}%
\pgfpathlineto{\pgfqpoint{2.203106in}{1.728586in}}%
\pgfpathlineto{\pgfqpoint{2.205624in}{1.712612in}}%
\pgfpathlineto{\pgfqpoint{2.208143in}{1.723929in}}%
\pgfpathlineto{\pgfqpoint{2.210662in}{1.727657in}}%
\pgfpathlineto{\pgfqpoint{2.213181in}{1.729474in}}%
\pgfpathlineto{\pgfqpoint{2.215699in}{1.738609in}}%
\pgfpathlineto{\pgfqpoint{2.218218in}{1.737809in}}%
\pgfpathlineto{\pgfqpoint{2.220737in}{1.736679in}}%
\pgfpathlineto{\pgfqpoint{2.223256in}{1.743551in}}%
\pgfpathlineto{\pgfqpoint{2.225774in}{1.716337in}}%
\pgfpathlineto{\pgfqpoint{2.228293in}{1.725132in}}%
\pgfpathlineto{\pgfqpoint{2.230812in}{1.744001in}}%
\pgfpathlineto{\pgfqpoint{2.233331in}{1.737026in}}%
\pgfpathlineto{\pgfqpoint{2.235849in}{1.736457in}}%
\pgfpathlineto{\pgfqpoint{2.238368in}{1.751065in}}%
\pgfpathlineto{\pgfqpoint{2.240887in}{1.774471in}}%
\pgfpathlineto{\pgfqpoint{2.243406in}{1.754299in}}%
\pgfpathlineto{\pgfqpoint{2.245924in}{1.747322in}}%
\pgfpathlineto{\pgfqpoint{2.248443in}{1.757659in}}%
\pgfpathlineto{\pgfqpoint{2.250962in}{1.755200in}}%
\pgfpathlineto{\pgfqpoint{2.253481in}{1.759503in}}%
\pgfpathlineto{\pgfqpoint{2.255999in}{1.753314in}}%
\pgfpathlineto{\pgfqpoint{2.258518in}{1.734312in}}%
\pgfpathlineto{\pgfqpoint{2.261037in}{1.713262in}}%
\pgfpathlineto{\pgfqpoint{2.263556in}{1.730191in}}%
\pgfpathlineto{\pgfqpoint{2.266074in}{1.768418in}}%
\pgfpathlineto{\pgfqpoint{2.268593in}{1.760426in}}%
\pgfpathlineto{\pgfqpoint{2.271112in}{1.766197in}}%
\pgfpathlineto{\pgfqpoint{2.273631in}{1.765928in}}%
\pgfpathlineto{\pgfqpoint{2.276149in}{1.750438in}}%
\pgfpathlineto{\pgfqpoint{2.278668in}{1.747762in}}%
\pgfpathlineto{\pgfqpoint{2.281187in}{1.731082in}}%
\pgfpathlineto{\pgfqpoint{2.283706in}{1.718980in}}%
\pgfpathlineto{\pgfqpoint{2.286224in}{1.732574in}}%
\pgfpathlineto{\pgfqpoint{2.288743in}{1.749455in}}%
\pgfpathlineto{\pgfqpoint{2.291262in}{1.762113in}}%
\pgfpathlineto{\pgfqpoint{2.293781in}{1.740677in}}%
\pgfpathlineto{\pgfqpoint{2.296299in}{1.752186in}}%
\pgfpathlineto{\pgfqpoint{2.298818in}{1.729350in}}%
\pgfpathlineto{\pgfqpoint{2.301337in}{1.719115in}}%
\pgfpathlineto{\pgfqpoint{2.303856in}{1.705080in}}%
\pgfpathlineto{\pgfqpoint{2.306374in}{1.710264in}}%
\pgfpathlineto{\pgfqpoint{2.308893in}{1.691287in}}%
\pgfpathlineto{\pgfqpoint{2.311412in}{1.692927in}}%
\pgfpathlineto{\pgfqpoint{2.313931in}{1.694040in}}%
\pgfpathlineto{\pgfqpoint{2.316449in}{1.689294in}}%
\pgfpathlineto{\pgfqpoint{2.318968in}{1.702283in}}%
\pgfpathlineto{\pgfqpoint{2.321487in}{1.675836in}}%
\pgfpathlineto{\pgfqpoint{2.324006in}{1.650523in}}%
\pgfpathlineto{\pgfqpoint{2.326524in}{1.661311in}}%
\pgfpathlineto{\pgfqpoint{2.329043in}{1.654261in}}%
\pgfpathlineto{\pgfqpoint{2.331562in}{1.671774in}}%
\pgfpathlineto{\pgfqpoint{2.334081in}{1.692556in}}%
\pgfpathlineto{\pgfqpoint{2.336599in}{1.685298in}}%
\pgfpathlineto{\pgfqpoint{2.339118in}{1.661008in}}%
\pgfpathlineto{\pgfqpoint{2.341637in}{1.649509in}}%
\pgfpathlineto{\pgfqpoint{2.344156in}{1.634686in}}%
\pgfpathlineto{\pgfqpoint{2.346674in}{1.636749in}}%
\pgfpathlineto{\pgfqpoint{2.349193in}{1.615891in}}%
\pgfpathlineto{\pgfqpoint{2.351712in}{1.620017in}}%
\pgfpathlineto{\pgfqpoint{2.354231in}{1.632803in}}%
\pgfpathlineto{\pgfqpoint{2.356749in}{1.613130in}}%
\pgfpathlineto{\pgfqpoint{2.359268in}{1.622002in}}%
\pgfpathlineto{\pgfqpoint{2.361787in}{1.612468in}}%
\pgfpathlineto{\pgfqpoint{2.364306in}{1.644810in}}%
\pgfpathlineto{\pgfqpoint{2.366824in}{1.661453in}}%
\pgfpathlineto{\pgfqpoint{2.369343in}{1.662004in}}%
\pgfpathlineto{\pgfqpoint{2.371862in}{1.665266in}}%
\pgfpathlineto{\pgfqpoint{2.374381in}{1.671283in}}%
\pgfpathlineto{\pgfqpoint{2.376899in}{1.679560in}}%
\pgfpathlineto{\pgfqpoint{2.379418in}{1.657394in}}%
\pgfpathlineto{\pgfqpoint{2.381937in}{1.640309in}}%
\pgfpathlineto{\pgfqpoint{2.384456in}{1.663394in}}%
\pgfpathlineto{\pgfqpoint{2.386974in}{1.644131in}}%
\pgfpathlineto{\pgfqpoint{2.389493in}{1.645342in}}%
\pgfpathlineto{\pgfqpoint{2.392012in}{1.637294in}}%
\pgfpathlineto{\pgfqpoint{2.394531in}{1.622561in}}%
\pgfpathlineto{\pgfqpoint{2.397049in}{1.622773in}}%
\pgfpathlineto{\pgfqpoint{2.399568in}{1.645232in}}%
\pgfpathlineto{\pgfqpoint{2.402087in}{1.648009in}}%
\pgfpathlineto{\pgfqpoint{2.404606in}{1.634195in}}%
\pgfpathlineto{\pgfqpoint{2.407124in}{1.621433in}}%
\pgfpathlineto{\pgfqpoint{2.409643in}{1.599337in}}%
\pgfpathlineto{\pgfqpoint{2.412162in}{1.603308in}}%
\pgfpathlineto{\pgfqpoint{2.414681in}{1.610955in}}%
\pgfpathlineto{\pgfqpoint{2.417199in}{1.630098in}}%
\pgfpathlineto{\pgfqpoint{2.419718in}{1.643559in}}%
\pgfpathlineto{\pgfqpoint{2.422237in}{1.642136in}}%
\pgfpathlineto{\pgfqpoint{2.424756in}{1.665982in}}%
\pgfpathlineto{\pgfqpoint{2.427274in}{1.662705in}}%
\pgfpathlineto{\pgfqpoint{2.429793in}{1.653928in}}%
\pgfpathlineto{\pgfqpoint{2.432312in}{1.646251in}}%
\pgfpathlineto{\pgfqpoint{2.434831in}{1.658705in}}%
\pgfpathlineto{\pgfqpoint{2.437349in}{1.652807in}}%
\pgfpathlineto{\pgfqpoint{2.439868in}{1.656614in}}%
\pgfpathlineto{\pgfqpoint{2.442387in}{1.670660in}}%
\pgfpathlineto{\pgfqpoint{2.444906in}{1.685955in}}%
\pgfpathlineto{\pgfqpoint{2.447424in}{1.688399in}}%
\pgfpathlineto{\pgfqpoint{2.449943in}{1.688458in}}%
\pgfpathlineto{\pgfqpoint{2.452462in}{1.697539in}}%
\pgfpathlineto{\pgfqpoint{2.454981in}{1.733663in}}%
\pgfpathlineto{\pgfqpoint{2.457499in}{1.746026in}}%
\pgfpathlineto{\pgfqpoint{2.460018in}{1.752009in}}%
\pgfpathlineto{\pgfqpoint{2.462537in}{1.749224in}}%
\pgfpathlineto{\pgfqpoint{2.465056in}{1.726971in}}%
\pgfpathlineto{\pgfqpoint{2.467574in}{1.746863in}}%
\pgfpathlineto{\pgfqpoint{2.470093in}{1.742021in}}%
\pgfpathlineto{\pgfqpoint{2.472612in}{1.745620in}}%
\pgfpathlineto{\pgfqpoint{2.475131in}{1.739439in}}%
\pgfpathlineto{\pgfqpoint{2.477649in}{1.738935in}}%
\pgfpathlineto{\pgfqpoint{2.480168in}{1.721172in}}%
\pgfpathlineto{\pgfqpoint{2.482687in}{1.712884in}}%
\pgfpathlineto{\pgfqpoint{2.485206in}{1.708462in}}%
\pgfpathlineto{\pgfqpoint{2.487724in}{1.700535in}}%
\pgfpathlineto{\pgfqpoint{2.490243in}{1.686789in}}%
\pgfpathlineto{\pgfqpoint{2.492762in}{1.679623in}}%
\pgfpathlineto{\pgfqpoint{2.495281in}{1.663358in}}%
\pgfpathlineto{\pgfqpoint{2.497799in}{1.656940in}}%
\pgfpathlineto{\pgfqpoint{2.500318in}{1.673255in}}%
\pgfpathlineto{\pgfqpoint{2.502837in}{1.684204in}}%
\pgfpathlineto{\pgfqpoint{2.505356in}{1.681145in}}%
\pgfpathlineto{\pgfqpoint{2.507874in}{1.685576in}}%
\pgfpathlineto{\pgfqpoint{2.510393in}{1.646045in}}%
\pgfpathlineto{\pgfqpoint{2.512912in}{1.640833in}}%
\pgfpathlineto{\pgfqpoint{2.515431in}{1.627516in}}%
\pgfpathlineto{\pgfqpoint{2.517949in}{1.610213in}}%
\pgfpathlineto{\pgfqpoint{2.520468in}{1.634240in}}%
\pgfpathlineto{\pgfqpoint{2.522987in}{1.646447in}}%
\pgfpathlineto{\pgfqpoint{2.525506in}{1.629852in}}%
\pgfpathlineto{\pgfqpoint{2.528024in}{1.647972in}}%
\pgfpathlineto{\pgfqpoint{2.530543in}{1.656505in}}%
\pgfpathlineto{\pgfqpoint{2.533062in}{1.658554in}}%
\pgfpathlineto{\pgfqpoint{2.535581in}{1.643906in}}%
\pgfpathlineto{\pgfqpoint{2.538099in}{1.664189in}}%
\pgfpathlineto{\pgfqpoint{2.540618in}{1.646572in}}%
\pgfpathlineto{\pgfqpoint{2.543137in}{1.651189in}}%
\pgfpathlineto{\pgfqpoint{2.545656in}{1.648160in}}%
\pgfpathlineto{\pgfqpoint{2.548174in}{1.662749in}}%
\pgfpathlineto{\pgfqpoint{2.550693in}{1.617686in}}%
\pgfpathlineto{\pgfqpoint{2.553212in}{1.626462in}}%
\pgfpathlineto{\pgfqpoint{2.555731in}{1.613512in}}%
\pgfpathlineto{\pgfqpoint{2.558249in}{1.591533in}}%
\pgfpathlineto{\pgfqpoint{2.560768in}{1.598795in}}%
\pgfpathlineto{\pgfqpoint{2.563287in}{1.582497in}}%
\pgfpathlineto{\pgfqpoint{2.565806in}{1.560611in}}%
\pgfpathlineto{\pgfqpoint{2.568324in}{1.551483in}}%
\pgfpathlineto{\pgfqpoint{2.570843in}{1.546416in}}%
\pgfpathlineto{\pgfqpoint{2.573362in}{1.571000in}}%
\pgfpathlineto{\pgfqpoint{2.575881in}{1.561913in}}%
\pgfpathlineto{\pgfqpoint{2.578399in}{1.585160in}}%
\pgfpathlineto{\pgfqpoint{2.580918in}{1.597303in}}%
\pgfpathlineto{\pgfqpoint{2.583437in}{1.607736in}}%
\pgfpathlineto{\pgfqpoint{2.585956in}{1.617198in}}%
\pgfpathlineto{\pgfqpoint{2.588474in}{1.630941in}}%
\pgfpathlineto{\pgfqpoint{2.590993in}{1.630199in}}%
\pgfpathlineto{\pgfqpoint{2.593512in}{1.594009in}}%
\pgfpathlineto{\pgfqpoint{2.596031in}{1.562916in}}%
\pgfpathlineto{\pgfqpoint{2.598549in}{1.573486in}}%
\pgfpathlineto{\pgfqpoint{2.601068in}{1.575131in}}%
\pgfpathlineto{\pgfqpoint{2.603587in}{1.568350in}}%
\pgfpathlineto{\pgfqpoint{2.606106in}{1.599248in}}%
\pgfpathlineto{\pgfqpoint{2.608624in}{1.615164in}}%
\pgfpathlineto{\pgfqpoint{2.611143in}{1.614047in}}%
\pgfpathlineto{\pgfqpoint{2.613662in}{1.595159in}}%
\pgfpathlineto{\pgfqpoint{2.616181in}{1.595056in}}%
\pgfpathlineto{\pgfqpoint{2.618699in}{1.605231in}}%
\pgfpathlineto{\pgfqpoint{2.621218in}{1.599568in}}%
\pgfpathlineto{\pgfqpoint{2.623737in}{1.587570in}}%
\pgfpathlineto{\pgfqpoint{2.626256in}{1.618964in}}%
\pgfpathlineto{\pgfqpoint{2.628774in}{1.594255in}}%
\pgfpathlineto{\pgfqpoint{2.631293in}{1.583037in}}%
\pgfpathlineto{\pgfqpoint{2.633812in}{1.591113in}}%
\pgfpathlineto{\pgfqpoint{2.636331in}{1.585937in}}%
\pgfpathlineto{\pgfqpoint{2.638849in}{1.590512in}}%
\pgfpathlineto{\pgfqpoint{2.641368in}{1.601070in}}%
\pgfpathlineto{\pgfqpoint{2.643887in}{1.603420in}}%
\pgfpathlineto{\pgfqpoint{2.646406in}{1.621630in}}%
\pgfpathlineto{\pgfqpoint{2.648924in}{1.619653in}}%
\pgfpathlineto{\pgfqpoint{2.651443in}{1.610570in}}%
\pgfpathlineto{\pgfqpoint{2.653962in}{1.621400in}}%
\pgfpathlineto{\pgfqpoint{2.656481in}{1.644968in}}%
\pgfpathlineto{\pgfqpoint{2.658999in}{1.655793in}}%
\pgfpathlineto{\pgfqpoint{2.661518in}{1.656438in}}%
\pgfpathlineto{\pgfqpoint{2.664037in}{1.681958in}}%
\pgfpathlineto{\pgfqpoint{2.666556in}{1.678628in}}%
\pgfpathlineto{\pgfqpoint{2.669074in}{1.692907in}}%
\pgfpathlineto{\pgfqpoint{2.671593in}{1.708436in}}%
\pgfpathlineto{\pgfqpoint{2.674112in}{1.725456in}}%
\pgfpathlineto{\pgfqpoint{2.676631in}{1.717312in}}%
\pgfpathlineto{\pgfqpoint{2.679149in}{1.730749in}}%
\pgfpathlineto{\pgfqpoint{2.681668in}{1.744778in}}%
\pgfpathlineto{\pgfqpoint{2.684187in}{1.746697in}}%
\pgfpathlineto{\pgfqpoint{2.686706in}{1.742697in}}%
\pgfpathlineto{\pgfqpoint{2.689224in}{1.738012in}}%
\pgfpathlineto{\pgfqpoint{2.691743in}{1.753020in}}%
\pgfpathlineto{\pgfqpoint{2.694262in}{1.758645in}}%
\pgfpathlineto{\pgfqpoint{2.696781in}{1.754917in}}%
\pgfpathlineto{\pgfqpoint{2.699299in}{1.769455in}}%
\pgfpathlineto{\pgfqpoint{2.701818in}{1.803655in}}%
\pgfpathlineto{\pgfqpoint{2.704337in}{1.792734in}}%
\pgfpathlineto{\pgfqpoint{2.706856in}{1.812118in}}%
\pgfpathlineto{\pgfqpoint{2.709374in}{1.804532in}}%
\pgfpathlineto{\pgfqpoint{2.711893in}{1.770665in}}%
\pgfpathlineto{\pgfqpoint{2.714412in}{1.763852in}}%
\pgfpathlineto{\pgfqpoint{2.716931in}{1.776228in}}%
\pgfpathlineto{\pgfqpoint{2.719449in}{1.776433in}}%
\pgfpathlineto{\pgfqpoint{2.721968in}{1.805804in}}%
\pgfpathlineto{\pgfqpoint{2.724487in}{1.795288in}}%
\pgfpathlineto{\pgfqpoint{2.727006in}{1.765002in}}%
\pgfpathlineto{\pgfqpoint{2.729524in}{1.785331in}}%
\pgfpathlineto{\pgfqpoint{2.732043in}{1.809264in}}%
\pgfpathlineto{\pgfqpoint{2.734562in}{1.796020in}}%
\pgfpathlineto{\pgfqpoint{2.737081in}{1.774055in}}%
\pgfpathlineto{\pgfqpoint{2.739599in}{1.784507in}}%
\pgfpathlineto{\pgfqpoint{2.742118in}{1.766458in}}%
\pgfpathlineto{\pgfqpoint{2.744637in}{1.765448in}}%
\pgfpathlineto{\pgfqpoint{2.747156in}{1.767123in}}%
\pgfpathlineto{\pgfqpoint{2.749674in}{1.756932in}}%
\pgfpathlineto{\pgfqpoint{2.752193in}{1.733809in}}%
\pgfpathlineto{\pgfqpoint{2.754712in}{1.708513in}}%
\pgfpathlineto{\pgfqpoint{2.757231in}{1.715181in}}%
\pgfpathlineto{\pgfqpoint{2.759749in}{1.695544in}}%
\pgfpathlineto{\pgfqpoint{2.762268in}{1.702889in}}%
\pgfpathlineto{\pgfqpoint{2.764787in}{1.718399in}}%
\pgfpathlineto{\pgfqpoint{2.767306in}{1.741437in}}%
\pgfpathlineto{\pgfqpoint{2.769824in}{1.725690in}}%
\pgfpathlineto{\pgfqpoint{2.772343in}{1.725433in}}%
\pgfpathlineto{\pgfqpoint{2.774862in}{1.703513in}}%
\pgfpathlineto{\pgfqpoint{2.777381in}{1.686154in}}%
\pgfpathlineto{\pgfqpoint{2.779899in}{1.664388in}}%
\pgfpathlineto{\pgfqpoint{2.782418in}{1.656620in}}%
\pgfpathlineto{\pgfqpoint{2.784937in}{1.656719in}}%
\pgfpathlineto{\pgfqpoint{2.787456in}{1.660115in}}%
\pgfpathlineto{\pgfqpoint{2.789974in}{1.690022in}}%
\pgfpathlineto{\pgfqpoint{2.792493in}{1.684318in}}%
\pgfpathlineto{\pgfqpoint{2.795012in}{1.707620in}}%
\pgfpathlineto{\pgfqpoint{2.797531in}{1.698040in}}%
\pgfpathlineto{\pgfqpoint{2.800049in}{1.679482in}}%
\pgfpathlineto{\pgfqpoint{2.802568in}{1.686397in}}%
\pgfpathlineto{\pgfqpoint{2.805087in}{1.678150in}}%
\pgfpathlineto{\pgfqpoint{2.807606in}{1.683037in}}%
\pgfpathlineto{\pgfqpoint{2.810124in}{1.679620in}}%
\pgfpathlineto{\pgfqpoint{2.812643in}{1.666745in}}%
\pgfpathlineto{\pgfqpoint{2.815162in}{1.655017in}}%
\pgfpathlineto{\pgfqpoint{2.817681in}{1.637000in}}%
\pgfpathlineto{\pgfqpoint{2.820199in}{1.620413in}}%
\pgfpathlineto{\pgfqpoint{2.822718in}{1.635855in}}%
\pgfpathlineto{\pgfqpoint{2.825237in}{1.590530in}}%
\pgfpathlineto{\pgfqpoint{2.827756in}{1.613968in}}%
\pgfpathlineto{\pgfqpoint{2.830274in}{1.602144in}}%
\pgfpathlineto{\pgfqpoint{2.832793in}{1.610468in}}%
\pgfpathlineto{\pgfqpoint{2.835312in}{1.610309in}}%
\pgfpathlineto{\pgfqpoint{2.837831in}{1.612447in}}%
\pgfpathlineto{\pgfqpoint{2.840349in}{1.601260in}}%
\pgfpathlineto{\pgfqpoint{2.842868in}{1.612867in}}%
\pgfpathlineto{\pgfqpoint{2.845387in}{1.602337in}}%
\pgfpathlineto{\pgfqpoint{2.847906in}{1.592294in}}%
\pgfpathlineto{\pgfqpoint{2.850424in}{1.561853in}}%
\pgfpathlineto{\pgfqpoint{2.852943in}{1.556782in}}%
\pgfpathlineto{\pgfqpoint{2.855462in}{1.555081in}}%
\pgfpathlineto{\pgfqpoint{2.857981in}{1.543590in}}%
\pgfpathlineto{\pgfqpoint{2.860499in}{1.548657in}}%
\pgfpathlineto{\pgfqpoint{2.863018in}{1.552854in}}%
\pgfpathlineto{\pgfqpoint{2.865537in}{1.546304in}}%
\pgfpathlineto{\pgfqpoint{2.868056in}{1.564057in}}%
\pgfpathlineto{\pgfqpoint{2.870574in}{1.559837in}}%
\pgfpathlineto{\pgfqpoint{2.873093in}{1.559947in}}%
\pgfpathlineto{\pgfqpoint{2.875612in}{1.608695in}}%
\pgfpathlineto{\pgfqpoint{2.878131in}{1.580000in}}%
\pgfpathlineto{\pgfqpoint{2.880649in}{1.565629in}}%
\pgfpathlineto{\pgfqpoint{2.883168in}{1.570235in}}%
\pgfpathlineto{\pgfqpoint{2.885687in}{1.546835in}}%
\pgfpathlineto{\pgfqpoint{2.888206in}{1.575279in}}%
\pgfpathlineto{\pgfqpoint{2.890724in}{1.572500in}}%
\pgfpathlineto{\pgfqpoint{2.893243in}{1.577587in}}%
\pgfpathlineto{\pgfqpoint{2.895762in}{1.574047in}}%
\pgfpathlineto{\pgfqpoint{2.898281in}{1.585601in}}%
\pgfpathlineto{\pgfqpoint{2.903318in}{1.595496in}}%
\pgfpathlineto{\pgfqpoint{2.905837in}{1.628127in}}%
\pgfpathlineto{\pgfqpoint{2.908356in}{1.616707in}}%
\pgfpathlineto{\pgfqpoint{2.910874in}{1.630472in}}%
\pgfpathlineto{\pgfqpoint{2.913393in}{1.636967in}}%
\pgfpathlineto{\pgfqpoint{2.915912in}{1.640279in}}%
\pgfpathlineto{\pgfqpoint{2.918431in}{1.660001in}}%
\pgfpathlineto{\pgfqpoint{2.920949in}{1.630679in}}%
\pgfpathlineto{\pgfqpoint{2.925987in}{1.683225in}}%
\pgfpathlineto{\pgfqpoint{2.928506in}{1.707194in}}%
\pgfpathlineto{\pgfqpoint{2.931024in}{1.680904in}}%
\pgfpathlineto{\pgfqpoint{2.933543in}{1.691379in}}%
\pgfpathlineto{\pgfqpoint{2.936062in}{1.686834in}}%
\pgfpathlineto{\pgfqpoint{2.938581in}{1.667630in}}%
\pgfpathlineto{\pgfqpoint{2.941099in}{1.651692in}}%
\pgfpathlineto{\pgfqpoint{2.943618in}{1.672173in}}%
\pgfpathlineto{\pgfqpoint{2.946137in}{1.667277in}}%
\pgfpathlineto{\pgfqpoint{2.948656in}{1.656008in}}%
\pgfpathlineto{\pgfqpoint{2.951174in}{1.615984in}}%
\pgfpathlineto{\pgfqpoint{2.953693in}{1.601612in}}%
\pgfpathlineto{\pgfqpoint{2.956212in}{1.602053in}}%
\pgfpathlineto{\pgfqpoint{2.958731in}{1.585762in}}%
\pgfpathlineto{\pgfqpoint{2.961249in}{1.589814in}}%
\pgfpathlineto{\pgfqpoint{2.963768in}{1.611534in}}%
\pgfpathlineto{\pgfqpoint{2.966287in}{1.609328in}}%
\pgfpathlineto{\pgfqpoint{2.968806in}{1.598187in}}%
\pgfpathlineto{\pgfqpoint{2.971324in}{1.561771in}}%
\pgfpathlineto{\pgfqpoint{2.973843in}{1.545921in}}%
\pgfpathlineto{\pgfqpoint{2.976362in}{1.552501in}}%
\pgfpathlineto{\pgfqpoint{2.978881in}{1.546844in}}%
\pgfpathlineto{\pgfqpoint{2.981399in}{1.540895in}}%
\pgfpathlineto{\pgfqpoint{2.983918in}{1.541714in}}%
\pgfpathlineto{\pgfqpoint{2.986437in}{1.533684in}}%
\pgfpathlineto{\pgfqpoint{2.988956in}{1.550709in}}%
\pgfpathlineto{\pgfqpoint{2.991474in}{1.568414in}}%
\pgfpathlineto{\pgfqpoint{2.993993in}{1.562888in}}%
\pgfpathlineto{\pgfqpoint{2.996512in}{1.587221in}}%
\pgfpathlineto{\pgfqpoint{2.999031in}{1.575825in}}%
\pgfpathlineto{\pgfqpoint{3.001549in}{1.589929in}}%
\pgfpathlineto{\pgfqpoint{3.004068in}{1.605528in}}%
\pgfpathlineto{\pgfqpoint{3.006587in}{1.599902in}}%
\pgfpathlineto{\pgfqpoint{3.009106in}{1.625856in}}%
\pgfpathlineto{\pgfqpoint{3.011624in}{1.626843in}}%
\pgfpathlineto{\pgfqpoint{3.014143in}{1.643475in}}%
\pgfpathlineto{\pgfqpoint{3.016662in}{1.674580in}}%
\pgfpathlineto{\pgfqpoint{3.019181in}{1.650860in}}%
\pgfpathlineto{\pgfqpoint{3.021699in}{1.619681in}}%
\pgfpathlineto{\pgfqpoint{3.026737in}{1.584978in}}%
\pgfpathlineto{\pgfqpoint{3.029256in}{1.578286in}}%
\pgfpathlineto{\pgfqpoint{3.031774in}{1.563415in}}%
\pgfpathlineto{\pgfqpoint{3.034293in}{1.574829in}}%
\pgfpathlineto{\pgfqpoint{3.036812in}{1.552771in}}%
\pgfpathlineto{\pgfqpoint{3.039331in}{1.536060in}}%
\pgfpathlineto{\pgfqpoint{3.041849in}{1.539472in}}%
\pgfpathlineto{\pgfqpoint{3.044368in}{1.541246in}}%
\pgfpathlineto{\pgfqpoint{3.046887in}{1.565512in}}%
\pgfpathlineto{\pgfqpoint{3.049406in}{1.545081in}}%
\pgfpathlineto{\pgfqpoint{3.051924in}{1.552280in}}%
\pgfpathlineto{\pgfqpoint{3.054443in}{1.539149in}}%
\pgfpathlineto{\pgfqpoint{3.056962in}{1.528641in}}%
\pgfpathlineto{\pgfqpoint{3.059481in}{1.533638in}}%
\pgfpathlineto{\pgfqpoint{3.061999in}{1.504839in}}%
\pgfpathlineto{\pgfqpoint{3.064518in}{1.498297in}}%
\pgfpathlineto{\pgfqpoint{3.067037in}{1.519149in}}%
\pgfpathlineto{\pgfqpoint{3.069556in}{1.495876in}}%
\pgfpathlineto{\pgfqpoint{3.072074in}{1.480114in}}%
\pgfpathlineto{\pgfqpoint{3.074593in}{1.477845in}}%
\pgfpathlineto{\pgfqpoint{3.077112in}{1.467822in}}%
\pgfpathlineto{\pgfqpoint{3.079631in}{1.465723in}}%
\pgfpathlineto{\pgfqpoint{3.082149in}{1.476470in}}%
\pgfpathlineto{\pgfqpoint{3.084668in}{1.462000in}}%
\pgfpathlineto{\pgfqpoint{3.087187in}{1.452773in}}%
\pgfpathlineto{\pgfqpoint{3.089706in}{1.439779in}}%
\pgfpathlineto{\pgfqpoint{3.092224in}{1.443915in}}%
\pgfpathlineto{\pgfqpoint{3.094743in}{1.434744in}}%
\pgfpathlineto{\pgfqpoint{3.097262in}{1.418132in}}%
\pgfpathlineto{\pgfqpoint{3.099781in}{1.420572in}}%
\pgfpathlineto{\pgfqpoint{3.102299in}{1.433206in}}%
\pgfpathlineto{\pgfqpoint{3.104818in}{1.446999in}}%
\pgfpathlineto{\pgfqpoint{3.109856in}{1.431769in}}%
\pgfpathlineto{\pgfqpoint{3.112374in}{1.455403in}}%
\pgfpathlineto{\pgfqpoint{3.114893in}{1.483637in}}%
\pgfpathlineto{\pgfqpoint{3.117412in}{1.475591in}}%
\pgfpathlineto{\pgfqpoint{3.119931in}{1.464526in}}%
\pgfpathlineto{\pgfqpoint{3.122449in}{1.452559in}}%
\pgfpathlineto{\pgfqpoint{3.124968in}{1.456612in}}%
\pgfpathlineto{\pgfqpoint{3.127487in}{1.437092in}}%
\pgfpathlineto{\pgfqpoint{3.130006in}{1.439515in}}%
\pgfpathlineto{\pgfqpoint{3.132524in}{1.448879in}}%
\pgfpathlineto{\pgfqpoint{3.135043in}{1.437500in}}%
\pgfpathlineto{\pgfqpoint{3.137562in}{1.448764in}}%
\pgfpathlineto{\pgfqpoint{3.140081in}{1.441290in}}%
\pgfpathlineto{\pgfqpoint{3.142599in}{1.438613in}}%
\pgfpathlineto{\pgfqpoint{3.145118in}{1.458773in}}%
\pgfpathlineto{\pgfqpoint{3.147637in}{1.450495in}}%
\pgfpathlineto{\pgfqpoint{3.150156in}{1.461078in}}%
\pgfpathlineto{\pgfqpoint{3.152674in}{1.470486in}}%
\pgfpathlineto{\pgfqpoint{3.155193in}{1.452642in}}%
\pgfpathlineto{\pgfqpoint{3.157712in}{1.465624in}}%
\pgfpathlineto{\pgfqpoint{3.160231in}{1.464974in}}%
\pgfpathlineto{\pgfqpoint{3.162749in}{1.463734in}}%
\pgfpathlineto{\pgfqpoint{3.165268in}{1.460388in}}%
\pgfpathlineto{\pgfqpoint{3.167787in}{1.447640in}}%
\pgfpathlineto{\pgfqpoint{3.170306in}{1.455057in}}%
\pgfpathlineto{\pgfqpoint{3.172824in}{1.441138in}}%
\pgfpathlineto{\pgfqpoint{3.175343in}{1.438277in}}%
\pgfpathlineto{\pgfqpoint{3.177862in}{1.457787in}}%
\pgfpathlineto{\pgfqpoint{3.180381in}{1.450908in}}%
\pgfpathlineto{\pgfqpoint{3.182899in}{1.428894in}}%
\pgfpathlineto{\pgfqpoint{3.185418in}{1.457531in}}%
\pgfpathlineto{\pgfqpoint{3.187937in}{1.465574in}}%
\pgfpathlineto{\pgfqpoint{3.190456in}{1.494673in}}%
\pgfpathlineto{\pgfqpoint{3.192974in}{1.514724in}}%
\pgfpathlineto{\pgfqpoint{3.195493in}{1.503072in}}%
\pgfpathlineto{\pgfqpoint{3.198012in}{1.508969in}}%
\pgfpathlineto{\pgfqpoint{3.200531in}{1.505603in}}%
\pgfpathlineto{\pgfqpoint{3.203049in}{1.500479in}}%
\pgfpathlineto{\pgfqpoint{3.205568in}{1.502348in}}%
\pgfpathlineto{\pgfqpoint{3.208087in}{1.512518in}}%
\pgfpathlineto{\pgfqpoint{3.210606in}{1.506074in}}%
\pgfpathlineto{\pgfqpoint{3.213124in}{1.508857in}}%
\pgfpathlineto{\pgfqpoint{3.215643in}{1.516093in}}%
\pgfpathlineto{\pgfqpoint{3.218162in}{1.514683in}}%
\pgfpathlineto{\pgfqpoint{3.220681in}{1.529589in}}%
\pgfpathlineto{\pgfqpoint{3.223199in}{1.535808in}}%
\pgfpathlineto{\pgfqpoint{3.225718in}{1.520741in}}%
\pgfpathlineto{\pgfqpoint{3.228237in}{1.524923in}}%
\pgfpathlineto{\pgfqpoint{3.230756in}{1.529507in}}%
\pgfpathlineto{\pgfqpoint{3.233274in}{1.576519in}}%
\pgfpathlineto{\pgfqpoint{3.235793in}{1.581513in}}%
\pgfpathlineto{\pgfqpoint{3.238312in}{1.574580in}}%
\pgfpathlineto{\pgfqpoint{3.240831in}{1.577479in}}%
\pgfpathlineto{\pgfqpoint{3.243349in}{1.566869in}}%
\pgfpathlineto{\pgfqpoint{3.245868in}{1.574667in}}%
\pgfpathlineto{\pgfqpoint{3.248387in}{1.563668in}}%
\pgfpathlineto{\pgfqpoint{3.250906in}{1.555521in}}%
\pgfpathlineto{\pgfqpoint{3.253424in}{1.567185in}}%
\pgfpathlineto{\pgfqpoint{3.255943in}{1.580292in}}%
\pgfpathlineto{\pgfqpoint{3.258462in}{1.591308in}}%
\pgfpathlineto{\pgfqpoint{3.260981in}{1.589121in}}%
\pgfpathlineto{\pgfqpoint{3.263499in}{1.566583in}}%
\pgfpathlineto{\pgfqpoint{3.266018in}{1.588719in}}%
\pgfpathlineto{\pgfqpoint{3.268537in}{1.606132in}}%
\pgfpathlineto{\pgfqpoint{3.271056in}{1.629106in}}%
\pgfpathlineto{\pgfqpoint{3.276093in}{1.591166in}}%
\pgfpathlineto{\pgfqpoint{3.278612in}{1.566456in}}%
\pgfpathlineto{\pgfqpoint{3.281131in}{1.552232in}}%
\pgfpathlineto{\pgfqpoint{3.283649in}{1.582695in}}%
\pgfpathlineto{\pgfqpoint{3.286168in}{1.571190in}}%
\pgfpathlineto{\pgfqpoint{3.288687in}{1.598028in}}%
\pgfpathlineto{\pgfqpoint{3.291206in}{1.641558in}}%
\pgfpathlineto{\pgfqpoint{3.293724in}{1.651769in}}%
\pgfpathlineto{\pgfqpoint{3.296243in}{1.635419in}}%
\pgfpathlineto{\pgfqpoint{3.298762in}{1.633069in}}%
\pgfpathlineto{\pgfqpoint{3.301281in}{1.641986in}}%
\pgfpathlineto{\pgfqpoint{3.303799in}{1.630212in}}%
\pgfpathlineto{\pgfqpoint{3.306318in}{1.609674in}}%
\pgfpathlineto{\pgfqpoint{3.308837in}{1.607608in}}%
\pgfpathlineto{\pgfqpoint{3.311356in}{1.643868in}}%
\pgfpathlineto{\pgfqpoint{3.313874in}{1.654431in}}%
\pgfpathlineto{\pgfqpoint{3.316393in}{1.678102in}}%
\pgfpathlineto{\pgfqpoint{3.318912in}{1.667850in}}%
\pgfpathlineto{\pgfqpoint{3.321431in}{1.678084in}}%
\pgfpathlineto{\pgfqpoint{3.323949in}{1.691085in}}%
\pgfpathlineto{\pgfqpoint{3.326468in}{1.663510in}}%
\pgfpathlineto{\pgfqpoint{3.328987in}{1.664890in}}%
\pgfpathlineto{\pgfqpoint{3.331506in}{1.668937in}}%
\pgfpathlineto{\pgfqpoint{3.334024in}{1.686338in}}%
\pgfpathlineto{\pgfqpoint{3.336543in}{1.691637in}}%
\pgfpathlineto{\pgfqpoint{3.339062in}{1.708006in}}%
\pgfpathlineto{\pgfqpoint{3.341581in}{1.703388in}}%
\pgfpathlineto{\pgfqpoint{3.344099in}{1.710617in}}%
\pgfpathlineto{\pgfqpoint{3.346618in}{1.705652in}}%
\pgfpathlineto{\pgfqpoint{3.349137in}{1.717534in}}%
\pgfpathlineto{\pgfqpoint{3.351656in}{1.707264in}}%
\pgfpathlineto{\pgfqpoint{3.354174in}{1.689444in}}%
\pgfpathlineto{\pgfqpoint{3.356693in}{1.701511in}}%
\pgfpathlineto{\pgfqpoint{3.359212in}{1.726716in}}%
\pgfpathlineto{\pgfqpoint{3.361731in}{1.726315in}}%
\pgfpathlineto{\pgfqpoint{3.364249in}{1.717565in}}%
\pgfpathlineto{\pgfqpoint{3.366768in}{1.680492in}}%
\pgfpathlineto{\pgfqpoint{3.369287in}{1.675543in}}%
\pgfpathlineto{\pgfqpoint{3.371806in}{1.696587in}}%
\pgfpathlineto{\pgfqpoint{3.374324in}{1.727665in}}%
\pgfpathlineto{\pgfqpoint{3.376843in}{1.732998in}}%
\pgfpathlineto{\pgfqpoint{3.379362in}{1.751472in}}%
\pgfpathlineto{\pgfqpoint{3.381881in}{1.719440in}}%
\pgfpathlineto{\pgfqpoint{3.384399in}{1.728804in}}%
\pgfpathlineto{\pgfqpoint{3.386918in}{1.735878in}}%
\pgfpathlineto{\pgfqpoint{3.389437in}{1.746087in}}%
\pgfpathlineto{\pgfqpoint{3.391956in}{1.764715in}}%
\pgfpathlineto{\pgfqpoint{3.394474in}{1.747694in}}%
\pgfpathlineto{\pgfqpoint{3.396993in}{1.760616in}}%
\pgfpathlineto{\pgfqpoint{3.399512in}{1.764240in}}%
\pgfpathlineto{\pgfqpoint{3.402031in}{1.780709in}}%
\pgfpathlineto{\pgfqpoint{3.404549in}{1.794922in}}%
\pgfpathlineto{\pgfqpoint{3.407068in}{1.774109in}}%
\pgfpathlineto{\pgfqpoint{3.409587in}{1.759460in}}%
\pgfpathlineto{\pgfqpoint{3.412106in}{1.781845in}}%
\pgfpathlineto{\pgfqpoint{3.414624in}{1.808652in}}%
\pgfpathlineto{\pgfqpoint{3.417143in}{1.843994in}}%
\pgfpathlineto{\pgfqpoint{3.419662in}{1.859652in}}%
\pgfpathlineto{\pgfqpoint{3.422181in}{1.850540in}}%
\pgfpathlineto{\pgfqpoint{3.424699in}{1.846095in}}%
\pgfpathlineto{\pgfqpoint{3.427218in}{1.850159in}}%
\pgfpathlineto{\pgfqpoint{3.429737in}{1.869148in}}%
\pgfpathlineto{\pgfqpoint{3.432256in}{1.861603in}}%
\pgfpathlineto{\pgfqpoint{3.434774in}{1.854761in}}%
\pgfpathlineto{\pgfqpoint{3.437293in}{1.870366in}}%
\pgfpathlineto{\pgfqpoint{3.439812in}{1.871540in}}%
\pgfpathlineto{\pgfqpoint{3.442331in}{1.864921in}}%
\pgfpathlineto{\pgfqpoint{3.444849in}{1.848853in}}%
\pgfpathlineto{\pgfqpoint{3.447368in}{1.837771in}}%
\pgfpathlineto{\pgfqpoint{3.449887in}{1.862547in}}%
\pgfpathlineto{\pgfqpoint{3.452406in}{1.838104in}}%
\pgfpathlineto{\pgfqpoint{3.454924in}{1.850443in}}%
\pgfpathlineto{\pgfqpoint{3.457443in}{1.868823in}}%
\pgfpathlineto{\pgfqpoint{3.459962in}{1.860581in}}%
\pgfpathlineto{\pgfqpoint{3.462481in}{1.891254in}}%
\pgfpathlineto{\pgfqpoint{3.464999in}{1.896433in}}%
\pgfpathlineto{\pgfqpoint{3.467518in}{1.898909in}}%
\pgfpathlineto{\pgfqpoint{3.470037in}{1.927461in}}%
\pgfpathlineto{\pgfqpoint{3.472556in}{1.910545in}}%
\pgfpathlineto{\pgfqpoint{3.475074in}{1.931360in}}%
\pgfpathlineto{\pgfqpoint{3.477593in}{1.934433in}}%
\pgfpathlineto{\pgfqpoint{3.480112in}{1.942032in}}%
\pgfpathlineto{\pgfqpoint{3.482631in}{1.925901in}}%
\pgfpathlineto{\pgfqpoint{3.485149in}{1.933868in}}%
\pgfpathlineto{\pgfqpoint{3.487668in}{1.938135in}}%
\pgfpathlineto{\pgfqpoint{3.490187in}{1.944018in}}%
\pgfpathlineto{\pgfqpoint{3.492706in}{1.973718in}}%
\pgfpathlineto{\pgfqpoint{3.495224in}{1.967686in}}%
\pgfpathlineto{\pgfqpoint{3.497743in}{1.932049in}}%
\pgfpathlineto{\pgfqpoint{3.500262in}{1.943085in}}%
\pgfpathlineto{\pgfqpoint{3.502781in}{1.942491in}}%
\pgfpathlineto{\pgfqpoint{3.505299in}{1.966651in}}%
\pgfpathlineto{\pgfqpoint{3.507818in}{1.951820in}}%
\pgfpathlineto{\pgfqpoint{3.510337in}{1.970933in}}%
\pgfpathlineto{\pgfqpoint{3.512856in}{1.987411in}}%
\pgfpathlineto{\pgfqpoint{3.515374in}{1.986212in}}%
\pgfpathlineto{\pgfqpoint{3.517893in}{1.938254in}}%
\pgfpathlineto{\pgfqpoint{3.520412in}{1.941163in}}%
\pgfpathlineto{\pgfqpoint{3.522931in}{1.905252in}}%
\pgfpathlineto{\pgfqpoint{3.525449in}{1.879985in}}%
\pgfpathlineto{\pgfqpoint{3.527968in}{1.890673in}}%
\pgfpathlineto{\pgfqpoint{3.530487in}{1.911605in}}%
\pgfpathlineto{\pgfqpoint{3.533006in}{1.894346in}}%
\pgfpathlineto{\pgfqpoint{3.535524in}{1.889127in}}%
\pgfpathlineto{\pgfqpoint{3.538043in}{1.872852in}}%
\pgfpathlineto{\pgfqpoint{3.540562in}{1.890552in}}%
\pgfpathlineto{\pgfqpoint{3.543081in}{1.881411in}}%
\pgfpathlineto{\pgfqpoint{3.545599in}{1.870589in}}%
\pgfpathlineto{\pgfqpoint{3.548118in}{1.889780in}}%
\pgfpathlineto{\pgfqpoint{3.550637in}{1.867336in}}%
\pgfpathlineto{\pgfqpoint{3.553156in}{1.899166in}}%
\pgfpathlineto{\pgfqpoint{3.555674in}{1.908408in}}%
\pgfpathlineto{\pgfqpoint{3.558193in}{1.896955in}}%
\pgfpathlineto{\pgfqpoint{3.560712in}{1.911247in}}%
\pgfpathlineto{\pgfqpoint{3.563231in}{1.882814in}}%
\pgfpathlineto{\pgfqpoint{3.565749in}{1.904430in}}%
\pgfpathlineto{\pgfqpoint{3.568268in}{1.919495in}}%
\pgfpathlineto{\pgfqpoint{3.570787in}{1.941636in}}%
\pgfpathlineto{\pgfqpoint{3.573306in}{1.942464in}}%
\pgfpathlineto{\pgfqpoint{3.575824in}{1.944561in}}%
\pgfpathlineto{\pgfqpoint{3.578343in}{1.935720in}}%
\pgfpathlineto{\pgfqpoint{3.580862in}{1.920691in}}%
\pgfpathlineto{\pgfqpoint{3.583381in}{1.912892in}}%
\pgfpathlineto{\pgfqpoint{3.585899in}{1.940566in}}%
\pgfpathlineto{\pgfqpoint{3.588418in}{1.915935in}}%
\pgfpathlineto{\pgfqpoint{3.590937in}{1.917689in}}%
\pgfpathlineto{\pgfqpoint{3.593456in}{1.934356in}}%
\pgfpathlineto{\pgfqpoint{3.595974in}{1.947868in}}%
\pgfpathlineto{\pgfqpoint{3.598493in}{1.935362in}}%
\pgfpathlineto{\pgfqpoint{3.601012in}{1.939093in}}%
\pgfpathlineto{\pgfqpoint{3.603531in}{1.928141in}}%
\pgfpathlineto{\pgfqpoint{3.606049in}{1.943949in}}%
\pgfpathlineto{\pgfqpoint{3.608568in}{1.984256in}}%
\pgfpathlineto{\pgfqpoint{3.611087in}{1.986781in}}%
\pgfpathlineto{\pgfqpoint{3.613606in}{1.983009in}}%
\pgfpathlineto{\pgfqpoint{3.616124in}{1.958165in}}%
\pgfpathlineto{\pgfqpoint{3.618643in}{1.969564in}}%
\pgfpathlineto{\pgfqpoint{3.621162in}{1.960310in}}%
\pgfpathlineto{\pgfqpoint{3.623681in}{1.986189in}}%
\pgfpathlineto{\pgfqpoint{3.626199in}{2.004520in}}%
\pgfpathlineto{\pgfqpoint{3.628718in}{1.988373in}}%
\pgfpathlineto{\pgfqpoint{3.631237in}{1.969750in}}%
\pgfpathlineto{\pgfqpoint{3.633756in}{1.989728in}}%
\pgfpathlineto{\pgfqpoint{3.636274in}{2.007668in}}%
\pgfpathlineto{\pgfqpoint{3.638793in}{2.016289in}}%
\pgfpathlineto{\pgfqpoint{3.641312in}{2.029245in}}%
\pgfpathlineto{\pgfqpoint{3.643831in}{2.049814in}}%
\pgfpathlineto{\pgfqpoint{3.646349in}{2.020022in}}%
\pgfpathlineto{\pgfqpoint{3.648868in}{2.034161in}}%
\pgfpathlineto{\pgfqpoint{3.651387in}{2.033205in}}%
\pgfpathlineto{\pgfqpoint{3.653906in}{2.039769in}}%
\pgfpathlineto{\pgfqpoint{3.656424in}{2.072654in}}%
\pgfpathlineto{\pgfqpoint{3.658943in}{2.061376in}}%
\pgfpathlineto{\pgfqpoint{3.661462in}{2.038958in}}%
\pgfpathlineto{\pgfqpoint{3.663981in}{2.054659in}}%
\pgfpathlineto{\pgfqpoint{3.666499in}{2.057848in}}%
\pgfpathlineto{\pgfqpoint{3.669018in}{2.048651in}}%
\pgfpathlineto{\pgfqpoint{3.671537in}{2.031182in}}%
\pgfpathlineto{\pgfqpoint{3.674056in}{2.026732in}}%
\pgfpathlineto{\pgfqpoint{3.676574in}{2.003294in}}%
\pgfpathlineto{\pgfqpoint{3.679093in}{2.048566in}}%
\pgfpathlineto{\pgfqpoint{3.681612in}{2.022360in}}%
\pgfpathlineto{\pgfqpoint{3.684131in}{2.024580in}}%
\pgfpathlineto{\pgfqpoint{3.686649in}{2.015417in}}%
\pgfpathlineto{\pgfqpoint{3.689168in}{2.036763in}}%
\pgfpathlineto{\pgfqpoint{3.691687in}{2.043470in}}%
\pgfpathlineto{\pgfqpoint{3.694206in}{2.041432in}}%
\pgfpathlineto{\pgfqpoint{3.696724in}{2.032547in}}%
\pgfpathlineto{\pgfqpoint{3.699243in}{2.053039in}}%
\pgfpathlineto{\pgfqpoint{3.701762in}{2.044560in}}%
\pgfpathlineto{\pgfqpoint{3.704281in}{2.040534in}}%
\pgfpathlineto{\pgfqpoint{3.706799in}{2.040405in}}%
\pgfpathlineto{\pgfqpoint{3.709318in}{2.032139in}}%
\pgfpathlineto{\pgfqpoint{3.711837in}{2.034782in}}%
\pgfpathlineto{\pgfqpoint{3.714356in}{2.017699in}}%
\pgfpathlineto{\pgfqpoint{3.716874in}{2.005150in}}%
\pgfpathlineto{\pgfqpoint{3.719393in}{1.987885in}}%
\pgfpathlineto{\pgfqpoint{3.721912in}{1.967213in}}%
\pgfpathlineto{\pgfqpoint{3.724431in}{1.968042in}}%
\pgfpathlineto{\pgfqpoint{3.726949in}{1.982813in}}%
\pgfpathlineto{\pgfqpoint{3.729468in}{1.983893in}}%
\pgfpathlineto{\pgfqpoint{3.731987in}{1.967603in}}%
\pgfpathlineto{\pgfqpoint{3.734506in}{1.985640in}}%
\pgfpathlineto{\pgfqpoint{3.737024in}{1.992417in}}%
\pgfpathlineto{\pgfqpoint{3.739543in}{2.003584in}}%
\pgfpathlineto{\pgfqpoint{3.742062in}{1.993452in}}%
\pgfpathlineto{\pgfqpoint{3.744581in}{1.969570in}}%
\pgfpathlineto{\pgfqpoint{3.747099in}{1.967792in}}%
\pgfpathlineto{\pgfqpoint{3.749618in}{1.973050in}}%
\pgfpathlineto{\pgfqpoint{3.752137in}{1.964422in}}%
\pgfpathlineto{\pgfqpoint{3.754656in}{1.956429in}}%
\pgfpathlineto{\pgfqpoint{3.757174in}{1.947775in}}%
\pgfpathlineto{\pgfqpoint{3.759693in}{1.974763in}}%
\pgfpathlineto{\pgfqpoint{3.762212in}{1.976691in}}%
\pgfpathlineto{\pgfqpoint{3.764731in}{1.967895in}}%
\pgfpathlineto{\pgfqpoint{3.767249in}{1.978836in}}%
\pgfpathlineto{\pgfqpoint{3.769768in}{1.973062in}}%
\pgfpathlineto{\pgfqpoint{3.772287in}{1.958151in}}%
\pgfpathlineto{\pgfqpoint{3.774806in}{1.952572in}}%
\pgfpathlineto{\pgfqpoint{3.777324in}{1.984581in}}%
\pgfpathlineto{\pgfqpoint{3.779843in}{1.979727in}}%
\pgfpathlineto{\pgfqpoint{3.782362in}{1.999100in}}%
\pgfpathlineto{\pgfqpoint{3.784881in}{2.016723in}}%
\pgfpathlineto{\pgfqpoint{3.787399in}{2.015327in}}%
\pgfpathlineto{\pgfqpoint{3.789918in}{2.014383in}}%
\pgfpathlineto{\pgfqpoint{3.792437in}{2.018548in}}%
\pgfpathlineto{\pgfqpoint{3.794956in}{2.025892in}}%
\pgfpathlineto{\pgfqpoint{3.797474in}{2.046877in}}%
\pgfpathlineto{\pgfqpoint{3.799993in}{2.061724in}}%
\pgfpathlineto{\pgfqpoint{3.802512in}{2.087195in}}%
\pgfpathlineto{\pgfqpoint{3.805031in}{2.062038in}}%
\pgfpathlineto{\pgfqpoint{3.810068in}{2.108160in}}%
\pgfpathlineto{\pgfqpoint{3.812587in}{2.109142in}}%
\pgfpathlineto{\pgfqpoint{3.815106in}{2.112940in}}%
\pgfpathlineto{\pgfqpoint{3.817624in}{2.141171in}}%
\pgfpathlineto{\pgfqpoint{3.820143in}{2.126938in}}%
\pgfpathlineto{\pgfqpoint{3.822662in}{2.126506in}}%
\pgfpathlineto{\pgfqpoint{3.825181in}{2.137076in}}%
\pgfpathlineto{\pgfqpoint{3.827699in}{2.144465in}}%
\pgfpathlineto{\pgfqpoint{3.830218in}{2.129328in}}%
\pgfpathlineto{\pgfqpoint{3.832737in}{2.131903in}}%
\pgfpathlineto{\pgfqpoint{3.835256in}{2.105249in}}%
\pgfpathlineto{\pgfqpoint{3.837774in}{2.125982in}}%
\pgfpathlineto{\pgfqpoint{3.840293in}{2.102153in}}%
\pgfpathlineto{\pgfqpoint{3.842812in}{2.093734in}}%
\pgfpathlineto{\pgfqpoint{3.845331in}{2.097043in}}%
\pgfpathlineto{\pgfqpoint{3.847849in}{2.087090in}}%
\pgfpathlineto{\pgfqpoint{3.850368in}{2.069068in}}%
\pgfpathlineto{\pgfqpoint{3.855406in}{2.061498in}}%
\pgfpathlineto{\pgfqpoint{3.857924in}{2.040869in}}%
\pgfpathlineto{\pgfqpoint{3.860443in}{2.017469in}}%
\pgfpathlineto{\pgfqpoint{3.862962in}{2.035857in}}%
\pgfpathlineto{\pgfqpoint{3.865481in}{2.061896in}}%
\pgfpathlineto{\pgfqpoint{3.867999in}{2.069872in}}%
\pgfpathlineto{\pgfqpoint{3.870518in}{2.065181in}}%
\pgfpathlineto{\pgfqpoint{3.873037in}{2.082545in}}%
\pgfpathlineto{\pgfqpoint{3.875556in}{2.101525in}}%
\pgfpathlineto{\pgfqpoint{3.878074in}{2.095796in}}%
\pgfpathlineto{\pgfqpoint{3.880593in}{2.108122in}}%
\pgfpathlineto{\pgfqpoint{3.883112in}{2.105668in}}%
\pgfpathlineto{\pgfqpoint{3.885631in}{2.126590in}}%
\pgfpathlineto{\pgfqpoint{3.888149in}{2.167416in}}%
\pgfpathlineto{\pgfqpoint{3.890668in}{2.153241in}}%
\pgfpathlineto{\pgfqpoint{3.893187in}{2.148261in}}%
\pgfpathlineto{\pgfqpoint{3.895706in}{2.176407in}}%
\pgfpathlineto{\pgfqpoint{3.898224in}{2.145309in}}%
\pgfpathlineto{\pgfqpoint{3.900743in}{2.131782in}}%
\pgfpathlineto{\pgfqpoint{3.903262in}{2.133669in}}%
\pgfpathlineto{\pgfqpoint{3.905781in}{2.118582in}}%
\pgfpathlineto{\pgfqpoint{3.908299in}{2.127064in}}%
\pgfpathlineto{\pgfqpoint{3.910818in}{2.098372in}}%
\pgfpathlineto{\pgfqpoint{3.913337in}{2.115793in}}%
\pgfpathlineto{\pgfqpoint{3.915856in}{2.129839in}}%
\pgfpathlineto{\pgfqpoint{3.918374in}{2.151056in}}%
\pgfpathlineto{\pgfqpoint{3.920893in}{2.143461in}}%
\pgfpathlineto{\pgfqpoint{3.923412in}{2.117780in}}%
\pgfpathlineto{\pgfqpoint{3.925931in}{2.130336in}}%
\pgfpathlineto{\pgfqpoint{3.928449in}{2.129644in}}%
\pgfpathlineto{\pgfqpoint{3.930968in}{2.159149in}}%
\pgfpathlineto{\pgfqpoint{3.933487in}{2.173276in}}%
\pgfpathlineto{\pgfqpoint{3.936006in}{2.141408in}}%
\pgfpathlineto{\pgfqpoint{3.938524in}{2.127644in}}%
\pgfpathlineto{\pgfqpoint{3.941043in}{2.126022in}}%
\pgfpathlineto{\pgfqpoint{3.943562in}{2.128539in}}%
\pgfpathlineto{\pgfqpoint{3.946081in}{2.124666in}}%
\pgfpathlineto{\pgfqpoint{3.948599in}{2.127422in}}%
\pgfpathlineto{\pgfqpoint{3.951118in}{2.119386in}}%
\pgfpathlineto{\pgfqpoint{3.956156in}{2.155618in}}%
\pgfpathlineto{\pgfqpoint{3.958674in}{2.152229in}}%
\pgfpathlineto{\pgfqpoint{3.961193in}{2.167096in}}%
\pgfpathlineto{\pgfqpoint{3.963712in}{2.178116in}}%
\pgfpathlineto{\pgfqpoint{3.966231in}{2.167428in}}%
\pgfpathlineto{\pgfqpoint{3.968749in}{2.166138in}}%
\pgfpathlineto{\pgfqpoint{3.971268in}{2.143204in}}%
\pgfpathlineto{\pgfqpoint{3.973787in}{2.131074in}}%
\pgfpathlineto{\pgfqpoint{3.976306in}{2.145403in}}%
\pgfpathlineto{\pgfqpoint{3.978824in}{2.136318in}}%
\pgfpathlineto{\pgfqpoint{3.981343in}{2.141102in}}%
\pgfpathlineto{\pgfqpoint{3.983862in}{2.114937in}}%
\pgfpathlineto{\pgfqpoint{3.986381in}{2.093491in}}%
\pgfpathlineto{\pgfqpoint{3.988899in}{2.100086in}}%
\pgfpathlineto{\pgfqpoint{3.991418in}{2.095882in}}%
\pgfpathlineto{\pgfqpoint{3.993937in}{2.063214in}}%
\pgfpathlineto{\pgfqpoint{3.996456in}{2.060892in}}%
\pgfpathlineto{\pgfqpoint{3.998974in}{2.059695in}}%
\pgfpathlineto{\pgfqpoint{4.001493in}{2.032208in}}%
\pgfpathlineto{\pgfqpoint{4.004012in}{2.009946in}}%
\pgfpathlineto{\pgfqpoint{4.006531in}{2.042816in}}%
\pgfpathlineto{\pgfqpoint{4.009049in}{2.030806in}}%
\pgfpathlineto{\pgfqpoint{4.011568in}{2.007245in}}%
\pgfpathlineto{\pgfqpoint{4.014087in}{2.006108in}}%
\pgfpathlineto{\pgfqpoint{4.016606in}{1.985906in}}%
\pgfpathlineto{\pgfqpoint{4.019124in}{1.990321in}}%
\pgfpathlineto{\pgfqpoint{4.021643in}{2.014930in}}%
\pgfpathlineto{\pgfqpoint{4.024162in}{2.028877in}}%
\pgfpathlineto{\pgfqpoint{4.026681in}{2.025096in}}%
\pgfpathlineto{\pgfqpoint{4.029199in}{1.986303in}}%
\pgfpathlineto{\pgfqpoint{4.031718in}{1.978018in}}%
\pgfpathlineto{\pgfqpoint{4.034237in}{1.977189in}}%
\pgfpathlineto{\pgfqpoint{4.036756in}{2.014544in}}%
\pgfpathlineto{\pgfqpoint{4.039274in}{1.992484in}}%
\pgfpathlineto{\pgfqpoint{4.041793in}{1.982426in}}%
\pgfpathlineto{\pgfqpoint{4.044312in}{1.995025in}}%
\pgfpathlineto{\pgfqpoint{4.046831in}{1.982269in}}%
\pgfpathlineto{\pgfqpoint{4.049349in}{2.014178in}}%
\pgfpathlineto{\pgfqpoint{4.051868in}{1.996001in}}%
\pgfpathlineto{\pgfqpoint{4.054387in}{2.006898in}}%
\pgfpathlineto{\pgfqpoint{4.056906in}{1.984540in}}%
\pgfpathlineto{\pgfqpoint{4.059424in}{1.990399in}}%
\pgfpathlineto{\pgfqpoint{4.061943in}{1.984358in}}%
\pgfpathlineto{\pgfqpoint{4.064462in}{1.981988in}}%
\pgfpathlineto{\pgfqpoint{4.066981in}{1.984468in}}%
\pgfpathlineto{\pgfqpoint{4.069499in}{1.980431in}}%
\pgfpathlineto{\pgfqpoint{4.072018in}{1.981691in}}%
\pgfpathlineto{\pgfqpoint{4.074537in}{1.987217in}}%
\pgfpathlineto{\pgfqpoint{4.077056in}{2.010475in}}%
\pgfpathlineto{\pgfqpoint{4.079574in}{2.014069in}}%
\pgfpathlineto{\pgfqpoint{4.082093in}{2.005580in}}%
\pgfpathlineto{\pgfqpoint{4.084612in}{2.001980in}}%
\pgfpathlineto{\pgfqpoint{4.087131in}{2.034461in}}%
\pgfpathlineto{\pgfqpoint{4.089649in}{1.997655in}}%
\pgfpathlineto{\pgfqpoint{4.092168in}{2.027184in}}%
\pgfpathlineto{\pgfqpoint{4.094687in}{1.997707in}}%
\pgfpathlineto{\pgfqpoint{4.097206in}{2.008454in}}%
\pgfpathlineto{\pgfqpoint{4.099724in}{2.011858in}}%
\pgfpathlineto{\pgfqpoint{4.102243in}{2.014307in}}%
\pgfpathlineto{\pgfqpoint{4.104762in}{2.030992in}}%
\pgfpathlineto{\pgfqpoint{4.107281in}{2.015739in}}%
\pgfpathlineto{\pgfqpoint{4.109799in}{2.009348in}}%
\pgfpathlineto{\pgfqpoint{4.112318in}{2.016491in}}%
\pgfpathlineto{\pgfqpoint{4.114837in}{1.993047in}}%
\pgfpathlineto{\pgfqpoint{4.117356in}{1.987941in}}%
\pgfpathlineto{\pgfqpoint{4.119874in}{2.007745in}}%
\pgfpathlineto{\pgfqpoint{4.122393in}{1.988397in}}%
\pgfpathlineto{\pgfqpoint{4.124912in}{1.973553in}}%
\pgfpathlineto{\pgfqpoint{4.127431in}{1.949191in}}%
\pgfpathlineto{\pgfqpoint{4.129949in}{1.947969in}}%
\pgfpathlineto{\pgfqpoint{4.132468in}{1.954002in}}%
\pgfpathlineto{\pgfqpoint{4.134987in}{1.938036in}}%
\pgfpathlineto{\pgfqpoint{4.137506in}{1.931920in}}%
\pgfpathlineto{\pgfqpoint{4.140024in}{1.901282in}}%
\pgfpathlineto{\pgfqpoint{4.142543in}{1.885301in}}%
\pgfpathlineto{\pgfqpoint{4.145062in}{1.904238in}}%
\pgfpathlineto{\pgfqpoint{4.147581in}{1.894679in}}%
\pgfpathlineto{\pgfqpoint{4.150099in}{1.871223in}}%
\pgfpathlineto{\pgfqpoint{4.152618in}{1.897662in}}%
\pgfpathlineto{\pgfqpoint{4.155137in}{1.889646in}}%
\pgfpathlineto{\pgfqpoint{4.157656in}{1.905952in}}%
\pgfpathlineto{\pgfqpoint{4.160174in}{1.935043in}}%
\pgfpathlineto{\pgfqpoint{4.162693in}{1.919228in}}%
\pgfpathlineto{\pgfqpoint{4.165212in}{1.940397in}}%
\pgfpathlineto{\pgfqpoint{4.167731in}{1.917044in}}%
\pgfpathlineto{\pgfqpoint{4.170249in}{1.920496in}}%
\pgfpathlineto{\pgfqpoint{4.172768in}{1.919775in}}%
\pgfpathlineto{\pgfqpoint{4.175287in}{1.932767in}}%
\pgfpathlineto{\pgfqpoint{4.177806in}{1.952563in}}%
\pgfpathlineto{\pgfqpoint{4.180324in}{1.965269in}}%
\pgfpathlineto{\pgfqpoint{4.182843in}{1.993331in}}%
\pgfpathlineto{\pgfqpoint{4.185362in}{1.997641in}}%
\pgfpathlineto{\pgfqpoint{4.187881in}{2.002912in}}%
\pgfpathlineto{\pgfqpoint{4.190399in}{2.001832in}}%
\pgfpathlineto{\pgfqpoint{4.192918in}{2.021669in}}%
\pgfpathlineto{\pgfqpoint{4.195437in}{2.038221in}}%
\pgfpathlineto{\pgfqpoint{4.197956in}{2.010295in}}%
\pgfpathlineto{\pgfqpoint{4.200474in}{2.018296in}}%
\pgfpathlineto{\pgfqpoint{4.202993in}{2.012431in}}%
\pgfpathlineto{\pgfqpoint{4.205512in}{2.019461in}}%
\pgfpathlineto{\pgfqpoint{4.208031in}{2.039325in}}%
\pgfpathlineto{\pgfqpoint{4.210549in}{2.052304in}}%
\pgfpathlineto{\pgfqpoint{4.213068in}{2.066122in}}%
\pgfpathlineto{\pgfqpoint{4.215587in}{2.041444in}}%
\pgfpathlineto{\pgfqpoint{4.218106in}{2.050836in}}%
\pgfpathlineto{\pgfqpoint{4.220624in}{2.053013in}}%
\pgfpathlineto{\pgfqpoint{4.223143in}{2.076301in}}%
\pgfpathlineto{\pgfqpoint{4.225662in}{2.082552in}}%
\pgfpathlineto{\pgfqpoint{4.228181in}{2.070966in}}%
\pgfpathlineto{\pgfqpoint{4.230699in}{2.088042in}}%
\pgfpathlineto{\pgfqpoint{4.233218in}{2.062591in}}%
\pgfpathlineto{\pgfqpoint{4.235737in}{2.065166in}}%
\pgfpathlineto{\pgfqpoint{4.238256in}{2.070606in}}%
\pgfpathlineto{\pgfqpoint{4.240774in}{2.046055in}}%
\pgfpathlineto{\pgfqpoint{4.243293in}{2.049046in}}%
\pgfpathlineto{\pgfqpoint{4.245812in}{2.043039in}}%
\pgfpathlineto{\pgfqpoint{4.248331in}{2.046621in}}%
\pgfpathlineto{\pgfqpoint{4.250849in}{2.040670in}}%
\pgfpathlineto{\pgfqpoint{4.253368in}{2.067111in}}%
\pgfpathlineto{\pgfqpoint{4.255887in}{2.056265in}}%
\pgfpathlineto{\pgfqpoint{4.258406in}{2.060066in}}%
\pgfpathlineto{\pgfqpoint{4.260924in}{2.055562in}}%
\pgfpathlineto{\pgfqpoint{4.263443in}{2.061542in}}%
\pgfpathlineto{\pgfqpoint{4.265962in}{2.034665in}}%
\pgfpathlineto{\pgfqpoint{4.268481in}{2.075143in}}%
\pgfpathlineto{\pgfqpoint{4.270999in}{2.095772in}}%
\pgfpathlineto{\pgfqpoint{4.273518in}{2.051113in}}%
\pgfpathlineto{\pgfqpoint{4.276037in}{2.073468in}}%
\pgfpathlineto{\pgfqpoint{4.278556in}{2.074061in}}%
\pgfpathlineto{\pgfqpoint{4.281074in}{2.061264in}}%
\pgfpathlineto{\pgfqpoint{4.283593in}{2.096781in}}%
\pgfpathlineto{\pgfqpoint{4.286112in}{2.086377in}}%
\pgfpathlineto{\pgfqpoint{4.288631in}{2.091929in}}%
\pgfpathlineto{\pgfqpoint{4.291149in}{2.109177in}}%
\pgfpathlineto{\pgfqpoint{4.293668in}{2.092834in}}%
\pgfpathlineto{\pgfqpoint{4.296187in}{2.079820in}}%
\pgfpathlineto{\pgfqpoint{4.298706in}{2.076661in}}%
\pgfpathlineto{\pgfqpoint{4.301224in}{2.070743in}}%
\pgfpathlineto{\pgfqpoint{4.303743in}{2.061210in}}%
\pgfpathlineto{\pgfqpoint{4.306262in}{2.114993in}}%
\pgfpathlineto{\pgfqpoint{4.308781in}{2.099057in}}%
\pgfpathlineto{\pgfqpoint{4.311299in}{2.106411in}}%
\pgfpathlineto{\pgfqpoint{4.313818in}{2.120201in}}%
\pgfpathlineto{\pgfqpoint{4.316337in}{2.128548in}}%
\pgfpathlineto{\pgfqpoint{4.318856in}{2.143360in}}%
\pgfpathlineto{\pgfqpoint{4.321374in}{2.135385in}}%
\pgfpathlineto{\pgfqpoint{4.323893in}{2.118155in}}%
\pgfpathlineto{\pgfqpoint{4.326412in}{2.132611in}}%
\pgfpathlineto{\pgfqpoint{4.328931in}{2.112579in}}%
\pgfpathlineto{\pgfqpoint{4.331449in}{2.105798in}}%
\pgfpathlineto{\pgfqpoint{4.333968in}{2.112206in}}%
\pgfpathlineto{\pgfqpoint{4.336487in}{2.105729in}}%
\pgfpathlineto{\pgfqpoint{4.339006in}{2.087376in}}%
\pgfpathlineto{\pgfqpoint{4.341524in}{2.075469in}}%
\pgfpathlineto{\pgfqpoint{4.344043in}{2.060423in}}%
\pgfpathlineto{\pgfqpoint{4.346562in}{2.052093in}}%
\pgfpathlineto{\pgfqpoint{4.349081in}{2.038392in}}%
\pgfpathlineto{\pgfqpoint{4.351599in}{2.028647in}}%
\pgfpathlineto{\pgfqpoint{4.354118in}{2.001212in}}%
\pgfpathlineto{\pgfqpoint{4.356637in}{2.003853in}}%
\pgfpathlineto{\pgfqpoint{4.359156in}{1.989034in}}%
\pgfpathlineto{\pgfqpoint{4.361674in}{1.992406in}}%
\pgfpathlineto{\pgfqpoint{4.364193in}{1.989301in}}%
\pgfpathlineto{\pgfqpoint{4.366712in}{1.990509in}}%
\pgfpathlineto{\pgfqpoint{4.369231in}{1.996616in}}%
\pgfpathlineto{\pgfqpoint{4.371749in}{1.990035in}}%
\pgfpathlineto{\pgfqpoint{4.374268in}{2.015193in}}%
\pgfpathlineto{\pgfqpoint{4.376787in}{1.979752in}}%
\pgfpathlineto{\pgfqpoint{4.379306in}{1.962872in}}%
\pgfpathlineto{\pgfqpoint{4.381824in}{1.955615in}}%
\pgfpathlineto{\pgfqpoint{4.384343in}{1.954072in}}%
\pgfpathlineto{\pgfqpoint{4.386862in}{1.925110in}}%
\pgfpathlineto{\pgfqpoint{4.389381in}{1.919368in}}%
\pgfpathlineto{\pgfqpoint{4.391899in}{1.910786in}}%
\pgfpathlineto{\pgfqpoint{4.394418in}{1.896325in}}%
\pgfpathlineto{\pgfqpoint{4.396937in}{1.875650in}}%
\pgfpathlineto{\pgfqpoint{4.399456in}{1.870877in}}%
\pgfpathlineto{\pgfqpoint{4.401974in}{1.890900in}}%
\pgfpathlineto{\pgfqpoint{4.404493in}{1.880256in}}%
\pgfpathlineto{\pgfqpoint{4.407012in}{1.845473in}}%
\pgfpathlineto{\pgfqpoint{4.409531in}{1.829015in}}%
\pgfpathlineto{\pgfqpoint{4.412049in}{1.824426in}}%
\pgfpathlineto{\pgfqpoint{4.414568in}{1.802389in}}%
\pgfpathlineto{\pgfqpoint{4.417087in}{1.813596in}}%
\pgfpathlineto{\pgfqpoint{4.419606in}{1.852716in}}%
\pgfpathlineto{\pgfqpoint{4.422124in}{1.832320in}}%
\pgfpathlineto{\pgfqpoint{4.424643in}{1.850036in}}%
\pgfpathlineto{\pgfqpoint{4.427162in}{1.851240in}}%
\pgfpathlineto{\pgfqpoint{4.429681in}{1.851663in}}%
\pgfpathlineto{\pgfqpoint{4.432199in}{1.863045in}}%
\pgfpathlineto{\pgfqpoint{4.434718in}{1.872208in}}%
\pgfpathlineto{\pgfqpoint{4.437237in}{1.884354in}}%
\pgfpathlineto{\pgfqpoint{4.439756in}{1.894448in}}%
\pgfpathlineto{\pgfqpoint{4.442274in}{1.891450in}}%
\pgfpathlineto{\pgfqpoint{4.444793in}{1.898954in}}%
\pgfpathlineto{\pgfqpoint{4.447312in}{1.917498in}}%
\pgfpathlineto{\pgfqpoint{4.449831in}{1.900226in}}%
\pgfpathlineto{\pgfqpoint{4.452349in}{1.872690in}}%
\pgfpathlineto{\pgfqpoint{4.454868in}{1.852559in}}%
\pgfpathlineto{\pgfqpoint{4.457387in}{1.831136in}}%
\pgfpathlineto{\pgfqpoint{4.459906in}{1.835352in}}%
\pgfpathlineto{\pgfqpoint{4.462424in}{1.826452in}}%
\pgfpathlineto{\pgfqpoint{4.464943in}{1.837111in}}%
\pgfpathlineto{\pgfqpoint{4.467462in}{1.838853in}}%
\pgfpathlineto{\pgfqpoint{4.469981in}{1.821836in}}%
\pgfpathlineto{\pgfqpoint{4.472499in}{1.817423in}}%
\pgfpathlineto{\pgfqpoint{4.475018in}{1.800678in}}%
\pgfpathlineto{\pgfqpoint{4.477537in}{1.800115in}}%
\pgfpathlineto{\pgfqpoint{4.480056in}{1.806924in}}%
\pgfpathlineto{\pgfqpoint{4.482574in}{1.792722in}}%
\pgfpathlineto{\pgfqpoint{4.485093in}{1.791621in}}%
\pgfpathlineto{\pgfqpoint{4.487612in}{1.805767in}}%
\pgfpathlineto{\pgfqpoint{4.490131in}{1.797551in}}%
\pgfpathlineto{\pgfqpoint{4.495168in}{1.815955in}}%
\pgfpathlineto{\pgfqpoint{4.497687in}{1.840795in}}%
\pgfpathlineto{\pgfqpoint{4.500206in}{1.847498in}}%
\pgfpathlineto{\pgfqpoint{4.502724in}{1.837916in}}%
\pgfpathlineto{\pgfqpoint{4.505243in}{1.830268in}}%
\pgfpathlineto{\pgfqpoint{4.507762in}{1.805003in}}%
\pgfpathlineto{\pgfqpoint{4.510281in}{1.816872in}}%
\pgfpathlineto{\pgfqpoint{4.512799in}{1.821754in}}%
\pgfpathlineto{\pgfqpoint{4.515318in}{1.835252in}}%
\pgfpathlineto{\pgfqpoint{4.517837in}{1.824014in}}%
\pgfpathlineto{\pgfqpoint{4.520356in}{1.826994in}}%
\pgfpathlineto{\pgfqpoint{4.522874in}{1.807140in}}%
\pgfpathlineto{\pgfqpoint{4.525393in}{1.791583in}}%
\pgfpathlineto{\pgfqpoint{4.527912in}{1.766214in}}%
\pgfpathlineto{\pgfqpoint{4.530431in}{1.781512in}}%
\pgfpathlineto{\pgfqpoint{4.532949in}{1.778874in}}%
\pgfpathlineto{\pgfqpoint{4.535468in}{1.775827in}}%
\pgfpathlineto{\pgfqpoint{4.537987in}{1.784866in}}%
\pgfpathlineto{\pgfqpoint{4.540506in}{1.780972in}}%
\pgfpathlineto{\pgfqpoint{4.543024in}{1.769860in}}%
\pgfpathlineto{\pgfqpoint{4.545543in}{1.727041in}}%
\pgfpathlineto{\pgfqpoint{4.548062in}{1.707439in}}%
\pgfpathlineto{\pgfqpoint{4.550581in}{1.707378in}}%
\pgfpathlineto{\pgfqpoint{4.553099in}{1.732188in}}%
\pgfpathlineto{\pgfqpoint{4.555618in}{1.730158in}}%
\pgfpathlineto{\pgfqpoint{4.558137in}{1.710327in}}%
\pgfpathlineto{\pgfqpoint{4.560656in}{1.719055in}}%
\pgfpathlineto{\pgfqpoint{4.563174in}{1.702955in}}%
\pgfpathlineto{\pgfqpoint{4.565693in}{1.684561in}}%
\pgfpathlineto{\pgfqpoint{4.568212in}{1.695210in}}%
\pgfpathlineto{\pgfqpoint{4.570731in}{1.671144in}}%
\pgfpathlineto{\pgfqpoint{4.573249in}{1.670423in}}%
\pgfpathlineto{\pgfqpoint{4.575768in}{1.666815in}}%
\pgfpathlineto{\pgfqpoint{4.578287in}{1.684246in}}%
\pgfpathlineto{\pgfqpoint{4.580806in}{1.721945in}}%
\pgfpathlineto{\pgfqpoint{4.583324in}{1.731836in}}%
\pgfpathlineto{\pgfqpoint{4.585843in}{1.729656in}}%
\pgfpathlineto{\pgfqpoint{4.588362in}{1.742058in}}%
\pgfpathlineto{\pgfqpoint{4.590881in}{1.747887in}}%
\pgfpathlineto{\pgfqpoint{4.593399in}{1.770347in}}%
\pgfpathlineto{\pgfqpoint{4.595918in}{1.769051in}}%
\pgfpathlineto{\pgfqpoint{4.598437in}{1.775615in}}%
\pgfpathlineto{\pgfqpoint{4.600956in}{1.750486in}}%
\pgfpathlineto{\pgfqpoint{4.603474in}{1.710926in}}%
\pgfpathlineto{\pgfqpoint{4.605993in}{1.715759in}}%
\pgfpathlineto{\pgfqpoint{4.608512in}{1.716462in}}%
\pgfpathlineto{\pgfqpoint{4.611031in}{1.713288in}}%
\pgfpathlineto{\pgfqpoint{4.613549in}{1.730403in}}%
\pgfpathlineto{\pgfqpoint{4.616068in}{1.737487in}}%
\pgfpathlineto{\pgfqpoint{4.618587in}{1.717954in}}%
\pgfpathlineto{\pgfqpoint{4.621106in}{1.712261in}}%
\pgfpathlineto{\pgfqpoint{4.623624in}{1.731829in}}%
\pgfpathlineto{\pgfqpoint{4.626143in}{1.748393in}}%
\pgfpathlineto{\pgfqpoint{4.628662in}{1.742380in}}%
\pgfpathlineto{\pgfqpoint{4.631181in}{1.754455in}}%
\pgfpathlineto{\pgfqpoint{4.633699in}{1.733002in}}%
\pgfpathlineto{\pgfqpoint{4.636218in}{1.695565in}}%
\pgfpathlineto{\pgfqpoint{4.638737in}{1.705194in}}%
\pgfpathlineto{\pgfqpoint{4.641256in}{1.675828in}}%
\pgfpathlineto{\pgfqpoint{4.643774in}{1.662272in}}%
\pgfpathlineto{\pgfqpoint{4.646293in}{1.681350in}}%
\pgfpathlineto{\pgfqpoint{4.648812in}{1.665131in}}%
\pgfpathlineto{\pgfqpoint{4.651331in}{1.675593in}}%
\pgfpathlineto{\pgfqpoint{4.653849in}{1.684888in}}%
\pgfpathlineto{\pgfqpoint{4.656368in}{1.653092in}}%
\pgfpathlineto{\pgfqpoint{4.658887in}{1.675220in}}%
\pgfpathlineto{\pgfqpoint{4.661406in}{1.683035in}}%
\pgfpathlineto{\pgfqpoint{4.663924in}{1.682434in}}%
\pgfpathlineto{\pgfqpoint{4.666443in}{1.663115in}}%
\pgfpathlineto{\pgfqpoint{4.668962in}{1.659478in}}%
\pgfpathlineto{\pgfqpoint{4.671481in}{1.669632in}}%
\pgfpathlineto{\pgfqpoint{4.673999in}{1.661709in}}%
\pgfpathlineto{\pgfqpoint{4.676518in}{1.686964in}}%
\pgfpathlineto{\pgfqpoint{4.679037in}{1.706478in}}%
\pgfpathlineto{\pgfqpoint{4.681556in}{1.703716in}}%
\pgfpathlineto{\pgfqpoint{4.684074in}{1.698409in}}%
\pgfpathlineto{\pgfqpoint{4.686593in}{1.680076in}}%
\pgfpathlineto{\pgfqpoint{4.689112in}{1.684443in}}%
\pgfpathlineto{\pgfqpoint{4.691631in}{1.682603in}}%
\pgfpathlineto{\pgfqpoint{4.694149in}{1.684151in}}%
\pgfpathlineto{\pgfqpoint{4.696668in}{1.702067in}}%
\pgfpathlineto{\pgfqpoint{4.699187in}{1.712930in}}%
\pgfpathlineto{\pgfqpoint{4.701706in}{1.710518in}}%
\pgfpathlineto{\pgfqpoint{4.704224in}{1.714407in}}%
\pgfpathlineto{\pgfqpoint{4.706743in}{1.727727in}}%
\pgfpathlineto{\pgfqpoint{4.709262in}{1.725099in}}%
\pgfpathlineto{\pgfqpoint{4.711781in}{1.726207in}}%
\pgfpathlineto{\pgfqpoint{4.714299in}{1.711983in}}%
\pgfpathlineto{\pgfqpoint{4.716818in}{1.725675in}}%
\pgfpathlineto{\pgfqpoint{4.719337in}{1.725976in}}%
\pgfpathlineto{\pgfqpoint{4.721856in}{1.728669in}}%
\pgfpathlineto{\pgfqpoint{4.724374in}{1.722447in}}%
\pgfpathlineto{\pgfqpoint{4.726893in}{1.695016in}}%
\pgfpathlineto{\pgfqpoint{4.729412in}{1.696710in}}%
\pgfpathlineto{\pgfqpoint{4.731931in}{1.723785in}}%
\pgfpathlineto{\pgfqpoint{4.734449in}{1.714303in}}%
\pgfpathlineto{\pgfqpoint{4.736968in}{1.730258in}}%
\pgfpathlineto{\pgfqpoint{4.739487in}{1.758203in}}%
\pgfpathlineto{\pgfqpoint{4.742006in}{1.752811in}}%
\pgfpathlineto{\pgfqpoint{4.744524in}{1.759582in}}%
\pgfpathlineto{\pgfqpoint{4.747043in}{1.748059in}}%
\pgfpathlineto{\pgfqpoint{4.749562in}{1.767289in}}%
\pgfpathlineto{\pgfqpoint{4.752081in}{1.764550in}}%
\pgfpathlineto{\pgfqpoint{4.754599in}{1.775328in}}%
\pgfpathlineto{\pgfqpoint{4.757118in}{1.764779in}}%
\pgfpathlineto{\pgfqpoint{4.759637in}{1.741817in}}%
\pgfpathlineto{\pgfqpoint{4.762156in}{1.768138in}}%
\pgfpathlineto{\pgfqpoint{4.764674in}{1.755936in}}%
\pgfpathlineto{\pgfqpoint{4.767193in}{1.739024in}}%
\pgfpathlineto{\pgfqpoint{4.769712in}{1.751564in}}%
\pgfpathlineto{\pgfqpoint{4.772231in}{1.758070in}}%
\pgfpathlineto{\pgfqpoint{4.774749in}{1.776342in}}%
\pgfpathlineto{\pgfqpoint{4.777268in}{1.785349in}}%
\pgfpathlineto{\pgfqpoint{4.779787in}{1.779130in}}%
\pgfpathlineto{\pgfqpoint{4.782306in}{1.785929in}}%
\pgfpathlineto{\pgfqpoint{4.784824in}{1.769359in}}%
\pgfpathlineto{\pgfqpoint{4.787343in}{1.769360in}}%
\pgfpathlineto{\pgfqpoint{4.789862in}{1.767153in}}%
\pgfpathlineto{\pgfqpoint{4.792381in}{1.779173in}}%
\pgfpathlineto{\pgfqpoint{4.794899in}{1.780037in}}%
\pgfpathlineto{\pgfqpoint{4.797418in}{1.808639in}}%
\pgfpathlineto{\pgfqpoint{4.799937in}{1.792076in}}%
\pgfpathlineto{\pgfqpoint{4.802456in}{1.796638in}}%
\pgfpathlineto{\pgfqpoint{4.804974in}{1.769683in}}%
\pgfpathlineto{\pgfqpoint{4.807493in}{1.750176in}}%
\pgfpathlineto{\pgfqpoint{4.810012in}{1.759628in}}%
\pgfpathlineto{\pgfqpoint{4.812531in}{1.801083in}}%
\pgfpathlineto{\pgfqpoint{4.815049in}{1.798859in}}%
\pgfpathlineto{\pgfqpoint{4.817568in}{1.753666in}}%
\pgfpathlineto{\pgfqpoint{4.820087in}{1.742523in}}%
\pgfpathlineto{\pgfqpoint{4.822606in}{1.756119in}}%
\pgfpathlineto{\pgfqpoint{4.825124in}{1.753401in}}%
\pgfpathlineto{\pgfqpoint{4.827643in}{1.771731in}}%
\pgfpathlineto{\pgfqpoint{4.830162in}{1.772173in}}%
\pgfpathlineto{\pgfqpoint{4.832681in}{1.748160in}}%
\pgfpathlineto{\pgfqpoint{4.835199in}{1.730523in}}%
\pgfpathlineto{\pgfqpoint{4.837718in}{1.748364in}}%
\pgfpathlineto{\pgfqpoint{4.840237in}{1.746468in}}%
\pgfpathlineto{\pgfqpoint{4.842756in}{1.740060in}}%
\pgfpathlineto{\pgfqpoint{4.845274in}{1.726929in}}%
\pgfpathlineto{\pgfqpoint{4.847793in}{1.698720in}}%
\pgfpathlineto{\pgfqpoint{4.850312in}{1.702258in}}%
\pgfpathlineto{\pgfqpoint{4.852831in}{1.691492in}}%
\pgfpathlineto{\pgfqpoint{4.855349in}{1.686607in}}%
\pgfpathlineto{\pgfqpoint{4.857868in}{1.698248in}}%
\pgfpathlineto{\pgfqpoint{4.860387in}{1.675645in}}%
\pgfpathlineto{\pgfqpoint{4.862906in}{1.673739in}}%
\pgfpathlineto{\pgfqpoint{4.865424in}{1.664795in}}%
\pgfpathlineto{\pgfqpoint{4.867943in}{1.686786in}}%
\pgfpathlineto{\pgfqpoint{4.870462in}{1.704248in}}%
\pgfpathlineto{\pgfqpoint{4.872981in}{1.712167in}}%
\pgfpathlineto{\pgfqpoint{4.875499in}{1.703583in}}%
\pgfpathlineto{\pgfqpoint{4.878018in}{1.699947in}}%
\pgfpathlineto{\pgfqpoint{4.880537in}{1.699414in}}%
\pgfpathlineto{\pgfqpoint{4.883056in}{1.699291in}}%
\pgfpathlineto{\pgfqpoint{4.885574in}{1.711369in}}%
\pgfpathlineto{\pgfqpoint{4.888093in}{1.685760in}}%
\pgfpathlineto{\pgfqpoint{4.890612in}{1.684088in}}%
\pgfpathlineto{\pgfqpoint{4.893131in}{1.692913in}}%
\pgfpathlineto{\pgfqpoint{4.895649in}{1.685363in}}%
\pgfpathlineto{\pgfqpoint{4.898168in}{1.668958in}}%
\pgfpathlineto{\pgfqpoint{4.900687in}{1.667523in}}%
\pgfpathlineto{\pgfqpoint{4.903206in}{1.650193in}}%
\pgfpathlineto{\pgfqpoint{4.905724in}{1.645672in}}%
\pgfpathlineto{\pgfqpoint{4.908243in}{1.675447in}}%
\pgfpathlineto{\pgfqpoint{4.910762in}{1.683916in}}%
\pgfpathlineto{\pgfqpoint{4.913281in}{1.701737in}}%
\pgfpathlineto{\pgfqpoint{4.915799in}{1.704440in}}%
\pgfpathlineto{\pgfqpoint{4.920837in}{1.727018in}}%
\pgfpathlineto{\pgfqpoint{4.923356in}{1.734205in}}%
\pgfpathlineto{\pgfqpoint{4.925874in}{1.746733in}}%
\pgfpathlineto{\pgfqpoint{4.928393in}{1.739980in}}%
\pgfpathlineto{\pgfqpoint{4.930912in}{1.751626in}}%
\pgfpathlineto{\pgfqpoint{4.933431in}{1.774603in}}%
\pgfpathlineto{\pgfqpoint{4.935949in}{1.780096in}}%
\pgfpathlineto{\pgfqpoint{4.938468in}{1.784547in}}%
\pgfpathlineto{\pgfqpoint{4.940987in}{1.784919in}}%
\pgfpathlineto{\pgfqpoint{4.943506in}{1.776464in}}%
\pgfpathlineto{\pgfqpoint{4.946024in}{1.764507in}}%
\pgfpathlineto{\pgfqpoint{4.948543in}{1.771462in}}%
\pgfpathlineto{\pgfqpoint{4.951062in}{1.781799in}}%
\pgfpathlineto{\pgfqpoint{4.953581in}{1.793725in}}%
\pgfpathlineto{\pgfqpoint{4.956099in}{1.795254in}}%
\pgfpathlineto{\pgfqpoint{4.958618in}{1.788565in}}%
\pgfpathlineto{\pgfqpoint{4.963656in}{1.750932in}}%
\pgfpathlineto{\pgfqpoint{4.966174in}{1.772225in}}%
\pgfpathlineto{\pgfqpoint{4.968693in}{1.765757in}}%
\pgfpathlineto{\pgfqpoint{4.971212in}{1.770632in}}%
\pgfpathlineto{\pgfqpoint{4.976249in}{1.746402in}}%
\pgfpathlineto{\pgfqpoint{4.978768in}{1.744147in}}%
\pgfpathlineto{\pgfqpoint{4.981287in}{1.743778in}}%
\pgfpathlineto{\pgfqpoint{4.983806in}{1.768161in}}%
\pgfpathlineto{\pgfqpoint{4.986324in}{1.772125in}}%
\pgfpathlineto{\pgfqpoint{4.988843in}{1.781371in}}%
\pgfpathlineto{\pgfqpoint{4.991362in}{1.771872in}}%
\pgfpathlineto{\pgfqpoint{4.993881in}{1.778083in}}%
\pgfpathlineto{\pgfqpoint{4.996399in}{1.796233in}}%
\pgfpathlineto{\pgfqpoint{4.998918in}{1.808610in}}%
\pgfpathlineto{\pgfqpoint{5.001437in}{1.819666in}}%
\pgfpathlineto{\pgfqpoint{5.003956in}{1.823809in}}%
\pgfpathlineto{\pgfqpoint{5.006474in}{1.798730in}}%
\pgfpathlineto{\pgfqpoint{5.008993in}{1.818494in}}%
\pgfpathlineto{\pgfqpoint{5.011512in}{1.844792in}}%
\pgfpathlineto{\pgfqpoint{5.014031in}{1.844788in}}%
\pgfpathlineto{\pgfqpoint{5.016549in}{1.829693in}}%
\pgfpathlineto{\pgfqpoint{5.019068in}{1.827419in}}%
\pgfpathlineto{\pgfqpoint{5.021587in}{1.826751in}}%
\pgfpathlineto{\pgfqpoint{5.024106in}{1.852987in}}%
\pgfpathlineto{\pgfqpoint{5.026624in}{1.847419in}}%
\pgfpathlineto{\pgfqpoint{5.029143in}{1.867274in}}%
\pgfpathlineto{\pgfqpoint{5.031662in}{1.849459in}}%
\pgfpathlineto{\pgfqpoint{5.034181in}{1.838310in}}%
\pgfpathlineto{\pgfqpoint{5.036699in}{1.867718in}}%
\pgfpathlineto{\pgfqpoint{5.039218in}{1.879540in}}%
\pgfpathlineto{\pgfqpoint{5.041737in}{1.869351in}}%
\pgfpathlineto{\pgfqpoint{5.044256in}{1.845167in}}%
\pgfpathlineto{\pgfqpoint{5.046774in}{1.832216in}}%
\pgfpathlineto{\pgfqpoint{5.049293in}{1.836873in}}%
\pgfpathlineto{\pgfqpoint{5.051812in}{1.836212in}}%
\pgfpathlineto{\pgfqpoint{5.054331in}{1.822360in}}%
\pgfpathlineto{\pgfqpoint{5.056849in}{1.850992in}}%
\pgfpathlineto{\pgfqpoint{5.059368in}{1.842833in}}%
\pgfpathlineto{\pgfqpoint{5.061887in}{1.853310in}}%
\pgfpathlineto{\pgfqpoint{5.064406in}{1.897138in}}%
\pgfpathlineto{\pgfqpoint{5.066924in}{1.881543in}}%
\pgfpathlineto{\pgfqpoint{5.069443in}{1.893370in}}%
\pgfpathlineto{\pgfqpoint{5.071962in}{1.896414in}}%
\pgfpathlineto{\pgfqpoint{5.074481in}{1.908411in}}%
\pgfpathlineto{\pgfqpoint{5.076999in}{1.858438in}}%
\pgfpathlineto{\pgfqpoint{5.079518in}{1.850666in}}%
\pgfpathlineto{\pgfqpoint{5.082037in}{1.840759in}}%
\pgfpathlineto{\pgfqpoint{5.084556in}{1.858433in}}%
\pgfpathlineto{\pgfqpoint{5.087074in}{1.857960in}}%
\pgfpathlineto{\pgfqpoint{5.089593in}{1.828542in}}%
\pgfpathlineto{\pgfqpoint{5.092112in}{1.824394in}}%
\pgfpathlineto{\pgfqpoint{5.094631in}{1.795534in}}%
\pgfpathlineto{\pgfqpoint{5.097149in}{1.790311in}}%
\pgfpathlineto{\pgfqpoint{5.099668in}{1.798548in}}%
\pgfpathlineto{\pgfqpoint{5.102187in}{1.803332in}}%
\pgfpathlineto{\pgfqpoint{5.104706in}{1.814937in}}%
\pgfpathlineto{\pgfqpoint{5.107224in}{1.863827in}}%
\pgfpathlineto{\pgfqpoint{5.109743in}{1.881075in}}%
\pgfpathlineto{\pgfqpoint{5.112262in}{1.896275in}}%
\pgfpathlineto{\pgfqpoint{5.114781in}{1.899891in}}%
\pgfpathlineto{\pgfqpoint{5.117299in}{1.928538in}}%
\pgfpathlineto{\pgfqpoint{5.119818in}{1.928554in}}%
\pgfpathlineto{\pgfqpoint{5.122337in}{1.911337in}}%
\pgfpathlineto{\pgfqpoint{5.124856in}{1.916782in}}%
\pgfpathlineto{\pgfqpoint{5.129893in}{1.915407in}}%
\pgfpathlineto{\pgfqpoint{5.132412in}{1.938517in}}%
\pgfpathlineto{\pgfqpoint{5.134931in}{1.948865in}}%
\pgfpathlineto{\pgfqpoint{5.137449in}{1.948140in}}%
\pgfpathlineto{\pgfqpoint{5.139968in}{1.975383in}}%
\pgfpathlineto{\pgfqpoint{5.142487in}{1.984249in}}%
\pgfpathlineto{\pgfqpoint{5.145006in}{1.969454in}}%
\pgfpathlineto{\pgfqpoint{5.147524in}{1.971874in}}%
\pgfpathlineto{\pgfqpoint{5.150043in}{1.968623in}}%
\pgfpathlineto{\pgfqpoint{5.152562in}{1.973871in}}%
\pgfpathlineto{\pgfqpoint{5.155081in}{1.979941in}}%
\pgfpathlineto{\pgfqpoint{5.157599in}{2.008377in}}%
\pgfpathlineto{\pgfqpoint{5.160118in}{2.008053in}}%
\pgfpathlineto{\pgfqpoint{5.162637in}{2.011975in}}%
\pgfpathlineto{\pgfqpoint{5.165156in}{2.012009in}}%
\pgfpathlineto{\pgfqpoint{5.167674in}{2.039687in}}%
\pgfpathlineto{\pgfqpoint{5.170193in}{2.061547in}}%
\pgfpathlineto{\pgfqpoint{5.172712in}{2.099339in}}%
\pgfpathlineto{\pgfqpoint{5.175231in}{2.079720in}}%
\pgfpathlineto{\pgfqpoint{5.177749in}{2.081741in}}%
\pgfpathlineto{\pgfqpoint{5.180268in}{2.095021in}}%
\pgfpathlineto{\pgfqpoint{5.182787in}{2.112609in}}%
\pgfpathlineto{\pgfqpoint{5.185306in}{2.131200in}}%
\pgfpathlineto{\pgfqpoint{5.187824in}{2.130075in}}%
\pgfpathlineto{\pgfqpoint{5.190343in}{2.151320in}}%
\pgfpathlineto{\pgfqpoint{5.192862in}{2.167403in}}%
\pgfpathlineto{\pgfqpoint{5.195381in}{2.178981in}}%
\pgfpathlineto{\pgfqpoint{5.197899in}{2.191817in}}%
\pgfpathlineto{\pgfqpoint{5.200418in}{2.198187in}}%
\pgfpathlineto{\pgfqpoint{5.202937in}{2.221912in}}%
\pgfpathlineto{\pgfqpoint{5.205456in}{2.206708in}}%
\pgfpathlineto{\pgfqpoint{5.207974in}{2.179551in}}%
\pgfpathlineto{\pgfqpoint{5.210493in}{2.183102in}}%
\pgfpathlineto{\pgfqpoint{5.213012in}{2.202236in}}%
\pgfpathlineto{\pgfqpoint{5.215531in}{2.200954in}}%
\pgfpathlineto{\pgfqpoint{5.218049in}{2.195191in}}%
\pgfpathlineto{\pgfqpoint{5.220568in}{2.206683in}}%
\pgfpathlineto{\pgfqpoint{5.223087in}{2.176390in}}%
\pgfpathlineto{\pgfqpoint{5.225606in}{2.178812in}}%
\pgfpathlineto{\pgfqpoint{5.228124in}{2.196205in}}%
\pgfpathlineto{\pgfqpoint{5.230643in}{2.189292in}}%
\pgfpathlineto{\pgfqpoint{5.233162in}{2.172054in}}%
\pgfpathlineto{\pgfqpoint{5.235681in}{2.187036in}}%
\pgfpathlineto{\pgfqpoint{5.238199in}{2.189580in}}%
\pgfpathlineto{\pgfqpoint{5.240718in}{2.182937in}}%
\pgfpathlineto{\pgfqpoint{5.243237in}{2.168421in}}%
\pgfpathlineto{\pgfqpoint{5.248274in}{2.150604in}}%
\pgfpathlineto{\pgfqpoint{5.250793in}{2.132836in}}%
\pgfpathlineto{\pgfqpoint{5.253312in}{2.139754in}}%
\pgfpathlineto{\pgfqpoint{5.255831in}{2.165899in}}%
\pgfpathlineto{\pgfqpoint{5.258349in}{2.165133in}}%
\pgfpathlineto{\pgfqpoint{5.260868in}{2.147917in}}%
\pgfpathlineto{\pgfqpoint{5.263387in}{2.127577in}}%
\pgfpathlineto{\pgfqpoint{5.265906in}{2.084208in}}%
\pgfpathlineto{\pgfqpoint{5.268424in}{2.068561in}}%
\pgfpathlineto{\pgfqpoint{5.270943in}{2.072364in}}%
\pgfpathlineto{\pgfqpoint{5.273462in}{2.066997in}}%
\pgfpathlineto{\pgfqpoint{5.275981in}{2.056233in}}%
\pgfpathlineto{\pgfqpoint{5.278499in}{2.055200in}}%
\pgfpathlineto{\pgfqpoint{5.281018in}{2.081200in}}%
\pgfpathlineto{\pgfqpoint{5.283537in}{2.077689in}}%
\pgfpathlineto{\pgfqpoint{5.286056in}{2.064796in}}%
\pgfpathlineto{\pgfqpoint{5.288574in}{2.071621in}}%
\pgfpathlineto{\pgfqpoint{5.291093in}{2.052092in}}%
\pgfpathlineto{\pgfqpoint{5.293612in}{2.035819in}}%
\pgfpathlineto{\pgfqpoint{5.296131in}{2.023538in}}%
\pgfpathlineto{\pgfqpoint{5.298649in}{2.029911in}}%
\pgfpathlineto{\pgfqpoint{5.301168in}{2.019559in}}%
\pgfpathlineto{\pgfqpoint{5.303687in}{2.047708in}}%
\pgfpathlineto{\pgfqpoint{5.306206in}{2.033047in}}%
\pgfpathlineto{\pgfqpoint{5.308724in}{2.033162in}}%
\pgfpathlineto{\pgfqpoint{5.311243in}{2.035016in}}%
\pgfpathlineto{\pgfqpoint{5.313762in}{2.037375in}}%
\pgfpathlineto{\pgfqpoint{5.316281in}{2.049730in}}%
\pgfpathlineto{\pgfqpoint{5.318799in}{2.042066in}}%
\pgfpathlineto{\pgfqpoint{5.321318in}{2.048900in}}%
\pgfpathlineto{\pgfqpoint{5.323837in}{2.045497in}}%
\pgfpathlineto{\pgfqpoint{5.326356in}{2.027214in}}%
\pgfpathlineto{\pgfqpoint{5.328874in}{2.034542in}}%
\pgfpathlineto{\pgfqpoint{5.331393in}{2.018610in}}%
\pgfpathlineto{\pgfqpoint{5.333912in}{2.009514in}}%
\pgfpathlineto{\pgfqpoint{5.336431in}{1.976611in}}%
\pgfpathlineto{\pgfqpoint{5.338949in}{1.977113in}}%
\pgfpathlineto{\pgfqpoint{5.341468in}{2.005247in}}%
\pgfpathlineto{\pgfqpoint{5.343987in}{2.022787in}}%
\pgfpathlineto{\pgfqpoint{5.346506in}{1.976773in}}%
\pgfpathlineto{\pgfqpoint{5.349024in}{1.985777in}}%
\pgfpathlineto{\pgfqpoint{5.351543in}{1.978612in}}%
\pgfpathlineto{\pgfqpoint{5.354062in}{1.992230in}}%
\pgfpathlineto{\pgfqpoint{5.356581in}{2.001813in}}%
\pgfpathlineto{\pgfqpoint{5.359099in}{2.010946in}}%
\pgfpathlineto{\pgfqpoint{5.361618in}{2.019464in}}%
\pgfpathlineto{\pgfqpoint{5.366656in}{1.998470in}}%
\pgfpathlineto{\pgfqpoint{5.369174in}{2.017528in}}%
\pgfpathlineto{\pgfqpoint{5.371693in}{2.048736in}}%
\pgfpathlineto{\pgfqpoint{5.374212in}{2.052406in}}%
\pgfpathlineto{\pgfqpoint{5.376731in}{2.041601in}}%
\pgfpathlineto{\pgfqpoint{5.379249in}{2.028314in}}%
\pgfpathlineto{\pgfqpoint{5.381768in}{2.016640in}}%
\pgfpathlineto{\pgfqpoint{5.384287in}{2.023346in}}%
\pgfpathlineto{\pgfqpoint{5.386806in}{2.044158in}}%
\pgfpathlineto{\pgfqpoint{5.389324in}{2.052242in}}%
\pgfpathlineto{\pgfqpoint{5.391843in}{2.064792in}}%
\pgfpathlineto{\pgfqpoint{5.394362in}{2.065658in}}%
\pgfpathlineto{\pgfqpoint{5.396881in}{2.086902in}}%
\pgfpathlineto{\pgfqpoint{5.399399in}{2.103172in}}%
\pgfpathlineto{\pgfqpoint{5.401918in}{2.101272in}}%
\pgfpathlineto{\pgfqpoint{5.404437in}{2.094408in}}%
\pgfpathlineto{\pgfqpoint{5.406956in}{2.075071in}}%
\pgfpathlineto{\pgfqpoint{5.409474in}{2.057277in}}%
\pgfpathlineto{\pgfqpoint{5.411993in}{2.086588in}}%
\pgfpathlineto{\pgfqpoint{5.414512in}{2.072537in}}%
\pgfpathlineto{\pgfqpoint{5.417031in}{2.030091in}}%
\pgfpathlineto{\pgfqpoint{5.419549in}{2.022135in}}%
\pgfpathlineto{\pgfqpoint{5.422068in}{2.044845in}}%
\pgfpathlineto{\pgfqpoint{5.424587in}{2.057933in}}%
\pgfpathlineto{\pgfqpoint{5.427106in}{2.052583in}}%
\pgfpathlineto{\pgfqpoint{5.429624in}{2.035398in}}%
\pgfpathlineto{\pgfqpoint{5.432143in}{2.037822in}}%
\pgfpathlineto{\pgfqpoint{5.434662in}{2.022035in}}%
\pgfpathlineto{\pgfqpoint{5.437181in}{2.007233in}}%
\pgfpathlineto{\pgfqpoint{5.439699in}{2.002974in}}%
\pgfpathlineto{\pgfqpoint{5.442218in}{2.003475in}}%
\pgfpathlineto{\pgfqpoint{5.444737in}{1.998320in}}%
\pgfpathlineto{\pgfqpoint{5.447256in}{1.982656in}}%
\pgfpathlineto{\pgfqpoint{5.449774in}{1.982439in}}%
\pgfpathlineto{\pgfqpoint{5.452293in}{1.975848in}}%
\pgfpathlineto{\pgfqpoint{5.454812in}{1.975055in}}%
\pgfpathlineto{\pgfqpoint{5.457331in}{1.959010in}}%
\pgfpathlineto{\pgfqpoint{5.459849in}{1.924351in}}%
\pgfpathlineto{\pgfqpoint{5.462368in}{1.936198in}}%
\pgfpathlineto{\pgfqpoint{5.464887in}{1.935814in}}%
\pgfpathlineto{\pgfqpoint{5.467406in}{1.946997in}}%
\pgfpathlineto{\pgfqpoint{5.472443in}{1.981776in}}%
\pgfpathlineto{\pgfqpoint{5.474962in}{1.941442in}}%
\pgfpathlineto{\pgfqpoint{5.477481in}{1.948946in}}%
\pgfpathlineto{\pgfqpoint{5.479999in}{1.972899in}}%
\pgfpathlineto{\pgfqpoint{5.482518in}{1.961656in}}%
\pgfpathlineto{\pgfqpoint{5.485037in}{1.970176in}}%
\pgfpathlineto{\pgfqpoint{5.487556in}{1.963294in}}%
\pgfpathlineto{\pgfqpoint{5.490074in}{1.972053in}}%
\pgfpathlineto{\pgfqpoint{5.492593in}{1.997315in}}%
\pgfpathlineto{\pgfqpoint{5.492593in}{1.997315in}}%
\pgfusepath{stroke}%
\end{pgfscope}%
\begin{pgfscope}%
\pgfpathrectangle{\pgfqpoint{0.455093in}{0.401080in}}{\pgfqpoint{5.037500in}{4.004000in}}%
\pgfusepath{clip}%
\pgfsetrectcap%
\pgfsetroundjoin%
\pgfsetlinewidth{0.501875pt}%
\definecolor{currentstroke}{rgb}{0.690196,0.188235,0.376471}%
\pgfsetstrokecolor{currentstroke}%
\pgfsetstrokeopacity{0.800000}%
\pgfsetdash{}{0pt}%
\pgfpathmoveto{\pgfqpoint{0.455093in}{1.021044in}}%
\pgfpathlineto{\pgfqpoint{0.457612in}{1.011585in}}%
\pgfpathlineto{\pgfqpoint{0.460131in}{1.020196in}}%
\pgfpathlineto{\pgfqpoint{0.462649in}{1.022793in}}%
\pgfpathlineto{\pgfqpoint{0.465168in}{1.009284in}}%
\pgfpathlineto{\pgfqpoint{0.467687in}{0.997475in}}%
\pgfpathlineto{\pgfqpoint{0.470206in}{0.998629in}}%
\pgfpathlineto{\pgfqpoint{0.472724in}{0.991943in}}%
\pgfpathlineto{\pgfqpoint{0.475243in}{0.997918in}}%
\pgfpathlineto{\pgfqpoint{0.477762in}{1.012556in}}%
\pgfpathlineto{\pgfqpoint{0.480281in}{1.020838in}}%
\pgfpathlineto{\pgfqpoint{0.482799in}{1.025401in}}%
\pgfpathlineto{\pgfqpoint{0.485318in}{1.018188in}}%
\pgfpathlineto{\pgfqpoint{0.487837in}{1.025055in}}%
\pgfpathlineto{\pgfqpoint{0.490356in}{1.027058in}}%
\pgfpathlineto{\pgfqpoint{0.492874in}{1.031991in}}%
\pgfpathlineto{\pgfqpoint{0.495393in}{1.037529in}}%
\pgfpathlineto{\pgfqpoint{0.497912in}{1.038085in}}%
\pgfpathlineto{\pgfqpoint{0.500431in}{1.052927in}}%
\pgfpathlineto{\pgfqpoint{0.502949in}{1.054804in}}%
\pgfpathlineto{\pgfqpoint{0.505468in}{1.049960in}}%
\pgfpathlineto{\pgfqpoint{0.507987in}{1.060910in}}%
\pgfpathlineto{\pgfqpoint{0.510506in}{1.044698in}}%
\pgfpathlineto{\pgfqpoint{0.513024in}{1.072945in}}%
\pgfpathlineto{\pgfqpoint{0.515543in}{1.086049in}}%
\pgfpathlineto{\pgfqpoint{0.518062in}{1.077399in}}%
\pgfpathlineto{\pgfqpoint{0.520581in}{1.108678in}}%
\pgfpathlineto{\pgfqpoint{0.523099in}{1.111499in}}%
\pgfpathlineto{\pgfqpoint{0.525618in}{1.080569in}}%
\pgfpathlineto{\pgfqpoint{0.528137in}{1.098155in}}%
\pgfpathlineto{\pgfqpoint{0.530656in}{1.091710in}}%
\pgfpathlineto{\pgfqpoint{0.533174in}{1.083603in}}%
\pgfpathlineto{\pgfqpoint{0.535693in}{1.085685in}}%
\pgfpathlineto{\pgfqpoint{0.538212in}{1.059642in}}%
\pgfpathlineto{\pgfqpoint{0.540731in}{1.075793in}}%
\pgfpathlineto{\pgfqpoint{0.543249in}{1.060062in}}%
\pgfpathlineto{\pgfqpoint{0.545768in}{1.052110in}}%
\pgfpathlineto{\pgfqpoint{0.548287in}{1.066488in}}%
\pgfpathlineto{\pgfqpoint{0.550806in}{1.057958in}}%
\pgfpathlineto{\pgfqpoint{0.553324in}{1.062343in}}%
\pgfpathlineto{\pgfqpoint{0.555843in}{1.047267in}}%
\pgfpathlineto{\pgfqpoint{0.558362in}{1.043680in}}%
\pgfpathlineto{\pgfqpoint{0.560881in}{1.052279in}}%
\pgfpathlineto{\pgfqpoint{0.563399in}{1.054213in}}%
\pgfpathlineto{\pgfqpoint{0.565918in}{1.061299in}}%
\pgfpathlineto{\pgfqpoint{0.568437in}{1.056961in}}%
\pgfpathlineto{\pgfqpoint{0.570956in}{1.050386in}}%
\pgfpathlineto{\pgfqpoint{0.573474in}{1.050475in}}%
\pgfpathlineto{\pgfqpoint{0.575993in}{1.031827in}}%
\pgfpathlineto{\pgfqpoint{0.578512in}{1.030829in}}%
\pgfpathlineto{\pgfqpoint{0.581031in}{1.038497in}}%
\pgfpathlineto{\pgfqpoint{0.583549in}{1.031736in}}%
\pgfpathlineto{\pgfqpoint{0.586068in}{1.036721in}}%
\pgfpathlineto{\pgfqpoint{0.588587in}{1.035499in}}%
\pgfpathlineto{\pgfqpoint{0.591106in}{1.046440in}}%
\pgfpathlineto{\pgfqpoint{0.593624in}{1.056507in}}%
\pgfpathlineto{\pgfqpoint{0.596143in}{1.065381in}}%
\pgfpathlineto{\pgfqpoint{0.598662in}{1.084670in}}%
\pgfpathlineto{\pgfqpoint{0.601181in}{1.090244in}}%
\pgfpathlineto{\pgfqpoint{0.603699in}{1.089681in}}%
\pgfpathlineto{\pgfqpoint{0.606218in}{1.066989in}}%
\pgfpathlineto{\pgfqpoint{0.608737in}{1.054486in}}%
\pgfpathlineto{\pgfqpoint{0.611256in}{1.062413in}}%
\pgfpathlineto{\pgfqpoint{0.613774in}{1.073952in}}%
\pgfpathlineto{\pgfqpoint{0.616293in}{1.068088in}}%
\pgfpathlineto{\pgfqpoint{0.618812in}{1.089227in}}%
\pgfpathlineto{\pgfqpoint{0.621331in}{1.093883in}}%
\pgfpathlineto{\pgfqpoint{0.623849in}{1.079739in}}%
\pgfpathlineto{\pgfqpoint{0.626368in}{1.072698in}}%
\pgfpathlineto{\pgfqpoint{0.628887in}{1.069886in}}%
\pgfpathlineto{\pgfqpoint{0.631406in}{1.076161in}}%
\pgfpathlineto{\pgfqpoint{0.633924in}{1.084831in}}%
\pgfpathlineto{\pgfqpoint{0.636443in}{1.071191in}}%
\pgfpathlineto{\pgfqpoint{0.638962in}{1.064372in}}%
\pgfpathlineto{\pgfqpoint{0.641481in}{1.064141in}}%
\pgfpathlineto{\pgfqpoint{0.643999in}{1.054226in}}%
\pgfpathlineto{\pgfqpoint{0.646518in}{1.066198in}}%
\pgfpathlineto{\pgfqpoint{0.649037in}{1.077532in}}%
\pgfpathlineto{\pgfqpoint{0.651556in}{1.060732in}}%
\pgfpathlineto{\pgfqpoint{0.656593in}{1.081183in}}%
\pgfpathlineto{\pgfqpoint{0.659112in}{1.085878in}}%
\pgfpathlineto{\pgfqpoint{0.661631in}{1.106901in}}%
\pgfpathlineto{\pgfqpoint{0.664149in}{1.096261in}}%
\pgfpathlineto{\pgfqpoint{0.666668in}{1.105564in}}%
\pgfpathlineto{\pgfqpoint{0.669187in}{1.128030in}}%
\pgfpathlineto{\pgfqpoint{0.671706in}{1.154722in}}%
\pgfpathlineto{\pgfqpoint{0.674224in}{1.177568in}}%
\pgfpathlineto{\pgfqpoint{0.676743in}{1.195793in}}%
\pgfpathlineto{\pgfqpoint{0.679262in}{1.208778in}}%
\pgfpathlineto{\pgfqpoint{0.681781in}{1.208807in}}%
\pgfpathlineto{\pgfqpoint{0.684299in}{1.208543in}}%
\pgfpathlineto{\pgfqpoint{0.686818in}{1.211166in}}%
\pgfpathlineto{\pgfqpoint{0.689337in}{1.236468in}}%
\pgfpathlineto{\pgfqpoint{0.691856in}{1.251683in}}%
\pgfpathlineto{\pgfqpoint{0.694374in}{1.231099in}}%
\pgfpathlineto{\pgfqpoint{0.696893in}{1.206700in}}%
\pgfpathlineto{\pgfqpoint{0.699412in}{1.203591in}}%
\pgfpathlineto{\pgfqpoint{0.701931in}{1.203925in}}%
\pgfpathlineto{\pgfqpoint{0.704449in}{1.192415in}}%
\pgfpathlineto{\pgfqpoint{0.706968in}{1.198860in}}%
\pgfpathlineto{\pgfqpoint{0.709487in}{1.201167in}}%
\pgfpathlineto{\pgfqpoint{0.712006in}{1.170643in}}%
\pgfpathlineto{\pgfqpoint{0.714524in}{1.168544in}}%
\pgfpathlineto{\pgfqpoint{0.717043in}{1.143840in}}%
\pgfpathlineto{\pgfqpoint{0.719562in}{1.138301in}}%
\pgfpathlineto{\pgfqpoint{0.722081in}{1.122812in}}%
\pgfpathlineto{\pgfqpoint{0.724599in}{1.122972in}}%
\pgfpathlineto{\pgfqpoint{0.727118in}{1.115454in}}%
\pgfpathlineto{\pgfqpoint{0.729637in}{1.114570in}}%
\pgfpathlineto{\pgfqpoint{0.732156in}{1.124198in}}%
\pgfpathlineto{\pgfqpoint{0.734674in}{1.108886in}}%
\pgfpathlineto{\pgfqpoint{0.737193in}{1.096592in}}%
\pgfpathlineto{\pgfqpoint{0.739712in}{1.092296in}}%
\pgfpathlineto{\pgfqpoint{0.742231in}{1.088206in}}%
\pgfpathlineto{\pgfqpoint{0.744749in}{1.096585in}}%
\pgfpathlineto{\pgfqpoint{0.747268in}{1.088572in}}%
\pgfpathlineto{\pgfqpoint{0.749787in}{1.071415in}}%
\pgfpathlineto{\pgfqpoint{0.752306in}{1.071051in}}%
\pgfpathlineto{\pgfqpoint{0.754824in}{1.046552in}}%
\pgfpathlineto{\pgfqpoint{0.757343in}{1.031285in}}%
\pgfpathlineto{\pgfqpoint{0.759862in}{1.045751in}}%
\pgfpathlineto{\pgfqpoint{0.762381in}{1.043674in}}%
\pgfpathlineto{\pgfqpoint{0.764899in}{1.042114in}}%
\pgfpathlineto{\pgfqpoint{0.767418in}{1.037084in}}%
\pgfpathlineto{\pgfqpoint{0.769937in}{1.028575in}}%
\pgfpathlineto{\pgfqpoint{0.772456in}{1.024670in}}%
\pgfpathlineto{\pgfqpoint{0.774974in}{1.014696in}}%
\pgfpathlineto{\pgfqpoint{0.777493in}{1.003607in}}%
\pgfpathlineto{\pgfqpoint{0.780012in}{1.004460in}}%
\pgfpathlineto{\pgfqpoint{0.782531in}{1.009936in}}%
\pgfpathlineto{\pgfqpoint{0.785049in}{1.017601in}}%
\pgfpathlineto{\pgfqpoint{0.787568in}{1.026030in}}%
\pgfpathlineto{\pgfqpoint{0.790087in}{1.036763in}}%
\pgfpathlineto{\pgfqpoint{0.792606in}{1.032821in}}%
\pgfpathlineto{\pgfqpoint{0.795124in}{1.047997in}}%
\pgfpathlineto{\pgfqpoint{0.797643in}{1.054343in}}%
\pgfpathlineto{\pgfqpoint{0.800162in}{1.045926in}}%
\pgfpathlineto{\pgfqpoint{0.802681in}{1.056526in}}%
\pgfpathlineto{\pgfqpoint{0.805199in}{1.058071in}}%
\pgfpathlineto{\pgfqpoint{0.807718in}{1.078685in}}%
\pgfpathlineto{\pgfqpoint{0.810237in}{1.067776in}}%
\pgfpathlineto{\pgfqpoint{0.812756in}{1.080529in}}%
\pgfpathlineto{\pgfqpoint{0.815274in}{1.088824in}}%
\pgfpathlineto{\pgfqpoint{0.817793in}{1.085998in}}%
\pgfpathlineto{\pgfqpoint{0.820312in}{1.098130in}}%
\pgfpathlineto{\pgfqpoint{0.822831in}{1.126667in}}%
\pgfpathlineto{\pgfqpoint{0.825349in}{1.095422in}}%
\pgfpathlineto{\pgfqpoint{0.827868in}{1.090370in}}%
\pgfpathlineto{\pgfqpoint{0.830387in}{1.084831in}}%
\pgfpathlineto{\pgfqpoint{0.832906in}{1.087249in}}%
\pgfpathlineto{\pgfqpoint{0.835424in}{1.081509in}}%
\pgfpathlineto{\pgfqpoint{0.837943in}{1.073099in}}%
\pgfpathlineto{\pgfqpoint{0.840462in}{1.079703in}}%
\pgfpathlineto{\pgfqpoint{0.842981in}{1.081481in}}%
\pgfpathlineto{\pgfqpoint{0.845499in}{1.068614in}}%
\pgfpathlineto{\pgfqpoint{0.848018in}{1.057669in}}%
\pgfpathlineto{\pgfqpoint{0.850537in}{1.095678in}}%
\pgfpathlineto{\pgfqpoint{0.853056in}{1.066558in}}%
\pgfpathlineto{\pgfqpoint{0.855574in}{1.079467in}}%
\pgfpathlineto{\pgfqpoint{0.858093in}{1.060237in}}%
\pgfpathlineto{\pgfqpoint{0.860612in}{1.046636in}}%
\pgfpathlineto{\pgfqpoint{0.863131in}{1.041958in}}%
\pgfpathlineto{\pgfqpoint{0.865649in}{1.063636in}}%
\pgfpathlineto{\pgfqpoint{0.868168in}{1.056977in}}%
\pgfpathlineto{\pgfqpoint{0.870687in}{1.032613in}}%
\pgfpathlineto{\pgfqpoint{0.873206in}{1.035431in}}%
\pgfpathlineto{\pgfqpoint{0.875724in}{1.034488in}}%
\pgfpathlineto{\pgfqpoint{0.878243in}{1.035268in}}%
\pgfpathlineto{\pgfqpoint{0.880762in}{1.047715in}}%
\pgfpathlineto{\pgfqpoint{0.883281in}{1.048785in}}%
\pgfpathlineto{\pgfqpoint{0.885799in}{1.060595in}}%
\pgfpathlineto{\pgfqpoint{0.888318in}{1.067765in}}%
\pgfpathlineto{\pgfqpoint{0.890837in}{1.059735in}}%
\pgfpathlineto{\pgfqpoint{0.893356in}{1.067682in}}%
\pgfpathlineto{\pgfqpoint{0.895874in}{1.067740in}}%
\pgfpathlineto{\pgfqpoint{0.898393in}{1.081595in}}%
\pgfpathlineto{\pgfqpoint{0.900912in}{1.067457in}}%
\pgfpathlineto{\pgfqpoint{0.903431in}{1.070120in}}%
\pgfpathlineto{\pgfqpoint{0.905949in}{1.105227in}}%
\pgfpathlineto{\pgfqpoint{0.908468in}{1.117005in}}%
\pgfpathlineto{\pgfqpoint{0.910987in}{1.145682in}}%
\pgfpathlineto{\pgfqpoint{0.913506in}{1.155915in}}%
\pgfpathlineto{\pgfqpoint{0.916024in}{1.177087in}}%
\pgfpathlineto{\pgfqpoint{0.918543in}{1.165573in}}%
\pgfpathlineto{\pgfqpoint{0.921062in}{1.158575in}}%
\pgfpathlineto{\pgfqpoint{0.923581in}{1.138605in}}%
\pgfpathlineto{\pgfqpoint{0.926099in}{1.127899in}}%
\pgfpathlineto{\pgfqpoint{0.928618in}{1.107954in}}%
\pgfpathlineto{\pgfqpoint{0.931137in}{1.116147in}}%
\pgfpathlineto{\pgfqpoint{0.933656in}{1.113282in}}%
\pgfpathlineto{\pgfqpoint{0.936174in}{1.129597in}}%
\pgfpathlineto{\pgfqpoint{0.938693in}{1.157974in}}%
\pgfpathlineto{\pgfqpoint{0.941212in}{1.184672in}}%
\pgfpathlineto{\pgfqpoint{0.943731in}{1.191927in}}%
\pgfpathlineto{\pgfqpoint{0.946249in}{1.224073in}}%
\pgfpathlineto{\pgfqpoint{0.948768in}{1.195972in}}%
\pgfpathlineto{\pgfqpoint{0.951287in}{1.193579in}}%
\pgfpathlineto{\pgfqpoint{0.953806in}{1.188586in}}%
\pgfpathlineto{\pgfqpoint{0.956324in}{1.188784in}}%
\pgfpathlineto{\pgfqpoint{0.958843in}{1.174563in}}%
\pgfpathlineto{\pgfqpoint{0.961362in}{1.153748in}}%
\pgfpathlineto{\pgfqpoint{0.963881in}{1.194804in}}%
\pgfpathlineto{\pgfqpoint{0.966399in}{1.192934in}}%
\pgfpathlineto{\pgfqpoint{0.968918in}{1.189224in}}%
\pgfpathlineto{\pgfqpoint{0.971437in}{1.184794in}}%
\pgfpathlineto{\pgfqpoint{0.973956in}{1.172856in}}%
\pgfpathlineto{\pgfqpoint{0.976474in}{1.173938in}}%
\pgfpathlineto{\pgfqpoint{0.978993in}{1.180243in}}%
\pgfpathlineto{\pgfqpoint{0.981512in}{1.170139in}}%
\pgfpathlineto{\pgfqpoint{0.984031in}{1.193751in}}%
\pgfpathlineto{\pgfqpoint{0.986549in}{1.181564in}}%
\pgfpathlineto{\pgfqpoint{0.989068in}{1.185329in}}%
\pgfpathlineto{\pgfqpoint{0.991587in}{1.190563in}}%
\pgfpathlineto{\pgfqpoint{0.994106in}{1.182919in}}%
\pgfpathlineto{\pgfqpoint{0.996624in}{1.176566in}}%
\pgfpathlineto{\pgfqpoint{0.999143in}{1.184767in}}%
\pgfpathlineto{\pgfqpoint{1.001662in}{1.162116in}}%
\pgfpathlineto{\pgfqpoint{1.004181in}{1.166452in}}%
\pgfpathlineto{\pgfqpoint{1.006699in}{1.144537in}}%
\pgfpathlineto{\pgfqpoint{1.009218in}{1.162473in}}%
\pgfpathlineto{\pgfqpoint{1.011737in}{1.146388in}}%
\pgfpathlineto{\pgfqpoint{1.014256in}{1.142326in}}%
\pgfpathlineto{\pgfqpoint{1.016774in}{1.121350in}}%
\pgfpathlineto{\pgfqpoint{1.019293in}{1.118698in}}%
\pgfpathlineto{\pgfqpoint{1.021812in}{1.100672in}}%
\pgfpathlineto{\pgfqpoint{1.024331in}{1.120497in}}%
\pgfpathlineto{\pgfqpoint{1.026849in}{1.113205in}}%
\pgfpathlineto{\pgfqpoint{1.029368in}{1.116015in}}%
\pgfpathlineto{\pgfqpoint{1.031887in}{1.091804in}}%
\pgfpathlineto{\pgfqpoint{1.034406in}{1.083433in}}%
\pgfpathlineto{\pgfqpoint{1.036924in}{1.095826in}}%
\pgfpathlineto{\pgfqpoint{1.039443in}{1.098997in}}%
\pgfpathlineto{\pgfqpoint{1.041962in}{1.093321in}}%
\pgfpathlineto{\pgfqpoint{1.044481in}{1.076593in}}%
\pgfpathlineto{\pgfqpoint{1.046999in}{1.064318in}}%
\pgfpathlineto{\pgfqpoint{1.049518in}{1.054009in}}%
\pgfpathlineto{\pgfqpoint{1.052037in}{1.054412in}}%
\pgfpathlineto{\pgfqpoint{1.054556in}{1.055363in}}%
\pgfpathlineto{\pgfqpoint{1.057074in}{1.051383in}}%
\pgfpathlineto{\pgfqpoint{1.059593in}{1.062087in}}%
\pgfpathlineto{\pgfqpoint{1.062112in}{1.064344in}}%
\pgfpathlineto{\pgfqpoint{1.064631in}{1.075177in}}%
\pgfpathlineto{\pgfqpoint{1.067149in}{1.080318in}}%
\pgfpathlineto{\pgfqpoint{1.069668in}{1.089430in}}%
\pgfpathlineto{\pgfqpoint{1.072187in}{1.101051in}}%
\pgfpathlineto{\pgfqpoint{1.074706in}{1.105457in}}%
\pgfpathlineto{\pgfqpoint{1.077224in}{1.100759in}}%
\pgfpathlineto{\pgfqpoint{1.079743in}{1.082203in}}%
\pgfpathlineto{\pgfqpoint{1.082262in}{1.091481in}}%
\pgfpathlineto{\pgfqpoint{1.084781in}{1.099728in}}%
\pgfpathlineto{\pgfqpoint{1.087299in}{1.085169in}}%
\pgfpathlineto{\pgfqpoint{1.089818in}{1.084412in}}%
\pgfpathlineto{\pgfqpoint{1.092337in}{1.083854in}}%
\pgfpathlineto{\pgfqpoint{1.094856in}{1.106956in}}%
\pgfpathlineto{\pgfqpoint{1.097374in}{1.133110in}}%
\pgfpathlineto{\pgfqpoint{1.099893in}{1.142623in}}%
\pgfpathlineto{\pgfqpoint{1.102412in}{1.150136in}}%
\pgfpathlineto{\pgfqpoint{1.104931in}{1.170086in}}%
\pgfpathlineto{\pgfqpoint{1.107449in}{1.156033in}}%
\pgfpathlineto{\pgfqpoint{1.109968in}{1.150670in}}%
\pgfpathlineto{\pgfqpoint{1.112487in}{1.164679in}}%
\pgfpathlineto{\pgfqpoint{1.115006in}{1.174727in}}%
\pgfpathlineto{\pgfqpoint{1.117524in}{1.180902in}}%
\pgfpathlineto{\pgfqpoint{1.120043in}{1.189958in}}%
\pgfpathlineto{\pgfqpoint{1.122562in}{1.181117in}}%
\pgfpathlineto{\pgfqpoint{1.125081in}{1.180021in}}%
\pgfpathlineto{\pgfqpoint{1.127599in}{1.197159in}}%
\pgfpathlineto{\pgfqpoint{1.130118in}{1.186959in}}%
\pgfpathlineto{\pgfqpoint{1.132637in}{1.174027in}}%
\pgfpathlineto{\pgfqpoint{1.135156in}{1.164912in}}%
\pgfpathlineto{\pgfqpoint{1.137674in}{1.166714in}}%
\pgfpathlineto{\pgfqpoint{1.142712in}{1.169366in}}%
\pgfpathlineto{\pgfqpoint{1.145231in}{1.184799in}}%
\pgfpathlineto{\pgfqpoint{1.147749in}{1.198743in}}%
\pgfpathlineto{\pgfqpoint{1.150268in}{1.224031in}}%
\pgfpathlineto{\pgfqpoint{1.152787in}{1.233509in}}%
\pgfpathlineto{\pgfqpoint{1.155306in}{1.239931in}}%
\pgfpathlineto{\pgfqpoint{1.157824in}{1.257975in}}%
\pgfpathlineto{\pgfqpoint{1.160343in}{1.274713in}}%
\pgfpathlineto{\pgfqpoint{1.162862in}{1.287737in}}%
\pgfpathlineto{\pgfqpoint{1.165381in}{1.299058in}}%
\pgfpathlineto{\pgfqpoint{1.167899in}{1.322869in}}%
\pgfpathlineto{\pgfqpoint{1.170418in}{1.324145in}}%
\pgfpathlineto{\pgfqpoint{1.172937in}{1.328848in}}%
\pgfpathlineto{\pgfqpoint{1.175456in}{1.337665in}}%
\pgfpathlineto{\pgfqpoint{1.177974in}{1.327735in}}%
\pgfpathlineto{\pgfqpoint{1.180493in}{1.324511in}}%
\pgfpathlineto{\pgfqpoint{1.183012in}{1.324259in}}%
\pgfpathlineto{\pgfqpoint{1.185531in}{1.321701in}}%
\pgfpathlineto{\pgfqpoint{1.188049in}{1.303182in}}%
\pgfpathlineto{\pgfqpoint{1.190568in}{1.299051in}}%
\pgfpathlineto{\pgfqpoint{1.193087in}{1.326093in}}%
\pgfpathlineto{\pgfqpoint{1.195606in}{1.318527in}}%
\pgfpathlineto{\pgfqpoint{1.198124in}{1.325979in}}%
\pgfpathlineto{\pgfqpoint{1.200643in}{1.309374in}}%
\pgfpathlineto{\pgfqpoint{1.203162in}{1.307646in}}%
\pgfpathlineto{\pgfqpoint{1.205681in}{1.312478in}}%
\pgfpathlineto{\pgfqpoint{1.208199in}{1.321333in}}%
\pgfpathlineto{\pgfqpoint{1.210718in}{1.336995in}}%
\pgfpathlineto{\pgfqpoint{1.213237in}{1.338973in}}%
\pgfpathlineto{\pgfqpoint{1.215756in}{1.348389in}}%
\pgfpathlineto{\pgfqpoint{1.218274in}{1.356507in}}%
\pgfpathlineto{\pgfqpoint{1.220793in}{1.353023in}}%
\pgfpathlineto{\pgfqpoint{1.223312in}{1.347970in}}%
\pgfpathlineto{\pgfqpoint{1.225831in}{1.358928in}}%
\pgfpathlineto{\pgfqpoint{1.228349in}{1.361391in}}%
\pgfpathlineto{\pgfqpoint{1.230868in}{1.341534in}}%
\pgfpathlineto{\pgfqpoint{1.233387in}{1.354845in}}%
\pgfpathlineto{\pgfqpoint{1.235906in}{1.343951in}}%
\pgfpathlineto{\pgfqpoint{1.238424in}{1.366893in}}%
\pgfpathlineto{\pgfqpoint{1.240943in}{1.354357in}}%
\pgfpathlineto{\pgfqpoint{1.243462in}{1.372817in}}%
\pgfpathlineto{\pgfqpoint{1.245981in}{1.359290in}}%
\pgfpathlineto{\pgfqpoint{1.248499in}{1.351140in}}%
\pgfpathlineto{\pgfqpoint{1.251018in}{1.348463in}}%
\pgfpathlineto{\pgfqpoint{1.253537in}{1.336672in}}%
\pgfpathlineto{\pgfqpoint{1.256056in}{1.357850in}}%
\pgfpathlineto{\pgfqpoint{1.258574in}{1.344861in}}%
\pgfpathlineto{\pgfqpoint{1.261093in}{1.340984in}}%
\pgfpathlineto{\pgfqpoint{1.263612in}{1.330239in}}%
\pgfpathlineto{\pgfqpoint{1.266131in}{1.352409in}}%
\pgfpathlineto{\pgfqpoint{1.268649in}{1.354340in}}%
\pgfpathlineto{\pgfqpoint{1.271168in}{1.356585in}}%
\pgfpathlineto{\pgfqpoint{1.273687in}{1.369428in}}%
\pgfpathlineto{\pgfqpoint{1.276206in}{1.363531in}}%
\pgfpathlineto{\pgfqpoint{1.278724in}{1.365145in}}%
\pgfpathlineto{\pgfqpoint{1.281243in}{1.388842in}}%
\pgfpathlineto{\pgfqpoint{1.283762in}{1.371429in}}%
\pgfpathlineto{\pgfqpoint{1.286281in}{1.382708in}}%
\pgfpathlineto{\pgfqpoint{1.288799in}{1.357302in}}%
\pgfpathlineto{\pgfqpoint{1.291318in}{1.360937in}}%
\pgfpathlineto{\pgfqpoint{1.293837in}{1.380323in}}%
\pgfpathlineto{\pgfqpoint{1.296356in}{1.384693in}}%
\pgfpathlineto{\pgfqpoint{1.298874in}{1.387110in}}%
\pgfpathlineto{\pgfqpoint{1.301393in}{1.374370in}}%
\pgfpathlineto{\pgfqpoint{1.303912in}{1.369009in}}%
\pgfpathlineto{\pgfqpoint{1.306431in}{1.370635in}}%
\pgfpathlineto{\pgfqpoint{1.308949in}{1.384249in}}%
\pgfpathlineto{\pgfqpoint{1.311468in}{1.392339in}}%
\pgfpathlineto{\pgfqpoint{1.313987in}{1.393258in}}%
\pgfpathlineto{\pgfqpoint{1.316506in}{1.392863in}}%
\pgfpathlineto{\pgfqpoint{1.321543in}{1.389173in}}%
\pgfpathlineto{\pgfqpoint{1.324062in}{1.397677in}}%
\pgfpathlineto{\pgfqpoint{1.326581in}{1.380655in}}%
\pgfpathlineto{\pgfqpoint{1.329099in}{1.393420in}}%
\pgfpathlineto{\pgfqpoint{1.331618in}{1.380426in}}%
\pgfpathlineto{\pgfqpoint{1.334137in}{1.373271in}}%
\pgfpathlineto{\pgfqpoint{1.336656in}{1.370216in}}%
\pgfpathlineto{\pgfqpoint{1.339174in}{1.357565in}}%
\pgfpathlineto{\pgfqpoint{1.341693in}{1.364164in}}%
\pgfpathlineto{\pgfqpoint{1.344212in}{1.362512in}}%
\pgfpathlineto{\pgfqpoint{1.346731in}{1.363306in}}%
\pgfpathlineto{\pgfqpoint{1.349249in}{1.367017in}}%
\pgfpathlineto{\pgfqpoint{1.351768in}{1.348184in}}%
\pgfpathlineto{\pgfqpoint{1.354287in}{1.348526in}}%
\pgfpathlineto{\pgfqpoint{1.356806in}{1.339934in}}%
\pgfpathlineto{\pgfqpoint{1.359324in}{1.337119in}}%
\pgfpathlineto{\pgfqpoint{1.361843in}{1.356067in}}%
\pgfpathlineto{\pgfqpoint{1.364362in}{1.362459in}}%
\pgfpathlineto{\pgfqpoint{1.366881in}{1.359440in}}%
\pgfpathlineto{\pgfqpoint{1.369399in}{1.353401in}}%
\pgfpathlineto{\pgfqpoint{1.371918in}{1.348517in}}%
\pgfpathlineto{\pgfqpoint{1.374437in}{1.374350in}}%
\pgfpathlineto{\pgfqpoint{1.376956in}{1.352087in}}%
\pgfpathlineto{\pgfqpoint{1.379474in}{1.358485in}}%
\pgfpathlineto{\pgfqpoint{1.381993in}{1.367502in}}%
\pgfpathlineto{\pgfqpoint{1.384512in}{1.373339in}}%
\pgfpathlineto{\pgfqpoint{1.387031in}{1.415411in}}%
\pgfpathlineto{\pgfqpoint{1.389549in}{1.413837in}}%
\pgfpathlineto{\pgfqpoint{1.392068in}{1.430826in}}%
\pgfpathlineto{\pgfqpoint{1.394587in}{1.438664in}}%
\pgfpathlineto{\pgfqpoint{1.397106in}{1.443871in}}%
\pgfpathlineto{\pgfqpoint{1.399624in}{1.476964in}}%
\pgfpathlineto{\pgfqpoint{1.402143in}{1.475997in}}%
\pgfpathlineto{\pgfqpoint{1.404662in}{1.442592in}}%
\pgfpathlineto{\pgfqpoint{1.407181in}{1.441119in}}%
\pgfpathlineto{\pgfqpoint{1.409699in}{1.466136in}}%
\pgfpathlineto{\pgfqpoint{1.412218in}{1.458396in}}%
\pgfpathlineto{\pgfqpoint{1.414737in}{1.466186in}}%
\pgfpathlineto{\pgfqpoint{1.417256in}{1.471271in}}%
\pgfpathlineto{\pgfqpoint{1.419774in}{1.477175in}}%
\pgfpathlineto{\pgfqpoint{1.422293in}{1.496244in}}%
\pgfpathlineto{\pgfqpoint{1.424812in}{1.464784in}}%
\pgfpathlineto{\pgfqpoint{1.427331in}{1.463291in}}%
\pgfpathlineto{\pgfqpoint{1.429849in}{1.478893in}}%
\pgfpathlineto{\pgfqpoint{1.432368in}{1.482305in}}%
\pgfpathlineto{\pgfqpoint{1.434887in}{1.448299in}}%
\pgfpathlineto{\pgfqpoint{1.437406in}{1.452749in}}%
\pgfpathlineto{\pgfqpoint{1.439924in}{1.449780in}}%
\pgfpathlineto{\pgfqpoint{1.442443in}{1.440837in}}%
\pgfpathlineto{\pgfqpoint{1.444962in}{1.458880in}}%
\pgfpathlineto{\pgfqpoint{1.447481in}{1.454639in}}%
\pgfpathlineto{\pgfqpoint{1.449999in}{1.438696in}}%
\pgfpathlineto{\pgfqpoint{1.452518in}{1.468417in}}%
\pgfpathlineto{\pgfqpoint{1.455037in}{1.433076in}}%
\pgfpathlineto{\pgfqpoint{1.457556in}{1.425778in}}%
\pgfpathlineto{\pgfqpoint{1.460074in}{1.446317in}}%
\pgfpathlineto{\pgfqpoint{1.462593in}{1.460805in}}%
\pgfpathlineto{\pgfqpoint{1.465112in}{1.451958in}}%
\pgfpathlineto{\pgfqpoint{1.467631in}{1.453246in}}%
\pgfpathlineto{\pgfqpoint{1.470149in}{1.445144in}}%
\pgfpathlineto{\pgfqpoint{1.472668in}{1.458060in}}%
\pgfpathlineto{\pgfqpoint{1.475187in}{1.451772in}}%
\pgfpathlineto{\pgfqpoint{1.477706in}{1.459469in}}%
\pgfpathlineto{\pgfqpoint{1.480224in}{1.448412in}}%
\pgfpathlineto{\pgfqpoint{1.482743in}{1.438425in}}%
\pgfpathlineto{\pgfqpoint{1.485262in}{1.437528in}}%
\pgfpathlineto{\pgfqpoint{1.487781in}{1.442898in}}%
\pgfpathlineto{\pgfqpoint{1.490299in}{1.445728in}}%
\pgfpathlineto{\pgfqpoint{1.492818in}{1.445189in}}%
\pgfpathlineto{\pgfqpoint{1.495337in}{1.408224in}}%
\pgfpathlineto{\pgfqpoint{1.497856in}{1.410505in}}%
\pgfpathlineto{\pgfqpoint{1.500374in}{1.408808in}}%
\pgfpathlineto{\pgfqpoint{1.502893in}{1.383055in}}%
\pgfpathlineto{\pgfqpoint{1.505412in}{1.383677in}}%
\pgfpathlineto{\pgfqpoint{1.507931in}{1.395404in}}%
\pgfpathlineto{\pgfqpoint{1.510449in}{1.397686in}}%
\pgfpathlineto{\pgfqpoint{1.512968in}{1.381028in}}%
\pgfpathlineto{\pgfqpoint{1.515487in}{1.395069in}}%
\pgfpathlineto{\pgfqpoint{1.518006in}{1.362007in}}%
\pgfpathlineto{\pgfqpoint{1.520524in}{1.400496in}}%
\pgfpathlineto{\pgfqpoint{1.523043in}{1.413604in}}%
\pgfpathlineto{\pgfqpoint{1.525562in}{1.405029in}}%
\pgfpathlineto{\pgfqpoint{1.528081in}{1.404916in}}%
\pgfpathlineto{\pgfqpoint{1.530599in}{1.407106in}}%
\pgfpathlineto{\pgfqpoint{1.533118in}{1.391190in}}%
\pgfpathlineto{\pgfqpoint{1.535637in}{1.358665in}}%
\pgfpathlineto{\pgfqpoint{1.538156in}{1.315934in}}%
\pgfpathlineto{\pgfqpoint{1.540674in}{1.312218in}}%
\pgfpathlineto{\pgfqpoint{1.543193in}{1.307734in}}%
\pgfpathlineto{\pgfqpoint{1.545712in}{1.315306in}}%
\pgfpathlineto{\pgfqpoint{1.548231in}{1.297592in}}%
\pgfpathlineto{\pgfqpoint{1.550749in}{1.305138in}}%
\pgfpathlineto{\pgfqpoint{1.553268in}{1.300190in}}%
\pgfpathlineto{\pgfqpoint{1.555787in}{1.293750in}}%
\pgfpathlineto{\pgfqpoint{1.558306in}{1.292590in}}%
\pgfpathlineto{\pgfqpoint{1.560824in}{1.298087in}}%
\pgfpathlineto{\pgfqpoint{1.563343in}{1.322762in}}%
\pgfpathlineto{\pgfqpoint{1.565862in}{1.321794in}}%
\pgfpathlineto{\pgfqpoint{1.568381in}{1.352640in}}%
\pgfpathlineto{\pgfqpoint{1.570899in}{1.367201in}}%
\pgfpathlineto{\pgfqpoint{1.573418in}{1.363212in}}%
\pgfpathlineto{\pgfqpoint{1.575937in}{1.381496in}}%
\pgfpathlineto{\pgfqpoint{1.578456in}{1.384612in}}%
\pgfpathlineto{\pgfqpoint{1.580974in}{1.370093in}}%
\pgfpathlineto{\pgfqpoint{1.583493in}{1.351036in}}%
\pgfpathlineto{\pgfqpoint{1.586012in}{1.368528in}}%
\pgfpathlineto{\pgfqpoint{1.588531in}{1.362211in}}%
\pgfpathlineto{\pgfqpoint{1.591049in}{1.361113in}}%
\pgfpathlineto{\pgfqpoint{1.593568in}{1.355532in}}%
\pgfpathlineto{\pgfqpoint{1.596087in}{1.357568in}}%
\pgfpathlineto{\pgfqpoint{1.598606in}{1.353805in}}%
\pgfpathlineto{\pgfqpoint{1.601124in}{1.352373in}}%
\pgfpathlineto{\pgfqpoint{1.603643in}{1.349436in}}%
\pgfpathlineto{\pgfqpoint{1.606162in}{1.344929in}}%
\pgfpathlineto{\pgfqpoint{1.608681in}{1.338669in}}%
\pgfpathlineto{\pgfqpoint{1.611199in}{1.345292in}}%
\pgfpathlineto{\pgfqpoint{1.613718in}{1.357492in}}%
\pgfpathlineto{\pgfqpoint{1.616237in}{1.368239in}}%
\pgfpathlineto{\pgfqpoint{1.618756in}{1.388076in}}%
\pgfpathlineto{\pgfqpoint{1.621274in}{1.361089in}}%
\pgfpathlineto{\pgfqpoint{1.623793in}{1.378428in}}%
\pgfpathlineto{\pgfqpoint{1.626312in}{1.373936in}}%
\pgfpathlineto{\pgfqpoint{1.628831in}{1.374585in}}%
\pgfpathlineto{\pgfqpoint{1.631349in}{1.356375in}}%
\pgfpathlineto{\pgfqpoint{1.633868in}{1.333692in}}%
\pgfpathlineto{\pgfqpoint{1.636387in}{1.308794in}}%
\pgfpathlineto{\pgfqpoint{1.638906in}{1.307453in}}%
\pgfpathlineto{\pgfqpoint{1.641424in}{1.301283in}}%
\pgfpathlineto{\pgfqpoint{1.643943in}{1.280011in}}%
\pgfpathlineto{\pgfqpoint{1.646462in}{1.299862in}}%
\pgfpathlineto{\pgfqpoint{1.648981in}{1.302117in}}%
\pgfpathlineto{\pgfqpoint{1.651499in}{1.319937in}}%
\pgfpathlineto{\pgfqpoint{1.654018in}{1.349468in}}%
\pgfpathlineto{\pgfqpoint{1.656537in}{1.358146in}}%
\pgfpathlineto{\pgfqpoint{1.659056in}{1.370176in}}%
\pgfpathlineto{\pgfqpoint{1.661574in}{1.376496in}}%
\pgfpathlineto{\pgfqpoint{1.664093in}{1.377565in}}%
\pgfpathlineto{\pgfqpoint{1.666612in}{1.377042in}}%
\pgfpathlineto{\pgfqpoint{1.669131in}{1.362449in}}%
\pgfpathlineto{\pgfqpoint{1.671649in}{1.358682in}}%
\pgfpathlineto{\pgfqpoint{1.674168in}{1.373753in}}%
\pgfpathlineto{\pgfqpoint{1.676687in}{1.380095in}}%
\pgfpathlineto{\pgfqpoint{1.679206in}{1.358176in}}%
\pgfpathlineto{\pgfqpoint{1.681724in}{1.357105in}}%
\pgfpathlineto{\pgfqpoint{1.684243in}{1.355035in}}%
\pgfpathlineto{\pgfqpoint{1.686762in}{1.364131in}}%
\pgfpathlineto{\pgfqpoint{1.689281in}{1.379621in}}%
\pgfpathlineto{\pgfqpoint{1.691799in}{1.392701in}}%
\pgfpathlineto{\pgfqpoint{1.694318in}{1.391078in}}%
\pgfpathlineto{\pgfqpoint{1.696837in}{1.405555in}}%
\pgfpathlineto{\pgfqpoint{1.699356in}{1.414473in}}%
\pgfpathlineto{\pgfqpoint{1.701874in}{1.419978in}}%
\pgfpathlineto{\pgfqpoint{1.704393in}{1.399880in}}%
\pgfpathlineto{\pgfqpoint{1.706912in}{1.381184in}}%
\pgfpathlineto{\pgfqpoint{1.709431in}{1.388828in}}%
\pgfpathlineto{\pgfqpoint{1.711949in}{1.390288in}}%
\pgfpathlineto{\pgfqpoint{1.714468in}{1.419924in}}%
\pgfpathlineto{\pgfqpoint{1.716987in}{1.409468in}}%
\pgfpathlineto{\pgfqpoint{1.719506in}{1.406233in}}%
\pgfpathlineto{\pgfqpoint{1.722024in}{1.407077in}}%
\pgfpathlineto{\pgfqpoint{1.724543in}{1.400608in}}%
\pgfpathlineto{\pgfqpoint{1.729581in}{1.392943in}}%
\pgfpathlineto{\pgfqpoint{1.732099in}{1.415774in}}%
\pgfpathlineto{\pgfqpoint{1.734618in}{1.389319in}}%
\pgfpathlineto{\pgfqpoint{1.737137in}{1.416508in}}%
\pgfpathlineto{\pgfqpoint{1.739656in}{1.414713in}}%
\pgfpathlineto{\pgfqpoint{1.742174in}{1.447291in}}%
\pgfpathlineto{\pgfqpoint{1.744693in}{1.462927in}}%
\pgfpathlineto{\pgfqpoint{1.747212in}{1.474646in}}%
\pgfpathlineto{\pgfqpoint{1.749731in}{1.446489in}}%
\pgfpathlineto{\pgfqpoint{1.752249in}{1.421288in}}%
\pgfpathlineto{\pgfqpoint{1.754768in}{1.414253in}}%
\pgfpathlineto{\pgfqpoint{1.757287in}{1.397434in}}%
\pgfpathlineto{\pgfqpoint{1.759806in}{1.395985in}}%
\pgfpathlineto{\pgfqpoint{1.762324in}{1.366294in}}%
\pgfpathlineto{\pgfqpoint{1.764843in}{1.361305in}}%
\pgfpathlineto{\pgfqpoint{1.767362in}{1.341528in}}%
\pgfpathlineto{\pgfqpoint{1.769881in}{1.358003in}}%
\pgfpathlineto{\pgfqpoint{1.772399in}{1.363856in}}%
\pgfpathlineto{\pgfqpoint{1.774918in}{1.365215in}}%
\pgfpathlineto{\pgfqpoint{1.777437in}{1.354709in}}%
\pgfpathlineto{\pgfqpoint{1.779956in}{1.350897in}}%
\pgfpathlineto{\pgfqpoint{1.782474in}{1.325519in}}%
\pgfpathlineto{\pgfqpoint{1.784993in}{1.326197in}}%
\pgfpathlineto{\pgfqpoint{1.787512in}{1.324909in}}%
\pgfpathlineto{\pgfqpoint{1.790031in}{1.336112in}}%
\pgfpathlineto{\pgfqpoint{1.792549in}{1.329390in}}%
\pgfpathlineto{\pgfqpoint{1.795068in}{1.347868in}}%
\pgfpathlineto{\pgfqpoint{1.797587in}{1.349142in}}%
\pgfpathlineto{\pgfqpoint{1.800106in}{1.338391in}}%
\pgfpathlineto{\pgfqpoint{1.802624in}{1.366083in}}%
\pgfpathlineto{\pgfqpoint{1.805143in}{1.338208in}}%
\pgfpathlineto{\pgfqpoint{1.807662in}{1.326670in}}%
\pgfpathlineto{\pgfqpoint{1.810181in}{1.333906in}}%
\pgfpathlineto{\pgfqpoint{1.812699in}{1.337473in}}%
\pgfpathlineto{\pgfqpoint{1.815218in}{1.336464in}}%
\pgfpathlineto{\pgfqpoint{1.817737in}{1.342561in}}%
\pgfpathlineto{\pgfqpoint{1.820256in}{1.347117in}}%
\pgfpathlineto{\pgfqpoint{1.822774in}{1.356956in}}%
\pgfpathlineto{\pgfqpoint{1.825293in}{1.341060in}}%
\pgfpathlineto{\pgfqpoint{1.827812in}{1.343545in}}%
\pgfpathlineto{\pgfqpoint{1.830331in}{1.343140in}}%
\pgfpathlineto{\pgfqpoint{1.832849in}{1.348877in}}%
\pgfpathlineto{\pgfqpoint{1.835368in}{1.337100in}}%
\pgfpathlineto{\pgfqpoint{1.837887in}{1.329029in}}%
\pgfpathlineto{\pgfqpoint{1.840406in}{1.311143in}}%
\pgfpathlineto{\pgfqpoint{1.842924in}{1.332171in}}%
\pgfpathlineto{\pgfqpoint{1.845443in}{1.337625in}}%
\pgfpathlineto{\pgfqpoint{1.847962in}{1.341870in}}%
\pgfpathlineto{\pgfqpoint{1.850481in}{1.344479in}}%
\pgfpathlineto{\pgfqpoint{1.852999in}{1.348464in}}%
\pgfpathlineto{\pgfqpoint{1.855518in}{1.356933in}}%
\pgfpathlineto{\pgfqpoint{1.858037in}{1.369747in}}%
\pgfpathlineto{\pgfqpoint{1.860556in}{1.352773in}}%
\pgfpathlineto{\pgfqpoint{1.863074in}{1.344905in}}%
\pgfpathlineto{\pgfqpoint{1.865593in}{1.335514in}}%
\pgfpathlineto{\pgfqpoint{1.868112in}{1.352638in}}%
\pgfpathlineto{\pgfqpoint{1.870631in}{1.356337in}}%
\pgfpathlineto{\pgfqpoint{1.873149in}{1.363347in}}%
\pgfpathlineto{\pgfqpoint{1.875668in}{1.378357in}}%
\pgfpathlineto{\pgfqpoint{1.878187in}{1.382513in}}%
\pgfpathlineto{\pgfqpoint{1.880706in}{1.385730in}}%
\pgfpathlineto{\pgfqpoint{1.883224in}{1.403829in}}%
\pgfpathlineto{\pgfqpoint{1.885743in}{1.409704in}}%
\pgfpathlineto{\pgfqpoint{1.888262in}{1.408875in}}%
\pgfpathlineto{\pgfqpoint{1.890781in}{1.415882in}}%
\pgfpathlineto{\pgfqpoint{1.893299in}{1.402657in}}%
\pgfpathlineto{\pgfqpoint{1.895818in}{1.417654in}}%
\pgfpathlineto{\pgfqpoint{1.898337in}{1.413772in}}%
\pgfpathlineto{\pgfqpoint{1.900856in}{1.404621in}}%
\pgfpathlineto{\pgfqpoint{1.903374in}{1.396199in}}%
\pgfpathlineto{\pgfqpoint{1.905893in}{1.422111in}}%
\pgfpathlineto{\pgfqpoint{1.908412in}{1.440197in}}%
\pgfpathlineto{\pgfqpoint{1.910931in}{1.460259in}}%
\pgfpathlineto{\pgfqpoint{1.913449in}{1.474493in}}%
\pgfpathlineto{\pgfqpoint{1.915968in}{1.473275in}}%
\pgfpathlineto{\pgfqpoint{1.918487in}{1.499690in}}%
\pgfpathlineto{\pgfqpoint{1.921006in}{1.495679in}}%
\pgfpathlineto{\pgfqpoint{1.923524in}{1.505168in}}%
\pgfpathlineto{\pgfqpoint{1.928562in}{1.516404in}}%
\pgfpathlineto{\pgfqpoint{1.931081in}{1.505176in}}%
\pgfpathlineto{\pgfqpoint{1.933599in}{1.485904in}}%
\pgfpathlineto{\pgfqpoint{1.936118in}{1.465506in}}%
\pgfpathlineto{\pgfqpoint{1.938637in}{1.464827in}}%
\pgfpathlineto{\pgfqpoint{1.941156in}{1.476950in}}%
\pgfpathlineto{\pgfqpoint{1.943674in}{1.474413in}}%
\pgfpathlineto{\pgfqpoint{1.946193in}{1.498104in}}%
\pgfpathlineto{\pgfqpoint{1.948712in}{1.491137in}}%
\pgfpathlineto{\pgfqpoint{1.951231in}{1.488377in}}%
\pgfpathlineto{\pgfqpoint{1.953749in}{1.509543in}}%
\pgfpathlineto{\pgfqpoint{1.956268in}{1.499946in}}%
\pgfpathlineto{\pgfqpoint{1.958787in}{1.528009in}}%
\pgfpathlineto{\pgfqpoint{1.961306in}{1.528236in}}%
\pgfpathlineto{\pgfqpoint{1.963824in}{1.543226in}}%
\pgfpathlineto{\pgfqpoint{1.966343in}{1.544158in}}%
\pgfpathlineto{\pgfqpoint{1.968862in}{1.547965in}}%
\pgfpathlineto{\pgfqpoint{1.971381in}{1.544294in}}%
\pgfpathlineto{\pgfqpoint{1.973899in}{1.544703in}}%
\pgfpathlineto{\pgfqpoint{1.976418in}{1.555417in}}%
\pgfpathlineto{\pgfqpoint{1.978937in}{1.563434in}}%
\pgfpathlineto{\pgfqpoint{1.981456in}{1.573283in}}%
\pgfpathlineto{\pgfqpoint{1.983974in}{1.593220in}}%
\pgfpathlineto{\pgfqpoint{1.986493in}{1.614349in}}%
\pgfpathlineto{\pgfqpoint{1.989012in}{1.618234in}}%
\pgfpathlineto{\pgfqpoint{1.991531in}{1.636142in}}%
\pgfpathlineto{\pgfqpoint{1.994049in}{1.656880in}}%
\pgfpathlineto{\pgfqpoint{1.996568in}{1.627932in}}%
\pgfpathlineto{\pgfqpoint{1.999087in}{1.608437in}}%
\pgfpathlineto{\pgfqpoint{2.001606in}{1.633276in}}%
\pgfpathlineto{\pgfqpoint{2.004124in}{1.624827in}}%
\pgfpathlineto{\pgfqpoint{2.006643in}{1.628489in}}%
\pgfpathlineto{\pgfqpoint{2.009162in}{1.612565in}}%
\pgfpathlineto{\pgfqpoint{2.011681in}{1.637352in}}%
\pgfpathlineto{\pgfqpoint{2.014199in}{1.665076in}}%
\pgfpathlineto{\pgfqpoint{2.016718in}{1.666768in}}%
\pgfpathlineto{\pgfqpoint{2.019237in}{1.656889in}}%
\pgfpathlineto{\pgfqpoint{2.021756in}{1.662371in}}%
\pgfpathlineto{\pgfqpoint{2.024274in}{1.657792in}}%
\pgfpathlineto{\pgfqpoint{2.026793in}{1.649646in}}%
\pgfpathlineto{\pgfqpoint{2.029312in}{1.668887in}}%
\pgfpathlineto{\pgfqpoint{2.031831in}{1.646504in}}%
\pgfpathlineto{\pgfqpoint{2.034349in}{1.642412in}}%
\pgfpathlineto{\pgfqpoint{2.036868in}{1.617797in}}%
\pgfpathlineto{\pgfqpoint{2.039387in}{1.599807in}}%
\pgfpathlineto{\pgfqpoint{2.041906in}{1.585085in}}%
\pgfpathlineto{\pgfqpoint{2.044424in}{1.551198in}}%
\pgfpathlineto{\pgfqpoint{2.046943in}{1.569971in}}%
\pgfpathlineto{\pgfqpoint{2.049462in}{1.559662in}}%
\pgfpathlineto{\pgfqpoint{2.051981in}{1.548076in}}%
\pgfpathlineto{\pgfqpoint{2.054499in}{1.545511in}}%
\pgfpathlineto{\pgfqpoint{2.057018in}{1.543755in}}%
\pgfpathlineto{\pgfqpoint{2.059537in}{1.538019in}}%
\pgfpathlineto{\pgfqpoint{2.062056in}{1.536835in}}%
\pgfpathlineto{\pgfqpoint{2.064574in}{1.524650in}}%
\pgfpathlineto{\pgfqpoint{2.067093in}{1.517166in}}%
\pgfpathlineto{\pgfqpoint{2.069612in}{1.498425in}}%
\pgfpathlineto{\pgfqpoint{2.072131in}{1.495801in}}%
\pgfpathlineto{\pgfqpoint{2.074649in}{1.474387in}}%
\pgfpathlineto{\pgfqpoint{2.077168in}{1.462816in}}%
\pgfpathlineto{\pgfqpoint{2.079687in}{1.471995in}}%
\pgfpathlineto{\pgfqpoint{2.082206in}{1.472845in}}%
\pgfpathlineto{\pgfqpoint{2.084724in}{1.469714in}}%
\pgfpathlineto{\pgfqpoint{2.087243in}{1.455053in}}%
\pgfpathlineto{\pgfqpoint{2.089762in}{1.457429in}}%
\pgfpathlineto{\pgfqpoint{2.092281in}{1.458571in}}%
\pgfpathlineto{\pgfqpoint{2.094799in}{1.474216in}}%
\pgfpathlineto{\pgfqpoint{2.097318in}{1.484027in}}%
\pgfpathlineto{\pgfqpoint{2.099837in}{1.516473in}}%
\pgfpathlineto{\pgfqpoint{2.102356in}{1.514948in}}%
\pgfpathlineto{\pgfqpoint{2.104874in}{1.532100in}}%
\pgfpathlineto{\pgfqpoint{2.107393in}{1.511966in}}%
\pgfpathlineto{\pgfqpoint{2.109912in}{1.522995in}}%
\pgfpathlineto{\pgfqpoint{2.112431in}{1.491840in}}%
\pgfpathlineto{\pgfqpoint{2.114949in}{1.497458in}}%
\pgfpathlineto{\pgfqpoint{2.117468in}{1.477673in}}%
\pgfpathlineto{\pgfqpoint{2.119987in}{1.478989in}}%
\pgfpathlineto{\pgfqpoint{2.122506in}{1.462772in}}%
\pgfpathlineto{\pgfqpoint{2.125024in}{1.461306in}}%
\pgfpathlineto{\pgfqpoint{2.127543in}{1.476412in}}%
\pgfpathlineto{\pgfqpoint{2.130062in}{1.471352in}}%
\pgfpathlineto{\pgfqpoint{2.132581in}{1.486388in}}%
\pgfpathlineto{\pgfqpoint{2.135099in}{1.480920in}}%
\pgfpathlineto{\pgfqpoint{2.137618in}{1.489504in}}%
\pgfpathlineto{\pgfqpoint{2.140137in}{1.493966in}}%
\pgfpathlineto{\pgfqpoint{2.142656in}{1.480576in}}%
\pgfpathlineto{\pgfqpoint{2.145174in}{1.460363in}}%
\pgfpathlineto{\pgfqpoint{2.147693in}{1.429294in}}%
\pgfpathlineto{\pgfqpoint{2.150212in}{1.456114in}}%
\pgfpathlineto{\pgfqpoint{2.152731in}{1.478720in}}%
\pgfpathlineto{\pgfqpoint{2.155249in}{1.473228in}}%
\pgfpathlineto{\pgfqpoint{2.157768in}{1.467185in}}%
\pgfpathlineto{\pgfqpoint{2.160287in}{1.473875in}}%
\pgfpathlineto{\pgfqpoint{2.162806in}{1.481760in}}%
\pgfpathlineto{\pgfqpoint{2.165324in}{1.473188in}}%
\pgfpathlineto{\pgfqpoint{2.167843in}{1.448267in}}%
\pgfpathlineto{\pgfqpoint{2.170362in}{1.455314in}}%
\pgfpathlineto{\pgfqpoint{2.172881in}{1.487461in}}%
\pgfpathlineto{\pgfqpoint{2.175399in}{1.466273in}}%
\pgfpathlineto{\pgfqpoint{2.177918in}{1.465894in}}%
\pgfpathlineto{\pgfqpoint{2.180437in}{1.460254in}}%
\pgfpathlineto{\pgfqpoint{2.182956in}{1.478324in}}%
\pgfpathlineto{\pgfqpoint{2.185474in}{1.483106in}}%
\pgfpathlineto{\pgfqpoint{2.187993in}{1.472796in}}%
\pgfpathlineto{\pgfqpoint{2.190512in}{1.463171in}}%
\pgfpathlineto{\pgfqpoint{2.193031in}{1.443223in}}%
\pgfpathlineto{\pgfqpoint{2.195549in}{1.444588in}}%
\pgfpathlineto{\pgfqpoint{2.198068in}{1.460684in}}%
\pgfpathlineto{\pgfqpoint{2.200587in}{1.463640in}}%
\pgfpathlineto{\pgfqpoint{2.203106in}{1.453046in}}%
\pgfpathlineto{\pgfqpoint{2.205624in}{1.423857in}}%
\pgfpathlineto{\pgfqpoint{2.208143in}{1.436369in}}%
\pgfpathlineto{\pgfqpoint{2.210662in}{1.426120in}}%
\pgfpathlineto{\pgfqpoint{2.213181in}{1.430647in}}%
\pgfpathlineto{\pgfqpoint{2.215699in}{1.437630in}}%
\pgfpathlineto{\pgfqpoint{2.218218in}{1.432260in}}%
\pgfpathlineto{\pgfqpoint{2.220737in}{1.444376in}}%
\pgfpathlineto{\pgfqpoint{2.223256in}{1.428963in}}%
\pgfpathlineto{\pgfqpoint{2.225774in}{1.437477in}}%
\pgfpathlineto{\pgfqpoint{2.228293in}{1.443197in}}%
\pgfpathlineto{\pgfqpoint{2.230812in}{1.433686in}}%
\pgfpathlineto{\pgfqpoint{2.233331in}{1.430617in}}%
\pgfpathlineto{\pgfqpoint{2.235849in}{1.431195in}}%
\pgfpathlineto{\pgfqpoint{2.238368in}{1.420770in}}%
\pgfpathlineto{\pgfqpoint{2.240887in}{1.421720in}}%
\pgfpathlineto{\pgfqpoint{2.243406in}{1.405592in}}%
\pgfpathlineto{\pgfqpoint{2.245924in}{1.384304in}}%
\pgfpathlineto{\pgfqpoint{2.248443in}{1.371240in}}%
\pgfpathlineto{\pgfqpoint{2.250962in}{1.364279in}}%
\pgfpathlineto{\pgfqpoint{2.253481in}{1.369538in}}%
\pgfpathlineto{\pgfqpoint{2.255999in}{1.352032in}}%
\pgfpathlineto{\pgfqpoint{2.258518in}{1.353484in}}%
\pgfpathlineto{\pgfqpoint{2.261037in}{1.322419in}}%
\pgfpathlineto{\pgfqpoint{2.263556in}{1.317677in}}%
\pgfpathlineto{\pgfqpoint{2.266074in}{1.334046in}}%
\pgfpathlineto{\pgfqpoint{2.268593in}{1.315329in}}%
\pgfpathlineto{\pgfqpoint{2.271112in}{1.333333in}}%
\pgfpathlineto{\pgfqpoint{2.273631in}{1.344540in}}%
\pgfpathlineto{\pgfqpoint{2.276149in}{1.350986in}}%
\pgfpathlineto{\pgfqpoint{2.278668in}{1.356496in}}%
\pgfpathlineto{\pgfqpoint{2.281187in}{1.370092in}}%
\pgfpathlineto{\pgfqpoint{2.283706in}{1.372411in}}%
\pgfpathlineto{\pgfqpoint{2.286224in}{1.402963in}}%
\pgfpathlineto{\pgfqpoint{2.288743in}{1.393260in}}%
\pgfpathlineto{\pgfqpoint{2.291262in}{1.378750in}}%
\pgfpathlineto{\pgfqpoint{2.293781in}{1.375815in}}%
\pgfpathlineto{\pgfqpoint{2.296299in}{1.364929in}}%
\pgfpathlineto{\pgfqpoint{2.298818in}{1.365079in}}%
\pgfpathlineto{\pgfqpoint{2.301337in}{1.366176in}}%
\pgfpathlineto{\pgfqpoint{2.303856in}{1.371731in}}%
\pgfpathlineto{\pgfqpoint{2.306374in}{1.373139in}}%
\pgfpathlineto{\pgfqpoint{2.308893in}{1.372367in}}%
\pgfpathlineto{\pgfqpoint{2.311412in}{1.377416in}}%
\pgfpathlineto{\pgfqpoint{2.313931in}{1.377135in}}%
\pgfpathlineto{\pgfqpoint{2.316449in}{1.373417in}}%
\pgfpathlineto{\pgfqpoint{2.318968in}{1.391873in}}%
\pgfpathlineto{\pgfqpoint{2.321487in}{1.405533in}}%
\pgfpathlineto{\pgfqpoint{2.324006in}{1.407335in}}%
\pgfpathlineto{\pgfqpoint{2.326524in}{1.398069in}}%
\pgfpathlineto{\pgfqpoint{2.329043in}{1.409367in}}%
\pgfpathlineto{\pgfqpoint{2.331562in}{1.397129in}}%
\pgfpathlineto{\pgfqpoint{2.334081in}{1.389911in}}%
\pgfpathlineto{\pgfqpoint{2.336599in}{1.390518in}}%
\pgfpathlineto{\pgfqpoint{2.339118in}{1.398967in}}%
\pgfpathlineto{\pgfqpoint{2.341637in}{1.420055in}}%
\pgfpathlineto{\pgfqpoint{2.344156in}{1.431841in}}%
\pgfpathlineto{\pgfqpoint{2.346674in}{1.436607in}}%
\pgfpathlineto{\pgfqpoint{2.349193in}{1.435826in}}%
\pgfpathlineto{\pgfqpoint{2.351712in}{1.445364in}}%
\pgfpathlineto{\pgfqpoint{2.354231in}{1.436256in}}%
\pgfpathlineto{\pgfqpoint{2.356749in}{1.441092in}}%
\pgfpathlineto{\pgfqpoint{2.359268in}{1.446328in}}%
\pgfpathlineto{\pgfqpoint{2.361787in}{1.454665in}}%
\pgfpathlineto{\pgfqpoint{2.364306in}{1.462447in}}%
\pgfpathlineto{\pgfqpoint{2.366824in}{1.485490in}}%
\pgfpathlineto{\pgfqpoint{2.369343in}{1.474804in}}%
\pgfpathlineto{\pgfqpoint{2.371862in}{1.480645in}}%
\pgfpathlineto{\pgfqpoint{2.374381in}{1.499953in}}%
\pgfpathlineto{\pgfqpoint{2.376899in}{1.511366in}}%
\pgfpathlineto{\pgfqpoint{2.379418in}{1.511204in}}%
\pgfpathlineto{\pgfqpoint{2.381937in}{1.508583in}}%
\pgfpathlineto{\pgfqpoint{2.384456in}{1.523898in}}%
\pgfpathlineto{\pgfqpoint{2.386974in}{1.508431in}}%
\pgfpathlineto{\pgfqpoint{2.389493in}{1.520259in}}%
\pgfpathlineto{\pgfqpoint{2.392012in}{1.522369in}}%
\pgfpathlineto{\pgfqpoint{2.394531in}{1.511231in}}%
\pgfpathlineto{\pgfqpoint{2.397049in}{1.493705in}}%
\pgfpathlineto{\pgfqpoint{2.399568in}{1.490246in}}%
\pgfpathlineto{\pgfqpoint{2.402087in}{1.496450in}}%
\pgfpathlineto{\pgfqpoint{2.404606in}{1.498550in}}%
\pgfpathlineto{\pgfqpoint{2.407124in}{1.500967in}}%
\pgfpathlineto{\pgfqpoint{2.409643in}{1.517150in}}%
\pgfpathlineto{\pgfqpoint{2.412162in}{1.499578in}}%
\pgfpathlineto{\pgfqpoint{2.414681in}{1.477780in}}%
\pgfpathlineto{\pgfqpoint{2.417199in}{1.460249in}}%
\pgfpathlineto{\pgfqpoint{2.419718in}{1.465785in}}%
\pgfpathlineto{\pgfqpoint{2.422237in}{1.472214in}}%
\pgfpathlineto{\pgfqpoint{2.424756in}{1.466907in}}%
\pgfpathlineto{\pgfqpoint{2.427274in}{1.457110in}}%
\pgfpathlineto{\pgfqpoint{2.429793in}{1.471713in}}%
\pgfpathlineto{\pgfqpoint{2.432312in}{1.447823in}}%
\pgfpathlineto{\pgfqpoint{2.434831in}{1.449071in}}%
\pgfpathlineto{\pgfqpoint{2.437349in}{1.421528in}}%
\pgfpathlineto{\pgfqpoint{2.439868in}{1.440520in}}%
\pgfpathlineto{\pgfqpoint{2.442387in}{1.426683in}}%
\pgfpathlineto{\pgfqpoint{2.444906in}{1.420256in}}%
\pgfpathlineto{\pgfqpoint{2.447424in}{1.445323in}}%
\pgfpathlineto{\pgfqpoint{2.449943in}{1.435578in}}%
\pgfpathlineto{\pgfqpoint{2.452462in}{1.436235in}}%
\pgfpathlineto{\pgfqpoint{2.454981in}{1.453284in}}%
\pgfpathlineto{\pgfqpoint{2.457499in}{1.432571in}}%
\pgfpathlineto{\pgfqpoint{2.460018in}{1.456993in}}%
\pgfpathlineto{\pgfqpoint{2.462537in}{1.446036in}}%
\pgfpathlineto{\pgfqpoint{2.465056in}{1.417469in}}%
\pgfpathlineto{\pgfqpoint{2.467574in}{1.402320in}}%
\pgfpathlineto{\pgfqpoint{2.470093in}{1.419269in}}%
\pgfpathlineto{\pgfqpoint{2.472612in}{1.421050in}}%
\pgfpathlineto{\pgfqpoint{2.475131in}{1.426245in}}%
\pgfpathlineto{\pgfqpoint{2.477649in}{1.419574in}}%
\pgfpathlineto{\pgfqpoint{2.480168in}{1.426739in}}%
\pgfpathlineto{\pgfqpoint{2.482687in}{1.426066in}}%
\pgfpathlineto{\pgfqpoint{2.485206in}{1.429088in}}%
\pgfpathlineto{\pgfqpoint{2.487724in}{1.457314in}}%
\pgfpathlineto{\pgfqpoint{2.490243in}{1.438295in}}%
\pgfpathlineto{\pgfqpoint{2.492762in}{1.429036in}}%
\pgfpathlineto{\pgfqpoint{2.495281in}{1.428607in}}%
\pgfpathlineto{\pgfqpoint{2.497799in}{1.440592in}}%
\pgfpathlineto{\pgfqpoint{2.500318in}{1.422053in}}%
\pgfpathlineto{\pgfqpoint{2.502837in}{1.408233in}}%
\pgfpathlineto{\pgfqpoint{2.505356in}{1.416493in}}%
\pgfpathlineto{\pgfqpoint{2.507874in}{1.418991in}}%
\pgfpathlineto{\pgfqpoint{2.510393in}{1.394842in}}%
\pgfpathlineto{\pgfqpoint{2.512912in}{1.382340in}}%
\pgfpathlineto{\pgfqpoint{2.515431in}{1.378013in}}%
\pgfpathlineto{\pgfqpoint{2.517949in}{1.372906in}}%
\pgfpathlineto{\pgfqpoint{2.520468in}{1.368177in}}%
\pgfpathlineto{\pgfqpoint{2.522987in}{1.374374in}}%
\pgfpathlineto{\pgfqpoint{2.525506in}{1.385596in}}%
\pgfpathlineto{\pgfqpoint{2.528024in}{1.402571in}}%
\pgfpathlineto{\pgfqpoint{2.530543in}{1.409162in}}%
\pgfpathlineto{\pgfqpoint{2.533062in}{1.401514in}}%
\pgfpathlineto{\pgfqpoint{2.535581in}{1.401361in}}%
\pgfpathlineto{\pgfqpoint{2.538099in}{1.431564in}}%
\pgfpathlineto{\pgfqpoint{2.540618in}{1.448585in}}%
\pgfpathlineto{\pgfqpoint{2.543137in}{1.436036in}}%
\pgfpathlineto{\pgfqpoint{2.545656in}{1.415788in}}%
\pgfpathlineto{\pgfqpoint{2.548174in}{1.412781in}}%
\pgfpathlineto{\pgfqpoint{2.550693in}{1.414280in}}%
\pgfpathlineto{\pgfqpoint{2.553212in}{1.404685in}}%
\pgfpathlineto{\pgfqpoint{2.555731in}{1.396441in}}%
\pgfpathlineto{\pgfqpoint{2.558249in}{1.399886in}}%
\pgfpathlineto{\pgfqpoint{2.560768in}{1.414153in}}%
\pgfpathlineto{\pgfqpoint{2.563287in}{1.412338in}}%
\pgfpathlineto{\pgfqpoint{2.565806in}{1.418358in}}%
\pgfpathlineto{\pgfqpoint{2.568324in}{1.419449in}}%
\pgfpathlineto{\pgfqpoint{2.570843in}{1.438456in}}%
\pgfpathlineto{\pgfqpoint{2.573362in}{1.448868in}}%
\pgfpathlineto{\pgfqpoint{2.575881in}{1.460348in}}%
\pgfpathlineto{\pgfqpoint{2.578399in}{1.465713in}}%
\pgfpathlineto{\pgfqpoint{2.580918in}{1.443895in}}%
\pgfpathlineto{\pgfqpoint{2.583437in}{1.440788in}}%
\pgfpathlineto{\pgfqpoint{2.585956in}{1.428609in}}%
\pgfpathlineto{\pgfqpoint{2.588474in}{1.445461in}}%
\pgfpathlineto{\pgfqpoint{2.593512in}{1.487170in}}%
\pgfpathlineto{\pgfqpoint{2.596031in}{1.508262in}}%
\pgfpathlineto{\pgfqpoint{2.598549in}{1.535710in}}%
\pgfpathlineto{\pgfqpoint{2.601068in}{1.537570in}}%
\pgfpathlineto{\pgfqpoint{2.603587in}{1.555187in}}%
\pgfpathlineto{\pgfqpoint{2.606106in}{1.533965in}}%
\pgfpathlineto{\pgfqpoint{2.608624in}{1.529784in}}%
\pgfpathlineto{\pgfqpoint{2.611143in}{1.528622in}}%
\pgfpathlineto{\pgfqpoint{2.613662in}{1.546577in}}%
\pgfpathlineto{\pgfqpoint{2.616181in}{1.572387in}}%
\pgfpathlineto{\pgfqpoint{2.618699in}{1.549650in}}%
\pgfpathlineto{\pgfqpoint{2.621218in}{1.580703in}}%
\pgfpathlineto{\pgfqpoint{2.623737in}{1.580664in}}%
\pgfpathlineto{\pgfqpoint{2.626256in}{1.559549in}}%
\pgfpathlineto{\pgfqpoint{2.628774in}{1.562370in}}%
\pgfpathlineto{\pgfqpoint{2.631293in}{1.558848in}}%
\pgfpathlineto{\pgfqpoint{2.636331in}{1.532730in}}%
\pgfpathlineto{\pgfqpoint{2.638849in}{1.538737in}}%
\pgfpathlineto{\pgfqpoint{2.641368in}{1.540633in}}%
\pgfpathlineto{\pgfqpoint{2.643887in}{1.547588in}}%
\pgfpathlineto{\pgfqpoint{2.646406in}{1.534124in}}%
\pgfpathlineto{\pgfqpoint{2.648924in}{1.548012in}}%
\pgfpathlineto{\pgfqpoint{2.651443in}{1.552204in}}%
\pgfpathlineto{\pgfqpoint{2.653962in}{1.571721in}}%
\pgfpathlineto{\pgfqpoint{2.656481in}{1.568617in}}%
\pgfpathlineto{\pgfqpoint{2.658999in}{1.577422in}}%
\pgfpathlineto{\pgfqpoint{2.661518in}{1.596329in}}%
\pgfpathlineto{\pgfqpoint{2.664037in}{1.607215in}}%
\pgfpathlineto{\pgfqpoint{2.666556in}{1.590664in}}%
\pgfpathlineto{\pgfqpoint{2.669074in}{1.637386in}}%
\pgfpathlineto{\pgfqpoint{2.671593in}{1.650339in}}%
\pgfpathlineto{\pgfqpoint{2.674112in}{1.661523in}}%
\pgfpathlineto{\pgfqpoint{2.676631in}{1.686767in}}%
\pgfpathlineto{\pgfqpoint{2.679149in}{1.678002in}}%
\pgfpathlineto{\pgfqpoint{2.681668in}{1.696350in}}%
\pgfpathlineto{\pgfqpoint{2.684187in}{1.712786in}}%
\pgfpathlineto{\pgfqpoint{2.686706in}{1.721455in}}%
\pgfpathlineto{\pgfqpoint{2.689224in}{1.739183in}}%
\pgfpathlineto{\pgfqpoint{2.691743in}{1.744322in}}%
\pgfpathlineto{\pgfqpoint{2.694262in}{1.760471in}}%
\pgfpathlineto{\pgfqpoint{2.696781in}{1.779850in}}%
\pgfpathlineto{\pgfqpoint{2.699299in}{1.789456in}}%
\pgfpathlineto{\pgfqpoint{2.701818in}{1.778295in}}%
\pgfpathlineto{\pgfqpoint{2.704337in}{1.757689in}}%
\pgfpathlineto{\pgfqpoint{2.706856in}{1.750908in}}%
\pgfpathlineto{\pgfqpoint{2.709374in}{1.734579in}}%
\pgfpathlineto{\pgfqpoint{2.711893in}{1.746532in}}%
\pgfpathlineto{\pgfqpoint{2.714412in}{1.760134in}}%
\pgfpathlineto{\pgfqpoint{2.716931in}{1.769565in}}%
\pgfpathlineto{\pgfqpoint{2.719449in}{1.753415in}}%
\pgfpathlineto{\pgfqpoint{2.721968in}{1.764022in}}%
\pgfpathlineto{\pgfqpoint{2.724487in}{1.776528in}}%
\pgfpathlineto{\pgfqpoint{2.727006in}{1.779315in}}%
\pgfpathlineto{\pgfqpoint{2.729524in}{1.794392in}}%
\pgfpathlineto{\pgfqpoint{2.732043in}{1.792939in}}%
\pgfpathlineto{\pgfqpoint{2.734562in}{1.813854in}}%
\pgfpathlineto{\pgfqpoint{2.737081in}{1.842542in}}%
\pgfpathlineto{\pgfqpoint{2.739599in}{1.833743in}}%
\pgfpathlineto{\pgfqpoint{2.742118in}{1.829631in}}%
\pgfpathlineto{\pgfqpoint{2.744637in}{1.832044in}}%
\pgfpathlineto{\pgfqpoint{2.747156in}{1.841192in}}%
\pgfpathlineto{\pgfqpoint{2.749674in}{1.829519in}}%
\pgfpathlineto{\pgfqpoint{2.752193in}{1.837054in}}%
\pgfpathlineto{\pgfqpoint{2.754712in}{1.833756in}}%
\pgfpathlineto{\pgfqpoint{2.757231in}{1.843304in}}%
\pgfpathlineto{\pgfqpoint{2.759749in}{1.841348in}}%
\pgfpathlineto{\pgfqpoint{2.762268in}{1.864402in}}%
\pgfpathlineto{\pgfqpoint{2.764787in}{1.859338in}}%
\pgfpathlineto{\pgfqpoint{2.767306in}{1.861940in}}%
\pgfpathlineto{\pgfqpoint{2.769824in}{1.859105in}}%
\pgfpathlineto{\pgfqpoint{2.772343in}{1.854812in}}%
\pgfpathlineto{\pgfqpoint{2.774862in}{1.858593in}}%
\pgfpathlineto{\pgfqpoint{2.777381in}{1.871311in}}%
\pgfpathlineto{\pgfqpoint{2.779899in}{1.880083in}}%
\pgfpathlineto{\pgfqpoint{2.782418in}{1.875001in}}%
\pgfpathlineto{\pgfqpoint{2.784937in}{1.897086in}}%
\pgfpathlineto{\pgfqpoint{2.787456in}{1.884345in}}%
\pgfpathlineto{\pgfqpoint{2.789974in}{1.892325in}}%
\pgfpathlineto{\pgfqpoint{2.792493in}{1.903109in}}%
\pgfpathlineto{\pgfqpoint{2.795012in}{1.890342in}}%
\pgfpathlineto{\pgfqpoint{2.797531in}{1.887981in}}%
\pgfpathlineto{\pgfqpoint{2.800049in}{1.899028in}}%
\pgfpathlineto{\pgfqpoint{2.802568in}{1.880427in}}%
\pgfpathlineto{\pgfqpoint{2.805087in}{1.909918in}}%
\pgfpathlineto{\pgfqpoint{2.807606in}{1.946286in}}%
\pgfpathlineto{\pgfqpoint{2.810124in}{1.927679in}}%
\pgfpathlineto{\pgfqpoint{2.812643in}{1.926519in}}%
\pgfpathlineto{\pgfqpoint{2.815162in}{1.912132in}}%
\pgfpathlineto{\pgfqpoint{2.817681in}{1.935245in}}%
\pgfpathlineto{\pgfqpoint{2.820199in}{1.964273in}}%
\pgfpathlineto{\pgfqpoint{2.822718in}{1.975999in}}%
\pgfpathlineto{\pgfqpoint{2.825237in}{1.979153in}}%
\pgfpathlineto{\pgfqpoint{2.827756in}{1.965183in}}%
\pgfpathlineto{\pgfqpoint{2.830274in}{1.955712in}}%
\pgfpathlineto{\pgfqpoint{2.832793in}{1.975089in}}%
\pgfpathlineto{\pgfqpoint{2.835312in}{1.973408in}}%
\pgfpathlineto{\pgfqpoint{2.837831in}{1.949978in}}%
\pgfpathlineto{\pgfqpoint{2.840349in}{1.966648in}}%
\pgfpathlineto{\pgfqpoint{2.842868in}{1.968815in}}%
\pgfpathlineto{\pgfqpoint{2.845387in}{1.990515in}}%
\pgfpathlineto{\pgfqpoint{2.847906in}{2.017785in}}%
\pgfpathlineto{\pgfqpoint{2.850424in}{2.040964in}}%
\pgfpathlineto{\pgfqpoint{2.852943in}{2.061748in}}%
\pgfpathlineto{\pgfqpoint{2.855462in}{2.055786in}}%
\pgfpathlineto{\pgfqpoint{2.857981in}{2.046381in}}%
\pgfpathlineto{\pgfqpoint{2.860499in}{2.028780in}}%
\pgfpathlineto{\pgfqpoint{2.863018in}{2.029102in}}%
\pgfpathlineto{\pgfqpoint{2.865537in}{2.030259in}}%
\pgfpathlineto{\pgfqpoint{2.868056in}{2.025383in}}%
\pgfpathlineto{\pgfqpoint{2.870574in}{2.024227in}}%
\pgfpathlineto{\pgfqpoint{2.873093in}{2.012423in}}%
\pgfpathlineto{\pgfqpoint{2.875612in}{2.004102in}}%
\pgfpathlineto{\pgfqpoint{2.878131in}{2.025055in}}%
\pgfpathlineto{\pgfqpoint{2.880649in}{2.024004in}}%
\pgfpathlineto{\pgfqpoint{2.883168in}{2.040286in}}%
\pgfpathlineto{\pgfqpoint{2.885687in}{2.026445in}}%
\pgfpathlineto{\pgfqpoint{2.888206in}{2.037754in}}%
\pgfpathlineto{\pgfqpoint{2.890724in}{2.013515in}}%
\pgfpathlineto{\pgfqpoint{2.893243in}{2.039485in}}%
\pgfpathlineto{\pgfqpoint{2.895762in}{2.032236in}}%
\pgfpathlineto{\pgfqpoint{2.898281in}{2.032516in}}%
\pgfpathlineto{\pgfqpoint{2.900799in}{2.030594in}}%
\pgfpathlineto{\pgfqpoint{2.903318in}{2.014979in}}%
\pgfpathlineto{\pgfqpoint{2.905837in}{2.012811in}}%
\pgfpathlineto{\pgfqpoint{2.908356in}{1.968813in}}%
\pgfpathlineto{\pgfqpoint{2.910874in}{1.969687in}}%
\pgfpathlineto{\pgfqpoint{2.913393in}{1.987461in}}%
\pgfpathlineto{\pgfqpoint{2.915912in}{1.950784in}}%
\pgfpathlineto{\pgfqpoint{2.918431in}{1.948334in}}%
\pgfpathlineto{\pgfqpoint{2.920949in}{1.917150in}}%
\pgfpathlineto{\pgfqpoint{2.923468in}{1.924289in}}%
\pgfpathlineto{\pgfqpoint{2.925987in}{1.914525in}}%
\pgfpathlineto{\pgfqpoint{2.928506in}{1.911107in}}%
\pgfpathlineto{\pgfqpoint{2.931024in}{1.887767in}}%
\pgfpathlineto{\pgfqpoint{2.933543in}{1.844779in}}%
\pgfpathlineto{\pgfqpoint{2.936062in}{1.814096in}}%
\pgfpathlineto{\pgfqpoint{2.938581in}{1.843545in}}%
\pgfpathlineto{\pgfqpoint{2.941099in}{1.845569in}}%
\pgfpathlineto{\pgfqpoint{2.943618in}{1.848873in}}%
\pgfpathlineto{\pgfqpoint{2.946137in}{1.839538in}}%
\pgfpathlineto{\pgfqpoint{2.948656in}{1.841533in}}%
\pgfpathlineto{\pgfqpoint{2.951174in}{1.835386in}}%
\pgfpathlineto{\pgfqpoint{2.953693in}{1.838317in}}%
\pgfpathlineto{\pgfqpoint{2.956212in}{1.840306in}}%
\pgfpathlineto{\pgfqpoint{2.958731in}{1.829711in}}%
\pgfpathlineto{\pgfqpoint{2.961249in}{1.845712in}}%
\pgfpathlineto{\pgfqpoint{2.963768in}{1.862916in}}%
\pgfpathlineto{\pgfqpoint{2.966287in}{1.842704in}}%
\pgfpathlineto{\pgfqpoint{2.968806in}{1.839529in}}%
\pgfpathlineto{\pgfqpoint{2.971324in}{1.837506in}}%
\pgfpathlineto{\pgfqpoint{2.973843in}{1.854776in}}%
\pgfpathlineto{\pgfqpoint{2.976362in}{1.850425in}}%
\pgfpathlineto{\pgfqpoint{2.978881in}{1.858757in}}%
\pgfpathlineto{\pgfqpoint{2.981399in}{1.859687in}}%
\pgfpathlineto{\pgfqpoint{2.983918in}{1.868056in}}%
\pgfpathlineto{\pgfqpoint{2.986437in}{1.843095in}}%
\pgfpathlineto{\pgfqpoint{2.988956in}{1.816481in}}%
\pgfpathlineto{\pgfqpoint{2.991474in}{1.827203in}}%
\pgfpathlineto{\pgfqpoint{2.993993in}{1.865030in}}%
\pgfpathlineto{\pgfqpoint{2.996512in}{1.854937in}}%
\pgfpathlineto{\pgfqpoint{2.999031in}{1.846236in}}%
\pgfpathlineto{\pgfqpoint{3.001549in}{1.862470in}}%
\pgfpathlineto{\pgfqpoint{3.004068in}{1.862313in}}%
\pgfpathlineto{\pgfqpoint{3.006587in}{1.876061in}}%
\pgfpathlineto{\pgfqpoint{3.009106in}{1.882842in}}%
\pgfpathlineto{\pgfqpoint{3.011624in}{1.880585in}}%
\pgfpathlineto{\pgfqpoint{3.014143in}{1.886973in}}%
\pgfpathlineto{\pgfqpoint{3.016662in}{1.926044in}}%
\pgfpathlineto{\pgfqpoint{3.019181in}{1.930527in}}%
\pgfpathlineto{\pgfqpoint{3.021699in}{1.917070in}}%
\pgfpathlineto{\pgfqpoint{3.024218in}{1.950817in}}%
\pgfpathlineto{\pgfqpoint{3.026737in}{1.957354in}}%
\pgfpathlineto{\pgfqpoint{3.029256in}{1.958699in}}%
\pgfpathlineto{\pgfqpoint{3.031774in}{1.967096in}}%
\pgfpathlineto{\pgfqpoint{3.034293in}{1.970942in}}%
\pgfpathlineto{\pgfqpoint{3.036812in}{1.968689in}}%
\pgfpathlineto{\pgfqpoint{3.039331in}{2.006593in}}%
\pgfpathlineto{\pgfqpoint{3.041849in}{1.997888in}}%
\pgfpathlineto{\pgfqpoint{3.044368in}{1.977220in}}%
\pgfpathlineto{\pgfqpoint{3.046887in}{1.983554in}}%
\pgfpathlineto{\pgfqpoint{3.049406in}{1.973996in}}%
\pgfpathlineto{\pgfqpoint{3.051924in}{1.960446in}}%
\pgfpathlineto{\pgfqpoint{3.054443in}{1.959719in}}%
\pgfpathlineto{\pgfqpoint{3.056962in}{1.941613in}}%
\pgfpathlineto{\pgfqpoint{3.059481in}{1.956483in}}%
\pgfpathlineto{\pgfqpoint{3.061999in}{1.949601in}}%
\pgfpathlineto{\pgfqpoint{3.064518in}{1.980667in}}%
\pgfpathlineto{\pgfqpoint{3.067037in}{1.965515in}}%
\pgfpathlineto{\pgfqpoint{3.069556in}{1.980809in}}%
\pgfpathlineto{\pgfqpoint{3.072074in}{1.976466in}}%
\pgfpathlineto{\pgfqpoint{3.074593in}{1.953025in}}%
\pgfpathlineto{\pgfqpoint{3.077112in}{1.940118in}}%
\pgfpathlineto{\pgfqpoint{3.079631in}{1.939067in}}%
\pgfpathlineto{\pgfqpoint{3.082149in}{1.962384in}}%
\pgfpathlineto{\pgfqpoint{3.084668in}{1.970535in}}%
\pgfpathlineto{\pgfqpoint{3.087187in}{1.943837in}}%
\pgfpathlineto{\pgfqpoint{3.089706in}{1.931810in}}%
\pgfpathlineto{\pgfqpoint{3.092224in}{1.937303in}}%
\pgfpathlineto{\pgfqpoint{3.094743in}{1.923054in}}%
\pgfpathlineto{\pgfqpoint{3.097262in}{1.986662in}}%
\pgfpathlineto{\pgfqpoint{3.099781in}{1.974835in}}%
\pgfpathlineto{\pgfqpoint{3.102299in}{1.980225in}}%
\pgfpathlineto{\pgfqpoint{3.104818in}{1.986572in}}%
\pgfpathlineto{\pgfqpoint{3.107337in}{1.991256in}}%
\pgfpathlineto{\pgfqpoint{3.109856in}{2.002053in}}%
\pgfpathlineto{\pgfqpoint{3.112374in}{2.041323in}}%
\pgfpathlineto{\pgfqpoint{3.114893in}{2.095972in}}%
\pgfpathlineto{\pgfqpoint{3.117412in}{2.096755in}}%
\pgfpathlineto{\pgfqpoint{3.119931in}{2.062213in}}%
\pgfpathlineto{\pgfqpoint{3.122449in}{2.030130in}}%
\pgfpathlineto{\pgfqpoint{3.124968in}{2.013763in}}%
\pgfpathlineto{\pgfqpoint{3.127487in}{2.043420in}}%
\pgfpathlineto{\pgfqpoint{3.130006in}{2.070193in}}%
\pgfpathlineto{\pgfqpoint{3.132524in}{2.030416in}}%
\pgfpathlineto{\pgfqpoint{3.135043in}{2.075555in}}%
\pgfpathlineto{\pgfqpoint{3.137562in}{2.072707in}}%
\pgfpathlineto{\pgfqpoint{3.140081in}{2.110174in}}%
\pgfpathlineto{\pgfqpoint{3.142599in}{2.119878in}}%
\pgfpathlineto{\pgfqpoint{3.145118in}{2.143818in}}%
\pgfpathlineto{\pgfqpoint{3.147637in}{2.137199in}}%
\pgfpathlineto{\pgfqpoint{3.150156in}{2.134865in}}%
\pgfpathlineto{\pgfqpoint{3.152674in}{2.166574in}}%
\pgfpathlineto{\pgfqpoint{3.155193in}{2.138247in}}%
\pgfpathlineto{\pgfqpoint{3.157712in}{2.176068in}}%
\pgfpathlineto{\pgfqpoint{3.160231in}{2.201840in}}%
\pgfpathlineto{\pgfqpoint{3.162749in}{2.190678in}}%
\pgfpathlineto{\pgfqpoint{3.165268in}{2.164445in}}%
\pgfpathlineto{\pgfqpoint{3.167787in}{2.155878in}}%
\pgfpathlineto{\pgfqpoint{3.170306in}{2.174672in}}%
\pgfpathlineto{\pgfqpoint{3.172824in}{2.175080in}}%
\pgfpathlineto{\pgfqpoint{3.175343in}{2.189354in}}%
\pgfpathlineto{\pgfqpoint{3.177862in}{2.192419in}}%
\pgfpathlineto{\pgfqpoint{3.180381in}{2.170156in}}%
\pgfpathlineto{\pgfqpoint{3.182899in}{2.160712in}}%
\pgfpathlineto{\pgfqpoint{3.185418in}{2.155016in}}%
\pgfpathlineto{\pgfqpoint{3.187937in}{2.169792in}}%
\pgfpathlineto{\pgfqpoint{3.190456in}{2.161587in}}%
\pgfpathlineto{\pgfqpoint{3.192974in}{2.158922in}}%
\pgfpathlineto{\pgfqpoint{3.195493in}{2.162060in}}%
\pgfpathlineto{\pgfqpoint{3.198012in}{2.153969in}}%
\pgfpathlineto{\pgfqpoint{3.200531in}{2.138723in}}%
\pgfpathlineto{\pgfqpoint{3.203049in}{2.129752in}}%
\pgfpathlineto{\pgfqpoint{3.205568in}{2.153419in}}%
\pgfpathlineto{\pgfqpoint{3.208087in}{2.155290in}}%
\pgfpathlineto{\pgfqpoint{3.210606in}{2.163130in}}%
\pgfpathlineto{\pgfqpoint{3.213124in}{2.144477in}}%
\pgfpathlineto{\pgfqpoint{3.215643in}{2.162831in}}%
\pgfpathlineto{\pgfqpoint{3.218162in}{2.185502in}}%
\pgfpathlineto{\pgfqpoint{3.220681in}{2.180949in}}%
\pgfpathlineto{\pgfqpoint{3.223199in}{2.187000in}}%
\pgfpathlineto{\pgfqpoint{3.225718in}{2.161703in}}%
\pgfpathlineto{\pgfqpoint{3.228237in}{2.144640in}}%
\pgfpathlineto{\pgfqpoint{3.230756in}{2.137729in}}%
\pgfpathlineto{\pgfqpoint{3.233274in}{2.163017in}}%
\pgfpathlineto{\pgfqpoint{3.235793in}{2.156941in}}%
\pgfpathlineto{\pgfqpoint{3.238312in}{2.131583in}}%
\pgfpathlineto{\pgfqpoint{3.240831in}{2.134546in}}%
\pgfpathlineto{\pgfqpoint{3.243349in}{2.113171in}}%
\pgfpathlineto{\pgfqpoint{3.245868in}{2.101051in}}%
\pgfpathlineto{\pgfqpoint{3.248387in}{2.115419in}}%
\pgfpathlineto{\pgfqpoint{3.250906in}{2.104326in}}%
\pgfpathlineto{\pgfqpoint{3.253424in}{2.117245in}}%
\pgfpathlineto{\pgfqpoint{3.255943in}{2.090906in}}%
\pgfpathlineto{\pgfqpoint{3.258462in}{2.115080in}}%
\pgfpathlineto{\pgfqpoint{3.260981in}{2.136132in}}%
\pgfpathlineto{\pgfqpoint{3.263499in}{2.123435in}}%
\pgfpathlineto{\pgfqpoint{3.266018in}{2.132210in}}%
\pgfpathlineto{\pgfqpoint{3.268537in}{2.120818in}}%
\pgfpathlineto{\pgfqpoint{3.271056in}{2.126768in}}%
\pgfpathlineto{\pgfqpoint{3.273574in}{2.124892in}}%
\pgfpathlineto{\pgfqpoint{3.276093in}{2.105471in}}%
\pgfpathlineto{\pgfqpoint{3.278612in}{2.098374in}}%
\pgfpathlineto{\pgfqpoint{3.281131in}{2.088732in}}%
\pgfpathlineto{\pgfqpoint{3.283649in}{2.111897in}}%
\pgfpathlineto{\pgfqpoint{3.286168in}{2.125028in}}%
\pgfpathlineto{\pgfqpoint{3.288687in}{2.158705in}}%
\pgfpathlineto{\pgfqpoint{3.291206in}{2.157932in}}%
\pgfpathlineto{\pgfqpoint{3.293724in}{2.163947in}}%
\pgfpathlineto{\pgfqpoint{3.296243in}{2.190128in}}%
\pgfpathlineto{\pgfqpoint{3.298762in}{2.180752in}}%
\pgfpathlineto{\pgfqpoint{3.301281in}{2.147408in}}%
\pgfpathlineto{\pgfqpoint{3.303799in}{2.145130in}}%
\pgfpathlineto{\pgfqpoint{3.306318in}{2.123329in}}%
\pgfpathlineto{\pgfqpoint{3.308837in}{2.127082in}}%
\pgfpathlineto{\pgfqpoint{3.311356in}{2.098376in}}%
\pgfpathlineto{\pgfqpoint{3.313874in}{2.101049in}}%
\pgfpathlineto{\pgfqpoint{3.316393in}{2.132614in}}%
\pgfpathlineto{\pgfqpoint{3.318912in}{2.147178in}}%
\pgfpathlineto{\pgfqpoint{3.321431in}{2.158372in}}%
\pgfpathlineto{\pgfqpoint{3.323949in}{2.154495in}}%
\pgfpathlineto{\pgfqpoint{3.326468in}{2.163719in}}%
\pgfpathlineto{\pgfqpoint{3.328987in}{2.165986in}}%
\pgfpathlineto{\pgfqpoint{3.331506in}{2.164421in}}%
\pgfpathlineto{\pgfqpoint{3.334024in}{2.173239in}}%
\pgfpathlineto{\pgfqpoint{3.336543in}{2.174991in}}%
\pgfpathlineto{\pgfqpoint{3.339062in}{2.184869in}}%
\pgfpathlineto{\pgfqpoint{3.341581in}{2.200296in}}%
\pgfpathlineto{\pgfqpoint{3.344099in}{2.188146in}}%
\pgfpathlineto{\pgfqpoint{3.346618in}{2.181253in}}%
\pgfpathlineto{\pgfqpoint{3.349137in}{2.179775in}}%
\pgfpathlineto{\pgfqpoint{3.351656in}{2.190792in}}%
\pgfpathlineto{\pgfqpoint{3.354174in}{2.180513in}}%
\pgfpathlineto{\pgfqpoint{3.356693in}{2.178584in}}%
\pgfpathlineto{\pgfqpoint{3.359212in}{2.201409in}}%
\pgfpathlineto{\pgfqpoint{3.361731in}{2.202443in}}%
\pgfpathlineto{\pgfqpoint{3.364249in}{2.212111in}}%
\pgfpathlineto{\pgfqpoint{3.366768in}{2.229469in}}%
\pgfpathlineto{\pgfqpoint{3.369287in}{2.251048in}}%
\pgfpathlineto{\pgfqpoint{3.371806in}{2.236863in}}%
\pgfpathlineto{\pgfqpoint{3.374324in}{2.238975in}}%
\pgfpathlineto{\pgfqpoint{3.376843in}{2.244938in}}%
\pgfpathlineto{\pgfqpoint{3.379362in}{2.253046in}}%
\pgfpathlineto{\pgfqpoint{3.381881in}{2.252839in}}%
\pgfpathlineto{\pgfqpoint{3.384399in}{2.218622in}}%
\pgfpathlineto{\pgfqpoint{3.386918in}{2.212569in}}%
\pgfpathlineto{\pgfqpoint{3.389437in}{2.218919in}}%
\pgfpathlineto{\pgfqpoint{3.391956in}{2.204341in}}%
\pgfpathlineto{\pgfqpoint{3.394474in}{2.223501in}}%
\pgfpathlineto{\pgfqpoint{3.396993in}{2.197555in}}%
\pgfpathlineto{\pgfqpoint{3.399512in}{2.216859in}}%
\pgfpathlineto{\pgfqpoint{3.402031in}{2.220420in}}%
\pgfpathlineto{\pgfqpoint{3.404549in}{2.231714in}}%
\pgfpathlineto{\pgfqpoint{3.407068in}{2.238872in}}%
\pgfpathlineto{\pgfqpoint{3.409587in}{2.260934in}}%
\pgfpathlineto{\pgfqpoint{3.412106in}{2.242360in}}%
\pgfpathlineto{\pgfqpoint{3.414624in}{2.230672in}}%
\pgfpathlineto{\pgfqpoint{3.417143in}{2.224328in}}%
\pgfpathlineto{\pgfqpoint{3.419662in}{2.267914in}}%
\pgfpathlineto{\pgfqpoint{3.422181in}{2.251054in}}%
\pgfpathlineto{\pgfqpoint{3.424699in}{2.246797in}}%
\pgfpathlineto{\pgfqpoint{3.427218in}{2.244923in}}%
\pgfpathlineto{\pgfqpoint{3.429737in}{2.226618in}}%
\pgfpathlineto{\pgfqpoint{3.432256in}{2.245417in}}%
\pgfpathlineto{\pgfqpoint{3.434774in}{2.229480in}}%
\pgfpathlineto{\pgfqpoint{3.437293in}{2.216018in}}%
\pgfpathlineto{\pgfqpoint{3.439812in}{2.195016in}}%
\pgfpathlineto{\pgfqpoint{3.442331in}{2.221990in}}%
\pgfpathlineto{\pgfqpoint{3.444849in}{2.234181in}}%
\pgfpathlineto{\pgfqpoint{3.447368in}{2.220027in}}%
\pgfpathlineto{\pgfqpoint{3.449887in}{2.191676in}}%
\pgfpathlineto{\pgfqpoint{3.452406in}{2.175912in}}%
\pgfpathlineto{\pgfqpoint{3.454924in}{2.150808in}}%
\pgfpathlineto{\pgfqpoint{3.457443in}{2.160657in}}%
\pgfpathlineto{\pgfqpoint{3.459962in}{2.159335in}}%
\pgfpathlineto{\pgfqpoint{3.464999in}{2.158966in}}%
\pgfpathlineto{\pgfqpoint{3.467518in}{2.144405in}}%
\pgfpathlineto{\pgfqpoint{3.470037in}{2.125118in}}%
\pgfpathlineto{\pgfqpoint{3.472556in}{2.138102in}}%
\pgfpathlineto{\pgfqpoint{3.475074in}{2.168304in}}%
\pgfpathlineto{\pgfqpoint{3.477593in}{2.184170in}}%
\pgfpathlineto{\pgfqpoint{3.480112in}{2.160451in}}%
\pgfpathlineto{\pgfqpoint{3.482631in}{2.156299in}}%
\pgfpathlineto{\pgfqpoint{3.485149in}{2.140394in}}%
\pgfpathlineto{\pgfqpoint{3.487668in}{2.138015in}}%
\pgfpathlineto{\pgfqpoint{3.490187in}{2.160416in}}%
\pgfpathlineto{\pgfqpoint{3.492706in}{2.149892in}}%
\pgfpathlineto{\pgfqpoint{3.495224in}{2.105966in}}%
\pgfpathlineto{\pgfqpoint{3.497743in}{2.079585in}}%
\pgfpathlineto{\pgfqpoint{3.500262in}{2.081069in}}%
\pgfpathlineto{\pgfqpoint{3.502781in}{2.079563in}}%
\pgfpathlineto{\pgfqpoint{3.505299in}{2.063800in}}%
\pgfpathlineto{\pgfqpoint{3.507818in}{2.065922in}}%
\pgfpathlineto{\pgfqpoint{3.510337in}{2.083382in}}%
\pgfpathlineto{\pgfqpoint{3.512856in}{2.077464in}}%
\pgfpathlineto{\pgfqpoint{3.515374in}{2.043274in}}%
\pgfpathlineto{\pgfqpoint{3.517893in}{2.050296in}}%
\pgfpathlineto{\pgfqpoint{3.520412in}{2.036198in}}%
\pgfpathlineto{\pgfqpoint{3.522931in}{2.055568in}}%
\pgfpathlineto{\pgfqpoint{3.525449in}{2.076019in}}%
\pgfpathlineto{\pgfqpoint{3.527968in}{2.092074in}}%
\pgfpathlineto{\pgfqpoint{3.530487in}{2.109247in}}%
\pgfpathlineto{\pgfqpoint{3.533006in}{2.106399in}}%
\pgfpathlineto{\pgfqpoint{3.535524in}{2.127985in}}%
\pgfpathlineto{\pgfqpoint{3.538043in}{2.125787in}}%
\pgfpathlineto{\pgfqpoint{3.540562in}{2.133410in}}%
\pgfpathlineto{\pgfqpoint{3.543081in}{2.161788in}}%
\pgfpathlineto{\pgfqpoint{3.545599in}{2.159150in}}%
\pgfpathlineto{\pgfqpoint{3.548118in}{2.133322in}}%
\pgfpathlineto{\pgfqpoint{3.550637in}{2.125445in}}%
\pgfpathlineto{\pgfqpoint{3.553156in}{2.125564in}}%
\pgfpathlineto{\pgfqpoint{3.555674in}{2.153874in}}%
\pgfpathlineto{\pgfqpoint{3.558193in}{2.186620in}}%
\pgfpathlineto{\pgfqpoint{3.560712in}{2.201008in}}%
\pgfpathlineto{\pgfqpoint{3.563231in}{2.173161in}}%
\pgfpathlineto{\pgfqpoint{3.565749in}{2.167631in}}%
\pgfpathlineto{\pgfqpoint{3.568268in}{2.156020in}}%
\pgfpathlineto{\pgfqpoint{3.570787in}{2.163339in}}%
\pgfpathlineto{\pgfqpoint{3.573306in}{2.142372in}}%
\pgfpathlineto{\pgfqpoint{3.575824in}{2.149508in}}%
\pgfpathlineto{\pgfqpoint{3.578343in}{2.168079in}}%
\pgfpathlineto{\pgfqpoint{3.580862in}{2.189680in}}%
\pgfpathlineto{\pgfqpoint{3.583381in}{2.189486in}}%
\pgfpathlineto{\pgfqpoint{3.585899in}{2.198610in}}%
\pgfpathlineto{\pgfqpoint{3.588418in}{2.216135in}}%
\pgfpathlineto{\pgfqpoint{3.590937in}{2.241099in}}%
\pgfpathlineto{\pgfqpoint{3.593456in}{2.237077in}}%
\pgfpathlineto{\pgfqpoint{3.595974in}{2.230268in}}%
\pgfpathlineto{\pgfqpoint{3.598493in}{2.234739in}}%
\pgfpathlineto{\pgfqpoint{3.601012in}{2.247488in}}%
\pgfpathlineto{\pgfqpoint{3.603531in}{2.229319in}}%
\pgfpathlineto{\pgfqpoint{3.606049in}{2.200755in}}%
\pgfpathlineto{\pgfqpoint{3.608568in}{2.217147in}}%
\pgfpathlineto{\pgfqpoint{3.611087in}{2.220745in}}%
\pgfpathlineto{\pgfqpoint{3.613606in}{2.223962in}}%
\pgfpathlineto{\pgfqpoint{3.616124in}{2.213827in}}%
\pgfpathlineto{\pgfqpoint{3.618643in}{2.241363in}}%
\pgfpathlineto{\pgfqpoint{3.621162in}{2.227754in}}%
\pgfpathlineto{\pgfqpoint{3.623681in}{2.237901in}}%
\pgfpathlineto{\pgfqpoint{3.626199in}{2.244216in}}%
\pgfpathlineto{\pgfqpoint{3.628718in}{2.252304in}}%
\pgfpathlineto{\pgfqpoint{3.631237in}{2.263907in}}%
\pgfpathlineto{\pgfqpoint{3.633756in}{2.280749in}}%
\pgfpathlineto{\pgfqpoint{3.636274in}{2.269967in}}%
\pgfpathlineto{\pgfqpoint{3.638793in}{2.274811in}}%
\pgfpathlineto{\pgfqpoint{3.641312in}{2.267572in}}%
\pgfpathlineto{\pgfqpoint{3.643831in}{2.238573in}}%
\pgfpathlineto{\pgfqpoint{3.646349in}{2.221095in}}%
\pgfpathlineto{\pgfqpoint{3.648868in}{2.192792in}}%
\pgfpathlineto{\pgfqpoint{3.651387in}{2.196443in}}%
\pgfpathlineto{\pgfqpoint{3.653906in}{2.185840in}}%
\pgfpathlineto{\pgfqpoint{3.656424in}{2.182932in}}%
\pgfpathlineto{\pgfqpoint{3.658943in}{2.171598in}}%
\pgfpathlineto{\pgfqpoint{3.661462in}{2.186278in}}%
\pgfpathlineto{\pgfqpoint{3.663981in}{2.159898in}}%
\pgfpathlineto{\pgfqpoint{3.666499in}{2.157137in}}%
\pgfpathlineto{\pgfqpoint{3.669018in}{2.162465in}}%
\pgfpathlineto{\pgfqpoint{3.671537in}{2.152093in}}%
\pgfpathlineto{\pgfqpoint{3.674056in}{2.163724in}}%
\pgfpathlineto{\pgfqpoint{3.676574in}{2.138926in}}%
\pgfpathlineto{\pgfqpoint{3.679093in}{2.174446in}}%
\pgfpathlineto{\pgfqpoint{3.681612in}{2.165360in}}%
\pgfpathlineto{\pgfqpoint{3.684131in}{2.137877in}}%
\pgfpathlineto{\pgfqpoint{3.686649in}{2.145234in}}%
\pgfpathlineto{\pgfqpoint{3.689168in}{2.146352in}}%
\pgfpathlineto{\pgfqpoint{3.691687in}{2.142776in}}%
\pgfpathlineto{\pgfqpoint{3.694206in}{2.150002in}}%
\pgfpathlineto{\pgfqpoint{3.696724in}{2.192670in}}%
\pgfpathlineto{\pgfqpoint{3.699243in}{2.174156in}}%
\pgfpathlineto{\pgfqpoint{3.701762in}{2.156837in}}%
\pgfpathlineto{\pgfqpoint{3.704281in}{2.137706in}}%
\pgfpathlineto{\pgfqpoint{3.706799in}{2.132915in}}%
\pgfpathlineto{\pgfqpoint{3.709318in}{2.125770in}}%
\pgfpathlineto{\pgfqpoint{3.711837in}{2.148228in}}%
\pgfpathlineto{\pgfqpoint{3.714356in}{2.177228in}}%
\pgfpathlineto{\pgfqpoint{3.716874in}{2.159915in}}%
\pgfpathlineto{\pgfqpoint{3.719393in}{2.157666in}}%
\pgfpathlineto{\pgfqpoint{3.721912in}{2.191927in}}%
\pgfpathlineto{\pgfqpoint{3.724431in}{2.207697in}}%
\pgfpathlineto{\pgfqpoint{3.726949in}{2.200633in}}%
\pgfpathlineto{\pgfqpoint{3.729468in}{2.210644in}}%
\pgfpathlineto{\pgfqpoint{3.731987in}{2.194692in}}%
\pgfpathlineto{\pgfqpoint{3.734506in}{2.215124in}}%
\pgfpathlineto{\pgfqpoint{3.737024in}{2.200729in}}%
\pgfpathlineto{\pgfqpoint{3.739543in}{2.196093in}}%
\pgfpathlineto{\pgfqpoint{3.742062in}{2.193139in}}%
\pgfpathlineto{\pgfqpoint{3.744581in}{2.188962in}}%
\pgfpathlineto{\pgfqpoint{3.747099in}{2.194679in}}%
\pgfpathlineto{\pgfqpoint{3.749618in}{2.196678in}}%
\pgfpathlineto{\pgfqpoint{3.752137in}{2.226273in}}%
\pgfpathlineto{\pgfqpoint{3.754656in}{2.210682in}}%
\pgfpathlineto{\pgfqpoint{3.757174in}{2.189475in}}%
\pgfpathlineto{\pgfqpoint{3.759693in}{2.189605in}}%
\pgfpathlineto{\pgfqpoint{3.762212in}{2.194938in}}%
\pgfpathlineto{\pgfqpoint{3.764731in}{2.174403in}}%
\pgfpathlineto{\pgfqpoint{3.767249in}{2.163903in}}%
\pgfpathlineto{\pgfqpoint{3.769768in}{2.184053in}}%
\pgfpathlineto{\pgfqpoint{3.772287in}{2.191594in}}%
\pgfpathlineto{\pgfqpoint{3.774806in}{2.197480in}}%
\pgfpathlineto{\pgfqpoint{3.777324in}{2.205834in}}%
\pgfpathlineto{\pgfqpoint{3.779843in}{2.240026in}}%
\pgfpathlineto{\pgfqpoint{3.782362in}{2.256620in}}%
\pgfpathlineto{\pgfqpoint{3.784881in}{2.263358in}}%
\pgfpathlineto{\pgfqpoint{3.787399in}{2.260548in}}%
\pgfpathlineto{\pgfqpoint{3.789918in}{2.265402in}}%
\pgfpathlineto{\pgfqpoint{3.792437in}{2.252356in}}%
\pgfpathlineto{\pgfqpoint{3.794956in}{2.277920in}}%
\pgfpathlineto{\pgfqpoint{3.797474in}{2.262782in}}%
\pgfpathlineto{\pgfqpoint{3.799993in}{2.262198in}}%
\pgfpathlineto{\pgfqpoint{3.802512in}{2.273680in}}%
\pgfpathlineto{\pgfqpoint{3.805031in}{2.312388in}}%
\pgfpathlineto{\pgfqpoint{3.807549in}{2.323911in}}%
\pgfpathlineto{\pgfqpoint{3.810068in}{2.316573in}}%
\pgfpathlineto{\pgfqpoint{3.812587in}{2.307485in}}%
\pgfpathlineto{\pgfqpoint{3.815106in}{2.303443in}}%
\pgfpathlineto{\pgfqpoint{3.817624in}{2.317936in}}%
\pgfpathlineto{\pgfqpoint{3.820143in}{2.340505in}}%
\pgfpathlineto{\pgfqpoint{3.822662in}{2.317842in}}%
\pgfpathlineto{\pgfqpoint{3.825181in}{2.305750in}}%
\pgfpathlineto{\pgfqpoint{3.827699in}{2.289694in}}%
\pgfpathlineto{\pgfqpoint{3.830218in}{2.309533in}}%
\pgfpathlineto{\pgfqpoint{3.832737in}{2.310269in}}%
\pgfpathlineto{\pgfqpoint{3.835256in}{2.323812in}}%
\pgfpathlineto{\pgfqpoint{3.837774in}{2.306165in}}%
\pgfpathlineto{\pgfqpoint{3.840293in}{2.318153in}}%
\pgfpathlineto{\pgfqpoint{3.842812in}{2.309899in}}%
\pgfpathlineto{\pgfqpoint{3.845331in}{2.316801in}}%
\pgfpathlineto{\pgfqpoint{3.847849in}{2.331282in}}%
\pgfpathlineto{\pgfqpoint{3.850368in}{2.318649in}}%
\pgfpathlineto{\pgfqpoint{3.852887in}{2.310043in}}%
\pgfpathlineto{\pgfqpoint{3.855406in}{2.304580in}}%
\pgfpathlineto{\pgfqpoint{3.857924in}{2.315474in}}%
\pgfpathlineto{\pgfqpoint{3.860443in}{2.297840in}}%
\pgfpathlineto{\pgfqpoint{3.862962in}{2.295759in}}%
\pgfpathlineto{\pgfqpoint{3.865481in}{2.286766in}}%
\pgfpathlineto{\pgfqpoint{3.867999in}{2.324856in}}%
\pgfpathlineto{\pgfqpoint{3.870518in}{2.302436in}}%
\pgfpathlineto{\pgfqpoint{3.873037in}{2.279040in}}%
\pgfpathlineto{\pgfqpoint{3.875556in}{2.294662in}}%
\pgfpathlineto{\pgfqpoint{3.878074in}{2.287327in}}%
\pgfpathlineto{\pgfqpoint{3.880593in}{2.297340in}}%
\pgfpathlineto{\pgfqpoint{3.883112in}{2.316435in}}%
\pgfpathlineto{\pgfqpoint{3.885631in}{2.306125in}}%
\pgfpathlineto{\pgfqpoint{3.888149in}{2.353566in}}%
\pgfpathlineto{\pgfqpoint{3.890668in}{2.346184in}}%
\pgfpathlineto{\pgfqpoint{3.893187in}{2.376431in}}%
\pgfpathlineto{\pgfqpoint{3.895706in}{2.334398in}}%
\pgfpathlineto{\pgfqpoint{3.898224in}{2.334740in}}%
\pgfpathlineto{\pgfqpoint{3.900743in}{2.328202in}}%
\pgfpathlineto{\pgfqpoint{3.903262in}{2.333913in}}%
\pgfpathlineto{\pgfqpoint{3.905781in}{2.324721in}}%
\pgfpathlineto{\pgfqpoint{3.908299in}{2.316127in}}%
\pgfpathlineto{\pgfqpoint{3.910818in}{2.337835in}}%
\pgfpathlineto{\pgfqpoint{3.913337in}{2.350906in}}%
\pgfpathlineto{\pgfqpoint{3.915856in}{2.334043in}}%
\pgfpathlineto{\pgfqpoint{3.918374in}{2.361623in}}%
\pgfpathlineto{\pgfqpoint{3.920893in}{2.390758in}}%
\pgfpathlineto{\pgfqpoint{3.923412in}{2.405716in}}%
\pgfpathlineto{\pgfqpoint{3.925931in}{2.425099in}}%
\pgfpathlineto{\pgfqpoint{3.928449in}{2.420740in}}%
\pgfpathlineto{\pgfqpoint{3.930968in}{2.427463in}}%
\pgfpathlineto{\pgfqpoint{3.933487in}{2.440156in}}%
\pgfpathlineto{\pgfqpoint{3.936006in}{2.450242in}}%
\pgfpathlineto{\pgfqpoint{3.938524in}{2.442468in}}%
\pgfpathlineto{\pgfqpoint{3.941043in}{2.426099in}}%
\pgfpathlineto{\pgfqpoint{3.943562in}{2.450536in}}%
\pgfpathlineto{\pgfqpoint{3.946081in}{2.449245in}}%
\pgfpathlineto{\pgfqpoint{3.948599in}{2.453681in}}%
\pgfpathlineto{\pgfqpoint{3.951118in}{2.493967in}}%
\pgfpathlineto{\pgfqpoint{3.953637in}{2.501357in}}%
\pgfpathlineto{\pgfqpoint{3.956156in}{2.489436in}}%
\pgfpathlineto{\pgfqpoint{3.958674in}{2.463882in}}%
\pgfpathlineto{\pgfqpoint{3.961193in}{2.455241in}}%
\pgfpathlineto{\pgfqpoint{3.963712in}{2.474498in}}%
\pgfpathlineto{\pgfqpoint{3.966231in}{2.480031in}}%
\pgfpathlineto{\pgfqpoint{3.968749in}{2.482657in}}%
\pgfpathlineto{\pgfqpoint{3.971268in}{2.473221in}}%
\pgfpathlineto{\pgfqpoint{3.973787in}{2.480380in}}%
\pgfpathlineto{\pgfqpoint{3.976306in}{2.443435in}}%
\pgfpathlineto{\pgfqpoint{3.978824in}{2.425102in}}%
\pgfpathlineto{\pgfqpoint{3.981343in}{2.436965in}}%
\pgfpathlineto{\pgfqpoint{3.983862in}{2.438341in}}%
\pgfpathlineto{\pgfqpoint{3.986381in}{2.453683in}}%
\pgfpathlineto{\pgfqpoint{3.988899in}{2.447485in}}%
\pgfpathlineto{\pgfqpoint{3.991418in}{2.432604in}}%
\pgfpathlineto{\pgfqpoint{3.993937in}{2.439947in}}%
\pgfpathlineto{\pgfqpoint{3.996456in}{2.462833in}}%
\pgfpathlineto{\pgfqpoint{3.998974in}{2.473873in}}%
\pgfpathlineto{\pgfqpoint{4.001493in}{2.468318in}}%
\pgfpathlineto{\pgfqpoint{4.004012in}{2.491160in}}%
\pgfpathlineto{\pgfqpoint{4.006531in}{2.493083in}}%
\pgfpathlineto{\pgfqpoint{4.009049in}{2.523987in}}%
\pgfpathlineto{\pgfqpoint{4.011568in}{2.535023in}}%
\pgfpathlineto{\pgfqpoint{4.014087in}{2.525667in}}%
\pgfpathlineto{\pgfqpoint{4.016606in}{2.465683in}}%
\pgfpathlineto{\pgfqpoint{4.019124in}{2.449832in}}%
\pgfpathlineto{\pgfqpoint{4.021643in}{2.454967in}}%
\pgfpathlineto{\pgfqpoint{4.026681in}{2.534866in}}%
\pgfpathlineto{\pgfqpoint{4.029199in}{2.552367in}}%
\pgfpathlineto{\pgfqpoint{4.031718in}{2.537106in}}%
\pgfpathlineto{\pgfqpoint{4.034237in}{2.538859in}}%
\pgfpathlineto{\pgfqpoint{4.036756in}{2.523190in}}%
\pgfpathlineto{\pgfqpoint{4.039274in}{2.536912in}}%
\pgfpathlineto{\pgfqpoint{4.041793in}{2.551698in}}%
\pgfpathlineto{\pgfqpoint{4.044312in}{2.544173in}}%
\pgfpathlineto{\pgfqpoint{4.046831in}{2.535378in}}%
\pgfpathlineto{\pgfqpoint{4.049349in}{2.541203in}}%
\pgfpathlineto{\pgfqpoint{4.051868in}{2.555577in}}%
\pgfpathlineto{\pgfqpoint{4.054387in}{2.556735in}}%
\pgfpathlineto{\pgfqpoint{4.056906in}{2.541046in}}%
\pgfpathlineto{\pgfqpoint{4.059424in}{2.560780in}}%
\pgfpathlineto{\pgfqpoint{4.061943in}{2.564015in}}%
\pgfpathlineto{\pgfqpoint{4.064462in}{2.566840in}}%
\pgfpathlineto{\pgfqpoint{4.066981in}{2.593731in}}%
\pgfpathlineto{\pgfqpoint{4.069499in}{2.617205in}}%
\pgfpathlineto{\pgfqpoint{4.072018in}{2.587740in}}%
\pgfpathlineto{\pgfqpoint{4.074537in}{2.550815in}}%
\pgfpathlineto{\pgfqpoint{4.077056in}{2.535286in}}%
\pgfpathlineto{\pgfqpoint{4.079574in}{2.504834in}}%
\pgfpathlineto{\pgfqpoint{4.082093in}{2.509026in}}%
\pgfpathlineto{\pgfqpoint{4.084612in}{2.525306in}}%
\pgfpathlineto{\pgfqpoint{4.087131in}{2.531189in}}%
\pgfpathlineto{\pgfqpoint{4.089649in}{2.525984in}}%
\pgfpathlineto{\pgfqpoint{4.092168in}{2.519789in}}%
\pgfpathlineto{\pgfqpoint{4.094687in}{2.512308in}}%
\pgfpathlineto{\pgfqpoint{4.097206in}{2.493685in}}%
\pgfpathlineto{\pgfqpoint{4.099724in}{2.495083in}}%
\pgfpathlineto{\pgfqpoint{4.102243in}{2.502846in}}%
\pgfpathlineto{\pgfqpoint{4.104762in}{2.497281in}}%
\pgfpathlineto{\pgfqpoint{4.107281in}{2.465588in}}%
\pgfpathlineto{\pgfqpoint{4.109799in}{2.439015in}}%
\pgfpathlineto{\pgfqpoint{4.112318in}{2.436358in}}%
\pgfpathlineto{\pgfqpoint{4.114837in}{2.456675in}}%
\pgfpathlineto{\pgfqpoint{4.117356in}{2.481737in}}%
\pgfpathlineto{\pgfqpoint{4.119874in}{2.484632in}}%
\pgfpathlineto{\pgfqpoint{4.122393in}{2.495984in}}%
\pgfpathlineto{\pgfqpoint{4.124912in}{2.478817in}}%
\pgfpathlineto{\pgfqpoint{4.127431in}{2.474591in}}%
\pgfpathlineto{\pgfqpoint{4.129949in}{2.495735in}}%
\pgfpathlineto{\pgfqpoint{4.132468in}{2.488959in}}%
\pgfpathlineto{\pgfqpoint{4.134987in}{2.507842in}}%
\pgfpathlineto{\pgfqpoint{4.137506in}{2.510001in}}%
\pgfpathlineto{\pgfqpoint{4.140024in}{2.521953in}}%
\pgfpathlineto{\pgfqpoint{4.142543in}{2.506150in}}%
\pgfpathlineto{\pgfqpoint{4.145062in}{2.498255in}}%
\pgfpathlineto{\pgfqpoint{4.147581in}{2.507579in}}%
\pgfpathlineto{\pgfqpoint{4.150099in}{2.510667in}}%
\pgfpathlineto{\pgfqpoint{4.152618in}{2.484902in}}%
\pgfpathlineto{\pgfqpoint{4.155137in}{2.494580in}}%
\pgfpathlineto{\pgfqpoint{4.157656in}{2.515147in}}%
\pgfpathlineto{\pgfqpoint{4.160174in}{2.555906in}}%
\pgfpathlineto{\pgfqpoint{4.162693in}{2.542692in}}%
\pgfpathlineto{\pgfqpoint{4.165212in}{2.544519in}}%
\pgfpathlineto{\pgfqpoint{4.167731in}{2.557664in}}%
\pgfpathlineto{\pgfqpoint{4.170249in}{2.517833in}}%
\pgfpathlineto{\pgfqpoint{4.172768in}{2.544550in}}%
\pgfpathlineto{\pgfqpoint{4.175287in}{2.563042in}}%
\pgfpathlineto{\pgfqpoint{4.177806in}{2.557937in}}%
\pgfpathlineto{\pgfqpoint{4.180324in}{2.542042in}}%
\pgfpathlineto{\pgfqpoint{4.182843in}{2.547063in}}%
\pgfpathlineto{\pgfqpoint{4.185362in}{2.524170in}}%
\pgfpathlineto{\pgfqpoint{4.187881in}{2.530707in}}%
\pgfpathlineto{\pgfqpoint{4.190399in}{2.523615in}}%
\pgfpathlineto{\pgfqpoint{4.192918in}{2.486336in}}%
\pgfpathlineto{\pgfqpoint{4.195437in}{2.495786in}}%
\pgfpathlineto{\pgfqpoint{4.197956in}{2.503769in}}%
\pgfpathlineto{\pgfqpoint{4.200474in}{2.525002in}}%
\pgfpathlineto{\pgfqpoint{4.202993in}{2.519466in}}%
\pgfpathlineto{\pgfqpoint{4.205512in}{2.507295in}}%
\pgfpathlineto{\pgfqpoint{4.208031in}{2.530131in}}%
\pgfpathlineto{\pgfqpoint{4.210549in}{2.527263in}}%
\pgfpathlineto{\pgfqpoint{4.213068in}{2.547621in}}%
\pgfpathlineto{\pgfqpoint{4.215587in}{2.558144in}}%
\pgfpathlineto{\pgfqpoint{4.218106in}{2.561035in}}%
\pgfpathlineto{\pgfqpoint{4.220624in}{2.554894in}}%
\pgfpathlineto{\pgfqpoint{4.223143in}{2.521399in}}%
\pgfpathlineto{\pgfqpoint{4.225662in}{2.537188in}}%
\pgfpathlineto{\pgfqpoint{4.228181in}{2.540699in}}%
\pgfpathlineto{\pgfqpoint{4.230699in}{2.558102in}}%
\pgfpathlineto{\pgfqpoint{4.233218in}{2.556404in}}%
\pgfpathlineto{\pgfqpoint{4.235737in}{2.519844in}}%
\pgfpathlineto{\pgfqpoint{4.238256in}{2.513732in}}%
\pgfpathlineto{\pgfqpoint{4.240774in}{2.521469in}}%
\pgfpathlineto{\pgfqpoint{4.243293in}{2.552382in}}%
\pgfpathlineto{\pgfqpoint{4.245812in}{2.549367in}}%
\pgfpathlineto{\pgfqpoint{4.248331in}{2.511706in}}%
\pgfpathlineto{\pgfqpoint{4.250849in}{2.555726in}}%
\pgfpathlineto{\pgfqpoint{4.253368in}{2.573485in}}%
\pgfpathlineto{\pgfqpoint{4.255887in}{2.587988in}}%
\pgfpathlineto{\pgfqpoint{4.258406in}{2.590207in}}%
\pgfpathlineto{\pgfqpoint{4.260924in}{2.603107in}}%
\pgfpathlineto{\pgfqpoint{4.263443in}{2.601982in}}%
\pgfpathlineto{\pgfqpoint{4.265962in}{2.608557in}}%
\pgfpathlineto{\pgfqpoint{4.268481in}{2.606882in}}%
\pgfpathlineto{\pgfqpoint{4.270999in}{2.596643in}}%
\pgfpathlineto{\pgfqpoint{4.273518in}{2.577361in}}%
\pgfpathlineto{\pgfqpoint{4.276037in}{2.567609in}}%
\pgfpathlineto{\pgfqpoint{4.278556in}{2.573267in}}%
\pgfpathlineto{\pgfqpoint{4.281074in}{2.592793in}}%
\pgfpathlineto{\pgfqpoint{4.283593in}{2.611551in}}%
\pgfpathlineto{\pgfqpoint{4.286112in}{2.585472in}}%
\pgfpathlineto{\pgfqpoint{4.288631in}{2.560711in}}%
\pgfpathlineto{\pgfqpoint{4.291149in}{2.541537in}}%
\pgfpathlineto{\pgfqpoint{4.293668in}{2.587950in}}%
\pgfpathlineto{\pgfqpoint{4.296187in}{2.613518in}}%
\pgfpathlineto{\pgfqpoint{4.298706in}{2.608541in}}%
\pgfpathlineto{\pgfqpoint{4.301224in}{2.622478in}}%
\pgfpathlineto{\pgfqpoint{4.303743in}{2.612812in}}%
\pgfpathlineto{\pgfqpoint{4.306262in}{2.633305in}}%
\pgfpathlineto{\pgfqpoint{4.308781in}{2.655258in}}%
\pgfpathlineto{\pgfqpoint{4.311299in}{2.642841in}}%
\pgfpathlineto{\pgfqpoint{4.313818in}{2.646769in}}%
\pgfpathlineto{\pgfqpoint{4.316337in}{2.682464in}}%
\pgfpathlineto{\pgfqpoint{4.318856in}{2.715558in}}%
\pgfpathlineto{\pgfqpoint{4.321374in}{2.700334in}}%
\pgfpathlineto{\pgfqpoint{4.323893in}{2.695837in}}%
\pgfpathlineto{\pgfqpoint{4.326412in}{2.679169in}}%
\pgfpathlineto{\pgfqpoint{4.328931in}{2.705175in}}%
\pgfpathlineto{\pgfqpoint{4.331449in}{2.726808in}}%
\pgfpathlineto{\pgfqpoint{4.333968in}{2.736865in}}%
\pgfpathlineto{\pgfqpoint{4.336487in}{2.730668in}}%
\pgfpathlineto{\pgfqpoint{4.339006in}{2.741525in}}%
\pgfpathlineto{\pgfqpoint{4.341524in}{2.741010in}}%
\pgfpathlineto{\pgfqpoint{4.344043in}{2.709827in}}%
\pgfpathlineto{\pgfqpoint{4.346562in}{2.697984in}}%
\pgfpathlineto{\pgfqpoint{4.349081in}{2.659591in}}%
\pgfpathlineto{\pgfqpoint{4.351599in}{2.681698in}}%
\pgfpathlineto{\pgfqpoint{4.354118in}{2.691665in}}%
\pgfpathlineto{\pgfqpoint{4.356637in}{2.696241in}}%
\pgfpathlineto{\pgfqpoint{4.359156in}{2.711831in}}%
\pgfpathlineto{\pgfqpoint{4.361674in}{2.730246in}}%
\pgfpathlineto{\pgfqpoint{4.364193in}{2.724644in}}%
\pgfpathlineto{\pgfqpoint{4.366712in}{2.751388in}}%
\pgfpathlineto{\pgfqpoint{4.369231in}{2.737233in}}%
\pgfpathlineto{\pgfqpoint{4.371749in}{2.730103in}}%
\pgfpathlineto{\pgfqpoint{4.374268in}{2.763195in}}%
\pgfpathlineto{\pgfqpoint{4.376787in}{2.779999in}}%
\pgfpathlineto{\pgfqpoint{4.379306in}{2.798682in}}%
\pgfpathlineto{\pgfqpoint{4.381824in}{2.803816in}}%
\pgfpathlineto{\pgfqpoint{4.384343in}{2.836929in}}%
\pgfpathlineto{\pgfqpoint{4.386862in}{2.889620in}}%
\pgfpathlineto{\pgfqpoint{4.389381in}{2.868173in}}%
\pgfpathlineto{\pgfqpoint{4.391899in}{2.839239in}}%
\pgfpathlineto{\pgfqpoint{4.394418in}{2.806256in}}%
\pgfpathlineto{\pgfqpoint{4.396937in}{2.814438in}}%
\pgfpathlineto{\pgfqpoint{4.399456in}{2.828323in}}%
\pgfpathlineto{\pgfqpoint{4.401974in}{2.844973in}}%
\pgfpathlineto{\pgfqpoint{4.404493in}{2.865503in}}%
\pgfpathlineto{\pgfqpoint{4.407012in}{2.860693in}}%
\pgfpathlineto{\pgfqpoint{4.409531in}{2.883273in}}%
\pgfpathlineto{\pgfqpoint{4.412049in}{2.922976in}}%
\pgfpathlineto{\pgfqpoint{4.414568in}{2.970241in}}%
\pgfpathlineto{\pgfqpoint{4.417087in}{3.008304in}}%
\pgfpathlineto{\pgfqpoint{4.419606in}{3.032494in}}%
\pgfpathlineto{\pgfqpoint{4.422124in}{3.026563in}}%
\pgfpathlineto{\pgfqpoint{4.424643in}{3.030111in}}%
\pgfpathlineto{\pgfqpoint{4.427162in}{3.022264in}}%
\pgfpathlineto{\pgfqpoint{4.429681in}{3.025612in}}%
\pgfpathlineto{\pgfqpoint{4.432199in}{3.044510in}}%
\pgfpathlineto{\pgfqpoint{4.434718in}{3.005274in}}%
\pgfpathlineto{\pgfqpoint{4.437237in}{3.022358in}}%
\pgfpathlineto{\pgfqpoint{4.439756in}{3.000428in}}%
\pgfpathlineto{\pgfqpoint{4.442274in}{2.988173in}}%
\pgfpathlineto{\pgfqpoint{4.444793in}{3.005526in}}%
\pgfpathlineto{\pgfqpoint{4.447312in}{3.008113in}}%
\pgfpathlineto{\pgfqpoint{4.449831in}{3.023736in}}%
\pgfpathlineto{\pgfqpoint{4.452349in}{2.999175in}}%
\pgfpathlineto{\pgfqpoint{4.454868in}{2.976222in}}%
\pgfpathlineto{\pgfqpoint{4.457387in}{2.961184in}}%
\pgfpathlineto{\pgfqpoint{4.459906in}{2.927250in}}%
\pgfpathlineto{\pgfqpoint{4.462424in}{2.950973in}}%
\pgfpathlineto{\pgfqpoint{4.464943in}{2.951938in}}%
\pgfpathlineto{\pgfqpoint{4.467462in}{2.962890in}}%
\pgfpathlineto{\pgfqpoint{4.469981in}{2.986175in}}%
\pgfpathlineto{\pgfqpoint{4.472499in}{2.972749in}}%
\pgfpathlineto{\pgfqpoint{4.475018in}{2.986568in}}%
\pgfpathlineto{\pgfqpoint{4.477537in}{2.981202in}}%
\pgfpathlineto{\pgfqpoint{4.480056in}{2.992297in}}%
\pgfpathlineto{\pgfqpoint{4.482574in}{3.007037in}}%
\pgfpathlineto{\pgfqpoint{4.485093in}{2.990784in}}%
\pgfpathlineto{\pgfqpoint{4.487612in}{2.971513in}}%
\pgfpathlineto{\pgfqpoint{4.490131in}{2.994369in}}%
\pgfpathlineto{\pgfqpoint{4.492649in}{3.018088in}}%
\pgfpathlineto{\pgfqpoint{4.495168in}{2.995647in}}%
\pgfpathlineto{\pgfqpoint{4.497687in}{2.955594in}}%
\pgfpathlineto{\pgfqpoint{4.500206in}{2.975787in}}%
\pgfpathlineto{\pgfqpoint{4.502724in}{2.957954in}}%
\pgfpathlineto{\pgfqpoint{4.505243in}{2.970991in}}%
\pgfpathlineto{\pgfqpoint{4.507762in}{2.987169in}}%
\pgfpathlineto{\pgfqpoint{4.510281in}{3.040168in}}%
\pgfpathlineto{\pgfqpoint{4.512799in}{3.059821in}}%
\pgfpathlineto{\pgfqpoint{4.515318in}{3.034888in}}%
\pgfpathlineto{\pgfqpoint{4.517837in}{3.049359in}}%
\pgfpathlineto{\pgfqpoint{4.520356in}{3.054358in}}%
\pgfpathlineto{\pgfqpoint{4.522874in}{3.035015in}}%
\pgfpathlineto{\pgfqpoint{4.525393in}{3.051048in}}%
\pgfpathlineto{\pgfqpoint{4.527912in}{3.032895in}}%
\pgfpathlineto{\pgfqpoint{4.530431in}{3.048873in}}%
\pgfpathlineto{\pgfqpoint{4.532949in}{2.998610in}}%
\pgfpathlineto{\pgfqpoint{4.535468in}{2.977791in}}%
\pgfpathlineto{\pgfqpoint{4.537987in}{2.966740in}}%
\pgfpathlineto{\pgfqpoint{4.540506in}{2.970515in}}%
\pgfpathlineto{\pgfqpoint{4.543024in}{2.980900in}}%
\pgfpathlineto{\pgfqpoint{4.545543in}{2.930252in}}%
\pgfpathlineto{\pgfqpoint{4.548062in}{2.903312in}}%
\pgfpathlineto{\pgfqpoint{4.550581in}{2.864265in}}%
\pgfpathlineto{\pgfqpoint{4.553099in}{2.853504in}}%
\pgfpathlineto{\pgfqpoint{4.555618in}{2.856120in}}%
\pgfpathlineto{\pgfqpoint{4.558137in}{2.854277in}}%
\pgfpathlineto{\pgfqpoint{4.560656in}{2.874709in}}%
\pgfpathlineto{\pgfqpoint{4.563174in}{2.882472in}}%
\pgfpathlineto{\pgfqpoint{4.565693in}{2.907956in}}%
\pgfpathlineto{\pgfqpoint{4.568212in}{2.912851in}}%
\pgfpathlineto{\pgfqpoint{4.570731in}{2.897709in}}%
\pgfpathlineto{\pgfqpoint{4.573249in}{2.909550in}}%
\pgfpathlineto{\pgfqpoint{4.575768in}{2.929940in}}%
\pgfpathlineto{\pgfqpoint{4.578287in}{2.929536in}}%
\pgfpathlineto{\pgfqpoint{4.580806in}{2.945869in}}%
\pgfpathlineto{\pgfqpoint{4.583324in}{2.910220in}}%
\pgfpathlineto{\pgfqpoint{4.585843in}{2.915170in}}%
\pgfpathlineto{\pgfqpoint{4.588362in}{2.966050in}}%
\pgfpathlineto{\pgfqpoint{4.590881in}{2.985417in}}%
\pgfpathlineto{\pgfqpoint{4.593399in}{3.018009in}}%
\pgfpathlineto{\pgfqpoint{4.595918in}{3.048689in}}%
\pgfpathlineto{\pgfqpoint{4.598437in}{3.045754in}}%
\pgfpathlineto{\pgfqpoint{4.600956in}{3.069724in}}%
\pgfpathlineto{\pgfqpoint{4.603474in}{3.078398in}}%
\pgfpathlineto{\pgfqpoint{4.605993in}{3.112229in}}%
\pgfpathlineto{\pgfqpoint{4.608512in}{3.137978in}}%
\pgfpathlineto{\pgfqpoint{4.611031in}{3.123075in}}%
\pgfpathlineto{\pgfqpoint{4.613549in}{3.148953in}}%
\pgfpathlineto{\pgfqpoint{4.616068in}{3.162685in}}%
\pgfpathlineto{\pgfqpoint{4.618587in}{3.172160in}}%
\pgfpathlineto{\pgfqpoint{4.621106in}{3.195045in}}%
\pgfpathlineto{\pgfqpoint{4.623624in}{3.174556in}}%
\pgfpathlineto{\pgfqpoint{4.626143in}{3.160587in}}%
\pgfpathlineto{\pgfqpoint{4.628662in}{3.122467in}}%
\pgfpathlineto{\pgfqpoint{4.631181in}{3.130920in}}%
\pgfpathlineto{\pgfqpoint{4.633699in}{3.191669in}}%
\pgfpathlineto{\pgfqpoint{4.636218in}{3.166915in}}%
\pgfpathlineto{\pgfqpoint{4.638737in}{3.167788in}}%
\pgfpathlineto{\pgfqpoint{4.641256in}{3.170504in}}%
\pgfpathlineto{\pgfqpoint{4.643774in}{3.177912in}}%
\pgfpathlineto{\pgfqpoint{4.646293in}{3.149225in}}%
\pgfpathlineto{\pgfqpoint{4.648812in}{3.137724in}}%
\pgfpathlineto{\pgfqpoint{4.651331in}{3.154154in}}%
\pgfpathlineto{\pgfqpoint{4.653849in}{3.128075in}}%
\pgfpathlineto{\pgfqpoint{4.656368in}{3.166851in}}%
\pgfpathlineto{\pgfqpoint{4.658887in}{3.198843in}}%
\pgfpathlineto{\pgfqpoint{4.661406in}{3.183521in}}%
\pgfpathlineto{\pgfqpoint{4.663924in}{3.198195in}}%
\pgfpathlineto{\pgfqpoint{4.666443in}{3.181811in}}%
\pgfpathlineto{\pgfqpoint{4.668962in}{3.185699in}}%
\pgfpathlineto{\pgfqpoint{4.671481in}{3.138525in}}%
\pgfpathlineto{\pgfqpoint{4.673999in}{3.119262in}}%
\pgfpathlineto{\pgfqpoint{4.676518in}{3.120011in}}%
\pgfpathlineto{\pgfqpoint{4.679037in}{3.145526in}}%
\pgfpathlineto{\pgfqpoint{4.681556in}{3.151925in}}%
\pgfpathlineto{\pgfqpoint{4.684074in}{3.144045in}}%
\pgfpathlineto{\pgfqpoint{4.686593in}{3.154286in}}%
\pgfpathlineto{\pgfqpoint{4.689112in}{3.146357in}}%
\pgfpathlineto{\pgfqpoint{4.691631in}{3.147046in}}%
\pgfpathlineto{\pgfqpoint{4.694149in}{3.175543in}}%
\pgfpathlineto{\pgfqpoint{4.696668in}{3.193752in}}%
\pgfpathlineto{\pgfqpoint{4.699187in}{3.214795in}}%
\pgfpathlineto{\pgfqpoint{4.701706in}{3.191188in}}%
\pgfpathlineto{\pgfqpoint{4.704224in}{3.189116in}}%
\pgfpathlineto{\pgfqpoint{4.706743in}{3.201443in}}%
\pgfpathlineto{\pgfqpoint{4.709262in}{3.160054in}}%
\pgfpathlineto{\pgfqpoint{4.711781in}{3.175361in}}%
\pgfpathlineto{\pgfqpoint{4.714299in}{3.172372in}}%
\pgfpathlineto{\pgfqpoint{4.716818in}{3.197590in}}%
\pgfpathlineto{\pgfqpoint{4.719337in}{3.167011in}}%
\pgfpathlineto{\pgfqpoint{4.721856in}{3.177339in}}%
\pgfpathlineto{\pgfqpoint{4.724374in}{3.235929in}}%
\pgfpathlineto{\pgfqpoint{4.726893in}{3.255327in}}%
\pgfpathlineto{\pgfqpoint{4.729412in}{3.285049in}}%
\pgfpathlineto{\pgfqpoint{4.731931in}{3.292855in}}%
\pgfpathlineto{\pgfqpoint{4.734449in}{3.333974in}}%
\pgfpathlineto{\pgfqpoint{4.736968in}{3.290043in}}%
\pgfpathlineto{\pgfqpoint{4.739487in}{3.292716in}}%
\pgfpathlineto{\pgfqpoint{4.742006in}{3.257557in}}%
\pgfpathlineto{\pgfqpoint{4.747043in}{3.243539in}}%
\pgfpathlineto{\pgfqpoint{4.749562in}{3.256374in}}%
\pgfpathlineto{\pgfqpoint{4.752081in}{3.229673in}}%
\pgfpathlineto{\pgfqpoint{4.754599in}{3.193653in}}%
\pgfpathlineto{\pgfqpoint{4.757118in}{3.163306in}}%
\pgfpathlineto{\pgfqpoint{4.759637in}{3.177186in}}%
\pgfpathlineto{\pgfqpoint{4.762156in}{3.205823in}}%
\pgfpathlineto{\pgfqpoint{4.764674in}{3.193266in}}%
\pgfpathlineto{\pgfqpoint{4.767193in}{3.202831in}}%
\pgfpathlineto{\pgfqpoint{4.769712in}{3.193786in}}%
\pgfpathlineto{\pgfqpoint{4.772231in}{3.213480in}}%
\pgfpathlineto{\pgfqpoint{4.774749in}{3.234592in}}%
\pgfpathlineto{\pgfqpoint{4.777268in}{3.241493in}}%
\pgfpathlineto{\pgfqpoint{4.779787in}{3.251955in}}%
\pgfpathlineto{\pgfqpoint{4.782306in}{3.252918in}}%
\pgfpathlineto{\pgfqpoint{4.784824in}{3.292480in}}%
\pgfpathlineto{\pgfqpoint{4.787343in}{3.267459in}}%
\pgfpathlineto{\pgfqpoint{4.789862in}{3.236784in}}%
\pgfpathlineto{\pgfqpoint{4.792381in}{3.208936in}}%
\pgfpathlineto{\pgfqpoint{4.794899in}{3.248320in}}%
\pgfpathlineto{\pgfqpoint{4.797418in}{3.235576in}}%
\pgfpathlineto{\pgfqpoint{4.799937in}{3.255137in}}%
\pgfpathlineto{\pgfqpoint{4.802456in}{3.244483in}}%
\pgfpathlineto{\pgfqpoint{4.804974in}{3.178166in}}%
\pgfpathlineto{\pgfqpoint{4.807493in}{3.190017in}}%
\pgfpathlineto{\pgfqpoint{4.810012in}{3.178757in}}%
\pgfpathlineto{\pgfqpoint{4.812531in}{3.177153in}}%
\pgfpathlineto{\pgfqpoint{4.815049in}{3.196185in}}%
\pgfpathlineto{\pgfqpoint{4.817568in}{3.194877in}}%
\pgfpathlineto{\pgfqpoint{4.822606in}{3.178035in}}%
\pgfpathlineto{\pgfqpoint{4.825124in}{3.183176in}}%
\pgfpathlineto{\pgfqpoint{4.827643in}{3.169708in}}%
\pgfpathlineto{\pgfqpoint{4.830162in}{3.171006in}}%
\pgfpathlineto{\pgfqpoint{4.832681in}{3.152692in}}%
\pgfpathlineto{\pgfqpoint{4.835199in}{3.130900in}}%
\pgfpathlineto{\pgfqpoint{4.837718in}{3.155583in}}%
\pgfpathlineto{\pgfqpoint{4.840237in}{3.181821in}}%
\pgfpathlineto{\pgfqpoint{4.842756in}{3.202073in}}%
\pgfpathlineto{\pgfqpoint{4.845274in}{3.205320in}}%
\pgfpathlineto{\pgfqpoint{4.847793in}{3.233824in}}%
\pgfpathlineto{\pgfqpoint{4.850312in}{3.209164in}}%
\pgfpathlineto{\pgfqpoint{4.852831in}{3.241755in}}%
\pgfpathlineto{\pgfqpoint{4.855349in}{3.226608in}}%
\pgfpathlineto{\pgfqpoint{4.857868in}{3.229489in}}%
\pgfpathlineto{\pgfqpoint{4.860387in}{3.214371in}}%
\pgfpathlineto{\pgfqpoint{4.862906in}{3.209804in}}%
\pgfpathlineto{\pgfqpoint{4.865424in}{3.223255in}}%
\pgfpathlineto{\pgfqpoint{4.867943in}{3.204490in}}%
\pgfpathlineto{\pgfqpoint{4.870462in}{3.224801in}}%
\pgfpathlineto{\pgfqpoint{4.872981in}{3.254297in}}%
\pgfpathlineto{\pgfqpoint{4.875499in}{3.261181in}}%
\pgfpathlineto{\pgfqpoint{4.878018in}{3.247130in}}%
\pgfpathlineto{\pgfqpoint{4.880537in}{3.204174in}}%
\pgfpathlineto{\pgfqpoint{4.883056in}{3.209532in}}%
\pgfpathlineto{\pgfqpoint{4.885574in}{3.215901in}}%
\pgfpathlineto{\pgfqpoint{4.888093in}{3.230302in}}%
\pgfpathlineto{\pgfqpoint{4.890612in}{3.279177in}}%
\pgfpathlineto{\pgfqpoint{4.893131in}{3.263458in}}%
\pgfpathlineto{\pgfqpoint{4.895649in}{3.294427in}}%
\pgfpathlineto{\pgfqpoint{4.898168in}{3.306380in}}%
\pgfpathlineto{\pgfqpoint{4.900687in}{3.316469in}}%
\pgfpathlineto{\pgfqpoint{4.903206in}{3.312364in}}%
\pgfpathlineto{\pgfqpoint{4.905724in}{3.306324in}}%
\pgfpathlineto{\pgfqpoint{4.908243in}{3.287612in}}%
\pgfpathlineto{\pgfqpoint{4.910762in}{3.296731in}}%
\pgfpathlineto{\pgfqpoint{4.913281in}{3.291241in}}%
\pgfpathlineto{\pgfqpoint{4.915799in}{3.292390in}}%
\pgfpathlineto{\pgfqpoint{4.918318in}{3.298611in}}%
\pgfpathlineto{\pgfqpoint{4.923356in}{3.287160in}}%
\pgfpathlineto{\pgfqpoint{4.925874in}{3.251420in}}%
\pgfpathlineto{\pgfqpoint{4.928393in}{3.235466in}}%
\pgfpathlineto{\pgfqpoint{4.930912in}{3.239781in}}%
\pgfpathlineto{\pgfqpoint{4.933431in}{3.260229in}}%
\pgfpathlineto{\pgfqpoint{4.935949in}{3.258825in}}%
\pgfpathlineto{\pgfqpoint{4.938468in}{3.289591in}}%
\pgfpathlineto{\pgfqpoint{4.940987in}{3.280302in}}%
\pgfpathlineto{\pgfqpoint{4.943506in}{3.270320in}}%
\pgfpathlineto{\pgfqpoint{4.946024in}{3.283745in}}%
\pgfpathlineto{\pgfqpoint{4.948543in}{3.266634in}}%
\pgfpathlineto{\pgfqpoint{4.951062in}{3.269893in}}%
\pgfpathlineto{\pgfqpoint{4.953581in}{3.291209in}}%
\pgfpathlineto{\pgfqpoint{4.956099in}{3.295590in}}%
\pgfpathlineto{\pgfqpoint{4.958618in}{3.282851in}}%
\pgfpathlineto{\pgfqpoint{4.961137in}{3.307341in}}%
\pgfpathlineto{\pgfqpoint{4.963656in}{3.290888in}}%
\pgfpathlineto{\pgfqpoint{4.966174in}{3.256450in}}%
\pgfpathlineto{\pgfqpoint{4.968693in}{3.228295in}}%
\pgfpathlineto{\pgfqpoint{4.971212in}{3.204684in}}%
\pgfpathlineto{\pgfqpoint{4.973731in}{3.192626in}}%
\pgfpathlineto{\pgfqpoint{4.976249in}{3.217099in}}%
\pgfpathlineto{\pgfqpoint{4.978768in}{3.219065in}}%
\pgfpathlineto{\pgfqpoint{4.981287in}{3.232650in}}%
\pgfpathlineto{\pgfqpoint{4.986324in}{3.187376in}}%
\pgfpathlineto{\pgfqpoint{4.988843in}{3.203868in}}%
\pgfpathlineto{\pgfqpoint{4.991362in}{3.210016in}}%
\pgfpathlineto{\pgfqpoint{4.993881in}{3.211003in}}%
\pgfpathlineto{\pgfqpoint{4.996399in}{3.235180in}}%
\pgfpathlineto{\pgfqpoint{4.998918in}{3.270157in}}%
\pgfpathlineto{\pgfqpoint{5.001437in}{3.299349in}}%
\pgfpathlineto{\pgfqpoint{5.003956in}{3.329807in}}%
\pgfpathlineto{\pgfqpoint{5.006474in}{3.353167in}}%
\pgfpathlineto{\pgfqpoint{5.008993in}{3.350289in}}%
\pgfpathlineto{\pgfqpoint{5.011512in}{3.366967in}}%
\pgfpathlineto{\pgfqpoint{5.014031in}{3.404046in}}%
\pgfpathlineto{\pgfqpoint{5.016549in}{3.378344in}}%
\pgfpathlineto{\pgfqpoint{5.019068in}{3.384716in}}%
\pgfpathlineto{\pgfqpoint{5.021587in}{3.392590in}}%
\pgfpathlineto{\pgfqpoint{5.024106in}{3.413477in}}%
\pgfpathlineto{\pgfqpoint{5.026624in}{3.364071in}}%
\pgfpathlineto{\pgfqpoint{5.029143in}{3.353123in}}%
\pgfpathlineto{\pgfqpoint{5.031662in}{3.404926in}}%
\pgfpathlineto{\pgfqpoint{5.034181in}{3.429865in}}%
\pgfpathlineto{\pgfqpoint{5.036699in}{3.443888in}}%
\pgfpathlineto{\pgfqpoint{5.039218in}{3.446008in}}%
\pgfpathlineto{\pgfqpoint{5.041737in}{3.452421in}}%
\pgfpathlineto{\pgfqpoint{5.044256in}{3.425427in}}%
\pgfpathlineto{\pgfqpoint{5.046774in}{3.431261in}}%
\pgfpathlineto{\pgfqpoint{5.049293in}{3.411520in}}%
\pgfpathlineto{\pgfqpoint{5.051812in}{3.449516in}}%
\pgfpathlineto{\pgfqpoint{5.054331in}{3.472453in}}%
\pgfpathlineto{\pgfqpoint{5.056849in}{3.480136in}}%
\pgfpathlineto{\pgfqpoint{5.059368in}{3.464519in}}%
\pgfpathlineto{\pgfqpoint{5.061887in}{3.501123in}}%
\pgfpathlineto{\pgfqpoint{5.064406in}{3.482960in}}%
\pgfpathlineto{\pgfqpoint{5.066924in}{3.486764in}}%
\pgfpathlineto{\pgfqpoint{5.069443in}{3.524001in}}%
\pgfpathlineto{\pgfqpoint{5.071962in}{3.511844in}}%
\pgfpathlineto{\pgfqpoint{5.074481in}{3.477316in}}%
\pgfpathlineto{\pgfqpoint{5.076999in}{3.484664in}}%
\pgfpathlineto{\pgfqpoint{5.079518in}{3.442043in}}%
\pgfpathlineto{\pgfqpoint{5.082037in}{3.461696in}}%
\pgfpathlineto{\pgfqpoint{5.084556in}{3.431472in}}%
\pgfpathlineto{\pgfqpoint{5.087074in}{3.453569in}}%
\pgfpathlineto{\pgfqpoint{5.089593in}{3.449385in}}%
\pgfpathlineto{\pgfqpoint{5.092112in}{3.456156in}}%
\pgfpathlineto{\pgfqpoint{5.094631in}{3.501016in}}%
\pgfpathlineto{\pgfqpoint{5.097149in}{3.524396in}}%
\pgfpathlineto{\pgfqpoint{5.099668in}{3.525737in}}%
\pgfpathlineto{\pgfqpoint{5.102187in}{3.539989in}}%
\pgfpathlineto{\pgfqpoint{5.104706in}{3.536004in}}%
\pgfpathlineto{\pgfqpoint{5.107224in}{3.541363in}}%
\pgfpathlineto{\pgfqpoint{5.109743in}{3.502883in}}%
\pgfpathlineto{\pgfqpoint{5.112262in}{3.506230in}}%
\pgfpathlineto{\pgfqpoint{5.114781in}{3.498653in}}%
\pgfpathlineto{\pgfqpoint{5.117299in}{3.521092in}}%
\pgfpathlineto{\pgfqpoint{5.119818in}{3.490194in}}%
\pgfpathlineto{\pgfqpoint{5.122337in}{3.504016in}}%
\pgfpathlineto{\pgfqpoint{5.124856in}{3.512890in}}%
\pgfpathlineto{\pgfqpoint{5.127374in}{3.495757in}}%
\pgfpathlineto{\pgfqpoint{5.129893in}{3.525297in}}%
\pgfpathlineto{\pgfqpoint{5.132412in}{3.542720in}}%
\pgfpathlineto{\pgfqpoint{5.134931in}{3.549990in}}%
\pgfpathlineto{\pgfqpoint{5.137449in}{3.562439in}}%
\pgfpathlineto{\pgfqpoint{5.139968in}{3.543244in}}%
\pgfpathlineto{\pgfqpoint{5.142487in}{3.530087in}}%
\pgfpathlineto{\pgfqpoint{5.145006in}{3.559587in}}%
\pgfpathlineto{\pgfqpoint{5.147524in}{3.549570in}}%
\pgfpathlineto{\pgfqpoint{5.150043in}{3.553025in}}%
\pgfpathlineto{\pgfqpoint{5.152562in}{3.523258in}}%
\pgfpathlineto{\pgfqpoint{5.155081in}{3.540495in}}%
\pgfpathlineto{\pgfqpoint{5.157599in}{3.540227in}}%
\pgfpathlineto{\pgfqpoint{5.160118in}{3.533284in}}%
\pgfpathlineto{\pgfqpoint{5.162637in}{3.559048in}}%
\pgfpathlineto{\pgfqpoint{5.165156in}{3.567205in}}%
\pgfpathlineto{\pgfqpoint{5.167674in}{3.542083in}}%
\pgfpathlineto{\pgfqpoint{5.170193in}{3.541041in}}%
\pgfpathlineto{\pgfqpoint{5.172712in}{3.536965in}}%
\pgfpathlineto{\pgfqpoint{5.175231in}{3.552915in}}%
\pgfpathlineto{\pgfqpoint{5.177749in}{3.557767in}}%
\pgfpathlineto{\pgfqpoint{5.180268in}{3.555967in}}%
\pgfpathlineto{\pgfqpoint{5.182787in}{3.561231in}}%
\pgfpathlineto{\pgfqpoint{5.185306in}{3.551042in}}%
\pgfpathlineto{\pgfqpoint{5.187824in}{3.528610in}}%
\pgfpathlineto{\pgfqpoint{5.190343in}{3.531193in}}%
\pgfpathlineto{\pgfqpoint{5.192862in}{3.541408in}}%
\pgfpathlineto{\pgfqpoint{5.195381in}{3.553329in}}%
\pgfpathlineto{\pgfqpoint{5.197899in}{3.530145in}}%
\pgfpathlineto{\pgfqpoint{5.200418in}{3.514562in}}%
\pgfpathlineto{\pgfqpoint{5.202937in}{3.545849in}}%
\pgfpathlineto{\pgfqpoint{5.205456in}{3.537122in}}%
\pgfpathlineto{\pgfqpoint{5.207974in}{3.516324in}}%
\pgfpathlineto{\pgfqpoint{5.210493in}{3.497727in}}%
\pgfpathlineto{\pgfqpoint{5.213012in}{3.510169in}}%
\pgfpathlineto{\pgfqpoint{5.215531in}{3.536451in}}%
\pgfpathlineto{\pgfqpoint{5.218049in}{3.554981in}}%
\pgfpathlineto{\pgfqpoint{5.220568in}{3.542731in}}%
\pgfpathlineto{\pgfqpoint{5.223087in}{3.560532in}}%
\pgfpathlineto{\pgfqpoint{5.225606in}{3.583396in}}%
\pgfpathlineto{\pgfqpoint{5.230643in}{3.607009in}}%
\pgfpathlineto{\pgfqpoint{5.233162in}{3.595186in}}%
\pgfpathlineto{\pgfqpoint{5.235681in}{3.560753in}}%
\pgfpathlineto{\pgfqpoint{5.238199in}{3.581356in}}%
\pgfpathlineto{\pgfqpoint{5.240718in}{3.611438in}}%
\pgfpathlineto{\pgfqpoint{5.245756in}{3.574374in}}%
\pgfpathlineto{\pgfqpoint{5.248274in}{3.578712in}}%
\pgfpathlineto{\pgfqpoint{5.250793in}{3.606172in}}%
\pgfpathlineto{\pgfqpoint{5.253312in}{3.617828in}}%
\pgfpathlineto{\pgfqpoint{5.255831in}{3.650999in}}%
\pgfpathlineto{\pgfqpoint{5.258349in}{3.678446in}}%
\pgfpathlineto{\pgfqpoint{5.260868in}{3.699849in}}%
\pgfpathlineto{\pgfqpoint{5.263387in}{3.704104in}}%
\pgfpathlineto{\pgfqpoint{5.265906in}{3.699970in}}%
\pgfpathlineto{\pgfqpoint{5.268424in}{3.696039in}}%
\pgfpathlineto{\pgfqpoint{5.270943in}{3.683136in}}%
\pgfpathlineto{\pgfqpoint{5.273462in}{3.679577in}}%
\pgfpathlineto{\pgfqpoint{5.275981in}{3.638378in}}%
\pgfpathlineto{\pgfqpoint{5.278499in}{3.630453in}}%
\pgfpathlineto{\pgfqpoint{5.281018in}{3.663410in}}%
\pgfpathlineto{\pgfqpoint{5.283537in}{3.654020in}}%
\pgfpathlineto{\pgfqpoint{5.286056in}{3.673190in}}%
\pgfpathlineto{\pgfqpoint{5.288574in}{3.655692in}}%
\pgfpathlineto{\pgfqpoint{5.291093in}{3.683147in}}%
\pgfpathlineto{\pgfqpoint{5.293612in}{3.714944in}}%
\pgfpathlineto{\pgfqpoint{5.296131in}{3.707281in}}%
\pgfpathlineto{\pgfqpoint{5.298649in}{3.705008in}}%
\pgfpathlineto{\pgfqpoint{5.301168in}{3.750178in}}%
\pgfpathlineto{\pgfqpoint{5.303687in}{3.713448in}}%
\pgfpathlineto{\pgfqpoint{5.306206in}{3.723046in}}%
\pgfpathlineto{\pgfqpoint{5.308724in}{3.718559in}}%
\pgfpathlineto{\pgfqpoint{5.311243in}{3.759313in}}%
\pgfpathlineto{\pgfqpoint{5.313762in}{3.839775in}}%
\pgfpathlineto{\pgfqpoint{5.316281in}{3.815230in}}%
\pgfpathlineto{\pgfqpoint{5.318799in}{3.786622in}}%
\pgfpathlineto{\pgfqpoint{5.321318in}{3.808463in}}%
\pgfpathlineto{\pgfqpoint{5.323837in}{3.793844in}}%
\pgfpathlineto{\pgfqpoint{5.326356in}{3.770434in}}%
\pgfpathlineto{\pgfqpoint{5.328874in}{3.771702in}}%
\pgfpathlineto{\pgfqpoint{5.331393in}{3.768757in}}%
\pgfpathlineto{\pgfqpoint{5.333912in}{3.752268in}}%
\pgfpathlineto{\pgfqpoint{5.336431in}{3.742418in}}%
\pgfpathlineto{\pgfqpoint{5.338949in}{3.734403in}}%
\pgfpathlineto{\pgfqpoint{5.341468in}{3.736270in}}%
\pgfpathlineto{\pgfqpoint{5.343987in}{3.718859in}}%
\pgfpathlineto{\pgfqpoint{5.346506in}{3.740260in}}%
\pgfpathlineto{\pgfqpoint{5.349024in}{3.748036in}}%
\pgfpathlineto{\pgfqpoint{5.351543in}{3.743288in}}%
\pgfpathlineto{\pgfqpoint{5.354062in}{3.748982in}}%
\pgfpathlineto{\pgfqpoint{5.356581in}{3.728600in}}%
\pgfpathlineto{\pgfqpoint{5.359099in}{3.718779in}}%
\pgfpathlineto{\pgfqpoint{5.361618in}{3.716489in}}%
\pgfpathlineto{\pgfqpoint{5.364137in}{3.734387in}}%
\pgfpathlineto{\pgfqpoint{5.366656in}{3.721361in}}%
\pgfpathlineto{\pgfqpoint{5.369174in}{3.704831in}}%
\pgfpathlineto{\pgfqpoint{5.371693in}{3.745184in}}%
\pgfpathlineto{\pgfqpoint{5.374212in}{3.756286in}}%
\pgfpathlineto{\pgfqpoint{5.376731in}{3.754991in}}%
\pgfpathlineto{\pgfqpoint{5.379249in}{3.794572in}}%
\pgfpathlineto{\pgfqpoint{5.381768in}{3.803965in}}%
\pgfpathlineto{\pgfqpoint{5.384287in}{3.767769in}}%
\pgfpathlineto{\pgfqpoint{5.386806in}{3.853251in}}%
\pgfpathlineto{\pgfqpoint{5.389324in}{3.870337in}}%
\pgfpathlineto{\pgfqpoint{5.391843in}{3.859863in}}%
\pgfpathlineto{\pgfqpoint{5.394362in}{3.882356in}}%
\pgfpathlineto{\pgfqpoint{5.396881in}{3.863866in}}%
\pgfpathlineto{\pgfqpoint{5.399399in}{3.915025in}}%
\pgfpathlineto{\pgfqpoint{5.401918in}{3.900967in}}%
\pgfpathlineto{\pgfqpoint{5.404437in}{3.947228in}}%
\pgfpathlineto{\pgfqpoint{5.406956in}{3.955174in}}%
\pgfpathlineto{\pgfqpoint{5.409474in}{3.946341in}}%
\pgfpathlineto{\pgfqpoint{5.411993in}{3.987493in}}%
\pgfpathlineto{\pgfqpoint{5.414512in}{3.950882in}}%
\pgfpathlineto{\pgfqpoint{5.417031in}{3.988414in}}%
\pgfpathlineto{\pgfqpoint{5.419549in}{3.925801in}}%
\pgfpathlineto{\pgfqpoint{5.422068in}{3.946183in}}%
\pgfpathlineto{\pgfqpoint{5.424587in}{3.950214in}}%
\pgfpathlineto{\pgfqpoint{5.427106in}{3.957518in}}%
\pgfpathlineto{\pgfqpoint{5.429624in}{3.962653in}}%
\pgfpathlineto{\pgfqpoint{5.432143in}{3.996533in}}%
\pgfpathlineto{\pgfqpoint{5.434662in}{4.040601in}}%
\pgfpathlineto{\pgfqpoint{5.437181in}{4.008182in}}%
\pgfpathlineto{\pgfqpoint{5.439699in}{3.991285in}}%
\pgfpathlineto{\pgfqpoint{5.442218in}{3.988153in}}%
\pgfpathlineto{\pgfqpoint{5.444737in}{3.994119in}}%
\pgfpathlineto{\pgfqpoint{5.447256in}{3.969178in}}%
\pgfpathlineto{\pgfqpoint{5.449774in}{3.961552in}}%
\pgfpathlineto{\pgfqpoint{5.452293in}{3.976155in}}%
\pgfpathlineto{\pgfqpoint{5.454812in}{4.014499in}}%
\pgfpathlineto{\pgfqpoint{5.457331in}{3.991107in}}%
\pgfpathlineto{\pgfqpoint{5.459849in}{3.955762in}}%
\pgfpathlineto{\pgfqpoint{5.462368in}{3.985628in}}%
\pgfpathlineto{\pgfqpoint{5.464887in}{3.970827in}}%
\pgfpathlineto{\pgfqpoint{5.467406in}{3.982350in}}%
\pgfpathlineto{\pgfqpoint{5.469924in}{3.967519in}}%
\pgfpathlineto{\pgfqpoint{5.472443in}{3.958092in}}%
\pgfpathlineto{\pgfqpoint{5.474962in}{3.956830in}}%
\pgfpathlineto{\pgfqpoint{5.477481in}{3.939747in}}%
\pgfpathlineto{\pgfqpoint{5.479999in}{3.952642in}}%
\pgfpathlineto{\pgfqpoint{5.482518in}{3.948172in}}%
\pgfpathlineto{\pgfqpoint{5.485037in}{3.956096in}}%
\pgfpathlineto{\pgfqpoint{5.487556in}{3.971985in}}%
\pgfpathlineto{\pgfqpoint{5.490074in}{3.948942in}}%
\pgfpathlineto{\pgfqpoint{5.492593in}{3.980544in}}%
\pgfpathlineto{\pgfqpoint{5.492593in}{3.980544in}}%
\pgfusepath{stroke}%
\end{pgfscope}%
\begin{pgfscope}%
\pgfpathrectangle{\pgfqpoint{0.455093in}{0.401080in}}{\pgfqpoint{5.037500in}{4.004000in}}%
\pgfusepath{clip}%
\pgfsetrectcap%
\pgfsetroundjoin%
\pgfsetlinewidth{0.501875pt}%
\definecolor{currentstroke}{rgb}{0.690196,0.188235,0.376471}%
\pgfsetstrokecolor{currentstroke}%
\pgfsetstrokeopacity{0.800000}%
\pgfsetdash{}{0pt}%
\pgfpathmoveto{\pgfqpoint{0.455093in}{1.021044in}}%
\pgfpathlineto{\pgfqpoint{0.457612in}{1.027051in}}%
\pgfpathlineto{\pgfqpoint{0.460131in}{1.016930in}}%
\pgfpathlineto{\pgfqpoint{0.462649in}{1.013788in}}%
\pgfpathlineto{\pgfqpoint{0.465168in}{0.997323in}}%
\pgfpathlineto{\pgfqpoint{0.467687in}{0.995272in}}%
\pgfpathlineto{\pgfqpoint{0.470206in}{0.989529in}}%
\pgfpathlineto{\pgfqpoint{0.472724in}{0.985867in}}%
\pgfpathlineto{\pgfqpoint{0.475243in}{0.977493in}}%
\pgfpathlineto{\pgfqpoint{0.477762in}{0.998727in}}%
\pgfpathlineto{\pgfqpoint{0.480281in}{0.998789in}}%
\pgfpathlineto{\pgfqpoint{0.482799in}{1.005496in}}%
\pgfpathlineto{\pgfqpoint{0.485318in}{0.974480in}}%
\pgfpathlineto{\pgfqpoint{0.487837in}{0.972288in}}%
\pgfpathlineto{\pgfqpoint{0.490356in}{0.991234in}}%
\pgfpathlineto{\pgfqpoint{0.492874in}{0.989176in}}%
\pgfpathlineto{\pgfqpoint{0.495393in}{0.993069in}}%
\pgfpathlineto{\pgfqpoint{0.497912in}{0.970259in}}%
\pgfpathlineto{\pgfqpoint{0.500431in}{0.978912in}}%
\pgfpathlineto{\pgfqpoint{0.502949in}{0.980365in}}%
\pgfpathlineto{\pgfqpoint{0.505468in}{0.978073in}}%
\pgfpathlineto{\pgfqpoint{0.507987in}{0.970864in}}%
\pgfpathlineto{\pgfqpoint{0.510506in}{0.971615in}}%
\pgfpathlineto{\pgfqpoint{0.513024in}{0.969338in}}%
\pgfpathlineto{\pgfqpoint{0.515543in}{0.937834in}}%
\pgfpathlineto{\pgfqpoint{0.518062in}{0.941528in}}%
\pgfpathlineto{\pgfqpoint{0.520581in}{0.967890in}}%
\pgfpathlineto{\pgfqpoint{0.523099in}{0.962027in}}%
\pgfpathlineto{\pgfqpoint{0.525618in}{0.960514in}}%
\pgfpathlineto{\pgfqpoint{0.528137in}{0.952465in}}%
\pgfpathlineto{\pgfqpoint{0.530656in}{0.923709in}}%
\pgfpathlineto{\pgfqpoint{0.533174in}{0.926259in}}%
\pgfpathlineto{\pgfqpoint{0.535693in}{0.921959in}}%
\pgfpathlineto{\pgfqpoint{0.538212in}{0.924458in}}%
\pgfpathlineto{\pgfqpoint{0.540731in}{0.923591in}}%
\pgfpathlineto{\pgfqpoint{0.543249in}{0.930946in}}%
\pgfpathlineto{\pgfqpoint{0.545768in}{0.928741in}}%
\pgfpathlineto{\pgfqpoint{0.550806in}{0.953988in}}%
\pgfpathlineto{\pgfqpoint{0.553324in}{0.978046in}}%
\pgfpathlineto{\pgfqpoint{0.555843in}{0.965658in}}%
\pgfpathlineto{\pgfqpoint{0.558362in}{0.973598in}}%
\pgfpathlineto{\pgfqpoint{0.560881in}{0.987287in}}%
\pgfpathlineto{\pgfqpoint{0.563399in}{0.971712in}}%
\pgfpathlineto{\pgfqpoint{0.565918in}{0.975005in}}%
\pgfpathlineto{\pgfqpoint{0.568437in}{0.982187in}}%
\pgfpathlineto{\pgfqpoint{0.570956in}{0.979770in}}%
\pgfpathlineto{\pgfqpoint{0.573474in}{0.992154in}}%
\pgfpathlineto{\pgfqpoint{0.575993in}{0.983217in}}%
\pgfpathlineto{\pgfqpoint{0.578512in}{0.986857in}}%
\pgfpathlineto{\pgfqpoint{0.581031in}{1.000025in}}%
\pgfpathlineto{\pgfqpoint{0.583549in}{1.030423in}}%
\pgfpathlineto{\pgfqpoint{0.586068in}{1.015553in}}%
\pgfpathlineto{\pgfqpoint{0.588587in}{1.038253in}}%
\pgfpathlineto{\pgfqpoint{0.591106in}{1.029612in}}%
\pgfpathlineto{\pgfqpoint{0.593624in}{1.036136in}}%
\pgfpathlineto{\pgfqpoint{0.596143in}{1.040330in}}%
\pgfpathlineto{\pgfqpoint{0.598662in}{1.026572in}}%
\pgfpathlineto{\pgfqpoint{0.601181in}{1.014623in}}%
\pgfpathlineto{\pgfqpoint{0.603699in}{1.030230in}}%
\pgfpathlineto{\pgfqpoint{0.606218in}{1.048606in}}%
\pgfpathlineto{\pgfqpoint{0.608737in}{1.031912in}}%
\pgfpathlineto{\pgfqpoint{0.611256in}{1.044350in}}%
\pgfpathlineto{\pgfqpoint{0.613774in}{1.039996in}}%
\pgfpathlineto{\pgfqpoint{0.616293in}{1.018575in}}%
\pgfpathlineto{\pgfqpoint{0.618812in}{1.016380in}}%
\pgfpathlineto{\pgfqpoint{0.621331in}{1.005439in}}%
\pgfpathlineto{\pgfqpoint{0.623849in}{0.995848in}}%
\pgfpathlineto{\pgfqpoint{0.626368in}{0.999511in}}%
\pgfpathlineto{\pgfqpoint{0.628887in}{1.019442in}}%
\pgfpathlineto{\pgfqpoint{0.631406in}{1.012899in}}%
\pgfpathlineto{\pgfqpoint{0.633924in}{0.999506in}}%
\pgfpathlineto{\pgfqpoint{0.636443in}{0.999029in}}%
\pgfpathlineto{\pgfqpoint{0.638962in}{0.987768in}}%
\pgfpathlineto{\pgfqpoint{0.641481in}{0.984902in}}%
\pgfpathlineto{\pgfqpoint{0.643999in}{0.985266in}}%
\pgfpathlineto{\pgfqpoint{0.646518in}{0.996791in}}%
\pgfpathlineto{\pgfqpoint{0.649037in}{0.976206in}}%
\pgfpathlineto{\pgfqpoint{0.651556in}{0.980400in}}%
\pgfpathlineto{\pgfqpoint{0.654074in}{0.980395in}}%
\pgfpathlineto{\pgfqpoint{0.656593in}{0.994170in}}%
\pgfpathlineto{\pgfqpoint{0.659112in}{1.003561in}}%
\pgfpathlineto{\pgfqpoint{0.661631in}{1.010046in}}%
\pgfpathlineto{\pgfqpoint{0.664149in}{0.998943in}}%
\pgfpathlineto{\pgfqpoint{0.666668in}{1.002431in}}%
\pgfpathlineto{\pgfqpoint{0.669187in}{1.021003in}}%
\pgfpathlineto{\pgfqpoint{0.671706in}{1.009895in}}%
\pgfpathlineto{\pgfqpoint{0.674224in}{1.024505in}}%
\pgfpathlineto{\pgfqpoint{0.676743in}{1.027955in}}%
\pgfpathlineto{\pgfqpoint{0.679262in}{1.039336in}}%
\pgfpathlineto{\pgfqpoint{0.681781in}{1.041636in}}%
\pgfpathlineto{\pgfqpoint{0.684299in}{1.052209in}}%
\pgfpathlineto{\pgfqpoint{0.686818in}{1.058621in}}%
\pgfpathlineto{\pgfqpoint{0.689337in}{1.041242in}}%
\pgfpathlineto{\pgfqpoint{0.691856in}{1.059753in}}%
\pgfpathlineto{\pgfqpoint{0.694374in}{1.082101in}}%
\pgfpathlineto{\pgfqpoint{0.696893in}{1.076660in}}%
\pgfpathlineto{\pgfqpoint{0.699412in}{1.077237in}}%
\pgfpathlineto{\pgfqpoint{0.701931in}{1.085637in}}%
\pgfpathlineto{\pgfqpoint{0.704449in}{1.073075in}}%
\pgfpathlineto{\pgfqpoint{0.706968in}{1.073724in}}%
\pgfpathlineto{\pgfqpoint{0.709487in}{1.076952in}}%
\pgfpathlineto{\pgfqpoint{0.712006in}{1.063912in}}%
\pgfpathlineto{\pgfqpoint{0.714524in}{1.065599in}}%
\pgfpathlineto{\pgfqpoint{0.717043in}{1.066238in}}%
\pgfpathlineto{\pgfqpoint{0.719562in}{1.065778in}}%
\pgfpathlineto{\pgfqpoint{0.722081in}{1.073019in}}%
\pgfpathlineto{\pgfqpoint{0.724599in}{1.068096in}}%
\pgfpathlineto{\pgfqpoint{0.727118in}{1.067559in}}%
\pgfpathlineto{\pgfqpoint{0.729637in}{1.061858in}}%
\pgfpathlineto{\pgfqpoint{0.732156in}{1.065727in}}%
\pgfpathlineto{\pgfqpoint{0.734674in}{1.053477in}}%
\pgfpathlineto{\pgfqpoint{0.737193in}{1.065664in}}%
\pgfpathlineto{\pgfqpoint{0.739712in}{1.060916in}}%
\pgfpathlineto{\pgfqpoint{0.742231in}{1.065523in}}%
\pgfpathlineto{\pgfqpoint{0.744749in}{1.062259in}}%
\pgfpathlineto{\pgfqpoint{0.747268in}{1.048145in}}%
\pgfpathlineto{\pgfqpoint{0.749787in}{1.038037in}}%
\pgfpathlineto{\pgfqpoint{0.752306in}{1.039800in}}%
\pgfpathlineto{\pgfqpoint{0.754824in}{1.042147in}}%
\pgfpathlineto{\pgfqpoint{0.757343in}{1.051845in}}%
\pgfpathlineto{\pgfqpoint{0.759862in}{1.049275in}}%
\pgfpathlineto{\pgfqpoint{0.762381in}{1.056595in}}%
\pgfpathlineto{\pgfqpoint{0.764899in}{1.064751in}}%
\pgfpathlineto{\pgfqpoint{0.767418in}{1.070208in}}%
\pgfpathlineto{\pgfqpoint{0.769937in}{1.076484in}}%
\pgfpathlineto{\pgfqpoint{0.772456in}{1.076552in}}%
\pgfpathlineto{\pgfqpoint{0.774974in}{1.079410in}}%
\pgfpathlineto{\pgfqpoint{0.777493in}{1.078779in}}%
\pgfpathlineto{\pgfqpoint{0.780012in}{1.061354in}}%
\pgfpathlineto{\pgfqpoint{0.782531in}{1.060630in}}%
\pgfpathlineto{\pgfqpoint{0.785049in}{1.068424in}}%
\pgfpathlineto{\pgfqpoint{0.787568in}{1.085621in}}%
\pgfpathlineto{\pgfqpoint{0.790087in}{1.119180in}}%
\pgfpathlineto{\pgfqpoint{0.792606in}{1.133970in}}%
\pgfpathlineto{\pgfqpoint{0.795124in}{1.136371in}}%
\pgfpathlineto{\pgfqpoint{0.797643in}{1.149280in}}%
\pgfpathlineto{\pgfqpoint{0.800162in}{1.163338in}}%
\pgfpathlineto{\pgfqpoint{0.802681in}{1.161445in}}%
\pgfpathlineto{\pgfqpoint{0.805199in}{1.171619in}}%
\pgfpathlineto{\pgfqpoint{0.807718in}{1.190011in}}%
\pgfpathlineto{\pgfqpoint{0.810237in}{1.183351in}}%
\pgfpathlineto{\pgfqpoint{0.812756in}{1.187587in}}%
\pgfpathlineto{\pgfqpoint{0.815274in}{1.192623in}}%
\pgfpathlineto{\pgfqpoint{0.817793in}{1.206562in}}%
\pgfpathlineto{\pgfqpoint{0.820312in}{1.218159in}}%
\pgfpathlineto{\pgfqpoint{0.822831in}{1.224929in}}%
\pgfpathlineto{\pgfqpoint{0.825349in}{1.228196in}}%
\pgfpathlineto{\pgfqpoint{0.827868in}{1.246465in}}%
\pgfpathlineto{\pgfqpoint{0.830387in}{1.256422in}}%
\pgfpathlineto{\pgfqpoint{0.832906in}{1.250057in}}%
\pgfpathlineto{\pgfqpoint{0.835424in}{1.274885in}}%
\pgfpathlineto{\pgfqpoint{0.837943in}{1.282438in}}%
\pgfpathlineto{\pgfqpoint{0.840462in}{1.281921in}}%
\pgfpathlineto{\pgfqpoint{0.842981in}{1.262058in}}%
\pgfpathlineto{\pgfqpoint{0.845499in}{1.244468in}}%
\pgfpathlineto{\pgfqpoint{0.848018in}{1.254615in}}%
\pgfpathlineto{\pgfqpoint{0.850537in}{1.289031in}}%
\pgfpathlineto{\pgfqpoint{0.853056in}{1.282481in}}%
\pgfpathlineto{\pgfqpoint{0.855574in}{1.283941in}}%
\pgfpathlineto{\pgfqpoint{0.858093in}{1.295815in}}%
\pgfpathlineto{\pgfqpoint{0.860612in}{1.302516in}}%
\pgfpathlineto{\pgfqpoint{0.863131in}{1.291303in}}%
\pgfpathlineto{\pgfqpoint{0.865649in}{1.286006in}}%
\pgfpathlineto{\pgfqpoint{0.868168in}{1.273988in}}%
\pgfpathlineto{\pgfqpoint{0.870687in}{1.286732in}}%
\pgfpathlineto{\pgfqpoint{0.873206in}{1.285746in}}%
\pgfpathlineto{\pgfqpoint{0.875724in}{1.268521in}}%
\pgfpathlineto{\pgfqpoint{0.878243in}{1.239871in}}%
\pgfpathlineto{\pgfqpoint{0.880762in}{1.243150in}}%
\pgfpathlineto{\pgfqpoint{0.883281in}{1.240283in}}%
\pgfpathlineto{\pgfqpoint{0.885799in}{1.245221in}}%
\pgfpathlineto{\pgfqpoint{0.888318in}{1.228372in}}%
\pgfpathlineto{\pgfqpoint{0.890837in}{1.229349in}}%
\pgfpathlineto{\pgfqpoint{0.893356in}{1.210983in}}%
\pgfpathlineto{\pgfqpoint{0.895874in}{1.216087in}}%
\pgfpathlineto{\pgfqpoint{0.898393in}{1.194616in}}%
\pgfpathlineto{\pgfqpoint{0.900912in}{1.197351in}}%
\pgfpathlineto{\pgfqpoint{0.903431in}{1.201211in}}%
\pgfpathlineto{\pgfqpoint{0.905949in}{1.213646in}}%
\pgfpathlineto{\pgfqpoint{0.908468in}{1.219961in}}%
\pgfpathlineto{\pgfqpoint{0.910987in}{1.218875in}}%
\pgfpathlineto{\pgfqpoint{0.913506in}{1.200750in}}%
\pgfpathlineto{\pgfqpoint{0.916024in}{1.199930in}}%
\pgfpathlineto{\pgfqpoint{0.918543in}{1.189218in}}%
\pgfpathlineto{\pgfqpoint{0.921062in}{1.175395in}}%
\pgfpathlineto{\pgfqpoint{0.923581in}{1.147757in}}%
\pgfpathlineto{\pgfqpoint{0.926099in}{1.161722in}}%
\pgfpathlineto{\pgfqpoint{0.928618in}{1.172829in}}%
\pgfpathlineto{\pgfqpoint{0.931137in}{1.167636in}}%
\pgfpathlineto{\pgfqpoint{0.933656in}{1.170963in}}%
\pgfpathlineto{\pgfqpoint{0.936174in}{1.168468in}}%
\pgfpathlineto{\pgfqpoint{0.938693in}{1.143470in}}%
\pgfpathlineto{\pgfqpoint{0.941212in}{1.133405in}}%
\pgfpathlineto{\pgfqpoint{0.943731in}{1.110734in}}%
\pgfpathlineto{\pgfqpoint{0.946249in}{1.101463in}}%
\pgfpathlineto{\pgfqpoint{0.948768in}{1.102284in}}%
\pgfpathlineto{\pgfqpoint{0.951287in}{1.076221in}}%
\pgfpathlineto{\pgfqpoint{0.953806in}{1.071535in}}%
\pgfpathlineto{\pgfqpoint{0.956324in}{1.065533in}}%
\pgfpathlineto{\pgfqpoint{0.958843in}{1.064045in}}%
\pgfpathlineto{\pgfqpoint{0.963881in}{1.087159in}}%
\pgfpathlineto{\pgfqpoint{0.966399in}{1.090726in}}%
\pgfpathlineto{\pgfqpoint{0.968918in}{1.100707in}}%
\pgfpathlineto{\pgfqpoint{0.971437in}{1.087252in}}%
\pgfpathlineto{\pgfqpoint{0.973956in}{1.072699in}}%
\pgfpathlineto{\pgfqpoint{0.976474in}{1.074576in}}%
\pgfpathlineto{\pgfqpoint{0.978993in}{1.079270in}}%
\pgfpathlineto{\pgfqpoint{0.981512in}{1.090411in}}%
\pgfpathlineto{\pgfqpoint{0.984031in}{1.067987in}}%
\pgfpathlineto{\pgfqpoint{0.986549in}{1.075021in}}%
\pgfpathlineto{\pgfqpoint{0.989068in}{1.083258in}}%
\pgfpathlineto{\pgfqpoint{0.991587in}{1.087173in}}%
\pgfpathlineto{\pgfqpoint{0.994106in}{1.107996in}}%
\pgfpathlineto{\pgfqpoint{0.999143in}{1.093848in}}%
\pgfpathlineto{\pgfqpoint{1.001662in}{1.093206in}}%
\pgfpathlineto{\pgfqpoint{1.004181in}{1.091828in}}%
\pgfpathlineto{\pgfqpoint{1.006699in}{1.096243in}}%
\pgfpathlineto{\pgfqpoint{1.009218in}{1.075780in}}%
\pgfpathlineto{\pgfqpoint{1.011737in}{1.087851in}}%
\pgfpathlineto{\pgfqpoint{1.014256in}{1.092326in}}%
\pgfpathlineto{\pgfqpoint{1.016774in}{1.098557in}}%
\pgfpathlineto{\pgfqpoint{1.019293in}{1.108250in}}%
\pgfpathlineto{\pgfqpoint{1.021812in}{1.107794in}}%
\pgfpathlineto{\pgfqpoint{1.024331in}{1.146151in}}%
\pgfpathlineto{\pgfqpoint{1.026849in}{1.124789in}}%
\pgfpathlineto{\pgfqpoint{1.029368in}{1.118651in}}%
\pgfpathlineto{\pgfqpoint{1.031887in}{1.116960in}}%
\pgfpathlineto{\pgfqpoint{1.034406in}{1.136761in}}%
\pgfpathlineto{\pgfqpoint{1.036924in}{1.128788in}}%
\pgfpathlineto{\pgfqpoint{1.039443in}{1.130515in}}%
\pgfpathlineto{\pgfqpoint{1.041962in}{1.101805in}}%
\pgfpathlineto{\pgfqpoint{1.044481in}{1.091816in}}%
\pgfpathlineto{\pgfqpoint{1.046999in}{1.090721in}}%
\pgfpathlineto{\pgfqpoint{1.049518in}{1.069198in}}%
\pgfpathlineto{\pgfqpoint{1.052037in}{1.096794in}}%
\pgfpathlineto{\pgfqpoint{1.054556in}{1.098784in}}%
\pgfpathlineto{\pgfqpoint{1.057074in}{1.106836in}}%
\pgfpathlineto{\pgfqpoint{1.059593in}{1.121742in}}%
\pgfpathlineto{\pgfqpoint{1.062112in}{1.113985in}}%
\pgfpathlineto{\pgfqpoint{1.064631in}{1.136575in}}%
\pgfpathlineto{\pgfqpoint{1.067149in}{1.151827in}}%
\pgfpathlineto{\pgfqpoint{1.069668in}{1.163535in}}%
\pgfpathlineto{\pgfqpoint{1.072187in}{1.166523in}}%
\pgfpathlineto{\pgfqpoint{1.074706in}{1.183847in}}%
\pgfpathlineto{\pgfqpoint{1.077224in}{1.187419in}}%
\pgfpathlineto{\pgfqpoint{1.079743in}{1.188672in}}%
\pgfpathlineto{\pgfqpoint{1.082262in}{1.219073in}}%
\pgfpathlineto{\pgfqpoint{1.084781in}{1.221086in}}%
\pgfpathlineto{\pgfqpoint{1.087299in}{1.220781in}}%
\pgfpathlineto{\pgfqpoint{1.089818in}{1.203733in}}%
\pgfpathlineto{\pgfqpoint{1.092337in}{1.178662in}}%
\pgfpathlineto{\pgfqpoint{1.094856in}{1.161514in}}%
\pgfpathlineto{\pgfqpoint{1.097374in}{1.158526in}}%
\pgfpathlineto{\pgfqpoint{1.099893in}{1.153076in}}%
\pgfpathlineto{\pgfqpoint{1.102412in}{1.156713in}}%
\pgfpathlineto{\pgfqpoint{1.104931in}{1.158248in}}%
\pgfpathlineto{\pgfqpoint{1.107449in}{1.155431in}}%
\pgfpathlineto{\pgfqpoint{1.109968in}{1.170847in}}%
\pgfpathlineto{\pgfqpoint{1.112487in}{1.154682in}}%
\pgfpathlineto{\pgfqpoint{1.115006in}{1.142939in}}%
\pgfpathlineto{\pgfqpoint{1.117524in}{1.141946in}}%
\pgfpathlineto{\pgfqpoint{1.120043in}{1.136703in}}%
\pgfpathlineto{\pgfqpoint{1.122562in}{1.123026in}}%
\pgfpathlineto{\pgfqpoint{1.125081in}{1.137380in}}%
\pgfpathlineto{\pgfqpoint{1.127599in}{1.171753in}}%
\pgfpathlineto{\pgfqpoint{1.130118in}{1.163927in}}%
\pgfpathlineto{\pgfqpoint{1.132637in}{1.147075in}}%
\pgfpathlineto{\pgfqpoint{1.135156in}{1.146349in}}%
\pgfpathlineto{\pgfqpoint{1.137674in}{1.166672in}}%
\pgfpathlineto{\pgfqpoint{1.140193in}{1.191063in}}%
\pgfpathlineto{\pgfqpoint{1.142712in}{1.200145in}}%
\pgfpathlineto{\pgfqpoint{1.145231in}{1.192373in}}%
\pgfpathlineto{\pgfqpoint{1.147749in}{1.196055in}}%
\pgfpathlineto{\pgfqpoint{1.150268in}{1.211406in}}%
\pgfpathlineto{\pgfqpoint{1.152787in}{1.220520in}}%
\pgfpathlineto{\pgfqpoint{1.155306in}{1.220236in}}%
\pgfpathlineto{\pgfqpoint{1.157824in}{1.232153in}}%
\pgfpathlineto{\pgfqpoint{1.162862in}{1.282504in}}%
\pgfpathlineto{\pgfqpoint{1.165381in}{1.295599in}}%
\pgfpathlineto{\pgfqpoint{1.167899in}{1.264981in}}%
\pgfpathlineto{\pgfqpoint{1.170418in}{1.241567in}}%
\pgfpathlineto{\pgfqpoint{1.172937in}{1.277985in}}%
\pgfpathlineto{\pgfqpoint{1.175456in}{1.287729in}}%
\pgfpathlineto{\pgfqpoint{1.177974in}{1.272672in}}%
\pgfpathlineto{\pgfqpoint{1.180493in}{1.274624in}}%
\pgfpathlineto{\pgfqpoint{1.183012in}{1.286300in}}%
\pgfpathlineto{\pgfqpoint{1.185531in}{1.270530in}}%
\pgfpathlineto{\pgfqpoint{1.188049in}{1.270231in}}%
\pgfpathlineto{\pgfqpoint{1.190568in}{1.280604in}}%
\pgfpathlineto{\pgfqpoint{1.193087in}{1.313772in}}%
\pgfpathlineto{\pgfqpoint{1.195606in}{1.299348in}}%
\pgfpathlineto{\pgfqpoint{1.198124in}{1.315723in}}%
\pgfpathlineto{\pgfqpoint{1.200643in}{1.332897in}}%
\pgfpathlineto{\pgfqpoint{1.203162in}{1.322879in}}%
\pgfpathlineto{\pgfqpoint{1.205681in}{1.318536in}}%
\pgfpathlineto{\pgfqpoint{1.208199in}{1.298520in}}%
\pgfpathlineto{\pgfqpoint{1.210718in}{1.295585in}}%
\pgfpathlineto{\pgfqpoint{1.213237in}{1.286705in}}%
\pgfpathlineto{\pgfqpoint{1.215756in}{1.269156in}}%
\pgfpathlineto{\pgfqpoint{1.218274in}{1.288535in}}%
\pgfpathlineto{\pgfqpoint{1.220793in}{1.286188in}}%
\pgfpathlineto{\pgfqpoint{1.223312in}{1.262842in}}%
\pgfpathlineto{\pgfqpoint{1.225831in}{1.265747in}}%
\pgfpathlineto{\pgfqpoint{1.228349in}{1.255644in}}%
\pgfpathlineto{\pgfqpoint{1.230868in}{1.250095in}}%
\pgfpathlineto{\pgfqpoint{1.233387in}{1.246057in}}%
\pgfpathlineto{\pgfqpoint{1.235906in}{1.242257in}}%
\pgfpathlineto{\pgfqpoint{1.238424in}{1.244009in}}%
\pgfpathlineto{\pgfqpoint{1.240943in}{1.255878in}}%
\pgfpathlineto{\pgfqpoint{1.243462in}{1.270907in}}%
\pgfpathlineto{\pgfqpoint{1.245981in}{1.277338in}}%
\pgfpathlineto{\pgfqpoint{1.251018in}{1.318976in}}%
\pgfpathlineto{\pgfqpoint{1.253537in}{1.328057in}}%
\pgfpathlineto{\pgfqpoint{1.256056in}{1.332101in}}%
\pgfpathlineto{\pgfqpoint{1.258574in}{1.338500in}}%
\pgfpathlineto{\pgfqpoint{1.261093in}{1.335809in}}%
\pgfpathlineto{\pgfqpoint{1.263612in}{1.331083in}}%
\pgfpathlineto{\pgfqpoint{1.266131in}{1.330485in}}%
\pgfpathlineto{\pgfqpoint{1.268649in}{1.361001in}}%
\pgfpathlineto{\pgfqpoint{1.271168in}{1.338035in}}%
\pgfpathlineto{\pgfqpoint{1.273687in}{1.357084in}}%
\pgfpathlineto{\pgfqpoint{1.276206in}{1.369360in}}%
\pgfpathlineto{\pgfqpoint{1.278724in}{1.388526in}}%
\pgfpathlineto{\pgfqpoint{1.281243in}{1.396612in}}%
\pgfpathlineto{\pgfqpoint{1.283762in}{1.420578in}}%
\pgfpathlineto{\pgfqpoint{1.286281in}{1.409932in}}%
\pgfpathlineto{\pgfqpoint{1.288799in}{1.447190in}}%
\pgfpathlineto{\pgfqpoint{1.291318in}{1.428166in}}%
\pgfpathlineto{\pgfqpoint{1.293837in}{1.444443in}}%
\pgfpathlineto{\pgfqpoint{1.296356in}{1.447675in}}%
\pgfpathlineto{\pgfqpoint{1.298874in}{1.479211in}}%
\pgfpathlineto{\pgfqpoint{1.301393in}{1.446833in}}%
\pgfpathlineto{\pgfqpoint{1.303912in}{1.432919in}}%
\pgfpathlineto{\pgfqpoint{1.306431in}{1.425429in}}%
\pgfpathlineto{\pgfqpoint{1.308949in}{1.435522in}}%
\pgfpathlineto{\pgfqpoint{1.311468in}{1.472624in}}%
\pgfpathlineto{\pgfqpoint{1.313987in}{1.447868in}}%
\pgfpathlineto{\pgfqpoint{1.316506in}{1.461004in}}%
\pgfpathlineto{\pgfqpoint{1.319024in}{1.458517in}}%
\pgfpathlineto{\pgfqpoint{1.321543in}{1.458508in}}%
\pgfpathlineto{\pgfqpoint{1.324062in}{1.467078in}}%
\pgfpathlineto{\pgfqpoint{1.326581in}{1.450566in}}%
\pgfpathlineto{\pgfqpoint{1.329099in}{1.461411in}}%
\pgfpathlineto{\pgfqpoint{1.331618in}{1.455140in}}%
\pgfpathlineto{\pgfqpoint{1.334137in}{1.462081in}}%
\pgfpathlineto{\pgfqpoint{1.336656in}{1.451776in}}%
\pgfpathlineto{\pgfqpoint{1.339174in}{1.425299in}}%
\pgfpathlineto{\pgfqpoint{1.341693in}{1.429132in}}%
\pgfpathlineto{\pgfqpoint{1.344212in}{1.423966in}}%
\pgfpathlineto{\pgfqpoint{1.346731in}{1.442797in}}%
\pgfpathlineto{\pgfqpoint{1.349249in}{1.450688in}}%
\pgfpathlineto{\pgfqpoint{1.351768in}{1.461230in}}%
\pgfpathlineto{\pgfqpoint{1.354287in}{1.459638in}}%
\pgfpathlineto{\pgfqpoint{1.356806in}{1.474924in}}%
\pgfpathlineto{\pgfqpoint{1.359324in}{1.450966in}}%
\pgfpathlineto{\pgfqpoint{1.361843in}{1.461009in}}%
\pgfpathlineto{\pgfqpoint{1.364362in}{1.448785in}}%
\pgfpathlineto{\pgfqpoint{1.366881in}{1.444851in}}%
\pgfpathlineto{\pgfqpoint{1.369399in}{1.416021in}}%
\pgfpathlineto{\pgfqpoint{1.371918in}{1.432019in}}%
\pgfpathlineto{\pgfqpoint{1.374437in}{1.441297in}}%
\pgfpathlineto{\pgfqpoint{1.376956in}{1.449167in}}%
\pgfpathlineto{\pgfqpoint{1.379474in}{1.444360in}}%
\pgfpathlineto{\pgfqpoint{1.381993in}{1.437860in}}%
\pgfpathlineto{\pgfqpoint{1.384512in}{1.427473in}}%
\pgfpathlineto{\pgfqpoint{1.387031in}{1.424012in}}%
\pgfpathlineto{\pgfqpoint{1.389549in}{1.419999in}}%
\pgfpathlineto{\pgfqpoint{1.392068in}{1.408404in}}%
\pgfpathlineto{\pgfqpoint{1.394587in}{1.376096in}}%
\pgfpathlineto{\pgfqpoint{1.397106in}{1.392699in}}%
\pgfpathlineto{\pgfqpoint{1.399624in}{1.367554in}}%
\pgfpathlineto{\pgfqpoint{1.402143in}{1.366892in}}%
\pgfpathlineto{\pgfqpoint{1.404662in}{1.367998in}}%
\pgfpathlineto{\pgfqpoint{1.407181in}{1.395153in}}%
\pgfpathlineto{\pgfqpoint{1.409699in}{1.432476in}}%
\pgfpathlineto{\pgfqpoint{1.412218in}{1.408385in}}%
\pgfpathlineto{\pgfqpoint{1.414737in}{1.412287in}}%
\pgfpathlineto{\pgfqpoint{1.417256in}{1.404226in}}%
\pgfpathlineto{\pgfqpoint{1.419774in}{1.396533in}}%
\pgfpathlineto{\pgfqpoint{1.422293in}{1.396794in}}%
\pgfpathlineto{\pgfqpoint{1.424812in}{1.367049in}}%
\pgfpathlineto{\pgfqpoint{1.427331in}{1.372556in}}%
\pgfpathlineto{\pgfqpoint{1.429849in}{1.366049in}}%
\pgfpathlineto{\pgfqpoint{1.432368in}{1.357852in}}%
\pgfpathlineto{\pgfqpoint{1.434887in}{1.373030in}}%
\pgfpathlineto{\pgfqpoint{1.437406in}{1.390704in}}%
\pgfpathlineto{\pgfqpoint{1.439924in}{1.369383in}}%
\pgfpathlineto{\pgfqpoint{1.442443in}{1.372153in}}%
\pgfpathlineto{\pgfqpoint{1.444962in}{1.388987in}}%
\pgfpathlineto{\pgfqpoint{1.447481in}{1.390267in}}%
\pgfpathlineto{\pgfqpoint{1.449999in}{1.377623in}}%
\pgfpathlineto{\pgfqpoint{1.452518in}{1.395992in}}%
\pgfpathlineto{\pgfqpoint{1.455037in}{1.412888in}}%
\pgfpathlineto{\pgfqpoint{1.457556in}{1.418909in}}%
\pgfpathlineto{\pgfqpoint{1.460074in}{1.412799in}}%
\pgfpathlineto{\pgfqpoint{1.462593in}{1.415449in}}%
\pgfpathlineto{\pgfqpoint{1.465112in}{1.410070in}}%
\pgfpathlineto{\pgfqpoint{1.467631in}{1.428410in}}%
\pgfpathlineto{\pgfqpoint{1.470149in}{1.436227in}}%
\pgfpathlineto{\pgfqpoint{1.472668in}{1.451499in}}%
\pgfpathlineto{\pgfqpoint{1.475187in}{1.456980in}}%
\pgfpathlineto{\pgfqpoint{1.477706in}{1.432179in}}%
\pgfpathlineto{\pgfqpoint{1.480224in}{1.431258in}}%
\pgfpathlineto{\pgfqpoint{1.482743in}{1.434554in}}%
\pgfpathlineto{\pgfqpoint{1.485262in}{1.414253in}}%
\pgfpathlineto{\pgfqpoint{1.487781in}{1.410014in}}%
\pgfpathlineto{\pgfqpoint{1.490299in}{1.407712in}}%
\pgfpathlineto{\pgfqpoint{1.492818in}{1.414266in}}%
\pgfpathlineto{\pgfqpoint{1.495337in}{1.417165in}}%
\pgfpathlineto{\pgfqpoint{1.497856in}{1.421239in}}%
\pgfpathlineto{\pgfqpoint{1.500374in}{1.415590in}}%
\pgfpathlineto{\pgfqpoint{1.502893in}{1.439907in}}%
\pgfpathlineto{\pgfqpoint{1.505412in}{1.436938in}}%
\pgfpathlineto{\pgfqpoint{1.507931in}{1.437270in}}%
\pgfpathlineto{\pgfqpoint{1.512968in}{1.406388in}}%
\pgfpathlineto{\pgfqpoint{1.515487in}{1.398670in}}%
\pgfpathlineto{\pgfqpoint{1.518006in}{1.376961in}}%
\pgfpathlineto{\pgfqpoint{1.520524in}{1.367562in}}%
\pgfpathlineto{\pgfqpoint{1.523043in}{1.333425in}}%
\pgfpathlineto{\pgfqpoint{1.525562in}{1.342177in}}%
\pgfpathlineto{\pgfqpoint{1.528081in}{1.342151in}}%
\pgfpathlineto{\pgfqpoint{1.530599in}{1.347693in}}%
\pgfpathlineto{\pgfqpoint{1.533118in}{1.377572in}}%
\pgfpathlineto{\pgfqpoint{1.535637in}{1.356787in}}%
\pgfpathlineto{\pgfqpoint{1.538156in}{1.348017in}}%
\pgfpathlineto{\pgfqpoint{1.540674in}{1.324903in}}%
\pgfpathlineto{\pgfqpoint{1.543193in}{1.335648in}}%
\pgfpathlineto{\pgfqpoint{1.545712in}{1.332393in}}%
\pgfpathlineto{\pgfqpoint{1.548231in}{1.331314in}}%
\pgfpathlineto{\pgfqpoint{1.550749in}{1.354293in}}%
\pgfpathlineto{\pgfqpoint{1.553268in}{1.382386in}}%
\pgfpathlineto{\pgfqpoint{1.555787in}{1.361202in}}%
\pgfpathlineto{\pgfqpoint{1.558306in}{1.373214in}}%
\pgfpathlineto{\pgfqpoint{1.560824in}{1.409293in}}%
\pgfpathlineto{\pgfqpoint{1.563343in}{1.404075in}}%
\pgfpathlineto{\pgfqpoint{1.565862in}{1.409304in}}%
\pgfpathlineto{\pgfqpoint{1.568381in}{1.407538in}}%
\pgfpathlineto{\pgfqpoint{1.570899in}{1.414567in}}%
\pgfpathlineto{\pgfqpoint{1.573418in}{1.436627in}}%
\pgfpathlineto{\pgfqpoint{1.575937in}{1.454629in}}%
\pgfpathlineto{\pgfqpoint{1.578456in}{1.481829in}}%
\pgfpathlineto{\pgfqpoint{1.580974in}{1.502535in}}%
\pgfpathlineto{\pgfqpoint{1.583493in}{1.500241in}}%
\pgfpathlineto{\pgfqpoint{1.586012in}{1.514173in}}%
\pgfpathlineto{\pgfqpoint{1.588531in}{1.501870in}}%
\pgfpathlineto{\pgfqpoint{1.591049in}{1.520072in}}%
\pgfpathlineto{\pgfqpoint{1.593568in}{1.526861in}}%
\pgfpathlineto{\pgfqpoint{1.596087in}{1.540739in}}%
\pgfpathlineto{\pgfqpoint{1.598606in}{1.544876in}}%
\pgfpathlineto{\pgfqpoint{1.601124in}{1.549789in}}%
\pgfpathlineto{\pgfqpoint{1.603643in}{1.545841in}}%
\pgfpathlineto{\pgfqpoint{1.606162in}{1.558667in}}%
\pgfpathlineto{\pgfqpoint{1.608681in}{1.552137in}}%
\pgfpathlineto{\pgfqpoint{1.611199in}{1.551596in}}%
\pgfpathlineto{\pgfqpoint{1.613718in}{1.567190in}}%
\pgfpathlineto{\pgfqpoint{1.616237in}{1.572286in}}%
\pgfpathlineto{\pgfqpoint{1.618756in}{1.555369in}}%
\pgfpathlineto{\pgfqpoint{1.621274in}{1.546835in}}%
\pgfpathlineto{\pgfqpoint{1.623793in}{1.534144in}}%
\pgfpathlineto{\pgfqpoint{1.626312in}{1.531905in}}%
\pgfpathlineto{\pgfqpoint{1.628831in}{1.529058in}}%
\pgfpathlineto{\pgfqpoint{1.631349in}{1.555850in}}%
\pgfpathlineto{\pgfqpoint{1.633868in}{1.545184in}}%
\pgfpathlineto{\pgfqpoint{1.636387in}{1.540249in}}%
\pgfpathlineto{\pgfqpoint{1.638906in}{1.551852in}}%
\pgfpathlineto{\pgfqpoint{1.641424in}{1.556051in}}%
\pgfpathlineto{\pgfqpoint{1.643943in}{1.563578in}}%
\pgfpathlineto{\pgfqpoint{1.646462in}{1.545410in}}%
\pgfpathlineto{\pgfqpoint{1.648981in}{1.538911in}}%
\pgfpathlineto{\pgfqpoint{1.651499in}{1.540005in}}%
\pgfpathlineto{\pgfqpoint{1.654018in}{1.545762in}}%
\pgfpathlineto{\pgfqpoint{1.656537in}{1.527127in}}%
\pgfpathlineto{\pgfqpoint{1.659056in}{1.537426in}}%
\pgfpathlineto{\pgfqpoint{1.661574in}{1.555333in}}%
\pgfpathlineto{\pgfqpoint{1.664093in}{1.557165in}}%
\pgfpathlineto{\pgfqpoint{1.666612in}{1.532837in}}%
\pgfpathlineto{\pgfqpoint{1.669131in}{1.539185in}}%
\pgfpathlineto{\pgfqpoint{1.671649in}{1.532736in}}%
\pgfpathlineto{\pgfqpoint{1.674168in}{1.511625in}}%
\pgfpathlineto{\pgfqpoint{1.676687in}{1.529927in}}%
\pgfpathlineto{\pgfqpoint{1.679206in}{1.563387in}}%
\pgfpathlineto{\pgfqpoint{1.681724in}{1.533367in}}%
\pgfpathlineto{\pgfqpoint{1.684243in}{1.522477in}}%
\pgfpathlineto{\pgfqpoint{1.686762in}{1.528915in}}%
\pgfpathlineto{\pgfqpoint{1.689281in}{1.527615in}}%
\pgfpathlineto{\pgfqpoint{1.691799in}{1.544839in}}%
\pgfpathlineto{\pgfqpoint{1.694318in}{1.541936in}}%
\pgfpathlineto{\pgfqpoint{1.696837in}{1.531265in}}%
\pgfpathlineto{\pgfqpoint{1.699356in}{1.527059in}}%
\pgfpathlineto{\pgfqpoint{1.701874in}{1.502524in}}%
\pgfpathlineto{\pgfqpoint{1.704393in}{1.527801in}}%
\pgfpathlineto{\pgfqpoint{1.709431in}{1.554774in}}%
\pgfpathlineto{\pgfqpoint{1.711949in}{1.537025in}}%
\pgfpathlineto{\pgfqpoint{1.714468in}{1.526192in}}%
\pgfpathlineto{\pgfqpoint{1.716987in}{1.534789in}}%
\pgfpathlineto{\pgfqpoint{1.719506in}{1.536551in}}%
\pgfpathlineto{\pgfqpoint{1.722024in}{1.526048in}}%
\pgfpathlineto{\pgfqpoint{1.724543in}{1.523484in}}%
\pgfpathlineto{\pgfqpoint{1.727062in}{1.533460in}}%
\pgfpathlineto{\pgfqpoint{1.729581in}{1.552227in}}%
\pgfpathlineto{\pgfqpoint{1.732099in}{1.558979in}}%
\pgfpathlineto{\pgfqpoint{1.734618in}{1.565231in}}%
\pgfpathlineto{\pgfqpoint{1.737137in}{1.547277in}}%
\pgfpathlineto{\pgfqpoint{1.739656in}{1.548040in}}%
\pgfpathlineto{\pgfqpoint{1.742174in}{1.554618in}}%
\pgfpathlineto{\pgfqpoint{1.744693in}{1.541182in}}%
\pgfpathlineto{\pgfqpoint{1.747212in}{1.519893in}}%
\pgfpathlineto{\pgfqpoint{1.749731in}{1.507634in}}%
\pgfpathlineto{\pgfqpoint{1.752249in}{1.502911in}}%
\pgfpathlineto{\pgfqpoint{1.754768in}{1.491574in}}%
\pgfpathlineto{\pgfqpoint{1.757287in}{1.472028in}}%
\pgfpathlineto{\pgfqpoint{1.759806in}{1.482896in}}%
\pgfpathlineto{\pgfqpoint{1.762324in}{1.497281in}}%
\pgfpathlineto{\pgfqpoint{1.764843in}{1.497174in}}%
\pgfpathlineto{\pgfqpoint{1.767362in}{1.505786in}}%
\pgfpathlineto{\pgfqpoint{1.769881in}{1.505777in}}%
\pgfpathlineto{\pgfqpoint{1.772399in}{1.508803in}}%
\pgfpathlineto{\pgfqpoint{1.774918in}{1.501877in}}%
\pgfpathlineto{\pgfqpoint{1.777437in}{1.474965in}}%
\pgfpathlineto{\pgfqpoint{1.779956in}{1.479181in}}%
\pgfpathlineto{\pgfqpoint{1.782474in}{1.497563in}}%
\pgfpathlineto{\pgfqpoint{1.784993in}{1.492119in}}%
\pgfpathlineto{\pgfqpoint{1.792549in}{1.511136in}}%
\pgfpathlineto{\pgfqpoint{1.795068in}{1.501167in}}%
\pgfpathlineto{\pgfqpoint{1.797587in}{1.502362in}}%
\pgfpathlineto{\pgfqpoint{1.800106in}{1.516563in}}%
\pgfpathlineto{\pgfqpoint{1.802624in}{1.542620in}}%
\pgfpathlineto{\pgfqpoint{1.805143in}{1.570425in}}%
\pgfpathlineto{\pgfqpoint{1.807662in}{1.566863in}}%
\pgfpathlineto{\pgfqpoint{1.810181in}{1.566459in}}%
\pgfpathlineto{\pgfqpoint{1.812699in}{1.568995in}}%
\pgfpathlineto{\pgfqpoint{1.815218in}{1.577206in}}%
\pgfpathlineto{\pgfqpoint{1.817737in}{1.560282in}}%
\pgfpathlineto{\pgfqpoint{1.820256in}{1.582749in}}%
\pgfpathlineto{\pgfqpoint{1.822774in}{1.586779in}}%
\pgfpathlineto{\pgfqpoint{1.825293in}{1.588446in}}%
\pgfpathlineto{\pgfqpoint{1.827812in}{1.585562in}}%
\pgfpathlineto{\pgfqpoint{1.830331in}{1.565395in}}%
\pgfpathlineto{\pgfqpoint{1.832849in}{1.573187in}}%
\pgfpathlineto{\pgfqpoint{1.835368in}{1.581350in}}%
\pgfpathlineto{\pgfqpoint{1.837887in}{1.598820in}}%
\pgfpathlineto{\pgfqpoint{1.840406in}{1.591003in}}%
\pgfpathlineto{\pgfqpoint{1.842924in}{1.587572in}}%
\pgfpathlineto{\pgfqpoint{1.845443in}{1.583470in}}%
\pgfpathlineto{\pgfqpoint{1.847962in}{1.585538in}}%
\pgfpathlineto{\pgfqpoint{1.850481in}{1.613045in}}%
\pgfpathlineto{\pgfqpoint{1.852999in}{1.602960in}}%
\pgfpathlineto{\pgfqpoint{1.855518in}{1.589466in}}%
\pgfpathlineto{\pgfqpoint{1.858037in}{1.591527in}}%
\pgfpathlineto{\pgfqpoint{1.860556in}{1.595323in}}%
\pgfpathlineto{\pgfqpoint{1.863074in}{1.590384in}}%
\pgfpathlineto{\pgfqpoint{1.868112in}{1.598880in}}%
\pgfpathlineto{\pgfqpoint{1.870631in}{1.591515in}}%
\pgfpathlineto{\pgfqpoint{1.873149in}{1.604165in}}%
\pgfpathlineto{\pgfqpoint{1.875668in}{1.605866in}}%
\pgfpathlineto{\pgfqpoint{1.878187in}{1.585561in}}%
\pgfpathlineto{\pgfqpoint{1.880706in}{1.607904in}}%
\pgfpathlineto{\pgfqpoint{1.883224in}{1.580438in}}%
\pgfpathlineto{\pgfqpoint{1.885743in}{1.555559in}}%
\pgfpathlineto{\pgfqpoint{1.888262in}{1.568978in}}%
\pgfpathlineto{\pgfqpoint{1.890781in}{1.555316in}}%
\pgfpathlineto{\pgfqpoint{1.893299in}{1.550258in}}%
\pgfpathlineto{\pgfqpoint{1.895818in}{1.540953in}}%
\pgfpathlineto{\pgfqpoint{1.898337in}{1.548595in}}%
\pgfpathlineto{\pgfqpoint{1.900856in}{1.543566in}}%
\pgfpathlineto{\pgfqpoint{1.903374in}{1.518772in}}%
\pgfpathlineto{\pgfqpoint{1.905893in}{1.502653in}}%
\pgfpathlineto{\pgfqpoint{1.908412in}{1.497910in}}%
\pgfpathlineto{\pgfqpoint{1.910931in}{1.499022in}}%
\pgfpathlineto{\pgfqpoint{1.913449in}{1.492888in}}%
\pgfpathlineto{\pgfqpoint{1.915968in}{1.508700in}}%
\pgfpathlineto{\pgfqpoint{1.918487in}{1.529346in}}%
\pgfpathlineto{\pgfqpoint{1.921006in}{1.527593in}}%
\pgfpathlineto{\pgfqpoint{1.923524in}{1.545137in}}%
\pgfpathlineto{\pgfqpoint{1.926043in}{1.534559in}}%
\pgfpathlineto{\pgfqpoint{1.928562in}{1.532194in}}%
\pgfpathlineto{\pgfqpoint{1.931081in}{1.521363in}}%
\pgfpathlineto{\pgfqpoint{1.933599in}{1.503966in}}%
\pgfpathlineto{\pgfqpoint{1.936118in}{1.496011in}}%
\pgfpathlineto{\pgfqpoint{1.938637in}{1.470219in}}%
\pgfpathlineto{\pgfqpoint{1.941156in}{1.496516in}}%
\pgfpathlineto{\pgfqpoint{1.943674in}{1.505088in}}%
\pgfpathlineto{\pgfqpoint{1.946193in}{1.508574in}}%
\pgfpathlineto{\pgfqpoint{1.948712in}{1.521454in}}%
\pgfpathlineto{\pgfqpoint{1.951231in}{1.510341in}}%
\pgfpathlineto{\pgfqpoint{1.953749in}{1.502411in}}%
\pgfpathlineto{\pgfqpoint{1.956268in}{1.503674in}}%
\pgfpathlineto{\pgfqpoint{1.958787in}{1.512729in}}%
\pgfpathlineto{\pgfqpoint{1.961306in}{1.538171in}}%
\pgfpathlineto{\pgfqpoint{1.963824in}{1.522203in}}%
\pgfpathlineto{\pgfqpoint{1.966343in}{1.502764in}}%
\pgfpathlineto{\pgfqpoint{1.968862in}{1.493440in}}%
\pgfpathlineto{\pgfqpoint{1.971381in}{1.481470in}}%
\pgfpathlineto{\pgfqpoint{1.973899in}{1.483181in}}%
\pgfpathlineto{\pgfqpoint{1.976418in}{1.459619in}}%
\pgfpathlineto{\pgfqpoint{1.978937in}{1.453797in}}%
\pgfpathlineto{\pgfqpoint{1.981456in}{1.464530in}}%
\pgfpathlineto{\pgfqpoint{1.983974in}{1.453525in}}%
\pgfpathlineto{\pgfqpoint{1.986493in}{1.445590in}}%
\pgfpathlineto{\pgfqpoint{1.989012in}{1.464921in}}%
\pgfpathlineto{\pgfqpoint{1.991531in}{1.480727in}}%
\pgfpathlineto{\pgfqpoint{1.994049in}{1.468373in}}%
\pgfpathlineto{\pgfqpoint{1.996568in}{1.450892in}}%
\pgfpathlineto{\pgfqpoint{1.999087in}{1.431222in}}%
\pgfpathlineto{\pgfqpoint{2.001606in}{1.425879in}}%
\pgfpathlineto{\pgfqpoint{2.004124in}{1.426419in}}%
\pgfpathlineto{\pgfqpoint{2.006643in}{1.439035in}}%
\pgfpathlineto{\pgfqpoint{2.009162in}{1.439177in}}%
\pgfpathlineto{\pgfqpoint{2.011681in}{1.444296in}}%
\pgfpathlineto{\pgfqpoint{2.014199in}{1.455230in}}%
\pgfpathlineto{\pgfqpoint{2.016718in}{1.452648in}}%
\pgfpathlineto{\pgfqpoint{2.019237in}{1.431633in}}%
\pgfpathlineto{\pgfqpoint{2.021756in}{1.416792in}}%
\pgfpathlineto{\pgfqpoint{2.024274in}{1.413608in}}%
\pgfpathlineto{\pgfqpoint{2.026793in}{1.430797in}}%
\pgfpathlineto{\pgfqpoint{2.029312in}{1.451520in}}%
\pgfpathlineto{\pgfqpoint{2.031831in}{1.468510in}}%
\pgfpathlineto{\pgfqpoint{2.034349in}{1.441197in}}%
\pgfpathlineto{\pgfqpoint{2.036868in}{1.472728in}}%
\pgfpathlineto{\pgfqpoint{2.039387in}{1.467336in}}%
\pgfpathlineto{\pgfqpoint{2.041906in}{1.462998in}}%
\pgfpathlineto{\pgfqpoint{2.044424in}{1.485159in}}%
\pgfpathlineto{\pgfqpoint{2.046943in}{1.484757in}}%
\pgfpathlineto{\pgfqpoint{2.049462in}{1.475058in}}%
\pgfpathlineto{\pgfqpoint{2.051981in}{1.480768in}}%
\pgfpathlineto{\pgfqpoint{2.054499in}{1.478063in}}%
\pgfpathlineto{\pgfqpoint{2.057018in}{1.449655in}}%
\pgfpathlineto{\pgfqpoint{2.059537in}{1.435591in}}%
\pgfpathlineto{\pgfqpoint{2.062056in}{1.442426in}}%
\pgfpathlineto{\pgfqpoint{2.064574in}{1.458447in}}%
\pgfpathlineto{\pgfqpoint{2.067093in}{1.437645in}}%
\pgfpathlineto{\pgfqpoint{2.069612in}{1.443505in}}%
\pgfpathlineto{\pgfqpoint{2.072131in}{1.456552in}}%
\pgfpathlineto{\pgfqpoint{2.074649in}{1.441597in}}%
\pgfpathlineto{\pgfqpoint{2.077168in}{1.415710in}}%
\pgfpathlineto{\pgfqpoint{2.079687in}{1.410170in}}%
\pgfpathlineto{\pgfqpoint{2.082206in}{1.413752in}}%
\pgfpathlineto{\pgfqpoint{2.084724in}{1.422299in}}%
\pgfpathlineto{\pgfqpoint{2.087243in}{1.407599in}}%
\pgfpathlineto{\pgfqpoint{2.089762in}{1.409458in}}%
\pgfpathlineto{\pgfqpoint{2.092281in}{1.410858in}}%
\pgfpathlineto{\pgfqpoint{2.094799in}{1.407711in}}%
\pgfpathlineto{\pgfqpoint{2.097318in}{1.408429in}}%
\pgfpathlineto{\pgfqpoint{2.099837in}{1.398153in}}%
\pgfpathlineto{\pgfqpoint{2.102356in}{1.403606in}}%
\pgfpathlineto{\pgfqpoint{2.104874in}{1.382724in}}%
\pgfpathlineto{\pgfqpoint{2.107393in}{1.399673in}}%
\pgfpathlineto{\pgfqpoint{2.109912in}{1.415168in}}%
\pgfpathlineto{\pgfqpoint{2.112431in}{1.420357in}}%
\pgfpathlineto{\pgfqpoint{2.114949in}{1.415416in}}%
\pgfpathlineto{\pgfqpoint{2.117468in}{1.436453in}}%
\pgfpathlineto{\pgfqpoint{2.119987in}{1.442331in}}%
\pgfpathlineto{\pgfqpoint{2.122506in}{1.420706in}}%
\pgfpathlineto{\pgfqpoint{2.125024in}{1.427082in}}%
\pgfpathlineto{\pgfqpoint{2.127543in}{1.401931in}}%
\pgfpathlineto{\pgfqpoint{2.130062in}{1.413866in}}%
\pgfpathlineto{\pgfqpoint{2.132581in}{1.422903in}}%
\pgfpathlineto{\pgfqpoint{2.135099in}{1.450549in}}%
\pgfpathlineto{\pgfqpoint{2.137618in}{1.438427in}}%
\pgfpathlineto{\pgfqpoint{2.140137in}{1.441241in}}%
\pgfpathlineto{\pgfqpoint{2.142656in}{1.439229in}}%
\pgfpathlineto{\pgfqpoint{2.145174in}{1.428950in}}%
\pgfpathlineto{\pgfqpoint{2.147693in}{1.419634in}}%
\pgfpathlineto{\pgfqpoint{2.150212in}{1.419793in}}%
\pgfpathlineto{\pgfqpoint{2.152731in}{1.412108in}}%
\pgfpathlineto{\pgfqpoint{2.155249in}{1.416423in}}%
\pgfpathlineto{\pgfqpoint{2.157768in}{1.434186in}}%
\pgfpathlineto{\pgfqpoint{2.160287in}{1.426159in}}%
\pgfpathlineto{\pgfqpoint{2.162806in}{1.416869in}}%
\pgfpathlineto{\pgfqpoint{2.165324in}{1.415008in}}%
\pgfpathlineto{\pgfqpoint{2.167843in}{1.406775in}}%
\pgfpathlineto{\pgfqpoint{2.170362in}{1.402237in}}%
\pgfpathlineto{\pgfqpoint{2.172881in}{1.417604in}}%
\pgfpathlineto{\pgfqpoint{2.175399in}{1.401959in}}%
\pgfpathlineto{\pgfqpoint{2.177918in}{1.398875in}}%
\pgfpathlineto{\pgfqpoint{2.180437in}{1.410150in}}%
\pgfpathlineto{\pgfqpoint{2.182956in}{1.397132in}}%
\pgfpathlineto{\pgfqpoint{2.185474in}{1.414981in}}%
\pgfpathlineto{\pgfqpoint{2.187993in}{1.448306in}}%
\pgfpathlineto{\pgfqpoint{2.190512in}{1.424474in}}%
\pgfpathlineto{\pgfqpoint{2.193031in}{1.411620in}}%
\pgfpathlineto{\pgfqpoint{2.195549in}{1.415493in}}%
\pgfpathlineto{\pgfqpoint{2.198068in}{1.399333in}}%
\pgfpathlineto{\pgfqpoint{2.200587in}{1.402697in}}%
\pgfpathlineto{\pgfqpoint{2.203106in}{1.397887in}}%
\pgfpathlineto{\pgfqpoint{2.205624in}{1.395797in}}%
\pgfpathlineto{\pgfqpoint{2.208143in}{1.402381in}}%
\pgfpathlineto{\pgfqpoint{2.210662in}{1.381706in}}%
\pgfpathlineto{\pgfqpoint{2.213181in}{1.377137in}}%
\pgfpathlineto{\pgfqpoint{2.215699in}{1.383605in}}%
\pgfpathlineto{\pgfqpoint{2.218218in}{1.379161in}}%
\pgfpathlineto{\pgfqpoint{2.220737in}{1.362588in}}%
\pgfpathlineto{\pgfqpoint{2.223256in}{1.355652in}}%
\pgfpathlineto{\pgfqpoint{2.225774in}{1.343730in}}%
\pgfpathlineto{\pgfqpoint{2.230812in}{1.375502in}}%
\pgfpathlineto{\pgfqpoint{2.233331in}{1.373409in}}%
\pgfpathlineto{\pgfqpoint{2.235849in}{1.373051in}}%
\pgfpathlineto{\pgfqpoint{2.238368in}{1.378924in}}%
\pgfpathlineto{\pgfqpoint{2.240887in}{1.386959in}}%
\pgfpathlineto{\pgfqpoint{2.243406in}{1.395915in}}%
\pgfpathlineto{\pgfqpoint{2.245924in}{1.402744in}}%
\pgfpathlineto{\pgfqpoint{2.248443in}{1.383392in}}%
\pgfpathlineto{\pgfqpoint{2.250962in}{1.377267in}}%
\pgfpathlineto{\pgfqpoint{2.253481in}{1.373864in}}%
\pgfpathlineto{\pgfqpoint{2.255999in}{1.367188in}}%
\pgfpathlineto{\pgfqpoint{2.258518in}{1.384804in}}%
\pgfpathlineto{\pgfqpoint{2.261037in}{1.364800in}}%
\pgfpathlineto{\pgfqpoint{2.263556in}{1.343228in}}%
\pgfpathlineto{\pgfqpoint{2.266074in}{1.327895in}}%
\pgfpathlineto{\pgfqpoint{2.268593in}{1.324582in}}%
\pgfpathlineto{\pgfqpoint{2.271112in}{1.338360in}}%
\pgfpathlineto{\pgfqpoint{2.273631in}{1.324492in}}%
\pgfpathlineto{\pgfqpoint{2.278668in}{1.345635in}}%
\pgfpathlineto{\pgfqpoint{2.281187in}{1.362242in}}%
\pgfpathlineto{\pgfqpoint{2.283706in}{1.364459in}}%
\pgfpathlineto{\pgfqpoint{2.286224in}{1.363900in}}%
\pgfpathlineto{\pgfqpoint{2.288743in}{1.342894in}}%
\pgfpathlineto{\pgfqpoint{2.291262in}{1.353518in}}%
\pgfpathlineto{\pgfqpoint{2.293781in}{1.356401in}}%
\pgfpathlineto{\pgfqpoint{2.296299in}{1.351911in}}%
\pgfpathlineto{\pgfqpoint{2.298818in}{1.348239in}}%
\pgfpathlineto{\pgfqpoint{2.301337in}{1.324003in}}%
\pgfpathlineto{\pgfqpoint{2.303856in}{1.326878in}}%
\pgfpathlineto{\pgfqpoint{2.306374in}{1.333962in}}%
\pgfpathlineto{\pgfqpoint{2.308893in}{1.343933in}}%
\pgfpathlineto{\pgfqpoint{2.311412in}{1.347192in}}%
\pgfpathlineto{\pgfqpoint{2.313931in}{1.345145in}}%
\pgfpathlineto{\pgfqpoint{2.316449in}{1.350036in}}%
\pgfpathlineto{\pgfqpoint{2.318968in}{1.358574in}}%
\pgfpathlineto{\pgfqpoint{2.321487in}{1.391490in}}%
\pgfpathlineto{\pgfqpoint{2.324006in}{1.398751in}}%
\pgfpathlineto{\pgfqpoint{2.326524in}{1.389868in}}%
\pgfpathlineto{\pgfqpoint{2.329043in}{1.392483in}}%
\pgfpathlineto{\pgfqpoint{2.331562in}{1.373100in}}%
\pgfpathlineto{\pgfqpoint{2.334081in}{1.392945in}}%
\pgfpathlineto{\pgfqpoint{2.336599in}{1.397400in}}%
\pgfpathlineto{\pgfqpoint{2.339118in}{1.413479in}}%
\pgfpathlineto{\pgfqpoint{2.341637in}{1.404593in}}%
\pgfpathlineto{\pgfqpoint{2.344156in}{1.398126in}}%
\pgfpathlineto{\pgfqpoint{2.346674in}{1.382813in}}%
\pgfpathlineto{\pgfqpoint{2.349193in}{1.384294in}}%
\pgfpathlineto{\pgfqpoint{2.351712in}{1.405395in}}%
\pgfpathlineto{\pgfqpoint{2.354231in}{1.424167in}}%
\pgfpathlineto{\pgfqpoint{2.356749in}{1.412231in}}%
\pgfpathlineto{\pgfqpoint{2.359268in}{1.397227in}}%
\pgfpathlineto{\pgfqpoint{2.361787in}{1.388753in}}%
\pgfpathlineto{\pgfqpoint{2.364306in}{1.383041in}}%
\pgfpathlineto{\pgfqpoint{2.366824in}{1.434403in}}%
\pgfpathlineto{\pgfqpoint{2.369343in}{1.424366in}}%
\pgfpathlineto{\pgfqpoint{2.371862in}{1.425404in}}%
\pgfpathlineto{\pgfqpoint{2.374381in}{1.430322in}}%
\pgfpathlineto{\pgfqpoint{2.376899in}{1.439631in}}%
\pgfpathlineto{\pgfqpoint{2.379418in}{1.406511in}}%
\pgfpathlineto{\pgfqpoint{2.381937in}{1.409220in}}%
\pgfpathlineto{\pgfqpoint{2.384456in}{1.414111in}}%
\pgfpathlineto{\pgfqpoint{2.386974in}{1.436663in}}%
\pgfpathlineto{\pgfqpoint{2.389493in}{1.443388in}}%
\pgfpathlineto{\pgfqpoint{2.392012in}{1.445809in}}%
\pgfpathlineto{\pgfqpoint{2.394531in}{1.445308in}}%
\pgfpathlineto{\pgfqpoint{2.397049in}{1.452981in}}%
\pgfpathlineto{\pgfqpoint{2.402087in}{1.419264in}}%
\pgfpathlineto{\pgfqpoint{2.404606in}{1.419592in}}%
\pgfpathlineto{\pgfqpoint{2.407124in}{1.408313in}}%
\pgfpathlineto{\pgfqpoint{2.409643in}{1.423410in}}%
\pgfpathlineto{\pgfqpoint{2.412162in}{1.416239in}}%
\pgfpathlineto{\pgfqpoint{2.414681in}{1.395946in}}%
\pgfpathlineto{\pgfqpoint{2.417199in}{1.405738in}}%
\pgfpathlineto{\pgfqpoint{2.419718in}{1.401888in}}%
\pgfpathlineto{\pgfqpoint{2.422237in}{1.411732in}}%
\pgfpathlineto{\pgfqpoint{2.424756in}{1.400666in}}%
\pgfpathlineto{\pgfqpoint{2.427274in}{1.398350in}}%
\pgfpathlineto{\pgfqpoint{2.429793in}{1.413217in}}%
\pgfpathlineto{\pgfqpoint{2.432312in}{1.419954in}}%
\pgfpathlineto{\pgfqpoint{2.434831in}{1.438374in}}%
\pgfpathlineto{\pgfqpoint{2.437349in}{1.462864in}}%
\pgfpathlineto{\pgfqpoint{2.439868in}{1.468611in}}%
\pgfpathlineto{\pgfqpoint{2.442387in}{1.480725in}}%
\pgfpathlineto{\pgfqpoint{2.444906in}{1.493616in}}%
\pgfpathlineto{\pgfqpoint{2.447424in}{1.489429in}}%
\pgfpathlineto{\pgfqpoint{2.449943in}{1.468263in}}%
\pgfpathlineto{\pgfqpoint{2.452462in}{1.488293in}}%
\pgfpathlineto{\pgfqpoint{2.454981in}{1.489693in}}%
\pgfpathlineto{\pgfqpoint{2.457499in}{1.481969in}}%
\pgfpathlineto{\pgfqpoint{2.460018in}{1.479046in}}%
\pgfpathlineto{\pgfqpoint{2.465056in}{1.532001in}}%
\pgfpathlineto{\pgfqpoint{2.467574in}{1.536190in}}%
\pgfpathlineto{\pgfqpoint{2.470093in}{1.547312in}}%
\pgfpathlineto{\pgfqpoint{2.472612in}{1.561906in}}%
\pgfpathlineto{\pgfqpoint{2.475131in}{1.560828in}}%
\pgfpathlineto{\pgfqpoint{2.477649in}{1.565136in}}%
\pgfpathlineto{\pgfqpoint{2.480168in}{1.571411in}}%
\pgfpathlineto{\pgfqpoint{2.482687in}{1.587034in}}%
\pgfpathlineto{\pgfqpoint{2.485206in}{1.590793in}}%
\pgfpathlineto{\pgfqpoint{2.487724in}{1.599602in}}%
\pgfpathlineto{\pgfqpoint{2.490243in}{1.578653in}}%
\pgfpathlineto{\pgfqpoint{2.492762in}{1.586162in}}%
\pgfpathlineto{\pgfqpoint{2.495281in}{1.576926in}}%
\pgfpathlineto{\pgfqpoint{2.497799in}{1.571376in}}%
\pgfpathlineto{\pgfqpoint{2.500318in}{1.595646in}}%
\pgfpathlineto{\pgfqpoint{2.502837in}{1.587930in}}%
\pgfpathlineto{\pgfqpoint{2.505356in}{1.609510in}}%
\pgfpathlineto{\pgfqpoint{2.507874in}{1.611781in}}%
\pgfpathlineto{\pgfqpoint{2.510393in}{1.638279in}}%
\pgfpathlineto{\pgfqpoint{2.512912in}{1.651096in}}%
\pgfpathlineto{\pgfqpoint{2.515431in}{1.653972in}}%
\pgfpathlineto{\pgfqpoint{2.517949in}{1.643541in}}%
\pgfpathlineto{\pgfqpoint{2.520468in}{1.624407in}}%
\pgfpathlineto{\pgfqpoint{2.522987in}{1.627714in}}%
\pgfpathlineto{\pgfqpoint{2.525506in}{1.669285in}}%
\pgfpathlineto{\pgfqpoint{2.528024in}{1.661934in}}%
\pgfpathlineto{\pgfqpoint{2.530543in}{1.680881in}}%
\pgfpathlineto{\pgfqpoint{2.533062in}{1.674444in}}%
\pgfpathlineto{\pgfqpoint{2.535581in}{1.675858in}}%
\pgfpathlineto{\pgfqpoint{2.538099in}{1.658455in}}%
\pgfpathlineto{\pgfqpoint{2.540618in}{1.673402in}}%
\pgfpathlineto{\pgfqpoint{2.543137in}{1.676834in}}%
\pgfpathlineto{\pgfqpoint{2.545656in}{1.677706in}}%
\pgfpathlineto{\pgfqpoint{2.548174in}{1.651773in}}%
\pgfpathlineto{\pgfqpoint{2.550693in}{1.645086in}}%
\pgfpathlineto{\pgfqpoint{2.553212in}{1.629836in}}%
\pgfpathlineto{\pgfqpoint{2.555731in}{1.619659in}}%
\pgfpathlineto{\pgfqpoint{2.558249in}{1.622884in}}%
\pgfpathlineto{\pgfqpoint{2.560768in}{1.629677in}}%
\pgfpathlineto{\pgfqpoint{2.563287in}{1.619571in}}%
\pgfpathlineto{\pgfqpoint{2.565806in}{1.614963in}}%
\pgfpathlineto{\pgfqpoint{2.568324in}{1.610086in}}%
\pgfpathlineto{\pgfqpoint{2.570843in}{1.612629in}}%
\pgfpathlineto{\pgfqpoint{2.573362in}{1.617932in}}%
\pgfpathlineto{\pgfqpoint{2.575881in}{1.604766in}}%
\pgfpathlineto{\pgfqpoint{2.578399in}{1.595685in}}%
\pgfpathlineto{\pgfqpoint{2.580918in}{1.589570in}}%
\pgfpathlineto{\pgfqpoint{2.583437in}{1.566306in}}%
\pgfpathlineto{\pgfqpoint{2.585956in}{1.581209in}}%
\pgfpathlineto{\pgfqpoint{2.588474in}{1.600146in}}%
\pgfpathlineto{\pgfqpoint{2.590993in}{1.615110in}}%
\pgfpathlineto{\pgfqpoint{2.593512in}{1.616751in}}%
\pgfpathlineto{\pgfqpoint{2.596031in}{1.634135in}}%
\pgfpathlineto{\pgfqpoint{2.598549in}{1.658442in}}%
\pgfpathlineto{\pgfqpoint{2.601068in}{1.662032in}}%
\pgfpathlineto{\pgfqpoint{2.603587in}{1.662044in}}%
\pgfpathlineto{\pgfqpoint{2.606106in}{1.632013in}}%
\pgfpathlineto{\pgfqpoint{2.608624in}{1.618689in}}%
\pgfpathlineto{\pgfqpoint{2.611143in}{1.623962in}}%
\pgfpathlineto{\pgfqpoint{2.613662in}{1.642967in}}%
\pgfpathlineto{\pgfqpoint{2.616181in}{1.653023in}}%
\pgfpathlineto{\pgfqpoint{2.618699in}{1.680597in}}%
\pgfpathlineto{\pgfqpoint{2.621218in}{1.687581in}}%
\pgfpathlineto{\pgfqpoint{2.623737in}{1.702651in}}%
\pgfpathlineto{\pgfqpoint{2.626256in}{1.682433in}}%
\pgfpathlineto{\pgfqpoint{2.628774in}{1.669756in}}%
\pgfpathlineto{\pgfqpoint{2.631293in}{1.672168in}}%
\pgfpathlineto{\pgfqpoint{2.633812in}{1.653913in}}%
\pgfpathlineto{\pgfqpoint{2.636331in}{1.647150in}}%
\pgfpathlineto{\pgfqpoint{2.638849in}{1.676807in}}%
\pgfpathlineto{\pgfqpoint{2.641368in}{1.692370in}}%
\pgfpathlineto{\pgfqpoint{2.643887in}{1.695881in}}%
\pgfpathlineto{\pgfqpoint{2.646406in}{1.681475in}}%
\pgfpathlineto{\pgfqpoint{2.648924in}{1.669315in}}%
\pgfpathlineto{\pgfqpoint{2.651443in}{1.663403in}}%
\pgfpathlineto{\pgfqpoint{2.653962in}{1.666006in}}%
\pgfpathlineto{\pgfqpoint{2.656481in}{1.636577in}}%
\pgfpathlineto{\pgfqpoint{2.658999in}{1.631682in}}%
\pgfpathlineto{\pgfqpoint{2.661518in}{1.639802in}}%
\pgfpathlineto{\pgfqpoint{2.664037in}{1.648462in}}%
\pgfpathlineto{\pgfqpoint{2.666556in}{1.642432in}}%
\pgfpathlineto{\pgfqpoint{2.669074in}{1.639176in}}%
\pgfpathlineto{\pgfqpoint{2.671593in}{1.611217in}}%
\pgfpathlineto{\pgfqpoint{2.674112in}{1.611847in}}%
\pgfpathlineto{\pgfqpoint{2.676631in}{1.611444in}}%
\pgfpathlineto{\pgfqpoint{2.679149in}{1.626868in}}%
\pgfpathlineto{\pgfqpoint{2.681668in}{1.617666in}}%
\pgfpathlineto{\pgfqpoint{2.684187in}{1.624583in}}%
\pgfpathlineto{\pgfqpoint{2.686706in}{1.604112in}}%
\pgfpathlineto{\pgfqpoint{2.689224in}{1.622152in}}%
\pgfpathlineto{\pgfqpoint{2.691743in}{1.632862in}}%
\pgfpathlineto{\pgfqpoint{2.694262in}{1.616995in}}%
\pgfpathlineto{\pgfqpoint{2.696781in}{1.632720in}}%
\pgfpathlineto{\pgfqpoint{2.701818in}{1.666623in}}%
\pgfpathlineto{\pgfqpoint{2.704337in}{1.677104in}}%
\pgfpathlineto{\pgfqpoint{2.706856in}{1.672360in}}%
\pgfpathlineto{\pgfqpoint{2.709374in}{1.670014in}}%
\pgfpathlineto{\pgfqpoint{2.711893in}{1.669385in}}%
\pgfpathlineto{\pgfqpoint{2.714412in}{1.694591in}}%
\pgfpathlineto{\pgfqpoint{2.716931in}{1.708617in}}%
\pgfpathlineto{\pgfqpoint{2.719449in}{1.724625in}}%
\pgfpathlineto{\pgfqpoint{2.721968in}{1.727294in}}%
\pgfpathlineto{\pgfqpoint{2.724487in}{1.700680in}}%
\pgfpathlineto{\pgfqpoint{2.727006in}{1.712152in}}%
\pgfpathlineto{\pgfqpoint{2.729524in}{1.733693in}}%
\pgfpathlineto{\pgfqpoint{2.732043in}{1.748172in}}%
\pgfpathlineto{\pgfqpoint{2.734562in}{1.756685in}}%
\pgfpathlineto{\pgfqpoint{2.737081in}{1.738551in}}%
\pgfpathlineto{\pgfqpoint{2.739599in}{1.744699in}}%
\pgfpathlineto{\pgfqpoint{2.742118in}{1.749659in}}%
\pgfpathlineto{\pgfqpoint{2.744637in}{1.735792in}}%
\pgfpathlineto{\pgfqpoint{2.747156in}{1.708655in}}%
\pgfpathlineto{\pgfqpoint{2.749674in}{1.720875in}}%
\pgfpathlineto{\pgfqpoint{2.752193in}{1.741189in}}%
\pgfpathlineto{\pgfqpoint{2.754712in}{1.712080in}}%
\pgfpathlineto{\pgfqpoint{2.757231in}{1.734088in}}%
\pgfpathlineto{\pgfqpoint{2.759749in}{1.716401in}}%
\pgfpathlineto{\pgfqpoint{2.762268in}{1.724944in}}%
\pgfpathlineto{\pgfqpoint{2.764787in}{1.760144in}}%
\pgfpathlineto{\pgfqpoint{2.767306in}{1.765638in}}%
\pgfpathlineto{\pgfqpoint{2.769824in}{1.759358in}}%
\pgfpathlineto{\pgfqpoint{2.772343in}{1.748520in}}%
\pgfpathlineto{\pgfqpoint{2.774862in}{1.756291in}}%
\pgfpathlineto{\pgfqpoint{2.777381in}{1.741746in}}%
\pgfpathlineto{\pgfqpoint{2.779899in}{1.715268in}}%
\pgfpathlineto{\pgfqpoint{2.782418in}{1.747609in}}%
\pgfpathlineto{\pgfqpoint{2.784937in}{1.753354in}}%
\pgfpathlineto{\pgfqpoint{2.787456in}{1.765203in}}%
\pgfpathlineto{\pgfqpoint{2.789974in}{1.790138in}}%
\pgfpathlineto{\pgfqpoint{2.792493in}{1.778119in}}%
\pgfpathlineto{\pgfqpoint{2.795012in}{1.778526in}}%
\pgfpathlineto{\pgfqpoint{2.797531in}{1.794929in}}%
\pgfpathlineto{\pgfqpoint{2.800049in}{1.793465in}}%
\pgfpathlineto{\pgfqpoint{2.802568in}{1.789498in}}%
\pgfpathlineto{\pgfqpoint{2.805087in}{1.794826in}}%
\pgfpathlineto{\pgfqpoint{2.807606in}{1.826249in}}%
\pgfpathlineto{\pgfqpoint{2.810124in}{1.818206in}}%
\pgfpathlineto{\pgfqpoint{2.812643in}{1.815549in}}%
\pgfpathlineto{\pgfqpoint{2.815162in}{1.812066in}}%
\pgfpathlineto{\pgfqpoint{2.817681in}{1.784974in}}%
\pgfpathlineto{\pgfqpoint{2.820199in}{1.755266in}}%
\pgfpathlineto{\pgfqpoint{2.822718in}{1.746159in}}%
\pgfpathlineto{\pgfqpoint{2.825237in}{1.773759in}}%
\pgfpathlineto{\pgfqpoint{2.827756in}{1.756689in}}%
\pgfpathlineto{\pgfqpoint{2.830274in}{1.772086in}}%
\pgfpathlineto{\pgfqpoint{2.832793in}{1.756439in}}%
\pgfpathlineto{\pgfqpoint{2.835312in}{1.722891in}}%
\pgfpathlineto{\pgfqpoint{2.837831in}{1.702724in}}%
\pgfpathlineto{\pgfqpoint{2.840349in}{1.694312in}}%
\pgfpathlineto{\pgfqpoint{2.842868in}{1.688809in}}%
\pgfpathlineto{\pgfqpoint{2.845387in}{1.696853in}}%
\pgfpathlineto{\pgfqpoint{2.847906in}{1.675885in}}%
\pgfpathlineto{\pgfqpoint{2.850424in}{1.676058in}}%
\pgfpathlineto{\pgfqpoint{2.852943in}{1.677934in}}%
\pgfpathlineto{\pgfqpoint{2.855462in}{1.662957in}}%
\pgfpathlineto{\pgfqpoint{2.857981in}{1.666638in}}%
\pgfpathlineto{\pgfqpoint{2.860499in}{1.677023in}}%
\pgfpathlineto{\pgfqpoint{2.863018in}{1.658788in}}%
\pgfpathlineto{\pgfqpoint{2.865537in}{1.687256in}}%
\pgfpathlineto{\pgfqpoint{2.868056in}{1.678186in}}%
\pgfpathlineto{\pgfqpoint{2.870574in}{1.684225in}}%
\pgfpathlineto{\pgfqpoint{2.873093in}{1.668265in}}%
\pgfpathlineto{\pgfqpoint{2.875612in}{1.663357in}}%
\pgfpathlineto{\pgfqpoint{2.878131in}{1.688317in}}%
\pgfpathlineto{\pgfqpoint{2.880649in}{1.703444in}}%
\pgfpathlineto{\pgfqpoint{2.883168in}{1.688961in}}%
\pgfpathlineto{\pgfqpoint{2.885687in}{1.702929in}}%
\pgfpathlineto{\pgfqpoint{2.888206in}{1.700538in}}%
\pgfpathlineto{\pgfqpoint{2.890724in}{1.712266in}}%
\pgfpathlineto{\pgfqpoint{2.893243in}{1.714318in}}%
\pgfpathlineto{\pgfqpoint{2.895762in}{1.747501in}}%
\pgfpathlineto{\pgfqpoint{2.898281in}{1.724209in}}%
\pgfpathlineto{\pgfqpoint{2.900799in}{1.737283in}}%
\pgfpathlineto{\pgfqpoint{2.903318in}{1.738499in}}%
\pgfpathlineto{\pgfqpoint{2.905837in}{1.760326in}}%
\pgfpathlineto{\pgfqpoint{2.908356in}{1.731999in}}%
\pgfpathlineto{\pgfqpoint{2.910874in}{1.745652in}}%
\pgfpathlineto{\pgfqpoint{2.913393in}{1.739380in}}%
\pgfpathlineto{\pgfqpoint{2.915912in}{1.733863in}}%
\pgfpathlineto{\pgfqpoint{2.918431in}{1.755252in}}%
\pgfpathlineto{\pgfqpoint{2.920949in}{1.746444in}}%
\pgfpathlineto{\pgfqpoint{2.923468in}{1.780389in}}%
\pgfpathlineto{\pgfqpoint{2.925987in}{1.798310in}}%
\pgfpathlineto{\pgfqpoint{2.928506in}{1.825318in}}%
\pgfpathlineto{\pgfqpoint{2.931024in}{1.818995in}}%
\pgfpathlineto{\pgfqpoint{2.933543in}{1.809715in}}%
\pgfpathlineto{\pgfqpoint{2.936062in}{1.831066in}}%
\pgfpathlineto{\pgfqpoint{2.938581in}{1.814104in}}%
\pgfpathlineto{\pgfqpoint{2.941099in}{1.816684in}}%
\pgfpathlineto{\pgfqpoint{2.943618in}{1.803417in}}%
\pgfpathlineto{\pgfqpoint{2.946137in}{1.793639in}}%
\pgfpathlineto{\pgfqpoint{2.948656in}{1.790535in}}%
\pgfpathlineto{\pgfqpoint{2.951174in}{1.769436in}}%
\pgfpathlineto{\pgfqpoint{2.953693in}{1.773368in}}%
\pgfpathlineto{\pgfqpoint{2.956212in}{1.791425in}}%
\pgfpathlineto{\pgfqpoint{2.958731in}{1.760180in}}%
\pgfpathlineto{\pgfqpoint{2.961249in}{1.746562in}}%
\pgfpathlineto{\pgfqpoint{2.963768in}{1.755266in}}%
\pgfpathlineto{\pgfqpoint{2.966287in}{1.744976in}}%
\pgfpathlineto{\pgfqpoint{2.968806in}{1.755879in}}%
\pgfpathlineto{\pgfqpoint{2.971324in}{1.751117in}}%
\pgfpathlineto{\pgfqpoint{2.973843in}{1.747377in}}%
\pgfpathlineto{\pgfqpoint{2.976362in}{1.761138in}}%
\pgfpathlineto{\pgfqpoint{2.978881in}{1.792126in}}%
\pgfpathlineto{\pgfqpoint{2.981399in}{1.798781in}}%
\pgfpathlineto{\pgfqpoint{2.983918in}{1.785034in}}%
\pgfpathlineto{\pgfqpoint{2.986437in}{1.763381in}}%
\pgfpathlineto{\pgfqpoint{2.988956in}{1.785436in}}%
\pgfpathlineto{\pgfqpoint{2.991474in}{1.773681in}}%
\pgfpathlineto{\pgfqpoint{2.993993in}{1.756498in}}%
\pgfpathlineto{\pgfqpoint{2.996512in}{1.758687in}}%
\pgfpathlineto{\pgfqpoint{2.999031in}{1.751848in}}%
\pgfpathlineto{\pgfqpoint{3.001549in}{1.723671in}}%
\pgfpathlineto{\pgfqpoint{3.004068in}{1.715915in}}%
\pgfpathlineto{\pgfqpoint{3.006587in}{1.715812in}}%
\pgfpathlineto{\pgfqpoint{3.009106in}{1.725794in}}%
\pgfpathlineto{\pgfqpoint{3.011624in}{1.714540in}}%
\pgfpathlineto{\pgfqpoint{3.014143in}{1.732308in}}%
\pgfpathlineto{\pgfqpoint{3.016662in}{1.740059in}}%
\pgfpathlineto{\pgfqpoint{3.019181in}{1.764046in}}%
\pgfpathlineto{\pgfqpoint{3.021699in}{1.776793in}}%
\pgfpathlineto{\pgfqpoint{3.024218in}{1.784744in}}%
\pgfpathlineto{\pgfqpoint{3.026737in}{1.787117in}}%
\pgfpathlineto{\pgfqpoint{3.029256in}{1.763687in}}%
\pgfpathlineto{\pgfqpoint{3.031774in}{1.747707in}}%
\pgfpathlineto{\pgfqpoint{3.034293in}{1.750680in}}%
\pgfpathlineto{\pgfqpoint{3.036812in}{1.723949in}}%
\pgfpathlineto{\pgfqpoint{3.039331in}{1.760856in}}%
\pgfpathlineto{\pgfqpoint{3.041849in}{1.761943in}}%
\pgfpathlineto{\pgfqpoint{3.044368in}{1.765209in}}%
\pgfpathlineto{\pgfqpoint{3.046887in}{1.754736in}}%
\pgfpathlineto{\pgfqpoint{3.049406in}{1.758446in}}%
\pgfpathlineto{\pgfqpoint{3.051924in}{1.744083in}}%
\pgfpathlineto{\pgfqpoint{3.054443in}{1.747980in}}%
\pgfpathlineto{\pgfqpoint{3.056962in}{1.730882in}}%
\pgfpathlineto{\pgfqpoint{3.059481in}{1.736209in}}%
\pgfpathlineto{\pgfqpoint{3.061999in}{1.728040in}}%
\pgfpathlineto{\pgfqpoint{3.064518in}{1.727860in}}%
\pgfpathlineto{\pgfqpoint{3.067037in}{1.765213in}}%
\pgfpathlineto{\pgfqpoint{3.069556in}{1.763827in}}%
\pgfpathlineto{\pgfqpoint{3.072074in}{1.755887in}}%
\pgfpathlineto{\pgfqpoint{3.074593in}{1.749057in}}%
\pgfpathlineto{\pgfqpoint{3.077112in}{1.721927in}}%
\pgfpathlineto{\pgfqpoint{3.079631in}{1.718696in}}%
\pgfpathlineto{\pgfqpoint{3.082149in}{1.711589in}}%
\pgfpathlineto{\pgfqpoint{3.087187in}{1.745992in}}%
\pgfpathlineto{\pgfqpoint{3.089706in}{1.738444in}}%
\pgfpathlineto{\pgfqpoint{3.094743in}{1.726907in}}%
\pgfpathlineto{\pgfqpoint{3.097262in}{1.736387in}}%
\pgfpathlineto{\pgfqpoint{3.099781in}{1.701463in}}%
\pgfpathlineto{\pgfqpoint{3.102299in}{1.688822in}}%
\pgfpathlineto{\pgfqpoint{3.104818in}{1.709976in}}%
\pgfpathlineto{\pgfqpoint{3.107337in}{1.673282in}}%
\pgfpathlineto{\pgfqpoint{3.109856in}{1.669058in}}%
\pgfpathlineto{\pgfqpoint{3.112374in}{1.683270in}}%
\pgfpathlineto{\pgfqpoint{3.114893in}{1.673810in}}%
\pgfpathlineto{\pgfqpoint{3.117412in}{1.641388in}}%
\pgfpathlineto{\pgfqpoint{3.119931in}{1.663579in}}%
\pgfpathlineto{\pgfqpoint{3.122449in}{1.655225in}}%
\pgfpathlineto{\pgfqpoint{3.124968in}{1.674581in}}%
\pgfpathlineto{\pgfqpoint{3.127487in}{1.704201in}}%
\pgfpathlineto{\pgfqpoint{3.130006in}{1.686350in}}%
\pgfpathlineto{\pgfqpoint{3.132524in}{1.690049in}}%
\pgfpathlineto{\pgfqpoint{3.135043in}{1.672307in}}%
\pgfpathlineto{\pgfqpoint{3.137562in}{1.692846in}}%
\pgfpathlineto{\pgfqpoint{3.140081in}{1.666288in}}%
\pgfpathlineto{\pgfqpoint{3.142599in}{1.667822in}}%
\pgfpathlineto{\pgfqpoint{3.145118in}{1.675555in}}%
\pgfpathlineto{\pgfqpoint{3.147637in}{1.685130in}}%
\pgfpathlineto{\pgfqpoint{3.150156in}{1.689356in}}%
\pgfpathlineto{\pgfqpoint{3.152674in}{1.737684in}}%
\pgfpathlineto{\pgfqpoint{3.155193in}{1.732419in}}%
\pgfpathlineto{\pgfqpoint{3.157712in}{1.740603in}}%
\pgfpathlineto{\pgfqpoint{3.160231in}{1.732754in}}%
\pgfpathlineto{\pgfqpoint{3.162749in}{1.713579in}}%
\pgfpathlineto{\pgfqpoint{3.165268in}{1.741461in}}%
\pgfpathlineto{\pgfqpoint{3.167787in}{1.686289in}}%
\pgfpathlineto{\pgfqpoint{3.170306in}{1.705991in}}%
\pgfpathlineto{\pgfqpoint{3.172824in}{1.710271in}}%
\pgfpathlineto{\pgfqpoint{3.175343in}{1.717581in}}%
\pgfpathlineto{\pgfqpoint{3.177862in}{1.708001in}}%
\pgfpathlineto{\pgfqpoint{3.180381in}{1.717518in}}%
\pgfpathlineto{\pgfqpoint{3.182899in}{1.733117in}}%
\pgfpathlineto{\pgfqpoint{3.185418in}{1.740831in}}%
\pgfpathlineto{\pgfqpoint{3.187937in}{1.755545in}}%
\pgfpathlineto{\pgfqpoint{3.190456in}{1.771217in}}%
\pgfpathlineto{\pgfqpoint{3.192974in}{1.811131in}}%
\pgfpathlineto{\pgfqpoint{3.195493in}{1.797469in}}%
\pgfpathlineto{\pgfqpoint{3.198012in}{1.797640in}}%
\pgfpathlineto{\pgfqpoint{3.200531in}{1.820841in}}%
\pgfpathlineto{\pgfqpoint{3.203049in}{1.832961in}}%
\pgfpathlineto{\pgfqpoint{3.205568in}{1.829952in}}%
\pgfpathlineto{\pgfqpoint{3.208087in}{1.827102in}}%
\pgfpathlineto{\pgfqpoint{3.210606in}{1.838091in}}%
\pgfpathlineto{\pgfqpoint{3.213124in}{1.837546in}}%
\pgfpathlineto{\pgfqpoint{3.215643in}{1.835471in}}%
\pgfpathlineto{\pgfqpoint{3.218162in}{1.839464in}}%
\pgfpathlineto{\pgfqpoint{3.220681in}{1.855837in}}%
\pgfpathlineto{\pgfqpoint{3.223199in}{1.847234in}}%
\pgfpathlineto{\pgfqpoint{3.225718in}{1.860070in}}%
\pgfpathlineto{\pgfqpoint{3.228237in}{1.871889in}}%
\pgfpathlineto{\pgfqpoint{3.230756in}{1.877926in}}%
\pgfpathlineto{\pgfqpoint{3.233274in}{1.898979in}}%
\pgfpathlineto{\pgfqpoint{3.235793in}{1.875447in}}%
\pgfpathlineto{\pgfqpoint{3.238312in}{1.887309in}}%
\pgfpathlineto{\pgfqpoint{3.240831in}{1.890756in}}%
\pgfpathlineto{\pgfqpoint{3.243349in}{1.888578in}}%
\pgfpathlineto{\pgfqpoint{3.245868in}{1.905044in}}%
\pgfpathlineto{\pgfqpoint{3.248387in}{1.914027in}}%
\pgfpathlineto{\pgfqpoint{3.250906in}{1.914994in}}%
\pgfpathlineto{\pgfqpoint{3.253424in}{1.903889in}}%
\pgfpathlineto{\pgfqpoint{3.255943in}{1.898767in}}%
\pgfpathlineto{\pgfqpoint{3.258462in}{1.902760in}}%
\pgfpathlineto{\pgfqpoint{3.260981in}{1.899306in}}%
\pgfpathlineto{\pgfqpoint{3.263499in}{1.893981in}}%
\pgfpathlineto{\pgfqpoint{3.266018in}{1.950545in}}%
\pgfpathlineto{\pgfqpoint{3.268537in}{1.956473in}}%
\pgfpathlineto{\pgfqpoint{3.271056in}{1.957930in}}%
\pgfpathlineto{\pgfqpoint{3.273574in}{1.949490in}}%
\pgfpathlineto{\pgfqpoint{3.276093in}{1.938821in}}%
\pgfpathlineto{\pgfqpoint{3.278612in}{1.948156in}}%
\pgfpathlineto{\pgfqpoint{3.281131in}{1.942976in}}%
\pgfpathlineto{\pgfqpoint{3.283649in}{1.965713in}}%
\pgfpathlineto{\pgfqpoint{3.286168in}{1.978235in}}%
\pgfpathlineto{\pgfqpoint{3.288687in}{1.986899in}}%
\pgfpathlineto{\pgfqpoint{3.291206in}{1.975631in}}%
\pgfpathlineto{\pgfqpoint{3.293724in}{1.983094in}}%
\pgfpathlineto{\pgfqpoint{3.296243in}{1.972144in}}%
\pgfpathlineto{\pgfqpoint{3.298762in}{1.988906in}}%
\pgfpathlineto{\pgfqpoint{3.301281in}{1.976556in}}%
\pgfpathlineto{\pgfqpoint{3.303799in}{1.962773in}}%
\pgfpathlineto{\pgfqpoint{3.306318in}{1.993531in}}%
\pgfpathlineto{\pgfqpoint{3.308837in}{2.000482in}}%
\pgfpathlineto{\pgfqpoint{3.311356in}{1.976255in}}%
\pgfpathlineto{\pgfqpoint{3.313874in}{1.984530in}}%
\pgfpathlineto{\pgfqpoint{3.316393in}{2.005779in}}%
\pgfpathlineto{\pgfqpoint{3.318912in}{1.979171in}}%
\pgfpathlineto{\pgfqpoint{3.321431in}{1.969481in}}%
\pgfpathlineto{\pgfqpoint{3.323949in}{1.967689in}}%
\pgfpathlineto{\pgfqpoint{3.326468in}{1.936998in}}%
\pgfpathlineto{\pgfqpoint{3.328987in}{1.923376in}}%
\pgfpathlineto{\pgfqpoint{3.331506in}{1.904770in}}%
\pgfpathlineto{\pgfqpoint{3.334024in}{1.875930in}}%
\pgfpathlineto{\pgfqpoint{3.336543in}{1.885830in}}%
\pgfpathlineto{\pgfqpoint{3.339062in}{1.871119in}}%
\pgfpathlineto{\pgfqpoint{3.341581in}{1.855181in}}%
\pgfpathlineto{\pgfqpoint{3.344099in}{1.862405in}}%
\pgfpathlineto{\pgfqpoint{3.346618in}{1.878565in}}%
\pgfpathlineto{\pgfqpoint{3.349137in}{1.878535in}}%
\pgfpathlineto{\pgfqpoint{3.351656in}{1.851055in}}%
\pgfpathlineto{\pgfqpoint{3.354174in}{1.843604in}}%
\pgfpathlineto{\pgfqpoint{3.356693in}{1.847438in}}%
\pgfpathlineto{\pgfqpoint{3.359212in}{1.866400in}}%
\pgfpathlineto{\pgfqpoint{3.361731in}{1.866408in}}%
\pgfpathlineto{\pgfqpoint{3.364249in}{1.894272in}}%
\pgfpathlineto{\pgfqpoint{3.366768in}{1.878689in}}%
\pgfpathlineto{\pgfqpoint{3.369287in}{1.884512in}}%
\pgfpathlineto{\pgfqpoint{3.371806in}{1.859098in}}%
\pgfpathlineto{\pgfqpoint{3.374324in}{1.858519in}}%
\pgfpathlineto{\pgfqpoint{3.376843in}{1.849482in}}%
\pgfpathlineto{\pgfqpoint{3.379362in}{1.836834in}}%
\pgfpathlineto{\pgfqpoint{3.381881in}{1.833111in}}%
\pgfpathlineto{\pgfqpoint{3.384399in}{1.845996in}}%
\pgfpathlineto{\pgfqpoint{3.386918in}{1.855787in}}%
\pgfpathlineto{\pgfqpoint{3.389437in}{1.848854in}}%
\pgfpathlineto{\pgfqpoint{3.391956in}{1.850540in}}%
\pgfpathlineto{\pgfqpoint{3.394474in}{1.855263in}}%
\pgfpathlineto{\pgfqpoint{3.396993in}{1.854467in}}%
\pgfpathlineto{\pgfqpoint{3.399512in}{1.890312in}}%
\pgfpathlineto{\pgfqpoint{3.402031in}{1.873071in}}%
\pgfpathlineto{\pgfqpoint{3.404549in}{1.850057in}}%
\pgfpathlineto{\pgfqpoint{3.407068in}{1.886022in}}%
\pgfpathlineto{\pgfqpoint{3.409587in}{1.890292in}}%
\pgfpathlineto{\pgfqpoint{3.412106in}{1.906699in}}%
\pgfpathlineto{\pgfqpoint{3.414624in}{1.927208in}}%
\pgfpathlineto{\pgfqpoint{3.417143in}{1.903030in}}%
\pgfpathlineto{\pgfqpoint{3.419662in}{1.912130in}}%
\pgfpathlineto{\pgfqpoint{3.422181in}{1.883278in}}%
\pgfpathlineto{\pgfqpoint{3.424699in}{1.916612in}}%
\pgfpathlineto{\pgfqpoint{3.427218in}{1.946218in}}%
\pgfpathlineto{\pgfqpoint{3.429737in}{1.951209in}}%
\pgfpathlineto{\pgfqpoint{3.432256in}{1.976533in}}%
\pgfpathlineto{\pgfqpoint{3.434774in}{1.976946in}}%
\pgfpathlineto{\pgfqpoint{3.437293in}{1.960495in}}%
\pgfpathlineto{\pgfqpoint{3.439812in}{1.963666in}}%
\pgfpathlineto{\pgfqpoint{3.442331in}{1.957665in}}%
\pgfpathlineto{\pgfqpoint{3.444849in}{1.941340in}}%
\pgfpathlineto{\pgfqpoint{3.447368in}{1.911462in}}%
\pgfpathlineto{\pgfqpoint{3.449887in}{1.955323in}}%
\pgfpathlineto{\pgfqpoint{3.452406in}{1.973144in}}%
\pgfpathlineto{\pgfqpoint{3.454924in}{1.981346in}}%
\pgfpathlineto{\pgfqpoint{3.457443in}{1.972755in}}%
\pgfpathlineto{\pgfqpoint{3.459962in}{1.971280in}}%
\pgfpathlineto{\pgfqpoint{3.462481in}{1.966803in}}%
\pgfpathlineto{\pgfqpoint{3.464999in}{1.960309in}}%
\pgfpathlineto{\pgfqpoint{3.467518in}{1.959042in}}%
\pgfpathlineto{\pgfqpoint{3.470037in}{1.983024in}}%
\pgfpathlineto{\pgfqpoint{3.472556in}{1.990821in}}%
\pgfpathlineto{\pgfqpoint{3.475074in}{1.993637in}}%
\pgfpathlineto{\pgfqpoint{3.477593in}{2.012875in}}%
\pgfpathlineto{\pgfqpoint{3.480112in}{2.006985in}}%
\pgfpathlineto{\pgfqpoint{3.482631in}{2.004465in}}%
\pgfpathlineto{\pgfqpoint{3.485149in}{2.011296in}}%
\pgfpathlineto{\pgfqpoint{3.487668in}{2.013930in}}%
\pgfpathlineto{\pgfqpoint{3.490187in}{1.998512in}}%
\pgfpathlineto{\pgfqpoint{3.492706in}{2.035787in}}%
\pgfpathlineto{\pgfqpoint{3.495224in}{2.034765in}}%
\pgfpathlineto{\pgfqpoint{3.497743in}{2.015282in}}%
\pgfpathlineto{\pgfqpoint{3.500262in}{2.022354in}}%
\pgfpathlineto{\pgfqpoint{3.502781in}{2.002344in}}%
\pgfpathlineto{\pgfqpoint{3.505299in}{2.000361in}}%
\pgfpathlineto{\pgfqpoint{3.507818in}{2.037810in}}%
\pgfpathlineto{\pgfqpoint{3.510337in}{2.035159in}}%
\pgfpathlineto{\pgfqpoint{3.512856in}{2.000961in}}%
\pgfpathlineto{\pgfqpoint{3.515374in}{1.969194in}}%
\pgfpathlineto{\pgfqpoint{3.517893in}{1.971055in}}%
\pgfpathlineto{\pgfqpoint{3.520412in}{1.977033in}}%
\pgfpathlineto{\pgfqpoint{3.522931in}{1.968744in}}%
\pgfpathlineto{\pgfqpoint{3.525449in}{1.957089in}}%
\pgfpathlineto{\pgfqpoint{3.527968in}{1.972692in}}%
\pgfpathlineto{\pgfqpoint{3.530487in}{1.993872in}}%
\pgfpathlineto{\pgfqpoint{3.533006in}{1.990952in}}%
\pgfpathlineto{\pgfqpoint{3.535524in}{1.977331in}}%
\pgfpathlineto{\pgfqpoint{3.538043in}{1.970769in}}%
\pgfpathlineto{\pgfqpoint{3.540562in}{1.980744in}}%
\pgfpathlineto{\pgfqpoint{3.543081in}{1.974133in}}%
\pgfpathlineto{\pgfqpoint{3.545599in}{1.969087in}}%
\pgfpathlineto{\pgfqpoint{3.548118in}{1.991028in}}%
\pgfpathlineto{\pgfqpoint{3.550637in}{2.007217in}}%
\pgfpathlineto{\pgfqpoint{3.553156in}{2.027164in}}%
\pgfpathlineto{\pgfqpoint{3.555674in}{2.018946in}}%
\pgfpathlineto{\pgfqpoint{3.558193in}{2.034354in}}%
\pgfpathlineto{\pgfqpoint{3.560712in}{2.066940in}}%
\pgfpathlineto{\pgfqpoint{3.563231in}{2.074960in}}%
\pgfpathlineto{\pgfqpoint{3.565749in}{2.088694in}}%
\pgfpathlineto{\pgfqpoint{3.568268in}{2.132547in}}%
\pgfpathlineto{\pgfqpoint{3.570787in}{2.107040in}}%
\pgfpathlineto{\pgfqpoint{3.573306in}{2.128239in}}%
\pgfpathlineto{\pgfqpoint{3.575824in}{2.136103in}}%
\pgfpathlineto{\pgfqpoint{3.578343in}{2.110501in}}%
\pgfpathlineto{\pgfqpoint{3.580862in}{2.073199in}}%
\pgfpathlineto{\pgfqpoint{3.583381in}{2.065370in}}%
\pgfpathlineto{\pgfqpoint{3.585899in}{2.100163in}}%
\pgfpathlineto{\pgfqpoint{3.588418in}{2.117790in}}%
\pgfpathlineto{\pgfqpoint{3.590937in}{2.137281in}}%
\pgfpathlineto{\pgfqpoint{3.593456in}{2.147859in}}%
\pgfpathlineto{\pgfqpoint{3.595974in}{2.145319in}}%
\pgfpathlineto{\pgfqpoint{3.598493in}{2.122878in}}%
\pgfpathlineto{\pgfqpoint{3.601012in}{2.148042in}}%
\pgfpathlineto{\pgfqpoint{3.603531in}{2.168808in}}%
\pgfpathlineto{\pgfqpoint{3.606049in}{2.167757in}}%
\pgfpathlineto{\pgfqpoint{3.608568in}{2.179465in}}%
\pgfpathlineto{\pgfqpoint{3.611087in}{2.176690in}}%
\pgfpathlineto{\pgfqpoint{3.613606in}{2.188378in}}%
\pgfpathlineto{\pgfqpoint{3.616124in}{2.159656in}}%
\pgfpathlineto{\pgfqpoint{3.618643in}{2.199141in}}%
\pgfpathlineto{\pgfqpoint{3.621162in}{2.179646in}}%
\pgfpathlineto{\pgfqpoint{3.623681in}{2.187482in}}%
\pgfpathlineto{\pgfqpoint{3.626199in}{2.191400in}}%
\pgfpathlineto{\pgfqpoint{3.628718in}{2.188068in}}%
\pgfpathlineto{\pgfqpoint{3.631237in}{2.194605in}}%
\pgfpathlineto{\pgfqpoint{3.633756in}{2.185706in}}%
\pgfpathlineto{\pgfqpoint{3.636274in}{2.179936in}}%
\pgfpathlineto{\pgfqpoint{3.638793in}{2.186880in}}%
\pgfpathlineto{\pgfqpoint{3.641312in}{2.205365in}}%
\pgfpathlineto{\pgfqpoint{3.643831in}{2.221583in}}%
\pgfpathlineto{\pgfqpoint{3.646349in}{2.263574in}}%
\pgfpathlineto{\pgfqpoint{3.648868in}{2.248235in}}%
\pgfpathlineto{\pgfqpoint{3.651387in}{2.266470in}}%
\pgfpathlineto{\pgfqpoint{3.653906in}{2.274049in}}%
\pgfpathlineto{\pgfqpoint{3.656424in}{2.293130in}}%
\pgfpathlineto{\pgfqpoint{3.658943in}{2.309920in}}%
\pgfpathlineto{\pgfqpoint{3.661462in}{2.310070in}}%
\pgfpathlineto{\pgfqpoint{3.663981in}{2.302151in}}%
\pgfpathlineto{\pgfqpoint{3.666499in}{2.284980in}}%
\pgfpathlineto{\pgfqpoint{3.669018in}{2.273353in}}%
\pgfpathlineto{\pgfqpoint{3.671537in}{2.278187in}}%
\pgfpathlineto{\pgfqpoint{3.674056in}{2.284177in}}%
\pgfpathlineto{\pgfqpoint{3.676574in}{2.264688in}}%
\pgfpathlineto{\pgfqpoint{3.679093in}{2.252855in}}%
\pgfpathlineto{\pgfqpoint{3.681612in}{2.249925in}}%
\pgfpathlineto{\pgfqpoint{3.684131in}{2.255909in}}%
\pgfpathlineto{\pgfqpoint{3.686649in}{2.261515in}}%
\pgfpathlineto{\pgfqpoint{3.689168in}{2.253137in}}%
\pgfpathlineto{\pgfqpoint{3.691687in}{2.267410in}}%
\pgfpathlineto{\pgfqpoint{3.694206in}{2.267983in}}%
\pgfpathlineto{\pgfqpoint{3.696724in}{2.238862in}}%
\pgfpathlineto{\pgfqpoint{3.701762in}{2.215330in}}%
\pgfpathlineto{\pgfqpoint{3.704281in}{2.198270in}}%
\pgfpathlineto{\pgfqpoint{3.706799in}{2.193259in}}%
\pgfpathlineto{\pgfqpoint{3.709318in}{2.132243in}}%
\pgfpathlineto{\pgfqpoint{3.711837in}{2.118624in}}%
\pgfpathlineto{\pgfqpoint{3.714356in}{2.120032in}}%
\pgfpathlineto{\pgfqpoint{3.716874in}{2.098939in}}%
\pgfpathlineto{\pgfqpoint{3.719393in}{2.110532in}}%
\pgfpathlineto{\pgfqpoint{3.721912in}{2.120064in}}%
\pgfpathlineto{\pgfqpoint{3.724431in}{2.120334in}}%
\pgfpathlineto{\pgfqpoint{3.726949in}{2.100993in}}%
\pgfpathlineto{\pgfqpoint{3.729468in}{2.128760in}}%
\pgfpathlineto{\pgfqpoint{3.731987in}{2.124315in}}%
\pgfpathlineto{\pgfqpoint{3.734506in}{2.120508in}}%
\pgfpathlineto{\pgfqpoint{3.737024in}{2.098264in}}%
\pgfpathlineto{\pgfqpoint{3.739543in}{2.114642in}}%
\pgfpathlineto{\pgfqpoint{3.742062in}{2.119801in}}%
\pgfpathlineto{\pgfqpoint{3.744581in}{2.099725in}}%
\pgfpathlineto{\pgfqpoint{3.747099in}{2.080621in}}%
\pgfpathlineto{\pgfqpoint{3.749618in}{2.076165in}}%
\pgfpathlineto{\pgfqpoint{3.752137in}{2.111127in}}%
\pgfpathlineto{\pgfqpoint{3.754656in}{2.106547in}}%
\pgfpathlineto{\pgfqpoint{3.757174in}{2.115331in}}%
\pgfpathlineto{\pgfqpoint{3.759693in}{2.112055in}}%
\pgfpathlineto{\pgfqpoint{3.762212in}{2.103829in}}%
\pgfpathlineto{\pgfqpoint{3.764731in}{2.120348in}}%
\pgfpathlineto{\pgfqpoint{3.767249in}{2.129592in}}%
\pgfpathlineto{\pgfqpoint{3.769768in}{2.153518in}}%
\pgfpathlineto{\pgfqpoint{3.772287in}{2.143814in}}%
\pgfpathlineto{\pgfqpoint{3.774806in}{2.129468in}}%
\pgfpathlineto{\pgfqpoint{3.777324in}{2.121081in}}%
\pgfpathlineto{\pgfqpoint{3.779843in}{2.128251in}}%
\pgfpathlineto{\pgfqpoint{3.782362in}{2.125115in}}%
\pgfpathlineto{\pgfqpoint{3.784881in}{2.097093in}}%
\pgfpathlineto{\pgfqpoint{3.787399in}{2.084867in}}%
\pgfpathlineto{\pgfqpoint{3.789918in}{2.099282in}}%
\pgfpathlineto{\pgfqpoint{3.792437in}{2.104296in}}%
\pgfpathlineto{\pgfqpoint{3.794956in}{2.121959in}}%
\pgfpathlineto{\pgfqpoint{3.797474in}{2.125978in}}%
\pgfpathlineto{\pgfqpoint{3.799993in}{2.125267in}}%
\pgfpathlineto{\pgfqpoint{3.802512in}{2.127705in}}%
\pgfpathlineto{\pgfqpoint{3.805031in}{2.138960in}}%
\pgfpathlineto{\pgfqpoint{3.807549in}{2.118979in}}%
\pgfpathlineto{\pgfqpoint{3.810068in}{2.108538in}}%
\pgfpathlineto{\pgfqpoint{3.812587in}{2.114003in}}%
\pgfpathlineto{\pgfqpoint{3.815106in}{2.122888in}}%
\pgfpathlineto{\pgfqpoint{3.817624in}{2.110959in}}%
\pgfpathlineto{\pgfqpoint{3.820143in}{2.134979in}}%
\pgfpathlineto{\pgfqpoint{3.822662in}{2.118134in}}%
\pgfpathlineto{\pgfqpoint{3.825181in}{2.149249in}}%
\pgfpathlineto{\pgfqpoint{3.827699in}{2.124476in}}%
\pgfpathlineto{\pgfqpoint{3.830218in}{2.119921in}}%
\pgfpathlineto{\pgfqpoint{3.832737in}{2.127361in}}%
\pgfpathlineto{\pgfqpoint{3.835256in}{2.130108in}}%
\pgfpathlineto{\pgfqpoint{3.837774in}{2.129303in}}%
\pgfpathlineto{\pgfqpoint{3.840293in}{2.145009in}}%
\pgfpathlineto{\pgfqpoint{3.842812in}{2.145984in}}%
\pgfpathlineto{\pgfqpoint{3.845331in}{2.135834in}}%
\pgfpathlineto{\pgfqpoint{3.847849in}{2.135564in}}%
\pgfpathlineto{\pgfqpoint{3.850368in}{2.148630in}}%
\pgfpathlineto{\pgfqpoint{3.852887in}{2.140216in}}%
\pgfpathlineto{\pgfqpoint{3.855406in}{2.144695in}}%
\pgfpathlineto{\pgfqpoint{3.857924in}{2.159968in}}%
\pgfpathlineto{\pgfqpoint{3.860443in}{2.143083in}}%
\pgfpathlineto{\pgfqpoint{3.862962in}{2.143527in}}%
\pgfpathlineto{\pgfqpoint{3.865481in}{2.127029in}}%
\pgfpathlineto{\pgfqpoint{3.867999in}{2.122606in}}%
\pgfpathlineto{\pgfqpoint{3.870518in}{2.121837in}}%
\pgfpathlineto{\pgfqpoint{3.873037in}{2.125493in}}%
\pgfpathlineto{\pgfqpoint{3.875556in}{2.105374in}}%
\pgfpathlineto{\pgfqpoint{3.878074in}{2.111951in}}%
\pgfpathlineto{\pgfqpoint{3.880593in}{2.128136in}}%
\pgfpathlineto{\pgfqpoint{3.883112in}{2.083209in}}%
\pgfpathlineto{\pgfqpoint{3.885631in}{2.069155in}}%
\pgfpathlineto{\pgfqpoint{3.888149in}{2.057902in}}%
\pgfpathlineto{\pgfqpoint{3.890668in}{2.060918in}}%
\pgfpathlineto{\pgfqpoint{3.893187in}{2.030806in}}%
\pgfpathlineto{\pgfqpoint{3.895706in}{2.041215in}}%
\pgfpathlineto{\pgfqpoint{3.898224in}{2.045935in}}%
\pgfpathlineto{\pgfqpoint{3.900743in}{2.018751in}}%
\pgfpathlineto{\pgfqpoint{3.903262in}{2.033716in}}%
\pgfpathlineto{\pgfqpoint{3.905781in}{2.043934in}}%
\pgfpathlineto{\pgfqpoint{3.908299in}{2.040615in}}%
\pgfpathlineto{\pgfqpoint{3.910818in}{2.055612in}}%
\pgfpathlineto{\pgfqpoint{3.913337in}{2.083321in}}%
\pgfpathlineto{\pgfqpoint{3.915856in}{2.096221in}}%
\pgfpathlineto{\pgfqpoint{3.918374in}{2.085489in}}%
\pgfpathlineto{\pgfqpoint{3.920893in}{2.077761in}}%
\pgfpathlineto{\pgfqpoint{3.923412in}{2.062041in}}%
\pgfpathlineto{\pgfqpoint{3.925931in}{2.057553in}}%
\pgfpathlineto{\pgfqpoint{3.928449in}{2.054812in}}%
\pgfpathlineto{\pgfqpoint{3.930968in}{2.065324in}}%
\pgfpathlineto{\pgfqpoint{3.933487in}{2.064843in}}%
\pgfpathlineto{\pgfqpoint{3.936006in}{2.056039in}}%
\pgfpathlineto{\pgfqpoint{3.938524in}{2.046185in}}%
\pgfpathlineto{\pgfqpoint{3.941043in}{2.025333in}}%
\pgfpathlineto{\pgfqpoint{3.943562in}{1.984360in}}%
\pgfpathlineto{\pgfqpoint{3.946081in}{1.984288in}}%
\pgfpathlineto{\pgfqpoint{3.948599in}{1.942337in}}%
\pgfpathlineto{\pgfqpoint{3.951118in}{1.930341in}}%
\pgfpathlineto{\pgfqpoint{3.953637in}{1.962441in}}%
\pgfpathlineto{\pgfqpoint{3.956156in}{1.982145in}}%
\pgfpathlineto{\pgfqpoint{3.958674in}{1.975099in}}%
\pgfpathlineto{\pgfqpoint{3.961193in}{1.967090in}}%
\pgfpathlineto{\pgfqpoint{3.963712in}{1.960800in}}%
\pgfpathlineto{\pgfqpoint{3.966231in}{1.966757in}}%
\pgfpathlineto{\pgfqpoint{3.968749in}{1.970449in}}%
\pgfpathlineto{\pgfqpoint{3.971268in}{1.984810in}}%
\pgfpathlineto{\pgfqpoint{3.973787in}{2.010877in}}%
\pgfpathlineto{\pgfqpoint{3.976306in}{1.996846in}}%
\pgfpathlineto{\pgfqpoint{3.978824in}{1.990919in}}%
\pgfpathlineto{\pgfqpoint{3.981343in}{1.983411in}}%
\pgfpathlineto{\pgfqpoint{3.983862in}{1.981570in}}%
\pgfpathlineto{\pgfqpoint{3.986381in}{1.972609in}}%
\pgfpathlineto{\pgfqpoint{3.988899in}{1.962921in}}%
\pgfpathlineto{\pgfqpoint{3.991418in}{1.970021in}}%
\pgfpathlineto{\pgfqpoint{3.993937in}{1.970500in}}%
\pgfpathlineto{\pgfqpoint{3.996456in}{1.937614in}}%
\pgfpathlineto{\pgfqpoint{3.998974in}{1.967866in}}%
\pgfpathlineto{\pgfqpoint{4.001493in}{2.015814in}}%
\pgfpathlineto{\pgfqpoint{4.004012in}{2.055750in}}%
\pgfpathlineto{\pgfqpoint{4.006531in}{2.039531in}}%
\pgfpathlineto{\pgfqpoint{4.009049in}{2.012107in}}%
\pgfpathlineto{\pgfqpoint{4.011568in}{2.015822in}}%
\pgfpathlineto{\pgfqpoint{4.014087in}{2.033096in}}%
\pgfpathlineto{\pgfqpoint{4.016606in}{2.059197in}}%
\pgfpathlineto{\pgfqpoint{4.019124in}{2.052227in}}%
\pgfpathlineto{\pgfqpoint{4.021643in}{2.053299in}}%
\pgfpathlineto{\pgfqpoint{4.024162in}{2.028194in}}%
\pgfpathlineto{\pgfqpoint{4.026681in}{2.022088in}}%
\pgfpathlineto{\pgfqpoint{4.029199in}{2.022901in}}%
\pgfpathlineto{\pgfqpoint{4.034237in}{2.061330in}}%
\pgfpathlineto{\pgfqpoint{4.036756in}{2.057460in}}%
\pgfpathlineto{\pgfqpoint{4.039274in}{2.037808in}}%
\pgfpathlineto{\pgfqpoint{4.041793in}{2.037734in}}%
\pgfpathlineto{\pgfqpoint{4.044312in}{2.038787in}}%
\pgfpathlineto{\pgfqpoint{4.046831in}{2.021169in}}%
\pgfpathlineto{\pgfqpoint{4.049349in}{2.030541in}}%
\pgfpathlineto{\pgfqpoint{4.051868in}{2.036270in}}%
\pgfpathlineto{\pgfqpoint{4.054387in}{2.033366in}}%
\pgfpathlineto{\pgfqpoint{4.056906in}{2.016045in}}%
\pgfpathlineto{\pgfqpoint{4.059424in}{1.996793in}}%
\pgfpathlineto{\pgfqpoint{4.061943in}{1.976385in}}%
\pgfpathlineto{\pgfqpoint{4.064462in}{1.978461in}}%
\pgfpathlineto{\pgfqpoint{4.066981in}{1.991662in}}%
\pgfpathlineto{\pgfqpoint{4.069499in}{1.987947in}}%
\pgfpathlineto{\pgfqpoint{4.072018in}{1.969764in}}%
\pgfpathlineto{\pgfqpoint{4.074537in}{1.992456in}}%
\pgfpathlineto{\pgfqpoint{4.077056in}{1.998010in}}%
\pgfpathlineto{\pgfqpoint{4.079574in}{2.009506in}}%
\pgfpathlineto{\pgfqpoint{4.082093in}{2.040093in}}%
\pgfpathlineto{\pgfqpoint{4.084612in}{2.067200in}}%
\pgfpathlineto{\pgfqpoint{4.087131in}{2.057836in}}%
\pgfpathlineto{\pgfqpoint{4.089649in}{2.019794in}}%
\pgfpathlineto{\pgfqpoint{4.092168in}{2.036761in}}%
\pgfpathlineto{\pgfqpoint{4.094687in}{2.031196in}}%
\pgfpathlineto{\pgfqpoint{4.097206in}{1.994743in}}%
\pgfpathlineto{\pgfqpoint{4.099724in}{1.963845in}}%
\pgfpathlineto{\pgfqpoint{4.102243in}{1.928401in}}%
\pgfpathlineto{\pgfqpoint{4.104762in}{1.916920in}}%
\pgfpathlineto{\pgfqpoint{4.107281in}{1.886961in}}%
\pgfpathlineto{\pgfqpoint{4.109799in}{1.905064in}}%
\pgfpathlineto{\pgfqpoint{4.112318in}{1.913361in}}%
\pgfpathlineto{\pgfqpoint{4.114837in}{1.920777in}}%
\pgfpathlineto{\pgfqpoint{4.117356in}{1.959058in}}%
\pgfpathlineto{\pgfqpoint{4.119874in}{1.987260in}}%
\pgfpathlineto{\pgfqpoint{4.122393in}{1.974996in}}%
\pgfpathlineto{\pgfqpoint{4.124912in}{1.981289in}}%
\pgfpathlineto{\pgfqpoint{4.127431in}{1.933603in}}%
\pgfpathlineto{\pgfqpoint{4.129949in}{1.968663in}}%
\pgfpathlineto{\pgfqpoint{4.132468in}{1.992259in}}%
\pgfpathlineto{\pgfqpoint{4.134987in}{1.989137in}}%
\pgfpathlineto{\pgfqpoint{4.137506in}{1.988202in}}%
\pgfpathlineto{\pgfqpoint{4.140024in}{1.980123in}}%
\pgfpathlineto{\pgfqpoint{4.142543in}{2.010862in}}%
\pgfpathlineto{\pgfqpoint{4.145062in}{1.996320in}}%
\pgfpathlineto{\pgfqpoint{4.147581in}{2.018848in}}%
\pgfpathlineto{\pgfqpoint{4.150099in}{2.010226in}}%
\pgfpathlineto{\pgfqpoint{4.152618in}{1.984305in}}%
\pgfpathlineto{\pgfqpoint{4.155137in}{1.987410in}}%
\pgfpathlineto{\pgfqpoint{4.157656in}{1.999837in}}%
\pgfpathlineto{\pgfqpoint{4.160174in}{2.006299in}}%
\pgfpathlineto{\pgfqpoint{4.162693in}{1.982785in}}%
\pgfpathlineto{\pgfqpoint{4.165212in}{1.986245in}}%
\pgfpathlineto{\pgfqpoint{4.167731in}{1.963493in}}%
\pgfpathlineto{\pgfqpoint{4.170249in}{1.938486in}}%
\pgfpathlineto{\pgfqpoint{4.172768in}{1.942833in}}%
\pgfpathlineto{\pgfqpoint{4.175287in}{1.936351in}}%
\pgfpathlineto{\pgfqpoint{4.177806in}{1.956152in}}%
\pgfpathlineto{\pgfqpoint{4.180324in}{1.969635in}}%
\pgfpathlineto{\pgfqpoint{4.182843in}{1.951215in}}%
\pgfpathlineto{\pgfqpoint{4.185362in}{1.954881in}}%
\pgfpathlineto{\pgfqpoint{4.187881in}{1.948551in}}%
\pgfpathlineto{\pgfqpoint{4.190399in}{1.921053in}}%
\pgfpathlineto{\pgfqpoint{4.192918in}{1.940930in}}%
\pgfpathlineto{\pgfqpoint{4.195437in}{1.957609in}}%
\pgfpathlineto{\pgfqpoint{4.197956in}{1.973293in}}%
\pgfpathlineto{\pgfqpoint{4.200474in}{1.993770in}}%
\pgfpathlineto{\pgfqpoint{4.202993in}{1.997354in}}%
\pgfpathlineto{\pgfqpoint{4.205512in}{1.983997in}}%
\pgfpathlineto{\pgfqpoint{4.208031in}{1.981227in}}%
\pgfpathlineto{\pgfqpoint{4.210549in}{1.953305in}}%
\pgfpathlineto{\pgfqpoint{4.213068in}{1.973947in}}%
\pgfpathlineto{\pgfqpoint{4.215587in}{1.976870in}}%
\pgfpathlineto{\pgfqpoint{4.218106in}{2.001191in}}%
\pgfpathlineto{\pgfqpoint{4.220624in}{2.015877in}}%
\pgfpathlineto{\pgfqpoint{4.223143in}{2.001111in}}%
\pgfpathlineto{\pgfqpoint{4.225662in}{1.990438in}}%
\pgfpathlineto{\pgfqpoint{4.228181in}{1.976448in}}%
\pgfpathlineto{\pgfqpoint{4.230699in}{1.984281in}}%
\pgfpathlineto{\pgfqpoint{4.233218in}{1.994607in}}%
\pgfpathlineto{\pgfqpoint{4.235737in}{1.959995in}}%
\pgfpathlineto{\pgfqpoint{4.238256in}{1.967071in}}%
\pgfpathlineto{\pgfqpoint{4.240774in}{1.960949in}}%
\pgfpathlineto{\pgfqpoint{4.243293in}{1.967485in}}%
\pgfpathlineto{\pgfqpoint{4.245812in}{1.924534in}}%
\pgfpathlineto{\pgfqpoint{4.248331in}{1.892899in}}%
\pgfpathlineto{\pgfqpoint{4.250849in}{1.907865in}}%
\pgfpathlineto{\pgfqpoint{4.253368in}{1.901290in}}%
\pgfpathlineto{\pgfqpoint{4.255887in}{1.896330in}}%
\pgfpathlineto{\pgfqpoint{4.258406in}{1.865657in}}%
\pgfpathlineto{\pgfqpoint{4.260924in}{1.883113in}}%
\pgfpathlineto{\pgfqpoint{4.263443in}{1.882143in}}%
\pgfpathlineto{\pgfqpoint{4.265962in}{1.878025in}}%
\pgfpathlineto{\pgfqpoint{4.268481in}{1.865404in}}%
\pgfpathlineto{\pgfqpoint{4.270999in}{1.858824in}}%
\pgfpathlineto{\pgfqpoint{4.273518in}{1.845802in}}%
\pgfpathlineto{\pgfqpoint{4.278556in}{1.897381in}}%
\pgfpathlineto{\pgfqpoint{4.281074in}{1.905885in}}%
\pgfpathlineto{\pgfqpoint{4.283593in}{1.920661in}}%
\pgfpathlineto{\pgfqpoint{4.286112in}{1.898623in}}%
\pgfpathlineto{\pgfqpoint{4.288631in}{1.898813in}}%
\pgfpathlineto{\pgfqpoint{4.291149in}{1.905628in}}%
\pgfpathlineto{\pgfqpoint{4.293668in}{1.901983in}}%
\pgfpathlineto{\pgfqpoint{4.296187in}{1.911551in}}%
\pgfpathlineto{\pgfqpoint{4.298706in}{1.903917in}}%
\pgfpathlineto{\pgfqpoint{4.301224in}{1.893359in}}%
\pgfpathlineto{\pgfqpoint{4.303743in}{1.893188in}}%
\pgfpathlineto{\pgfqpoint{4.306262in}{1.908394in}}%
\pgfpathlineto{\pgfqpoint{4.308781in}{1.914640in}}%
\pgfpathlineto{\pgfqpoint{4.311299in}{1.912305in}}%
\pgfpathlineto{\pgfqpoint{4.313818in}{1.898107in}}%
\pgfpathlineto{\pgfqpoint{4.316337in}{1.922054in}}%
\pgfpathlineto{\pgfqpoint{4.318856in}{1.898604in}}%
\pgfpathlineto{\pgfqpoint{4.321374in}{1.883397in}}%
\pgfpathlineto{\pgfqpoint{4.323893in}{1.899899in}}%
\pgfpathlineto{\pgfqpoint{4.326412in}{1.889325in}}%
\pgfpathlineto{\pgfqpoint{4.328931in}{1.891079in}}%
\pgfpathlineto{\pgfqpoint{4.331449in}{1.879111in}}%
\pgfpathlineto{\pgfqpoint{4.333968in}{1.861383in}}%
\pgfpathlineto{\pgfqpoint{4.336487in}{1.886483in}}%
\pgfpathlineto{\pgfqpoint{4.339006in}{1.884556in}}%
\pgfpathlineto{\pgfqpoint{4.341524in}{1.864755in}}%
\pgfpathlineto{\pgfqpoint{4.344043in}{1.888035in}}%
\pgfpathlineto{\pgfqpoint{4.346562in}{1.885196in}}%
\pgfpathlineto{\pgfqpoint{4.349081in}{1.907622in}}%
\pgfpathlineto{\pgfqpoint{4.351599in}{1.909859in}}%
\pgfpathlineto{\pgfqpoint{4.354118in}{1.912710in}}%
\pgfpathlineto{\pgfqpoint{4.356637in}{1.916058in}}%
\pgfpathlineto{\pgfqpoint{4.359156in}{1.914641in}}%
\pgfpathlineto{\pgfqpoint{4.361674in}{1.935656in}}%
\pgfpathlineto{\pgfqpoint{4.364193in}{1.954421in}}%
\pgfpathlineto{\pgfqpoint{4.366712in}{1.969157in}}%
\pgfpathlineto{\pgfqpoint{4.369231in}{1.958968in}}%
\pgfpathlineto{\pgfqpoint{4.371749in}{1.955034in}}%
\pgfpathlineto{\pgfqpoint{4.374268in}{1.973982in}}%
\pgfpathlineto{\pgfqpoint{4.376787in}{1.939071in}}%
\pgfpathlineto{\pgfqpoint{4.379306in}{1.936664in}}%
\pgfpathlineto{\pgfqpoint{4.381824in}{1.906358in}}%
\pgfpathlineto{\pgfqpoint{4.384343in}{1.914188in}}%
\pgfpathlineto{\pgfqpoint{4.386862in}{1.921329in}}%
\pgfpathlineto{\pgfqpoint{4.389381in}{1.941438in}}%
\pgfpathlineto{\pgfqpoint{4.391899in}{1.953948in}}%
\pgfpathlineto{\pgfqpoint{4.394418in}{1.927637in}}%
\pgfpathlineto{\pgfqpoint{4.396937in}{1.954312in}}%
\pgfpathlineto{\pgfqpoint{4.399456in}{1.968618in}}%
\pgfpathlineto{\pgfqpoint{4.401974in}{1.974004in}}%
\pgfpathlineto{\pgfqpoint{4.404493in}{1.981802in}}%
\pgfpathlineto{\pgfqpoint{4.407012in}{2.006984in}}%
\pgfpathlineto{\pgfqpoint{4.409531in}{1.994637in}}%
\pgfpathlineto{\pgfqpoint{4.412049in}{1.996778in}}%
\pgfpathlineto{\pgfqpoint{4.414568in}{1.960404in}}%
\pgfpathlineto{\pgfqpoint{4.417087in}{1.963970in}}%
\pgfpathlineto{\pgfqpoint{4.422124in}{1.933107in}}%
\pgfpathlineto{\pgfqpoint{4.424643in}{1.927485in}}%
\pgfpathlineto{\pgfqpoint{4.427162in}{1.912955in}}%
\pgfpathlineto{\pgfqpoint{4.429681in}{1.932692in}}%
\pgfpathlineto{\pgfqpoint{4.432199in}{1.925449in}}%
\pgfpathlineto{\pgfqpoint{4.434718in}{1.930691in}}%
\pgfpathlineto{\pgfqpoint{4.437237in}{1.920205in}}%
\pgfpathlineto{\pgfqpoint{4.439756in}{1.929463in}}%
\pgfpathlineto{\pgfqpoint{4.442274in}{1.914690in}}%
\pgfpathlineto{\pgfqpoint{4.444793in}{1.918873in}}%
\pgfpathlineto{\pgfqpoint{4.447312in}{1.925650in}}%
\pgfpathlineto{\pgfqpoint{4.449831in}{1.946978in}}%
\pgfpathlineto{\pgfqpoint{4.452349in}{1.944882in}}%
\pgfpathlineto{\pgfqpoint{4.454868in}{1.959081in}}%
\pgfpathlineto{\pgfqpoint{4.457387in}{1.978531in}}%
\pgfpathlineto{\pgfqpoint{4.459906in}{2.015835in}}%
\pgfpathlineto{\pgfqpoint{4.462424in}{2.005365in}}%
\pgfpathlineto{\pgfqpoint{4.464943in}{2.004499in}}%
\pgfpathlineto{\pgfqpoint{4.467462in}{1.981755in}}%
\pgfpathlineto{\pgfqpoint{4.469981in}{1.982010in}}%
\pgfpathlineto{\pgfqpoint{4.472499in}{1.981347in}}%
\pgfpathlineto{\pgfqpoint{4.475018in}{1.997505in}}%
\pgfpathlineto{\pgfqpoint{4.477537in}{2.009951in}}%
\pgfpathlineto{\pgfqpoint{4.480056in}{1.987327in}}%
\pgfpathlineto{\pgfqpoint{4.482574in}{1.988683in}}%
\pgfpathlineto{\pgfqpoint{4.485093in}{1.999166in}}%
\pgfpathlineto{\pgfqpoint{4.487612in}{1.982663in}}%
\pgfpathlineto{\pgfqpoint{4.490131in}{1.964994in}}%
\pgfpathlineto{\pgfqpoint{4.492649in}{2.011152in}}%
\pgfpathlineto{\pgfqpoint{4.495168in}{2.015578in}}%
\pgfpathlineto{\pgfqpoint{4.497687in}{2.026142in}}%
\pgfpathlineto{\pgfqpoint{4.500206in}{2.008414in}}%
\pgfpathlineto{\pgfqpoint{4.502724in}{2.034902in}}%
\pgfpathlineto{\pgfqpoint{4.505243in}{2.011445in}}%
\pgfpathlineto{\pgfqpoint{4.507762in}{2.009089in}}%
\pgfpathlineto{\pgfqpoint{4.510281in}{2.005229in}}%
\pgfpathlineto{\pgfqpoint{4.512799in}{2.002858in}}%
\pgfpathlineto{\pgfqpoint{4.515318in}{1.997752in}}%
\pgfpathlineto{\pgfqpoint{4.517837in}{2.044788in}}%
\pgfpathlineto{\pgfqpoint{4.520356in}{2.017392in}}%
\pgfpathlineto{\pgfqpoint{4.522874in}{2.023363in}}%
\pgfpathlineto{\pgfqpoint{4.525393in}{2.032591in}}%
\pgfpathlineto{\pgfqpoint{4.527912in}{2.036564in}}%
\pgfpathlineto{\pgfqpoint{4.530431in}{2.039464in}}%
\pgfpathlineto{\pgfqpoint{4.532949in}{2.055849in}}%
\pgfpathlineto{\pgfqpoint{4.535468in}{2.062457in}}%
\pgfpathlineto{\pgfqpoint{4.537987in}{2.069437in}}%
\pgfpathlineto{\pgfqpoint{4.540506in}{2.021206in}}%
\pgfpathlineto{\pgfqpoint{4.543024in}{2.019747in}}%
\pgfpathlineto{\pgfqpoint{4.545543in}{2.016235in}}%
\pgfpathlineto{\pgfqpoint{4.548062in}{2.020270in}}%
\pgfpathlineto{\pgfqpoint{4.550581in}{2.036387in}}%
\pgfpathlineto{\pgfqpoint{4.553099in}{2.022211in}}%
\pgfpathlineto{\pgfqpoint{4.555618in}{2.013880in}}%
\pgfpathlineto{\pgfqpoint{4.558137in}{2.001965in}}%
\pgfpathlineto{\pgfqpoint{4.560656in}{2.013846in}}%
\pgfpathlineto{\pgfqpoint{4.563174in}{2.010830in}}%
\pgfpathlineto{\pgfqpoint{4.565693in}{2.010885in}}%
\pgfpathlineto{\pgfqpoint{4.568212in}{2.001962in}}%
\pgfpathlineto{\pgfqpoint{4.570731in}{2.013845in}}%
\pgfpathlineto{\pgfqpoint{4.573249in}{2.040172in}}%
\pgfpathlineto{\pgfqpoint{4.575768in}{2.029555in}}%
\pgfpathlineto{\pgfqpoint{4.578287in}{2.037886in}}%
\pgfpathlineto{\pgfqpoint{4.580806in}{2.039551in}}%
\pgfpathlineto{\pgfqpoint{4.583324in}{2.022898in}}%
\pgfpathlineto{\pgfqpoint{4.585843in}{1.995429in}}%
\pgfpathlineto{\pgfqpoint{4.588362in}{1.979354in}}%
\pgfpathlineto{\pgfqpoint{4.590881in}{1.990276in}}%
\pgfpathlineto{\pgfqpoint{4.593399in}{1.970026in}}%
\pgfpathlineto{\pgfqpoint{4.595918in}{1.974534in}}%
\pgfpathlineto{\pgfqpoint{4.598437in}{2.008553in}}%
\pgfpathlineto{\pgfqpoint{4.600956in}{2.009584in}}%
\pgfpathlineto{\pgfqpoint{4.603474in}{2.009238in}}%
\pgfpathlineto{\pgfqpoint{4.605993in}{1.994355in}}%
\pgfpathlineto{\pgfqpoint{4.608512in}{1.990899in}}%
\pgfpathlineto{\pgfqpoint{4.611031in}{2.012173in}}%
\pgfpathlineto{\pgfqpoint{4.613549in}{2.008527in}}%
\pgfpathlineto{\pgfqpoint{4.616068in}{1.996228in}}%
\pgfpathlineto{\pgfqpoint{4.621106in}{2.016588in}}%
\pgfpathlineto{\pgfqpoint{4.623624in}{2.024606in}}%
\pgfpathlineto{\pgfqpoint{4.626143in}{2.020828in}}%
\pgfpathlineto{\pgfqpoint{4.628662in}{1.995819in}}%
\pgfpathlineto{\pgfqpoint{4.633699in}{1.959511in}}%
\pgfpathlineto{\pgfqpoint{4.636218in}{1.943780in}}%
\pgfpathlineto{\pgfqpoint{4.638737in}{1.967846in}}%
\pgfpathlineto{\pgfqpoint{4.641256in}{1.947770in}}%
\pgfpathlineto{\pgfqpoint{4.643774in}{1.939607in}}%
\pgfpathlineto{\pgfqpoint{4.646293in}{1.956705in}}%
\pgfpathlineto{\pgfqpoint{4.648812in}{1.975230in}}%
\pgfpathlineto{\pgfqpoint{4.651331in}{1.977174in}}%
\pgfpathlineto{\pgfqpoint{4.653849in}{1.990647in}}%
\pgfpathlineto{\pgfqpoint{4.656368in}{1.943700in}}%
\pgfpathlineto{\pgfqpoint{4.658887in}{1.961507in}}%
\pgfpathlineto{\pgfqpoint{4.661406in}{1.946487in}}%
\pgfpathlineto{\pgfqpoint{4.663924in}{1.929702in}}%
\pgfpathlineto{\pgfqpoint{4.666443in}{1.963944in}}%
\pgfpathlineto{\pgfqpoint{4.668962in}{1.985018in}}%
\pgfpathlineto{\pgfqpoint{4.671481in}{1.972442in}}%
\pgfpathlineto{\pgfqpoint{4.673999in}{2.005893in}}%
\pgfpathlineto{\pgfqpoint{4.676518in}{2.008598in}}%
\pgfpathlineto{\pgfqpoint{4.679037in}{2.040576in}}%
\pgfpathlineto{\pgfqpoint{4.681556in}{2.023786in}}%
\pgfpathlineto{\pgfqpoint{4.684074in}{2.029663in}}%
\pgfpathlineto{\pgfqpoint{4.686593in}{2.039857in}}%
\pgfpathlineto{\pgfqpoint{4.689112in}{2.021719in}}%
\pgfpathlineto{\pgfqpoint{4.691631in}{2.032826in}}%
\pgfpathlineto{\pgfqpoint{4.694149in}{2.014431in}}%
\pgfpathlineto{\pgfqpoint{4.696668in}{1.980111in}}%
\pgfpathlineto{\pgfqpoint{4.699187in}{1.921953in}}%
\pgfpathlineto{\pgfqpoint{4.701706in}{1.923822in}}%
\pgfpathlineto{\pgfqpoint{4.704224in}{1.955845in}}%
\pgfpathlineto{\pgfqpoint{4.706743in}{2.001895in}}%
\pgfpathlineto{\pgfqpoint{4.709262in}{2.012387in}}%
\pgfpathlineto{\pgfqpoint{4.711781in}{2.000593in}}%
\pgfpathlineto{\pgfqpoint{4.714299in}{2.013333in}}%
\pgfpathlineto{\pgfqpoint{4.716818in}{2.014548in}}%
\pgfpathlineto{\pgfqpoint{4.719337in}{1.989459in}}%
\pgfpathlineto{\pgfqpoint{4.721856in}{1.987357in}}%
\pgfpathlineto{\pgfqpoint{4.726893in}{2.024925in}}%
\pgfpathlineto{\pgfqpoint{4.729412in}{2.006487in}}%
\pgfpathlineto{\pgfqpoint{4.731931in}{1.986207in}}%
\pgfpathlineto{\pgfqpoint{4.734449in}{1.959106in}}%
\pgfpathlineto{\pgfqpoint{4.736968in}{1.949408in}}%
\pgfpathlineto{\pgfqpoint{4.739487in}{1.950604in}}%
\pgfpathlineto{\pgfqpoint{4.742006in}{1.966146in}}%
\pgfpathlineto{\pgfqpoint{4.744524in}{1.980482in}}%
\pgfpathlineto{\pgfqpoint{4.747043in}{2.002832in}}%
\pgfpathlineto{\pgfqpoint{4.749562in}{1.976210in}}%
\pgfpathlineto{\pgfqpoint{4.752081in}{1.966515in}}%
\pgfpathlineto{\pgfqpoint{4.754599in}{1.962939in}}%
\pgfpathlineto{\pgfqpoint{4.757118in}{1.937748in}}%
\pgfpathlineto{\pgfqpoint{4.759637in}{1.939069in}}%
\pgfpathlineto{\pgfqpoint{4.762156in}{1.951724in}}%
\pgfpathlineto{\pgfqpoint{4.764674in}{1.958184in}}%
\pgfpathlineto{\pgfqpoint{4.767193in}{1.942814in}}%
\pgfpathlineto{\pgfqpoint{4.769712in}{1.913675in}}%
\pgfpathlineto{\pgfqpoint{4.772231in}{1.893576in}}%
\pgfpathlineto{\pgfqpoint{4.774749in}{1.868627in}}%
\pgfpathlineto{\pgfqpoint{4.777268in}{1.874962in}}%
\pgfpathlineto{\pgfqpoint{4.779787in}{1.872542in}}%
\pgfpathlineto{\pgfqpoint{4.782306in}{1.905369in}}%
\pgfpathlineto{\pgfqpoint{4.784824in}{1.908647in}}%
\pgfpathlineto{\pgfqpoint{4.787343in}{1.925274in}}%
\pgfpathlineto{\pgfqpoint{4.789862in}{1.952953in}}%
\pgfpathlineto{\pgfqpoint{4.792381in}{1.967420in}}%
\pgfpathlineto{\pgfqpoint{4.794899in}{1.966718in}}%
\pgfpathlineto{\pgfqpoint{4.797418in}{1.966196in}}%
\pgfpathlineto{\pgfqpoint{4.799937in}{1.970372in}}%
\pgfpathlineto{\pgfqpoint{4.802456in}{1.976390in}}%
\pgfpathlineto{\pgfqpoint{4.804974in}{1.986147in}}%
\pgfpathlineto{\pgfqpoint{4.807493in}{1.981008in}}%
\pgfpathlineto{\pgfqpoint{4.810012in}{1.991425in}}%
\pgfpathlineto{\pgfqpoint{4.812531in}{1.994484in}}%
\pgfpathlineto{\pgfqpoint{4.815049in}{1.999781in}}%
\pgfpathlineto{\pgfqpoint{4.817568in}{2.008331in}}%
\pgfpathlineto{\pgfqpoint{4.820087in}{1.975781in}}%
\pgfpathlineto{\pgfqpoint{4.822606in}{1.956226in}}%
\pgfpathlineto{\pgfqpoint{4.825124in}{1.958928in}}%
\pgfpathlineto{\pgfqpoint{4.827643in}{1.953529in}}%
\pgfpathlineto{\pgfqpoint{4.830162in}{1.986002in}}%
\pgfpathlineto{\pgfqpoint{4.832681in}{1.981542in}}%
\pgfpathlineto{\pgfqpoint{4.835199in}{1.988944in}}%
\pgfpathlineto{\pgfqpoint{4.837718in}{1.997128in}}%
\pgfpathlineto{\pgfqpoint{4.840237in}{1.987250in}}%
\pgfpathlineto{\pgfqpoint{4.842756in}{1.979762in}}%
\pgfpathlineto{\pgfqpoint{4.845274in}{1.977060in}}%
\pgfpathlineto{\pgfqpoint{4.847793in}{1.962226in}}%
\pgfpathlineto{\pgfqpoint{4.850312in}{1.973293in}}%
\pgfpathlineto{\pgfqpoint{4.852831in}{1.987485in}}%
\pgfpathlineto{\pgfqpoint{4.855349in}{2.008882in}}%
\pgfpathlineto{\pgfqpoint{4.857868in}{1.990962in}}%
\pgfpathlineto{\pgfqpoint{4.860387in}{1.979121in}}%
\pgfpathlineto{\pgfqpoint{4.862906in}{1.970239in}}%
\pgfpathlineto{\pgfqpoint{4.865424in}{1.973921in}}%
\pgfpathlineto{\pgfqpoint{4.867943in}{1.969426in}}%
\pgfpathlineto{\pgfqpoint{4.870462in}{1.958680in}}%
\pgfpathlineto{\pgfqpoint{4.872981in}{1.965475in}}%
\pgfpathlineto{\pgfqpoint{4.875499in}{1.927137in}}%
\pgfpathlineto{\pgfqpoint{4.878018in}{1.909731in}}%
\pgfpathlineto{\pgfqpoint{4.880537in}{1.904558in}}%
\pgfpathlineto{\pgfqpoint{4.883056in}{1.919367in}}%
\pgfpathlineto{\pgfqpoint{4.885574in}{1.898792in}}%
\pgfpathlineto{\pgfqpoint{4.888093in}{1.912257in}}%
\pgfpathlineto{\pgfqpoint{4.890612in}{1.901831in}}%
\pgfpathlineto{\pgfqpoint{4.893131in}{1.915920in}}%
\pgfpathlineto{\pgfqpoint{4.895649in}{1.932576in}}%
\pgfpathlineto{\pgfqpoint{4.898168in}{1.921363in}}%
\pgfpathlineto{\pgfqpoint{4.900687in}{1.915288in}}%
\pgfpathlineto{\pgfqpoint{4.903206in}{1.902293in}}%
\pgfpathlineto{\pgfqpoint{4.905724in}{1.915427in}}%
\pgfpathlineto{\pgfqpoint{4.908243in}{1.910649in}}%
\pgfpathlineto{\pgfqpoint{4.910762in}{1.932475in}}%
\pgfpathlineto{\pgfqpoint{4.913281in}{1.931392in}}%
\pgfpathlineto{\pgfqpoint{4.915799in}{1.979508in}}%
\pgfpathlineto{\pgfqpoint{4.918318in}{1.948474in}}%
\pgfpathlineto{\pgfqpoint{4.920837in}{1.969258in}}%
\pgfpathlineto{\pgfqpoint{4.923356in}{1.967742in}}%
\pgfpathlineto{\pgfqpoint{4.925874in}{1.977153in}}%
\pgfpathlineto{\pgfqpoint{4.928393in}{1.953624in}}%
\pgfpathlineto{\pgfqpoint{4.930912in}{1.945902in}}%
\pgfpathlineto{\pgfqpoint{4.933431in}{1.913562in}}%
\pgfpathlineto{\pgfqpoint{4.935949in}{1.917828in}}%
\pgfpathlineto{\pgfqpoint{4.938468in}{1.938141in}}%
\pgfpathlineto{\pgfqpoint{4.940987in}{1.969681in}}%
\pgfpathlineto{\pgfqpoint{4.943506in}{1.970581in}}%
\pgfpathlineto{\pgfqpoint{4.946024in}{1.943025in}}%
\pgfpathlineto{\pgfqpoint{4.948543in}{1.912271in}}%
\pgfpathlineto{\pgfqpoint{4.951062in}{1.859896in}}%
\pgfpathlineto{\pgfqpoint{4.953581in}{1.877081in}}%
\pgfpathlineto{\pgfqpoint{4.956099in}{1.869225in}}%
\pgfpathlineto{\pgfqpoint{4.958618in}{1.873365in}}%
\pgfpathlineto{\pgfqpoint{4.961137in}{1.869728in}}%
\pgfpathlineto{\pgfqpoint{4.963656in}{1.875682in}}%
\pgfpathlineto{\pgfqpoint{4.966174in}{1.856571in}}%
\pgfpathlineto{\pgfqpoint{4.968693in}{1.828791in}}%
\pgfpathlineto{\pgfqpoint{4.971212in}{1.828671in}}%
\pgfpathlineto{\pgfqpoint{4.973731in}{1.807073in}}%
\pgfpathlineto{\pgfqpoint{4.976249in}{1.815315in}}%
\pgfpathlineto{\pgfqpoint{4.978768in}{1.826439in}}%
\pgfpathlineto{\pgfqpoint{4.981287in}{1.832111in}}%
\pgfpathlineto{\pgfqpoint{4.983806in}{1.833106in}}%
\pgfpathlineto{\pgfqpoint{4.986324in}{1.843052in}}%
\pgfpathlineto{\pgfqpoint{4.991362in}{1.856408in}}%
\pgfpathlineto{\pgfqpoint{4.993881in}{1.873806in}}%
\pgfpathlineto{\pgfqpoint{4.996399in}{1.854785in}}%
\pgfpathlineto{\pgfqpoint{4.998918in}{1.840715in}}%
\pgfpathlineto{\pgfqpoint{5.001437in}{1.838993in}}%
\pgfpathlineto{\pgfqpoint{5.003956in}{1.873362in}}%
\pgfpathlineto{\pgfqpoint{5.006474in}{1.873170in}}%
\pgfpathlineto{\pgfqpoint{5.008993in}{1.875994in}}%
\pgfpathlineto{\pgfqpoint{5.011512in}{1.874096in}}%
\pgfpathlineto{\pgfqpoint{5.014031in}{1.888867in}}%
\pgfpathlineto{\pgfqpoint{5.016549in}{1.898286in}}%
\pgfpathlineto{\pgfqpoint{5.019068in}{1.917754in}}%
\pgfpathlineto{\pgfqpoint{5.021587in}{1.900611in}}%
\pgfpathlineto{\pgfqpoint{5.024106in}{1.887377in}}%
\pgfpathlineto{\pgfqpoint{5.026624in}{1.876706in}}%
\pgfpathlineto{\pgfqpoint{5.029143in}{1.862237in}}%
\pgfpathlineto{\pgfqpoint{5.031662in}{1.841805in}}%
\pgfpathlineto{\pgfqpoint{5.034181in}{1.858155in}}%
\pgfpathlineto{\pgfqpoint{5.036699in}{1.879873in}}%
\pgfpathlineto{\pgfqpoint{5.039218in}{1.915428in}}%
\pgfpathlineto{\pgfqpoint{5.041737in}{1.916717in}}%
\pgfpathlineto{\pgfqpoint{5.044256in}{1.929126in}}%
\pgfpathlineto{\pgfqpoint{5.046774in}{1.916466in}}%
\pgfpathlineto{\pgfqpoint{5.049293in}{1.939793in}}%
\pgfpathlineto{\pgfqpoint{5.051812in}{1.926328in}}%
\pgfpathlineto{\pgfqpoint{5.054331in}{1.902120in}}%
\pgfpathlineto{\pgfqpoint{5.056849in}{1.893308in}}%
\pgfpathlineto{\pgfqpoint{5.059368in}{1.913776in}}%
\pgfpathlineto{\pgfqpoint{5.061887in}{1.915609in}}%
\pgfpathlineto{\pgfqpoint{5.064406in}{1.911252in}}%
\pgfpathlineto{\pgfqpoint{5.066924in}{1.918982in}}%
\pgfpathlineto{\pgfqpoint{5.069443in}{1.925374in}}%
\pgfpathlineto{\pgfqpoint{5.071962in}{1.909165in}}%
\pgfpathlineto{\pgfqpoint{5.074481in}{1.881938in}}%
\pgfpathlineto{\pgfqpoint{5.076999in}{1.904881in}}%
\pgfpathlineto{\pgfqpoint{5.079518in}{1.923875in}}%
\pgfpathlineto{\pgfqpoint{5.082037in}{1.925855in}}%
\pgfpathlineto{\pgfqpoint{5.084556in}{1.937631in}}%
\pgfpathlineto{\pgfqpoint{5.087074in}{1.915616in}}%
\pgfpathlineto{\pgfqpoint{5.089593in}{1.891132in}}%
\pgfpathlineto{\pgfqpoint{5.092112in}{1.897694in}}%
\pgfpathlineto{\pgfqpoint{5.094631in}{1.899047in}}%
\pgfpathlineto{\pgfqpoint{5.097149in}{1.896348in}}%
\pgfpathlineto{\pgfqpoint{5.099668in}{1.891858in}}%
\pgfpathlineto{\pgfqpoint{5.102187in}{1.888780in}}%
\pgfpathlineto{\pgfqpoint{5.104706in}{1.882833in}}%
\pgfpathlineto{\pgfqpoint{5.107224in}{1.881521in}}%
\pgfpathlineto{\pgfqpoint{5.109743in}{1.898106in}}%
\pgfpathlineto{\pgfqpoint{5.112262in}{1.874833in}}%
\pgfpathlineto{\pgfqpoint{5.114781in}{1.862800in}}%
\pgfpathlineto{\pgfqpoint{5.117299in}{1.844279in}}%
\pgfpathlineto{\pgfqpoint{5.119818in}{1.845409in}}%
\pgfpathlineto{\pgfqpoint{5.122337in}{1.846429in}}%
\pgfpathlineto{\pgfqpoint{5.124856in}{1.889506in}}%
\pgfpathlineto{\pgfqpoint{5.127374in}{1.898289in}}%
\pgfpathlineto{\pgfqpoint{5.129893in}{1.902683in}}%
\pgfpathlineto{\pgfqpoint{5.132412in}{1.907270in}}%
\pgfpathlineto{\pgfqpoint{5.134931in}{1.908185in}}%
\pgfpathlineto{\pgfqpoint{5.137449in}{1.922203in}}%
\pgfpathlineto{\pgfqpoint{5.139968in}{1.907624in}}%
\pgfpathlineto{\pgfqpoint{5.142487in}{1.934437in}}%
\pgfpathlineto{\pgfqpoint{5.145006in}{1.933975in}}%
\pgfpathlineto{\pgfqpoint{5.147524in}{1.947938in}}%
\pgfpathlineto{\pgfqpoint{5.150043in}{1.970686in}}%
\pgfpathlineto{\pgfqpoint{5.152562in}{1.981836in}}%
\pgfpathlineto{\pgfqpoint{5.155081in}{1.986957in}}%
\pgfpathlineto{\pgfqpoint{5.157599in}{2.012768in}}%
\pgfpathlineto{\pgfqpoint{5.160118in}{2.037183in}}%
\pgfpathlineto{\pgfqpoint{5.162637in}{2.074391in}}%
\pgfpathlineto{\pgfqpoint{5.165156in}{2.058713in}}%
\pgfpathlineto{\pgfqpoint{5.167674in}{2.060816in}}%
\pgfpathlineto{\pgfqpoint{5.170193in}{2.094889in}}%
\pgfpathlineto{\pgfqpoint{5.172712in}{2.082573in}}%
\pgfpathlineto{\pgfqpoint{5.175231in}{2.081198in}}%
\pgfpathlineto{\pgfqpoint{5.177749in}{2.085655in}}%
\pgfpathlineto{\pgfqpoint{5.180268in}{2.068608in}}%
\pgfpathlineto{\pgfqpoint{5.182787in}{2.095156in}}%
\pgfpathlineto{\pgfqpoint{5.185306in}{2.089804in}}%
\pgfpathlineto{\pgfqpoint{5.187824in}{2.097909in}}%
\pgfpathlineto{\pgfqpoint{5.190343in}{2.081988in}}%
\pgfpathlineto{\pgfqpoint{5.192862in}{2.084179in}}%
\pgfpathlineto{\pgfqpoint{5.195381in}{2.092700in}}%
\pgfpathlineto{\pgfqpoint{5.197899in}{2.071271in}}%
\pgfpathlineto{\pgfqpoint{5.200418in}{2.053615in}}%
\pgfpathlineto{\pgfqpoint{5.202937in}{2.070336in}}%
\pgfpathlineto{\pgfqpoint{5.205456in}{2.078886in}}%
\pgfpathlineto{\pgfqpoint{5.207974in}{2.080214in}}%
\pgfpathlineto{\pgfqpoint{5.210493in}{2.096579in}}%
\pgfpathlineto{\pgfqpoint{5.213012in}{2.098959in}}%
\pgfpathlineto{\pgfqpoint{5.215531in}{2.111820in}}%
\pgfpathlineto{\pgfqpoint{5.218049in}{2.097036in}}%
\pgfpathlineto{\pgfqpoint{5.220568in}{2.097897in}}%
\pgfpathlineto{\pgfqpoint{5.223087in}{2.091922in}}%
\pgfpathlineto{\pgfqpoint{5.225606in}{2.118737in}}%
\pgfpathlineto{\pgfqpoint{5.228124in}{2.130637in}}%
\pgfpathlineto{\pgfqpoint{5.230643in}{2.164584in}}%
\pgfpathlineto{\pgfqpoint{5.233162in}{2.141917in}}%
\pgfpathlineto{\pgfqpoint{5.235681in}{2.115224in}}%
\pgfpathlineto{\pgfqpoint{5.238199in}{2.118221in}}%
\pgfpathlineto{\pgfqpoint{5.240718in}{2.125821in}}%
\pgfpathlineto{\pgfqpoint{5.243237in}{2.114652in}}%
\pgfpathlineto{\pgfqpoint{5.245756in}{2.126294in}}%
\pgfpathlineto{\pgfqpoint{5.248274in}{2.117977in}}%
\pgfpathlineto{\pgfqpoint{5.250793in}{2.112211in}}%
\pgfpathlineto{\pgfqpoint{5.253312in}{2.135063in}}%
\pgfpathlineto{\pgfqpoint{5.255831in}{2.112505in}}%
\pgfpathlineto{\pgfqpoint{5.258349in}{2.079697in}}%
\pgfpathlineto{\pgfqpoint{5.260868in}{2.073959in}}%
\pgfpathlineto{\pgfqpoint{5.263387in}{2.046172in}}%
\pgfpathlineto{\pgfqpoint{5.265906in}{2.044617in}}%
\pgfpathlineto{\pgfqpoint{5.268424in}{2.077404in}}%
\pgfpathlineto{\pgfqpoint{5.270943in}{2.060445in}}%
\pgfpathlineto{\pgfqpoint{5.273462in}{2.068296in}}%
\pgfpathlineto{\pgfqpoint{5.275981in}{2.071952in}}%
\pgfpathlineto{\pgfqpoint{5.278499in}{2.053188in}}%
\pgfpathlineto{\pgfqpoint{5.281018in}{2.077466in}}%
\pgfpathlineto{\pgfqpoint{5.283537in}{2.063730in}}%
\pgfpathlineto{\pgfqpoint{5.286056in}{2.056045in}}%
\pgfpathlineto{\pgfqpoint{5.288574in}{2.036531in}}%
\pgfpathlineto{\pgfqpoint{5.291093in}{2.031126in}}%
\pgfpathlineto{\pgfqpoint{5.293612in}{2.018444in}}%
\pgfpathlineto{\pgfqpoint{5.296131in}{2.062192in}}%
\pgfpathlineto{\pgfqpoint{5.298649in}{2.070095in}}%
\pgfpathlineto{\pgfqpoint{5.301168in}{2.063734in}}%
\pgfpathlineto{\pgfqpoint{5.303687in}{2.079228in}}%
\pgfpathlineto{\pgfqpoint{5.306206in}{2.074987in}}%
\pgfpathlineto{\pgfqpoint{5.308724in}{2.083986in}}%
\pgfpathlineto{\pgfqpoint{5.311243in}{2.069957in}}%
\pgfpathlineto{\pgfqpoint{5.313762in}{2.058955in}}%
\pgfpathlineto{\pgfqpoint{5.316281in}{2.057612in}}%
\pgfpathlineto{\pgfqpoint{5.318799in}{2.056055in}}%
\pgfpathlineto{\pgfqpoint{5.321318in}{2.055717in}}%
\pgfpathlineto{\pgfqpoint{5.323837in}{2.037163in}}%
\pgfpathlineto{\pgfqpoint{5.326356in}{2.012021in}}%
\pgfpathlineto{\pgfqpoint{5.328874in}{1.984057in}}%
\pgfpathlineto{\pgfqpoint{5.331393in}{1.995339in}}%
\pgfpathlineto{\pgfqpoint{5.333912in}{2.001354in}}%
\pgfpathlineto{\pgfqpoint{5.336431in}{2.031891in}}%
\pgfpathlineto{\pgfqpoint{5.338949in}{2.013163in}}%
\pgfpathlineto{\pgfqpoint{5.341468in}{2.010035in}}%
\pgfpathlineto{\pgfqpoint{5.343987in}{2.006699in}}%
\pgfpathlineto{\pgfqpoint{5.346506in}{2.008686in}}%
\pgfpathlineto{\pgfqpoint{5.349024in}{2.028984in}}%
\pgfpathlineto{\pgfqpoint{5.351543in}{2.025546in}}%
\pgfpathlineto{\pgfqpoint{5.354062in}{2.028412in}}%
\pgfpathlineto{\pgfqpoint{5.356581in}{2.032450in}}%
\pgfpathlineto{\pgfqpoint{5.359099in}{2.041349in}}%
\pgfpathlineto{\pgfqpoint{5.361618in}{2.033160in}}%
\pgfpathlineto{\pgfqpoint{5.364137in}{2.039836in}}%
\pgfpathlineto{\pgfqpoint{5.366656in}{2.032349in}}%
\pgfpathlineto{\pgfqpoint{5.369174in}{2.036447in}}%
\pgfpathlineto{\pgfqpoint{5.371693in}{2.024030in}}%
\pgfpathlineto{\pgfqpoint{5.374212in}{2.052562in}}%
\pgfpathlineto{\pgfqpoint{5.376731in}{2.062224in}}%
\pgfpathlineto{\pgfqpoint{5.379249in}{2.087052in}}%
\pgfpathlineto{\pgfqpoint{5.381768in}{2.116986in}}%
\pgfpathlineto{\pgfqpoint{5.384287in}{2.113591in}}%
\pgfpathlineto{\pgfqpoint{5.386806in}{2.092116in}}%
\pgfpathlineto{\pgfqpoint{5.389324in}{2.071882in}}%
\pgfpathlineto{\pgfqpoint{5.391843in}{2.035605in}}%
\pgfpathlineto{\pgfqpoint{5.394362in}{2.020923in}}%
\pgfpathlineto{\pgfqpoint{5.396881in}{2.027913in}}%
\pgfpathlineto{\pgfqpoint{5.399399in}{2.024294in}}%
\pgfpathlineto{\pgfqpoint{5.401918in}{2.009829in}}%
\pgfpathlineto{\pgfqpoint{5.404437in}{2.019214in}}%
\pgfpathlineto{\pgfqpoint{5.406956in}{2.001681in}}%
\pgfpathlineto{\pgfqpoint{5.409474in}{1.989269in}}%
\pgfpathlineto{\pgfqpoint{5.411993in}{1.966934in}}%
\pgfpathlineto{\pgfqpoint{5.414512in}{1.989376in}}%
\pgfpathlineto{\pgfqpoint{5.417031in}{2.005008in}}%
\pgfpathlineto{\pgfqpoint{5.419549in}{1.991698in}}%
\pgfpathlineto{\pgfqpoint{5.422068in}{2.013288in}}%
\pgfpathlineto{\pgfqpoint{5.424587in}{2.019701in}}%
\pgfpathlineto{\pgfqpoint{5.427106in}{2.033036in}}%
\pgfpathlineto{\pgfqpoint{5.429624in}{2.022867in}}%
\pgfpathlineto{\pgfqpoint{5.432143in}{2.020572in}}%
\pgfpathlineto{\pgfqpoint{5.434662in}{2.033390in}}%
\pgfpathlineto{\pgfqpoint{5.437181in}{2.029373in}}%
\pgfpathlineto{\pgfqpoint{5.439699in}{2.036874in}}%
\pgfpathlineto{\pgfqpoint{5.442218in}{2.041702in}}%
\pgfpathlineto{\pgfqpoint{5.444737in}{2.026243in}}%
\pgfpathlineto{\pgfqpoint{5.447256in}{2.058874in}}%
\pgfpathlineto{\pgfqpoint{5.449774in}{2.057369in}}%
\pgfpathlineto{\pgfqpoint{5.452293in}{2.085872in}}%
\pgfpathlineto{\pgfqpoint{5.454812in}{2.072051in}}%
\pgfpathlineto{\pgfqpoint{5.457331in}{2.068263in}}%
\pgfpathlineto{\pgfqpoint{5.459849in}{2.067344in}}%
\pgfpathlineto{\pgfqpoint{5.462368in}{2.046282in}}%
\pgfpathlineto{\pgfqpoint{5.464887in}{2.024017in}}%
\pgfpathlineto{\pgfqpoint{5.467406in}{2.010685in}}%
\pgfpathlineto{\pgfqpoint{5.469924in}{1.985993in}}%
\pgfpathlineto{\pgfqpoint{5.472443in}{1.956741in}}%
\pgfpathlineto{\pgfqpoint{5.474962in}{1.982947in}}%
\pgfpathlineto{\pgfqpoint{5.477481in}{1.951084in}}%
\pgfpathlineto{\pgfqpoint{5.479999in}{1.943107in}}%
\pgfpathlineto{\pgfqpoint{5.482518in}{1.948584in}}%
\pgfpathlineto{\pgfqpoint{5.485037in}{1.952415in}}%
\pgfpathlineto{\pgfqpoint{5.487556in}{1.966038in}}%
\pgfpathlineto{\pgfqpoint{5.490074in}{1.969549in}}%
\pgfpathlineto{\pgfqpoint{5.492593in}{1.977039in}}%
\pgfpathlineto{\pgfqpoint{5.492593in}{1.977039in}}%
\pgfusepath{stroke}%
\end{pgfscope}%
\begin{pgfscope}%
\pgfpathrectangle{\pgfqpoint{0.455093in}{0.401080in}}{\pgfqpoint{5.037500in}{4.004000in}}%
\pgfusepath{clip}%
\pgfsetrectcap%
\pgfsetroundjoin%
\pgfsetlinewidth{0.501875pt}%
\definecolor{currentstroke}{rgb}{0.690196,0.188235,0.376471}%
\pgfsetstrokecolor{currentstroke}%
\pgfsetstrokeopacity{0.800000}%
\pgfsetdash{}{0pt}%
\pgfpathmoveto{\pgfqpoint{0.455093in}{1.021044in}}%
\pgfpathlineto{\pgfqpoint{0.457612in}{1.029679in}}%
\pgfpathlineto{\pgfqpoint{0.460131in}{1.015164in}}%
\pgfpathlineto{\pgfqpoint{0.462649in}{1.017856in}}%
\pgfpathlineto{\pgfqpoint{0.465168in}{1.020285in}}%
\pgfpathlineto{\pgfqpoint{0.467687in}{1.015400in}}%
\pgfpathlineto{\pgfqpoint{0.470206in}{1.012432in}}%
\pgfpathlineto{\pgfqpoint{0.472724in}{1.003410in}}%
\pgfpathlineto{\pgfqpoint{0.475243in}{0.996584in}}%
\pgfpathlineto{\pgfqpoint{0.477762in}{0.997810in}}%
\pgfpathlineto{\pgfqpoint{0.480281in}{0.990664in}}%
\pgfpathlineto{\pgfqpoint{0.482799in}{0.998388in}}%
\pgfpathlineto{\pgfqpoint{0.485318in}{0.975999in}}%
\pgfpathlineto{\pgfqpoint{0.487837in}{0.968162in}}%
\pgfpathlineto{\pgfqpoint{0.490356in}{0.954787in}}%
\pgfpathlineto{\pgfqpoint{0.492874in}{0.955257in}}%
\pgfpathlineto{\pgfqpoint{0.495393in}{0.936363in}}%
\pgfpathlineto{\pgfqpoint{0.497912in}{0.929289in}}%
\pgfpathlineto{\pgfqpoint{0.500431in}{0.920482in}}%
\pgfpathlineto{\pgfqpoint{0.502949in}{0.925836in}}%
\pgfpathlineto{\pgfqpoint{0.505468in}{0.915238in}}%
\pgfpathlineto{\pgfqpoint{0.507987in}{0.914121in}}%
\pgfpathlineto{\pgfqpoint{0.510506in}{0.908933in}}%
\pgfpathlineto{\pgfqpoint{0.513024in}{0.919763in}}%
\pgfpathlineto{\pgfqpoint{0.515543in}{0.932413in}}%
\pgfpathlineto{\pgfqpoint{0.518062in}{0.921397in}}%
\pgfpathlineto{\pgfqpoint{0.520581in}{0.932672in}}%
\pgfpathlineto{\pgfqpoint{0.523099in}{0.921939in}}%
\pgfpathlineto{\pgfqpoint{0.525618in}{0.931553in}}%
\pgfpathlineto{\pgfqpoint{0.528137in}{0.944289in}}%
\pgfpathlineto{\pgfqpoint{0.530656in}{0.948403in}}%
\pgfpathlineto{\pgfqpoint{0.533174in}{0.945547in}}%
\pgfpathlineto{\pgfqpoint{0.535693in}{0.954342in}}%
\pgfpathlineto{\pgfqpoint{0.538212in}{0.941920in}}%
\pgfpathlineto{\pgfqpoint{0.540731in}{0.942358in}}%
\pgfpathlineto{\pgfqpoint{0.543249in}{0.963708in}}%
\pgfpathlineto{\pgfqpoint{0.545768in}{0.969889in}}%
\pgfpathlineto{\pgfqpoint{0.548287in}{0.976705in}}%
\pgfpathlineto{\pgfqpoint{0.550806in}{0.970600in}}%
\pgfpathlineto{\pgfqpoint{0.553324in}{0.962956in}}%
\pgfpathlineto{\pgfqpoint{0.555843in}{0.975441in}}%
\pgfpathlineto{\pgfqpoint{0.558362in}{1.005274in}}%
\pgfpathlineto{\pgfqpoint{0.560881in}{0.983083in}}%
\pgfpathlineto{\pgfqpoint{0.563399in}{0.991666in}}%
\pgfpathlineto{\pgfqpoint{0.565918in}{0.993195in}}%
\pgfpathlineto{\pgfqpoint{0.568437in}{0.991730in}}%
\pgfpathlineto{\pgfqpoint{0.570956in}{1.007881in}}%
\pgfpathlineto{\pgfqpoint{0.573474in}{1.005808in}}%
\pgfpathlineto{\pgfqpoint{0.575993in}{1.020402in}}%
\pgfpathlineto{\pgfqpoint{0.578512in}{1.005412in}}%
\pgfpathlineto{\pgfqpoint{0.581031in}{1.002739in}}%
\pgfpathlineto{\pgfqpoint{0.583549in}{0.985020in}}%
\pgfpathlineto{\pgfqpoint{0.586068in}{1.003920in}}%
\pgfpathlineto{\pgfqpoint{0.588587in}{1.002741in}}%
\pgfpathlineto{\pgfqpoint{0.591106in}{1.007470in}}%
\pgfpathlineto{\pgfqpoint{0.593624in}{1.016177in}}%
\pgfpathlineto{\pgfqpoint{0.596143in}{1.029568in}}%
\pgfpathlineto{\pgfqpoint{0.598662in}{1.032213in}}%
\pgfpathlineto{\pgfqpoint{0.601181in}{1.023539in}}%
\pgfpathlineto{\pgfqpoint{0.603699in}{1.007807in}}%
\pgfpathlineto{\pgfqpoint{0.606218in}{1.012579in}}%
\pgfpathlineto{\pgfqpoint{0.608737in}{0.994136in}}%
\pgfpathlineto{\pgfqpoint{0.611256in}{1.027962in}}%
\pgfpathlineto{\pgfqpoint{0.613774in}{1.029038in}}%
\pgfpathlineto{\pgfqpoint{0.616293in}{1.039233in}}%
\pgfpathlineto{\pgfqpoint{0.618812in}{1.045004in}}%
\pgfpathlineto{\pgfqpoint{0.621331in}{1.077108in}}%
\pgfpathlineto{\pgfqpoint{0.623849in}{1.081592in}}%
\pgfpathlineto{\pgfqpoint{0.626368in}{1.073384in}}%
\pgfpathlineto{\pgfqpoint{0.628887in}{1.061129in}}%
\pgfpathlineto{\pgfqpoint{0.631406in}{1.038714in}}%
\pgfpathlineto{\pgfqpoint{0.633924in}{1.036843in}}%
\pgfpathlineto{\pgfqpoint{0.636443in}{1.022727in}}%
\pgfpathlineto{\pgfqpoint{0.638962in}{1.028853in}}%
\pgfpathlineto{\pgfqpoint{0.641481in}{1.023117in}}%
\pgfpathlineto{\pgfqpoint{0.643999in}{1.022520in}}%
\pgfpathlineto{\pgfqpoint{0.646518in}{1.002442in}}%
\pgfpathlineto{\pgfqpoint{0.649037in}{0.997254in}}%
\pgfpathlineto{\pgfqpoint{0.651556in}{0.991586in}}%
\pgfpathlineto{\pgfqpoint{0.654074in}{1.002169in}}%
\pgfpathlineto{\pgfqpoint{0.656593in}{0.981488in}}%
\pgfpathlineto{\pgfqpoint{0.659112in}{0.987159in}}%
\pgfpathlineto{\pgfqpoint{0.661631in}{0.978120in}}%
\pgfpathlineto{\pgfqpoint{0.664149in}{0.988318in}}%
\pgfpathlineto{\pgfqpoint{0.666668in}{0.983201in}}%
\pgfpathlineto{\pgfqpoint{0.669187in}{0.988706in}}%
\pgfpathlineto{\pgfqpoint{0.671706in}{0.980843in}}%
\pgfpathlineto{\pgfqpoint{0.674224in}{0.948943in}}%
\pgfpathlineto{\pgfqpoint{0.676743in}{0.934679in}}%
\pgfpathlineto{\pgfqpoint{0.679262in}{0.935851in}}%
\pgfpathlineto{\pgfqpoint{0.681781in}{0.929555in}}%
\pgfpathlineto{\pgfqpoint{0.684299in}{0.891862in}}%
\pgfpathlineto{\pgfqpoint{0.686818in}{0.857266in}}%
\pgfpathlineto{\pgfqpoint{0.689337in}{0.845272in}}%
\pgfpathlineto{\pgfqpoint{0.691856in}{0.875686in}}%
\pgfpathlineto{\pgfqpoint{0.694374in}{0.876390in}}%
\pgfpathlineto{\pgfqpoint{0.696893in}{0.868857in}}%
\pgfpathlineto{\pgfqpoint{0.699412in}{0.886078in}}%
\pgfpathlineto{\pgfqpoint{0.701931in}{0.900915in}}%
\pgfpathlineto{\pgfqpoint{0.704449in}{0.909775in}}%
\pgfpathlineto{\pgfqpoint{0.706968in}{0.901316in}}%
\pgfpathlineto{\pgfqpoint{0.709487in}{0.910730in}}%
\pgfpathlineto{\pgfqpoint{0.712006in}{0.911718in}}%
\pgfpathlineto{\pgfqpoint{0.714524in}{0.892206in}}%
\pgfpathlineto{\pgfqpoint{0.717043in}{0.916070in}}%
\pgfpathlineto{\pgfqpoint{0.719562in}{0.915513in}}%
\pgfpathlineto{\pgfqpoint{0.722081in}{0.926737in}}%
\pgfpathlineto{\pgfqpoint{0.724599in}{0.906659in}}%
\pgfpathlineto{\pgfqpoint{0.727118in}{0.931321in}}%
\pgfpathlineto{\pgfqpoint{0.729637in}{0.919443in}}%
\pgfpathlineto{\pgfqpoint{0.732156in}{0.925992in}}%
\pgfpathlineto{\pgfqpoint{0.734674in}{0.916836in}}%
\pgfpathlineto{\pgfqpoint{0.737193in}{0.890477in}}%
\pgfpathlineto{\pgfqpoint{0.739712in}{0.884395in}}%
\pgfpathlineto{\pgfqpoint{0.744749in}{0.941230in}}%
\pgfpathlineto{\pgfqpoint{0.747268in}{0.959898in}}%
\pgfpathlineto{\pgfqpoint{0.749787in}{0.987161in}}%
\pgfpathlineto{\pgfqpoint{0.752306in}{0.980346in}}%
\pgfpathlineto{\pgfqpoint{0.754824in}{0.980904in}}%
\pgfpathlineto{\pgfqpoint{0.757343in}{0.974786in}}%
\pgfpathlineto{\pgfqpoint{0.759862in}{0.983920in}}%
\pgfpathlineto{\pgfqpoint{0.762381in}{0.972617in}}%
\pgfpathlineto{\pgfqpoint{0.764899in}{0.968626in}}%
\pgfpathlineto{\pgfqpoint{0.767418in}{0.960999in}}%
\pgfpathlineto{\pgfqpoint{0.769937in}{0.944440in}}%
\pgfpathlineto{\pgfqpoint{0.772456in}{0.952653in}}%
\pgfpathlineto{\pgfqpoint{0.774974in}{0.950967in}}%
\pgfpathlineto{\pgfqpoint{0.777493in}{0.947803in}}%
\pgfpathlineto{\pgfqpoint{0.780012in}{0.930490in}}%
\pgfpathlineto{\pgfqpoint{0.782531in}{0.926700in}}%
\pgfpathlineto{\pgfqpoint{0.785049in}{0.944744in}}%
\pgfpathlineto{\pgfqpoint{0.787568in}{0.944938in}}%
\pgfpathlineto{\pgfqpoint{0.790087in}{0.953756in}}%
\pgfpathlineto{\pgfqpoint{0.792606in}{0.969838in}}%
\pgfpathlineto{\pgfqpoint{0.795124in}{0.976158in}}%
\pgfpathlineto{\pgfqpoint{0.797643in}{0.985279in}}%
\pgfpathlineto{\pgfqpoint{0.802681in}{0.988606in}}%
\pgfpathlineto{\pgfqpoint{0.805199in}{0.960231in}}%
\pgfpathlineto{\pgfqpoint{0.807718in}{0.962408in}}%
\pgfpathlineto{\pgfqpoint{0.810237in}{0.939114in}}%
\pgfpathlineto{\pgfqpoint{0.812756in}{0.928497in}}%
\pgfpathlineto{\pgfqpoint{0.815274in}{0.947028in}}%
\pgfpathlineto{\pgfqpoint{0.817793in}{0.947186in}}%
\pgfpathlineto{\pgfqpoint{0.820312in}{0.944499in}}%
\pgfpathlineto{\pgfqpoint{0.822831in}{0.969070in}}%
\pgfpathlineto{\pgfqpoint{0.825349in}{0.979524in}}%
\pgfpathlineto{\pgfqpoint{0.827868in}{0.967023in}}%
\pgfpathlineto{\pgfqpoint{0.830387in}{0.985367in}}%
\pgfpathlineto{\pgfqpoint{0.832906in}{0.985035in}}%
\pgfpathlineto{\pgfqpoint{0.835424in}{0.982582in}}%
\pgfpathlineto{\pgfqpoint{0.837943in}{0.992819in}}%
\pgfpathlineto{\pgfqpoint{0.840462in}{0.986035in}}%
\pgfpathlineto{\pgfqpoint{0.842981in}{0.988064in}}%
\pgfpathlineto{\pgfqpoint{0.845499in}{0.989027in}}%
\pgfpathlineto{\pgfqpoint{0.848018in}{1.006900in}}%
\pgfpathlineto{\pgfqpoint{0.850537in}{1.012952in}}%
\pgfpathlineto{\pgfqpoint{0.853056in}{1.015104in}}%
\pgfpathlineto{\pgfqpoint{0.855574in}{1.016943in}}%
\pgfpathlineto{\pgfqpoint{0.858093in}{1.044282in}}%
\pgfpathlineto{\pgfqpoint{0.860612in}{1.051707in}}%
\pgfpathlineto{\pgfqpoint{0.863131in}{1.034963in}}%
\pgfpathlineto{\pgfqpoint{0.865649in}{1.027621in}}%
\pgfpathlineto{\pgfqpoint{0.868168in}{0.989131in}}%
\pgfpathlineto{\pgfqpoint{0.870687in}{1.012970in}}%
\pgfpathlineto{\pgfqpoint{0.873206in}{1.021191in}}%
\pgfpathlineto{\pgfqpoint{0.875724in}{1.017254in}}%
\pgfpathlineto{\pgfqpoint{0.878243in}{1.019379in}}%
\pgfpathlineto{\pgfqpoint{0.880762in}{1.030313in}}%
\pgfpathlineto{\pgfqpoint{0.883281in}{1.027218in}}%
\pgfpathlineto{\pgfqpoint{0.885799in}{1.038202in}}%
\pgfpathlineto{\pgfqpoint{0.888318in}{1.023539in}}%
\pgfpathlineto{\pgfqpoint{0.890837in}{1.032562in}}%
\pgfpathlineto{\pgfqpoint{0.893356in}{1.011258in}}%
\pgfpathlineto{\pgfqpoint{0.895874in}{1.032894in}}%
\pgfpathlineto{\pgfqpoint{0.898393in}{1.021973in}}%
\pgfpathlineto{\pgfqpoint{0.900912in}{1.003585in}}%
\pgfpathlineto{\pgfqpoint{0.903431in}{1.003516in}}%
\pgfpathlineto{\pgfqpoint{0.905949in}{1.010242in}}%
\pgfpathlineto{\pgfqpoint{0.908468in}{1.010956in}}%
\pgfpathlineto{\pgfqpoint{0.910987in}{1.013478in}}%
\pgfpathlineto{\pgfqpoint{0.913506in}{1.007660in}}%
\pgfpathlineto{\pgfqpoint{0.916024in}{0.995130in}}%
\pgfpathlineto{\pgfqpoint{0.918543in}{0.984657in}}%
\pgfpathlineto{\pgfqpoint{0.921062in}{0.975430in}}%
\pgfpathlineto{\pgfqpoint{0.923581in}{1.002755in}}%
\pgfpathlineto{\pgfqpoint{0.926099in}{1.007066in}}%
\pgfpathlineto{\pgfqpoint{0.928618in}{1.023165in}}%
\pgfpathlineto{\pgfqpoint{0.931137in}{1.019029in}}%
\pgfpathlineto{\pgfqpoint{0.933656in}{1.004564in}}%
\pgfpathlineto{\pgfqpoint{0.936174in}{0.998228in}}%
\pgfpathlineto{\pgfqpoint{0.938693in}{1.019229in}}%
\pgfpathlineto{\pgfqpoint{0.941212in}{1.026012in}}%
\pgfpathlineto{\pgfqpoint{0.943731in}{1.054659in}}%
\pgfpathlineto{\pgfqpoint{0.946249in}{1.051116in}}%
\pgfpathlineto{\pgfqpoint{0.948768in}{1.056986in}}%
\pgfpathlineto{\pgfqpoint{0.951287in}{1.030593in}}%
\pgfpathlineto{\pgfqpoint{0.953806in}{1.052850in}}%
\pgfpathlineto{\pgfqpoint{0.956324in}{1.076458in}}%
\pgfpathlineto{\pgfqpoint{0.958843in}{1.088357in}}%
\pgfpathlineto{\pgfqpoint{0.961362in}{1.090371in}}%
\pgfpathlineto{\pgfqpoint{0.963881in}{1.093702in}}%
\pgfpathlineto{\pgfqpoint{0.966399in}{1.111253in}}%
\pgfpathlineto{\pgfqpoint{0.968918in}{1.090379in}}%
\pgfpathlineto{\pgfqpoint{0.971437in}{1.079897in}}%
\pgfpathlineto{\pgfqpoint{0.973956in}{1.060386in}}%
\pgfpathlineto{\pgfqpoint{0.976474in}{1.058310in}}%
\pgfpathlineto{\pgfqpoint{0.978993in}{1.050045in}}%
\pgfpathlineto{\pgfqpoint{0.981512in}{1.056512in}}%
\pgfpathlineto{\pgfqpoint{0.984031in}{1.058729in}}%
\pgfpathlineto{\pgfqpoint{0.986549in}{1.044547in}}%
\pgfpathlineto{\pgfqpoint{0.989068in}{1.058364in}}%
\pgfpathlineto{\pgfqpoint{0.991587in}{1.061985in}}%
\pgfpathlineto{\pgfqpoint{0.994106in}{1.073090in}}%
\pgfpathlineto{\pgfqpoint{0.996624in}{1.073367in}}%
\pgfpathlineto{\pgfqpoint{0.999143in}{1.073782in}}%
\pgfpathlineto{\pgfqpoint{1.001662in}{1.105704in}}%
\pgfpathlineto{\pgfqpoint{1.004181in}{1.102628in}}%
\pgfpathlineto{\pgfqpoint{1.006699in}{1.117322in}}%
\pgfpathlineto{\pgfqpoint{1.009218in}{1.105245in}}%
\pgfpathlineto{\pgfqpoint{1.011737in}{1.099745in}}%
\pgfpathlineto{\pgfqpoint{1.014256in}{1.116981in}}%
\pgfpathlineto{\pgfqpoint{1.016774in}{1.121232in}}%
\pgfpathlineto{\pgfqpoint{1.019293in}{1.118114in}}%
\pgfpathlineto{\pgfqpoint{1.021812in}{1.118800in}}%
\pgfpathlineto{\pgfqpoint{1.024331in}{1.123640in}}%
\pgfpathlineto{\pgfqpoint{1.026849in}{1.146170in}}%
\pgfpathlineto{\pgfqpoint{1.029368in}{1.150454in}}%
\pgfpathlineto{\pgfqpoint{1.031887in}{1.155838in}}%
\pgfpathlineto{\pgfqpoint{1.034406in}{1.155508in}}%
\pgfpathlineto{\pgfqpoint{1.036924in}{1.164391in}}%
\pgfpathlineto{\pgfqpoint{1.039443in}{1.163524in}}%
\pgfpathlineto{\pgfqpoint{1.041962in}{1.136416in}}%
\pgfpathlineto{\pgfqpoint{1.044481in}{1.153789in}}%
\pgfpathlineto{\pgfqpoint{1.046999in}{1.167430in}}%
\pgfpathlineto{\pgfqpoint{1.049518in}{1.151547in}}%
\pgfpathlineto{\pgfqpoint{1.052037in}{1.172183in}}%
\pgfpathlineto{\pgfqpoint{1.054556in}{1.174991in}}%
\pgfpathlineto{\pgfqpoint{1.059593in}{1.179482in}}%
\pgfpathlineto{\pgfqpoint{1.062112in}{1.187636in}}%
\pgfpathlineto{\pgfqpoint{1.064631in}{1.197759in}}%
\pgfpathlineto{\pgfqpoint{1.067149in}{1.192587in}}%
\pgfpathlineto{\pgfqpoint{1.069668in}{1.177598in}}%
\pgfpathlineto{\pgfqpoint{1.072187in}{1.157065in}}%
\pgfpathlineto{\pgfqpoint{1.074706in}{1.169522in}}%
\pgfpathlineto{\pgfqpoint{1.077224in}{1.175024in}}%
\pgfpathlineto{\pgfqpoint{1.079743in}{1.185129in}}%
\pgfpathlineto{\pgfqpoint{1.084781in}{1.155469in}}%
\pgfpathlineto{\pgfqpoint{1.087299in}{1.153842in}}%
\pgfpathlineto{\pgfqpoint{1.089818in}{1.151261in}}%
\pgfpathlineto{\pgfqpoint{1.092337in}{1.141937in}}%
\pgfpathlineto{\pgfqpoint{1.094856in}{1.116302in}}%
\pgfpathlineto{\pgfqpoint{1.097374in}{1.123425in}}%
\pgfpathlineto{\pgfqpoint{1.099893in}{1.131050in}}%
\pgfpathlineto{\pgfqpoint{1.102412in}{1.119689in}}%
\pgfpathlineto{\pgfqpoint{1.104931in}{1.090351in}}%
\pgfpathlineto{\pgfqpoint{1.107449in}{1.095529in}}%
\pgfpathlineto{\pgfqpoint{1.109968in}{1.080712in}}%
\pgfpathlineto{\pgfqpoint{1.112487in}{1.094158in}}%
\pgfpathlineto{\pgfqpoint{1.115006in}{1.110113in}}%
\pgfpathlineto{\pgfqpoint{1.117524in}{1.103797in}}%
\pgfpathlineto{\pgfqpoint{1.120043in}{1.094791in}}%
\pgfpathlineto{\pgfqpoint{1.122562in}{1.111923in}}%
\pgfpathlineto{\pgfqpoint{1.125081in}{1.099423in}}%
\pgfpathlineto{\pgfqpoint{1.127599in}{1.119006in}}%
\pgfpathlineto{\pgfqpoint{1.130118in}{1.139527in}}%
\pgfpathlineto{\pgfqpoint{1.132637in}{1.143456in}}%
\pgfpathlineto{\pgfqpoint{1.135156in}{1.154002in}}%
\pgfpathlineto{\pgfqpoint{1.137674in}{1.142483in}}%
\pgfpathlineto{\pgfqpoint{1.140193in}{1.120044in}}%
\pgfpathlineto{\pgfqpoint{1.142712in}{1.148396in}}%
\pgfpathlineto{\pgfqpoint{1.145231in}{1.156341in}}%
\pgfpathlineto{\pgfqpoint{1.147749in}{1.145808in}}%
\pgfpathlineto{\pgfqpoint{1.150268in}{1.149934in}}%
\pgfpathlineto{\pgfqpoint{1.152787in}{1.114126in}}%
\pgfpathlineto{\pgfqpoint{1.155306in}{1.099501in}}%
\pgfpathlineto{\pgfqpoint{1.157824in}{1.079465in}}%
\pgfpathlineto{\pgfqpoint{1.160343in}{1.084520in}}%
\pgfpathlineto{\pgfqpoint{1.162862in}{1.105966in}}%
\pgfpathlineto{\pgfqpoint{1.165381in}{1.108268in}}%
\pgfpathlineto{\pgfqpoint{1.167899in}{1.109718in}}%
\pgfpathlineto{\pgfqpoint{1.170418in}{1.121476in}}%
\pgfpathlineto{\pgfqpoint{1.172937in}{1.115734in}}%
\pgfpathlineto{\pgfqpoint{1.175456in}{1.124914in}}%
\pgfpathlineto{\pgfqpoint{1.177974in}{1.119640in}}%
\pgfpathlineto{\pgfqpoint{1.180493in}{1.097403in}}%
\pgfpathlineto{\pgfqpoint{1.183012in}{1.106975in}}%
\pgfpathlineto{\pgfqpoint{1.185531in}{1.107987in}}%
\pgfpathlineto{\pgfqpoint{1.188049in}{1.110466in}}%
\pgfpathlineto{\pgfqpoint{1.190568in}{1.111957in}}%
\pgfpathlineto{\pgfqpoint{1.193087in}{1.114568in}}%
\pgfpathlineto{\pgfqpoint{1.195606in}{1.100033in}}%
\pgfpathlineto{\pgfqpoint{1.198124in}{1.090042in}}%
\pgfpathlineto{\pgfqpoint{1.200643in}{1.070048in}}%
\pgfpathlineto{\pgfqpoint{1.203162in}{1.074811in}}%
\pgfpathlineto{\pgfqpoint{1.205681in}{1.055949in}}%
\pgfpathlineto{\pgfqpoint{1.208199in}{1.057351in}}%
\pgfpathlineto{\pgfqpoint{1.210718in}{1.075277in}}%
\pgfpathlineto{\pgfqpoint{1.213237in}{1.083740in}}%
\pgfpathlineto{\pgfqpoint{1.215756in}{1.100955in}}%
\pgfpathlineto{\pgfqpoint{1.218274in}{1.103291in}}%
\pgfpathlineto{\pgfqpoint{1.220793in}{1.109649in}}%
\pgfpathlineto{\pgfqpoint{1.223312in}{1.102821in}}%
\pgfpathlineto{\pgfqpoint{1.225831in}{1.095288in}}%
\pgfpathlineto{\pgfqpoint{1.228349in}{1.098054in}}%
\pgfpathlineto{\pgfqpoint{1.230868in}{1.118525in}}%
\pgfpathlineto{\pgfqpoint{1.233387in}{1.106996in}}%
\pgfpathlineto{\pgfqpoint{1.235906in}{1.104686in}}%
\pgfpathlineto{\pgfqpoint{1.238424in}{1.103848in}}%
\pgfpathlineto{\pgfqpoint{1.240943in}{1.119087in}}%
\pgfpathlineto{\pgfqpoint{1.243462in}{1.120621in}}%
\pgfpathlineto{\pgfqpoint{1.245981in}{1.093357in}}%
\pgfpathlineto{\pgfqpoint{1.248499in}{1.075728in}}%
\pgfpathlineto{\pgfqpoint{1.251018in}{1.082733in}}%
\pgfpathlineto{\pgfqpoint{1.253537in}{1.058083in}}%
\pgfpathlineto{\pgfqpoint{1.256056in}{1.061589in}}%
\pgfpathlineto{\pgfqpoint{1.258574in}{1.079862in}}%
\pgfpathlineto{\pgfqpoint{1.261093in}{1.058751in}}%
\pgfpathlineto{\pgfqpoint{1.263612in}{1.055311in}}%
\pgfpathlineto{\pgfqpoint{1.266131in}{1.047468in}}%
\pgfpathlineto{\pgfqpoint{1.268649in}{1.038104in}}%
\pgfpathlineto{\pgfqpoint{1.271168in}{1.032167in}}%
\pgfpathlineto{\pgfqpoint{1.273687in}{1.052034in}}%
\pgfpathlineto{\pgfqpoint{1.276206in}{1.031417in}}%
\pgfpathlineto{\pgfqpoint{1.278724in}{1.025323in}}%
\pgfpathlineto{\pgfqpoint{1.281243in}{1.027689in}}%
\pgfpathlineto{\pgfqpoint{1.283762in}{1.023988in}}%
\pgfpathlineto{\pgfqpoint{1.286281in}{1.032017in}}%
\pgfpathlineto{\pgfqpoint{1.288799in}{1.028711in}}%
\pgfpathlineto{\pgfqpoint{1.291318in}{1.033486in}}%
\pgfpathlineto{\pgfqpoint{1.293837in}{1.034493in}}%
\pgfpathlineto{\pgfqpoint{1.296356in}{1.033397in}}%
\pgfpathlineto{\pgfqpoint{1.298874in}{1.035669in}}%
\pgfpathlineto{\pgfqpoint{1.301393in}{1.029723in}}%
\pgfpathlineto{\pgfqpoint{1.303912in}{1.025234in}}%
\pgfpathlineto{\pgfqpoint{1.306431in}{1.004552in}}%
\pgfpathlineto{\pgfqpoint{1.308949in}{0.995579in}}%
\pgfpathlineto{\pgfqpoint{1.311468in}{0.995095in}}%
\pgfpathlineto{\pgfqpoint{1.313987in}{0.979796in}}%
\pgfpathlineto{\pgfqpoint{1.316506in}{0.980160in}}%
\pgfpathlineto{\pgfqpoint{1.319024in}{0.983156in}}%
\pgfpathlineto{\pgfqpoint{1.321543in}{0.969098in}}%
\pgfpathlineto{\pgfqpoint{1.324062in}{0.973609in}}%
\pgfpathlineto{\pgfqpoint{1.326581in}{0.996070in}}%
\pgfpathlineto{\pgfqpoint{1.329099in}{0.978663in}}%
\pgfpathlineto{\pgfqpoint{1.331618in}{0.987799in}}%
\pgfpathlineto{\pgfqpoint{1.334137in}{0.999597in}}%
\pgfpathlineto{\pgfqpoint{1.336656in}{0.993851in}}%
\pgfpathlineto{\pgfqpoint{1.339174in}{0.997488in}}%
\pgfpathlineto{\pgfqpoint{1.344212in}{1.019080in}}%
\pgfpathlineto{\pgfqpoint{1.346731in}{1.005762in}}%
\pgfpathlineto{\pgfqpoint{1.349249in}{0.999328in}}%
\pgfpathlineto{\pgfqpoint{1.351768in}{1.028976in}}%
\pgfpathlineto{\pgfqpoint{1.354287in}{1.006481in}}%
\pgfpathlineto{\pgfqpoint{1.356806in}{1.001576in}}%
\pgfpathlineto{\pgfqpoint{1.359324in}{1.009157in}}%
\pgfpathlineto{\pgfqpoint{1.361843in}{1.018166in}}%
\pgfpathlineto{\pgfqpoint{1.364362in}{1.038598in}}%
\pgfpathlineto{\pgfqpoint{1.366881in}{1.032444in}}%
\pgfpathlineto{\pgfqpoint{1.369399in}{1.056498in}}%
\pgfpathlineto{\pgfqpoint{1.371918in}{1.062988in}}%
\pgfpathlineto{\pgfqpoint{1.374437in}{1.073506in}}%
\pgfpathlineto{\pgfqpoint{1.376956in}{1.102378in}}%
\pgfpathlineto{\pgfqpoint{1.379474in}{1.111928in}}%
\pgfpathlineto{\pgfqpoint{1.381993in}{1.108193in}}%
\pgfpathlineto{\pgfqpoint{1.384512in}{1.135436in}}%
\pgfpathlineto{\pgfqpoint{1.387031in}{1.145006in}}%
\pgfpathlineto{\pgfqpoint{1.389549in}{1.142270in}}%
\pgfpathlineto{\pgfqpoint{1.392068in}{1.157411in}}%
\pgfpathlineto{\pgfqpoint{1.394587in}{1.161368in}}%
\pgfpathlineto{\pgfqpoint{1.397106in}{1.184722in}}%
\pgfpathlineto{\pgfqpoint{1.399624in}{1.197139in}}%
\pgfpathlineto{\pgfqpoint{1.402143in}{1.225085in}}%
\pgfpathlineto{\pgfqpoint{1.404662in}{1.208840in}}%
\pgfpathlineto{\pgfqpoint{1.407181in}{1.203809in}}%
\pgfpathlineto{\pgfqpoint{1.409699in}{1.195432in}}%
\pgfpathlineto{\pgfqpoint{1.412218in}{1.190398in}}%
\pgfpathlineto{\pgfqpoint{1.414737in}{1.172700in}}%
\pgfpathlineto{\pgfqpoint{1.417256in}{1.181183in}}%
\pgfpathlineto{\pgfqpoint{1.419774in}{1.183591in}}%
\pgfpathlineto{\pgfqpoint{1.422293in}{1.167802in}}%
\pgfpathlineto{\pgfqpoint{1.424812in}{1.194459in}}%
\pgfpathlineto{\pgfqpoint{1.427331in}{1.194290in}}%
\pgfpathlineto{\pgfqpoint{1.429849in}{1.198626in}}%
\pgfpathlineto{\pgfqpoint{1.432368in}{1.213320in}}%
\pgfpathlineto{\pgfqpoint{1.434887in}{1.226021in}}%
\pgfpathlineto{\pgfqpoint{1.437406in}{1.237747in}}%
\pgfpathlineto{\pgfqpoint{1.439924in}{1.229069in}}%
\pgfpathlineto{\pgfqpoint{1.442443in}{1.221532in}}%
\pgfpathlineto{\pgfqpoint{1.444962in}{1.224542in}}%
\pgfpathlineto{\pgfqpoint{1.447481in}{1.250398in}}%
\pgfpathlineto{\pgfqpoint{1.449999in}{1.238409in}}%
\pgfpathlineto{\pgfqpoint{1.452518in}{1.228467in}}%
\pgfpathlineto{\pgfqpoint{1.455037in}{1.205907in}}%
\pgfpathlineto{\pgfqpoint{1.457556in}{1.216719in}}%
\pgfpathlineto{\pgfqpoint{1.460074in}{1.223969in}}%
\pgfpathlineto{\pgfqpoint{1.462593in}{1.209183in}}%
\pgfpathlineto{\pgfqpoint{1.465112in}{1.203762in}}%
\pgfpathlineto{\pgfqpoint{1.467631in}{1.206179in}}%
\pgfpathlineto{\pgfqpoint{1.470149in}{1.199804in}}%
\pgfpathlineto{\pgfqpoint{1.472668in}{1.212756in}}%
\pgfpathlineto{\pgfqpoint{1.475187in}{1.209098in}}%
\pgfpathlineto{\pgfqpoint{1.477706in}{1.220009in}}%
\pgfpathlineto{\pgfqpoint{1.480224in}{1.240314in}}%
\pgfpathlineto{\pgfqpoint{1.482743in}{1.269273in}}%
\pgfpathlineto{\pgfqpoint{1.485262in}{1.268287in}}%
\pgfpathlineto{\pgfqpoint{1.487781in}{1.277700in}}%
\pgfpathlineto{\pgfqpoint{1.490299in}{1.282674in}}%
\pgfpathlineto{\pgfqpoint{1.492818in}{1.286863in}}%
\pgfpathlineto{\pgfqpoint{1.497856in}{1.306302in}}%
\pgfpathlineto{\pgfqpoint{1.500374in}{1.314618in}}%
\pgfpathlineto{\pgfqpoint{1.502893in}{1.300459in}}%
\pgfpathlineto{\pgfqpoint{1.505412in}{1.284977in}}%
\pgfpathlineto{\pgfqpoint{1.507931in}{1.298515in}}%
\pgfpathlineto{\pgfqpoint{1.510449in}{1.286444in}}%
\pgfpathlineto{\pgfqpoint{1.512968in}{1.293813in}}%
\pgfpathlineto{\pgfqpoint{1.515487in}{1.294648in}}%
\pgfpathlineto{\pgfqpoint{1.518006in}{1.294679in}}%
\pgfpathlineto{\pgfqpoint{1.520524in}{1.290279in}}%
\pgfpathlineto{\pgfqpoint{1.523043in}{1.311634in}}%
\pgfpathlineto{\pgfqpoint{1.525562in}{1.330669in}}%
\pgfpathlineto{\pgfqpoint{1.528081in}{1.313188in}}%
\pgfpathlineto{\pgfqpoint{1.530599in}{1.326240in}}%
\pgfpathlineto{\pgfqpoint{1.533118in}{1.348777in}}%
\pgfpathlineto{\pgfqpoint{1.535637in}{1.350083in}}%
\pgfpathlineto{\pgfqpoint{1.538156in}{1.353402in}}%
\pgfpathlineto{\pgfqpoint{1.540674in}{1.342809in}}%
\pgfpathlineto{\pgfqpoint{1.543193in}{1.336781in}}%
\pgfpathlineto{\pgfqpoint{1.545712in}{1.355911in}}%
\pgfpathlineto{\pgfqpoint{1.548231in}{1.344309in}}%
\pgfpathlineto{\pgfqpoint{1.550749in}{1.318469in}}%
\pgfpathlineto{\pgfqpoint{1.553268in}{1.341272in}}%
\pgfpathlineto{\pgfqpoint{1.555787in}{1.313800in}}%
\pgfpathlineto{\pgfqpoint{1.558306in}{1.317816in}}%
\pgfpathlineto{\pgfqpoint{1.560824in}{1.327102in}}%
\pgfpathlineto{\pgfqpoint{1.563343in}{1.317111in}}%
\pgfpathlineto{\pgfqpoint{1.565862in}{1.319885in}}%
\pgfpathlineto{\pgfqpoint{1.568381in}{1.330252in}}%
\pgfpathlineto{\pgfqpoint{1.570899in}{1.350078in}}%
\pgfpathlineto{\pgfqpoint{1.573418in}{1.353698in}}%
\pgfpathlineto{\pgfqpoint{1.575937in}{1.364564in}}%
\pgfpathlineto{\pgfqpoint{1.578456in}{1.366328in}}%
\pgfpathlineto{\pgfqpoint{1.580974in}{1.340865in}}%
\pgfpathlineto{\pgfqpoint{1.583493in}{1.329161in}}%
\pgfpathlineto{\pgfqpoint{1.586012in}{1.326269in}}%
\pgfpathlineto{\pgfqpoint{1.588531in}{1.293039in}}%
\pgfpathlineto{\pgfqpoint{1.591049in}{1.289768in}}%
\pgfpathlineto{\pgfqpoint{1.593568in}{1.283291in}}%
\pgfpathlineto{\pgfqpoint{1.596087in}{1.276011in}}%
\pgfpathlineto{\pgfqpoint{1.598606in}{1.278055in}}%
\pgfpathlineto{\pgfqpoint{1.601124in}{1.237344in}}%
\pgfpathlineto{\pgfqpoint{1.603643in}{1.236210in}}%
\pgfpathlineto{\pgfqpoint{1.606162in}{1.221093in}}%
\pgfpathlineto{\pgfqpoint{1.608681in}{1.214103in}}%
\pgfpathlineto{\pgfqpoint{1.611199in}{1.216151in}}%
\pgfpathlineto{\pgfqpoint{1.613718in}{1.209931in}}%
\pgfpathlineto{\pgfqpoint{1.616237in}{1.200793in}}%
\pgfpathlineto{\pgfqpoint{1.618756in}{1.195457in}}%
\pgfpathlineto{\pgfqpoint{1.621274in}{1.194757in}}%
\pgfpathlineto{\pgfqpoint{1.623793in}{1.195723in}}%
\pgfpathlineto{\pgfqpoint{1.626312in}{1.181107in}}%
\pgfpathlineto{\pgfqpoint{1.628831in}{1.203429in}}%
\pgfpathlineto{\pgfqpoint{1.631349in}{1.218134in}}%
\pgfpathlineto{\pgfqpoint{1.633868in}{1.222051in}}%
\pgfpathlineto{\pgfqpoint{1.636387in}{1.229855in}}%
\pgfpathlineto{\pgfqpoint{1.638906in}{1.229274in}}%
\pgfpathlineto{\pgfqpoint{1.641424in}{1.233151in}}%
\pgfpathlineto{\pgfqpoint{1.643943in}{1.238153in}}%
\pgfpathlineto{\pgfqpoint{1.646462in}{1.230467in}}%
\pgfpathlineto{\pgfqpoint{1.648981in}{1.235318in}}%
\pgfpathlineto{\pgfqpoint{1.651499in}{1.252782in}}%
\pgfpathlineto{\pgfqpoint{1.654018in}{1.245051in}}%
\pgfpathlineto{\pgfqpoint{1.656537in}{1.262534in}}%
\pgfpathlineto{\pgfqpoint{1.659056in}{1.238314in}}%
\pgfpathlineto{\pgfqpoint{1.661574in}{1.224921in}}%
\pgfpathlineto{\pgfqpoint{1.664093in}{1.238042in}}%
\pgfpathlineto{\pgfqpoint{1.666612in}{1.244550in}}%
\pgfpathlineto{\pgfqpoint{1.669131in}{1.227725in}}%
\pgfpathlineto{\pgfqpoint{1.671649in}{1.208588in}}%
\pgfpathlineto{\pgfqpoint{1.674168in}{1.214873in}}%
\pgfpathlineto{\pgfqpoint{1.676687in}{1.229479in}}%
\pgfpathlineto{\pgfqpoint{1.679206in}{1.251194in}}%
\pgfpathlineto{\pgfqpoint{1.681724in}{1.254950in}}%
\pgfpathlineto{\pgfqpoint{1.684243in}{1.243108in}}%
\pgfpathlineto{\pgfqpoint{1.686762in}{1.238318in}}%
\pgfpathlineto{\pgfqpoint{1.689281in}{1.241133in}}%
\pgfpathlineto{\pgfqpoint{1.691799in}{1.268717in}}%
\pgfpathlineto{\pgfqpoint{1.694318in}{1.261415in}}%
\pgfpathlineto{\pgfqpoint{1.696837in}{1.280965in}}%
\pgfpathlineto{\pgfqpoint{1.699356in}{1.294433in}}%
\pgfpathlineto{\pgfqpoint{1.701874in}{1.301420in}}%
\pgfpathlineto{\pgfqpoint{1.704393in}{1.307588in}}%
\pgfpathlineto{\pgfqpoint{1.706912in}{1.317577in}}%
\pgfpathlineto{\pgfqpoint{1.709431in}{1.329308in}}%
\pgfpathlineto{\pgfqpoint{1.711949in}{1.333732in}}%
\pgfpathlineto{\pgfqpoint{1.714468in}{1.329399in}}%
\pgfpathlineto{\pgfqpoint{1.716987in}{1.322205in}}%
\pgfpathlineto{\pgfqpoint{1.719506in}{1.303244in}}%
\pgfpathlineto{\pgfqpoint{1.722024in}{1.282358in}}%
\pgfpathlineto{\pgfqpoint{1.724543in}{1.272785in}}%
\pgfpathlineto{\pgfqpoint{1.727062in}{1.275436in}}%
\pgfpathlineto{\pgfqpoint{1.729581in}{1.260233in}}%
\pgfpathlineto{\pgfqpoint{1.732099in}{1.269872in}}%
\pgfpathlineto{\pgfqpoint{1.734618in}{1.273920in}}%
\pgfpathlineto{\pgfqpoint{1.737137in}{1.295576in}}%
\pgfpathlineto{\pgfqpoint{1.739656in}{1.297619in}}%
\pgfpathlineto{\pgfqpoint{1.742174in}{1.283012in}}%
\pgfpathlineto{\pgfqpoint{1.744693in}{1.286804in}}%
\pgfpathlineto{\pgfqpoint{1.747212in}{1.289147in}}%
\pgfpathlineto{\pgfqpoint{1.749731in}{1.302415in}}%
\pgfpathlineto{\pgfqpoint{1.752249in}{1.294280in}}%
\pgfpathlineto{\pgfqpoint{1.754768in}{1.296398in}}%
\pgfpathlineto{\pgfqpoint{1.757287in}{1.304924in}}%
\pgfpathlineto{\pgfqpoint{1.759806in}{1.297626in}}%
\pgfpathlineto{\pgfqpoint{1.762324in}{1.303219in}}%
\pgfpathlineto{\pgfqpoint{1.764843in}{1.330895in}}%
\pgfpathlineto{\pgfqpoint{1.769881in}{1.361078in}}%
\pgfpathlineto{\pgfqpoint{1.772399in}{1.389800in}}%
\pgfpathlineto{\pgfqpoint{1.774918in}{1.393612in}}%
\pgfpathlineto{\pgfqpoint{1.777437in}{1.367873in}}%
\pgfpathlineto{\pgfqpoint{1.779956in}{1.374380in}}%
\pgfpathlineto{\pgfqpoint{1.782474in}{1.382194in}}%
\pgfpathlineto{\pgfqpoint{1.784993in}{1.391259in}}%
\pgfpathlineto{\pgfqpoint{1.787512in}{1.410018in}}%
\pgfpathlineto{\pgfqpoint{1.790031in}{1.375538in}}%
\pgfpathlineto{\pgfqpoint{1.792549in}{1.374441in}}%
\pgfpathlineto{\pgfqpoint{1.795068in}{1.372472in}}%
\pgfpathlineto{\pgfqpoint{1.797587in}{1.369985in}}%
\pgfpathlineto{\pgfqpoint{1.800106in}{1.373024in}}%
\pgfpathlineto{\pgfqpoint{1.802624in}{1.361004in}}%
\pgfpathlineto{\pgfqpoint{1.805143in}{1.343532in}}%
\pgfpathlineto{\pgfqpoint{1.807662in}{1.330308in}}%
\pgfpathlineto{\pgfqpoint{1.810181in}{1.350347in}}%
\pgfpathlineto{\pgfqpoint{1.812699in}{1.333991in}}%
\pgfpathlineto{\pgfqpoint{1.815218in}{1.325804in}}%
\pgfpathlineto{\pgfqpoint{1.817737in}{1.333442in}}%
\pgfpathlineto{\pgfqpoint{1.820256in}{1.319135in}}%
\pgfpathlineto{\pgfqpoint{1.822774in}{1.307051in}}%
\pgfpathlineto{\pgfqpoint{1.825293in}{1.288933in}}%
\pgfpathlineto{\pgfqpoint{1.827812in}{1.288668in}}%
\pgfpathlineto{\pgfqpoint{1.830331in}{1.284842in}}%
\pgfpathlineto{\pgfqpoint{1.832849in}{1.285908in}}%
\pgfpathlineto{\pgfqpoint{1.835368in}{1.290513in}}%
\pgfpathlineto{\pgfqpoint{1.837887in}{1.278698in}}%
\pgfpathlineto{\pgfqpoint{1.840406in}{1.307434in}}%
\pgfpathlineto{\pgfqpoint{1.842924in}{1.303945in}}%
\pgfpathlineto{\pgfqpoint{1.845443in}{1.322837in}}%
\pgfpathlineto{\pgfqpoint{1.847962in}{1.329103in}}%
\pgfpathlineto{\pgfqpoint{1.850481in}{1.340312in}}%
\pgfpathlineto{\pgfqpoint{1.852999in}{1.317799in}}%
\pgfpathlineto{\pgfqpoint{1.855518in}{1.326272in}}%
\pgfpathlineto{\pgfqpoint{1.858037in}{1.320251in}}%
\pgfpathlineto{\pgfqpoint{1.860556in}{1.313692in}}%
\pgfpathlineto{\pgfqpoint{1.863074in}{1.314508in}}%
\pgfpathlineto{\pgfqpoint{1.865593in}{1.321576in}}%
\pgfpathlineto{\pgfqpoint{1.868112in}{1.336733in}}%
\pgfpathlineto{\pgfqpoint{1.870631in}{1.358188in}}%
\pgfpathlineto{\pgfqpoint{1.873149in}{1.340981in}}%
\pgfpathlineto{\pgfqpoint{1.875668in}{1.327656in}}%
\pgfpathlineto{\pgfqpoint{1.878187in}{1.323509in}}%
\pgfpathlineto{\pgfqpoint{1.880706in}{1.307176in}}%
\pgfpathlineto{\pgfqpoint{1.883224in}{1.338622in}}%
\pgfpathlineto{\pgfqpoint{1.885743in}{1.340955in}}%
\pgfpathlineto{\pgfqpoint{1.888262in}{1.332470in}}%
\pgfpathlineto{\pgfqpoint{1.890781in}{1.338612in}}%
\pgfpathlineto{\pgfqpoint{1.893299in}{1.329161in}}%
\pgfpathlineto{\pgfqpoint{1.895818in}{1.344713in}}%
\pgfpathlineto{\pgfqpoint{1.898337in}{1.347416in}}%
\pgfpathlineto{\pgfqpoint{1.900856in}{1.360697in}}%
\pgfpathlineto{\pgfqpoint{1.903374in}{1.379314in}}%
\pgfpathlineto{\pgfqpoint{1.905893in}{1.411518in}}%
\pgfpathlineto{\pgfqpoint{1.908412in}{1.402130in}}%
\pgfpathlineto{\pgfqpoint{1.910931in}{1.397714in}}%
\pgfpathlineto{\pgfqpoint{1.913449in}{1.390380in}}%
\pgfpathlineto{\pgfqpoint{1.915968in}{1.390327in}}%
\pgfpathlineto{\pgfqpoint{1.918487in}{1.383170in}}%
\pgfpathlineto{\pgfqpoint{1.921006in}{1.396514in}}%
\pgfpathlineto{\pgfqpoint{1.923524in}{1.392335in}}%
\pgfpathlineto{\pgfqpoint{1.926043in}{1.410174in}}%
\pgfpathlineto{\pgfqpoint{1.928562in}{1.406144in}}%
\pgfpathlineto{\pgfqpoint{1.931081in}{1.400241in}}%
\pgfpathlineto{\pgfqpoint{1.933599in}{1.409427in}}%
\pgfpathlineto{\pgfqpoint{1.936118in}{1.404745in}}%
\pgfpathlineto{\pgfqpoint{1.938637in}{1.395975in}}%
\pgfpathlineto{\pgfqpoint{1.941156in}{1.399959in}}%
\pgfpathlineto{\pgfqpoint{1.943674in}{1.397562in}}%
\pgfpathlineto{\pgfqpoint{1.946193in}{1.396234in}}%
\pgfpathlineto{\pgfqpoint{1.948712in}{1.420418in}}%
\pgfpathlineto{\pgfqpoint{1.951231in}{1.410787in}}%
\pgfpathlineto{\pgfqpoint{1.953749in}{1.394222in}}%
\pgfpathlineto{\pgfqpoint{1.956268in}{1.397224in}}%
\pgfpathlineto{\pgfqpoint{1.958787in}{1.389281in}}%
\pgfpathlineto{\pgfqpoint{1.961306in}{1.400717in}}%
\pgfpathlineto{\pgfqpoint{1.963824in}{1.409389in}}%
\pgfpathlineto{\pgfqpoint{1.966343in}{1.402299in}}%
\pgfpathlineto{\pgfqpoint{1.968862in}{1.411075in}}%
\pgfpathlineto{\pgfqpoint{1.971381in}{1.409967in}}%
\pgfpathlineto{\pgfqpoint{1.973899in}{1.402246in}}%
\pgfpathlineto{\pgfqpoint{1.976418in}{1.419787in}}%
\pgfpathlineto{\pgfqpoint{1.978937in}{1.419019in}}%
\pgfpathlineto{\pgfqpoint{1.981456in}{1.442136in}}%
\pgfpathlineto{\pgfqpoint{1.983974in}{1.463105in}}%
\pgfpathlineto{\pgfqpoint{1.986493in}{1.457315in}}%
\pgfpathlineto{\pgfqpoint{1.989012in}{1.462589in}}%
\pgfpathlineto{\pgfqpoint{1.991531in}{1.462147in}}%
\pgfpathlineto{\pgfqpoint{1.994049in}{1.483220in}}%
\pgfpathlineto{\pgfqpoint{1.996568in}{1.472430in}}%
\pgfpathlineto{\pgfqpoint{1.999087in}{1.442568in}}%
\pgfpathlineto{\pgfqpoint{2.001606in}{1.445077in}}%
\pgfpathlineto{\pgfqpoint{2.004124in}{1.451360in}}%
\pgfpathlineto{\pgfqpoint{2.006643in}{1.452726in}}%
\pgfpathlineto{\pgfqpoint{2.009162in}{1.465708in}}%
\pgfpathlineto{\pgfqpoint{2.011681in}{1.489016in}}%
\pgfpathlineto{\pgfqpoint{2.014199in}{1.500923in}}%
\pgfpathlineto{\pgfqpoint{2.016718in}{1.493397in}}%
\pgfpathlineto{\pgfqpoint{2.019237in}{1.477176in}}%
\pgfpathlineto{\pgfqpoint{2.021756in}{1.475882in}}%
\pgfpathlineto{\pgfqpoint{2.024274in}{1.477170in}}%
\pgfpathlineto{\pgfqpoint{2.026793in}{1.466159in}}%
\pgfpathlineto{\pgfqpoint{2.029312in}{1.456268in}}%
\pgfpathlineto{\pgfqpoint{2.031831in}{1.439973in}}%
\pgfpathlineto{\pgfqpoint{2.034349in}{1.430278in}}%
\pgfpathlineto{\pgfqpoint{2.036868in}{1.437926in}}%
\pgfpathlineto{\pgfqpoint{2.039387in}{1.455085in}}%
\pgfpathlineto{\pgfqpoint{2.041906in}{1.456300in}}%
\pgfpathlineto{\pgfqpoint{2.044424in}{1.461353in}}%
\pgfpathlineto{\pgfqpoint{2.046943in}{1.437234in}}%
\pgfpathlineto{\pgfqpoint{2.049462in}{1.437107in}}%
\pgfpathlineto{\pgfqpoint{2.051981in}{1.443646in}}%
\pgfpathlineto{\pgfqpoint{2.054499in}{1.449635in}}%
\pgfpathlineto{\pgfqpoint{2.057018in}{1.453838in}}%
\pgfpathlineto{\pgfqpoint{2.059537in}{1.436610in}}%
\pgfpathlineto{\pgfqpoint{2.062056in}{1.456118in}}%
\pgfpathlineto{\pgfqpoint{2.064574in}{1.438191in}}%
\pgfpathlineto{\pgfqpoint{2.067093in}{1.432834in}}%
\pgfpathlineto{\pgfqpoint{2.069612in}{1.403345in}}%
\pgfpathlineto{\pgfqpoint{2.072131in}{1.394399in}}%
\pgfpathlineto{\pgfqpoint{2.074649in}{1.407287in}}%
\pgfpathlineto{\pgfqpoint{2.077168in}{1.407219in}}%
\pgfpathlineto{\pgfqpoint{2.079687in}{1.400356in}}%
\pgfpathlineto{\pgfqpoint{2.082206in}{1.386314in}}%
\pgfpathlineto{\pgfqpoint{2.084724in}{1.408873in}}%
\pgfpathlineto{\pgfqpoint{2.087243in}{1.395035in}}%
\pgfpathlineto{\pgfqpoint{2.089762in}{1.394459in}}%
\pgfpathlineto{\pgfqpoint{2.092281in}{1.384747in}}%
\pgfpathlineto{\pgfqpoint{2.094799in}{1.392489in}}%
\pgfpathlineto{\pgfqpoint{2.097318in}{1.385413in}}%
\pgfpathlineto{\pgfqpoint{2.099837in}{1.391139in}}%
\pgfpathlineto{\pgfqpoint{2.102356in}{1.412375in}}%
\pgfpathlineto{\pgfqpoint{2.104874in}{1.399566in}}%
\pgfpathlineto{\pgfqpoint{2.107393in}{1.380617in}}%
\pgfpathlineto{\pgfqpoint{2.109912in}{1.380419in}}%
\pgfpathlineto{\pgfqpoint{2.112431in}{1.386111in}}%
\pgfpathlineto{\pgfqpoint{2.114949in}{1.407393in}}%
\pgfpathlineto{\pgfqpoint{2.117468in}{1.394098in}}%
\pgfpathlineto{\pgfqpoint{2.119987in}{1.420140in}}%
\pgfpathlineto{\pgfqpoint{2.122506in}{1.449333in}}%
\pgfpathlineto{\pgfqpoint{2.125024in}{1.471443in}}%
\pgfpathlineto{\pgfqpoint{2.127543in}{1.480072in}}%
\pgfpathlineto{\pgfqpoint{2.130062in}{1.491663in}}%
\pgfpathlineto{\pgfqpoint{2.132581in}{1.508719in}}%
\pgfpathlineto{\pgfqpoint{2.135099in}{1.502808in}}%
\pgfpathlineto{\pgfqpoint{2.137618in}{1.482106in}}%
\pgfpathlineto{\pgfqpoint{2.140137in}{1.490583in}}%
\pgfpathlineto{\pgfqpoint{2.142656in}{1.496485in}}%
\pgfpathlineto{\pgfqpoint{2.145174in}{1.504417in}}%
\pgfpathlineto{\pgfqpoint{2.147693in}{1.530456in}}%
\pgfpathlineto{\pgfqpoint{2.150212in}{1.549459in}}%
\pgfpathlineto{\pgfqpoint{2.152731in}{1.537349in}}%
\pgfpathlineto{\pgfqpoint{2.155249in}{1.535532in}}%
\pgfpathlineto{\pgfqpoint{2.157768in}{1.499434in}}%
\pgfpathlineto{\pgfqpoint{2.160287in}{1.485041in}}%
\pgfpathlineto{\pgfqpoint{2.162806in}{1.514864in}}%
\pgfpathlineto{\pgfqpoint{2.165324in}{1.499824in}}%
\pgfpathlineto{\pgfqpoint{2.167843in}{1.500659in}}%
\pgfpathlineto{\pgfqpoint{2.170362in}{1.496875in}}%
\pgfpathlineto{\pgfqpoint{2.172881in}{1.490124in}}%
\pgfpathlineto{\pgfqpoint{2.175399in}{1.497348in}}%
\pgfpathlineto{\pgfqpoint{2.177918in}{1.508517in}}%
\pgfpathlineto{\pgfqpoint{2.180437in}{1.512318in}}%
\pgfpathlineto{\pgfqpoint{2.182956in}{1.494350in}}%
\pgfpathlineto{\pgfqpoint{2.185474in}{1.500838in}}%
\pgfpathlineto{\pgfqpoint{2.187993in}{1.499055in}}%
\pgfpathlineto{\pgfqpoint{2.190512in}{1.491426in}}%
\pgfpathlineto{\pgfqpoint{2.193031in}{1.472032in}}%
\pgfpathlineto{\pgfqpoint{2.195549in}{1.465309in}}%
\pgfpathlineto{\pgfqpoint{2.198068in}{1.465774in}}%
\pgfpathlineto{\pgfqpoint{2.200587in}{1.484485in}}%
\pgfpathlineto{\pgfqpoint{2.203106in}{1.476422in}}%
\pgfpathlineto{\pgfqpoint{2.205624in}{1.480167in}}%
\pgfpathlineto{\pgfqpoint{2.208143in}{1.460144in}}%
\pgfpathlineto{\pgfqpoint{2.210662in}{1.466282in}}%
\pgfpathlineto{\pgfqpoint{2.213181in}{1.484674in}}%
\pgfpathlineto{\pgfqpoint{2.215699in}{1.509055in}}%
\pgfpathlineto{\pgfqpoint{2.218218in}{1.498828in}}%
\pgfpathlineto{\pgfqpoint{2.220737in}{1.509095in}}%
\pgfpathlineto{\pgfqpoint{2.223256in}{1.513498in}}%
\pgfpathlineto{\pgfqpoint{2.225774in}{1.500583in}}%
\pgfpathlineto{\pgfqpoint{2.228293in}{1.495789in}}%
\pgfpathlineto{\pgfqpoint{2.230812in}{1.483686in}}%
\pgfpathlineto{\pgfqpoint{2.233331in}{1.482146in}}%
\pgfpathlineto{\pgfqpoint{2.235849in}{1.472608in}}%
\pgfpathlineto{\pgfqpoint{2.238368in}{1.456096in}}%
\pgfpathlineto{\pgfqpoint{2.240887in}{1.472714in}}%
\pgfpathlineto{\pgfqpoint{2.243406in}{1.455578in}}%
\pgfpathlineto{\pgfqpoint{2.245924in}{1.455344in}}%
\pgfpathlineto{\pgfqpoint{2.248443in}{1.459504in}}%
\pgfpathlineto{\pgfqpoint{2.250962in}{1.445581in}}%
\pgfpathlineto{\pgfqpoint{2.258518in}{1.494492in}}%
\pgfpathlineto{\pgfqpoint{2.261037in}{1.492463in}}%
\pgfpathlineto{\pgfqpoint{2.263556in}{1.497340in}}%
\pgfpathlineto{\pgfqpoint{2.266074in}{1.511211in}}%
\pgfpathlineto{\pgfqpoint{2.268593in}{1.519193in}}%
\pgfpathlineto{\pgfqpoint{2.271112in}{1.503845in}}%
\pgfpathlineto{\pgfqpoint{2.273631in}{1.499425in}}%
\pgfpathlineto{\pgfqpoint{2.276149in}{1.513009in}}%
\pgfpathlineto{\pgfqpoint{2.278668in}{1.513995in}}%
\pgfpathlineto{\pgfqpoint{2.281187in}{1.528156in}}%
\pgfpathlineto{\pgfqpoint{2.283706in}{1.552285in}}%
\pgfpathlineto{\pgfqpoint{2.286224in}{1.573417in}}%
\pgfpathlineto{\pgfqpoint{2.288743in}{1.590953in}}%
\pgfpathlineto{\pgfqpoint{2.291262in}{1.605340in}}%
\pgfpathlineto{\pgfqpoint{2.293781in}{1.621040in}}%
\pgfpathlineto{\pgfqpoint{2.296299in}{1.612566in}}%
\pgfpathlineto{\pgfqpoint{2.298818in}{1.597605in}}%
\pgfpathlineto{\pgfqpoint{2.301337in}{1.590633in}}%
\pgfpathlineto{\pgfqpoint{2.303856in}{1.625882in}}%
\pgfpathlineto{\pgfqpoint{2.306374in}{1.637432in}}%
\pgfpathlineto{\pgfqpoint{2.308893in}{1.637143in}}%
\pgfpathlineto{\pgfqpoint{2.311412in}{1.640087in}}%
\pgfpathlineto{\pgfqpoint{2.313931in}{1.622796in}}%
\pgfpathlineto{\pgfqpoint{2.316449in}{1.609165in}}%
\pgfpathlineto{\pgfqpoint{2.318968in}{1.604741in}}%
\pgfpathlineto{\pgfqpoint{2.321487in}{1.619461in}}%
\pgfpathlineto{\pgfqpoint{2.324006in}{1.618517in}}%
\pgfpathlineto{\pgfqpoint{2.326524in}{1.616643in}}%
\pgfpathlineto{\pgfqpoint{2.329043in}{1.605716in}}%
\pgfpathlineto{\pgfqpoint{2.331562in}{1.605752in}}%
\pgfpathlineto{\pgfqpoint{2.334081in}{1.603834in}}%
\pgfpathlineto{\pgfqpoint{2.336599in}{1.582783in}}%
\pgfpathlineto{\pgfqpoint{2.339118in}{1.606537in}}%
\pgfpathlineto{\pgfqpoint{2.341637in}{1.625942in}}%
\pgfpathlineto{\pgfqpoint{2.344156in}{1.631822in}}%
\pgfpathlineto{\pgfqpoint{2.346674in}{1.633243in}}%
\pgfpathlineto{\pgfqpoint{2.349193in}{1.641232in}}%
\pgfpathlineto{\pgfqpoint{2.354231in}{1.667301in}}%
\pgfpathlineto{\pgfqpoint{2.356749in}{1.682851in}}%
\pgfpathlineto{\pgfqpoint{2.359268in}{1.682849in}}%
\pgfpathlineto{\pgfqpoint{2.361787in}{1.693164in}}%
\pgfpathlineto{\pgfqpoint{2.364306in}{1.683170in}}%
\pgfpathlineto{\pgfqpoint{2.366824in}{1.714686in}}%
\pgfpathlineto{\pgfqpoint{2.369343in}{1.716671in}}%
\pgfpathlineto{\pgfqpoint{2.371862in}{1.695876in}}%
\pgfpathlineto{\pgfqpoint{2.374381in}{1.697086in}}%
\pgfpathlineto{\pgfqpoint{2.376899in}{1.688626in}}%
\pgfpathlineto{\pgfqpoint{2.379418in}{1.677474in}}%
\pgfpathlineto{\pgfqpoint{2.381937in}{1.682735in}}%
\pgfpathlineto{\pgfqpoint{2.384456in}{1.692647in}}%
\pgfpathlineto{\pgfqpoint{2.386974in}{1.676845in}}%
\pgfpathlineto{\pgfqpoint{2.389493in}{1.678943in}}%
\pgfpathlineto{\pgfqpoint{2.392012in}{1.719638in}}%
\pgfpathlineto{\pgfqpoint{2.394531in}{1.732805in}}%
\pgfpathlineto{\pgfqpoint{2.399568in}{1.713135in}}%
\pgfpathlineto{\pgfqpoint{2.402087in}{1.719247in}}%
\pgfpathlineto{\pgfqpoint{2.404606in}{1.711334in}}%
\pgfpathlineto{\pgfqpoint{2.407124in}{1.719216in}}%
\pgfpathlineto{\pgfqpoint{2.409643in}{1.718255in}}%
\pgfpathlineto{\pgfqpoint{2.412162in}{1.719507in}}%
\pgfpathlineto{\pgfqpoint{2.414681in}{1.710996in}}%
\pgfpathlineto{\pgfqpoint{2.417199in}{1.701017in}}%
\pgfpathlineto{\pgfqpoint{2.419718in}{1.691683in}}%
\pgfpathlineto{\pgfqpoint{2.422237in}{1.726691in}}%
\pgfpathlineto{\pgfqpoint{2.424756in}{1.714226in}}%
\pgfpathlineto{\pgfqpoint{2.427274in}{1.730450in}}%
\pgfpathlineto{\pgfqpoint{2.429793in}{1.749641in}}%
\pgfpathlineto{\pgfqpoint{2.432312in}{1.759970in}}%
\pgfpathlineto{\pgfqpoint{2.434831in}{1.753552in}}%
\pgfpathlineto{\pgfqpoint{2.437349in}{1.756952in}}%
\pgfpathlineto{\pgfqpoint{2.439868in}{1.761205in}}%
\pgfpathlineto{\pgfqpoint{2.442387in}{1.755074in}}%
\pgfpathlineto{\pgfqpoint{2.444906in}{1.735120in}}%
\pgfpathlineto{\pgfqpoint{2.447424in}{1.752267in}}%
\pgfpathlineto{\pgfqpoint{2.449943in}{1.727757in}}%
\pgfpathlineto{\pgfqpoint{2.452462in}{1.741090in}}%
\pgfpathlineto{\pgfqpoint{2.454981in}{1.742973in}}%
\pgfpathlineto{\pgfqpoint{2.457499in}{1.731594in}}%
\pgfpathlineto{\pgfqpoint{2.460018in}{1.736613in}}%
\pgfpathlineto{\pgfqpoint{2.462537in}{1.738206in}}%
\pgfpathlineto{\pgfqpoint{2.465056in}{1.734236in}}%
\pgfpathlineto{\pgfqpoint{2.467574in}{1.719280in}}%
\pgfpathlineto{\pgfqpoint{2.470093in}{1.723278in}}%
\pgfpathlineto{\pgfqpoint{2.472612in}{1.775363in}}%
\pgfpathlineto{\pgfqpoint{2.475131in}{1.776949in}}%
\pgfpathlineto{\pgfqpoint{2.477649in}{1.775952in}}%
\pgfpathlineto{\pgfqpoint{2.480168in}{1.765888in}}%
\pgfpathlineto{\pgfqpoint{2.482687in}{1.763295in}}%
\pgfpathlineto{\pgfqpoint{2.485206in}{1.731067in}}%
\pgfpathlineto{\pgfqpoint{2.487724in}{1.731709in}}%
\pgfpathlineto{\pgfqpoint{2.490243in}{1.740690in}}%
\pgfpathlineto{\pgfqpoint{2.492762in}{1.722165in}}%
\pgfpathlineto{\pgfqpoint{2.495281in}{1.736802in}}%
\pgfpathlineto{\pgfqpoint{2.497799in}{1.733187in}}%
\pgfpathlineto{\pgfqpoint{2.500318in}{1.740686in}}%
\pgfpathlineto{\pgfqpoint{2.502837in}{1.709302in}}%
\pgfpathlineto{\pgfqpoint{2.505356in}{1.703534in}}%
\pgfpathlineto{\pgfqpoint{2.507874in}{1.665511in}}%
\pgfpathlineto{\pgfqpoint{2.510393in}{1.677472in}}%
\pgfpathlineto{\pgfqpoint{2.512912in}{1.669092in}}%
\pgfpathlineto{\pgfqpoint{2.515431in}{1.660107in}}%
\pgfpathlineto{\pgfqpoint{2.517949in}{1.648843in}}%
\pgfpathlineto{\pgfqpoint{2.520468in}{1.627705in}}%
\pgfpathlineto{\pgfqpoint{2.522987in}{1.638269in}}%
\pgfpathlineto{\pgfqpoint{2.525506in}{1.626957in}}%
\pgfpathlineto{\pgfqpoint{2.528024in}{1.616984in}}%
\pgfpathlineto{\pgfqpoint{2.530543in}{1.626805in}}%
\pgfpathlineto{\pgfqpoint{2.533062in}{1.618693in}}%
\pgfpathlineto{\pgfqpoint{2.535581in}{1.632284in}}%
\pgfpathlineto{\pgfqpoint{2.538099in}{1.661769in}}%
\pgfpathlineto{\pgfqpoint{2.540618in}{1.630655in}}%
\pgfpathlineto{\pgfqpoint{2.543137in}{1.614580in}}%
\pgfpathlineto{\pgfqpoint{2.545656in}{1.610395in}}%
\pgfpathlineto{\pgfqpoint{2.548174in}{1.620073in}}%
\pgfpathlineto{\pgfqpoint{2.550693in}{1.611138in}}%
\pgfpathlineto{\pgfqpoint{2.553212in}{1.643602in}}%
\pgfpathlineto{\pgfqpoint{2.555731in}{1.630323in}}%
\pgfpathlineto{\pgfqpoint{2.558249in}{1.599354in}}%
\pgfpathlineto{\pgfqpoint{2.560768in}{1.582034in}}%
\pgfpathlineto{\pgfqpoint{2.563287in}{1.588777in}}%
\pgfpathlineto{\pgfqpoint{2.565806in}{1.573090in}}%
\pgfpathlineto{\pgfqpoint{2.568324in}{1.561393in}}%
\pgfpathlineto{\pgfqpoint{2.570843in}{1.558468in}}%
\pgfpathlineto{\pgfqpoint{2.573362in}{1.547017in}}%
\pgfpathlineto{\pgfqpoint{2.575881in}{1.550650in}}%
\pgfpathlineto{\pgfqpoint{2.578399in}{1.564819in}}%
\pgfpathlineto{\pgfqpoint{2.580918in}{1.552185in}}%
\pgfpathlineto{\pgfqpoint{2.583437in}{1.576848in}}%
\pgfpathlineto{\pgfqpoint{2.585956in}{1.571706in}}%
\pgfpathlineto{\pgfqpoint{2.588474in}{1.570472in}}%
\pgfpathlineto{\pgfqpoint{2.590993in}{1.586348in}}%
\pgfpathlineto{\pgfqpoint{2.593512in}{1.579893in}}%
\pgfpathlineto{\pgfqpoint{2.596031in}{1.569397in}}%
\pgfpathlineto{\pgfqpoint{2.598549in}{1.568921in}}%
\pgfpathlineto{\pgfqpoint{2.601068in}{1.563287in}}%
\pgfpathlineto{\pgfqpoint{2.606106in}{1.531792in}}%
\pgfpathlineto{\pgfqpoint{2.608624in}{1.538845in}}%
\pgfpathlineto{\pgfqpoint{2.611143in}{1.533536in}}%
\pgfpathlineto{\pgfqpoint{2.613662in}{1.569274in}}%
\pgfpathlineto{\pgfqpoint{2.616181in}{1.571461in}}%
\pgfpathlineto{\pgfqpoint{2.618699in}{1.576566in}}%
\pgfpathlineto{\pgfqpoint{2.621218in}{1.542551in}}%
\pgfpathlineto{\pgfqpoint{2.623737in}{1.532384in}}%
\pgfpathlineto{\pgfqpoint{2.626256in}{1.539493in}}%
\pgfpathlineto{\pgfqpoint{2.628774in}{1.536890in}}%
\pgfpathlineto{\pgfqpoint{2.631293in}{1.533019in}}%
\pgfpathlineto{\pgfqpoint{2.633812in}{1.534496in}}%
\pgfpathlineto{\pgfqpoint{2.636331in}{1.525406in}}%
\pgfpathlineto{\pgfqpoint{2.638849in}{1.496506in}}%
\pgfpathlineto{\pgfqpoint{2.641368in}{1.503462in}}%
\pgfpathlineto{\pgfqpoint{2.643887in}{1.513895in}}%
\pgfpathlineto{\pgfqpoint{2.646406in}{1.533720in}}%
\pgfpathlineto{\pgfqpoint{2.648924in}{1.548654in}}%
\pgfpathlineto{\pgfqpoint{2.653962in}{1.601251in}}%
\pgfpathlineto{\pgfqpoint{2.656481in}{1.614301in}}%
\pgfpathlineto{\pgfqpoint{2.658999in}{1.633610in}}%
\pgfpathlineto{\pgfqpoint{2.661518in}{1.646814in}}%
\pgfpathlineto{\pgfqpoint{2.664037in}{1.657616in}}%
\pgfpathlineto{\pgfqpoint{2.666556in}{1.656247in}}%
\pgfpathlineto{\pgfqpoint{2.669074in}{1.671285in}}%
\pgfpathlineto{\pgfqpoint{2.671593in}{1.676787in}}%
\pgfpathlineto{\pgfqpoint{2.674112in}{1.699928in}}%
\pgfpathlineto{\pgfqpoint{2.676631in}{1.698019in}}%
\pgfpathlineto{\pgfqpoint{2.679149in}{1.708474in}}%
\pgfpathlineto{\pgfqpoint{2.681668in}{1.709089in}}%
\pgfpathlineto{\pgfqpoint{2.684187in}{1.698726in}}%
\pgfpathlineto{\pgfqpoint{2.686706in}{1.705409in}}%
\pgfpathlineto{\pgfqpoint{2.689224in}{1.694547in}}%
\pgfpathlineto{\pgfqpoint{2.691743in}{1.722358in}}%
\pgfpathlineto{\pgfqpoint{2.694262in}{1.729343in}}%
\pgfpathlineto{\pgfqpoint{2.696781in}{1.716227in}}%
\pgfpathlineto{\pgfqpoint{2.699299in}{1.712845in}}%
\pgfpathlineto{\pgfqpoint{2.701818in}{1.718213in}}%
\pgfpathlineto{\pgfqpoint{2.704337in}{1.708697in}}%
\pgfpathlineto{\pgfqpoint{2.706856in}{1.711040in}}%
\pgfpathlineto{\pgfqpoint{2.709374in}{1.724474in}}%
\pgfpathlineto{\pgfqpoint{2.711893in}{1.733711in}}%
\pgfpathlineto{\pgfqpoint{2.714412in}{1.747758in}}%
\pgfpathlineto{\pgfqpoint{2.716931in}{1.753029in}}%
\pgfpathlineto{\pgfqpoint{2.719449in}{1.756831in}}%
\pgfpathlineto{\pgfqpoint{2.721968in}{1.761237in}}%
\pgfpathlineto{\pgfqpoint{2.724487in}{1.754682in}}%
\pgfpathlineto{\pgfqpoint{2.727006in}{1.757829in}}%
\pgfpathlineto{\pgfqpoint{2.729524in}{1.749992in}}%
\pgfpathlineto{\pgfqpoint{2.732043in}{1.760379in}}%
\pgfpathlineto{\pgfqpoint{2.734562in}{1.795217in}}%
\pgfpathlineto{\pgfqpoint{2.737081in}{1.831741in}}%
\pgfpathlineto{\pgfqpoint{2.739599in}{1.839965in}}%
\pgfpathlineto{\pgfqpoint{2.742118in}{1.834773in}}%
\pgfpathlineto{\pgfqpoint{2.744637in}{1.832721in}}%
\pgfpathlineto{\pgfqpoint{2.747156in}{1.830468in}}%
\pgfpathlineto{\pgfqpoint{2.749674in}{1.809275in}}%
\pgfpathlineto{\pgfqpoint{2.752193in}{1.813103in}}%
\pgfpathlineto{\pgfqpoint{2.754712in}{1.810232in}}%
\pgfpathlineto{\pgfqpoint{2.757231in}{1.803742in}}%
\pgfpathlineto{\pgfqpoint{2.759749in}{1.813858in}}%
\pgfpathlineto{\pgfqpoint{2.762268in}{1.812483in}}%
\pgfpathlineto{\pgfqpoint{2.764787in}{1.822555in}}%
\pgfpathlineto{\pgfqpoint{2.767306in}{1.826687in}}%
\pgfpathlineto{\pgfqpoint{2.769824in}{1.809401in}}%
\pgfpathlineto{\pgfqpoint{2.772343in}{1.804892in}}%
\pgfpathlineto{\pgfqpoint{2.774862in}{1.820744in}}%
\pgfpathlineto{\pgfqpoint{2.777381in}{1.800324in}}%
\pgfpathlineto{\pgfqpoint{2.779899in}{1.781883in}}%
\pgfpathlineto{\pgfqpoint{2.782418in}{1.782755in}}%
\pgfpathlineto{\pgfqpoint{2.784937in}{1.782401in}}%
\pgfpathlineto{\pgfqpoint{2.787456in}{1.792077in}}%
\pgfpathlineto{\pgfqpoint{2.789974in}{1.776012in}}%
\pgfpathlineto{\pgfqpoint{2.792493in}{1.770868in}}%
\pgfpathlineto{\pgfqpoint{2.795012in}{1.763202in}}%
\pgfpathlineto{\pgfqpoint{2.797531in}{1.763920in}}%
\pgfpathlineto{\pgfqpoint{2.800049in}{1.766226in}}%
\pgfpathlineto{\pgfqpoint{2.802568in}{1.776494in}}%
\pgfpathlineto{\pgfqpoint{2.805087in}{1.800906in}}%
\pgfpathlineto{\pgfqpoint{2.807606in}{1.789605in}}%
\pgfpathlineto{\pgfqpoint{2.810124in}{1.813326in}}%
\pgfpathlineto{\pgfqpoint{2.812643in}{1.850111in}}%
\pgfpathlineto{\pgfqpoint{2.815162in}{1.880475in}}%
\pgfpathlineto{\pgfqpoint{2.817681in}{1.876027in}}%
\pgfpathlineto{\pgfqpoint{2.820199in}{1.921364in}}%
\pgfpathlineto{\pgfqpoint{2.822718in}{1.916282in}}%
\pgfpathlineto{\pgfqpoint{2.825237in}{1.923633in}}%
\pgfpathlineto{\pgfqpoint{2.827756in}{1.948097in}}%
\pgfpathlineto{\pgfqpoint{2.830274in}{1.955777in}}%
\pgfpathlineto{\pgfqpoint{2.832793in}{1.954958in}}%
\pgfpathlineto{\pgfqpoint{2.835312in}{1.962357in}}%
\pgfpathlineto{\pgfqpoint{2.837831in}{1.970735in}}%
\pgfpathlineto{\pgfqpoint{2.840349in}{1.973002in}}%
\pgfpathlineto{\pgfqpoint{2.842868in}{1.994917in}}%
\pgfpathlineto{\pgfqpoint{2.845387in}{1.969038in}}%
\pgfpathlineto{\pgfqpoint{2.847906in}{1.990524in}}%
\pgfpathlineto{\pgfqpoint{2.852943in}{1.958815in}}%
\pgfpathlineto{\pgfqpoint{2.855462in}{1.949927in}}%
\pgfpathlineto{\pgfqpoint{2.857981in}{1.947785in}}%
\pgfpathlineto{\pgfqpoint{2.860499in}{1.914081in}}%
\pgfpathlineto{\pgfqpoint{2.863018in}{1.865596in}}%
\pgfpathlineto{\pgfqpoint{2.865537in}{1.840946in}}%
\pgfpathlineto{\pgfqpoint{2.868056in}{1.830803in}}%
\pgfpathlineto{\pgfqpoint{2.870574in}{1.831865in}}%
\pgfpathlineto{\pgfqpoint{2.873093in}{1.813287in}}%
\pgfpathlineto{\pgfqpoint{2.875612in}{1.847328in}}%
\pgfpathlineto{\pgfqpoint{2.878131in}{1.839422in}}%
\pgfpathlineto{\pgfqpoint{2.880649in}{1.873131in}}%
\pgfpathlineto{\pgfqpoint{2.883168in}{1.901521in}}%
\pgfpathlineto{\pgfqpoint{2.885687in}{1.921367in}}%
\pgfpathlineto{\pgfqpoint{2.888206in}{1.921860in}}%
\pgfpathlineto{\pgfqpoint{2.890724in}{1.947829in}}%
\pgfpathlineto{\pgfqpoint{2.893243in}{1.955925in}}%
\pgfpathlineto{\pgfqpoint{2.895762in}{1.964649in}}%
\pgfpathlineto{\pgfqpoint{2.898281in}{1.977040in}}%
\pgfpathlineto{\pgfqpoint{2.900799in}{1.983376in}}%
\pgfpathlineto{\pgfqpoint{2.903318in}{1.968329in}}%
\pgfpathlineto{\pgfqpoint{2.905837in}{1.961160in}}%
\pgfpathlineto{\pgfqpoint{2.910874in}{1.937653in}}%
\pgfpathlineto{\pgfqpoint{2.913393in}{1.936894in}}%
\pgfpathlineto{\pgfqpoint{2.915912in}{1.954503in}}%
\pgfpathlineto{\pgfqpoint{2.918431in}{1.948405in}}%
\pgfpathlineto{\pgfqpoint{2.920949in}{1.982506in}}%
\pgfpathlineto{\pgfqpoint{2.923468in}{1.998527in}}%
\pgfpathlineto{\pgfqpoint{2.925987in}{1.992905in}}%
\pgfpathlineto{\pgfqpoint{2.928506in}{1.979690in}}%
\pgfpathlineto{\pgfqpoint{2.931024in}{1.996056in}}%
\pgfpathlineto{\pgfqpoint{2.933543in}{1.993711in}}%
\pgfpathlineto{\pgfqpoint{2.936062in}{1.992871in}}%
\pgfpathlineto{\pgfqpoint{2.938581in}{2.013807in}}%
\pgfpathlineto{\pgfqpoint{2.941099in}{2.007337in}}%
\pgfpathlineto{\pgfqpoint{2.943618in}{1.974522in}}%
\pgfpathlineto{\pgfqpoint{2.946137in}{1.976600in}}%
\pgfpathlineto{\pgfqpoint{2.948656in}{1.976105in}}%
\pgfpathlineto{\pgfqpoint{2.951174in}{1.967114in}}%
\pgfpathlineto{\pgfqpoint{2.953693in}{2.000131in}}%
\pgfpathlineto{\pgfqpoint{2.956212in}{1.987542in}}%
\pgfpathlineto{\pgfqpoint{2.958731in}{1.960539in}}%
\pgfpathlineto{\pgfqpoint{2.961249in}{1.965710in}}%
\pgfpathlineto{\pgfqpoint{2.963768in}{1.970098in}}%
\pgfpathlineto{\pgfqpoint{2.966287in}{1.979593in}}%
\pgfpathlineto{\pgfqpoint{2.968806in}{1.974824in}}%
\pgfpathlineto{\pgfqpoint{2.971324in}{1.968399in}}%
\pgfpathlineto{\pgfqpoint{2.973843in}{1.961569in}}%
\pgfpathlineto{\pgfqpoint{2.976362in}{1.957503in}}%
\pgfpathlineto{\pgfqpoint{2.978881in}{1.964586in}}%
\pgfpathlineto{\pgfqpoint{2.981399in}{1.969805in}}%
\pgfpathlineto{\pgfqpoint{2.983918in}{1.951880in}}%
\pgfpathlineto{\pgfqpoint{2.986437in}{1.952516in}}%
\pgfpathlineto{\pgfqpoint{2.988956in}{1.953798in}}%
\pgfpathlineto{\pgfqpoint{2.991474in}{1.966642in}}%
\pgfpathlineto{\pgfqpoint{2.993993in}{1.950791in}}%
\pgfpathlineto{\pgfqpoint{2.996512in}{1.972930in}}%
\pgfpathlineto{\pgfqpoint{2.999031in}{1.952262in}}%
\pgfpathlineto{\pgfqpoint{3.001549in}{1.964086in}}%
\pgfpathlineto{\pgfqpoint{3.004068in}{1.981464in}}%
\pgfpathlineto{\pgfqpoint{3.006587in}{1.979055in}}%
\pgfpathlineto{\pgfqpoint{3.009106in}{1.978037in}}%
\pgfpathlineto{\pgfqpoint{3.011624in}{1.960755in}}%
\pgfpathlineto{\pgfqpoint{3.014143in}{1.973409in}}%
\pgfpathlineto{\pgfqpoint{3.016662in}{1.987470in}}%
\pgfpathlineto{\pgfqpoint{3.019181in}{1.996728in}}%
\pgfpathlineto{\pgfqpoint{3.021699in}{1.990126in}}%
\pgfpathlineto{\pgfqpoint{3.024218in}{2.012505in}}%
\pgfpathlineto{\pgfqpoint{3.026737in}{2.011687in}}%
\pgfpathlineto{\pgfqpoint{3.029256in}{2.013441in}}%
\pgfpathlineto{\pgfqpoint{3.031774in}{2.067583in}}%
\pgfpathlineto{\pgfqpoint{3.034293in}{2.063518in}}%
\pgfpathlineto{\pgfqpoint{3.036812in}{2.083186in}}%
\pgfpathlineto{\pgfqpoint{3.039331in}{2.068995in}}%
\pgfpathlineto{\pgfqpoint{3.041849in}{2.081271in}}%
\pgfpathlineto{\pgfqpoint{3.044368in}{2.082381in}}%
\pgfpathlineto{\pgfqpoint{3.046887in}{2.096537in}}%
\pgfpathlineto{\pgfqpoint{3.049406in}{2.097041in}}%
\pgfpathlineto{\pgfqpoint{3.051924in}{2.109086in}}%
\pgfpathlineto{\pgfqpoint{3.054443in}{2.132576in}}%
\pgfpathlineto{\pgfqpoint{3.056962in}{2.103080in}}%
\pgfpathlineto{\pgfqpoint{3.059481in}{2.115419in}}%
\pgfpathlineto{\pgfqpoint{3.061999in}{2.113098in}}%
\pgfpathlineto{\pgfqpoint{3.064518in}{2.117867in}}%
\pgfpathlineto{\pgfqpoint{3.067037in}{2.100407in}}%
\pgfpathlineto{\pgfqpoint{3.069556in}{2.096329in}}%
\pgfpathlineto{\pgfqpoint{3.072074in}{2.107022in}}%
\pgfpathlineto{\pgfqpoint{3.074593in}{2.100426in}}%
\pgfpathlineto{\pgfqpoint{3.079631in}{2.081219in}}%
\pgfpathlineto{\pgfqpoint{3.082149in}{2.076941in}}%
\pgfpathlineto{\pgfqpoint{3.084668in}{2.040683in}}%
\pgfpathlineto{\pgfqpoint{3.087187in}{2.029099in}}%
\pgfpathlineto{\pgfqpoint{3.089706in}{2.033303in}}%
\pgfpathlineto{\pgfqpoint{3.092224in}{2.033384in}}%
\pgfpathlineto{\pgfqpoint{3.094743in}{2.033603in}}%
\pgfpathlineto{\pgfqpoint{3.097262in}{2.039940in}}%
\pgfpathlineto{\pgfqpoint{3.099781in}{2.072694in}}%
\pgfpathlineto{\pgfqpoint{3.102299in}{2.090158in}}%
\pgfpathlineto{\pgfqpoint{3.104818in}{2.069739in}}%
\pgfpathlineto{\pgfqpoint{3.107337in}{2.062288in}}%
\pgfpathlineto{\pgfqpoint{3.109856in}{2.064315in}}%
\pgfpathlineto{\pgfqpoint{3.112374in}{2.058235in}}%
\pgfpathlineto{\pgfqpoint{3.114893in}{2.057224in}}%
\pgfpathlineto{\pgfqpoint{3.117412in}{2.105326in}}%
\pgfpathlineto{\pgfqpoint{3.119931in}{2.105110in}}%
\pgfpathlineto{\pgfqpoint{3.122449in}{2.088840in}}%
\pgfpathlineto{\pgfqpoint{3.124968in}{2.094390in}}%
\pgfpathlineto{\pgfqpoint{3.127487in}{2.087902in}}%
\pgfpathlineto{\pgfqpoint{3.130006in}{2.086664in}}%
\pgfpathlineto{\pgfqpoint{3.132524in}{2.067292in}}%
\pgfpathlineto{\pgfqpoint{3.135043in}{2.081776in}}%
\pgfpathlineto{\pgfqpoint{3.137562in}{2.095706in}}%
\pgfpathlineto{\pgfqpoint{3.140081in}{2.101257in}}%
\pgfpathlineto{\pgfqpoint{3.142599in}{2.139643in}}%
\pgfpathlineto{\pgfqpoint{3.145118in}{2.108264in}}%
\pgfpathlineto{\pgfqpoint{3.147637in}{2.128676in}}%
\pgfpathlineto{\pgfqpoint{3.150156in}{2.121195in}}%
\pgfpathlineto{\pgfqpoint{3.152674in}{2.144079in}}%
\pgfpathlineto{\pgfqpoint{3.155193in}{2.160567in}}%
\pgfpathlineto{\pgfqpoint{3.157712in}{2.164450in}}%
\pgfpathlineto{\pgfqpoint{3.160231in}{2.157620in}}%
\pgfpathlineto{\pgfqpoint{3.162749in}{2.169934in}}%
\pgfpathlineto{\pgfqpoint{3.165268in}{2.155853in}}%
\pgfpathlineto{\pgfqpoint{3.167787in}{2.167558in}}%
\pgfpathlineto{\pgfqpoint{3.170306in}{2.176901in}}%
\pgfpathlineto{\pgfqpoint{3.172824in}{2.173311in}}%
\pgfpathlineto{\pgfqpoint{3.175343in}{2.169997in}}%
\pgfpathlineto{\pgfqpoint{3.177862in}{2.167819in}}%
\pgfpathlineto{\pgfqpoint{3.180381in}{2.172730in}}%
\pgfpathlineto{\pgfqpoint{3.182899in}{2.144832in}}%
\pgfpathlineto{\pgfqpoint{3.185418in}{2.129827in}}%
\pgfpathlineto{\pgfqpoint{3.187937in}{2.148933in}}%
\pgfpathlineto{\pgfqpoint{3.190456in}{2.157324in}}%
\pgfpathlineto{\pgfqpoint{3.195493in}{2.192282in}}%
\pgfpathlineto{\pgfqpoint{3.198012in}{2.202078in}}%
\pgfpathlineto{\pgfqpoint{3.200531in}{2.188337in}}%
\pgfpathlineto{\pgfqpoint{3.203049in}{2.204666in}}%
\pgfpathlineto{\pgfqpoint{3.205568in}{2.233085in}}%
\pgfpathlineto{\pgfqpoint{3.208087in}{2.209381in}}%
\pgfpathlineto{\pgfqpoint{3.210606in}{2.199424in}}%
\pgfpathlineto{\pgfqpoint{3.213124in}{2.182160in}}%
\pgfpathlineto{\pgfqpoint{3.215643in}{2.178570in}}%
\pgfpathlineto{\pgfqpoint{3.218162in}{2.170160in}}%
\pgfpathlineto{\pgfqpoint{3.220681in}{2.189444in}}%
\pgfpathlineto{\pgfqpoint{3.223199in}{2.210065in}}%
\pgfpathlineto{\pgfqpoint{3.225718in}{2.213932in}}%
\pgfpathlineto{\pgfqpoint{3.228237in}{2.246865in}}%
\pgfpathlineto{\pgfqpoint{3.230756in}{2.243973in}}%
\pgfpathlineto{\pgfqpoint{3.233274in}{2.231422in}}%
\pgfpathlineto{\pgfqpoint{3.235793in}{2.208510in}}%
\pgfpathlineto{\pgfqpoint{3.238312in}{2.198272in}}%
\pgfpathlineto{\pgfqpoint{3.240831in}{2.217596in}}%
\pgfpathlineto{\pgfqpoint{3.243349in}{2.227738in}}%
\pgfpathlineto{\pgfqpoint{3.245868in}{2.195190in}}%
\pgfpathlineto{\pgfqpoint{3.248387in}{2.229136in}}%
\pgfpathlineto{\pgfqpoint{3.250906in}{2.241940in}}%
\pgfpathlineto{\pgfqpoint{3.253424in}{2.244289in}}%
\pgfpathlineto{\pgfqpoint{3.255943in}{2.240807in}}%
\pgfpathlineto{\pgfqpoint{3.258462in}{2.244469in}}%
\pgfpathlineto{\pgfqpoint{3.260981in}{2.265948in}}%
\pgfpathlineto{\pgfqpoint{3.263499in}{2.261827in}}%
\pgfpathlineto{\pgfqpoint{3.266018in}{2.269377in}}%
\pgfpathlineto{\pgfqpoint{3.268537in}{2.234706in}}%
\pgfpathlineto{\pgfqpoint{3.271056in}{2.217364in}}%
\pgfpathlineto{\pgfqpoint{3.273574in}{2.212685in}}%
\pgfpathlineto{\pgfqpoint{3.276093in}{2.204575in}}%
\pgfpathlineto{\pgfqpoint{3.278612in}{2.209340in}}%
\pgfpathlineto{\pgfqpoint{3.281131in}{2.215365in}}%
\pgfpathlineto{\pgfqpoint{3.283649in}{2.207759in}}%
\pgfpathlineto{\pgfqpoint{3.286168in}{2.195370in}}%
\pgfpathlineto{\pgfqpoint{3.288687in}{2.202241in}}%
\pgfpathlineto{\pgfqpoint{3.291206in}{2.180328in}}%
\pgfpathlineto{\pgfqpoint{3.293724in}{2.178997in}}%
\pgfpathlineto{\pgfqpoint{3.296243in}{2.183509in}}%
\pgfpathlineto{\pgfqpoint{3.298762in}{2.217005in}}%
\pgfpathlineto{\pgfqpoint{3.301281in}{2.212809in}}%
\pgfpathlineto{\pgfqpoint{3.303799in}{2.178604in}}%
\pgfpathlineto{\pgfqpoint{3.306318in}{2.201320in}}%
\pgfpathlineto{\pgfqpoint{3.308837in}{2.194776in}}%
\pgfpathlineto{\pgfqpoint{3.311356in}{2.221225in}}%
\pgfpathlineto{\pgfqpoint{3.313874in}{2.243756in}}%
\pgfpathlineto{\pgfqpoint{3.316393in}{2.254032in}}%
\pgfpathlineto{\pgfqpoint{3.318912in}{2.234007in}}%
\pgfpathlineto{\pgfqpoint{3.321431in}{2.240903in}}%
\pgfpathlineto{\pgfqpoint{3.323949in}{2.266674in}}%
\pgfpathlineto{\pgfqpoint{3.326468in}{2.273684in}}%
\pgfpathlineto{\pgfqpoint{3.328987in}{2.282001in}}%
\pgfpathlineto{\pgfqpoint{3.331506in}{2.308002in}}%
\pgfpathlineto{\pgfqpoint{3.334024in}{2.329488in}}%
\pgfpathlineto{\pgfqpoint{3.336543in}{2.306785in}}%
\pgfpathlineto{\pgfqpoint{3.339062in}{2.308757in}}%
\pgfpathlineto{\pgfqpoint{3.341581in}{2.305656in}}%
\pgfpathlineto{\pgfqpoint{3.344099in}{2.293296in}}%
\pgfpathlineto{\pgfqpoint{3.346618in}{2.288283in}}%
\pgfpathlineto{\pgfqpoint{3.349137in}{2.299674in}}%
\pgfpathlineto{\pgfqpoint{3.351656in}{2.287988in}}%
\pgfpathlineto{\pgfqpoint{3.354174in}{2.283625in}}%
\pgfpathlineto{\pgfqpoint{3.356693in}{2.292420in}}%
\pgfpathlineto{\pgfqpoint{3.359212in}{2.295566in}}%
\pgfpathlineto{\pgfqpoint{3.361731in}{2.316032in}}%
\pgfpathlineto{\pgfqpoint{3.364249in}{2.333987in}}%
\pgfpathlineto{\pgfqpoint{3.366768in}{2.344017in}}%
\pgfpathlineto{\pgfqpoint{3.369287in}{2.357804in}}%
\pgfpathlineto{\pgfqpoint{3.371806in}{2.380580in}}%
\pgfpathlineto{\pgfqpoint{3.374324in}{2.388399in}}%
\pgfpathlineto{\pgfqpoint{3.376843in}{2.378539in}}%
\pgfpathlineto{\pgfqpoint{3.379362in}{2.382360in}}%
\pgfpathlineto{\pgfqpoint{3.381881in}{2.354140in}}%
\pgfpathlineto{\pgfqpoint{3.384399in}{2.339889in}}%
\pgfpathlineto{\pgfqpoint{3.386918in}{2.324712in}}%
\pgfpathlineto{\pgfqpoint{3.389437in}{2.321868in}}%
\pgfpathlineto{\pgfqpoint{3.391956in}{2.334473in}}%
\pgfpathlineto{\pgfqpoint{3.394474in}{2.332274in}}%
\pgfpathlineto{\pgfqpoint{3.396993in}{2.324436in}}%
\pgfpathlineto{\pgfqpoint{3.399512in}{2.323024in}}%
\pgfpathlineto{\pgfqpoint{3.402031in}{2.330697in}}%
\pgfpathlineto{\pgfqpoint{3.404549in}{2.351132in}}%
\pgfpathlineto{\pgfqpoint{3.407068in}{2.404374in}}%
\pgfpathlineto{\pgfqpoint{3.409587in}{2.416562in}}%
\pgfpathlineto{\pgfqpoint{3.412106in}{2.415825in}}%
\pgfpathlineto{\pgfqpoint{3.414624in}{2.392825in}}%
\pgfpathlineto{\pgfqpoint{3.417143in}{2.379091in}}%
\pgfpathlineto{\pgfqpoint{3.419662in}{2.376771in}}%
\pgfpathlineto{\pgfqpoint{3.422181in}{2.410111in}}%
\pgfpathlineto{\pgfqpoint{3.424699in}{2.404562in}}%
\pgfpathlineto{\pgfqpoint{3.427218in}{2.414556in}}%
\pgfpathlineto{\pgfqpoint{3.429737in}{2.434710in}}%
\pgfpathlineto{\pgfqpoint{3.432256in}{2.453473in}}%
\pgfpathlineto{\pgfqpoint{3.434774in}{2.421522in}}%
\pgfpathlineto{\pgfqpoint{3.437293in}{2.399407in}}%
\pgfpathlineto{\pgfqpoint{3.439812in}{2.417258in}}%
\pgfpathlineto{\pgfqpoint{3.442331in}{2.414912in}}%
\pgfpathlineto{\pgfqpoint{3.444849in}{2.422660in}}%
\pgfpathlineto{\pgfqpoint{3.447368in}{2.413990in}}%
\pgfpathlineto{\pgfqpoint{3.449887in}{2.410138in}}%
\pgfpathlineto{\pgfqpoint{3.452406in}{2.371495in}}%
\pgfpathlineto{\pgfqpoint{3.454924in}{2.338659in}}%
\pgfpathlineto{\pgfqpoint{3.457443in}{2.317421in}}%
\pgfpathlineto{\pgfqpoint{3.459962in}{2.307626in}}%
\pgfpathlineto{\pgfqpoint{3.464999in}{2.245359in}}%
\pgfpathlineto{\pgfqpoint{3.467518in}{2.212613in}}%
\pgfpathlineto{\pgfqpoint{3.470037in}{2.193581in}}%
\pgfpathlineto{\pgfqpoint{3.472556in}{2.204165in}}%
\pgfpathlineto{\pgfqpoint{3.475074in}{2.234141in}}%
\pgfpathlineto{\pgfqpoint{3.477593in}{2.245456in}}%
\pgfpathlineto{\pgfqpoint{3.480112in}{2.240863in}}%
\pgfpathlineto{\pgfqpoint{3.482631in}{2.248166in}}%
\pgfpathlineto{\pgfqpoint{3.485149in}{2.228456in}}%
\pgfpathlineto{\pgfqpoint{3.487668in}{2.234497in}}%
\pgfpathlineto{\pgfqpoint{3.490187in}{2.250106in}}%
\pgfpathlineto{\pgfqpoint{3.492706in}{2.205148in}}%
\pgfpathlineto{\pgfqpoint{3.495224in}{2.224915in}}%
\pgfpathlineto{\pgfqpoint{3.497743in}{2.197461in}}%
\pgfpathlineto{\pgfqpoint{3.500262in}{2.171559in}}%
\pgfpathlineto{\pgfqpoint{3.502781in}{2.177795in}}%
\pgfpathlineto{\pgfqpoint{3.505299in}{2.171584in}}%
\pgfpathlineto{\pgfqpoint{3.507818in}{2.189640in}}%
\pgfpathlineto{\pgfqpoint{3.510337in}{2.187882in}}%
\pgfpathlineto{\pgfqpoint{3.512856in}{2.162899in}}%
\pgfpathlineto{\pgfqpoint{3.515374in}{2.173594in}}%
\pgfpathlineto{\pgfqpoint{3.517893in}{2.159607in}}%
\pgfpathlineto{\pgfqpoint{3.520412in}{2.161640in}}%
\pgfpathlineto{\pgfqpoint{3.522931in}{2.150703in}}%
\pgfpathlineto{\pgfqpoint{3.525449in}{2.145225in}}%
\pgfpathlineto{\pgfqpoint{3.527968in}{2.142393in}}%
\pgfpathlineto{\pgfqpoint{3.530487in}{2.151399in}}%
\pgfpathlineto{\pgfqpoint{3.533006in}{2.185036in}}%
\pgfpathlineto{\pgfqpoint{3.535524in}{2.209755in}}%
\pgfpathlineto{\pgfqpoint{3.538043in}{2.211948in}}%
\pgfpathlineto{\pgfqpoint{3.540562in}{2.211665in}}%
\pgfpathlineto{\pgfqpoint{3.543081in}{2.202448in}}%
\pgfpathlineto{\pgfqpoint{3.545599in}{2.210111in}}%
\pgfpathlineto{\pgfqpoint{3.548118in}{2.224193in}}%
\pgfpathlineto{\pgfqpoint{3.550637in}{2.220359in}}%
\pgfpathlineto{\pgfqpoint{3.553156in}{2.247595in}}%
\pgfpathlineto{\pgfqpoint{3.555674in}{2.225817in}}%
\pgfpathlineto{\pgfqpoint{3.558193in}{2.240338in}}%
\pgfpathlineto{\pgfqpoint{3.560712in}{2.237822in}}%
\pgfpathlineto{\pgfqpoint{3.563231in}{2.257549in}}%
\pgfpathlineto{\pgfqpoint{3.565749in}{2.281231in}}%
\pgfpathlineto{\pgfqpoint{3.568268in}{2.286816in}}%
\pgfpathlineto{\pgfqpoint{3.570787in}{2.292811in}}%
\pgfpathlineto{\pgfqpoint{3.573306in}{2.295072in}}%
\pgfpathlineto{\pgfqpoint{3.575824in}{2.309612in}}%
\pgfpathlineto{\pgfqpoint{3.578343in}{2.327762in}}%
\pgfpathlineto{\pgfqpoint{3.580862in}{2.331391in}}%
\pgfpathlineto{\pgfqpoint{3.583381in}{2.323610in}}%
\pgfpathlineto{\pgfqpoint{3.585899in}{2.302899in}}%
\pgfpathlineto{\pgfqpoint{3.588418in}{2.287541in}}%
\pgfpathlineto{\pgfqpoint{3.590937in}{2.294880in}}%
\pgfpathlineto{\pgfqpoint{3.593456in}{2.276023in}}%
\pgfpathlineto{\pgfqpoint{3.595974in}{2.272620in}}%
\pgfpathlineto{\pgfqpoint{3.598493in}{2.254275in}}%
\pgfpathlineto{\pgfqpoint{3.601012in}{2.287335in}}%
\pgfpathlineto{\pgfqpoint{3.603531in}{2.283288in}}%
\pgfpathlineto{\pgfqpoint{3.606049in}{2.289177in}}%
\pgfpathlineto{\pgfqpoint{3.608568in}{2.299027in}}%
\pgfpathlineto{\pgfqpoint{3.611087in}{2.304392in}}%
\pgfpathlineto{\pgfqpoint{3.613606in}{2.303767in}}%
\pgfpathlineto{\pgfqpoint{3.616124in}{2.330622in}}%
\pgfpathlineto{\pgfqpoint{3.618643in}{2.318043in}}%
\pgfpathlineto{\pgfqpoint{3.621162in}{2.318298in}}%
\pgfpathlineto{\pgfqpoint{3.623681in}{2.304564in}}%
\pgfpathlineto{\pgfqpoint{3.626199in}{2.295335in}}%
\pgfpathlineto{\pgfqpoint{3.628718in}{2.295661in}}%
\pgfpathlineto{\pgfqpoint{3.631237in}{2.286042in}}%
\pgfpathlineto{\pgfqpoint{3.633756in}{2.293828in}}%
\pgfpathlineto{\pgfqpoint{3.636274in}{2.267404in}}%
\pgfpathlineto{\pgfqpoint{3.638793in}{2.271550in}}%
\pgfpathlineto{\pgfqpoint{3.646349in}{2.222155in}}%
\pgfpathlineto{\pgfqpoint{3.648868in}{2.219380in}}%
\pgfpathlineto{\pgfqpoint{3.651387in}{2.262101in}}%
\pgfpathlineto{\pgfqpoint{3.653906in}{2.246752in}}%
\pgfpathlineto{\pgfqpoint{3.656424in}{2.258928in}}%
\pgfpathlineto{\pgfqpoint{3.658943in}{2.278431in}}%
\pgfpathlineto{\pgfqpoint{3.661462in}{2.258641in}}%
\pgfpathlineto{\pgfqpoint{3.663981in}{2.283792in}}%
\pgfpathlineto{\pgfqpoint{3.669018in}{2.259420in}}%
\pgfpathlineto{\pgfqpoint{3.671537in}{2.240775in}}%
\pgfpathlineto{\pgfqpoint{3.674056in}{2.235212in}}%
\pgfpathlineto{\pgfqpoint{3.676574in}{2.206037in}}%
\pgfpathlineto{\pgfqpoint{3.679093in}{2.189642in}}%
\pgfpathlineto{\pgfqpoint{3.681612in}{2.181054in}}%
\pgfpathlineto{\pgfqpoint{3.684131in}{2.196501in}}%
\pgfpathlineto{\pgfqpoint{3.686649in}{2.199034in}}%
\pgfpathlineto{\pgfqpoint{3.689168in}{2.176792in}}%
\pgfpathlineto{\pgfqpoint{3.691687in}{2.180475in}}%
\pgfpathlineto{\pgfqpoint{3.694206in}{2.199947in}}%
\pgfpathlineto{\pgfqpoint{3.696724in}{2.216214in}}%
\pgfpathlineto{\pgfqpoint{3.699243in}{2.238437in}}%
\pgfpathlineto{\pgfqpoint{3.701762in}{2.252006in}}%
\pgfpathlineto{\pgfqpoint{3.704281in}{2.236360in}}%
\pgfpathlineto{\pgfqpoint{3.706799in}{2.244776in}}%
\pgfpathlineto{\pgfqpoint{3.709318in}{2.228287in}}%
\pgfpathlineto{\pgfqpoint{3.711837in}{2.195257in}}%
\pgfpathlineto{\pgfqpoint{3.714356in}{2.195273in}}%
\pgfpathlineto{\pgfqpoint{3.716874in}{2.183354in}}%
\pgfpathlineto{\pgfqpoint{3.719393in}{2.165276in}}%
\pgfpathlineto{\pgfqpoint{3.721912in}{2.161146in}}%
\pgfpathlineto{\pgfqpoint{3.724431in}{2.194025in}}%
\pgfpathlineto{\pgfqpoint{3.726949in}{2.193684in}}%
\pgfpathlineto{\pgfqpoint{3.729468in}{2.165801in}}%
\pgfpathlineto{\pgfqpoint{3.731987in}{2.144427in}}%
\pgfpathlineto{\pgfqpoint{3.734506in}{2.140042in}}%
\pgfpathlineto{\pgfqpoint{3.737024in}{2.191150in}}%
\pgfpathlineto{\pgfqpoint{3.739543in}{2.197410in}}%
\pgfpathlineto{\pgfqpoint{3.742062in}{2.213891in}}%
\pgfpathlineto{\pgfqpoint{3.744581in}{2.213131in}}%
\pgfpathlineto{\pgfqpoint{3.747099in}{2.186222in}}%
\pgfpathlineto{\pgfqpoint{3.749618in}{2.195368in}}%
\pgfpathlineto{\pgfqpoint{3.752137in}{2.183371in}}%
\pgfpathlineto{\pgfqpoint{3.754656in}{2.160076in}}%
\pgfpathlineto{\pgfqpoint{3.757174in}{2.121616in}}%
\pgfpathlineto{\pgfqpoint{3.759693in}{2.148184in}}%
\pgfpathlineto{\pgfqpoint{3.762212in}{2.133145in}}%
\pgfpathlineto{\pgfqpoint{3.764731in}{2.101324in}}%
\pgfpathlineto{\pgfqpoint{3.767249in}{2.109035in}}%
\pgfpathlineto{\pgfqpoint{3.769768in}{2.080160in}}%
\pgfpathlineto{\pgfqpoint{3.772287in}{2.109226in}}%
\pgfpathlineto{\pgfqpoint{3.774806in}{2.113711in}}%
\pgfpathlineto{\pgfqpoint{3.777324in}{2.116906in}}%
\pgfpathlineto{\pgfqpoint{3.779843in}{2.122822in}}%
\pgfpathlineto{\pgfqpoint{3.782362in}{2.136786in}}%
\pgfpathlineto{\pgfqpoint{3.784881in}{2.140551in}}%
\pgfpathlineto{\pgfqpoint{3.787399in}{2.143396in}}%
\pgfpathlineto{\pgfqpoint{3.789918in}{2.142178in}}%
\pgfpathlineto{\pgfqpoint{3.792437in}{2.137637in}}%
\pgfpathlineto{\pgfqpoint{3.794956in}{2.118909in}}%
\pgfpathlineto{\pgfqpoint{3.797474in}{2.135049in}}%
\pgfpathlineto{\pgfqpoint{3.799993in}{2.129495in}}%
\pgfpathlineto{\pgfqpoint{3.802512in}{2.137512in}}%
\pgfpathlineto{\pgfqpoint{3.805031in}{2.138136in}}%
\pgfpathlineto{\pgfqpoint{3.807549in}{2.188002in}}%
\pgfpathlineto{\pgfqpoint{3.810068in}{2.201757in}}%
\pgfpathlineto{\pgfqpoint{3.812587in}{2.212875in}}%
\pgfpathlineto{\pgfqpoint{3.817624in}{2.210250in}}%
\pgfpathlineto{\pgfqpoint{3.820143in}{2.215453in}}%
\pgfpathlineto{\pgfqpoint{3.822662in}{2.212160in}}%
\pgfpathlineto{\pgfqpoint{3.825181in}{2.196879in}}%
\pgfpathlineto{\pgfqpoint{3.827699in}{2.203320in}}%
\pgfpathlineto{\pgfqpoint{3.830218in}{2.201023in}}%
\pgfpathlineto{\pgfqpoint{3.832737in}{2.167830in}}%
\pgfpathlineto{\pgfqpoint{3.835256in}{2.162670in}}%
\pgfpathlineto{\pgfqpoint{3.837774in}{2.160736in}}%
\pgfpathlineto{\pgfqpoint{3.840293in}{2.151834in}}%
\pgfpathlineto{\pgfqpoint{3.842812in}{2.119933in}}%
\pgfpathlineto{\pgfqpoint{3.845331in}{2.120118in}}%
\pgfpathlineto{\pgfqpoint{3.847849in}{2.084614in}}%
\pgfpathlineto{\pgfqpoint{3.850368in}{2.124100in}}%
\pgfpathlineto{\pgfqpoint{3.852887in}{2.128479in}}%
\pgfpathlineto{\pgfqpoint{3.855406in}{2.152788in}}%
\pgfpathlineto{\pgfqpoint{3.857924in}{2.182629in}}%
\pgfpathlineto{\pgfqpoint{3.860443in}{2.187840in}}%
\pgfpathlineto{\pgfqpoint{3.862962in}{2.172224in}}%
\pgfpathlineto{\pgfqpoint{3.865481in}{2.176916in}}%
\pgfpathlineto{\pgfqpoint{3.867999in}{2.177485in}}%
\pgfpathlineto{\pgfqpoint{3.870518in}{2.200515in}}%
\pgfpathlineto{\pgfqpoint{3.873037in}{2.222684in}}%
\pgfpathlineto{\pgfqpoint{3.875556in}{2.198710in}}%
\pgfpathlineto{\pgfqpoint{3.878074in}{2.201659in}}%
\pgfpathlineto{\pgfqpoint{3.880593in}{2.208718in}}%
\pgfpathlineto{\pgfqpoint{3.883112in}{2.201618in}}%
\pgfpathlineto{\pgfqpoint{3.885631in}{2.233726in}}%
\pgfpathlineto{\pgfqpoint{3.888149in}{2.231903in}}%
\pgfpathlineto{\pgfqpoint{3.890668in}{2.208515in}}%
\pgfpathlineto{\pgfqpoint{3.893187in}{2.191957in}}%
\pgfpathlineto{\pgfqpoint{3.895706in}{2.215154in}}%
\pgfpathlineto{\pgfqpoint{3.898224in}{2.220077in}}%
\pgfpathlineto{\pgfqpoint{3.900743in}{2.250746in}}%
\pgfpathlineto{\pgfqpoint{3.903262in}{2.243614in}}%
\pgfpathlineto{\pgfqpoint{3.905781in}{2.253317in}}%
\pgfpathlineto{\pgfqpoint{3.908299in}{2.254705in}}%
\pgfpathlineto{\pgfqpoint{3.910818in}{2.301926in}}%
\pgfpathlineto{\pgfqpoint{3.913337in}{2.311241in}}%
\pgfpathlineto{\pgfqpoint{3.915856in}{2.277582in}}%
\pgfpathlineto{\pgfqpoint{3.918374in}{2.295975in}}%
\pgfpathlineto{\pgfqpoint{3.920893in}{2.295718in}}%
\pgfpathlineto{\pgfqpoint{3.923412in}{2.297974in}}%
\pgfpathlineto{\pgfqpoint{3.925931in}{2.326237in}}%
\pgfpathlineto{\pgfqpoint{3.928449in}{2.328920in}}%
\pgfpathlineto{\pgfqpoint{3.930968in}{2.339366in}}%
\pgfpathlineto{\pgfqpoint{3.933487in}{2.314716in}}%
\pgfpathlineto{\pgfqpoint{3.936006in}{2.318702in}}%
\pgfpathlineto{\pgfqpoint{3.938524in}{2.326947in}}%
\pgfpathlineto{\pgfqpoint{3.941043in}{2.343664in}}%
\pgfpathlineto{\pgfqpoint{3.943562in}{2.352740in}}%
\pgfpathlineto{\pgfqpoint{3.946081in}{2.343434in}}%
\pgfpathlineto{\pgfqpoint{3.948599in}{2.339398in}}%
\pgfpathlineto{\pgfqpoint{3.951118in}{2.348893in}}%
\pgfpathlineto{\pgfqpoint{3.953637in}{2.348778in}}%
\pgfpathlineto{\pgfqpoint{3.956156in}{2.327201in}}%
\pgfpathlineto{\pgfqpoint{3.958674in}{2.329241in}}%
\pgfpathlineto{\pgfqpoint{3.961193in}{2.334807in}}%
\pgfpathlineto{\pgfqpoint{3.963712in}{2.298690in}}%
\pgfpathlineto{\pgfqpoint{3.966231in}{2.324925in}}%
\pgfpathlineto{\pgfqpoint{3.968749in}{2.320379in}}%
\pgfpathlineto{\pgfqpoint{3.971268in}{2.313168in}}%
\pgfpathlineto{\pgfqpoint{3.973787in}{2.312598in}}%
\pgfpathlineto{\pgfqpoint{3.976306in}{2.274402in}}%
\pgfpathlineto{\pgfqpoint{3.978824in}{2.263823in}}%
\pgfpathlineto{\pgfqpoint{3.981343in}{2.263499in}}%
\pgfpathlineto{\pgfqpoint{3.983862in}{2.224943in}}%
\pgfpathlineto{\pgfqpoint{3.986381in}{2.224957in}}%
\pgfpathlineto{\pgfqpoint{3.988899in}{2.248439in}}%
\pgfpathlineto{\pgfqpoint{3.991418in}{2.260854in}}%
\pgfpathlineto{\pgfqpoint{3.993937in}{2.292065in}}%
\pgfpathlineto{\pgfqpoint{3.996456in}{2.273807in}}%
\pgfpathlineto{\pgfqpoint{3.998974in}{2.262414in}}%
\pgfpathlineto{\pgfqpoint{4.001493in}{2.269889in}}%
\pgfpathlineto{\pgfqpoint{4.004012in}{2.243419in}}%
\pgfpathlineto{\pgfqpoint{4.006531in}{2.258991in}}%
\pgfpathlineto{\pgfqpoint{4.009049in}{2.248244in}}%
\pgfpathlineto{\pgfqpoint{4.011568in}{2.285145in}}%
\pgfpathlineto{\pgfqpoint{4.014087in}{2.316408in}}%
\pgfpathlineto{\pgfqpoint{4.016606in}{2.310009in}}%
\pgfpathlineto{\pgfqpoint{4.019124in}{2.327356in}}%
\pgfpathlineto{\pgfqpoint{4.021643in}{2.360671in}}%
\pgfpathlineto{\pgfqpoint{4.024162in}{2.362759in}}%
\pgfpathlineto{\pgfqpoint{4.026681in}{2.364599in}}%
\pgfpathlineto{\pgfqpoint{4.029199in}{2.362357in}}%
\pgfpathlineto{\pgfqpoint{4.031718in}{2.379039in}}%
\pgfpathlineto{\pgfqpoint{4.034237in}{2.394533in}}%
\pgfpathlineto{\pgfqpoint{4.036756in}{2.394214in}}%
\pgfpathlineto{\pgfqpoint{4.039274in}{2.395678in}}%
\pgfpathlineto{\pgfqpoint{4.041793in}{2.358948in}}%
\pgfpathlineto{\pgfqpoint{4.044312in}{2.353144in}}%
\pgfpathlineto{\pgfqpoint{4.046831in}{2.391219in}}%
\pgfpathlineto{\pgfqpoint{4.049349in}{2.411958in}}%
\pgfpathlineto{\pgfqpoint{4.051868in}{2.438109in}}%
\pgfpathlineto{\pgfqpoint{4.054387in}{2.472923in}}%
\pgfpathlineto{\pgfqpoint{4.056906in}{2.485752in}}%
\pgfpathlineto{\pgfqpoint{4.059424in}{2.458961in}}%
\pgfpathlineto{\pgfqpoint{4.061943in}{2.461553in}}%
\pgfpathlineto{\pgfqpoint{4.064462in}{2.445748in}}%
\pgfpathlineto{\pgfqpoint{4.066981in}{2.476963in}}%
\pgfpathlineto{\pgfqpoint{4.072018in}{2.465401in}}%
\pgfpathlineto{\pgfqpoint{4.074537in}{2.428975in}}%
\pgfpathlineto{\pgfqpoint{4.077056in}{2.448966in}}%
\pgfpathlineto{\pgfqpoint{4.079574in}{2.432253in}}%
\pgfpathlineto{\pgfqpoint{4.082093in}{2.413702in}}%
\pgfpathlineto{\pgfqpoint{4.084612in}{2.434120in}}%
\pgfpathlineto{\pgfqpoint{4.087131in}{2.449179in}}%
\pgfpathlineto{\pgfqpoint{4.089649in}{2.420067in}}%
\pgfpathlineto{\pgfqpoint{4.094687in}{2.460147in}}%
\pgfpathlineto{\pgfqpoint{4.097206in}{2.433809in}}%
\pgfpathlineto{\pgfqpoint{4.099724in}{2.428439in}}%
\pgfpathlineto{\pgfqpoint{4.102243in}{2.424344in}}%
\pgfpathlineto{\pgfqpoint{4.104762in}{2.416509in}}%
\pgfpathlineto{\pgfqpoint{4.107281in}{2.412449in}}%
\pgfpathlineto{\pgfqpoint{4.109799in}{2.415075in}}%
\pgfpathlineto{\pgfqpoint{4.112318in}{2.419090in}}%
\pgfpathlineto{\pgfqpoint{4.114837in}{2.422144in}}%
\pgfpathlineto{\pgfqpoint{4.117356in}{2.414807in}}%
\pgfpathlineto{\pgfqpoint{4.119874in}{2.404619in}}%
\pgfpathlineto{\pgfqpoint{4.122393in}{2.436674in}}%
\pgfpathlineto{\pgfqpoint{4.124912in}{2.483731in}}%
\pgfpathlineto{\pgfqpoint{4.127431in}{2.478840in}}%
\pgfpathlineto{\pgfqpoint{4.129949in}{2.535769in}}%
\pgfpathlineto{\pgfqpoint{4.132468in}{2.574868in}}%
\pgfpathlineto{\pgfqpoint{4.134987in}{2.584329in}}%
\pgfpathlineto{\pgfqpoint{4.137506in}{2.610545in}}%
\pgfpathlineto{\pgfqpoint{4.140024in}{2.635250in}}%
\pgfpathlineto{\pgfqpoint{4.142543in}{2.615930in}}%
\pgfpathlineto{\pgfqpoint{4.145062in}{2.606803in}}%
\pgfpathlineto{\pgfqpoint{4.147581in}{2.637830in}}%
\pgfpathlineto{\pgfqpoint{4.150099in}{2.644762in}}%
\pgfpathlineto{\pgfqpoint{4.152618in}{2.669723in}}%
\pgfpathlineto{\pgfqpoint{4.155137in}{2.643694in}}%
\pgfpathlineto{\pgfqpoint{4.157656in}{2.648615in}}%
\pgfpathlineto{\pgfqpoint{4.160174in}{2.667255in}}%
\pgfpathlineto{\pgfqpoint{4.162693in}{2.674408in}}%
\pgfpathlineto{\pgfqpoint{4.165212in}{2.683119in}}%
\pgfpathlineto{\pgfqpoint{4.167731in}{2.700179in}}%
\pgfpathlineto{\pgfqpoint{4.170249in}{2.724543in}}%
\pgfpathlineto{\pgfqpoint{4.172768in}{2.718209in}}%
\pgfpathlineto{\pgfqpoint{4.175287in}{2.722272in}}%
\pgfpathlineto{\pgfqpoint{4.177806in}{2.695353in}}%
\pgfpathlineto{\pgfqpoint{4.180324in}{2.636869in}}%
\pgfpathlineto{\pgfqpoint{4.182843in}{2.632465in}}%
\pgfpathlineto{\pgfqpoint{4.185362in}{2.642334in}}%
\pgfpathlineto{\pgfqpoint{4.187881in}{2.656049in}}%
\pgfpathlineto{\pgfqpoint{4.190399in}{2.709365in}}%
\pgfpathlineto{\pgfqpoint{4.192918in}{2.696268in}}%
\pgfpathlineto{\pgfqpoint{4.195437in}{2.690947in}}%
\pgfpathlineto{\pgfqpoint{4.197956in}{2.669209in}}%
\pgfpathlineto{\pgfqpoint{4.200474in}{2.662620in}}%
\pgfpathlineto{\pgfqpoint{4.202993in}{2.689525in}}%
\pgfpathlineto{\pgfqpoint{4.205512in}{2.699724in}}%
\pgfpathlineto{\pgfqpoint{4.208031in}{2.708069in}}%
\pgfpathlineto{\pgfqpoint{4.210549in}{2.730900in}}%
\pgfpathlineto{\pgfqpoint{4.213068in}{2.713424in}}%
\pgfpathlineto{\pgfqpoint{4.215587in}{2.681174in}}%
\pgfpathlineto{\pgfqpoint{4.218106in}{2.708150in}}%
\pgfpathlineto{\pgfqpoint{4.220624in}{2.693173in}}%
\pgfpathlineto{\pgfqpoint{4.223143in}{2.693000in}}%
\pgfpathlineto{\pgfqpoint{4.225662in}{2.680600in}}%
\pgfpathlineto{\pgfqpoint{4.228181in}{2.694535in}}%
\pgfpathlineto{\pgfqpoint{4.230699in}{2.687414in}}%
\pgfpathlineto{\pgfqpoint{4.233218in}{2.669343in}}%
\pgfpathlineto{\pgfqpoint{4.235737in}{2.701084in}}%
\pgfpathlineto{\pgfqpoint{4.238256in}{2.731357in}}%
\pgfpathlineto{\pgfqpoint{4.240774in}{2.744522in}}%
\pgfpathlineto{\pgfqpoint{4.243293in}{2.684976in}}%
\pgfpathlineto{\pgfqpoint{4.245812in}{2.633648in}}%
\pgfpathlineto{\pgfqpoint{4.248331in}{2.639486in}}%
\pgfpathlineto{\pgfqpoint{4.250849in}{2.611980in}}%
\pgfpathlineto{\pgfqpoint{4.253368in}{2.609346in}}%
\pgfpathlineto{\pgfqpoint{4.255887in}{2.611401in}}%
\pgfpathlineto{\pgfqpoint{4.258406in}{2.622964in}}%
\pgfpathlineto{\pgfqpoint{4.260924in}{2.638483in}}%
\pgfpathlineto{\pgfqpoint{4.263443in}{2.626211in}}%
\pgfpathlineto{\pgfqpoint{4.265962in}{2.632001in}}%
\pgfpathlineto{\pgfqpoint{4.268481in}{2.608079in}}%
\pgfpathlineto{\pgfqpoint{4.270999in}{2.575370in}}%
\pgfpathlineto{\pgfqpoint{4.273518in}{2.585047in}}%
\pgfpathlineto{\pgfqpoint{4.276037in}{2.545691in}}%
\pgfpathlineto{\pgfqpoint{4.278556in}{2.542454in}}%
\pgfpathlineto{\pgfqpoint{4.281074in}{2.570372in}}%
\pgfpathlineto{\pgfqpoint{4.283593in}{2.542733in}}%
\pgfpathlineto{\pgfqpoint{4.286112in}{2.531215in}}%
\pgfpathlineto{\pgfqpoint{4.288631in}{2.556899in}}%
\pgfpathlineto{\pgfqpoint{4.291149in}{2.527373in}}%
\pgfpathlineto{\pgfqpoint{4.293668in}{2.513054in}}%
\pgfpathlineto{\pgfqpoint{4.296187in}{2.475560in}}%
\pgfpathlineto{\pgfqpoint{4.298706in}{2.481948in}}%
\pgfpathlineto{\pgfqpoint{4.301224in}{2.483584in}}%
\pgfpathlineto{\pgfqpoint{4.303743in}{2.472057in}}%
\pgfpathlineto{\pgfqpoint{4.306262in}{2.476949in}}%
\pgfpathlineto{\pgfqpoint{4.308781in}{2.472860in}}%
\pgfpathlineto{\pgfqpoint{4.311299in}{2.456993in}}%
\pgfpathlineto{\pgfqpoint{4.313818in}{2.480565in}}%
\pgfpathlineto{\pgfqpoint{4.316337in}{2.486678in}}%
\pgfpathlineto{\pgfqpoint{4.318856in}{2.497725in}}%
\pgfpathlineto{\pgfqpoint{4.321374in}{2.520978in}}%
\pgfpathlineto{\pgfqpoint{4.323893in}{2.528821in}}%
\pgfpathlineto{\pgfqpoint{4.326412in}{2.528587in}}%
\pgfpathlineto{\pgfqpoint{4.328931in}{2.528858in}}%
\pgfpathlineto{\pgfqpoint{4.331449in}{2.527631in}}%
\pgfpathlineto{\pgfqpoint{4.333968in}{2.547311in}}%
\pgfpathlineto{\pgfqpoint{4.336487in}{2.532696in}}%
\pgfpathlineto{\pgfqpoint{4.339006in}{2.537524in}}%
\pgfpathlineto{\pgfqpoint{4.341524in}{2.510276in}}%
\pgfpathlineto{\pgfqpoint{4.344043in}{2.544790in}}%
\pgfpathlineto{\pgfqpoint{4.346562in}{2.555969in}}%
\pgfpathlineto{\pgfqpoint{4.349081in}{2.536948in}}%
\pgfpathlineto{\pgfqpoint{4.351599in}{2.523175in}}%
\pgfpathlineto{\pgfqpoint{4.354118in}{2.523697in}}%
\pgfpathlineto{\pgfqpoint{4.356637in}{2.532675in}}%
\pgfpathlineto{\pgfqpoint{4.359156in}{2.557866in}}%
\pgfpathlineto{\pgfqpoint{4.361674in}{2.543791in}}%
\pgfpathlineto{\pgfqpoint{4.364193in}{2.545137in}}%
\pgfpathlineto{\pgfqpoint{4.366712in}{2.544252in}}%
\pgfpathlineto{\pgfqpoint{4.371749in}{2.504741in}}%
\pgfpathlineto{\pgfqpoint{4.374268in}{2.504163in}}%
\pgfpathlineto{\pgfqpoint{4.376787in}{2.521524in}}%
\pgfpathlineto{\pgfqpoint{4.379306in}{2.494379in}}%
\pgfpathlineto{\pgfqpoint{4.381824in}{2.515920in}}%
\pgfpathlineto{\pgfqpoint{4.384343in}{2.535760in}}%
\pgfpathlineto{\pgfqpoint{4.386862in}{2.485446in}}%
\pgfpathlineto{\pgfqpoint{4.389381in}{2.522763in}}%
\pgfpathlineto{\pgfqpoint{4.391899in}{2.526806in}}%
\pgfpathlineto{\pgfqpoint{4.394418in}{2.557842in}}%
\pgfpathlineto{\pgfqpoint{4.396937in}{2.548013in}}%
\pgfpathlineto{\pgfqpoint{4.399456in}{2.565421in}}%
\pgfpathlineto{\pgfqpoint{4.401974in}{2.558580in}}%
\pgfpathlineto{\pgfqpoint{4.404493in}{2.546830in}}%
\pgfpathlineto{\pgfqpoint{4.407012in}{2.509369in}}%
\pgfpathlineto{\pgfqpoint{4.409531in}{2.530538in}}%
\pgfpathlineto{\pgfqpoint{4.412049in}{2.550023in}}%
\pgfpathlineto{\pgfqpoint{4.414568in}{2.580923in}}%
\pgfpathlineto{\pgfqpoint{4.417087in}{2.573908in}}%
\pgfpathlineto{\pgfqpoint{4.419606in}{2.568686in}}%
\pgfpathlineto{\pgfqpoint{4.422124in}{2.545819in}}%
\pgfpathlineto{\pgfqpoint{4.424643in}{2.550311in}}%
\pgfpathlineto{\pgfqpoint{4.427162in}{2.584009in}}%
\pgfpathlineto{\pgfqpoint{4.429681in}{2.611021in}}%
\pgfpathlineto{\pgfqpoint{4.432199in}{2.629196in}}%
\pgfpathlineto{\pgfqpoint{4.434718in}{2.608467in}}%
\pgfpathlineto{\pgfqpoint{4.437237in}{2.607955in}}%
\pgfpathlineto{\pgfqpoint{4.439756in}{2.606836in}}%
\pgfpathlineto{\pgfqpoint{4.442274in}{2.610478in}}%
\pgfpathlineto{\pgfqpoint{4.444793in}{2.639769in}}%
\pgfpathlineto{\pgfqpoint{4.447312in}{2.621962in}}%
\pgfpathlineto{\pgfqpoint{4.449831in}{2.653010in}}%
\pgfpathlineto{\pgfqpoint{4.452349in}{2.685872in}}%
\pgfpathlineto{\pgfqpoint{4.454868in}{2.682325in}}%
\pgfpathlineto{\pgfqpoint{4.459906in}{2.641164in}}%
\pgfpathlineto{\pgfqpoint{4.462424in}{2.629618in}}%
\pgfpathlineto{\pgfqpoint{4.464943in}{2.633817in}}%
\pgfpathlineto{\pgfqpoint{4.467462in}{2.593653in}}%
\pgfpathlineto{\pgfqpoint{4.469981in}{2.596150in}}%
\pgfpathlineto{\pgfqpoint{4.472499in}{2.569521in}}%
\pgfpathlineto{\pgfqpoint{4.475018in}{2.577001in}}%
\pgfpathlineto{\pgfqpoint{4.477537in}{2.589801in}}%
\pgfpathlineto{\pgfqpoint{4.480056in}{2.617185in}}%
\pgfpathlineto{\pgfqpoint{4.482574in}{2.633226in}}%
\pgfpathlineto{\pgfqpoint{4.485093in}{2.625247in}}%
\pgfpathlineto{\pgfqpoint{4.487612in}{2.585823in}}%
\pgfpathlineto{\pgfqpoint{4.490131in}{2.606215in}}%
\pgfpathlineto{\pgfqpoint{4.492649in}{2.612182in}}%
\pgfpathlineto{\pgfqpoint{4.495168in}{2.615158in}}%
\pgfpathlineto{\pgfqpoint{4.497687in}{2.603620in}}%
\pgfpathlineto{\pgfqpoint{4.500206in}{2.604889in}}%
\pgfpathlineto{\pgfqpoint{4.502724in}{2.614995in}}%
\pgfpathlineto{\pgfqpoint{4.505243in}{2.593215in}}%
\pgfpathlineto{\pgfqpoint{4.507762in}{2.606398in}}%
\pgfpathlineto{\pgfqpoint{4.510281in}{2.591849in}}%
\pgfpathlineto{\pgfqpoint{4.512799in}{2.578448in}}%
\pgfpathlineto{\pgfqpoint{4.515318in}{2.636150in}}%
\pgfpathlineto{\pgfqpoint{4.517837in}{2.631242in}}%
\pgfpathlineto{\pgfqpoint{4.520356in}{2.652440in}}%
\pgfpathlineto{\pgfqpoint{4.522874in}{2.652004in}}%
\pgfpathlineto{\pgfqpoint{4.525393in}{2.664278in}}%
\pgfpathlineto{\pgfqpoint{4.527912in}{2.657421in}}%
\pgfpathlineto{\pgfqpoint{4.530431in}{2.628221in}}%
\pgfpathlineto{\pgfqpoint{4.532949in}{2.638054in}}%
\pgfpathlineto{\pgfqpoint{4.535468in}{2.630495in}}%
\pgfpathlineto{\pgfqpoint{4.537987in}{2.650305in}}%
\pgfpathlineto{\pgfqpoint{4.540506in}{2.644560in}}%
\pgfpathlineto{\pgfqpoint{4.543024in}{2.676060in}}%
\pgfpathlineto{\pgfqpoint{4.545543in}{2.685213in}}%
\pgfpathlineto{\pgfqpoint{4.548062in}{2.685599in}}%
\pgfpathlineto{\pgfqpoint{4.550581in}{2.677190in}}%
\pgfpathlineto{\pgfqpoint{4.553099in}{2.676344in}}%
\pgfpathlineto{\pgfqpoint{4.555618in}{2.684604in}}%
\pgfpathlineto{\pgfqpoint{4.558137in}{2.640188in}}%
\pgfpathlineto{\pgfqpoint{4.560656in}{2.647567in}}%
\pgfpathlineto{\pgfqpoint{4.563174in}{2.700048in}}%
\pgfpathlineto{\pgfqpoint{4.565693in}{2.682210in}}%
\pgfpathlineto{\pgfqpoint{4.568212in}{2.673063in}}%
\pgfpathlineto{\pgfqpoint{4.570731in}{2.665384in}}%
\pgfpathlineto{\pgfqpoint{4.573249in}{2.688407in}}%
\pgfpathlineto{\pgfqpoint{4.575768in}{2.695272in}}%
\pgfpathlineto{\pgfqpoint{4.578287in}{2.721636in}}%
\pgfpathlineto{\pgfqpoint{4.580806in}{2.724344in}}%
\pgfpathlineto{\pgfqpoint{4.583324in}{2.718003in}}%
\pgfpathlineto{\pgfqpoint{4.585843in}{2.728897in}}%
\pgfpathlineto{\pgfqpoint{4.588362in}{2.764982in}}%
\pgfpathlineto{\pgfqpoint{4.590881in}{2.743619in}}%
\pgfpathlineto{\pgfqpoint{4.595918in}{2.708462in}}%
\pgfpathlineto{\pgfqpoint{4.598437in}{2.762194in}}%
\pgfpathlineto{\pgfqpoint{4.600956in}{2.763414in}}%
\pgfpathlineto{\pgfqpoint{4.603474in}{2.722783in}}%
\pgfpathlineto{\pgfqpoint{4.605993in}{2.738691in}}%
\pgfpathlineto{\pgfqpoint{4.608512in}{2.757530in}}%
\pgfpathlineto{\pgfqpoint{4.611031in}{2.763634in}}%
\pgfpathlineto{\pgfqpoint{4.613549in}{2.750169in}}%
\pgfpathlineto{\pgfqpoint{4.616068in}{2.752196in}}%
\pgfpathlineto{\pgfqpoint{4.618587in}{2.760131in}}%
\pgfpathlineto{\pgfqpoint{4.621106in}{2.737611in}}%
\pgfpathlineto{\pgfqpoint{4.623624in}{2.732489in}}%
\pgfpathlineto{\pgfqpoint{4.626143in}{2.731491in}}%
\pgfpathlineto{\pgfqpoint{4.628662in}{2.738587in}}%
\pgfpathlineto{\pgfqpoint{4.631181in}{2.756178in}}%
\pgfpathlineto{\pgfqpoint{4.633699in}{2.737981in}}%
\pgfpathlineto{\pgfqpoint{4.636218in}{2.716329in}}%
\pgfpathlineto{\pgfqpoint{4.638737in}{2.726791in}}%
\pgfpathlineto{\pgfqpoint{4.641256in}{2.739035in}}%
\pgfpathlineto{\pgfqpoint{4.643774in}{2.740250in}}%
\pgfpathlineto{\pgfqpoint{4.646293in}{2.719333in}}%
\pgfpathlineto{\pgfqpoint{4.648812in}{2.738428in}}%
\pgfpathlineto{\pgfqpoint{4.651331in}{2.732384in}}%
\pgfpathlineto{\pgfqpoint{4.653849in}{2.721584in}}%
\pgfpathlineto{\pgfqpoint{4.656368in}{2.760057in}}%
\pgfpathlineto{\pgfqpoint{4.658887in}{2.759063in}}%
\pgfpathlineto{\pgfqpoint{4.661406in}{2.755599in}}%
\pgfpathlineto{\pgfqpoint{4.663924in}{2.799607in}}%
\pgfpathlineto{\pgfqpoint{4.666443in}{2.768489in}}%
\pgfpathlineto{\pgfqpoint{4.668962in}{2.743571in}}%
\pgfpathlineto{\pgfqpoint{4.671481in}{2.740570in}}%
\pgfpathlineto{\pgfqpoint{4.673999in}{2.758956in}}%
\pgfpathlineto{\pgfqpoint{4.676518in}{2.739456in}}%
\pgfpathlineto{\pgfqpoint{4.679037in}{2.714456in}}%
\pgfpathlineto{\pgfqpoint{4.681556in}{2.694603in}}%
\pgfpathlineto{\pgfqpoint{4.684074in}{2.698847in}}%
\pgfpathlineto{\pgfqpoint{4.686593in}{2.718453in}}%
\pgfpathlineto{\pgfqpoint{4.689112in}{2.746024in}}%
\pgfpathlineto{\pgfqpoint{4.691631in}{2.724176in}}%
\pgfpathlineto{\pgfqpoint{4.694149in}{2.706891in}}%
\pgfpathlineto{\pgfqpoint{4.696668in}{2.704657in}}%
\pgfpathlineto{\pgfqpoint{4.699187in}{2.717779in}}%
\pgfpathlineto{\pgfqpoint{4.701706in}{2.694210in}}%
\pgfpathlineto{\pgfqpoint{4.704224in}{2.688197in}}%
\pgfpathlineto{\pgfqpoint{4.706743in}{2.652338in}}%
\pgfpathlineto{\pgfqpoint{4.709262in}{2.656825in}}%
\pgfpathlineto{\pgfqpoint{4.711781in}{2.683634in}}%
\pgfpathlineto{\pgfqpoint{4.714299in}{2.629358in}}%
\pgfpathlineto{\pgfqpoint{4.716818in}{2.605114in}}%
\pgfpathlineto{\pgfqpoint{4.719337in}{2.606542in}}%
\pgfpathlineto{\pgfqpoint{4.721856in}{2.613685in}}%
\pgfpathlineto{\pgfqpoint{4.724374in}{2.637031in}}%
\pgfpathlineto{\pgfqpoint{4.726893in}{2.596228in}}%
\pgfpathlineto{\pgfqpoint{4.729412in}{2.606931in}}%
\pgfpathlineto{\pgfqpoint{4.731931in}{2.608846in}}%
\pgfpathlineto{\pgfqpoint{4.734449in}{2.603808in}}%
\pgfpathlineto{\pgfqpoint{4.736968in}{2.581803in}}%
\pgfpathlineto{\pgfqpoint{4.739487in}{2.594158in}}%
\pgfpathlineto{\pgfqpoint{4.742006in}{2.626830in}}%
\pgfpathlineto{\pgfqpoint{4.744524in}{2.661658in}}%
\pgfpathlineto{\pgfqpoint{4.747043in}{2.655740in}}%
\pgfpathlineto{\pgfqpoint{4.749562in}{2.625722in}}%
\pgfpathlineto{\pgfqpoint{4.752081in}{2.629886in}}%
\pgfpathlineto{\pgfqpoint{4.757118in}{2.593329in}}%
\pgfpathlineto{\pgfqpoint{4.759637in}{2.593638in}}%
\pgfpathlineto{\pgfqpoint{4.762156in}{2.602787in}}%
\pgfpathlineto{\pgfqpoint{4.764674in}{2.590567in}}%
\pgfpathlineto{\pgfqpoint{4.767193in}{2.625604in}}%
\pgfpathlineto{\pgfqpoint{4.769712in}{2.622521in}}%
\pgfpathlineto{\pgfqpoint{4.772231in}{2.636974in}}%
\pgfpathlineto{\pgfqpoint{4.774749in}{2.636408in}}%
\pgfpathlineto{\pgfqpoint{4.777268in}{2.644118in}}%
\pgfpathlineto{\pgfqpoint{4.779787in}{2.627833in}}%
\pgfpathlineto{\pgfqpoint{4.782306in}{2.624026in}}%
\pgfpathlineto{\pgfqpoint{4.784824in}{2.629927in}}%
\pgfpathlineto{\pgfqpoint{4.787343in}{2.657296in}}%
\pgfpathlineto{\pgfqpoint{4.789862in}{2.707148in}}%
\pgfpathlineto{\pgfqpoint{4.792381in}{2.721965in}}%
\pgfpathlineto{\pgfqpoint{4.794899in}{2.720516in}}%
\pgfpathlineto{\pgfqpoint{4.797418in}{2.716607in}}%
\pgfpathlineto{\pgfqpoint{4.799937in}{2.699462in}}%
\pgfpathlineto{\pgfqpoint{4.802456in}{2.710065in}}%
\pgfpathlineto{\pgfqpoint{4.804974in}{2.703179in}}%
\pgfpathlineto{\pgfqpoint{4.810012in}{2.752516in}}%
\pgfpathlineto{\pgfqpoint{4.812531in}{2.732483in}}%
\pgfpathlineto{\pgfqpoint{4.815049in}{2.744393in}}%
\pgfpathlineto{\pgfqpoint{4.817568in}{2.752939in}}%
\pgfpathlineto{\pgfqpoint{4.820087in}{2.754607in}}%
\pgfpathlineto{\pgfqpoint{4.822606in}{2.723446in}}%
\pgfpathlineto{\pgfqpoint{4.825124in}{2.721630in}}%
\pgfpathlineto{\pgfqpoint{4.827643in}{2.702388in}}%
\pgfpathlineto{\pgfqpoint{4.830162in}{2.702873in}}%
\pgfpathlineto{\pgfqpoint{4.832681in}{2.701820in}}%
\pgfpathlineto{\pgfqpoint{4.835199in}{2.704565in}}%
\pgfpathlineto{\pgfqpoint{4.837718in}{2.710862in}}%
\pgfpathlineto{\pgfqpoint{4.840237in}{2.695437in}}%
\pgfpathlineto{\pgfqpoint{4.842756in}{2.716946in}}%
\pgfpathlineto{\pgfqpoint{4.845274in}{2.736677in}}%
\pgfpathlineto{\pgfqpoint{4.847793in}{2.724733in}}%
\pgfpathlineto{\pgfqpoint{4.850312in}{2.759113in}}%
\pgfpathlineto{\pgfqpoint{4.852831in}{2.767754in}}%
\pgfpathlineto{\pgfqpoint{4.855349in}{2.778869in}}%
\pgfpathlineto{\pgfqpoint{4.857868in}{2.785693in}}%
\pgfpathlineto{\pgfqpoint{4.860387in}{2.786571in}}%
\pgfpathlineto{\pgfqpoint{4.862906in}{2.767816in}}%
\pgfpathlineto{\pgfqpoint{4.865424in}{2.758040in}}%
\pgfpathlineto{\pgfqpoint{4.867943in}{2.760407in}}%
\pgfpathlineto{\pgfqpoint{4.870462in}{2.756412in}}%
\pgfpathlineto{\pgfqpoint{4.872981in}{2.756778in}}%
\pgfpathlineto{\pgfqpoint{4.875499in}{2.788777in}}%
\pgfpathlineto{\pgfqpoint{4.878018in}{2.777781in}}%
\pgfpathlineto{\pgfqpoint{4.880537in}{2.772048in}}%
\pgfpathlineto{\pgfqpoint{4.883056in}{2.796049in}}%
\pgfpathlineto{\pgfqpoint{4.888093in}{2.827344in}}%
\pgfpathlineto{\pgfqpoint{4.890612in}{2.832004in}}%
\pgfpathlineto{\pgfqpoint{4.893131in}{2.850524in}}%
\pgfpathlineto{\pgfqpoint{4.895649in}{2.895945in}}%
\pgfpathlineto{\pgfqpoint{4.898168in}{2.882880in}}%
\pgfpathlineto{\pgfqpoint{4.900687in}{2.911681in}}%
\pgfpathlineto{\pgfqpoint{4.903206in}{2.920775in}}%
\pgfpathlineto{\pgfqpoint{4.905724in}{2.901951in}}%
\pgfpathlineto{\pgfqpoint{4.908243in}{2.889463in}}%
\pgfpathlineto{\pgfqpoint{4.910762in}{2.910366in}}%
\pgfpathlineto{\pgfqpoint{4.913281in}{2.943049in}}%
\pgfpathlineto{\pgfqpoint{4.915799in}{2.924097in}}%
\pgfpathlineto{\pgfqpoint{4.918318in}{2.953985in}}%
\pgfpathlineto{\pgfqpoint{4.920837in}{2.941930in}}%
\pgfpathlineto{\pgfqpoint{4.923356in}{2.928273in}}%
\pgfpathlineto{\pgfqpoint{4.925874in}{2.905761in}}%
\pgfpathlineto{\pgfqpoint{4.928393in}{2.904486in}}%
\pgfpathlineto{\pgfqpoint{4.930912in}{2.949667in}}%
\pgfpathlineto{\pgfqpoint{4.933431in}{2.953181in}}%
\pgfpathlineto{\pgfqpoint{4.935949in}{2.989373in}}%
\pgfpathlineto{\pgfqpoint{4.938468in}{3.006187in}}%
\pgfpathlineto{\pgfqpoint{4.940987in}{3.027960in}}%
\pgfpathlineto{\pgfqpoint{4.943506in}{3.030830in}}%
\pgfpathlineto{\pgfqpoint{4.946024in}{3.044493in}}%
\pgfpathlineto{\pgfqpoint{4.948543in}{3.068705in}}%
\pgfpathlineto{\pgfqpoint{4.951062in}{3.061632in}}%
\pgfpathlineto{\pgfqpoint{4.953581in}{3.081191in}}%
\pgfpathlineto{\pgfqpoint{4.956099in}{3.079904in}}%
\pgfpathlineto{\pgfqpoint{4.958618in}{3.074871in}}%
\pgfpathlineto{\pgfqpoint{4.961137in}{3.087093in}}%
\pgfpathlineto{\pgfqpoint{4.963656in}{3.099941in}}%
\pgfpathlineto{\pgfqpoint{4.966174in}{3.106665in}}%
\pgfpathlineto{\pgfqpoint{4.968693in}{3.099057in}}%
\pgfpathlineto{\pgfqpoint{4.971212in}{3.111114in}}%
\pgfpathlineto{\pgfqpoint{4.973731in}{3.154073in}}%
\pgfpathlineto{\pgfqpoint{4.976249in}{3.189792in}}%
\pgfpathlineto{\pgfqpoint{4.978768in}{3.172070in}}%
\pgfpathlineto{\pgfqpoint{4.981287in}{3.160296in}}%
\pgfpathlineto{\pgfqpoint{4.983806in}{3.180173in}}%
\pgfpathlineto{\pgfqpoint{4.986324in}{3.155379in}}%
\pgfpathlineto{\pgfqpoint{4.988843in}{3.184989in}}%
\pgfpathlineto{\pgfqpoint{4.991362in}{3.197416in}}%
\pgfpathlineto{\pgfqpoint{4.993881in}{3.200702in}}%
\pgfpathlineto{\pgfqpoint{4.996399in}{3.221160in}}%
\pgfpathlineto{\pgfqpoint{4.998918in}{3.245680in}}%
\pgfpathlineto{\pgfqpoint{5.001437in}{3.222820in}}%
\pgfpathlineto{\pgfqpoint{5.003956in}{3.198761in}}%
\pgfpathlineto{\pgfqpoint{5.006474in}{3.177188in}}%
\pgfpathlineto{\pgfqpoint{5.008993in}{3.197086in}}%
\pgfpathlineto{\pgfqpoint{5.011512in}{3.212251in}}%
\pgfpathlineto{\pgfqpoint{5.014031in}{3.209507in}}%
\pgfpathlineto{\pgfqpoint{5.016549in}{3.199405in}}%
\pgfpathlineto{\pgfqpoint{5.019068in}{3.177124in}}%
\pgfpathlineto{\pgfqpoint{5.021587in}{3.165358in}}%
\pgfpathlineto{\pgfqpoint{5.024106in}{3.201514in}}%
\pgfpathlineto{\pgfqpoint{5.026624in}{3.231293in}}%
\pgfpathlineto{\pgfqpoint{5.029143in}{3.273494in}}%
\pgfpathlineto{\pgfqpoint{5.031662in}{3.299955in}}%
\pgfpathlineto{\pgfqpoint{5.036699in}{3.274756in}}%
\pgfpathlineto{\pgfqpoint{5.039218in}{3.290578in}}%
\pgfpathlineto{\pgfqpoint{5.041737in}{3.283880in}}%
\pgfpathlineto{\pgfqpoint{5.044256in}{3.252840in}}%
\pgfpathlineto{\pgfqpoint{5.046774in}{3.265493in}}%
\pgfpathlineto{\pgfqpoint{5.049293in}{3.291493in}}%
\pgfpathlineto{\pgfqpoint{5.051812in}{3.266238in}}%
\pgfpathlineto{\pgfqpoint{5.054331in}{3.275127in}}%
\pgfpathlineto{\pgfqpoint{5.056849in}{3.272689in}}%
\pgfpathlineto{\pgfqpoint{5.059368in}{3.260178in}}%
\pgfpathlineto{\pgfqpoint{5.061887in}{3.254958in}}%
\pgfpathlineto{\pgfqpoint{5.064406in}{3.259255in}}%
\pgfpathlineto{\pgfqpoint{5.066924in}{3.251465in}}%
\pgfpathlineto{\pgfqpoint{5.069443in}{3.228243in}}%
\pgfpathlineto{\pgfqpoint{5.071962in}{3.185543in}}%
\pgfpathlineto{\pgfqpoint{5.074481in}{3.180128in}}%
\pgfpathlineto{\pgfqpoint{5.076999in}{3.145351in}}%
\pgfpathlineto{\pgfqpoint{5.079518in}{3.113164in}}%
\pgfpathlineto{\pgfqpoint{5.082037in}{3.088041in}}%
\pgfpathlineto{\pgfqpoint{5.084556in}{3.077619in}}%
\pgfpathlineto{\pgfqpoint{5.087074in}{3.080735in}}%
\pgfpathlineto{\pgfqpoint{5.089593in}{3.067090in}}%
\pgfpathlineto{\pgfqpoint{5.092112in}{3.078916in}}%
\pgfpathlineto{\pgfqpoint{5.094631in}{3.077049in}}%
\pgfpathlineto{\pgfqpoint{5.097149in}{3.114487in}}%
\pgfpathlineto{\pgfqpoint{5.099668in}{3.143908in}}%
\pgfpathlineto{\pgfqpoint{5.102187in}{3.138266in}}%
\pgfpathlineto{\pgfqpoint{5.104706in}{3.120233in}}%
\pgfpathlineto{\pgfqpoint{5.107224in}{3.132076in}}%
\pgfpathlineto{\pgfqpoint{5.109743in}{3.085003in}}%
\pgfpathlineto{\pgfqpoint{5.112262in}{3.097238in}}%
\pgfpathlineto{\pgfqpoint{5.114781in}{3.083473in}}%
\pgfpathlineto{\pgfqpoint{5.117299in}{3.062883in}}%
\pgfpathlineto{\pgfqpoint{5.119818in}{3.076928in}}%
\pgfpathlineto{\pgfqpoint{5.122337in}{3.065188in}}%
\pgfpathlineto{\pgfqpoint{5.124856in}{3.071252in}}%
\pgfpathlineto{\pgfqpoint{5.127374in}{3.041121in}}%
\pgfpathlineto{\pgfqpoint{5.129893in}{3.023500in}}%
\pgfpathlineto{\pgfqpoint{5.132412in}{3.023590in}}%
\pgfpathlineto{\pgfqpoint{5.134931in}{3.024040in}}%
\pgfpathlineto{\pgfqpoint{5.137449in}{3.052095in}}%
\pgfpathlineto{\pgfqpoint{5.139968in}{3.007066in}}%
\pgfpathlineto{\pgfqpoint{5.142487in}{2.981239in}}%
\pgfpathlineto{\pgfqpoint{5.145006in}{2.966315in}}%
\pgfpathlineto{\pgfqpoint{5.147524in}{2.988644in}}%
\pgfpathlineto{\pgfqpoint{5.150043in}{2.986714in}}%
\pgfpathlineto{\pgfqpoint{5.152562in}{2.993804in}}%
\pgfpathlineto{\pgfqpoint{5.155081in}{2.977939in}}%
\pgfpathlineto{\pgfqpoint{5.157599in}{2.984040in}}%
\pgfpathlineto{\pgfqpoint{5.160118in}{2.956288in}}%
\pgfpathlineto{\pgfqpoint{5.162637in}{2.956518in}}%
\pgfpathlineto{\pgfqpoint{5.165156in}{2.927636in}}%
\pgfpathlineto{\pgfqpoint{5.167674in}{2.912818in}}%
\pgfpathlineto{\pgfqpoint{5.170193in}{2.911426in}}%
\pgfpathlineto{\pgfqpoint{5.172712in}{2.898146in}}%
\pgfpathlineto{\pgfqpoint{5.175231in}{2.887703in}}%
\pgfpathlineto{\pgfqpoint{5.177749in}{2.875843in}}%
\pgfpathlineto{\pgfqpoint{5.180268in}{2.889245in}}%
\pgfpathlineto{\pgfqpoint{5.182787in}{2.917259in}}%
\pgfpathlineto{\pgfqpoint{5.185306in}{2.899276in}}%
\pgfpathlineto{\pgfqpoint{5.187824in}{2.921777in}}%
\pgfpathlineto{\pgfqpoint{5.190343in}{2.908788in}}%
\pgfpathlineto{\pgfqpoint{5.192862in}{2.891801in}}%
\pgfpathlineto{\pgfqpoint{5.197899in}{2.907893in}}%
\pgfpathlineto{\pgfqpoint{5.200418in}{2.903903in}}%
\pgfpathlineto{\pgfqpoint{5.202937in}{2.897498in}}%
\pgfpathlineto{\pgfqpoint{5.205456in}{2.929488in}}%
\pgfpathlineto{\pgfqpoint{5.207974in}{2.919007in}}%
\pgfpathlineto{\pgfqpoint{5.210493in}{2.943042in}}%
\pgfpathlineto{\pgfqpoint{5.213012in}{2.937993in}}%
\pgfpathlineto{\pgfqpoint{5.215531in}{2.948906in}}%
\pgfpathlineto{\pgfqpoint{5.218049in}{2.926061in}}%
\pgfpathlineto{\pgfqpoint{5.220568in}{2.913339in}}%
\pgfpathlineto{\pgfqpoint{5.223087in}{2.888137in}}%
\pgfpathlineto{\pgfqpoint{5.225606in}{2.866646in}}%
\pgfpathlineto{\pgfqpoint{5.228124in}{2.855850in}}%
\pgfpathlineto{\pgfqpoint{5.230643in}{2.878265in}}%
\pgfpathlineto{\pgfqpoint{5.233162in}{2.843991in}}%
\pgfpathlineto{\pgfqpoint{5.235681in}{2.850439in}}%
\pgfpathlineto{\pgfqpoint{5.238199in}{2.851975in}}%
\pgfpathlineto{\pgfqpoint{5.240718in}{2.838133in}}%
\pgfpathlineto{\pgfqpoint{5.243237in}{2.849512in}}%
\pgfpathlineto{\pgfqpoint{5.245756in}{2.820883in}}%
\pgfpathlineto{\pgfqpoint{5.248274in}{2.824054in}}%
\pgfpathlineto{\pgfqpoint{5.250793in}{2.847508in}}%
\pgfpathlineto{\pgfqpoint{5.253312in}{2.840766in}}%
\pgfpathlineto{\pgfqpoint{5.255831in}{2.861730in}}%
\pgfpathlineto{\pgfqpoint{5.258349in}{2.896909in}}%
\pgfpathlineto{\pgfqpoint{5.260868in}{2.896367in}}%
\pgfpathlineto{\pgfqpoint{5.263387in}{2.913093in}}%
\pgfpathlineto{\pgfqpoint{5.265906in}{2.904316in}}%
\pgfpathlineto{\pgfqpoint{5.268424in}{2.908613in}}%
\pgfpathlineto{\pgfqpoint{5.270943in}{2.922343in}}%
\pgfpathlineto{\pgfqpoint{5.273462in}{2.911039in}}%
\pgfpathlineto{\pgfqpoint{5.275981in}{2.940139in}}%
\pgfpathlineto{\pgfqpoint{5.278499in}{2.903595in}}%
\pgfpathlineto{\pgfqpoint{5.281018in}{2.923330in}}%
\pgfpathlineto{\pgfqpoint{5.283537in}{2.922019in}}%
\pgfpathlineto{\pgfqpoint{5.286056in}{2.936150in}}%
\pgfpathlineto{\pgfqpoint{5.288574in}{2.874604in}}%
\pgfpathlineto{\pgfqpoint{5.291093in}{2.840948in}}%
\pgfpathlineto{\pgfqpoint{5.293612in}{2.847705in}}%
\pgfpathlineto{\pgfqpoint{5.296131in}{2.844749in}}%
\pgfpathlineto{\pgfqpoint{5.298649in}{2.855053in}}%
\pgfpathlineto{\pgfqpoint{5.301168in}{2.871835in}}%
\pgfpathlineto{\pgfqpoint{5.303687in}{2.905473in}}%
\pgfpathlineto{\pgfqpoint{5.306206in}{2.912737in}}%
\pgfpathlineto{\pgfqpoint{5.308724in}{2.909154in}}%
\pgfpathlineto{\pgfqpoint{5.311243in}{2.907498in}}%
\pgfpathlineto{\pgfqpoint{5.313762in}{2.914816in}}%
\pgfpathlineto{\pgfqpoint{5.316281in}{2.906428in}}%
\pgfpathlineto{\pgfqpoint{5.318799in}{2.898626in}}%
\pgfpathlineto{\pgfqpoint{5.321318in}{2.900108in}}%
\pgfpathlineto{\pgfqpoint{5.323837in}{2.910125in}}%
\pgfpathlineto{\pgfqpoint{5.326356in}{2.877480in}}%
\pgfpathlineto{\pgfqpoint{5.328874in}{2.901691in}}%
\pgfpathlineto{\pgfqpoint{5.331393in}{2.904501in}}%
\pgfpathlineto{\pgfqpoint{5.333912in}{2.928344in}}%
\pgfpathlineto{\pgfqpoint{5.336431in}{2.933344in}}%
\pgfpathlineto{\pgfqpoint{5.341468in}{2.942523in}}%
\pgfpathlineto{\pgfqpoint{5.343987in}{2.956916in}}%
\pgfpathlineto{\pgfqpoint{5.346506in}{2.934294in}}%
\pgfpathlineto{\pgfqpoint{5.349024in}{2.953890in}}%
\pgfpathlineto{\pgfqpoint{5.351543in}{2.928428in}}%
\pgfpathlineto{\pgfqpoint{5.354062in}{2.905011in}}%
\pgfpathlineto{\pgfqpoint{5.356581in}{2.916665in}}%
\pgfpathlineto{\pgfqpoint{5.359099in}{2.909462in}}%
\pgfpathlineto{\pgfqpoint{5.361618in}{2.909074in}}%
\pgfpathlineto{\pgfqpoint{5.364137in}{2.903780in}}%
\pgfpathlineto{\pgfqpoint{5.366656in}{2.925319in}}%
\pgfpathlineto{\pgfqpoint{5.369174in}{2.903624in}}%
\pgfpathlineto{\pgfqpoint{5.371693in}{2.893021in}}%
\pgfpathlineto{\pgfqpoint{5.374212in}{2.861134in}}%
\pgfpathlineto{\pgfqpoint{5.376731in}{2.843202in}}%
\pgfpathlineto{\pgfqpoint{5.379249in}{2.823196in}}%
\pgfpathlineto{\pgfqpoint{5.381768in}{2.822556in}}%
\pgfpathlineto{\pgfqpoint{5.384287in}{2.806786in}}%
\pgfpathlineto{\pgfqpoint{5.386806in}{2.783042in}}%
\pgfpathlineto{\pgfqpoint{5.389324in}{2.801372in}}%
\pgfpathlineto{\pgfqpoint{5.394362in}{2.833075in}}%
\pgfpathlineto{\pgfqpoint{5.396881in}{2.833020in}}%
\pgfpathlineto{\pgfqpoint{5.399399in}{2.853776in}}%
\pgfpathlineto{\pgfqpoint{5.401918in}{2.865923in}}%
\pgfpathlineto{\pgfqpoint{5.404437in}{2.856684in}}%
\pgfpathlineto{\pgfqpoint{5.406956in}{2.837219in}}%
\pgfpathlineto{\pgfqpoint{5.409474in}{2.829729in}}%
\pgfpathlineto{\pgfqpoint{5.411993in}{2.759313in}}%
\pgfpathlineto{\pgfqpoint{5.414512in}{2.747501in}}%
\pgfpathlineto{\pgfqpoint{5.417031in}{2.753045in}}%
\pgfpathlineto{\pgfqpoint{5.419549in}{2.755620in}}%
\pgfpathlineto{\pgfqpoint{5.422068in}{2.763768in}}%
\pgfpathlineto{\pgfqpoint{5.424587in}{2.734242in}}%
\pgfpathlineto{\pgfqpoint{5.427106in}{2.725019in}}%
\pgfpathlineto{\pgfqpoint{5.429624in}{2.757897in}}%
\pgfpathlineto{\pgfqpoint{5.432143in}{2.782070in}}%
\pgfpathlineto{\pgfqpoint{5.434662in}{2.755412in}}%
\pgfpathlineto{\pgfqpoint{5.437181in}{2.772979in}}%
\pgfpathlineto{\pgfqpoint{5.439699in}{2.777279in}}%
\pgfpathlineto{\pgfqpoint{5.442218in}{2.757255in}}%
\pgfpathlineto{\pgfqpoint{5.444737in}{2.764842in}}%
\pgfpathlineto{\pgfqpoint{5.447256in}{2.762517in}}%
\pgfpathlineto{\pgfqpoint{5.449774in}{2.753971in}}%
\pgfpathlineto{\pgfqpoint{5.452293in}{2.730926in}}%
\pgfpathlineto{\pgfqpoint{5.454812in}{2.793531in}}%
\pgfpathlineto{\pgfqpoint{5.457331in}{2.819068in}}%
\pgfpathlineto{\pgfqpoint{5.459849in}{2.854379in}}%
\pgfpathlineto{\pgfqpoint{5.462368in}{2.837964in}}%
\pgfpathlineto{\pgfqpoint{5.464887in}{2.850635in}}%
\pgfpathlineto{\pgfqpoint{5.467406in}{2.856375in}}%
\pgfpathlineto{\pgfqpoint{5.469924in}{2.876085in}}%
\pgfpathlineto{\pgfqpoint{5.472443in}{2.865309in}}%
\pgfpathlineto{\pgfqpoint{5.474962in}{2.872158in}}%
\pgfpathlineto{\pgfqpoint{5.477481in}{2.886254in}}%
\pgfpathlineto{\pgfqpoint{5.479999in}{2.880158in}}%
\pgfpathlineto{\pgfqpoint{5.482518in}{2.877645in}}%
\pgfpathlineto{\pgfqpoint{5.485037in}{2.881072in}}%
\pgfpathlineto{\pgfqpoint{5.487556in}{2.844139in}}%
\pgfpathlineto{\pgfqpoint{5.490074in}{2.810948in}}%
\pgfpathlineto{\pgfqpoint{5.492593in}{2.867447in}}%
\pgfpathlineto{\pgfqpoint{5.492593in}{2.867447in}}%
\pgfusepath{stroke}%
\end{pgfscope}%
\begin{pgfscope}%
\pgfpathrectangle{\pgfqpoint{0.455093in}{0.401080in}}{\pgfqpoint{5.037500in}{4.004000in}}%
\pgfusepath{clip}%
\pgfsetrectcap%
\pgfsetroundjoin%
\pgfsetlinewidth{0.501875pt}%
\definecolor{currentstroke}{rgb}{0.690196,0.188235,0.376471}%
\pgfsetstrokecolor{currentstroke}%
\pgfsetstrokeopacity{0.800000}%
\pgfsetdash{}{0pt}%
\pgfpathmoveto{\pgfqpoint{0.455093in}{1.021044in}}%
\pgfpathlineto{\pgfqpoint{0.457612in}{1.047069in}}%
\pgfpathlineto{\pgfqpoint{0.460131in}{1.062107in}}%
\pgfpathlineto{\pgfqpoint{0.462649in}{1.091579in}}%
\pgfpathlineto{\pgfqpoint{0.465168in}{1.101060in}}%
\pgfpathlineto{\pgfqpoint{0.467687in}{1.085864in}}%
\pgfpathlineto{\pgfqpoint{0.470206in}{1.082999in}}%
\pgfpathlineto{\pgfqpoint{0.472724in}{1.093602in}}%
\pgfpathlineto{\pgfqpoint{0.475243in}{1.109780in}}%
\pgfpathlineto{\pgfqpoint{0.477762in}{1.087863in}}%
\pgfpathlineto{\pgfqpoint{0.480281in}{1.085411in}}%
\pgfpathlineto{\pgfqpoint{0.482799in}{1.097190in}}%
\pgfpathlineto{\pgfqpoint{0.485318in}{1.093129in}}%
\pgfpathlineto{\pgfqpoint{0.487837in}{1.105311in}}%
\pgfpathlineto{\pgfqpoint{0.490356in}{1.112536in}}%
\pgfpathlineto{\pgfqpoint{0.492874in}{1.114115in}}%
\pgfpathlineto{\pgfqpoint{0.495393in}{1.108969in}}%
\pgfpathlineto{\pgfqpoint{0.497912in}{1.098409in}}%
\pgfpathlineto{\pgfqpoint{0.500431in}{1.104755in}}%
\pgfpathlineto{\pgfqpoint{0.502949in}{1.117573in}}%
\pgfpathlineto{\pgfqpoint{0.505468in}{1.134369in}}%
\pgfpathlineto{\pgfqpoint{0.507987in}{1.124502in}}%
\pgfpathlineto{\pgfqpoint{0.510506in}{1.116330in}}%
\pgfpathlineto{\pgfqpoint{0.513024in}{1.122817in}}%
\pgfpathlineto{\pgfqpoint{0.515543in}{1.148160in}}%
\pgfpathlineto{\pgfqpoint{0.518062in}{1.161032in}}%
\pgfpathlineto{\pgfqpoint{0.520581in}{1.166535in}}%
\pgfpathlineto{\pgfqpoint{0.523099in}{1.160876in}}%
\pgfpathlineto{\pgfqpoint{0.525618in}{1.163124in}}%
\pgfpathlineto{\pgfqpoint{0.528137in}{1.175285in}}%
\pgfpathlineto{\pgfqpoint{0.530656in}{1.177311in}}%
\pgfpathlineto{\pgfqpoint{0.533174in}{1.177448in}}%
\pgfpathlineto{\pgfqpoint{0.535693in}{1.172805in}}%
\pgfpathlineto{\pgfqpoint{0.538212in}{1.183283in}}%
\pgfpathlineto{\pgfqpoint{0.540731in}{1.196874in}}%
\pgfpathlineto{\pgfqpoint{0.543249in}{1.194209in}}%
\pgfpathlineto{\pgfqpoint{0.545768in}{1.213141in}}%
\pgfpathlineto{\pgfqpoint{0.548287in}{1.215992in}}%
\pgfpathlineto{\pgfqpoint{0.550806in}{1.231293in}}%
\pgfpathlineto{\pgfqpoint{0.553324in}{1.258431in}}%
\pgfpathlineto{\pgfqpoint{0.555843in}{1.250425in}}%
\pgfpathlineto{\pgfqpoint{0.558362in}{1.235228in}}%
\pgfpathlineto{\pgfqpoint{0.560881in}{1.242053in}}%
\pgfpathlineto{\pgfqpoint{0.563399in}{1.265746in}}%
\pgfpathlineto{\pgfqpoint{0.565918in}{1.249865in}}%
\pgfpathlineto{\pgfqpoint{0.568437in}{1.243540in}}%
\pgfpathlineto{\pgfqpoint{0.570956in}{1.235203in}}%
\pgfpathlineto{\pgfqpoint{0.573474in}{1.227482in}}%
\pgfpathlineto{\pgfqpoint{0.575993in}{1.248970in}}%
\pgfpathlineto{\pgfqpoint{0.578512in}{1.239092in}}%
\pgfpathlineto{\pgfqpoint{0.581031in}{1.267120in}}%
\pgfpathlineto{\pgfqpoint{0.583549in}{1.247201in}}%
\pgfpathlineto{\pgfqpoint{0.586068in}{1.257071in}}%
\pgfpathlineto{\pgfqpoint{0.588587in}{1.280566in}}%
\pgfpathlineto{\pgfqpoint{0.591106in}{1.262589in}}%
\pgfpathlineto{\pgfqpoint{0.593624in}{1.255440in}}%
\pgfpathlineto{\pgfqpoint{0.596143in}{1.267115in}}%
\pgfpathlineto{\pgfqpoint{0.598662in}{1.268860in}}%
\pgfpathlineto{\pgfqpoint{0.601181in}{1.278967in}}%
\pgfpathlineto{\pgfqpoint{0.603699in}{1.295218in}}%
\pgfpathlineto{\pgfqpoint{0.606218in}{1.308996in}}%
\pgfpathlineto{\pgfqpoint{0.608737in}{1.305745in}}%
\pgfpathlineto{\pgfqpoint{0.611256in}{1.322921in}}%
\pgfpathlineto{\pgfqpoint{0.613774in}{1.334818in}}%
\pgfpathlineto{\pgfqpoint{0.616293in}{1.338150in}}%
\pgfpathlineto{\pgfqpoint{0.618812in}{1.329491in}}%
\pgfpathlineto{\pgfqpoint{0.621331in}{1.327243in}}%
\pgfpathlineto{\pgfqpoint{0.623849in}{1.336492in}}%
\pgfpathlineto{\pgfqpoint{0.626368in}{1.331971in}}%
\pgfpathlineto{\pgfqpoint{0.628887in}{1.358646in}}%
\pgfpathlineto{\pgfqpoint{0.631406in}{1.347234in}}%
\pgfpathlineto{\pgfqpoint{0.633924in}{1.326144in}}%
\pgfpathlineto{\pgfqpoint{0.636443in}{1.321909in}}%
\pgfpathlineto{\pgfqpoint{0.638962in}{1.333194in}}%
\pgfpathlineto{\pgfqpoint{0.641481in}{1.334641in}}%
\pgfpathlineto{\pgfqpoint{0.643999in}{1.297233in}}%
\pgfpathlineto{\pgfqpoint{0.646518in}{1.312675in}}%
\pgfpathlineto{\pgfqpoint{0.649037in}{1.299937in}}%
\pgfpathlineto{\pgfqpoint{0.651556in}{1.298528in}}%
\pgfpathlineto{\pgfqpoint{0.654074in}{1.313105in}}%
\pgfpathlineto{\pgfqpoint{0.656593in}{1.323389in}}%
\pgfpathlineto{\pgfqpoint{0.659112in}{1.357813in}}%
\pgfpathlineto{\pgfqpoint{0.661631in}{1.365132in}}%
\pgfpathlineto{\pgfqpoint{0.664149in}{1.338872in}}%
\pgfpathlineto{\pgfqpoint{0.666668in}{1.337296in}}%
\pgfpathlineto{\pgfqpoint{0.669187in}{1.332662in}}%
\pgfpathlineto{\pgfqpoint{0.671706in}{1.342289in}}%
\pgfpathlineto{\pgfqpoint{0.674224in}{1.377055in}}%
\pgfpathlineto{\pgfqpoint{0.676743in}{1.393588in}}%
\pgfpathlineto{\pgfqpoint{0.679262in}{1.389456in}}%
\pgfpathlineto{\pgfqpoint{0.681781in}{1.392894in}}%
\pgfpathlineto{\pgfqpoint{0.684299in}{1.407137in}}%
\pgfpathlineto{\pgfqpoint{0.686818in}{1.401228in}}%
\pgfpathlineto{\pgfqpoint{0.689337in}{1.389202in}}%
\pgfpathlineto{\pgfqpoint{0.691856in}{1.392826in}}%
\pgfpathlineto{\pgfqpoint{0.694374in}{1.370852in}}%
\pgfpathlineto{\pgfqpoint{0.696893in}{1.377476in}}%
\pgfpathlineto{\pgfqpoint{0.699412in}{1.368721in}}%
\pgfpathlineto{\pgfqpoint{0.701931in}{1.365642in}}%
\pgfpathlineto{\pgfqpoint{0.704449in}{1.379635in}}%
\pgfpathlineto{\pgfqpoint{0.706968in}{1.387088in}}%
\pgfpathlineto{\pgfqpoint{0.709487in}{1.385580in}}%
\pgfpathlineto{\pgfqpoint{0.712006in}{1.404000in}}%
\pgfpathlineto{\pgfqpoint{0.714524in}{1.392907in}}%
\pgfpathlineto{\pgfqpoint{0.717043in}{1.392458in}}%
\pgfpathlineto{\pgfqpoint{0.719562in}{1.395507in}}%
\pgfpathlineto{\pgfqpoint{0.722081in}{1.409525in}}%
\pgfpathlineto{\pgfqpoint{0.724599in}{1.409287in}}%
\pgfpathlineto{\pgfqpoint{0.727118in}{1.385946in}}%
\pgfpathlineto{\pgfqpoint{0.729637in}{1.438329in}}%
\pgfpathlineto{\pgfqpoint{0.732156in}{1.421001in}}%
\pgfpathlineto{\pgfqpoint{0.734674in}{1.418202in}}%
\pgfpathlineto{\pgfqpoint{0.737193in}{1.422632in}}%
\pgfpathlineto{\pgfqpoint{0.739712in}{1.427987in}}%
\pgfpathlineto{\pgfqpoint{0.742231in}{1.432932in}}%
\pgfpathlineto{\pgfqpoint{0.744749in}{1.445318in}}%
\pgfpathlineto{\pgfqpoint{0.747268in}{1.452375in}}%
\pgfpathlineto{\pgfqpoint{0.749787in}{1.462192in}}%
\pgfpathlineto{\pgfqpoint{0.752306in}{1.464660in}}%
\pgfpathlineto{\pgfqpoint{0.754824in}{1.455514in}}%
\pgfpathlineto{\pgfqpoint{0.757343in}{1.483808in}}%
\pgfpathlineto{\pgfqpoint{0.759862in}{1.500033in}}%
\pgfpathlineto{\pgfqpoint{0.762381in}{1.504914in}}%
\pgfpathlineto{\pgfqpoint{0.764899in}{1.519842in}}%
\pgfpathlineto{\pgfqpoint{0.767418in}{1.520238in}}%
\pgfpathlineto{\pgfqpoint{0.772456in}{1.540571in}}%
\pgfpathlineto{\pgfqpoint{0.774974in}{1.561470in}}%
\pgfpathlineto{\pgfqpoint{0.777493in}{1.578713in}}%
\pgfpathlineto{\pgfqpoint{0.780012in}{1.581629in}}%
\pgfpathlineto{\pgfqpoint{0.782531in}{1.574562in}}%
\pgfpathlineto{\pgfqpoint{0.785049in}{1.566160in}}%
\pgfpathlineto{\pgfqpoint{0.787568in}{1.583076in}}%
\pgfpathlineto{\pgfqpoint{0.790087in}{1.580088in}}%
\pgfpathlineto{\pgfqpoint{0.792606in}{1.578834in}}%
\pgfpathlineto{\pgfqpoint{0.795124in}{1.591362in}}%
\pgfpathlineto{\pgfqpoint{0.797643in}{1.600833in}}%
\pgfpathlineto{\pgfqpoint{0.800162in}{1.605360in}}%
\pgfpathlineto{\pgfqpoint{0.802681in}{1.602759in}}%
\pgfpathlineto{\pgfqpoint{0.805199in}{1.586049in}}%
\pgfpathlineto{\pgfqpoint{0.807718in}{1.589894in}}%
\pgfpathlineto{\pgfqpoint{0.810237in}{1.580696in}}%
\pgfpathlineto{\pgfqpoint{0.812756in}{1.555254in}}%
\pgfpathlineto{\pgfqpoint{0.815274in}{1.541776in}}%
\pgfpathlineto{\pgfqpoint{0.817793in}{1.524156in}}%
\pgfpathlineto{\pgfqpoint{0.820312in}{1.524515in}}%
\pgfpathlineto{\pgfqpoint{0.822831in}{1.526492in}}%
\pgfpathlineto{\pgfqpoint{0.825349in}{1.539342in}}%
\pgfpathlineto{\pgfqpoint{0.827868in}{1.527589in}}%
\pgfpathlineto{\pgfqpoint{0.830387in}{1.518761in}}%
\pgfpathlineto{\pgfqpoint{0.832906in}{1.503728in}}%
\pgfpathlineto{\pgfqpoint{0.835424in}{1.498141in}}%
\pgfpathlineto{\pgfqpoint{0.837943in}{1.498533in}}%
\pgfpathlineto{\pgfqpoint{0.840462in}{1.480528in}}%
\pgfpathlineto{\pgfqpoint{0.842981in}{1.481967in}}%
\pgfpathlineto{\pgfqpoint{0.845499in}{1.499534in}}%
\pgfpathlineto{\pgfqpoint{0.848018in}{1.522557in}}%
\pgfpathlineto{\pgfqpoint{0.850537in}{1.504573in}}%
\pgfpathlineto{\pgfqpoint{0.853056in}{1.498360in}}%
\pgfpathlineto{\pgfqpoint{0.855574in}{1.506314in}}%
\pgfpathlineto{\pgfqpoint{0.858093in}{1.503201in}}%
\pgfpathlineto{\pgfqpoint{0.860612in}{1.473628in}}%
\pgfpathlineto{\pgfqpoint{0.863131in}{1.464665in}}%
\pgfpathlineto{\pgfqpoint{0.865649in}{1.470166in}}%
\pgfpathlineto{\pgfqpoint{0.868168in}{1.490442in}}%
\pgfpathlineto{\pgfqpoint{0.870687in}{1.498196in}}%
\pgfpathlineto{\pgfqpoint{0.873206in}{1.482301in}}%
\pgfpathlineto{\pgfqpoint{0.875724in}{1.473096in}}%
\pgfpathlineto{\pgfqpoint{0.878243in}{1.475675in}}%
\pgfpathlineto{\pgfqpoint{0.880762in}{1.482820in}}%
\pgfpathlineto{\pgfqpoint{0.883281in}{1.465303in}}%
\pgfpathlineto{\pgfqpoint{0.885799in}{1.474569in}}%
\pgfpathlineto{\pgfqpoint{0.888318in}{1.466102in}}%
\pgfpathlineto{\pgfqpoint{0.890837in}{1.467523in}}%
\pgfpathlineto{\pgfqpoint{0.893356in}{1.445539in}}%
\pgfpathlineto{\pgfqpoint{0.895874in}{1.454063in}}%
\pgfpathlineto{\pgfqpoint{0.898393in}{1.470680in}}%
\pgfpathlineto{\pgfqpoint{0.900912in}{1.460510in}}%
\pgfpathlineto{\pgfqpoint{0.903431in}{1.462517in}}%
\pgfpathlineto{\pgfqpoint{0.905949in}{1.485977in}}%
\pgfpathlineto{\pgfqpoint{0.908468in}{1.498512in}}%
\pgfpathlineto{\pgfqpoint{0.910987in}{1.505221in}}%
\pgfpathlineto{\pgfqpoint{0.913506in}{1.502703in}}%
\pgfpathlineto{\pgfqpoint{0.916024in}{1.490416in}}%
\pgfpathlineto{\pgfqpoint{0.918543in}{1.475556in}}%
\pgfpathlineto{\pgfqpoint{0.921062in}{1.503166in}}%
\pgfpathlineto{\pgfqpoint{0.923581in}{1.487942in}}%
\pgfpathlineto{\pgfqpoint{0.926099in}{1.478923in}}%
\pgfpathlineto{\pgfqpoint{0.928618in}{1.497540in}}%
\pgfpathlineto{\pgfqpoint{0.931137in}{1.496388in}}%
\pgfpathlineto{\pgfqpoint{0.933656in}{1.509155in}}%
\pgfpathlineto{\pgfqpoint{0.936174in}{1.499628in}}%
\pgfpathlineto{\pgfqpoint{0.938693in}{1.517360in}}%
\pgfpathlineto{\pgfqpoint{0.941212in}{1.514566in}}%
\pgfpathlineto{\pgfqpoint{0.943731in}{1.509578in}}%
\pgfpathlineto{\pgfqpoint{0.946249in}{1.484423in}}%
\pgfpathlineto{\pgfqpoint{0.948768in}{1.494079in}}%
\pgfpathlineto{\pgfqpoint{0.951287in}{1.477303in}}%
\pgfpathlineto{\pgfqpoint{0.953806in}{1.463898in}}%
\pgfpathlineto{\pgfqpoint{0.956324in}{1.437649in}}%
\pgfpathlineto{\pgfqpoint{0.958843in}{1.449135in}}%
\pgfpathlineto{\pgfqpoint{0.961362in}{1.454522in}}%
\pgfpathlineto{\pgfqpoint{0.963881in}{1.443325in}}%
\pgfpathlineto{\pgfqpoint{0.966399in}{1.438289in}}%
\pgfpathlineto{\pgfqpoint{0.968918in}{1.451332in}}%
\pgfpathlineto{\pgfqpoint{0.971437in}{1.453093in}}%
\pgfpathlineto{\pgfqpoint{0.973956in}{1.424273in}}%
\pgfpathlineto{\pgfqpoint{0.976474in}{1.449295in}}%
\pgfpathlineto{\pgfqpoint{0.978993in}{1.449849in}}%
\pgfpathlineto{\pgfqpoint{0.981512in}{1.432728in}}%
\pgfpathlineto{\pgfqpoint{0.984031in}{1.383466in}}%
\pgfpathlineto{\pgfqpoint{0.986549in}{1.386408in}}%
\pgfpathlineto{\pgfqpoint{0.989068in}{1.388715in}}%
\pgfpathlineto{\pgfqpoint{0.991587in}{1.399122in}}%
\pgfpathlineto{\pgfqpoint{0.994106in}{1.394769in}}%
\pgfpathlineto{\pgfqpoint{0.996624in}{1.402291in}}%
\pgfpathlineto{\pgfqpoint{0.999143in}{1.405094in}}%
\pgfpathlineto{\pgfqpoint{1.001662in}{1.421410in}}%
\pgfpathlineto{\pgfqpoint{1.004181in}{1.429740in}}%
\pgfpathlineto{\pgfqpoint{1.006699in}{1.416661in}}%
\pgfpathlineto{\pgfqpoint{1.009218in}{1.425875in}}%
\pgfpathlineto{\pgfqpoint{1.011737in}{1.446450in}}%
\pgfpathlineto{\pgfqpoint{1.014256in}{1.421485in}}%
\pgfpathlineto{\pgfqpoint{1.016774in}{1.419271in}}%
\pgfpathlineto{\pgfqpoint{1.019293in}{1.418628in}}%
\pgfpathlineto{\pgfqpoint{1.021812in}{1.416913in}}%
\pgfpathlineto{\pgfqpoint{1.024331in}{1.441920in}}%
\pgfpathlineto{\pgfqpoint{1.026849in}{1.457830in}}%
\pgfpathlineto{\pgfqpoint{1.029368in}{1.456478in}}%
\pgfpathlineto{\pgfqpoint{1.031887in}{1.458865in}}%
\pgfpathlineto{\pgfqpoint{1.034406in}{1.446047in}}%
\pgfpathlineto{\pgfqpoint{1.036924in}{1.465842in}}%
\pgfpathlineto{\pgfqpoint{1.039443in}{1.467382in}}%
\pgfpathlineto{\pgfqpoint{1.041962in}{1.465016in}}%
\pgfpathlineto{\pgfqpoint{1.044481in}{1.472517in}}%
\pgfpathlineto{\pgfqpoint{1.046999in}{1.482893in}}%
\pgfpathlineto{\pgfqpoint{1.049518in}{1.495127in}}%
\pgfpathlineto{\pgfqpoint{1.052037in}{1.498329in}}%
\pgfpathlineto{\pgfqpoint{1.054556in}{1.490285in}}%
\pgfpathlineto{\pgfqpoint{1.057074in}{1.500827in}}%
\pgfpathlineto{\pgfqpoint{1.059593in}{1.504034in}}%
\pgfpathlineto{\pgfqpoint{1.062112in}{1.510142in}}%
\pgfpathlineto{\pgfqpoint{1.064631in}{1.522736in}}%
\pgfpathlineto{\pgfqpoint{1.067149in}{1.521977in}}%
\pgfpathlineto{\pgfqpoint{1.069668in}{1.526528in}}%
\pgfpathlineto{\pgfqpoint{1.072187in}{1.555315in}}%
\pgfpathlineto{\pgfqpoint{1.074706in}{1.536423in}}%
\pgfpathlineto{\pgfqpoint{1.077224in}{1.532243in}}%
\pgfpathlineto{\pgfqpoint{1.079743in}{1.545795in}}%
\pgfpathlineto{\pgfqpoint{1.082262in}{1.514153in}}%
\pgfpathlineto{\pgfqpoint{1.084781in}{1.505641in}}%
\pgfpathlineto{\pgfqpoint{1.087299in}{1.487720in}}%
\pgfpathlineto{\pgfqpoint{1.089818in}{1.500869in}}%
\pgfpathlineto{\pgfqpoint{1.092337in}{1.495778in}}%
\pgfpathlineto{\pgfqpoint{1.094856in}{1.480462in}}%
\pgfpathlineto{\pgfqpoint{1.097374in}{1.497045in}}%
\pgfpathlineto{\pgfqpoint{1.099893in}{1.507620in}}%
\pgfpathlineto{\pgfqpoint{1.102412in}{1.519271in}}%
\pgfpathlineto{\pgfqpoint{1.104931in}{1.533757in}}%
\pgfpathlineto{\pgfqpoint{1.107449in}{1.538159in}}%
\pgfpathlineto{\pgfqpoint{1.109968in}{1.554056in}}%
\pgfpathlineto{\pgfqpoint{1.112487in}{1.584499in}}%
\pgfpathlineto{\pgfqpoint{1.115006in}{1.565418in}}%
\pgfpathlineto{\pgfqpoint{1.117524in}{1.529258in}}%
\pgfpathlineto{\pgfqpoint{1.120043in}{1.522956in}}%
\pgfpathlineto{\pgfqpoint{1.122562in}{1.510733in}}%
\pgfpathlineto{\pgfqpoint{1.125081in}{1.520345in}}%
\pgfpathlineto{\pgfqpoint{1.127599in}{1.499785in}}%
\pgfpathlineto{\pgfqpoint{1.130118in}{1.501210in}}%
\pgfpathlineto{\pgfqpoint{1.132637in}{1.478141in}}%
\pgfpathlineto{\pgfqpoint{1.135156in}{1.474997in}}%
\pgfpathlineto{\pgfqpoint{1.137674in}{1.473624in}}%
\pgfpathlineto{\pgfqpoint{1.140193in}{1.476536in}}%
\pgfpathlineto{\pgfqpoint{1.142712in}{1.505156in}}%
\pgfpathlineto{\pgfqpoint{1.145231in}{1.506243in}}%
\pgfpathlineto{\pgfqpoint{1.147749in}{1.499802in}}%
\pgfpathlineto{\pgfqpoint{1.150268in}{1.490208in}}%
\pgfpathlineto{\pgfqpoint{1.152787in}{1.498227in}}%
\pgfpathlineto{\pgfqpoint{1.155306in}{1.513844in}}%
\pgfpathlineto{\pgfqpoint{1.157824in}{1.498463in}}%
\pgfpathlineto{\pgfqpoint{1.160343in}{1.481023in}}%
\pgfpathlineto{\pgfqpoint{1.162862in}{1.464751in}}%
\pgfpathlineto{\pgfqpoint{1.165381in}{1.470286in}}%
\pgfpathlineto{\pgfqpoint{1.167899in}{1.475246in}}%
\pgfpathlineto{\pgfqpoint{1.170418in}{1.479023in}}%
\pgfpathlineto{\pgfqpoint{1.172937in}{1.483980in}}%
\pgfpathlineto{\pgfqpoint{1.175456in}{1.494822in}}%
\pgfpathlineto{\pgfqpoint{1.177974in}{1.524499in}}%
\pgfpathlineto{\pgfqpoint{1.180493in}{1.509591in}}%
\pgfpathlineto{\pgfqpoint{1.183012in}{1.505699in}}%
\pgfpathlineto{\pgfqpoint{1.185531in}{1.486302in}}%
\pgfpathlineto{\pgfqpoint{1.188049in}{1.534569in}}%
\pgfpathlineto{\pgfqpoint{1.190568in}{1.566198in}}%
\pgfpathlineto{\pgfqpoint{1.193087in}{1.565637in}}%
\pgfpathlineto{\pgfqpoint{1.195606in}{1.571777in}}%
\pgfpathlineto{\pgfqpoint{1.198124in}{1.571771in}}%
\pgfpathlineto{\pgfqpoint{1.200643in}{1.557187in}}%
\pgfpathlineto{\pgfqpoint{1.203162in}{1.569595in}}%
\pgfpathlineto{\pgfqpoint{1.205681in}{1.563706in}}%
\pgfpathlineto{\pgfqpoint{1.208199in}{1.560896in}}%
\pgfpathlineto{\pgfqpoint{1.210718in}{1.566003in}}%
\pgfpathlineto{\pgfqpoint{1.213237in}{1.576089in}}%
\pgfpathlineto{\pgfqpoint{1.215756in}{1.602424in}}%
\pgfpathlineto{\pgfqpoint{1.218274in}{1.608543in}}%
\pgfpathlineto{\pgfqpoint{1.220793in}{1.593407in}}%
\pgfpathlineto{\pgfqpoint{1.223312in}{1.590790in}}%
\pgfpathlineto{\pgfqpoint{1.225831in}{1.574915in}}%
\pgfpathlineto{\pgfqpoint{1.228349in}{1.579290in}}%
\pgfpathlineto{\pgfqpoint{1.230868in}{1.575838in}}%
\pgfpathlineto{\pgfqpoint{1.233387in}{1.570014in}}%
\pgfpathlineto{\pgfqpoint{1.235906in}{1.576447in}}%
\pgfpathlineto{\pgfqpoint{1.238424in}{1.597110in}}%
\pgfpathlineto{\pgfqpoint{1.240943in}{1.607084in}}%
\pgfpathlineto{\pgfqpoint{1.243462in}{1.607181in}}%
\pgfpathlineto{\pgfqpoint{1.245981in}{1.615041in}}%
\pgfpathlineto{\pgfqpoint{1.248499in}{1.607578in}}%
\pgfpathlineto{\pgfqpoint{1.251018in}{1.613138in}}%
\pgfpathlineto{\pgfqpoint{1.253537in}{1.612615in}}%
\pgfpathlineto{\pgfqpoint{1.256056in}{1.611995in}}%
\pgfpathlineto{\pgfqpoint{1.258574in}{1.622104in}}%
\pgfpathlineto{\pgfqpoint{1.261093in}{1.628030in}}%
\pgfpathlineto{\pgfqpoint{1.263612in}{1.624871in}}%
\pgfpathlineto{\pgfqpoint{1.266131in}{1.645226in}}%
\pgfpathlineto{\pgfqpoint{1.268649in}{1.668282in}}%
\pgfpathlineto{\pgfqpoint{1.271168in}{1.669277in}}%
\pgfpathlineto{\pgfqpoint{1.273687in}{1.688348in}}%
\pgfpathlineto{\pgfqpoint{1.276206in}{1.696315in}}%
\pgfpathlineto{\pgfqpoint{1.278724in}{1.689429in}}%
\pgfpathlineto{\pgfqpoint{1.281243in}{1.681392in}}%
\pgfpathlineto{\pgfqpoint{1.283762in}{1.691350in}}%
\pgfpathlineto{\pgfqpoint{1.286281in}{1.727883in}}%
\pgfpathlineto{\pgfqpoint{1.288799in}{1.720124in}}%
\pgfpathlineto{\pgfqpoint{1.291318in}{1.699904in}}%
\pgfpathlineto{\pgfqpoint{1.293837in}{1.708552in}}%
\pgfpathlineto{\pgfqpoint{1.296356in}{1.730392in}}%
\pgfpathlineto{\pgfqpoint{1.298874in}{1.730255in}}%
\pgfpathlineto{\pgfqpoint{1.301393in}{1.720170in}}%
\pgfpathlineto{\pgfqpoint{1.303912in}{1.732771in}}%
\pgfpathlineto{\pgfqpoint{1.306431in}{1.744742in}}%
\pgfpathlineto{\pgfqpoint{1.308949in}{1.763311in}}%
\pgfpathlineto{\pgfqpoint{1.311468in}{1.739909in}}%
\pgfpathlineto{\pgfqpoint{1.313987in}{1.732741in}}%
\pgfpathlineto{\pgfqpoint{1.316506in}{1.725183in}}%
\pgfpathlineto{\pgfqpoint{1.319024in}{1.721712in}}%
\pgfpathlineto{\pgfqpoint{1.321543in}{1.719184in}}%
\pgfpathlineto{\pgfqpoint{1.324062in}{1.749969in}}%
\pgfpathlineto{\pgfqpoint{1.326581in}{1.758154in}}%
\pgfpathlineto{\pgfqpoint{1.329099in}{1.758732in}}%
\pgfpathlineto{\pgfqpoint{1.331618in}{1.745920in}}%
\pgfpathlineto{\pgfqpoint{1.334137in}{1.737680in}}%
\pgfpathlineto{\pgfqpoint{1.336656in}{1.754149in}}%
\pgfpathlineto{\pgfqpoint{1.339174in}{1.740605in}}%
\pgfpathlineto{\pgfqpoint{1.341693in}{1.735630in}}%
\pgfpathlineto{\pgfqpoint{1.344212in}{1.723908in}}%
\pgfpathlineto{\pgfqpoint{1.346731in}{1.696201in}}%
\pgfpathlineto{\pgfqpoint{1.349249in}{1.699577in}}%
\pgfpathlineto{\pgfqpoint{1.351768in}{1.705510in}}%
\pgfpathlineto{\pgfqpoint{1.354287in}{1.677326in}}%
\pgfpathlineto{\pgfqpoint{1.356806in}{1.702135in}}%
\pgfpathlineto{\pgfqpoint{1.359324in}{1.701758in}}%
\pgfpathlineto{\pgfqpoint{1.361843in}{1.712965in}}%
\pgfpathlineto{\pgfqpoint{1.364362in}{1.726316in}}%
\pgfpathlineto{\pgfqpoint{1.366881in}{1.734199in}}%
\pgfpathlineto{\pgfqpoint{1.369399in}{1.754954in}}%
\pgfpathlineto{\pgfqpoint{1.371918in}{1.757255in}}%
\pgfpathlineto{\pgfqpoint{1.374437in}{1.749450in}}%
\pgfpathlineto{\pgfqpoint{1.376956in}{1.757325in}}%
\pgfpathlineto{\pgfqpoint{1.379474in}{1.776407in}}%
\pgfpathlineto{\pgfqpoint{1.381993in}{1.797628in}}%
\pgfpathlineto{\pgfqpoint{1.384512in}{1.764667in}}%
\pgfpathlineto{\pgfqpoint{1.387031in}{1.769023in}}%
\pgfpathlineto{\pgfqpoint{1.389549in}{1.764918in}}%
\pgfpathlineto{\pgfqpoint{1.392068in}{1.776463in}}%
\pgfpathlineto{\pgfqpoint{1.394587in}{1.762544in}}%
\pgfpathlineto{\pgfqpoint{1.397106in}{1.746493in}}%
\pgfpathlineto{\pgfqpoint{1.399624in}{1.732374in}}%
\pgfpathlineto{\pgfqpoint{1.402143in}{1.694872in}}%
\pgfpathlineto{\pgfqpoint{1.404662in}{1.708847in}}%
\pgfpathlineto{\pgfqpoint{1.407181in}{1.684441in}}%
\pgfpathlineto{\pgfqpoint{1.409699in}{1.690823in}}%
\pgfpathlineto{\pgfqpoint{1.412218in}{1.695287in}}%
\pgfpathlineto{\pgfqpoint{1.414737in}{1.686693in}}%
\pgfpathlineto{\pgfqpoint{1.417256in}{1.672748in}}%
\pgfpathlineto{\pgfqpoint{1.419774in}{1.688829in}}%
\pgfpathlineto{\pgfqpoint{1.422293in}{1.683949in}}%
\pgfpathlineto{\pgfqpoint{1.424812in}{1.682413in}}%
\pgfpathlineto{\pgfqpoint{1.427331in}{1.675392in}}%
\pgfpathlineto{\pgfqpoint{1.429849in}{1.687603in}}%
\pgfpathlineto{\pgfqpoint{1.432368in}{1.693772in}}%
\pgfpathlineto{\pgfqpoint{1.434887in}{1.706741in}}%
\pgfpathlineto{\pgfqpoint{1.437406in}{1.690160in}}%
\pgfpathlineto{\pgfqpoint{1.439924in}{1.705088in}}%
\pgfpathlineto{\pgfqpoint{1.442443in}{1.698339in}}%
\pgfpathlineto{\pgfqpoint{1.444962in}{1.694840in}}%
\pgfpathlineto{\pgfqpoint{1.447481in}{1.696172in}}%
\pgfpathlineto{\pgfqpoint{1.449999in}{1.713917in}}%
\pgfpathlineto{\pgfqpoint{1.452518in}{1.721173in}}%
\pgfpathlineto{\pgfqpoint{1.457556in}{1.713251in}}%
\pgfpathlineto{\pgfqpoint{1.460074in}{1.696801in}}%
\pgfpathlineto{\pgfqpoint{1.462593in}{1.715745in}}%
\pgfpathlineto{\pgfqpoint{1.465112in}{1.679333in}}%
\pgfpathlineto{\pgfqpoint{1.467631in}{1.687904in}}%
\pgfpathlineto{\pgfqpoint{1.470149in}{1.671702in}}%
\pgfpathlineto{\pgfqpoint{1.472668in}{1.668562in}}%
\pgfpathlineto{\pgfqpoint{1.475187in}{1.685303in}}%
\pgfpathlineto{\pgfqpoint{1.477706in}{1.693681in}}%
\pgfpathlineto{\pgfqpoint{1.480224in}{1.680588in}}%
\pgfpathlineto{\pgfqpoint{1.482743in}{1.703969in}}%
\pgfpathlineto{\pgfqpoint{1.485262in}{1.710762in}}%
\pgfpathlineto{\pgfqpoint{1.487781in}{1.734368in}}%
\pgfpathlineto{\pgfqpoint{1.490299in}{1.722292in}}%
\pgfpathlineto{\pgfqpoint{1.492818in}{1.733971in}}%
\pgfpathlineto{\pgfqpoint{1.495337in}{1.718612in}}%
\pgfpathlineto{\pgfqpoint{1.497856in}{1.718489in}}%
\pgfpathlineto{\pgfqpoint{1.500374in}{1.709894in}}%
\pgfpathlineto{\pgfqpoint{1.502893in}{1.714260in}}%
\pgfpathlineto{\pgfqpoint{1.505412in}{1.705899in}}%
\pgfpathlineto{\pgfqpoint{1.507931in}{1.713929in}}%
\pgfpathlineto{\pgfqpoint{1.510449in}{1.734072in}}%
\pgfpathlineto{\pgfqpoint{1.515487in}{1.763106in}}%
\pgfpathlineto{\pgfqpoint{1.518006in}{1.781089in}}%
\pgfpathlineto{\pgfqpoint{1.520524in}{1.767210in}}%
\pgfpathlineto{\pgfqpoint{1.523043in}{1.742208in}}%
\pgfpathlineto{\pgfqpoint{1.525562in}{1.780033in}}%
\pgfpathlineto{\pgfqpoint{1.528081in}{1.778556in}}%
\pgfpathlineto{\pgfqpoint{1.530599in}{1.786743in}}%
\pgfpathlineto{\pgfqpoint{1.533118in}{1.782920in}}%
\pgfpathlineto{\pgfqpoint{1.535637in}{1.774734in}}%
\pgfpathlineto{\pgfqpoint{1.538156in}{1.776362in}}%
\pgfpathlineto{\pgfqpoint{1.540674in}{1.756693in}}%
\pgfpathlineto{\pgfqpoint{1.543193in}{1.775274in}}%
\pgfpathlineto{\pgfqpoint{1.545712in}{1.777366in}}%
\pgfpathlineto{\pgfqpoint{1.548231in}{1.776720in}}%
\pgfpathlineto{\pgfqpoint{1.550749in}{1.813721in}}%
\pgfpathlineto{\pgfqpoint{1.553268in}{1.837197in}}%
\pgfpathlineto{\pgfqpoint{1.555787in}{1.822944in}}%
\pgfpathlineto{\pgfqpoint{1.558306in}{1.799006in}}%
\pgfpathlineto{\pgfqpoint{1.560824in}{1.808527in}}%
\pgfpathlineto{\pgfqpoint{1.563343in}{1.817627in}}%
\pgfpathlineto{\pgfqpoint{1.565862in}{1.814402in}}%
\pgfpathlineto{\pgfqpoint{1.568381in}{1.794496in}}%
\pgfpathlineto{\pgfqpoint{1.570899in}{1.799361in}}%
\pgfpathlineto{\pgfqpoint{1.573418in}{1.791027in}}%
\pgfpathlineto{\pgfqpoint{1.575937in}{1.786569in}}%
\pgfpathlineto{\pgfqpoint{1.578456in}{1.770967in}}%
\pgfpathlineto{\pgfqpoint{1.580974in}{1.748253in}}%
\pgfpathlineto{\pgfqpoint{1.583493in}{1.784716in}}%
\pgfpathlineto{\pgfqpoint{1.586012in}{1.772544in}}%
\pgfpathlineto{\pgfqpoint{1.588531in}{1.772212in}}%
\pgfpathlineto{\pgfqpoint{1.591049in}{1.775941in}}%
\pgfpathlineto{\pgfqpoint{1.593568in}{1.770655in}}%
\pgfpathlineto{\pgfqpoint{1.596087in}{1.761163in}}%
\pgfpathlineto{\pgfqpoint{1.598606in}{1.742814in}}%
\pgfpathlineto{\pgfqpoint{1.601124in}{1.726009in}}%
\pgfpathlineto{\pgfqpoint{1.603643in}{1.716301in}}%
\pgfpathlineto{\pgfqpoint{1.606162in}{1.716380in}}%
\pgfpathlineto{\pgfqpoint{1.608681in}{1.698746in}}%
\pgfpathlineto{\pgfqpoint{1.611199in}{1.694991in}}%
\pgfpathlineto{\pgfqpoint{1.613718in}{1.708607in}}%
\pgfpathlineto{\pgfqpoint{1.616237in}{1.690007in}}%
\pgfpathlineto{\pgfqpoint{1.618756in}{1.698859in}}%
\pgfpathlineto{\pgfqpoint{1.621274in}{1.696775in}}%
\pgfpathlineto{\pgfqpoint{1.623793in}{1.678574in}}%
\pgfpathlineto{\pgfqpoint{1.626312in}{1.709310in}}%
\pgfpathlineto{\pgfqpoint{1.628831in}{1.709192in}}%
\pgfpathlineto{\pgfqpoint{1.631349in}{1.716144in}}%
\pgfpathlineto{\pgfqpoint{1.633868in}{1.725800in}}%
\pgfpathlineto{\pgfqpoint{1.636387in}{1.737707in}}%
\pgfpathlineto{\pgfqpoint{1.638906in}{1.699234in}}%
\pgfpathlineto{\pgfqpoint{1.641424in}{1.693600in}}%
\pgfpathlineto{\pgfqpoint{1.643943in}{1.685402in}}%
\pgfpathlineto{\pgfqpoint{1.646462in}{1.689714in}}%
\pgfpathlineto{\pgfqpoint{1.648981in}{1.666324in}}%
\pgfpathlineto{\pgfqpoint{1.651499in}{1.674357in}}%
\pgfpathlineto{\pgfqpoint{1.654018in}{1.682927in}}%
\pgfpathlineto{\pgfqpoint{1.656537in}{1.695766in}}%
\pgfpathlineto{\pgfqpoint{1.659056in}{1.675742in}}%
\pgfpathlineto{\pgfqpoint{1.661574in}{1.671613in}}%
\pgfpathlineto{\pgfqpoint{1.664093in}{1.647033in}}%
\pgfpathlineto{\pgfqpoint{1.666612in}{1.666368in}}%
\pgfpathlineto{\pgfqpoint{1.669131in}{1.673408in}}%
\pgfpathlineto{\pgfqpoint{1.671649in}{1.710046in}}%
\pgfpathlineto{\pgfqpoint{1.674168in}{1.730934in}}%
\pgfpathlineto{\pgfqpoint{1.676687in}{1.738495in}}%
\pgfpathlineto{\pgfqpoint{1.679206in}{1.754089in}}%
\pgfpathlineto{\pgfqpoint{1.681724in}{1.746611in}}%
\pgfpathlineto{\pgfqpoint{1.684243in}{1.729885in}}%
\pgfpathlineto{\pgfqpoint{1.686762in}{1.737860in}}%
\pgfpathlineto{\pgfqpoint{1.689281in}{1.748130in}}%
\pgfpathlineto{\pgfqpoint{1.691799in}{1.756788in}}%
\pgfpathlineto{\pgfqpoint{1.694318in}{1.750540in}}%
\pgfpathlineto{\pgfqpoint{1.696837in}{1.743193in}}%
\pgfpathlineto{\pgfqpoint{1.699356in}{1.709428in}}%
\pgfpathlineto{\pgfqpoint{1.701874in}{1.739812in}}%
\pgfpathlineto{\pgfqpoint{1.704393in}{1.744064in}}%
\pgfpathlineto{\pgfqpoint{1.706912in}{1.754088in}}%
\pgfpathlineto{\pgfqpoint{1.709431in}{1.733739in}}%
\pgfpathlineto{\pgfqpoint{1.711949in}{1.730015in}}%
\pgfpathlineto{\pgfqpoint{1.714468in}{1.708377in}}%
\pgfpathlineto{\pgfqpoint{1.716987in}{1.717726in}}%
\pgfpathlineto{\pgfqpoint{1.719506in}{1.720723in}}%
\pgfpathlineto{\pgfqpoint{1.722024in}{1.736493in}}%
\pgfpathlineto{\pgfqpoint{1.724543in}{1.768656in}}%
\pgfpathlineto{\pgfqpoint{1.727062in}{1.760696in}}%
\pgfpathlineto{\pgfqpoint{1.729581in}{1.766870in}}%
\pgfpathlineto{\pgfqpoint{1.732099in}{1.764780in}}%
\pgfpathlineto{\pgfqpoint{1.734618in}{1.765254in}}%
\pgfpathlineto{\pgfqpoint{1.737137in}{1.781390in}}%
\pgfpathlineto{\pgfqpoint{1.739656in}{1.773867in}}%
\pgfpathlineto{\pgfqpoint{1.742174in}{1.743273in}}%
\pgfpathlineto{\pgfqpoint{1.744693in}{1.718693in}}%
\pgfpathlineto{\pgfqpoint{1.747212in}{1.718459in}}%
\pgfpathlineto{\pgfqpoint{1.749731in}{1.721634in}}%
\pgfpathlineto{\pgfqpoint{1.752249in}{1.752376in}}%
\pgfpathlineto{\pgfqpoint{1.754768in}{1.719138in}}%
\pgfpathlineto{\pgfqpoint{1.757287in}{1.713591in}}%
\pgfpathlineto{\pgfqpoint{1.759806in}{1.712494in}}%
\pgfpathlineto{\pgfqpoint{1.762324in}{1.703426in}}%
\pgfpathlineto{\pgfqpoint{1.764843in}{1.724467in}}%
\pgfpathlineto{\pgfqpoint{1.767362in}{1.711623in}}%
\pgfpathlineto{\pgfqpoint{1.769881in}{1.700761in}}%
\pgfpathlineto{\pgfqpoint{1.772399in}{1.725549in}}%
\pgfpathlineto{\pgfqpoint{1.774918in}{1.709333in}}%
\pgfpathlineto{\pgfqpoint{1.777437in}{1.688465in}}%
\pgfpathlineto{\pgfqpoint{1.779956in}{1.697588in}}%
\pgfpathlineto{\pgfqpoint{1.782474in}{1.692323in}}%
\pgfpathlineto{\pgfqpoint{1.784993in}{1.676931in}}%
\pgfpathlineto{\pgfqpoint{1.787512in}{1.668363in}}%
\pgfpathlineto{\pgfqpoint{1.790031in}{1.652613in}}%
\pgfpathlineto{\pgfqpoint{1.792549in}{1.656843in}}%
\pgfpathlineto{\pgfqpoint{1.795068in}{1.642874in}}%
\pgfpathlineto{\pgfqpoint{1.797587in}{1.678732in}}%
\pgfpathlineto{\pgfqpoint{1.800106in}{1.703095in}}%
\pgfpathlineto{\pgfqpoint{1.802624in}{1.689132in}}%
\pgfpathlineto{\pgfqpoint{1.805143in}{1.682902in}}%
\pgfpathlineto{\pgfqpoint{1.807662in}{1.656514in}}%
\pgfpathlineto{\pgfqpoint{1.810181in}{1.656868in}}%
\pgfpathlineto{\pgfqpoint{1.812699in}{1.652282in}}%
\pgfpathlineto{\pgfqpoint{1.815218in}{1.649394in}}%
\pgfpathlineto{\pgfqpoint{1.817737in}{1.642342in}}%
\pgfpathlineto{\pgfqpoint{1.820256in}{1.663277in}}%
\pgfpathlineto{\pgfqpoint{1.822774in}{1.690121in}}%
\pgfpathlineto{\pgfqpoint{1.825293in}{1.705915in}}%
\pgfpathlineto{\pgfqpoint{1.827812in}{1.701508in}}%
\pgfpathlineto{\pgfqpoint{1.830331in}{1.677094in}}%
\pgfpathlineto{\pgfqpoint{1.832849in}{1.686520in}}%
\pgfpathlineto{\pgfqpoint{1.835368in}{1.684139in}}%
\pgfpathlineto{\pgfqpoint{1.837887in}{1.679808in}}%
\pgfpathlineto{\pgfqpoint{1.840406in}{1.671942in}}%
\pgfpathlineto{\pgfqpoint{1.842924in}{1.679504in}}%
\pgfpathlineto{\pgfqpoint{1.845443in}{1.677243in}}%
\pgfpathlineto{\pgfqpoint{1.847962in}{1.669128in}}%
\pgfpathlineto{\pgfqpoint{1.850481in}{1.665725in}}%
\pgfpathlineto{\pgfqpoint{1.852999in}{1.662833in}}%
\pgfpathlineto{\pgfqpoint{1.855518in}{1.673079in}}%
\pgfpathlineto{\pgfqpoint{1.858037in}{1.666012in}}%
\pgfpathlineto{\pgfqpoint{1.860556in}{1.686351in}}%
\pgfpathlineto{\pgfqpoint{1.863074in}{1.708830in}}%
\pgfpathlineto{\pgfqpoint{1.865593in}{1.717271in}}%
\pgfpathlineto{\pgfqpoint{1.868112in}{1.715621in}}%
\pgfpathlineto{\pgfqpoint{1.870631in}{1.703588in}}%
\pgfpathlineto{\pgfqpoint{1.873149in}{1.703874in}}%
\pgfpathlineto{\pgfqpoint{1.875668in}{1.711734in}}%
\pgfpathlineto{\pgfqpoint{1.878187in}{1.700144in}}%
\pgfpathlineto{\pgfqpoint{1.880706in}{1.685046in}}%
\pgfpathlineto{\pgfqpoint{1.883224in}{1.679795in}}%
\pgfpathlineto{\pgfqpoint{1.885743in}{1.670934in}}%
\pgfpathlineto{\pgfqpoint{1.888262in}{1.660432in}}%
\pgfpathlineto{\pgfqpoint{1.890781in}{1.660628in}}%
\pgfpathlineto{\pgfqpoint{1.893299in}{1.626839in}}%
\pgfpathlineto{\pgfqpoint{1.895818in}{1.615874in}}%
\pgfpathlineto{\pgfqpoint{1.898337in}{1.627810in}}%
\pgfpathlineto{\pgfqpoint{1.900856in}{1.644343in}}%
\pgfpathlineto{\pgfqpoint{1.903374in}{1.637192in}}%
\pgfpathlineto{\pgfqpoint{1.905893in}{1.627881in}}%
\pgfpathlineto{\pgfqpoint{1.908412in}{1.608934in}}%
\pgfpathlineto{\pgfqpoint{1.910931in}{1.621462in}}%
\pgfpathlineto{\pgfqpoint{1.913449in}{1.626716in}}%
\pgfpathlineto{\pgfqpoint{1.915968in}{1.655951in}}%
\pgfpathlineto{\pgfqpoint{1.918487in}{1.651668in}}%
\pgfpathlineto{\pgfqpoint{1.921006in}{1.632936in}}%
\pgfpathlineto{\pgfqpoint{1.926043in}{1.678847in}}%
\pgfpathlineto{\pgfqpoint{1.928562in}{1.665911in}}%
\pgfpathlineto{\pgfqpoint{1.931081in}{1.676232in}}%
\pgfpathlineto{\pgfqpoint{1.933599in}{1.665490in}}%
\pgfpathlineto{\pgfqpoint{1.936118in}{1.629078in}}%
\pgfpathlineto{\pgfqpoint{1.938637in}{1.634294in}}%
\pgfpathlineto{\pgfqpoint{1.941156in}{1.644760in}}%
\pgfpathlineto{\pgfqpoint{1.943674in}{1.667075in}}%
\pgfpathlineto{\pgfqpoint{1.948712in}{1.619671in}}%
\pgfpathlineto{\pgfqpoint{1.951231in}{1.644457in}}%
\pgfpathlineto{\pgfqpoint{1.953749in}{1.635944in}}%
\pgfpathlineto{\pgfqpoint{1.956268in}{1.644196in}}%
\pgfpathlineto{\pgfqpoint{1.958787in}{1.617723in}}%
\pgfpathlineto{\pgfqpoint{1.961306in}{1.605524in}}%
\pgfpathlineto{\pgfqpoint{1.963824in}{1.604280in}}%
\pgfpathlineto{\pgfqpoint{1.966343in}{1.642939in}}%
\pgfpathlineto{\pgfqpoint{1.968862in}{1.657778in}}%
\pgfpathlineto{\pgfqpoint{1.971381in}{1.642328in}}%
\pgfpathlineto{\pgfqpoint{1.973899in}{1.624404in}}%
\pgfpathlineto{\pgfqpoint{1.976418in}{1.602296in}}%
\pgfpathlineto{\pgfqpoint{1.978937in}{1.609169in}}%
\pgfpathlineto{\pgfqpoint{1.981456in}{1.606965in}}%
\pgfpathlineto{\pgfqpoint{1.983974in}{1.598884in}}%
\pgfpathlineto{\pgfqpoint{1.986493in}{1.585617in}}%
\pgfpathlineto{\pgfqpoint{1.991531in}{1.587886in}}%
\pgfpathlineto{\pgfqpoint{1.994049in}{1.602224in}}%
\pgfpathlineto{\pgfqpoint{1.996568in}{1.567191in}}%
\pgfpathlineto{\pgfqpoint{1.999087in}{1.542275in}}%
\pgfpathlineto{\pgfqpoint{2.001606in}{1.559229in}}%
\pgfpathlineto{\pgfqpoint{2.004124in}{1.575221in}}%
\pgfpathlineto{\pgfqpoint{2.006643in}{1.553814in}}%
\pgfpathlineto{\pgfqpoint{2.009162in}{1.540285in}}%
\pgfpathlineto{\pgfqpoint{2.011681in}{1.535631in}}%
\pgfpathlineto{\pgfqpoint{2.016718in}{1.557335in}}%
\pgfpathlineto{\pgfqpoint{2.019237in}{1.551396in}}%
\pgfpathlineto{\pgfqpoint{2.021756in}{1.562308in}}%
\pgfpathlineto{\pgfqpoint{2.024274in}{1.565126in}}%
\pgfpathlineto{\pgfqpoint{2.026793in}{1.555247in}}%
\pgfpathlineto{\pgfqpoint{2.029312in}{1.552998in}}%
\pgfpathlineto{\pgfqpoint{2.031831in}{1.554966in}}%
\pgfpathlineto{\pgfqpoint{2.034349in}{1.546138in}}%
\pgfpathlineto{\pgfqpoint{2.036868in}{1.571539in}}%
\pgfpathlineto{\pgfqpoint{2.039387in}{1.585508in}}%
\pgfpathlineto{\pgfqpoint{2.041906in}{1.585990in}}%
\pgfpathlineto{\pgfqpoint{2.044424in}{1.618772in}}%
\pgfpathlineto{\pgfqpoint{2.046943in}{1.613745in}}%
\pgfpathlineto{\pgfqpoint{2.049462in}{1.618755in}}%
\pgfpathlineto{\pgfqpoint{2.051981in}{1.642038in}}%
\pgfpathlineto{\pgfqpoint{2.054499in}{1.661502in}}%
\pgfpathlineto{\pgfqpoint{2.057018in}{1.651543in}}%
\pgfpathlineto{\pgfqpoint{2.062056in}{1.661918in}}%
\pgfpathlineto{\pgfqpoint{2.064574in}{1.656296in}}%
\pgfpathlineto{\pgfqpoint{2.067093in}{1.653646in}}%
\pgfpathlineto{\pgfqpoint{2.069612in}{1.637309in}}%
\pgfpathlineto{\pgfqpoint{2.072131in}{1.629884in}}%
\pgfpathlineto{\pgfqpoint{2.074649in}{1.637673in}}%
\pgfpathlineto{\pgfqpoint{2.077168in}{1.637840in}}%
\pgfpathlineto{\pgfqpoint{2.079687in}{1.608877in}}%
\pgfpathlineto{\pgfqpoint{2.082206in}{1.606976in}}%
\pgfpathlineto{\pgfqpoint{2.084724in}{1.606965in}}%
\pgfpathlineto{\pgfqpoint{2.087243in}{1.595068in}}%
\pgfpathlineto{\pgfqpoint{2.089762in}{1.585170in}}%
\pgfpathlineto{\pgfqpoint{2.092281in}{1.612980in}}%
\pgfpathlineto{\pgfqpoint{2.094799in}{1.600289in}}%
\pgfpathlineto{\pgfqpoint{2.097318in}{1.597036in}}%
\pgfpathlineto{\pgfqpoint{2.099837in}{1.599683in}}%
\pgfpathlineto{\pgfqpoint{2.102356in}{1.618304in}}%
\pgfpathlineto{\pgfqpoint{2.104874in}{1.625273in}}%
\pgfpathlineto{\pgfqpoint{2.107393in}{1.635482in}}%
\pgfpathlineto{\pgfqpoint{2.109912in}{1.627408in}}%
\pgfpathlineto{\pgfqpoint{2.112431in}{1.609987in}}%
\pgfpathlineto{\pgfqpoint{2.114949in}{1.606848in}}%
\pgfpathlineto{\pgfqpoint{2.117468in}{1.633356in}}%
\pgfpathlineto{\pgfqpoint{2.119987in}{1.626240in}}%
\pgfpathlineto{\pgfqpoint{2.122506in}{1.638363in}}%
\pgfpathlineto{\pgfqpoint{2.125024in}{1.633392in}}%
\pgfpathlineto{\pgfqpoint{2.127543in}{1.633759in}}%
\pgfpathlineto{\pgfqpoint{2.130062in}{1.652084in}}%
\pgfpathlineto{\pgfqpoint{2.132581in}{1.655592in}}%
\pgfpathlineto{\pgfqpoint{2.135099in}{1.652651in}}%
\pgfpathlineto{\pgfqpoint{2.137618in}{1.646169in}}%
\pgfpathlineto{\pgfqpoint{2.140137in}{1.642396in}}%
\pgfpathlineto{\pgfqpoint{2.142656in}{1.635300in}}%
\pgfpathlineto{\pgfqpoint{2.145174in}{1.641226in}}%
\pgfpathlineto{\pgfqpoint{2.147693in}{1.684342in}}%
\pgfpathlineto{\pgfqpoint{2.150212in}{1.664140in}}%
\pgfpathlineto{\pgfqpoint{2.152731in}{1.654092in}}%
\pgfpathlineto{\pgfqpoint{2.155249in}{1.662929in}}%
\pgfpathlineto{\pgfqpoint{2.157768in}{1.653641in}}%
\pgfpathlineto{\pgfqpoint{2.160287in}{1.665795in}}%
\pgfpathlineto{\pgfqpoint{2.162806in}{1.693028in}}%
\pgfpathlineto{\pgfqpoint{2.165324in}{1.712446in}}%
\pgfpathlineto{\pgfqpoint{2.167843in}{1.721034in}}%
\pgfpathlineto{\pgfqpoint{2.170362in}{1.717871in}}%
\pgfpathlineto{\pgfqpoint{2.172881in}{1.705654in}}%
\pgfpathlineto{\pgfqpoint{2.175399in}{1.728946in}}%
\pgfpathlineto{\pgfqpoint{2.177918in}{1.708166in}}%
\pgfpathlineto{\pgfqpoint{2.180437in}{1.700534in}}%
\pgfpathlineto{\pgfqpoint{2.182956in}{1.715313in}}%
\pgfpathlineto{\pgfqpoint{2.185474in}{1.727434in}}%
\pgfpathlineto{\pgfqpoint{2.187993in}{1.734203in}}%
\pgfpathlineto{\pgfqpoint{2.190512in}{1.756926in}}%
\pgfpathlineto{\pgfqpoint{2.193031in}{1.740296in}}%
\pgfpathlineto{\pgfqpoint{2.195549in}{1.757202in}}%
\pgfpathlineto{\pgfqpoint{2.198068in}{1.767111in}}%
\pgfpathlineto{\pgfqpoint{2.200587in}{1.748962in}}%
\pgfpathlineto{\pgfqpoint{2.205624in}{1.765398in}}%
\pgfpathlineto{\pgfqpoint{2.208143in}{1.748343in}}%
\pgfpathlineto{\pgfqpoint{2.210662in}{1.707255in}}%
\pgfpathlineto{\pgfqpoint{2.213181in}{1.724624in}}%
\pgfpathlineto{\pgfqpoint{2.215699in}{1.737188in}}%
\pgfpathlineto{\pgfqpoint{2.218218in}{1.764637in}}%
\pgfpathlineto{\pgfqpoint{2.220737in}{1.777885in}}%
\pgfpathlineto{\pgfqpoint{2.223256in}{1.783808in}}%
\pgfpathlineto{\pgfqpoint{2.225774in}{1.791013in}}%
\pgfpathlineto{\pgfqpoint{2.228293in}{1.786318in}}%
\pgfpathlineto{\pgfqpoint{2.230812in}{1.744731in}}%
\pgfpathlineto{\pgfqpoint{2.233331in}{1.755750in}}%
\pgfpathlineto{\pgfqpoint{2.235849in}{1.765818in}}%
\pgfpathlineto{\pgfqpoint{2.238368in}{1.744978in}}%
\pgfpathlineto{\pgfqpoint{2.240887in}{1.755777in}}%
\pgfpathlineto{\pgfqpoint{2.243406in}{1.777737in}}%
\pgfpathlineto{\pgfqpoint{2.245924in}{1.792188in}}%
\pgfpathlineto{\pgfqpoint{2.248443in}{1.797265in}}%
\pgfpathlineto{\pgfqpoint{2.250962in}{1.812069in}}%
\pgfpathlineto{\pgfqpoint{2.253481in}{1.786492in}}%
\pgfpathlineto{\pgfqpoint{2.255999in}{1.777935in}}%
\pgfpathlineto{\pgfqpoint{2.258518in}{1.802017in}}%
\pgfpathlineto{\pgfqpoint{2.261037in}{1.801330in}}%
\pgfpathlineto{\pgfqpoint{2.263556in}{1.812613in}}%
\pgfpathlineto{\pgfqpoint{2.266074in}{1.849168in}}%
\pgfpathlineto{\pgfqpoint{2.268593in}{1.867694in}}%
\pgfpathlineto{\pgfqpoint{2.271112in}{1.870476in}}%
\pgfpathlineto{\pgfqpoint{2.273631in}{1.889875in}}%
\pgfpathlineto{\pgfqpoint{2.276149in}{1.889384in}}%
\pgfpathlineto{\pgfqpoint{2.278668in}{1.917059in}}%
\pgfpathlineto{\pgfqpoint{2.281187in}{1.928917in}}%
\pgfpathlineto{\pgfqpoint{2.283706in}{1.944161in}}%
\pgfpathlineto{\pgfqpoint{2.286224in}{1.931490in}}%
\pgfpathlineto{\pgfqpoint{2.288743in}{1.922062in}}%
\pgfpathlineto{\pgfqpoint{2.291262in}{1.918208in}}%
\pgfpathlineto{\pgfqpoint{2.293781in}{1.922035in}}%
\pgfpathlineto{\pgfqpoint{2.296299in}{1.951476in}}%
\pgfpathlineto{\pgfqpoint{2.298818in}{1.942670in}}%
\pgfpathlineto{\pgfqpoint{2.301337in}{1.928350in}}%
\pgfpathlineto{\pgfqpoint{2.303856in}{1.928528in}}%
\pgfpathlineto{\pgfqpoint{2.306374in}{1.935004in}}%
\pgfpathlineto{\pgfqpoint{2.308893in}{1.936590in}}%
\pgfpathlineto{\pgfqpoint{2.311412in}{1.916626in}}%
\pgfpathlineto{\pgfqpoint{2.313931in}{1.938161in}}%
\pgfpathlineto{\pgfqpoint{2.316449in}{1.923608in}}%
\pgfpathlineto{\pgfqpoint{2.318968in}{1.921908in}}%
\pgfpathlineto{\pgfqpoint{2.321487in}{1.936433in}}%
\pgfpathlineto{\pgfqpoint{2.324006in}{1.906790in}}%
\pgfpathlineto{\pgfqpoint{2.326524in}{1.920865in}}%
\pgfpathlineto{\pgfqpoint{2.329043in}{1.904683in}}%
\pgfpathlineto{\pgfqpoint{2.331562in}{1.872803in}}%
\pgfpathlineto{\pgfqpoint{2.334081in}{1.870672in}}%
\pgfpathlineto{\pgfqpoint{2.336599in}{1.848168in}}%
\pgfpathlineto{\pgfqpoint{2.339118in}{1.855226in}}%
\pgfpathlineto{\pgfqpoint{2.341637in}{1.846737in}}%
\pgfpathlineto{\pgfqpoint{2.344156in}{1.858561in}}%
\pgfpathlineto{\pgfqpoint{2.346674in}{1.858923in}}%
\pgfpathlineto{\pgfqpoint{2.349193in}{1.855129in}}%
\pgfpathlineto{\pgfqpoint{2.351712in}{1.856177in}}%
\pgfpathlineto{\pgfqpoint{2.354231in}{1.840369in}}%
\pgfpathlineto{\pgfqpoint{2.356749in}{1.812088in}}%
\pgfpathlineto{\pgfqpoint{2.359268in}{1.837787in}}%
\pgfpathlineto{\pgfqpoint{2.361787in}{1.842683in}}%
\pgfpathlineto{\pgfqpoint{2.364306in}{1.845157in}}%
\pgfpathlineto{\pgfqpoint{2.366824in}{1.854449in}}%
\pgfpathlineto{\pgfqpoint{2.369343in}{1.857696in}}%
\pgfpathlineto{\pgfqpoint{2.371862in}{1.870456in}}%
\pgfpathlineto{\pgfqpoint{2.374381in}{1.880661in}}%
\pgfpathlineto{\pgfqpoint{2.376899in}{1.878190in}}%
\pgfpathlineto{\pgfqpoint{2.379418in}{1.866945in}}%
\pgfpathlineto{\pgfqpoint{2.381937in}{1.865501in}}%
\pgfpathlineto{\pgfqpoint{2.384456in}{1.850566in}}%
\pgfpathlineto{\pgfqpoint{2.386974in}{1.843340in}}%
\pgfpathlineto{\pgfqpoint{2.389493in}{1.895374in}}%
\pgfpathlineto{\pgfqpoint{2.392012in}{1.915839in}}%
\pgfpathlineto{\pgfqpoint{2.394531in}{1.903996in}}%
\pgfpathlineto{\pgfqpoint{2.397049in}{1.894864in}}%
\pgfpathlineto{\pgfqpoint{2.399568in}{1.879584in}}%
\pgfpathlineto{\pgfqpoint{2.402087in}{1.871358in}}%
\pgfpathlineto{\pgfqpoint{2.404606in}{1.880834in}}%
\pgfpathlineto{\pgfqpoint{2.407124in}{1.860988in}}%
\pgfpathlineto{\pgfqpoint{2.409643in}{1.858201in}}%
\pgfpathlineto{\pgfqpoint{2.412162in}{1.883491in}}%
\pgfpathlineto{\pgfqpoint{2.414681in}{1.922185in}}%
\pgfpathlineto{\pgfqpoint{2.417199in}{1.926166in}}%
\pgfpathlineto{\pgfqpoint{2.419718in}{1.919766in}}%
\pgfpathlineto{\pgfqpoint{2.422237in}{1.960173in}}%
\pgfpathlineto{\pgfqpoint{2.424756in}{1.960180in}}%
\pgfpathlineto{\pgfqpoint{2.427274in}{1.934643in}}%
\pgfpathlineto{\pgfqpoint{2.429793in}{1.945471in}}%
\pgfpathlineto{\pgfqpoint{2.432312in}{1.929964in}}%
\pgfpathlineto{\pgfqpoint{2.434831in}{1.934875in}}%
\pgfpathlineto{\pgfqpoint{2.437349in}{1.921511in}}%
\pgfpathlineto{\pgfqpoint{2.439868in}{1.925833in}}%
\pgfpathlineto{\pgfqpoint{2.442387in}{1.935130in}}%
\pgfpathlineto{\pgfqpoint{2.444906in}{1.947978in}}%
\pgfpathlineto{\pgfqpoint{2.447424in}{1.943351in}}%
\pgfpathlineto{\pgfqpoint{2.449943in}{1.959654in}}%
\pgfpathlineto{\pgfqpoint{2.452462in}{1.967246in}}%
\pgfpathlineto{\pgfqpoint{2.454981in}{1.997391in}}%
\pgfpathlineto{\pgfqpoint{2.457499in}{1.988181in}}%
\pgfpathlineto{\pgfqpoint{2.460018in}{1.981337in}}%
\pgfpathlineto{\pgfqpoint{2.462537in}{1.986883in}}%
\pgfpathlineto{\pgfqpoint{2.465056in}{2.012941in}}%
\pgfpathlineto{\pgfqpoint{2.467574in}{2.007650in}}%
\pgfpathlineto{\pgfqpoint{2.470093in}{2.017203in}}%
\pgfpathlineto{\pgfqpoint{2.472612in}{2.004698in}}%
\pgfpathlineto{\pgfqpoint{2.475131in}{1.999774in}}%
\pgfpathlineto{\pgfqpoint{2.477649in}{1.971921in}}%
\pgfpathlineto{\pgfqpoint{2.480168in}{1.963659in}}%
\pgfpathlineto{\pgfqpoint{2.482687in}{1.976604in}}%
\pgfpathlineto{\pgfqpoint{2.485206in}{1.957884in}}%
\pgfpathlineto{\pgfqpoint{2.487724in}{1.961903in}}%
\pgfpathlineto{\pgfqpoint{2.490243in}{1.962484in}}%
\pgfpathlineto{\pgfqpoint{2.492762in}{1.978093in}}%
\pgfpathlineto{\pgfqpoint{2.495281in}{1.976171in}}%
\pgfpathlineto{\pgfqpoint{2.497799in}{1.957083in}}%
\pgfpathlineto{\pgfqpoint{2.500318in}{1.949243in}}%
\pgfpathlineto{\pgfqpoint{2.502837in}{1.964486in}}%
\pgfpathlineto{\pgfqpoint{2.505356in}{1.963224in}}%
\pgfpathlineto{\pgfqpoint{2.507874in}{1.980872in}}%
\pgfpathlineto{\pgfqpoint{2.510393in}{1.990178in}}%
\pgfpathlineto{\pgfqpoint{2.512912in}{2.049398in}}%
\pgfpathlineto{\pgfqpoint{2.515431in}{2.022691in}}%
\pgfpathlineto{\pgfqpoint{2.517949in}{2.022794in}}%
\pgfpathlineto{\pgfqpoint{2.520468in}{2.044377in}}%
\pgfpathlineto{\pgfqpoint{2.522987in}{2.053305in}}%
\pgfpathlineto{\pgfqpoint{2.525506in}{2.064299in}}%
\pgfpathlineto{\pgfqpoint{2.528024in}{2.094161in}}%
\pgfpathlineto{\pgfqpoint{2.530543in}{2.059132in}}%
\pgfpathlineto{\pgfqpoint{2.533062in}{2.069632in}}%
\pgfpathlineto{\pgfqpoint{2.535581in}{2.041916in}}%
\pgfpathlineto{\pgfqpoint{2.538099in}{2.053051in}}%
\pgfpathlineto{\pgfqpoint{2.540618in}{2.075788in}}%
\pgfpathlineto{\pgfqpoint{2.543137in}{2.053979in}}%
\pgfpathlineto{\pgfqpoint{2.545656in}{2.046036in}}%
\pgfpathlineto{\pgfqpoint{2.548174in}{2.052876in}}%
\pgfpathlineto{\pgfqpoint{2.550693in}{2.079013in}}%
\pgfpathlineto{\pgfqpoint{2.553212in}{2.077001in}}%
\pgfpathlineto{\pgfqpoint{2.555731in}{2.059538in}}%
\pgfpathlineto{\pgfqpoint{2.558249in}{2.104629in}}%
\pgfpathlineto{\pgfqpoint{2.560768in}{2.102106in}}%
\pgfpathlineto{\pgfqpoint{2.563287in}{2.092085in}}%
\pgfpathlineto{\pgfqpoint{2.565806in}{2.104264in}}%
\pgfpathlineto{\pgfqpoint{2.568324in}{2.125629in}}%
\pgfpathlineto{\pgfqpoint{2.570843in}{2.082661in}}%
\pgfpathlineto{\pgfqpoint{2.573362in}{2.090673in}}%
\pgfpathlineto{\pgfqpoint{2.575881in}{2.095654in}}%
\pgfpathlineto{\pgfqpoint{2.578399in}{2.110446in}}%
\pgfpathlineto{\pgfqpoint{2.580918in}{2.108607in}}%
\pgfpathlineto{\pgfqpoint{2.583437in}{2.118652in}}%
\pgfpathlineto{\pgfqpoint{2.585956in}{2.147275in}}%
\pgfpathlineto{\pgfqpoint{2.588474in}{2.193541in}}%
\pgfpathlineto{\pgfqpoint{2.590993in}{2.196862in}}%
\pgfpathlineto{\pgfqpoint{2.593512in}{2.198990in}}%
\pgfpathlineto{\pgfqpoint{2.596031in}{2.183721in}}%
\pgfpathlineto{\pgfqpoint{2.598549in}{2.201103in}}%
\pgfpathlineto{\pgfqpoint{2.601068in}{2.206502in}}%
\pgfpathlineto{\pgfqpoint{2.603587in}{2.208771in}}%
\pgfpathlineto{\pgfqpoint{2.606106in}{2.245299in}}%
\pgfpathlineto{\pgfqpoint{2.608624in}{2.263286in}}%
\pgfpathlineto{\pgfqpoint{2.611143in}{2.310827in}}%
\pgfpathlineto{\pgfqpoint{2.613662in}{2.316762in}}%
\pgfpathlineto{\pgfqpoint{2.616181in}{2.352655in}}%
\pgfpathlineto{\pgfqpoint{2.618699in}{2.339576in}}%
\pgfpathlineto{\pgfqpoint{2.623737in}{2.293168in}}%
\pgfpathlineto{\pgfqpoint{2.626256in}{2.363921in}}%
\pgfpathlineto{\pgfqpoint{2.628774in}{2.373330in}}%
\pgfpathlineto{\pgfqpoint{2.631293in}{2.388048in}}%
\pgfpathlineto{\pgfqpoint{2.633812in}{2.395886in}}%
\pgfpathlineto{\pgfqpoint{2.636331in}{2.437585in}}%
\pgfpathlineto{\pgfqpoint{2.638849in}{2.418676in}}%
\pgfpathlineto{\pgfqpoint{2.641368in}{2.479269in}}%
\pgfpathlineto{\pgfqpoint{2.643887in}{2.487717in}}%
\pgfpathlineto{\pgfqpoint{2.646406in}{2.445925in}}%
\pgfpathlineto{\pgfqpoint{2.648924in}{2.444598in}}%
\pgfpathlineto{\pgfqpoint{2.651443in}{2.436851in}}%
\pgfpathlineto{\pgfqpoint{2.653962in}{2.441345in}}%
\pgfpathlineto{\pgfqpoint{2.656481in}{2.444208in}}%
\pgfpathlineto{\pgfqpoint{2.658999in}{2.431258in}}%
\pgfpathlineto{\pgfqpoint{2.661518in}{2.441688in}}%
\pgfpathlineto{\pgfqpoint{2.664037in}{2.412510in}}%
\pgfpathlineto{\pgfqpoint{2.666556in}{2.442476in}}%
\pgfpathlineto{\pgfqpoint{2.669074in}{2.455924in}}%
\pgfpathlineto{\pgfqpoint{2.671593in}{2.464036in}}%
\pgfpathlineto{\pgfqpoint{2.674112in}{2.473379in}}%
\pgfpathlineto{\pgfqpoint{2.676631in}{2.441981in}}%
\pgfpathlineto{\pgfqpoint{2.681668in}{2.496227in}}%
\pgfpathlineto{\pgfqpoint{2.684187in}{2.508096in}}%
\pgfpathlineto{\pgfqpoint{2.686706in}{2.497639in}}%
\pgfpathlineto{\pgfqpoint{2.689224in}{2.494796in}}%
\pgfpathlineto{\pgfqpoint{2.691743in}{2.458491in}}%
\pgfpathlineto{\pgfqpoint{2.694262in}{2.459727in}}%
\pgfpathlineto{\pgfqpoint{2.696781in}{2.465661in}}%
\pgfpathlineto{\pgfqpoint{2.699299in}{2.450341in}}%
\pgfpathlineto{\pgfqpoint{2.701818in}{2.491822in}}%
\pgfpathlineto{\pgfqpoint{2.704337in}{2.470746in}}%
\pgfpathlineto{\pgfqpoint{2.706856in}{2.459337in}}%
\pgfpathlineto{\pgfqpoint{2.709374in}{2.464204in}}%
\pgfpathlineto{\pgfqpoint{2.711893in}{2.443688in}}%
\pgfpathlineto{\pgfqpoint{2.714412in}{2.408419in}}%
\pgfpathlineto{\pgfqpoint{2.716931in}{2.427152in}}%
\pgfpathlineto{\pgfqpoint{2.719449in}{2.443420in}}%
\pgfpathlineto{\pgfqpoint{2.721968in}{2.454458in}}%
\pgfpathlineto{\pgfqpoint{2.724487in}{2.448245in}}%
\pgfpathlineto{\pgfqpoint{2.727006in}{2.451584in}}%
\pgfpathlineto{\pgfqpoint{2.729524in}{2.471809in}}%
\pgfpathlineto{\pgfqpoint{2.732043in}{2.504456in}}%
\pgfpathlineto{\pgfqpoint{2.734562in}{2.460757in}}%
\pgfpathlineto{\pgfqpoint{2.737081in}{2.460849in}}%
\pgfpathlineto{\pgfqpoint{2.739599in}{2.479389in}}%
\pgfpathlineto{\pgfqpoint{2.742118in}{2.506767in}}%
\pgfpathlineto{\pgfqpoint{2.744637in}{2.488606in}}%
\pgfpathlineto{\pgfqpoint{2.747156in}{2.474033in}}%
\pgfpathlineto{\pgfqpoint{2.749674in}{2.479369in}}%
\pgfpathlineto{\pgfqpoint{2.752193in}{2.456384in}}%
\pgfpathlineto{\pgfqpoint{2.754712in}{2.454655in}}%
\pgfpathlineto{\pgfqpoint{2.757231in}{2.468413in}}%
\pgfpathlineto{\pgfqpoint{2.759749in}{2.475249in}}%
\pgfpathlineto{\pgfqpoint{2.762268in}{2.508968in}}%
\pgfpathlineto{\pgfqpoint{2.764787in}{2.483232in}}%
\pgfpathlineto{\pgfqpoint{2.767306in}{2.484240in}}%
\pgfpathlineto{\pgfqpoint{2.769824in}{2.492030in}}%
\pgfpathlineto{\pgfqpoint{2.772343in}{2.507638in}}%
\pgfpathlineto{\pgfqpoint{2.774862in}{2.489423in}}%
\pgfpathlineto{\pgfqpoint{2.777381in}{2.460668in}}%
\pgfpathlineto{\pgfqpoint{2.779899in}{2.427893in}}%
\pgfpathlineto{\pgfqpoint{2.782418in}{2.429036in}}%
\pgfpathlineto{\pgfqpoint{2.784937in}{2.406970in}}%
\pgfpathlineto{\pgfqpoint{2.787456in}{2.429328in}}%
\pgfpathlineto{\pgfqpoint{2.789974in}{2.413633in}}%
\pgfpathlineto{\pgfqpoint{2.792493in}{2.425983in}}%
\pgfpathlineto{\pgfqpoint{2.795012in}{2.462410in}}%
\pgfpathlineto{\pgfqpoint{2.797531in}{2.506319in}}%
\pgfpathlineto{\pgfqpoint{2.802568in}{2.528448in}}%
\pgfpathlineto{\pgfqpoint{2.805087in}{2.540898in}}%
\pgfpathlineto{\pgfqpoint{2.807606in}{2.552097in}}%
\pgfpathlineto{\pgfqpoint{2.810124in}{2.562869in}}%
\pgfpathlineto{\pgfqpoint{2.812643in}{2.582122in}}%
\pgfpathlineto{\pgfqpoint{2.815162in}{2.581848in}}%
\pgfpathlineto{\pgfqpoint{2.817681in}{2.595720in}}%
\pgfpathlineto{\pgfqpoint{2.820199in}{2.593733in}}%
\pgfpathlineto{\pgfqpoint{2.822718in}{2.576838in}}%
\pgfpathlineto{\pgfqpoint{2.825237in}{2.590787in}}%
\pgfpathlineto{\pgfqpoint{2.827756in}{2.616537in}}%
\pgfpathlineto{\pgfqpoint{2.830274in}{2.611140in}}%
\pgfpathlineto{\pgfqpoint{2.832793in}{2.632631in}}%
\pgfpathlineto{\pgfqpoint{2.835312in}{2.627118in}}%
\pgfpathlineto{\pgfqpoint{2.837831in}{2.607134in}}%
\pgfpathlineto{\pgfqpoint{2.840349in}{2.585825in}}%
\pgfpathlineto{\pgfqpoint{2.842868in}{2.598600in}}%
\pgfpathlineto{\pgfqpoint{2.845387in}{2.578254in}}%
\pgfpathlineto{\pgfqpoint{2.847906in}{2.562453in}}%
\pgfpathlineto{\pgfqpoint{2.850424in}{2.581212in}}%
\pgfpathlineto{\pgfqpoint{2.852943in}{2.574482in}}%
\pgfpathlineto{\pgfqpoint{2.855462in}{2.577197in}}%
\pgfpathlineto{\pgfqpoint{2.857981in}{2.560139in}}%
\pgfpathlineto{\pgfqpoint{2.860499in}{2.566160in}}%
\pgfpathlineto{\pgfqpoint{2.863018in}{2.583508in}}%
\pgfpathlineto{\pgfqpoint{2.868056in}{2.626124in}}%
\pgfpathlineto{\pgfqpoint{2.870574in}{2.590878in}}%
\pgfpathlineto{\pgfqpoint{2.873093in}{2.629273in}}%
\pgfpathlineto{\pgfqpoint{2.875612in}{2.655810in}}%
\pgfpathlineto{\pgfqpoint{2.878131in}{2.635150in}}%
\pgfpathlineto{\pgfqpoint{2.880649in}{2.594623in}}%
\pgfpathlineto{\pgfqpoint{2.883168in}{2.570829in}}%
\pgfpathlineto{\pgfqpoint{2.885687in}{2.543068in}}%
\pgfpathlineto{\pgfqpoint{2.888206in}{2.551900in}}%
\pgfpathlineto{\pgfqpoint{2.890724in}{2.576911in}}%
\pgfpathlineto{\pgfqpoint{2.893243in}{2.595350in}}%
\pgfpathlineto{\pgfqpoint{2.895762in}{2.552487in}}%
\pgfpathlineto{\pgfqpoint{2.898281in}{2.599178in}}%
\pgfpathlineto{\pgfqpoint{2.900799in}{2.607981in}}%
\pgfpathlineto{\pgfqpoint{2.903318in}{2.555851in}}%
\pgfpathlineto{\pgfqpoint{2.905837in}{2.589099in}}%
\pgfpathlineto{\pgfqpoint{2.908356in}{2.560819in}}%
\pgfpathlineto{\pgfqpoint{2.910874in}{2.528017in}}%
\pgfpathlineto{\pgfqpoint{2.913393in}{2.553169in}}%
\pgfpathlineto{\pgfqpoint{2.915912in}{2.595199in}}%
\pgfpathlineto{\pgfqpoint{2.918431in}{2.593539in}}%
\pgfpathlineto{\pgfqpoint{2.920949in}{2.664621in}}%
\pgfpathlineto{\pgfqpoint{2.923468in}{2.672121in}}%
\pgfpathlineto{\pgfqpoint{2.925987in}{2.664655in}}%
\pgfpathlineto{\pgfqpoint{2.928506in}{2.682371in}}%
\pgfpathlineto{\pgfqpoint{2.931024in}{2.728997in}}%
\pgfpathlineto{\pgfqpoint{2.933543in}{2.732939in}}%
\pgfpathlineto{\pgfqpoint{2.936062in}{2.753395in}}%
\pgfpathlineto{\pgfqpoint{2.938581in}{2.749454in}}%
\pgfpathlineto{\pgfqpoint{2.941099in}{2.736043in}}%
\pgfpathlineto{\pgfqpoint{2.943618in}{2.761637in}}%
\pgfpathlineto{\pgfqpoint{2.946137in}{2.764479in}}%
\pgfpathlineto{\pgfqpoint{2.948656in}{2.785375in}}%
\pgfpathlineto{\pgfqpoint{2.951174in}{2.763427in}}%
\pgfpathlineto{\pgfqpoint{2.953693in}{2.761572in}}%
\pgfpathlineto{\pgfqpoint{2.956212in}{2.769697in}}%
\pgfpathlineto{\pgfqpoint{2.958731in}{2.792787in}}%
\pgfpathlineto{\pgfqpoint{2.961249in}{2.746791in}}%
\pgfpathlineto{\pgfqpoint{2.963768in}{2.715436in}}%
\pgfpathlineto{\pgfqpoint{2.966287in}{2.689817in}}%
\pgfpathlineto{\pgfqpoint{2.968806in}{2.695300in}}%
\pgfpathlineto{\pgfqpoint{2.971324in}{2.708306in}}%
\pgfpathlineto{\pgfqpoint{2.973843in}{2.712128in}}%
\pgfpathlineto{\pgfqpoint{2.976362in}{2.752902in}}%
\pgfpathlineto{\pgfqpoint{2.978881in}{2.716962in}}%
\pgfpathlineto{\pgfqpoint{2.981399in}{2.730149in}}%
\pgfpathlineto{\pgfqpoint{2.983918in}{2.721668in}}%
\pgfpathlineto{\pgfqpoint{2.986437in}{2.707709in}}%
\pgfpathlineto{\pgfqpoint{2.988956in}{2.715678in}}%
\pgfpathlineto{\pgfqpoint{2.991474in}{2.725241in}}%
\pgfpathlineto{\pgfqpoint{2.993993in}{2.704023in}}%
\pgfpathlineto{\pgfqpoint{2.996512in}{2.708360in}}%
\pgfpathlineto{\pgfqpoint{2.999031in}{2.682013in}}%
\pgfpathlineto{\pgfqpoint{3.001549in}{2.666985in}}%
\pgfpathlineto{\pgfqpoint{3.004068in}{2.698455in}}%
\pgfpathlineto{\pgfqpoint{3.006587in}{2.739719in}}%
\pgfpathlineto{\pgfqpoint{3.009106in}{2.752679in}}%
\pgfpathlineto{\pgfqpoint{3.011624in}{2.785975in}}%
\pgfpathlineto{\pgfqpoint{3.014143in}{2.759560in}}%
\pgfpathlineto{\pgfqpoint{3.016662in}{2.775528in}}%
\pgfpathlineto{\pgfqpoint{3.019181in}{2.779171in}}%
\pgfpathlineto{\pgfqpoint{3.024218in}{2.786046in}}%
\pgfpathlineto{\pgfqpoint{3.026737in}{2.784666in}}%
\pgfpathlineto{\pgfqpoint{3.029256in}{2.767668in}}%
\pgfpathlineto{\pgfqpoint{3.031774in}{2.787459in}}%
\pgfpathlineto{\pgfqpoint{3.034293in}{2.783694in}}%
\pgfpathlineto{\pgfqpoint{3.036812in}{2.750341in}}%
\pgfpathlineto{\pgfqpoint{3.039331in}{2.769990in}}%
\pgfpathlineto{\pgfqpoint{3.041849in}{2.783167in}}%
\pgfpathlineto{\pgfqpoint{3.044368in}{2.776968in}}%
\pgfpathlineto{\pgfqpoint{3.046887in}{2.800052in}}%
\pgfpathlineto{\pgfqpoint{3.049406in}{2.804554in}}%
\pgfpathlineto{\pgfqpoint{3.051924in}{2.787414in}}%
\pgfpathlineto{\pgfqpoint{3.054443in}{2.812215in}}%
\pgfpathlineto{\pgfqpoint{3.056962in}{2.817117in}}%
\pgfpathlineto{\pgfqpoint{3.059481in}{2.842883in}}%
\pgfpathlineto{\pgfqpoint{3.061999in}{2.831355in}}%
\pgfpathlineto{\pgfqpoint{3.064518in}{2.866396in}}%
\pgfpathlineto{\pgfqpoint{3.067037in}{2.840570in}}%
\pgfpathlineto{\pgfqpoint{3.069556in}{2.870337in}}%
\pgfpathlineto{\pgfqpoint{3.072074in}{2.887001in}}%
\pgfpathlineto{\pgfqpoint{3.074593in}{2.878297in}}%
\pgfpathlineto{\pgfqpoint{3.077112in}{2.906306in}}%
\pgfpathlineto{\pgfqpoint{3.079631in}{2.960958in}}%
\pgfpathlineto{\pgfqpoint{3.082149in}{2.979184in}}%
\pgfpathlineto{\pgfqpoint{3.084668in}{2.973546in}}%
\pgfpathlineto{\pgfqpoint{3.087187in}{2.946185in}}%
\pgfpathlineto{\pgfqpoint{3.089706in}{2.933230in}}%
\pgfpathlineto{\pgfqpoint{3.092224in}{2.901961in}}%
\pgfpathlineto{\pgfqpoint{3.094743in}{2.877331in}}%
\pgfpathlineto{\pgfqpoint{3.097262in}{2.874929in}}%
\pgfpathlineto{\pgfqpoint{3.099781in}{2.870699in}}%
\pgfpathlineto{\pgfqpoint{3.102299in}{2.888794in}}%
\pgfpathlineto{\pgfqpoint{3.104818in}{2.909598in}}%
\pgfpathlineto{\pgfqpoint{3.107337in}{2.873212in}}%
\pgfpathlineto{\pgfqpoint{3.109856in}{2.867197in}}%
\pgfpathlineto{\pgfqpoint{3.112374in}{2.851020in}}%
\pgfpathlineto{\pgfqpoint{3.114893in}{2.842998in}}%
\pgfpathlineto{\pgfqpoint{3.117412in}{2.835441in}}%
\pgfpathlineto{\pgfqpoint{3.119931in}{2.844761in}}%
\pgfpathlineto{\pgfqpoint{3.122449in}{2.836200in}}%
\pgfpathlineto{\pgfqpoint{3.124968in}{2.886753in}}%
\pgfpathlineto{\pgfqpoint{3.127487in}{2.866182in}}%
\pgfpathlineto{\pgfqpoint{3.130006in}{2.872079in}}%
\pgfpathlineto{\pgfqpoint{3.132524in}{2.874503in}}%
\pgfpathlineto{\pgfqpoint{3.135043in}{2.867791in}}%
\pgfpathlineto{\pgfqpoint{3.137562in}{2.895797in}}%
\pgfpathlineto{\pgfqpoint{3.140081in}{2.884318in}}%
\pgfpathlineto{\pgfqpoint{3.142599in}{2.871444in}}%
\pgfpathlineto{\pgfqpoint{3.145118in}{2.869640in}}%
\pgfpathlineto{\pgfqpoint{3.147637in}{2.856237in}}%
\pgfpathlineto{\pgfqpoint{3.150156in}{2.835949in}}%
\pgfpathlineto{\pgfqpoint{3.152674in}{2.858930in}}%
\pgfpathlineto{\pgfqpoint{3.155193in}{2.876742in}}%
\pgfpathlineto{\pgfqpoint{3.157712in}{2.861203in}}%
\pgfpathlineto{\pgfqpoint{3.160231in}{2.870886in}}%
\pgfpathlineto{\pgfqpoint{3.162749in}{2.921813in}}%
\pgfpathlineto{\pgfqpoint{3.165268in}{2.939014in}}%
\pgfpathlineto{\pgfqpoint{3.167787in}{2.920759in}}%
\pgfpathlineto{\pgfqpoint{3.170306in}{2.893915in}}%
\pgfpathlineto{\pgfqpoint{3.172824in}{2.875003in}}%
\pgfpathlineto{\pgfqpoint{3.175343in}{2.904732in}}%
\pgfpathlineto{\pgfqpoint{3.177862in}{2.900237in}}%
\pgfpathlineto{\pgfqpoint{3.180381in}{2.901423in}}%
\pgfpathlineto{\pgfqpoint{3.182899in}{2.888246in}}%
\pgfpathlineto{\pgfqpoint{3.185418in}{2.855243in}}%
\pgfpathlineto{\pgfqpoint{3.187937in}{2.876195in}}%
\pgfpathlineto{\pgfqpoint{3.190456in}{2.911322in}}%
\pgfpathlineto{\pgfqpoint{3.192974in}{2.928278in}}%
\pgfpathlineto{\pgfqpoint{3.195493in}{2.937026in}}%
\pgfpathlineto{\pgfqpoint{3.198012in}{2.923912in}}%
\pgfpathlineto{\pgfqpoint{3.200531in}{2.919005in}}%
\pgfpathlineto{\pgfqpoint{3.203049in}{2.928574in}}%
\pgfpathlineto{\pgfqpoint{3.205568in}{2.939442in}}%
\pgfpathlineto{\pgfqpoint{3.208087in}{2.902028in}}%
\pgfpathlineto{\pgfqpoint{3.210606in}{2.914863in}}%
\pgfpathlineto{\pgfqpoint{3.213124in}{2.919339in}}%
\pgfpathlineto{\pgfqpoint{3.215643in}{2.916833in}}%
\pgfpathlineto{\pgfqpoint{3.218162in}{2.921243in}}%
\pgfpathlineto{\pgfqpoint{3.220681in}{2.886166in}}%
\pgfpathlineto{\pgfqpoint{3.225718in}{2.954835in}}%
\pgfpathlineto{\pgfqpoint{3.228237in}{2.947084in}}%
\pgfpathlineto{\pgfqpoint{3.230756in}{2.930707in}}%
\pgfpathlineto{\pgfqpoint{3.233274in}{2.895630in}}%
\pgfpathlineto{\pgfqpoint{3.235793in}{2.896701in}}%
\pgfpathlineto{\pgfqpoint{3.238312in}{2.894289in}}%
\pgfpathlineto{\pgfqpoint{3.240831in}{2.900112in}}%
\pgfpathlineto{\pgfqpoint{3.243349in}{2.861270in}}%
\pgfpathlineto{\pgfqpoint{3.245868in}{2.864570in}}%
\pgfpathlineto{\pgfqpoint{3.248387in}{2.849789in}}%
\pgfpathlineto{\pgfqpoint{3.250906in}{2.846727in}}%
\pgfpathlineto{\pgfqpoint{3.253424in}{2.828127in}}%
\pgfpathlineto{\pgfqpoint{3.255943in}{2.819120in}}%
\pgfpathlineto{\pgfqpoint{3.258462in}{2.832664in}}%
\pgfpathlineto{\pgfqpoint{3.260981in}{2.875760in}}%
\pgfpathlineto{\pgfqpoint{3.263499in}{2.903076in}}%
\pgfpathlineto{\pgfqpoint{3.266018in}{2.869689in}}%
\pgfpathlineto{\pgfqpoint{3.268537in}{2.829296in}}%
\pgfpathlineto{\pgfqpoint{3.271056in}{2.846498in}}%
\pgfpathlineto{\pgfqpoint{3.273574in}{2.892724in}}%
\pgfpathlineto{\pgfqpoint{3.276093in}{2.892268in}}%
\pgfpathlineto{\pgfqpoint{3.278612in}{2.956996in}}%
\pgfpathlineto{\pgfqpoint{3.281131in}{2.992086in}}%
\pgfpathlineto{\pgfqpoint{3.283649in}{2.993659in}}%
\pgfpathlineto{\pgfqpoint{3.286168in}{2.987962in}}%
\pgfpathlineto{\pgfqpoint{3.288687in}{2.983682in}}%
\pgfpathlineto{\pgfqpoint{3.291206in}{2.994044in}}%
\pgfpathlineto{\pgfqpoint{3.293724in}{2.999341in}}%
\pgfpathlineto{\pgfqpoint{3.296243in}{3.001837in}}%
\pgfpathlineto{\pgfqpoint{3.298762in}{2.978331in}}%
\pgfpathlineto{\pgfqpoint{3.301281in}{2.986744in}}%
\pgfpathlineto{\pgfqpoint{3.303799in}{2.975025in}}%
\pgfpathlineto{\pgfqpoint{3.306318in}{2.977170in}}%
\pgfpathlineto{\pgfqpoint{3.308837in}{3.001058in}}%
\pgfpathlineto{\pgfqpoint{3.311356in}{2.986209in}}%
\pgfpathlineto{\pgfqpoint{3.313874in}{2.980714in}}%
\pgfpathlineto{\pgfqpoint{3.316393in}{2.987366in}}%
\pgfpathlineto{\pgfqpoint{3.318912in}{2.984512in}}%
\pgfpathlineto{\pgfqpoint{3.321431in}{2.928985in}}%
\pgfpathlineto{\pgfqpoint{3.323949in}{2.925902in}}%
\pgfpathlineto{\pgfqpoint{3.326468in}{2.938753in}}%
\pgfpathlineto{\pgfqpoint{3.331506in}{2.883181in}}%
\pgfpathlineto{\pgfqpoint{3.334024in}{2.929474in}}%
\pgfpathlineto{\pgfqpoint{3.336543in}{2.938176in}}%
\pgfpathlineto{\pgfqpoint{3.339062in}{2.950331in}}%
\pgfpathlineto{\pgfqpoint{3.341581in}{2.981531in}}%
\pgfpathlineto{\pgfqpoint{3.344099in}{2.995514in}}%
\pgfpathlineto{\pgfqpoint{3.346618in}{3.016908in}}%
\pgfpathlineto{\pgfqpoint{3.349137in}{3.015771in}}%
\pgfpathlineto{\pgfqpoint{3.351656in}{3.033095in}}%
\pgfpathlineto{\pgfqpoint{3.354174in}{3.044091in}}%
\pgfpathlineto{\pgfqpoint{3.356693in}{3.021931in}}%
\pgfpathlineto{\pgfqpoint{3.359212in}{3.016175in}}%
\pgfpathlineto{\pgfqpoint{3.361731in}{3.051555in}}%
\pgfpathlineto{\pgfqpoint{3.364249in}{3.019692in}}%
\pgfpathlineto{\pgfqpoint{3.366768in}{3.018595in}}%
\pgfpathlineto{\pgfqpoint{3.369287in}{3.014370in}}%
\pgfpathlineto{\pgfqpoint{3.371806in}{3.020984in}}%
\pgfpathlineto{\pgfqpoint{3.374324in}{3.009473in}}%
\pgfpathlineto{\pgfqpoint{3.376843in}{2.994270in}}%
\pgfpathlineto{\pgfqpoint{3.379362in}{2.984313in}}%
\pgfpathlineto{\pgfqpoint{3.381881in}{2.988271in}}%
\pgfpathlineto{\pgfqpoint{3.384399in}{3.003070in}}%
\pgfpathlineto{\pgfqpoint{3.386918in}{3.034225in}}%
\pgfpathlineto{\pgfqpoint{3.389437in}{3.051173in}}%
\pgfpathlineto{\pgfqpoint{3.391956in}{3.055896in}}%
\pgfpathlineto{\pgfqpoint{3.394474in}{3.038962in}}%
\pgfpathlineto{\pgfqpoint{3.396993in}{3.077877in}}%
\pgfpathlineto{\pgfqpoint{3.399512in}{3.073275in}}%
\pgfpathlineto{\pgfqpoint{3.402031in}{3.068920in}}%
\pgfpathlineto{\pgfqpoint{3.404549in}{3.046985in}}%
\pgfpathlineto{\pgfqpoint{3.409587in}{3.010572in}}%
\pgfpathlineto{\pgfqpoint{3.414624in}{3.069391in}}%
\pgfpathlineto{\pgfqpoint{3.417143in}{3.056201in}}%
\pgfpathlineto{\pgfqpoint{3.419662in}{3.063073in}}%
\pgfpathlineto{\pgfqpoint{3.422181in}{3.042613in}}%
\pgfpathlineto{\pgfqpoint{3.424699in}{3.052818in}}%
\pgfpathlineto{\pgfqpoint{3.427218in}{3.058467in}}%
\pgfpathlineto{\pgfqpoint{3.429737in}{3.063358in}}%
\pgfpathlineto{\pgfqpoint{3.432256in}{3.084626in}}%
\pgfpathlineto{\pgfqpoint{3.434774in}{3.093972in}}%
\pgfpathlineto{\pgfqpoint{3.437293in}{3.072760in}}%
\pgfpathlineto{\pgfqpoint{3.439812in}{3.108869in}}%
\pgfpathlineto{\pgfqpoint{3.442331in}{3.121753in}}%
\pgfpathlineto{\pgfqpoint{3.444849in}{3.108927in}}%
\pgfpathlineto{\pgfqpoint{3.447368in}{3.125466in}}%
\pgfpathlineto{\pgfqpoint{3.449887in}{3.115064in}}%
\pgfpathlineto{\pgfqpoint{3.452406in}{3.147006in}}%
\pgfpathlineto{\pgfqpoint{3.454924in}{3.096006in}}%
\pgfpathlineto{\pgfqpoint{3.457443in}{3.115949in}}%
\pgfpathlineto{\pgfqpoint{3.459962in}{3.122646in}}%
\pgfpathlineto{\pgfqpoint{3.462481in}{3.087701in}}%
\pgfpathlineto{\pgfqpoint{3.464999in}{3.082066in}}%
\pgfpathlineto{\pgfqpoint{3.470037in}{3.104824in}}%
\pgfpathlineto{\pgfqpoint{3.472556in}{3.096940in}}%
\pgfpathlineto{\pgfqpoint{3.475074in}{3.119099in}}%
\pgfpathlineto{\pgfqpoint{3.477593in}{3.140165in}}%
\pgfpathlineto{\pgfqpoint{3.480112in}{3.132585in}}%
\pgfpathlineto{\pgfqpoint{3.482631in}{3.104349in}}%
\pgfpathlineto{\pgfqpoint{3.485149in}{3.130391in}}%
\pgfpathlineto{\pgfqpoint{3.487668in}{3.138759in}}%
\pgfpathlineto{\pgfqpoint{3.490187in}{3.112521in}}%
\pgfpathlineto{\pgfqpoint{3.492706in}{3.102210in}}%
\pgfpathlineto{\pgfqpoint{3.495224in}{3.113690in}}%
\pgfpathlineto{\pgfqpoint{3.497743in}{3.096847in}}%
\pgfpathlineto{\pgfqpoint{3.500262in}{3.076463in}}%
\pgfpathlineto{\pgfqpoint{3.502781in}{3.082282in}}%
\pgfpathlineto{\pgfqpoint{3.505299in}{3.072876in}}%
\pgfpathlineto{\pgfqpoint{3.507818in}{3.076837in}}%
\pgfpathlineto{\pgfqpoint{3.510337in}{3.049598in}}%
\pgfpathlineto{\pgfqpoint{3.512856in}{3.031934in}}%
\pgfpathlineto{\pgfqpoint{3.515374in}{3.033150in}}%
\pgfpathlineto{\pgfqpoint{3.517893in}{3.010212in}}%
\pgfpathlineto{\pgfqpoint{3.520412in}{2.997570in}}%
\pgfpathlineto{\pgfqpoint{3.522931in}{2.976339in}}%
\pgfpathlineto{\pgfqpoint{3.525449in}{3.004993in}}%
\pgfpathlineto{\pgfqpoint{3.527968in}{3.004324in}}%
\pgfpathlineto{\pgfqpoint{3.530487in}{3.028656in}}%
\pgfpathlineto{\pgfqpoint{3.533006in}{3.050100in}}%
\pgfpathlineto{\pgfqpoint{3.535524in}{3.078180in}}%
\pgfpathlineto{\pgfqpoint{3.538043in}{3.080694in}}%
\pgfpathlineto{\pgfqpoint{3.540562in}{3.079085in}}%
\pgfpathlineto{\pgfqpoint{3.543081in}{3.075808in}}%
\pgfpathlineto{\pgfqpoint{3.545599in}{3.078605in}}%
\pgfpathlineto{\pgfqpoint{3.548118in}{3.112072in}}%
\pgfpathlineto{\pgfqpoint{3.550637in}{3.090268in}}%
\pgfpathlineto{\pgfqpoint{3.553156in}{3.075340in}}%
\pgfpathlineto{\pgfqpoint{3.555674in}{3.078180in}}%
\pgfpathlineto{\pgfqpoint{3.558193in}{3.066155in}}%
\pgfpathlineto{\pgfqpoint{3.560712in}{3.095159in}}%
\pgfpathlineto{\pgfqpoint{3.563231in}{3.136521in}}%
\pgfpathlineto{\pgfqpoint{3.565749in}{3.135821in}}%
\pgfpathlineto{\pgfqpoint{3.568268in}{3.164307in}}%
\pgfpathlineto{\pgfqpoint{3.570787in}{3.154300in}}%
\pgfpathlineto{\pgfqpoint{3.573306in}{3.108333in}}%
\pgfpathlineto{\pgfqpoint{3.575824in}{3.102004in}}%
\pgfpathlineto{\pgfqpoint{3.578343in}{3.099417in}}%
\pgfpathlineto{\pgfqpoint{3.580862in}{3.060038in}}%
\pgfpathlineto{\pgfqpoint{3.583381in}{3.057167in}}%
\pgfpathlineto{\pgfqpoint{3.585899in}{3.051966in}}%
\pgfpathlineto{\pgfqpoint{3.588418in}{3.049471in}}%
\pgfpathlineto{\pgfqpoint{3.590937in}{3.046170in}}%
\pgfpathlineto{\pgfqpoint{3.593456in}{3.064568in}}%
\pgfpathlineto{\pgfqpoint{3.595974in}{3.074735in}}%
\pgfpathlineto{\pgfqpoint{3.598493in}{3.046120in}}%
\pgfpathlineto{\pgfqpoint{3.601012in}{3.083008in}}%
\pgfpathlineto{\pgfqpoint{3.603531in}{3.088377in}}%
\pgfpathlineto{\pgfqpoint{3.606049in}{3.069620in}}%
\pgfpathlineto{\pgfqpoint{3.608568in}{3.060318in}}%
\pgfpathlineto{\pgfqpoint{3.611087in}{3.068093in}}%
\pgfpathlineto{\pgfqpoint{3.613606in}{3.084529in}}%
\pgfpathlineto{\pgfqpoint{3.616124in}{3.071820in}}%
\pgfpathlineto{\pgfqpoint{3.618643in}{3.068319in}}%
\pgfpathlineto{\pgfqpoint{3.621162in}{3.097892in}}%
\pgfpathlineto{\pgfqpoint{3.623681in}{3.108565in}}%
\pgfpathlineto{\pgfqpoint{3.626199in}{3.116422in}}%
\pgfpathlineto{\pgfqpoint{3.628718in}{3.114241in}}%
\pgfpathlineto{\pgfqpoint{3.631237in}{3.089599in}}%
\pgfpathlineto{\pgfqpoint{3.633756in}{3.089078in}}%
\pgfpathlineto{\pgfqpoint{3.636274in}{3.052508in}}%
\pgfpathlineto{\pgfqpoint{3.638793in}{3.075431in}}%
\pgfpathlineto{\pgfqpoint{3.641312in}{3.086019in}}%
\pgfpathlineto{\pgfqpoint{3.643831in}{3.083757in}}%
\pgfpathlineto{\pgfqpoint{3.646349in}{3.097430in}}%
\pgfpathlineto{\pgfqpoint{3.648868in}{3.074099in}}%
\pgfpathlineto{\pgfqpoint{3.651387in}{3.037608in}}%
\pgfpathlineto{\pgfqpoint{3.656424in}{2.992973in}}%
\pgfpathlineto{\pgfqpoint{3.658943in}{3.002377in}}%
\pgfpathlineto{\pgfqpoint{3.661462in}{2.992820in}}%
\pgfpathlineto{\pgfqpoint{3.663981in}{2.973742in}}%
\pgfpathlineto{\pgfqpoint{3.666499in}{2.927637in}}%
\pgfpathlineto{\pgfqpoint{3.669018in}{2.948237in}}%
\pgfpathlineto{\pgfqpoint{3.671537in}{2.964036in}}%
\pgfpathlineto{\pgfqpoint{3.674056in}{2.937749in}}%
\pgfpathlineto{\pgfqpoint{3.676574in}{2.926155in}}%
\pgfpathlineto{\pgfqpoint{3.679093in}{2.916310in}}%
\pgfpathlineto{\pgfqpoint{3.681612in}{2.956189in}}%
\pgfpathlineto{\pgfqpoint{3.684131in}{2.927929in}}%
\pgfpathlineto{\pgfqpoint{3.686649in}{2.913625in}}%
\pgfpathlineto{\pgfqpoint{3.689168in}{2.888351in}}%
\pgfpathlineto{\pgfqpoint{3.691687in}{2.884426in}}%
\pgfpathlineto{\pgfqpoint{3.694206in}{2.928506in}}%
\pgfpathlineto{\pgfqpoint{3.696724in}{2.922130in}}%
\pgfpathlineto{\pgfqpoint{3.699243in}{2.908439in}}%
\pgfpathlineto{\pgfqpoint{3.701762in}{2.913285in}}%
\pgfpathlineto{\pgfqpoint{3.704281in}{2.863730in}}%
\pgfpathlineto{\pgfqpoint{3.706799in}{2.852816in}}%
\pgfpathlineto{\pgfqpoint{3.709318in}{2.835868in}}%
\pgfpathlineto{\pgfqpoint{3.711837in}{2.827984in}}%
\pgfpathlineto{\pgfqpoint{3.714356in}{2.817817in}}%
\pgfpathlineto{\pgfqpoint{3.716874in}{2.817034in}}%
\pgfpathlineto{\pgfqpoint{3.719393in}{2.859738in}}%
\pgfpathlineto{\pgfqpoint{3.721912in}{2.867514in}}%
\pgfpathlineto{\pgfqpoint{3.724431in}{2.859734in}}%
\pgfpathlineto{\pgfqpoint{3.726949in}{2.837241in}}%
\pgfpathlineto{\pgfqpoint{3.729468in}{2.819250in}}%
\pgfpathlineto{\pgfqpoint{3.731987in}{2.791331in}}%
\pgfpathlineto{\pgfqpoint{3.734506in}{2.797901in}}%
\pgfpathlineto{\pgfqpoint{3.737024in}{2.798956in}}%
\pgfpathlineto{\pgfqpoint{3.739543in}{2.795808in}}%
\pgfpathlineto{\pgfqpoint{3.742062in}{2.768775in}}%
\pgfpathlineto{\pgfqpoint{3.744581in}{2.762131in}}%
\pgfpathlineto{\pgfqpoint{3.747099in}{2.746063in}}%
\pgfpathlineto{\pgfqpoint{3.749618in}{2.741543in}}%
\pgfpathlineto{\pgfqpoint{3.752137in}{2.718323in}}%
\pgfpathlineto{\pgfqpoint{3.754656in}{2.728349in}}%
\pgfpathlineto{\pgfqpoint{3.757174in}{2.774235in}}%
\pgfpathlineto{\pgfqpoint{3.759693in}{2.790804in}}%
\pgfpathlineto{\pgfqpoint{3.762212in}{2.772619in}}%
\pgfpathlineto{\pgfqpoint{3.764731in}{2.786659in}}%
\pgfpathlineto{\pgfqpoint{3.767249in}{2.794601in}}%
\pgfpathlineto{\pgfqpoint{3.769768in}{2.748284in}}%
\pgfpathlineto{\pgfqpoint{3.772287in}{2.729183in}}%
\pgfpathlineto{\pgfqpoint{3.774806in}{2.722069in}}%
\pgfpathlineto{\pgfqpoint{3.777324in}{2.749130in}}%
\pgfpathlineto{\pgfqpoint{3.779843in}{2.746629in}}%
\pgfpathlineto{\pgfqpoint{3.782362in}{2.753270in}}%
\pgfpathlineto{\pgfqpoint{3.784881in}{2.734102in}}%
\pgfpathlineto{\pgfqpoint{3.787399in}{2.752839in}}%
\pgfpathlineto{\pgfqpoint{3.789918in}{2.768106in}}%
\pgfpathlineto{\pgfqpoint{3.792437in}{2.773536in}}%
\pgfpathlineto{\pgfqpoint{3.794956in}{2.737666in}}%
\pgfpathlineto{\pgfqpoint{3.797474in}{2.735071in}}%
\pgfpathlineto{\pgfqpoint{3.799993in}{2.751672in}}%
\pgfpathlineto{\pgfqpoint{3.802512in}{2.740593in}}%
\pgfpathlineto{\pgfqpoint{3.805031in}{2.738301in}}%
\pgfpathlineto{\pgfqpoint{3.807549in}{2.736582in}}%
\pgfpathlineto{\pgfqpoint{3.810068in}{2.744764in}}%
\pgfpathlineto{\pgfqpoint{3.812587in}{2.768394in}}%
\pgfpathlineto{\pgfqpoint{3.817624in}{2.794262in}}%
\pgfpathlineto{\pgfqpoint{3.820143in}{2.793890in}}%
\pgfpathlineto{\pgfqpoint{3.822662in}{2.810716in}}%
\pgfpathlineto{\pgfqpoint{3.825181in}{2.790698in}}%
\pgfpathlineto{\pgfqpoint{3.827699in}{2.792167in}}%
\pgfpathlineto{\pgfqpoint{3.830218in}{2.781981in}}%
\pgfpathlineto{\pgfqpoint{3.832737in}{2.776378in}}%
\pgfpathlineto{\pgfqpoint{3.835256in}{2.803835in}}%
\pgfpathlineto{\pgfqpoint{3.837774in}{2.801966in}}%
\pgfpathlineto{\pgfqpoint{3.840293in}{2.803964in}}%
\pgfpathlineto{\pgfqpoint{3.842812in}{2.819180in}}%
\pgfpathlineto{\pgfqpoint{3.845331in}{2.826312in}}%
\pgfpathlineto{\pgfqpoint{3.847849in}{2.857073in}}%
\pgfpathlineto{\pgfqpoint{3.850368in}{2.837622in}}%
\pgfpathlineto{\pgfqpoint{3.852887in}{2.865719in}}%
\pgfpathlineto{\pgfqpoint{3.855406in}{2.847398in}}%
\pgfpathlineto{\pgfqpoint{3.857924in}{2.861072in}}%
\pgfpathlineto{\pgfqpoint{3.860443in}{2.880644in}}%
\pgfpathlineto{\pgfqpoint{3.862962in}{2.894811in}}%
\pgfpathlineto{\pgfqpoint{3.865481in}{2.895656in}}%
\pgfpathlineto{\pgfqpoint{3.867999in}{2.853381in}}%
\pgfpathlineto{\pgfqpoint{3.870518in}{2.798333in}}%
\pgfpathlineto{\pgfqpoint{3.873037in}{2.761072in}}%
\pgfpathlineto{\pgfqpoint{3.875556in}{2.760587in}}%
\pgfpathlineto{\pgfqpoint{3.878074in}{2.767330in}}%
\pgfpathlineto{\pgfqpoint{3.880593in}{2.766419in}}%
\pgfpathlineto{\pgfqpoint{3.883112in}{2.760404in}}%
\pgfpathlineto{\pgfqpoint{3.885631in}{2.761053in}}%
\pgfpathlineto{\pgfqpoint{3.888149in}{2.748188in}}%
\pgfpathlineto{\pgfqpoint{3.890668in}{2.762389in}}%
\pgfpathlineto{\pgfqpoint{3.893187in}{2.796984in}}%
\pgfpathlineto{\pgfqpoint{3.895706in}{2.778949in}}%
\pgfpathlineto{\pgfqpoint{3.898224in}{2.802381in}}%
\pgfpathlineto{\pgfqpoint{3.900743in}{2.798127in}}%
\pgfpathlineto{\pgfqpoint{3.903262in}{2.792131in}}%
\pgfpathlineto{\pgfqpoint{3.905781in}{2.769494in}}%
\pgfpathlineto{\pgfqpoint{3.908299in}{2.734110in}}%
\pgfpathlineto{\pgfqpoint{3.910818in}{2.737211in}}%
\pgfpathlineto{\pgfqpoint{3.915856in}{2.683134in}}%
\pgfpathlineto{\pgfqpoint{3.918374in}{2.699722in}}%
\pgfpathlineto{\pgfqpoint{3.920893in}{2.706561in}}%
\pgfpathlineto{\pgfqpoint{3.923412in}{2.704847in}}%
\pgfpathlineto{\pgfqpoint{3.925931in}{2.712522in}}%
\pgfpathlineto{\pgfqpoint{3.928449in}{2.711316in}}%
\pgfpathlineto{\pgfqpoint{3.930968in}{2.690968in}}%
\pgfpathlineto{\pgfqpoint{3.933487in}{2.696729in}}%
\pgfpathlineto{\pgfqpoint{3.936006in}{2.686807in}}%
\pgfpathlineto{\pgfqpoint{3.938524in}{2.665775in}}%
\pgfpathlineto{\pgfqpoint{3.941043in}{2.727168in}}%
\pgfpathlineto{\pgfqpoint{3.943562in}{2.727339in}}%
\pgfpathlineto{\pgfqpoint{3.946081in}{2.703463in}}%
\pgfpathlineto{\pgfqpoint{3.948599in}{2.700011in}}%
\pgfpathlineto{\pgfqpoint{3.951118in}{2.695035in}}%
\pgfpathlineto{\pgfqpoint{3.953637in}{2.674300in}}%
\pgfpathlineto{\pgfqpoint{3.956156in}{2.701718in}}%
\pgfpathlineto{\pgfqpoint{3.958674in}{2.701129in}}%
\pgfpathlineto{\pgfqpoint{3.961193in}{2.691165in}}%
\pgfpathlineto{\pgfqpoint{3.963712in}{2.693311in}}%
\pgfpathlineto{\pgfqpoint{3.966231in}{2.720104in}}%
\pgfpathlineto{\pgfqpoint{3.968749in}{2.726374in}}%
\pgfpathlineto{\pgfqpoint{3.971268in}{2.715914in}}%
\pgfpathlineto{\pgfqpoint{3.973787in}{2.703653in}}%
\pgfpathlineto{\pgfqpoint{3.976306in}{2.689922in}}%
\pgfpathlineto{\pgfqpoint{3.978824in}{2.686552in}}%
\pgfpathlineto{\pgfqpoint{3.981343in}{2.693602in}}%
\pgfpathlineto{\pgfqpoint{3.983862in}{2.702383in}}%
\pgfpathlineto{\pgfqpoint{3.986381in}{2.731587in}}%
\pgfpathlineto{\pgfqpoint{3.988899in}{2.703448in}}%
\pgfpathlineto{\pgfqpoint{3.991418in}{2.718884in}}%
\pgfpathlineto{\pgfqpoint{3.993937in}{2.747131in}}%
\pgfpathlineto{\pgfqpoint{3.996456in}{2.757676in}}%
\pgfpathlineto{\pgfqpoint{3.998974in}{2.724813in}}%
\pgfpathlineto{\pgfqpoint{4.001493in}{2.725403in}}%
\pgfpathlineto{\pgfqpoint{4.004012in}{2.731231in}}%
\pgfpathlineto{\pgfqpoint{4.006531in}{2.729601in}}%
\pgfpathlineto{\pgfqpoint{4.009049in}{2.737164in}}%
\pgfpathlineto{\pgfqpoint{4.011568in}{2.760343in}}%
\pgfpathlineto{\pgfqpoint{4.014087in}{2.724076in}}%
\pgfpathlineto{\pgfqpoint{4.016606in}{2.705966in}}%
\pgfpathlineto{\pgfqpoint{4.019124in}{2.722203in}}%
\pgfpathlineto{\pgfqpoint{4.021643in}{2.775224in}}%
\pgfpathlineto{\pgfqpoint{4.024162in}{2.773926in}}%
\pgfpathlineto{\pgfqpoint{4.026681in}{2.773177in}}%
\pgfpathlineto{\pgfqpoint{4.029199in}{2.799189in}}%
\pgfpathlineto{\pgfqpoint{4.031718in}{2.789373in}}%
\pgfpathlineto{\pgfqpoint{4.036756in}{2.770868in}}%
\pgfpathlineto{\pgfqpoint{4.041793in}{2.722482in}}%
\pgfpathlineto{\pgfqpoint{4.044312in}{2.654103in}}%
\pgfpathlineto{\pgfqpoint{4.046831in}{2.684979in}}%
\pgfpathlineto{\pgfqpoint{4.049349in}{2.713130in}}%
\pgfpathlineto{\pgfqpoint{4.051868in}{2.677723in}}%
\pgfpathlineto{\pgfqpoint{4.054387in}{2.671317in}}%
\pgfpathlineto{\pgfqpoint{4.056906in}{2.670191in}}%
\pgfpathlineto{\pgfqpoint{4.059424in}{2.648629in}}%
\pgfpathlineto{\pgfqpoint{4.061943in}{2.659460in}}%
\pgfpathlineto{\pgfqpoint{4.064462in}{2.659014in}}%
\pgfpathlineto{\pgfqpoint{4.066981in}{2.678863in}}%
\pgfpathlineto{\pgfqpoint{4.069499in}{2.678035in}}%
\pgfpathlineto{\pgfqpoint{4.072018in}{2.646992in}}%
\pgfpathlineto{\pgfqpoint{4.074537in}{2.651796in}}%
\pgfpathlineto{\pgfqpoint{4.077056in}{2.686358in}}%
\pgfpathlineto{\pgfqpoint{4.079574in}{2.697470in}}%
\pgfpathlineto{\pgfqpoint{4.082093in}{2.690660in}}%
\pgfpathlineto{\pgfqpoint{4.084612in}{2.666921in}}%
\pgfpathlineto{\pgfqpoint{4.087131in}{2.669009in}}%
\pgfpathlineto{\pgfqpoint{4.089649in}{2.643383in}}%
\pgfpathlineto{\pgfqpoint{4.092168in}{2.627782in}}%
\pgfpathlineto{\pgfqpoint{4.094687in}{2.636732in}}%
\pgfpathlineto{\pgfqpoint{4.097206in}{2.661460in}}%
\pgfpathlineto{\pgfqpoint{4.099724in}{2.669910in}}%
\pgfpathlineto{\pgfqpoint{4.102243in}{2.679012in}}%
\pgfpathlineto{\pgfqpoint{4.104762in}{2.673100in}}%
\pgfpathlineto{\pgfqpoint{4.107281in}{2.671595in}}%
\pgfpathlineto{\pgfqpoint{4.109799in}{2.664502in}}%
\pgfpathlineto{\pgfqpoint{4.112318in}{2.665922in}}%
\pgfpathlineto{\pgfqpoint{4.114837in}{2.705865in}}%
\pgfpathlineto{\pgfqpoint{4.117356in}{2.674461in}}%
\pgfpathlineto{\pgfqpoint{4.119874in}{2.683830in}}%
\pgfpathlineto{\pgfqpoint{4.122393in}{2.707640in}}%
\pgfpathlineto{\pgfqpoint{4.124912in}{2.696253in}}%
\pgfpathlineto{\pgfqpoint{4.127431in}{2.683787in}}%
\pgfpathlineto{\pgfqpoint{4.129949in}{2.687819in}}%
\pgfpathlineto{\pgfqpoint{4.132468in}{2.694467in}}%
\pgfpathlineto{\pgfqpoint{4.134987in}{2.717828in}}%
\pgfpathlineto{\pgfqpoint{4.137506in}{2.758474in}}%
\pgfpathlineto{\pgfqpoint{4.140024in}{2.766584in}}%
\pgfpathlineto{\pgfqpoint{4.142543in}{2.783764in}}%
\pgfpathlineto{\pgfqpoint{4.145062in}{2.781247in}}%
\pgfpathlineto{\pgfqpoint{4.147581in}{2.794765in}}%
\pgfpathlineto{\pgfqpoint{4.150099in}{2.784875in}}%
\pgfpathlineto{\pgfqpoint{4.152618in}{2.792329in}}%
\pgfpathlineto{\pgfqpoint{4.155137in}{2.819798in}}%
\pgfpathlineto{\pgfqpoint{4.157656in}{2.833434in}}%
\pgfpathlineto{\pgfqpoint{4.160174in}{2.806176in}}%
\pgfpathlineto{\pgfqpoint{4.162693in}{2.786430in}}%
\pgfpathlineto{\pgfqpoint{4.165212in}{2.765022in}}%
\pgfpathlineto{\pgfqpoint{4.167731in}{2.750420in}}%
\pgfpathlineto{\pgfqpoint{4.170249in}{2.739127in}}%
\pgfpathlineto{\pgfqpoint{4.172768in}{2.754139in}}%
\pgfpathlineto{\pgfqpoint{4.175287in}{2.741418in}}%
\pgfpathlineto{\pgfqpoint{4.177806in}{2.752237in}}%
\pgfpathlineto{\pgfqpoint{4.180324in}{2.742267in}}%
\pgfpathlineto{\pgfqpoint{4.182843in}{2.677950in}}%
\pgfpathlineto{\pgfqpoint{4.185362in}{2.632822in}}%
\pgfpathlineto{\pgfqpoint{4.187881in}{2.622127in}}%
\pgfpathlineto{\pgfqpoint{4.190399in}{2.621137in}}%
\pgfpathlineto{\pgfqpoint{4.192918in}{2.628112in}}%
\pgfpathlineto{\pgfqpoint{4.195437in}{2.655001in}}%
\pgfpathlineto{\pgfqpoint{4.197956in}{2.667488in}}%
\pgfpathlineto{\pgfqpoint{4.200474in}{2.631350in}}%
\pgfpathlineto{\pgfqpoint{4.202993in}{2.634167in}}%
\pgfpathlineto{\pgfqpoint{4.205512in}{2.644310in}}%
\pgfpathlineto{\pgfqpoint{4.208031in}{2.684077in}}%
\pgfpathlineto{\pgfqpoint{4.210549in}{2.675102in}}%
\pgfpathlineto{\pgfqpoint{4.213068in}{2.679724in}}%
\pgfpathlineto{\pgfqpoint{4.215587in}{2.703363in}}%
\pgfpathlineto{\pgfqpoint{4.218106in}{2.722963in}}%
\pgfpathlineto{\pgfqpoint{4.220624in}{2.709478in}}%
\pgfpathlineto{\pgfqpoint{4.223143in}{2.728207in}}%
\pgfpathlineto{\pgfqpoint{4.225662in}{2.710172in}}%
\pgfpathlineto{\pgfqpoint{4.228181in}{2.744078in}}%
\pgfpathlineto{\pgfqpoint{4.230699in}{2.759935in}}%
\pgfpathlineto{\pgfqpoint{4.233218in}{2.757528in}}%
\pgfpathlineto{\pgfqpoint{4.235737in}{2.789503in}}%
\pgfpathlineto{\pgfqpoint{4.238256in}{2.819002in}}%
\pgfpathlineto{\pgfqpoint{4.240774in}{2.837635in}}%
\pgfpathlineto{\pgfqpoint{4.243293in}{2.832605in}}%
\pgfpathlineto{\pgfqpoint{4.245812in}{2.815827in}}%
\pgfpathlineto{\pgfqpoint{4.248331in}{2.755317in}}%
\pgfpathlineto{\pgfqpoint{4.250849in}{2.725501in}}%
\pgfpathlineto{\pgfqpoint{4.253368in}{2.721135in}}%
\pgfpathlineto{\pgfqpoint{4.255887in}{2.702270in}}%
\pgfpathlineto{\pgfqpoint{4.258406in}{2.725088in}}%
\pgfpathlineto{\pgfqpoint{4.260924in}{2.721118in}}%
\pgfpathlineto{\pgfqpoint{4.263443in}{2.727139in}}%
\pgfpathlineto{\pgfqpoint{4.265962in}{2.706689in}}%
\pgfpathlineto{\pgfqpoint{4.268481in}{2.706017in}}%
\pgfpathlineto{\pgfqpoint{4.270999in}{2.680121in}}%
\pgfpathlineto{\pgfqpoint{4.273518in}{2.685020in}}%
\pgfpathlineto{\pgfqpoint{4.276037in}{2.688618in}}%
\pgfpathlineto{\pgfqpoint{4.278556in}{2.663251in}}%
\pgfpathlineto{\pgfqpoint{4.281074in}{2.666788in}}%
\pgfpathlineto{\pgfqpoint{4.283593in}{2.668178in}}%
\pgfpathlineto{\pgfqpoint{4.286112in}{2.663148in}}%
\pgfpathlineto{\pgfqpoint{4.288631in}{2.639791in}}%
\pgfpathlineto{\pgfqpoint{4.291149in}{2.612043in}}%
\pgfpathlineto{\pgfqpoint{4.293668in}{2.621192in}}%
\pgfpathlineto{\pgfqpoint{4.296187in}{2.648736in}}%
\pgfpathlineto{\pgfqpoint{4.298706in}{2.645118in}}%
\pgfpathlineto{\pgfqpoint{4.301224in}{2.625723in}}%
\pgfpathlineto{\pgfqpoint{4.303743in}{2.641184in}}%
\pgfpathlineto{\pgfqpoint{4.306262in}{2.672155in}}%
\pgfpathlineto{\pgfqpoint{4.308781in}{2.641811in}}%
\pgfpathlineto{\pgfqpoint{4.311299in}{2.659647in}}%
\pgfpathlineto{\pgfqpoint{4.313818in}{2.657329in}}%
\pgfpathlineto{\pgfqpoint{4.316337in}{2.653999in}}%
\pgfpathlineto{\pgfqpoint{4.318856in}{2.649857in}}%
\pgfpathlineto{\pgfqpoint{4.321374in}{2.698310in}}%
\pgfpathlineto{\pgfqpoint{4.323893in}{2.685879in}}%
\pgfpathlineto{\pgfqpoint{4.326412in}{2.675627in}}%
\pgfpathlineto{\pgfqpoint{4.328931in}{2.641108in}}%
\pgfpathlineto{\pgfqpoint{4.331449in}{2.661618in}}%
\pgfpathlineto{\pgfqpoint{4.333968in}{2.623562in}}%
\pgfpathlineto{\pgfqpoint{4.336487in}{2.607588in}}%
\pgfpathlineto{\pgfqpoint{4.339006in}{2.576660in}}%
\pgfpathlineto{\pgfqpoint{4.341524in}{2.590032in}}%
\pgfpathlineto{\pgfqpoint{4.344043in}{2.611243in}}%
\pgfpathlineto{\pgfqpoint{4.346562in}{2.613041in}}%
\pgfpathlineto{\pgfqpoint{4.349081in}{2.615862in}}%
\pgfpathlineto{\pgfqpoint{4.351599in}{2.621897in}}%
\pgfpathlineto{\pgfqpoint{4.354118in}{2.606914in}}%
\pgfpathlineto{\pgfqpoint{4.356637in}{2.590787in}}%
\pgfpathlineto{\pgfqpoint{4.359156in}{2.549518in}}%
\pgfpathlineto{\pgfqpoint{4.361674in}{2.541203in}}%
\pgfpathlineto{\pgfqpoint{4.364193in}{2.550985in}}%
\pgfpathlineto{\pgfqpoint{4.366712in}{2.565737in}}%
\pgfpathlineto{\pgfqpoint{4.369231in}{2.592330in}}%
\pgfpathlineto{\pgfqpoint{4.371749in}{2.602096in}}%
\pgfpathlineto{\pgfqpoint{4.374268in}{2.602556in}}%
\pgfpathlineto{\pgfqpoint{4.376787in}{2.559261in}}%
\pgfpathlineto{\pgfqpoint{4.379306in}{2.552064in}}%
\pgfpathlineto{\pgfqpoint{4.381824in}{2.583081in}}%
\pgfpathlineto{\pgfqpoint{4.384343in}{2.626876in}}%
\pgfpathlineto{\pgfqpoint{4.386862in}{2.648841in}}%
\pgfpathlineto{\pgfqpoint{4.389381in}{2.677471in}}%
\pgfpathlineto{\pgfqpoint{4.391899in}{2.654654in}}%
\pgfpathlineto{\pgfqpoint{4.394418in}{2.648305in}}%
\pgfpathlineto{\pgfqpoint{4.396937in}{2.665544in}}%
\pgfpathlineto{\pgfqpoint{4.399456in}{2.629392in}}%
\pgfpathlineto{\pgfqpoint{4.401974in}{2.616189in}}%
\pgfpathlineto{\pgfqpoint{4.404493in}{2.592947in}}%
\pgfpathlineto{\pgfqpoint{4.407012in}{2.620094in}}%
\pgfpathlineto{\pgfqpoint{4.409531in}{2.605707in}}%
\pgfpathlineto{\pgfqpoint{4.412049in}{2.631622in}}%
\pgfpathlineto{\pgfqpoint{4.414568in}{2.670927in}}%
\pgfpathlineto{\pgfqpoint{4.417087in}{2.686002in}}%
\pgfpathlineto{\pgfqpoint{4.419606in}{2.697319in}}%
\pgfpathlineto{\pgfqpoint{4.422124in}{2.700400in}}%
\pgfpathlineto{\pgfqpoint{4.424643in}{2.738733in}}%
\pgfpathlineto{\pgfqpoint{4.429681in}{2.761291in}}%
\pgfpathlineto{\pgfqpoint{4.432199in}{2.759946in}}%
\pgfpathlineto{\pgfqpoint{4.434718in}{2.717443in}}%
\pgfpathlineto{\pgfqpoint{4.437237in}{2.733539in}}%
\pgfpathlineto{\pgfqpoint{4.439756in}{2.734193in}}%
\pgfpathlineto{\pgfqpoint{4.442274in}{2.731748in}}%
\pgfpathlineto{\pgfqpoint{4.444793in}{2.755656in}}%
\pgfpathlineto{\pgfqpoint{4.447312in}{2.730524in}}%
\pgfpathlineto{\pgfqpoint{4.449831in}{2.715602in}}%
\pgfpathlineto{\pgfqpoint{4.452349in}{2.717217in}}%
\pgfpathlineto{\pgfqpoint{4.454868in}{2.728525in}}%
\pgfpathlineto{\pgfqpoint{4.457387in}{2.728267in}}%
\pgfpathlineto{\pgfqpoint{4.459906in}{2.713044in}}%
\pgfpathlineto{\pgfqpoint{4.462424in}{2.740906in}}%
\pgfpathlineto{\pgfqpoint{4.464943in}{2.743479in}}%
\pgfpathlineto{\pgfqpoint{4.467462in}{2.764592in}}%
\pgfpathlineto{\pgfqpoint{4.469981in}{2.788235in}}%
\pgfpathlineto{\pgfqpoint{4.472499in}{2.775618in}}%
\pgfpathlineto{\pgfqpoint{4.477537in}{2.770244in}}%
\pgfpathlineto{\pgfqpoint{4.480056in}{2.794094in}}%
\pgfpathlineto{\pgfqpoint{4.482574in}{2.759546in}}%
\pgfpathlineto{\pgfqpoint{4.485093in}{2.762042in}}%
\pgfpathlineto{\pgfqpoint{4.487612in}{2.729983in}}%
\pgfpathlineto{\pgfqpoint{4.490131in}{2.759155in}}%
\pgfpathlineto{\pgfqpoint{4.492649in}{2.768731in}}%
\pgfpathlineto{\pgfqpoint{4.495168in}{2.784597in}}%
\pgfpathlineto{\pgfqpoint{4.497687in}{2.766048in}}%
\pgfpathlineto{\pgfqpoint{4.500206in}{2.778797in}}%
\pgfpathlineto{\pgfqpoint{4.502724in}{2.799003in}}%
\pgfpathlineto{\pgfqpoint{4.505243in}{2.827911in}}%
\pgfpathlineto{\pgfqpoint{4.507762in}{2.825203in}}%
\pgfpathlineto{\pgfqpoint{4.510281in}{2.810365in}}%
\pgfpathlineto{\pgfqpoint{4.512799in}{2.776435in}}%
\pgfpathlineto{\pgfqpoint{4.515318in}{2.815383in}}%
\pgfpathlineto{\pgfqpoint{4.517837in}{2.842622in}}%
\pgfpathlineto{\pgfqpoint{4.520356in}{2.880463in}}%
\pgfpathlineto{\pgfqpoint{4.522874in}{2.892547in}}%
\pgfpathlineto{\pgfqpoint{4.525393in}{2.894216in}}%
\pgfpathlineto{\pgfqpoint{4.527912in}{2.940014in}}%
\pgfpathlineto{\pgfqpoint{4.530431in}{2.936190in}}%
\pgfpathlineto{\pgfqpoint{4.532949in}{2.920214in}}%
\pgfpathlineto{\pgfqpoint{4.535468in}{2.930871in}}%
\pgfpathlineto{\pgfqpoint{4.540506in}{2.911938in}}%
\pgfpathlineto{\pgfqpoint{4.543024in}{2.946280in}}%
\pgfpathlineto{\pgfqpoint{4.545543in}{2.954193in}}%
\pgfpathlineto{\pgfqpoint{4.548062in}{2.969071in}}%
\pgfpathlineto{\pgfqpoint{4.550581in}{2.953361in}}%
\pgfpathlineto{\pgfqpoint{4.553099in}{2.985809in}}%
\pgfpathlineto{\pgfqpoint{4.555618in}{2.970627in}}%
\pgfpathlineto{\pgfqpoint{4.558137in}{2.990629in}}%
\pgfpathlineto{\pgfqpoint{4.560656in}{3.009693in}}%
\pgfpathlineto{\pgfqpoint{4.563174in}{2.993389in}}%
\pgfpathlineto{\pgfqpoint{4.565693in}{2.936429in}}%
\pgfpathlineto{\pgfqpoint{4.568212in}{2.958138in}}%
\pgfpathlineto{\pgfqpoint{4.570731in}{2.950541in}}%
\pgfpathlineto{\pgfqpoint{4.573249in}{2.948215in}}%
\pgfpathlineto{\pgfqpoint{4.575768in}{2.967510in}}%
\pgfpathlineto{\pgfqpoint{4.578287in}{2.998350in}}%
\pgfpathlineto{\pgfqpoint{4.580806in}{3.022385in}}%
\pgfpathlineto{\pgfqpoint{4.583324in}{3.042387in}}%
\pgfpathlineto{\pgfqpoint{4.585843in}{3.078180in}}%
\pgfpathlineto{\pgfqpoint{4.588362in}{3.065343in}}%
\pgfpathlineto{\pgfqpoint{4.590881in}{3.082836in}}%
\pgfpathlineto{\pgfqpoint{4.593399in}{3.055566in}}%
\pgfpathlineto{\pgfqpoint{4.595918in}{3.081572in}}%
\pgfpathlineto{\pgfqpoint{4.598437in}{3.098773in}}%
\pgfpathlineto{\pgfqpoint{4.600956in}{3.100414in}}%
\pgfpathlineto{\pgfqpoint{4.603474in}{3.128318in}}%
\pgfpathlineto{\pgfqpoint{4.605993in}{3.119870in}}%
\pgfpathlineto{\pgfqpoint{4.608512in}{3.181130in}}%
\pgfpathlineto{\pgfqpoint{4.611031in}{3.227548in}}%
\pgfpathlineto{\pgfqpoint{4.613549in}{3.244565in}}%
\pgfpathlineto{\pgfqpoint{4.616068in}{3.223892in}}%
\pgfpathlineto{\pgfqpoint{4.618587in}{3.247622in}}%
\pgfpathlineto{\pgfqpoint{4.621106in}{3.257344in}}%
\pgfpathlineto{\pgfqpoint{4.623624in}{3.260270in}}%
\pgfpathlineto{\pgfqpoint{4.626143in}{3.255358in}}%
\pgfpathlineto{\pgfqpoint{4.631181in}{3.227625in}}%
\pgfpathlineto{\pgfqpoint{4.633699in}{3.240215in}}%
\pgfpathlineto{\pgfqpoint{4.636218in}{3.203549in}}%
\pgfpathlineto{\pgfqpoint{4.638737in}{3.195826in}}%
\pgfpathlineto{\pgfqpoint{4.641256in}{3.221477in}}%
\pgfpathlineto{\pgfqpoint{4.643774in}{3.249941in}}%
\pgfpathlineto{\pgfqpoint{4.646293in}{3.254225in}}%
\pgfpathlineto{\pgfqpoint{4.648812in}{3.280982in}}%
\pgfpathlineto{\pgfqpoint{4.651331in}{3.263763in}}%
\pgfpathlineto{\pgfqpoint{4.653849in}{3.290845in}}%
\pgfpathlineto{\pgfqpoint{4.656368in}{3.299768in}}%
\pgfpathlineto{\pgfqpoint{4.658887in}{3.272984in}}%
\pgfpathlineto{\pgfqpoint{4.661406in}{3.312780in}}%
\pgfpathlineto{\pgfqpoint{4.663924in}{3.316816in}}%
\pgfpathlineto{\pgfqpoint{4.666443in}{3.294254in}}%
\pgfpathlineto{\pgfqpoint{4.668962in}{3.305793in}}%
\pgfpathlineto{\pgfqpoint{4.671481in}{3.308886in}}%
\pgfpathlineto{\pgfqpoint{4.673999in}{3.308851in}}%
\pgfpathlineto{\pgfqpoint{4.676518in}{3.355260in}}%
\pgfpathlineto{\pgfqpoint{4.679037in}{3.368754in}}%
\pgfpathlineto{\pgfqpoint{4.681556in}{3.303784in}}%
\pgfpathlineto{\pgfqpoint{4.684074in}{3.335153in}}%
\pgfpathlineto{\pgfqpoint{4.686593in}{3.299590in}}%
\pgfpathlineto{\pgfqpoint{4.689112in}{3.269315in}}%
\pgfpathlineto{\pgfqpoint{4.691631in}{3.255189in}}%
\pgfpathlineto{\pgfqpoint{4.694149in}{3.268792in}}%
\pgfpathlineto{\pgfqpoint{4.696668in}{3.270676in}}%
\pgfpathlineto{\pgfqpoint{4.699187in}{3.273723in}}%
\pgfpathlineto{\pgfqpoint{4.701706in}{3.258208in}}%
\pgfpathlineto{\pgfqpoint{4.704224in}{3.268471in}}%
\pgfpathlineto{\pgfqpoint{4.706743in}{3.222073in}}%
\pgfpathlineto{\pgfqpoint{4.709262in}{3.192635in}}%
\pgfpathlineto{\pgfqpoint{4.711781in}{3.187833in}}%
\pgfpathlineto{\pgfqpoint{4.714299in}{3.204055in}}%
\pgfpathlineto{\pgfqpoint{4.716818in}{3.201837in}}%
\pgfpathlineto{\pgfqpoint{4.719337in}{3.203067in}}%
\pgfpathlineto{\pgfqpoint{4.721856in}{3.196415in}}%
\pgfpathlineto{\pgfqpoint{4.724374in}{3.173138in}}%
\pgfpathlineto{\pgfqpoint{4.726893in}{3.166175in}}%
\pgfpathlineto{\pgfqpoint{4.729412in}{3.146912in}}%
\pgfpathlineto{\pgfqpoint{4.731931in}{3.132887in}}%
\pgfpathlineto{\pgfqpoint{4.734449in}{3.173541in}}%
\pgfpathlineto{\pgfqpoint{4.736968in}{3.156163in}}%
\pgfpathlineto{\pgfqpoint{4.739487in}{3.171168in}}%
\pgfpathlineto{\pgfqpoint{4.742006in}{3.167590in}}%
\pgfpathlineto{\pgfqpoint{4.744524in}{3.168243in}}%
\pgfpathlineto{\pgfqpoint{4.747043in}{3.170218in}}%
\pgfpathlineto{\pgfqpoint{4.749562in}{3.206265in}}%
\pgfpathlineto{\pgfqpoint{4.752081in}{3.196000in}}%
\pgfpathlineto{\pgfqpoint{4.754599in}{3.217914in}}%
\pgfpathlineto{\pgfqpoint{4.757118in}{3.222528in}}%
\pgfpathlineto{\pgfqpoint{4.759637in}{3.225877in}}%
\pgfpathlineto{\pgfqpoint{4.762156in}{3.186112in}}%
\pgfpathlineto{\pgfqpoint{4.764674in}{3.201696in}}%
\pgfpathlineto{\pgfqpoint{4.767193in}{3.203968in}}%
\pgfpathlineto{\pgfqpoint{4.769712in}{3.207629in}}%
\pgfpathlineto{\pgfqpoint{4.772231in}{3.198938in}}%
\pgfpathlineto{\pgfqpoint{4.774749in}{3.224295in}}%
\pgfpathlineto{\pgfqpoint{4.777268in}{3.225650in}}%
\pgfpathlineto{\pgfqpoint{4.779787in}{3.214570in}}%
\pgfpathlineto{\pgfqpoint{4.782306in}{3.217992in}}%
\pgfpathlineto{\pgfqpoint{4.784824in}{3.258829in}}%
\pgfpathlineto{\pgfqpoint{4.787343in}{3.262983in}}%
\pgfpathlineto{\pgfqpoint{4.789862in}{3.303811in}}%
\pgfpathlineto{\pgfqpoint{4.792381in}{3.307825in}}%
\pgfpathlineto{\pgfqpoint{4.794899in}{3.331996in}}%
\pgfpathlineto{\pgfqpoint{4.797418in}{3.354094in}}%
\pgfpathlineto{\pgfqpoint{4.799937in}{3.396002in}}%
\pgfpathlineto{\pgfqpoint{4.802456in}{3.395417in}}%
\pgfpathlineto{\pgfqpoint{4.804974in}{3.432139in}}%
\pgfpathlineto{\pgfqpoint{4.807493in}{3.450373in}}%
\pgfpathlineto{\pgfqpoint{4.810012in}{3.406939in}}%
\pgfpathlineto{\pgfqpoint{4.812531in}{3.414033in}}%
\pgfpathlineto{\pgfqpoint{4.815049in}{3.384484in}}%
\pgfpathlineto{\pgfqpoint{4.817568in}{3.391297in}}%
\pgfpathlineto{\pgfqpoint{4.820087in}{3.364648in}}%
\pgfpathlineto{\pgfqpoint{4.822606in}{3.354894in}}%
\pgfpathlineto{\pgfqpoint{4.825124in}{3.365480in}}%
\pgfpathlineto{\pgfqpoint{4.827643in}{3.369277in}}%
\pgfpathlineto{\pgfqpoint{4.830162in}{3.361388in}}%
\pgfpathlineto{\pgfqpoint{4.832681in}{3.368855in}}%
\pgfpathlineto{\pgfqpoint{4.835199in}{3.324335in}}%
\pgfpathlineto{\pgfqpoint{4.837718in}{3.328707in}}%
\pgfpathlineto{\pgfqpoint{4.840237in}{3.312561in}}%
\pgfpathlineto{\pgfqpoint{4.842756in}{3.306713in}}%
\pgfpathlineto{\pgfqpoint{4.845274in}{3.296683in}}%
\pgfpathlineto{\pgfqpoint{4.847793in}{3.340544in}}%
\pgfpathlineto{\pgfqpoint{4.850312in}{3.331730in}}%
\pgfpathlineto{\pgfqpoint{4.852831in}{3.345877in}}%
\pgfpathlineto{\pgfqpoint{4.855349in}{3.377361in}}%
\pgfpathlineto{\pgfqpoint{4.857868in}{3.387152in}}%
\pgfpathlineto{\pgfqpoint{4.860387in}{3.388749in}}%
\pgfpathlineto{\pgfqpoint{4.862906in}{3.439750in}}%
\pgfpathlineto{\pgfqpoint{4.865424in}{3.445102in}}%
\pgfpathlineto{\pgfqpoint{4.867943in}{3.460250in}}%
\pgfpathlineto{\pgfqpoint{4.870462in}{3.466421in}}%
\pgfpathlineto{\pgfqpoint{4.872981in}{3.411580in}}%
\pgfpathlineto{\pgfqpoint{4.875499in}{3.372153in}}%
\pgfpathlineto{\pgfqpoint{4.878018in}{3.352237in}}%
\pgfpathlineto{\pgfqpoint{4.880537in}{3.383935in}}%
\pgfpathlineto{\pgfqpoint{4.883056in}{3.352214in}}%
\pgfpathlineto{\pgfqpoint{4.885574in}{3.336552in}}%
\pgfpathlineto{\pgfqpoint{4.888093in}{3.308540in}}%
\pgfpathlineto{\pgfqpoint{4.890612in}{3.285750in}}%
\pgfpathlineto{\pgfqpoint{4.893131in}{3.286621in}}%
\pgfpathlineto{\pgfqpoint{4.895649in}{3.294828in}}%
\pgfpathlineto{\pgfqpoint{4.898168in}{3.313596in}}%
\pgfpathlineto{\pgfqpoint{4.900687in}{3.343191in}}%
\pgfpathlineto{\pgfqpoint{4.903206in}{3.355015in}}%
\pgfpathlineto{\pgfqpoint{4.905724in}{3.403001in}}%
\pgfpathlineto{\pgfqpoint{4.908243in}{3.402417in}}%
\pgfpathlineto{\pgfqpoint{4.910762in}{3.411532in}}%
\pgfpathlineto{\pgfqpoint{4.913281in}{3.442232in}}%
\pgfpathlineto{\pgfqpoint{4.915799in}{3.444433in}}%
\pgfpathlineto{\pgfqpoint{4.918318in}{3.418570in}}%
\pgfpathlineto{\pgfqpoint{4.920837in}{3.406401in}}%
\pgfpathlineto{\pgfqpoint{4.923356in}{3.407413in}}%
\pgfpathlineto{\pgfqpoint{4.925874in}{3.395982in}}%
\pgfpathlineto{\pgfqpoint{4.928393in}{3.412451in}}%
\pgfpathlineto{\pgfqpoint{4.930912in}{3.440473in}}%
\pgfpathlineto{\pgfqpoint{4.933431in}{3.479909in}}%
\pgfpathlineto{\pgfqpoint{4.935949in}{3.457462in}}%
\pgfpathlineto{\pgfqpoint{4.938468in}{3.446552in}}%
\pgfpathlineto{\pgfqpoint{4.940987in}{3.474880in}}%
\pgfpathlineto{\pgfqpoint{4.943506in}{3.532569in}}%
\pgfpathlineto{\pgfqpoint{4.946024in}{3.542364in}}%
\pgfpathlineto{\pgfqpoint{4.948543in}{3.534736in}}%
\pgfpathlineto{\pgfqpoint{4.951062in}{3.566592in}}%
\pgfpathlineto{\pgfqpoint{4.953581in}{3.542687in}}%
\pgfpathlineto{\pgfqpoint{4.956099in}{3.559227in}}%
\pgfpathlineto{\pgfqpoint{4.958618in}{3.560456in}}%
\pgfpathlineto{\pgfqpoint{4.961137in}{3.526856in}}%
\pgfpathlineto{\pgfqpoint{4.963656in}{3.528442in}}%
\pgfpathlineto{\pgfqpoint{4.966174in}{3.557368in}}%
\pgfpathlineto{\pgfqpoint{4.968693in}{3.559285in}}%
\pgfpathlineto{\pgfqpoint{4.971212in}{3.540648in}}%
\pgfpathlineto{\pgfqpoint{4.973731in}{3.567932in}}%
\pgfpathlineto{\pgfqpoint{4.976249in}{3.519438in}}%
\pgfpathlineto{\pgfqpoint{4.978768in}{3.553821in}}%
\pgfpathlineto{\pgfqpoint{4.981287in}{3.507409in}}%
\pgfpathlineto{\pgfqpoint{4.983806in}{3.505193in}}%
\pgfpathlineto{\pgfqpoint{4.986324in}{3.479610in}}%
\pgfpathlineto{\pgfqpoint{4.988843in}{3.520798in}}%
\pgfpathlineto{\pgfqpoint{4.991362in}{3.522048in}}%
\pgfpathlineto{\pgfqpoint{4.993881in}{3.483118in}}%
\pgfpathlineto{\pgfqpoint{4.996399in}{3.435746in}}%
\pgfpathlineto{\pgfqpoint{4.998918in}{3.475327in}}%
\pgfpathlineto{\pgfqpoint{5.001437in}{3.479938in}}%
\pgfpathlineto{\pgfqpoint{5.003956in}{3.485780in}}%
\pgfpathlineto{\pgfqpoint{5.006474in}{3.498107in}}%
\pgfpathlineto{\pgfqpoint{5.008993in}{3.526586in}}%
\pgfpathlineto{\pgfqpoint{5.011512in}{3.552403in}}%
\pgfpathlineto{\pgfqpoint{5.014031in}{3.594492in}}%
\pgfpathlineto{\pgfqpoint{5.016549in}{3.587459in}}%
\pgfpathlineto{\pgfqpoint{5.019068in}{3.613306in}}%
\pgfpathlineto{\pgfqpoint{5.021587in}{3.647398in}}%
\pgfpathlineto{\pgfqpoint{5.024106in}{3.659583in}}%
\pgfpathlineto{\pgfqpoint{5.026624in}{3.694722in}}%
\pgfpathlineto{\pgfqpoint{5.029143in}{3.688093in}}%
\pgfpathlineto{\pgfqpoint{5.031662in}{3.634930in}}%
\pgfpathlineto{\pgfqpoint{5.034181in}{3.643051in}}%
\pgfpathlineto{\pgfqpoint{5.036699in}{3.610870in}}%
\pgfpathlineto{\pgfqpoint{5.039218in}{3.624616in}}%
\pgfpathlineto{\pgfqpoint{5.041737in}{3.620534in}}%
\pgfpathlineto{\pgfqpoint{5.044256in}{3.637761in}}%
\pgfpathlineto{\pgfqpoint{5.046774in}{3.616046in}}%
\pgfpathlineto{\pgfqpoint{5.049293in}{3.630947in}}%
\pgfpathlineto{\pgfqpoint{5.051812in}{3.613923in}}%
\pgfpathlineto{\pgfqpoint{5.054331in}{3.627601in}}%
\pgfpathlineto{\pgfqpoint{5.056849in}{3.676214in}}%
\pgfpathlineto{\pgfqpoint{5.059368in}{3.672111in}}%
\pgfpathlineto{\pgfqpoint{5.061887in}{3.726049in}}%
\pgfpathlineto{\pgfqpoint{5.064406in}{3.725411in}}%
\pgfpathlineto{\pgfqpoint{5.066924in}{3.756671in}}%
\pgfpathlineto{\pgfqpoint{5.069443in}{3.741029in}}%
\pgfpathlineto{\pgfqpoint{5.071962in}{3.705540in}}%
\pgfpathlineto{\pgfqpoint{5.074481in}{3.707992in}}%
\pgfpathlineto{\pgfqpoint{5.076999in}{3.695715in}}%
\pgfpathlineto{\pgfqpoint{5.079518in}{3.717859in}}%
\pgfpathlineto{\pgfqpoint{5.082037in}{3.677598in}}%
\pgfpathlineto{\pgfqpoint{5.084556in}{3.667052in}}%
\pgfpathlineto{\pgfqpoint{5.087074in}{3.661574in}}%
\pgfpathlineto{\pgfqpoint{5.089593in}{3.661039in}}%
\pgfpathlineto{\pgfqpoint{5.092112in}{3.614885in}}%
\pgfpathlineto{\pgfqpoint{5.094631in}{3.658483in}}%
\pgfpathlineto{\pgfqpoint{5.097149in}{3.620225in}}%
\pgfpathlineto{\pgfqpoint{5.099668in}{3.672857in}}%
\pgfpathlineto{\pgfqpoint{5.102187in}{3.658769in}}%
\pgfpathlineto{\pgfqpoint{5.104706in}{3.671178in}}%
\pgfpathlineto{\pgfqpoint{5.107224in}{3.664117in}}%
\pgfpathlineto{\pgfqpoint{5.109743in}{3.684326in}}%
\pgfpathlineto{\pgfqpoint{5.112262in}{3.723717in}}%
\pgfpathlineto{\pgfqpoint{5.114781in}{3.723401in}}%
\pgfpathlineto{\pgfqpoint{5.117299in}{3.710787in}}%
\pgfpathlineto{\pgfqpoint{5.119818in}{3.708535in}}%
\pgfpathlineto{\pgfqpoint{5.122337in}{3.704796in}}%
\pgfpathlineto{\pgfqpoint{5.124856in}{3.706333in}}%
\pgfpathlineto{\pgfqpoint{5.127374in}{3.702819in}}%
\pgfpathlineto{\pgfqpoint{5.129893in}{3.701129in}}%
\pgfpathlineto{\pgfqpoint{5.132412in}{3.677912in}}%
\pgfpathlineto{\pgfqpoint{5.134931in}{3.671066in}}%
\pgfpathlineto{\pgfqpoint{5.137449in}{3.673509in}}%
\pgfpathlineto{\pgfqpoint{5.139968in}{3.654661in}}%
\pgfpathlineto{\pgfqpoint{5.142487in}{3.652508in}}%
\pgfpathlineto{\pgfqpoint{5.145006in}{3.621504in}}%
\pgfpathlineto{\pgfqpoint{5.147524in}{3.628293in}}%
\pgfpathlineto{\pgfqpoint{5.150043in}{3.645582in}}%
\pgfpathlineto{\pgfqpoint{5.152562in}{3.665529in}}%
\pgfpathlineto{\pgfqpoint{5.155081in}{3.637758in}}%
\pgfpathlineto{\pgfqpoint{5.157599in}{3.601283in}}%
\pgfpathlineto{\pgfqpoint{5.160118in}{3.596388in}}%
\pgfpathlineto{\pgfqpoint{5.162637in}{3.598593in}}%
\pgfpathlineto{\pgfqpoint{5.165156in}{3.608535in}}%
\pgfpathlineto{\pgfqpoint{5.167674in}{3.583413in}}%
\pgfpathlineto{\pgfqpoint{5.170193in}{3.556380in}}%
\pgfpathlineto{\pgfqpoint{5.172712in}{3.523337in}}%
\pgfpathlineto{\pgfqpoint{5.175231in}{3.509121in}}%
\pgfpathlineto{\pgfqpoint{5.177749in}{3.474178in}}%
\pgfpathlineto{\pgfqpoint{5.180268in}{3.485251in}}%
\pgfpathlineto{\pgfqpoint{5.182787in}{3.476575in}}%
\pgfpathlineto{\pgfqpoint{5.185306in}{3.446174in}}%
\pgfpathlineto{\pgfqpoint{5.187824in}{3.435991in}}%
\pgfpathlineto{\pgfqpoint{5.192862in}{3.463528in}}%
\pgfpathlineto{\pgfqpoint{5.195381in}{3.459732in}}%
\pgfpathlineto{\pgfqpoint{5.197899in}{3.484089in}}%
\pgfpathlineto{\pgfqpoint{5.200418in}{3.487782in}}%
\pgfpathlineto{\pgfqpoint{5.202937in}{3.504041in}}%
\pgfpathlineto{\pgfqpoint{5.205456in}{3.517937in}}%
\pgfpathlineto{\pgfqpoint{5.207974in}{3.546054in}}%
\pgfpathlineto{\pgfqpoint{5.210493in}{3.556515in}}%
\pgfpathlineto{\pgfqpoint{5.213012in}{3.569883in}}%
\pgfpathlineto{\pgfqpoint{5.215531in}{3.543863in}}%
\pgfpathlineto{\pgfqpoint{5.218049in}{3.557600in}}%
\pgfpathlineto{\pgfqpoint{5.220568in}{3.545570in}}%
\pgfpathlineto{\pgfqpoint{5.223087in}{3.500643in}}%
\pgfpathlineto{\pgfqpoint{5.225606in}{3.504945in}}%
\pgfpathlineto{\pgfqpoint{5.228124in}{3.501226in}}%
\pgfpathlineto{\pgfqpoint{5.230643in}{3.504585in}}%
\pgfpathlineto{\pgfqpoint{5.233162in}{3.526837in}}%
\pgfpathlineto{\pgfqpoint{5.235681in}{3.505625in}}%
\pgfpathlineto{\pgfqpoint{5.238199in}{3.470942in}}%
\pgfpathlineto{\pgfqpoint{5.240718in}{3.471193in}}%
\pgfpathlineto{\pgfqpoint{5.243237in}{3.484966in}}%
\pgfpathlineto{\pgfqpoint{5.245756in}{3.492491in}}%
\pgfpathlineto{\pgfqpoint{5.248274in}{3.472994in}}%
\pgfpathlineto{\pgfqpoint{5.250793in}{3.459766in}}%
\pgfpathlineto{\pgfqpoint{5.253312in}{3.455275in}}%
\pgfpathlineto{\pgfqpoint{5.255831in}{3.416960in}}%
\pgfpathlineto{\pgfqpoint{5.258349in}{3.439711in}}%
\pgfpathlineto{\pgfqpoint{5.260868in}{3.442045in}}%
\pgfpathlineto{\pgfqpoint{5.263387in}{3.394247in}}%
\pgfpathlineto{\pgfqpoint{5.265906in}{3.395093in}}%
\pgfpathlineto{\pgfqpoint{5.268424in}{3.412261in}}%
\pgfpathlineto{\pgfqpoint{5.270943in}{3.438251in}}%
\pgfpathlineto{\pgfqpoint{5.273462in}{3.424416in}}%
\pgfpathlineto{\pgfqpoint{5.275981in}{3.439160in}}%
\pgfpathlineto{\pgfqpoint{5.278499in}{3.462817in}}%
\pgfpathlineto{\pgfqpoint{5.281018in}{3.500582in}}%
\pgfpathlineto{\pgfqpoint{5.283537in}{3.471667in}}%
\pgfpathlineto{\pgfqpoint{5.286056in}{3.495289in}}%
\pgfpathlineto{\pgfqpoint{5.288574in}{3.490505in}}%
\pgfpathlineto{\pgfqpoint{5.291093in}{3.534080in}}%
\pgfpathlineto{\pgfqpoint{5.293612in}{3.574309in}}%
\pgfpathlineto{\pgfqpoint{5.296131in}{3.587429in}}%
\pgfpathlineto{\pgfqpoint{5.298649in}{3.558572in}}%
\pgfpathlineto{\pgfqpoint{5.301168in}{3.561834in}}%
\pgfpathlineto{\pgfqpoint{5.303687in}{3.579456in}}%
\pgfpathlineto{\pgfqpoint{5.306206in}{3.563675in}}%
\pgfpathlineto{\pgfqpoint{5.308724in}{3.512350in}}%
\pgfpathlineto{\pgfqpoint{5.311243in}{3.500970in}}%
\pgfpathlineto{\pgfqpoint{5.313762in}{3.500830in}}%
\pgfpathlineto{\pgfqpoint{5.316281in}{3.542241in}}%
\pgfpathlineto{\pgfqpoint{5.318799in}{3.554011in}}%
\pgfpathlineto{\pgfqpoint{5.321318in}{3.564645in}}%
\pgfpathlineto{\pgfqpoint{5.323837in}{3.603209in}}%
\pgfpathlineto{\pgfqpoint{5.326356in}{3.580636in}}%
\pgfpathlineto{\pgfqpoint{5.328874in}{3.570149in}}%
\pgfpathlineto{\pgfqpoint{5.331393in}{3.582924in}}%
\pgfpathlineto{\pgfqpoint{5.333912in}{3.566624in}}%
\pgfpathlineto{\pgfqpoint{5.336431in}{3.567836in}}%
\pgfpathlineto{\pgfqpoint{5.338949in}{3.576189in}}%
\pgfpathlineto{\pgfqpoint{5.341468in}{3.579892in}}%
\pgfpathlineto{\pgfqpoint{5.343987in}{3.571970in}}%
\pgfpathlineto{\pgfqpoint{5.346506in}{3.591938in}}%
\pgfpathlineto{\pgfqpoint{5.349024in}{3.579777in}}%
\pgfpathlineto{\pgfqpoint{5.351543in}{3.515588in}}%
\pgfpathlineto{\pgfqpoint{5.354062in}{3.493870in}}%
\pgfpathlineto{\pgfqpoint{5.359099in}{3.478188in}}%
\pgfpathlineto{\pgfqpoint{5.361618in}{3.453727in}}%
\pgfpathlineto{\pgfqpoint{5.364137in}{3.458873in}}%
\pgfpathlineto{\pgfqpoint{5.366656in}{3.483303in}}%
\pgfpathlineto{\pgfqpoint{5.369174in}{3.472326in}}%
\pgfpathlineto{\pgfqpoint{5.371693in}{3.455198in}}%
\pgfpathlineto{\pgfqpoint{5.374212in}{3.485257in}}%
\pgfpathlineto{\pgfqpoint{5.376731in}{3.487103in}}%
\pgfpathlineto{\pgfqpoint{5.379249in}{3.463273in}}%
\pgfpathlineto{\pgfqpoint{5.381768in}{3.501858in}}%
\pgfpathlineto{\pgfqpoint{5.384287in}{3.503776in}}%
\pgfpathlineto{\pgfqpoint{5.386806in}{3.458736in}}%
\pgfpathlineto{\pgfqpoint{5.389324in}{3.478090in}}%
\pgfpathlineto{\pgfqpoint{5.391843in}{3.468536in}}%
\pgfpathlineto{\pgfqpoint{5.394362in}{3.489040in}}%
\pgfpathlineto{\pgfqpoint{5.396881in}{3.487818in}}%
\pgfpathlineto{\pgfqpoint{5.399399in}{3.507671in}}%
\pgfpathlineto{\pgfqpoint{5.401918in}{3.545649in}}%
\pgfpathlineto{\pgfqpoint{5.404437in}{3.536427in}}%
\pgfpathlineto{\pgfqpoint{5.406956in}{3.559060in}}%
\pgfpathlineto{\pgfqpoint{5.409474in}{3.586002in}}%
\pgfpathlineto{\pgfqpoint{5.411993in}{3.566339in}}%
\pgfpathlineto{\pgfqpoint{5.414512in}{3.544505in}}%
\pgfpathlineto{\pgfqpoint{5.417031in}{3.557580in}}%
\pgfpathlineto{\pgfqpoint{5.419549in}{3.588041in}}%
\pgfpathlineto{\pgfqpoint{5.422068in}{3.581364in}}%
\pgfpathlineto{\pgfqpoint{5.424587in}{3.577874in}}%
\pgfpathlineto{\pgfqpoint{5.427106in}{3.628421in}}%
\pgfpathlineto{\pgfqpoint{5.429624in}{3.653202in}}%
\pgfpathlineto{\pgfqpoint{5.432143in}{3.682356in}}%
\pgfpathlineto{\pgfqpoint{5.434662in}{3.707150in}}%
\pgfpathlineto{\pgfqpoint{5.437181in}{3.738255in}}%
\pgfpathlineto{\pgfqpoint{5.439699in}{3.779397in}}%
\pgfpathlineto{\pgfqpoint{5.442218in}{3.794215in}}%
\pgfpathlineto{\pgfqpoint{5.444737in}{3.848607in}}%
\pgfpathlineto{\pgfqpoint{5.447256in}{3.861362in}}%
\pgfpathlineto{\pgfqpoint{5.449774in}{3.886149in}}%
\pgfpathlineto{\pgfqpoint{5.452293in}{3.901956in}}%
\pgfpathlineto{\pgfqpoint{5.454812in}{3.902036in}}%
\pgfpathlineto{\pgfqpoint{5.457331in}{3.943055in}}%
\pgfpathlineto{\pgfqpoint{5.459849in}{3.949338in}}%
\pgfpathlineto{\pgfqpoint{5.462368in}{3.891519in}}%
\pgfpathlineto{\pgfqpoint{5.464887in}{3.932322in}}%
\pgfpathlineto{\pgfqpoint{5.467406in}{3.948631in}}%
\pgfpathlineto{\pgfqpoint{5.469924in}{3.988649in}}%
\pgfpathlineto{\pgfqpoint{5.472443in}{3.973229in}}%
\pgfpathlineto{\pgfqpoint{5.474962in}{3.942736in}}%
\pgfpathlineto{\pgfqpoint{5.477481in}{3.945564in}}%
\pgfpathlineto{\pgfqpoint{5.479999in}{3.960524in}}%
\pgfpathlineto{\pgfqpoint{5.482518in}{3.949778in}}%
\pgfpathlineto{\pgfqpoint{5.485037in}{3.922975in}}%
\pgfpathlineto{\pgfqpoint{5.487556in}{3.863934in}}%
\pgfpathlineto{\pgfqpoint{5.490074in}{3.900189in}}%
\pgfpathlineto{\pgfqpoint{5.492593in}{3.890996in}}%
\pgfpathlineto{\pgfqpoint{5.492593in}{3.890996in}}%
\pgfusepath{stroke}%
\end{pgfscope}%
\begin{pgfscope}%
\pgfpathrectangle{\pgfqpoint{0.455093in}{0.401080in}}{\pgfqpoint{5.037500in}{4.004000in}}%
\pgfusepath{clip}%
\pgfsetrectcap%
\pgfsetroundjoin%
\pgfsetlinewidth{0.501875pt}%
\definecolor{currentstroke}{rgb}{0.690196,0.188235,0.376471}%
\pgfsetstrokecolor{currentstroke}%
\pgfsetstrokeopacity{0.800000}%
\pgfsetdash{}{0pt}%
\pgfpathmoveto{\pgfqpoint{0.455093in}{1.021044in}}%
\pgfpathlineto{\pgfqpoint{0.457612in}{1.011511in}}%
\pgfpathlineto{\pgfqpoint{0.460131in}{1.013365in}}%
\pgfpathlineto{\pgfqpoint{0.462649in}{1.007434in}}%
\pgfpathlineto{\pgfqpoint{0.465168in}{1.004596in}}%
\pgfpathlineto{\pgfqpoint{0.467687in}{0.971475in}}%
\pgfpathlineto{\pgfqpoint{0.470206in}{0.981668in}}%
\pgfpathlineto{\pgfqpoint{0.472724in}{0.982317in}}%
\pgfpathlineto{\pgfqpoint{0.475243in}{0.983688in}}%
\pgfpathlineto{\pgfqpoint{0.477762in}{0.974390in}}%
\pgfpathlineto{\pgfqpoint{0.480281in}{0.966089in}}%
\pgfpathlineto{\pgfqpoint{0.482799in}{0.956495in}}%
\pgfpathlineto{\pgfqpoint{0.485318in}{0.956180in}}%
\pgfpathlineto{\pgfqpoint{0.487837in}{0.954302in}}%
\pgfpathlineto{\pgfqpoint{0.490356in}{0.939990in}}%
\pgfpathlineto{\pgfqpoint{0.492874in}{0.935923in}}%
\pgfpathlineto{\pgfqpoint{0.495393in}{0.958865in}}%
\pgfpathlineto{\pgfqpoint{0.497912in}{0.939795in}}%
\pgfpathlineto{\pgfqpoint{0.500431in}{0.954620in}}%
\pgfpathlineto{\pgfqpoint{0.502949in}{0.938903in}}%
\pgfpathlineto{\pgfqpoint{0.505468in}{0.960529in}}%
\pgfpathlineto{\pgfqpoint{0.507987in}{0.986886in}}%
\pgfpathlineto{\pgfqpoint{0.510506in}{0.979962in}}%
\pgfpathlineto{\pgfqpoint{0.513024in}{0.983186in}}%
\pgfpathlineto{\pgfqpoint{0.515543in}{0.980939in}}%
\pgfpathlineto{\pgfqpoint{0.518062in}{0.967521in}}%
\pgfpathlineto{\pgfqpoint{0.520581in}{0.966719in}}%
\pgfpathlineto{\pgfqpoint{0.523099in}{0.965791in}}%
\pgfpathlineto{\pgfqpoint{0.525618in}{0.972050in}}%
\pgfpathlineto{\pgfqpoint{0.528137in}{0.958146in}}%
\pgfpathlineto{\pgfqpoint{0.530656in}{0.953964in}}%
\pgfpathlineto{\pgfqpoint{0.533174in}{0.946587in}}%
\pgfpathlineto{\pgfqpoint{0.535693in}{0.952217in}}%
\pgfpathlineto{\pgfqpoint{0.538212in}{0.947503in}}%
\pgfpathlineto{\pgfqpoint{0.540731in}{0.983057in}}%
\pgfpathlineto{\pgfqpoint{0.543249in}{0.953156in}}%
\pgfpathlineto{\pgfqpoint{0.545768in}{0.960360in}}%
\pgfpathlineto{\pgfqpoint{0.548287in}{0.937481in}}%
\pgfpathlineto{\pgfqpoint{0.550806in}{0.947388in}}%
\pgfpathlineto{\pgfqpoint{0.553324in}{0.980887in}}%
\pgfpathlineto{\pgfqpoint{0.555843in}{0.965709in}}%
\pgfpathlineto{\pgfqpoint{0.558362in}{0.970501in}}%
\pgfpathlineto{\pgfqpoint{0.560881in}{0.972298in}}%
\pgfpathlineto{\pgfqpoint{0.563399in}{0.982493in}}%
\pgfpathlineto{\pgfqpoint{0.565918in}{1.006600in}}%
\pgfpathlineto{\pgfqpoint{0.568437in}{1.008105in}}%
\pgfpathlineto{\pgfqpoint{0.570956in}{0.998190in}}%
\pgfpathlineto{\pgfqpoint{0.573474in}{0.998448in}}%
\pgfpathlineto{\pgfqpoint{0.575993in}{0.994213in}}%
\pgfpathlineto{\pgfqpoint{0.578512in}{1.007678in}}%
\pgfpathlineto{\pgfqpoint{0.581031in}{1.025236in}}%
\pgfpathlineto{\pgfqpoint{0.583549in}{1.027491in}}%
\pgfpathlineto{\pgfqpoint{0.586068in}{1.029264in}}%
\pgfpathlineto{\pgfqpoint{0.588587in}{1.018982in}}%
\pgfpathlineto{\pgfqpoint{0.591106in}{1.024363in}}%
\pgfpathlineto{\pgfqpoint{0.593624in}{1.006862in}}%
\pgfpathlineto{\pgfqpoint{0.596143in}{1.005791in}}%
\pgfpathlineto{\pgfqpoint{0.598662in}{0.999494in}}%
\pgfpathlineto{\pgfqpoint{0.601181in}{0.993560in}}%
\pgfpathlineto{\pgfqpoint{0.603699in}{0.988738in}}%
\pgfpathlineto{\pgfqpoint{0.606218in}{0.987926in}}%
\pgfpathlineto{\pgfqpoint{0.608737in}{1.013179in}}%
\pgfpathlineto{\pgfqpoint{0.611256in}{0.998662in}}%
\pgfpathlineto{\pgfqpoint{0.613774in}{1.001389in}}%
\pgfpathlineto{\pgfqpoint{0.616293in}{1.000658in}}%
\pgfpathlineto{\pgfqpoint{0.618812in}{1.001013in}}%
\pgfpathlineto{\pgfqpoint{0.621331in}{0.998395in}}%
\pgfpathlineto{\pgfqpoint{0.623849in}{0.992656in}}%
\pgfpathlineto{\pgfqpoint{0.626368in}{0.993565in}}%
\pgfpathlineto{\pgfqpoint{0.628887in}{0.994027in}}%
\pgfpathlineto{\pgfqpoint{0.631406in}{0.992209in}}%
\pgfpathlineto{\pgfqpoint{0.633924in}{0.978428in}}%
\pgfpathlineto{\pgfqpoint{0.636443in}{0.962206in}}%
\pgfpathlineto{\pgfqpoint{0.638962in}{0.962914in}}%
\pgfpathlineto{\pgfqpoint{0.641481in}{0.944849in}}%
\pgfpathlineto{\pgfqpoint{0.643999in}{0.943859in}}%
\pgfpathlineto{\pgfqpoint{0.646518in}{0.933857in}}%
\pgfpathlineto{\pgfqpoint{0.649037in}{0.938717in}}%
\pgfpathlineto{\pgfqpoint{0.651556in}{0.946928in}}%
\pgfpathlineto{\pgfqpoint{0.654074in}{0.950450in}}%
\pgfpathlineto{\pgfqpoint{0.656593in}{0.956995in}}%
\pgfpathlineto{\pgfqpoint{0.659112in}{0.937383in}}%
\pgfpathlineto{\pgfqpoint{0.661631in}{0.937851in}}%
\pgfpathlineto{\pgfqpoint{0.664149in}{0.940622in}}%
\pgfpathlineto{\pgfqpoint{0.666668in}{0.928511in}}%
\pgfpathlineto{\pgfqpoint{0.669187in}{0.934819in}}%
\pgfpathlineto{\pgfqpoint{0.671706in}{0.922782in}}%
\pgfpathlineto{\pgfqpoint{0.674224in}{0.921528in}}%
\pgfpathlineto{\pgfqpoint{0.676743in}{0.921412in}}%
\pgfpathlineto{\pgfqpoint{0.679262in}{0.916861in}}%
\pgfpathlineto{\pgfqpoint{0.681781in}{0.899197in}}%
\pgfpathlineto{\pgfqpoint{0.684299in}{0.899564in}}%
\pgfpathlineto{\pgfqpoint{0.686818in}{0.914656in}}%
\pgfpathlineto{\pgfqpoint{0.689337in}{0.921658in}}%
\pgfpathlineto{\pgfqpoint{0.691856in}{0.926163in}}%
\pgfpathlineto{\pgfqpoint{0.694374in}{0.903769in}}%
\pgfpathlineto{\pgfqpoint{0.696893in}{0.922634in}}%
\pgfpathlineto{\pgfqpoint{0.699412in}{0.924369in}}%
\pgfpathlineto{\pgfqpoint{0.701931in}{0.938361in}}%
\pgfpathlineto{\pgfqpoint{0.704449in}{0.948369in}}%
\pgfpathlineto{\pgfqpoint{0.706968in}{0.956124in}}%
\pgfpathlineto{\pgfqpoint{0.709487in}{0.961025in}}%
\pgfpathlineto{\pgfqpoint{0.712006in}{0.964520in}}%
\pgfpathlineto{\pgfqpoint{0.714524in}{0.972825in}}%
\pgfpathlineto{\pgfqpoint{0.717043in}{0.979631in}}%
\pgfpathlineto{\pgfqpoint{0.719562in}{0.968286in}}%
\pgfpathlineto{\pgfqpoint{0.722081in}{0.975145in}}%
\pgfpathlineto{\pgfqpoint{0.724599in}{0.977666in}}%
\pgfpathlineto{\pgfqpoint{0.727118in}{0.988496in}}%
\pgfpathlineto{\pgfqpoint{0.729637in}{1.000680in}}%
\pgfpathlineto{\pgfqpoint{0.732156in}{0.973189in}}%
\pgfpathlineto{\pgfqpoint{0.734674in}{0.971426in}}%
\pgfpathlineto{\pgfqpoint{0.737193in}{0.978298in}}%
\pgfpathlineto{\pgfqpoint{0.739712in}{0.998321in}}%
\pgfpathlineto{\pgfqpoint{0.742231in}{1.004841in}}%
\pgfpathlineto{\pgfqpoint{0.744749in}{0.994578in}}%
\pgfpathlineto{\pgfqpoint{0.747268in}{0.969796in}}%
\pgfpathlineto{\pgfqpoint{0.749787in}{0.950119in}}%
\pgfpathlineto{\pgfqpoint{0.752306in}{0.932933in}}%
\pgfpathlineto{\pgfqpoint{0.754824in}{0.926272in}}%
\pgfpathlineto{\pgfqpoint{0.757343in}{0.912963in}}%
\pgfpathlineto{\pgfqpoint{0.759862in}{0.922908in}}%
\pgfpathlineto{\pgfqpoint{0.762381in}{0.920696in}}%
\pgfpathlineto{\pgfqpoint{0.764899in}{0.929308in}}%
\pgfpathlineto{\pgfqpoint{0.767418in}{0.925617in}}%
\pgfpathlineto{\pgfqpoint{0.769937in}{0.936703in}}%
\pgfpathlineto{\pgfqpoint{0.772456in}{0.958997in}}%
\pgfpathlineto{\pgfqpoint{0.774974in}{0.957236in}}%
\pgfpathlineto{\pgfqpoint{0.777493in}{0.949777in}}%
\pgfpathlineto{\pgfqpoint{0.780012in}{0.954640in}}%
\pgfpathlineto{\pgfqpoint{0.782531in}{0.960555in}}%
\pgfpathlineto{\pgfqpoint{0.785049in}{0.948159in}}%
\pgfpathlineto{\pgfqpoint{0.787568in}{0.950350in}}%
\pgfpathlineto{\pgfqpoint{0.790087in}{0.969021in}}%
\pgfpathlineto{\pgfqpoint{0.792606in}{0.975488in}}%
\pgfpathlineto{\pgfqpoint{0.797643in}{0.986435in}}%
\pgfpathlineto{\pgfqpoint{0.800162in}{0.987401in}}%
\pgfpathlineto{\pgfqpoint{0.802681in}{0.982507in}}%
\pgfpathlineto{\pgfqpoint{0.805199in}{0.970383in}}%
\pgfpathlineto{\pgfqpoint{0.807718in}{0.958960in}}%
\pgfpathlineto{\pgfqpoint{0.810237in}{0.956889in}}%
\pgfpathlineto{\pgfqpoint{0.812756in}{0.938578in}}%
\pgfpathlineto{\pgfqpoint{0.815274in}{0.942356in}}%
\pgfpathlineto{\pgfqpoint{0.817793in}{0.928258in}}%
\pgfpathlineto{\pgfqpoint{0.820312in}{0.901951in}}%
\pgfpathlineto{\pgfqpoint{0.822831in}{0.892281in}}%
\pgfpathlineto{\pgfqpoint{0.825349in}{0.879409in}}%
\pgfpathlineto{\pgfqpoint{0.827868in}{0.893575in}}%
\pgfpathlineto{\pgfqpoint{0.830387in}{0.903195in}}%
\pgfpathlineto{\pgfqpoint{0.832906in}{0.910126in}}%
\pgfpathlineto{\pgfqpoint{0.835424in}{0.910977in}}%
\pgfpathlineto{\pgfqpoint{0.837943in}{0.926247in}}%
\pgfpathlineto{\pgfqpoint{0.840462in}{0.918178in}}%
\pgfpathlineto{\pgfqpoint{0.842981in}{0.914970in}}%
\pgfpathlineto{\pgfqpoint{0.845499in}{0.912184in}}%
\pgfpathlineto{\pgfqpoint{0.848018in}{0.883055in}}%
\pgfpathlineto{\pgfqpoint{0.850537in}{0.897727in}}%
\pgfpathlineto{\pgfqpoint{0.853056in}{0.890545in}}%
\pgfpathlineto{\pgfqpoint{0.855574in}{0.881242in}}%
\pgfpathlineto{\pgfqpoint{0.858093in}{0.880826in}}%
\pgfpathlineto{\pgfqpoint{0.860612in}{0.881524in}}%
\pgfpathlineto{\pgfqpoint{0.863131in}{0.891225in}}%
\pgfpathlineto{\pgfqpoint{0.865649in}{0.906737in}}%
\pgfpathlineto{\pgfqpoint{0.868168in}{0.929125in}}%
\pgfpathlineto{\pgfqpoint{0.870687in}{0.939157in}}%
\pgfpathlineto{\pgfqpoint{0.873206in}{0.940979in}}%
\pgfpathlineto{\pgfqpoint{0.875724in}{0.942936in}}%
\pgfpathlineto{\pgfqpoint{0.878243in}{0.946985in}}%
\pgfpathlineto{\pgfqpoint{0.880762in}{0.974620in}}%
\pgfpathlineto{\pgfqpoint{0.883281in}{0.997815in}}%
\pgfpathlineto{\pgfqpoint{0.885799in}{1.015175in}}%
\pgfpathlineto{\pgfqpoint{0.888318in}{1.025161in}}%
\pgfpathlineto{\pgfqpoint{0.890837in}{1.045721in}}%
\pgfpathlineto{\pgfqpoint{0.893356in}{1.046688in}}%
\pgfpathlineto{\pgfqpoint{0.895874in}{1.066892in}}%
\pgfpathlineto{\pgfqpoint{0.898393in}{1.046736in}}%
\pgfpathlineto{\pgfqpoint{0.900912in}{1.050051in}}%
\pgfpathlineto{\pgfqpoint{0.903431in}{1.071858in}}%
\pgfpathlineto{\pgfqpoint{0.905949in}{1.081275in}}%
\pgfpathlineto{\pgfqpoint{0.908468in}{1.087229in}}%
\pgfpathlineto{\pgfqpoint{0.910987in}{1.079724in}}%
\pgfpathlineto{\pgfqpoint{0.913506in}{1.055801in}}%
\pgfpathlineto{\pgfqpoint{0.916024in}{1.067463in}}%
\pgfpathlineto{\pgfqpoint{0.918543in}{1.053951in}}%
\pgfpathlineto{\pgfqpoint{0.921062in}{1.042360in}}%
\pgfpathlineto{\pgfqpoint{0.923581in}{1.054619in}}%
\pgfpathlineto{\pgfqpoint{0.926099in}{1.054550in}}%
\pgfpathlineto{\pgfqpoint{0.928618in}{1.064365in}}%
\pgfpathlineto{\pgfqpoint{0.931137in}{1.068139in}}%
\pgfpathlineto{\pgfqpoint{0.933656in}{1.075559in}}%
\pgfpathlineto{\pgfqpoint{0.936174in}{1.081130in}}%
\pgfpathlineto{\pgfqpoint{0.938693in}{1.095396in}}%
\pgfpathlineto{\pgfqpoint{0.941212in}{1.122853in}}%
\pgfpathlineto{\pgfqpoint{0.943731in}{1.116623in}}%
\pgfpathlineto{\pgfqpoint{0.946249in}{1.144589in}}%
\pgfpathlineto{\pgfqpoint{0.948768in}{1.132334in}}%
\pgfpathlineto{\pgfqpoint{0.951287in}{1.136392in}}%
\pgfpathlineto{\pgfqpoint{0.953806in}{1.135919in}}%
\pgfpathlineto{\pgfqpoint{0.956324in}{1.124269in}}%
\pgfpathlineto{\pgfqpoint{0.958843in}{1.130388in}}%
\pgfpathlineto{\pgfqpoint{0.961362in}{1.113113in}}%
\pgfpathlineto{\pgfqpoint{0.963881in}{1.136371in}}%
\pgfpathlineto{\pgfqpoint{0.966399in}{1.138043in}}%
\pgfpathlineto{\pgfqpoint{0.968918in}{1.134445in}}%
\pgfpathlineto{\pgfqpoint{0.971437in}{1.140940in}}%
\pgfpathlineto{\pgfqpoint{0.973956in}{1.137550in}}%
\pgfpathlineto{\pgfqpoint{0.976474in}{1.153357in}}%
\pgfpathlineto{\pgfqpoint{0.978993in}{1.160052in}}%
\pgfpathlineto{\pgfqpoint{0.981512in}{1.178490in}}%
\pgfpathlineto{\pgfqpoint{0.984031in}{1.183071in}}%
\pgfpathlineto{\pgfqpoint{0.986549in}{1.190957in}}%
\pgfpathlineto{\pgfqpoint{0.989068in}{1.175955in}}%
\pgfpathlineto{\pgfqpoint{0.991587in}{1.147602in}}%
\pgfpathlineto{\pgfqpoint{0.994106in}{1.142571in}}%
\pgfpathlineto{\pgfqpoint{0.996624in}{1.126962in}}%
\pgfpathlineto{\pgfqpoint{0.999143in}{1.136107in}}%
\pgfpathlineto{\pgfqpoint{1.001662in}{1.139369in}}%
\pgfpathlineto{\pgfqpoint{1.004181in}{1.137899in}}%
\pgfpathlineto{\pgfqpoint{1.006699in}{1.143751in}}%
\pgfpathlineto{\pgfqpoint{1.009218in}{1.134537in}}%
\pgfpathlineto{\pgfqpoint{1.011737in}{1.132105in}}%
\pgfpathlineto{\pgfqpoint{1.014256in}{1.127032in}}%
\pgfpathlineto{\pgfqpoint{1.016774in}{1.123519in}}%
\pgfpathlineto{\pgfqpoint{1.019293in}{1.114850in}}%
\pgfpathlineto{\pgfqpoint{1.021812in}{1.112502in}}%
\pgfpathlineto{\pgfqpoint{1.024331in}{1.120890in}}%
\pgfpathlineto{\pgfqpoint{1.026849in}{1.132212in}}%
\pgfpathlineto{\pgfqpoint{1.029368in}{1.136382in}}%
\pgfpathlineto{\pgfqpoint{1.031887in}{1.153592in}}%
\pgfpathlineto{\pgfqpoint{1.034406in}{1.138338in}}%
\pgfpathlineto{\pgfqpoint{1.036924in}{1.130145in}}%
\pgfpathlineto{\pgfqpoint{1.039443in}{1.146674in}}%
\pgfpathlineto{\pgfqpoint{1.041962in}{1.164011in}}%
\pgfpathlineto{\pgfqpoint{1.044481in}{1.186506in}}%
\pgfpathlineto{\pgfqpoint{1.046999in}{1.214192in}}%
\pgfpathlineto{\pgfqpoint{1.049518in}{1.215116in}}%
\pgfpathlineto{\pgfqpoint{1.052037in}{1.215337in}}%
\pgfpathlineto{\pgfqpoint{1.054556in}{1.219422in}}%
\pgfpathlineto{\pgfqpoint{1.057074in}{1.206661in}}%
\pgfpathlineto{\pgfqpoint{1.059593in}{1.196059in}}%
\pgfpathlineto{\pgfqpoint{1.062112in}{1.204658in}}%
\pgfpathlineto{\pgfqpoint{1.064631in}{1.201070in}}%
\pgfpathlineto{\pgfqpoint{1.067149in}{1.186052in}}%
\pgfpathlineto{\pgfqpoint{1.069668in}{1.206917in}}%
\pgfpathlineto{\pgfqpoint{1.072187in}{1.208441in}}%
\pgfpathlineto{\pgfqpoint{1.074706in}{1.220105in}}%
\pgfpathlineto{\pgfqpoint{1.079743in}{1.247410in}}%
\pgfpathlineto{\pgfqpoint{1.082262in}{1.241804in}}%
\pgfpathlineto{\pgfqpoint{1.084781in}{1.261022in}}%
\pgfpathlineto{\pgfqpoint{1.087299in}{1.262115in}}%
\pgfpathlineto{\pgfqpoint{1.089818in}{1.281957in}}%
\pgfpathlineto{\pgfqpoint{1.092337in}{1.267817in}}%
\pgfpathlineto{\pgfqpoint{1.094856in}{1.298161in}}%
\pgfpathlineto{\pgfqpoint{1.097374in}{1.291394in}}%
\pgfpathlineto{\pgfqpoint{1.099893in}{1.306503in}}%
\pgfpathlineto{\pgfqpoint{1.102412in}{1.336505in}}%
\pgfpathlineto{\pgfqpoint{1.104931in}{1.311845in}}%
\pgfpathlineto{\pgfqpoint{1.107449in}{1.295325in}}%
\pgfpathlineto{\pgfqpoint{1.109968in}{1.282520in}}%
\pgfpathlineto{\pgfqpoint{1.112487in}{1.272917in}}%
\pgfpathlineto{\pgfqpoint{1.115006in}{1.267187in}}%
\pgfpathlineto{\pgfqpoint{1.117524in}{1.282125in}}%
\pgfpathlineto{\pgfqpoint{1.120043in}{1.269147in}}%
\pgfpathlineto{\pgfqpoint{1.122562in}{1.286764in}}%
\pgfpathlineto{\pgfqpoint{1.125081in}{1.278593in}}%
\pgfpathlineto{\pgfqpoint{1.127599in}{1.286663in}}%
\pgfpathlineto{\pgfqpoint{1.130118in}{1.293023in}}%
\pgfpathlineto{\pgfqpoint{1.132637in}{1.291851in}}%
\pgfpathlineto{\pgfqpoint{1.135156in}{1.289485in}}%
\pgfpathlineto{\pgfqpoint{1.137674in}{1.275560in}}%
\pgfpathlineto{\pgfqpoint{1.140193in}{1.255312in}}%
\pgfpathlineto{\pgfqpoint{1.142712in}{1.277221in}}%
\pgfpathlineto{\pgfqpoint{1.145231in}{1.274983in}}%
\pgfpathlineto{\pgfqpoint{1.147749in}{1.264544in}}%
\pgfpathlineto{\pgfqpoint{1.155306in}{1.325428in}}%
\pgfpathlineto{\pgfqpoint{1.157824in}{1.325563in}}%
\pgfpathlineto{\pgfqpoint{1.160343in}{1.319040in}}%
\pgfpathlineto{\pgfqpoint{1.162862in}{1.307069in}}%
\pgfpathlineto{\pgfqpoint{1.165381in}{1.319473in}}%
\pgfpathlineto{\pgfqpoint{1.167899in}{1.321115in}}%
\pgfpathlineto{\pgfqpoint{1.170418in}{1.308511in}}%
\pgfpathlineto{\pgfqpoint{1.172937in}{1.330620in}}%
\pgfpathlineto{\pgfqpoint{1.175456in}{1.316451in}}%
\pgfpathlineto{\pgfqpoint{1.177974in}{1.326951in}}%
\pgfpathlineto{\pgfqpoint{1.180493in}{1.338974in}}%
\pgfpathlineto{\pgfqpoint{1.183012in}{1.329239in}}%
\pgfpathlineto{\pgfqpoint{1.185531in}{1.310144in}}%
\pgfpathlineto{\pgfqpoint{1.188049in}{1.325071in}}%
\pgfpathlineto{\pgfqpoint{1.190568in}{1.331178in}}%
\pgfpathlineto{\pgfqpoint{1.193087in}{1.338926in}}%
\pgfpathlineto{\pgfqpoint{1.195606in}{1.330731in}}%
\pgfpathlineto{\pgfqpoint{1.198124in}{1.334999in}}%
\pgfpathlineto{\pgfqpoint{1.200643in}{1.324401in}}%
\pgfpathlineto{\pgfqpoint{1.203162in}{1.309773in}}%
\pgfpathlineto{\pgfqpoint{1.205681in}{1.313890in}}%
\pgfpathlineto{\pgfqpoint{1.208199in}{1.314222in}}%
\pgfpathlineto{\pgfqpoint{1.210718in}{1.316874in}}%
\pgfpathlineto{\pgfqpoint{1.213237in}{1.312296in}}%
\pgfpathlineto{\pgfqpoint{1.215756in}{1.328496in}}%
\pgfpathlineto{\pgfqpoint{1.218274in}{1.341499in}}%
\pgfpathlineto{\pgfqpoint{1.220793in}{1.326434in}}%
\pgfpathlineto{\pgfqpoint{1.223312in}{1.333470in}}%
\pgfpathlineto{\pgfqpoint{1.225831in}{1.322747in}}%
\pgfpathlineto{\pgfqpoint{1.228349in}{1.322532in}}%
\pgfpathlineto{\pgfqpoint{1.230868in}{1.276302in}}%
\pgfpathlineto{\pgfqpoint{1.233387in}{1.287737in}}%
\pgfpathlineto{\pgfqpoint{1.235906in}{1.296648in}}%
\pgfpathlineto{\pgfqpoint{1.238424in}{1.290380in}}%
\pgfpathlineto{\pgfqpoint{1.240943in}{1.314287in}}%
\pgfpathlineto{\pgfqpoint{1.243462in}{1.306459in}}%
\pgfpathlineto{\pgfqpoint{1.245981in}{1.337528in}}%
\pgfpathlineto{\pgfqpoint{1.248499in}{1.334643in}}%
\pgfpathlineto{\pgfqpoint{1.251018in}{1.321548in}}%
\pgfpathlineto{\pgfqpoint{1.253537in}{1.330745in}}%
\pgfpathlineto{\pgfqpoint{1.256056in}{1.321459in}}%
\pgfpathlineto{\pgfqpoint{1.258574in}{1.316447in}}%
\pgfpathlineto{\pgfqpoint{1.261093in}{1.320880in}}%
\pgfpathlineto{\pgfqpoint{1.263612in}{1.338424in}}%
\pgfpathlineto{\pgfqpoint{1.266131in}{1.344975in}}%
\pgfpathlineto{\pgfqpoint{1.268649in}{1.335585in}}%
\pgfpathlineto{\pgfqpoint{1.271168in}{1.366927in}}%
\pgfpathlineto{\pgfqpoint{1.273687in}{1.357141in}}%
\pgfpathlineto{\pgfqpoint{1.276206in}{1.359808in}}%
\pgfpathlineto{\pgfqpoint{1.278724in}{1.375281in}}%
\pgfpathlineto{\pgfqpoint{1.281243in}{1.385213in}}%
\pgfpathlineto{\pgfqpoint{1.283762in}{1.425166in}}%
\pgfpathlineto{\pgfqpoint{1.286281in}{1.444024in}}%
\pgfpathlineto{\pgfqpoint{1.291318in}{1.473283in}}%
\pgfpathlineto{\pgfqpoint{1.293837in}{1.485280in}}%
\pgfpathlineto{\pgfqpoint{1.296356in}{1.478654in}}%
\pgfpathlineto{\pgfqpoint{1.298874in}{1.487574in}}%
\pgfpathlineto{\pgfqpoint{1.301393in}{1.515642in}}%
\pgfpathlineto{\pgfqpoint{1.303912in}{1.519405in}}%
\pgfpathlineto{\pgfqpoint{1.306431in}{1.532947in}}%
\pgfpathlineto{\pgfqpoint{1.308949in}{1.532844in}}%
\pgfpathlineto{\pgfqpoint{1.311468in}{1.535650in}}%
\pgfpathlineto{\pgfqpoint{1.313987in}{1.530267in}}%
\pgfpathlineto{\pgfqpoint{1.316506in}{1.549967in}}%
\pgfpathlineto{\pgfqpoint{1.319024in}{1.571075in}}%
\pgfpathlineto{\pgfqpoint{1.321543in}{1.561016in}}%
\pgfpathlineto{\pgfqpoint{1.324062in}{1.557642in}}%
\pgfpathlineto{\pgfqpoint{1.326581in}{1.574968in}}%
\pgfpathlineto{\pgfqpoint{1.329099in}{1.572007in}}%
\pgfpathlineto{\pgfqpoint{1.331618in}{1.559851in}}%
\pgfpathlineto{\pgfqpoint{1.334137in}{1.536518in}}%
\pgfpathlineto{\pgfqpoint{1.336656in}{1.551237in}}%
\pgfpathlineto{\pgfqpoint{1.339174in}{1.531947in}}%
\pgfpathlineto{\pgfqpoint{1.341693in}{1.545704in}}%
\pgfpathlineto{\pgfqpoint{1.344212in}{1.549796in}}%
\pgfpathlineto{\pgfqpoint{1.346731in}{1.544796in}}%
\pgfpathlineto{\pgfqpoint{1.349249in}{1.547595in}}%
\pgfpathlineto{\pgfqpoint{1.351768in}{1.557738in}}%
\pgfpathlineto{\pgfqpoint{1.354287in}{1.536864in}}%
\pgfpathlineto{\pgfqpoint{1.356806in}{1.559405in}}%
\pgfpathlineto{\pgfqpoint{1.359324in}{1.593574in}}%
\pgfpathlineto{\pgfqpoint{1.361843in}{1.597905in}}%
\pgfpathlineto{\pgfqpoint{1.364362in}{1.609756in}}%
\pgfpathlineto{\pgfqpoint{1.366881in}{1.613007in}}%
\pgfpathlineto{\pgfqpoint{1.369399in}{1.598731in}}%
\pgfpathlineto{\pgfqpoint{1.371918in}{1.576467in}}%
\pgfpathlineto{\pgfqpoint{1.374437in}{1.589281in}}%
\pgfpathlineto{\pgfqpoint{1.376956in}{1.611335in}}%
\pgfpathlineto{\pgfqpoint{1.379474in}{1.600901in}}%
\pgfpathlineto{\pgfqpoint{1.381993in}{1.621242in}}%
\pgfpathlineto{\pgfqpoint{1.384512in}{1.608720in}}%
\pgfpathlineto{\pgfqpoint{1.387031in}{1.632622in}}%
\pgfpathlineto{\pgfqpoint{1.389549in}{1.628168in}}%
\pgfpathlineto{\pgfqpoint{1.392068in}{1.642686in}}%
\pgfpathlineto{\pgfqpoint{1.394587in}{1.615642in}}%
\pgfpathlineto{\pgfqpoint{1.397106in}{1.593369in}}%
\pgfpathlineto{\pgfqpoint{1.399624in}{1.592042in}}%
\pgfpathlineto{\pgfqpoint{1.402143in}{1.595008in}}%
\pgfpathlineto{\pgfqpoint{1.404662in}{1.590502in}}%
\pgfpathlineto{\pgfqpoint{1.407181in}{1.586589in}}%
\pgfpathlineto{\pgfqpoint{1.409699in}{1.573109in}}%
\pgfpathlineto{\pgfqpoint{1.412218in}{1.581248in}}%
\pgfpathlineto{\pgfqpoint{1.414737in}{1.606884in}}%
\pgfpathlineto{\pgfqpoint{1.417256in}{1.625592in}}%
\pgfpathlineto{\pgfqpoint{1.419774in}{1.610317in}}%
\pgfpathlineto{\pgfqpoint{1.422293in}{1.608147in}}%
\pgfpathlineto{\pgfqpoint{1.424812in}{1.614498in}}%
\pgfpathlineto{\pgfqpoint{1.427331in}{1.598836in}}%
\pgfpathlineto{\pgfqpoint{1.429849in}{1.614014in}}%
\pgfpathlineto{\pgfqpoint{1.432368in}{1.577827in}}%
\pgfpathlineto{\pgfqpoint{1.434887in}{1.582277in}}%
\pgfpathlineto{\pgfqpoint{1.437406in}{1.581938in}}%
\pgfpathlineto{\pgfqpoint{1.439924in}{1.548768in}}%
\pgfpathlineto{\pgfqpoint{1.442443in}{1.544906in}}%
\pgfpathlineto{\pgfqpoint{1.444962in}{1.523909in}}%
\pgfpathlineto{\pgfqpoint{1.447481in}{1.547541in}}%
\pgfpathlineto{\pgfqpoint{1.449999in}{1.544357in}}%
\pgfpathlineto{\pgfqpoint{1.452518in}{1.535915in}}%
\pgfpathlineto{\pgfqpoint{1.455037in}{1.554144in}}%
\pgfpathlineto{\pgfqpoint{1.457556in}{1.551892in}}%
\pgfpathlineto{\pgfqpoint{1.460074in}{1.552593in}}%
\pgfpathlineto{\pgfqpoint{1.462593in}{1.552937in}}%
\pgfpathlineto{\pgfqpoint{1.465112in}{1.545929in}}%
\pgfpathlineto{\pgfqpoint{1.467631in}{1.512204in}}%
\pgfpathlineto{\pgfqpoint{1.470149in}{1.531753in}}%
\pgfpathlineto{\pgfqpoint{1.472668in}{1.516998in}}%
\pgfpathlineto{\pgfqpoint{1.475187in}{1.504874in}}%
\pgfpathlineto{\pgfqpoint{1.477706in}{1.518949in}}%
\pgfpathlineto{\pgfqpoint{1.480224in}{1.538616in}}%
\pgfpathlineto{\pgfqpoint{1.482743in}{1.538584in}}%
\pgfpathlineto{\pgfqpoint{1.485262in}{1.575700in}}%
\pgfpathlineto{\pgfqpoint{1.487781in}{1.563942in}}%
\pgfpathlineto{\pgfqpoint{1.490299in}{1.562918in}}%
\pgfpathlineto{\pgfqpoint{1.492818in}{1.550560in}}%
\pgfpathlineto{\pgfqpoint{1.495337in}{1.554338in}}%
\pgfpathlineto{\pgfqpoint{1.497856in}{1.539823in}}%
\pgfpathlineto{\pgfqpoint{1.500374in}{1.523578in}}%
\pgfpathlineto{\pgfqpoint{1.502893in}{1.535053in}}%
\pgfpathlineto{\pgfqpoint{1.505412in}{1.520846in}}%
\pgfpathlineto{\pgfqpoint{1.507931in}{1.521014in}}%
\pgfpathlineto{\pgfqpoint{1.510449in}{1.525108in}}%
\pgfpathlineto{\pgfqpoint{1.512968in}{1.530550in}}%
\pgfpathlineto{\pgfqpoint{1.515487in}{1.506235in}}%
\pgfpathlineto{\pgfqpoint{1.518006in}{1.507369in}}%
\pgfpathlineto{\pgfqpoint{1.520524in}{1.513398in}}%
\pgfpathlineto{\pgfqpoint{1.523043in}{1.504282in}}%
\pgfpathlineto{\pgfqpoint{1.525562in}{1.516303in}}%
\pgfpathlineto{\pgfqpoint{1.528081in}{1.523382in}}%
\pgfpathlineto{\pgfqpoint{1.530599in}{1.536355in}}%
\pgfpathlineto{\pgfqpoint{1.533118in}{1.531912in}}%
\pgfpathlineto{\pgfqpoint{1.535637in}{1.538204in}}%
\pgfpathlineto{\pgfqpoint{1.538156in}{1.550981in}}%
\pgfpathlineto{\pgfqpoint{1.540674in}{1.587857in}}%
\pgfpathlineto{\pgfqpoint{1.543193in}{1.571510in}}%
\pgfpathlineto{\pgfqpoint{1.545712in}{1.580382in}}%
\pgfpathlineto{\pgfqpoint{1.548231in}{1.594661in}}%
\pgfpathlineto{\pgfqpoint{1.550749in}{1.572080in}}%
\pgfpathlineto{\pgfqpoint{1.553268in}{1.563951in}}%
\pgfpathlineto{\pgfqpoint{1.555787in}{1.558659in}}%
\pgfpathlineto{\pgfqpoint{1.558306in}{1.567324in}}%
\pgfpathlineto{\pgfqpoint{1.560824in}{1.559661in}}%
\pgfpathlineto{\pgfqpoint{1.563343in}{1.549165in}}%
\pgfpathlineto{\pgfqpoint{1.565862in}{1.547724in}}%
\pgfpathlineto{\pgfqpoint{1.568381in}{1.544716in}}%
\pgfpathlineto{\pgfqpoint{1.570899in}{1.564960in}}%
\pgfpathlineto{\pgfqpoint{1.573418in}{1.579732in}}%
\pgfpathlineto{\pgfqpoint{1.575937in}{1.584046in}}%
\pgfpathlineto{\pgfqpoint{1.578456in}{1.597917in}}%
\pgfpathlineto{\pgfqpoint{1.580974in}{1.589838in}}%
\pgfpathlineto{\pgfqpoint{1.583493in}{1.582695in}}%
\pgfpathlineto{\pgfqpoint{1.586012in}{1.588236in}}%
\pgfpathlineto{\pgfqpoint{1.588531in}{1.577080in}}%
\pgfpathlineto{\pgfqpoint{1.591049in}{1.572257in}}%
\pgfpathlineto{\pgfqpoint{1.593568in}{1.591262in}}%
\pgfpathlineto{\pgfqpoint{1.596087in}{1.626288in}}%
\pgfpathlineto{\pgfqpoint{1.598606in}{1.644570in}}%
\pgfpathlineto{\pgfqpoint{1.601124in}{1.670648in}}%
\pgfpathlineto{\pgfqpoint{1.603643in}{1.645070in}}%
\pgfpathlineto{\pgfqpoint{1.606162in}{1.664123in}}%
\pgfpathlineto{\pgfqpoint{1.608681in}{1.667652in}}%
\pgfpathlineto{\pgfqpoint{1.611199in}{1.649241in}}%
\pgfpathlineto{\pgfqpoint{1.613718in}{1.651773in}}%
\pgfpathlineto{\pgfqpoint{1.616237in}{1.623550in}}%
\pgfpathlineto{\pgfqpoint{1.621274in}{1.609052in}}%
\pgfpathlineto{\pgfqpoint{1.626312in}{1.637342in}}%
\pgfpathlineto{\pgfqpoint{1.628831in}{1.656274in}}%
\pgfpathlineto{\pgfqpoint{1.631349in}{1.668833in}}%
\pgfpathlineto{\pgfqpoint{1.633868in}{1.629810in}}%
\pgfpathlineto{\pgfqpoint{1.636387in}{1.644268in}}%
\pgfpathlineto{\pgfqpoint{1.638906in}{1.655932in}}%
\pgfpathlineto{\pgfqpoint{1.641424in}{1.653116in}}%
\pgfpathlineto{\pgfqpoint{1.643943in}{1.656047in}}%
\pgfpathlineto{\pgfqpoint{1.646462in}{1.663356in}}%
\pgfpathlineto{\pgfqpoint{1.648981in}{1.655875in}}%
\pgfpathlineto{\pgfqpoint{1.651499in}{1.644191in}}%
\pgfpathlineto{\pgfqpoint{1.654018in}{1.623582in}}%
\pgfpathlineto{\pgfqpoint{1.656537in}{1.617870in}}%
\pgfpathlineto{\pgfqpoint{1.659056in}{1.607847in}}%
\pgfpathlineto{\pgfqpoint{1.661574in}{1.613287in}}%
\pgfpathlineto{\pgfqpoint{1.664093in}{1.603241in}}%
\pgfpathlineto{\pgfqpoint{1.666612in}{1.562733in}}%
\pgfpathlineto{\pgfqpoint{1.669131in}{1.578227in}}%
\pgfpathlineto{\pgfqpoint{1.671649in}{1.553309in}}%
\pgfpathlineto{\pgfqpoint{1.674168in}{1.537646in}}%
\pgfpathlineto{\pgfqpoint{1.676687in}{1.556266in}}%
\pgfpathlineto{\pgfqpoint{1.679206in}{1.532468in}}%
\pgfpathlineto{\pgfqpoint{1.681724in}{1.515877in}}%
\pgfpathlineto{\pgfqpoint{1.684243in}{1.529334in}}%
\pgfpathlineto{\pgfqpoint{1.686762in}{1.535972in}}%
\pgfpathlineto{\pgfqpoint{1.689281in}{1.570035in}}%
\pgfpathlineto{\pgfqpoint{1.691799in}{1.572672in}}%
\pgfpathlineto{\pgfqpoint{1.694318in}{1.576053in}}%
\pgfpathlineto{\pgfqpoint{1.696837in}{1.572483in}}%
\pgfpathlineto{\pgfqpoint{1.699356in}{1.569817in}}%
\pgfpathlineto{\pgfqpoint{1.701874in}{1.550203in}}%
\pgfpathlineto{\pgfqpoint{1.704393in}{1.578680in}}%
\pgfpathlineto{\pgfqpoint{1.706912in}{1.559590in}}%
\pgfpathlineto{\pgfqpoint{1.709431in}{1.577650in}}%
\pgfpathlineto{\pgfqpoint{1.711949in}{1.573823in}}%
\pgfpathlineto{\pgfqpoint{1.714468in}{1.578529in}}%
\pgfpathlineto{\pgfqpoint{1.716987in}{1.570983in}}%
\pgfpathlineto{\pgfqpoint{1.719506in}{1.576864in}}%
\pgfpathlineto{\pgfqpoint{1.722024in}{1.569414in}}%
\pgfpathlineto{\pgfqpoint{1.724543in}{1.531615in}}%
\pgfpathlineto{\pgfqpoint{1.727062in}{1.523751in}}%
\pgfpathlineto{\pgfqpoint{1.732099in}{1.525215in}}%
\pgfpathlineto{\pgfqpoint{1.734618in}{1.542625in}}%
\pgfpathlineto{\pgfqpoint{1.737137in}{1.533257in}}%
\pgfpathlineto{\pgfqpoint{1.739656in}{1.538799in}}%
\pgfpathlineto{\pgfqpoint{1.742174in}{1.541103in}}%
\pgfpathlineto{\pgfqpoint{1.744693in}{1.573553in}}%
\pgfpathlineto{\pgfqpoint{1.747212in}{1.575847in}}%
\pgfpathlineto{\pgfqpoint{1.749731in}{1.571504in}}%
\pgfpathlineto{\pgfqpoint{1.752249in}{1.556322in}}%
\pgfpathlineto{\pgfqpoint{1.754768in}{1.560621in}}%
\pgfpathlineto{\pgfqpoint{1.757287in}{1.561486in}}%
\pgfpathlineto{\pgfqpoint{1.759806in}{1.553467in}}%
\pgfpathlineto{\pgfqpoint{1.762324in}{1.559317in}}%
\pgfpathlineto{\pgfqpoint{1.764843in}{1.548448in}}%
\pgfpathlineto{\pgfqpoint{1.767362in}{1.560779in}}%
\pgfpathlineto{\pgfqpoint{1.769881in}{1.553027in}}%
\pgfpathlineto{\pgfqpoint{1.772399in}{1.526192in}}%
\pgfpathlineto{\pgfqpoint{1.774918in}{1.556945in}}%
\pgfpathlineto{\pgfqpoint{1.777437in}{1.559668in}}%
\pgfpathlineto{\pgfqpoint{1.779956in}{1.545234in}}%
\pgfpathlineto{\pgfqpoint{1.782474in}{1.558955in}}%
\pgfpathlineto{\pgfqpoint{1.784993in}{1.556848in}}%
\pgfpathlineto{\pgfqpoint{1.787512in}{1.550840in}}%
\pgfpathlineto{\pgfqpoint{1.792549in}{1.570871in}}%
\pgfpathlineto{\pgfqpoint{1.795068in}{1.586092in}}%
\pgfpathlineto{\pgfqpoint{1.797587in}{1.600225in}}%
\pgfpathlineto{\pgfqpoint{1.800106in}{1.602283in}}%
\pgfpathlineto{\pgfqpoint{1.802624in}{1.584919in}}%
\pgfpathlineto{\pgfqpoint{1.805143in}{1.553675in}}%
\pgfpathlineto{\pgfqpoint{1.807662in}{1.568517in}}%
\pgfpathlineto{\pgfqpoint{1.810181in}{1.564614in}}%
\pgfpathlineto{\pgfqpoint{1.812699in}{1.544024in}}%
\pgfpathlineto{\pgfqpoint{1.815218in}{1.516854in}}%
\pgfpathlineto{\pgfqpoint{1.817737in}{1.514071in}}%
\pgfpathlineto{\pgfqpoint{1.820256in}{1.504499in}}%
\pgfpathlineto{\pgfqpoint{1.822774in}{1.501234in}}%
\pgfpathlineto{\pgfqpoint{1.825293in}{1.478669in}}%
\pgfpathlineto{\pgfqpoint{1.827812in}{1.447067in}}%
\pgfpathlineto{\pgfqpoint{1.830331in}{1.453595in}}%
\pgfpathlineto{\pgfqpoint{1.832849in}{1.455455in}}%
\pgfpathlineto{\pgfqpoint{1.835368in}{1.453015in}}%
\pgfpathlineto{\pgfqpoint{1.837887in}{1.446028in}}%
\pgfpathlineto{\pgfqpoint{1.840406in}{1.451208in}}%
\pgfpathlineto{\pgfqpoint{1.842924in}{1.423969in}}%
\pgfpathlineto{\pgfqpoint{1.845443in}{1.418254in}}%
\pgfpathlineto{\pgfqpoint{1.847962in}{1.413119in}}%
\pgfpathlineto{\pgfqpoint{1.850481in}{1.419666in}}%
\pgfpathlineto{\pgfqpoint{1.852999in}{1.444343in}}%
\pgfpathlineto{\pgfqpoint{1.855518in}{1.429195in}}%
\pgfpathlineto{\pgfqpoint{1.858037in}{1.439927in}}%
\pgfpathlineto{\pgfqpoint{1.860556in}{1.422400in}}%
\pgfpathlineto{\pgfqpoint{1.863074in}{1.427954in}}%
\pgfpathlineto{\pgfqpoint{1.865593in}{1.434039in}}%
\pgfpathlineto{\pgfqpoint{1.868112in}{1.430064in}}%
\pgfpathlineto{\pgfqpoint{1.870631in}{1.450036in}}%
\pgfpathlineto{\pgfqpoint{1.873149in}{1.448177in}}%
\pgfpathlineto{\pgfqpoint{1.875668in}{1.446065in}}%
\pgfpathlineto{\pgfqpoint{1.878187in}{1.447760in}}%
\pgfpathlineto{\pgfqpoint{1.880706in}{1.455764in}}%
\pgfpathlineto{\pgfqpoint{1.883224in}{1.471272in}}%
\pgfpathlineto{\pgfqpoint{1.885743in}{1.458774in}}%
\pgfpathlineto{\pgfqpoint{1.888262in}{1.471203in}}%
\pgfpathlineto{\pgfqpoint{1.890781in}{1.467885in}}%
\pgfpathlineto{\pgfqpoint{1.893299in}{1.473933in}}%
\pgfpathlineto{\pgfqpoint{1.895818in}{1.468406in}}%
\pgfpathlineto{\pgfqpoint{1.898337in}{1.480607in}}%
\pgfpathlineto{\pgfqpoint{1.900856in}{1.466403in}}%
\pgfpathlineto{\pgfqpoint{1.903374in}{1.466033in}}%
\pgfpathlineto{\pgfqpoint{1.905893in}{1.471166in}}%
\pgfpathlineto{\pgfqpoint{1.908412in}{1.470432in}}%
\pgfpathlineto{\pgfqpoint{1.910931in}{1.443011in}}%
\pgfpathlineto{\pgfqpoint{1.913449in}{1.451324in}}%
\pgfpathlineto{\pgfqpoint{1.915968in}{1.463147in}}%
\pgfpathlineto{\pgfqpoint{1.918487in}{1.441344in}}%
\pgfpathlineto{\pgfqpoint{1.921006in}{1.464045in}}%
\pgfpathlineto{\pgfqpoint{1.923524in}{1.467864in}}%
\pgfpathlineto{\pgfqpoint{1.926043in}{1.457412in}}%
\pgfpathlineto{\pgfqpoint{1.928562in}{1.459152in}}%
\pgfpathlineto{\pgfqpoint{1.931081in}{1.478399in}}%
\pgfpathlineto{\pgfqpoint{1.933599in}{1.474401in}}%
\pgfpathlineto{\pgfqpoint{1.936118in}{1.475959in}}%
\pgfpathlineto{\pgfqpoint{1.938637in}{1.499190in}}%
\pgfpathlineto{\pgfqpoint{1.941156in}{1.480453in}}%
\pgfpathlineto{\pgfqpoint{1.943674in}{1.484481in}}%
\pgfpathlineto{\pgfqpoint{1.946193in}{1.486945in}}%
\pgfpathlineto{\pgfqpoint{1.948712in}{1.479410in}}%
\pgfpathlineto{\pgfqpoint{1.951231in}{1.461864in}}%
\pgfpathlineto{\pgfqpoint{1.953749in}{1.481380in}}%
\pgfpathlineto{\pgfqpoint{1.956268in}{1.491271in}}%
\pgfpathlineto{\pgfqpoint{1.958787in}{1.467351in}}%
\pgfpathlineto{\pgfqpoint{1.961306in}{1.483534in}}%
\pgfpathlineto{\pgfqpoint{1.963824in}{1.490102in}}%
\pgfpathlineto{\pgfqpoint{1.968862in}{1.457731in}}%
\pgfpathlineto{\pgfqpoint{1.971381in}{1.431206in}}%
\pgfpathlineto{\pgfqpoint{1.973899in}{1.427879in}}%
\pgfpathlineto{\pgfqpoint{1.976418in}{1.421475in}}%
\pgfpathlineto{\pgfqpoint{1.978937in}{1.419084in}}%
\pgfpathlineto{\pgfqpoint{1.981456in}{1.397830in}}%
\pgfpathlineto{\pgfqpoint{1.983974in}{1.391242in}}%
\pgfpathlineto{\pgfqpoint{1.986493in}{1.398964in}}%
\pgfpathlineto{\pgfqpoint{1.989012in}{1.397347in}}%
\pgfpathlineto{\pgfqpoint{1.991531in}{1.399565in}}%
\pgfpathlineto{\pgfqpoint{1.994049in}{1.387721in}}%
\pgfpathlineto{\pgfqpoint{1.996568in}{1.389624in}}%
\pgfpathlineto{\pgfqpoint{1.999087in}{1.392538in}}%
\pgfpathlineto{\pgfqpoint{2.001606in}{1.384803in}}%
\pgfpathlineto{\pgfqpoint{2.004124in}{1.389912in}}%
\pgfpathlineto{\pgfqpoint{2.006643in}{1.398719in}}%
\pgfpathlineto{\pgfqpoint{2.009162in}{1.386714in}}%
\pgfpathlineto{\pgfqpoint{2.011681in}{1.367773in}}%
\pgfpathlineto{\pgfqpoint{2.014199in}{1.325951in}}%
\pgfpathlineto{\pgfqpoint{2.016718in}{1.323204in}}%
\pgfpathlineto{\pgfqpoint{2.019237in}{1.332205in}}%
\pgfpathlineto{\pgfqpoint{2.021756in}{1.337842in}}%
\pgfpathlineto{\pgfqpoint{2.024274in}{1.335380in}}%
\pgfpathlineto{\pgfqpoint{2.026793in}{1.345667in}}%
\pgfpathlineto{\pgfqpoint{2.029312in}{1.352780in}}%
\pgfpathlineto{\pgfqpoint{2.031831in}{1.372689in}}%
\pgfpathlineto{\pgfqpoint{2.034349in}{1.385083in}}%
\pgfpathlineto{\pgfqpoint{2.036868in}{1.351037in}}%
\pgfpathlineto{\pgfqpoint{2.039387in}{1.346773in}}%
\pgfpathlineto{\pgfqpoint{2.041906in}{1.378164in}}%
\pgfpathlineto{\pgfqpoint{2.044424in}{1.390574in}}%
\pgfpathlineto{\pgfqpoint{2.046943in}{1.398377in}}%
\pgfpathlineto{\pgfqpoint{2.049462in}{1.413124in}}%
\pgfpathlineto{\pgfqpoint{2.051981in}{1.439182in}}%
\pgfpathlineto{\pgfqpoint{2.054499in}{1.446429in}}%
\pgfpathlineto{\pgfqpoint{2.057018in}{1.437799in}}%
\pgfpathlineto{\pgfqpoint{2.059537in}{1.437853in}}%
\pgfpathlineto{\pgfqpoint{2.062056in}{1.452350in}}%
\pgfpathlineto{\pgfqpoint{2.064574in}{1.462048in}}%
\pgfpathlineto{\pgfqpoint{2.067093in}{1.442504in}}%
\pgfpathlineto{\pgfqpoint{2.069612in}{1.410685in}}%
\pgfpathlineto{\pgfqpoint{2.072131in}{1.411634in}}%
\pgfpathlineto{\pgfqpoint{2.074649in}{1.426822in}}%
\pgfpathlineto{\pgfqpoint{2.077168in}{1.419120in}}%
\pgfpathlineto{\pgfqpoint{2.079687in}{1.422326in}}%
\pgfpathlineto{\pgfqpoint{2.082206in}{1.424402in}}%
\pgfpathlineto{\pgfqpoint{2.084724in}{1.417758in}}%
\pgfpathlineto{\pgfqpoint{2.087243in}{1.468597in}}%
\pgfpathlineto{\pgfqpoint{2.089762in}{1.454953in}}%
\pgfpathlineto{\pgfqpoint{2.092281in}{1.467224in}}%
\pgfpathlineto{\pgfqpoint{2.094799in}{1.483343in}}%
\pgfpathlineto{\pgfqpoint{2.097318in}{1.466564in}}%
\pgfpathlineto{\pgfqpoint{2.099837in}{1.461501in}}%
\pgfpathlineto{\pgfqpoint{2.102356in}{1.450195in}}%
\pgfpathlineto{\pgfqpoint{2.104874in}{1.432217in}}%
\pgfpathlineto{\pgfqpoint{2.107393in}{1.459841in}}%
\pgfpathlineto{\pgfqpoint{2.109912in}{1.438806in}}%
\pgfpathlineto{\pgfqpoint{2.114949in}{1.471696in}}%
\pgfpathlineto{\pgfqpoint{2.117468in}{1.474498in}}%
\pgfpathlineto{\pgfqpoint{2.119987in}{1.470886in}}%
\pgfpathlineto{\pgfqpoint{2.122506in}{1.473438in}}%
\pgfpathlineto{\pgfqpoint{2.125024in}{1.483088in}}%
\pgfpathlineto{\pgfqpoint{2.127543in}{1.486840in}}%
\pgfpathlineto{\pgfqpoint{2.130062in}{1.483274in}}%
\pgfpathlineto{\pgfqpoint{2.132581in}{1.472890in}}%
\pgfpathlineto{\pgfqpoint{2.135099in}{1.477905in}}%
\pgfpathlineto{\pgfqpoint{2.137618in}{1.460090in}}%
\pgfpathlineto{\pgfqpoint{2.140137in}{1.463252in}}%
\pgfpathlineto{\pgfqpoint{2.142656in}{1.485518in}}%
\pgfpathlineto{\pgfqpoint{2.145174in}{1.480379in}}%
\pgfpathlineto{\pgfqpoint{2.147693in}{1.475018in}}%
\pgfpathlineto{\pgfqpoint{2.150212in}{1.485131in}}%
\pgfpathlineto{\pgfqpoint{2.152731in}{1.492533in}}%
\pgfpathlineto{\pgfqpoint{2.155249in}{1.472476in}}%
\pgfpathlineto{\pgfqpoint{2.157768in}{1.465986in}}%
\pgfpathlineto{\pgfqpoint{2.160287in}{1.454448in}}%
\pgfpathlineto{\pgfqpoint{2.162806in}{1.494059in}}%
\pgfpathlineto{\pgfqpoint{2.165324in}{1.494750in}}%
\pgfpathlineto{\pgfqpoint{2.167843in}{1.460469in}}%
\pgfpathlineto{\pgfqpoint{2.170362in}{1.462536in}}%
\pgfpathlineto{\pgfqpoint{2.172881in}{1.459888in}}%
\pgfpathlineto{\pgfqpoint{2.175399in}{1.431151in}}%
\pgfpathlineto{\pgfqpoint{2.177918in}{1.441090in}}%
\pgfpathlineto{\pgfqpoint{2.180437in}{1.447190in}}%
\pgfpathlineto{\pgfqpoint{2.182956in}{1.439664in}}%
\pgfpathlineto{\pgfqpoint{2.185474in}{1.439033in}}%
\pgfpathlineto{\pgfqpoint{2.187993in}{1.440740in}}%
\pgfpathlineto{\pgfqpoint{2.190512in}{1.462786in}}%
\pgfpathlineto{\pgfqpoint{2.193031in}{1.456599in}}%
\pgfpathlineto{\pgfqpoint{2.195549in}{1.451091in}}%
\pgfpathlineto{\pgfqpoint{2.198068in}{1.453798in}}%
\pgfpathlineto{\pgfqpoint{2.200587in}{1.459815in}}%
\pgfpathlineto{\pgfqpoint{2.203106in}{1.479262in}}%
\pgfpathlineto{\pgfqpoint{2.205624in}{1.449650in}}%
\pgfpathlineto{\pgfqpoint{2.208143in}{1.426615in}}%
\pgfpathlineto{\pgfqpoint{2.210662in}{1.445729in}}%
\pgfpathlineto{\pgfqpoint{2.213181in}{1.440782in}}%
\pgfpathlineto{\pgfqpoint{2.215699in}{1.429081in}}%
\pgfpathlineto{\pgfqpoint{2.218218in}{1.416862in}}%
\pgfpathlineto{\pgfqpoint{2.220737in}{1.402756in}}%
\pgfpathlineto{\pgfqpoint{2.223256in}{1.418190in}}%
\pgfpathlineto{\pgfqpoint{2.225774in}{1.400403in}}%
\pgfpathlineto{\pgfqpoint{2.228293in}{1.391120in}}%
\pgfpathlineto{\pgfqpoint{2.230812in}{1.389491in}}%
\pgfpathlineto{\pgfqpoint{2.233331in}{1.401431in}}%
\pgfpathlineto{\pgfqpoint{2.235849in}{1.388843in}}%
\pgfpathlineto{\pgfqpoint{2.238368in}{1.391544in}}%
\pgfpathlineto{\pgfqpoint{2.240887in}{1.405379in}}%
\pgfpathlineto{\pgfqpoint{2.243406in}{1.413014in}}%
\pgfpathlineto{\pgfqpoint{2.245924in}{1.415255in}}%
\pgfpathlineto{\pgfqpoint{2.248443in}{1.418084in}}%
\pgfpathlineto{\pgfqpoint{2.250962in}{1.422085in}}%
\pgfpathlineto{\pgfqpoint{2.253481in}{1.408898in}}%
\pgfpathlineto{\pgfqpoint{2.255999in}{1.419823in}}%
\pgfpathlineto{\pgfqpoint{2.258518in}{1.405978in}}%
\pgfpathlineto{\pgfqpoint{2.261037in}{1.384895in}}%
\pgfpathlineto{\pgfqpoint{2.263556in}{1.400608in}}%
\pgfpathlineto{\pgfqpoint{2.266074in}{1.407670in}}%
\pgfpathlineto{\pgfqpoint{2.268593in}{1.425765in}}%
\pgfpathlineto{\pgfqpoint{2.271112in}{1.422741in}}%
\pgfpathlineto{\pgfqpoint{2.273631in}{1.416601in}}%
\pgfpathlineto{\pgfqpoint{2.276149in}{1.455598in}}%
\pgfpathlineto{\pgfqpoint{2.278668in}{1.448273in}}%
\pgfpathlineto{\pgfqpoint{2.281187in}{1.461566in}}%
\pgfpathlineto{\pgfqpoint{2.283706in}{1.469623in}}%
\pgfpathlineto{\pgfqpoint{2.286224in}{1.472237in}}%
\pgfpathlineto{\pgfqpoint{2.288743in}{1.468508in}}%
\pgfpathlineto{\pgfqpoint{2.291262in}{1.461796in}}%
\pgfpathlineto{\pgfqpoint{2.293781in}{1.454318in}}%
\pgfpathlineto{\pgfqpoint{2.296299in}{1.473100in}}%
\pgfpathlineto{\pgfqpoint{2.298818in}{1.455581in}}%
\pgfpathlineto{\pgfqpoint{2.301337in}{1.468544in}}%
\pgfpathlineto{\pgfqpoint{2.303856in}{1.468307in}}%
\pgfpathlineto{\pgfqpoint{2.306374in}{1.469348in}}%
\pgfpathlineto{\pgfqpoint{2.308893in}{1.491294in}}%
\pgfpathlineto{\pgfqpoint{2.311412in}{1.496585in}}%
\pgfpathlineto{\pgfqpoint{2.313931in}{1.485787in}}%
\pgfpathlineto{\pgfqpoint{2.316449in}{1.483345in}}%
\pgfpathlineto{\pgfqpoint{2.318968in}{1.481767in}}%
\pgfpathlineto{\pgfqpoint{2.321487in}{1.479691in}}%
\pgfpathlineto{\pgfqpoint{2.324006in}{1.478618in}}%
\pgfpathlineto{\pgfqpoint{2.326524in}{1.458889in}}%
\pgfpathlineto{\pgfqpoint{2.329043in}{1.453681in}}%
\pgfpathlineto{\pgfqpoint{2.331562in}{1.443153in}}%
\pgfpathlineto{\pgfqpoint{2.334081in}{1.449678in}}%
\pgfpathlineto{\pgfqpoint{2.336599in}{1.444340in}}%
\pgfpathlineto{\pgfqpoint{2.339118in}{1.454688in}}%
\pgfpathlineto{\pgfqpoint{2.341637in}{1.477582in}}%
\pgfpathlineto{\pgfqpoint{2.344156in}{1.480505in}}%
\pgfpathlineto{\pgfqpoint{2.346674in}{1.464185in}}%
\pgfpathlineto{\pgfqpoint{2.349193in}{1.470116in}}%
\pgfpathlineto{\pgfqpoint{2.351712in}{1.495735in}}%
\pgfpathlineto{\pgfqpoint{2.354231in}{1.498410in}}%
\pgfpathlineto{\pgfqpoint{2.356749in}{1.518387in}}%
\pgfpathlineto{\pgfqpoint{2.359268in}{1.515496in}}%
\pgfpathlineto{\pgfqpoint{2.361787in}{1.520834in}}%
\pgfpathlineto{\pgfqpoint{2.364306in}{1.524278in}}%
\pgfpathlineto{\pgfqpoint{2.366824in}{1.504822in}}%
\pgfpathlineto{\pgfqpoint{2.369343in}{1.508767in}}%
\pgfpathlineto{\pgfqpoint{2.371862in}{1.506472in}}%
\pgfpathlineto{\pgfqpoint{2.374381in}{1.524783in}}%
\pgfpathlineto{\pgfqpoint{2.376899in}{1.528784in}}%
\pgfpathlineto{\pgfqpoint{2.379418in}{1.539392in}}%
\pgfpathlineto{\pgfqpoint{2.381937in}{1.528498in}}%
\pgfpathlineto{\pgfqpoint{2.384456in}{1.534960in}}%
\pgfpathlineto{\pgfqpoint{2.386974in}{1.533323in}}%
\pgfpathlineto{\pgfqpoint{2.389493in}{1.525090in}}%
\pgfpathlineto{\pgfqpoint{2.392012in}{1.536400in}}%
\pgfpathlineto{\pgfqpoint{2.394531in}{1.528062in}}%
\pgfpathlineto{\pgfqpoint{2.397049in}{1.570164in}}%
\pgfpathlineto{\pgfqpoint{2.399568in}{1.552382in}}%
\pgfpathlineto{\pgfqpoint{2.402087in}{1.575030in}}%
\pgfpathlineto{\pgfqpoint{2.404606in}{1.565540in}}%
\pgfpathlineto{\pgfqpoint{2.407124in}{1.571497in}}%
\pgfpathlineto{\pgfqpoint{2.412162in}{1.617236in}}%
\pgfpathlineto{\pgfqpoint{2.414681in}{1.612337in}}%
\pgfpathlineto{\pgfqpoint{2.417199in}{1.601804in}}%
\pgfpathlineto{\pgfqpoint{2.419718in}{1.611292in}}%
\pgfpathlineto{\pgfqpoint{2.422237in}{1.645129in}}%
\pgfpathlineto{\pgfqpoint{2.424756in}{1.652858in}}%
\pgfpathlineto{\pgfqpoint{2.427274in}{1.656368in}}%
\pgfpathlineto{\pgfqpoint{2.429793in}{1.666715in}}%
\pgfpathlineto{\pgfqpoint{2.432312in}{1.687539in}}%
\pgfpathlineto{\pgfqpoint{2.434831in}{1.670828in}}%
\pgfpathlineto{\pgfqpoint{2.437349in}{1.682311in}}%
\pgfpathlineto{\pgfqpoint{2.439868in}{1.686334in}}%
\pgfpathlineto{\pgfqpoint{2.442387in}{1.695412in}}%
\pgfpathlineto{\pgfqpoint{2.444906in}{1.713954in}}%
\pgfpathlineto{\pgfqpoint{2.447424in}{1.711271in}}%
\pgfpathlineto{\pgfqpoint{2.449943in}{1.717307in}}%
\pgfpathlineto{\pgfqpoint{2.452462in}{1.739717in}}%
\pgfpathlineto{\pgfqpoint{2.454981in}{1.737083in}}%
\pgfpathlineto{\pgfqpoint{2.457499in}{1.739643in}}%
\pgfpathlineto{\pgfqpoint{2.460018in}{1.747351in}}%
\pgfpathlineto{\pgfqpoint{2.462537in}{1.722986in}}%
\pgfpathlineto{\pgfqpoint{2.465056in}{1.724096in}}%
\pgfpathlineto{\pgfqpoint{2.467574in}{1.716806in}}%
\pgfpathlineto{\pgfqpoint{2.470093in}{1.723595in}}%
\pgfpathlineto{\pgfqpoint{2.472612in}{1.713263in}}%
\pgfpathlineto{\pgfqpoint{2.475131in}{1.725084in}}%
\pgfpathlineto{\pgfqpoint{2.477649in}{1.729924in}}%
\pgfpathlineto{\pgfqpoint{2.480168in}{1.760695in}}%
\pgfpathlineto{\pgfqpoint{2.482687in}{1.751709in}}%
\pgfpathlineto{\pgfqpoint{2.485206in}{1.762411in}}%
\pgfpathlineto{\pgfqpoint{2.487724in}{1.734772in}}%
\pgfpathlineto{\pgfqpoint{2.490243in}{1.740505in}}%
\pgfpathlineto{\pgfqpoint{2.492762in}{1.737246in}}%
\pgfpathlineto{\pgfqpoint{2.495281in}{1.736566in}}%
\pgfpathlineto{\pgfqpoint{2.497799in}{1.741983in}}%
\pgfpathlineto{\pgfqpoint{2.500318in}{1.742452in}}%
\pgfpathlineto{\pgfqpoint{2.502837in}{1.712450in}}%
\pgfpathlineto{\pgfqpoint{2.505356in}{1.714374in}}%
\pgfpathlineto{\pgfqpoint{2.507874in}{1.688103in}}%
\pgfpathlineto{\pgfqpoint{2.510393in}{1.704295in}}%
\pgfpathlineto{\pgfqpoint{2.512912in}{1.697853in}}%
\pgfpathlineto{\pgfqpoint{2.515431in}{1.697542in}}%
\pgfpathlineto{\pgfqpoint{2.517949in}{1.726805in}}%
\pgfpathlineto{\pgfqpoint{2.520468in}{1.701255in}}%
\pgfpathlineto{\pgfqpoint{2.522987in}{1.691573in}}%
\pgfpathlineto{\pgfqpoint{2.525506in}{1.705892in}}%
\pgfpathlineto{\pgfqpoint{2.528024in}{1.769221in}}%
\pgfpathlineto{\pgfqpoint{2.530543in}{1.753523in}}%
\pgfpathlineto{\pgfqpoint{2.533062in}{1.768321in}}%
\pgfpathlineto{\pgfqpoint{2.535581in}{1.770871in}}%
\pgfpathlineto{\pgfqpoint{2.538099in}{1.736246in}}%
\pgfpathlineto{\pgfqpoint{2.540618in}{1.735938in}}%
\pgfpathlineto{\pgfqpoint{2.543137in}{1.710091in}}%
\pgfpathlineto{\pgfqpoint{2.545656in}{1.730547in}}%
\pgfpathlineto{\pgfqpoint{2.548174in}{1.724673in}}%
\pgfpathlineto{\pgfqpoint{2.550693in}{1.712790in}}%
\pgfpathlineto{\pgfqpoint{2.553212in}{1.737850in}}%
\pgfpathlineto{\pgfqpoint{2.555731in}{1.737136in}}%
\pgfpathlineto{\pgfqpoint{2.558249in}{1.747994in}}%
\pgfpathlineto{\pgfqpoint{2.560768in}{1.750938in}}%
\pgfpathlineto{\pgfqpoint{2.563287in}{1.749163in}}%
\pgfpathlineto{\pgfqpoint{2.565806in}{1.764615in}}%
\pgfpathlineto{\pgfqpoint{2.568324in}{1.767019in}}%
\pgfpathlineto{\pgfqpoint{2.570843in}{1.758012in}}%
\pgfpathlineto{\pgfqpoint{2.573362in}{1.770376in}}%
\pgfpathlineto{\pgfqpoint{2.575881in}{1.798670in}}%
\pgfpathlineto{\pgfqpoint{2.578399in}{1.791913in}}%
\pgfpathlineto{\pgfqpoint{2.580918in}{1.804141in}}%
\pgfpathlineto{\pgfqpoint{2.583437in}{1.805437in}}%
\pgfpathlineto{\pgfqpoint{2.585956in}{1.793449in}}%
\pgfpathlineto{\pgfqpoint{2.588474in}{1.788860in}}%
\pgfpathlineto{\pgfqpoint{2.590993in}{1.779708in}}%
\pgfpathlineto{\pgfqpoint{2.593512in}{1.758230in}}%
\pgfpathlineto{\pgfqpoint{2.596031in}{1.739175in}}%
\pgfpathlineto{\pgfqpoint{2.598549in}{1.726432in}}%
\pgfpathlineto{\pgfqpoint{2.601068in}{1.739876in}}%
\pgfpathlineto{\pgfqpoint{2.603587in}{1.763362in}}%
\pgfpathlineto{\pgfqpoint{2.606106in}{1.766657in}}%
\pgfpathlineto{\pgfqpoint{2.608624in}{1.774451in}}%
\pgfpathlineto{\pgfqpoint{2.611143in}{1.785415in}}%
\pgfpathlineto{\pgfqpoint{2.613662in}{1.775141in}}%
\pgfpathlineto{\pgfqpoint{2.616181in}{1.769921in}}%
\pgfpathlineto{\pgfqpoint{2.618699in}{1.760357in}}%
\pgfpathlineto{\pgfqpoint{2.621218in}{1.771823in}}%
\pgfpathlineto{\pgfqpoint{2.623737in}{1.775812in}}%
\pgfpathlineto{\pgfqpoint{2.626256in}{1.779199in}}%
\pgfpathlineto{\pgfqpoint{2.628774in}{1.740323in}}%
\pgfpathlineto{\pgfqpoint{2.631293in}{1.724122in}}%
\pgfpathlineto{\pgfqpoint{2.633812in}{1.723922in}}%
\pgfpathlineto{\pgfqpoint{2.636331in}{1.743685in}}%
\pgfpathlineto{\pgfqpoint{2.638849in}{1.744516in}}%
\pgfpathlineto{\pgfqpoint{2.641368in}{1.724092in}}%
\pgfpathlineto{\pgfqpoint{2.643887in}{1.735670in}}%
\pgfpathlineto{\pgfqpoint{2.646406in}{1.716059in}}%
\pgfpathlineto{\pgfqpoint{2.648924in}{1.703336in}}%
\pgfpathlineto{\pgfqpoint{2.651443in}{1.685722in}}%
\pgfpathlineto{\pgfqpoint{2.653962in}{1.689310in}}%
\pgfpathlineto{\pgfqpoint{2.656481in}{1.697750in}}%
\pgfpathlineto{\pgfqpoint{2.658999in}{1.717309in}}%
\pgfpathlineto{\pgfqpoint{2.661518in}{1.721002in}}%
\pgfpathlineto{\pgfqpoint{2.664037in}{1.706518in}}%
\pgfpathlineto{\pgfqpoint{2.666556in}{1.713092in}}%
\pgfpathlineto{\pgfqpoint{2.669074in}{1.700753in}}%
\pgfpathlineto{\pgfqpoint{2.671593in}{1.711950in}}%
\pgfpathlineto{\pgfqpoint{2.674112in}{1.670754in}}%
\pgfpathlineto{\pgfqpoint{2.676631in}{1.648256in}}%
\pgfpathlineto{\pgfqpoint{2.679149in}{1.645169in}}%
\pgfpathlineto{\pgfqpoint{2.681668in}{1.644878in}}%
\pgfpathlineto{\pgfqpoint{2.684187in}{1.611245in}}%
\pgfpathlineto{\pgfqpoint{2.686706in}{1.614731in}}%
\pgfpathlineto{\pgfqpoint{2.689224in}{1.623688in}}%
\pgfpathlineto{\pgfqpoint{2.691743in}{1.630613in}}%
\pgfpathlineto{\pgfqpoint{2.694262in}{1.628849in}}%
\pgfpathlineto{\pgfqpoint{2.696781in}{1.603562in}}%
\pgfpathlineto{\pgfqpoint{2.699299in}{1.618960in}}%
\pgfpathlineto{\pgfqpoint{2.701818in}{1.627252in}}%
\pgfpathlineto{\pgfqpoint{2.704337in}{1.630826in}}%
\pgfpathlineto{\pgfqpoint{2.706856in}{1.600492in}}%
\pgfpathlineto{\pgfqpoint{2.709374in}{1.599226in}}%
\pgfpathlineto{\pgfqpoint{2.711893in}{1.598749in}}%
\pgfpathlineto{\pgfqpoint{2.714412in}{1.603687in}}%
\pgfpathlineto{\pgfqpoint{2.716931in}{1.617379in}}%
\pgfpathlineto{\pgfqpoint{2.719449in}{1.645961in}}%
\pgfpathlineto{\pgfqpoint{2.721968in}{1.660084in}}%
\pgfpathlineto{\pgfqpoint{2.724487in}{1.652402in}}%
\pgfpathlineto{\pgfqpoint{2.727006in}{1.638833in}}%
\pgfpathlineto{\pgfqpoint{2.729524in}{1.657061in}}%
\pgfpathlineto{\pgfqpoint{2.732043in}{1.652085in}}%
\pgfpathlineto{\pgfqpoint{2.734562in}{1.625471in}}%
\pgfpathlineto{\pgfqpoint{2.737081in}{1.619388in}}%
\pgfpathlineto{\pgfqpoint{2.739599in}{1.576916in}}%
\pgfpathlineto{\pgfqpoint{2.742118in}{1.572247in}}%
\pgfpathlineto{\pgfqpoint{2.744637in}{1.582069in}}%
\pgfpathlineto{\pgfqpoint{2.747156in}{1.584586in}}%
\pgfpathlineto{\pgfqpoint{2.749674in}{1.551567in}}%
\pgfpathlineto{\pgfqpoint{2.752193in}{1.541538in}}%
\pgfpathlineto{\pgfqpoint{2.754712in}{1.551275in}}%
\pgfpathlineto{\pgfqpoint{2.757231in}{1.526459in}}%
\pgfpathlineto{\pgfqpoint{2.759749in}{1.533441in}}%
\pgfpathlineto{\pgfqpoint{2.762268in}{1.497319in}}%
\pgfpathlineto{\pgfqpoint{2.764787in}{1.484514in}}%
\pgfpathlineto{\pgfqpoint{2.767306in}{1.476608in}}%
\pgfpathlineto{\pgfqpoint{2.769824in}{1.483380in}}%
\pgfpathlineto{\pgfqpoint{2.772343in}{1.471777in}}%
\pgfpathlineto{\pgfqpoint{2.774862in}{1.454834in}}%
\pgfpathlineto{\pgfqpoint{2.777381in}{1.459930in}}%
\pgfpathlineto{\pgfqpoint{2.779899in}{1.447957in}}%
\pgfpathlineto{\pgfqpoint{2.782418in}{1.449251in}}%
\pgfpathlineto{\pgfqpoint{2.784937in}{1.453545in}}%
\pgfpathlineto{\pgfqpoint{2.787456in}{1.438276in}}%
\pgfpathlineto{\pgfqpoint{2.789974in}{1.426367in}}%
\pgfpathlineto{\pgfqpoint{2.792493in}{1.399253in}}%
\pgfpathlineto{\pgfqpoint{2.795012in}{1.394398in}}%
\pgfpathlineto{\pgfqpoint{2.797531in}{1.346595in}}%
\pgfpathlineto{\pgfqpoint{2.800049in}{1.309906in}}%
\pgfpathlineto{\pgfqpoint{2.802568in}{1.325275in}}%
\pgfpathlineto{\pgfqpoint{2.805087in}{1.334575in}}%
\pgfpathlineto{\pgfqpoint{2.807606in}{1.337320in}}%
\pgfpathlineto{\pgfqpoint{2.810124in}{1.318938in}}%
\pgfpathlineto{\pgfqpoint{2.812643in}{1.320468in}}%
\pgfpathlineto{\pgfqpoint{2.815162in}{1.310351in}}%
\pgfpathlineto{\pgfqpoint{2.817681in}{1.297576in}}%
\pgfpathlineto{\pgfqpoint{2.820199in}{1.290444in}}%
\pgfpathlineto{\pgfqpoint{2.822718in}{1.302834in}}%
\pgfpathlineto{\pgfqpoint{2.825237in}{1.282638in}}%
\pgfpathlineto{\pgfqpoint{2.827756in}{1.271486in}}%
\pgfpathlineto{\pgfqpoint{2.830274in}{1.277683in}}%
\pgfpathlineto{\pgfqpoint{2.832793in}{1.281572in}}%
\pgfpathlineto{\pgfqpoint{2.835312in}{1.290682in}}%
\pgfpathlineto{\pgfqpoint{2.837831in}{1.293146in}}%
\pgfpathlineto{\pgfqpoint{2.840349in}{1.283509in}}%
\pgfpathlineto{\pgfqpoint{2.842868in}{1.305826in}}%
\pgfpathlineto{\pgfqpoint{2.845387in}{1.289663in}}%
\pgfpathlineto{\pgfqpoint{2.847906in}{1.311439in}}%
\pgfpathlineto{\pgfqpoint{2.850424in}{1.302021in}}%
\pgfpathlineto{\pgfqpoint{2.852943in}{1.313842in}}%
\pgfpathlineto{\pgfqpoint{2.855462in}{1.310945in}}%
\pgfpathlineto{\pgfqpoint{2.857981in}{1.285391in}}%
\pgfpathlineto{\pgfqpoint{2.860499in}{1.280374in}}%
\pgfpathlineto{\pgfqpoint{2.863018in}{1.300357in}}%
\pgfpathlineto{\pgfqpoint{2.865537in}{1.324949in}}%
\pgfpathlineto{\pgfqpoint{2.868056in}{1.335816in}}%
\pgfpathlineto{\pgfqpoint{2.870574in}{1.328443in}}%
\pgfpathlineto{\pgfqpoint{2.873093in}{1.327587in}}%
\pgfpathlineto{\pgfqpoint{2.878131in}{1.345658in}}%
\pgfpathlineto{\pgfqpoint{2.880649in}{1.339237in}}%
\pgfpathlineto{\pgfqpoint{2.883168in}{1.354105in}}%
\pgfpathlineto{\pgfqpoint{2.885687in}{1.335801in}}%
\pgfpathlineto{\pgfqpoint{2.888206in}{1.323682in}}%
\pgfpathlineto{\pgfqpoint{2.890724in}{1.341270in}}%
\pgfpathlineto{\pgfqpoint{2.893243in}{1.339031in}}%
\pgfpathlineto{\pgfqpoint{2.895762in}{1.353346in}}%
\pgfpathlineto{\pgfqpoint{2.898281in}{1.355592in}}%
\pgfpathlineto{\pgfqpoint{2.900799in}{1.325045in}}%
\pgfpathlineto{\pgfqpoint{2.903318in}{1.317664in}}%
\pgfpathlineto{\pgfqpoint{2.905837in}{1.301723in}}%
\pgfpathlineto{\pgfqpoint{2.908356in}{1.325717in}}%
\pgfpathlineto{\pgfqpoint{2.910874in}{1.341261in}}%
\pgfpathlineto{\pgfqpoint{2.913393in}{1.331917in}}%
\pgfpathlineto{\pgfqpoint{2.915912in}{1.333435in}}%
\pgfpathlineto{\pgfqpoint{2.918431in}{1.330928in}}%
\pgfpathlineto{\pgfqpoint{2.920949in}{1.315163in}}%
\pgfpathlineto{\pgfqpoint{2.923468in}{1.313004in}}%
\pgfpathlineto{\pgfqpoint{2.925987in}{1.299101in}}%
\pgfpathlineto{\pgfqpoint{2.928506in}{1.296477in}}%
\pgfpathlineto{\pgfqpoint{2.931024in}{1.282915in}}%
\pgfpathlineto{\pgfqpoint{2.933543in}{1.283191in}}%
\pgfpathlineto{\pgfqpoint{2.936062in}{1.278891in}}%
\pgfpathlineto{\pgfqpoint{2.938581in}{1.281864in}}%
\pgfpathlineto{\pgfqpoint{2.941099in}{1.292114in}}%
\pgfpathlineto{\pgfqpoint{2.943618in}{1.291971in}}%
\pgfpathlineto{\pgfqpoint{2.946137in}{1.289201in}}%
\pgfpathlineto{\pgfqpoint{2.948656in}{1.275726in}}%
\pgfpathlineto{\pgfqpoint{2.951174in}{1.256669in}}%
\pgfpathlineto{\pgfqpoint{2.953693in}{1.268763in}}%
\pgfpathlineto{\pgfqpoint{2.956212in}{1.287540in}}%
\pgfpathlineto{\pgfqpoint{2.958731in}{1.318126in}}%
\pgfpathlineto{\pgfqpoint{2.961249in}{1.317115in}}%
\pgfpathlineto{\pgfqpoint{2.963768in}{1.295499in}}%
\pgfpathlineto{\pgfqpoint{2.966287in}{1.284195in}}%
\pgfpathlineto{\pgfqpoint{2.968806in}{1.336564in}}%
\pgfpathlineto{\pgfqpoint{2.971324in}{1.335248in}}%
\pgfpathlineto{\pgfqpoint{2.973843in}{1.338020in}}%
\pgfpathlineto{\pgfqpoint{2.976362in}{1.357570in}}%
\pgfpathlineto{\pgfqpoint{2.978881in}{1.368354in}}%
\pgfpathlineto{\pgfqpoint{2.981399in}{1.358464in}}%
\pgfpathlineto{\pgfqpoint{2.983918in}{1.403011in}}%
\pgfpathlineto{\pgfqpoint{2.986437in}{1.414519in}}%
\pgfpathlineto{\pgfqpoint{2.988956in}{1.423852in}}%
\pgfpathlineto{\pgfqpoint{2.991474in}{1.408880in}}%
\pgfpathlineto{\pgfqpoint{2.993993in}{1.383772in}}%
\pgfpathlineto{\pgfqpoint{2.996512in}{1.403854in}}%
\pgfpathlineto{\pgfqpoint{2.999031in}{1.406893in}}%
\pgfpathlineto{\pgfqpoint{3.001549in}{1.383030in}}%
\pgfpathlineto{\pgfqpoint{3.004068in}{1.385553in}}%
\pgfpathlineto{\pgfqpoint{3.006587in}{1.396868in}}%
\pgfpathlineto{\pgfqpoint{3.009106in}{1.419433in}}%
\pgfpathlineto{\pgfqpoint{3.011624in}{1.419626in}}%
\pgfpathlineto{\pgfqpoint{3.014143in}{1.436770in}}%
\pgfpathlineto{\pgfqpoint{3.016662in}{1.456528in}}%
\pgfpathlineto{\pgfqpoint{3.019181in}{1.480195in}}%
\pgfpathlineto{\pgfqpoint{3.021699in}{1.474755in}}%
\pgfpathlineto{\pgfqpoint{3.024218in}{1.471472in}}%
\pgfpathlineto{\pgfqpoint{3.026737in}{1.458996in}}%
\pgfpathlineto{\pgfqpoint{3.029256in}{1.484063in}}%
\pgfpathlineto{\pgfqpoint{3.031774in}{1.486073in}}%
\pgfpathlineto{\pgfqpoint{3.034293in}{1.512981in}}%
\pgfpathlineto{\pgfqpoint{3.036812in}{1.517497in}}%
\pgfpathlineto{\pgfqpoint{3.039331in}{1.519429in}}%
\pgfpathlineto{\pgfqpoint{3.041849in}{1.492967in}}%
\pgfpathlineto{\pgfqpoint{3.044368in}{1.502240in}}%
\pgfpathlineto{\pgfqpoint{3.046887in}{1.504290in}}%
\pgfpathlineto{\pgfqpoint{3.049406in}{1.530997in}}%
\pgfpathlineto{\pgfqpoint{3.051924in}{1.520333in}}%
\pgfpathlineto{\pgfqpoint{3.054443in}{1.518909in}}%
\pgfpathlineto{\pgfqpoint{3.056962in}{1.502939in}}%
\pgfpathlineto{\pgfqpoint{3.059481in}{1.500853in}}%
\pgfpathlineto{\pgfqpoint{3.061999in}{1.499365in}}%
\pgfpathlineto{\pgfqpoint{3.064518in}{1.504182in}}%
\pgfpathlineto{\pgfqpoint{3.067037in}{1.512513in}}%
\pgfpathlineto{\pgfqpoint{3.069556in}{1.493955in}}%
\pgfpathlineto{\pgfqpoint{3.072074in}{1.495515in}}%
\pgfpathlineto{\pgfqpoint{3.074593in}{1.503554in}}%
\pgfpathlineto{\pgfqpoint{3.077112in}{1.521986in}}%
\pgfpathlineto{\pgfqpoint{3.079631in}{1.524191in}}%
\pgfpathlineto{\pgfqpoint{3.082149in}{1.515661in}}%
\pgfpathlineto{\pgfqpoint{3.084668in}{1.541321in}}%
\pgfpathlineto{\pgfqpoint{3.087187in}{1.529961in}}%
\pgfpathlineto{\pgfqpoint{3.089706in}{1.509205in}}%
\pgfpathlineto{\pgfqpoint{3.092224in}{1.510775in}}%
\pgfpathlineto{\pgfqpoint{3.094743in}{1.501398in}}%
\pgfpathlineto{\pgfqpoint{3.097262in}{1.498275in}}%
\pgfpathlineto{\pgfqpoint{3.099781in}{1.487080in}}%
\pgfpathlineto{\pgfqpoint{3.102299in}{1.494272in}}%
\pgfpathlineto{\pgfqpoint{3.104818in}{1.477308in}}%
\pgfpathlineto{\pgfqpoint{3.107337in}{1.481070in}}%
\pgfpathlineto{\pgfqpoint{3.109856in}{1.486278in}}%
\pgfpathlineto{\pgfqpoint{3.112374in}{1.478020in}}%
\pgfpathlineto{\pgfqpoint{3.114893in}{1.442648in}}%
\pgfpathlineto{\pgfqpoint{3.117412in}{1.461872in}}%
\pgfpathlineto{\pgfqpoint{3.119931in}{1.472212in}}%
\pgfpathlineto{\pgfqpoint{3.122449in}{1.481558in}}%
\pgfpathlineto{\pgfqpoint{3.124968in}{1.495361in}}%
\pgfpathlineto{\pgfqpoint{3.127487in}{1.510597in}}%
\pgfpathlineto{\pgfqpoint{3.130006in}{1.528068in}}%
\pgfpathlineto{\pgfqpoint{3.132524in}{1.523271in}}%
\pgfpathlineto{\pgfqpoint{3.135043in}{1.522430in}}%
\pgfpathlineto{\pgfqpoint{3.137562in}{1.523504in}}%
\pgfpathlineto{\pgfqpoint{3.140081in}{1.547714in}}%
\pgfpathlineto{\pgfqpoint{3.142599in}{1.548131in}}%
\pgfpathlineto{\pgfqpoint{3.145118in}{1.521140in}}%
\pgfpathlineto{\pgfqpoint{3.147637in}{1.536502in}}%
\pgfpathlineto{\pgfqpoint{3.150156in}{1.561904in}}%
\pgfpathlineto{\pgfqpoint{3.152674in}{1.551827in}}%
\pgfpathlineto{\pgfqpoint{3.155193in}{1.536624in}}%
\pgfpathlineto{\pgfqpoint{3.157712in}{1.557632in}}%
\pgfpathlineto{\pgfqpoint{3.160231in}{1.570388in}}%
\pgfpathlineto{\pgfqpoint{3.162749in}{1.565763in}}%
\pgfpathlineto{\pgfqpoint{3.165268in}{1.566261in}}%
\pgfpathlineto{\pgfqpoint{3.167787in}{1.568595in}}%
\pgfpathlineto{\pgfqpoint{3.170306in}{1.578235in}}%
\pgfpathlineto{\pgfqpoint{3.172824in}{1.577331in}}%
\pgfpathlineto{\pgfqpoint{3.175343in}{1.552189in}}%
\pgfpathlineto{\pgfqpoint{3.177862in}{1.558107in}}%
\pgfpathlineto{\pgfqpoint{3.180381in}{1.576771in}}%
\pgfpathlineto{\pgfqpoint{3.182899in}{1.577543in}}%
\pgfpathlineto{\pgfqpoint{3.185418in}{1.581980in}}%
\pgfpathlineto{\pgfqpoint{3.187937in}{1.578964in}}%
\pgfpathlineto{\pgfqpoint{3.190456in}{1.583329in}}%
\pgfpathlineto{\pgfqpoint{3.192974in}{1.578111in}}%
\pgfpathlineto{\pgfqpoint{3.195493in}{1.584529in}}%
\pgfpathlineto{\pgfqpoint{3.198012in}{1.616099in}}%
\pgfpathlineto{\pgfqpoint{3.200531in}{1.589809in}}%
\pgfpathlineto{\pgfqpoint{3.203049in}{1.603974in}}%
\pgfpathlineto{\pgfqpoint{3.205568in}{1.608406in}}%
\pgfpathlineto{\pgfqpoint{3.208087in}{1.609931in}}%
\pgfpathlineto{\pgfqpoint{3.210606in}{1.609420in}}%
\pgfpathlineto{\pgfqpoint{3.213124in}{1.602769in}}%
\pgfpathlineto{\pgfqpoint{3.215643in}{1.621355in}}%
\pgfpathlineto{\pgfqpoint{3.218162in}{1.627138in}}%
\pgfpathlineto{\pgfqpoint{3.220681in}{1.619328in}}%
\pgfpathlineto{\pgfqpoint{3.223199in}{1.635730in}}%
\pgfpathlineto{\pgfqpoint{3.225718in}{1.648968in}}%
\pgfpathlineto{\pgfqpoint{3.228237in}{1.657297in}}%
\pgfpathlineto{\pgfqpoint{3.230756in}{1.656953in}}%
\pgfpathlineto{\pgfqpoint{3.233274in}{1.654292in}}%
\pgfpathlineto{\pgfqpoint{3.235793in}{1.647655in}}%
\pgfpathlineto{\pgfqpoint{3.238312in}{1.650689in}}%
\pgfpathlineto{\pgfqpoint{3.240831in}{1.637427in}}%
\pgfpathlineto{\pgfqpoint{3.243349in}{1.629997in}}%
\pgfpathlineto{\pgfqpoint{3.245868in}{1.623236in}}%
\pgfpathlineto{\pgfqpoint{3.248387in}{1.619834in}}%
\pgfpathlineto{\pgfqpoint{3.250906in}{1.590217in}}%
\pgfpathlineto{\pgfqpoint{3.253424in}{1.572477in}}%
\pgfpathlineto{\pgfqpoint{3.255943in}{1.570832in}}%
\pgfpathlineto{\pgfqpoint{3.258462in}{1.586439in}}%
\pgfpathlineto{\pgfqpoint{3.260981in}{1.596030in}}%
\pgfpathlineto{\pgfqpoint{3.263499in}{1.577838in}}%
\pgfpathlineto{\pgfqpoint{3.266018in}{1.577013in}}%
\pgfpathlineto{\pgfqpoint{3.268537in}{1.569739in}}%
\pgfpathlineto{\pgfqpoint{3.271056in}{1.561383in}}%
\pgfpathlineto{\pgfqpoint{3.273574in}{1.564666in}}%
\pgfpathlineto{\pgfqpoint{3.276093in}{1.547319in}}%
\pgfpathlineto{\pgfqpoint{3.278612in}{1.557333in}}%
\pgfpathlineto{\pgfqpoint{3.281131in}{1.555065in}}%
\pgfpathlineto{\pgfqpoint{3.283649in}{1.562565in}}%
\pgfpathlineto{\pgfqpoint{3.286168in}{1.560881in}}%
\pgfpathlineto{\pgfqpoint{3.288687in}{1.565657in}}%
\pgfpathlineto{\pgfqpoint{3.291206in}{1.580933in}}%
\pgfpathlineto{\pgfqpoint{3.293724in}{1.603488in}}%
\pgfpathlineto{\pgfqpoint{3.296243in}{1.605355in}}%
\pgfpathlineto{\pgfqpoint{3.298762in}{1.609104in}}%
\pgfpathlineto{\pgfqpoint{3.301281in}{1.586612in}}%
\pgfpathlineto{\pgfqpoint{3.303799in}{1.579736in}}%
\pgfpathlineto{\pgfqpoint{3.306318in}{1.595361in}}%
\pgfpathlineto{\pgfqpoint{3.308837in}{1.574574in}}%
\pgfpathlineto{\pgfqpoint{3.311356in}{1.583954in}}%
\pgfpathlineto{\pgfqpoint{3.313874in}{1.569916in}}%
\pgfpathlineto{\pgfqpoint{3.316393in}{1.554284in}}%
\pgfpathlineto{\pgfqpoint{3.318912in}{1.563664in}}%
\pgfpathlineto{\pgfqpoint{3.321431in}{1.539080in}}%
\pgfpathlineto{\pgfqpoint{3.323949in}{1.550597in}}%
\pgfpathlineto{\pgfqpoint{3.326468in}{1.584636in}}%
\pgfpathlineto{\pgfqpoint{3.328987in}{1.595402in}}%
\pgfpathlineto{\pgfqpoint{3.331506in}{1.597876in}}%
\pgfpathlineto{\pgfqpoint{3.334024in}{1.585718in}}%
\pgfpathlineto{\pgfqpoint{3.336543in}{1.581666in}}%
\pgfpathlineto{\pgfqpoint{3.339062in}{1.532047in}}%
\pgfpathlineto{\pgfqpoint{3.341581in}{1.513365in}}%
\pgfpathlineto{\pgfqpoint{3.344099in}{1.504469in}}%
\pgfpathlineto{\pgfqpoint{3.346618in}{1.506987in}}%
\pgfpathlineto{\pgfqpoint{3.349137in}{1.508189in}}%
\pgfpathlineto{\pgfqpoint{3.351656in}{1.468547in}}%
\pgfpathlineto{\pgfqpoint{3.354174in}{1.474725in}}%
\pgfpathlineto{\pgfqpoint{3.356693in}{1.468829in}}%
\pgfpathlineto{\pgfqpoint{3.359212in}{1.469489in}}%
\pgfpathlineto{\pgfqpoint{3.364249in}{1.527791in}}%
\pgfpathlineto{\pgfqpoint{3.366768in}{1.498759in}}%
\pgfpathlineto{\pgfqpoint{3.369287in}{1.535336in}}%
\pgfpathlineto{\pgfqpoint{3.371806in}{1.522613in}}%
\pgfpathlineto{\pgfqpoint{3.374324in}{1.516337in}}%
\pgfpathlineto{\pgfqpoint{3.376843in}{1.524054in}}%
\pgfpathlineto{\pgfqpoint{3.379362in}{1.517060in}}%
\pgfpathlineto{\pgfqpoint{3.381881in}{1.535820in}}%
\pgfpathlineto{\pgfqpoint{3.384399in}{1.555589in}}%
\pgfpathlineto{\pgfqpoint{3.386918in}{1.550188in}}%
\pgfpathlineto{\pgfqpoint{3.389437in}{1.566843in}}%
\pgfpathlineto{\pgfqpoint{3.391956in}{1.577120in}}%
\pgfpathlineto{\pgfqpoint{3.394474in}{1.574884in}}%
\pgfpathlineto{\pgfqpoint{3.396993in}{1.565324in}}%
\pgfpathlineto{\pgfqpoint{3.399512in}{1.565895in}}%
\pgfpathlineto{\pgfqpoint{3.402031in}{1.561921in}}%
\pgfpathlineto{\pgfqpoint{3.404549in}{1.558158in}}%
\pgfpathlineto{\pgfqpoint{3.407068in}{1.564725in}}%
\pgfpathlineto{\pgfqpoint{3.412106in}{1.590784in}}%
\pgfpathlineto{\pgfqpoint{3.414624in}{1.598334in}}%
\pgfpathlineto{\pgfqpoint{3.417143in}{1.591853in}}%
\pgfpathlineto{\pgfqpoint{3.419662in}{1.559584in}}%
\pgfpathlineto{\pgfqpoint{3.422181in}{1.559121in}}%
\pgfpathlineto{\pgfqpoint{3.424699in}{1.584314in}}%
\pgfpathlineto{\pgfqpoint{3.429737in}{1.638157in}}%
\pgfpathlineto{\pgfqpoint{3.432256in}{1.671280in}}%
\pgfpathlineto{\pgfqpoint{3.434774in}{1.669090in}}%
\pgfpathlineto{\pgfqpoint{3.437293in}{1.697699in}}%
\pgfpathlineto{\pgfqpoint{3.439812in}{1.688497in}}%
\pgfpathlineto{\pgfqpoint{3.442331in}{1.688037in}}%
\pgfpathlineto{\pgfqpoint{3.444849in}{1.700619in}}%
\pgfpathlineto{\pgfqpoint{3.447368in}{1.693857in}}%
\pgfpathlineto{\pgfqpoint{3.449887in}{1.706524in}}%
\pgfpathlineto{\pgfqpoint{3.452406in}{1.706968in}}%
\pgfpathlineto{\pgfqpoint{3.454924in}{1.686193in}}%
\pgfpathlineto{\pgfqpoint{3.457443in}{1.678635in}}%
\pgfpathlineto{\pgfqpoint{3.459962in}{1.659527in}}%
\pgfpathlineto{\pgfqpoint{3.462481in}{1.649446in}}%
\pgfpathlineto{\pgfqpoint{3.464999in}{1.652922in}}%
\pgfpathlineto{\pgfqpoint{3.467518in}{1.656798in}}%
\pgfpathlineto{\pgfqpoint{3.470037in}{1.694305in}}%
\pgfpathlineto{\pgfqpoint{3.472556in}{1.676339in}}%
\pgfpathlineto{\pgfqpoint{3.475074in}{1.660034in}}%
\pgfpathlineto{\pgfqpoint{3.477593in}{1.677821in}}%
\pgfpathlineto{\pgfqpoint{3.480112in}{1.697367in}}%
\pgfpathlineto{\pgfqpoint{3.482631in}{1.718660in}}%
\pgfpathlineto{\pgfqpoint{3.485149in}{1.708670in}}%
\pgfpathlineto{\pgfqpoint{3.487668in}{1.722295in}}%
\pgfpathlineto{\pgfqpoint{3.490187in}{1.721777in}}%
\pgfpathlineto{\pgfqpoint{3.492706in}{1.730847in}}%
\pgfpathlineto{\pgfqpoint{3.495224in}{1.756239in}}%
\pgfpathlineto{\pgfqpoint{3.497743in}{1.764964in}}%
\pgfpathlineto{\pgfqpoint{3.500262in}{1.784814in}}%
\pgfpathlineto{\pgfqpoint{3.502781in}{1.799118in}}%
\pgfpathlineto{\pgfqpoint{3.505299in}{1.809403in}}%
\pgfpathlineto{\pgfqpoint{3.507818in}{1.823803in}}%
\pgfpathlineto{\pgfqpoint{3.510337in}{1.807535in}}%
\pgfpathlineto{\pgfqpoint{3.512856in}{1.812379in}}%
\pgfpathlineto{\pgfqpoint{3.515374in}{1.793144in}}%
\pgfpathlineto{\pgfqpoint{3.517893in}{1.772935in}}%
\pgfpathlineto{\pgfqpoint{3.520412in}{1.768695in}}%
\pgfpathlineto{\pgfqpoint{3.522931in}{1.794689in}}%
\pgfpathlineto{\pgfqpoint{3.525449in}{1.804664in}}%
\pgfpathlineto{\pgfqpoint{3.527968in}{1.794213in}}%
\pgfpathlineto{\pgfqpoint{3.530487in}{1.814139in}}%
\pgfpathlineto{\pgfqpoint{3.533006in}{1.819477in}}%
\pgfpathlineto{\pgfqpoint{3.535524in}{1.832296in}}%
\pgfpathlineto{\pgfqpoint{3.538043in}{1.813496in}}%
\pgfpathlineto{\pgfqpoint{3.540562in}{1.777245in}}%
\pgfpathlineto{\pgfqpoint{3.543081in}{1.775796in}}%
\pgfpathlineto{\pgfqpoint{3.545599in}{1.768554in}}%
\pgfpathlineto{\pgfqpoint{3.548118in}{1.778233in}}%
\pgfpathlineto{\pgfqpoint{3.550637in}{1.756938in}}%
\pgfpathlineto{\pgfqpoint{3.553156in}{1.751290in}}%
\pgfpathlineto{\pgfqpoint{3.555674in}{1.751369in}}%
\pgfpathlineto{\pgfqpoint{3.558193in}{1.753632in}}%
\pgfpathlineto{\pgfqpoint{3.560712in}{1.766681in}}%
\pgfpathlineto{\pgfqpoint{3.563231in}{1.760348in}}%
\pgfpathlineto{\pgfqpoint{3.565749in}{1.739757in}}%
\pgfpathlineto{\pgfqpoint{3.568268in}{1.765975in}}%
\pgfpathlineto{\pgfqpoint{3.570787in}{1.767623in}}%
\pgfpathlineto{\pgfqpoint{3.573306in}{1.754951in}}%
\pgfpathlineto{\pgfqpoint{3.575824in}{1.758034in}}%
\pgfpathlineto{\pgfqpoint{3.578343in}{1.773664in}}%
\pgfpathlineto{\pgfqpoint{3.580862in}{1.765787in}}%
\pgfpathlineto{\pgfqpoint{3.583381in}{1.772352in}}%
\pgfpathlineto{\pgfqpoint{3.585899in}{1.770477in}}%
\pgfpathlineto{\pgfqpoint{3.588418in}{1.776133in}}%
\pgfpathlineto{\pgfqpoint{3.590937in}{1.789260in}}%
\pgfpathlineto{\pgfqpoint{3.593456in}{1.811038in}}%
\pgfpathlineto{\pgfqpoint{3.595974in}{1.835697in}}%
\pgfpathlineto{\pgfqpoint{3.598493in}{1.850562in}}%
\pgfpathlineto{\pgfqpoint{3.601012in}{1.839028in}}%
\pgfpathlineto{\pgfqpoint{3.603531in}{1.821493in}}%
\pgfpathlineto{\pgfqpoint{3.606049in}{1.824388in}}%
\pgfpathlineto{\pgfqpoint{3.608568in}{1.816140in}}%
\pgfpathlineto{\pgfqpoint{3.611087in}{1.806158in}}%
\pgfpathlineto{\pgfqpoint{3.613606in}{1.802213in}}%
\pgfpathlineto{\pgfqpoint{3.616124in}{1.782955in}}%
\pgfpathlineto{\pgfqpoint{3.618643in}{1.773647in}}%
\pgfpathlineto{\pgfqpoint{3.621162in}{1.779568in}}%
\pgfpathlineto{\pgfqpoint{3.623681in}{1.747179in}}%
\pgfpathlineto{\pgfqpoint{3.626199in}{1.741611in}}%
\pgfpathlineto{\pgfqpoint{3.628718in}{1.738110in}}%
\pgfpathlineto{\pgfqpoint{3.631237in}{1.756895in}}%
\pgfpathlineto{\pgfqpoint{3.633756in}{1.748558in}}%
\pgfpathlineto{\pgfqpoint{3.636274in}{1.790912in}}%
\pgfpathlineto{\pgfqpoint{3.638793in}{1.778032in}}%
\pgfpathlineto{\pgfqpoint{3.641312in}{1.783749in}}%
\pgfpathlineto{\pgfqpoint{3.643831in}{1.787539in}}%
\pgfpathlineto{\pgfqpoint{3.646349in}{1.784298in}}%
\pgfpathlineto{\pgfqpoint{3.648868in}{1.777635in}}%
\pgfpathlineto{\pgfqpoint{3.651387in}{1.790566in}}%
\pgfpathlineto{\pgfqpoint{3.653906in}{1.791653in}}%
\pgfpathlineto{\pgfqpoint{3.656424in}{1.769759in}}%
\pgfpathlineto{\pgfqpoint{3.658943in}{1.787059in}}%
\pgfpathlineto{\pgfqpoint{3.661462in}{1.791214in}}%
\pgfpathlineto{\pgfqpoint{3.663981in}{1.803407in}}%
\pgfpathlineto{\pgfqpoint{3.666499in}{1.798389in}}%
\pgfpathlineto{\pgfqpoint{3.669018in}{1.781684in}}%
\pgfpathlineto{\pgfqpoint{3.671537in}{1.806899in}}%
\pgfpathlineto{\pgfqpoint{3.674056in}{1.813497in}}%
\pgfpathlineto{\pgfqpoint{3.676574in}{1.833285in}}%
\pgfpathlineto{\pgfqpoint{3.679093in}{1.833938in}}%
\pgfpathlineto{\pgfqpoint{3.681612in}{1.821270in}}%
\pgfpathlineto{\pgfqpoint{3.684131in}{1.826877in}}%
\pgfpathlineto{\pgfqpoint{3.686649in}{1.812126in}}%
\pgfpathlineto{\pgfqpoint{3.689168in}{1.790077in}}%
\pgfpathlineto{\pgfqpoint{3.691687in}{1.787653in}}%
\pgfpathlineto{\pgfqpoint{3.694206in}{1.772832in}}%
\pgfpathlineto{\pgfqpoint{3.696724in}{1.762831in}}%
\pgfpathlineto{\pgfqpoint{3.699243in}{1.775528in}}%
\pgfpathlineto{\pgfqpoint{3.701762in}{1.767623in}}%
\pgfpathlineto{\pgfqpoint{3.704281in}{1.766193in}}%
\pgfpathlineto{\pgfqpoint{3.706799in}{1.789996in}}%
\pgfpathlineto{\pgfqpoint{3.709318in}{1.811999in}}%
\pgfpathlineto{\pgfqpoint{3.711837in}{1.797170in}}%
\pgfpathlineto{\pgfqpoint{3.714356in}{1.811981in}}%
\pgfpathlineto{\pgfqpoint{3.716874in}{1.797215in}}%
\pgfpathlineto{\pgfqpoint{3.719393in}{1.779195in}}%
\pgfpathlineto{\pgfqpoint{3.721912in}{1.808072in}}%
\pgfpathlineto{\pgfqpoint{3.724431in}{1.802492in}}%
\pgfpathlineto{\pgfqpoint{3.726949in}{1.819117in}}%
\pgfpathlineto{\pgfqpoint{3.729468in}{1.820408in}}%
\pgfpathlineto{\pgfqpoint{3.731987in}{1.826817in}}%
\pgfpathlineto{\pgfqpoint{3.734506in}{1.828144in}}%
\pgfpathlineto{\pgfqpoint{3.737024in}{1.867453in}}%
\pgfpathlineto{\pgfqpoint{3.739543in}{1.845736in}}%
\pgfpathlineto{\pgfqpoint{3.742062in}{1.843729in}}%
\pgfpathlineto{\pgfqpoint{3.744581in}{1.851382in}}%
\pgfpathlineto{\pgfqpoint{3.747099in}{1.829060in}}%
\pgfpathlineto{\pgfqpoint{3.749618in}{1.828470in}}%
\pgfpathlineto{\pgfqpoint{3.752137in}{1.830685in}}%
\pgfpathlineto{\pgfqpoint{3.754656in}{1.857676in}}%
\pgfpathlineto{\pgfqpoint{3.757174in}{1.857090in}}%
\pgfpathlineto{\pgfqpoint{3.759693in}{1.822417in}}%
\pgfpathlineto{\pgfqpoint{3.762212in}{1.838329in}}%
\pgfpathlineto{\pgfqpoint{3.764731in}{1.845052in}}%
\pgfpathlineto{\pgfqpoint{3.767249in}{1.832224in}}%
\pgfpathlineto{\pgfqpoint{3.769768in}{1.830917in}}%
\pgfpathlineto{\pgfqpoint{3.772287in}{1.837309in}}%
\pgfpathlineto{\pgfqpoint{3.774806in}{1.831514in}}%
\pgfpathlineto{\pgfqpoint{3.777324in}{1.833489in}}%
\pgfpathlineto{\pgfqpoint{3.779843in}{1.803921in}}%
\pgfpathlineto{\pgfqpoint{3.782362in}{1.795503in}}%
\pgfpathlineto{\pgfqpoint{3.784881in}{1.794429in}}%
\pgfpathlineto{\pgfqpoint{3.787399in}{1.797699in}}%
\pgfpathlineto{\pgfqpoint{3.789918in}{1.815230in}}%
\pgfpathlineto{\pgfqpoint{3.792437in}{1.804572in}}%
\pgfpathlineto{\pgfqpoint{3.794956in}{1.799851in}}%
\pgfpathlineto{\pgfqpoint{3.797474in}{1.791209in}}%
\pgfpathlineto{\pgfqpoint{3.799993in}{1.805784in}}%
\pgfpathlineto{\pgfqpoint{3.802512in}{1.799615in}}%
\pgfpathlineto{\pgfqpoint{3.805031in}{1.807080in}}%
\pgfpathlineto{\pgfqpoint{3.807549in}{1.795588in}}%
\pgfpathlineto{\pgfqpoint{3.810068in}{1.787128in}}%
\pgfpathlineto{\pgfqpoint{3.812587in}{1.783936in}}%
\pgfpathlineto{\pgfqpoint{3.815106in}{1.791656in}}%
\pgfpathlineto{\pgfqpoint{3.817624in}{1.790531in}}%
\pgfpathlineto{\pgfqpoint{3.820143in}{1.771331in}}%
\pgfpathlineto{\pgfqpoint{3.822662in}{1.758255in}}%
\pgfpathlineto{\pgfqpoint{3.825181in}{1.748123in}}%
\pgfpathlineto{\pgfqpoint{3.830218in}{1.729771in}}%
\pgfpathlineto{\pgfqpoint{3.832737in}{1.745097in}}%
\pgfpathlineto{\pgfqpoint{3.835256in}{1.756912in}}%
\pgfpathlineto{\pgfqpoint{3.837774in}{1.771434in}}%
\pgfpathlineto{\pgfqpoint{3.840293in}{1.761904in}}%
\pgfpathlineto{\pgfqpoint{3.842812in}{1.764678in}}%
\pgfpathlineto{\pgfqpoint{3.845331in}{1.763383in}}%
\pgfpathlineto{\pgfqpoint{3.847849in}{1.755679in}}%
\pgfpathlineto{\pgfqpoint{3.850368in}{1.758580in}}%
\pgfpathlineto{\pgfqpoint{3.852887in}{1.758813in}}%
\pgfpathlineto{\pgfqpoint{3.855406in}{1.746987in}}%
\pgfpathlineto{\pgfqpoint{3.857924in}{1.758370in}}%
\pgfpathlineto{\pgfqpoint{3.860443in}{1.767844in}}%
\pgfpathlineto{\pgfqpoint{3.862962in}{1.777932in}}%
\pgfpathlineto{\pgfqpoint{3.865481in}{1.759703in}}%
\pgfpathlineto{\pgfqpoint{3.867999in}{1.759343in}}%
\pgfpathlineto{\pgfqpoint{3.870518in}{1.745041in}}%
\pgfpathlineto{\pgfqpoint{3.873037in}{1.749015in}}%
\pgfpathlineto{\pgfqpoint{3.875556in}{1.749777in}}%
\pgfpathlineto{\pgfqpoint{3.878074in}{1.772350in}}%
\pgfpathlineto{\pgfqpoint{3.880593in}{1.796388in}}%
\pgfpathlineto{\pgfqpoint{3.883112in}{1.809969in}}%
\pgfpathlineto{\pgfqpoint{3.885631in}{1.848182in}}%
\pgfpathlineto{\pgfqpoint{3.888149in}{1.857639in}}%
\pgfpathlineto{\pgfqpoint{3.890668in}{1.833618in}}%
\pgfpathlineto{\pgfqpoint{3.893187in}{1.817184in}}%
\pgfpathlineto{\pgfqpoint{3.895706in}{1.808141in}}%
\pgfpathlineto{\pgfqpoint{3.898224in}{1.801383in}}%
\pgfpathlineto{\pgfqpoint{3.900743in}{1.797001in}}%
\pgfpathlineto{\pgfqpoint{3.903262in}{1.796785in}}%
\pgfpathlineto{\pgfqpoint{3.905781in}{1.805855in}}%
\pgfpathlineto{\pgfqpoint{3.908299in}{1.828854in}}%
\pgfpathlineto{\pgfqpoint{3.910818in}{1.834373in}}%
\pgfpathlineto{\pgfqpoint{3.913337in}{1.836674in}}%
\pgfpathlineto{\pgfqpoint{3.915856in}{1.831989in}}%
\pgfpathlineto{\pgfqpoint{3.918374in}{1.850222in}}%
\pgfpathlineto{\pgfqpoint{3.920893in}{1.814685in}}%
\pgfpathlineto{\pgfqpoint{3.923412in}{1.809282in}}%
\pgfpathlineto{\pgfqpoint{3.925931in}{1.795192in}}%
\pgfpathlineto{\pgfqpoint{3.928449in}{1.779963in}}%
\pgfpathlineto{\pgfqpoint{3.930968in}{1.797837in}}%
\pgfpathlineto{\pgfqpoint{3.933487in}{1.804048in}}%
\pgfpathlineto{\pgfqpoint{3.936006in}{1.799438in}}%
\pgfpathlineto{\pgfqpoint{3.938524in}{1.802217in}}%
\pgfpathlineto{\pgfqpoint{3.941043in}{1.815556in}}%
\pgfpathlineto{\pgfqpoint{3.943562in}{1.800980in}}%
\pgfpathlineto{\pgfqpoint{3.946081in}{1.768468in}}%
\pgfpathlineto{\pgfqpoint{3.948599in}{1.770463in}}%
\pgfpathlineto{\pgfqpoint{3.951118in}{1.769250in}}%
\pgfpathlineto{\pgfqpoint{3.953637in}{1.731964in}}%
\pgfpathlineto{\pgfqpoint{3.956156in}{1.741630in}}%
\pgfpathlineto{\pgfqpoint{3.958674in}{1.758371in}}%
\pgfpathlineto{\pgfqpoint{3.961193in}{1.758003in}}%
\pgfpathlineto{\pgfqpoint{3.963712in}{1.732869in}}%
\pgfpathlineto{\pgfqpoint{3.966231in}{1.733896in}}%
\pgfpathlineto{\pgfqpoint{3.968749in}{1.757622in}}%
\pgfpathlineto{\pgfqpoint{3.971268in}{1.722476in}}%
\pgfpathlineto{\pgfqpoint{3.973787in}{1.723460in}}%
\pgfpathlineto{\pgfqpoint{3.976306in}{1.720946in}}%
\pgfpathlineto{\pgfqpoint{3.978824in}{1.689280in}}%
\pgfpathlineto{\pgfqpoint{3.981343in}{1.690114in}}%
\pgfpathlineto{\pgfqpoint{3.983862in}{1.700076in}}%
\pgfpathlineto{\pgfqpoint{3.986381in}{1.716501in}}%
\pgfpathlineto{\pgfqpoint{3.988899in}{1.736128in}}%
\pgfpathlineto{\pgfqpoint{3.991418in}{1.759224in}}%
\pgfpathlineto{\pgfqpoint{3.993937in}{1.761832in}}%
\pgfpathlineto{\pgfqpoint{3.996456in}{1.769015in}}%
\pgfpathlineto{\pgfqpoint{3.998974in}{1.764246in}}%
\pgfpathlineto{\pgfqpoint{4.001493in}{1.784393in}}%
\pgfpathlineto{\pgfqpoint{4.004012in}{1.796610in}}%
\pgfpathlineto{\pgfqpoint{4.006531in}{1.791751in}}%
\pgfpathlineto{\pgfqpoint{4.009049in}{1.778144in}}%
\pgfpathlineto{\pgfqpoint{4.011568in}{1.747715in}}%
\pgfpathlineto{\pgfqpoint{4.014087in}{1.752698in}}%
\pgfpathlineto{\pgfqpoint{4.016606in}{1.767276in}}%
\pgfpathlineto{\pgfqpoint{4.019124in}{1.751444in}}%
\pgfpathlineto{\pgfqpoint{4.021643in}{1.760567in}}%
\pgfpathlineto{\pgfqpoint{4.024162in}{1.788907in}}%
\pgfpathlineto{\pgfqpoint{4.026681in}{1.815571in}}%
\pgfpathlineto{\pgfqpoint{4.029199in}{1.823746in}}%
\pgfpathlineto{\pgfqpoint{4.031718in}{1.829583in}}%
\pgfpathlineto{\pgfqpoint{4.034237in}{1.826633in}}%
\pgfpathlineto{\pgfqpoint{4.036756in}{1.818024in}}%
\pgfpathlineto{\pgfqpoint{4.039274in}{1.807351in}}%
\pgfpathlineto{\pgfqpoint{4.041793in}{1.827267in}}%
\pgfpathlineto{\pgfqpoint{4.044312in}{1.792987in}}%
\pgfpathlineto{\pgfqpoint{4.046831in}{1.818899in}}%
\pgfpathlineto{\pgfqpoint{4.049349in}{1.824963in}}%
\pgfpathlineto{\pgfqpoint{4.051868in}{1.841810in}}%
\pgfpathlineto{\pgfqpoint{4.054387in}{1.842952in}}%
\pgfpathlineto{\pgfqpoint{4.056906in}{1.847742in}}%
\pgfpathlineto{\pgfqpoint{4.059424in}{1.879831in}}%
\pgfpathlineto{\pgfqpoint{4.061943in}{1.866593in}}%
\pgfpathlineto{\pgfqpoint{4.064462in}{1.873338in}}%
\pgfpathlineto{\pgfqpoint{4.066981in}{1.872796in}}%
\pgfpathlineto{\pgfqpoint{4.069499in}{1.880342in}}%
\pgfpathlineto{\pgfqpoint{4.072018in}{1.881431in}}%
\pgfpathlineto{\pgfqpoint{4.074537in}{1.865274in}}%
\pgfpathlineto{\pgfqpoint{4.077056in}{1.883646in}}%
\pgfpathlineto{\pgfqpoint{4.079574in}{1.888864in}}%
\pgfpathlineto{\pgfqpoint{4.082093in}{1.879080in}}%
\pgfpathlineto{\pgfqpoint{4.084612in}{1.901289in}}%
\pgfpathlineto{\pgfqpoint{4.087131in}{1.907256in}}%
\pgfpathlineto{\pgfqpoint{4.089649in}{1.910391in}}%
\pgfpathlineto{\pgfqpoint{4.092168in}{1.928242in}}%
\pgfpathlineto{\pgfqpoint{4.094687in}{1.902125in}}%
\pgfpathlineto{\pgfqpoint{4.097206in}{1.906790in}}%
\pgfpathlineto{\pgfqpoint{4.099724in}{1.912150in}}%
\pgfpathlineto{\pgfqpoint{4.102243in}{1.879815in}}%
\pgfpathlineto{\pgfqpoint{4.104762in}{1.873570in}}%
\pgfpathlineto{\pgfqpoint{4.107281in}{1.864827in}}%
\pgfpathlineto{\pgfqpoint{4.109799in}{1.874731in}}%
\pgfpathlineto{\pgfqpoint{4.112318in}{1.898634in}}%
\pgfpathlineto{\pgfqpoint{4.114837in}{1.897021in}}%
\pgfpathlineto{\pgfqpoint{4.117356in}{1.904291in}}%
\pgfpathlineto{\pgfqpoint{4.119874in}{1.887350in}}%
\pgfpathlineto{\pgfqpoint{4.122393in}{1.902522in}}%
\pgfpathlineto{\pgfqpoint{4.124912in}{1.887756in}}%
\pgfpathlineto{\pgfqpoint{4.127431in}{1.880743in}}%
\pgfpathlineto{\pgfqpoint{4.129949in}{1.879314in}}%
\pgfpathlineto{\pgfqpoint{4.132468in}{1.903072in}}%
\pgfpathlineto{\pgfqpoint{4.134987in}{1.921083in}}%
\pgfpathlineto{\pgfqpoint{4.137506in}{1.926856in}}%
\pgfpathlineto{\pgfqpoint{4.140024in}{1.913917in}}%
\pgfpathlineto{\pgfqpoint{4.142543in}{1.916224in}}%
\pgfpathlineto{\pgfqpoint{4.147581in}{1.900563in}}%
\pgfpathlineto{\pgfqpoint{4.150099in}{1.897450in}}%
\pgfpathlineto{\pgfqpoint{4.152618in}{1.876083in}}%
\pgfpathlineto{\pgfqpoint{4.155137in}{1.887759in}}%
\pgfpathlineto{\pgfqpoint{4.157656in}{1.890303in}}%
\pgfpathlineto{\pgfqpoint{4.160174in}{1.911856in}}%
\pgfpathlineto{\pgfqpoint{4.162693in}{1.891156in}}%
\pgfpathlineto{\pgfqpoint{4.165212in}{1.917674in}}%
\pgfpathlineto{\pgfqpoint{4.167731in}{1.904182in}}%
\pgfpathlineto{\pgfqpoint{4.170249in}{1.885902in}}%
\pgfpathlineto{\pgfqpoint{4.172768in}{1.891510in}}%
\pgfpathlineto{\pgfqpoint{4.175287in}{1.908080in}}%
\pgfpathlineto{\pgfqpoint{4.177806in}{1.929448in}}%
\pgfpathlineto{\pgfqpoint{4.180324in}{1.933829in}}%
\pgfpathlineto{\pgfqpoint{4.182843in}{1.918396in}}%
\pgfpathlineto{\pgfqpoint{4.185362in}{1.912483in}}%
\pgfpathlineto{\pgfqpoint{4.187881in}{1.933001in}}%
\pgfpathlineto{\pgfqpoint{4.190399in}{1.911565in}}%
\pgfpathlineto{\pgfqpoint{4.192918in}{1.880171in}}%
\pgfpathlineto{\pgfqpoint{4.195437in}{1.880117in}}%
\pgfpathlineto{\pgfqpoint{4.197956in}{1.899395in}}%
\pgfpathlineto{\pgfqpoint{4.200474in}{1.886909in}}%
\pgfpathlineto{\pgfqpoint{4.202993in}{1.880419in}}%
\pgfpathlineto{\pgfqpoint{4.205512in}{1.872029in}}%
\pgfpathlineto{\pgfqpoint{4.208031in}{1.864958in}}%
\pgfpathlineto{\pgfqpoint{4.210549in}{1.857579in}}%
\pgfpathlineto{\pgfqpoint{4.213068in}{1.872466in}}%
\pgfpathlineto{\pgfqpoint{4.215587in}{1.898207in}}%
\pgfpathlineto{\pgfqpoint{4.218106in}{1.901218in}}%
\pgfpathlineto{\pgfqpoint{4.220624in}{1.919729in}}%
\pgfpathlineto{\pgfqpoint{4.223143in}{1.908367in}}%
\pgfpathlineto{\pgfqpoint{4.225662in}{1.916455in}}%
\pgfpathlineto{\pgfqpoint{4.228181in}{1.906543in}}%
\pgfpathlineto{\pgfqpoint{4.230699in}{1.886762in}}%
\pgfpathlineto{\pgfqpoint{4.233218in}{1.883276in}}%
\pgfpathlineto{\pgfqpoint{4.235737in}{1.915293in}}%
\pgfpathlineto{\pgfqpoint{4.238256in}{1.891399in}}%
\pgfpathlineto{\pgfqpoint{4.240774in}{1.883299in}}%
\pgfpathlineto{\pgfqpoint{4.243293in}{1.874006in}}%
\pgfpathlineto{\pgfqpoint{4.245812in}{1.895272in}}%
\pgfpathlineto{\pgfqpoint{4.248331in}{1.932660in}}%
\pgfpathlineto{\pgfqpoint{4.250849in}{1.946423in}}%
\pgfpathlineto{\pgfqpoint{4.253368in}{1.962944in}}%
\pgfpathlineto{\pgfqpoint{4.255887in}{1.995137in}}%
\pgfpathlineto{\pgfqpoint{4.258406in}{1.998915in}}%
\pgfpathlineto{\pgfqpoint{4.260924in}{1.982856in}}%
\pgfpathlineto{\pgfqpoint{4.263443in}{1.971036in}}%
\pgfpathlineto{\pgfqpoint{4.265962in}{1.965049in}}%
\pgfpathlineto{\pgfqpoint{4.268481in}{1.962908in}}%
\pgfpathlineto{\pgfqpoint{4.270999in}{1.951939in}}%
\pgfpathlineto{\pgfqpoint{4.273518in}{1.935391in}}%
\pgfpathlineto{\pgfqpoint{4.276037in}{1.960809in}}%
\pgfpathlineto{\pgfqpoint{4.278556in}{1.997809in}}%
\pgfpathlineto{\pgfqpoint{4.281074in}{1.985127in}}%
\pgfpathlineto{\pgfqpoint{4.283593in}{2.006234in}}%
\pgfpathlineto{\pgfqpoint{4.286112in}{1.999107in}}%
\pgfpathlineto{\pgfqpoint{4.288631in}{2.003197in}}%
\pgfpathlineto{\pgfqpoint{4.291149in}{1.982664in}}%
\pgfpathlineto{\pgfqpoint{4.293668in}{1.971997in}}%
\pgfpathlineto{\pgfqpoint{4.296187in}{1.968787in}}%
\pgfpathlineto{\pgfqpoint{4.298706in}{1.973482in}}%
\pgfpathlineto{\pgfqpoint{4.301224in}{1.987668in}}%
\pgfpathlineto{\pgfqpoint{4.303743in}{2.006120in}}%
\pgfpathlineto{\pgfqpoint{4.306262in}{1.980495in}}%
\pgfpathlineto{\pgfqpoint{4.308781in}{1.982904in}}%
\pgfpathlineto{\pgfqpoint{4.311299in}{1.991044in}}%
\pgfpathlineto{\pgfqpoint{4.313818in}{2.019215in}}%
\pgfpathlineto{\pgfqpoint{4.316337in}{2.003371in}}%
\pgfpathlineto{\pgfqpoint{4.318856in}{2.023886in}}%
\pgfpathlineto{\pgfqpoint{4.321374in}{2.015565in}}%
\pgfpathlineto{\pgfqpoint{4.323893in}{2.012664in}}%
\pgfpathlineto{\pgfqpoint{4.326412in}{1.978992in}}%
\pgfpathlineto{\pgfqpoint{4.328931in}{2.005542in}}%
\pgfpathlineto{\pgfqpoint{4.331449in}{1.997463in}}%
\pgfpathlineto{\pgfqpoint{4.333968in}{1.999350in}}%
\pgfpathlineto{\pgfqpoint{4.336487in}{1.979042in}}%
\pgfpathlineto{\pgfqpoint{4.339006in}{1.974462in}}%
\pgfpathlineto{\pgfqpoint{4.341524in}{2.005851in}}%
\pgfpathlineto{\pgfqpoint{4.344043in}{1.993666in}}%
\pgfpathlineto{\pgfqpoint{4.346562in}{2.009094in}}%
\pgfpathlineto{\pgfqpoint{4.349081in}{1.971509in}}%
\pgfpathlineto{\pgfqpoint{4.351599in}{1.960917in}}%
\pgfpathlineto{\pgfqpoint{4.354118in}{1.961052in}}%
\pgfpathlineto{\pgfqpoint{4.356637in}{1.935923in}}%
\pgfpathlineto{\pgfqpoint{4.359156in}{1.918852in}}%
\pgfpathlineto{\pgfqpoint{4.361674in}{1.909190in}}%
\pgfpathlineto{\pgfqpoint{4.364193in}{1.910146in}}%
\pgfpathlineto{\pgfqpoint{4.366712in}{1.871332in}}%
\pgfpathlineto{\pgfqpoint{4.369231in}{1.873029in}}%
\pgfpathlineto{\pgfqpoint{4.371749in}{1.860488in}}%
\pgfpathlineto{\pgfqpoint{4.374268in}{1.855488in}}%
\pgfpathlineto{\pgfqpoint{4.376787in}{1.862528in}}%
\pgfpathlineto{\pgfqpoint{4.379306in}{1.866243in}}%
\pgfpathlineto{\pgfqpoint{4.381824in}{1.862202in}}%
\pgfpathlineto{\pgfqpoint{4.384343in}{1.888675in}}%
\pgfpathlineto{\pgfqpoint{4.386862in}{1.900012in}}%
\pgfpathlineto{\pgfqpoint{4.389381in}{1.887104in}}%
\pgfpathlineto{\pgfqpoint{4.391899in}{1.892958in}}%
\pgfpathlineto{\pgfqpoint{4.394418in}{1.886027in}}%
\pgfpathlineto{\pgfqpoint{4.396937in}{1.889577in}}%
\pgfpathlineto{\pgfqpoint{4.399456in}{1.890625in}}%
\pgfpathlineto{\pgfqpoint{4.401974in}{1.880224in}}%
\pgfpathlineto{\pgfqpoint{4.404493in}{1.880223in}}%
\pgfpathlineto{\pgfqpoint{4.407012in}{1.857197in}}%
\pgfpathlineto{\pgfqpoint{4.409531in}{1.867883in}}%
\pgfpathlineto{\pgfqpoint{4.412049in}{1.856004in}}%
\pgfpathlineto{\pgfqpoint{4.414568in}{1.834928in}}%
\pgfpathlineto{\pgfqpoint{4.417087in}{1.831954in}}%
\pgfpathlineto{\pgfqpoint{4.419606in}{1.801068in}}%
\pgfpathlineto{\pgfqpoint{4.422124in}{1.785369in}}%
\pgfpathlineto{\pgfqpoint{4.424643in}{1.745143in}}%
\pgfpathlineto{\pgfqpoint{4.427162in}{1.781347in}}%
\pgfpathlineto{\pgfqpoint{4.429681in}{1.804387in}}%
\pgfpathlineto{\pgfqpoint{4.432199in}{1.809035in}}%
\pgfpathlineto{\pgfqpoint{4.434718in}{1.811249in}}%
\pgfpathlineto{\pgfqpoint{4.437237in}{1.822951in}}%
\pgfpathlineto{\pgfqpoint{4.439756in}{1.848430in}}%
\pgfpathlineto{\pgfqpoint{4.442274in}{1.825572in}}%
\pgfpathlineto{\pgfqpoint{4.444793in}{1.839581in}}%
\pgfpathlineto{\pgfqpoint{4.447312in}{1.840144in}}%
\pgfpathlineto{\pgfqpoint{4.449831in}{1.872287in}}%
\pgfpathlineto{\pgfqpoint{4.452349in}{1.854968in}}%
\pgfpathlineto{\pgfqpoint{4.454868in}{1.862809in}}%
\pgfpathlineto{\pgfqpoint{4.457387in}{1.860649in}}%
\pgfpathlineto{\pgfqpoint{4.459906in}{1.882924in}}%
\pgfpathlineto{\pgfqpoint{4.462424in}{1.900002in}}%
\pgfpathlineto{\pgfqpoint{4.464943in}{1.898509in}}%
\pgfpathlineto{\pgfqpoint{4.467462in}{1.914066in}}%
\pgfpathlineto{\pgfqpoint{4.469981in}{1.919941in}}%
\pgfpathlineto{\pgfqpoint{4.472499in}{1.900598in}}%
\pgfpathlineto{\pgfqpoint{4.475018in}{1.882753in}}%
\pgfpathlineto{\pgfqpoint{4.477537in}{1.887307in}}%
\pgfpathlineto{\pgfqpoint{4.480056in}{1.885138in}}%
\pgfpathlineto{\pgfqpoint{4.482574in}{1.899244in}}%
\pgfpathlineto{\pgfqpoint{4.485093in}{1.888209in}}%
\pgfpathlineto{\pgfqpoint{4.487612in}{1.868948in}}%
\pgfpathlineto{\pgfqpoint{4.490131in}{1.879042in}}%
\pgfpathlineto{\pgfqpoint{4.492649in}{1.882962in}}%
\pgfpathlineto{\pgfqpoint{4.495168in}{1.864163in}}%
\pgfpathlineto{\pgfqpoint{4.497687in}{1.893470in}}%
\pgfpathlineto{\pgfqpoint{4.500206in}{1.886588in}}%
\pgfpathlineto{\pgfqpoint{4.502724in}{1.900010in}}%
\pgfpathlineto{\pgfqpoint{4.505243in}{1.918920in}}%
\pgfpathlineto{\pgfqpoint{4.507762in}{1.876331in}}%
\pgfpathlineto{\pgfqpoint{4.510281in}{1.884623in}}%
\pgfpathlineto{\pgfqpoint{4.512799in}{1.879902in}}%
\pgfpathlineto{\pgfqpoint{4.515318in}{1.878724in}}%
\pgfpathlineto{\pgfqpoint{4.517837in}{1.877802in}}%
\pgfpathlineto{\pgfqpoint{4.520356in}{1.877732in}}%
\pgfpathlineto{\pgfqpoint{4.522874in}{1.900317in}}%
\pgfpathlineto{\pgfqpoint{4.525393in}{1.857658in}}%
\pgfpathlineto{\pgfqpoint{4.527912in}{1.849401in}}%
\pgfpathlineto{\pgfqpoint{4.530431in}{1.844097in}}%
\pgfpathlineto{\pgfqpoint{4.532949in}{1.867084in}}%
\pgfpathlineto{\pgfqpoint{4.535468in}{1.880248in}}%
\pgfpathlineto{\pgfqpoint{4.537987in}{1.872610in}}%
\pgfpathlineto{\pgfqpoint{4.540506in}{1.875036in}}%
\pgfpathlineto{\pgfqpoint{4.543024in}{1.874779in}}%
\pgfpathlineto{\pgfqpoint{4.545543in}{1.864492in}}%
\pgfpathlineto{\pgfqpoint{4.548062in}{1.861678in}}%
\pgfpathlineto{\pgfqpoint{4.550581in}{1.854377in}}%
\pgfpathlineto{\pgfqpoint{4.553099in}{1.873798in}}%
\pgfpathlineto{\pgfqpoint{4.555618in}{1.902127in}}%
\pgfpathlineto{\pgfqpoint{4.558137in}{1.889729in}}%
\pgfpathlineto{\pgfqpoint{4.560656in}{1.868861in}}%
\pgfpathlineto{\pgfqpoint{4.563174in}{1.873468in}}%
\pgfpathlineto{\pgfqpoint{4.565693in}{1.885711in}}%
\pgfpathlineto{\pgfqpoint{4.568212in}{1.886395in}}%
\pgfpathlineto{\pgfqpoint{4.570731in}{1.872559in}}%
\pgfpathlineto{\pgfqpoint{4.573249in}{1.874944in}}%
\pgfpathlineto{\pgfqpoint{4.575768in}{1.867776in}}%
\pgfpathlineto{\pgfqpoint{4.578287in}{1.852249in}}%
\pgfpathlineto{\pgfqpoint{4.580806in}{1.874825in}}%
\pgfpathlineto{\pgfqpoint{4.583324in}{1.879932in}}%
\pgfpathlineto{\pgfqpoint{4.585843in}{1.877873in}}%
\pgfpathlineto{\pgfqpoint{4.588362in}{1.883509in}}%
\pgfpathlineto{\pgfqpoint{4.590881in}{1.890124in}}%
\pgfpathlineto{\pgfqpoint{4.593399in}{1.915302in}}%
\pgfpathlineto{\pgfqpoint{4.595918in}{1.892369in}}%
\pgfpathlineto{\pgfqpoint{4.598437in}{1.885286in}}%
\pgfpathlineto{\pgfqpoint{4.600956in}{1.905769in}}%
\pgfpathlineto{\pgfqpoint{4.603474in}{1.923731in}}%
\pgfpathlineto{\pgfqpoint{4.605993in}{1.933187in}}%
\pgfpathlineto{\pgfqpoint{4.608512in}{1.946282in}}%
\pgfpathlineto{\pgfqpoint{4.611031in}{1.950855in}}%
\pgfpathlineto{\pgfqpoint{4.613549in}{1.964120in}}%
\pgfpathlineto{\pgfqpoint{4.616068in}{1.989372in}}%
\pgfpathlineto{\pgfqpoint{4.618587in}{1.997571in}}%
\pgfpathlineto{\pgfqpoint{4.621106in}{2.017839in}}%
\pgfpathlineto{\pgfqpoint{4.623624in}{2.019805in}}%
\pgfpathlineto{\pgfqpoint{4.626143in}{2.030775in}}%
\pgfpathlineto{\pgfqpoint{4.628662in}{2.034437in}}%
\pgfpathlineto{\pgfqpoint{4.631181in}{2.056873in}}%
\pgfpathlineto{\pgfqpoint{4.633699in}{2.059531in}}%
\pgfpathlineto{\pgfqpoint{4.636218in}{2.058619in}}%
\pgfpathlineto{\pgfqpoint{4.638737in}{2.082049in}}%
\pgfpathlineto{\pgfqpoint{4.641256in}{2.095925in}}%
\pgfpathlineto{\pgfqpoint{4.646293in}{2.044996in}}%
\pgfpathlineto{\pgfqpoint{4.648812in}{2.071213in}}%
\pgfpathlineto{\pgfqpoint{4.651331in}{2.109594in}}%
\pgfpathlineto{\pgfqpoint{4.653849in}{2.114577in}}%
\pgfpathlineto{\pgfqpoint{4.656368in}{2.087579in}}%
\pgfpathlineto{\pgfqpoint{4.658887in}{2.099811in}}%
\pgfpathlineto{\pgfqpoint{4.661406in}{2.112612in}}%
\pgfpathlineto{\pgfqpoint{4.663924in}{2.138807in}}%
\pgfpathlineto{\pgfqpoint{4.666443in}{2.147550in}}%
\pgfpathlineto{\pgfqpoint{4.668962in}{2.161163in}}%
\pgfpathlineto{\pgfqpoint{4.671481in}{2.169520in}}%
\pgfpathlineto{\pgfqpoint{4.673999in}{2.193526in}}%
\pgfpathlineto{\pgfqpoint{4.676518in}{2.192591in}}%
\pgfpathlineto{\pgfqpoint{4.679037in}{2.184093in}}%
\pgfpathlineto{\pgfqpoint{4.681556in}{2.167181in}}%
\pgfpathlineto{\pgfqpoint{4.684074in}{2.161843in}}%
\pgfpathlineto{\pgfqpoint{4.686593in}{2.149067in}}%
\pgfpathlineto{\pgfqpoint{4.689112in}{2.120703in}}%
\pgfpathlineto{\pgfqpoint{4.691631in}{2.127376in}}%
\pgfpathlineto{\pgfqpoint{4.694149in}{2.143185in}}%
\pgfpathlineto{\pgfqpoint{4.696668in}{2.130570in}}%
\pgfpathlineto{\pgfqpoint{4.699187in}{2.160135in}}%
\pgfpathlineto{\pgfqpoint{4.701706in}{2.178448in}}%
\pgfpathlineto{\pgfqpoint{4.704224in}{2.154527in}}%
\pgfpathlineto{\pgfqpoint{4.706743in}{2.169940in}}%
\pgfpathlineto{\pgfqpoint{4.709262in}{2.132874in}}%
\pgfpathlineto{\pgfqpoint{4.711781in}{2.121200in}}%
\pgfpathlineto{\pgfqpoint{4.714299in}{2.124631in}}%
\pgfpathlineto{\pgfqpoint{4.716818in}{2.139501in}}%
\pgfpathlineto{\pgfqpoint{4.719337in}{2.132101in}}%
\pgfpathlineto{\pgfqpoint{4.721856in}{2.128903in}}%
\pgfpathlineto{\pgfqpoint{4.724374in}{2.113921in}}%
\pgfpathlineto{\pgfqpoint{4.726893in}{2.123496in}}%
\pgfpathlineto{\pgfqpoint{4.729412in}{2.138766in}}%
\pgfpathlineto{\pgfqpoint{4.731931in}{2.147615in}}%
\pgfpathlineto{\pgfqpoint{4.734449in}{2.151111in}}%
\pgfpathlineto{\pgfqpoint{4.736968in}{2.145367in}}%
\pgfpathlineto{\pgfqpoint{4.739487in}{2.115138in}}%
\pgfpathlineto{\pgfqpoint{4.742006in}{2.091008in}}%
\pgfpathlineto{\pgfqpoint{4.744524in}{2.100545in}}%
\pgfpathlineto{\pgfqpoint{4.747043in}{2.097260in}}%
\pgfpathlineto{\pgfqpoint{4.749562in}{2.096548in}}%
\pgfpathlineto{\pgfqpoint{4.752081in}{2.082617in}}%
\pgfpathlineto{\pgfqpoint{4.754599in}{2.085813in}}%
\pgfpathlineto{\pgfqpoint{4.757118in}{2.093460in}}%
\pgfpathlineto{\pgfqpoint{4.759637in}{2.098809in}}%
\pgfpathlineto{\pgfqpoint{4.762156in}{2.122025in}}%
\pgfpathlineto{\pgfqpoint{4.764674in}{2.135245in}}%
\pgfpathlineto{\pgfqpoint{4.767193in}{2.130864in}}%
\pgfpathlineto{\pgfqpoint{4.769712in}{2.154357in}}%
\pgfpathlineto{\pgfqpoint{4.772231in}{2.162116in}}%
\pgfpathlineto{\pgfqpoint{4.774749in}{2.166566in}}%
\pgfpathlineto{\pgfqpoint{4.777268in}{2.161670in}}%
\pgfpathlineto{\pgfqpoint{4.779787in}{2.149037in}}%
\pgfpathlineto{\pgfqpoint{4.782306in}{2.151151in}}%
\pgfpathlineto{\pgfqpoint{4.784824in}{2.160389in}}%
\pgfpathlineto{\pgfqpoint{4.787343in}{2.169153in}}%
\pgfpathlineto{\pgfqpoint{4.789862in}{2.187021in}}%
\pgfpathlineto{\pgfqpoint{4.792381in}{2.200880in}}%
\pgfpathlineto{\pgfqpoint{4.794899in}{2.176342in}}%
\pgfpathlineto{\pgfqpoint{4.797418in}{2.198508in}}%
\pgfpathlineto{\pgfqpoint{4.799937in}{2.182769in}}%
\pgfpathlineto{\pgfqpoint{4.802456in}{2.174929in}}%
\pgfpathlineto{\pgfqpoint{4.804974in}{2.183934in}}%
\pgfpathlineto{\pgfqpoint{4.807493in}{2.187115in}}%
\pgfpathlineto{\pgfqpoint{4.810012in}{2.212984in}}%
\pgfpathlineto{\pgfqpoint{4.812531in}{2.268383in}}%
\pgfpathlineto{\pgfqpoint{4.815049in}{2.298535in}}%
\pgfpathlineto{\pgfqpoint{4.817568in}{2.291060in}}%
\pgfpathlineto{\pgfqpoint{4.820087in}{2.277333in}}%
\pgfpathlineto{\pgfqpoint{4.822606in}{2.259576in}}%
\pgfpathlineto{\pgfqpoint{4.825124in}{2.280682in}}%
\pgfpathlineto{\pgfqpoint{4.827643in}{2.292717in}}%
\pgfpathlineto{\pgfqpoint{4.830162in}{2.290495in}}%
\pgfpathlineto{\pgfqpoint{4.832681in}{2.277430in}}%
\pgfpathlineto{\pgfqpoint{4.835199in}{2.320197in}}%
\pgfpathlineto{\pgfqpoint{4.837718in}{2.297216in}}%
\pgfpathlineto{\pgfqpoint{4.840237in}{2.318260in}}%
\pgfpathlineto{\pgfqpoint{4.842756in}{2.337575in}}%
\pgfpathlineto{\pgfqpoint{4.845274in}{2.369431in}}%
\pgfpathlineto{\pgfqpoint{4.847793in}{2.382632in}}%
\pgfpathlineto{\pgfqpoint{4.850312in}{2.393748in}}%
\pgfpathlineto{\pgfqpoint{4.855349in}{2.361771in}}%
\pgfpathlineto{\pgfqpoint{4.857868in}{2.371467in}}%
\pgfpathlineto{\pgfqpoint{4.860387in}{2.372532in}}%
\pgfpathlineto{\pgfqpoint{4.862906in}{2.348102in}}%
\pgfpathlineto{\pgfqpoint{4.865424in}{2.356725in}}%
\pgfpathlineto{\pgfqpoint{4.867943in}{2.338479in}}%
\pgfpathlineto{\pgfqpoint{4.870462in}{2.329942in}}%
\pgfpathlineto{\pgfqpoint{4.872981in}{2.340861in}}%
\pgfpathlineto{\pgfqpoint{4.875499in}{2.332887in}}%
\pgfpathlineto{\pgfqpoint{4.878018in}{2.309442in}}%
\pgfpathlineto{\pgfqpoint{4.880537in}{2.305097in}}%
\pgfpathlineto{\pgfqpoint{4.883056in}{2.341921in}}%
\pgfpathlineto{\pgfqpoint{4.885574in}{2.363645in}}%
\pgfpathlineto{\pgfqpoint{4.888093in}{2.354226in}}%
\pgfpathlineto{\pgfqpoint{4.890612in}{2.340766in}}%
\pgfpathlineto{\pgfqpoint{4.893131in}{2.326356in}}%
\pgfpathlineto{\pgfqpoint{4.895649in}{2.332729in}}%
\pgfpathlineto{\pgfqpoint{4.898168in}{2.282306in}}%
\pgfpathlineto{\pgfqpoint{4.900687in}{2.303572in}}%
\pgfpathlineto{\pgfqpoint{4.903206in}{2.293173in}}%
\pgfpathlineto{\pgfqpoint{4.905724in}{2.285901in}}%
\pgfpathlineto{\pgfqpoint{4.908243in}{2.281541in}}%
\pgfpathlineto{\pgfqpoint{4.910762in}{2.277777in}}%
\pgfpathlineto{\pgfqpoint{4.913281in}{2.251585in}}%
\pgfpathlineto{\pgfqpoint{4.915799in}{2.287743in}}%
\pgfpathlineto{\pgfqpoint{4.918318in}{2.298530in}}%
\pgfpathlineto{\pgfqpoint{4.920837in}{2.318174in}}%
\pgfpathlineto{\pgfqpoint{4.923356in}{2.329048in}}%
\pgfpathlineto{\pgfqpoint{4.925874in}{2.306073in}}%
\pgfpathlineto{\pgfqpoint{4.928393in}{2.311981in}}%
\pgfpathlineto{\pgfqpoint{4.930912in}{2.252549in}}%
\pgfpathlineto{\pgfqpoint{4.933431in}{2.252014in}}%
\pgfpathlineto{\pgfqpoint{4.935949in}{2.262137in}}%
\pgfpathlineto{\pgfqpoint{4.938468in}{2.256116in}}%
\pgfpathlineto{\pgfqpoint{4.940987in}{2.276309in}}%
\pgfpathlineto{\pgfqpoint{4.943506in}{2.323450in}}%
\pgfpathlineto{\pgfqpoint{4.946024in}{2.348704in}}%
\pgfpathlineto{\pgfqpoint{4.948543in}{2.362641in}}%
\pgfpathlineto{\pgfqpoint{4.951062in}{2.346286in}}%
\pgfpathlineto{\pgfqpoint{4.953581in}{2.346912in}}%
\pgfpathlineto{\pgfqpoint{4.956099in}{2.358394in}}%
\pgfpathlineto{\pgfqpoint{4.958618in}{2.363375in}}%
\pgfpathlineto{\pgfqpoint{4.961137in}{2.386006in}}%
\pgfpathlineto{\pgfqpoint{4.963656in}{2.419852in}}%
\pgfpathlineto{\pgfqpoint{4.966174in}{2.436532in}}%
\pgfpathlineto{\pgfqpoint{4.968693in}{2.430165in}}%
\pgfpathlineto{\pgfqpoint{4.971212in}{2.406514in}}%
\pgfpathlineto{\pgfqpoint{4.973731in}{2.397754in}}%
\pgfpathlineto{\pgfqpoint{4.976249in}{2.375222in}}%
\pgfpathlineto{\pgfqpoint{4.978768in}{2.382823in}}%
\pgfpathlineto{\pgfqpoint{4.981287in}{2.364116in}}%
\pgfpathlineto{\pgfqpoint{4.983806in}{2.384950in}}%
\pgfpathlineto{\pgfqpoint{4.986324in}{2.362039in}}%
\pgfpathlineto{\pgfqpoint{4.988843in}{2.372865in}}%
\pgfpathlineto{\pgfqpoint{4.991362in}{2.347529in}}%
\pgfpathlineto{\pgfqpoint{4.993881in}{2.347926in}}%
\pgfpathlineto{\pgfqpoint{4.996399in}{2.321757in}}%
\pgfpathlineto{\pgfqpoint{4.998918in}{2.327882in}}%
\pgfpathlineto{\pgfqpoint{5.001437in}{2.320240in}}%
\pgfpathlineto{\pgfqpoint{5.003956in}{2.348062in}}%
\pgfpathlineto{\pgfqpoint{5.006474in}{2.337149in}}%
\pgfpathlineto{\pgfqpoint{5.008993in}{2.330846in}}%
\pgfpathlineto{\pgfqpoint{5.011512in}{2.332506in}}%
\pgfpathlineto{\pgfqpoint{5.014031in}{2.335499in}}%
\pgfpathlineto{\pgfqpoint{5.016549in}{2.339289in}}%
\pgfpathlineto{\pgfqpoint{5.019068in}{2.335178in}}%
\pgfpathlineto{\pgfqpoint{5.021587in}{2.305778in}}%
\pgfpathlineto{\pgfqpoint{5.024106in}{2.260288in}}%
\pgfpathlineto{\pgfqpoint{5.026624in}{2.277645in}}%
\pgfpathlineto{\pgfqpoint{5.029143in}{2.272780in}}%
\pgfpathlineto{\pgfqpoint{5.031662in}{2.275088in}}%
\pgfpathlineto{\pgfqpoint{5.034181in}{2.227157in}}%
\pgfpathlineto{\pgfqpoint{5.036699in}{2.227449in}}%
\pgfpathlineto{\pgfqpoint{5.039218in}{2.209542in}}%
\pgfpathlineto{\pgfqpoint{5.041737in}{2.234271in}}%
\pgfpathlineto{\pgfqpoint{5.044256in}{2.242209in}}%
\pgfpathlineto{\pgfqpoint{5.046774in}{2.230123in}}%
\pgfpathlineto{\pgfqpoint{5.049293in}{2.244928in}}%
\pgfpathlineto{\pgfqpoint{5.051812in}{2.256482in}}%
\pgfpathlineto{\pgfqpoint{5.054331in}{2.257641in}}%
\pgfpathlineto{\pgfqpoint{5.056849in}{2.267829in}}%
\pgfpathlineto{\pgfqpoint{5.059368in}{2.267882in}}%
\pgfpathlineto{\pgfqpoint{5.061887in}{2.272017in}}%
\pgfpathlineto{\pgfqpoint{5.064406in}{2.306581in}}%
\pgfpathlineto{\pgfqpoint{5.066924in}{2.290442in}}%
\pgfpathlineto{\pgfqpoint{5.069443in}{2.296942in}}%
\pgfpathlineto{\pgfqpoint{5.071962in}{2.302892in}}%
\pgfpathlineto{\pgfqpoint{5.074481in}{2.290461in}}%
\pgfpathlineto{\pgfqpoint{5.076999in}{2.283885in}}%
\pgfpathlineto{\pgfqpoint{5.079518in}{2.285571in}}%
\pgfpathlineto{\pgfqpoint{5.082037in}{2.296529in}}%
\pgfpathlineto{\pgfqpoint{5.084556in}{2.301486in}}%
\pgfpathlineto{\pgfqpoint{5.087074in}{2.309600in}}%
\pgfpathlineto{\pgfqpoint{5.089593in}{2.310402in}}%
\pgfpathlineto{\pgfqpoint{5.092112in}{2.300247in}}%
\pgfpathlineto{\pgfqpoint{5.094631in}{2.305098in}}%
\pgfpathlineto{\pgfqpoint{5.097149in}{2.307503in}}%
\pgfpathlineto{\pgfqpoint{5.099668in}{2.273938in}}%
\pgfpathlineto{\pgfqpoint{5.102187in}{2.270298in}}%
\pgfpathlineto{\pgfqpoint{5.104706in}{2.299181in}}%
\pgfpathlineto{\pgfqpoint{5.107224in}{2.291701in}}%
\pgfpathlineto{\pgfqpoint{5.109743in}{2.277480in}}%
\pgfpathlineto{\pgfqpoint{5.112262in}{2.275172in}}%
\pgfpathlineto{\pgfqpoint{5.114781in}{2.252350in}}%
\pgfpathlineto{\pgfqpoint{5.117299in}{2.215477in}}%
\pgfpathlineto{\pgfqpoint{5.119818in}{2.258390in}}%
\pgfpathlineto{\pgfqpoint{5.122337in}{2.266413in}}%
\pgfpathlineto{\pgfqpoint{5.124856in}{2.269087in}}%
\pgfpathlineto{\pgfqpoint{5.127374in}{2.272414in}}%
\pgfpathlineto{\pgfqpoint{5.129893in}{2.224264in}}%
\pgfpathlineto{\pgfqpoint{5.132412in}{2.234296in}}%
\pgfpathlineto{\pgfqpoint{5.134931in}{2.232137in}}%
\pgfpathlineto{\pgfqpoint{5.137449in}{2.221098in}}%
\pgfpathlineto{\pgfqpoint{5.139968in}{2.224937in}}%
\pgfpathlineto{\pgfqpoint{5.142487in}{2.224635in}}%
\pgfpathlineto{\pgfqpoint{5.145006in}{2.210662in}}%
\pgfpathlineto{\pgfqpoint{5.147524in}{2.222805in}}%
\pgfpathlineto{\pgfqpoint{5.150043in}{2.218754in}}%
\pgfpathlineto{\pgfqpoint{5.152562in}{2.230412in}}%
\pgfpathlineto{\pgfqpoint{5.155081in}{2.241137in}}%
\pgfpathlineto{\pgfqpoint{5.157599in}{2.237933in}}%
\pgfpathlineto{\pgfqpoint{5.160118in}{2.272307in}}%
\pgfpathlineto{\pgfqpoint{5.162637in}{2.291818in}}%
\pgfpathlineto{\pgfqpoint{5.165156in}{2.305542in}}%
\pgfpathlineto{\pgfqpoint{5.167674in}{2.298186in}}%
\pgfpathlineto{\pgfqpoint{5.170193in}{2.286070in}}%
\pgfpathlineto{\pgfqpoint{5.172712in}{2.286003in}}%
\pgfpathlineto{\pgfqpoint{5.175231in}{2.285313in}}%
\pgfpathlineto{\pgfqpoint{5.177749in}{2.317717in}}%
\pgfpathlineto{\pgfqpoint{5.180268in}{2.324025in}}%
\pgfpathlineto{\pgfqpoint{5.182787in}{2.325522in}}%
\pgfpathlineto{\pgfqpoint{5.185306in}{2.335253in}}%
\pgfpathlineto{\pgfqpoint{5.187824in}{2.335858in}}%
\pgfpathlineto{\pgfqpoint{5.190343in}{2.361362in}}%
\pgfpathlineto{\pgfqpoint{5.192862in}{2.368317in}}%
\pgfpathlineto{\pgfqpoint{5.195381in}{2.388526in}}%
\pgfpathlineto{\pgfqpoint{5.197899in}{2.437880in}}%
\pgfpathlineto{\pgfqpoint{5.200418in}{2.457451in}}%
\pgfpathlineto{\pgfqpoint{5.202937in}{2.432070in}}%
\pgfpathlineto{\pgfqpoint{5.205456in}{2.432893in}}%
\pgfpathlineto{\pgfqpoint{5.207974in}{2.418395in}}%
\pgfpathlineto{\pgfqpoint{5.210493in}{2.461722in}}%
\pgfpathlineto{\pgfqpoint{5.213012in}{2.475360in}}%
\pgfpathlineto{\pgfqpoint{5.215531in}{2.480220in}}%
\pgfpathlineto{\pgfqpoint{5.218049in}{2.463614in}}%
\pgfpathlineto{\pgfqpoint{5.220568in}{2.463588in}}%
\pgfpathlineto{\pgfqpoint{5.223087in}{2.483949in}}%
\pgfpathlineto{\pgfqpoint{5.225606in}{2.448951in}}%
\pgfpathlineto{\pgfqpoint{5.228124in}{2.438357in}}%
\pgfpathlineto{\pgfqpoint{5.230643in}{2.437416in}}%
\pgfpathlineto{\pgfqpoint{5.233162in}{2.447267in}}%
\pgfpathlineto{\pgfqpoint{5.235681in}{2.428462in}}%
\pgfpathlineto{\pgfqpoint{5.238199in}{2.460945in}}%
\pgfpathlineto{\pgfqpoint{5.240718in}{2.427554in}}%
\pgfpathlineto{\pgfqpoint{5.243237in}{2.436157in}}%
\pgfpathlineto{\pgfqpoint{5.245756in}{2.423986in}}%
\pgfpathlineto{\pgfqpoint{5.248274in}{2.421268in}}%
\pgfpathlineto{\pgfqpoint{5.250793in}{2.424820in}}%
\pgfpathlineto{\pgfqpoint{5.253312in}{2.392632in}}%
\pgfpathlineto{\pgfqpoint{5.255831in}{2.388997in}}%
\pgfpathlineto{\pgfqpoint{5.258349in}{2.373639in}}%
\pgfpathlineto{\pgfqpoint{5.260868in}{2.339934in}}%
\pgfpathlineto{\pgfqpoint{5.263387in}{2.350269in}}%
\pgfpathlineto{\pgfqpoint{5.265906in}{2.351535in}}%
\pgfpathlineto{\pgfqpoint{5.268424in}{2.369938in}}%
\pgfpathlineto{\pgfqpoint{5.270943in}{2.416377in}}%
\pgfpathlineto{\pgfqpoint{5.273462in}{2.416149in}}%
\pgfpathlineto{\pgfqpoint{5.275981in}{2.407299in}}%
\pgfpathlineto{\pgfqpoint{5.278499in}{2.428528in}}%
\pgfpathlineto{\pgfqpoint{5.281018in}{2.407008in}}%
\pgfpathlineto{\pgfqpoint{5.283537in}{2.392045in}}%
\pgfpathlineto{\pgfqpoint{5.286056in}{2.391533in}}%
\pgfpathlineto{\pgfqpoint{5.288574in}{2.403608in}}%
\pgfpathlineto{\pgfqpoint{5.291093in}{2.420820in}}%
\pgfpathlineto{\pgfqpoint{5.293612in}{2.423963in}}%
\pgfpathlineto{\pgfqpoint{5.296131in}{2.413261in}}%
\pgfpathlineto{\pgfqpoint{5.298649in}{2.445472in}}%
\pgfpathlineto{\pgfqpoint{5.301168in}{2.436576in}}%
\pgfpathlineto{\pgfqpoint{5.303687in}{2.379490in}}%
\pgfpathlineto{\pgfqpoint{5.306206in}{2.405411in}}%
\pgfpathlineto{\pgfqpoint{5.308724in}{2.411020in}}%
\pgfpathlineto{\pgfqpoint{5.311243in}{2.404312in}}%
\pgfpathlineto{\pgfqpoint{5.313762in}{2.414137in}}%
\pgfpathlineto{\pgfqpoint{5.316281in}{2.422210in}}%
\pgfpathlineto{\pgfqpoint{5.318799in}{2.417966in}}%
\pgfpathlineto{\pgfqpoint{5.321318in}{2.453025in}}%
\pgfpathlineto{\pgfqpoint{5.323837in}{2.463436in}}%
\pgfpathlineto{\pgfqpoint{5.326356in}{2.438974in}}%
\pgfpathlineto{\pgfqpoint{5.328874in}{2.469480in}}%
\pgfpathlineto{\pgfqpoint{5.331393in}{2.470886in}}%
\pgfpathlineto{\pgfqpoint{5.333912in}{2.475455in}}%
\pgfpathlineto{\pgfqpoint{5.336431in}{2.479755in}}%
\pgfpathlineto{\pgfqpoint{5.338949in}{2.470677in}}%
\pgfpathlineto{\pgfqpoint{5.341468in}{2.468210in}}%
\pgfpathlineto{\pgfqpoint{5.343987in}{2.469258in}}%
\pgfpathlineto{\pgfqpoint{5.346506in}{2.456350in}}%
\pgfpathlineto{\pgfqpoint{5.349024in}{2.439566in}}%
\pgfpathlineto{\pgfqpoint{5.351543in}{2.450247in}}%
\pgfpathlineto{\pgfqpoint{5.354062in}{2.430744in}}%
\pgfpathlineto{\pgfqpoint{5.356581in}{2.431258in}}%
\pgfpathlineto{\pgfqpoint{5.359099in}{2.449432in}}%
\pgfpathlineto{\pgfqpoint{5.361618in}{2.446205in}}%
\pgfpathlineto{\pgfqpoint{5.364137in}{2.424474in}}%
\pgfpathlineto{\pgfqpoint{5.366656in}{2.419588in}}%
\pgfpathlineto{\pgfqpoint{5.369174in}{2.434484in}}%
\pgfpathlineto{\pgfqpoint{5.371693in}{2.433170in}}%
\pgfpathlineto{\pgfqpoint{5.374212in}{2.428392in}}%
\pgfpathlineto{\pgfqpoint{5.376731in}{2.444843in}}%
\pgfpathlineto{\pgfqpoint{5.379249in}{2.440469in}}%
\pgfpathlineto{\pgfqpoint{5.381768in}{2.437778in}}%
\pgfpathlineto{\pgfqpoint{5.384287in}{2.453158in}}%
\pgfpathlineto{\pgfqpoint{5.386806in}{2.466567in}}%
\pgfpathlineto{\pgfqpoint{5.389324in}{2.455430in}}%
\pgfpathlineto{\pgfqpoint{5.391843in}{2.434546in}}%
\pgfpathlineto{\pgfqpoint{5.394362in}{2.423996in}}%
\pgfpathlineto{\pgfqpoint{5.396881in}{2.414969in}}%
\pgfpathlineto{\pgfqpoint{5.399399in}{2.462852in}}%
\pgfpathlineto{\pgfqpoint{5.401918in}{2.437384in}}%
\pgfpathlineto{\pgfqpoint{5.404437in}{2.425087in}}%
\pgfpathlineto{\pgfqpoint{5.406956in}{2.429819in}}%
\pgfpathlineto{\pgfqpoint{5.409474in}{2.426125in}}%
\pgfpathlineto{\pgfqpoint{5.411993in}{2.425808in}}%
\pgfpathlineto{\pgfqpoint{5.414512in}{2.415741in}}%
\pgfpathlineto{\pgfqpoint{5.417031in}{2.384039in}}%
\pgfpathlineto{\pgfqpoint{5.419549in}{2.407569in}}%
\pgfpathlineto{\pgfqpoint{5.422068in}{2.415420in}}%
\pgfpathlineto{\pgfqpoint{5.424587in}{2.389984in}}%
\pgfpathlineto{\pgfqpoint{5.427106in}{2.395802in}}%
\pgfpathlineto{\pgfqpoint{5.429624in}{2.422368in}}%
\pgfpathlineto{\pgfqpoint{5.434662in}{2.421686in}}%
\pgfpathlineto{\pgfqpoint{5.437181in}{2.422807in}}%
\pgfpathlineto{\pgfqpoint{5.439699in}{2.417657in}}%
\pgfpathlineto{\pgfqpoint{5.442218in}{2.434830in}}%
\pgfpathlineto{\pgfqpoint{5.444737in}{2.424727in}}%
\pgfpathlineto{\pgfqpoint{5.447256in}{2.423076in}}%
\pgfpathlineto{\pgfqpoint{5.449774in}{2.441194in}}%
\pgfpathlineto{\pgfqpoint{5.452293in}{2.445514in}}%
\pgfpathlineto{\pgfqpoint{5.454812in}{2.461025in}}%
\pgfpathlineto{\pgfqpoint{5.457331in}{2.488656in}}%
\pgfpathlineto{\pgfqpoint{5.459849in}{2.509044in}}%
\pgfpathlineto{\pgfqpoint{5.462368in}{2.516492in}}%
\pgfpathlineto{\pgfqpoint{5.464887in}{2.520438in}}%
\pgfpathlineto{\pgfqpoint{5.467406in}{2.493133in}}%
\pgfpathlineto{\pgfqpoint{5.469924in}{2.477036in}}%
\pgfpathlineto{\pgfqpoint{5.472443in}{2.477127in}}%
\pgfpathlineto{\pgfqpoint{5.474962in}{2.485694in}}%
\pgfpathlineto{\pgfqpoint{5.477481in}{2.498618in}}%
\pgfpathlineto{\pgfqpoint{5.479999in}{2.528400in}}%
\pgfpathlineto{\pgfqpoint{5.482518in}{2.521159in}}%
\pgfpathlineto{\pgfqpoint{5.485037in}{2.500073in}}%
\pgfpathlineto{\pgfqpoint{5.487556in}{2.540818in}}%
\pgfpathlineto{\pgfqpoint{5.490074in}{2.551794in}}%
\pgfpathlineto{\pgfqpoint{5.492593in}{2.548498in}}%
\pgfpathlineto{\pgfqpoint{5.492593in}{2.548498in}}%
\pgfusepath{stroke}%
\end{pgfscope}%
\begin{pgfscope}%
\pgfpathrectangle{\pgfqpoint{0.455093in}{0.401080in}}{\pgfqpoint{5.037500in}{4.004000in}}%
\pgfusepath{clip}%
\pgfsetrectcap%
\pgfsetroundjoin%
\pgfsetlinewidth{0.501875pt}%
\definecolor{currentstroke}{rgb}{0.690196,0.188235,0.376471}%
\pgfsetstrokecolor{currentstroke}%
\pgfsetstrokeopacity{0.800000}%
\pgfsetdash{}{0pt}%
\pgfpathmoveto{\pgfqpoint{0.455093in}{1.021044in}}%
\pgfpathlineto{\pgfqpoint{0.457612in}{1.020288in}}%
\pgfpathlineto{\pgfqpoint{0.460131in}{1.039779in}}%
\pgfpathlineto{\pgfqpoint{0.462649in}{1.045156in}}%
\pgfpathlineto{\pgfqpoint{0.465168in}{1.031182in}}%
\pgfpathlineto{\pgfqpoint{0.467687in}{1.015642in}}%
\pgfpathlineto{\pgfqpoint{0.470206in}{1.021806in}}%
\pgfpathlineto{\pgfqpoint{0.472724in}{1.034780in}}%
\pgfpathlineto{\pgfqpoint{0.475243in}{1.027109in}}%
\pgfpathlineto{\pgfqpoint{0.477762in}{1.016814in}}%
\pgfpathlineto{\pgfqpoint{0.480281in}{1.029938in}}%
\pgfpathlineto{\pgfqpoint{0.482799in}{1.042160in}}%
\pgfpathlineto{\pgfqpoint{0.485318in}{1.039528in}}%
\pgfpathlineto{\pgfqpoint{0.487837in}{1.038476in}}%
\pgfpathlineto{\pgfqpoint{0.490356in}{1.054486in}}%
\pgfpathlineto{\pgfqpoint{0.492874in}{1.062256in}}%
\pgfpathlineto{\pgfqpoint{0.495393in}{1.051092in}}%
\pgfpathlineto{\pgfqpoint{0.497912in}{1.029584in}}%
\pgfpathlineto{\pgfqpoint{0.500431in}{1.055708in}}%
\pgfpathlineto{\pgfqpoint{0.502949in}{1.058923in}}%
\pgfpathlineto{\pgfqpoint{0.505468in}{1.068067in}}%
\pgfpathlineto{\pgfqpoint{0.507987in}{1.078342in}}%
\pgfpathlineto{\pgfqpoint{0.510506in}{1.085599in}}%
\pgfpathlineto{\pgfqpoint{0.513024in}{1.061642in}}%
\pgfpathlineto{\pgfqpoint{0.515543in}{1.062741in}}%
\pgfpathlineto{\pgfqpoint{0.518062in}{1.073836in}}%
\pgfpathlineto{\pgfqpoint{0.520581in}{1.058800in}}%
\pgfpathlineto{\pgfqpoint{0.523099in}{1.073374in}}%
\pgfpathlineto{\pgfqpoint{0.525618in}{1.104722in}}%
\pgfpathlineto{\pgfqpoint{0.528137in}{1.087263in}}%
\pgfpathlineto{\pgfqpoint{0.530656in}{1.073477in}}%
\pgfpathlineto{\pgfqpoint{0.533174in}{1.066706in}}%
\pgfpathlineto{\pgfqpoint{0.535693in}{1.058172in}}%
\pgfpathlineto{\pgfqpoint{0.538212in}{1.067688in}}%
\pgfpathlineto{\pgfqpoint{0.540731in}{1.076245in}}%
\pgfpathlineto{\pgfqpoint{0.543249in}{1.094580in}}%
\pgfpathlineto{\pgfqpoint{0.545768in}{1.081528in}}%
\pgfpathlineto{\pgfqpoint{0.548287in}{1.087245in}}%
\pgfpathlineto{\pgfqpoint{0.550806in}{1.094400in}}%
\pgfpathlineto{\pgfqpoint{0.553324in}{1.074851in}}%
\pgfpathlineto{\pgfqpoint{0.555843in}{1.047638in}}%
\pgfpathlineto{\pgfqpoint{0.558362in}{1.046510in}}%
\pgfpathlineto{\pgfqpoint{0.560881in}{1.048125in}}%
\pgfpathlineto{\pgfqpoint{0.563399in}{1.062639in}}%
\pgfpathlineto{\pgfqpoint{0.565918in}{1.052438in}}%
\pgfpathlineto{\pgfqpoint{0.568437in}{1.063402in}}%
\pgfpathlineto{\pgfqpoint{0.570956in}{1.085423in}}%
\pgfpathlineto{\pgfqpoint{0.573474in}{1.073507in}}%
\pgfpathlineto{\pgfqpoint{0.575993in}{1.071021in}}%
\pgfpathlineto{\pgfqpoint{0.578512in}{1.058099in}}%
\pgfpathlineto{\pgfqpoint{0.581031in}{1.079796in}}%
\pgfpathlineto{\pgfqpoint{0.583549in}{1.061688in}}%
\pgfpathlineto{\pgfqpoint{0.586068in}{1.033140in}}%
\pgfpathlineto{\pgfqpoint{0.588587in}{1.045492in}}%
\pgfpathlineto{\pgfqpoint{0.591106in}{1.054568in}}%
\pgfpathlineto{\pgfqpoint{0.593624in}{1.045182in}}%
\pgfpathlineto{\pgfqpoint{0.596143in}{1.058925in}}%
\pgfpathlineto{\pgfqpoint{0.598662in}{1.038940in}}%
\pgfpathlineto{\pgfqpoint{0.601181in}{1.037079in}}%
\pgfpathlineto{\pgfqpoint{0.603699in}{1.025650in}}%
\pgfpathlineto{\pgfqpoint{0.606218in}{1.019038in}}%
\pgfpathlineto{\pgfqpoint{0.608737in}{1.039751in}}%
\pgfpathlineto{\pgfqpoint{0.611256in}{1.038852in}}%
\pgfpathlineto{\pgfqpoint{0.613774in}{1.042502in}}%
\pgfpathlineto{\pgfqpoint{0.616293in}{1.052772in}}%
\pgfpathlineto{\pgfqpoint{0.618812in}{1.056042in}}%
\pgfpathlineto{\pgfqpoint{0.621331in}{1.046626in}}%
\pgfpathlineto{\pgfqpoint{0.623849in}{1.021514in}}%
\pgfpathlineto{\pgfqpoint{0.626368in}{0.991288in}}%
\pgfpathlineto{\pgfqpoint{0.628887in}{0.999462in}}%
\pgfpathlineto{\pgfqpoint{0.631406in}{1.005886in}}%
\pgfpathlineto{\pgfqpoint{0.633924in}{1.009013in}}%
\pgfpathlineto{\pgfqpoint{0.636443in}{1.010328in}}%
\pgfpathlineto{\pgfqpoint{0.638962in}{1.019519in}}%
\pgfpathlineto{\pgfqpoint{0.641481in}{1.015494in}}%
\pgfpathlineto{\pgfqpoint{0.643999in}{1.028492in}}%
\pgfpathlineto{\pgfqpoint{0.646518in}{1.036710in}}%
\pgfpathlineto{\pgfqpoint{0.649037in}{1.027656in}}%
\pgfpathlineto{\pgfqpoint{0.651556in}{1.043307in}}%
\pgfpathlineto{\pgfqpoint{0.654074in}{1.029779in}}%
\pgfpathlineto{\pgfqpoint{0.656593in}{1.023669in}}%
\pgfpathlineto{\pgfqpoint{0.659112in}{1.028592in}}%
\pgfpathlineto{\pgfqpoint{0.661631in}{1.045575in}}%
\pgfpathlineto{\pgfqpoint{0.664149in}{1.058432in}}%
\pgfpathlineto{\pgfqpoint{0.666668in}{1.057516in}}%
\pgfpathlineto{\pgfqpoint{0.669187in}{1.029940in}}%
\pgfpathlineto{\pgfqpoint{0.671706in}{1.037755in}}%
\pgfpathlineto{\pgfqpoint{0.674224in}{1.031837in}}%
\pgfpathlineto{\pgfqpoint{0.676743in}{1.028353in}}%
\pgfpathlineto{\pgfqpoint{0.679262in}{1.057232in}}%
\pgfpathlineto{\pgfqpoint{0.681781in}{1.066836in}}%
\pgfpathlineto{\pgfqpoint{0.684299in}{1.061083in}}%
\pgfpathlineto{\pgfqpoint{0.686818in}{1.047626in}}%
\pgfpathlineto{\pgfqpoint{0.689337in}{1.063142in}}%
\pgfpathlineto{\pgfqpoint{0.691856in}{1.073141in}}%
\pgfpathlineto{\pgfqpoint{0.694374in}{1.059766in}}%
\pgfpathlineto{\pgfqpoint{0.696893in}{1.075844in}}%
\pgfpathlineto{\pgfqpoint{0.699412in}{1.088890in}}%
\pgfpathlineto{\pgfqpoint{0.701931in}{1.076520in}}%
\pgfpathlineto{\pgfqpoint{0.704449in}{1.059975in}}%
\pgfpathlineto{\pgfqpoint{0.706968in}{1.086107in}}%
\pgfpathlineto{\pgfqpoint{0.709487in}{1.078529in}}%
\pgfpathlineto{\pgfqpoint{0.712006in}{1.078608in}}%
\pgfpathlineto{\pgfqpoint{0.714524in}{1.084042in}}%
\pgfpathlineto{\pgfqpoint{0.717043in}{1.082718in}}%
\pgfpathlineto{\pgfqpoint{0.719562in}{1.089580in}}%
\pgfpathlineto{\pgfqpoint{0.722081in}{1.103426in}}%
\pgfpathlineto{\pgfqpoint{0.724599in}{1.093395in}}%
\pgfpathlineto{\pgfqpoint{0.727118in}{1.102163in}}%
\pgfpathlineto{\pgfqpoint{0.729637in}{1.091425in}}%
\pgfpathlineto{\pgfqpoint{0.732156in}{1.087368in}}%
\pgfpathlineto{\pgfqpoint{0.734674in}{1.106691in}}%
\pgfpathlineto{\pgfqpoint{0.737193in}{1.082932in}}%
\pgfpathlineto{\pgfqpoint{0.739712in}{1.083199in}}%
\pgfpathlineto{\pgfqpoint{0.742231in}{1.082019in}}%
\pgfpathlineto{\pgfqpoint{0.744749in}{1.091935in}}%
\pgfpathlineto{\pgfqpoint{0.747268in}{1.102919in}}%
\pgfpathlineto{\pgfqpoint{0.749787in}{1.092835in}}%
\pgfpathlineto{\pgfqpoint{0.752306in}{1.084098in}}%
\pgfpathlineto{\pgfqpoint{0.754824in}{1.079065in}}%
\pgfpathlineto{\pgfqpoint{0.757343in}{1.082672in}}%
\pgfpathlineto{\pgfqpoint{0.759862in}{1.113091in}}%
\pgfpathlineto{\pgfqpoint{0.762381in}{1.100621in}}%
\pgfpathlineto{\pgfqpoint{0.764899in}{1.089335in}}%
\pgfpathlineto{\pgfqpoint{0.767418in}{1.070672in}}%
\pgfpathlineto{\pgfqpoint{0.769937in}{1.063134in}}%
\pgfpathlineto{\pgfqpoint{0.772456in}{1.064152in}}%
\pgfpathlineto{\pgfqpoint{0.774974in}{1.069109in}}%
\pgfpathlineto{\pgfqpoint{0.777493in}{1.059927in}}%
\pgfpathlineto{\pgfqpoint{0.780012in}{1.060554in}}%
\pgfpathlineto{\pgfqpoint{0.782531in}{1.059652in}}%
\pgfpathlineto{\pgfqpoint{0.785049in}{1.052723in}}%
\pgfpathlineto{\pgfqpoint{0.787568in}{1.068512in}}%
\pgfpathlineto{\pgfqpoint{0.790087in}{1.075204in}}%
\pgfpathlineto{\pgfqpoint{0.792606in}{1.095872in}}%
\pgfpathlineto{\pgfqpoint{0.795124in}{1.100216in}}%
\pgfpathlineto{\pgfqpoint{0.797643in}{1.127929in}}%
\pgfpathlineto{\pgfqpoint{0.800162in}{1.141512in}}%
\pgfpathlineto{\pgfqpoint{0.802681in}{1.150303in}}%
\pgfpathlineto{\pgfqpoint{0.805199in}{1.184906in}}%
\pgfpathlineto{\pgfqpoint{0.807718in}{1.179271in}}%
\pgfpathlineto{\pgfqpoint{0.810237in}{1.178274in}}%
\pgfpathlineto{\pgfqpoint{0.812756in}{1.187491in}}%
\pgfpathlineto{\pgfqpoint{0.815274in}{1.174133in}}%
\pgfpathlineto{\pgfqpoint{0.817793in}{1.168582in}}%
\pgfpathlineto{\pgfqpoint{0.820312in}{1.147109in}}%
\pgfpathlineto{\pgfqpoint{0.822831in}{1.147626in}}%
\pgfpathlineto{\pgfqpoint{0.825349in}{1.123095in}}%
\pgfpathlineto{\pgfqpoint{0.827868in}{1.140682in}}%
\pgfpathlineto{\pgfqpoint{0.830387in}{1.144808in}}%
\pgfpathlineto{\pgfqpoint{0.832906in}{1.159820in}}%
\pgfpathlineto{\pgfqpoint{0.835424in}{1.165056in}}%
\pgfpathlineto{\pgfqpoint{0.837943in}{1.190716in}}%
\pgfpathlineto{\pgfqpoint{0.840462in}{1.169922in}}%
\pgfpathlineto{\pgfqpoint{0.842981in}{1.172157in}}%
\pgfpathlineto{\pgfqpoint{0.845499in}{1.150991in}}%
\pgfpathlineto{\pgfqpoint{0.848018in}{1.142871in}}%
\pgfpathlineto{\pgfqpoint{0.850537in}{1.135470in}}%
\pgfpathlineto{\pgfqpoint{0.853056in}{1.157330in}}%
\pgfpathlineto{\pgfqpoint{0.855574in}{1.149786in}}%
\pgfpathlineto{\pgfqpoint{0.858093in}{1.161595in}}%
\pgfpathlineto{\pgfqpoint{0.860612in}{1.154589in}}%
\pgfpathlineto{\pgfqpoint{0.863131in}{1.140758in}}%
\pgfpathlineto{\pgfqpoint{0.865649in}{1.129896in}}%
\pgfpathlineto{\pgfqpoint{0.868168in}{1.132815in}}%
\pgfpathlineto{\pgfqpoint{0.870687in}{1.129474in}}%
\pgfpathlineto{\pgfqpoint{0.873206in}{1.145528in}}%
\pgfpathlineto{\pgfqpoint{0.875724in}{1.157385in}}%
\pgfpathlineto{\pgfqpoint{0.878243in}{1.138373in}}%
\pgfpathlineto{\pgfqpoint{0.880762in}{1.146215in}}%
\pgfpathlineto{\pgfqpoint{0.883281in}{1.138147in}}%
\pgfpathlineto{\pgfqpoint{0.885799in}{1.146628in}}%
\pgfpathlineto{\pgfqpoint{0.888318in}{1.139289in}}%
\pgfpathlineto{\pgfqpoint{0.890837in}{1.148515in}}%
\pgfpathlineto{\pgfqpoint{0.893356in}{1.169488in}}%
\pgfpathlineto{\pgfqpoint{0.895874in}{1.181105in}}%
\pgfpathlineto{\pgfqpoint{0.898393in}{1.161723in}}%
\pgfpathlineto{\pgfqpoint{0.900912in}{1.151938in}}%
\pgfpathlineto{\pgfqpoint{0.903431in}{1.146476in}}%
\pgfpathlineto{\pgfqpoint{0.905949in}{1.134716in}}%
\pgfpathlineto{\pgfqpoint{0.908468in}{1.123402in}}%
\pgfpathlineto{\pgfqpoint{0.910987in}{1.155985in}}%
\pgfpathlineto{\pgfqpoint{0.913506in}{1.173336in}}%
\pgfpathlineto{\pgfqpoint{0.916024in}{1.158881in}}%
\pgfpathlineto{\pgfqpoint{0.918543in}{1.147899in}}%
\pgfpathlineto{\pgfqpoint{0.921062in}{1.143020in}}%
\pgfpathlineto{\pgfqpoint{0.923581in}{1.158113in}}%
\pgfpathlineto{\pgfqpoint{0.926099in}{1.170936in}}%
\pgfpathlineto{\pgfqpoint{0.928618in}{1.175358in}}%
\pgfpathlineto{\pgfqpoint{0.931137in}{1.172596in}}%
\pgfpathlineto{\pgfqpoint{0.933656in}{1.169124in}}%
\pgfpathlineto{\pgfqpoint{0.936174in}{1.167452in}}%
\pgfpathlineto{\pgfqpoint{0.938693in}{1.182284in}}%
\pgfpathlineto{\pgfqpoint{0.941212in}{1.159006in}}%
\pgfpathlineto{\pgfqpoint{0.943731in}{1.163205in}}%
\pgfpathlineto{\pgfqpoint{0.946249in}{1.145085in}}%
\pgfpathlineto{\pgfqpoint{0.948768in}{1.138578in}}%
\pgfpathlineto{\pgfqpoint{0.951287in}{1.180012in}}%
\pgfpathlineto{\pgfqpoint{0.953806in}{1.189903in}}%
\pgfpathlineto{\pgfqpoint{0.956324in}{1.194869in}}%
\pgfpathlineto{\pgfqpoint{0.958843in}{1.198288in}}%
\pgfpathlineto{\pgfqpoint{0.961362in}{1.213774in}}%
\pgfpathlineto{\pgfqpoint{0.963881in}{1.217021in}}%
\pgfpathlineto{\pgfqpoint{0.966399in}{1.203458in}}%
\pgfpathlineto{\pgfqpoint{0.968918in}{1.207532in}}%
\pgfpathlineto{\pgfqpoint{0.971437in}{1.193935in}}%
\pgfpathlineto{\pgfqpoint{0.973956in}{1.200583in}}%
\pgfpathlineto{\pgfqpoint{0.976474in}{1.197997in}}%
\pgfpathlineto{\pgfqpoint{0.978993in}{1.194511in}}%
\pgfpathlineto{\pgfqpoint{0.981512in}{1.201612in}}%
\pgfpathlineto{\pgfqpoint{0.984031in}{1.197841in}}%
\pgfpathlineto{\pgfqpoint{0.986549in}{1.204022in}}%
\pgfpathlineto{\pgfqpoint{0.989068in}{1.218211in}}%
\pgfpathlineto{\pgfqpoint{0.991587in}{1.219492in}}%
\pgfpathlineto{\pgfqpoint{0.994106in}{1.216550in}}%
\pgfpathlineto{\pgfqpoint{0.996624in}{1.229739in}}%
\pgfpathlineto{\pgfqpoint{0.999143in}{1.216237in}}%
\pgfpathlineto{\pgfqpoint{1.001662in}{1.227536in}}%
\pgfpathlineto{\pgfqpoint{1.004181in}{1.228857in}}%
\pgfpathlineto{\pgfqpoint{1.006699in}{1.228703in}}%
\pgfpathlineto{\pgfqpoint{1.009218in}{1.236680in}}%
\pgfpathlineto{\pgfqpoint{1.011737in}{1.237633in}}%
\pgfpathlineto{\pgfqpoint{1.014256in}{1.244298in}}%
\pgfpathlineto{\pgfqpoint{1.016774in}{1.240303in}}%
\pgfpathlineto{\pgfqpoint{1.019293in}{1.246940in}}%
\pgfpathlineto{\pgfqpoint{1.021812in}{1.275353in}}%
\pgfpathlineto{\pgfqpoint{1.024331in}{1.278315in}}%
\pgfpathlineto{\pgfqpoint{1.026849in}{1.273082in}}%
\pgfpathlineto{\pgfqpoint{1.029368in}{1.272132in}}%
\pgfpathlineto{\pgfqpoint{1.031887in}{1.281753in}}%
\pgfpathlineto{\pgfqpoint{1.034406in}{1.276864in}}%
\pgfpathlineto{\pgfqpoint{1.036924in}{1.291299in}}%
\pgfpathlineto{\pgfqpoint{1.039443in}{1.310951in}}%
\pgfpathlineto{\pgfqpoint{1.041962in}{1.316207in}}%
\pgfpathlineto{\pgfqpoint{1.044481in}{1.304640in}}%
\pgfpathlineto{\pgfqpoint{1.046999in}{1.315706in}}%
\pgfpathlineto{\pgfqpoint{1.049518in}{1.341546in}}%
\pgfpathlineto{\pgfqpoint{1.052037in}{1.330370in}}%
\pgfpathlineto{\pgfqpoint{1.054556in}{1.302924in}}%
\pgfpathlineto{\pgfqpoint{1.057074in}{1.336253in}}%
\pgfpathlineto{\pgfqpoint{1.059593in}{1.364513in}}%
\pgfpathlineto{\pgfqpoint{1.062112in}{1.341439in}}%
\pgfpathlineto{\pgfqpoint{1.064631in}{1.355614in}}%
\pgfpathlineto{\pgfqpoint{1.067149in}{1.361747in}}%
\pgfpathlineto{\pgfqpoint{1.069668in}{1.357679in}}%
\pgfpathlineto{\pgfqpoint{1.072187in}{1.367232in}}%
\pgfpathlineto{\pgfqpoint{1.074706in}{1.379154in}}%
\pgfpathlineto{\pgfqpoint{1.077224in}{1.366683in}}%
\pgfpathlineto{\pgfqpoint{1.079743in}{1.364654in}}%
\pgfpathlineto{\pgfqpoint{1.082262in}{1.322970in}}%
\pgfpathlineto{\pgfqpoint{1.084781in}{1.315828in}}%
\pgfpathlineto{\pgfqpoint{1.087299in}{1.322031in}}%
\pgfpathlineto{\pgfqpoint{1.092337in}{1.357573in}}%
\pgfpathlineto{\pgfqpoint{1.094856in}{1.346561in}}%
\pgfpathlineto{\pgfqpoint{1.097374in}{1.357751in}}%
\pgfpathlineto{\pgfqpoint{1.099893in}{1.351108in}}%
\pgfpathlineto{\pgfqpoint{1.102412in}{1.340358in}}%
\pgfpathlineto{\pgfqpoint{1.104931in}{1.321284in}}%
\pgfpathlineto{\pgfqpoint{1.107449in}{1.313133in}}%
\pgfpathlineto{\pgfqpoint{1.109968in}{1.320552in}}%
\pgfpathlineto{\pgfqpoint{1.112487in}{1.330599in}}%
\pgfpathlineto{\pgfqpoint{1.115006in}{1.314613in}}%
\pgfpathlineto{\pgfqpoint{1.117524in}{1.300895in}}%
\pgfpathlineto{\pgfqpoint{1.120043in}{1.300526in}}%
\pgfpathlineto{\pgfqpoint{1.122562in}{1.312329in}}%
\pgfpathlineto{\pgfqpoint{1.125081in}{1.346825in}}%
\pgfpathlineto{\pgfqpoint{1.127599in}{1.359277in}}%
\pgfpathlineto{\pgfqpoint{1.130118in}{1.351664in}}%
\pgfpathlineto{\pgfqpoint{1.132637in}{1.367839in}}%
\pgfpathlineto{\pgfqpoint{1.135156in}{1.372113in}}%
\pgfpathlineto{\pgfqpoint{1.137674in}{1.364345in}}%
\pgfpathlineto{\pgfqpoint{1.140193in}{1.359056in}}%
\pgfpathlineto{\pgfqpoint{1.142712in}{1.362189in}}%
\pgfpathlineto{\pgfqpoint{1.145231in}{1.383889in}}%
\pgfpathlineto{\pgfqpoint{1.147749in}{1.381665in}}%
\pgfpathlineto{\pgfqpoint{1.150268in}{1.406529in}}%
\pgfpathlineto{\pgfqpoint{1.152787in}{1.410621in}}%
\pgfpathlineto{\pgfqpoint{1.155306in}{1.422185in}}%
\pgfpathlineto{\pgfqpoint{1.157824in}{1.418024in}}%
\pgfpathlineto{\pgfqpoint{1.160343in}{1.429443in}}%
\pgfpathlineto{\pgfqpoint{1.162862in}{1.420648in}}%
\pgfpathlineto{\pgfqpoint{1.165381in}{1.417477in}}%
\pgfpathlineto{\pgfqpoint{1.167899in}{1.435098in}}%
\pgfpathlineto{\pgfqpoint{1.170418in}{1.449204in}}%
\pgfpathlineto{\pgfqpoint{1.172937in}{1.454810in}}%
\pgfpathlineto{\pgfqpoint{1.175456in}{1.456511in}}%
\pgfpathlineto{\pgfqpoint{1.177974in}{1.463930in}}%
\pgfpathlineto{\pgfqpoint{1.180493in}{1.473709in}}%
\pgfpathlineto{\pgfqpoint{1.183012in}{1.480880in}}%
\pgfpathlineto{\pgfqpoint{1.185531in}{1.444092in}}%
\pgfpathlineto{\pgfqpoint{1.188049in}{1.477773in}}%
\pgfpathlineto{\pgfqpoint{1.190568in}{1.479455in}}%
\pgfpathlineto{\pgfqpoint{1.193087in}{1.472669in}}%
\pgfpathlineto{\pgfqpoint{1.195606in}{1.484042in}}%
\pgfpathlineto{\pgfqpoint{1.198124in}{1.492542in}}%
\pgfpathlineto{\pgfqpoint{1.200643in}{1.483587in}}%
\pgfpathlineto{\pgfqpoint{1.203162in}{1.477263in}}%
\pgfpathlineto{\pgfqpoint{1.205681in}{1.465946in}}%
\pgfpathlineto{\pgfqpoint{1.208199in}{1.455395in}}%
\pgfpathlineto{\pgfqpoint{1.210718in}{1.473969in}}%
\pgfpathlineto{\pgfqpoint{1.213237in}{1.483787in}}%
\pgfpathlineto{\pgfqpoint{1.215756in}{1.494280in}}%
\pgfpathlineto{\pgfqpoint{1.218274in}{1.509950in}}%
\pgfpathlineto{\pgfqpoint{1.220793in}{1.506253in}}%
\pgfpathlineto{\pgfqpoint{1.223312in}{1.488998in}}%
\pgfpathlineto{\pgfqpoint{1.225831in}{1.492010in}}%
\pgfpathlineto{\pgfqpoint{1.228349in}{1.495330in}}%
\pgfpathlineto{\pgfqpoint{1.230868in}{1.493935in}}%
\pgfpathlineto{\pgfqpoint{1.233387in}{1.469465in}}%
\pgfpathlineto{\pgfqpoint{1.235906in}{1.469775in}}%
\pgfpathlineto{\pgfqpoint{1.238424in}{1.494321in}}%
\pgfpathlineto{\pgfqpoint{1.240943in}{1.500044in}}%
\pgfpathlineto{\pgfqpoint{1.243462in}{1.502245in}}%
\pgfpathlineto{\pgfqpoint{1.245981in}{1.464069in}}%
\pgfpathlineto{\pgfqpoint{1.248499in}{1.454069in}}%
\pgfpathlineto{\pgfqpoint{1.251018in}{1.466681in}}%
\pgfpathlineto{\pgfqpoint{1.253537in}{1.448098in}}%
\pgfpathlineto{\pgfqpoint{1.256056in}{1.445651in}}%
\pgfpathlineto{\pgfqpoint{1.258574in}{1.460183in}}%
\pgfpathlineto{\pgfqpoint{1.261093in}{1.459169in}}%
\pgfpathlineto{\pgfqpoint{1.263612in}{1.441121in}}%
\pgfpathlineto{\pgfqpoint{1.266131in}{1.454435in}}%
\pgfpathlineto{\pgfqpoint{1.268649in}{1.444190in}}%
\pgfpathlineto{\pgfqpoint{1.271168in}{1.446847in}}%
\pgfpathlineto{\pgfqpoint{1.273687in}{1.474172in}}%
\pgfpathlineto{\pgfqpoint{1.276206in}{1.490184in}}%
\pgfpathlineto{\pgfqpoint{1.278724in}{1.488370in}}%
\pgfpathlineto{\pgfqpoint{1.281243in}{1.487820in}}%
\pgfpathlineto{\pgfqpoint{1.283762in}{1.465698in}}%
\pgfpathlineto{\pgfqpoint{1.286281in}{1.493377in}}%
\pgfpathlineto{\pgfqpoint{1.288799in}{1.506873in}}%
\pgfpathlineto{\pgfqpoint{1.291318in}{1.504246in}}%
\pgfpathlineto{\pgfqpoint{1.293837in}{1.511374in}}%
\pgfpathlineto{\pgfqpoint{1.296356in}{1.522055in}}%
\pgfpathlineto{\pgfqpoint{1.298874in}{1.521946in}}%
\pgfpathlineto{\pgfqpoint{1.301393in}{1.516647in}}%
\pgfpathlineto{\pgfqpoint{1.303912in}{1.499548in}}%
\pgfpathlineto{\pgfqpoint{1.306431in}{1.489136in}}%
\pgfpathlineto{\pgfqpoint{1.308949in}{1.472915in}}%
\pgfpathlineto{\pgfqpoint{1.311468in}{1.495954in}}%
\pgfpathlineto{\pgfqpoint{1.313987in}{1.512619in}}%
\pgfpathlineto{\pgfqpoint{1.316506in}{1.531470in}}%
\pgfpathlineto{\pgfqpoint{1.319024in}{1.537959in}}%
\pgfpathlineto{\pgfqpoint{1.321543in}{1.555695in}}%
\pgfpathlineto{\pgfqpoint{1.324062in}{1.587433in}}%
\pgfpathlineto{\pgfqpoint{1.326581in}{1.603691in}}%
\pgfpathlineto{\pgfqpoint{1.329099in}{1.604012in}}%
\pgfpathlineto{\pgfqpoint{1.331618in}{1.613957in}}%
\pgfpathlineto{\pgfqpoint{1.334137in}{1.606206in}}%
\pgfpathlineto{\pgfqpoint{1.336656in}{1.612147in}}%
\pgfpathlineto{\pgfqpoint{1.339174in}{1.591729in}}%
\pgfpathlineto{\pgfqpoint{1.341693in}{1.587192in}}%
\pgfpathlineto{\pgfqpoint{1.344212in}{1.609846in}}%
\pgfpathlineto{\pgfqpoint{1.346731in}{1.590589in}}%
\pgfpathlineto{\pgfqpoint{1.349249in}{1.610839in}}%
\pgfpathlineto{\pgfqpoint{1.351768in}{1.603115in}}%
\pgfpathlineto{\pgfqpoint{1.354287in}{1.613939in}}%
\pgfpathlineto{\pgfqpoint{1.356806in}{1.603838in}}%
\pgfpathlineto{\pgfqpoint{1.359324in}{1.584802in}}%
\pgfpathlineto{\pgfqpoint{1.361843in}{1.575731in}}%
\pgfpathlineto{\pgfqpoint{1.364362in}{1.570279in}}%
\pgfpathlineto{\pgfqpoint{1.366881in}{1.582795in}}%
\pgfpathlineto{\pgfqpoint{1.369399in}{1.553054in}}%
\pgfpathlineto{\pgfqpoint{1.371918in}{1.537259in}}%
\pgfpathlineto{\pgfqpoint{1.374437in}{1.526802in}}%
\pgfpathlineto{\pgfqpoint{1.376956in}{1.535969in}}%
\pgfpathlineto{\pgfqpoint{1.379474in}{1.535493in}}%
\pgfpathlineto{\pgfqpoint{1.381993in}{1.538123in}}%
\pgfpathlineto{\pgfqpoint{1.384512in}{1.571131in}}%
\pgfpathlineto{\pgfqpoint{1.387031in}{1.589864in}}%
\pgfpathlineto{\pgfqpoint{1.389549in}{1.577923in}}%
\pgfpathlineto{\pgfqpoint{1.392068in}{1.578797in}}%
\pgfpathlineto{\pgfqpoint{1.394587in}{1.588654in}}%
\pgfpathlineto{\pgfqpoint{1.397106in}{1.570843in}}%
\pgfpathlineto{\pgfqpoint{1.399624in}{1.556115in}}%
\pgfpathlineto{\pgfqpoint{1.402143in}{1.572057in}}%
\pgfpathlineto{\pgfqpoint{1.404662in}{1.587241in}}%
\pgfpathlineto{\pgfqpoint{1.407181in}{1.564132in}}%
\pgfpathlineto{\pgfqpoint{1.409699in}{1.585166in}}%
\pgfpathlineto{\pgfqpoint{1.412218in}{1.598677in}}%
\pgfpathlineto{\pgfqpoint{1.414737in}{1.596052in}}%
\pgfpathlineto{\pgfqpoint{1.417256in}{1.586787in}}%
\pgfpathlineto{\pgfqpoint{1.419774in}{1.617133in}}%
\pgfpathlineto{\pgfqpoint{1.422293in}{1.620119in}}%
\pgfpathlineto{\pgfqpoint{1.424812in}{1.620840in}}%
\pgfpathlineto{\pgfqpoint{1.427331in}{1.621973in}}%
\pgfpathlineto{\pgfqpoint{1.429849in}{1.613772in}}%
\pgfpathlineto{\pgfqpoint{1.432368in}{1.603528in}}%
\pgfpathlineto{\pgfqpoint{1.434887in}{1.606338in}}%
\pgfpathlineto{\pgfqpoint{1.437406in}{1.627577in}}%
\pgfpathlineto{\pgfqpoint{1.442443in}{1.685499in}}%
\pgfpathlineto{\pgfqpoint{1.444962in}{1.649869in}}%
\pgfpathlineto{\pgfqpoint{1.447481in}{1.668889in}}%
\pgfpathlineto{\pgfqpoint{1.449999in}{1.655842in}}%
\pgfpathlineto{\pgfqpoint{1.452518in}{1.657441in}}%
\pgfpathlineto{\pgfqpoint{1.455037in}{1.664558in}}%
\pgfpathlineto{\pgfqpoint{1.457556in}{1.662463in}}%
\pgfpathlineto{\pgfqpoint{1.460074in}{1.660154in}}%
\pgfpathlineto{\pgfqpoint{1.462593in}{1.660961in}}%
\pgfpathlineto{\pgfqpoint{1.465112in}{1.638938in}}%
\pgfpathlineto{\pgfqpoint{1.467631in}{1.610621in}}%
\pgfpathlineto{\pgfqpoint{1.470149in}{1.614447in}}%
\pgfpathlineto{\pgfqpoint{1.472668in}{1.628478in}}%
\pgfpathlineto{\pgfqpoint{1.475187in}{1.607079in}}%
\pgfpathlineto{\pgfqpoint{1.477706in}{1.609234in}}%
\pgfpathlineto{\pgfqpoint{1.480224in}{1.623644in}}%
\pgfpathlineto{\pgfqpoint{1.482743in}{1.616978in}}%
\pgfpathlineto{\pgfqpoint{1.485262in}{1.625777in}}%
\pgfpathlineto{\pgfqpoint{1.487781in}{1.668134in}}%
\pgfpathlineto{\pgfqpoint{1.490299in}{1.670217in}}%
\pgfpathlineto{\pgfqpoint{1.492818in}{1.697132in}}%
\pgfpathlineto{\pgfqpoint{1.495337in}{1.686089in}}%
\pgfpathlineto{\pgfqpoint{1.497856in}{1.698629in}}%
\pgfpathlineto{\pgfqpoint{1.500374in}{1.715261in}}%
\pgfpathlineto{\pgfqpoint{1.502893in}{1.709010in}}%
\pgfpathlineto{\pgfqpoint{1.505412in}{1.721305in}}%
\pgfpathlineto{\pgfqpoint{1.507931in}{1.716212in}}%
\pgfpathlineto{\pgfqpoint{1.510449in}{1.701747in}}%
\pgfpathlineto{\pgfqpoint{1.512968in}{1.701913in}}%
\pgfpathlineto{\pgfqpoint{1.515487in}{1.682649in}}%
\pgfpathlineto{\pgfqpoint{1.518006in}{1.677228in}}%
\pgfpathlineto{\pgfqpoint{1.520524in}{1.685940in}}%
\pgfpathlineto{\pgfqpoint{1.523043in}{1.686631in}}%
\pgfpathlineto{\pgfqpoint{1.525562in}{1.685385in}}%
\pgfpathlineto{\pgfqpoint{1.528081in}{1.682238in}}%
\pgfpathlineto{\pgfqpoint{1.530599in}{1.676004in}}%
\pgfpathlineto{\pgfqpoint{1.533118in}{1.667826in}}%
\pgfpathlineto{\pgfqpoint{1.535637in}{1.657103in}}%
\pgfpathlineto{\pgfqpoint{1.538156in}{1.636233in}}%
\pgfpathlineto{\pgfqpoint{1.540674in}{1.623254in}}%
\pgfpathlineto{\pgfqpoint{1.543193in}{1.640893in}}%
\pgfpathlineto{\pgfqpoint{1.545712in}{1.632918in}}%
\pgfpathlineto{\pgfqpoint{1.548231in}{1.679103in}}%
\pgfpathlineto{\pgfqpoint{1.550749in}{1.690851in}}%
\pgfpathlineto{\pgfqpoint{1.553268in}{1.681478in}}%
\pgfpathlineto{\pgfqpoint{1.555787in}{1.714002in}}%
\pgfpathlineto{\pgfqpoint{1.558306in}{1.719863in}}%
\pgfpathlineto{\pgfqpoint{1.560824in}{1.745573in}}%
\pgfpathlineto{\pgfqpoint{1.563343in}{1.766287in}}%
\pgfpathlineto{\pgfqpoint{1.565862in}{1.744990in}}%
\pgfpathlineto{\pgfqpoint{1.568381in}{1.719263in}}%
\pgfpathlineto{\pgfqpoint{1.570899in}{1.703838in}}%
\pgfpathlineto{\pgfqpoint{1.573418in}{1.676325in}}%
\pgfpathlineto{\pgfqpoint{1.575937in}{1.696499in}}%
\pgfpathlineto{\pgfqpoint{1.578456in}{1.713340in}}%
\pgfpathlineto{\pgfqpoint{1.580974in}{1.740424in}}%
\pgfpathlineto{\pgfqpoint{1.583493in}{1.757700in}}%
\pgfpathlineto{\pgfqpoint{1.586012in}{1.785223in}}%
\pgfpathlineto{\pgfqpoint{1.588531in}{1.767228in}}%
\pgfpathlineto{\pgfqpoint{1.591049in}{1.764880in}}%
\pgfpathlineto{\pgfqpoint{1.593568in}{1.767133in}}%
\pgfpathlineto{\pgfqpoint{1.596087in}{1.744300in}}%
\pgfpathlineto{\pgfqpoint{1.598606in}{1.743377in}}%
\pgfpathlineto{\pgfqpoint{1.601124in}{1.761240in}}%
\pgfpathlineto{\pgfqpoint{1.603643in}{1.760860in}}%
\pgfpathlineto{\pgfqpoint{1.606162in}{1.758588in}}%
\pgfpathlineto{\pgfqpoint{1.608681in}{1.757523in}}%
\pgfpathlineto{\pgfqpoint{1.611199in}{1.759563in}}%
\pgfpathlineto{\pgfqpoint{1.613718in}{1.743112in}}%
\pgfpathlineto{\pgfqpoint{1.616237in}{1.744586in}}%
\pgfpathlineto{\pgfqpoint{1.618756in}{1.782916in}}%
\pgfpathlineto{\pgfqpoint{1.621274in}{1.788691in}}%
\pgfpathlineto{\pgfqpoint{1.623793in}{1.775555in}}%
\pgfpathlineto{\pgfqpoint{1.626312in}{1.785198in}}%
\pgfpathlineto{\pgfqpoint{1.628831in}{1.786987in}}%
\pgfpathlineto{\pgfqpoint{1.631349in}{1.790159in}}%
\pgfpathlineto{\pgfqpoint{1.633868in}{1.810429in}}%
\pgfpathlineto{\pgfqpoint{1.636387in}{1.811588in}}%
\pgfpathlineto{\pgfqpoint{1.638906in}{1.833605in}}%
\pgfpathlineto{\pgfqpoint{1.641424in}{1.809985in}}%
\pgfpathlineto{\pgfqpoint{1.643943in}{1.800578in}}%
\pgfpathlineto{\pgfqpoint{1.646462in}{1.780366in}}%
\pgfpathlineto{\pgfqpoint{1.648981in}{1.797036in}}%
\pgfpathlineto{\pgfqpoint{1.651499in}{1.780589in}}%
\pgfpathlineto{\pgfqpoint{1.654018in}{1.784289in}}%
\pgfpathlineto{\pgfqpoint{1.656537in}{1.826916in}}%
\pgfpathlineto{\pgfqpoint{1.659056in}{1.874636in}}%
\pgfpathlineto{\pgfqpoint{1.661574in}{1.880493in}}%
\pgfpathlineto{\pgfqpoint{1.664093in}{1.846830in}}%
\pgfpathlineto{\pgfqpoint{1.666612in}{1.840243in}}%
\pgfpathlineto{\pgfqpoint{1.669131in}{1.821423in}}%
\pgfpathlineto{\pgfqpoint{1.671649in}{1.813897in}}%
\pgfpathlineto{\pgfqpoint{1.674168in}{1.807904in}}%
\pgfpathlineto{\pgfqpoint{1.676687in}{1.817890in}}%
\pgfpathlineto{\pgfqpoint{1.679206in}{1.820521in}}%
\pgfpathlineto{\pgfqpoint{1.681724in}{1.844995in}}%
\pgfpathlineto{\pgfqpoint{1.684243in}{1.825085in}}%
\pgfpathlineto{\pgfqpoint{1.686762in}{1.849433in}}%
\pgfpathlineto{\pgfqpoint{1.689281in}{1.851218in}}%
\pgfpathlineto{\pgfqpoint{1.691799in}{1.846597in}}%
\pgfpathlineto{\pgfqpoint{1.694318in}{1.837035in}}%
\pgfpathlineto{\pgfqpoint{1.696837in}{1.837616in}}%
\pgfpathlineto{\pgfqpoint{1.699356in}{1.836350in}}%
\pgfpathlineto{\pgfqpoint{1.701874in}{1.830318in}}%
\pgfpathlineto{\pgfqpoint{1.704393in}{1.827687in}}%
\pgfpathlineto{\pgfqpoint{1.706912in}{1.816458in}}%
\pgfpathlineto{\pgfqpoint{1.709431in}{1.809643in}}%
\pgfpathlineto{\pgfqpoint{1.711949in}{1.846091in}}%
\pgfpathlineto{\pgfqpoint{1.714468in}{1.836608in}}%
\pgfpathlineto{\pgfqpoint{1.716987in}{1.835557in}}%
\pgfpathlineto{\pgfqpoint{1.719506in}{1.822140in}}%
\pgfpathlineto{\pgfqpoint{1.722024in}{1.820975in}}%
\pgfpathlineto{\pgfqpoint{1.724543in}{1.841643in}}%
\pgfpathlineto{\pgfqpoint{1.727062in}{1.842784in}}%
\pgfpathlineto{\pgfqpoint{1.729581in}{1.836396in}}%
\pgfpathlineto{\pgfqpoint{1.732099in}{1.812989in}}%
\pgfpathlineto{\pgfqpoint{1.734618in}{1.785323in}}%
\pgfpathlineto{\pgfqpoint{1.737137in}{1.794572in}}%
\pgfpathlineto{\pgfqpoint{1.739656in}{1.794501in}}%
\pgfpathlineto{\pgfqpoint{1.742174in}{1.769887in}}%
\pgfpathlineto{\pgfqpoint{1.744693in}{1.773193in}}%
\pgfpathlineto{\pgfqpoint{1.747212in}{1.777523in}}%
\pgfpathlineto{\pgfqpoint{1.749731in}{1.797426in}}%
\pgfpathlineto{\pgfqpoint{1.752249in}{1.770165in}}%
\pgfpathlineto{\pgfqpoint{1.754768in}{1.769865in}}%
\pgfpathlineto{\pgfqpoint{1.757287in}{1.772393in}}%
\pgfpathlineto{\pgfqpoint{1.759806in}{1.772601in}}%
\pgfpathlineto{\pgfqpoint{1.762324in}{1.782289in}}%
\pgfpathlineto{\pgfqpoint{1.764843in}{1.788423in}}%
\pgfpathlineto{\pgfqpoint{1.767362in}{1.795582in}}%
\pgfpathlineto{\pgfqpoint{1.769881in}{1.820110in}}%
\pgfpathlineto{\pgfqpoint{1.772399in}{1.838239in}}%
\pgfpathlineto{\pgfqpoint{1.774918in}{1.846366in}}%
\pgfpathlineto{\pgfqpoint{1.777437in}{1.859556in}}%
\pgfpathlineto{\pgfqpoint{1.779956in}{1.863712in}}%
\pgfpathlineto{\pgfqpoint{1.782474in}{1.869297in}}%
\pgfpathlineto{\pgfqpoint{1.784993in}{1.879357in}}%
\pgfpathlineto{\pgfqpoint{1.787512in}{1.864648in}}%
\pgfpathlineto{\pgfqpoint{1.790031in}{1.865417in}}%
\pgfpathlineto{\pgfqpoint{1.792549in}{1.890117in}}%
\pgfpathlineto{\pgfqpoint{1.795068in}{1.888592in}}%
\pgfpathlineto{\pgfqpoint{1.797587in}{1.920751in}}%
\pgfpathlineto{\pgfqpoint{1.800106in}{1.903788in}}%
\pgfpathlineto{\pgfqpoint{1.802624in}{1.894394in}}%
\pgfpathlineto{\pgfqpoint{1.805143in}{1.914377in}}%
\pgfpathlineto{\pgfqpoint{1.807662in}{1.893937in}}%
\pgfpathlineto{\pgfqpoint{1.810181in}{1.880790in}}%
\pgfpathlineto{\pgfqpoint{1.812699in}{1.860108in}}%
\pgfpathlineto{\pgfqpoint{1.815218in}{1.856207in}}%
\pgfpathlineto{\pgfqpoint{1.817737in}{1.856719in}}%
\pgfpathlineto{\pgfqpoint{1.820256in}{1.861449in}}%
\pgfpathlineto{\pgfqpoint{1.822774in}{1.843215in}}%
\pgfpathlineto{\pgfqpoint{1.825293in}{1.832965in}}%
\pgfpathlineto{\pgfqpoint{1.827812in}{1.850853in}}%
\pgfpathlineto{\pgfqpoint{1.830331in}{1.825389in}}%
\pgfpathlineto{\pgfqpoint{1.832849in}{1.819406in}}%
\pgfpathlineto{\pgfqpoint{1.835368in}{1.818488in}}%
\pgfpathlineto{\pgfqpoint{1.837887in}{1.826356in}}%
\pgfpathlineto{\pgfqpoint{1.840406in}{1.832703in}}%
\pgfpathlineto{\pgfqpoint{1.842924in}{1.835897in}}%
\pgfpathlineto{\pgfqpoint{1.845443in}{1.822990in}}%
\pgfpathlineto{\pgfqpoint{1.847962in}{1.818309in}}%
\pgfpathlineto{\pgfqpoint{1.850481in}{1.826599in}}%
\pgfpathlineto{\pgfqpoint{1.852999in}{1.830070in}}%
\pgfpathlineto{\pgfqpoint{1.855518in}{1.819021in}}%
\pgfpathlineto{\pgfqpoint{1.858037in}{1.796578in}}%
\pgfpathlineto{\pgfqpoint{1.860556in}{1.787538in}}%
\pgfpathlineto{\pgfqpoint{1.863074in}{1.797861in}}%
\pgfpathlineto{\pgfqpoint{1.865593in}{1.793495in}}%
\pgfpathlineto{\pgfqpoint{1.868112in}{1.830118in}}%
\pgfpathlineto{\pgfqpoint{1.870631in}{1.830794in}}%
\pgfpathlineto{\pgfqpoint{1.873149in}{1.829259in}}%
\pgfpathlineto{\pgfqpoint{1.875668in}{1.791445in}}%
\pgfpathlineto{\pgfqpoint{1.878187in}{1.799612in}}%
\pgfpathlineto{\pgfqpoint{1.880706in}{1.813535in}}%
\pgfpathlineto{\pgfqpoint{1.883224in}{1.795171in}}%
\pgfpathlineto{\pgfqpoint{1.885743in}{1.780460in}}%
\pgfpathlineto{\pgfqpoint{1.888262in}{1.781647in}}%
\pgfpathlineto{\pgfqpoint{1.890781in}{1.772124in}}%
\pgfpathlineto{\pgfqpoint{1.893299in}{1.780261in}}%
\pgfpathlineto{\pgfqpoint{1.895818in}{1.779766in}}%
\pgfpathlineto{\pgfqpoint{1.898337in}{1.770020in}}%
\pgfpathlineto{\pgfqpoint{1.900856in}{1.771880in}}%
\pgfpathlineto{\pgfqpoint{1.903374in}{1.749323in}}%
\pgfpathlineto{\pgfqpoint{1.905893in}{1.758387in}}%
\pgfpathlineto{\pgfqpoint{1.908412in}{1.764531in}}%
\pgfpathlineto{\pgfqpoint{1.910931in}{1.758849in}}%
\pgfpathlineto{\pgfqpoint{1.913449in}{1.758295in}}%
\pgfpathlineto{\pgfqpoint{1.915968in}{1.770099in}}%
\pgfpathlineto{\pgfqpoint{1.918487in}{1.747878in}}%
\pgfpathlineto{\pgfqpoint{1.921006in}{1.744432in}}%
\pgfpathlineto{\pgfqpoint{1.923524in}{1.761812in}}%
\pgfpathlineto{\pgfqpoint{1.926043in}{1.775668in}}%
\pgfpathlineto{\pgfqpoint{1.928562in}{1.775165in}}%
\pgfpathlineto{\pgfqpoint{1.931081in}{1.777153in}}%
\pgfpathlineto{\pgfqpoint{1.933599in}{1.791529in}}%
\pgfpathlineto{\pgfqpoint{1.936118in}{1.775268in}}%
\pgfpathlineto{\pgfqpoint{1.938637in}{1.782094in}}%
\pgfpathlineto{\pgfqpoint{1.941156in}{1.787777in}}%
\pgfpathlineto{\pgfqpoint{1.943674in}{1.764470in}}%
\pgfpathlineto{\pgfqpoint{1.946193in}{1.778162in}}%
\pgfpathlineto{\pgfqpoint{1.948712in}{1.788940in}}%
\pgfpathlineto{\pgfqpoint{1.951231in}{1.788337in}}%
\pgfpathlineto{\pgfqpoint{1.953749in}{1.793297in}}%
\pgfpathlineto{\pgfqpoint{1.956268in}{1.790906in}}%
\pgfpathlineto{\pgfqpoint{1.958787in}{1.798995in}}%
\pgfpathlineto{\pgfqpoint{1.961306in}{1.793559in}}%
\pgfpathlineto{\pgfqpoint{1.963824in}{1.818686in}}%
\pgfpathlineto{\pgfqpoint{1.966343in}{1.834716in}}%
\pgfpathlineto{\pgfqpoint{1.968862in}{1.815809in}}%
\pgfpathlineto{\pgfqpoint{1.971381in}{1.800250in}}%
\pgfpathlineto{\pgfqpoint{1.973899in}{1.826913in}}%
\pgfpathlineto{\pgfqpoint{1.976418in}{1.832865in}}%
\pgfpathlineto{\pgfqpoint{1.978937in}{1.842346in}}%
\pgfpathlineto{\pgfqpoint{1.981456in}{1.863955in}}%
\pgfpathlineto{\pgfqpoint{1.983974in}{1.865777in}}%
\pgfpathlineto{\pgfqpoint{1.986493in}{1.888706in}}%
\pgfpathlineto{\pgfqpoint{1.989012in}{1.859793in}}%
\pgfpathlineto{\pgfqpoint{1.991531in}{1.847203in}}%
\pgfpathlineto{\pgfqpoint{1.994049in}{1.845587in}}%
\pgfpathlineto{\pgfqpoint{1.996568in}{1.850314in}}%
\pgfpathlineto{\pgfqpoint{1.999087in}{1.842381in}}%
\pgfpathlineto{\pgfqpoint{2.001606in}{1.852761in}}%
\pgfpathlineto{\pgfqpoint{2.004124in}{1.869951in}}%
\pgfpathlineto{\pgfqpoint{2.006643in}{1.894334in}}%
\pgfpathlineto{\pgfqpoint{2.009162in}{1.896903in}}%
\pgfpathlineto{\pgfqpoint{2.011681in}{1.886087in}}%
\pgfpathlineto{\pgfqpoint{2.014199in}{1.889545in}}%
\pgfpathlineto{\pgfqpoint{2.016718in}{1.890985in}}%
\pgfpathlineto{\pgfqpoint{2.019237in}{1.898569in}}%
\pgfpathlineto{\pgfqpoint{2.021756in}{1.892826in}}%
\pgfpathlineto{\pgfqpoint{2.024274in}{1.898536in}}%
\pgfpathlineto{\pgfqpoint{2.026793in}{1.902353in}}%
\pgfpathlineto{\pgfqpoint{2.029312in}{1.914175in}}%
\pgfpathlineto{\pgfqpoint{2.031831in}{1.918636in}}%
\pgfpathlineto{\pgfqpoint{2.034349in}{1.876893in}}%
\pgfpathlineto{\pgfqpoint{2.036868in}{1.890666in}}%
\pgfpathlineto{\pgfqpoint{2.039387in}{1.873327in}}%
\pgfpathlineto{\pgfqpoint{2.041906in}{1.879309in}}%
\pgfpathlineto{\pgfqpoint{2.044424in}{1.845145in}}%
\pgfpathlineto{\pgfqpoint{2.046943in}{1.853854in}}%
\pgfpathlineto{\pgfqpoint{2.049462in}{1.881133in}}%
\pgfpathlineto{\pgfqpoint{2.051981in}{1.882080in}}%
\pgfpathlineto{\pgfqpoint{2.054499in}{1.896851in}}%
\pgfpathlineto{\pgfqpoint{2.057018in}{1.934222in}}%
\pgfpathlineto{\pgfqpoint{2.059537in}{1.915289in}}%
\pgfpathlineto{\pgfqpoint{2.062056in}{1.951049in}}%
\pgfpathlineto{\pgfqpoint{2.064574in}{1.953919in}}%
\pgfpathlineto{\pgfqpoint{2.067093in}{1.953182in}}%
\pgfpathlineto{\pgfqpoint{2.069612in}{1.953275in}}%
\pgfpathlineto{\pgfqpoint{2.072131in}{1.953161in}}%
\pgfpathlineto{\pgfqpoint{2.074649in}{1.948901in}}%
\pgfpathlineto{\pgfqpoint{2.077168in}{1.923057in}}%
\pgfpathlineto{\pgfqpoint{2.079687in}{1.924203in}}%
\pgfpathlineto{\pgfqpoint{2.082206in}{1.901645in}}%
\pgfpathlineto{\pgfqpoint{2.084724in}{1.912680in}}%
\pgfpathlineto{\pgfqpoint{2.087243in}{1.938978in}}%
\pgfpathlineto{\pgfqpoint{2.089762in}{1.908796in}}%
\pgfpathlineto{\pgfqpoint{2.092281in}{1.894989in}}%
\pgfpathlineto{\pgfqpoint{2.094799in}{1.900688in}}%
\pgfpathlineto{\pgfqpoint{2.097318in}{1.876247in}}%
\pgfpathlineto{\pgfqpoint{2.099837in}{1.867471in}}%
\pgfpathlineto{\pgfqpoint{2.102356in}{1.850056in}}%
\pgfpathlineto{\pgfqpoint{2.104874in}{1.862960in}}%
\pgfpathlineto{\pgfqpoint{2.107393in}{1.835762in}}%
\pgfpathlineto{\pgfqpoint{2.109912in}{1.821708in}}%
\pgfpathlineto{\pgfqpoint{2.112431in}{1.817189in}}%
\pgfpathlineto{\pgfqpoint{2.114949in}{1.824002in}}%
\pgfpathlineto{\pgfqpoint{2.117468in}{1.822282in}}%
\pgfpathlineto{\pgfqpoint{2.119987in}{1.829296in}}%
\pgfpathlineto{\pgfqpoint{2.122506in}{1.821189in}}%
\pgfpathlineto{\pgfqpoint{2.125024in}{1.822857in}}%
\pgfpathlineto{\pgfqpoint{2.127543in}{1.832457in}}%
\pgfpathlineto{\pgfqpoint{2.130062in}{1.827211in}}%
\pgfpathlineto{\pgfqpoint{2.132581in}{1.818225in}}%
\pgfpathlineto{\pgfqpoint{2.135099in}{1.814004in}}%
\pgfpathlineto{\pgfqpoint{2.137618in}{1.803293in}}%
\pgfpathlineto{\pgfqpoint{2.140137in}{1.860174in}}%
\pgfpathlineto{\pgfqpoint{2.142656in}{1.889438in}}%
\pgfpathlineto{\pgfqpoint{2.145174in}{1.883644in}}%
\pgfpathlineto{\pgfqpoint{2.147693in}{1.852103in}}%
\pgfpathlineto{\pgfqpoint{2.150212in}{1.850471in}}%
\pgfpathlineto{\pgfqpoint{2.152731in}{1.837494in}}%
\pgfpathlineto{\pgfqpoint{2.155249in}{1.826462in}}%
\pgfpathlineto{\pgfqpoint{2.157768in}{1.851311in}}%
\pgfpathlineto{\pgfqpoint{2.160287in}{1.862185in}}%
\pgfpathlineto{\pgfqpoint{2.162806in}{1.875038in}}%
\pgfpathlineto{\pgfqpoint{2.165324in}{1.871305in}}%
\pgfpathlineto{\pgfqpoint{2.167843in}{1.877573in}}%
\pgfpathlineto{\pgfqpoint{2.170362in}{1.909046in}}%
\pgfpathlineto{\pgfqpoint{2.172881in}{1.919757in}}%
\pgfpathlineto{\pgfqpoint{2.175399in}{1.910038in}}%
\pgfpathlineto{\pgfqpoint{2.177918in}{1.871700in}}%
\pgfpathlineto{\pgfqpoint{2.180437in}{1.848774in}}%
\pgfpathlineto{\pgfqpoint{2.182956in}{1.849162in}}%
\pgfpathlineto{\pgfqpoint{2.185474in}{1.857160in}}%
\pgfpathlineto{\pgfqpoint{2.187993in}{1.889778in}}%
\pgfpathlineto{\pgfqpoint{2.190512in}{1.896058in}}%
\pgfpathlineto{\pgfqpoint{2.193031in}{1.872028in}}%
\pgfpathlineto{\pgfqpoint{2.195549in}{1.895439in}}%
\pgfpathlineto{\pgfqpoint{2.198068in}{1.880504in}}%
\pgfpathlineto{\pgfqpoint{2.200587in}{1.866780in}}%
\pgfpathlineto{\pgfqpoint{2.203106in}{1.897889in}}%
\pgfpathlineto{\pgfqpoint{2.205624in}{1.923963in}}%
\pgfpathlineto{\pgfqpoint{2.208143in}{1.905775in}}%
\pgfpathlineto{\pgfqpoint{2.210662in}{1.873171in}}%
\pgfpathlineto{\pgfqpoint{2.213181in}{1.886283in}}%
\pgfpathlineto{\pgfqpoint{2.215699in}{1.902115in}}%
\pgfpathlineto{\pgfqpoint{2.218218in}{1.912594in}}%
\pgfpathlineto{\pgfqpoint{2.220737in}{1.911625in}}%
\pgfpathlineto{\pgfqpoint{2.223256in}{1.935694in}}%
\pgfpathlineto{\pgfqpoint{2.225774in}{1.945845in}}%
\pgfpathlineto{\pgfqpoint{2.228293in}{1.928847in}}%
\pgfpathlineto{\pgfqpoint{2.230812in}{1.910476in}}%
\pgfpathlineto{\pgfqpoint{2.233331in}{1.902510in}}%
\pgfpathlineto{\pgfqpoint{2.235849in}{1.884176in}}%
\pgfpathlineto{\pgfqpoint{2.238368in}{1.893427in}}%
\pgfpathlineto{\pgfqpoint{2.240887in}{1.893794in}}%
\pgfpathlineto{\pgfqpoint{2.243406in}{1.864599in}}%
\pgfpathlineto{\pgfqpoint{2.245924in}{1.865017in}}%
\pgfpathlineto{\pgfqpoint{2.248443in}{1.865586in}}%
\pgfpathlineto{\pgfqpoint{2.250962in}{1.866778in}}%
\pgfpathlineto{\pgfqpoint{2.253481in}{1.877607in}}%
\pgfpathlineto{\pgfqpoint{2.255999in}{1.870266in}}%
\pgfpathlineto{\pgfqpoint{2.258518in}{1.866259in}}%
\pgfpathlineto{\pgfqpoint{2.261037in}{1.853452in}}%
\pgfpathlineto{\pgfqpoint{2.263556in}{1.870734in}}%
\pgfpathlineto{\pgfqpoint{2.266074in}{1.871259in}}%
\pgfpathlineto{\pgfqpoint{2.268593in}{1.902801in}}%
\pgfpathlineto{\pgfqpoint{2.271112in}{1.893376in}}%
\pgfpathlineto{\pgfqpoint{2.273631in}{1.898574in}}%
\pgfpathlineto{\pgfqpoint{2.276149in}{1.915541in}}%
\pgfpathlineto{\pgfqpoint{2.278668in}{1.907622in}}%
\pgfpathlineto{\pgfqpoint{2.281187in}{1.896822in}}%
\pgfpathlineto{\pgfqpoint{2.283706in}{1.886854in}}%
\pgfpathlineto{\pgfqpoint{2.286224in}{1.890333in}}%
\pgfpathlineto{\pgfqpoint{2.288743in}{1.910598in}}%
\pgfpathlineto{\pgfqpoint{2.291262in}{1.919730in}}%
\pgfpathlineto{\pgfqpoint{2.293781in}{1.910596in}}%
\pgfpathlineto{\pgfqpoint{2.296299in}{1.940177in}}%
\pgfpathlineto{\pgfqpoint{2.298818in}{1.952148in}}%
\pgfpathlineto{\pgfqpoint{2.301337in}{1.931299in}}%
\pgfpathlineto{\pgfqpoint{2.303856in}{1.902158in}}%
\pgfpathlineto{\pgfqpoint{2.306374in}{1.903436in}}%
\pgfpathlineto{\pgfqpoint{2.308893in}{1.905066in}}%
\pgfpathlineto{\pgfqpoint{2.311412in}{1.915125in}}%
\pgfpathlineto{\pgfqpoint{2.313931in}{1.901347in}}%
\pgfpathlineto{\pgfqpoint{2.316449in}{1.874439in}}%
\pgfpathlineto{\pgfqpoint{2.318968in}{1.863668in}}%
\pgfpathlineto{\pgfqpoint{2.321487in}{1.832530in}}%
\pgfpathlineto{\pgfqpoint{2.324006in}{1.838307in}}%
\pgfpathlineto{\pgfqpoint{2.326524in}{1.842546in}}%
\pgfpathlineto{\pgfqpoint{2.329043in}{1.846564in}}%
\pgfpathlineto{\pgfqpoint{2.331562in}{1.829907in}}%
\pgfpathlineto{\pgfqpoint{2.334081in}{1.829612in}}%
\pgfpathlineto{\pgfqpoint{2.336599in}{1.789876in}}%
\pgfpathlineto{\pgfqpoint{2.339118in}{1.780233in}}%
\pgfpathlineto{\pgfqpoint{2.341637in}{1.788633in}}%
\pgfpathlineto{\pgfqpoint{2.344156in}{1.796475in}}%
\pgfpathlineto{\pgfqpoint{2.346674in}{1.809187in}}%
\pgfpathlineto{\pgfqpoint{2.349193in}{1.828788in}}%
\pgfpathlineto{\pgfqpoint{2.351712in}{1.806581in}}%
\pgfpathlineto{\pgfqpoint{2.354231in}{1.802745in}}%
\pgfpathlineto{\pgfqpoint{2.356749in}{1.791846in}}%
\pgfpathlineto{\pgfqpoint{2.359268in}{1.807128in}}%
\pgfpathlineto{\pgfqpoint{2.361787in}{1.794948in}}%
\pgfpathlineto{\pgfqpoint{2.364306in}{1.788758in}}%
\pgfpathlineto{\pgfqpoint{2.366824in}{1.781971in}}%
\pgfpathlineto{\pgfqpoint{2.369343in}{1.819398in}}%
\pgfpathlineto{\pgfqpoint{2.371862in}{1.803173in}}%
\pgfpathlineto{\pgfqpoint{2.374381in}{1.803217in}}%
\pgfpathlineto{\pgfqpoint{2.376899in}{1.822291in}}%
\pgfpathlineto{\pgfqpoint{2.379418in}{1.801352in}}%
\pgfpathlineto{\pgfqpoint{2.381937in}{1.799691in}}%
\pgfpathlineto{\pgfqpoint{2.384456in}{1.815806in}}%
\pgfpathlineto{\pgfqpoint{2.386974in}{1.797824in}}%
\pgfpathlineto{\pgfqpoint{2.389493in}{1.795053in}}%
\pgfpathlineto{\pgfqpoint{2.392012in}{1.805407in}}%
\pgfpathlineto{\pgfqpoint{2.394531in}{1.801048in}}%
\pgfpathlineto{\pgfqpoint{2.397049in}{1.839396in}}%
\pgfpathlineto{\pgfqpoint{2.399568in}{1.842818in}}%
\pgfpathlineto{\pgfqpoint{2.402087in}{1.830740in}}%
\pgfpathlineto{\pgfqpoint{2.404606in}{1.828068in}}%
\pgfpathlineto{\pgfqpoint{2.407124in}{1.837665in}}%
\pgfpathlineto{\pgfqpoint{2.409643in}{1.840555in}}%
\pgfpathlineto{\pgfqpoint{2.412162in}{1.833214in}}%
\pgfpathlineto{\pgfqpoint{2.414681in}{1.830236in}}%
\pgfpathlineto{\pgfqpoint{2.417199in}{1.841486in}}%
\pgfpathlineto{\pgfqpoint{2.419718in}{1.869026in}}%
\pgfpathlineto{\pgfqpoint{2.422237in}{1.854390in}}%
\pgfpathlineto{\pgfqpoint{2.424756in}{1.854402in}}%
\pgfpathlineto{\pgfqpoint{2.427274in}{1.829135in}}%
\pgfpathlineto{\pgfqpoint{2.429793in}{1.821292in}}%
\pgfpathlineto{\pgfqpoint{2.432312in}{1.831223in}}%
\pgfpathlineto{\pgfqpoint{2.434831in}{1.829618in}}%
\pgfpathlineto{\pgfqpoint{2.437349in}{1.831146in}}%
\pgfpathlineto{\pgfqpoint{2.439868in}{1.833702in}}%
\pgfpathlineto{\pgfqpoint{2.442387in}{1.842323in}}%
\pgfpathlineto{\pgfqpoint{2.444906in}{1.842153in}}%
\pgfpathlineto{\pgfqpoint{2.447424in}{1.860429in}}%
\pgfpathlineto{\pgfqpoint{2.449943in}{1.876641in}}%
\pgfpathlineto{\pgfqpoint{2.452462in}{1.893764in}}%
\pgfpathlineto{\pgfqpoint{2.454981in}{1.916635in}}%
\pgfpathlineto{\pgfqpoint{2.457499in}{1.899423in}}%
\pgfpathlineto{\pgfqpoint{2.460018in}{1.894623in}}%
\pgfpathlineto{\pgfqpoint{2.462537in}{1.857186in}}%
\pgfpathlineto{\pgfqpoint{2.465056in}{1.874653in}}%
\pgfpathlineto{\pgfqpoint{2.467574in}{1.898422in}}%
\pgfpathlineto{\pgfqpoint{2.470093in}{1.897941in}}%
\pgfpathlineto{\pgfqpoint{2.472612in}{1.894463in}}%
\pgfpathlineto{\pgfqpoint{2.475131in}{1.886326in}}%
\pgfpathlineto{\pgfqpoint{2.477649in}{1.890514in}}%
\pgfpathlineto{\pgfqpoint{2.480168in}{1.912751in}}%
\pgfpathlineto{\pgfqpoint{2.482687in}{1.912486in}}%
\pgfpathlineto{\pgfqpoint{2.485206in}{1.917750in}}%
\pgfpathlineto{\pgfqpoint{2.487724in}{1.912057in}}%
\pgfpathlineto{\pgfqpoint{2.490243in}{1.916589in}}%
\pgfpathlineto{\pgfqpoint{2.492762in}{1.917172in}}%
\pgfpathlineto{\pgfqpoint{2.495281in}{1.916404in}}%
\pgfpathlineto{\pgfqpoint{2.497799in}{1.897622in}}%
\pgfpathlineto{\pgfqpoint{2.500318in}{1.926865in}}%
\pgfpathlineto{\pgfqpoint{2.502837in}{1.935801in}}%
\pgfpathlineto{\pgfqpoint{2.505356in}{1.936426in}}%
\pgfpathlineto{\pgfqpoint{2.507874in}{1.960142in}}%
\pgfpathlineto{\pgfqpoint{2.510393in}{1.952455in}}%
\pgfpathlineto{\pgfqpoint{2.512912in}{1.958650in}}%
\pgfpathlineto{\pgfqpoint{2.515431in}{1.943407in}}%
\pgfpathlineto{\pgfqpoint{2.517949in}{1.970405in}}%
\pgfpathlineto{\pgfqpoint{2.520468in}{1.973672in}}%
\pgfpathlineto{\pgfqpoint{2.522987in}{1.998107in}}%
\pgfpathlineto{\pgfqpoint{2.525506in}{1.987080in}}%
\pgfpathlineto{\pgfqpoint{2.528024in}{1.979977in}}%
\pgfpathlineto{\pgfqpoint{2.530543in}{1.976651in}}%
\pgfpathlineto{\pgfqpoint{2.533062in}{2.008449in}}%
\pgfpathlineto{\pgfqpoint{2.535581in}{2.020787in}}%
\pgfpathlineto{\pgfqpoint{2.538099in}{2.027002in}}%
\pgfpathlineto{\pgfqpoint{2.540618in}{2.017884in}}%
\pgfpathlineto{\pgfqpoint{2.543137in}{2.038193in}}%
\pgfpathlineto{\pgfqpoint{2.545656in}{2.034107in}}%
\pgfpathlineto{\pgfqpoint{2.548174in}{2.027296in}}%
\pgfpathlineto{\pgfqpoint{2.550693in}{2.052993in}}%
\pgfpathlineto{\pgfqpoint{2.553212in}{2.041215in}}%
\pgfpathlineto{\pgfqpoint{2.555731in}{1.996970in}}%
\pgfpathlineto{\pgfqpoint{2.558249in}{1.992247in}}%
\pgfpathlineto{\pgfqpoint{2.560768in}{2.023390in}}%
\pgfpathlineto{\pgfqpoint{2.563287in}{2.035858in}}%
\pgfpathlineto{\pgfqpoint{2.565806in}{2.039968in}}%
\pgfpathlineto{\pgfqpoint{2.568324in}{2.050040in}}%
\pgfpathlineto{\pgfqpoint{2.570843in}{2.050161in}}%
\pgfpathlineto{\pgfqpoint{2.573362in}{2.063573in}}%
\pgfpathlineto{\pgfqpoint{2.575881in}{2.063996in}}%
\pgfpathlineto{\pgfqpoint{2.578399in}{2.044762in}}%
\pgfpathlineto{\pgfqpoint{2.580918in}{2.030399in}}%
\pgfpathlineto{\pgfqpoint{2.583437in}{2.033961in}}%
\pgfpathlineto{\pgfqpoint{2.585956in}{2.042204in}}%
\pgfpathlineto{\pgfqpoint{2.588474in}{2.069202in}}%
\pgfpathlineto{\pgfqpoint{2.590993in}{2.091542in}}%
\pgfpathlineto{\pgfqpoint{2.593512in}{2.077495in}}%
\pgfpathlineto{\pgfqpoint{2.596031in}{2.086469in}}%
\pgfpathlineto{\pgfqpoint{2.598549in}{2.077137in}}%
\pgfpathlineto{\pgfqpoint{2.601068in}{2.072259in}}%
\pgfpathlineto{\pgfqpoint{2.603587in}{2.082573in}}%
\pgfpathlineto{\pgfqpoint{2.606106in}{2.085304in}}%
\pgfpathlineto{\pgfqpoint{2.608624in}{2.089267in}}%
\pgfpathlineto{\pgfqpoint{2.611143in}{2.093454in}}%
\pgfpathlineto{\pgfqpoint{2.613662in}{2.112747in}}%
\pgfpathlineto{\pgfqpoint{2.616181in}{2.102787in}}%
\pgfpathlineto{\pgfqpoint{2.618699in}{2.099206in}}%
\pgfpathlineto{\pgfqpoint{2.621218in}{2.111608in}}%
\pgfpathlineto{\pgfqpoint{2.623737in}{2.090073in}}%
\pgfpathlineto{\pgfqpoint{2.626256in}{2.076346in}}%
\pgfpathlineto{\pgfqpoint{2.628774in}{2.077862in}}%
\pgfpathlineto{\pgfqpoint{2.631293in}{2.086966in}}%
\pgfpathlineto{\pgfqpoint{2.633812in}{2.093517in}}%
\pgfpathlineto{\pgfqpoint{2.636331in}{2.109467in}}%
\pgfpathlineto{\pgfqpoint{2.638849in}{2.111239in}}%
\pgfpathlineto{\pgfqpoint{2.641368in}{2.136524in}}%
\pgfpathlineto{\pgfqpoint{2.643887in}{2.099951in}}%
\pgfpathlineto{\pgfqpoint{2.646406in}{2.074269in}}%
\pgfpathlineto{\pgfqpoint{2.648924in}{2.073572in}}%
\pgfpathlineto{\pgfqpoint{2.651443in}{2.067340in}}%
\pgfpathlineto{\pgfqpoint{2.653962in}{2.067294in}}%
\pgfpathlineto{\pgfqpoint{2.656481in}{2.037491in}}%
\pgfpathlineto{\pgfqpoint{2.658999in}{2.013366in}}%
\pgfpathlineto{\pgfqpoint{2.661518in}{2.027200in}}%
\pgfpathlineto{\pgfqpoint{2.664037in}{2.004638in}}%
\pgfpathlineto{\pgfqpoint{2.666556in}{1.942604in}}%
\pgfpathlineto{\pgfqpoint{2.669074in}{1.963256in}}%
\pgfpathlineto{\pgfqpoint{2.671593in}{1.990804in}}%
\pgfpathlineto{\pgfqpoint{2.674112in}{2.005443in}}%
\pgfpathlineto{\pgfqpoint{2.676631in}{2.026087in}}%
\pgfpathlineto{\pgfqpoint{2.679149in}{2.023470in}}%
\pgfpathlineto{\pgfqpoint{2.681668in}{2.029137in}}%
\pgfpathlineto{\pgfqpoint{2.684187in}{2.019499in}}%
\pgfpathlineto{\pgfqpoint{2.686706in}{2.007801in}}%
\pgfpathlineto{\pgfqpoint{2.689224in}{2.009297in}}%
\pgfpathlineto{\pgfqpoint{2.691743in}{1.999302in}}%
\pgfpathlineto{\pgfqpoint{2.694262in}{1.974126in}}%
\pgfpathlineto{\pgfqpoint{2.696781in}{1.968526in}}%
\pgfpathlineto{\pgfqpoint{2.699299in}{1.979476in}}%
\pgfpathlineto{\pgfqpoint{2.701818in}{1.960029in}}%
\pgfpathlineto{\pgfqpoint{2.704337in}{1.930642in}}%
\pgfpathlineto{\pgfqpoint{2.706856in}{1.936319in}}%
\pgfpathlineto{\pgfqpoint{2.709374in}{1.948504in}}%
\pgfpathlineto{\pgfqpoint{2.711893in}{1.952793in}}%
\pgfpathlineto{\pgfqpoint{2.714412in}{1.975377in}}%
\pgfpathlineto{\pgfqpoint{2.716931in}{2.006850in}}%
\pgfpathlineto{\pgfqpoint{2.719449in}{2.015538in}}%
\pgfpathlineto{\pgfqpoint{2.721968in}{2.021542in}}%
\pgfpathlineto{\pgfqpoint{2.724487in}{2.027080in}}%
\pgfpathlineto{\pgfqpoint{2.727006in}{2.026546in}}%
\pgfpathlineto{\pgfqpoint{2.729524in}{1.982405in}}%
\pgfpathlineto{\pgfqpoint{2.732043in}{1.980139in}}%
\pgfpathlineto{\pgfqpoint{2.734562in}{1.973925in}}%
\pgfpathlineto{\pgfqpoint{2.737081in}{1.985611in}}%
\pgfpathlineto{\pgfqpoint{2.739599in}{1.977508in}}%
\pgfpathlineto{\pgfqpoint{2.742118in}{1.980944in}}%
\pgfpathlineto{\pgfqpoint{2.744637in}{1.981021in}}%
\pgfpathlineto{\pgfqpoint{2.747156in}{1.989186in}}%
\pgfpathlineto{\pgfqpoint{2.749674in}{2.010457in}}%
\pgfpathlineto{\pgfqpoint{2.752193in}{1.976190in}}%
\pgfpathlineto{\pgfqpoint{2.754712in}{1.979865in}}%
\pgfpathlineto{\pgfqpoint{2.757231in}{1.983781in}}%
\pgfpathlineto{\pgfqpoint{2.759749in}{1.960458in}}%
\pgfpathlineto{\pgfqpoint{2.762268in}{1.930903in}}%
\pgfpathlineto{\pgfqpoint{2.764787in}{1.933586in}}%
\pgfpathlineto{\pgfqpoint{2.767306in}{1.959764in}}%
\pgfpathlineto{\pgfqpoint{2.769824in}{1.949171in}}%
\pgfpathlineto{\pgfqpoint{2.772343in}{1.943125in}}%
\pgfpathlineto{\pgfqpoint{2.774862in}{1.949866in}}%
\pgfpathlineto{\pgfqpoint{2.777381in}{1.971005in}}%
\pgfpathlineto{\pgfqpoint{2.779899in}{1.956724in}}%
\pgfpathlineto{\pgfqpoint{2.782418in}{1.946933in}}%
\pgfpathlineto{\pgfqpoint{2.784937in}{1.943838in}}%
\pgfpathlineto{\pgfqpoint{2.787456in}{1.935592in}}%
\pgfpathlineto{\pgfqpoint{2.789974in}{1.946660in}}%
\pgfpathlineto{\pgfqpoint{2.792493in}{1.942832in}}%
\pgfpathlineto{\pgfqpoint{2.795012in}{1.941391in}}%
\pgfpathlineto{\pgfqpoint{2.797531in}{1.969109in}}%
\pgfpathlineto{\pgfqpoint{2.800049in}{1.953757in}}%
\pgfpathlineto{\pgfqpoint{2.802568in}{1.973474in}}%
\pgfpathlineto{\pgfqpoint{2.805087in}{1.975506in}}%
\pgfpathlineto{\pgfqpoint{2.807606in}{2.003948in}}%
\pgfpathlineto{\pgfqpoint{2.810124in}{1.995997in}}%
\pgfpathlineto{\pgfqpoint{2.812643in}{1.995267in}}%
\pgfpathlineto{\pgfqpoint{2.815162in}{1.970690in}}%
\pgfpathlineto{\pgfqpoint{2.817681in}{1.968600in}}%
\pgfpathlineto{\pgfqpoint{2.820199in}{1.985497in}}%
\pgfpathlineto{\pgfqpoint{2.822718in}{1.969867in}}%
\pgfpathlineto{\pgfqpoint{2.825237in}{1.925469in}}%
\pgfpathlineto{\pgfqpoint{2.827756in}{1.920382in}}%
\pgfpathlineto{\pgfqpoint{2.830274in}{1.900811in}}%
\pgfpathlineto{\pgfqpoint{2.832793in}{1.906109in}}%
\pgfpathlineto{\pgfqpoint{2.835312in}{1.938074in}}%
\pgfpathlineto{\pgfqpoint{2.837831in}{1.938769in}}%
\pgfpathlineto{\pgfqpoint{2.840349in}{1.944755in}}%
\pgfpathlineto{\pgfqpoint{2.842868in}{1.942074in}}%
\pgfpathlineto{\pgfqpoint{2.845387in}{1.947262in}}%
\pgfpathlineto{\pgfqpoint{2.847906in}{1.931658in}}%
\pgfpathlineto{\pgfqpoint{2.850424in}{1.954276in}}%
\pgfpathlineto{\pgfqpoint{2.852943in}{1.972631in}}%
\pgfpathlineto{\pgfqpoint{2.855462in}{1.969564in}}%
\pgfpathlineto{\pgfqpoint{2.857981in}{1.984804in}}%
\pgfpathlineto{\pgfqpoint{2.860499in}{2.013017in}}%
\pgfpathlineto{\pgfqpoint{2.863018in}{1.998064in}}%
\pgfpathlineto{\pgfqpoint{2.865537in}{1.977480in}}%
\pgfpathlineto{\pgfqpoint{2.868056in}{1.962033in}}%
\pgfpathlineto{\pgfqpoint{2.870574in}{1.951605in}}%
\pgfpathlineto{\pgfqpoint{2.873093in}{1.945637in}}%
\pgfpathlineto{\pgfqpoint{2.875612in}{1.940234in}}%
\pgfpathlineto{\pgfqpoint{2.878131in}{1.964305in}}%
\pgfpathlineto{\pgfqpoint{2.880649in}{1.953861in}}%
\pgfpathlineto{\pgfqpoint{2.883168in}{1.977070in}}%
\pgfpathlineto{\pgfqpoint{2.885687in}{1.998436in}}%
\pgfpathlineto{\pgfqpoint{2.888206in}{1.985998in}}%
\pgfpathlineto{\pgfqpoint{2.890724in}{1.994449in}}%
\pgfpathlineto{\pgfqpoint{2.893243in}{2.003458in}}%
\pgfpathlineto{\pgfqpoint{2.895762in}{1.997599in}}%
\pgfpathlineto{\pgfqpoint{2.898281in}{1.972402in}}%
\pgfpathlineto{\pgfqpoint{2.900799in}{1.959622in}}%
\pgfpathlineto{\pgfqpoint{2.903318in}{1.972130in}}%
\pgfpathlineto{\pgfqpoint{2.905837in}{1.976887in}}%
\pgfpathlineto{\pgfqpoint{2.908356in}{1.991545in}}%
\pgfpathlineto{\pgfqpoint{2.910874in}{1.977856in}}%
\pgfpathlineto{\pgfqpoint{2.913393in}{1.958552in}}%
\pgfpathlineto{\pgfqpoint{2.915912in}{1.949019in}}%
\pgfpathlineto{\pgfqpoint{2.918431in}{1.940194in}}%
\pgfpathlineto{\pgfqpoint{2.920949in}{1.945247in}}%
\pgfpathlineto{\pgfqpoint{2.923468in}{1.936624in}}%
\pgfpathlineto{\pgfqpoint{2.925987in}{1.961258in}}%
\pgfpathlineto{\pgfqpoint{2.928506in}{1.933092in}}%
\pgfpathlineto{\pgfqpoint{2.931024in}{1.935161in}}%
\pgfpathlineto{\pgfqpoint{2.933543in}{1.923890in}}%
\pgfpathlineto{\pgfqpoint{2.936062in}{1.909915in}}%
\pgfpathlineto{\pgfqpoint{2.938581in}{1.916097in}}%
\pgfpathlineto{\pgfqpoint{2.941099in}{1.908127in}}%
\pgfpathlineto{\pgfqpoint{2.943618in}{1.917391in}}%
\pgfpathlineto{\pgfqpoint{2.946137in}{1.933244in}}%
\pgfpathlineto{\pgfqpoint{2.948656in}{1.920768in}}%
\pgfpathlineto{\pgfqpoint{2.951174in}{1.902402in}}%
\pgfpathlineto{\pgfqpoint{2.953693in}{1.912878in}}%
\pgfpathlineto{\pgfqpoint{2.956212in}{1.908272in}}%
\pgfpathlineto{\pgfqpoint{2.958731in}{1.926014in}}%
\pgfpathlineto{\pgfqpoint{2.961249in}{1.930878in}}%
\pgfpathlineto{\pgfqpoint{2.963768in}{1.948133in}}%
\pgfpathlineto{\pgfqpoint{2.966287in}{1.937370in}}%
\pgfpathlineto{\pgfqpoint{2.968806in}{1.943217in}}%
\pgfpathlineto{\pgfqpoint{2.971324in}{1.912761in}}%
\pgfpathlineto{\pgfqpoint{2.973843in}{1.896268in}}%
\pgfpathlineto{\pgfqpoint{2.976362in}{1.919953in}}%
\pgfpathlineto{\pgfqpoint{2.978881in}{1.925277in}}%
\pgfpathlineto{\pgfqpoint{2.981399in}{1.920145in}}%
\pgfpathlineto{\pgfqpoint{2.983918in}{1.938918in}}%
\pgfpathlineto{\pgfqpoint{2.986437in}{1.971272in}}%
\pgfpathlineto{\pgfqpoint{2.991474in}{1.967753in}}%
\pgfpathlineto{\pgfqpoint{2.993993in}{1.969918in}}%
\pgfpathlineto{\pgfqpoint{2.996512in}{1.965536in}}%
\pgfpathlineto{\pgfqpoint{2.999031in}{1.991445in}}%
\pgfpathlineto{\pgfqpoint{3.001549in}{1.965615in}}%
\pgfpathlineto{\pgfqpoint{3.004068in}{1.932853in}}%
\pgfpathlineto{\pgfqpoint{3.006587in}{1.954537in}}%
\pgfpathlineto{\pgfqpoint{3.011624in}{1.969369in}}%
\pgfpathlineto{\pgfqpoint{3.014143in}{1.973316in}}%
\pgfpathlineto{\pgfqpoint{3.016662in}{1.981843in}}%
\pgfpathlineto{\pgfqpoint{3.019181in}{2.006805in}}%
\pgfpathlineto{\pgfqpoint{3.021699in}{2.011191in}}%
\pgfpathlineto{\pgfqpoint{3.024218in}{2.007083in}}%
\pgfpathlineto{\pgfqpoint{3.026737in}{2.011686in}}%
\pgfpathlineto{\pgfqpoint{3.029256in}{1.992895in}}%
\pgfpathlineto{\pgfqpoint{3.031774in}{2.028330in}}%
\pgfpathlineto{\pgfqpoint{3.034293in}{2.017530in}}%
\pgfpathlineto{\pgfqpoint{3.036812in}{2.030721in}}%
\pgfpathlineto{\pgfqpoint{3.039331in}{2.058099in}}%
\pgfpathlineto{\pgfqpoint{3.041849in}{2.040329in}}%
\pgfpathlineto{\pgfqpoint{3.044368in}{2.028137in}}%
\pgfpathlineto{\pgfqpoint{3.046887in}{2.051568in}}%
\pgfpathlineto{\pgfqpoint{3.049406in}{2.062824in}}%
\pgfpathlineto{\pgfqpoint{3.051924in}{2.067454in}}%
\pgfpathlineto{\pgfqpoint{3.054443in}{2.048817in}}%
\pgfpathlineto{\pgfqpoint{3.056962in}{2.008776in}}%
\pgfpathlineto{\pgfqpoint{3.059481in}{2.035253in}}%
\pgfpathlineto{\pgfqpoint{3.061999in}{2.021081in}}%
\pgfpathlineto{\pgfqpoint{3.064518in}{2.059163in}}%
\pgfpathlineto{\pgfqpoint{3.067037in}{2.083220in}}%
\pgfpathlineto{\pgfqpoint{3.069556in}{2.117177in}}%
\pgfpathlineto{\pgfqpoint{3.072074in}{2.100675in}}%
\pgfpathlineto{\pgfqpoint{3.074593in}{2.092281in}}%
\pgfpathlineto{\pgfqpoint{3.077112in}{2.070196in}}%
\pgfpathlineto{\pgfqpoint{3.079631in}{2.066397in}}%
\pgfpathlineto{\pgfqpoint{3.082149in}{2.081909in}}%
\pgfpathlineto{\pgfqpoint{3.084668in}{2.112294in}}%
\pgfpathlineto{\pgfqpoint{3.087187in}{2.099241in}}%
\pgfpathlineto{\pgfqpoint{3.089706in}{2.105653in}}%
\pgfpathlineto{\pgfqpoint{3.094743in}{2.090901in}}%
\pgfpathlineto{\pgfqpoint{3.097262in}{2.097498in}}%
\pgfpathlineto{\pgfqpoint{3.099781in}{2.117374in}}%
\pgfpathlineto{\pgfqpoint{3.102299in}{2.122129in}}%
\pgfpathlineto{\pgfqpoint{3.104818in}{2.111454in}}%
\pgfpathlineto{\pgfqpoint{3.107337in}{2.143090in}}%
\pgfpathlineto{\pgfqpoint{3.109856in}{2.110876in}}%
\pgfpathlineto{\pgfqpoint{3.112374in}{2.095258in}}%
\pgfpathlineto{\pgfqpoint{3.114893in}{2.072743in}}%
\pgfpathlineto{\pgfqpoint{3.117412in}{2.090070in}}%
\pgfpathlineto{\pgfqpoint{3.119931in}{2.059461in}}%
\pgfpathlineto{\pgfqpoint{3.122449in}{2.030095in}}%
\pgfpathlineto{\pgfqpoint{3.124968in}{2.023373in}}%
\pgfpathlineto{\pgfqpoint{3.127487in}{2.037508in}}%
\pgfpathlineto{\pgfqpoint{3.130006in}{2.033150in}}%
\pgfpathlineto{\pgfqpoint{3.132524in}{2.051764in}}%
\pgfpathlineto{\pgfqpoint{3.135043in}{2.055279in}}%
\pgfpathlineto{\pgfqpoint{3.137562in}{2.021941in}}%
\pgfpathlineto{\pgfqpoint{3.140081in}{2.011375in}}%
\pgfpathlineto{\pgfqpoint{3.142599in}{2.027280in}}%
\pgfpathlineto{\pgfqpoint{3.145118in}{2.035309in}}%
\pgfpathlineto{\pgfqpoint{3.147637in}{2.040676in}}%
\pgfpathlineto{\pgfqpoint{3.150156in}{2.040707in}}%
\pgfpathlineto{\pgfqpoint{3.152674in}{2.047886in}}%
\pgfpathlineto{\pgfqpoint{3.155193in}{2.047856in}}%
\pgfpathlineto{\pgfqpoint{3.157712in}{2.057508in}}%
\pgfpathlineto{\pgfqpoint{3.160231in}{2.071533in}}%
\pgfpathlineto{\pgfqpoint{3.162749in}{2.045403in}}%
\pgfpathlineto{\pgfqpoint{3.165268in}{2.056229in}}%
\pgfpathlineto{\pgfqpoint{3.167787in}{2.095230in}}%
\pgfpathlineto{\pgfqpoint{3.170306in}{2.090995in}}%
\pgfpathlineto{\pgfqpoint{3.172824in}{2.093243in}}%
\pgfpathlineto{\pgfqpoint{3.175343in}{2.069204in}}%
\pgfpathlineto{\pgfqpoint{3.177862in}{2.080472in}}%
\pgfpathlineto{\pgfqpoint{3.180381in}{2.117158in}}%
\pgfpathlineto{\pgfqpoint{3.182899in}{2.115726in}}%
\pgfpathlineto{\pgfqpoint{3.185418in}{2.111977in}}%
\pgfpathlineto{\pgfqpoint{3.187937in}{2.127821in}}%
\pgfpathlineto{\pgfqpoint{3.190456in}{2.115844in}}%
\pgfpathlineto{\pgfqpoint{3.192974in}{2.100859in}}%
\pgfpathlineto{\pgfqpoint{3.195493in}{2.090065in}}%
\pgfpathlineto{\pgfqpoint{3.198012in}{2.115654in}}%
\pgfpathlineto{\pgfqpoint{3.200531in}{2.085311in}}%
\pgfpathlineto{\pgfqpoint{3.203049in}{2.088687in}}%
\pgfpathlineto{\pgfqpoint{3.205568in}{2.118590in}}%
\pgfpathlineto{\pgfqpoint{3.208087in}{2.126655in}}%
\pgfpathlineto{\pgfqpoint{3.210606in}{2.133728in}}%
\pgfpathlineto{\pgfqpoint{3.213124in}{2.114599in}}%
\pgfpathlineto{\pgfqpoint{3.215643in}{2.111065in}}%
\pgfpathlineto{\pgfqpoint{3.218162in}{2.095172in}}%
\pgfpathlineto{\pgfqpoint{3.220681in}{2.097470in}}%
\pgfpathlineto{\pgfqpoint{3.223199in}{2.092382in}}%
\pgfpathlineto{\pgfqpoint{3.225718in}{2.080274in}}%
\pgfpathlineto{\pgfqpoint{3.228237in}{2.069838in}}%
\pgfpathlineto{\pgfqpoint{3.230756in}{2.065056in}}%
\pgfpathlineto{\pgfqpoint{3.233274in}{2.037782in}}%
\pgfpathlineto{\pgfqpoint{3.235793in}{2.069712in}}%
\pgfpathlineto{\pgfqpoint{3.238312in}{2.048144in}}%
\pgfpathlineto{\pgfqpoint{3.240831in}{2.057881in}}%
\pgfpathlineto{\pgfqpoint{3.243349in}{2.083073in}}%
\pgfpathlineto{\pgfqpoint{3.245868in}{2.082602in}}%
\pgfpathlineto{\pgfqpoint{3.248387in}{2.086647in}}%
\pgfpathlineto{\pgfqpoint{3.250906in}{2.112345in}}%
\pgfpathlineto{\pgfqpoint{3.253424in}{2.063312in}}%
\pgfpathlineto{\pgfqpoint{3.255943in}{2.066544in}}%
\pgfpathlineto{\pgfqpoint{3.258462in}{2.053586in}}%
\pgfpathlineto{\pgfqpoint{3.260981in}{2.038397in}}%
\pgfpathlineto{\pgfqpoint{3.263499in}{2.035490in}}%
\pgfpathlineto{\pgfqpoint{3.266018in}{2.025298in}}%
\pgfpathlineto{\pgfqpoint{3.268537in}{2.014670in}}%
\pgfpathlineto{\pgfqpoint{3.271056in}{1.994532in}}%
\pgfpathlineto{\pgfqpoint{3.273574in}{1.966784in}}%
\pgfpathlineto{\pgfqpoint{3.276093in}{1.975770in}}%
\pgfpathlineto{\pgfqpoint{3.278612in}{1.977735in}}%
\pgfpathlineto{\pgfqpoint{3.281131in}{1.999310in}}%
\pgfpathlineto{\pgfqpoint{3.283649in}{1.975597in}}%
\pgfpathlineto{\pgfqpoint{3.286168in}{1.982602in}}%
\pgfpathlineto{\pgfqpoint{3.288687in}{1.965236in}}%
\pgfpathlineto{\pgfqpoint{3.291206in}{1.967383in}}%
\pgfpathlineto{\pgfqpoint{3.293724in}{1.962617in}}%
\pgfpathlineto{\pgfqpoint{3.296243in}{1.973726in}}%
\pgfpathlineto{\pgfqpoint{3.298762in}{1.955492in}}%
\pgfpathlineto{\pgfqpoint{3.301281in}{1.923471in}}%
\pgfpathlineto{\pgfqpoint{3.303799in}{1.910810in}}%
\pgfpathlineto{\pgfqpoint{3.306318in}{1.916330in}}%
\pgfpathlineto{\pgfqpoint{3.308837in}{1.941330in}}%
\pgfpathlineto{\pgfqpoint{3.311356in}{1.967659in}}%
\pgfpathlineto{\pgfqpoint{3.313874in}{1.970729in}}%
\pgfpathlineto{\pgfqpoint{3.316393in}{1.991554in}}%
\pgfpathlineto{\pgfqpoint{3.318912in}{1.995948in}}%
\pgfpathlineto{\pgfqpoint{3.323949in}{1.968164in}}%
\pgfpathlineto{\pgfqpoint{3.326468in}{1.985925in}}%
\pgfpathlineto{\pgfqpoint{3.328987in}{1.963020in}}%
\pgfpathlineto{\pgfqpoint{3.331506in}{1.946651in}}%
\pgfpathlineto{\pgfqpoint{3.334024in}{1.957745in}}%
\pgfpathlineto{\pgfqpoint{3.336543in}{1.976983in}}%
\pgfpathlineto{\pgfqpoint{3.339062in}{1.976024in}}%
\pgfpathlineto{\pgfqpoint{3.341581in}{1.967741in}}%
\pgfpathlineto{\pgfqpoint{3.344099in}{1.981498in}}%
\pgfpathlineto{\pgfqpoint{3.346618in}{1.963012in}}%
\pgfpathlineto{\pgfqpoint{3.349137in}{1.940172in}}%
\pgfpathlineto{\pgfqpoint{3.351656in}{1.930897in}}%
\pgfpathlineto{\pgfqpoint{3.354174in}{1.920321in}}%
\pgfpathlineto{\pgfqpoint{3.356693in}{1.929456in}}%
\pgfpathlineto{\pgfqpoint{3.359212in}{1.946422in}}%
\pgfpathlineto{\pgfqpoint{3.361731in}{1.987923in}}%
\pgfpathlineto{\pgfqpoint{3.366768in}{1.970104in}}%
\pgfpathlineto{\pgfqpoint{3.369287in}{1.998177in}}%
\pgfpathlineto{\pgfqpoint{3.371806in}{2.012562in}}%
\pgfpathlineto{\pgfqpoint{3.374324in}{2.018275in}}%
\pgfpathlineto{\pgfqpoint{3.376843in}{2.008971in}}%
\pgfpathlineto{\pgfqpoint{3.379362in}{2.023278in}}%
\pgfpathlineto{\pgfqpoint{3.381881in}{2.019117in}}%
\pgfpathlineto{\pgfqpoint{3.384399in}{2.020109in}}%
\pgfpathlineto{\pgfqpoint{3.386918in}{2.014986in}}%
\pgfpathlineto{\pgfqpoint{3.389437in}{2.034374in}}%
\pgfpathlineto{\pgfqpoint{3.391956in}{2.029076in}}%
\pgfpathlineto{\pgfqpoint{3.394474in}{2.021463in}}%
\pgfpathlineto{\pgfqpoint{3.396993in}{2.034373in}}%
\pgfpathlineto{\pgfqpoint{3.399512in}{2.033657in}}%
\pgfpathlineto{\pgfqpoint{3.402031in}{2.033981in}}%
\pgfpathlineto{\pgfqpoint{3.404549in}{2.032512in}}%
\pgfpathlineto{\pgfqpoint{3.407068in}{2.015667in}}%
\pgfpathlineto{\pgfqpoint{3.409587in}{2.020088in}}%
\pgfpathlineto{\pgfqpoint{3.412106in}{2.019504in}}%
\pgfpathlineto{\pgfqpoint{3.414624in}{2.020257in}}%
\pgfpathlineto{\pgfqpoint{3.417143in}{2.061643in}}%
\pgfpathlineto{\pgfqpoint{3.419662in}{2.065907in}}%
\pgfpathlineto{\pgfqpoint{3.422181in}{2.085818in}}%
\pgfpathlineto{\pgfqpoint{3.424699in}{2.121069in}}%
\pgfpathlineto{\pgfqpoint{3.427218in}{2.121603in}}%
\pgfpathlineto{\pgfqpoint{3.429737in}{2.095191in}}%
\pgfpathlineto{\pgfqpoint{3.432256in}{2.086594in}}%
\pgfpathlineto{\pgfqpoint{3.434774in}{2.070568in}}%
\pgfpathlineto{\pgfqpoint{3.437293in}{2.062741in}}%
\pgfpathlineto{\pgfqpoint{3.439812in}{2.027357in}}%
\pgfpathlineto{\pgfqpoint{3.442331in}{1.995387in}}%
\pgfpathlineto{\pgfqpoint{3.444849in}{2.005729in}}%
\pgfpathlineto{\pgfqpoint{3.447368in}{1.987573in}}%
\pgfpathlineto{\pgfqpoint{3.449887in}{2.007992in}}%
\pgfpathlineto{\pgfqpoint{3.452406in}{2.021741in}}%
\pgfpathlineto{\pgfqpoint{3.454924in}{2.018417in}}%
\pgfpathlineto{\pgfqpoint{3.457443in}{2.020641in}}%
\pgfpathlineto{\pgfqpoint{3.459962in}{2.021727in}}%
\pgfpathlineto{\pgfqpoint{3.462481in}{2.008278in}}%
\pgfpathlineto{\pgfqpoint{3.464999in}{1.994280in}}%
\pgfpathlineto{\pgfqpoint{3.467518in}{2.014809in}}%
\pgfpathlineto{\pgfqpoint{3.470037in}{2.010613in}}%
\pgfpathlineto{\pgfqpoint{3.472556in}{1.999684in}}%
\pgfpathlineto{\pgfqpoint{3.475074in}{2.020615in}}%
\pgfpathlineto{\pgfqpoint{3.477593in}{2.057054in}}%
\pgfpathlineto{\pgfqpoint{3.480112in}{2.026874in}}%
\pgfpathlineto{\pgfqpoint{3.482631in}{2.028719in}}%
\pgfpathlineto{\pgfqpoint{3.485149in}{2.014995in}}%
\pgfpathlineto{\pgfqpoint{3.487668in}{2.036393in}}%
\pgfpathlineto{\pgfqpoint{3.490187in}{2.021017in}}%
\pgfpathlineto{\pgfqpoint{3.492706in}{2.021458in}}%
\pgfpathlineto{\pgfqpoint{3.495224in}{2.013069in}}%
\pgfpathlineto{\pgfqpoint{3.497743in}{1.998556in}}%
\pgfpathlineto{\pgfqpoint{3.500262in}{2.011995in}}%
\pgfpathlineto{\pgfqpoint{3.502781in}{2.019819in}}%
\pgfpathlineto{\pgfqpoint{3.505299in}{1.988856in}}%
\pgfpathlineto{\pgfqpoint{3.507818in}{2.011597in}}%
\pgfpathlineto{\pgfqpoint{3.510337in}{2.011214in}}%
\pgfpathlineto{\pgfqpoint{3.512856in}{1.988692in}}%
\pgfpathlineto{\pgfqpoint{3.515374in}{2.000554in}}%
\pgfpathlineto{\pgfqpoint{3.517893in}{1.986694in}}%
\pgfpathlineto{\pgfqpoint{3.520412in}{1.994920in}}%
\pgfpathlineto{\pgfqpoint{3.522931in}{1.992005in}}%
\pgfpathlineto{\pgfqpoint{3.525449in}{1.956773in}}%
\pgfpathlineto{\pgfqpoint{3.527968in}{1.970415in}}%
\pgfpathlineto{\pgfqpoint{3.530487in}{1.966361in}}%
\pgfpathlineto{\pgfqpoint{3.533006in}{1.957188in}}%
\pgfpathlineto{\pgfqpoint{3.535524in}{1.957753in}}%
\pgfpathlineto{\pgfqpoint{3.538043in}{1.982682in}}%
\pgfpathlineto{\pgfqpoint{3.540562in}{1.964681in}}%
\pgfpathlineto{\pgfqpoint{3.543081in}{1.951098in}}%
\pgfpathlineto{\pgfqpoint{3.545599in}{1.945957in}}%
\pgfpathlineto{\pgfqpoint{3.548118in}{1.944814in}}%
\pgfpathlineto{\pgfqpoint{3.550637in}{1.944705in}}%
\pgfpathlineto{\pgfqpoint{3.553156in}{1.938836in}}%
\pgfpathlineto{\pgfqpoint{3.555674in}{1.942560in}}%
\pgfpathlineto{\pgfqpoint{3.558193in}{1.957687in}}%
\pgfpathlineto{\pgfqpoint{3.560712in}{1.965787in}}%
\pgfpathlineto{\pgfqpoint{3.563231in}{1.972116in}}%
\pgfpathlineto{\pgfqpoint{3.565749in}{1.981708in}}%
\pgfpathlineto{\pgfqpoint{3.568268in}{1.967252in}}%
\pgfpathlineto{\pgfqpoint{3.570787in}{1.978196in}}%
\pgfpathlineto{\pgfqpoint{3.573306in}{1.984043in}}%
\pgfpathlineto{\pgfqpoint{3.575824in}{1.962868in}}%
\pgfpathlineto{\pgfqpoint{3.578343in}{1.964502in}}%
\pgfpathlineto{\pgfqpoint{3.580862in}{1.993216in}}%
\pgfpathlineto{\pgfqpoint{3.583381in}{1.974045in}}%
\pgfpathlineto{\pgfqpoint{3.585899in}{1.971986in}}%
\pgfpathlineto{\pgfqpoint{3.588418in}{1.990389in}}%
\pgfpathlineto{\pgfqpoint{3.590937in}{1.973329in}}%
\pgfpathlineto{\pgfqpoint{3.593456in}{1.975034in}}%
\pgfpathlineto{\pgfqpoint{3.595974in}{2.000238in}}%
\pgfpathlineto{\pgfqpoint{3.598493in}{2.016131in}}%
\pgfpathlineto{\pgfqpoint{3.601012in}{2.003846in}}%
\pgfpathlineto{\pgfqpoint{3.603531in}{2.020585in}}%
\pgfpathlineto{\pgfqpoint{3.606049in}{2.002103in}}%
\pgfpathlineto{\pgfqpoint{3.608568in}{2.002880in}}%
\pgfpathlineto{\pgfqpoint{3.611087in}{1.991798in}}%
\pgfpathlineto{\pgfqpoint{3.613606in}{1.994717in}}%
\pgfpathlineto{\pgfqpoint{3.616124in}{2.001328in}}%
\pgfpathlineto{\pgfqpoint{3.618643in}{2.008477in}}%
\pgfpathlineto{\pgfqpoint{3.621162in}{1.993423in}}%
\pgfpathlineto{\pgfqpoint{3.623681in}{2.002867in}}%
\pgfpathlineto{\pgfqpoint{3.626199in}{1.976959in}}%
\pgfpathlineto{\pgfqpoint{3.628718in}{1.969629in}}%
\pgfpathlineto{\pgfqpoint{3.631237in}{1.951527in}}%
\pgfpathlineto{\pgfqpoint{3.633756in}{1.969252in}}%
\pgfpathlineto{\pgfqpoint{3.636274in}{1.985735in}}%
\pgfpathlineto{\pgfqpoint{3.638793in}{2.013954in}}%
\pgfpathlineto{\pgfqpoint{3.641312in}{1.992140in}}%
\pgfpathlineto{\pgfqpoint{3.643831in}{1.967852in}}%
\pgfpathlineto{\pgfqpoint{3.646349in}{1.979306in}}%
\pgfpathlineto{\pgfqpoint{3.648868in}{1.982259in}}%
\pgfpathlineto{\pgfqpoint{3.651387in}{1.995614in}}%
\pgfpathlineto{\pgfqpoint{3.653906in}{2.002723in}}%
\pgfpathlineto{\pgfqpoint{3.656424in}{2.015332in}}%
\pgfpathlineto{\pgfqpoint{3.658943in}{2.015675in}}%
\pgfpathlineto{\pgfqpoint{3.661462in}{2.034232in}}%
\pgfpathlineto{\pgfqpoint{3.663981in}{2.047789in}}%
\pgfpathlineto{\pgfqpoint{3.666499in}{2.040495in}}%
\pgfpathlineto{\pgfqpoint{3.669018in}{2.039571in}}%
\pgfpathlineto{\pgfqpoint{3.671537in}{2.048127in}}%
\pgfpathlineto{\pgfqpoint{3.674056in}{2.054004in}}%
\pgfpathlineto{\pgfqpoint{3.676574in}{2.038418in}}%
\pgfpathlineto{\pgfqpoint{3.679093in}{2.025329in}}%
\pgfpathlineto{\pgfqpoint{3.681612in}{2.005360in}}%
\pgfpathlineto{\pgfqpoint{3.684131in}{2.023162in}}%
\pgfpathlineto{\pgfqpoint{3.686649in}{2.025995in}}%
\pgfpathlineto{\pgfqpoint{3.689168in}{2.032494in}}%
\pgfpathlineto{\pgfqpoint{3.691687in}{2.030502in}}%
\pgfpathlineto{\pgfqpoint{3.694206in}{2.031841in}}%
\pgfpathlineto{\pgfqpoint{3.696724in}{2.023069in}}%
\pgfpathlineto{\pgfqpoint{3.699243in}{2.022771in}}%
\pgfpathlineto{\pgfqpoint{3.701762in}{2.053439in}}%
\pgfpathlineto{\pgfqpoint{3.704281in}{2.059622in}}%
\pgfpathlineto{\pgfqpoint{3.706799in}{2.058915in}}%
\pgfpathlineto{\pgfqpoint{3.709318in}{2.061683in}}%
\pgfpathlineto{\pgfqpoint{3.711837in}{2.065538in}}%
\pgfpathlineto{\pgfqpoint{3.714356in}{2.081639in}}%
\pgfpathlineto{\pgfqpoint{3.716874in}{2.084572in}}%
\pgfpathlineto{\pgfqpoint{3.719393in}{2.104008in}}%
\pgfpathlineto{\pgfqpoint{3.721912in}{2.096259in}}%
\pgfpathlineto{\pgfqpoint{3.724431in}{2.106854in}}%
\pgfpathlineto{\pgfqpoint{3.726949in}{2.097894in}}%
\pgfpathlineto{\pgfqpoint{3.729468in}{2.091540in}}%
\pgfpathlineto{\pgfqpoint{3.731987in}{2.102747in}}%
\pgfpathlineto{\pgfqpoint{3.734506in}{2.152126in}}%
\pgfpathlineto{\pgfqpoint{3.737024in}{2.158321in}}%
\pgfpathlineto{\pgfqpoint{3.739543in}{2.112835in}}%
\pgfpathlineto{\pgfqpoint{3.742062in}{2.109077in}}%
\pgfpathlineto{\pgfqpoint{3.744581in}{2.139729in}}%
\pgfpathlineto{\pgfqpoint{3.747099in}{2.108477in}}%
\pgfpathlineto{\pgfqpoint{3.749618in}{2.119234in}}%
\pgfpathlineto{\pgfqpoint{3.752137in}{2.100652in}}%
\pgfpathlineto{\pgfqpoint{3.754656in}{2.101947in}}%
\pgfpathlineto{\pgfqpoint{3.757174in}{2.143875in}}%
\pgfpathlineto{\pgfqpoint{3.759693in}{2.133195in}}%
\pgfpathlineto{\pgfqpoint{3.762212in}{2.134387in}}%
\pgfpathlineto{\pgfqpoint{3.764731in}{2.146172in}}%
\pgfpathlineto{\pgfqpoint{3.767249in}{2.135618in}}%
\pgfpathlineto{\pgfqpoint{3.769768in}{2.129557in}}%
\pgfpathlineto{\pgfqpoint{3.772287in}{2.165810in}}%
\pgfpathlineto{\pgfqpoint{3.774806in}{2.153839in}}%
\pgfpathlineto{\pgfqpoint{3.777324in}{2.117123in}}%
\pgfpathlineto{\pgfqpoint{3.779843in}{2.122406in}}%
\pgfpathlineto{\pgfqpoint{3.782362in}{2.124882in}}%
\pgfpathlineto{\pgfqpoint{3.784881in}{2.126273in}}%
\pgfpathlineto{\pgfqpoint{3.787399in}{2.149357in}}%
\pgfpathlineto{\pgfqpoint{3.789918in}{2.183072in}}%
\pgfpathlineto{\pgfqpoint{3.792437in}{2.205554in}}%
\pgfpathlineto{\pgfqpoint{3.794956in}{2.226745in}}%
\pgfpathlineto{\pgfqpoint{3.797474in}{2.256943in}}%
\pgfpathlineto{\pgfqpoint{3.799993in}{2.268887in}}%
\pgfpathlineto{\pgfqpoint{3.802512in}{2.295835in}}%
\pgfpathlineto{\pgfqpoint{3.805031in}{2.265960in}}%
\pgfpathlineto{\pgfqpoint{3.807549in}{2.272775in}}%
\pgfpathlineto{\pgfqpoint{3.810068in}{2.227824in}}%
\pgfpathlineto{\pgfqpoint{3.812587in}{2.221603in}}%
\pgfpathlineto{\pgfqpoint{3.815106in}{2.228361in}}%
\pgfpathlineto{\pgfqpoint{3.817624in}{2.220515in}}%
\pgfpathlineto{\pgfqpoint{3.820143in}{2.234004in}}%
\pgfpathlineto{\pgfqpoint{3.822662in}{2.249528in}}%
\pgfpathlineto{\pgfqpoint{3.825181in}{2.246321in}}%
\pgfpathlineto{\pgfqpoint{3.827699in}{2.280825in}}%
\pgfpathlineto{\pgfqpoint{3.830218in}{2.304369in}}%
\pgfpathlineto{\pgfqpoint{3.832737in}{2.306144in}}%
\pgfpathlineto{\pgfqpoint{3.835256in}{2.295320in}}%
\pgfpathlineto{\pgfqpoint{3.837774in}{2.304084in}}%
\pgfpathlineto{\pgfqpoint{3.840293in}{2.281704in}}%
\pgfpathlineto{\pgfqpoint{3.842812in}{2.303393in}}%
\pgfpathlineto{\pgfqpoint{3.845331in}{2.317571in}}%
\pgfpathlineto{\pgfqpoint{3.847849in}{2.324267in}}%
\pgfpathlineto{\pgfqpoint{3.850368in}{2.332833in}}%
\pgfpathlineto{\pgfqpoint{3.852887in}{2.347911in}}%
\pgfpathlineto{\pgfqpoint{3.855406in}{2.340290in}}%
\pgfpathlineto{\pgfqpoint{3.857924in}{2.316475in}}%
\pgfpathlineto{\pgfqpoint{3.860443in}{2.298862in}}%
\pgfpathlineto{\pgfqpoint{3.862962in}{2.289542in}}%
\pgfpathlineto{\pgfqpoint{3.865481in}{2.274087in}}%
\pgfpathlineto{\pgfqpoint{3.867999in}{2.264649in}}%
\pgfpathlineto{\pgfqpoint{3.870518in}{2.278873in}}%
\pgfpathlineto{\pgfqpoint{3.873037in}{2.262333in}}%
\pgfpathlineto{\pgfqpoint{3.875556in}{2.297217in}}%
\pgfpathlineto{\pgfqpoint{3.878074in}{2.298203in}}%
\pgfpathlineto{\pgfqpoint{3.880593in}{2.299974in}}%
\pgfpathlineto{\pgfqpoint{3.883112in}{2.295374in}}%
\pgfpathlineto{\pgfqpoint{3.885631in}{2.293250in}}%
\pgfpathlineto{\pgfqpoint{3.888149in}{2.307633in}}%
\pgfpathlineto{\pgfqpoint{3.890668in}{2.316727in}}%
\pgfpathlineto{\pgfqpoint{3.893187in}{2.331812in}}%
\pgfpathlineto{\pgfqpoint{3.895706in}{2.337216in}}%
\pgfpathlineto{\pgfqpoint{3.898224in}{2.347236in}}%
\pgfpathlineto{\pgfqpoint{3.900743in}{2.354461in}}%
\pgfpathlineto{\pgfqpoint{3.903262in}{2.365101in}}%
\pgfpathlineto{\pgfqpoint{3.905781in}{2.374898in}}%
\pgfpathlineto{\pgfqpoint{3.908299in}{2.380697in}}%
\pgfpathlineto{\pgfqpoint{3.910818in}{2.399953in}}%
\pgfpathlineto{\pgfqpoint{3.915856in}{2.346128in}}%
\pgfpathlineto{\pgfqpoint{3.918374in}{2.324195in}}%
\pgfpathlineto{\pgfqpoint{3.920893in}{2.354584in}}%
\pgfpathlineto{\pgfqpoint{3.923412in}{2.344341in}}%
\pgfpathlineto{\pgfqpoint{3.925931in}{2.318190in}}%
\pgfpathlineto{\pgfqpoint{3.928449in}{2.356432in}}%
\pgfpathlineto{\pgfqpoint{3.930968in}{2.364531in}}%
\pgfpathlineto{\pgfqpoint{3.936006in}{2.415062in}}%
\pgfpathlineto{\pgfqpoint{3.938524in}{2.447471in}}%
\pgfpathlineto{\pgfqpoint{3.941043in}{2.459162in}}%
\pgfpathlineto{\pgfqpoint{3.943562in}{2.492706in}}%
\pgfpathlineto{\pgfqpoint{3.946081in}{2.489742in}}%
\pgfpathlineto{\pgfqpoint{3.948599in}{2.519026in}}%
\pgfpathlineto{\pgfqpoint{3.951118in}{2.516540in}}%
\pgfpathlineto{\pgfqpoint{3.953637in}{2.483878in}}%
\pgfpathlineto{\pgfqpoint{3.956156in}{2.500468in}}%
\pgfpathlineto{\pgfqpoint{3.958674in}{2.486675in}}%
\pgfpathlineto{\pgfqpoint{3.961193in}{2.492025in}}%
\pgfpathlineto{\pgfqpoint{3.963712in}{2.510851in}}%
\pgfpathlineto{\pgfqpoint{3.966231in}{2.518892in}}%
\pgfpathlineto{\pgfqpoint{3.968749in}{2.517374in}}%
\pgfpathlineto{\pgfqpoint{3.971268in}{2.538191in}}%
\pgfpathlineto{\pgfqpoint{3.973787in}{2.548636in}}%
\pgfpathlineto{\pgfqpoint{3.976306in}{2.590998in}}%
\pgfpathlineto{\pgfqpoint{3.978824in}{2.583539in}}%
\pgfpathlineto{\pgfqpoint{3.981343in}{2.615418in}}%
\pgfpathlineto{\pgfqpoint{3.983862in}{2.629190in}}%
\pgfpathlineto{\pgfqpoint{3.986381in}{2.667657in}}%
\pgfpathlineto{\pgfqpoint{3.988899in}{2.690359in}}%
\pgfpathlineto{\pgfqpoint{3.991418in}{2.669447in}}%
\pgfpathlineto{\pgfqpoint{3.993937in}{2.697235in}}%
\pgfpathlineto{\pgfqpoint{3.996456in}{2.685719in}}%
\pgfpathlineto{\pgfqpoint{3.998974in}{2.685322in}}%
\pgfpathlineto{\pgfqpoint{4.001493in}{2.650043in}}%
\pgfpathlineto{\pgfqpoint{4.004012in}{2.622617in}}%
\pgfpathlineto{\pgfqpoint{4.006531in}{2.641375in}}%
\pgfpathlineto{\pgfqpoint{4.009049in}{2.673986in}}%
\pgfpathlineto{\pgfqpoint{4.011568in}{2.700041in}}%
\pgfpathlineto{\pgfqpoint{4.014087in}{2.687985in}}%
\pgfpathlineto{\pgfqpoint{4.016606in}{2.697488in}}%
\pgfpathlineto{\pgfqpoint{4.019124in}{2.701746in}}%
\pgfpathlineto{\pgfqpoint{4.021643in}{2.700183in}}%
\pgfpathlineto{\pgfqpoint{4.024162in}{2.665699in}}%
\pgfpathlineto{\pgfqpoint{4.029199in}{2.744057in}}%
\pgfpathlineto{\pgfqpoint{4.031718in}{2.759007in}}%
\pgfpathlineto{\pgfqpoint{4.034237in}{2.742743in}}%
\pgfpathlineto{\pgfqpoint{4.036756in}{2.769734in}}%
\pgfpathlineto{\pgfqpoint{4.039274in}{2.776604in}}%
\pgfpathlineto{\pgfqpoint{4.041793in}{2.759901in}}%
\pgfpathlineto{\pgfqpoint{4.044312in}{2.752357in}}%
\pgfpathlineto{\pgfqpoint{4.046831in}{2.772559in}}%
\pgfpathlineto{\pgfqpoint{4.049349in}{2.809266in}}%
\pgfpathlineto{\pgfqpoint{4.051868in}{2.794966in}}%
\pgfpathlineto{\pgfqpoint{4.054387in}{2.797937in}}%
\pgfpathlineto{\pgfqpoint{4.056906in}{2.774685in}}%
\pgfpathlineto{\pgfqpoint{4.059424in}{2.767606in}}%
\pgfpathlineto{\pgfqpoint{4.061943in}{2.740389in}}%
\pgfpathlineto{\pgfqpoint{4.064462in}{2.716733in}}%
\pgfpathlineto{\pgfqpoint{4.066981in}{2.720992in}}%
\pgfpathlineto{\pgfqpoint{4.069499in}{2.710205in}}%
\pgfpathlineto{\pgfqpoint{4.072018in}{2.667348in}}%
\pgfpathlineto{\pgfqpoint{4.074537in}{2.648988in}}%
\pgfpathlineto{\pgfqpoint{4.077056in}{2.668891in}}%
\pgfpathlineto{\pgfqpoint{4.079574in}{2.689608in}}%
\pgfpathlineto{\pgfqpoint{4.082093in}{2.693914in}}%
\pgfpathlineto{\pgfqpoint{4.084612in}{2.644349in}}%
\pgfpathlineto{\pgfqpoint{4.087131in}{2.611659in}}%
\pgfpathlineto{\pgfqpoint{4.089649in}{2.638509in}}%
\pgfpathlineto{\pgfqpoint{4.092168in}{2.627550in}}%
\pgfpathlineto{\pgfqpoint{4.094687in}{2.629692in}}%
\pgfpathlineto{\pgfqpoint{4.097206in}{2.673859in}}%
\pgfpathlineto{\pgfqpoint{4.099724in}{2.674617in}}%
\pgfpathlineto{\pgfqpoint{4.102243in}{2.677224in}}%
\pgfpathlineto{\pgfqpoint{4.104762in}{2.661281in}}%
\pgfpathlineto{\pgfqpoint{4.107281in}{2.644365in}}%
\pgfpathlineto{\pgfqpoint{4.109799in}{2.609958in}}%
\pgfpathlineto{\pgfqpoint{4.112318in}{2.587032in}}%
\pgfpathlineto{\pgfqpoint{4.114837in}{2.581548in}}%
\pgfpathlineto{\pgfqpoint{4.117356in}{2.583499in}}%
\pgfpathlineto{\pgfqpoint{4.119874in}{2.578651in}}%
\pgfpathlineto{\pgfqpoint{4.122393in}{2.580661in}}%
\pgfpathlineto{\pgfqpoint{4.124912in}{2.590591in}}%
\pgfpathlineto{\pgfqpoint{4.127431in}{2.589586in}}%
\pgfpathlineto{\pgfqpoint{4.129949in}{2.602474in}}%
\pgfpathlineto{\pgfqpoint{4.132468in}{2.605842in}}%
\pgfpathlineto{\pgfqpoint{4.134987in}{2.581292in}}%
\pgfpathlineto{\pgfqpoint{4.137506in}{2.531344in}}%
\pgfpathlineto{\pgfqpoint{4.140024in}{2.520896in}}%
\pgfpathlineto{\pgfqpoint{4.142543in}{2.542818in}}%
\pgfpathlineto{\pgfqpoint{4.145062in}{2.558698in}}%
\pgfpathlineto{\pgfqpoint{4.147581in}{2.557704in}}%
\pgfpathlineto{\pgfqpoint{4.150099in}{2.558309in}}%
\pgfpathlineto{\pgfqpoint{4.152618in}{2.560121in}}%
\pgfpathlineto{\pgfqpoint{4.155137in}{2.566940in}}%
\pgfpathlineto{\pgfqpoint{4.157656in}{2.607006in}}%
\pgfpathlineto{\pgfqpoint{4.160174in}{2.598935in}}%
\pgfpathlineto{\pgfqpoint{4.162693in}{2.579274in}}%
\pgfpathlineto{\pgfqpoint{4.165212in}{2.601756in}}%
\pgfpathlineto{\pgfqpoint{4.167731in}{2.625151in}}%
\pgfpathlineto{\pgfqpoint{4.170249in}{2.643776in}}%
\pgfpathlineto{\pgfqpoint{4.172768in}{2.630855in}}%
\pgfpathlineto{\pgfqpoint{4.175287in}{2.598994in}}%
\pgfpathlineto{\pgfqpoint{4.177806in}{2.607831in}}%
\pgfpathlineto{\pgfqpoint{4.180324in}{2.581862in}}%
\pgfpathlineto{\pgfqpoint{4.182843in}{2.590433in}}%
\pgfpathlineto{\pgfqpoint{4.185362in}{2.598255in}}%
\pgfpathlineto{\pgfqpoint{4.187881in}{2.591102in}}%
\pgfpathlineto{\pgfqpoint{4.190399in}{2.594834in}}%
\pgfpathlineto{\pgfqpoint{4.192918in}{2.605199in}}%
\pgfpathlineto{\pgfqpoint{4.195437in}{2.612226in}}%
\pgfpathlineto{\pgfqpoint{4.197956in}{2.621706in}}%
\pgfpathlineto{\pgfqpoint{4.200474in}{2.636459in}}%
\pgfpathlineto{\pgfqpoint{4.202993in}{2.633842in}}%
\pgfpathlineto{\pgfqpoint{4.205512in}{2.639709in}}%
\pgfpathlineto{\pgfqpoint{4.208031in}{2.656091in}}%
\pgfpathlineto{\pgfqpoint{4.210549in}{2.668767in}}%
\pgfpathlineto{\pgfqpoint{4.213068in}{2.662630in}}%
\pgfpathlineto{\pgfqpoint{4.215587in}{2.661379in}}%
\pgfpathlineto{\pgfqpoint{4.218106in}{2.667368in}}%
\pgfpathlineto{\pgfqpoint{4.220624in}{2.657460in}}%
\pgfpathlineto{\pgfqpoint{4.223143in}{2.668378in}}%
\pgfpathlineto{\pgfqpoint{4.225662in}{2.681878in}}%
\pgfpathlineto{\pgfqpoint{4.228181in}{2.682020in}}%
\pgfpathlineto{\pgfqpoint{4.230699in}{2.661789in}}%
\pgfpathlineto{\pgfqpoint{4.233218in}{2.652435in}}%
\pgfpathlineto{\pgfqpoint{4.235737in}{2.650520in}}%
\pgfpathlineto{\pgfqpoint{4.238256in}{2.621094in}}%
\pgfpathlineto{\pgfqpoint{4.240774in}{2.668680in}}%
\pgfpathlineto{\pgfqpoint{4.245812in}{2.696304in}}%
\pgfpathlineto{\pgfqpoint{4.248331in}{2.694736in}}%
\pgfpathlineto{\pgfqpoint{4.250849in}{2.662375in}}%
\pgfpathlineto{\pgfqpoint{4.253368in}{2.650477in}}%
\pgfpathlineto{\pgfqpoint{4.255887in}{2.646887in}}%
\pgfpathlineto{\pgfqpoint{4.258406in}{2.652408in}}%
\pgfpathlineto{\pgfqpoint{4.260924in}{2.663113in}}%
\pgfpathlineto{\pgfqpoint{4.263443in}{2.693263in}}%
\pgfpathlineto{\pgfqpoint{4.265962in}{2.707728in}}%
\pgfpathlineto{\pgfqpoint{4.268481in}{2.674324in}}%
\pgfpathlineto{\pgfqpoint{4.273518in}{2.703247in}}%
\pgfpathlineto{\pgfqpoint{4.276037in}{2.743736in}}%
\pgfpathlineto{\pgfqpoint{4.278556in}{2.758644in}}%
\pgfpathlineto{\pgfqpoint{4.281074in}{2.799686in}}%
\pgfpathlineto{\pgfqpoint{4.283593in}{2.820410in}}%
\pgfpathlineto{\pgfqpoint{4.286112in}{2.788612in}}%
\pgfpathlineto{\pgfqpoint{4.288631in}{2.776758in}}%
\pgfpathlineto{\pgfqpoint{4.291149in}{2.757802in}}%
\pgfpathlineto{\pgfqpoint{4.293668in}{2.752171in}}%
\pgfpathlineto{\pgfqpoint{4.296187in}{2.755684in}}%
\pgfpathlineto{\pgfqpoint{4.298706in}{2.780306in}}%
\pgfpathlineto{\pgfqpoint{4.301224in}{2.776120in}}%
\pgfpathlineto{\pgfqpoint{4.303743in}{2.780042in}}%
\pgfpathlineto{\pgfqpoint{4.306262in}{2.791932in}}%
\pgfpathlineto{\pgfqpoint{4.308781in}{2.794140in}}%
\pgfpathlineto{\pgfqpoint{4.311299in}{2.803506in}}%
\pgfpathlineto{\pgfqpoint{4.313818in}{2.804320in}}%
\pgfpathlineto{\pgfqpoint{4.316337in}{2.779110in}}%
\pgfpathlineto{\pgfqpoint{4.318856in}{2.781877in}}%
\pgfpathlineto{\pgfqpoint{4.321374in}{2.732588in}}%
\pgfpathlineto{\pgfqpoint{4.323893in}{2.748098in}}%
\pgfpathlineto{\pgfqpoint{4.326412in}{2.713504in}}%
\pgfpathlineto{\pgfqpoint{4.328931in}{2.749359in}}%
\pgfpathlineto{\pgfqpoint{4.331449in}{2.748935in}}%
\pgfpathlineto{\pgfqpoint{4.333968in}{2.725509in}}%
\pgfpathlineto{\pgfqpoint{4.336487in}{2.731391in}}%
\pgfpathlineto{\pgfqpoint{4.339006in}{2.751869in}}%
\pgfpathlineto{\pgfqpoint{4.341524in}{2.770342in}}%
\pgfpathlineto{\pgfqpoint{4.344043in}{2.752374in}}%
\pgfpathlineto{\pgfqpoint{4.346562in}{2.771480in}}%
\pgfpathlineto{\pgfqpoint{4.349081in}{2.760456in}}%
\pgfpathlineto{\pgfqpoint{4.351599in}{2.823674in}}%
\pgfpathlineto{\pgfqpoint{4.354118in}{2.790410in}}%
\pgfpathlineto{\pgfqpoint{4.356637in}{2.805380in}}%
\pgfpathlineto{\pgfqpoint{4.359156in}{2.809265in}}%
\pgfpathlineto{\pgfqpoint{4.361674in}{2.803787in}}%
\pgfpathlineto{\pgfqpoint{4.364193in}{2.839956in}}%
\pgfpathlineto{\pgfqpoint{4.366712in}{2.819040in}}%
\pgfpathlineto{\pgfqpoint{4.369231in}{2.815265in}}%
\pgfpathlineto{\pgfqpoint{4.371749in}{2.822966in}}%
\pgfpathlineto{\pgfqpoint{4.374268in}{2.785734in}}%
\pgfpathlineto{\pgfqpoint{4.376787in}{2.795329in}}%
\pgfpathlineto{\pgfqpoint{4.379306in}{2.786214in}}%
\pgfpathlineto{\pgfqpoint{4.381824in}{2.762993in}}%
\pgfpathlineto{\pgfqpoint{4.384343in}{2.795275in}}%
\pgfpathlineto{\pgfqpoint{4.386862in}{2.790739in}}%
\pgfpathlineto{\pgfqpoint{4.389381in}{2.795030in}}%
\pgfpathlineto{\pgfqpoint{4.391899in}{2.782810in}}%
\pgfpathlineto{\pgfqpoint{4.394418in}{2.810551in}}%
\pgfpathlineto{\pgfqpoint{4.396937in}{2.799681in}}%
\pgfpathlineto{\pgfqpoint{4.399456in}{2.756740in}}%
\pgfpathlineto{\pgfqpoint{4.404493in}{2.832304in}}%
\pgfpathlineto{\pgfqpoint{4.407012in}{2.819311in}}%
\pgfpathlineto{\pgfqpoint{4.409531in}{2.835958in}}%
\pgfpathlineto{\pgfqpoint{4.412049in}{2.853295in}}%
\pgfpathlineto{\pgfqpoint{4.414568in}{2.882861in}}%
\pgfpathlineto{\pgfqpoint{4.417087in}{2.885106in}}%
\pgfpathlineto{\pgfqpoint{4.419606in}{2.879393in}}%
\pgfpathlineto{\pgfqpoint{4.422124in}{2.901226in}}%
\pgfpathlineto{\pgfqpoint{4.424643in}{2.898147in}}%
\pgfpathlineto{\pgfqpoint{4.427162in}{2.858353in}}%
\pgfpathlineto{\pgfqpoint{4.429681in}{2.856413in}}%
\pgfpathlineto{\pgfqpoint{4.432199in}{2.859658in}}%
\pgfpathlineto{\pgfqpoint{4.434718in}{2.888642in}}%
\pgfpathlineto{\pgfqpoint{4.437237in}{2.925878in}}%
\pgfpathlineto{\pgfqpoint{4.439756in}{2.912187in}}%
\pgfpathlineto{\pgfqpoint{4.442274in}{2.921576in}}%
\pgfpathlineto{\pgfqpoint{4.444793in}{2.913426in}}%
\pgfpathlineto{\pgfqpoint{4.447312in}{2.962214in}}%
\pgfpathlineto{\pgfqpoint{4.449831in}{2.981052in}}%
\pgfpathlineto{\pgfqpoint{4.452349in}{3.007158in}}%
\pgfpathlineto{\pgfqpoint{4.454868in}{2.988463in}}%
\pgfpathlineto{\pgfqpoint{4.457387in}{2.981404in}}%
\pgfpathlineto{\pgfqpoint{4.459906in}{2.991587in}}%
\pgfpathlineto{\pgfqpoint{4.462424in}{2.980046in}}%
\pgfpathlineto{\pgfqpoint{4.464943in}{2.986645in}}%
\pgfpathlineto{\pgfqpoint{4.467462in}{2.982187in}}%
\pgfpathlineto{\pgfqpoint{4.469981in}{2.990711in}}%
\pgfpathlineto{\pgfqpoint{4.472499in}{2.980263in}}%
\pgfpathlineto{\pgfqpoint{4.475018in}{2.947015in}}%
\pgfpathlineto{\pgfqpoint{4.477537in}{2.962627in}}%
\pgfpathlineto{\pgfqpoint{4.480056in}{2.955800in}}%
\pgfpathlineto{\pgfqpoint{4.482574in}{2.936377in}}%
\pgfpathlineto{\pgfqpoint{4.485093in}{2.931301in}}%
\pgfpathlineto{\pgfqpoint{4.487612in}{2.917986in}}%
\pgfpathlineto{\pgfqpoint{4.490131in}{2.885218in}}%
\pgfpathlineto{\pgfqpoint{4.492649in}{2.880355in}}%
\pgfpathlineto{\pgfqpoint{4.495168in}{2.914059in}}%
\pgfpathlineto{\pgfqpoint{4.497687in}{2.892982in}}%
\pgfpathlineto{\pgfqpoint{4.500206in}{2.891889in}}%
\pgfpathlineto{\pgfqpoint{4.502724in}{2.877001in}}%
\pgfpathlineto{\pgfqpoint{4.505243in}{2.879658in}}%
\pgfpathlineto{\pgfqpoint{4.507762in}{2.896259in}}%
\pgfpathlineto{\pgfqpoint{4.510281in}{2.886875in}}%
\pgfpathlineto{\pgfqpoint{4.512799in}{2.917627in}}%
\pgfpathlineto{\pgfqpoint{4.515318in}{2.928316in}}%
\pgfpathlineto{\pgfqpoint{4.517837in}{2.922457in}}%
\pgfpathlineto{\pgfqpoint{4.520356in}{2.896446in}}%
\pgfpathlineto{\pgfqpoint{4.522874in}{2.880862in}}%
\pgfpathlineto{\pgfqpoint{4.525393in}{2.898121in}}%
\pgfpathlineto{\pgfqpoint{4.527912in}{2.892961in}}%
\pgfpathlineto{\pgfqpoint{4.530431in}{2.970230in}}%
\pgfpathlineto{\pgfqpoint{4.532949in}{2.986271in}}%
\pgfpathlineto{\pgfqpoint{4.535468in}{2.967071in}}%
\pgfpathlineto{\pgfqpoint{4.537987in}{2.964332in}}%
\pgfpathlineto{\pgfqpoint{4.540506in}{2.944237in}}%
\pgfpathlineto{\pgfqpoint{4.543024in}{2.959263in}}%
\pgfpathlineto{\pgfqpoint{4.545543in}{2.949434in}}%
\pgfpathlineto{\pgfqpoint{4.548062in}{2.922878in}}%
\pgfpathlineto{\pgfqpoint{4.550581in}{2.941757in}}%
\pgfpathlineto{\pgfqpoint{4.553099in}{2.950068in}}%
\pgfpathlineto{\pgfqpoint{4.555618in}{2.950388in}}%
\pgfpathlineto{\pgfqpoint{4.558137in}{2.955302in}}%
\pgfpathlineto{\pgfqpoint{4.560656in}{2.939358in}}%
\pgfpathlineto{\pgfqpoint{4.563174in}{2.946938in}}%
\pgfpathlineto{\pgfqpoint{4.565693in}{2.939336in}}%
\pgfpathlineto{\pgfqpoint{4.568212in}{2.945527in}}%
\pgfpathlineto{\pgfqpoint{4.570731in}{2.954748in}}%
\pgfpathlineto{\pgfqpoint{4.573249in}{3.014633in}}%
\pgfpathlineto{\pgfqpoint{4.575768in}{3.018486in}}%
\pgfpathlineto{\pgfqpoint{4.578287in}{3.065280in}}%
\pgfpathlineto{\pgfqpoint{4.580806in}{3.030053in}}%
\pgfpathlineto{\pgfqpoint{4.583324in}{3.048193in}}%
\pgfpathlineto{\pgfqpoint{4.585843in}{2.999873in}}%
\pgfpathlineto{\pgfqpoint{4.588362in}{3.010584in}}%
\pgfpathlineto{\pgfqpoint{4.590881in}{3.018475in}}%
\pgfpathlineto{\pgfqpoint{4.593399in}{3.020936in}}%
\pgfpathlineto{\pgfqpoint{4.595918in}{3.013472in}}%
\pgfpathlineto{\pgfqpoint{4.598437in}{3.046063in}}%
\pgfpathlineto{\pgfqpoint{4.600956in}{3.032751in}}%
\pgfpathlineto{\pgfqpoint{4.603474in}{3.038907in}}%
\pgfpathlineto{\pgfqpoint{4.605993in}{3.042492in}}%
\pgfpathlineto{\pgfqpoint{4.608512in}{3.047664in}}%
\pgfpathlineto{\pgfqpoint{4.611031in}{3.074673in}}%
\pgfpathlineto{\pgfqpoint{4.613549in}{3.077484in}}%
\pgfpathlineto{\pgfqpoint{4.616068in}{3.089914in}}%
\pgfpathlineto{\pgfqpoint{4.618587in}{3.065768in}}%
\pgfpathlineto{\pgfqpoint{4.621106in}{3.036533in}}%
\pgfpathlineto{\pgfqpoint{4.623624in}{3.042365in}}%
\pgfpathlineto{\pgfqpoint{4.626143in}{3.096272in}}%
\pgfpathlineto{\pgfqpoint{4.628662in}{3.152885in}}%
\pgfpathlineto{\pgfqpoint{4.631181in}{3.110406in}}%
\pgfpathlineto{\pgfqpoint{4.633699in}{3.114807in}}%
\pgfpathlineto{\pgfqpoint{4.636218in}{3.146278in}}%
\pgfpathlineto{\pgfqpoint{4.638737in}{3.151164in}}%
\pgfpathlineto{\pgfqpoint{4.641256in}{3.182121in}}%
\pgfpathlineto{\pgfqpoint{4.643774in}{3.195407in}}%
\pgfpathlineto{\pgfqpoint{4.646293in}{3.176104in}}%
\pgfpathlineto{\pgfqpoint{4.651331in}{3.262227in}}%
\pgfpathlineto{\pgfqpoint{4.653849in}{3.288274in}}%
\pgfpathlineto{\pgfqpoint{4.656368in}{3.304592in}}%
\pgfpathlineto{\pgfqpoint{4.658887in}{3.275437in}}%
\pgfpathlineto{\pgfqpoint{4.661406in}{3.259092in}}%
\pgfpathlineto{\pgfqpoint{4.663924in}{3.255546in}}%
\pgfpathlineto{\pgfqpoint{4.666443in}{3.224950in}}%
\pgfpathlineto{\pgfqpoint{4.668962in}{3.224797in}}%
\pgfpathlineto{\pgfqpoint{4.671481in}{3.252317in}}%
\pgfpathlineto{\pgfqpoint{4.673999in}{3.289502in}}%
\pgfpathlineto{\pgfqpoint{4.676518in}{3.286584in}}%
\pgfpathlineto{\pgfqpoint{4.679037in}{3.226512in}}%
\pgfpathlineto{\pgfqpoint{4.681556in}{3.237156in}}%
\pgfpathlineto{\pgfqpoint{4.684074in}{3.239731in}}%
\pgfpathlineto{\pgfqpoint{4.686593in}{3.206924in}}%
\pgfpathlineto{\pgfqpoint{4.689112in}{3.203225in}}%
\pgfpathlineto{\pgfqpoint{4.691631in}{3.216063in}}%
\pgfpathlineto{\pgfqpoint{4.694149in}{3.253751in}}%
\pgfpathlineto{\pgfqpoint{4.696668in}{3.281000in}}%
\pgfpathlineto{\pgfqpoint{4.699187in}{3.293135in}}%
\pgfpathlineto{\pgfqpoint{4.701706in}{3.279558in}}%
\pgfpathlineto{\pgfqpoint{4.704224in}{3.308003in}}%
\pgfpathlineto{\pgfqpoint{4.706743in}{3.304361in}}%
\pgfpathlineto{\pgfqpoint{4.709262in}{3.311039in}}%
\pgfpathlineto{\pgfqpoint{4.711781in}{3.296292in}}%
\pgfpathlineto{\pgfqpoint{4.714299in}{3.280691in}}%
\pgfpathlineto{\pgfqpoint{4.716818in}{3.295147in}}%
\pgfpathlineto{\pgfqpoint{4.719337in}{3.269443in}}%
\pgfpathlineto{\pgfqpoint{4.721856in}{3.256459in}}%
\pgfpathlineto{\pgfqpoint{4.724374in}{3.252630in}}%
\pgfpathlineto{\pgfqpoint{4.726893in}{3.258080in}}%
\pgfpathlineto{\pgfqpoint{4.729412in}{3.280261in}}%
\pgfpathlineto{\pgfqpoint{4.731931in}{3.260004in}}%
\pgfpathlineto{\pgfqpoint{4.734449in}{3.269716in}}%
\pgfpathlineto{\pgfqpoint{4.736968in}{3.312923in}}%
\pgfpathlineto{\pgfqpoint{4.742006in}{3.274774in}}%
\pgfpathlineto{\pgfqpoint{4.744524in}{3.287961in}}%
\pgfpathlineto{\pgfqpoint{4.747043in}{3.314018in}}%
\pgfpathlineto{\pgfqpoint{4.749562in}{3.344987in}}%
\pgfpathlineto{\pgfqpoint{4.752081in}{3.334428in}}%
\pgfpathlineto{\pgfqpoint{4.754599in}{3.333810in}}%
\pgfpathlineto{\pgfqpoint{4.757118in}{3.320172in}}%
\pgfpathlineto{\pgfqpoint{4.759637in}{3.310473in}}%
\pgfpathlineto{\pgfqpoint{4.762156in}{3.321510in}}%
\pgfpathlineto{\pgfqpoint{4.764674in}{3.330528in}}%
\pgfpathlineto{\pgfqpoint{4.767193in}{3.302268in}}%
\pgfpathlineto{\pgfqpoint{4.769712in}{3.258167in}}%
\pgfpathlineto{\pgfqpoint{4.772231in}{3.279059in}}%
\pgfpathlineto{\pgfqpoint{4.774749in}{3.253115in}}%
\pgfpathlineto{\pgfqpoint{4.777268in}{3.262562in}}%
\pgfpathlineto{\pgfqpoint{4.779787in}{3.281337in}}%
\pgfpathlineto{\pgfqpoint{4.782306in}{3.285869in}}%
\pgfpathlineto{\pgfqpoint{4.784824in}{3.277349in}}%
\pgfpathlineto{\pgfqpoint{4.787343in}{3.254760in}}%
\pgfpathlineto{\pgfqpoint{4.789862in}{3.296211in}}%
\pgfpathlineto{\pgfqpoint{4.792381in}{3.298466in}}%
\pgfpathlineto{\pgfqpoint{4.794899in}{3.311130in}}%
\pgfpathlineto{\pgfqpoint{4.797418in}{3.330157in}}%
\pgfpathlineto{\pgfqpoint{4.799937in}{3.341229in}}%
\pgfpathlineto{\pgfqpoint{4.802456in}{3.338724in}}%
\pgfpathlineto{\pgfqpoint{4.804974in}{3.353297in}}%
\pgfpathlineto{\pgfqpoint{4.807493in}{3.341974in}}%
\pgfpathlineto{\pgfqpoint{4.810012in}{3.346931in}}%
\pgfpathlineto{\pgfqpoint{4.812531in}{3.371650in}}%
\pgfpathlineto{\pgfqpoint{4.815049in}{3.411980in}}%
\pgfpathlineto{\pgfqpoint{4.817568in}{3.388244in}}%
\pgfpathlineto{\pgfqpoint{4.820087in}{3.379952in}}%
\pgfpathlineto{\pgfqpoint{4.822606in}{3.412054in}}%
\pgfpathlineto{\pgfqpoint{4.825124in}{3.435158in}}%
\pgfpathlineto{\pgfqpoint{4.827643in}{3.438243in}}%
\pgfpathlineto{\pgfqpoint{4.830162in}{3.446546in}}%
\pgfpathlineto{\pgfqpoint{4.832681in}{3.450969in}}%
\pgfpathlineto{\pgfqpoint{4.835199in}{3.469864in}}%
\pgfpathlineto{\pgfqpoint{4.837718in}{3.477590in}}%
\pgfpathlineto{\pgfqpoint{4.840237in}{3.456969in}}%
\pgfpathlineto{\pgfqpoint{4.842756in}{3.439371in}}%
\pgfpathlineto{\pgfqpoint{4.845274in}{3.451512in}}%
\pgfpathlineto{\pgfqpoint{4.847793in}{3.443936in}}%
\pgfpathlineto{\pgfqpoint{4.850312in}{3.446528in}}%
\pgfpathlineto{\pgfqpoint{4.852831in}{3.455503in}}%
\pgfpathlineto{\pgfqpoint{4.855349in}{3.435174in}}%
\pgfpathlineto{\pgfqpoint{4.857868in}{3.473392in}}%
\pgfpathlineto{\pgfqpoint{4.860387in}{3.454008in}}%
\pgfpathlineto{\pgfqpoint{4.862906in}{3.416708in}}%
\pgfpathlineto{\pgfqpoint{4.865424in}{3.421221in}}%
\pgfpathlineto{\pgfqpoint{4.867943in}{3.427457in}}%
\pgfpathlineto{\pgfqpoint{4.870462in}{3.436568in}}%
\pgfpathlineto{\pgfqpoint{4.872981in}{3.451225in}}%
\pgfpathlineto{\pgfqpoint{4.875499in}{3.446038in}}%
\pgfpathlineto{\pgfqpoint{4.878018in}{3.452999in}}%
\pgfpathlineto{\pgfqpoint{4.880537in}{3.469766in}}%
\pgfpathlineto{\pgfqpoint{4.883056in}{3.500963in}}%
\pgfpathlineto{\pgfqpoint{4.885574in}{3.465251in}}%
\pgfpathlineto{\pgfqpoint{4.888093in}{3.475218in}}%
\pgfpathlineto{\pgfqpoint{4.890612in}{3.469936in}}%
\pgfpathlineto{\pgfqpoint{4.893131in}{3.510001in}}%
\pgfpathlineto{\pgfqpoint{4.895649in}{3.512830in}}%
\pgfpathlineto{\pgfqpoint{4.898168in}{3.508379in}}%
\pgfpathlineto{\pgfqpoint{4.900687in}{3.508375in}}%
\pgfpathlineto{\pgfqpoint{4.903206in}{3.551190in}}%
\pgfpathlineto{\pgfqpoint{4.905724in}{3.544086in}}%
\pgfpathlineto{\pgfqpoint{4.908243in}{3.571105in}}%
\pgfpathlineto{\pgfqpoint{4.910762in}{3.601514in}}%
\pgfpathlineto{\pgfqpoint{4.913281in}{3.548489in}}%
\pgfpathlineto{\pgfqpoint{4.915799in}{3.537716in}}%
\pgfpathlineto{\pgfqpoint{4.918318in}{3.537738in}}%
\pgfpathlineto{\pgfqpoint{4.920837in}{3.542042in}}%
\pgfpathlineto{\pgfqpoint{4.923356in}{3.561365in}}%
\pgfpathlineto{\pgfqpoint{4.925874in}{3.579572in}}%
\pgfpathlineto{\pgfqpoint{4.928393in}{3.580559in}}%
\pgfpathlineto{\pgfqpoint{4.930912in}{3.600528in}}%
\pgfpathlineto{\pgfqpoint{4.933431in}{3.594798in}}%
\pgfpathlineto{\pgfqpoint{4.935949in}{3.609187in}}%
\pgfpathlineto{\pgfqpoint{4.938468in}{3.647650in}}%
\pgfpathlineto{\pgfqpoint{4.940987in}{3.584943in}}%
\pgfpathlineto{\pgfqpoint{4.943506in}{3.547673in}}%
\pgfpathlineto{\pgfqpoint{4.946024in}{3.539731in}}%
\pgfpathlineto{\pgfqpoint{4.948543in}{3.544675in}}%
\pgfpathlineto{\pgfqpoint{4.951062in}{3.517185in}}%
\pgfpathlineto{\pgfqpoint{4.953581in}{3.516613in}}%
\pgfpathlineto{\pgfqpoint{4.956099in}{3.542192in}}%
\pgfpathlineto{\pgfqpoint{4.958618in}{3.489824in}}%
\pgfpathlineto{\pgfqpoint{4.961137in}{3.450733in}}%
\pgfpathlineto{\pgfqpoint{4.963656in}{3.430268in}}%
\pgfpathlineto{\pgfqpoint{4.966174in}{3.445993in}}%
\pgfpathlineto{\pgfqpoint{4.968693in}{3.425383in}}%
\pgfpathlineto{\pgfqpoint{4.971212in}{3.448719in}}%
\pgfpathlineto{\pgfqpoint{4.973731in}{3.410557in}}%
\pgfpathlineto{\pgfqpoint{4.976249in}{3.451720in}}%
\pgfpathlineto{\pgfqpoint{4.978768in}{3.450647in}}%
\pgfpathlineto{\pgfqpoint{4.981287in}{3.439202in}}%
\pgfpathlineto{\pgfqpoint{4.983806in}{3.437040in}}%
\pgfpathlineto{\pgfqpoint{4.986324in}{3.442867in}}%
\pgfpathlineto{\pgfqpoint{4.988843in}{3.449237in}}%
\pgfpathlineto{\pgfqpoint{4.991362in}{3.454189in}}%
\pgfpathlineto{\pgfqpoint{4.993881in}{3.433774in}}%
\pgfpathlineto{\pgfqpoint{4.996399in}{3.360281in}}%
\pgfpathlineto{\pgfqpoint{4.998918in}{3.366995in}}%
\pgfpathlineto{\pgfqpoint{5.001437in}{3.348629in}}%
\pgfpathlineto{\pgfqpoint{5.003956in}{3.373246in}}%
\pgfpathlineto{\pgfqpoint{5.006474in}{3.376709in}}%
\pgfpathlineto{\pgfqpoint{5.008993in}{3.402896in}}%
\pgfpathlineto{\pgfqpoint{5.011512in}{3.360716in}}%
\pgfpathlineto{\pgfqpoint{5.014031in}{3.364192in}}%
\pgfpathlineto{\pgfqpoint{5.016549in}{3.371116in}}%
\pgfpathlineto{\pgfqpoint{5.019068in}{3.357340in}}%
\pgfpathlineto{\pgfqpoint{5.021587in}{3.330320in}}%
\pgfpathlineto{\pgfqpoint{5.024106in}{3.300825in}}%
\pgfpathlineto{\pgfqpoint{5.026624in}{3.299352in}}%
\pgfpathlineto{\pgfqpoint{5.031662in}{3.281900in}}%
\pgfpathlineto{\pgfqpoint{5.034181in}{3.291458in}}%
\pgfpathlineto{\pgfqpoint{5.036699in}{3.306679in}}%
\pgfpathlineto{\pgfqpoint{5.039218in}{3.295483in}}%
\pgfpathlineto{\pgfqpoint{5.041737in}{3.289734in}}%
\pgfpathlineto{\pgfqpoint{5.044256in}{3.343342in}}%
\pgfpathlineto{\pgfqpoint{5.046774in}{3.362183in}}%
\pgfpathlineto{\pgfqpoint{5.051812in}{3.288098in}}%
\pgfpathlineto{\pgfqpoint{5.054331in}{3.289473in}}%
\pgfpathlineto{\pgfqpoint{5.056849in}{3.296364in}}%
\pgfpathlineto{\pgfqpoint{5.059368in}{3.307275in}}%
\pgfpathlineto{\pgfqpoint{5.061887in}{3.319781in}}%
\pgfpathlineto{\pgfqpoint{5.064406in}{3.365055in}}%
\pgfpathlineto{\pgfqpoint{5.066924in}{3.340342in}}%
\pgfpathlineto{\pgfqpoint{5.069443in}{3.342042in}}%
\pgfpathlineto{\pgfqpoint{5.071962in}{3.329147in}}%
\pgfpathlineto{\pgfqpoint{5.074481in}{3.310933in}}%
\pgfpathlineto{\pgfqpoint{5.076999in}{3.303599in}}%
\pgfpathlineto{\pgfqpoint{5.079518in}{3.297698in}}%
\pgfpathlineto{\pgfqpoint{5.082037in}{3.308443in}}%
\pgfpathlineto{\pgfqpoint{5.084556in}{3.307127in}}%
\pgfpathlineto{\pgfqpoint{5.089593in}{3.269586in}}%
\pgfpathlineto{\pgfqpoint{5.092112in}{3.269895in}}%
\pgfpathlineto{\pgfqpoint{5.094631in}{3.291051in}}%
\pgfpathlineto{\pgfqpoint{5.097149in}{3.307975in}}%
\pgfpathlineto{\pgfqpoint{5.099668in}{3.302888in}}%
\pgfpathlineto{\pgfqpoint{5.102187in}{3.282989in}}%
\pgfpathlineto{\pgfqpoint{5.104706in}{3.271513in}}%
\pgfpathlineto{\pgfqpoint{5.107224in}{3.297496in}}%
\pgfpathlineto{\pgfqpoint{5.109743in}{3.319851in}}%
\pgfpathlineto{\pgfqpoint{5.112262in}{3.322479in}}%
\pgfpathlineto{\pgfqpoint{5.114781in}{3.336850in}}%
\pgfpathlineto{\pgfqpoint{5.117299in}{3.373676in}}%
\pgfpathlineto{\pgfqpoint{5.119818in}{3.396690in}}%
\pgfpathlineto{\pgfqpoint{5.122337in}{3.399447in}}%
\pgfpathlineto{\pgfqpoint{5.124856in}{3.394688in}}%
\pgfpathlineto{\pgfqpoint{5.127374in}{3.362295in}}%
\pgfpathlineto{\pgfqpoint{5.129893in}{3.348938in}}%
\pgfpathlineto{\pgfqpoint{5.132412in}{3.337470in}}%
\pgfpathlineto{\pgfqpoint{5.134931in}{3.338370in}}%
\pgfpathlineto{\pgfqpoint{5.137449in}{3.299198in}}%
\pgfpathlineto{\pgfqpoint{5.139968in}{3.285350in}}%
\pgfpathlineto{\pgfqpoint{5.142487in}{3.278226in}}%
\pgfpathlineto{\pgfqpoint{5.145006in}{3.284650in}}%
\pgfpathlineto{\pgfqpoint{5.147524in}{3.315721in}}%
\pgfpathlineto{\pgfqpoint{5.150043in}{3.318994in}}%
\pgfpathlineto{\pgfqpoint{5.152562in}{3.349855in}}%
\pgfpathlineto{\pgfqpoint{5.155081in}{3.367314in}}%
\pgfpathlineto{\pgfqpoint{5.157599in}{3.378642in}}%
\pgfpathlineto{\pgfqpoint{5.160118in}{3.402096in}}%
\pgfpathlineto{\pgfqpoint{5.162637in}{3.428602in}}%
\pgfpathlineto{\pgfqpoint{5.165156in}{3.419244in}}%
\pgfpathlineto{\pgfqpoint{5.167674in}{3.409088in}}%
\pgfpathlineto{\pgfqpoint{5.170193in}{3.394639in}}%
\pgfpathlineto{\pgfqpoint{5.172712in}{3.395916in}}%
\pgfpathlineto{\pgfqpoint{5.175231in}{3.440810in}}%
\pgfpathlineto{\pgfqpoint{5.177749in}{3.468912in}}%
\pgfpathlineto{\pgfqpoint{5.180268in}{3.482995in}}%
\pgfpathlineto{\pgfqpoint{5.182787in}{3.468703in}}%
\pgfpathlineto{\pgfqpoint{5.185306in}{3.477484in}}%
\pgfpathlineto{\pgfqpoint{5.187824in}{3.480245in}}%
\pgfpathlineto{\pgfqpoint{5.190343in}{3.510570in}}%
\pgfpathlineto{\pgfqpoint{5.192862in}{3.547578in}}%
\pgfpathlineto{\pgfqpoint{5.195381in}{3.530938in}}%
\pgfpathlineto{\pgfqpoint{5.197899in}{3.534491in}}%
\pgfpathlineto{\pgfqpoint{5.200418in}{3.545283in}}%
\pgfpathlineto{\pgfqpoint{5.202937in}{3.530112in}}%
\pgfpathlineto{\pgfqpoint{5.205456in}{3.553526in}}%
\pgfpathlineto{\pgfqpoint{5.207974in}{3.545036in}}%
\pgfpathlineto{\pgfqpoint{5.210493in}{3.554359in}}%
\pgfpathlineto{\pgfqpoint{5.213012in}{3.557065in}}%
\pgfpathlineto{\pgfqpoint{5.215531in}{3.569134in}}%
\pgfpathlineto{\pgfqpoint{5.218049in}{3.597560in}}%
\pgfpathlineto{\pgfqpoint{5.220568in}{3.607895in}}%
\pgfpathlineto{\pgfqpoint{5.223087in}{3.561955in}}%
\pgfpathlineto{\pgfqpoint{5.225606in}{3.585869in}}%
\pgfpathlineto{\pgfqpoint{5.228124in}{3.576289in}}%
\pgfpathlineto{\pgfqpoint{5.230643in}{3.580486in}}%
\pgfpathlineto{\pgfqpoint{5.233162in}{3.571774in}}%
\pgfpathlineto{\pgfqpoint{5.235681in}{3.576095in}}%
\pgfpathlineto{\pgfqpoint{5.238199in}{3.580222in}}%
\pgfpathlineto{\pgfqpoint{5.240718in}{3.564933in}}%
\pgfpathlineto{\pgfqpoint{5.243237in}{3.563132in}}%
\pgfpathlineto{\pgfqpoint{5.245756in}{3.568687in}}%
\pgfpathlineto{\pgfqpoint{5.248274in}{3.560107in}}%
\pgfpathlineto{\pgfqpoint{5.250793in}{3.582939in}}%
\pgfpathlineto{\pgfqpoint{5.253312in}{3.590102in}}%
\pgfpathlineto{\pgfqpoint{5.255831in}{3.592577in}}%
\pgfpathlineto{\pgfqpoint{5.258349in}{3.619917in}}%
\pgfpathlineto{\pgfqpoint{5.260868in}{3.637405in}}%
\pgfpathlineto{\pgfqpoint{5.263387in}{3.671934in}}%
\pgfpathlineto{\pgfqpoint{5.268424in}{3.703406in}}%
\pgfpathlineto{\pgfqpoint{5.270943in}{3.699984in}}%
\pgfpathlineto{\pgfqpoint{5.275981in}{3.669732in}}%
\pgfpathlineto{\pgfqpoint{5.278499in}{3.629037in}}%
\pgfpathlineto{\pgfqpoint{5.281018in}{3.580800in}}%
\pgfpathlineto{\pgfqpoint{5.283537in}{3.544987in}}%
\pgfpathlineto{\pgfqpoint{5.286056in}{3.543873in}}%
\pgfpathlineto{\pgfqpoint{5.288574in}{3.516554in}}%
\pgfpathlineto{\pgfqpoint{5.291093in}{3.498356in}}%
\pgfpathlineto{\pgfqpoint{5.293612in}{3.502776in}}%
\pgfpathlineto{\pgfqpoint{5.296131in}{3.491421in}}%
\pgfpathlineto{\pgfqpoint{5.298649in}{3.465209in}}%
\pgfpathlineto{\pgfqpoint{5.301168in}{3.458092in}}%
\pgfpathlineto{\pgfqpoint{5.303687in}{3.477819in}}%
\pgfpathlineto{\pgfqpoint{5.306206in}{3.427303in}}%
\pgfpathlineto{\pgfqpoint{5.308724in}{3.447927in}}%
\pgfpathlineto{\pgfqpoint{5.311243in}{3.417342in}}%
\pgfpathlineto{\pgfqpoint{5.313762in}{3.419286in}}%
\pgfpathlineto{\pgfqpoint{5.316281in}{3.430022in}}%
\pgfpathlineto{\pgfqpoint{5.318799in}{3.422405in}}%
\pgfpathlineto{\pgfqpoint{5.321318in}{3.462239in}}%
\pgfpathlineto{\pgfqpoint{5.323837in}{3.439068in}}%
\pgfpathlineto{\pgfqpoint{5.326356in}{3.460009in}}%
\pgfpathlineto{\pgfqpoint{5.328874in}{3.484332in}}%
\pgfpathlineto{\pgfqpoint{5.331393in}{3.505219in}}%
\pgfpathlineto{\pgfqpoint{5.333912in}{3.476138in}}%
\pgfpathlineto{\pgfqpoint{5.336431in}{3.474808in}}%
\pgfpathlineto{\pgfqpoint{5.338949in}{3.477319in}}%
\pgfpathlineto{\pgfqpoint{5.341468in}{3.520298in}}%
\pgfpathlineto{\pgfqpoint{5.343987in}{3.507958in}}%
\pgfpathlineto{\pgfqpoint{5.346506in}{3.523800in}}%
\pgfpathlineto{\pgfqpoint{5.349024in}{3.548762in}}%
\pgfpathlineto{\pgfqpoint{5.351543in}{3.561291in}}%
\pgfpathlineto{\pgfqpoint{5.354062in}{3.596724in}}%
\pgfpathlineto{\pgfqpoint{5.356581in}{3.622597in}}%
\pgfpathlineto{\pgfqpoint{5.359099in}{3.651540in}}%
\pgfpathlineto{\pgfqpoint{5.361618in}{3.661955in}}%
\pgfpathlineto{\pgfqpoint{5.364137in}{3.656284in}}%
\pgfpathlineto{\pgfqpoint{5.366656in}{3.688208in}}%
\pgfpathlineto{\pgfqpoint{5.369174in}{3.663906in}}%
\pgfpathlineto{\pgfqpoint{5.371693in}{3.646701in}}%
\pgfpathlineto{\pgfqpoint{5.374212in}{3.679283in}}%
\pgfpathlineto{\pgfqpoint{5.376731in}{3.697588in}}%
\pgfpathlineto{\pgfqpoint{5.379249in}{3.699497in}}%
\pgfpathlineto{\pgfqpoint{5.381768in}{3.707505in}}%
\pgfpathlineto{\pgfqpoint{5.384287in}{3.739851in}}%
\pgfpathlineto{\pgfqpoint{5.386806in}{3.738161in}}%
\pgfpathlineto{\pgfqpoint{5.389324in}{3.699335in}}%
\pgfpathlineto{\pgfqpoint{5.391843in}{3.680390in}}%
\pgfpathlineto{\pgfqpoint{5.394362in}{3.686197in}}%
\pgfpathlineto{\pgfqpoint{5.396881in}{3.711995in}}%
\pgfpathlineto{\pgfqpoint{5.399399in}{3.705895in}}%
\pgfpathlineto{\pgfqpoint{5.401918in}{3.695222in}}%
\pgfpathlineto{\pgfqpoint{5.404437in}{3.661496in}}%
\pgfpathlineto{\pgfqpoint{5.406956in}{3.657245in}}%
\pgfpathlineto{\pgfqpoint{5.409474in}{3.661430in}}%
\pgfpathlineto{\pgfqpoint{5.411993in}{3.695145in}}%
\pgfpathlineto{\pgfqpoint{5.414512in}{3.719347in}}%
\pgfpathlineto{\pgfqpoint{5.417031in}{3.689344in}}%
\pgfpathlineto{\pgfqpoint{5.419549in}{3.679340in}}%
\pgfpathlineto{\pgfqpoint{5.422068in}{3.693777in}}%
\pgfpathlineto{\pgfqpoint{5.424587in}{3.711021in}}%
\pgfpathlineto{\pgfqpoint{5.427106in}{3.697249in}}%
\pgfpathlineto{\pgfqpoint{5.429624in}{3.661030in}}%
\pgfpathlineto{\pgfqpoint{5.432143in}{3.697589in}}%
\pgfpathlineto{\pgfqpoint{5.434662in}{3.666062in}}%
\pgfpathlineto{\pgfqpoint{5.437181in}{3.695502in}}%
\pgfpathlineto{\pgfqpoint{5.439699in}{3.666864in}}%
\pgfpathlineto{\pgfqpoint{5.442218in}{3.625397in}}%
\pgfpathlineto{\pgfqpoint{5.444737in}{3.629232in}}%
\pgfpathlineto{\pgfqpoint{5.447256in}{3.662764in}}%
\pgfpathlineto{\pgfqpoint{5.449774in}{3.637589in}}%
\pgfpathlineto{\pgfqpoint{5.452293in}{3.629767in}}%
\pgfpathlineto{\pgfqpoint{5.454812in}{3.637196in}}%
\pgfpathlineto{\pgfqpoint{5.457331in}{3.650185in}}%
\pgfpathlineto{\pgfqpoint{5.459849in}{3.637525in}}%
\pgfpathlineto{\pgfqpoint{5.462368in}{3.642220in}}%
\pgfpathlineto{\pgfqpoint{5.464887in}{3.661588in}}%
\pgfpathlineto{\pgfqpoint{5.467406in}{3.668439in}}%
\pgfpathlineto{\pgfqpoint{5.469924in}{3.680095in}}%
\pgfpathlineto{\pgfqpoint{5.472443in}{3.719818in}}%
\pgfpathlineto{\pgfqpoint{5.474962in}{3.741527in}}%
\pgfpathlineto{\pgfqpoint{5.477481in}{3.732643in}}%
\pgfpathlineto{\pgfqpoint{5.479999in}{3.692199in}}%
\pgfpathlineto{\pgfqpoint{5.482518in}{3.729781in}}%
\pgfpathlineto{\pgfqpoint{5.485037in}{3.753676in}}%
\pgfpathlineto{\pgfqpoint{5.487556in}{3.727511in}}%
\pgfpathlineto{\pgfqpoint{5.490074in}{3.725834in}}%
\pgfpathlineto{\pgfqpoint{5.492593in}{3.745908in}}%
\pgfpathlineto{\pgfqpoint{5.492593in}{3.745908in}}%
\pgfusepath{stroke}%
\end{pgfscope}%
\begin{pgfscope}%
\pgfpathrectangle{\pgfqpoint{0.455093in}{0.401080in}}{\pgfqpoint{5.037500in}{4.004000in}}%
\pgfusepath{clip}%
\pgfsetrectcap%
\pgfsetroundjoin%
\pgfsetlinewidth{0.501875pt}%
\definecolor{currentstroke}{rgb}{0.690196,0.188235,0.376471}%
\pgfsetstrokecolor{currentstroke}%
\pgfsetstrokeopacity{0.800000}%
\pgfsetdash{}{0pt}%
\pgfpathmoveto{\pgfqpoint{0.455093in}{1.021044in}}%
\pgfpathlineto{\pgfqpoint{0.457612in}{1.001588in}}%
\pgfpathlineto{\pgfqpoint{0.460131in}{1.005684in}}%
\pgfpathlineto{\pgfqpoint{0.462649in}{1.000133in}}%
\pgfpathlineto{\pgfqpoint{0.465168in}{1.015457in}}%
\pgfpathlineto{\pgfqpoint{0.467687in}{1.025728in}}%
\pgfpathlineto{\pgfqpoint{0.470206in}{1.020092in}}%
\pgfpathlineto{\pgfqpoint{0.472724in}{1.005098in}}%
\pgfpathlineto{\pgfqpoint{0.475243in}{0.991539in}}%
\pgfpathlineto{\pgfqpoint{0.477762in}{0.969811in}}%
\pgfpathlineto{\pgfqpoint{0.480281in}{0.950466in}}%
\pgfpathlineto{\pgfqpoint{0.482799in}{0.951294in}}%
\pgfpathlineto{\pgfqpoint{0.485318in}{0.930976in}}%
\pgfpathlineto{\pgfqpoint{0.487837in}{0.928514in}}%
\pgfpathlineto{\pgfqpoint{0.490356in}{0.925493in}}%
\pgfpathlineto{\pgfqpoint{0.492874in}{0.924851in}}%
\pgfpathlineto{\pgfqpoint{0.495393in}{0.922947in}}%
\pgfpathlineto{\pgfqpoint{0.497912in}{0.922149in}}%
\pgfpathlineto{\pgfqpoint{0.500431in}{0.933256in}}%
\pgfpathlineto{\pgfqpoint{0.505468in}{0.905493in}}%
\pgfpathlineto{\pgfqpoint{0.507987in}{0.912536in}}%
\pgfpathlineto{\pgfqpoint{0.510506in}{0.895726in}}%
\pgfpathlineto{\pgfqpoint{0.513024in}{0.905698in}}%
\pgfpathlineto{\pgfqpoint{0.515543in}{0.900165in}}%
\pgfpathlineto{\pgfqpoint{0.518062in}{0.908935in}}%
\pgfpathlineto{\pgfqpoint{0.520581in}{0.906255in}}%
\pgfpathlineto{\pgfqpoint{0.523099in}{0.917956in}}%
\pgfpathlineto{\pgfqpoint{0.525618in}{0.906481in}}%
\pgfpathlineto{\pgfqpoint{0.528137in}{0.908607in}}%
\pgfpathlineto{\pgfqpoint{0.530656in}{0.904055in}}%
\pgfpathlineto{\pgfqpoint{0.533174in}{0.892218in}}%
\pgfpathlineto{\pgfqpoint{0.535693in}{0.899592in}}%
\pgfpathlineto{\pgfqpoint{0.538212in}{0.881206in}}%
\pgfpathlineto{\pgfqpoint{0.540731in}{0.894427in}}%
\pgfpathlineto{\pgfqpoint{0.543249in}{0.883352in}}%
\pgfpathlineto{\pgfqpoint{0.545768in}{0.884323in}}%
\pgfpathlineto{\pgfqpoint{0.548287in}{0.879142in}}%
\pgfpathlineto{\pgfqpoint{0.550806in}{0.878069in}}%
\pgfpathlineto{\pgfqpoint{0.553324in}{0.904159in}}%
\pgfpathlineto{\pgfqpoint{0.555843in}{0.908724in}}%
\pgfpathlineto{\pgfqpoint{0.558362in}{0.918346in}}%
\pgfpathlineto{\pgfqpoint{0.560881in}{0.954909in}}%
\pgfpathlineto{\pgfqpoint{0.563399in}{0.971596in}}%
\pgfpathlineto{\pgfqpoint{0.565918in}{0.969415in}}%
\pgfpathlineto{\pgfqpoint{0.568437in}{0.970000in}}%
\pgfpathlineto{\pgfqpoint{0.570956in}{0.962594in}}%
\pgfpathlineto{\pgfqpoint{0.573474in}{0.958047in}}%
\pgfpathlineto{\pgfqpoint{0.575993in}{0.963183in}}%
\pgfpathlineto{\pgfqpoint{0.578512in}{0.969824in}}%
\pgfpathlineto{\pgfqpoint{0.581031in}{0.979348in}}%
\pgfpathlineto{\pgfqpoint{0.583549in}{0.956077in}}%
\pgfpathlineto{\pgfqpoint{0.586068in}{0.946648in}}%
\pgfpathlineto{\pgfqpoint{0.588587in}{0.946262in}}%
\pgfpathlineto{\pgfqpoint{0.591106in}{0.947793in}}%
\pgfpathlineto{\pgfqpoint{0.593624in}{0.935522in}}%
\pgfpathlineto{\pgfqpoint{0.596143in}{0.957296in}}%
\pgfpathlineto{\pgfqpoint{0.598662in}{0.974825in}}%
\pgfpathlineto{\pgfqpoint{0.601181in}{0.975583in}}%
\pgfpathlineto{\pgfqpoint{0.603699in}{0.977058in}}%
\pgfpathlineto{\pgfqpoint{0.606218in}{0.982253in}}%
\pgfpathlineto{\pgfqpoint{0.608737in}{0.988818in}}%
\pgfpathlineto{\pgfqpoint{0.611256in}{1.004465in}}%
\pgfpathlineto{\pgfqpoint{0.613774in}{1.005036in}}%
\pgfpathlineto{\pgfqpoint{0.616293in}{1.009581in}}%
\pgfpathlineto{\pgfqpoint{0.618812in}{1.013699in}}%
\pgfpathlineto{\pgfqpoint{0.621331in}{1.029347in}}%
\pgfpathlineto{\pgfqpoint{0.623849in}{1.054384in}}%
\pgfpathlineto{\pgfqpoint{0.626368in}{1.031366in}}%
\pgfpathlineto{\pgfqpoint{0.628887in}{1.003101in}}%
\pgfpathlineto{\pgfqpoint{0.631406in}{1.003937in}}%
\pgfpathlineto{\pgfqpoint{0.633924in}{1.020885in}}%
\pgfpathlineto{\pgfqpoint{0.636443in}{1.016418in}}%
\pgfpathlineto{\pgfqpoint{0.638962in}{1.007204in}}%
\pgfpathlineto{\pgfqpoint{0.641481in}{0.991211in}}%
\pgfpathlineto{\pgfqpoint{0.643999in}{1.001561in}}%
\pgfpathlineto{\pgfqpoint{0.646518in}{0.996238in}}%
\pgfpathlineto{\pgfqpoint{0.649037in}{1.006290in}}%
\pgfpathlineto{\pgfqpoint{0.651556in}{0.999559in}}%
\pgfpathlineto{\pgfqpoint{0.654074in}{0.986153in}}%
\pgfpathlineto{\pgfqpoint{0.656593in}{0.993370in}}%
\pgfpathlineto{\pgfqpoint{0.659112in}{1.004854in}}%
\pgfpathlineto{\pgfqpoint{0.661631in}{0.999603in}}%
\pgfpathlineto{\pgfqpoint{0.664149in}{1.007938in}}%
\pgfpathlineto{\pgfqpoint{0.666668in}{1.008343in}}%
\pgfpathlineto{\pgfqpoint{0.669187in}{1.013802in}}%
\pgfpathlineto{\pgfqpoint{0.671706in}{1.008285in}}%
\pgfpathlineto{\pgfqpoint{0.674224in}{1.013538in}}%
\pgfpathlineto{\pgfqpoint{0.676743in}{1.009906in}}%
\pgfpathlineto{\pgfqpoint{0.679262in}{1.024651in}}%
\pgfpathlineto{\pgfqpoint{0.681781in}{1.019192in}}%
\pgfpathlineto{\pgfqpoint{0.684299in}{0.993749in}}%
\pgfpathlineto{\pgfqpoint{0.686818in}{0.973520in}}%
\pgfpathlineto{\pgfqpoint{0.689337in}{0.984478in}}%
\pgfpathlineto{\pgfqpoint{0.691856in}{1.000766in}}%
\pgfpathlineto{\pgfqpoint{0.694374in}{1.009337in}}%
\pgfpathlineto{\pgfqpoint{0.696893in}{1.004825in}}%
\pgfpathlineto{\pgfqpoint{0.699412in}{1.007786in}}%
\pgfpathlineto{\pgfqpoint{0.701931in}{0.992367in}}%
\pgfpathlineto{\pgfqpoint{0.704449in}{0.987501in}}%
\pgfpathlineto{\pgfqpoint{0.706968in}{0.967551in}}%
\pgfpathlineto{\pgfqpoint{0.709487in}{0.942350in}}%
\pgfpathlineto{\pgfqpoint{0.712006in}{0.964594in}}%
\pgfpathlineto{\pgfqpoint{0.714524in}{0.966455in}}%
\pgfpathlineto{\pgfqpoint{0.717043in}{0.949117in}}%
\pgfpathlineto{\pgfqpoint{0.719562in}{0.963071in}}%
\pgfpathlineto{\pgfqpoint{0.722081in}{0.941799in}}%
\pgfpathlineto{\pgfqpoint{0.724599in}{0.950659in}}%
\pgfpathlineto{\pgfqpoint{0.727118in}{0.976734in}}%
\pgfpathlineto{\pgfqpoint{0.729637in}{0.984576in}}%
\pgfpathlineto{\pgfqpoint{0.732156in}{0.977123in}}%
\pgfpathlineto{\pgfqpoint{0.734674in}{0.975913in}}%
\pgfpathlineto{\pgfqpoint{0.737193in}{0.978566in}}%
\pgfpathlineto{\pgfqpoint{0.739712in}{0.991528in}}%
\pgfpathlineto{\pgfqpoint{0.742231in}{0.984451in}}%
\pgfpathlineto{\pgfqpoint{0.744749in}{0.994702in}}%
\pgfpathlineto{\pgfqpoint{0.747268in}{0.971280in}}%
\pgfpathlineto{\pgfqpoint{0.749787in}{0.954586in}}%
\pgfpathlineto{\pgfqpoint{0.752306in}{0.960895in}}%
\pgfpathlineto{\pgfqpoint{0.754824in}{0.959309in}}%
\pgfpathlineto{\pgfqpoint{0.757343in}{0.965563in}}%
\pgfpathlineto{\pgfqpoint{0.759862in}{0.994331in}}%
\pgfpathlineto{\pgfqpoint{0.762381in}{0.996264in}}%
\pgfpathlineto{\pgfqpoint{0.764899in}{1.004300in}}%
\pgfpathlineto{\pgfqpoint{0.767418in}{0.989161in}}%
\pgfpathlineto{\pgfqpoint{0.769937in}{0.979921in}}%
\pgfpathlineto{\pgfqpoint{0.772456in}{0.999209in}}%
\pgfpathlineto{\pgfqpoint{0.774974in}{1.007301in}}%
\pgfpathlineto{\pgfqpoint{0.777493in}{0.992964in}}%
\pgfpathlineto{\pgfqpoint{0.780012in}{1.009819in}}%
\pgfpathlineto{\pgfqpoint{0.782531in}{1.014905in}}%
\pgfpathlineto{\pgfqpoint{0.785049in}{1.009082in}}%
\pgfpathlineto{\pgfqpoint{0.787568in}{1.015962in}}%
\pgfpathlineto{\pgfqpoint{0.790087in}{1.014099in}}%
\pgfpathlineto{\pgfqpoint{0.792606in}{1.017947in}}%
\pgfpathlineto{\pgfqpoint{0.795124in}{1.027283in}}%
\pgfpathlineto{\pgfqpoint{0.797643in}{1.040104in}}%
\pgfpathlineto{\pgfqpoint{0.800162in}{1.056100in}}%
\pgfpathlineto{\pgfqpoint{0.802681in}{1.055540in}}%
\pgfpathlineto{\pgfqpoint{0.805199in}{1.062971in}}%
\pgfpathlineto{\pgfqpoint{0.807718in}{1.061851in}}%
\pgfpathlineto{\pgfqpoint{0.810237in}{1.078871in}}%
\pgfpathlineto{\pgfqpoint{0.812756in}{1.069561in}}%
\pgfpathlineto{\pgfqpoint{0.815274in}{1.092277in}}%
\pgfpathlineto{\pgfqpoint{0.817793in}{1.102299in}}%
\pgfpathlineto{\pgfqpoint{0.820312in}{1.094756in}}%
\pgfpathlineto{\pgfqpoint{0.822831in}{1.090117in}}%
\pgfpathlineto{\pgfqpoint{0.825349in}{1.090073in}}%
\pgfpathlineto{\pgfqpoint{0.827868in}{1.056577in}}%
\pgfpathlineto{\pgfqpoint{0.832906in}{1.012517in}}%
\pgfpathlineto{\pgfqpoint{0.835424in}{1.031716in}}%
\pgfpathlineto{\pgfqpoint{0.837943in}{1.020089in}}%
\pgfpathlineto{\pgfqpoint{0.842981in}{1.068652in}}%
\pgfpathlineto{\pgfqpoint{0.845499in}{1.057902in}}%
\pgfpathlineto{\pgfqpoint{0.848018in}{1.085548in}}%
\pgfpathlineto{\pgfqpoint{0.850537in}{1.094048in}}%
\pgfpathlineto{\pgfqpoint{0.853056in}{1.076925in}}%
\pgfpathlineto{\pgfqpoint{0.855574in}{1.081108in}}%
\pgfpathlineto{\pgfqpoint{0.858093in}{1.075450in}}%
\pgfpathlineto{\pgfqpoint{0.860612in}{1.066705in}}%
\pgfpathlineto{\pgfqpoint{0.863131in}{1.065421in}}%
\pgfpathlineto{\pgfqpoint{0.865649in}{1.070834in}}%
\pgfpathlineto{\pgfqpoint{0.868168in}{1.039438in}}%
\pgfpathlineto{\pgfqpoint{0.870687in}{1.027850in}}%
\pgfpathlineto{\pgfqpoint{0.873206in}{1.029277in}}%
\pgfpathlineto{\pgfqpoint{0.875724in}{1.033536in}}%
\pgfpathlineto{\pgfqpoint{0.880762in}{1.043742in}}%
\pgfpathlineto{\pgfqpoint{0.883281in}{1.023370in}}%
\pgfpathlineto{\pgfqpoint{0.885799in}{1.021147in}}%
\pgfpathlineto{\pgfqpoint{0.888318in}{1.008227in}}%
\pgfpathlineto{\pgfqpoint{0.890837in}{1.007421in}}%
\pgfpathlineto{\pgfqpoint{0.893356in}{1.008276in}}%
\pgfpathlineto{\pgfqpoint{0.895874in}{1.008824in}}%
\pgfpathlineto{\pgfqpoint{0.898393in}{0.988802in}}%
\pgfpathlineto{\pgfqpoint{0.900912in}{0.983118in}}%
\pgfpathlineto{\pgfqpoint{0.903431in}{0.984354in}}%
\pgfpathlineto{\pgfqpoint{0.905949in}{0.994892in}}%
\pgfpathlineto{\pgfqpoint{0.908468in}{0.992084in}}%
\pgfpathlineto{\pgfqpoint{0.910987in}{0.995500in}}%
\pgfpathlineto{\pgfqpoint{0.913506in}{0.992252in}}%
\pgfpathlineto{\pgfqpoint{0.916024in}{0.979735in}}%
\pgfpathlineto{\pgfqpoint{0.918543in}{0.970311in}}%
\pgfpathlineto{\pgfqpoint{0.921062in}{0.946714in}}%
\pgfpathlineto{\pgfqpoint{0.923581in}{0.961321in}}%
\pgfpathlineto{\pgfqpoint{0.926099in}{0.946013in}}%
\pgfpathlineto{\pgfqpoint{0.928618in}{0.949080in}}%
\pgfpathlineto{\pgfqpoint{0.931137in}{0.948355in}}%
\pgfpathlineto{\pgfqpoint{0.933656in}{0.963610in}}%
\pgfpathlineto{\pgfqpoint{0.936174in}{0.981423in}}%
\pgfpathlineto{\pgfqpoint{0.938693in}{0.977295in}}%
\pgfpathlineto{\pgfqpoint{0.941212in}{0.998388in}}%
\pgfpathlineto{\pgfqpoint{0.943731in}{0.994533in}}%
\pgfpathlineto{\pgfqpoint{0.946249in}{0.987801in}}%
\pgfpathlineto{\pgfqpoint{0.948768in}{0.982080in}}%
\pgfpathlineto{\pgfqpoint{0.951287in}{0.964463in}}%
\pgfpathlineto{\pgfqpoint{0.953806in}{1.001641in}}%
\pgfpathlineto{\pgfqpoint{0.956324in}{1.008536in}}%
\pgfpathlineto{\pgfqpoint{0.958843in}{0.998591in}}%
\pgfpathlineto{\pgfqpoint{0.961362in}{1.006551in}}%
\pgfpathlineto{\pgfqpoint{0.963881in}{1.017151in}}%
\pgfpathlineto{\pgfqpoint{0.966399in}{1.026163in}}%
\pgfpathlineto{\pgfqpoint{0.968918in}{1.026298in}}%
\pgfpathlineto{\pgfqpoint{0.971437in}{1.045473in}}%
\pgfpathlineto{\pgfqpoint{0.973956in}{1.038306in}}%
\pgfpathlineto{\pgfqpoint{0.976474in}{1.040860in}}%
\pgfpathlineto{\pgfqpoint{0.978993in}{1.029500in}}%
\pgfpathlineto{\pgfqpoint{0.981512in}{1.040793in}}%
\pgfpathlineto{\pgfqpoint{0.984031in}{1.050195in}}%
\pgfpathlineto{\pgfqpoint{0.986549in}{1.067348in}}%
\pgfpathlineto{\pgfqpoint{0.989068in}{1.074342in}}%
\pgfpathlineto{\pgfqpoint{0.991587in}{1.092397in}}%
\pgfpathlineto{\pgfqpoint{0.994106in}{1.076939in}}%
\pgfpathlineto{\pgfqpoint{0.996624in}{1.091060in}}%
\pgfpathlineto{\pgfqpoint{0.999143in}{1.082340in}}%
\pgfpathlineto{\pgfqpoint{1.001662in}{1.071345in}}%
\pgfpathlineto{\pgfqpoint{1.004181in}{1.058102in}}%
\pgfpathlineto{\pgfqpoint{1.006699in}{1.061169in}}%
\pgfpathlineto{\pgfqpoint{1.009218in}{1.049885in}}%
\pgfpathlineto{\pgfqpoint{1.011737in}{1.046296in}}%
\pgfpathlineto{\pgfqpoint{1.014256in}{1.059009in}}%
\pgfpathlineto{\pgfqpoint{1.016774in}{1.063912in}}%
\pgfpathlineto{\pgfqpoint{1.019293in}{1.067366in}}%
\pgfpathlineto{\pgfqpoint{1.021812in}{1.074949in}}%
\pgfpathlineto{\pgfqpoint{1.024331in}{1.085085in}}%
\pgfpathlineto{\pgfqpoint{1.026849in}{1.089188in}}%
\pgfpathlineto{\pgfqpoint{1.029368in}{1.082253in}}%
\pgfpathlineto{\pgfqpoint{1.031887in}{1.080201in}}%
\pgfpathlineto{\pgfqpoint{1.034406in}{1.067288in}}%
\pgfpathlineto{\pgfqpoint{1.036924in}{1.057971in}}%
\pgfpathlineto{\pgfqpoint{1.039443in}{1.056248in}}%
\pgfpathlineto{\pgfqpoint{1.041962in}{1.062735in}}%
\pgfpathlineto{\pgfqpoint{1.044481in}{1.044542in}}%
\pgfpathlineto{\pgfqpoint{1.046999in}{1.050145in}}%
\pgfpathlineto{\pgfqpoint{1.049518in}{1.023863in}}%
\pgfpathlineto{\pgfqpoint{1.052037in}{1.057159in}}%
\pgfpathlineto{\pgfqpoint{1.054556in}{1.038090in}}%
\pgfpathlineto{\pgfqpoint{1.057074in}{1.056965in}}%
\pgfpathlineto{\pgfqpoint{1.059593in}{1.042501in}}%
\pgfpathlineto{\pgfqpoint{1.062112in}{1.030729in}}%
\pgfpathlineto{\pgfqpoint{1.064631in}{1.046523in}}%
\pgfpathlineto{\pgfqpoint{1.067149in}{1.039909in}}%
\pgfpathlineto{\pgfqpoint{1.069668in}{1.028139in}}%
\pgfpathlineto{\pgfqpoint{1.072187in}{1.046640in}}%
\pgfpathlineto{\pgfqpoint{1.074706in}{1.050083in}}%
\pgfpathlineto{\pgfqpoint{1.077224in}{1.057409in}}%
\pgfpathlineto{\pgfqpoint{1.079743in}{1.051917in}}%
\pgfpathlineto{\pgfqpoint{1.082262in}{1.057995in}}%
\pgfpathlineto{\pgfqpoint{1.084781in}{1.075915in}}%
\pgfpathlineto{\pgfqpoint{1.087299in}{1.080289in}}%
\pgfpathlineto{\pgfqpoint{1.089818in}{1.064830in}}%
\pgfpathlineto{\pgfqpoint{1.092337in}{1.079745in}}%
\pgfpathlineto{\pgfqpoint{1.094856in}{1.087897in}}%
\pgfpathlineto{\pgfqpoint{1.097374in}{1.091692in}}%
\pgfpathlineto{\pgfqpoint{1.099893in}{1.077857in}}%
\pgfpathlineto{\pgfqpoint{1.102412in}{1.093932in}}%
\pgfpathlineto{\pgfqpoint{1.107449in}{1.068566in}}%
\pgfpathlineto{\pgfqpoint{1.109968in}{1.086957in}}%
\pgfpathlineto{\pgfqpoint{1.112487in}{1.101911in}}%
\pgfpathlineto{\pgfqpoint{1.115006in}{1.092687in}}%
\pgfpathlineto{\pgfqpoint{1.117524in}{1.102720in}}%
\pgfpathlineto{\pgfqpoint{1.120043in}{1.099614in}}%
\pgfpathlineto{\pgfqpoint{1.122562in}{1.121034in}}%
\pgfpathlineto{\pgfqpoint{1.125081in}{1.111187in}}%
\pgfpathlineto{\pgfqpoint{1.127599in}{1.105014in}}%
\pgfpathlineto{\pgfqpoint{1.130118in}{1.120140in}}%
\pgfpathlineto{\pgfqpoint{1.132637in}{1.137973in}}%
\pgfpathlineto{\pgfqpoint{1.135156in}{1.157039in}}%
\pgfpathlineto{\pgfqpoint{1.137674in}{1.157561in}}%
\pgfpathlineto{\pgfqpoint{1.140193in}{1.146930in}}%
\pgfpathlineto{\pgfqpoint{1.142712in}{1.151705in}}%
\pgfpathlineto{\pgfqpoint{1.145231in}{1.145195in}}%
\pgfpathlineto{\pgfqpoint{1.147749in}{1.165813in}}%
\pgfpathlineto{\pgfqpoint{1.150268in}{1.167943in}}%
\pgfpathlineto{\pgfqpoint{1.152787in}{1.183398in}}%
\pgfpathlineto{\pgfqpoint{1.155306in}{1.212436in}}%
\pgfpathlineto{\pgfqpoint{1.160343in}{1.267624in}}%
\pgfpathlineto{\pgfqpoint{1.162862in}{1.265494in}}%
\pgfpathlineto{\pgfqpoint{1.165381in}{1.270428in}}%
\pgfpathlineto{\pgfqpoint{1.167899in}{1.272238in}}%
\pgfpathlineto{\pgfqpoint{1.170418in}{1.252165in}}%
\pgfpathlineto{\pgfqpoint{1.172937in}{1.245489in}}%
\pgfpathlineto{\pgfqpoint{1.175456in}{1.254702in}}%
\pgfpathlineto{\pgfqpoint{1.177974in}{1.285565in}}%
\pgfpathlineto{\pgfqpoint{1.180493in}{1.275593in}}%
\pgfpathlineto{\pgfqpoint{1.183012in}{1.271341in}}%
\pgfpathlineto{\pgfqpoint{1.185531in}{1.263331in}}%
\pgfpathlineto{\pgfqpoint{1.188049in}{1.277076in}}%
\pgfpathlineto{\pgfqpoint{1.190568in}{1.285234in}}%
\pgfpathlineto{\pgfqpoint{1.193087in}{1.309341in}}%
\pgfpathlineto{\pgfqpoint{1.195606in}{1.303635in}}%
\pgfpathlineto{\pgfqpoint{1.198124in}{1.318746in}}%
\pgfpathlineto{\pgfqpoint{1.200643in}{1.307026in}}%
\pgfpathlineto{\pgfqpoint{1.203162in}{1.309462in}}%
\pgfpathlineto{\pgfqpoint{1.205681in}{1.293851in}}%
\pgfpathlineto{\pgfqpoint{1.208199in}{1.279429in}}%
\pgfpathlineto{\pgfqpoint{1.210718in}{1.272690in}}%
\pgfpathlineto{\pgfqpoint{1.213237in}{1.266329in}}%
\pgfpathlineto{\pgfqpoint{1.215756in}{1.254578in}}%
\pgfpathlineto{\pgfqpoint{1.218274in}{1.232829in}}%
\pgfpathlineto{\pgfqpoint{1.220793in}{1.238255in}}%
\pgfpathlineto{\pgfqpoint{1.223312in}{1.205710in}}%
\pgfpathlineto{\pgfqpoint{1.225831in}{1.225616in}}%
\pgfpathlineto{\pgfqpoint{1.228349in}{1.241782in}}%
\pgfpathlineto{\pgfqpoint{1.230868in}{1.236152in}}%
\pgfpathlineto{\pgfqpoint{1.233387in}{1.225714in}}%
\pgfpathlineto{\pgfqpoint{1.235906in}{1.211162in}}%
\pgfpathlineto{\pgfqpoint{1.238424in}{1.213538in}}%
\pgfpathlineto{\pgfqpoint{1.240943in}{1.227780in}}%
\pgfpathlineto{\pgfqpoint{1.243462in}{1.229468in}}%
\pgfpathlineto{\pgfqpoint{1.245981in}{1.235066in}}%
\pgfpathlineto{\pgfqpoint{1.248499in}{1.226623in}}%
\pgfpathlineto{\pgfqpoint{1.251018in}{1.236718in}}%
\pgfpathlineto{\pgfqpoint{1.253537in}{1.263012in}}%
\pgfpathlineto{\pgfqpoint{1.256056in}{1.270126in}}%
\pgfpathlineto{\pgfqpoint{1.258574in}{1.273518in}}%
\pgfpathlineto{\pgfqpoint{1.261093in}{1.267165in}}%
\pgfpathlineto{\pgfqpoint{1.263612in}{1.273365in}}%
\pgfpathlineto{\pgfqpoint{1.266131in}{1.277349in}}%
\pgfpathlineto{\pgfqpoint{1.268649in}{1.282876in}}%
\pgfpathlineto{\pgfqpoint{1.271168in}{1.264677in}}%
\pgfpathlineto{\pgfqpoint{1.273687in}{1.262843in}}%
\pgfpathlineto{\pgfqpoint{1.276206in}{1.271355in}}%
\pgfpathlineto{\pgfqpoint{1.278724in}{1.264763in}}%
\pgfpathlineto{\pgfqpoint{1.281243in}{1.280510in}}%
\pgfpathlineto{\pgfqpoint{1.283762in}{1.281650in}}%
\pgfpathlineto{\pgfqpoint{1.286281in}{1.309362in}}%
\pgfpathlineto{\pgfqpoint{1.288799in}{1.303819in}}%
\pgfpathlineto{\pgfqpoint{1.291318in}{1.283537in}}%
\pgfpathlineto{\pgfqpoint{1.293837in}{1.279880in}}%
\pgfpathlineto{\pgfqpoint{1.296356in}{1.273757in}}%
\pgfpathlineto{\pgfqpoint{1.298874in}{1.279722in}}%
\pgfpathlineto{\pgfqpoint{1.301393in}{1.271377in}}%
\pgfpathlineto{\pgfqpoint{1.303912in}{1.279638in}}%
\pgfpathlineto{\pgfqpoint{1.306431in}{1.277404in}}%
\pgfpathlineto{\pgfqpoint{1.308949in}{1.288241in}}%
\pgfpathlineto{\pgfqpoint{1.311468in}{1.306054in}}%
\pgfpathlineto{\pgfqpoint{1.313987in}{1.284060in}}%
\pgfpathlineto{\pgfqpoint{1.316506in}{1.272548in}}%
\pgfpathlineto{\pgfqpoint{1.319024in}{1.267138in}}%
\pgfpathlineto{\pgfqpoint{1.321543in}{1.257584in}}%
\pgfpathlineto{\pgfqpoint{1.324062in}{1.266637in}}%
\pgfpathlineto{\pgfqpoint{1.326581in}{1.272169in}}%
\pgfpathlineto{\pgfqpoint{1.329099in}{1.286342in}}%
\pgfpathlineto{\pgfqpoint{1.331618in}{1.304628in}}%
\pgfpathlineto{\pgfqpoint{1.334137in}{1.298698in}}%
\pgfpathlineto{\pgfqpoint{1.336656in}{1.291223in}}%
\pgfpathlineto{\pgfqpoint{1.339174in}{1.294402in}}%
\pgfpathlineto{\pgfqpoint{1.341693in}{1.307099in}}%
\pgfpathlineto{\pgfqpoint{1.344212in}{1.284944in}}%
\pgfpathlineto{\pgfqpoint{1.346731in}{1.254048in}}%
\pgfpathlineto{\pgfqpoint{1.349249in}{1.255049in}}%
\pgfpathlineto{\pgfqpoint{1.351768in}{1.238004in}}%
\pgfpathlineto{\pgfqpoint{1.354287in}{1.247523in}}%
\pgfpathlineto{\pgfqpoint{1.356806in}{1.250413in}}%
\pgfpathlineto{\pgfqpoint{1.359324in}{1.243002in}}%
\pgfpathlineto{\pgfqpoint{1.361843in}{1.233271in}}%
\pgfpathlineto{\pgfqpoint{1.364362in}{1.223953in}}%
\pgfpathlineto{\pgfqpoint{1.366881in}{1.218155in}}%
\pgfpathlineto{\pgfqpoint{1.369399in}{1.219330in}}%
\pgfpathlineto{\pgfqpoint{1.371918in}{1.204161in}}%
\pgfpathlineto{\pgfqpoint{1.374437in}{1.211027in}}%
\pgfpathlineto{\pgfqpoint{1.376956in}{1.209917in}}%
\pgfpathlineto{\pgfqpoint{1.379474in}{1.219910in}}%
\pgfpathlineto{\pgfqpoint{1.381993in}{1.217938in}}%
\pgfpathlineto{\pgfqpoint{1.384512in}{1.213952in}}%
\pgfpathlineto{\pgfqpoint{1.387031in}{1.196910in}}%
\pgfpathlineto{\pgfqpoint{1.389549in}{1.218189in}}%
\pgfpathlineto{\pgfqpoint{1.392068in}{1.211068in}}%
\pgfpathlineto{\pgfqpoint{1.394587in}{1.197394in}}%
\pgfpathlineto{\pgfqpoint{1.397106in}{1.198713in}}%
\pgfpathlineto{\pgfqpoint{1.399624in}{1.207636in}}%
\pgfpathlineto{\pgfqpoint{1.402143in}{1.212688in}}%
\pgfpathlineto{\pgfqpoint{1.404662in}{1.207804in}}%
\pgfpathlineto{\pgfqpoint{1.407181in}{1.205964in}}%
\pgfpathlineto{\pgfqpoint{1.409699in}{1.214822in}}%
\pgfpathlineto{\pgfqpoint{1.412218in}{1.232993in}}%
\pgfpathlineto{\pgfqpoint{1.414737in}{1.248662in}}%
\pgfpathlineto{\pgfqpoint{1.417256in}{1.238715in}}%
\pgfpathlineto{\pgfqpoint{1.419774in}{1.251018in}}%
\pgfpathlineto{\pgfqpoint{1.422293in}{1.258682in}}%
\pgfpathlineto{\pgfqpoint{1.424812in}{1.240163in}}%
\pgfpathlineto{\pgfqpoint{1.427331in}{1.237166in}}%
\pgfpathlineto{\pgfqpoint{1.429849in}{1.226314in}}%
\pgfpathlineto{\pgfqpoint{1.432368in}{1.246771in}}%
\pgfpathlineto{\pgfqpoint{1.434887in}{1.230854in}}%
\pgfpathlineto{\pgfqpoint{1.437406in}{1.236714in}}%
\pgfpathlineto{\pgfqpoint{1.439924in}{1.223034in}}%
\pgfpathlineto{\pgfqpoint{1.442443in}{1.212471in}}%
\pgfpathlineto{\pgfqpoint{1.444962in}{1.210445in}}%
\pgfpathlineto{\pgfqpoint{1.447481in}{1.195828in}}%
\pgfpathlineto{\pgfqpoint{1.449999in}{1.206987in}}%
\pgfpathlineto{\pgfqpoint{1.452518in}{1.220197in}}%
\pgfpathlineto{\pgfqpoint{1.455037in}{1.238663in}}%
\pgfpathlineto{\pgfqpoint{1.457556in}{1.243118in}}%
\pgfpathlineto{\pgfqpoint{1.460074in}{1.250783in}}%
\pgfpathlineto{\pgfqpoint{1.462593in}{1.245686in}}%
\pgfpathlineto{\pgfqpoint{1.465112in}{1.250765in}}%
\pgfpathlineto{\pgfqpoint{1.467631in}{1.246200in}}%
\pgfpathlineto{\pgfqpoint{1.470149in}{1.227918in}}%
\pgfpathlineto{\pgfqpoint{1.472668in}{1.212141in}}%
\pgfpathlineto{\pgfqpoint{1.475187in}{1.246759in}}%
\pgfpathlineto{\pgfqpoint{1.477706in}{1.256131in}}%
\pgfpathlineto{\pgfqpoint{1.480224in}{1.252100in}}%
\pgfpathlineto{\pgfqpoint{1.482743in}{1.261030in}}%
\pgfpathlineto{\pgfqpoint{1.485262in}{1.249393in}}%
\pgfpathlineto{\pgfqpoint{1.487781in}{1.263521in}}%
\pgfpathlineto{\pgfqpoint{1.490299in}{1.281143in}}%
\pgfpathlineto{\pgfqpoint{1.492818in}{1.271267in}}%
\pgfpathlineto{\pgfqpoint{1.495337in}{1.276718in}}%
\pgfpathlineto{\pgfqpoint{1.497856in}{1.267044in}}%
\pgfpathlineto{\pgfqpoint{1.500374in}{1.270607in}}%
\pgfpathlineto{\pgfqpoint{1.502893in}{1.267356in}}%
\pgfpathlineto{\pgfqpoint{1.505412in}{1.243987in}}%
\pgfpathlineto{\pgfqpoint{1.507931in}{1.239383in}}%
\pgfpathlineto{\pgfqpoint{1.510449in}{1.217213in}}%
\pgfpathlineto{\pgfqpoint{1.512968in}{1.231236in}}%
\pgfpathlineto{\pgfqpoint{1.515487in}{1.233548in}}%
\pgfpathlineto{\pgfqpoint{1.518006in}{1.228349in}}%
\pgfpathlineto{\pgfqpoint{1.520524in}{1.247619in}}%
\pgfpathlineto{\pgfqpoint{1.523043in}{1.254237in}}%
\pgfpathlineto{\pgfqpoint{1.525562in}{1.237533in}}%
\pgfpathlineto{\pgfqpoint{1.528081in}{1.251815in}}%
\pgfpathlineto{\pgfqpoint{1.530599in}{1.230546in}}%
\pgfpathlineto{\pgfqpoint{1.533118in}{1.242874in}}%
\pgfpathlineto{\pgfqpoint{1.535637in}{1.246112in}}%
\pgfpathlineto{\pgfqpoint{1.538156in}{1.234019in}}%
\pgfpathlineto{\pgfqpoint{1.540674in}{1.243242in}}%
\pgfpathlineto{\pgfqpoint{1.543193in}{1.244848in}}%
\pgfpathlineto{\pgfqpoint{1.545712in}{1.239242in}}%
\pgfpathlineto{\pgfqpoint{1.548231in}{1.232162in}}%
\pgfpathlineto{\pgfqpoint{1.550749in}{1.235136in}}%
\pgfpathlineto{\pgfqpoint{1.553268in}{1.228310in}}%
\pgfpathlineto{\pgfqpoint{1.555787in}{1.241978in}}%
\pgfpathlineto{\pgfqpoint{1.558306in}{1.249943in}}%
\pgfpathlineto{\pgfqpoint{1.560824in}{1.240685in}}%
\pgfpathlineto{\pgfqpoint{1.563343in}{1.228822in}}%
\pgfpathlineto{\pgfqpoint{1.565862in}{1.223789in}}%
\pgfpathlineto{\pgfqpoint{1.568381in}{1.223545in}}%
\pgfpathlineto{\pgfqpoint{1.570899in}{1.254031in}}%
\pgfpathlineto{\pgfqpoint{1.573418in}{1.267073in}}%
\pgfpathlineto{\pgfqpoint{1.575937in}{1.267461in}}%
\pgfpathlineto{\pgfqpoint{1.578456in}{1.275294in}}%
\pgfpathlineto{\pgfqpoint{1.580974in}{1.292865in}}%
\pgfpathlineto{\pgfqpoint{1.583493in}{1.309512in}}%
\pgfpathlineto{\pgfqpoint{1.586012in}{1.313992in}}%
\pgfpathlineto{\pgfqpoint{1.588531in}{1.285519in}}%
\pgfpathlineto{\pgfqpoint{1.591049in}{1.283194in}}%
\pgfpathlineto{\pgfqpoint{1.593568in}{1.292812in}}%
\pgfpathlineto{\pgfqpoint{1.596087in}{1.256198in}}%
\pgfpathlineto{\pgfqpoint{1.598606in}{1.250012in}}%
\pgfpathlineto{\pgfqpoint{1.601124in}{1.262910in}}%
\pgfpathlineto{\pgfqpoint{1.603643in}{1.258961in}}%
\pgfpathlineto{\pgfqpoint{1.606162in}{1.241963in}}%
\pgfpathlineto{\pgfqpoint{1.608681in}{1.226194in}}%
\pgfpathlineto{\pgfqpoint{1.611199in}{1.229276in}}%
\pgfpathlineto{\pgfqpoint{1.613718in}{1.207364in}}%
\pgfpathlineto{\pgfqpoint{1.616237in}{1.194630in}}%
\pgfpathlineto{\pgfqpoint{1.618756in}{1.228149in}}%
\pgfpathlineto{\pgfqpoint{1.621274in}{1.209091in}}%
\pgfpathlineto{\pgfqpoint{1.623793in}{1.208199in}}%
\pgfpathlineto{\pgfqpoint{1.626312in}{1.219499in}}%
\pgfpathlineto{\pgfqpoint{1.628831in}{1.253466in}}%
\pgfpathlineto{\pgfqpoint{1.631349in}{1.251820in}}%
\pgfpathlineto{\pgfqpoint{1.633868in}{1.240655in}}%
\pgfpathlineto{\pgfqpoint{1.636387in}{1.220010in}}%
\pgfpathlineto{\pgfqpoint{1.638906in}{1.224644in}}%
\pgfpathlineto{\pgfqpoint{1.641424in}{1.222097in}}%
\pgfpathlineto{\pgfqpoint{1.643943in}{1.233997in}}%
\pgfpathlineto{\pgfqpoint{1.646462in}{1.227809in}}%
\pgfpathlineto{\pgfqpoint{1.648981in}{1.233690in}}%
\pgfpathlineto{\pgfqpoint{1.651499in}{1.223801in}}%
\pgfpathlineto{\pgfqpoint{1.654018in}{1.230199in}}%
\pgfpathlineto{\pgfqpoint{1.656537in}{1.230224in}}%
\pgfpathlineto{\pgfqpoint{1.659056in}{1.220379in}}%
\pgfpathlineto{\pgfqpoint{1.661574in}{1.229488in}}%
\pgfpathlineto{\pgfqpoint{1.664093in}{1.231061in}}%
\pgfpathlineto{\pgfqpoint{1.666612in}{1.249857in}}%
\pgfpathlineto{\pgfqpoint{1.669131in}{1.254895in}}%
\pgfpathlineto{\pgfqpoint{1.671649in}{1.243089in}}%
\pgfpathlineto{\pgfqpoint{1.674168in}{1.233373in}}%
\pgfpathlineto{\pgfqpoint{1.676687in}{1.227083in}}%
\pgfpathlineto{\pgfqpoint{1.679206in}{1.229307in}}%
\pgfpathlineto{\pgfqpoint{1.681724in}{1.204440in}}%
\pgfpathlineto{\pgfqpoint{1.684243in}{1.184366in}}%
\pgfpathlineto{\pgfqpoint{1.686762in}{1.192045in}}%
\pgfpathlineto{\pgfqpoint{1.689281in}{1.181029in}}%
\pgfpathlineto{\pgfqpoint{1.691799in}{1.206081in}}%
\pgfpathlineto{\pgfqpoint{1.694318in}{1.202869in}}%
\pgfpathlineto{\pgfqpoint{1.696837in}{1.228504in}}%
\pgfpathlineto{\pgfqpoint{1.699356in}{1.213114in}}%
\pgfpathlineto{\pgfqpoint{1.701874in}{1.204048in}}%
\pgfpathlineto{\pgfqpoint{1.704393in}{1.213865in}}%
\pgfpathlineto{\pgfqpoint{1.706912in}{1.207464in}}%
\pgfpathlineto{\pgfqpoint{1.709431in}{1.215864in}}%
\pgfpathlineto{\pgfqpoint{1.711949in}{1.223919in}}%
\pgfpathlineto{\pgfqpoint{1.714468in}{1.209917in}}%
\pgfpathlineto{\pgfqpoint{1.716987in}{1.225335in}}%
\pgfpathlineto{\pgfqpoint{1.719506in}{1.218443in}}%
\pgfpathlineto{\pgfqpoint{1.722024in}{1.234432in}}%
\pgfpathlineto{\pgfqpoint{1.724543in}{1.251422in}}%
\pgfpathlineto{\pgfqpoint{1.727062in}{1.281882in}}%
\pgfpathlineto{\pgfqpoint{1.729581in}{1.293816in}}%
\pgfpathlineto{\pgfqpoint{1.732099in}{1.309725in}}%
\pgfpathlineto{\pgfqpoint{1.734618in}{1.312698in}}%
\pgfpathlineto{\pgfqpoint{1.737137in}{1.286479in}}%
\pgfpathlineto{\pgfqpoint{1.739656in}{1.274221in}}%
\pgfpathlineto{\pgfqpoint{1.742174in}{1.302917in}}%
\pgfpathlineto{\pgfqpoint{1.744693in}{1.294046in}}%
\pgfpathlineto{\pgfqpoint{1.747212in}{1.311558in}}%
\pgfpathlineto{\pgfqpoint{1.749731in}{1.304250in}}%
\pgfpathlineto{\pgfqpoint{1.752249in}{1.307821in}}%
\pgfpathlineto{\pgfqpoint{1.754768in}{1.326348in}}%
\pgfpathlineto{\pgfqpoint{1.757287in}{1.331862in}}%
\pgfpathlineto{\pgfqpoint{1.759806in}{1.334623in}}%
\pgfpathlineto{\pgfqpoint{1.762324in}{1.345266in}}%
\pgfpathlineto{\pgfqpoint{1.764843in}{1.352761in}}%
\pgfpathlineto{\pgfqpoint{1.767362in}{1.348856in}}%
\pgfpathlineto{\pgfqpoint{1.769881in}{1.371077in}}%
\pgfpathlineto{\pgfqpoint{1.772399in}{1.376417in}}%
\pgfpathlineto{\pgfqpoint{1.774918in}{1.394959in}}%
\pgfpathlineto{\pgfqpoint{1.777437in}{1.384787in}}%
\pgfpathlineto{\pgfqpoint{1.779956in}{1.395314in}}%
\pgfpathlineto{\pgfqpoint{1.782474in}{1.407898in}}%
\pgfpathlineto{\pgfqpoint{1.784993in}{1.413937in}}%
\pgfpathlineto{\pgfqpoint{1.787512in}{1.430623in}}%
\pgfpathlineto{\pgfqpoint{1.790031in}{1.451243in}}%
\pgfpathlineto{\pgfqpoint{1.792549in}{1.470342in}}%
\pgfpathlineto{\pgfqpoint{1.795068in}{1.484520in}}%
\pgfpathlineto{\pgfqpoint{1.797587in}{1.485907in}}%
\pgfpathlineto{\pgfqpoint{1.802624in}{1.470511in}}%
\pgfpathlineto{\pgfqpoint{1.805143in}{1.484435in}}%
\pgfpathlineto{\pgfqpoint{1.807662in}{1.480194in}}%
\pgfpathlineto{\pgfqpoint{1.810181in}{1.478661in}}%
\pgfpathlineto{\pgfqpoint{1.812699in}{1.459853in}}%
\pgfpathlineto{\pgfqpoint{1.815218in}{1.438878in}}%
\pgfpathlineto{\pgfqpoint{1.817737in}{1.445082in}}%
\pgfpathlineto{\pgfqpoint{1.820256in}{1.447700in}}%
\pgfpathlineto{\pgfqpoint{1.822774in}{1.433730in}}%
\pgfpathlineto{\pgfqpoint{1.825293in}{1.464181in}}%
\pgfpathlineto{\pgfqpoint{1.827812in}{1.444286in}}%
\pgfpathlineto{\pgfqpoint{1.830331in}{1.437514in}}%
\pgfpathlineto{\pgfqpoint{1.832849in}{1.445543in}}%
\pgfpathlineto{\pgfqpoint{1.835368in}{1.440766in}}%
\pgfpathlineto{\pgfqpoint{1.837887in}{1.432739in}}%
\pgfpathlineto{\pgfqpoint{1.840406in}{1.461712in}}%
\pgfpathlineto{\pgfqpoint{1.842924in}{1.472664in}}%
\pgfpathlineto{\pgfqpoint{1.845443in}{1.470641in}}%
\pgfpathlineto{\pgfqpoint{1.847962in}{1.461611in}}%
\pgfpathlineto{\pgfqpoint{1.850481in}{1.470433in}}%
\pgfpathlineto{\pgfqpoint{1.852999in}{1.451103in}}%
\pgfpathlineto{\pgfqpoint{1.855518in}{1.445523in}}%
\pgfpathlineto{\pgfqpoint{1.858037in}{1.421943in}}%
\pgfpathlineto{\pgfqpoint{1.860556in}{1.422282in}}%
\pgfpathlineto{\pgfqpoint{1.863074in}{1.425782in}}%
\pgfpathlineto{\pgfqpoint{1.865593in}{1.451900in}}%
\pgfpathlineto{\pgfqpoint{1.868112in}{1.448976in}}%
\pgfpathlineto{\pgfqpoint{1.870631in}{1.445189in}}%
\pgfpathlineto{\pgfqpoint{1.873149in}{1.470038in}}%
\pgfpathlineto{\pgfqpoint{1.875668in}{1.454803in}}%
\pgfpathlineto{\pgfqpoint{1.878187in}{1.461239in}}%
\pgfpathlineto{\pgfqpoint{1.880706in}{1.483082in}}%
\pgfpathlineto{\pgfqpoint{1.883224in}{1.496219in}}%
\pgfpathlineto{\pgfqpoint{1.885743in}{1.477538in}}%
\pgfpathlineto{\pgfqpoint{1.888262in}{1.482411in}}%
\pgfpathlineto{\pgfqpoint{1.890781in}{1.482276in}}%
\pgfpathlineto{\pgfqpoint{1.893299in}{1.481914in}}%
\pgfpathlineto{\pgfqpoint{1.895818in}{1.457090in}}%
\pgfpathlineto{\pgfqpoint{1.898337in}{1.456184in}}%
\pgfpathlineto{\pgfqpoint{1.900856in}{1.446620in}}%
\pgfpathlineto{\pgfqpoint{1.903374in}{1.468810in}}%
\pgfpathlineto{\pgfqpoint{1.905893in}{1.452555in}}%
\pgfpathlineto{\pgfqpoint{1.908412in}{1.460070in}}%
\pgfpathlineto{\pgfqpoint{1.910931in}{1.484435in}}%
\pgfpathlineto{\pgfqpoint{1.913449in}{1.460259in}}%
\pgfpathlineto{\pgfqpoint{1.915968in}{1.486213in}}%
\pgfpathlineto{\pgfqpoint{1.918487in}{1.489429in}}%
\pgfpathlineto{\pgfqpoint{1.921006in}{1.486337in}}%
\pgfpathlineto{\pgfqpoint{1.923524in}{1.490354in}}%
\pgfpathlineto{\pgfqpoint{1.926043in}{1.493984in}}%
\pgfpathlineto{\pgfqpoint{1.928562in}{1.501830in}}%
\pgfpathlineto{\pgfqpoint{1.931081in}{1.505157in}}%
\pgfpathlineto{\pgfqpoint{1.933599in}{1.487846in}}%
\pgfpathlineto{\pgfqpoint{1.936118in}{1.511965in}}%
\pgfpathlineto{\pgfqpoint{1.938637in}{1.533229in}}%
\pgfpathlineto{\pgfqpoint{1.941156in}{1.524654in}}%
\pgfpathlineto{\pgfqpoint{1.943674in}{1.539206in}}%
\pgfpathlineto{\pgfqpoint{1.946193in}{1.533063in}}%
\pgfpathlineto{\pgfqpoint{1.948712in}{1.560778in}}%
\pgfpathlineto{\pgfqpoint{1.951231in}{1.575873in}}%
\pgfpathlineto{\pgfqpoint{1.953749in}{1.596750in}}%
\pgfpathlineto{\pgfqpoint{1.956268in}{1.576143in}}%
\pgfpathlineto{\pgfqpoint{1.958787in}{1.559940in}}%
\pgfpathlineto{\pgfqpoint{1.961306in}{1.547151in}}%
\pgfpathlineto{\pgfqpoint{1.963824in}{1.554553in}}%
\pgfpathlineto{\pgfqpoint{1.966343in}{1.553070in}}%
\pgfpathlineto{\pgfqpoint{1.968862in}{1.567449in}}%
\pgfpathlineto{\pgfqpoint{1.971381in}{1.537979in}}%
\pgfpathlineto{\pgfqpoint{1.973899in}{1.536073in}}%
\pgfpathlineto{\pgfqpoint{1.976418in}{1.543362in}}%
\pgfpathlineto{\pgfqpoint{1.978937in}{1.574661in}}%
\pgfpathlineto{\pgfqpoint{1.981456in}{1.564945in}}%
\pgfpathlineto{\pgfqpoint{1.989012in}{1.537622in}}%
\pgfpathlineto{\pgfqpoint{1.991531in}{1.515970in}}%
\pgfpathlineto{\pgfqpoint{1.994049in}{1.536406in}}%
\pgfpathlineto{\pgfqpoint{1.996568in}{1.549645in}}%
\pgfpathlineto{\pgfqpoint{1.999087in}{1.522425in}}%
\pgfpathlineto{\pgfqpoint{2.001606in}{1.521595in}}%
\pgfpathlineto{\pgfqpoint{2.004124in}{1.512884in}}%
\pgfpathlineto{\pgfqpoint{2.006643in}{1.505925in}}%
\pgfpathlineto{\pgfqpoint{2.009162in}{1.524610in}}%
\pgfpathlineto{\pgfqpoint{2.011681in}{1.547407in}}%
\pgfpathlineto{\pgfqpoint{2.014199in}{1.551369in}}%
\pgfpathlineto{\pgfqpoint{2.016718in}{1.565380in}}%
\pgfpathlineto{\pgfqpoint{2.019237in}{1.580351in}}%
\pgfpathlineto{\pgfqpoint{2.021756in}{1.603415in}}%
\pgfpathlineto{\pgfqpoint{2.024274in}{1.605008in}}%
\pgfpathlineto{\pgfqpoint{2.026793in}{1.605113in}}%
\pgfpathlineto{\pgfqpoint{2.031831in}{1.634673in}}%
\pgfpathlineto{\pgfqpoint{2.034349in}{1.647973in}}%
\pgfpathlineto{\pgfqpoint{2.036868in}{1.651174in}}%
\pgfpathlineto{\pgfqpoint{2.039387in}{1.655545in}}%
\pgfpathlineto{\pgfqpoint{2.041906in}{1.659162in}}%
\pgfpathlineto{\pgfqpoint{2.044424in}{1.663552in}}%
\pgfpathlineto{\pgfqpoint{2.046943in}{1.656849in}}%
\pgfpathlineto{\pgfqpoint{2.049462in}{1.627607in}}%
\pgfpathlineto{\pgfqpoint{2.051981in}{1.617692in}}%
\pgfpathlineto{\pgfqpoint{2.054499in}{1.591727in}}%
\pgfpathlineto{\pgfqpoint{2.057018in}{1.601257in}}%
\pgfpathlineto{\pgfqpoint{2.059537in}{1.577778in}}%
\pgfpathlineto{\pgfqpoint{2.062056in}{1.559947in}}%
\pgfpathlineto{\pgfqpoint{2.064574in}{1.561555in}}%
\pgfpathlineto{\pgfqpoint{2.067093in}{1.534953in}}%
\pgfpathlineto{\pgfqpoint{2.069612in}{1.526426in}}%
\pgfpathlineto{\pgfqpoint{2.072131in}{1.527311in}}%
\pgfpathlineto{\pgfqpoint{2.074649in}{1.519519in}}%
\pgfpathlineto{\pgfqpoint{2.077168in}{1.508092in}}%
\pgfpathlineto{\pgfqpoint{2.079687in}{1.477873in}}%
\pgfpathlineto{\pgfqpoint{2.082206in}{1.470103in}}%
\pgfpathlineto{\pgfqpoint{2.084724in}{1.473380in}}%
\pgfpathlineto{\pgfqpoint{2.087243in}{1.481434in}}%
\pgfpathlineto{\pgfqpoint{2.089762in}{1.479822in}}%
\pgfpathlineto{\pgfqpoint{2.092281in}{1.481028in}}%
\pgfpathlineto{\pgfqpoint{2.094799in}{1.462437in}}%
\pgfpathlineto{\pgfqpoint{2.097318in}{1.458061in}}%
\pgfpathlineto{\pgfqpoint{2.099837in}{1.470067in}}%
\pgfpathlineto{\pgfqpoint{2.102356in}{1.478053in}}%
\pgfpathlineto{\pgfqpoint{2.104874in}{1.472522in}}%
\pgfpathlineto{\pgfqpoint{2.107393in}{1.472484in}}%
\pgfpathlineto{\pgfqpoint{2.109912in}{1.482127in}}%
\pgfpathlineto{\pgfqpoint{2.112431in}{1.481811in}}%
\pgfpathlineto{\pgfqpoint{2.114949in}{1.474735in}}%
\pgfpathlineto{\pgfqpoint{2.117468in}{1.451106in}}%
\pgfpathlineto{\pgfqpoint{2.119987in}{1.456468in}}%
\pgfpathlineto{\pgfqpoint{2.122506in}{1.484861in}}%
\pgfpathlineto{\pgfqpoint{2.125024in}{1.458853in}}%
\pgfpathlineto{\pgfqpoint{2.127543in}{1.522485in}}%
\pgfpathlineto{\pgfqpoint{2.130062in}{1.518381in}}%
\pgfpathlineto{\pgfqpoint{2.132581in}{1.509684in}}%
\pgfpathlineto{\pgfqpoint{2.135099in}{1.512045in}}%
\pgfpathlineto{\pgfqpoint{2.137618in}{1.500487in}}%
\pgfpathlineto{\pgfqpoint{2.140137in}{1.498406in}}%
\pgfpathlineto{\pgfqpoint{2.142656in}{1.513599in}}%
\pgfpathlineto{\pgfqpoint{2.145174in}{1.514615in}}%
\pgfpathlineto{\pgfqpoint{2.147693in}{1.501301in}}%
\pgfpathlineto{\pgfqpoint{2.150212in}{1.499967in}}%
\pgfpathlineto{\pgfqpoint{2.152731in}{1.510614in}}%
\pgfpathlineto{\pgfqpoint{2.155249in}{1.524791in}}%
\pgfpathlineto{\pgfqpoint{2.157768in}{1.522587in}}%
\pgfpathlineto{\pgfqpoint{2.160287in}{1.533232in}}%
\pgfpathlineto{\pgfqpoint{2.162806in}{1.531276in}}%
\pgfpathlineto{\pgfqpoint{2.165324in}{1.520857in}}%
\pgfpathlineto{\pgfqpoint{2.167843in}{1.517535in}}%
\pgfpathlineto{\pgfqpoint{2.170362in}{1.513167in}}%
\pgfpathlineto{\pgfqpoint{2.172881in}{1.534762in}}%
\pgfpathlineto{\pgfqpoint{2.175399in}{1.521374in}}%
\pgfpathlineto{\pgfqpoint{2.177918in}{1.537463in}}%
\pgfpathlineto{\pgfqpoint{2.180437in}{1.539582in}}%
\pgfpathlineto{\pgfqpoint{2.182956in}{1.540054in}}%
\pgfpathlineto{\pgfqpoint{2.185474in}{1.547311in}}%
\pgfpathlineto{\pgfqpoint{2.187993in}{1.549304in}}%
\pgfpathlineto{\pgfqpoint{2.190512in}{1.564783in}}%
\pgfpathlineto{\pgfqpoint{2.193031in}{1.532970in}}%
\pgfpathlineto{\pgfqpoint{2.195549in}{1.527932in}}%
\pgfpathlineto{\pgfqpoint{2.198068in}{1.548098in}}%
\pgfpathlineto{\pgfqpoint{2.200587in}{1.582428in}}%
\pgfpathlineto{\pgfqpoint{2.203106in}{1.566133in}}%
\pgfpathlineto{\pgfqpoint{2.205624in}{1.545234in}}%
\pgfpathlineto{\pgfqpoint{2.208143in}{1.537921in}}%
\pgfpathlineto{\pgfqpoint{2.210662in}{1.570036in}}%
\pgfpathlineto{\pgfqpoint{2.213181in}{1.562402in}}%
\pgfpathlineto{\pgfqpoint{2.215699in}{1.575780in}}%
\pgfpathlineto{\pgfqpoint{2.218218in}{1.556243in}}%
\pgfpathlineto{\pgfqpoint{2.220737in}{1.553781in}}%
\pgfpathlineto{\pgfqpoint{2.223256in}{1.548477in}}%
\pgfpathlineto{\pgfqpoint{2.225774in}{1.557814in}}%
\pgfpathlineto{\pgfqpoint{2.228293in}{1.554872in}}%
\pgfpathlineto{\pgfqpoint{2.230812in}{1.544528in}}%
\pgfpathlineto{\pgfqpoint{2.233331in}{1.544835in}}%
\pgfpathlineto{\pgfqpoint{2.235849in}{1.579864in}}%
\pgfpathlineto{\pgfqpoint{2.238368in}{1.600413in}}%
\pgfpathlineto{\pgfqpoint{2.240887in}{1.588505in}}%
\pgfpathlineto{\pgfqpoint{2.243406in}{1.602840in}}%
\pgfpathlineto{\pgfqpoint{2.245924in}{1.565479in}}%
\pgfpathlineto{\pgfqpoint{2.248443in}{1.584205in}}%
\pgfpathlineto{\pgfqpoint{2.250962in}{1.588843in}}%
\pgfpathlineto{\pgfqpoint{2.253481in}{1.564066in}}%
\pgfpathlineto{\pgfqpoint{2.255999in}{1.571750in}}%
\pgfpathlineto{\pgfqpoint{2.258518in}{1.583671in}}%
\pgfpathlineto{\pgfqpoint{2.261037in}{1.578578in}}%
\pgfpathlineto{\pgfqpoint{2.266074in}{1.543275in}}%
\pgfpathlineto{\pgfqpoint{2.268593in}{1.510711in}}%
\pgfpathlineto{\pgfqpoint{2.271112in}{1.502820in}}%
\pgfpathlineto{\pgfqpoint{2.273631in}{1.485384in}}%
\pgfpathlineto{\pgfqpoint{2.276149in}{1.500078in}}%
\pgfpathlineto{\pgfqpoint{2.278668in}{1.481377in}}%
\pgfpathlineto{\pgfqpoint{2.281187in}{1.488122in}}%
\pgfpathlineto{\pgfqpoint{2.283706in}{1.501768in}}%
\pgfpathlineto{\pgfqpoint{2.286224in}{1.514532in}}%
\pgfpathlineto{\pgfqpoint{2.288743in}{1.499300in}}%
\pgfpathlineto{\pgfqpoint{2.291262in}{1.489121in}}%
\pgfpathlineto{\pgfqpoint{2.293781in}{1.485104in}}%
\pgfpathlineto{\pgfqpoint{2.298818in}{1.483615in}}%
\pgfpathlineto{\pgfqpoint{2.301337in}{1.474668in}}%
\pgfpathlineto{\pgfqpoint{2.303856in}{1.455810in}}%
\pgfpathlineto{\pgfqpoint{2.306374in}{1.464797in}}%
\pgfpathlineto{\pgfqpoint{2.308893in}{1.442770in}}%
\pgfpathlineto{\pgfqpoint{2.311412in}{1.442666in}}%
\pgfpathlineto{\pgfqpoint{2.313931in}{1.439671in}}%
\pgfpathlineto{\pgfqpoint{2.316449in}{1.434787in}}%
\pgfpathlineto{\pgfqpoint{2.318968in}{1.423909in}}%
\pgfpathlineto{\pgfqpoint{2.324006in}{1.485312in}}%
\pgfpathlineto{\pgfqpoint{2.326524in}{1.482733in}}%
\pgfpathlineto{\pgfqpoint{2.329043in}{1.497723in}}%
\pgfpathlineto{\pgfqpoint{2.331562in}{1.478289in}}%
\pgfpathlineto{\pgfqpoint{2.334081in}{1.473563in}}%
\pgfpathlineto{\pgfqpoint{2.336599in}{1.468206in}}%
\pgfpathlineto{\pgfqpoint{2.339118in}{1.461523in}}%
\pgfpathlineto{\pgfqpoint{2.341637in}{1.484127in}}%
\pgfpathlineto{\pgfqpoint{2.344156in}{1.488344in}}%
\pgfpathlineto{\pgfqpoint{2.346674in}{1.506317in}}%
\pgfpathlineto{\pgfqpoint{2.349193in}{1.485433in}}%
\pgfpathlineto{\pgfqpoint{2.351712in}{1.515543in}}%
\pgfpathlineto{\pgfqpoint{2.354231in}{1.522426in}}%
\pgfpathlineto{\pgfqpoint{2.356749in}{1.509280in}}%
\pgfpathlineto{\pgfqpoint{2.359268in}{1.526208in}}%
\pgfpathlineto{\pgfqpoint{2.361787in}{1.515932in}}%
\pgfpathlineto{\pgfqpoint{2.364306in}{1.534810in}}%
\pgfpathlineto{\pgfqpoint{2.366824in}{1.533611in}}%
\pgfpathlineto{\pgfqpoint{2.369343in}{1.517728in}}%
\pgfpathlineto{\pgfqpoint{2.371862in}{1.506974in}}%
\pgfpathlineto{\pgfqpoint{2.374381in}{1.516869in}}%
\pgfpathlineto{\pgfqpoint{2.376899in}{1.534198in}}%
\pgfpathlineto{\pgfqpoint{2.379418in}{1.518548in}}%
\pgfpathlineto{\pgfqpoint{2.381937in}{1.516738in}}%
\pgfpathlineto{\pgfqpoint{2.384456in}{1.541397in}}%
\pgfpathlineto{\pgfqpoint{2.386974in}{1.556807in}}%
\pgfpathlineto{\pgfqpoint{2.389493in}{1.539382in}}%
\pgfpathlineto{\pgfqpoint{2.392012in}{1.516716in}}%
\pgfpathlineto{\pgfqpoint{2.394531in}{1.530437in}}%
\pgfpathlineto{\pgfqpoint{2.397049in}{1.538692in}}%
\pgfpathlineto{\pgfqpoint{2.399568in}{1.541284in}}%
\pgfpathlineto{\pgfqpoint{2.402087in}{1.542713in}}%
\pgfpathlineto{\pgfqpoint{2.404606in}{1.515816in}}%
\pgfpathlineto{\pgfqpoint{2.407124in}{1.531152in}}%
\pgfpathlineto{\pgfqpoint{2.409643in}{1.540536in}}%
\pgfpathlineto{\pgfqpoint{2.412162in}{1.527213in}}%
\pgfpathlineto{\pgfqpoint{2.414681in}{1.519782in}}%
\pgfpathlineto{\pgfqpoint{2.417199in}{1.543223in}}%
\pgfpathlineto{\pgfqpoint{2.419718in}{1.530075in}}%
\pgfpathlineto{\pgfqpoint{2.422237in}{1.533839in}}%
\pgfpathlineto{\pgfqpoint{2.424756in}{1.560886in}}%
\pgfpathlineto{\pgfqpoint{2.427274in}{1.552665in}}%
\pgfpathlineto{\pgfqpoint{2.429793in}{1.573783in}}%
\pgfpathlineto{\pgfqpoint{2.432312in}{1.556939in}}%
\pgfpathlineto{\pgfqpoint{2.434831in}{1.544669in}}%
\pgfpathlineto{\pgfqpoint{2.437349in}{1.536268in}}%
\pgfpathlineto{\pgfqpoint{2.439868in}{1.563982in}}%
\pgfpathlineto{\pgfqpoint{2.442387in}{1.545907in}}%
\pgfpathlineto{\pgfqpoint{2.444906in}{1.543616in}}%
\pgfpathlineto{\pgfqpoint{2.447424in}{1.546657in}}%
\pgfpathlineto{\pgfqpoint{2.449943in}{1.517128in}}%
\pgfpathlineto{\pgfqpoint{2.452462in}{1.498860in}}%
\pgfpathlineto{\pgfqpoint{2.454981in}{1.506518in}}%
\pgfpathlineto{\pgfqpoint{2.457499in}{1.519264in}}%
\pgfpathlineto{\pgfqpoint{2.460018in}{1.525009in}}%
\pgfpathlineto{\pgfqpoint{2.462537in}{1.519214in}}%
\pgfpathlineto{\pgfqpoint{2.465056in}{1.505324in}}%
\pgfpathlineto{\pgfqpoint{2.467574in}{1.532624in}}%
\pgfpathlineto{\pgfqpoint{2.470093in}{1.524153in}}%
\pgfpathlineto{\pgfqpoint{2.472612in}{1.553172in}}%
\pgfpathlineto{\pgfqpoint{2.475131in}{1.561814in}}%
\pgfpathlineto{\pgfqpoint{2.477649in}{1.571379in}}%
\pgfpathlineto{\pgfqpoint{2.480168in}{1.558057in}}%
\pgfpathlineto{\pgfqpoint{2.482687in}{1.543486in}}%
\pgfpathlineto{\pgfqpoint{2.485206in}{1.514281in}}%
\pgfpathlineto{\pgfqpoint{2.487724in}{1.487868in}}%
\pgfpathlineto{\pgfqpoint{2.490243in}{1.491485in}}%
\pgfpathlineto{\pgfqpoint{2.492762in}{1.499532in}}%
\pgfpathlineto{\pgfqpoint{2.495281in}{1.504478in}}%
\pgfpathlineto{\pgfqpoint{2.497799in}{1.507118in}}%
\pgfpathlineto{\pgfqpoint{2.500318in}{1.513021in}}%
\pgfpathlineto{\pgfqpoint{2.502837in}{1.516449in}}%
\pgfpathlineto{\pgfqpoint{2.505356in}{1.503153in}}%
\pgfpathlineto{\pgfqpoint{2.507874in}{1.525347in}}%
\pgfpathlineto{\pgfqpoint{2.510393in}{1.539272in}}%
\pgfpathlineto{\pgfqpoint{2.512912in}{1.539028in}}%
\pgfpathlineto{\pgfqpoint{2.515431in}{1.530805in}}%
\pgfpathlineto{\pgfqpoint{2.517949in}{1.507524in}}%
\pgfpathlineto{\pgfqpoint{2.520468in}{1.503827in}}%
\pgfpathlineto{\pgfqpoint{2.522987in}{1.511224in}}%
\pgfpathlineto{\pgfqpoint{2.525506in}{1.490679in}}%
\pgfpathlineto{\pgfqpoint{2.528024in}{1.509192in}}%
\pgfpathlineto{\pgfqpoint{2.530543in}{1.524481in}}%
\pgfpathlineto{\pgfqpoint{2.533062in}{1.516555in}}%
\pgfpathlineto{\pgfqpoint{2.535581in}{1.521427in}}%
\pgfpathlineto{\pgfqpoint{2.538099in}{1.520687in}}%
\pgfpathlineto{\pgfqpoint{2.540618in}{1.518088in}}%
\pgfpathlineto{\pgfqpoint{2.543137in}{1.532193in}}%
\pgfpathlineto{\pgfqpoint{2.545656in}{1.523896in}}%
\pgfpathlineto{\pgfqpoint{2.548174in}{1.508437in}}%
\pgfpathlineto{\pgfqpoint{2.550693in}{1.513199in}}%
\pgfpathlineto{\pgfqpoint{2.553212in}{1.507448in}}%
\pgfpathlineto{\pgfqpoint{2.555731in}{1.511644in}}%
\pgfpathlineto{\pgfqpoint{2.558249in}{1.485390in}}%
\pgfpathlineto{\pgfqpoint{2.560768in}{1.485949in}}%
\pgfpathlineto{\pgfqpoint{2.563287in}{1.480858in}}%
\pgfpathlineto{\pgfqpoint{2.565806in}{1.476917in}}%
\pgfpathlineto{\pgfqpoint{2.568324in}{1.471702in}}%
\pgfpathlineto{\pgfqpoint{2.570843in}{1.471175in}}%
\pgfpathlineto{\pgfqpoint{2.573362in}{1.454342in}}%
\pgfpathlineto{\pgfqpoint{2.575881in}{1.483702in}}%
\pgfpathlineto{\pgfqpoint{2.578399in}{1.474845in}}%
\pgfpathlineto{\pgfqpoint{2.580918in}{1.475102in}}%
\pgfpathlineto{\pgfqpoint{2.585956in}{1.438698in}}%
\pgfpathlineto{\pgfqpoint{2.588474in}{1.462466in}}%
\pgfpathlineto{\pgfqpoint{2.590993in}{1.452738in}}%
\pgfpathlineto{\pgfqpoint{2.593512in}{1.461543in}}%
\pgfpathlineto{\pgfqpoint{2.596031in}{1.472242in}}%
\pgfpathlineto{\pgfqpoint{2.598549in}{1.456680in}}%
\pgfpathlineto{\pgfqpoint{2.601068in}{1.466160in}}%
\pgfpathlineto{\pgfqpoint{2.603587in}{1.481134in}}%
\pgfpathlineto{\pgfqpoint{2.606106in}{1.474653in}}%
\pgfpathlineto{\pgfqpoint{2.608624in}{1.469547in}}%
\pgfpathlineto{\pgfqpoint{2.611143in}{1.461674in}}%
\pgfpathlineto{\pgfqpoint{2.613662in}{1.449338in}}%
\pgfpathlineto{\pgfqpoint{2.616181in}{1.443880in}}%
\pgfpathlineto{\pgfqpoint{2.618699in}{1.444517in}}%
\pgfpathlineto{\pgfqpoint{2.621218in}{1.461347in}}%
\pgfpathlineto{\pgfqpoint{2.623737in}{1.454888in}}%
\pgfpathlineto{\pgfqpoint{2.626256in}{1.432689in}}%
\pgfpathlineto{\pgfqpoint{2.628774in}{1.460569in}}%
\pgfpathlineto{\pgfqpoint{2.631293in}{1.433834in}}%
\pgfpathlineto{\pgfqpoint{2.633812in}{1.438776in}}%
\pgfpathlineto{\pgfqpoint{2.636331in}{1.449052in}}%
\pgfpathlineto{\pgfqpoint{2.638849in}{1.445235in}}%
\pgfpathlineto{\pgfqpoint{2.641368in}{1.439692in}}%
\pgfpathlineto{\pgfqpoint{2.643887in}{1.442908in}}%
\pgfpathlineto{\pgfqpoint{2.646406in}{1.434014in}}%
\pgfpathlineto{\pgfqpoint{2.648924in}{1.412943in}}%
\pgfpathlineto{\pgfqpoint{2.651443in}{1.404641in}}%
\pgfpathlineto{\pgfqpoint{2.653962in}{1.403415in}}%
\pgfpathlineto{\pgfqpoint{2.656481in}{1.398161in}}%
\pgfpathlineto{\pgfqpoint{2.658999in}{1.382193in}}%
\pgfpathlineto{\pgfqpoint{2.661518in}{1.380464in}}%
\pgfpathlineto{\pgfqpoint{2.664037in}{1.404943in}}%
\pgfpathlineto{\pgfqpoint{2.666556in}{1.413679in}}%
\pgfpathlineto{\pgfqpoint{2.669074in}{1.387074in}}%
\pgfpathlineto{\pgfqpoint{2.671593in}{1.378357in}}%
\pgfpathlineto{\pgfqpoint{2.674112in}{1.362360in}}%
\pgfpathlineto{\pgfqpoint{2.676631in}{1.359878in}}%
\pgfpathlineto{\pgfqpoint{2.679149in}{1.348442in}}%
\pgfpathlineto{\pgfqpoint{2.681668in}{1.346739in}}%
\pgfpathlineto{\pgfqpoint{2.684187in}{1.325455in}}%
\pgfpathlineto{\pgfqpoint{2.686706in}{1.325672in}}%
\pgfpathlineto{\pgfqpoint{2.689224in}{1.307122in}}%
\pgfpathlineto{\pgfqpoint{2.691743in}{1.322164in}}%
\pgfpathlineto{\pgfqpoint{2.694262in}{1.316556in}}%
\pgfpathlineto{\pgfqpoint{2.696781in}{1.333218in}}%
\pgfpathlineto{\pgfqpoint{2.699299in}{1.313954in}}%
\pgfpathlineto{\pgfqpoint{2.701818in}{1.305114in}}%
\pgfpathlineto{\pgfqpoint{2.704337in}{1.303173in}}%
\pgfpathlineto{\pgfqpoint{2.706856in}{1.285853in}}%
\pgfpathlineto{\pgfqpoint{2.709374in}{1.288356in}}%
\pgfpathlineto{\pgfqpoint{2.711893in}{1.255686in}}%
\pgfpathlineto{\pgfqpoint{2.714412in}{1.244539in}}%
\pgfpathlineto{\pgfqpoint{2.716931in}{1.232938in}}%
\pgfpathlineto{\pgfqpoint{2.719449in}{1.215513in}}%
\pgfpathlineto{\pgfqpoint{2.721968in}{1.228316in}}%
\pgfpathlineto{\pgfqpoint{2.724487in}{1.253421in}}%
\pgfpathlineto{\pgfqpoint{2.727006in}{1.237388in}}%
\pgfpathlineto{\pgfqpoint{2.729524in}{1.254404in}}%
\pgfpathlineto{\pgfqpoint{2.732043in}{1.253714in}}%
\pgfpathlineto{\pgfqpoint{2.734562in}{1.238557in}}%
\pgfpathlineto{\pgfqpoint{2.737081in}{1.238706in}}%
\pgfpathlineto{\pgfqpoint{2.739599in}{1.233336in}}%
\pgfpathlineto{\pgfqpoint{2.742118in}{1.223984in}}%
\pgfpathlineto{\pgfqpoint{2.744637in}{1.232757in}}%
\pgfpathlineto{\pgfqpoint{2.747156in}{1.221218in}}%
\pgfpathlineto{\pgfqpoint{2.749674in}{1.213267in}}%
\pgfpathlineto{\pgfqpoint{2.752193in}{1.199352in}}%
\pgfpathlineto{\pgfqpoint{2.754712in}{1.198865in}}%
\pgfpathlineto{\pgfqpoint{2.757231in}{1.196500in}}%
\pgfpathlineto{\pgfqpoint{2.759749in}{1.204999in}}%
\pgfpathlineto{\pgfqpoint{2.762268in}{1.189623in}}%
\pgfpathlineto{\pgfqpoint{2.764787in}{1.191346in}}%
\pgfpathlineto{\pgfqpoint{2.767306in}{1.196810in}}%
\pgfpathlineto{\pgfqpoint{2.769824in}{1.197240in}}%
\pgfpathlineto{\pgfqpoint{2.772343in}{1.168145in}}%
\pgfpathlineto{\pgfqpoint{2.774862in}{1.177620in}}%
\pgfpathlineto{\pgfqpoint{2.777381in}{1.164337in}}%
\pgfpathlineto{\pgfqpoint{2.779899in}{1.165009in}}%
\pgfpathlineto{\pgfqpoint{2.782418in}{1.176679in}}%
\pgfpathlineto{\pgfqpoint{2.784937in}{1.174107in}}%
\pgfpathlineto{\pgfqpoint{2.789974in}{1.138307in}}%
\pgfpathlineto{\pgfqpoint{2.792493in}{1.118806in}}%
\pgfpathlineto{\pgfqpoint{2.795012in}{1.120931in}}%
\pgfpathlineto{\pgfqpoint{2.797531in}{1.113580in}}%
\pgfpathlineto{\pgfqpoint{2.800049in}{1.123250in}}%
\pgfpathlineto{\pgfqpoint{2.802568in}{1.105477in}}%
\pgfpathlineto{\pgfqpoint{2.805087in}{1.113335in}}%
\pgfpathlineto{\pgfqpoint{2.807606in}{1.128181in}}%
\pgfpathlineto{\pgfqpoint{2.810124in}{1.129089in}}%
\pgfpathlineto{\pgfqpoint{2.812643in}{1.129292in}}%
\pgfpathlineto{\pgfqpoint{2.815162in}{1.130347in}}%
\pgfpathlineto{\pgfqpoint{2.817681in}{1.165216in}}%
\pgfpathlineto{\pgfqpoint{2.820199in}{1.161729in}}%
\pgfpathlineto{\pgfqpoint{2.825237in}{1.139151in}}%
\pgfpathlineto{\pgfqpoint{2.827756in}{1.164087in}}%
\pgfpathlineto{\pgfqpoint{2.830274in}{1.149619in}}%
\pgfpathlineto{\pgfqpoint{2.832793in}{1.158419in}}%
\pgfpathlineto{\pgfqpoint{2.835312in}{1.148007in}}%
\pgfpathlineto{\pgfqpoint{2.837831in}{1.130658in}}%
\pgfpathlineto{\pgfqpoint{2.840349in}{1.148981in}}%
\pgfpathlineto{\pgfqpoint{2.842868in}{1.150751in}}%
\pgfpathlineto{\pgfqpoint{2.845387in}{1.162534in}}%
\pgfpathlineto{\pgfqpoint{2.847906in}{1.181907in}}%
\pgfpathlineto{\pgfqpoint{2.850424in}{1.169447in}}%
\pgfpathlineto{\pgfqpoint{2.852943in}{1.153220in}}%
\pgfpathlineto{\pgfqpoint{2.855462in}{1.127907in}}%
\pgfpathlineto{\pgfqpoint{2.857981in}{1.152269in}}%
\pgfpathlineto{\pgfqpoint{2.860499in}{1.173504in}}%
\pgfpathlineto{\pgfqpoint{2.863018in}{1.199833in}}%
\pgfpathlineto{\pgfqpoint{2.865537in}{1.193512in}}%
\pgfpathlineto{\pgfqpoint{2.868056in}{1.196997in}}%
\pgfpathlineto{\pgfqpoint{2.870574in}{1.232785in}}%
\pgfpathlineto{\pgfqpoint{2.873093in}{1.230287in}}%
\pgfpathlineto{\pgfqpoint{2.875612in}{1.246388in}}%
\pgfpathlineto{\pgfqpoint{2.878131in}{1.250255in}}%
\pgfpathlineto{\pgfqpoint{2.880649in}{1.255903in}}%
\pgfpathlineto{\pgfqpoint{2.883168in}{1.253416in}}%
\pgfpathlineto{\pgfqpoint{2.885687in}{1.228078in}}%
\pgfpathlineto{\pgfqpoint{2.888206in}{1.241987in}}%
\pgfpathlineto{\pgfqpoint{2.890724in}{1.248046in}}%
\pgfpathlineto{\pgfqpoint{2.893243in}{1.260059in}}%
\pgfpathlineto{\pgfqpoint{2.895762in}{1.267592in}}%
\pgfpathlineto{\pgfqpoint{2.898281in}{1.285197in}}%
\pgfpathlineto{\pgfqpoint{2.900799in}{1.297556in}}%
\pgfpathlineto{\pgfqpoint{2.903318in}{1.309252in}}%
\pgfpathlineto{\pgfqpoint{2.905837in}{1.327714in}}%
\pgfpathlineto{\pgfqpoint{2.910874in}{1.348027in}}%
\pgfpathlineto{\pgfqpoint{2.913393in}{1.327722in}}%
\pgfpathlineto{\pgfqpoint{2.915912in}{1.320440in}}%
\pgfpathlineto{\pgfqpoint{2.918431in}{1.350493in}}%
\pgfpathlineto{\pgfqpoint{2.920949in}{1.336543in}}%
\pgfpathlineto{\pgfqpoint{2.923468in}{1.340117in}}%
\pgfpathlineto{\pgfqpoint{2.925987in}{1.303038in}}%
\pgfpathlineto{\pgfqpoint{2.928506in}{1.310914in}}%
\pgfpathlineto{\pgfqpoint{2.931024in}{1.312035in}}%
\pgfpathlineto{\pgfqpoint{2.933543in}{1.291846in}}%
\pgfpathlineto{\pgfqpoint{2.936062in}{1.250660in}}%
\pgfpathlineto{\pgfqpoint{2.938581in}{1.223695in}}%
\pgfpathlineto{\pgfqpoint{2.941099in}{1.234358in}}%
\pgfpathlineto{\pgfqpoint{2.943618in}{1.231425in}}%
\pgfpathlineto{\pgfqpoint{2.946137in}{1.214494in}}%
\pgfpathlineto{\pgfqpoint{2.948656in}{1.212200in}}%
\pgfpathlineto{\pgfqpoint{2.951174in}{1.216763in}}%
\pgfpathlineto{\pgfqpoint{2.953693in}{1.227185in}}%
\pgfpathlineto{\pgfqpoint{2.956212in}{1.241441in}}%
\pgfpathlineto{\pgfqpoint{2.958731in}{1.242407in}}%
\pgfpathlineto{\pgfqpoint{2.961249in}{1.279954in}}%
\pgfpathlineto{\pgfqpoint{2.963768in}{1.276952in}}%
\pgfpathlineto{\pgfqpoint{2.966287in}{1.287089in}}%
\pgfpathlineto{\pgfqpoint{2.968806in}{1.283046in}}%
\pgfpathlineto{\pgfqpoint{2.971324in}{1.279242in}}%
\pgfpathlineto{\pgfqpoint{2.973843in}{1.302219in}}%
\pgfpathlineto{\pgfqpoint{2.976362in}{1.287285in}}%
\pgfpathlineto{\pgfqpoint{2.978881in}{1.304490in}}%
\pgfpathlineto{\pgfqpoint{2.981399in}{1.293630in}}%
\pgfpathlineto{\pgfqpoint{2.983918in}{1.295536in}}%
\pgfpathlineto{\pgfqpoint{2.986437in}{1.289849in}}%
\pgfpathlineto{\pgfqpoint{2.988956in}{1.293428in}}%
\pgfpathlineto{\pgfqpoint{2.991474in}{1.299778in}}%
\pgfpathlineto{\pgfqpoint{2.993993in}{1.300180in}}%
\pgfpathlineto{\pgfqpoint{2.996512in}{1.284717in}}%
\pgfpathlineto{\pgfqpoint{2.999031in}{1.292759in}}%
\pgfpathlineto{\pgfqpoint{3.001549in}{1.324446in}}%
\pgfpathlineto{\pgfqpoint{3.004068in}{1.306323in}}%
\pgfpathlineto{\pgfqpoint{3.006587in}{1.340483in}}%
\pgfpathlineto{\pgfqpoint{3.009106in}{1.333516in}}%
\pgfpathlineto{\pgfqpoint{3.011624in}{1.342054in}}%
\pgfpathlineto{\pgfqpoint{3.014143in}{1.355091in}}%
\pgfpathlineto{\pgfqpoint{3.016662in}{1.361877in}}%
\pgfpathlineto{\pgfqpoint{3.019181in}{1.344678in}}%
\pgfpathlineto{\pgfqpoint{3.021699in}{1.329917in}}%
\pgfpathlineto{\pgfqpoint{3.024218in}{1.316732in}}%
\pgfpathlineto{\pgfqpoint{3.026737in}{1.292684in}}%
\pgfpathlineto{\pgfqpoint{3.029256in}{1.270611in}}%
\pgfpathlineto{\pgfqpoint{3.031774in}{1.284995in}}%
\pgfpathlineto{\pgfqpoint{3.034293in}{1.314655in}}%
\pgfpathlineto{\pgfqpoint{3.036812in}{1.317594in}}%
\pgfpathlineto{\pgfqpoint{3.039331in}{1.308188in}}%
\pgfpathlineto{\pgfqpoint{3.041849in}{1.299433in}}%
\pgfpathlineto{\pgfqpoint{3.044368in}{1.306376in}}%
\pgfpathlineto{\pgfqpoint{3.046887in}{1.300549in}}%
\pgfpathlineto{\pgfqpoint{3.049406in}{1.306030in}}%
\pgfpathlineto{\pgfqpoint{3.051924in}{1.315931in}}%
\pgfpathlineto{\pgfqpoint{3.054443in}{1.335640in}}%
\pgfpathlineto{\pgfqpoint{3.056962in}{1.326675in}}%
\pgfpathlineto{\pgfqpoint{3.059481in}{1.321989in}}%
\pgfpathlineto{\pgfqpoint{3.061999in}{1.361599in}}%
\pgfpathlineto{\pgfqpoint{3.064518in}{1.347820in}}%
\pgfpathlineto{\pgfqpoint{3.067037in}{1.340542in}}%
\pgfpathlineto{\pgfqpoint{3.069556in}{1.324298in}}%
\pgfpathlineto{\pgfqpoint{3.072074in}{1.340673in}}%
\pgfpathlineto{\pgfqpoint{3.074593in}{1.351136in}}%
\pgfpathlineto{\pgfqpoint{3.077112in}{1.349202in}}%
\pgfpathlineto{\pgfqpoint{3.079631in}{1.331421in}}%
\pgfpathlineto{\pgfqpoint{3.082149in}{1.359803in}}%
\pgfpathlineto{\pgfqpoint{3.084668in}{1.385504in}}%
\pgfpathlineto{\pgfqpoint{3.087187in}{1.372506in}}%
\pgfpathlineto{\pgfqpoint{3.089706in}{1.393735in}}%
\pgfpathlineto{\pgfqpoint{3.092224in}{1.410252in}}%
\pgfpathlineto{\pgfqpoint{3.094743in}{1.419919in}}%
\pgfpathlineto{\pgfqpoint{3.097262in}{1.434797in}}%
\pgfpathlineto{\pgfqpoint{3.099781in}{1.437900in}}%
\pgfpathlineto{\pgfqpoint{3.102299in}{1.432040in}}%
\pgfpathlineto{\pgfqpoint{3.104818in}{1.401510in}}%
\pgfpathlineto{\pgfqpoint{3.107337in}{1.392608in}}%
\pgfpathlineto{\pgfqpoint{3.109856in}{1.399659in}}%
\pgfpathlineto{\pgfqpoint{3.112374in}{1.402777in}}%
\pgfpathlineto{\pgfqpoint{3.114893in}{1.408141in}}%
\pgfpathlineto{\pgfqpoint{3.117412in}{1.400713in}}%
\pgfpathlineto{\pgfqpoint{3.119931in}{1.399462in}}%
\pgfpathlineto{\pgfqpoint{3.122449in}{1.414457in}}%
\pgfpathlineto{\pgfqpoint{3.124968in}{1.405890in}}%
\pgfpathlineto{\pgfqpoint{3.127487in}{1.413141in}}%
\pgfpathlineto{\pgfqpoint{3.130006in}{1.432987in}}%
\pgfpathlineto{\pgfqpoint{3.132524in}{1.423394in}}%
\pgfpathlineto{\pgfqpoint{3.135043in}{1.435084in}}%
\pgfpathlineto{\pgfqpoint{3.137562in}{1.405077in}}%
\pgfpathlineto{\pgfqpoint{3.140081in}{1.423576in}}%
\pgfpathlineto{\pgfqpoint{3.142599in}{1.441184in}}%
\pgfpathlineto{\pgfqpoint{3.145118in}{1.425988in}}%
\pgfpathlineto{\pgfqpoint{3.147637in}{1.413821in}}%
\pgfpathlineto{\pgfqpoint{3.150156in}{1.419974in}}%
\pgfpathlineto{\pgfqpoint{3.152674in}{1.415455in}}%
\pgfpathlineto{\pgfqpoint{3.155193in}{1.405067in}}%
\pgfpathlineto{\pgfqpoint{3.157712in}{1.410754in}}%
\pgfpathlineto{\pgfqpoint{3.160231in}{1.415888in}}%
\pgfpathlineto{\pgfqpoint{3.162749in}{1.428706in}}%
\pgfpathlineto{\pgfqpoint{3.165268in}{1.442669in}}%
\pgfpathlineto{\pgfqpoint{3.167787in}{1.441750in}}%
\pgfpathlineto{\pgfqpoint{3.170306in}{1.453462in}}%
\pgfpathlineto{\pgfqpoint{3.172824in}{1.465807in}}%
\pgfpathlineto{\pgfqpoint{3.175343in}{1.466564in}}%
\pgfpathlineto{\pgfqpoint{3.177862in}{1.448661in}}%
\pgfpathlineto{\pgfqpoint{3.180381in}{1.452261in}}%
\pgfpathlineto{\pgfqpoint{3.182899in}{1.474655in}}%
\pgfpathlineto{\pgfqpoint{3.185418in}{1.474557in}}%
\pgfpathlineto{\pgfqpoint{3.187937in}{1.497253in}}%
\pgfpathlineto{\pgfqpoint{3.190456in}{1.501503in}}%
\pgfpathlineto{\pgfqpoint{3.192974in}{1.511831in}}%
\pgfpathlineto{\pgfqpoint{3.195493in}{1.514387in}}%
\pgfpathlineto{\pgfqpoint{3.198012in}{1.497648in}}%
\pgfpathlineto{\pgfqpoint{3.200531in}{1.493047in}}%
\pgfpathlineto{\pgfqpoint{3.203049in}{1.467129in}}%
\pgfpathlineto{\pgfqpoint{3.205568in}{1.450607in}}%
\pgfpathlineto{\pgfqpoint{3.208087in}{1.464497in}}%
\pgfpathlineto{\pgfqpoint{3.210606in}{1.467187in}}%
\pgfpathlineto{\pgfqpoint{3.213124in}{1.464625in}}%
\pgfpathlineto{\pgfqpoint{3.215643in}{1.502284in}}%
\pgfpathlineto{\pgfqpoint{3.218162in}{1.499218in}}%
\pgfpathlineto{\pgfqpoint{3.220681in}{1.493668in}}%
\pgfpathlineto{\pgfqpoint{3.223199in}{1.502489in}}%
\pgfpathlineto{\pgfqpoint{3.225718in}{1.536890in}}%
\pgfpathlineto{\pgfqpoint{3.228237in}{1.566495in}}%
\pgfpathlineto{\pgfqpoint{3.230756in}{1.581642in}}%
\pgfpathlineto{\pgfqpoint{3.233274in}{1.586098in}}%
\pgfpathlineto{\pgfqpoint{3.235793in}{1.594971in}}%
\pgfpathlineto{\pgfqpoint{3.238312in}{1.586540in}}%
\pgfpathlineto{\pgfqpoint{3.240831in}{1.580669in}}%
\pgfpathlineto{\pgfqpoint{3.243349in}{1.582165in}}%
\pgfpathlineto{\pgfqpoint{3.245868in}{1.566842in}}%
\pgfpathlineto{\pgfqpoint{3.248387in}{1.584588in}}%
\pgfpathlineto{\pgfqpoint{3.250906in}{1.599798in}}%
\pgfpathlineto{\pgfqpoint{3.253424in}{1.608427in}}%
\pgfpathlineto{\pgfqpoint{3.255943in}{1.586997in}}%
\pgfpathlineto{\pgfqpoint{3.258462in}{1.569739in}}%
\pgfpathlineto{\pgfqpoint{3.260981in}{1.576334in}}%
\pgfpathlineto{\pgfqpoint{3.263499in}{1.580113in}}%
\pgfpathlineto{\pgfqpoint{3.266018in}{1.590152in}}%
\pgfpathlineto{\pgfqpoint{3.268537in}{1.607408in}}%
\pgfpathlineto{\pgfqpoint{3.271056in}{1.608783in}}%
\pgfpathlineto{\pgfqpoint{3.273574in}{1.616163in}}%
\pgfpathlineto{\pgfqpoint{3.276093in}{1.636608in}}%
\pgfpathlineto{\pgfqpoint{3.278612in}{1.644268in}}%
\pgfpathlineto{\pgfqpoint{3.281131in}{1.682118in}}%
\pgfpathlineto{\pgfqpoint{3.283649in}{1.668485in}}%
\pgfpathlineto{\pgfqpoint{3.286168in}{1.664479in}}%
\pgfpathlineto{\pgfqpoint{3.288687in}{1.682494in}}%
\pgfpathlineto{\pgfqpoint{3.291206in}{1.691196in}}%
\pgfpathlineto{\pgfqpoint{3.293724in}{1.671951in}}%
\pgfpathlineto{\pgfqpoint{3.296243in}{1.669626in}}%
\pgfpathlineto{\pgfqpoint{3.298762in}{1.681922in}}%
\pgfpathlineto{\pgfqpoint{3.301281in}{1.685266in}}%
\pgfpathlineto{\pgfqpoint{3.303799in}{1.702225in}}%
\pgfpathlineto{\pgfqpoint{3.306318in}{1.698670in}}%
\pgfpathlineto{\pgfqpoint{3.308837in}{1.693232in}}%
\pgfpathlineto{\pgfqpoint{3.311356in}{1.696030in}}%
\pgfpathlineto{\pgfqpoint{3.313874in}{1.691493in}}%
\pgfpathlineto{\pgfqpoint{3.316393in}{1.687661in}}%
\pgfpathlineto{\pgfqpoint{3.318912in}{1.677510in}}%
\pgfpathlineto{\pgfqpoint{3.321431in}{1.708092in}}%
\pgfpathlineto{\pgfqpoint{3.323949in}{1.734090in}}%
\pgfpathlineto{\pgfqpoint{3.326468in}{1.740774in}}%
\pgfpathlineto{\pgfqpoint{3.328987in}{1.741604in}}%
\pgfpathlineto{\pgfqpoint{3.331506in}{1.750670in}}%
\pgfpathlineto{\pgfqpoint{3.334024in}{1.755503in}}%
\pgfpathlineto{\pgfqpoint{3.336543in}{1.745522in}}%
\pgfpathlineto{\pgfqpoint{3.339062in}{1.771789in}}%
\pgfpathlineto{\pgfqpoint{3.341581in}{1.756224in}}%
\pgfpathlineto{\pgfqpoint{3.344099in}{1.716038in}}%
\pgfpathlineto{\pgfqpoint{3.346618in}{1.719774in}}%
\pgfpathlineto{\pgfqpoint{3.349137in}{1.733626in}}%
\pgfpathlineto{\pgfqpoint{3.351656in}{1.738152in}}%
\pgfpathlineto{\pgfqpoint{3.354174in}{1.757798in}}%
\pgfpathlineto{\pgfqpoint{3.356693in}{1.743183in}}%
\pgfpathlineto{\pgfqpoint{3.359212in}{1.721529in}}%
\pgfpathlineto{\pgfqpoint{3.361731in}{1.718212in}}%
\pgfpathlineto{\pgfqpoint{3.364249in}{1.700667in}}%
\pgfpathlineto{\pgfqpoint{3.366768in}{1.727357in}}%
\pgfpathlineto{\pgfqpoint{3.369287in}{1.720397in}}%
\pgfpathlineto{\pgfqpoint{3.371806in}{1.718938in}}%
\pgfpathlineto{\pgfqpoint{3.374324in}{1.748526in}}%
\pgfpathlineto{\pgfqpoint{3.376843in}{1.761858in}}%
\pgfpathlineto{\pgfqpoint{3.379362in}{1.772947in}}%
\pgfpathlineto{\pgfqpoint{3.381881in}{1.764549in}}%
\pgfpathlineto{\pgfqpoint{3.384399in}{1.772668in}}%
\pgfpathlineto{\pgfqpoint{3.386918in}{1.775805in}}%
\pgfpathlineto{\pgfqpoint{3.389437in}{1.774149in}}%
\pgfpathlineto{\pgfqpoint{3.391956in}{1.749611in}}%
\pgfpathlineto{\pgfqpoint{3.394474in}{1.751477in}}%
\pgfpathlineto{\pgfqpoint{3.396993in}{1.749743in}}%
\pgfpathlineto{\pgfqpoint{3.399512in}{1.752524in}}%
\pgfpathlineto{\pgfqpoint{3.402031in}{1.757711in}}%
\pgfpathlineto{\pgfqpoint{3.404549in}{1.751495in}}%
\pgfpathlineto{\pgfqpoint{3.407068in}{1.733018in}}%
\pgfpathlineto{\pgfqpoint{3.412106in}{1.765054in}}%
\pgfpathlineto{\pgfqpoint{3.414624in}{1.759145in}}%
\pgfpathlineto{\pgfqpoint{3.417143in}{1.782016in}}%
\pgfpathlineto{\pgfqpoint{3.419662in}{1.785233in}}%
\pgfpathlineto{\pgfqpoint{3.422181in}{1.776010in}}%
\pgfpathlineto{\pgfqpoint{3.424699in}{1.786324in}}%
\pgfpathlineto{\pgfqpoint{3.427218in}{1.761289in}}%
\pgfpathlineto{\pgfqpoint{3.429737in}{1.795031in}}%
\pgfpathlineto{\pgfqpoint{3.432256in}{1.794683in}}%
\pgfpathlineto{\pgfqpoint{3.434774in}{1.783008in}}%
\pgfpathlineto{\pgfqpoint{3.437293in}{1.763708in}}%
\pgfpathlineto{\pgfqpoint{3.439812in}{1.758611in}}%
\pgfpathlineto{\pgfqpoint{3.442331in}{1.785339in}}%
\pgfpathlineto{\pgfqpoint{3.444849in}{1.794977in}}%
\pgfpathlineto{\pgfqpoint{3.447368in}{1.779895in}}%
\pgfpathlineto{\pgfqpoint{3.449887in}{1.789186in}}%
\pgfpathlineto{\pgfqpoint{3.452406in}{1.792928in}}%
\pgfpathlineto{\pgfqpoint{3.454924in}{1.787745in}}%
\pgfpathlineto{\pgfqpoint{3.457443in}{1.770258in}}%
\pgfpathlineto{\pgfqpoint{3.459962in}{1.746191in}}%
\pgfpathlineto{\pgfqpoint{3.462481in}{1.745877in}}%
\pgfpathlineto{\pgfqpoint{3.464999in}{1.735793in}}%
\pgfpathlineto{\pgfqpoint{3.467518in}{1.768634in}}%
\pgfpathlineto{\pgfqpoint{3.470037in}{1.791662in}}%
\pgfpathlineto{\pgfqpoint{3.472556in}{1.790141in}}%
\pgfpathlineto{\pgfqpoint{3.475074in}{1.780886in}}%
\pgfpathlineto{\pgfqpoint{3.477593in}{1.775617in}}%
\pgfpathlineto{\pgfqpoint{3.480112in}{1.809385in}}%
\pgfpathlineto{\pgfqpoint{3.482631in}{1.819761in}}%
\pgfpathlineto{\pgfqpoint{3.485149in}{1.831644in}}%
\pgfpathlineto{\pgfqpoint{3.487668in}{1.853337in}}%
\pgfpathlineto{\pgfqpoint{3.490187in}{1.839607in}}%
\pgfpathlineto{\pgfqpoint{3.492706in}{1.849774in}}%
\pgfpathlineto{\pgfqpoint{3.495224in}{1.864847in}}%
\pgfpathlineto{\pgfqpoint{3.497743in}{1.825094in}}%
\pgfpathlineto{\pgfqpoint{3.500262in}{1.807657in}}%
\pgfpathlineto{\pgfqpoint{3.502781in}{1.813009in}}%
\pgfpathlineto{\pgfqpoint{3.505299in}{1.806121in}}%
\pgfpathlineto{\pgfqpoint{3.507818in}{1.776757in}}%
\pgfpathlineto{\pgfqpoint{3.510337in}{1.765842in}}%
\pgfpathlineto{\pgfqpoint{3.512856in}{1.753726in}}%
\pgfpathlineto{\pgfqpoint{3.515374in}{1.748931in}}%
\pgfpathlineto{\pgfqpoint{3.517893in}{1.758608in}}%
\pgfpathlineto{\pgfqpoint{3.520412in}{1.757780in}}%
\pgfpathlineto{\pgfqpoint{3.522931in}{1.747542in}}%
\pgfpathlineto{\pgfqpoint{3.525449in}{1.747446in}}%
\pgfpathlineto{\pgfqpoint{3.527968in}{1.733468in}}%
\pgfpathlineto{\pgfqpoint{3.530487in}{1.720883in}}%
\pgfpathlineto{\pgfqpoint{3.533006in}{1.723707in}}%
\pgfpathlineto{\pgfqpoint{3.535524in}{1.725940in}}%
\pgfpathlineto{\pgfqpoint{3.538043in}{1.747430in}}%
\pgfpathlineto{\pgfqpoint{3.540562in}{1.748435in}}%
\pgfpathlineto{\pgfqpoint{3.543081in}{1.754541in}}%
\pgfpathlineto{\pgfqpoint{3.545599in}{1.754220in}}%
\pgfpathlineto{\pgfqpoint{3.548118in}{1.757005in}}%
\pgfpathlineto{\pgfqpoint{3.550637in}{1.788151in}}%
\pgfpathlineto{\pgfqpoint{3.553156in}{1.813938in}}%
\pgfpathlineto{\pgfqpoint{3.555674in}{1.812908in}}%
\pgfpathlineto{\pgfqpoint{3.558193in}{1.808729in}}%
\pgfpathlineto{\pgfqpoint{3.560712in}{1.783301in}}%
\pgfpathlineto{\pgfqpoint{3.563231in}{1.784740in}}%
\pgfpathlineto{\pgfqpoint{3.565749in}{1.813874in}}%
\pgfpathlineto{\pgfqpoint{3.568268in}{1.795691in}}%
\pgfpathlineto{\pgfqpoint{3.570787in}{1.765809in}}%
\pgfpathlineto{\pgfqpoint{3.573306in}{1.760679in}}%
\pgfpathlineto{\pgfqpoint{3.575824in}{1.761712in}}%
\pgfpathlineto{\pgfqpoint{3.578343in}{1.754093in}}%
\pgfpathlineto{\pgfqpoint{3.580862in}{1.763595in}}%
\pgfpathlineto{\pgfqpoint{3.583381in}{1.757379in}}%
\pgfpathlineto{\pgfqpoint{3.585899in}{1.763915in}}%
\pgfpathlineto{\pgfqpoint{3.588418in}{1.767996in}}%
\pgfpathlineto{\pgfqpoint{3.590937in}{1.766560in}}%
\pgfpathlineto{\pgfqpoint{3.593456in}{1.753121in}}%
\pgfpathlineto{\pgfqpoint{3.595974in}{1.764347in}}%
\pgfpathlineto{\pgfqpoint{3.598493in}{1.772364in}}%
\pgfpathlineto{\pgfqpoint{3.601012in}{1.771861in}}%
\pgfpathlineto{\pgfqpoint{3.603531in}{1.779062in}}%
\pgfpathlineto{\pgfqpoint{3.606049in}{1.774239in}}%
\pgfpathlineto{\pgfqpoint{3.608568in}{1.779996in}}%
\pgfpathlineto{\pgfqpoint{3.611087in}{1.798856in}}%
\pgfpathlineto{\pgfqpoint{3.613606in}{1.804757in}}%
\pgfpathlineto{\pgfqpoint{3.616124in}{1.760629in}}%
\pgfpathlineto{\pgfqpoint{3.618643in}{1.748553in}}%
\pgfpathlineto{\pgfqpoint{3.621162in}{1.757912in}}%
\pgfpathlineto{\pgfqpoint{3.623681in}{1.761277in}}%
\pgfpathlineto{\pgfqpoint{3.626199in}{1.771440in}}%
\pgfpathlineto{\pgfqpoint{3.628718in}{1.755220in}}%
\pgfpathlineto{\pgfqpoint{3.631237in}{1.754581in}}%
\pgfpathlineto{\pgfqpoint{3.633756in}{1.756099in}}%
\pgfpathlineto{\pgfqpoint{3.636274in}{1.782913in}}%
\pgfpathlineto{\pgfqpoint{3.638793in}{1.803131in}}%
\pgfpathlineto{\pgfqpoint{3.641312in}{1.814451in}}%
\pgfpathlineto{\pgfqpoint{3.643831in}{1.844484in}}%
\pgfpathlineto{\pgfqpoint{3.646349in}{1.839786in}}%
\pgfpathlineto{\pgfqpoint{3.648868in}{1.823747in}}%
\pgfpathlineto{\pgfqpoint{3.651387in}{1.821564in}}%
\pgfpathlineto{\pgfqpoint{3.653906in}{1.833544in}}%
\pgfpathlineto{\pgfqpoint{3.656424in}{1.838234in}}%
\pgfpathlineto{\pgfqpoint{3.658943in}{1.843274in}}%
\pgfpathlineto{\pgfqpoint{3.661462in}{1.799820in}}%
\pgfpathlineto{\pgfqpoint{3.663981in}{1.821099in}}%
\pgfpathlineto{\pgfqpoint{3.666499in}{1.808908in}}%
\pgfpathlineto{\pgfqpoint{3.669018in}{1.831922in}}%
\pgfpathlineto{\pgfqpoint{3.671537in}{1.832677in}}%
\pgfpathlineto{\pgfqpoint{3.674056in}{1.837882in}}%
\pgfpathlineto{\pgfqpoint{3.676574in}{1.839975in}}%
\pgfpathlineto{\pgfqpoint{3.679093in}{1.851729in}}%
\pgfpathlineto{\pgfqpoint{3.681612in}{1.841712in}}%
\pgfpathlineto{\pgfqpoint{3.684131in}{1.854061in}}%
\pgfpathlineto{\pgfqpoint{3.686649in}{1.859040in}}%
\pgfpathlineto{\pgfqpoint{3.689168in}{1.861011in}}%
\pgfpathlineto{\pgfqpoint{3.691687in}{1.858947in}}%
\pgfpathlineto{\pgfqpoint{3.694206in}{1.865092in}}%
\pgfpathlineto{\pgfqpoint{3.696724in}{1.849428in}}%
\pgfpathlineto{\pgfqpoint{3.699243in}{1.872028in}}%
\pgfpathlineto{\pgfqpoint{3.701762in}{1.854035in}}%
\pgfpathlineto{\pgfqpoint{3.704281in}{1.852158in}}%
\pgfpathlineto{\pgfqpoint{3.706799in}{1.848069in}}%
\pgfpathlineto{\pgfqpoint{3.709318in}{1.850693in}}%
\pgfpathlineto{\pgfqpoint{3.711837in}{1.845381in}}%
\pgfpathlineto{\pgfqpoint{3.714356in}{1.840676in}}%
\pgfpathlineto{\pgfqpoint{3.716874in}{1.819231in}}%
\pgfpathlineto{\pgfqpoint{3.719393in}{1.825073in}}%
\pgfpathlineto{\pgfqpoint{3.721912in}{1.797046in}}%
\pgfpathlineto{\pgfqpoint{3.724431in}{1.788386in}}%
\pgfpathlineto{\pgfqpoint{3.726949in}{1.767239in}}%
\pgfpathlineto{\pgfqpoint{3.729468in}{1.762074in}}%
\pgfpathlineto{\pgfqpoint{3.731987in}{1.735388in}}%
\pgfpathlineto{\pgfqpoint{3.734506in}{1.726565in}}%
\pgfpathlineto{\pgfqpoint{3.737024in}{1.727563in}}%
\pgfpathlineto{\pgfqpoint{3.739543in}{1.712468in}}%
\pgfpathlineto{\pgfqpoint{3.742062in}{1.711618in}}%
\pgfpathlineto{\pgfqpoint{3.744581in}{1.735589in}}%
\pgfpathlineto{\pgfqpoint{3.747099in}{1.732972in}}%
\pgfpathlineto{\pgfqpoint{3.749618in}{1.731370in}}%
\pgfpathlineto{\pgfqpoint{3.752137in}{1.723417in}}%
\pgfpathlineto{\pgfqpoint{3.754656in}{1.758400in}}%
\pgfpathlineto{\pgfqpoint{3.757174in}{1.772990in}}%
\pgfpathlineto{\pgfqpoint{3.759693in}{1.750948in}}%
\pgfpathlineto{\pgfqpoint{3.762212in}{1.739671in}}%
\pgfpathlineto{\pgfqpoint{3.764731in}{1.731101in}}%
\pgfpathlineto{\pgfqpoint{3.767249in}{1.738462in}}%
\pgfpathlineto{\pgfqpoint{3.769768in}{1.762882in}}%
\pgfpathlineto{\pgfqpoint{3.772287in}{1.753465in}}%
\pgfpathlineto{\pgfqpoint{3.774806in}{1.744953in}}%
\pgfpathlineto{\pgfqpoint{3.777324in}{1.755497in}}%
\pgfpathlineto{\pgfqpoint{3.779843in}{1.727879in}}%
\pgfpathlineto{\pgfqpoint{3.782362in}{1.739220in}}%
\pgfpathlineto{\pgfqpoint{3.784881in}{1.715003in}}%
\pgfpathlineto{\pgfqpoint{3.787399in}{1.687936in}}%
\pgfpathlineto{\pgfqpoint{3.789918in}{1.669456in}}%
\pgfpathlineto{\pgfqpoint{3.792437in}{1.676117in}}%
\pgfpathlineto{\pgfqpoint{3.794956in}{1.669509in}}%
\pgfpathlineto{\pgfqpoint{3.797474in}{1.677349in}}%
\pgfpathlineto{\pgfqpoint{3.799993in}{1.681558in}}%
\pgfpathlineto{\pgfqpoint{3.802512in}{1.696022in}}%
\pgfpathlineto{\pgfqpoint{3.805031in}{1.724391in}}%
\pgfpathlineto{\pgfqpoint{3.807549in}{1.720412in}}%
\pgfpathlineto{\pgfqpoint{3.810068in}{1.730555in}}%
\pgfpathlineto{\pgfqpoint{3.812587in}{1.692622in}}%
\pgfpathlineto{\pgfqpoint{3.815106in}{1.672344in}}%
\pgfpathlineto{\pgfqpoint{3.817624in}{1.689936in}}%
\pgfpathlineto{\pgfqpoint{3.820143in}{1.691967in}}%
\pgfpathlineto{\pgfqpoint{3.822662in}{1.685563in}}%
\pgfpathlineto{\pgfqpoint{3.825181in}{1.695430in}}%
\pgfpathlineto{\pgfqpoint{3.827699in}{1.696648in}}%
\pgfpathlineto{\pgfqpoint{3.830218in}{1.695533in}}%
\pgfpathlineto{\pgfqpoint{3.832737in}{1.702408in}}%
\pgfpathlineto{\pgfqpoint{3.835256in}{1.688937in}}%
\pgfpathlineto{\pgfqpoint{3.837774in}{1.684568in}}%
\pgfpathlineto{\pgfqpoint{3.840293in}{1.676048in}}%
\pgfpathlineto{\pgfqpoint{3.842812in}{1.704255in}}%
\pgfpathlineto{\pgfqpoint{3.845331in}{1.706537in}}%
\pgfpathlineto{\pgfqpoint{3.847849in}{1.698960in}}%
\pgfpathlineto{\pgfqpoint{3.850368in}{1.715599in}}%
\pgfpathlineto{\pgfqpoint{3.852887in}{1.739713in}}%
\pgfpathlineto{\pgfqpoint{3.855406in}{1.722035in}}%
\pgfpathlineto{\pgfqpoint{3.857924in}{1.743740in}}%
\pgfpathlineto{\pgfqpoint{3.860443in}{1.753749in}}%
\pgfpathlineto{\pgfqpoint{3.862962in}{1.782961in}}%
\pgfpathlineto{\pgfqpoint{3.865481in}{1.794412in}}%
\pgfpathlineto{\pgfqpoint{3.867999in}{1.805262in}}%
\pgfpathlineto{\pgfqpoint{3.870518in}{1.816528in}}%
\pgfpathlineto{\pgfqpoint{3.873037in}{1.799805in}}%
\pgfpathlineto{\pgfqpoint{3.875556in}{1.789616in}}%
\pgfpathlineto{\pgfqpoint{3.878074in}{1.786196in}}%
\pgfpathlineto{\pgfqpoint{3.880593in}{1.774116in}}%
\pgfpathlineto{\pgfqpoint{3.883112in}{1.807476in}}%
\pgfpathlineto{\pgfqpoint{3.885631in}{1.809326in}}%
\pgfpathlineto{\pgfqpoint{3.888149in}{1.801756in}}%
\pgfpathlineto{\pgfqpoint{3.890668in}{1.805719in}}%
\pgfpathlineto{\pgfqpoint{3.893187in}{1.814289in}}%
\pgfpathlineto{\pgfqpoint{3.895706in}{1.792621in}}%
\pgfpathlineto{\pgfqpoint{3.898224in}{1.780383in}}%
\pgfpathlineto{\pgfqpoint{3.900743in}{1.783537in}}%
\pgfpathlineto{\pgfqpoint{3.903262in}{1.808220in}}%
\pgfpathlineto{\pgfqpoint{3.905781in}{1.793797in}}%
\pgfpathlineto{\pgfqpoint{3.908299in}{1.819377in}}%
\pgfpathlineto{\pgfqpoint{3.910818in}{1.798231in}}%
\pgfpathlineto{\pgfqpoint{3.913337in}{1.785661in}}%
\pgfpathlineto{\pgfqpoint{3.915856in}{1.762658in}}%
\pgfpathlineto{\pgfqpoint{3.918374in}{1.758505in}}%
\pgfpathlineto{\pgfqpoint{3.920893in}{1.727124in}}%
\pgfpathlineto{\pgfqpoint{3.923412in}{1.706882in}}%
\pgfpathlineto{\pgfqpoint{3.925931in}{1.733953in}}%
\pgfpathlineto{\pgfqpoint{3.928449in}{1.737545in}}%
\pgfpathlineto{\pgfqpoint{3.933487in}{1.707717in}}%
\pgfpathlineto{\pgfqpoint{3.936006in}{1.710474in}}%
\pgfpathlineto{\pgfqpoint{3.938524in}{1.713541in}}%
\pgfpathlineto{\pgfqpoint{3.941043in}{1.704705in}}%
\pgfpathlineto{\pgfqpoint{3.943562in}{1.703945in}}%
\pgfpathlineto{\pgfqpoint{3.946081in}{1.704870in}}%
\pgfpathlineto{\pgfqpoint{3.948599in}{1.679002in}}%
\pgfpathlineto{\pgfqpoint{3.951118in}{1.686940in}}%
\pgfpathlineto{\pgfqpoint{3.953637in}{1.669276in}}%
\pgfpathlineto{\pgfqpoint{3.956156in}{1.653460in}}%
\pgfpathlineto{\pgfqpoint{3.958674in}{1.644102in}}%
\pgfpathlineto{\pgfqpoint{3.961193in}{1.651769in}}%
\pgfpathlineto{\pgfqpoint{3.963712in}{1.637100in}}%
\pgfpathlineto{\pgfqpoint{3.966231in}{1.653558in}}%
\pgfpathlineto{\pgfqpoint{3.968749in}{1.655493in}}%
\pgfpathlineto{\pgfqpoint{3.971268in}{1.634311in}}%
\pgfpathlineto{\pgfqpoint{3.973787in}{1.627473in}}%
\pgfpathlineto{\pgfqpoint{3.976306in}{1.637591in}}%
\pgfpathlineto{\pgfqpoint{3.978824in}{1.640002in}}%
\pgfpathlineto{\pgfqpoint{3.981343in}{1.628553in}}%
\pgfpathlineto{\pgfqpoint{3.983862in}{1.599131in}}%
\pgfpathlineto{\pgfqpoint{3.986381in}{1.601143in}}%
\pgfpathlineto{\pgfqpoint{3.988899in}{1.600567in}}%
\pgfpathlineto{\pgfqpoint{3.991418in}{1.602388in}}%
\pgfpathlineto{\pgfqpoint{3.993937in}{1.577817in}}%
\pgfpathlineto{\pgfqpoint{3.996456in}{1.596551in}}%
\pgfpathlineto{\pgfqpoint{3.998974in}{1.612137in}}%
\pgfpathlineto{\pgfqpoint{4.001493in}{1.597198in}}%
\pgfpathlineto{\pgfqpoint{4.004012in}{1.584949in}}%
\pgfpathlineto{\pgfqpoint{4.006531in}{1.623329in}}%
\pgfpathlineto{\pgfqpoint{4.009049in}{1.610019in}}%
\pgfpathlineto{\pgfqpoint{4.011568in}{1.635150in}}%
\pgfpathlineto{\pgfqpoint{4.014087in}{1.656720in}}%
\pgfpathlineto{\pgfqpoint{4.016606in}{1.642878in}}%
\pgfpathlineto{\pgfqpoint{4.019124in}{1.641673in}}%
\pgfpathlineto{\pgfqpoint{4.021643in}{1.632103in}}%
\pgfpathlineto{\pgfqpoint{4.024162in}{1.605742in}}%
\pgfpathlineto{\pgfqpoint{4.026681in}{1.588931in}}%
\pgfpathlineto{\pgfqpoint{4.029199in}{1.591410in}}%
\pgfpathlineto{\pgfqpoint{4.031718in}{1.563649in}}%
\pgfpathlineto{\pgfqpoint{4.034237in}{1.564644in}}%
\pgfpathlineto{\pgfqpoint{4.036756in}{1.553706in}}%
\pgfpathlineto{\pgfqpoint{4.039274in}{1.526551in}}%
\pgfpathlineto{\pgfqpoint{4.041793in}{1.529230in}}%
\pgfpathlineto{\pgfqpoint{4.044312in}{1.524117in}}%
\pgfpathlineto{\pgfqpoint{4.046831in}{1.509217in}}%
\pgfpathlineto{\pgfqpoint{4.049349in}{1.494936in}}%
\pgfpathlineto{\pgfqpoint{4.051868in}{1.485153in}}%
\pgfpathlineto{\pgfqpoint{4.054387in}{1.488056in}}%
\pgfpathlineto{\pgfqpoint{4.056906in}{1.480364in}}%
\pgfpathlineto{\pgfqpoint{4.059424in}{1.482128in}}%
\pgfpathlineto{\pgfqpoint{4.061943in}{1.481735in}}%
\pgfpathlineto{\pgfqpoint{4.064462in}{1.498790in}}%
\pgfpathlineto{\pgfqpoint{4.066981in}{1.499954in}}%
\pgfpathlineto{\pgfqpoint{4.069499in}{1.476292in}}%
\pgfpathlineto{\pgfqpoint{4.072018in}{1.496196in}}%
\pgfpathlineto{\pgfqpoint{4.074537in}{1.480370in}}%
\pgfpathlineto{\pgfqpoint{4.077056in}{1.461440in}}%
\pgfpathlineto{\pgfqpoint{4.079574in}{1.451549in}}%
\pgfpathlineto{\pgfqpoint{4.082093in}{1.455167in}}%
\pgfpathlineto{\pgfqpoint{4.084612in}{1.446343in}}%
\pgfpathlineto{\pgfqpoint{4.087131in}{1.451856in}}%
\pgfpathlineto{\pgfqpoint{4.089649in}{1.446529in}}%
\pgfpathlineto{\pgfqpoint{4.092168in}{1.466608in}}%
\pgfpathlineto{\pgfqpoint{4.094687in}{1.453855in}}%
\pgfpathlineto{\pgfqpoint{4.097206in}{1.464446in}}%
\pgfpathlineto{\pgfqpoint{4.099724in}{1.447625in}}%
\pgfpathlineto{\pgfqpoint{4.102243in}{1.468582in}}%
\pgfpathlineto{\pgfqpoint{4.104762in}{1.451700in}}%
\pgfpathlineto{\pgfqpoint{4.107281in}{1.456073in}}%
\pgfpathlineto{\pgfqpoint{4.109799in}{1.461046in}}%
\pgfpathlineto{\pgfqpoint{4.112318in}{1.461154in}}%
\pgfpathlineto{\pgfqpoint{4.114837in}{1.437412in}}%
\pgfpathlineto{\pgfqpoint{4.117356in}{1.450164in}}%
\pgfpathlineto{\pgfqpoint{4.119874in}{1.445200in}}%
\pgfpathlineto{\pgfqpoint{4.122393in}{1.437641in}}%
\pgfpathlineto{\pgfqpoint{4.124912in}{1.442492in}}%
\pgfpathlineto{\pgfqpoint{4.127431in}{1.442907in}}%
\pgfpathlineto{\pgfqpoint{4.129949in}{1.436450in}}%
\pgfpathlineto{\pgfqpoint{4.132468in}{1.415340in}}%
\pgfpathlineto{\pgfqpoint{4.134987in}{1.424204in}}%
\pgfpathlineto{\pgfqpoint{4.137506in}{1.407350in}}%
\pgfpathlineto{\pgfqpoint{4.140024in}{1.387473in}}%
\pgfpathlineto{\pgfqpoint{4.142543in}{1.390751in}}%
\pgfpathlineto{\pgfqpoint{4.145062in}{1.397992in}}%
\pgfpathlineto{\pgfqpoint{4.147581in}{1.396360in}}%
\pgfpathlineto{\pgfqpoint{4.150099in}{1.391731in}}%
\pgfpathlineto{\pgfqpoint{4.152618in}{1.386508in}}%
\pgfpathlineto{\pgfqpoint{4.155137in}{1.387908in}}%
\pgfpathlineto{\pgfqpoint{4.157656in}{1.393931in}}%
\pgfpathlineto{\pgfqpoint{4.160174in}{1.372714in}}%
\pgfpathlineto{\pgfqpoint{4.162693in}{1.397135in}}%
\pgfpathlineto{\pgfqpoint{4.165212in}{1.389810in}}%
\pgfpathlineto{\pgfqpoint{4.167731in}{1.389868in}}%
\pgfpathlineto{\pgfqpoint{4.170249in}{1.401143in}}%
\pgfpathlineto{\pgfqpoint{4.172768in}{1.387747in}}%
\pgfpathlineto{\pgfqpoint{4.175287in}{1.412320in}}%
\pgfpathlineto{\pgfqpoint{4.177806in}{1.410332in}}%
\pgfpathlineto{\pgfqpoint{4.180324in}{1.402967in}}%
\pgfpathlineto{\pgfqpoint{4.182843in}{1.427671in}}%
\pgfpathlineto{\pgfqpoint{4.185362in}{1.414899in}}%
\pgfpathlineto{\pgfqpoint{4.187881in}{1.403811in}}%
\pgfpathlineto{\pgfqpoint{4.190399in}{1.423749in}}%
\pgfpathlineto{\pgfqpoint{4.192918in}{1.419759in}}%
\pgfpathlineto{\pgfqpoint{4.195437in}{1.420778in}}%
\pgfpathlineto{\pgfqpoint{4.197956in}{1.455924in}}%
\pgfpathlineto{\pgfqpoint{4.200474in}{1.457070in}}%
\pgfpathlineto{\pgfqpoint{4.202993in}{1.466341in}}%
\pgfpathlineto{\pgfqpoint{4.205512in}{1.482182in}}%
\pgfpathlineto{\pgfqpoint{4.208031in}{1.481706in}}%
\pgfpathlineto{\pgfqpoint{4.210549in}{1.472292in}}%
\pgfpathlineto{\pgfqpoint{4.213068in}{1.473171in}}%
\pgfpathlineto{\pgfqpoint{4.215587in}{1.483370in}}%
\pgfpathlineto{\pgfqpoint{4.218106in}{1.482917in}}%
\pgfpathlineto{\pgfqpoint{4.220624in}{1.494948in}}%
\pgfpathlineto{\pgfqpoint{4.223143in}{1.516432in}}%
\pgfpathlineto{\pgfqpoint{4.225662in}{1.511413in}}%
\pgfpathlineto{\pgfqpoint{4.228181in}{1.533441in}}%
\pgfpathlineto{\pgfqpoint{4.230699in}{1.547367in}}%
\pgfpathlineto{\pgfqpoint{4.233218in}{1.563917in}}%
\pgfpathlineto{\pgfqpoint{4.235737in}{1.570202in}}%
\pgfpathlineto{\pgfqpoint{4.238256in}{1.565747in}}%
\pgfpathlineto{\pgfqpoint{4.240774in}{1.589074in}}%
\pgfpathlineto{\pgfqpoint{4.243293in}{1.592979in}}%
\pgfpathlineto{\pgfqpoint{4.245812in}{1.600452in}}%
\pgfpathlineto{\pgfqpoint{4.248331in}{1.649809in}}%
\pgfpathlineto{\pgfqpoint{4.253368in}{1.678480in}}%
\pgfpathlineto{\pgfqpoint{4.255887in}{1.661863in}}%
\pgfpathlineto{\pgfqpoint{4.258406in}{1.673779in}}%
\pgfpathlineto{\pgfqpoint{4.260924in}{1.664043in}}%
\pgfpathlineto{\pgfqpoint{4.263443in}{1.640227in}}%
\pgfpathlineto{\pgfqpoint{4.265962in}{1.659743in}}%
\pgfpathlineto{\pgfqpoint{4.268481in}{1.659578in}}%
\pgfpathlineto{\pgfqpoint{4.270999in}{1.651161in}}%
\pgfpathlineto{\pgfqpoint{4.273518in}{1.658637in}}%
\pgfpathlineto{\pgfqpoint{4.276037in}{1.643984in}}%
\pgfpathlineto{\pgfqpoint{4.278556in}{1.650431in}}%
\pgfpathlineto{\pgfqpoint{4.281074in}{1.636168in}}%
\pgfpathlineto{\pgfqpoint{4.283593in}{1.628351in}}%
\pgfpathlineto{\pgfqpoint{4.286112in}{1.643600in}}%
\pgfpathlineto{\pgfqpoint{4.288631in}{1.631977in}}%
\pgfpathlineto{\pgfqpoint{4.291149in}{1.624358in}}%
\pgfpathlineto{\pgfqpoint{4.293668in}{1.619839in}}%
\pgfpathlineto{\pgfqpoint{4.296187in}{1.631588in}}%
\pgfpathlineto{\pgfqpoint{4.298706in}{1.611914in}}%
\pgfpathlineto{\pgfqpoint{4.301224in}{1.583563in}}%
\pgfpathlineto{\pgfqpoint{4.303743in}{1.591415in}}%
\pgfpathlineto{\pgfqpoint{4.306262in}{1.578812in}}%
\pgfpathlineto{\pgfqpoint{4.308781in}{1.582854in}}%
\pgfpathlineto{\pgfqpoint{4.311299in}{1.570530in}}%
\pgfpathlineto{\pgfqpoint{4.313818in}{1.582629in}}%
\pgfpathlineto{\pgfqpoint{4.316337in}{1.587998in}}%
\pgfpathlineto{\pgfqpoint{4.318856in}{1.592914in}}%
\pgfpathlineto{\pgfqpoint{4.321374in}{1.626481in}}%
\pgfpathlineto{\pgfqpoint{4.323893in}{1.626712in}}%
\pgfpathlineto{\pgfqpoint{4.326412in}{1.634490in}}%
\pgfpathlineto{\pgfqpoint{4.328931in}{1.632100in}}%
\pgfpathlineto{\pgfqpoint{4.331449in}{1.635192in}}%
\pgfpathlineto{\pgfqpoint{4.333968in}{1.628619in}}%
\pgfpathlineto{\pgfqpoint{4.336487in}{1.635239in}}%
\pgfpathlineto{\pgfqpoint{4.339006in}{1.633869in}}%
\pgfpathlineto{\pgfqpoint{4.341524in}{1.641990in}}%
\pgfpathlineto{\pgfqpoint{4.344043in}{1.675077in}}%
\pgfpathlineto{\pgfqpoint{4.346562in}{1.677718in}}%
\pgfpathlineto{\pgfqpoint{4.349081in}{1.646076in}}%
\pgfpathlineto{\pgfqpoint{4.351599in}{1.641767in}}%
\pgfpathlineto{\pgfqpoint{4.354118in}{1.626373in}}%
\pgfpathlineto{\pgfqpoint{4.356637in}{1.603073in}}%
\pgfpathlineto{\pgfqpoint{4.359156in}{1.562168in}}%
\pgfpathlineto{\pgfqpoint{4.364193in}{1.562114in}}%
\pgfpathlineto{\pgfqpoint{4.366712in}{1.584372in}}%
\pgfpathlineto{\pgfqpoint{4.369231in}{1.576884in}}%
\pgfpathlineto{\pgfqpoint{4.371749in}{1.561034in}}%
\pgfpathlineto{\pgfqpoint{4.374268in}{1.570690in}}%
\pgfpathlineto{\pgfqpoint{4.376787in}{1.569658in}}%
\pgfpathlineto{\pgfqpoint{4.379306in}{1.589004in}}%
\pgfpathlineto{\pgfqpoint{4.381824in}{1.577741in}}%
\pgfpathlineto{\pgfqpoint{4.384343in}{1.570527in}}%
\pgfpathlineto{\pgfqpoint{4.386862in}{1.579095in}}%
\pgfpathlineto{\pgfqpoint{4.389381in}{1.564090in}}%
\pgfpathlineto{\pgfqpoint{4.391899in}{1.577340in}}%
\pgfpathlineto{\pgfqpoint{4.394418in}{1.573998in}}%
\pgfpathlineto{\pgfqpoint{4.396937in}{1.580395in}}%
\pgfpathlineto{\pgfqpoint{4.399456in}{1.564759in}}%
\pgfpathlineto{\pgfqpoint{4.401974in}{1.572789in}}%
\pgfpathlineto{\pgfqpoint{4.404493in}{1.582669in}}%
\pgfpathlineto{\pgfqpoint{4.407012in}{1.583646in}}%
\pgfpathlineto{\pgfqpoint{4.409531in}{1.589243in}}%
\pgfpathlineto{\pgfqpoint{4.412049in}{1.603951in}}%
\pgfpathlineto{\pgfqpoint{4.414568in}{1.638075in}}%
\pgfpathlineto{\pgfqpoint{4.417087in}{1.639788in}}%
\pgfpathlineto{\pgfqpoint{4.419606in}{1.643413in}}%
\pgfpathlineto{\pgfqpoint{4.422124in}{1.647457in}}%
\pgfpathlineto{\pgfqpoint{4.424643in}{1.657174in}}%
\pgfpathlineto{\pgfqpoint{4.427162in}{1.675494in}}%
\pgfpathlineto{\pgfqpoint{4.429681in}{1.655299in}}%
\pgfpathlineto{\pgfqpoint{4.432199in}{1.715818in}}%
\pgfpathlineto{\pgfqpoint{4.434718in}{1.688670in}}%
\pgfpathlineto{\pgfqpoint{4.437237in}{1.697143in}}%
\pgfpathlineto{\pgfqpoint{4.439756in}{1.716095in}}%
\pgfpathlineto{\pgfqpoint{4.442274in}{1.743485in}}%
\pgfpathlineto{\pgfqpoint{4.444793in}{1.732903in}}%
\pgfpathlineto{\pgfqpoint{4.447312in}{1.726002in}}%
\pgfpathlineto{\pgfqpoint{4.449831in}{1.748447in}}%
\pgfpathlineto{\pgfqpoint{4.452349in}{1.741276in}}%
\pgfpathlineto{\pgfqpoint{4.454868in}{1.731259in}}%
\pgfpathlineto{\pgfqpoint{4.457387in}{1.734653in}}%
\pgfpathlineto{\pgfqpoint{4.459906in}{1.724684in}}%
\pgfpathlineto{\pgfqpoint{4.462424in}{1.678041in}}%
\pgfpathlineto{\pgfqpoint{4.464943in}{1.675751in}}%
\pgfpathlineto{\pgfqpoint{4.467462in}{1.691034in}}%
\pgfpathlineto{\pgfqpoint{4.469981in}{1.697708in}}%
\pgfpathlineto{\pgfqpoint{4.472499in}{1.717762in}}%
\pgfpathlineto{\pgfqpoint{4.475018in}{1.705141in}}%
\pgfpathlineto{\pgfqpoint{4.477537in}{1.703545in}}%
\pgfpathlineto{\pgfqpoint{4.480056in}{1.707134in}}%
\pgfpathlineto{\pgfqpoint{4.482574in}{1.654428in}}%
\pgfpathlineto{\pgfqpoint{4.485093in}{1.666818in}}%
\pgfpathlineto{\pgfqpoint{4.487612in}{1.672084in}}%
\pgfpathlineto{\pgfqpoint{4.490131in}{1.678453in}}%
\pgfpathlineto{\pgfqpoint{4.492649in}{1.683831in}}%
\pgfpathlineto{\pgfqpoint{4.495168in}{1.687480in}}%
\pgfpathlineto{\pgfqpoint{4.497687in}{1.689745in}}%
\pgfpathlineto{\pgfqpoint{4.500206in}{1.691121in}}%
\pgfpathlineto{\pgfqpoint{4.502724in}{1.705474in}}%
\pgfpathlineto{\pgfqpoint{4.505243in}{1.675492in}}%
\pgfpathlineto{\pgfqpoint{4.507762in}{1.660224in}}%
\pgfpathlineto{\pgfqpoint{4.510281in}{1.638675in}}%
\pgfpathlineto{\pgfqpoint{4.512799in}{1.599077in}}%
\pgfpathlineto{\pgfqpoint{4.515318in}{1.598922in}}%
\pgfpathlineto{\pgfqpoint{4.517837in}{1.596187in}}%
\pgfpathlineto{\pgfqpoint{4.520356in}{1.599042in}}%
\pgfpathlineto{\pgfqpoint{4.522874in}{1.600521in}}%
\pgfpathlineto{\pgfqpoint{4.525393in}{1.616806in}}%
\pgfpathlineto{\pgfqpoint{4.527912in}{1.621014in}}%
\pgfpathlineto{\pgfqpoint{4.530431in}{1.638132in}}%
\pgfpathlineto{\pgfqpoint{4.532949in}{1.634512in}}%
\pgfpathlineto{\pgfqpoint{4.535468in}{1.626555in}}%
\pgfpathlineto{\pgfqpoint{4.537987in}{1.633381in}}%
\pgfpathlineto{\pgfqpoint{4.540506in}{1.621344in}}%
\pgfpathlineto{\pgfqpoint{4.543024in}{1.607389in}}%
\pgfpathlineto{\pgfqpoint{4.545543in}{1.615108in}}%
\pgfpathlineto{\pgfqpoint{4.548062in}{1.616290in}}%
\pgfpathlineto{\pgfqpoint{4.550581in}{1.616033in}}%
\pgfpathlineto{\pgfqpoint{4.553099in}{1.618615in}}%
\pgfpathlineto{\pgfqpoint{4.555618in}{1.618968in}}%
\pgfpathlineto{\pgfqpoint{4.558137in}{1.634671in}}%
\pgfpathlineto{\pgfqpoint{4.560656in}{1.644142in}}%
\pgfpathlineto{\pgfqpoint{4.563174in}{1.643298in}}%
\pgfpathlineto{\pgfqpoint{4.565693in}{1.620928in}}%
\pgfpathlineto{\pgfqpoint{4.568212in}{1.628966in}}%
\pgfpathlineto{\pgfqpoint{4.570731in}{1.614751in}}%
\pgfpathlineto{\pgfqpoint{4.573249in}{1.621123in}}%
\pgfpathlineto{\pgfqpoint{4.575768in}{1.631371in}}%
\pgfpathlineto{\pgfqpoint{4.578287in}{1.625089in}}%
\pgfpathlineto{\pgfqpoint{4.580806in}{1.641533in}}%
\pgfpathlineto{\pgfqpoint{4.583324in}{1.628451in}}%
\pgfpathlineto{\pgfqpoint{4.585843in}{1.630703in}}%
\pgfpathlineto{\pgfqpoint{4.588362in}{1.628519in}}%
\pgfpathlineto{\pgfqpoint{4.590881in}{1.641170in}}%
\pgfpathlineto{\pgfqpoint{4.593399in}{1.665989in}}%
\pgfpathlineto{\pgfqpoint{4.595918in}{1.670451in}}%
\pgfpathlineto{\pgfqpoint{4.598437in}{1.668696in}}%
\pgfpathlineto{\pgfqpoint{4.600956in}{1.668527in}}%
\pgfpathlineto{\pgfqpoint{4.603474in}{1.707509in}}%
\pgfpathlineto{\pgfqpoint{4.605993in}{1.703855in}}%
\pgfpathlineto{\pgfqpoint{4.608512in}{1.702905in}}%
\pgfpathlineto{\pgfqpoint{4.611031in}{1.744415in}}%
\pgfpathlineto{\pgfqpoint{4.613549in}{1.716217in}}%
\pgfpathlineto{\pgfqpoint{4.616068in}{1.735570in}}%
\pgfpathlineto{\pgfqpoint{4.618587in}{1.749153in}}%
\pgfpathlineto{\pgfqpoint{4.621106in}{1.759771in}}%
\pgfpathlineto{\pgfqpoint{4.623624in}{1.751115in}}%
\pgfpathlineto{\pgfqpoint{4.626143in}{1.761838in}}%
\pgfpathlineto{\pgfqpoint{4.628662in}{1.756261in}}%
\pgfpathlineto{\pgfqpoint{4.631181in}{1.758877in}}%
\pgfpathlineto{\pgfqpoint{4.633699in}{1.763127in}}%
\pgfpathlineto{\pgfqpoint{4.636218in}{1.769797in}}%
\pgfpathlineto{\pgfqpoint{4.638737in}{1.782623in}}%
\pgfpathlineto{\pgfqpoint{4.641256in}{1.772843in}}%
\pgfpathlineto{\pgfqpoint{4.643774in}{1.789920in}}%
\pgfpathlineto{\pgfqpoint{4.646293in}{1.809779in}}%
\pgfpathlineto{\pgfqpoint{4.648812in}{1.807125in}}%
\pgfpathlineto{\pgfqpoint{4.651331in}{1.816934in}}%
\pgfpathlineto{\pgfqpoint{4.653849in}{1.833990in}}%
\pgfpathlineto{\pgfqpoint{4.656368in}{1.813661in}}%
\pgfpathlineto{\pgfqpoint{4.658887in}{1.834819in}}%
\pgfpathlineto{\pgfqpoint{4.661406in}{1.839522in}}%
\pgfpathlineto{\pgfqpoint{4.663924in}{1.846344in}}%
\pgfpathlineto{\pgfqpoint{4.666443in}{1.844526in}}%
\pgfpathlineto{\pgfqpoint{4.668962in}{1.856874in}}%
\pgfpathlineto{\pgfqpoint{4.671481in}{1.858441in}}%
\pgfpathlineto{\pgfqpoint{4.673999in}{1.866355in}}%
\pgfpathlineto{\pgfqpoint{4.676518in}{1.853574in}}%
\pgfpathlineto{\pgfqpoint{4.679037in}{1.832572in}}%
\pgfpathlineto{\pgfqpoint{4.681556in}{1.816563in}}%
\pgfpathlineto{\pgfqpoint{4.684074in}{1.819683in}}%
\pgfpathlineto{\pgfqpoint{4.686593in}{1.826052in}}%
\pgfpathlineto{\pgfqpoint{4.689112in}{1.849412in}}%
\pgfpathlineto{\pgfqpoint{4.691631in}{1.840072in}}%
\pgfpathlineto{\pgfqpoint{4.694149in}{1.819988in}}%
\pgfpathlineto{\pgfqpoint{4.696668in}{1.820903in}}%
\pgfpathlineto{\pgfqpoint{4.699187in}{1.847262in}}%
\pgfpathlineto{\pgfqpoint{4.701706in}{1.870104in}}%
\pgfpathlineto{\pgfqpoint{4.704224in}{1.894386in}}%
\pgfpathlineto{\pgfqpoint{4.706743in}{1.919576in}}%
\pgfpathlineto{\pgfqpoint{4.709262in}{1.918119in}}%
\pgfpathlineto{\pgfqpoint{4.711781in}{1.928408in}}%
\pgfpathlineto{\pgfqpoint{4.714299in}{1.941272in}}%
\pgfpathlineto{\pgfqpoint{4.716818in}{1.920836in}}%
\pgfpathlineto{\pgfqpoint{4.719337in}{1.890209in}}%
\pgfpathlineto{\pgfqpoint{4.721856in}{1.889416in}}%
\pgfpathlineto{\pgfqpoint{4.724374in}{1.872459in}}%
\pgfpathlineto{\pgfqpoint{4.726893in}{1.853785in}}%
\pgfpathlineto{\pgfqpoint{4.729412in}{1.876045in}}%
\pgfpathlineto{\pgfqpoint{4.731931in}{1.871380in}}%
\pgfpathlineto{\pgfqpoint{4.734449in}{1.850669in}}%
\pgfpathlineto{\pgfqpoint{4.736968in}{1.849450in}}%
\pgfpathlineto{\pgfqpoint{4.739487in}{1.859933in}}%
\pgfpathlineto{\pgfqpoint{4.742006in}{1.852740in}}%
\pgfpathlineto{\pgfqpoint{4.744524in}{1.867360in}}%
\pgfpathlineto{\pgfqpoint{4.747043in}{1.855777in}}%
\pgfpathlineto{\pgfqpoint{4.749562in}{1.843294in}}%
\pgfpathlineto{\pgfqpoint{4.752081in}{1.854656in}}%
\pgfpathlineto{\pgfqpoint{4.754599in}{1.858439in}}%
\pgfpathlineto{\pgfqpoint{4.757118in}{1.867029in}}%
\pgfpathlineto{\pgfqpoint{4.759637in}{1.886851in}}%
\pgfpathlineto{\pgfqpoint{4.762156in}{1.900104in}}%
\pgfpathlineto{\pgfqpoint{4.764674in}{1.881573in}}%
\pgfpathlineto{\pgfqpoint{4.767193in}{1.909169in}}%
\pgfpathlineto{\pgfqpoint{4.769712in}{1.900797in}}%
\pgfpathlineto{\pgfqpoint{4.772231in}{1.896356in}}%
\pgfpathlineto{\pgfqpoint{4.774749in}{1.914849in}}%
\pgfpathlineto{\pgfqpoint{4.777268in}{1.909463in}}%
\pgfpathlineto{\pgfqpoint{4.779787in}{1.908249in}}%
\pgfpathlineto{\pgfqpoint{4.782306in}{1.891464in}}%
\pgfpathlineto{\pgfqpoint{4.784824in}{1.898225in}}%
\pgfpathlineto{\pgfqpoint{4.787343in}{1.902786in}}%
\pgfpathlineto{\pgfqpoint{4.789862in}{1.912040in}}%
\pgfpathlineto{\pgfqpoint{4.792381in}{1.909240in}}%
\pgfpathlineto{\pgfqpoint{4.794899in}{1.896412in}}%
\pgfpathlineto{\pgfqpoint{4.797418in}{1.892570in}}%
\pgfpathlineto{\pgfqpoint{4.799937in}{1.902888in}}%
\pgfpathlineto{\pgfqpoint{4.802456in}{1.899032in}}%
\pgfpathlineto{\pgfqpoint{4.804974in}{1.925428in}}%
\pgfpathlineto{\pgfqpoint{4.807493in}{1.938281in}}%
\pgfpathlineto{\pgfqpoint{4.810012in}{1.918499in}}%
\pgfpathlineto{\pgfqpoint{4.812531in}{1.937698in}}%
\pgfpathlineto{\pgfqpoint{4.815049in}{1.931300in}}%
\pgfpathlineto{\pgfqpoint{4.817568in}{1.917344in}}%
\pgfpathlineto{\pgfqpoint{4.820087in}{1.938648in}}%
\pgfpathlineto{\pgfqpoint{4.822606in}{1.976795in}}%
\pgfpathlineto{\pgfqpoint{4.825124in}{1.969105in}}%
\pgfpathlineto{\pgfqpoint{4.827643in}{1.978718in}}%
\pgfpathlineto{\pgfqpoint{4.830162in}{1.969759in}}%
\pgfpathlineto{\pgfqpoint{4.832681in}{1.942705in}}%
\pgfpathlineto{\pgfqpoint{4.835199in}{1.938770in}}%
\pgfpathlineto{\pgfqpoint{4.837718in}{1.926029in}}%
\pgfpathlineto{\pgfqpoint{4.840237in}{1.902820in}}%
\pgfpathlineto{\pgfqpoint{4.842756in}{1.906497in}}%
\pgfpathlineto{\pgfqpoint{4.845274in}{1.922987in}}%
\pgfpathlineto{\pgfqpoint{4.847793in}{1.947200in}}%
\pgfpathlineto{\pgfqpoint{4.850312in}{1.944921in}}%
\pgfpathlineto{\pgfqpoint{4.852831in}{1.933145in}}%
\pgfpathlineto{\pgfqpoint{4.855349in}{1.916687in}}%
\pgfpathlineto{\pgfqpoint{4.857868in}{1.924776in}}%
\pgfpathlineto{\pgfqpoint{4.860387in}{1.924292in}}%
\pgfpathlineto{\pgfqpoint{4.862906in}{1.926427in}}%
\pgfpathlineto{\pgfqpoint{4.865424in}{1.949439in}}%
\pgfpathlineto{\pgfqpoint{4.867943in}{1.944614in}}%
\pgfpathlineto{\pgfqpoint{4.870462in}{1.938361in}}%
\pgfpathlineto{\pgfqpoint{4.872981in}{1.914954in}}%
\pgfpathlineto{\pgfqpoint{4.875499in}{1.905911in}}%
\pgfpathlineto{\pgfqpoint{4.878018in}{1.892047in}}%
\pgfpathlineto{\pgfqpoint{4.880537in}{1.882272in}}%
\pgfpathlineto{\pgfqpoint{4.883056in}{1.882177in}}%
\pgfpathlineto{\pgfqpoint{4.885574in}{1.874102in}}%
\pgfpathlineto{\pgfqpoint{4.888093in}{1.865662in}}%
\pgfpathlineto{\pgfqpoint{4.890612in}{1.851452in}}%
\pgfpathlineto{\pgfqpoint{4.893131in}{1.845903in}}%
\pgfpathlineto{\pgfqpoint{4.895649in}{1.841967in}}%
\pgfpathlineto{\pgfqpoint{4.898168in}{1.862181in}}%
\pgfpathlineto{\pgfqpoint{4.900687in}{1.861270in}}%
\pgfpathlineto{\pgfqpoint{4.903206in}{1.838268in}}%
\pgfpathlineto{\pgfqpoint{4.905724in}{1.820534in}}%
\pgfpathlineto{\pgfqpoint{4.908243in}{1.831698in}}%
\pgfpathlineto{\pgfqpoint{4.910762in}{1.846485in}}%
\pgfpathlineto{\pgfqpoint{4.913281in}{1.831668in}}%
\pgfpathlineto{\pgfqpoint{4.915799in}{1.868656in}}%
\pgfpathlineto{\pgfqpoint{4.918318in}{1.892838in}}%
\pgfpathlineto{\pgfqpoint{4.920837in}{1.892713in}}%
\pgfpathlineto{\pgfqpoint{4.923356in}{1.922918in}}%
\pgfpathlineto{\pgfqpoint{4.925874in}{1.930501in}}%
\pgfpathlineto{\pgfqpoint{4.928393in}{1.935339in}}%
\pgfpathlineto{\pgfqpoint{4.930912in}{1.949613in}}%
\pgfpathlineto{\pgfqpoint{4.933431in}{1.938164in}}%
\pgfpathlineto{\pgfqpoint{4.935949in}{1.924654in}}%
\pgfpathlineto{\pgfqpoint{4.938468in}{1.957078in}}%
\pgfpathlineto{\pgfqpoint{4.940987in}{1.977978in}}%
\pgfpathlineto{\pgfqpoint{4.943506in}{1.966309in}}%
\pgfpathlineto{\pgfqpoint{4.946024in}{1.971502in}}%
\pgfpathlineto{\pgfqpoint{4.948543in}{1.975622in}}%
\pgfpathlineto{\pgfqpoint{4.951062in}{1.987184in}}%
\pgfpathlineto{\pgfqpoint{4.953581in}{1.999911in}}%
\pgfpathlineto{\pgfqpoint{4.956099in}{2.007956in}}%
\pgfpathlineto{\pgfqpoint{4.958618in}{2.004112in}}%
\pgfpathlineto{\pgfqpoint{4.961137in}{1.985321in}}%
\pgfpathlineto{\pgfqpoint{4.966174in}{1.942613in}}%
\pgfpathlineto{\pgfqpoint{4.968693in}{1.943195in}}%
\pgfpathlineto{\pgfqpoint{4.971212in}{1.950092in}}%
\pgfpathlineto{\pgfqpoint{4.973731in}{1.940751in}}%
\pgfpathlineto{\pgfqpoint{4.976249in}{1.938842in}}%
\pgfpathlineto{\pgfqpoint{4.978768in}{1.924344in}}%
\pgfpathlineto{\pgfqpoint{4.981287in}{1.896543in}}%
\pgfpathlineto{\pgfqpoint{4.983806in}{1.896283in}}%
\pgfpathlineto{\pgfqpoint{4.986324in}{1.904557in}}%
\pgfpathlineto{\pgfqpoint{4.988843in}{1.900939in}}%
\pgfpathlineto{\pgfqpoint{4.991362in}{1.910557in}}%
\pgfpathlineto{\pgfqpoint{4.993881in}{1.921994in}}%
\pgfpathlineto{\pgfqpoint{4.996399in}{1.924140in}}%
\pgfpathlineto{\pgfqpoint{4.998918in}{1.899820in}}%
\pgfpathlineto{\pgfqpoint{5.001437in}{1.902674in}}%
\pgfpathlineto{\pgfqpoint{5.003956in}{1.915007in}}%
\pgfpathlineto{\pgfqpoint{5.006474in}{1.909289in}}%
\pgfpathlineto{\pgfqpoint{5.008993in}{1.888758in}}%
\pgfpathlineto{\pgfqpoint{5.011512in}{1.861864in}}%
\pgfpathlineto{\pgfqpoint{5.014031in}{1.891680in}}%
\pgfpathlineto{\pgfqpoint{5.016549in}{1.893338in}}%
\pgfpathlineto{\pgfqpoint{5.019068in}{1.904559in}}%
\pgfpathlineto{\pgfqpoint{5.021587in}{1.897667in}}%
\pgfpathlineto{\pgfqpoint{5.024106in}{1.883012in}}%
\pgfpathlineto{\pgfqpoint{5.026624in}{1.883764in}}%
\pgfpathlineto{\pgfqpoint{5.029143in}{1.858670in}}%
\pgfpathlineto{\pgfqpoint{5.031662in}{1.848886in}}%
\pgfpathlineto{\pgfqpoint{5.034181in}{1.851473in}}%
\pgfpathlineto{\pgfqpoint{5.036699in}{1.855075in}}%
\pgfpathlineto{\pgfqpoint{5.039218in}{1.860053in}}%
\pgfpathlineto{\pgfqpoint{5.041737in}{1.866259in}}%
\pgfpathlineto{\pgfqpoint{5.044256in}{1.849882in}}%
\pgfpathlineto{\pgfqpoint{5.046774in}{1.834490in}}%
\pgfpathlineto{\pgfqpoint{5.049293in}{1.856460in}}%
\pgfpathlineto{\pgfqpoint{5.051812in}{1.854864in}}%
\pgfpathlineto{\pgfqpoint{5.054331in}{1.831688in}}%
\pgfpathlineto{\pgfqpoint{5.056849in}{1.839459in}}%
\pgfpathlineto{\pgfqpoint{5.059368in}{1.850573in}}%
\pgfpathlineto{\pgfqpoint{5.061887in}{1.838485in}}%
\pgfpathlineto{\pgfqpoint{5.064406in}{1.838908in}}%
\pgfpathlineto{\pgfqpoint{5.066924in}{1.864581in}}%
\pgfpathlineto{\pgfqpoint{5.069443in}{1.854925in}}%
\pgfpathlineto{\pgfqpoint{5.071962in}{1.823230in}}%
\pgfpathlineto{\pgfqpoint{5.074481in}{1.821414in}}%
\pgfpathlineto{\pgfqpoint{5.076999in}{1.839650in}}%
\pgfpathlineto{\pgfqpoint{5.079518in}{1.827433in}}%
\pgfpathlineto{\pgfqpoint{5.082037in}{1.835330in}}%
\pgfpathlineto{\pgfqpoint{5.084556in}{1.836806in}}%
\pgfpathlineto{\pgfqpoint{5.087074in}{1.820728in}}%
\pgfpathlineto{\pgfqpoint{5.089593in}{1.835897in}}%
\pgfpathlineto{\pgfqpoint{5.092112in}{1.848399in}}%
\pgfpathlineto{\pgfqpoint{5.094631in}{1.828776in}}%
\pgfpathlineto{\pgfqpoint{5.097149in}{1.786992in}}%
\pgfpathlineto{\pgfqpoint{5.099668in}{1.782301in}}%
\pgfpathlineto{\pgfqpoint{5.102187in}{1.787462in}}%
\pgfpathlineto{\pgfqpoint{5.107224in}{1.765262in}}%
\pgfpathlineto{\pgfqpoint{5.109743in}{1.766465in}}%
\pgfpathlineto{\pgfqpoint{5.112262in}{1.791615in}}%
\pgfpathlineto{\pgfqpoint{5.114781in}{1.752584in}}%
\pgfpathlineto{\pgfqpoint{5.117299in}{1.735799in}}%
\pgfpathlineto{\pgfqpoint{5.119818in}{1.702766in}}%
\pgfpathlineto{\pgfqpoint{5.122337in}{1.717607in}}%
\pgfpathlineto{\pgfqpoint{5.124856in}{1.738447in}}%
\pgfpathlineto{\pgfqpoint{5.127374in}{1.726084in}}%
\pgfpathlineto{\pgfqpoint{5.129893in}{1.707565in}}%
\pgfpathlineto{\pgfqpoint{5.132412in}{1.682665in}}%
\pgfpathlineto{\pgfqpoint{5.134931in}{1.694493in}}%
\pgfpathlineto{\pgfqpoint{5.137449in}{1.722519in}}%
\pgfpathlineto{\pgfqpoint{5.139968in}{1.691778in}}%
\pgfpathlineto{\pgfqpoint{5.142487in}{1.707752in}}%
\pgfpathlineto{\pgfqpoint{5.145006in}{1.686264in}}%
\pgfpathlineto{\pgfqpoint{5.147524in}{1.674378in}}%
\pgfpathlineto{\pgfqpoint{5.150043in}{1.665214in}}%
\pgfpathlineto{\pgfqpoint{5.152562in}{1.638735in}}%
\pgfpathlineto{\pgfqpoint{5.155081in}{1.645784in}}%
\pgfpathlineto{\pgfqpoint{5.157599in}{1.640838in}}%
\pgfpathlineto{\pgfqpoint{5.160118in}{1.640142in}}%
\pgfpathlineto{\pgfqpoint{5.162637in}{1.657532in}}%
\pgfpathlineto{\pgfqpoint{5.165156in}{1.678479in}}%
\pgfpathlineto{\pgfqpoint{5.167674in}{1.693747in}}%
\pgfpathlineto{\pgfqpoint{5.170193in}{1.686823in}}%
\pgfpathlineto{\pgfqpoint{5.172712in}{1.673701in}}%
\pgfpathlineto{\pgfqpoint{5.175231in}{1.668853in}}%
\pgfpathlineto{\pgfqpoint{5.177749in}{1.676857in}}%
\pgfpathlineto{\pgfqpoint{5.180268in}{1.665205in}}%
\pgfpathlineto{\pgfqpoint{5.182787in}{1.672638in}}%
\pgfpathlineto{\pgfqpoint{5.185306in}{1.683065in}}%
\pgfpathlineto{\pgfqpoint{5.187824in}{1.695280in}}%
\pgfpathlineto{\pgfqpoint{5.190343in}{1.674606in}}%
\pgfpathlineto{\pgfqpoint{5.192862in}{1.701776in}}%
\pgfpathlineto{\pgfqpoint{5.195381in}{1.681299in}}%
\pgfpathlineto{\pgfqpoint{5.197899in}{1.683580in}}%
\pgfpathlineto{\pgfqpoint{5.200418in}{1.670752in}}%
\pgfpathlineto{\pgfqpoint{5.202937in}{1.673043in}}%
\pgfpathlineto{\pgfqpoint{5.205456in}{1.679991in}}%
\pgfpathlineto{\pgfqpoint{5.207974in}{1.675389in}}%
\pgfpathlineto{\pgfqpoint{5.210493in}{1.701089in}}%
\pgfpathlineto{\pgfqpoint{5.213012in}{1.691571in}}%
\pgfpathlineto{\pgfqpoint{5.215531in}{1.694746in}}%
\pgfpathlineto{\pgfqpoint{5.218049in}{1.682871in}}%
\pgfpathlineto{\pgfqpoint{5.220568in}{1.694502in}}%
\pgfpathlineto{\pgfqpoint{5.223087in}{1.695272in}}%
\pgfpathlineto{\pgfqpoint{5.225606in}{1.683981in}}%
\pgfpathlineto{\pgfqpoint{5.228124in}{1.692508in}}%
\pgfpathlineto{\pgfqpoint{5.230643in}{1.681984in}}%
\pgfpathlineto{\pgfqpoint{5.233162in}{1.671956in}}%
\pgfpathlineto{\pgfqpoint{5.235681in}{1.673006in}}%
\pgfpathlineto{\pgfqpoint{5.238199in}{1.693427in}}%
\pgfpathlineto{\pgfqpoint{5.240718in}{1.688965in}}%
\pgfpathlineto{\pgfqpoint{5.243237in}{1.708092in}}%
\pgfpathlineto{\pgfqpoint{5.245756in}{1.702427in}}%
\pgfpathlineto{\pgfqpoint{5.248274in}{1.715548in}}%
\pgfpathlineto{\pgfqpoint{5.250793in}{1.714027in}}%
\pgfpathlineto{\pgfqpoint{5.253312in}{1.699657in}}%
\pgfpathlineto{\pgfqpoint{5.255831in}{1.664950in}}%
\pgfpathlineto{\pgfqpoint{5.258349in}{1.689959in}}%
\pgfpathlineto{\pgfqpoint{5.260868in}{1.689121in}}%
\pgfpathlineto{\pgfqpoint{5.263387in}{1.695603in}}%
\pgfpathlineto{\pgfqpoint{5.265906in}{1.686837in}}%
\pgfpathlineto{\pgfqpoint{5.268424in}{1.670336in}}%
\pgfpathlineto{\pgfqpoint{5.270943in}{1.685403in}}%
\pgfpathlineto{\pgfqpoint{5.273462in}{1.707064in}}%
\pgfpathlineto{\pgfqpoint{5.275981in}{1.731712in}}%
\pgfpathlineto{\pgfqpoint{5.278499in}{1.723966in}}%
\pgfpathlineto{\pgfqpoint{5.281018in}{1.719941in}}%
\pgfpathlineto{\pgfqpoint{5.283537in}{1.738241in}}%
\pgfpathlineto{\pgfqpoint{5.286056in}{1.749635in}}%
\pgfpathlineto{\pgfqpoint{5.288574in}{1.765468in}}%
\pgfpathlineto{\pgfqpoint{5.291093in}{1.783511in}}%
\pgfpathlineto{\pgfqpoint{5.293612in}{1.810418in}}%
\pgfpathlineto{\pgfqpoint{5.296131in}{1.815652in}}%
\pgfpathlineto{\pgfqpoint{5.298649in}{1.808175in}}%
\pgfpathlineto{\pgfqpoint{5.301168in}{1.831649in}}%
\pgfpathlineto{\pgfqpoint{5.303687in}{1.874316in}}%
\pgfpathlineto{\pgfqpoint{5.306206in}{1.884986in}}%
\pgfpathlineto{\pgfqpoint{5.308724in}{1.906907in}}%
\pgfpathlineto{\pgfqpoint{5.311243in}{1.918153in}}%
\pgfpathlineto{\pgfqpoint{5.313762in}{1.924767in}}%
\pgfpathlineto{\pgfqpoint{5.318799in}{1.942064in}}%
\pgfpathlineto{\pgfqpoint{5.321318in}{1.964056in}}%
\pgfpathlineto{\pgfqpoint{5.323837in}{1.995375in}}%
\pgfpathlineto{\pgfqpoint{5.326356in}{1.989453in}}%
\pgfpathlineto{\pgfqpoint{5.328874in}{1.986809in}}%
\pgfpathlineto{\pgfqpoint{5.331393in}{2.024526in}}%
\pgfpathlineto{\pgfqpoint{5.333912in}{2.043887in}}%
\pgfpathlineto{\pgfqpoint{5.336431in}{2.051447in}}%
\pgfpathlineto{\pgfqpoint{5.338949in}{2.052906in}}%
\pgfpathlineto{\pgfqpoint{5.341468in}{2.044746in}}%
\pgfpathlineto{\pgfqpoint{5.343987in}{2.027223in}}%
\pgfpathlineto{\pgfqpoint{5.346506in}{2.031524in}}%
\pgfpathlineto{\pgfqpoint{5.349024in}{2.038108in}}%
\pgfpathlineto{\pgfqpoint{5.351543in}{2.014794in}}%
\pgfpathlineto{\pgfqpoint{5.354062in}{2.003531in}}%
\pgfpathlineto{\pgfqpoint{5.356581in}{2.025149in}}%
\pgfpathlineto{\pgfqpoint{5.359099in}{2.035512in}}%
\pgfpathlineto{\pgfqpoint{5.361618in}{2.056971in}}%
\pgfpathlineto{\pgfqpoint{5.364137in}{2.064330in}}%
\pgfpathlineto{\pgfqpoint{5.366656in}{2.030808in}}%
\pgfpathlineto{\pgfqpoint{5.369174in}{2.035952in}}%
\pgfpathlineto{\pgfqpoint{5.371693in}{2.020188in}}%
\pgfpathlineto{\pgfqpoint{5.374212in}{2.020656in}}%
\pgfpathlineto{\pgfqpoint{5.376731in}{2.021614in}}%
\pgfpathlineto{\pgfqpoint{5.379249in}{2.009344in}}%
\pgfpathlineto{\pgfqpoint{5.381768in}{2.017018in}}%
\pgfpathlineto{\pgfqpoint{5.384287in}{2.026750in}}%
\pgfpathlineto{\pgfqpoint{5.386806in}{2.046693in}}%
\pgfpathlineto{\pgfqpoint{5.389324in}{2.018536in}}%
\pgfpathlineto{\pgfqpoint{5.391843in}{2.007240in}}%
\pgfpathlineto{\pgfqpoint{5.394362in}{2.030038in}}%
\pgfpathlineto{\pgfqpoint{5.396881in}{2.016596in}}%
\pgfpathlineto{\pgfqpoint{5.399399in}{2.023987in}}%
\pgfpathlineto{\pgfqpoint{5.401918in}{2.010047in}}%
\pgfpathlineto{\pgfqpoint{5.404437in}{1.994597in}}%
\pgfpathlineto{\pgfqpoint{5.406956in}{2.009498in}}%
\pgfpathlineto{\pgfqpoint{5.409474in}{2.002303in}}%
\pgfpathlineto{\pgfqpoint{5.411993in}{1.962360in}}%
\pgfpathlineto{\pgfqpoint{5.414512in}{1.949820in}}%
\pgfpathlineto{\pgfqpoint{5.417031in}{1.929348in}}%
\pgfpathlineto{\pgfqpoint{5.419549in}{1.921639in}}%
\pgfpathlineto{\pgfqpoint{5.422068in}{1.911776in}}%
\pgfpathlineto{\pgfqpoint{5.424587in}{1.925493in}}%
\pgfpathlineto{\pgfqpoint{5.427106in}{1.918633in}}%
\pgfpathlineto{\pgfqpoint{5.429624in}{1.930506in}}%
\pgfpathlineto{\pgfqpoint{5.432143in}{1.940860in}}%
\pgfpathlineto{\pgfqpoint{5.434662in}{1.965314in}}%
\pgfpathlineto{\pgfqpoint{5.437181in}{1.956883in}}%
\pgfpathlineto{\pgfqpoint{5.439699in}{2.006310in}}%
\pgfpathlineto{\pgfqpoint{5.442218in}{2.014249in}}%
\pgfpathlineto{\pgfqpoint{5.444737in}{2.005001in}}%
\pgfpathlineto{\pgfqpoint{5.447256in}{1.988814in}}%
\pgfpathlineto{\pgfqpoint{5.449774in}{1.986904in}}%
\pgfpathlineto{\pgfqpoint{5.452293in}{1.947620in}}%
\pgfpathlineto{\pgfqpoint{5.454812in}{1.971419in}}%
\pgfpathlineto{\pgfqpoint{5.457331in}{1.971437in}}%
\pgfpathlineto{\pgfqpoint{5.459849in}{1.964058in}}%
\pgfpathlineto{\pgfqpoint{5.462368in}{1.967044in}}%
\pgfpathlineto{\pgfqpoint{5.464887in}{1.995650in}}%
\pgfpathlineto{\pgfqpoint{5.467406in}{1.998723in}}%
\pgfpathlineto{\pgfqpoint{5.469924in}{2.003389in}}%
\pgfpathlineto{\pgfqpoint{5.472443in}{1.992061in}}%
\pgfpathlineto{\pgfqpoint{5.474962in}{1.985580in}}%
\pgfpathlineto{\pgfqpoint{5.477481in}{1.966760in}}%
\pgfpathlineto{\pgfqpoint{5.479999in}{1.992981in}}%
\pgfpathlineto{\pgfqpoint{5.482518in}{2.001544in}}%
\pgfpathlineto{\pgfqpoint{5.485037in}{2.012280in}}%
\pgfpathlineto{\pgfqpoint{5.487556in}{2.017532in}}%
\pgfpathlineto{\pgfqpoint{5.490074in}{2.025200in}}%
\pgfpathlineto{\pgfqpoint{5.492593in}{2.041757in}}%
\pgfpathlineto{\pgfqpoint{5.492593in}{2.041757in}}%
\pgfusepath{stroke}%
\end{pgfscope}%
\begin{pgfscope}%
\pgfsetrectcap%
\pgfsetmiterjoin%
\pgfsetlinewidth{0.501875pt}%
\definecolor{currentstroke}{rgb}{0.000000,0.000000,0.000000}%
\pgfsetstrokecolor{currentstroke}%
\pgfsetdash{}{0pt}%
\pgfpathmoveto{\pgfqpoint{0.455093in}{0.401080in}}%
\pgfpathlineto{\pgfqpoint{0.455093in}{4.405080in}}%
\pgfusepath{stroke}%
\end{pgfscope}%
\begin{pgfscope}%
\pgfsetrectcap%
\pgfsetmiterjoin%
\pgfsetlinewidth{0.501875pt}%
\definecolor{currentstroke}{rgb}{0.000000,0.000000,0.000000}%
\pgfsetstrokecolor{currentstroke}%
\pgfsetdash{}{0pt}%
\pgfpathmoveto{\pgfqpoint{5.492593in}{0.401080in}}%
\pgfpathlineto{\pgfqpoint{5.492593in}{4.405080in}}%
\pgfusepath{stroke}%
\end{pgfscope}%
\begin{pgfscope}%
\pgfsetrectcap%
\pgfsetmiterjoin%
\pgfsetlinewidth{0.501875pt}%
\definecolor{currentstroke}{rgb}{0.000000,0.000000,0.000000}%
\pgfsetstrokecolor{currentstroke}%
\pgfsetdash{}{0pt}%
\pgfpathmoveto{\pgfqpoint{0.455093in}{0.401080in}}%
\pgfpathlineto{\pgfqpoint{5.492593in}{0.401080in}}%
\pgfusepath{stroke}%
\end{pgfscope}%
\begin{pgfscope}%
\pgfsetrectcap%
\pgfsetmiterjoin%
\pgfsetlinewidth{0.501875pt}%
\definecolor{currentstroke}{rgb}{0.000000,0.000000,0.000000}%
\pgfsetstrokecolor{currentstroke}%
\pgfsetdash{}{0pt}%
\pgfpathmoveto{\pgfqpoint{0.455093in}{4.405080in}}%
\pgfpathlineto{\pgfqpoint{5.492593in}{4.405080in}}%
\pgfusepath{stroke}%
\end{pgfscope}%
\begin{pgfscope}%
\pgfsetrectcap%
\pgfsetroundjoin%
\pgfsetlinewidth{1.505625pt}%
\definecolor{currentstroke}{rgb}{0.000000,0.000000,0.000000}%
\pgfsetstrokecolor{currentstroke}%
\pgfsetdash{}{0pt}%
\pgfpathmoveto{\pgfqpoint{4.243748in}{4.231469in}}%
\pgfpathlineto{\pgfqpoint{4.382637in}{4.231469in}}%
\pgfpathlineto{\pgfqpoint{4.521526in}{4.231469in}}%
\pgfusepath{stroke}%
\end{pgfscope}%
\begin{pgfscope}%
\definecolor{textcolor}{rgb}{0.000000,0.000000,0.000000}%
\pgfsetstrokecolor{textcolor}%
\pgfsetfillcolor{textcolor}%
\pgftext[x=4.632637in,y=4.182858in,left,base]{\color{textcolor}{\rmfamily\fontsize{10.000000}{12.000000}\selectfont\catcode`\^=\active\def^{\ifmmode\sp\else\^{}\fi}\catcode`\%=\active\def%{\%}Expectation}}%
\end{pgfscope}%
\begin{pgfscope}%
\pgfpathrectangle{\pgfqpoint{0.455093in}{0.401080in}}{\pgfqpoint{5.037500in}{4.004000in}}%
\pgfusepath{clip}%
\pgfsetrectcap%
\pgfsetroundjoin%
\pgfsetlinewidth{1.505625pt}%
\definecolor{currentstroke}{rgb}{0.000000,0.000000,0.000000}%
\pgfsetstrokecolor{currentstroke}%
\pgfsetdash{}{0pt}%
\pgfpathmoveto{\pgfqpoint{0.455093in}{1.021044in}}%
\pgfpathlineto{\pgfqpoint{0.535693in}{1.044902in}}%
\pgfpathlineto{\pgfqpoint{0.616293in}{1.068951in}}%
\pgfpathlineto{\pgfqpoint{0.696893in}{1.093194in}}%
\pgfpathlineto{\pgfqpoint{0.777493in}{1.117631in}}%
\pgfpathlineto{\pgfqpoint{0.858093in}{1.142264in}}%
\pgfpathlineto{\pgfqpoint{0.938693in}{1.167096in}}%
\pgfpathlineto{\pgfqpoint{1.019293in}{1.192127in}}%
\pgfpathlineto{\pgfqpoint{1.097374in}{1.216567in}}%
\pgfpathlineto{\pgfqpoint{1.175456in}{1.241197in}}%
\pgfpathlineto{\pgfqpoint{1.253537in}{1.266020in}}%
\pgfpathlineto{\pgfqpoint{1.331618in}{1.291035in}}%
\pgfpathlineto{\pgfqpoint{1.409699in}{1.316245in}}%
\pgfpathlineto{\pgfqpoint{1.487781in}{1.341651in}}%
\pgfpathlineto{\pgfqpoint{1.565862in}{1.367254in}}%
\pgfpathlineto{\pgfqpoint{1.643943in}{1.393057in}}%
\pgfpathlineto{\pgfqpoint{1.722024in}{1.419061in}}%
\pgfpathlineto{\pgfqpoint{1.797587in}{1.444418in}}%
\pgfpathlineto{\pgfqpoint{1.873149in}{1.469967in}}%
\pgfpathlineto{\pgfqpoint{1.948712in}{1.495707in}}%
\pgfpathlineto{\pgfqpoint{2.024274in}{1.521642in}}%
\pgfpathlineto{\pgfqpoint{2.099837in}{1.547771in}}%
\pgfpathlineto{\pgfqpoint{2.175399in}{1.574098in}}%
\pgfpathlineto{\pgfqpoint{2.250962in}{1.600622in}}%
\pgfpathlineto{\pgfqpoint{2.326524in}{1.627346in}}%
\pgfpathlineto{\pgfqpoint{2.402087in}{1.654272in}}%
\pgfpathlineto{\pgfqpoint{2.477649in}{1.681400in}}%
\pgfpathlineto{\pgfqpoint{2.550693in}{1.707818in}}%
\pgfpathlineto{\pgfqpoint{2.623737in}{1.734428in}}%
\pgfpathlineto{\pgfqpoint{2.696781in}{1.761232in}}%
\pgfpathlineto{\pgfqpoint{2.769824in}{1.788231in}}%
\pgfpathlineto{\pgfqpoint{2.842868in}{1.815426in}}%
\pgfpathlineto{\pgfqpoint{2.915912in}{1.842819in}}%
\pgfpathlineto{\pgfqpoint{2.988956in}{1.870412in}}%
\pgfpathlineto{\pgfqpoint{3.061999in}{1.898205in}}%
\pgfpathlineto{\pgfqpoint{3.135043in}{1.926200in}}%
\pgfpathlineto{\pgfqpoint{3.208087in}{1.954400in}}%
\pgfpathlineto{\pgfqpoint{3.278612in}{1.981821in}}%
\pgfpathlineto{\pgfqpoint{3.349137in}{2.009435in}}%
\pgfpathlineto{\pgfqpoint{3.419662in}{2.037243in}}%
\pgfpathlineto{\pgfqpoint{3.490187in}{2.065247in}}%
\pgfpathlineto{\pgfqpoint{3.560712in}{2.093447in}}%
\pgfpathlineto{\pgfqpoint{3.631237in}{2.121845in}}%
\pgfpathlineto{\pgfqpoint{3.701762in}{2.150443in}}%
\pgfpathlineto{\pgfqpoint{3.772287in}{2.179241in}}%
\pgfpathlineto{\pgfqpoint{3.842812in}{2.208242in}}%
\pgfpathlineto{\pgfqpoint{3.913337in}{2.237447in}}%
\pgfpathlineto{\pgfqpoint{3.983862in}{2.266857in}}%
\pgfpathlineto{\pgfqpoint{4.054387in}{2.296473in}}%
\pgfpathlineto{\pgfqpoint{4.122393in}{2.325229in}}%
\pgfpathlineto{\pgfqpoint{4.190399in}{2.354179in}}%
\pgfpathlineto{\pgfqpoint{4.258406in}{2.383326in}}%
\pgfpathlineto{\pgfqpoint{4.326412in}{2.412670in}}%
\pgfpathlineto{\pgfqpoint{4.394418in}{2.442212in}}%
\pgfpathlineto{\pgfqpoint{4.462424in}{2.471955in}}%
\pgfpathlineto{\pgfqpoint{4.530431in}{2.501899in}}%
\pgfpathlineto{\pgfqpoint{4.598437in}{2.532046in}}%
\pgfpathlineto{\pgfqpoint{4.666443in}{2.562397in}}%
\pgfpathlineto{\pgfqpoint{4.734449in}{2.592954in}}%
\pgfpathlineto{\pgfqpoint{4.802456in}{2.623718in}}%
\pgfpathlineto{\pgfqpoint{4.870462in}{2.654690in}}%
\pgfpathlineto{\pgfqpoint{4.938468in}{2.685871in}}%
\pgfpathlineto{\pgfqpoint{5.006474in}{2.717264in}}%
\pgfpathlineto{\pgfqpoint{5.071962in}{2.747696in}}%
\pgfpathlineto{\pgfqpoint{5.137449in}{2.778325in}}%
\pgfpathlineto{\pgfqpoint{5.202937in}{2.809155in}}%
\pgfpathlineto{\pgfqpoint{5.268424in}{2.840185in}}%
\pgfpathlineto{\pgfqpoint{5.333912in}{2.871418in}}%
\pgfpathlineto{\pgfqpoint{5.399399in}{2.902854in}}%
\pgfpathlineto{\pgfqpoint{5.464887in}{2.934496in}}%
\pgfpathlineto{\pgfqpoint{5.492593in}{2.947945in}}%
\pgfpathlineto{\pgfqpoint{5.492593in}{2.947945in}}%
\pgfusepath{stroke}%
\end{pgfscope}%
\end{pgfpicture}%
\makeatother%
\endgroup%

    \caption{Geometric Brownian motion sample path and expectation for $\mu=0.5$, $\sigma=0.2$ and $S_0=1$.}
    \label{fig:geometricBrownianMotion}
\end{figure}